\documentclass[10pt,a4paper,oneside,openany]{book}
\usepackage[margin=19mm,top=18mm,bottom=19mm,headheight=14pt]{geometry}
\usepackage[T1]{fontenc}
\usepackage[utf8]{inputenc}
\usepackage{lmodern,microtype,amsmath,amssymb,amsthm,mathtools,mathrsfs}
\usepackage{booktabs,array,longtable,graphicx,xcolor,tikz,xurl}
\usepackage{titlesec,fancyhdr}
\usepackage[hypertexnames=false,hidelinks]{hyperref}
\newtheorem{theorem}{Theorem}[chapter]
\newtheorem{lemma}[theorem]{Lemma}
\newtheorem{proposition}[theorem]{Proposition}
\newtheorem{corollary}[theorem]{Corollary}
\newtheorem{claim}[theorem]{Claim}
\theoremstyle{definition}
\newtheorem{definition}[theorem]{Definition}
\theoremstyle{remark}
\newtheorem{remark}[theorem]{Remark}
\providecommand{\R}{\mathbb R}
\providecommand{\Q}{\mathbb Q}
\providecommand{\Rc}{\mathbb Q[\delta,k,b,a,t]}
\providecommand{\I}{\mathsf I}
\providecommand{\U}{\mathsf U}
\providecommand{\Otype}{\mathsf O}
\providecommand{\SOS}{\operatorname{SOS}}
\providecommand{\Inf}{\operatorname{Inf}}
\providecommand{\res}{\operatorname{res}}
\providecommand{\Sym}{\operatorname{Sym}}
\providecommand{\Ow}{O_{\mathrm{wt}}}
\providecommand{\wt}{\operatorname{wt}}
\providecommand{\norm}[1]{\left\lVert #1\right\rVert}
\providecommand{\status}[1]{\par\smallskip\noindent\textit{Scope. }#1\par\smallskip}
\providecommand{\tightlist}{\setlength{\itemsep}{0pt}\setlength{\parskip}{0pt}}
\DeclareUnicodeCharacter{2264}{\ensuremath{\leq}}
\DeclareUnicodeCharacter{2265}{\ensuremath{\geq}}
\DeclareUnicodeCharacter{221E}{\ensuremath{\infty}}
\DeclareUnicodeCharacter{2212}{-}
\DeclareUnicodeCharacter{03B1}{\ensuremath{\alpha}}
\DeclareUnicodeCharacter{03B2}{\ensuremath{\beta}}
\DeclareUnicodeCharacter{03B4}{\ensuremath{\delta}}
\DeclareUnicodeCharacter{03BB}{\ensuremath{\lambda}}
\DeclareUnicodeCharacter{03C1}{\ensuremath{\rho}}
\DeclareUnicodeCharacter{03C3}{\ensuremath{\sigma}}
\DeclareUnicodeCharacter{03C4}{\ensuremath{\tau}}
\DeclareUnicodeCharacter{03A3}{\ensuremath{\Sigma}}
\DeclareUnicodeCharacter{00D7}{\ensuremath{\times}}
\DeclareUnicodeCharacter{2192}{\ensuremath{\to}}
\DeclareUnicodeCharacter{21D4}{\ensuremath{\Longleftrightarrow}}
\allowdisplaybreaks[2]
\setlength{\emergencystretch}{2em}
\setlength{\parskip}{2pt}
\setlength{\parindent}{1em}
\titleformat{\chapter}[display]{\normalfont\Large\bfseries}{\chaptername\ \thechapter}{4pt}{\LARGE}
\titlespacing*{\chapter}{0pt}{0pt}{16pt}
\titlespacing*{\section}{0pt}{14pt plus 3pt minus 2pt}{6pt}
\titlespacing*{\subsection}{0pt}{10pt plus 2pt minus 2pt}{4pt}
\titlespacing*{\subsubsection}{0pt}{8pt plus 2pt minus 2pt}{3pt}
\pagestyle{fancy}
\fancyhf{}
\fancyhead[L]{\small Q222: finite semidefinite representability}
\fancyhead[R]{\small\thepage}
\renewcommand{\headrulewidth}{0.3pt}
\fancypagestyle{plain}{\fancyhf{}\fancyfoot[C]{\thepage}\renewcommand{\headrulewidth}{0pt}}
\setcounter{tocdepth}{1}
\hypersetup{bookmarksdepth=2,pdftitle={Complete proofs: finite semidefinite representation of the simplest quantum correlation set},pdfsubject={Complete qualitative existence argument and explicit full-slack certificates},pdfauthor={}}
\begin{document}
\frontmatter
\begin{titlepage}
\centering
\vspace*{25mm}
{\LARGE\bfseries Finite semidefinite representation\\[4pt]
of the simplest quantum correlation set\par}
\vspace{16mm}
{\Huge Complete proofs\par}
\vspace{12mm}
{\large 23 October 2026\par}
\vfill
\begin{minipage}{0.88\textwidth}
A Bell contact balances eight angular monomials. Their elementary
square identities determine the angular positivity. The remaining
argument keeps every state direction and rescales its coefficients
until a positive real Gram matrix controls the full polynomial.

\medskip
Chapters 2--6 prove the physical reduction and the algebra used in
that construction. Appendix A treats bounded source coefficients.
Appendices B and C organize the singular contacts by their source
scales, with the local identities and complementary measurement
charts proved together.

\medskip
Chapter 7 combines these certificates and proves the universal
sum-of-squares statement that gives finite semidefinite
representability. Every singular case remains part of this volume.
\end{minipage}
\end{titlepage}
\tableofcontents
\mainmatter
\chapter{Finite semidefinite representation}
\label{ch:theorem}

Consider two separated quantum systems, with two binary measurements
on each. Their behaviour is specified by the eight expectations
\[
 q=(a_0,a_1,b_0,b_1,E_{00},E_{01},E_{10},E_{11}),\qquad
 p(\alpha,\beta\mid x,y)=
 \frac{1+\alpha a_x+\beta b_y+\alpha\beta E_{xy}}4.
\]
Let $\mathcal Q_{222}$ contain all such behaviours obtained from
finite-dimensional quantum states and binary positive-operator-valued
measurements. Local dimensions are unrestricted.

\begin{theorem}\label{thm:main}
There are finite integers $N,m$ and real symmetric matrices
$A_0,A_1,\ldots,A_8,B_1,\ldots,B_m$ of order $N$ such that
\[
 \mathcal Q_{222}=\left\{q\in\mathbb R^8:\ \exists y\in\mathbb R^m,\quad
 A_0+\sum_{\nu=1}^8q_\nu A_\nu+\sum_{j=1}^my_jB_j\succeq0\right\}.
\]
\end{theorem}

Thus every Bell functional in this scenario can be optimized over
one fixed finite semidefinite representation. We prove existence by
certifying all supporting inequalities. The certificate must remain
valid when the inequality and the quantum realization are perturbed
independently. This independence is expressed by a tensor product.

\begin{theorem}[Separated sums of squares]\label{thm:slack}
Let $K_c,K_t$ be labelled real closed extensions of $\mathbb R$.
If $f=f_0+\sum_{\nu=1}^8f_\nu q_\nu$ is nonnegative on
$(\mathcal Q_{222})_{K_c}$ and $q\in(\mathcal Q_{222})_{K_t}$, then
\[
 [f,q]:=f_0\otimes1+\sum_{\nu=1}^8f_\nu\otimes q_\nu
       \in\Sigma(K_c\otimes_{\mathbb R}K_t)^2.
\]
Here $\Sigma T^2$ denotes finite sums of squares in $T$.
The two fields may be labelled copies of the same extension.
\end{theorem}

The proof proceeds from the physical model to the supporting
inequalities, and then to their singular contacts.

Chapter~\ref{ch:model} first reduces the quantum set to
$\operatorname{conv}\Phi(S^1\times S^1\times S^3)$, where
$\Phi$ is the polynomial expectation map for two real qubits.
At most nine such points are needed in a mixture. This gives a
compact convex semialgebraic set and a finite formula for its
extension to any real closed field.

Chapter~\ref{ch:geometry} then reduces a valid inequality to
nonnegative combinations of events and canonical supporting
inequalities. Two closed families have explicit small matrix
representations. Every other extreme inequality has a stationary
zero and, after relabelling, one of two event supports:
\[
\begin{array}{ll}
\mathsf C_1:&\lambda=n=0,\quad d,m,C,A,\alpha>0,
                       \quad D,\beta\ge0,\quad D+\beta>0,\\
\mathsf C_2:&\lambda=\beta=0,\quad d,m,C,A,D,\alpha>0,
                       \quad n\ge0.
\end{array}
\]
Chapter~\ref{ch:entry} derives the signs and rational coordinates
of these contacts from the original stationary equations.

It remains to certify the corresponding quadratic form in the
state amplitudes. Chapter~\ref{ch:tensor} provides the common
algebraic step: a matrix with bounded tensor coefficients and
positive-definite real residue gives an exact sum of squares.
The local calculations express the full quadratic form in roots
for which this criterion applies. Chapter~\ref{ch:cover} gives
an exhaustive source partition; the subsequent proofs establish
its identities at every boundary. Vanishing event weights and
commuting measurement limits remain part of the argument.

Finally, finite mixtures preserve the certificates. If
$f=\sum_j a_jf_j$, $q=\sum_iw_iq_i$, $a_j,w_i\ge0$ and
$\sum_iw_i=1$, and if $[f_j,q_i]=\sum_k h_{ijk}^2$, then
\[
 [f,q]=\sum_{i,j,k}
       \bigl[(\sqrt{a_j}\otimes\sqrt{w_i})h_{ijk}\bigr]^2.
 \tag{M2}
\]
The normalization of the testing weights preserves the constant
term of $f$. Theorem~\ref{thm:slack} follows. Applying
Scheiderer's tensor-evaluation criterion
\cite[Corollary 3.19]{scheiderer} to two labelled copies of each
real closed field gives Theorem~\ref{thm:main}.

Throughout, source coefficients and contact coordinates belong
to $K_c$; testing measurements and states belong to $K_t$.
The symbols $\mathrm I,\mathrm U,\mathrm O$ distinguish positive
infinitesimals, positive bounded elements with bounded inverse,
and positive unbounded elements. Formula tags are local to each
section. Every sum and algebraic operation in a certificate is finite.

\chapter{The finite real-qubit model and exact normalization}
\label{ch:model}
\label{component:p9b00698ff36a}
\label{p9b00698ff36a:sec:fresh-r21-conditional-assembly}
\section{An exact finite-dimensional physical model}

Let $Q\subset\mathbb R^8$ be the set of correlations obtained from arbitrary
finite-dimensional complex bipartite tensor-product states and two binary
measurements per party. The coordinates are
\[
 (a_0,a_1,b_0,b_1,E_{00},E_{01},E_{10},E_{11}).
 \tag{RFA1}
\]
A binary measurement is represented by an effect $0\preceq E\preceq I$,
with observable $2E-I$. Put
$Z=\operatorname{diag}(1,-1)$ and
$X=\left(\begin{smallmatrix}0&1\\1&0\end{smallmatrix}\right)$.
For $(c_A,s_A),(c_B,s_B)\in S^1$ and $\psi\in S^3\subset\mathbb R^4$,
let $\Phi$ be the eight expectations of
\[
 A_0=B_0=Z,\qquad A_1=c_AZ+s_AX,\qquad B_1=c_BZ+s_BX
 \tag{RFA2}
\]
in the state $\psi$. Every coordinate of $\Phi$ is polynomial in these
eight real parameters.

\begin{theorem}[Real-qubit representation]
With $\mathcal Z=S^1\times S^1\times S^3$,
\[
 Q=\operatorname{conv}\Phi(\mathcal Z)
   =\left\{\sum_{i=1}^{9}w_i\Phi(z_i):
       z_i\in\mathcal Z,\ w_i\geq0,\ \sum_iw_i=1\right\}.
 \tag{RFA3}
\]
In particular $Q$ is compact, convex and semialgebraic over the real field.
\end{theorem}

\begin{proof}
First, both effects of a party can be dilated with the same embedding.
For an effect $E$ on $\mathcal H$ define, on $\mathcal H\oplus\mathcal H$,
\[
 \widehat P_E=
 \begin{pmatrix}E&\sqrt{E(I-E)}\\
                 \sqrt{E(I-E)}&I-E\end{pmatrix},
 \qquad Vh=(h,0).
 \tag{RFA4}
\]
The square root commutes with $E$ by its spectral construction. Direct
block multiplication gives $\widehat P_E^2=\widehat P_E$ and
$V^*\widehat P_EV=E$. Thus for both measurement settings the same $V$
compresses the two projections to the original effects. Applying this
construction to each party and replacing $\rho$ by
$(V_A\otimes V_B)\rho(V_A\otimes V_B)^*$ preserves every marginal and
joint expectation. We may therefore assume the two observables of each
party are Hermitian involutions.

Here is the needed two-dimensional block decomposition. If $U,V$ are
Hermitian involutions, $C=UV$ is unitary and $UCU=C^*$. For an eigenvector
$e$ of $C$ with nonreal eigenvalue $\lambda$, $Ue$ has eigenvalue
$\overline\lambda$. The vectors are orthogonal and their span is invariant
under $U$ and $V=UC$. Its orthogonal complement is invariant too. Repeating
on orthogonal eigenvectors gives two-dimensional blocks. On the real
eigenvalue spaces $\lambda=\pm1$ one has $V=\pm U$; diagonalizing $U$ gives
one-dimensional blocks. This proves the decomposition in finite dimension
without an assumption on multiplicities.

On a two-dimensional block $U$ has eigenvalues $1,-1$. In its eigenbasis,
the other nontrivial involution has the form
$\left(\begin{smallmatrix}c&\overline z\\z&-c\end{smallmatrix}\right)$,
where $c$ is real and $c^2+|z|^2=1$. A diagonal change of phase leaves
$U=Z$ and makes $z$ real. Thus the pair has the form $Z,cZ+sX$.
A one-dimensional block with outcomes $a_0,a_1\in\{\pm1\}$ embeds into
this form: use the $a_0$ eigenvector of $Z$ and take
$c=a_0a_1,s=0$. The unused basis vector does not affect the experiment.

Apply both local decompositions. All measurement operators are block
diagonal, so their expectations depend only on the diagonal blocks of
the state. The traces of these positive blocks are nonnegative weights
summing to one. Every nonzero block, after normalization and the harmless
one-dimensional embedding, is a two-qubit state with real measurement
matrices. Decompose that state into finitely many complex pure states.
For a real symmetric matrix $M$ and a vector $r+is$,
\[
 (r+is)^*M(r+is)=r^TMr+s^TMs,\qquad
 \|r+is\|^2=\|r\|^2+\|s\|^2.
 \tag{RFA5}
\]
The cross terms cancel by symmetry. The nonzero real and imaginary parts,
normalized separately, are real pure states with their squared norms as
weights. This proves $Q\subseteq\operatorname{conv}\Phi(\mathcal Z)$.

Conversely each atom is a physical two-qubit experiment. A finite convex
combination is realized by a finite direct sum on each party and a state
with the corresponding classical diagonal weights in matching summands.
Therefore the reverse inclusion holds.

Any convex combination of more than nine points in $\mathbb R^8$ can be
shortened. The augmented vectors $(1,q_i)\in\mathbb R^9$ are linearly
dependent, so there are $\alpha_i$, not all zero, with
$\sum_i\alpha_i=0$ and $\sum_i\alpha_iq_i=0$. Change their common sign if
needed so that some $\alpha_i>0$, and replace the positive weights $w_i$
by $w_i-\tau\alpha_i$, where
$\tau=\min_{\alpha_i>0}w_i/\alpha_i$. They remain nonnegative, keep the
same sum and point, and at least one vanishes. Iterate. Zero weights
allow a fixed list of nine atoms. The right side of (RFA3) is the
continuous image of the compact product of a simplex and nine copies of
$\mathcal Z$, hence is compact. It is semialgebraic by the projection
theorem for real closed fields applied to its displayed finite polynomial
formula. Convexity has already been proved.
\end{proof}

For any real closed extension $K\supset\mathbb R$, define $Q_K$ by
(RFA3) with all variables in $K$. Define
$P=\{f\in\mathbb R[q]_{\leq1}:f\geq0\text{ on }Q\}$.
The defining formula
$\forall q\,(q\in Q\Rightarrow f(q)\geq0)$, together with quantifier
elimination and elementary transfer, gives exactly
\[
 P_K=\{f\in K[q]_{\leq1}:f\geq0\text{ on }Q_K\}.
 \tag{RFA6}
\]
Compactness was used only over $\mathbb R$. No compactness assertion for
the order topology of a non-Archimedean field is needed.

\section{Separated fields and the full homogeneous parent}

Fix arbitrary real closed extensions $K_c,K_t\supset\mathbb R$ and put
$T=K_c\otimes_{\mathbb R}K_t$. The factors are labelled; they may be two
copies of the same extension. There is no assumption of algebraic
disjointness and no chosen ordering on $T$. For
$f=f_0+\sum_{\nu=1}^8f_\nu q_\nu$ define
\[
 \mathscr S(f,q)=f_0\otimes1+\sum_{\nu=1}^8f_\nu\otimes q_\nu.
 \tag{RFA7}
\]

The event and CHSH convention is
\[
 S=E_{00}-E_{01}+E_{10}+E_{11},\quad H=1-S/2,\quad
 f=dH+4\sum_{\xi}\xi\,p_\xi,
 \quad \xi\in\{\lambda,n,m,C,A,D,\alpha,\beta\}.
 \tag{RFA8}
\]
The homogeneous parent is explicit.

Choose projective angle pairs $(a_0,a_1)$ and $(b_0,b_1)$ in $K_t^2$,
each nonzero, and set
$\Delta_A=a_0^2+a_1^2>0$, $\Delta_B=b_0^2+b_1^2>0$.
The corresponding circle coordinates are
\[
 c_A=(a_0^2-a_1^2)/\Delta_A,\quad s_A=2a_0a_1/\Delta_A,
 \tag{RFA9}
\]
and similarly for Bob. These cover every circle point over $K_t$:
if $c\ne-1$ use $(1+c,s)$, and if $c=-1$ use $(0,1)$.
No affine half-angle chart is therefore needed at infinity.
For $\Psi=(W,X,Y,Z)$ put
\[
 \begin{split}
 L_m&=a_0b_1W-a_0b_0X+a_1b_1Y-a_1b_0Z,\\
 L_D&=a_1b_0W+a_1b_1X-a_0b_0Y-a_0b_1Z.
 \end{split}
\]
Products of a coefficient-field element with a testing-field expression
in the next identity mean their separated tensor product. The complete
homogeneous parent is
\[
 \begin{split}
 \mathcal F={}&4\lambda\Delta_A(b_0W+b_1X)^2
       +4n\Delta_A\Delta_BX^2+4mL_m^2\\
 &+4C\Delta_B(a_1W-a_0Y)^2+4A\Delta_A\Delta_BY^2\\
 &+(4D+2d)L_D^2
       +(4\alpha+2d)\Delta_A(b_1Y-b_0Z)^2\\
 &+(4\beta+2d)\Delta_B(a_0X+a_1Z)^2
       -2d\Delta_A\Delta_BZ^2.
 \end{split}                                                    \tag{RFA10}
\]
For a nonzero state let $N=\|\Psi\|^2>0$ and let $z$ be the atom
with the circle coordinates (RFA9) and unit state $\Psi/\sqrt N$.
Then
\[
 \mathcal F=(1\otimes\Delta_A\Delta_BN)\,
                 \mathscr S(f,\Phi(z)).
 \tag{RFA11}
\]

To check the normalization, first take $a_0=b_0=1$, $u=a_1$ and
$v=b_1$. The plus/minus measurement eigenvectors are
$(1,u)/\sqrt{1+u^2}$ and $(u,-1)/\sqrt{1+u^2}$.
Their product amplitudes give, in the order of (RFA8),
\[
 \frac{(W+vX)^2}{1+v^2},\quad X^2,\quad
 \frac{(vW-X+uvY-uZ)^2}{(1+u^2)(1+v^2)},\quad
 \frac{(uW-Y)^2}{1+u^2},\quad Y^2,
\]
\[
 \frac{(uW+uvX-Y-vZ)^2}{(1+u^2)(1+v^2)},\quad
 \frac{(vY-Z)^2}{1+v^2},\quad
 \frac{(X+uZ)^2}{1+u^2}.
 \tag{RFA12}
\]
For a normalized state these are exactly the eight selected event
probabilities. For a general state they are multiplied by $N$.
Substitution of the matrices (RFA2) into $H$ gives its numerator
\[
 2\{(uW+uvX-Y-vZ)^2+(1+u^2)(vY-Z)^2
       +(1+v^2)(X+uZ)^2-(1+u^2)(1+v^2)Z^2\}.
\]
This proves (RFA11) in the affine chart. Homogenization gives (RFA10);
it is a polynomial identity and therefore also gives the zero and
infinite projective coordinates without dividing by $a_0$ or $b_0$.
The zero state has zero parent separately.

If $\mathcal F$ or $(s\otimes1)\mathcal F$ has an exact finite SOS
certificate, with $s\in K_c^{>0}$, the original slack is restored by
multiplying each root by
\[
 s^{-1/2}\otimes(\Delta_A\Delta_BN)^{-1/2}.
 \tag{RFA13}
\]
All these inverses are taken within one field before forming a tensor.
The identity does not permit cancellation of an arbitrary mixed tensor.
Additional positive source gauges are handled in the same way. A state
chart or a change of frame must first prove its own denominator nonzero.

\paragraph{Literal commuting tests.}
These require no limiting local certificate. If one party's two
involutions commute, on a simultaneous eigenspace the CHSH operator
$S=A_0(B_0-B_1)+A_1(B_0+B_1)$ is either $\pm2B_0$ or $\pm2B_1$.
Thus $H=I-S/2$ is positive semidefinite in the testing field. For
$d\geq0$, its contribution is a pure-source square times a
pure-testing square. Every nonnegative event packet is also separated
SOS. This covers the actual commuting angle axes and infinities for
every nonnegative canonical coefficient tuple. It does not cover merely
infinitesimal nonzero angles, which are covered by the canonical certificate theorem.


\chapter{Exact sums of squares in the separated tensor ring}
\label{ch:tensor}
\section{Two labelled fields and their residues}
Let $K_c,K_t$ be real closed extensions of $\mathbb R$ and put
\[
 T=K_c\otimes_{\mathbb R}K_t,\qquad
 B_i=\{x\in K_i: |x|<M\text{ for some }M\in\mathbb R_{>0}\},
 \quad i\in\{c,t\}.
\]
Write $\mathfrak m_i$ for the elements whose absolute value is smaller
than every positive real number. The order on $K_i$ is its unique field
order. No order, embedding into a field, or algebraic-disjointness
assumption is imposed on $T$.

Every $x\in B_i$ has a unique real standard part $\bar x$ with
$x-\bar x\in\mathfrak m_i$. Indeed, the real cut
$\{r\in\mathbb R:r\le x\}$ is nonempty and bounded; its supremum
has the claimed property. Two such real numbers differ by a real
infinitesimal and hence agree. Infinitesimals form an ideal of $B_i$,
so $x\mapsto\bar x$ is an $\mathbb R$-algebra homomorphism with
kernel $\mathfrak m_i$. In particular
$B_i=\mathbb R\oplus\mathfrak m_i$ as real vector spaces.

Since tensoring real vector spaces preserves injections,
$B:=B_c\otimes_{\mathbb R}B_t$ embeds into $T$. Its residue map is
\[
 \rho:B\longrightarrow\mathbb R,\qquad
 \rho\Bigl(\sum_i a_i\otimes b_i\Bigr)=\sum_i\bar a_i\bar b_i.
 \tag{TF1}
\]
The direct-sum decompositions above show that
\[
 \ker\rho=\mathfrak m_c\otimes B_t+B_c\otimes\mathfrak m_t.
 \tag{TF2}
\]
Thus every tensor of zero residue has a finite expression in which
each elementary term has at least one infinitesimal factor. None of
these observations requires the two field extensions to be disjoint.

\section{A strict real reserve gives an exact tensor sum of squares}
\label{component:p25af2c58da7e}
\label{p25af2c58da7e:sec:mechanism}
\label{p25af2c58da7e:lem:residue}
\paragraph{Lemma TF1 (positive residue).}
If $a\in B$ and $\rho(a)>0$, then $a$ is a finite sum of squares in
$T$. In fact its square roots may all be chosen to be elementary
tensors.

\begin{proof}
By (TF2), write
\[
 a=\varepsilon+\sum_{j=1}^N x_j\otimes y_j,\qquad
 \varepsilon=\rho(a)>0,
 \tag{TF3}
\]
with bounded pure factors and an infinitesimal in at least one factor
of each term. If $N=0$, the assertion is immediate. Otherwise choose
positive real numbers $\varepsilon_j$ summing to $\varepsilon$.
For a term with $x_j$ infinitesimal, choose a positive real
$M_j>|y_j|$. Direct multiplication gives
\[
 \varepsilon_j+x_j\otimes y_j
 =\frac12\left[
  (\varepsilon_j/M_j+x_j)\otimes(M_j+y_j)
 +(\varepsilon_j/M_j-x_j)\otimes(M_j-y_j)\right].
 \tag{TF4}
\]
All four pure factors are strictly positive in their respective
fields. Each elementary summand therefore is a square: for example
$\tfrac12u\otimes v=(\sqrt{u/2}\otimes\sqrt v)^2$.
If instead $y_j$ is infinitesimal, interchange the roles of the
two factors in this calculation. Sum (TF4) over $j$. Every operation
is a field operation in one labelled factor; no mixed radical or
mixed inverse occurs.
\end{proof}

\label{p25af2c58da7e:lem:gram}
\paragraph{Lemma TF2 (strict bounded Gram).}
Let $Q\in\operatorname{Sym}_r(B)$ and suppose its entrywise real
residue $\rho(Q)$ is positive definite. Then
\[
                 g^{\mathsf T}Qg\in\Sigma T^2
                 \qquad\text{for every }g\in T^r.
 \tag{TF5}
\]
The entries of $g$ need not be bounded, independent, or polynomial.

\begin{proof}
Choose a real invertible matrix $C$ with
$C^{\mathsf T}\rho(Q)C=I$. Replacing $Q$ by
$C^{\mathsf T}QC$ and $g$ by $C^{-1}g$ preserves the quadratic
expression, so assume $\rho(Q)=I$. Choose a real $\delta>0$ such
that $(r-1)\delta<1$. With $e_i$ the real coordinate vectors,
\[
\begin{split}
 Q={}&\sum_i\bigl(Q_{ii}-(r-1)\delta\bigr)e_ie_i^{\mathsf T}\\
 &+\sum_{i<j}\frac{\delta+Q_{ij}}2
                (e_i+e_j)(e_i+e_j)^{\mathsf T}\\
 &+\sum_{i<j}\frac{\delta-Q_{ij}}2
                (e_i-e_j)(e_i-e_j)^{\mathsf T}.
\end{split}
 \tag{TF6}
\]
The off-diagonal entries agree because the last two contributions
have difference $Q_{ij}$; their diagonal contributions sum to
$\delta$ for every incident pair. The residues of all scalar
coefficients in (TF6) are positive. Lemma TF1 expresses each
coefficient as a finite sum of squares in $T$. Evaluation on $g$
then expresses (TF5) as a finite sum of squares of those scalar roots
times $g_i$ or $g_i\pm g_j$. This is a polynomial identity in the
commutative ring $T$, so no size condition on $g$ is involved.
\end{proof}

\section{Legal normalization and finite assembly}
\paragraph{Lemma TF3 (pure-factor cancellation).}
For $s\in K_c^{>0}$, $t\in K_t^{>0}$, and $F\in T$,
\[
 F\in\Sigma T^2\quad\Longleftrightarrow\quad
                  (s\otimes t)F\in\Sigma T^2.
 \tag{TF8}
\]
Indeed multiply each root by $\sqrt s\otimes\sqrt t$ in one
direction and by $s^{-1/2}\otimes t^{-1/2}$ in the other.
These are inverse units of $T$. Nonvanishing must be proved before
this lemma is used. A homogeneous coordinate that is literally zero
cannot be cancelled, and the lemma does not authorize cancellation
of a general mixed tensor.

\paragraph{Lemma TF4 (positive matrix pairing).}
If $X\in\operatorname{Sym}_r(K_c)$ and
$Y\in\operatorname{Sym}_r(K_t)$ are positive semidefinite in their
own fields, then
\[
 \sum_{i,j}X_{ij}\otimes Y_{ij}
 =\sum_{a,b}\left(\sum_i U_{ia}\otimes V_{ib}\right)^2,
 \qquad X=UU^{\mathsf T},\quad Y=VV^{\mathsf T}.
 \tag{TF9}
\]
To obtain the factorizations over a real closed field, pivot on a
positive diagonal entry and complete its square. The remaining
Schur complement is positive semidefinite. If every diagonal entry
is zero, every off-diagonal entry is zero by its two-dimensional
quadratic restriction. Induction, and positive square roots of the
pivots, finish the factorization even for singular matrices.
Expansion of the right side proves (TF9).

\paragraph{Lemma TF5 (conic and convex closure).}
For $f=f_0+\sum_{\nu=1}^n f_\nu q_\nu$ set
$[f,q]=f_0\otimes1+\sum_\nu f_\nu\otimes q_\nu$.
Suppose in the source field
$f=\sum_j a_j f_j$ with $a_j\ge0$, and in the testing field
$q=\sum_i w_i q_i$ with $w_i\ge0$ and $\sum_iw_i=1$.
If $[f_j,q_i]=\sum_k h_{ijk}^2$ in $T$, then
\[
             [f,q]=\sum_{i,j,k}
                \left[(\sqrt{a_j}\otimes\sqrt{w_i})h_{ijk}\right]^2.
 \tag{TF10}
\]
This follows from linearity in $f$ and affineness in $q$; the equation
$\sum_iw_i=1$ preserves the constant coefficient. All sums are finite.
Zero weights give zero roots and are never inverted.

\section{From separated squares to a finite representation}
Let $S\subset\mathbb R^n$ be closed, convex, and semialgebraic, and
let $P_S$ be its cone of nonnegative affine functions. Fix a finite
real formula defining $S$. Over any real closed extension $K$, its
transfer $S_K$ is the set defined by that same formula. The formula
\[
          \forall q\ (q\in S\ \Longrightarrow\ f(q)\ge0)
 \tag{TF11}
\]
shows, by quantifier elimination and elementary transfer, that
$(P_S)_K$ consists exactly of the $K$-affine functions nonnegative
on $S_K$. This is a statement about formulas; it makes no claim
that $S_K$ is compact in the order topology of $K$.

Scheiderer's Corollary 3.19~\cite{scheiderer} says that $S$ has a finite semidefinite
lift if and only if, for every real closed extension $K/\mathbb R$,
every $f\in(P_S)_K$, and every $q\in S_K$, the point-first tensor
evaluation is a finite sum of squares in $K\otimes_{\mathbb R}K$.
The coefficient-first convention used here is equivalent by the
factor-flip automorphism, which preserves squares.
Therefore a proof of
\[
       [f,q]\in\Sigma(K\otimes_{\mathbb R}K)^2
       \quad\text{for every }K,\ f\in(P_S)_K,\ q\in S_K
 \tag{TF12}
\]
is sufficient. Two labelled copies of each one arbitrary $K$ suffice;
a theorem for all unequal pairs $K_c,K_t$ is stronger than necessary.
The labels must still remain separate in every identity. In particular,
algebraic disjointness of the copies must not be imposed.

The number and tensor ranks of the roots in (TF12) may depend on
$K,f,q$. Corollary 3.19 converts this universal assertion into one
finite semidefinite representation by a finite-dimensional algebraic
construction and compactness. This is the implication used here.

\chapter{The source cone: two closed faces and two contact families}
\label{ch:geometry}
\label{component:p0c8f420148b7}
\label{p0c8f420148b7:sec:physical}
\label{p0c8f420148b7:eq:parent}
\label{p0c8f420148b7:prop:compat}
\label{component:pd9a6195f0af3}
\label{pd9a6195f0af3:sec:closed-faces-short}
\label{pd9a6195f0af3:lem:closed-face-orbits}
\label{pd9a6195f0af3:lem:quadratic-matrix-lift}
\label{pd9a6195f0af3:thm:closed-face-adjacent}
\label{pd9a6195f0af3:thm:closed-face-opposite}
\label{pd9a6195f0af3:cor:all-closed-faces}
\section{The exact physical polynomial}

Use the affine functional $H=1-(E_{00}-E_{01}+E_{10}+E_{11})/2$.
Order the eight losing events as
\[
 (\lambda,n,m,C,A,D,\alpha,\beta)
 =\bigl((++|01),(+-|00),(+-|11),(-+|10),
             (-+|00),(-+|11),(--|01),(+-|10)\bigr).
\]
Every event vanishes at the chosen PR point.  Each is one at its own
deterministic CHSH-facet neighbour and zero at the other seven.
Consequently every affine functional has the unique expansion
\[
 f=dH+4\sum_i a_i p_i,\qquad
 d=-f(\mathrm{PR}),\qquad a_i=f(v_i)/4.                 \tag{SG1}
\]
For a quantum-valid functional negative at that PR point, $d>0$ and
$a_i\geq0$.  If a quantum-valid functional is nonnegative at every PR
point, it is nonnegative on the entire nonsignalling polytope and is a
nonnegative combination of its probability inequalities.  The latter
sources have immediate separated square certificates.  Here is a complete short vertex argument. The nonsignalling affine
space has dimension eight. If no local marginal is deterministic, each
two-by-two outcome table has positive row and column sums, hence at
most two zero entries. A vertex requires at least eight active
probability inequalities. Thus every table has exactly two zeros,
necessarily in opposite corners. Write its perfect correlation or
anticorrelation as $a_x=s_{xy}b_y$, with $s_{xy}=\pm1$.
If $\prod_{x,y}s_{xy}=-1$, these equations force all four marginals
to vanish and give a PR vertex. If the product is $+1$, one marginal
remains free and the table is a mixture of two deterministic points,
so an interior marginal value is not a vertex. If instead one
marginal is deterministic, say $A_0$, a local model is obtained by
sampling $A_1$ and then sampling $B_0,B_1$ independently conditional
on $A_1$, using the two given joint tables. Nonsignalling makes their
$A_1$ marginals agree; the tables involving fixed $A_0$ are reproduced
automatically. Such a vertex is therefore deterministic. This proves
that the only vertices are the sixteen deterministic points and the
eight PR points. Linear separation (equivalently finite-dimensional
Farkas duality) now gives the asserted event representation.

Put $D_u=1+u^2$, $D_v=1+v^2$, and write the real state as
$\psi=(W,X,Y,Z)$.  The entire cleared parent is
\[
\begin{split}
 F={}&4\lambda D_u(W+vX)^2+4nD_uD_vX^2
       +4m(vW-X+uvY-uZ)^2\\
 &+4CD_v(uW-Y)^2+4AD_uD_vY^2
 +(4D+2d)(uW+uvX-Y-vZ)^2\\
 &+(4\alpha+2d)D_u(vY-Z)^2
 +(4\beta+2d)D_v(X+uZ)^2-2dD_uD_vZ^2.
\end{split}                                                   \tag{SG2}
\]
It is $D_uD_v\|\psi\|^2f(q)$ for a nonzero unnormalised state.
All these factors are pure testing factors.  Write $P(u,v)$ for its
symmetric state matrix.  If $d>0$ and $uv\ne0$, its $Z=0$ principal
block is positive definite, since the signed part alone is
\[
 2d\{(uW+uvX-Y)^2+D_uv^2Y^2+D_vX^2\}.                 \tag{SG3}
\]
Thus $P$ has at least three positive eigenvalues, its kernel at a
zero is one-dimensional with $Z\ne0$, and
\[
                 P\succeq0\quad\Longleftrightarrow\quad\det P\geq0.
                                                               \tag{SG4}
\]

Here is the entire determinant datum needed for the geometric
reduction.  Set
\[
\begin{array}{ll}
t_1=d+2(\lambda+n+D),&t_2=d+2(\alpha+m+A),\\
t_3=d+2(\lambda+C+D),&t_4=d+2(\alpha+\beta+m),\\
t_5=d+2(\lambda+\beta+n),&t_6=d+2(\alpha+C+A),\\
t_7=d+2(\beta+n+m),&t_8=d+2(C+A+D).
\end{array}
\]
All $t_i>0$.  Directly expanding the four-dimensional determinant of
SG2, or applying Cauchy--Binet to its eight positive and one negative
rank-one packets, gives
\[
 \det P=16D_u^2D_v^2u^2v^2L(u^2,v^2),\qquad
 L(U,V)=h+\sum_{(i,j)\ne(1,1)}q_{ij}U^{i-1}V^{j-1},       \tag{SG5}
\]
where
\[
\begin{array}{ll}
q_{00}=4\lambda\alpha t_7t_8,&
q_{01}=4(\lambda Dt_7t_6+\alpha mt_5t_8),\\
q_{02}=4mDt_5t_6,&
q_{10}=4(\alpha At_7t_3+\lambda nt_8t_4),\\
q_{12}=4(\beta mt_1t_6+CDt_5t_2),&
q_{20}=4nAt_3t_4,\\
q_{21}=4(\beta nt_3t_2+CAt_1t_4),&
q_{22}=4\beta Ct_1t_2.
\end{array}                                                   \tag{SG6}
\]
The central coefficient is fixed without a long expansion by
$h=\det P(1,1)/256-\sum_{(i,j)\ne(1,1)}q_{ij}$.
The eight nonnegative outer coefficients determine the possible
contact geometries.

\section{The twelve elementary angular identities}
\label{sec:angular-circuits}

The determinant of the state matrix has eight nonconstant Laurent
monomials. Their coefficients are nonnegative. A stationary zero
therefore has a particularly simple algebraic decomposition, which
also explains the cubic and quartic identities used below.

Let $E=\{-1,0,1\}^2\setminus\{(0,0)\}$ and write
\[
 L(U,V)=h+\sum_{e\in E}q_eU^{e_1}V^{e_2},\qquad q_e\ge0.
\]
Suppose $U_c,V_c>0$ satisfy $L=L_U=L_V=0$. Put
$w_e=q_eU_c^{e_1}V_c^{e_2}$. The value and derivative equations are
exactly
\[
 h=-\sum_e w_e,\qquad \sum_e w_e e=0.
 \label{eq:angular-balance}
\]
Thus the positive coefficients balance around the origin in the
exponent plane.

\begin{lemma}[Balanced Laurent polynomial]
Let $K_c,K_t$ be labelled real closed extensions of $\mathbb R$.
If the coefficients and contact above belong to $K_c$, then for every
$U,V>0$ in $K_t$ the separated evaluation of $L$ is a finite sum of
squares in $K_c\otimes_{\mathbb R}K_t$.
\end{lemma}
\begin{proof}
First decompose the coefficient vector $w$ using its balance equation.
If $w\ne0$, the origin lies in the convex hull of its active exponents.
A minimal subset containing the origin has two or three elements.
For the eight exponents in $E$, the possibilities, up to rotation
and reflection of the square, are
\[
\begin{array}{c|c|c}
\text{exponents}&\text{positive integer weights}&\text{number}\\ \hline
\{e,-e\}&(1,1)&4\\
\{(1,1),(-1,0),(0,-1)\}&(1,1,1)&4\\
\{(1,1),(-1,1),(0,-1)\}&(1,1,2)&4.
\end{array}
\]
Indeed a two-element set must be opposite. A triangle containing the
origin contains either one corner and its two opposite axial points,
or two adjacent corners and their opposite axial point. These are
the last two rows. For its integer weight vector $b$, subtract
$\lambda b$, where $\lambda=\min_{b_e>0}w_e/b_e$.
The remainder is nonnegative and balanced, and has fewer active
entries. Repeating at most eight times gives
$w=\sum_j\lambda_jb^{(j)}$ with $\lambda_j>0$ in $K_c$.

Set $X=U_c^{-1}\otimes U$ and $Y=V_c^{-1}\otimes V$.
Their rational powers mean separate positive roots in the two fields.
For each row of the table the monomials $p_i=X^{e_{i1}}Y^{e_{i2}}$
satisfy $\prod_i p_i^{b_i}=1$. The required elementary identities are
\begin{align*}
 p+p^{-1}-2&=(\sqrt p-p^{-1/2})^2,\\
 a^3+b^3+c^3-3abc
   &=\tfrac12(a+b+c)\bigl[(a-b)^2+(b-c)^2+(c-a)^2\bigr],\\
 a^4+b^4+c^4+d^4-4abcd
   &=(a^2-b^2)^2+(c^2-d^2)^2+2(ab-cd)^2.
\end{align*}
For weights $(1,1,1)$ take $a=p_1^{1/3}$, $b=p_2^{1/3}$,
$c=p_3^{1/3}$; then $abc=1$. Each of $a,b,c$ is an elementary
tensor square, so the middle line is a sum of tensor squares.
For weights $(1,1,2)$ take $a=p_1^{1/4}$, $b=p_2^{1/4}$ and
$c=d=p_3^{1/4}$; then $abcd=1$.
Finally multiply the squares by $\sqrt{\lambda_j}\otimes1$ and sum.
Equation~\eqref{eq:angular-balance} makes the result exactly $L$.
\end{proof}

In a single real closed field this proves that a stationary zero of
$L$ is a global minimum, using only the three displayed identities.
For the Bell state matrix, its positive three-dimensional principal
block then turns determinant nonnegativity into matrix positivity.
The same calculation proves a separated sum-of-squares identity for
the determinant at independent tests. To obtain a certificate for the
whole state polynomial we must also control the state-normal terms;
the following matrix lemmas perform that step.


\section{Stationary zeros and physicality}

\paragraph{Contact lemma.}
Suppose $d>0$, all event weights are nonnegative, and there exist
positive $u_c,v_c$ and a nonzero real state $\psi_c$ such that
\[
 P(u_c,v_c)\psi_c=0,
 \qquad\psi_c^TP_u(u_c,v_c)\psi_c=0,
 \qquad\psi_c^TP_v(u_c,v_c)\psi_c=0.                    \tag{SG7}
\]
Then the source is physical at every angle and every state.  This
statement holds over every real closed field.

\emph{Proof.}
SG3 gives three positive complementary state eigenvalues. The adjugate
at contact is a positive multiple of $\psi_c\psi_c^T$, so SG7 and SG5
give $L=L_U=L_V=0$ there. The balanced Laurent identity just proved
makes $L$ nonnegative at every positive pair in the same field.
Equation SG4 therefore gives physicality for $uv\ne0$.
At a commuting or projective endpoint, simultaneous diagonalization
of the commuting local pair reduces the CHSH operator to $\pm2$
times a remote reflection, so $H\succeq0$. The event terms are
nonnegative there as well. Homogenizing SG2 includes all endpoints.
This proves the claim over every real closed field.\hfill$\square$

This lemma is also the justification for deleting a positive event
whose amplitude vanishes at a contact.  Its square and all first
derivatives vanish; deleting it preserves SG7.  It is not legitimate
to delete an event whose amplitude is nonzero.

\section{The geometry of a stationary zero}

\paragraph{Coercivity at every physical finite contact.}
Every physical source with $d>0$ and a finite noncommuting zero has
coercive Newton support.  Indeed, put
$a_e=q_eU_c^{e_1}V_c^{e_2}>0$ for the active exponent vectors.
The two contact derivatives give
\[
                         \sum_e a_e e=0.                \tag{SG17}
\]
If zero is not in the interior of their convex hull, a separating
linear functional nonnegative on every active vector must vanish
on all of them, by SG17.  Thus the active support is empty or lies
on one of the four lines through the origin in the square grid.
A nonempty one-sided subset cannot satisfy SG17.  The horizontal
opposite pair $q_{01},q_{21}>0$ has four possible product pairings
in SG6.  These respectively force $\lambda n$, $CD$, $\beta m$
or $\alpha A$ positive, activating an off-line coefficient
$q_{10}$ or $q_{12}$.  The vertical pair similarly forces
$\alpha m$, $CA$, $\beta n$ or $\lambda D$ and activates
$q_{01}$ or $q_{21}$.

For the diagonal pair $q_{00},q_{22}>0$, all off-line coefficients
can vanish only on the alternating matching
$\lambda,\alpha,\beta,C>0$, $n=m=A=D=0$.
For the other diagonal they can vanish only on the opposite matching
$n,m,A,D>0$, $\lambda=C=\alpha=\beta=0$.
The latter is not physical: SG2 at $u=v=1$ and
$\psi=(1,-\varepsilon,\varepsilon,1)$ gives
\[
 F=16\varepsilon\{(d+n+m+A+D)\varepsilon-d\}<0
 \quad\text{for }\quad
 \varepsilon=\frac{d}{2(d+n+m+A+D)}.
\]
The first matching is its physical one-step cycle rotation.
If every outer coefficient is zero, SG6 forbids every adjacent or
opposite positive pair in the cycle
$(\lambda,D,C,A,\alpha,m,\beta,n)$.
Any maximal remaining set is a triple with cyclic gaps $(2,3,3)$;
smaller sets extend to such a triple.  The eight triples are rotations
of $\{\lambda,m,A\}$.  At $u=v=1$, $\psi=(1,-1,0,2)$
annihilates precisely the needed three event rows and has signed-core
value $-4$, giving $F=-4d<0$.  This excludes empty support too.

For completeness, the physical cycle symmetries used here are the
CHSH-preserving maps on event coordinates $(x,y,a,b)$
\[
\begin{split}
g_A:(x,y,a,b)&\mapsto(1-x,y,a,(-1)^yb),\\
g_B:(x,y,a,b)&\mapsto(x,1-y,(-1)^{1-x}a,b),\\
g_T:(x,y,a,b)&\mapsto(y,x,(-1)^yb,(-1)^xa).
\end{split}
\]
In the displayed cycle, indexed from zero, they act as
$i\mapsto5-i,3-i,2-i$ modulo eight.  Their compositions give
a one-step rotation and a reflection.  The assertions follow
by inserting the eight event definitions in SG1.

\paragraph{Strict single-field contact Hessian.}
The logarithmic angular Hessian is now
\[
                 \mathcal H_L=\sum_e a_e ee^T\succ0,       \tag{SG18}
\]
since the active vectors span the plane.  Write the state matrix
as $P=\left(\begin{smallmatrix}H_3&b\\b^T&c\end{smallmatrix}\right)$
in the coordinates $(W,X,Y;Z)$ and set $Z=1$.
The positive definite state block SG3 gives, within the single
source field,
\[
 F=(z+H_3^{-1}b)^TH_3(z+H_3^{-1}b)+s(u,v),
 \qquad s=\det P/\det H_3,\quad z=(W,X,Y)^T.             \tag{SG19}
\]
At contact $s=0$ and its gradient vanishes.  Its angular Hessian
is the positive scalar $16D_u^2D_v^2u^2v^2/\det H_3$ times
the congruence of SG18 by $\operatorname{diag}(2/u,2/v)$.
Consequently the full Hessian of $F$ in $(W,X,Y,u,v)$ has
positive state block $2H_3$ and a positive angular Schur complement;
it is positive definite.  All formulas for derivatives and
positive definiteness are algebraic.  The finite real assertion
therefore transfers to every real closed field; logarithms are
only notation for $U\partial_U,V\partial_V$.
The inverse in SG19 belongs to the source field.

\paragraph{Positive contact state orientation.}
Let $F_H$ denote SG2 with $d=1$ and all eight event weights zero.
Its state matrix satisfies
\[
              \det P_H=-16D_u^2D_v^2u^2v^2\ne0.          \tag{SG20}
\]
This follows by a four-dimensional determinant expansion of the
three signed-core squares and the negative $Z^2$ term in SG2.
If all positive-weight event rows vanished at a physical contact,
their state matrices would annihilate the contact state.  The
kernel equation would then give $P_H\psi_c=0$, contradicting
SG20.  Hence at least one event contributes positively.  Since
$F(\psi_c)=0$ and $d>0$, the core value $F_H(\psi_c)$ is strictly
negative.

For positive $u,v$, the $X,Y$ block of $F_H$ is
$2\left(\begin{smallmatrix}h_0&-uv\\-uv&h_0\end{smallmatrix}\right)$,
where $h_0=1+v^2+u^2v^2$.  Its Schur form is exactly SG15.
It is nonnegative when $WZ\le0$.  Thus the negative core value
forces
\[
                  W_cZ_c>0,\qquad Z_c=1\Longrightarrow W_c>0.
                                                               \tag{SG21}
\]

\paragraph{The forbidden-deletion implication.}
If a positive event row vanished at a contact, deleting that
coefficient would preserve the four kernel equations and both
angular derivatives.  SG7 would make the resulting nonnegative
source physical, with the same finite contact.  SG17--SG18
therefore make its support coercive.  We obtain the exact rule
\[
 \text{a positive contact event cannot vanish if deleting it
 makes the Newton support noncoercive}.                       \tag{SG22}
\]
In particular every deletion introducing an exclusive omitted pair
is forbidden by SG9.  This rule will fix the signs of the contact coordinates in
Chapter~\ref{ch:entry}.

\section{Reduction to two contact families}

The active exponent vectors in SG5 lie on the boundary of
$\{-1,0,1\}^2$.  The exponential sum is coercive exactly when zero
is in the interior of their convex hull.  One direction follows from
uniform positive support in every unit direction; the other follows
from a separating direction along which every active exponential is
nonincreasing.

For this eight-point configuration a separating direction can be
chosen from $(\pm1,0),(0,\pm1),(\pm1,\pm1)$.  For example the axis
condition in SG6 gives
$\lambda\alpha=\lambda D=\alpha m=mD=0$, whose minimal omitted
event pairs are $\{\lambda,m\}$ and $\{\alpha,D\}$.
For a diagonal the additional minimal pair is
$\{\lambda,\alpha\}$; the remaining minimal triples already
contain an axis pair.  Rotating the event cycle
\[
                 (\lambda,D,C,A,\alpha,m,\beta,n)          \tag{SG8}
\]
gives exactly the twelve pairs of exclusive events.  Therefore
\[
 \text{coercive exponent support}\quad\Longleftrightarrow\quad
 \text{every two omitted losing events are compatible}.       \tag{SG9}
\]
Compatibility means agreement on each shared local measurement.
Pairwise compatible events extend to a single deterministic
assignment.  Such an assignment loses one or three CHSH inputs;
there are consequently at most three compatible omissions.
In SG8 compatible pairs have cyclic distance one or two, and a
compatible triple is three consecutive vertices.

\paragraph{Extreme-source lemma.}
Normalise the canonical source by $d+\sum a_i=1$ and impose
physicality and nonnegativity.  This is a compact convex subset of
an eight-dimensional simplex.  Every extreme point with $d>0$ and
coercive support has a finite stationary zero and exactly five or
six positive event weights.

\emph{Proof.}
Coercivity gives a finite minimum $\mu\geq0$ of $L$.  If $\mu>0$,
take sufficiently small opposite nonradial perturbations inside the
exact active coefficient face, preserving the normalisation.  Every
positive $q_e$ changes by relative error at most $\varepsilon$,
every zero $q_e$ remains zero, and $h$ changes by absolute error at
most $\varepsilon$.  Thus, uniformly at all positive angles,
\[
 L'\geq(1-\varepsilon)L+\varepsilon h-\varepsilon
       \geq(1-\varepsilon)\mu+\varepsilon h-\varepsilon>0.
\]
The perturbations remain physical by SG4 and the direct endpoint
argument, contradicting extremality.  The perturbation exists:
SG9 permits at most three omissions, hence at least six active
coordinates including $d$.  Thus $\mu=0$ and there is a contact.

Fix its angles and state.  SG7 consists of six homogeneous linear
equations in the source coefficients.  If at least seven events
were positive, there would be at least eight active coordinates;
their solution space has dimension at least two and contains the
source.  Intersecting its directions with coefficient sum zero
leaves a nonzero direction.  Small opposite perturbations preserve
all signs and SG7, and are physical by the contact lemma.  This again
contradicts extremality.  Hence at most six events are positive.
The lower bound five follows from SG9. \hfill$\square$

The CHSH-preserving event relabellings generate the dihedral group
of SG8.  Thus the coercive extreme sources, including every allowed
third omission, reduce to the two patterns
\[
\begin{array}{ll}
\mathsf C_1:&\lambda=n=0,\quad d,m,C,A,\alpha>0,
                  \quad D,\beta\geq0,\quad D+\beta>0,\\
\mathsf C_2:&\lambda=\beta=0,\quad d,m,C,A,D,\alpha>0,
                  \quad n\geq0.
\end{array}                                                     \tag{SG10}
\]
The C1 third omission is $D=0$ or $\beta=0$; the C2 third omission
is $n=0$, already a C1 orbit.  This overlap is harmless.  It is not
legitimate to replace positive infinitesimal coefficients by zeros.

Every point of the normalised source body is a convex combination
of at most nine extreme points.  The fixed nine-term assertion,
the midpoint definition of extremality, SG7, and physicality all
have first-order polynomial formulations.  Their real proof
therefore transfers to a real closed coefficient field, without
using topological compactness or arbitrary infinite decompositions
in that field.  A nonnegative decomposition preserves each literal
zero coefficient term by term.

\section{Noncoercive sources have explicit lifts}

Every noncoercive source omits an exclusive pair.  Up to the same
finite symmetries there are only two possibilities:
$\lambda=m=0$ and $\lambda=\alpha=0$.

We first record the elementary quadratic matrix lemma.  A real
symmetric matrix polynomial $K(t)=K_0+tK_1+t^2K_2$ is PSD for every
real $t$ exactly when there is a skew-symmetric $S$ such that
\[
 \begin{pmatrix}K_0&K_1/2+S\\K_1/2-S&K_2\end{pmatrix}\succeq0.
                                                               \tag{SG11}
\]
For completeness, the cone of Gram-representable coefficient
triples is closed because the trace of a representing PSD Gram is
$\operatorname{tr}(K_0+K_2)$.  A nonnegative dual functional has
block Hankel matrix $\left(\begin{smallmatrix}U&V\\V&W\end{smallmatrix}\right)
\succeq0$, with $U,V,W$ symmetric.  Diagonalising
$U^{+/2}VU^{+/2}$ on $\operatorname{ran}U$ writes its moment part as
$\sum_j(z_j,t_jz_j)(z_j,t_jz_j)^T$; the remainder is
$\operatorname{diag}(0,W-VU^+V)\succeq0$.
Its value on $K$ is consequently
$\sum_jz_j^TK(t_j)z_j+\sum_iw_i^TK_2w_i\geq0$.
Separation proves SG11, including singular $U$.  The fixed-size
statement transfers to real closed fields.

For $\lambda=m=0$, use SG11 with
\[
\begin{split}
 K_0&=\operatorname{diag}(4C+4D+2d,4A,4\alpha),\\
 K_2&=\operatorname{diag}(4C,4A+4\alpha+2d,4D),\\
 K_1&=\begin{pmatrix}0&0&-(4D+2d)\\
                    0&0&-(4\alpha+2d)\\
                    -(4D+2d)&-(4\alpha+2d)&0\end{pmatrix}.
\end{split}                                                     \tag{SG12}
\]
Necessity follows by setting $W=a/u$ and letting $u\to0$ in the
homogeneous state form.  The surviving form, after $R=a-Y$, is
$2t_5D_vX^2+(R,Y,Z)^TK(v)(R,Y,Z)$.
For sufficiency the determinant identity gives
\[
 \det K(v)=\frac{8D_v}{t_5}(q_{10}+hv^2+q_{12}v^4),
\]
whereas
$L=h+q_{10}/V+q_{12}V+U(q_{20}/V+q_{21}+q_{22}V)$.
PSD of $K$ makes the first three terms nonnegative; all remaining
terms are nonnegative.  SG4 and the endpoint argument finish the
equivalence.  When $d=0$ all forms are event-square sums.
The Gram is strictly feasible at
$C=A=2$, $D=\alpha=d=n=\beta=1$, $S=0$:
after permutation its two Schur margins are $125/56>0$.

For $\lambda=\alpha=0$, put $p=4C+4D+2d$,
$q=4m+4\beta+2d$, $r=4D+2d$.  The exact lift is
\[
 B_1=\begin{pmatrix}4m&-4m&\xi\\-4m&q+4n&q\\\xi&q&q\end{pmatrix}
 \succeq0,\quad
 B_2=\begin{pmatrix}p+4A&4D&-p\\
                 4D&4D&-(4m+r)-\xi\\
                 -p&-(4m+r)-\xi&p\end{pmatrix}\succeq0.          \tag{SG13}
\]
Indeed set $u=\varepsilon s$, $v=\varepsilon$,
$(W,Z)=(w,z)/\varepsilon$, $X=x$, $Y=y$.
The surviving form is
\[
\begin{split}
 G_s={}&4nx^2+4m(w-x-sz)^2+4C(sw-y)^2+4Ay^2\\
       &+r(sw-y-z)^2+(4\beta+2d)(x+sz)^2-4dyz-2dz^2.
\end{split}
\]
Its determinant is $16(q_{02}+hs^2+q_{20}s^4)$.  The same
nonnegative-remainder argument proves equivalence with physicality.
Applying SG11 to its state matrix, the zero diagonal entries on
$sx,sy$ force their Gram rows to vanish.  Coefficient matching leaves
only $\xi$, and the remaining Gram splits into the two blocks SG13.
Equivalently
$G_s=(w,x,sz)B_1(w,x,sz)^T+(y,z,sw)B_2(y,z,sw)^T$;
the $\xi$ terms cancel because $w(sz)=z(sw)$.
At $C=A=n=\beta=2$, $m=D=d=1$, $\xi=-5$, the successive
principal minors are $(4,72,234)$ and $(22,72,234)$.
Both lifts retain all nonnegative coefficient constraints.

\paragraph{Why these source lifts give separated certificates.}
Represent either source cone as $a\geq0$ and
$\exists z\,M(a,z)\succeq0$, absorbing the signs into diagonal
blocks.  The displayed points make this entire homogeneous lift
strictly feasible.  The image of its dual PSD cone is closed:
pairing with a fixed strict feasible matrix bounds the trace of
every PSD preimage of a convergent dual sequence.  Separation
therefore proves attained dual representation for every functional
nonnegative on the source cone.  This is a finite-dimensional
argument, and its fixed matrix-size quantified conclusion transfers.

Given a source in $K_c$, choose its lift $A_c\succeq0$ there.
Given a physical test in $K_t$, its evaluation functional belongs
to the dual source cone in $K_t$; choose its attained dual matrix
$B_t\succeq0$ there.  Their separated trace pairing equals the
slack because all auxiliary-coordinate coefficients cancel.
Factoring $A_c=RR^T$ and $B_t=SS^T$ in their respective fields gives
\[
 \langle A_c,B_t\rangle_{\otimes}
       =\sum_{i,j}\left(\sum_k R_{ki}\otimes S_{kj}\right)^2.
                                                               \tag{SG14}
\]
No mixed inverse or limiting certificate is involved.  Literal
zero coefficients and singular matrices are included.

\section{The remaining quadratic form}

Let $K_c,K_t$ be separately labelled real closed extensions of
$\mathbb R$ and $T=K_c\otimes_{\mathbb R}K_t$.
The source partition in Chapter~\ref{ch:cover} reduces the following
assertion to explicit local calculations.

\paragraph{Canonical assertion $\mathcal C$.}
Take $d=1$ and source coefficients in either pattern SG10.
Assume only that there exist $u_c,v_c>0$, $W_c>0$, and
$X_c,Y_c\in K_c$ for which SG7 holds with
$\psi_c=(W_c,X_c,Y_c,1)$.  For every independent
$u,v,W>0$ and $X,Y\in K_t$, the full expression SG2 with
$Z=1$ belongs to $\Sigma T^2$.

Physicality follows from SG7. Every source produced by the extreme
reduction has this form after the angle signs and state normalization
are fixed in its own field. The independent testing state remains
arbitrary in $\mathcal C$.

\paragraph{Assembly lemma.}
Once $\mathcal C$ is proved, every transferred quantum-valid affine
functional and every independent transferred quantum behaviour
have separated slack in $\Sigma T^2$.  Therefore the universal
tensor-SOS finite-lift criterion applies to the full eight-coordinate
quantum set.

\emph{Proof.}
The source reduction gives at most nine canonical source atoms,
each of which is an event atom, a closed-face atom, or SG10 with
an actual contact.  Event atoms and closed faces were handled above.
Restore $d$ by multiplying each root by $\sqrt d\otimes1$.

For a canonical contact source, finite angle signs are harmless.
The polynomial is unchanged by
\[
 (u;W,X,Y,Z)\mapsto(-u;W,X,-Y,-Z),\qquad
 (v;W,X,Y,Z)\mapsto(-v;W,-X,Y,-Z).
\]
Literal zero/infinite angles have nonnegative pure-testing $H$ and
nonnegative event squares, as already proved.  For $u,v>0$, the
signed core is nonnegative whenever $WZ\leq0$.  To check this,
put $h_0=1+v^2+u^2v^2$; its $X,Y$ block is
$2\left(\begin{smallmatrix}h_0&-uv\\-uv&h_0\end{smallmatrix}\right)
\succ0$, and the remaining Schur form on $W,Z$ is
\[
 \frac{2D_uD_v}{h_0^2-u^2v^2}
 \{u^2v^2(W^2+Z^2)-2uvh_0WZ\}\geq0.                 \tag{SG15}
\]
This is a pure testing calculation.  Its square root in $K_t$,
together with $\sqrt d$ in $K_c$, certifies the signed core;
every event remains a square.  It includes $Z=0$.
Otherwise $Z\ne0$ and $W/Z>0$; divide the state by $Z$ inside
$K_t$ and apply $\mathcal C$.  Restore $Z^2$ afterwards.
Real and imaginary state parts add because the measurements are
real.  All homogeneous angle normalisations are pure testing
squares and have nonzero denominators on their stated charts.

Finally every testing behaviour is a mixture of at most nine real
two-qubit behaviours.  If source and testing weights are $a_j,w_i$
and atom certificates are $\sum_kh_{ijk}^2$, the entire slack is
\[
       \sum_{i,j,k}
           [ (\sqrt{a_j}\otimes\sqrt{w_i})h_{ijk}]^2.
\]
The affine constant is preserved because $\sum_iw_i=1$.
Zero weights are never inverted.  This proves the universal tensor
assertion; applying the criterion in Chapter~\ref{ch:tensor} is the final
qualitative existence step. \hfill$\square$


\chapter{Coordinates of the contact families}
\label{ch:entry}
\label{component:p477c0ac30750}
\label{p477c0ac30750:sec:rederived-contact-entry}
The source support, physicality, curvature and deletion facts used below were proved in Chapter~\ref{ch:geometry}, including (SG17)--(SG22).
\section{Signs and stationary coordinates}
\label{component:paf13bbf8fd32}

\label{paf13bbf8fd32:sec:global-native-entry-rederived}

The stationary equations give rational formulas for the event weights.
To use those formulas at a singular limit, we first determine the signs
of every denominator. The contact geometry of Chapter~\ref{ch:geometry}
provides the required strict inequalities.

\subsubsection*{Statement and the real maximal-contact graph}

Use the raw full parent \(F\), with \(d>0\).
Every source below is physical and has an actual finite noncommuting
zero.  The two allowed coefficient domains are
\[
\begin{array}{ll}
 {\rm C1}:&\lambda=n=0,\quad m,C,A,\alpha>0,\quad
 D,\beta\ge0,\quad D+\beta>0,\\
 {\rm C2}:&\lambda=\beta=0,\quad m,C,A,D,\alpha>0,\quad n\ge0.
\end{array}\tag{EN1}
\]
Make the actual angle gauges positive and put \(Z_c=1\), as permitted
by SG3--SG4.  Then \(W_c>0\) by SG15 and SG21.  We prove the strict signs
\[
\begin{gathered}
 X_c<0,\quad Y_c>0,\quad u_cW_c-Y_c>0,\quad
 1-v_cY_c>0,\\
 R_{m,c}:=v_cW_c-X_c+u_cv_cY_c-u_c>0,\\
 R_{D,c}:=u_cW_c+u_cv_cX_c-Y_c-v_c<0,\\
 X_c+u_c>0,\qquad W_c+v_cX_c>0.
\end{gathered}\tag{EN2}
\]
These inequalities hold over every real closed coefficient field.
In particular they are strict ordered-field statements, with no
assertion that a gap has positive real residue.

Here is the common connected-graph argument, proved over \(\mathbb R\)
before first-order transfer.  In either exact two-omission family,
initially require all six retained events positive, normalize \(C=2\),
and let their other five weights vary over the positive orthant.  Put
\[
 d_*(w)=\max\{d\ge0:(d,w)\text{ is physical}\}.
 \tag{EN3}
\]
The admissible set is a closed interval by convexity.  It has positive
length: at \(d=0\) the central Laurent coefficient is positive, its
nonnegative Cauchy--Binet expansion containing \(16mCAD\); the outer
coefficients remain nonnegative for \(d\ge0\).  Thus sufficiently small
positive \(d\) is physical by SG3--SG6 and the homogeneous endpoint
limits.  A fixed strictly CHSH-violating real test bounds \(d\) above.

For these positive six-event supports the active outer exponents are
constant and have zero in the interior of their convex hull.  The
logarithmic Laurent sum is strictly convex and coercive.  At \(d_*\)
its unique minimum is zero.  Indeed if its minimum were \(\mu>0\),
a small increase of \(d\) preserves every active outer coefficient
within a relative error \(\varepsilon\) and the central coefficient
within absolute error \(\varepsilon\).  As in the extreme-source argument, this preserves a
positive uniform lower bound, contradicting maximality.
There is therefore one positive squared-angle contact, and SG3--SG4 gives
one state with \(Z=1\).

At that contact the core value is strictly negative (SG20--SG21).
The Jacobian of the six equations
\(\nabla_zF=0,F=0\), in \(z=(W,X,Y,u,v)\) and \(d\), is
\[
 \begin{pmatrix}\nabla_z^2F&\partial_d\nabla_zF\\
 0&\partial_dF\end{pmatrix},
 \qquad \nabla_z^2F\succ0,\quad\partial_dF<0 .
 \tag{EN4}
\]
The real analytic implicit-function theorem gives a local analytic
contact branch.  SG7 makes it physical and the negative core value
makes its \(d\) maximal.  Uniqueness makes these branches agree.
Thus the normalized contact is analytic on the entire connected
positive orthant.  Every physical finite zero with those event weights
lies on this graph, because its negative core value prevents increasing
\(d\) at that same test.

\subsubsection*{C1: strict signs including the two literal boundary faces}

Write
\[
 (W,X,Y,Z)=(S,u(\theta-1),uY_0,1).
 \tag{EN5}
\]
We first determine the sign of $\theta$ from an adjugate entry. With \(B=\operatorname{adj}P\),
\[
\begin{gathered}
 uB_{44}+B_{24}
 =16uv^2D_u^2D_v\,\mathcal F,\\
\begin{aligned}
 \mathcal F={}&(d+2D)\{(d+2\alpha)(mD_v+Cu^2v^2)
                       +2AD_v(m+Cu^2)\}\\
 &+2mC\{d(u^2v^2+D_v)+2\alpha+2Dv^2D_u\}>0 .
\end{aligned}
\end{gathered}\tag{EN6}
\]
Since \(P\) has one-dimensional kernel and three positive eigenvalues,
\(B=\gamma\psi\psi^T\) with \(\gamma>0\).  Equation EN6 gives
\(\theta>0\).

Neither \(Y_0=0\) nor \(S-Y_0=0\) is possible.  These walls would delete
the positive \(A\)-event or \(C\)-event, respectively, while preserving
the kernel and both angle jets.  The resulting omission set contains
an exclusive pair, so has noncoercive support.  SG7 and SG22 rule this out.
The remaining cross-multiplied contact identity is
\[
 m\,[S+uv(\theta-1)]\,
       [v(S+u^2Y_0)-u\theta]
   =2u^3(S-Y_0)(1-\theta)D_v .
 \tag{EN7}
\]
For precision about exceptional source minors, use the four raw kernel
equations and the two full angular derivatives.  Their selected six-column
system, for \((d,m,A,D,\alpha,\beta)\) with \(C=2\), has
determinant
\[
 -8192\theta u^2vD_u^2D_v^2\,\mathcal D,\qquad
 \mathcal D=Y_0D_u(uY_0+v)(1-uvY_0)
 [S+uv(\theta-1)]R_m(-R_D).
 \tag{EN8}
\]
EN7 follows first from its rational solution wherever this selected
minor is nonzero: the Cramer determinant for the \(m\)-column,
multiplied by \([S+uv(\theta-1)]R_m\), equals this determinant
times \(2u^3(S-Y_0)(1-\theta)D_v\).  This is an exact
polynomial determinant identity.  The equivalent native formula is
\(m=2ae^2c_2/(Mc_3)\), with the definitions in EN13 below:
\(c_2/c_3=(1-\theta)D_v/E\),
\(M=R_m/(Sv)\), \(E=[S+uv(\theta-1)]/S\),
\(a=u(S-Y_0)/v\), \(e=u/S\).  These substitutions give EN7
exactly.  The native contact equations are polynomial identities, so
this argument does not assume the signs being proved.

At the exact positive seed
\[
 (u,v;S,Y_0,\theta)=(1,1;1,\tfrac14,\tfrac12),\qquad
 (d,m,C,A,D,\alpha,\beta)
 =(\tfrac{27}2,4,2,9,\tfrac54,\tfrac{13}4,\tfrac94),
 \tag{EN9}
\]
all six stationary equations hold and the selected minor is nonzero.
Consequently that analytic minor is not identically zero on the
connected contact graph.  Its nonzero set is dense.  Continuity extends
EN7 to every contact on the graph.  At \(\theta=1\), EN7 and \(S,m>0\)
would imply \(R_m=0\), permitting deletion of \(m\) and creating an
exclusive omitted pair, again impossible.  The three continuous
nonzero functions \(Y_0,S-Y_0,\theta-1\) therefore have constant signs,
fixed by EN9:
\[
 S>0,\qquad 0<Y_0<S,\qquad 0<\theta<1.
 \tag{EN10}
\]

The full-parent C1 reflection swaps
\((\lambda,n),(m,C),(A,\alpha),(D,\beta)\), sends \(u\) to \(1/u\),
and sends the state to
\[
 (vS-u(\theta-1),\,-S-uv(\theta-1),\,
                   1-uvY_0,\,uY_0+v).
\]
Its state matrix \(G\) has \(G^TG=D_vI\); the full core and every event
obey \(u^2F_{\sigma_1c}(1/u,v;G\psi)=D_vF_c(u,v;\psi)\).
Its last coordinate is positive.
Applying EN10 to the transformed normalized state proves successively
\(1-uvY_0>0,\ S+uv(\theta-1)>0,\ R_D<0,\ R_m>0\).
Together with EN10 this is precisely EN2.

For a literal third omission \(D=0\) or \(\beta=0\), start at the
given real contact and open that coefficient to \(\varepsilon>0\).
EN4 is still invertible by the strict contact Hessian (SG18--SG19).  The local analytic branch stays at
positive \(d\), has nonnegative event weights, and is physical by SG7.
The already proved inequalities hold on it.  At the limit, EN6 and
SG15 and SG21 stay strict; the still positive \(A,C,m\) and the same deletion
arguments make the other weak inequalities strict.  EN7 passes to
the limit, and so does the reflection.  Thus EN2 holds on both
boundary faces.  This is an implicit-function limit at a fixed
finite real contact, not an assertion that an optimizer cannot escape
to infinity.  First-order transfer now gives the entire implication
over every real closed field.

\subsubsection*{C2: all signs and a positive omitted-row cofactor}

Use the same positive six-event graph.  Five retained-amplitude walls
are excluded directly by SG22:
\[
 Y=0,\quad uW-Y=0,\quad R_m=0,\quad R_D=0,\quad vY-1=0.
 \tag{EN11}
\]
Deleting \(A,C,m,D,\alpha\), respectively, introduces an exclusive
omitted pair in addition to \(\lambda=\beta=0\).
The wall \(X=0\) deletes \(n\); the resulting source is on the C1
three-consecutive-omission face \(\lambda=n=\beta=0\).
It is physical by SG7, so EN10 gives \(X<0\), a contradiction.
Thus all six retained event rows have fixed signs on the connected
C2 graph.  The exact seed
\[
 (n,m,C,A,D,\alpha)=(\tfrac18,\tfrac38,\tfrac16,
                              \tfrac58,\tfrac38,\tfrac16),\quad d=1,
 \qquad (u,v;W,X,Y,Z)=(1,1;1,-\tfrac14,\tfrac14,1)
\]
has every required positive weight and every stationary equation.
It fixes their signs to the first six signs in EN2.

Here is a direct proof of the omitted-row sign, including \(n=0\).
Let \(P_0\) be this raw C2 matrix with \(n\) set to zero.
Changing \(n\) adds only \(4nD_uD_v\) to its \(XX\) entry.
The cofactor \(B_{24}\) does not contain that column, whereas
\(B_{44}\) changes by \(4nD_uD_v\det(P_0|_{W,Y})\).
Consequently
\[
\begin{split}
 uB_{44}+B_{24}
  ={}&16uv^2D_u^2D_v\,\mathcal F\\
    &+4unD_uD_v\det(P_0|_{W,Y})>0 .
\end{split}\tag{EN12}
\]
Here \(\mathcal F\) is EN6 with \(\beta=0\).  The determinant in
the second line is positive: the \(Z=0\) state block already has
the strictly positive signed-core form SG3--SG4.  Thus \(X+u>0\).
The C2 full-parent reflection
\[
 (u,v;W,X,Y,Z)\mapsto(1/v,1/u;Z,X,Y,W),
 \quad (m,D),(C,\alpha),(\lambda,\beta)
\]
has the exact homogeneous identity
\(u^2v^2F_{\sigma_2c}(1/v,1/u;Z,X,Y,W)=F_c(u,v;W,X,Y,Z)\).
The transformed contact
has positive last coordinate \(W\).  Applying the just-proved
\(X'+u'Z'>0\) gives \(W+vX>0\).
This proves EN2 for positive \(n\).  For literal \(n=0\), EN1 is
already a C1 triple face and its signs apply directly; EN12 remains
strict as well.

\subsubsection*{Exact native-coordinate transport and uniqueness}

For C1, reflect if needed so \(0<k=u_c\le1\), and define
\[
\begin{gathered}
 S=W_c,\quad r=Y_c/(kS),\quad
 \theta=1+X_c/k,\quad h=1-\theta,\\
 \eta=v_c/(kS),\quad \zeta=Sv_c/k,\quad
 f=\sqrt\eta,\quad g=\sqrt\zeta,\quad e=k/S=kf/g,\\
 a=(1-r)/\eta,\quad t=\theta/\zeta,\quad x=k^2,\\
 R=1+xr,\quad T=1+xh,\quad M=R-t,\quad J=T-a,\\
 E=1-x\eta h,\qquad Q=1-x\zeta r .
\end{gathered}\tag{EN13}
\]
By EN2, \(0<r,\theta<1\), all \(\eta,\zeta,a,t\) are positive,
and
\[
 R_m=k\zeta M>0,\quad R_D=-v_cJ<0,\quad
 W_c+v_cX_c=SE>0,\quad 1-v_cY_c=Q>0.
\]
Thus \(M,J,E,Q>0\), and \(0<a,t<2\).
The selected determinant EN8 is now strictly nonzero.  The native
seven coquotients and source coefficients SD1 in Section~\ref{component:p0acdb66bba2f}, with their exact stationarity identities SD9--SD11,
satisfy exactly the six stationary equations at this same contact.
Uniqueness therefore identifies them with the original normalized
source, including \(D=0\) or \(\beta=0\).  No asymptotic hypothesis
is used for these polynomial identities.

For C2, normalize \(d=1\), and set
\[
\begin{gathered}
 c=u_c,\quad a=v_c/u_c,\quad b=aW_c,\quad r=-X_c/u_c,\quad
 t=(u_cW_c-Y_c)/v_c,\quad z=c^2,\\
 h=b-a^2t=aY_c/u_c,\quad
 g=(1+z)b-(1-r+a^2zt)=R_m/u_c,\quad
 j=1+zr-t=-R_D/v_c .
\end{gathered}\tag{EN14}
\]
The complete strict native domain follows:
\[
 a,b,c,t>0,\quad 0<r<1,\quad h,g,j>0,\quad
 1-zh=1-v_cY_c>0,\quad b-a^2zr=a(W_c+v_cX_c)>0.
 \tag{EN15}
\]
The actual six-column stationary determinant is, in the unweighted
event convention of the retained determinant proof,
\[
 -64ac^6rgt\,hj(1-zh)(1+z)^2(1+a^2z)^2\Delta,
\]
\[
 \Delta=1+(1+z)b+z(1+z)b^2+2z(1+z)br+z^2r^2+2zr
        +a^2z\{(1+z)b+1+zr(2-r)\}\ge1 .
 \tag{EN16}
\]
That body provides its explicit reduction to a three-by-three
determinant.  The six rational coefficients in
Section~\ref{component:p6e5d41ed80b3} solve these six equations at the
contact EN14.  EN16 makes the solution unique, so they are exactly
the original source coefficients.  This remains valid at \(n=0\);
one does not delete its column from the square system.

All gauge actions, positivity statements and source equations used
above are finite polynomial or rational ordered-field statements
with denominators explicitly nonzero.  The real connected-graph
proof therefore transfers as one fixed first-order implication.
Together, EN1--EN16 discharge the native sign and positive-gap
interface for both families in SG10.  They do not prove an
all-tests tensor assertion; that assertion is supplied by the
source partitions and their proofs in Appendices B and C.


\chapter{Algebraic tools for singular contacts}
\label{ch:mechanisms}
A singular contact has vanishing quadratic terms. The following identities
express its higher-order terms using roots with a positive residue Gram.
Chapter~\ref{ch:tensor} then turns that Gram into an exact sum of squares.
\section{Finite algebra and contact completion}
\paragraph{Finite penalty lemma.}
Let $R$ be a real symmetric quadratic form and $L$ a real linear
map.  If $R$ is positive definite on $\ker L$, then
$R+M\|Lz\|^2$ is positive definite for some real $M>0$.

\emph{Proof.}
Choose a linear complement of $\ker L$.  The restriction of $L$
to that complement has a positive least singular value.  In these
coordinates the kernel block of $R$ is positive definite, and
adding $M L^TL$ makes the other Schur complement positive for
sufficiently large $M$. \hfill$\square$

For an opening argument one must additionally exhibit a bounded
remainder after subtracting the divergent squares, and a bounded
right inverse for their independent hard rows.  If a family has a
uniform positive lower bound before subtraction, its limit on the
limiting kernel is positive: lift a kernel vector to exact kernels
with that right inverse and take the real limit.  The finite penalty
lemma then supplies finite caps.  The applications in Appendices B and C give the
required bounded minimizers and independent rows.

\paragraph{Small-cost helper lemma.}
Suppose the main roots have a strict real residue Gram and an exact
identity has the form $A(z,H)-H^2=\sum R_i^2$, where every
coefficient of the quadratic form $A$ on the main roots $z$ and
$H$ is infinitesimal.  Then any prescribed finite positive diagonal
budget $cH^2$ can be introduced by the exact replacement
\[
          F=(F-cA+cH^2)+c\sum R_i^2.
\]
It changes the main residue Gram by zero and retains the displayed
positive squares.  Finitely
many helper roots can be treated in a stated order, after which a
real Schur complement handles their bounded cross coefficients.
When the identities are triangular, introduce the auxiliary roots in
that order. In the double-infinitesimal C1 case this means obtaining
the radial pivot before the angular roots.


\section{Finite polynomial division at a contact}
\label{component:pfabc29e28ab3}
% Optional module: the source notation is that of (UZ4)--(UZ6).
% All statements below are exact; no localized ring is used.


Write $d=\delta$ in this section; this local parameter is distinct from
the global Bell coefficient $d_{\rm B}$. Put
\[
 R_0=\mathbb Q[d,k,b,a,t],\qquad R=R_0[u,w,v],\qquad
 J=(u-k,w-dk,v-dk^2b^2).
\]
For $F\in R$, write $F_c=F(k,dk,dk^2b^2)$. Define
\[
 A=wv,\quad B=dv,\quad C=db^2uw,\quad D_0=d^2k^3b^2,
 \qquad M=(A,B,C),
\]
\[
 x=B-C,\qquad y=A-kC,\qquad z=uB-kC,\qquad h=kC-D_0.
\]
Equivalently, write $g_0=x$, $h_{13}=y$, $h_{23}=z$, and $h_{34}=h$. Then
$D_0=a_4^2$, whereas the source monomial in $g_4$ is $kD_0$.
Adjoin the separated radicals used in the four-monomial identity,
\[
 a_1=\sqrt{wv},\quad a_2=\sqrt{duv},\quad
 a_3=b\sqrt{dkuw},\quad a_4=dk^{3/2}b,
 \qquad q=a_1a_2-a_3a_4.
\]
Only $q^2\in R$ is needed in the scalar lift. The following three identities
are the entire conductor calculation:
\begin{align}
 C(u-k)&=z-ux,\label{pfabc29e28ab3:CC-u}\\
 C(w-dk)&=dy-wx,\label{pfabc29e28ab3:CC-w}\\
 C(v-dk^2b^2)&=b^2(wz+dky-kwx).\label{pfabc29e28ab3:CC-v}
\end{align}
Thus $CJ\subset(x,y,z)$ without inverting any parameter.
For
\[
 r_*=y+z-h-2kx
\]
we also have the two exact identities
\begin{align}
 r_*&=(v-dk^2b^2)\bigl((w-dk)+d(u-k)\bigr)
       -dkb^2(u-k)(w-dk),\label{pfabc29e28ab3:CC-star}\\
 kCr_*&=q^2-yz.\label{pfabc29e28ab3:CC-repair}
\end{align}
The first says $r_*\in J^2$. The second supplies the otherwise missing
conductor factor through one of the original circuit squares.

\paragraph{Explicit monic division.}
Every $S\in R$ has a polynomial decomposition
\[
 S=S_c+(u-k)T_u+(w-dk)T_w+(v-dk^2b^2)T_v.                 \tag{CC1}
\]
One obtains it by substituting $u=k$, then $w=dk$, then $v=dk^2b^2$,
and using $X^n-Y^n=(X-Y)\sum_{j=0}^{n-1}X^{n-1-j}Y^j$ at each step.
There are no parameter denominators. If $S_c=0$, equations
\eqref{pfabc29e28ab3:CC-u}--\eqref{pfabc29e28ab3:CC-v} give
\[
 CS=x(-uT_u-wT_w-b^2kwT_v)
     +y(dT_w+dkb^2T_v)+z(T_u+b^2wT_v).                 \tag{CC2}
\]

\begin{lemma}[First-order and restricted second-order contact lifts]
In the above ring, the following assertions hold, with polynomial
coefficients in $R$:
\[
 M\cap J\subset (x,y,z),\qquad
 \bigl(M^2+(D_0^2)\bigr)\cap J^2
       \subset (x,y,z,h,q)^2.                         \tag{CC3}
\]
Here $(D_0^2)$ is the principal ideal generated by $D_0^2$; the second
assertion does not replace it by all mixed products involving $D_0$.
\end{lemma}
\begin{proof}
For the first assertion, write $F=R_AA+R_BB+R_CC$. Since
$A=y+kC$ and $B=x+C$,
\[
 F=R_Ay+R_Bx+CS,\qquad S=kR_A+R_B+R_C.
\]
The contact value is $F_c=C_cS_c$, where $C_c=d^2k^2b^2$ is a nonzero
polynomial. Since $R_0$ is a domain, $F_c=0$ implies $S_c=0$.
Apply (CC2). This cancellation proves a zero identity; it introduces no
coefficient $1/C_c$.

For the second assertion choose any displayed representation
$F=\sum_{i\le j}f_{ij}m_im_j+\rho D_0^2$, with
$(m_1,m_2,m_3)=(A,B,C)$. Substitute
\[
 A=y+kC,\qquad B=x+C,\qquad D_0=kC-h.
\]
This gives
\[
 F=Q_2+C(xH_x+yH_y+zH_z+hH_h)+C^2S,\qquad H_h=-2k\rho,       \tag{CC4}
\]
where $Q_2\in(x,y,h)^2$ and initially $H_z=0$.
The contact value again implies $S_c=0$. Use (CC2) on $CS$ and absorb
$C^2S$ into the four coefficients in (CC4). This leaves $H_h=-2k\rho$
unchanged and gives the same formula without $C^2S$.

The first derivatives of $F\in J^2$ vanish at contact. The first three
root gradients have determinant
\[
 \det\left(\frac{\partial(x,y,z)}{\partial(u,w,v)}\right)_c
       =-d^4k^4b^4\ne0.                              \tag{CC5}
\]
Equation~\eqref{pfabc29e28ab3:CC-star} gives a gradient relation with coefficient vector
$(-2k,1,1,-1)$. Its coefficient at $h$ is $-1$.
Set $\lambda=2k\rho_c$. Subtract
$\lambda(-2k,1,1,-1)$ from $(H_x,H_y,H_z,H_h)$.
The last resulting contact value is zero, and the first three satisfy the
homogeneous linear system with determinant (CC5). Multiplication by its
adjugate and cancellation in the domain $R_0$ show that those three
contact values are zero too. No inverse of the determinant is used in a
coefficient: each resulting polynomial $J_i$ lies in $J$ and is divided
only by the monic procedure (CC1).
Consequently
\[
 F=Q_2+2k\rho_cCr_*+C\sum_{i=1}^4r_iJ_i
   =Q_2+2\rho_c(q^2-yz)+\sum_{i=1}^4r_i(CJ_i),          \tag{CC6}
\]
where $(r_1,r_2,r_3,r_4)=(x,y,z,h)$. Each $CJ_i$ is explicitly a
combination of $x,y,z$ by (CC2), which proves the assertion in the
unlocalized polynomial ring.
\end{proof}

The explicit scalar and both state-normal coefficients in
Section~\ref{sec:c1-bounded-reciprocal} satisfy these ideal conditions.
There the same monic divisions also preserve a positive coefficient
weight. This supplies bounded corrections with zero residue, which
can be added to the strict Gram matrix.

\section{A real two-block allocation}
\label{component:p6b22e944bd7c}
\paragraph{A rank-one allocation lemma.}
Let $P$ be a positive definite real symmetric matrix, let $\ell,b$ be
real columns, and let $H,c_0>0$. Set
\[
 d=\ell^TP^{-1}\ell,\qquad e=\ell^TP^{-1}b,\qquad
 h=c_2-b^TP^{-1}b.
\]
Suppose
\[
 p(T)=c_0+(h+4Hdc_0)T+4H(dh+e^2)T^2>0\quad(T\ge0),
 \qquad dh+e^2>0.                                      \tag{UA1}
\]
There exists a real $\lambda$ such that
\[
 |\lambda|<\sqrt{4Hc_0},\qquad
 h+2e\lambda-d\lambda^2>0.                             \tag{UA2}
\]
Indeed put $s=\sqrt{4Hc_0}$. If $d>0$ and $|e|<ds$, take
$\lambda=e/d$. If $|e|\ge ds$, the value at the appropriate endpoint
is $h+2|e|s-ds^2$. Were this nonpositive, writing
$x=|e|-ds\ge0$ would give
$h+ds^2\le-2sx$ and $dh+e^2\le x^2$; the quadratic in (UA1)
would then have a nonpositive value at a positive argument. Thus a
nearby interior $\lambda$ works. This includes equality $|e|=ds$.
The scalar interval argument also holds in a limiting case $d=0$;
then $e\ne0$ follows from $dh+e^2>0$. (For a finite positive $P$
itself, $d=0$ would force $\ell=e=0$, so this case occurs only
when using the continuous scalar extension at a boundary.)

Consequently the two real matrices
\[
 \begin{pmatrix}P&b-\lambda\ell\\(b-\lambda\ell)^T&c_2\end{pmatrix},
 \qquad
 \begin{pmatrix}4H&\lambda\\\lambda&c_0\end{pmatrix}       \tag{UA3}
\]
are positive definite. This is an elementary real lemma. Its inverse
and square root are never taken in the tensor ring.


\section{Noncommuting testing residues}
\label{component:pe25c50379365}
\label{component:p76173e99815f}
\label{component:p24f8a75e651e}
\label{component:p4e9ae18b7af4}
\label{component:pa2cbe6ecbeb4}
\label{component:pb16e5066334b}
\label{component:pf47f6a4e3a60}
\label{component:p94dfbc8a5970}
\label{component:p2c3fc22497a5}
% Standalone replacement for the NONCOMMUTING portions of components4--12.
% Uses only the tensor foundation (positive residue, strict bounded Gram,
% pure-factor cancellation), standard real algebra and the displayed parent.

\label{sec:noncommuting-refined}

Use the two labelled fields and the bounded tensor ring of
Chapter~\ref{ch:tensor}. The strict-Gram lemma applies to arbitrary
root columns, including unbounded rows.

\subsection{Parent, physicality and contact curvature}

For real homogeneous state coordinates $\psi=(W,X,Y,Z)$ and half angles
$u,v$, set
\[
 D_u=1+u^2,\quad D_v=1+v^2,\quad
 M=vW-X+uvY-uZ,\quad J=uW+uvX-Y-vZ.
\]
All nine source coefficients $d,\lambda,n,m,c,A,D,\alpha,\beta$ are
nonnegative. The entire parent is
\begin{align}
 F={}&4\lambda D_u(W+vX)^2+4nD_uD_vX^2+4mM^2
       +4cD_v(uW-Y)^2+4AD_uD_vY^2\notag\\
 &+(4D+2d)J^2+(4\alpha+2d)D_u(vY-Z)^2
       +(4\beta+2d)D_v(X+uZ)^2-2dD_uD_vZ^2.
 \tag{NC1}
\end{align}
Write $C$ for the coefficient of $d$ in this expression. A source is
\emph{physical} if its state matrix is positive semidefinite at every
measurement pair over its own field $K_c$. The claim below concerns
evaluation of (NC1) at an independent test over $K_t$.

\begin{theorem}\label{thm:noncommuting-refined}
For every physical source in (NC1), its tensor evaluation is a finite
sum of squares at every independent test for which both measurement
pairs are noncommuting in residue, and every homogeneous real or
complex state. No relation between source scales or between source
and testing scales is assumed.
\end{theorem}

Local frame changes and a projective state chart are pure-testing
operations. At a noncommuting residue one may take $u,v$ to be bounded
units, and normalize one state coordinate to a bounded unit. The
original slack is restored by positive pure-testing norm factors.
For complex states, apply the real result separately to real and
imaginary parts, since the state matrix is real symmetric. We prove
the real case.

If $d=0$, the parent is already a sum of event squares.
For the remaining $d>0$ case, use the eight outer coefficients
$q_{ij}$ and the positive quantities $t_i$ in (SG5)--(SG6), put
$q_{11}=h$, and denote the event coefficient by $c$ here. Equivalently
\[
 \det P=16D_u^2D_v^2\sum_{i,j=0}^2q_{ij}u^{2i}v^{2j}.
 \tag{NC2}
\]
Thus every coefficient in the constrained determinant calculations
below is specified by that displayed finite dictionary.

For $U=u^2,V=v^2>0$, let
\[
 L(U,V)=q_{11}+\sum_{(i,j)\ne(1,1)}q_{ij}U^{i-1}V^{j-1}.
 \tag{NC4}
\]
Every nonconstant coefficient is nonnegative. Over the real residue
field, a stationary zero of such a Laurent polynomial satisfies
\[
 \sum_{e\in E}w_e e=0,\qquad
 \nabla^2_{\log U,\log V}L=\sum_{e\in E}w_e ee^T,
 \quad w_e=q_eU^{e_1}V^{e_2}>0.                       \tag{NC5}
\]
Consequently its angular Hessian is positive definite exactly when
the active exponent vectors span $\mathbb R^2$. Logarithms are used
only for this real calculation; the displayed formulas are algebraic.

When $d>0$ and $uv\ne0$, the core restricted to $Z=0$ is
\[
 C|_{Z=0}=2\{(uW+uvX-Y)^2+D_uv^2Y^2+D_vX^2\},       \tag{NC6}
\]
which is positive definite in $(W,X,Y)$. Thus a physical state matrix
at a nonzero kernel has exactly three positive transverse directions.
The same observation supplies two or one positive transverse
directions after imposing one or two independent state-linear event
constraints, respectively.

\subsection{Exact local lifting with unequal hard weights}

We record the local argument that will turn real curvature into an
exact tensor certificate. Suppose, in a real coordinate neighborhood,
the complete physical parent has the form
\[
 F(x,y)=\sum_{i=1}^k a_i b_i(y)x_i^2+\sigma G(x,y),
 \qquad \varepsilon_i=\sigma/a_i>0\text{ infinitesimal},
 \tag{NC7}
\]
where $a_i,\sigma$ are pure-source, all coefficients of $b_i,G$ are
bounded, $\bar x=0$, $\bar y=y_0$, and $\bar b_i(y_0)>0$.
The case $k=0$ is allowed. Rational expressions are also allowed when
their denominators are products of pure-source units and real testing
polynomials nonzero at the residue. Multiply first by their squares;
the resulting polynomial has the same physical sign and the same
strictness at a zero, and the factors are canceled afterward.

Let $g(y)=\overline G(0,y)$. Physicality on the exact hard kernel
implies $g\ge0$ in its real neighborhood. If $g(y_0)>0$, the roots
$(x_i/\sqrt{\varepsilon_i},1)$ give an exact bounded Gram with residue
$\operatorname{diag}(\bar b_i(y_0),g(y_0))\succ0$:
every term of $G-G(0,y)$ containing $x_i$ acquires a factor
$\sqrt{\varepsilon_i}$ in this frame.

If $g(y_0)=0$ and $\nabla^2g(y_0)\succ0$, physicality also gives
$\nabla g(y_0)=0$. Put $x_i=\varepsilon_i s_i$ in the stationary
equations for $F/\sigma$ and divide the $x_i$ equations by their
nonzero factors. The exact equations are
\[
 2b_i(y)s_i+G_{x_i}(\varepsilon s,y)=0,\qquad
 \sum_i\varepsilon_i s_i^2\nabla b_i(y)
                       +G_y(\varepsilon s,y)=0.       \tag{NC8}
\]
At residue they have a block upper triangular Jacobian with diagonal
blocks $2\operatorname{diag}(\bar b_i(y_0))$ and
$\nabla^2g(y_0)$. The real implicit-function theorem gives a unique
local solution in fixed semialgebraic neighborhoods. Existence,
uniqueness and continuity bounds in those neighborhoods are fixed
first-order statements and transfer to $K_c$. They give an exact
source-field stationary point $(x_c,y_c)$ with
$x_{ci}=\varepsilon_i s_{ci}$ and bounded $s_c$, of the prescribed
residue. No center depends on a testing infinitesimal.

Finite polynomial Taylor expansion at this center, using the rows
\[
 h=\big((x_i-x_{ci})/\sqrt{\varepsilon_i}\big)_i
                    \mathbin{\oplus}(y-y_c),
\]
gives exactly
\[
 F/\sigma=a+h^TQh,\qquad
 a=F(x_c,y_c)/\sigma\ge0,\qquad
 \bar Q=\operatorname{diag}
   (\bar b_i(y_0),\tfrac12\nabla^2g(y_0))\succ0.      \tag{NC9}
\]
Indeed two hard deviations are allocated to their two scaled rows;
one hard deviation and a tangent deviation leave
$\sqrt{\varepsilon_i}$ in the coefficient; a hard packet's pure
tangent term leaves $\varepsilon_i$, because
$x_{ci}/\varepsilon_i$ is bounded. The $G$ terms have the same or
smaller factors. All higher Taylor terms leave an infinitesimal
deviation in the coefficient. The exact linear terms vanish at the
center. This proves the full-parent certificate, including all state
normals and higher terms. In particular, $k=0$ proves the ordinary
regular-contact case.

\subsection{Geometry of the event constraints}

The losing events form one bipartite cycle of local outcome rays,
in edge order
\[
 (\lambda,D,c,A,\alpha,m,\beta,n).                   \tag{NC10}
\]
For completeness its consecutive vertices may be taken as
$A_0^+,B_1^+,A_1^-,B_0^+,A_0^-,B_1^-,A_1^+,B_0^-$.
The four rays at either party are distinct at noncommuting residue
angles. CHSH-preserving relabelings act on $(x,y,a,b)$ as
\[
\begin{split}
g_A:(x,y,a,b)&\mapsto(1-x,y,a,(-1)^yb),\\
g_B:(x,y,a,b)&\mapsto(x,1-y,(-1)^{1-x}a,b),\\
g_T:(x,y,a,b)&\mapsto(y,x,(-1)^yb,(-1)^xa).
\end{split}
\]
They act on the edge indices of (NC10) by $i\mapsto5-i,3-i,2-i$
modulo eight. Hence they generate a cycle rotation and reflection.
They preserve the parent interpretation; chart changes contribute
only positive pure-testing homogeneous factors.

Product residue states need no further argument. For a product state
the CHSH expression satisfies
\[
 S=a_0(b_0-b_1)+a_1(b_0+b_1)
 \le |b_0-b_1|+|b_0+b_1|=2\max(|b_0|,|b_1|)\le2.
\]
The inequality is strict when both measurement pairs are
noncommuting. If a Bob modulus is one, the other is strictly less
than one, and equality would require Alice to be a common eigenstate
of her two noncommuting reflections. Thus $C$ has positive testing
residue and its actual pure-testing value is positive. The term $dC$
and all event packets are then sums of squares using separate roots.

For an entangled real residue state, write its amplitudes as an
invertible two-by-two matrix $\Psi$. A fixed Alice ray $a$ has
exactly one Bob projective ray satisfying $a^T\Psi b=0$, and
conversely. The vanishing events therefore share no endpoint in
(NC10): they form a matching and there are at most four of them.
Any three matching event vectors are state-linearly independent.
To see this, the span of two rank-one matrices with independent
left and right factors contains no other rank-one projective point;
three dependent rank-one matrices must therefore share a left or a
right factor, contrary to the matching condition. For $k\le3$, the
event amplitudes are consequently exact local projective-state
coordinates. Their coordinate inverse has only pure-testing unit
denominators and their packets have the diagonal form (NC7).

There are two perfect matchings. For $\{n,m,A,D\}$ their equations
are $X=Y=0$, $vW=uZ$, $uW=vZ$. Thus $u^2=v^2$; on either sign
branch $W\ne0$ and, after a sign gauge, $W=Z=1$, $v=u=t\ne0$.
With $t=u$ and $\delta=v-u$, the independent event-normal differentials are
\[
 dX,\quad dY,\quad
 dR_m=d\delta-dX+t^2dY-t\,dZ,\quad
 dR_D=t^2dX-dY-d\delta-t\,dZ.
\]
Their rank is four; the remaining tangent is $t$. The core is zero
on the curve and has differential
\[
 dC=2t\{(1+t^2)dX-(1+t^2)dY-dR_m+dR_D\}.           \tag{NC11}
\]
Every coefficient here is nonzero. The other perfect matching is
$\{\lambda,c,\alpha,\beta\}$; with $W=1$ its equations give
$X=-1/v$, $Y=u$, $Z=uv$, $u^2v^2=1$. It is the symmetry image
of the first and has the same rank and differential properties.

\subsection{Curvature when the signed coefficient is visible}

All statements in this subsection are over $\mathbb R$ with $d>0$.
We prove that every physical zero of the full parent, or of its
restriction to at most three hard matching events, has strictly
positive Hessian in its remaining angle and projective-state
coordinates. This supplies precisely the soft Hessian in (NC8).

\paragraph{No hard event.}
At a zero the transverse state block is positive by (NC6). If the
active exponents in (NC4) have rank one, stationarity (NC5) requires
an opposite pair on one of the four lines through the Newton square.
A horizontal opposite pair combines $\lambda D$ or $\alpha m$
with $\beta n$ or $cA$; in the four cases it also activates
$\lambda n,cD,\beta m,\alpha A$, respectively, off that line.
The vertical case is identical: combining $\alpha A$ or $\lambda n$
with $\beta m$ or $cD$ activates $\alpha m,cA,\beta n,\lambda D$.
The two diagonal pairs force exactly the supports
$\{n,m,A,D\}$ and $\{\lambda,c,\alpha,\beta\}$.
Any further event activates an exponent off the diagonal.

Neither diagonal support is physical. At $u=v=1$ the core matrix is
\[
 \begin{pmatrix}
 2&2&-2&-2\\2&6&-2&2\\-2&-2&6&-2\\-2&2&-2&2
 \end{pmatrix}.                                      \tag{NC12}
\]
The corresponding vectors $(1,0,0,1)$ and $(1,-1,1,1)$ are
annihilated by all selected events, have zero core value, and have
nonzero core image. A positive semidefinite matrix cannot have zero
quadratic value and nonzero image at the same vector.

If no outer exponent is active, the event support contains neither
neighboring nor opposite pairs in (NC10); these are exactly the
twelve activating pairs in (SG6). Such a set extends to a triple
with cyclic gaps $(2,3,3)$. Indeed an independent four-set in the
cycle would be alternating and contain an opposite pair, and the
only possible independent triple gap patterns are $(2,3,3)$ and
$(2,2,4)$, of which the latter contains an opposite pair. All the
allowed triples are cyclic images of $\{\lambda,m,c\}$.
The state
\[
 \psi_H=(uv,-u,u^2v,1+v^2+u^2v^2)
\]
annihilates this triple and has
$C(\psi_H)=-2u^2v^2D_uD_v<0$. This contradicts physicality.
Thus the active exponents span the plane.

\paragraph{One hard event: explicit coefficients.}
By symmetry take the hard event to be $\lambda$. Its kernel is
$W=-vX$. The constrained matrix in $(X,Y,Z)$ is positive definite
on $Z=0$, so it has two positive state pivots at a zero. Extracting
the coefficient of $\lambda$ from (NC2) gives its determinant as
$4D_uD_v^2UV L_\lambda(U,V)$, whose outer coefficients are
\[
\begin{array}{ll}
a_{00}=4\alpha t_7t_8,&
a_{01}=4(Dt_7t_6+2\alpha mt_8),\\
a_{02}=8mDt_6,&
a_{10}=4(2\alpha At_7+nt_8t_4),\\
a_{12}=8(\beta mt_6+cDt_2),&
a_{20}=8nAt_4,\\
a_{21}=8(\beta nt_2+cAt_4),&
a_{22}=8\beta ct_2.
\end{array}                                                    \tag{NC13}
\]
The factor follows directly by using first the normal coordinate
$W+vX$, whose packet is $4\lambda D_u(W+vX)^2$, and then the
three kernel coordinates. All $t_i$ in (NC13) are strictly positive.

The four rank-one possibilities are exhausted as follows. On the
horizontal line, off-line vanishing forces $\alpha=n=0$; its negative
end requires $D>0$, forcing $m=c=0$, so its positive end is absent.
On the vertical line, off-line vanishing gives $\alpha=D=0$; its
negative end requires $n>0$, forcing $\beta=0$, so its positive end
is absent. The slope-minus-one pair requires $m,D,n,A>0$, which
activates $a_{01}$ off the line. The slope-plus-one pair forces
$\alpha,\beta,c>0$ and $m=D=n=A=0$: it completes the perfect
matching containing the hard $\lambda$ event.

This last possibility is not nonnegative on the hard kernel. On the
matching curve every event and the core vanish, but (NC11) has a
nonzero derivative after only one event-normal direction is fixed.
All event squares have zero derivative there, so the constrained
parent changes sign. If no exponent is active, (NC13) gives
$\alpha=D=n=0$ and $\beta m=cA=\beta c=0$. Its support, including
the hard event, is contained in one of
\[
 \{\lambda,A,\beta\},\quad\{\lambda,m,c\},
 \quad\{\lambda,m,A\}.
\]
These are cyclic images of the negative-core triple above, so each
has a kernel witness contradicting constrained physicality. The
active exponents of $L_\lambda$ therefore span the plane.

\paragraph{Two hard events: exhaustive positivity table.}
Matching pairs have cyclic distances two, three or four, represented
by $\{\lambda,c\}$, $\{\lambda,m\}$ and
$\{\lambda,\alpha\}$. Extract the coefficient of the two hard
weights from (SG6). An adapted state basis shows that it is a
positive angle-dependent factor times the determinant of the
constrained two-by-two state form. After the monomial $UV$ is removed,
the exact positivity conditions of its eight outer coefficients are
the following. Juxtaposition means both entries are positive;
``or'' means at least one condition holds.
\[
\begin{array}{c|c|c|c}
 e&\{\lambda,m\}&\{\lambda,c\}&\{\lambda,\alpha\}\\\hline
(-1,-1)&\alpha&\alpha&\text{always}\\
(-1,0)&D\text{ or }\alpha&D\text{ or }\alpha m&D\text{ or }m\\
(-1,1)&D&Dm&Dm\\
(0,-1)&n\text{ or }A\alpha&n&A\text{ or }n\\
(0,1)&\beta\text{ or }Dc&D\text{ or }\beta m&Dc\text{ or }\beta m\\
(1,-1)&An&0&An\\
(1,0)&Ac\text{ or }\beta n&A&Ac\text{ or }\beta n\\
(1,1)&\beta c&\beta&\beta c
\end{array}                                                    \tag{NC14}
\]
This table is obtained by differentiating (SG6), not by a sign
approximation. For example the first two entries in the
$\{\lambda,m\}$ column, apart from the positive factor four, are
$2\alpha(2A+2D+2c+d)$ and
$2(2AD+2A\alpha+4D\alpha+2Dc+Dd+2\alpha c+\alpha d)$.

For $\{\lambda,m\}$, either diagonal pair activates the off-line
exponent $(-1,0)$. A horizontal pair combines $D$ or $\alpha$ with
$Ac$ or $\beta n$ and activates an off-line exponent in all four
cases. A vertical pair combines $n$ or $A\alpha$ with $\beta$ or
$Dc$ and does the same. Rank zero forces
$\alpha=D=n=\beta=0$ and $Ac=0$, so only $A$ or only $c$ may remain.

For $\{\lambda,c\}$ the slope-minus-one pair is impossible because
one end is identically zero. A horizontal pair requires $A>0$ and
$D>0$ or $\alpha m>0$, activating $(0,1)$ or $(-1,-1)$ off the
line. A vertical pair requires $n>0$ and either $D>0$ or
$\beta m>0$, activating $(-1,0)$ or $(1,1)$ off the line.
The slope-plus-one pair forces $\alpha,\beta>0$ and
$D=n=A=m=0$. Rank zero leaves only $m$.

For $\{\lambda,\alpha\}$ the exponent $(-1,-1)$ is always active.
A rank-one stationary support must be the slope-plus-one pair,
forcing $c,\beta>0$ and $D=m=A=n=0$. These last two rank-one
possibilities complete $\{\lambda,c,\alpha,\beta\}$. Imposing
only two hard events still leaves a nonzero core derivative in
(NC11), so the constrained parent changes sign on the matching
curve. The rank-zero possibilities are contained in
$\{\lambda,m,A\}$ or $\{\lambda,m,c\}$. At $u=v=1$, their kernel
vectors $(1,-1,0,2)$ and $(1,-1,1,3)$ have core values $-4,-8$.
Thus every obstruction is excluded by constrained physicality.
The active exponents again span the plane.

\paragraph{Three hard events: two complete Laurent restrictions.}
A three-edge matching has cyclic gaps $(2,2,4)$ or $(2,3,3)$;
rotations and reflections are transitive on each pattern. The two
orbits have representatives and kernel states
\[
 M_N=\{\lambda,m,A\},\quad\psi_N=(uv,-u,0,1+v^2);
 \qquad
 M_E=\{n,m,A\},\quad\psi_E=(u,0,0,v).
\]
Direct substitution into the entire parent gives
\begin{align}
 {F(\psi_N)\over2UV(1+V)}
 ={}&{2n(1+U)\over V}+2cU+{2D(1+V)\over U}
       +{2\alpha(1+U)(1+V)\over UV}+2\beta V-d,
                                                        \tag{NC15}\\
 {F(\psi_E)\over2UV}
 ={}&d(U/V-1)+{2\lambda(1+U)\over V}
       +{2cU(1+V)\over V}+{2D(U-V)^2\over UV}
       +{2\alpha(1+U)\over U}+2\beta(1+V).            \tag{NC16}
\end{align}
For (NC15), any positive $n,D$ or $\alpha$ supplies two independent
exponents. If they vanish, $2cU+2\beta V-d$ cannot have a stationary
zero: stationarity forces $c=\beta=0$ and the value is $-d<0$.
For (NC16), $d>0$ supplies $(1,-1)$. Rank less than two forces
$\lambda=c=\alpha=\beta=0$. The remaining function of $x=U/V$
has infimum
\[
 2\sqrt{2D(d+2D)}-d-4D<0,
 \qquad (d+4D)^2-8D(d+2D)=d^2>0.
\]
The formula also gives a negative infimum when $D=0$. This contradicts
physicality. Both constrained Laurent Hessians are therefore strict
at every zero.

\paragraph{From angular determinants to the full soft Hessian.}
For zero, one or two hard events, (NC6) supplies a positive state
block on the remaining projective-state coordinates. In a fixed real
kernel chart write the constrained state form as
$a+2b^Ty+y^TC_0y$, with $C_0\succ0$ at contact. Its state Schur
complement is $a-b^TC_0^{-1}b$, and its determinant equals this
scalar times $\det C_0$. At a stationary zero, differentiating a
positive prefactor does not change the Hessian sign. Equations
(NC5), (NC12)--(NC14) therefore give a positive angular Schur
complement and hence the full positive soft Hessian. For three hard
events only the two angles remain and (NC15)--(NC16) apply directly.
All these matrix inverses are ordinary real residue calculations.
The exact certificate itself is (NC7)--(NC9), with no tensor inverse.

\subsection{Four matching weights, with no comparability assumption}

Suppose the four positive weights of a perfect matching dominate
$d$ and every complementary event. By symmetry use
$\{n,m,A,D\}$ and arrange that $D=\mu=\min(n,m,A,D)$.
Set
\[
 \eta_1=\mu/n,\quad\eta_2=\mu/A,\quad\eta_3=\mu/m,
 \qquad0<\eta_i\le1.
\]
Their residues may vanish independently. Work near the matching
curve in $W=1$ and introduce the exact state coordinates
\[
 R=v-X+uvY-uZ,\qquad
 Z={v-X+uvY-R\over u},\qquad\delta=v-u,
\]
using only the testing unit $u^{-1}$. The fourth event is exactly
$f+L$, where
\[
 f={u^2-v^2\over u}=-2\delta-\delta^2/u,\qquad
 L=v(u+u^{-1})X-(1+v^2)Y+(v/u)R.                    \tag{NC17}
\]
The first three hard packets are
$4nD_uD_vX^2+4AD_uD_vY^2+4mR^2$.
On $X=Y=R=0$ the core is exactly $2uf$, so
\[
 4\mu f^2+2duf
 =4\mu\left(f+{du\over4\mu}\right)^2
                         -{d^2u^2\over4\mu}.         \tag{NC18}
\]
Let $\tau=\max(\lambda,c,\alpha,\beta)$. Testing (NC18) at
$u=1$, $v=\sqrt{1+d/(4\mu)}$ proves
\[
 \tau>0,\qquad \rho={d^2\over\mu\tau}\text{ bounded}.
 \tag{NC19}
\]
Indeed all complementary amplitudes are bounded at this source-field
test; if their scale were infinitesimal relative to $d^2/\mu$, the
complete parent would be negative. The case $d=0$ was already an
event-square sum. Put $\varepsilon=d/\mu$, which is infinitesimal.

If $\bar\rho=0$, use the five rows
\[
 (\sqrt{n/\tau}X,\sqrt{A/\tau}Y,
       \sqrt{m/\tau}R,\sqrt{\mu/\tau}\delta,1).
\]
Their four-normal residue is, with $B_0=1+u_0^2$,
\[
4B_0^2z_1^2+4B_0^2z_2^2+4z_3^2
 +4(-2z_4+B_0\sqrt{\bar\eta_1}z_1
       -B_0\sqrt{\bar\eta_2}z_2+\sqrt{\bar\eta_3}z_3)^2,
 \tag{NC20}
\]
which is positive definite. The constant residue is
$4(1+u_0^2)(P+Qu_0^2)>0$, where
$P=\overline{(\lambda+\alpha)/\tau}$,
$Q=\overline{(c+\beta)/\tau}$ and $P+Q>0$.
The signed linear coefficients are $\sqrt{\rho\eta_i}$ and
$\sqrt\rho$, so vanish. All other normalized coefficients are
bounded by $\varepsilon,\eta_i,\tau/\mu$ and bounded testing
units. This is an exact strict bounded Gram of the whole parent.

Now suppose $\bar\rho>0$. In the source-field stationary equations
put
\[
 X=\varepsilon\eta_1s_1,\quad
 Y=\varepsilon\eta_2s_2,\quad
 R=\varepsilon\eta_3s_3,\quad\delta=\varepsilon r,
 \quad u=t.                                          \tag{NC21}
\]
Divide the three state equations and the $\delta$ equation by $d$,
and the tangent equation by $\tau$. All coefficients are bounded
after clearing powers of the unit $t$: the factors $\eta_i$ cancel
the large weights $\mu/\eta_i$, and
$\tau=\mu\varepsilon^2/\rho$. Put, in residue,
\[
 B_t=1+t^2,\quad\ell=(B_t,-B_t,1),\quad
 b=(B_t^2,B_t^2,1),\quad
 k=-2r+\sum_i\bar\eta_i\ell_i s_i.
\]
Direct differentiation of (NC1), (NC17) gives the four normal equations
\[
 8b_i s_i+8k\ell_i+(4tB_t,-4tB_t,0)_i=0,
 \qquad-16k-4t=0.                                   \tag{NC22}
\]
Eliminating the last row from the first three gives diagonal pivots
$8b_i$, with the last pivot $32$. The determinant is
$32\cdot8^3B_t^4>0$, even if some $\bar\eta_i=0$.
Their solution is
\[
 s_1^*=-{t\over4B_t},\quad s_2^*={t\over4B_t},\quad
 s_3^*={t\over4},\quad
 r^*={t\over8}(1-\bar\eta_1-\bar\eta_2+\bar\eta_3).
 \tag{NC23}
\]
The complete normalized residue objective is
\[
\begin{split}
 \Psi={}&\bar\rho\{4\sum_i\bar\eta_i b_i s_i^2+4k^2
           +4tB_t(\bar\eta_1s_1-\bar\eta_2s_2)-4tr\}\\
 &+4(1+t^2)(P+Qt^2).
\end{split}
\]
Completing its four normal squares by (NC23) leaves exactly
\[
 g(t)=4(1+t^2)(P+Qt^2)
       -{\bar\rho t^2\over4}
          (1+\bar\eta_1+\bar\eta_2+\bar\eta_3).     \tag{NC24}
\]
This completion is polynomial in the $\bar\eta_i$ and remains
valid when one vanishes; its scaled equation in (NC22) still has
the displayed positive pivot. Substituting the corresponding normal
values into the physical source and taking residue proves $g(t)\ge0$
for every real $t\ne0$.

If $g(t_0)>0$, shifting by (NC23) gives the positive normal form
(NC20) and a positive constant, hence a strict bounded Gram. At a
zero, divide (NC24) by $4t^2$:
\[
 P(1+t^{-2})+Q(1+t^2)
       -{\bar\rho\over16}(1+\bar\eta_1+\bar\eta_2+\bar\eta_3).
\]
Stationarity gives $P/t_0^2=Qt_0^2$. Since $P+Q>0$, both are
positive; the logarithmic second derivative is positive, and thus
$g''(t_0)>0$. Append the tangent equation to (NC22). Its Schur
complement is $g''(t_0)$, so its real Jacobian is invertible. The
semialgebraic implicit-function argument gives an exact source
stationary center of the form (NC21), with bounded $s_c,r_c,t_c$.
Its normalized value is nonnegative by physicality.

The exact full Taylor rows are
\[
\big(\sqrt{n/\tau}(X-X_c),\sqrt{A/\tau}(Y-Y_c),
\sqrt{m/\tau}(R-R_c),\sqrt{\mu/\tau}(\delta-\delta_c),t-t_c\big).
 \tag{NC25}
\]
Here is the required coefficient control, including all normal terms.
For $a_i=\mu/\eta_i$ the two center factors in a hard packet are
\[
 {a_ix_{ci}^2\over\tau}=\rho\eta_i s_{ci}^2,
 \qquad {a_ix_{ci}\over\sqrt{a_i\tau}}
                    =\sqrt{\rho\eta_i}\,s_{ci}.      \tag{NC26}
\]
The fourth packet has the analogous factors $\rho r_c^2$ and
$\sqrt\rho r_c$. The core belongs to the ideal $(X,Y,R,\delta)$
because it vanishes identically on the matching curve. Its center-linear
factors after division by $\tau$ are therefore
$d x_{ci}/\tau=\rho\eta_i s_{ci}$ and
$d\delta_c/\tau=\rho r_c$; its normalized normal factors are
$d/\sqrt{a_i\tau}=\sqrt{\rho\eta_i}$ and
$d/\sqrt{\mu\tau}=\sqrt\rho$. Its quadratic normal factors are
bounded by $d/\mu=\varepsilon$. Complementary coefficients are at
most $\tau$ and give smaller factors. Each packet is at most quadratic in the
state coordinates, and all remaining higher Taylor factors are
actual angle or state deviations with zero residue. The only testing
denominators are powers of the unit $t$. Thus exact packetwise Taylor
allocation is a bounded Gram on (NC25). Its normal residue is (NC20)
and its tangent Schur complement is $g''(t_0)/2>0$, as follows by
differentiating the completed squares (NC23)--(NC24). This assertion
also holds when $\bar\eta_i=0$: the normal form (NC20), its cross
entries and the completed Schur expression are continuous polynomial
expressions in $\sqrt{\bar\eta_i}$, with a strict normal block.
The strict-Gram lemma, the nonnegative center value and pure-factor
restoration prove the four-event case without any restriction
on its internal coefficient ratios.

\subsection{A terminating induction on event constraints}

Normalize the nonzero source by $d+\lambda+n+m+c+A+D+\alpha+\beta=1$.
The factor is positive and pure-source. Literal $d=0$ is already an
event-square sum; positive slack residue is covered by the positive
residue lemma. If $\bar d>0$ and the residue slack is zero, the
unconstrained curvature argument and (NC9) with $k=0$ apply.

It remains that $d>0$, $\bar d=0$ and the residue slack is zero.
Product residue states were settled above. At an entangled residue,
let $S$ be all events with positive coefficient residue. This is a
nonempty matching. If it has four edges, every complementary
coefficient and $d$ is infinitesimal relative to its smallest weight,
and (NC17)--(NC26) apply. Otherwise use its at most three independent
event amplitudes as the exact hard state coordinates. Set
\[
 \sigma=\max\bigl(\{d\}\cup\{c_e:e\notin S\}\bigr)>0.
\]
Every hard weight $a_i$ satisfies $\sigma/a_i$ infinitesimal. The
parent has the diagonal form (NC7), and its constrained next-scale
residue $g=\overline G(0,y)$ is nonnegative by physicality.

If $g(y_0)>0$, (NC9)'s positive-value argument finishes. If $g(y_0)=0$
and $\overline{d/\sigma}>0$, the one-, two- or three-hard curvature
calculation above gives a strict full soft Hessian; (NC8)--(NC9)
finish without assuming comparability of the hard weights.

The only remaining case has $g(y_0)=0$ and
$\overline{d/\sigma}=0$. Then $g$ is a nonnegative event-square
sum, and at least one remaining event coefficient has positive
normalized residue. All such event amplitudes vanish at the same
entangled residue state. Add them all to $S$. The enlarged set is
still a matching; every coefficient left outside it is infinitesimal
relative to each of its weights. If it has at most three edges, repeat.
If it has four, apply the four-event proof above, which allows arbitrary source scales.
Each repetition adds an edge, and a matching has at most four edges.
There are therefore at most three enlargements. Every step retains
the entire original parent. This proves
Theorem~\ref{thm:noncommuting-refined}.




\chapter{Assembly of the local certificates}
\label{ch:cover}
\label{component:pe1088004af65}
\label{pe1088004af65:sec:complete-source-routing}

The preceding chapters reduce every valid inequality to events, the
closed faces, and the two contact families C1 and C2. Events and closed
faces already have separated certificates. We now prove the same
assertion for the contact families.

\begin{theorem}[Canonical certificate]
For every source in either family SG10 with an actual finite
noncommuting stationary zero, the complete homogeneous state polynomial
is a finite sum of squares in $K_c\otimes_{\mathbb R}K_t$ at every
independent projective testing pair and every real or complex state.
\end{theorem}

Here and below, a source coefficient times a testing expression denotes
its separated tensor product. With $D_u=1+u^2$, $D_v=1+v^2$, the full
polynomial to be certified is
\[
\begin{split}
 F={}&4\lambda D_u(W+vX)^2+4nD_uD_vX^2
       +4m(vW-X+uvY-uZ)^2\\
 &+4CD_v(uW-Y)^2+4AD_uD_vY^2
       +(4D+2d)(uW+uvX-Y-vZ)^2\\
 &+(4\alpha+2d)D_u(vY-Z)^2
       +(4\beta+2d)D_v(X+uZ)^2-2dD_uD_vZ^2.
\end{split}\tag{LC2}
\]
Its homogeneous form defines the literal projective endpoints.

\begin{proof}
Normalize $d=1$. Every original coefficient can be restored by
multiplying the roots by $\sqrt d\otimes1$. The source parameter
\[
 \delta_{\rm tot}=\frac d{d+\lambda+n+m+C+A+D+\alpha+\beta}
 \tag{LC3}
\]
separates two possibilities. If it is a positive unit, all coefficients
after normalization are bounded. The proof in
Chapter~\ref{ch:cert-2} covers this domain. Otherwise
$\delta_{\rm tot}$ is a positive infinitesimal, and the contact
coordinates from Chapter~\ref{ch:entry} give the source partitions
of Chapters~\ref{ch:c1-short} and \ref{ch:c2-short}.

For reference, the C1 coordinates satisfy
\[
\begin{gathered}
 0<k\le1,\quad \eta,\zeta>0,\quad x=k^2,\\
 r=1-a\eta\in(0,1),\quad\theta=t\zeta\in(0,1),\quad h=1-\theta,\\
 0<a,t<2,\quad E=1-x\eta h>0,\quad Q=1-x\zeta r>0,\\
 M=1+xr-t>0,\quad J=1+xh-a>0,\quad
 e=k\sqrt{\eta/\zeta}>0.
\end{gathered}\tag{LC4}
\]
The C1 partition first separates bounded and infinitesimal $k$, then
the relative source scales $\eta,\zeta,e$. When $\zeta$ is
noninfinitesimal, the product $L=x\zeta$ gives the remaining three
alternatives. Chapter~\ref{ch:c1-short} proves their source inequalities
and constructs the corresponding full quadratic forms. The literal
faces $D=0$ and $\beta=0$ are included wherever the stationary equations
permit them.

For C2, the reflection
$(u,v;W,X,Y,Z)\mapsto(1/v,1/u;Z,X,Y,W)$ exchanges
$(m,D)$ and $(C,\alpha)$. Chapter~\ref{ch:c2-short} uses it to reduce
the source angle types. The case of two weak complementary events
has a separate radial calculation. With one weak event, the directed
source partition distinguishes bounded and unbounded ratios; the
bounded ratios share one full matrix completion. The remaining
unbounded cases and every $n=0$ boundary are proved there. All these
partitions concern coordinates in $K_c$ alone.

Each local identity certifies the full state quadratic form, including
its two first-normal terms and its state Hessian. Positive source
excesses are retained as whole squares. The remaining coefficient
matrix is bounded and has positive-definite real residue, so the
lemma of Chapter~\ref{ch:tensor} gives a finite tensor sum of squares.
The explicit complementary charts cover all independent testing
angles and state ratios.

The sign changes
$(u,Y,Z)\mapsto(-u,-Y,-Z)$ and
$(v,X,Z)\mapsto(-v,-X,-Z)$ preserve LC2. Its homogeneous identities
include zero and infinite half-angles. State normalizations use a
nonzero testing coordinate and are restored by its square. Finally,
real and imaginary state parts add, since the measurements are real.
This proves the theorem.
\end{proof}

To finish the global argument, write an arbitrary valid source as a
nonnegative combination of at most nine extreme sources, and an
arbitrary testing behaviour as a mixture of at most nine real-qubit
behaviours. Events, closed faces and the theorem above certify every
atom pair. Equation M2 combines their squares and preserves the affine
constant. This proves the separated-slack theorem, and hence the finite
semidefinite representation.

\appendix
\chapter{Bounded source coefficients}
\label{ch:cert-2}

\section{Two quartic contacts}
\label{component:p4ec1a51262f9}
\label{p4ec1a51262f9:a-full-parent-case-for-a-finite-positive-d-commuting-edge}

When one measurement pair becomes commuting, a quadratic angular
coefficient can vanish. The next coefficient is then quartic. We treat
the two such contacts together: the source determines a critical centre,
and a finite polynomial expansion about that centre gives the required
Gram matrix. All coefficients in this section are bounded, $d$ is a
positive unit, and $\lambda=0$. The source is physical.
Write $(n,m,C,A,D)$ for the coefficients denoted
$(e_1,e_2,e_3,e_4,e_5)$ in the event order of SG1.

\paragraph{A finite commuting edge.}
Suppose the testing residues are $u=0$, $v=v_0\ne0$ and
$[W:X:Y:Z]=[1:0:0:0]$, and $m$ is infinitesimal.
On the residue source the determinant numerator is
\[
 Q(U,V)=UE(V)+U^2B(V),\quad
 E(V)=q_{10}+hV+q_{12}V^2,\quad
 B(V)=q_{20}+q_{21}V+q_{22}V^2.
\]
Assume $q_{10},q_{12},B(v_0^2)>0$ and $E(v_0^2)=0$.
The unit event amplitude $vW-X+uvY-uZ$ can be set to one,
so $W=(1+X-uvY+uZ)/v$. The opened event contributes the
constant $4m$. The state block at $u=0$ is
\[
 M_0(v)=\begin{pmatrix}
 2(d+2n+2\beta)(1+v^2)&0&0\\
 0&2d(1+v^2)+4(C+A)(1+v^2)+4D+4\alpha v^2&4v(D-\alpha)\\
 0&4v(D-\alpha)&4\alpha+4Dv^2
 \end{pmatrix}.
\]
Assume $M_0(v_0)\succ0$. In the applications below its lower
minor is positive because
$\alpha d+2\alpha C+2\alpha A+2\alpha D(1+v_0^2)
+Dv_0^2(d+2C+2A)>0$.

The substitution $X=u^2x$, $Y=ua$, $Z=ub$, $U=u^2$ defines
the polynomial $K$ by the exact identity
\[
 F(u^2x,ua,ub,u,v)=4m+u^2K(x,a,b,v,U).
 \tag{FE1}
\]
Its coefficients are Laurent polynomials in $v$. Substitution in
SG2 shows that $K_x=UE_x$ and
\[
 (E_x)_x\big|_{U=0}=4(d+2n+2\beta)(1+v^2)>0.
 \tag{FE2}
\]
The residue of $K$ at $U=0$ is independent of $x$.
Minimizing its positive $(a,b)$ block leaves
\[
 r(v)=\frac{16(1+v^2)^2E(v^2)}{v^2\det M_0(v)}.
\]
Since $E\ge0$ and has a zero at $v_0^2>0$,
$E(V)=q_{12}(V-v_0^2)^2$. Thus the Hessian in $(a,b,v)$
is positive definite at the residue zero. Equation FE2 then
determines $x$. The equations
$\mathcal E=(E_x,K_a,K_b,K_v)=0$ have an invertible
Jacobian in $q=(x,a,b,v)$. They define a source-field branch
$q_*(c,U)$. Its minimized profile $k(c,U)=K(q_*(c,U),U)$ satisfies
\[
 k(c_0,0)=0,\qquad
 \gamma=k_U(c_0,0)=
 \frac{16(1+v_0^2)^2B(v_0^2)}{v_0^2\det M_0(v_0)}>0.
 \tag{FE3}
\]
The last equality follows by taking the Schur complement of the
three-state block in the event-normalized parent. Its determinant
is $\det P/v^2$; the coefficients of $U$ and $U^2$ give
$r(v)$ and FE3 respectively.

\paragraph{A finite-ratio corner.}
\label{component:p03b551f2a9a0}
\label{p03b551f2a9a0:a-positive-d-case-at-a-finite-ratio-distance-four-corner}
Now let $\alpha$ be infinitesimal. At the residue source suppose
\[
 Q(U,V)=q_{20}U^2+hUV+q_{02}V^2+
 q_{21}U^2V+q_{12}UV^2+q_{22}U^2V^2,
\]
where $q_{20},q_{02}>0$. Choose a nonzero real $t_0$ such that
\[
 q_{20}t_0^4+ht_0^2+q_{02}=0,
 \qquad q_{21}t_0^4+q_{12}t_0^2>0.
\]
Normalize the unit amplitude $vY-Z$ to $-1$, giving $Z=1+vY$.
The substitution $X=vx$, $Y=vy$, $u=vt$, $W=w$, $V=v^2$
defines a polynomial $K$ by
\[
 F=4\alpha+v^2K(x,y,t,w,V).
 \tag{FR1}
\]
Its constant coefficient in $V$ is
\[
\begin{split}
 K_0={}&4nx^2+4m(w-x-t)^2+4C(tw-y)^2+4Ay^2\\
 &+4D(tw-y-1)^2+4\beta(x+t)^2\\
 &+2d\{(tw-y)^2+(x+t)^2-2tw\}+4\alpha t^2.
\end{split}
\]
Assume the residue quadratic form at $t_0$ has kernel
$(w_0,x_0,y_0,1)$ and a positive $(w,x,y)$ principal block
$M_J(t_0)$. Its minimized scalar is
$16(q_{20}t^4+ht^2+q_{02})/\det M_J(t)$.
Physicality makes the numerator $q_{20}(t^2-t_0^2)^2$.
The full Hessian in $q=(x,y,t,w)$ is therefore positive definite.
The source equations $\nabla_qK=0$ define $q_*(c,V)$, and
\[
 \gamma=\partial_VK(q_0,0)=
 \frac{16(q_{21}t_0^4+q_{12}t_0^2)}{\det M_J(t_0)}>0.
 \tag{FR2}
\]
This is again the next coefficient of the full determinant after
the positive state block is minimized. The testing chart has
$u,v$ infinitesimal and $[W:X:Y:Z]$ of residue $[w_0:0:0:1]$.
The actual testing ratio $u/v$ is unrestricted.

\paragraph{The common choice of centre.}
In either case write $S=U$ or $V$, and put $\kappa=k(c,0)$.
If $\kappa\ge0$, choose $S_c=0$ and $q_c=q_*(c,0)$.
If $\kappa<0$, solve
\[
 \mathcal E(q_c,S_c)=0,\qquad
 K(q_c,S_c)+S_cK_S(q_c,S_c)=0,
 \tag{FC1}
\]
where $\mathcal E=\nabla_qK$ in the corner case. The Jacobian
has the invertible normal block above and last pivot $2\gamma>0$.
The semialgebraic implicit-function theorem and real closed transfer
give a unique nearby source solution. The function
$k(c,S)+Sk_S(c,S)$ increases from the negative value $\kappa$,
so $S_c>0$. Realizing the corresponding source measurement and
state shows that
\[
 \mu=c_*+S_cK(q_c,S_c)\ge0,
 \qquad c_*=4m\ \hbox{or}\ 4\alpha.
 \tag{FC2}
\]
Every use of the implicit-function theorem here is within the
source field. The tensor identity uses only the resulting source
coordinates and finite polynomial divisions.

\paragraph{Full state expansion.}
At the edge set
\[
 N_s=(X-u^2x_c,Y-ua_c,Z-ub_c),\qquad N_v=u(v-v_c).
\]
Quadratic state dependence gives exactly
\[
 F=4m+UK(x_c,a_c,b_c,v,U)
       +UE_xN_x+uK_aN_y+uK_bN_z+N_s^TMN_s.
\]
The centre equations make each derivative vanish at $(v_c,U_c)$.
Difference division in $v$ and $U$ therefore gives only the
cross terms
\[
 N_xU(U-U_c),\quad N_yu(U-U_c),\quad N_zu(U-U_c),
 \quad N_xuN_v,\quad N_yN_v,\quad N_zN_v,
\]
and $N_v^2$ and $N_vu(U-U_c)$ from the scalar profile.
All coefficients are bounded. For negative powers of the unit
$v$, use $(v^{-j}-v_c^{-j})/(v-v_c)
=-\sum_{i=0}^{j-1}v^{-i-1}v_c^{-(j-i)}$.

At the corner set
\[
 N=(X-vx_c,Y-vy_c,u-vt_c,v(W-w_c)),\qquad z=W-w_c.
\]
The linear part is $v\nabla_qK(q_c,V)\cdot N$.
The rest belongs to $(N_1,N_2,N_3,vz)^2$: in SG2, a term
containing one $z$ and only one of the first three normals carries
$v$, and a term containing $z^2$ and no such normal carries $v^2$.
Terms with two of the first three normals already belong to this
ideal, with bounded $z$ left in their coefficients. This gives the
whole quadratic remainder without dividing by $v$.

In both cases the derivative coefficients are divisible by $S-S_c$.
Their cross coefficients with that last generator are infinitesimal.
The fixed-centre scalar profile has a double zero at $S_c$ when
$\kappa<0$. Finite division thus yields
\[
 \begin{array}{ll}
 \kappa\ge0:&F=c_*+\kappa S+G^THG,\quad G=(N,S),\\
 \kappa<0:&F=\mu+G^THG,\quad G=(N,S-S_c).
 \end{array}
 \tag{FC3}
\]
At the edge $N=(N_s,N_v)$ and the normal residue block is the
positive $X$ entry together with half the positive $(a,b,v)$
Hessian. At the corner it is half the positive four-variable
Hessian. The last diagonal is $\gamma>0$, and its cross entries
vanish in residue. The strict bounded-Gram lemma applies.
The remaining scalars in FC3 are nonnegative in the source field;
$S$ is a testing square. Restoring the event normalization completes
both certificates, including arbitrary relative infinitesimal rates.

For the edge with $E(v_0^2)>0$, the four-normal quadratic form in
$(X,Y,Z,u)$ is already strict, with $4m$ retained separately.
This supplies the other finite-edge alternative.


\section{\texorpdfstring{Positive total scale: the finite-pitchfork case}{Positive total scale: the finite-pitchfork case}}
\label{component:pfb0cf454f653}


The finite-edge calculation applies whenever $d,C,A,D,\alpha$
are positive units. We establish its hypotheses and then treat the
remaining commuting and projective limits.

\begin{theorem}
Let the original normalized C1 or C2 source be bounded, nonnegative and
physical over a real closed coefficient extension of $\mathbb R$.
Assume $d,C,A,D,\alpha$ are positive bounded units. Retain its exact
canonical zero coefficients, in particular $\lambda=0$, and every
other nonnegative original coefficient $n,m,\beta$. Then its entire
parent is tensor-SOS at every independent projective testing pair and
every real or complex testing state.
\end{theorem}

Noncommuting testing residues are covered by
Theorem~\ref{thm:noncommuting-refined}.

\subsection*{The finite edge}

The only possible zero at $u=0$, $v=v_0\ne0$ is the $W$ state
ray when $m$ is infinitesimal. We verify the conditions of
Section~\ref{component:p4ec1a51262f9} on that domain.

At $\lambda=m=0$, the exact determinant polynomial is
\[
 Q(U,V)=U E(V)+U^2 B(V),\quad
 E(V)=q_{10}+q_{11}V+q_{12}V^2,\quad
 B(V)=q_{20}+q_{21}V+q_{22}V^2.
\]
Use the positive quantities $t_1,\ldots,t_8$ in SG6.

All eight have positive bounded-unit residues. At the residue $m=0$,
the required coefficients are
\[
 q_{10}=4\alpha A t_7t_3>0,\quad
 q_{12}=4CDt_5t_2>0,\quad
 q_{21}=4(\beta n t_3t_2+CA t_1t_4)>0,
\]
and $q_{20},q_{22}\geq0$. Therefore $B(v_0^2)>0$ for every
$v_0\ne0$. Physicality implies $E(V)\geq0$ for all real $V>0$:
divide $Q(U,V)\geq0$ by $U>0$ and let the real variable $U$ tend to
zero. Consequently a zero $E(v_0^2)=0$ is a double zero,
$E(V)=q_{12}(V-v_0^2)^2$, and otherwise $E(v_0^2)>0$.

The state complement $M_0(v_0)$ in
Section~\ref{component:p4ec1a51262f9} is positive definite: its
lower determinant has the positive factor
$\alpha d+2\alpha C+2\alpha A+2\alpha D(1+v_0^2)
+Dv_0^2(d+2C+2A)$. Thus that section proves the full state
certificate at this edge, including both possible signs of the
optimized source value.



\subsection*{The other measurement limits}

Noncommuting residue tests use A6. It remains that at least one angle
has residue zero or infinity. For finite $u_0\ne0$ and $v_0=0$,
the real residue parent is a sum of nonnegative event terms and the
core $2d[(uW-Y)^2+(X+uZ)^2]$. The unit $A,\alpha,C$ terms
force $Y=Z=W=0$, and the core then forces $X=0$.
Hence the residue state matrix is positive definite.

For $u_0=0$ and finite $v_0\ne0$, the unit $C,A,D,\alpha$ terms
and $d>0$ leave only the $W$ ray if $m$ is weak, while unit $m$
removes that ray too. In the weak case, positive residue slack is
immediate; at the $W$ ray, the previous section applies for either
$E(v_0^2)>0$ or $E(v_0^2)=0$.

For $u=\infty$ and finite nonzero $v$, set $p=1/u$ and homogenize
as $F_p=p^2F(1/p,v)$. Its real residue has core
\[
 2d[(W+vX)^2+(vY-Z)^2]
\]
and positive $4C(1+v^2)W^2+4A(1+v^2)Y^2$ squares. These
force every state coordinate to vanish at a zero, so the matrix is
positive definite. For $v=\infty$ and finite nonzero $u$, the
homogenized residue has core $2d(1+u^2)(X^2+Y^2)$ together with
$4C(uW-Y)^2+4D(uX-Z)^2$. Again it is positive definite.
These are exact residue identities and use no density argument.

It remains to cover the four double corners. The following displayed
quadratic forms are complete, with coefficient letters interpreted in
real residue. All coefficients that are merely nonnegative remain so.
At a unit state coordinate, normalize it in the testing field.

At $(u,v)=(0,0)$ the only possible zero state is the $W$ ray.
Put $W=1$ and retain the entire $m$ square. The remainder lies
exactly in $(u,X,Y,Z)^2$, and its residue quadratic is
\[
 2(d+2n+2\beta)X^2+(4C+4D+2d)(u-Y)^2+4AY^2+4\alpha Z^2.
\]
It is strict. The testing $v$ is bounded infinitesimal in its
coefficients and is never divided out.

At $(u,v)=(\infty,0)$ the only possible zero state is the $X$ ray
when $n$ is weak; if $n$ is a unit there is no zero. Put $X=1$,
retain the entire $n$ square in $F_p$, and use roots $(p,v,W,Y,Z)$.
The whole remainder belongs to their square ideal and its residue is
\[
 4CW^2+4AY^2+(4D+2d)(W+v)^2
                 +4\alpha Z^2+[4(m+\beta)+2d](p+Z)^2\succ0.
\]

At $(u,v)=(0,\infty)$ the only possible zero state is the $W$ ray
when $m$ is weak. Put $q=1/v$, $W=1$, and retain the entire $m$
square in $F_q=q^2F(u,1/q)$. The whole remainder is in
$(u,X,Y,Z)^2$ and its residue quadratic is
\[
 (4n+4\beta+2d)X^2+4C(u-Y)^2+[4(A+\alpha)+2d]Y^2+4DZ^2\succ0.
\]
The testing $q$ remains a bounded infinitesimal coefficient.

At $(u,v)=(\infty,\infty)$ the only possible zero state is the $Z$
ray when $\beta$ is weak. Put $Z=1$ and retain the entire $\beta$
square in $F_{pq}=p^2q^2F(1/p,1/q)$. The complete remainder lies
in $(p,q,W,X,Y)^2$ and its residue quadratic is
\[
 4CW^2+4AY^2+4D(X-p)^2+(2d+4n)X^2
                            +[2d+4(m+\alpha)](Y-q)^2\succ0.
\]

For each corner, allocate two normal factors from every monomial to
the Gram roots; every remaining factor is bounded. Terms above
quadratic order have infinitesimal coefficients. The displayed full
quadratic residues are strict, so the separated strict-Gram lemma
applies. Restore each retained square. Homogeneous angle factors and
state normalizations are pure testing squares. The identities include
literal zero and infinity and every original nonnegative opening rate.
For nonzero state residues away from the listed rays, positivity gives
the scalar strict-residue exit. The zero state is immediate and complex
states split into real and imaginary evaluations. This proves the theorem.



\section{\texorpdfstring{Bounded C2 sources with infinitesimal alpha}{Bounded C2 sources with infinitesimal alpha}}
\label{component:pf31e38136e03}



\begin{theorem}
Let the original bounded normalized C2 source be physical, with
$\lambda=\beta=0$, all event coefficients nonnegative,
$d,n,m,C,A,D$ positive bounded units, and $\alpha$ weak, meaning
$\bar\alpha=0$. Then the complete original parent is tensor-SOS
for every independent projective testing pair and every real or
complex testing state.
\end{theorem}

The finite-ratio calculation of
Section~\ref{component:p03b551f2a9a0} treats the remaining corner.
We first identify the zero states to which it applies.

\subsection*{Only the zero-zero corner needs a new argument}

A6 covers noncommuting testing residues. On finite $u=0$, $v\ne0$,
the unit $m$ removes the $W$ ray, and the $(X,Y,Z)$ complement is
strict because $d,D,C,A$ are units and $v\ne0$. On finite $v=0$,
$u\ne0$, the units $n,A$ force $X=Y=0$ at a zero; the $C$ square
forces $W=0$, and the signed core $(X+uZ)^2$ forces $Z=0$.

The exact single-infinity core identities in
Section~\ref{component:pfb0cf454f653} require no unit $\alpha$: at
$u=\infty$, $v\ne0$ use
$2d[(W+vX)^2+(vY-Z)^2]$ and the $C,A$ squares; at $v=\infty$,
$u\ne0$ use $2d(1+u^2)(X^2+Y^2)$ and the $C,D$ squares.
Both state matrices are strictly positive.

At $(\infty,0)$ the unit $n$ removes the only possible $X$ zero
ray, and at $(0,\infty)$ the unit $m$ removes the $W$ ray.
At $(\infty,\infty)$, put $p=1/u,q=1/v$, normalize the possible
$Z$ zero ray to $Z=1$, and use $(p,q,W,X,Y)$ as roots. The exact
full-parent square-ideal allocation given in
Section~\ref{component:pfb0cf454f653} has residue
\[
 4CW^2+4AY^2+4D(X-p)^2+(2d+4n)X^2
                         +[2d+4(m+\alpha)](Y-q)^2\succ0.
\]
This form remains strict for weak $\alpha$. Thus every projective
test except $u,v$ infinitesimal is already covered, including literal
zero/infinite endpoints and all signs.

\subsection*{The exact two state charts at zero-zero}

Write $F$ for the original parent, with $\lambda=\beta=0$, and retain
the whole nonnegative $\alpha$ square
$4\alpha(1+u^2)(vY-Z)^2$ whenever using an ordinary strict jet.
At $u=v=0$ the residue is strictly positive on the $(X,Y)$ state
subspace. Hence $X$- or $Y$-unit states have positive residue slack.
For all other states, normalize a maximal coordinate. Then either
$W=1$ with $Z$ infinitesimal, or $Z=1$ with $W$ bounded.

In the first chart the entire remaining parent belongs to
$(X,Y,u,v)^2$. Its quadratic residue is
\begin{equation}
 (4n+2d)X^2+4m(v-X)^2+[4(C+D)+2d](u-Y)^2+4AY^2\succ0. \tag{1}
\end{equation}
Consequently the finite normal allocation and restoration of the
complete $\alpha$ square give tensor-SOS. No testing ratio is used.

In the second chart, set $w=\bar W\in\mathbb R$. The full residue
quadratic jet on $(X,Y,u,v)$ is
\begin{align}
 Q_w={}&4nX^2+4m(vw-X-u)^2+4C(uw-Y)^2+4AY^2\notag\\
      &+4D(uw-Y-v)^2
        +2d[(uw-Y)^2+(X+u)^2-2uvw].                    \tag{2}
\end{align}
All coefficient letters in (1)--(2) now denote their real residues.
Physicality implies $Q_w\succeq0$: evaluate the physical residue
parent at $(X,Y,u,v)=s(x,y,a,b)$ and let the real $s$ tend to zero.
The $(X,Y)$ principal block is strictly positive, already from the
$nX^2$ and $AY^2$ terms. If $Q_w$ is strict, the same ordinary
four-normal allocation applies, with $W-w$ infinitesimal in its
coefficients, and the whole $\alpha$ square restored.

Suppose instead that $Q_w$ is singular. It has no nonzero kernel
vector with $v=0$: the nonnegative sum at $v=0$ includes
$4nX^2+(4m+2d)(X+u)^2+4AY^2$, so $X=u=Y=0$.
It also has no nonzero kernel vector with $u=0$: the sum at $u=0$
includes $4nX^2+4AY^2+4D(Y+v)^2$, so $X=Y=v=0$.
Its real angular Schur complement therefore has positive diagonal
entries. Being singular and positive semidefinite, it has rank one
and a kernel with finite nonzero ratio $t_0=u/v$. Normalize a real
kernel vector by $v=1$ and write it as $(x_0,y_0,t_0,1)$.
This constructs a real coefficient-side jet center; it imposes no
condition on the actual independent testing value of $u/v$.

\subsection*{Application of the finite-ratio contact calculation}

The homogeneous state jet at a real ratio $t$ is
\begin{align*}
 J_t(W,x,y,Z)={}&4nx^2+4m(W-x-tZ)^2+4C(tW-y)^2+4Ay^2\\
 &+4D(tW-y-Z)^2
   +2d[(tW-y)^2+(x+tZ)^2-2tWZ].
\end{align*}
Its $(W,x,y)$ complement at $Z=0$ is strict: it contains
$4nx^2+4m(W-x)^2+4Ay^2$, with the remaining terms nonnegative.
At $t=t_0$, its kernel is $(w,x_0,y_0,1)$ by (2).

In the original determinant notation, at $\lambda=\beta=\alpha=0$,
\[
 q_{20}=4nAt_3t_4>0,\qquad q_{02}=4mDt_5t_6>0,
 \qquad q_{21}=4CAt_1t_4>0,\qquad q_{12}=4CDt_5t_2>0.
\]
Here the $t_i$ are the original positive sums of bounded coefficients
printed in the finite-pitchfork fit; all are units since $d$ is a
unit. The full jet determinant is
\[
 \det J_t=16(q_{20}t^4+q_{11}t^2+q_{02}).
\]
Thus the singular jet gives precisely the finite nonzero root required
by Section~\ref{component:p03b551f2a9a0}, and its next coefficient satisfies
$q_{21}t_0^4+q_{12}t_0^2>0$. The state complement, both endpoint
coefficients, and the positive next coefficient are all verified.

That calculation normalizes the pure-testing event
$vY-Z$ to $-1$, setting $Z=1+vY$. It keeps the complete opened
$\alpha$ square, including its $4\alpha u^2$ term, in the exact
profile
\[
 F=4\alpha+v^2K(x,y,t,w,V),\qquad
 (X,Y,u,W,V)=(vx,vy,vt,w,v^2)
\]
used only to construct coefficient-field critical centers. The positive
state complement and the double finite root give a strict full Hessian
in $(x,y,t,w)$, with positive next coefficient
\[
 \frac{16(q_{21}t_0^4+q_{12}t_0^2)}{\det M_J(t_0)}.
\]
The two signs of the optimized constant are handled by coefficient-only
implicit centers exactly as in the finite-edge proof; physicality makes
the remaining scalar nonnegative.

Its actual tensor roots are the un-divided expressions
\[
 X-vx_c,\quad Y-vy_c,\quad u-vt_c,\quad v(W-w_c),
 \quad v^2-V_c.
\]
All four full first-normal coefficients are $v\nabla K$ on this
frame. Monic difference division makes their coupling to $v^2-V_c$
infinitesimal. All 55 full-tail monomials have an exact allocation to
the first four roots; factors of $W-w_c$ that remain in coefficients
are infinitesimal. There is no substitution dividing $v(W-w_c)$ by
$v$. The full residue Gram is the positive four-variable Hessian and
the positive last scalar pivot. Hence the theorem applies to every
testing point with the stated residue state, regardless of whether its
actual tested $u/v$ is bounded or approaches $t_0$. This includes literal
$v=0$ and arbitrary independent opening rates.

The two zero-corner state charts and the strict/singular jet split
are exhaustive. Together with the complementary projective cases, this
proves the theorem. The zero state is direct, and complex states split
into real and imaginary quadratic evaluations.



\section{\texorpdfstring{The complete positive-scale C1 reciprocal case}{The complete positive-scale C1 reciprocal case}}
\label{component:p98f24864b071}

\label{p98f24864b071:sec:c1-reciprocal-rederived}

\paragraph{The source domain.}
Let the canonical C1 source be physical, with an actual finite
noncommuting zero contact, and normalize $d=1$. Suppose every source
coefficient is bounded, $\lambda=n=0$, $C,A,\alpha$ are positive
bounded units, $m>0$ and $D\geq0$ are infinitesimal, and
$\beta\geq0$ is arbitrary bounded. Retain the original canonical
source gaps and intersections. Then the entire parent has a finite
sum-of-squares representation in $K_c\otimes_{\mathbb R}K_t$ for
every independent projective test and every real or complex state.
The assertion includes literal $D=0$ and literal $\beta=0$ when
permitted by the actual contact; they cannot both vanish.

We use $O(1)$ and $o(1)$ for bounded and infinitesimal elements of
one of the two fields, and $a\asymp b$ for positive source quantities
whose ratio is a bounded unit. A bounded tensor Gram with a strictly
positive real residue is a sum of squares by the separated strict-Gram lemma. Its rows need not be bounded. No order on the tensor
product, mixed inverse, or cancellation of a determinant multiplier is
used below.

\subsection*{Source balances at the reciprocal corner}
\label{lem:reciprocal-source-balance}
Both C1 and C2 have a reciprocal corner when $m,D$ are
infinitesimal and $C,A,\alpha$ are positive units. All coefficients
are bounded, $d=1$, $\lambda=0$, and $n\beta=0$.
At a stationary source contact put $U=u_c^2$, $P=v_c^{-2}$.
The determinant formula SG5--SG6 becomes
\[
\begin{split}
 L(U,P^{-1})={}&h+\frac{ma_1}{U}+\frac{mDa_2}{UP}+a_3P\\
 &+\frac{\beta ma_4+Da_5}{P}+a_6U+na_8UP
 +\frac{\beta a_7U}{P},
\end{split}
\]
where the positive units are explicitly
\[
\begin{array}{llll}
a_1=4\alpha t_5t_8,&a_2=4t_5t_6,&a_3=4\alpha A t_7t_3,&a_4=4t_1t_6,\\
a_5=4Ct_5t_2,&a_6=4CAt_1t_4,&a_7=4Ct_1t_2,&a_8=4At_3t_4.
\end{array}
\]
Here $h$ is bounded. At a zero, every positive Laurent term is
bounded by $-h$, so $U,P$ are bounded. Differentiating at the
stationary contact gives the exact balances
\begin{equation}
\begin{aligned}
 U^2(a_6+na_8P+\beta a_7/P)&=m(a_1+Da_2/P),\\
 P^2(a_3+na_8U)&=mDa_2/U+\beta ma_4+Da_5+\beta a_7U.
\end{aligned}
\label{eq:reciprocal-source-balances}
\end{equation}
These equations force both $U$ and $P$ to be infinitesimal.
Indeed, if $U$ were a unit and $P$ a unit, the first balance
would have a positive-unit left side and an infinitesimal right
side. If $U$ were a unit and $P$ infinitesimal, the second balance
would give $mDa_2/P\le UP(a_3+na_8U)$, again making the first
right side infinitesimal. Thus $U$ is infinitesimal. If $P$ were
a unit, the first balance would then give
\[
 \frac{mDa_2}{U}
 =UP(a_6+na_8P)+\beta a_7U-ma_1P/U
 \le UP(a_6+na_8P)+\beta a_7U,
\]
which is infinitesimal. Every term on the second right side would
be infinitesimal, contradicting its positive-unit left side.

The nonnegative terms in the second balance now give
$D/P^2=O(1)$ and $\beta U/P^2=O(1)$. The first consequently gives
\begin{equation}
 \frac m{U^2}\asymp1+\frac\beta P,\qquad
 \frac mU,\ \frac DP,\ \frac{\beta U}P\in\operatorname{Inf}.
\label{eq:reciprocal-source-ratios}
\end{equation}
In C2, $\beta=0$, so $m/U^2$ is a positive unit and
\[
 \frac D{P^2}=\frac{a_3+na_8U}{a_5+ma_2/U}
\]
is a positive unit as well. Thus $m\asymp U^2$, $D\asymp P^2$.
For C1 the preceding bounded estimates also allow literal $D=0$
or $\beta=0$ whenever the source contact permits it.


\subsubsection*{Source entry before any state-ratio division}

Write the actual contact angles as $u_c=k>0$, $b_c=1/v_c=p>0$,
and put $U=k^2$, $P=p^2$. Equation~\eqref{eq:reciprocal-source-balances}, with $n=0$,
gives $U,P$ positive infinitesimals and
\[
 D/P^2=O(1),\qquad \beta U/P^2=O(1),\qquad
 m/U^2\asymp1+\beta/P.
 \tag{CR1}
\]
The same balances give
\[
 m/U=o(1),\qquad D/P=o(1),\qquad \beta U/P=o(1).
 \tag{CR2}
\]

At a physical noncommuting contact the signed core is negative, and
the exact core Schur identity SE9 gives $W_cZ_c>0$. Normalize $W_c=1$
and write $Z_c=q>0$. Define
\[
 Q=q^2,\quad L=PQ,\qquad
 r=\frac{k}{pq},\quad y=\frac{Y_c}{pq},\quad
 x=\frac{X_c}{pq^2}.
 \tag{CR3}
\]
We now prove these divisions are bounded wherever they are used.
In the reciprocal quarter-parent CR9 below the hard $X,Y$ matrix at
the source has residue
$\operatorname{diag}(\bar\beta+1/2,H_0)$, where
$H=C+A+\alpha+1/2$ is a positive unit. Eliminate this two-by-two
block inside the source field. Its Schur block on $W,Z$ satisfies
\[
 \frac{S_{WW}}{U}\sim a:=C-C^2/H>0,\qquad
 -\frac{S_{WZ}}{kp}\sim b_*:=1/2+C\alpha/H>0.
 \tag{CR4}
\]
Here and below a displayed positive expression in an asymptotic
formula is evaluated in residue. To check CR4 without a shape
assumption, the original hard couplings are
\[
\begin{aligned}
 P_{XW}&=p\{-m+(D+1/2)U\},&
 P_{YW}&=k\{m-C(1+P)-(D+1/2)P\},\\
 P_{XZ}&=k\{\beta-D+P(m+\beta+1/2)\},&
 P_{YZ}&=p\{D-\alpha-U(m+\alpha+1/2)\},\\
 P_{WZ}&=-(m+D+1/2)kp.
\end{aligned}
\]
CR2 shows that the $X$ corrections in CR4 vanish in residue and
the $Y$ corrections are respectively $C^2/H$ and $C\alpha/H$.
The kernel equation gives $q=S_{WW}/(-S_{WZ})$, so $r$ is a
positive bounded unit.

The scaled $X,Y$ kernel equations in the unknowns $x,y$ have
bounded entries and a diagonal positive residue. For example their
$X$ equation is exactly
\[
\begin{split}
 0={}&\{(\beta+1/2)(1+P)+Pm+(D+1/2)r^2PQ\}x\\
 &+r\{\beta-D+P[\beta+1/2-m(y-1)+(D+1/2)(r-y)]\}-m/Q.
\end{split}\tag{CR5}
\]
The quantity $m/Q=(m/U)r^2P$ is bounded before this equation is
used. The $y$ coefficient in CR5 is infinitesimal; the $x$
coefficient in the scaled $Y$ equation is $O(PQ)$ and its $y$
pivot is $H+o(1)$. Thus $x,y$ are bounded. Taking the $Y$ and
$Z$ kernel equations in residue, using CR1--CR2, yields
\[
 Hy\sim Cr+\alpha,\qquad C(r-y)\sim1/2,
 \qquad \alpha(1-y)\sim r/2.
 \tag{CR6}
\]
Consequently $y,r-y,1-y$ are positive bounded units. Equivalently
their real residues satisfy
\[
 C_0=\frac1{2(r_0-y_0)},\quad
 A_0=\frac{1+r_0-y_0}{2y_0},\quad
 \alpha_0=\frac{r_0}{2(1-y_0)}.
 \tag{CR7}
\]
For the last equation in CR6, divide the original $Z$ kernel
equation by $Pq$. Its extra $X$ contribution is
$O(x\beta Q+xPQ)$; CR1 and the already proved unit $r$ give
$\beta Q=O(P)$, while $PQ=U/r^2$ is infinitesimal. All the
omitted terms therefore vanish in residue. This proves CR6 in the
unbounded-$q$ case as well.

In CR5, $D=o(P)$ and $m/Q=o(P)$. Its numerator for $-x$ is
the sum of $r\beta$ and a positive-unit multiple of $P$, up to
$o(P)$. Hence the following bounds hold on the whole source domain:
\[
\begin{gathered}
 L=PQ=o(1),\qquad \beta Q/P=O(1),\qquad
 -x\asymp P+\beta,\qquad r+x\text{ is a positive unit},\\
 m\asymp PQ^2(P+\beta),\qquad
 m/(PQ)=O(PQ+\beta Q)=o(1),\qquad xQ=o(1).
\end{gathered}\tag{CR8}
\]
The statement about $r+x$ follows directly from CR5:
$x/r=-\beta/(\beta+1/2)+o(1)$.

\subsubsection*{The entire scalar and all normal coefficients}

Use $b=1/v$, $D_u=1+u^2$, $D_b=1+b^2$. The full reciprocal
quarter-parent is
\[
\begin{aligned}
 F={}&m(W-bX+uY-ubZ)^2+CD_b(uW-Y)^2+AD_uD_bY^2\\
 &+(D+1/2)(ubW+uX-bY-Z)^2
 +(\alpha+1/2)D_u(Y-bZ)^2\\
 &+(\beta+1/2)D_b(X+uZ)^2-D_uD_bZ^2/2.
\end{aligned}\tag{CR9}
\]
All further local identities are expansions of this full expression.
In the chart $W=1$, $Z=z$ put
\[
 u_g=rbz,\qquad X_g=xbz^2,\qquad Y_g=ybz,
 \qquad V=b^2,\quad T=z^2.
 \tag{CR10}
\]
Its entire scalar restriction is
\[
 \Phi=m+DT+hVT+\mathcal C V^2T+EVT^2+
 \mathcal F V^2T^2+GV^3T^2+JV^2T^3,
 \tag{CR11}
\]
where
\[
\begin{aligned}
 \mathcal C&=Ay^2+(C+D+1/2)(r-y)^2,\\
 E&=\beta(r+x)^2+x^2/2-2Drx,\qquad
 G=Ar^2y^2,\qquad J=(D+1/2)r^2x^2,\\
 \mathcal F&=Ar^2y^2+(\alpha+1/2)r^2(y-1)^2
       +\beta(r+x)^2+m[r(y-1)-x]^2\\
 &\hspace{12mm}+(2D+1)rx(r-y)+rx+x^2/2,\\
 h&=Ay^2+C(r-y)^2+\alpha(y-1)^2+y^2/2-r
          -2D(r-y)+2m[r(y-1)-x].
\end{aligned}\tag{CR12}
\]
In particular $\mathcal C,G$ are positive units,
$E\asymp\beta+P^2>0$, $J\geq0$, and all these coefficients
are bounded. If $\beta$ is infinitesimal, $\mathcal F$ is a
positive unit, with residue $r_0^2(1+r_0-y_0)/2$.
For unit $\beta$, $\mathcal F$ can have either sign; it is not
discarded or assumed positive.

Actual contact and its two angular/state jets give
$\Phi(P,Q)=\Phi_V(P,Q)=\Phi_T(P,Q)=0$. Therefore
\[
\begin{aligned}
 D&=P^2\mathcal C-PQE+P^3QG-P^2Q^2J,\\
 m&=PQ^2E+P^2Q^2\mathcal F+P^3Q^2G+2P^2Q^3J,\\
 h&=-2P\mathcal C-QE-2PQ\mathcal F-3P^2QG-2PQ^2J.
\end{aligned}\tag{CR13}
\]
Set
\[
 \delta=D/P^2\geq0,\qquad \theta=EQ/P>0.
 \quad \delta,\theta=O(1),\qquad
 \delta+\theta=\mathcal C+GL-JQ^2=\mathcal C+o(1),
 \tag{CR14}
\]
where $L=PQ$. The bounds follow either from CR1 and CR8 or from CR13;
$JQ^2=O((xQ)^2)=o(1)$. At least one of $\delta,\theta$
is a unit. Define
\[
 R=z(V-P),\quad K=VT-PQ,\quad
 M=V^2T+(P/Q)VT^2+P^2Q-3PVT.
\]
The exact complete balanced scalar is
\[
\begin{aligned}
 \Phi={}&\delta R^2+\theta M+\mathcal F K^2
       +G\{PK^2+VzRK\}\\
       &+J\{(QR-zK)^2+2QK^2\}.
\end{aligned}\tag{CR15}
\]
The entire $J$ line is retained as positive squares.

To retain the complete first-normal terms, Let $L_u,L_X,L_Y$
be half the first derivatives of CR9 at CR10. Then
\[
 L_u=bz A_u(V,T),\quad L_Y=bz A_y(V,T),\quad
 L_X=b\{-m+TB_X(V,T)\},
 \tag{CR16}
\]
where
\[
\begin{aligned}
 A_u&=a_{00}+a_{10}V+a_{01}T+a_{11}VT+a_{21}V^2T+a_{12}VT^2,\\
 A_y&=e_{00}+e_{10}V+e_{11}VT+e_{21}V^2T,\\
 B_X&=c_0+c_1V+c_2VT,\\
 c_0&=-Dr+\beta(r+x)+x/2,\\
 c_1&=(D+1/2)r(r-y)+(\beta+1/2)(r+x)
                        +m[x-r(y-1)],\\
 c_2&=(D+1/2)r^2x.
\end{aligned}\tag{CR17}
\]
Every displayed coefficient is bounded. The only bounds beyond this
needed below are
\[
 a_{01}=\beta(r+x)-Dx=O(E),\quad
 a_{12}=(D+1/2)rx^2,\quad a_{12}Q=O(1),
 \qquad c_2=O(P+E).
 \tag{CR18}
\]
They follow from CR8 and CR12. Every other coefficient is the corresponding monomial coefficient
of the explicit polynomial CR9.

State and angle stationarity give
$A_u(P,Q)=A_y(P,Q)=0$ and $B_X(P,Q)=m/Q$.
Put $E_T=bz(T-Q)$. Finite monic division therefore gives exactly
\[
\begin{aligned}
 L_i&=A_{i,V}(V,T)bR+B_i(T)E_T,\
 B_u(T)&=a_{01}+P\{a_{11}+Pa_{21}+a_{12}(T+Q)\},\\
 B_y(T)&=P(e_{11}+Pe_{21}),\\
 L_X&=b\{[m/(PQ)+c_2T]K+z[c_1-m/(PQ)]R\}.
\end{aligned}\tag{CR19}
\]
The first line is for $i=u,y$ and
$A_{i,V}=[A_i(V,T)-A_i(P,T)]/(V-P)$, an exact polynomial.
In particular, using $P(T-Q)=K-zR$, each of $L_u,L_Y,L_X$
is a bounded linear combination of
\[
 bR,\qquad bzK,\qquad b[m/(PQ)+c_2T]K,
 \qquad E E_T.
 \tag{CR20}
\]
Every occurrence of $K$ outside the last product has an
infinitesimal coefficient. CR19 remains literal at $b=0$ or $z=0$.

\subsubsection*{The cubic conductors and the bounded source-ratio chart}

All radicals in this subsection are products of radicals in the
separate fields. Set
\[
 \mathsf A=(V^2T)^{1/3},\quad
 \mathsf B=((P/Q)VT^2)^{1/3},\quad
 \mathsf C=(P^2Q)^{1/3},\qquad \rho=(P/Q)^{1/3}.
\]
Because $\mathsf A\mathsf B\mathsf C=PVT$,
\[
 M=\tfrac12(\mathsf A+\mathsf B+\mathsf C)
 [ (\mathsf A-\mathsf B)^2+(\mathsf A-\mathsf C)^2
                         +(\mathsf B-\mathsf C)^2].
 \tag{CR21}
\]
Let $\mathcal Z$ be its nine roots multiplied by $\sqrt\theta$,
so that $\theta M=\frac12\sum_{j=1}^9\mathcal Z_j^2$.
For nonnegative testing $b,z$, the following exact identities hold:
\[
\begin{aligned}
 bR={}&b\sqrt{\mathsf A}(\mathsf A-\mathsf B)
       +\rho b^{1/3}z^{2/3}\sqrt{\mathsf A}(\mathsf A-\mathsf C),\\
 E E_T={}&\sqrt E\,z\,[\sqrt\theta\sqrt{\mathsf B}
                                      (\mathsf B-\mathsf A)]\\
 &+\theta^{1/3}E^{1/6}b^{2/3}z^{1/3}
       [\sqrt\theta\sqrt{\mathsf B}(\mathsf B-\mathsf C)].
\end{aligned}\tag{CR22}
\]
Negative signs are handled by the exact testing gauges; absolute
values give the same monomial identities if desired. If $\theta$
is a unit, the first line has infinitesimal coefficients on
$\mathcal Z$: $\rho=(E/\theta)^{1/3}$ is bounded and $b$ is
infinitesimal. If $\delta$ is a unit, $bR$ instead has the
infinitesimal coefficient $b$ on the strict row $R$. The second
line has infinitesimal coefficients whenever $E$ is infinitesimal
or $z$ is infinitesimal. Its displayed coefficients use no
comparison of a source rate with a testing rate.

Suppose first that actual $q$ is bounded and $\beta$ is
infinitesimal. Use $K$ and $\mathcal Z$, with $R$ added when
$\delta$ is a unit. If only $\theta$ is a unit, retain
$\delta R^2$ as a separate positive square. In CR15 the $K$
pivot $\mathcal F$ is a positive unit. The only signed extra
term is
$GVzRK=Gbz(bR)K$, an infinitesimal Gram correction by CR22.
The term $GPK^2$ is nonnegative. Thus CR15 has a strict residue
Gram on the stated roots. CR19--CR22 give infinitesimal lifts of
all three complete first normals to these same roots.

If $\beta$ is a unit, CR8 forces $Q$ infinitesimal and $E$ a
unit. The only chart still requiring CR10 will be $z$ infinitesimal.
The second line of CR22 still has infinitesimal coefficients.
When $\delta$ is a unit, use the exact identity
\[
 2QM+TR^2-K^2=(K-zR)^2+2QR^2.
 \tag{CR23}
\]
Choose a positive real $a$ with $\overline{\mathcal F}+a>1$.
Subtracting $a$ times the corresponding zero identity changes the
first three terms of CR15 into
\[
 (\delta-aT+2aQ)R^2+(1-2aP/E)\theta M
       +(\mathcal F+a)K^2+a(K-zR)^2.
 \tag{CR24}
\]
The first three pivots have strictly positive residues; both costs
$T$ and $P/E$ are infinitesimal. This handles the signed
$\mathcal F$ exactly.

When instead $\theta$ is a unit, $\rho$ is also a unit and
\[
 K=\rho^{-1}\{\sqrt{\mathsf A}\,
       [\sqrt{\mathsf A}(\mathsf B-\mathsf C)]
       +\sqrt{\mathsf C}\,
       [\sqrt{\mathsf C}(\mathsf A-\mathsf C)]\}.
 \tag{CR25}
\]
Both coefficients on $\mathcal Z$ are infinitesimal. Thus the
entire signed $\mathcal F K^2$ is an infinitesimal correction on
the strict circuit Gram. The $G$ cross term and every first normal
are controlled by CR22 and CR25; $\delta R^2$ and the $J$ line
remain positive squares. No inverse of $\delta$ is used.

\subsubsection*{An unbounded actual state ratio in the same testing chart}

Now let actual $q$ be unbounded. CR8 and CR12 imply
\[
 \beta/P=O(1/Q)=o(1),\quad -x\asymp P,\quad
 m\asymp L^2,\quad E/P=o(1),\quad JQ^2=O(L^2)=o(1).
 \tag{CR26}
\]
Here $\beta$ is infinitesimal and $\mathcal F$ is a positive
unit. In the testing chart $W=1$, $z$ bounded, the exact relation
\[
 \theta M=\theta R^2+(E/P)(K^2-zRK)
 \tag{CR27}
\]
makes CR15 a strict two-row Gram on $R,K$, plus its whole positive
$J$ line. Its residue is
$\operatorname{diag}(\overline{\delta+\theta},\overline{\mathcal F})$.
Every correction coefficient is infinitesimal: $E/P$, $P$ and
$Vz$ are infinitesimal. In CR19, $B_u/P$ and $B_y/P$ are bounded,
because $a_{01}/P=O(\beta/P+Dx/P)$ and $a_{12}Q$ are bounded.
Using $P E_T=bzK-bz^2R$ shows that all three full normals are
infinitesimal linear combinations of $R,K$.

\subsubsection*{The complementary testing state chart}

Use $Z=1$, $W=w$ infinitesimal, and write $a=u/r$. The graph is
\[
 b=aw,\qquad X=xa,\qquad Y=yaw.
 \tag{CR28}
\]
The entire scalar is
\[
 \Phi'=D+mw^2+Ea^2+ha^2w^2+\mathcal C a^4w^4
       +\mathcal F a^4w^2+Ga^6w^4+Ja^4.
 \tag{CR29}
\]
First assume actual $q$ bounded and $\beta$ infinitesimal. If
$\delta$ is a unit, use the three strict rows
$P,a^2w,\sqrt E\,a$. The sole signed cross in CR29 is
$h a^2w^2=(h/P)w\,P(a^2w)$, and $h/P$ is bounded by CR13.
Every other term is positive or has the positive-unit $\mathcal F$
pivot. If $\theta$ is a unit, $P/E=Q/\theta$ is bounded;
use instead $\sqrt E\,a,a^2w$ and retain $D,mw^2$ whole.
The same signed term is the infinitesimal correction
$(h/E)w^2(\sqrt E\,a)^2$.

Here are the complete complementary first normals, which also prove
the bounds for unbounded actual $q$. Define
\[
 R'=a^2w^2-P,\qquad K'=w(a^2-L).
 \tag{CR30}
\]
Half the $b,Y$ derivatives of the original full CR9 at CR28 have
the respective forms
\[
 L_i'=aw(d_{i0}+d_{i1}a^2+d_{i2}a^2w^2+d_{i3}a^4w^2),
 \qquad i=b,Y,
 \tag{CR31}
\]
with bounded coefficients. At the actual contact
$a^2=L,w^2=1/Q$, they satisfy
$d_{i0}+Ld_{i1}+Pd_{i2}+PLd_{i3}=0$.
Consequently the entire, exact divisions are
\[
 L_i'=(d_{i1}+Pd_{i3})aK'
                   +(d_{i2}+a^2d_{i3})awR'.
 \tag{CR32}
\]
The full $X$ derivative is
\[
 L_X'=a(c_0+c_2a^2-mw^2+c_1a^2w^2),\qquad
 c_0+c_2L-m/Q+c_1P=0.
 \tag{CR33}
\]
For bounded $Q$, $c_0/P$ and $m/P$ are bounded, and
$c_2\in P B_c+E B_c$ by CR18. The products $aK'$ and
$awR'$ are infinitesimal linear combinations of $P,a^2w$.
Expanding CR33 gives infinitesimal coefficients on
$P,a^2w,\sqrt E\,a$. If $\theta$ is a unit, use $P/E=O(1)$:
$c_0/\sqrt E=o(1)$, $c_2=O(E)$ and $m/E=O(1)$, giving the
same assertion on the two rows $\sqrt E\,a,a^2w$.
These are pure-source bounded quotients; no testing $w$ is inverted.

For unbounded actual $q$, CR29 has the exact alternative form
\[
\begin{aligned}
 \Phi'={}&\delta R'^2+\theta M'+\mathcal F K'^2
       +G\{PK'^2+a^2wR'K'\}\\
       &+J\{a^4(Qw^2-1)^2+2QK'^2\},\\
 M'={}&a^4w^4+(P/Q)a^2+P^2Qw^2-3Pa^2w^2.
\end{aligned}\tag{CR34}
\]
Define nonnegative elementary radicals
\[
 \mathsf A'=|aw|^{4/3},\quad
 \mathsf B'=(P/Q)^{1/3}|a|^{2/3},\quad
 \mathsf C'=(P^2Q)^{1/3}|w|^{2/3}.
\]
They give the same nine-square cubic identity for $M'$. The exact
conductor needed for the complete normals is
\[
 |a|R'=|a|\sqrt{\mathsf A'}(\mathsf A'-\mathsf C')
          +\sqrt L\sqrt{\mathsf B'}(\mathsf A'-\mathsf B').
 \tag{CR35}
\]
Its coefficients are infinitesimal. If $\delta$ is a unit, use
$R',K'$ and retain the circuit as positive squares. If $\theta$
is a unit, use its scaled nine roots and $K'$, retaining
$\delta R'^2$ whole. In either case CR34 is strict: its sole
signed cross is $Gaw(aR')K'$, controlled by CR35, and its last
line is retained as squares.

CR26 implies $c_2/P=O(1)$ and
$c_0/P=m/(PQ)-Qc_2-c_1=O(1)$. The full exceptional normal
has the exact division
\[
\begin{aligned}
 L_X'={}&-[c_0/P+(c_2/P)a^2]\,aR'\\
 &+[(c_2/P)(a^2+L)+c_0/P+c_1]\,awK'.
\end{aligned}\tag{CR36}
\]
Together with CR32, this uses only bounded multiples of
$aR',awR',aK'$. Each has infinitesimal coefficients on the strict
soft rows by CR35. This also proves the result at literal $w=0$.

\subsubsection*{All higher normals and the remaining unit-beta states}

The preceding scalar and first-normal formulas certify the entire
parent. Here are the required hard blocks and the justification for
every higher term. Normalize the testing state so its coordinates
are bounded with at least one unit. At the reciprocal angle corner
$u,b=o(1)$ the state residue hard block is
$(\bar\beta+1/2)X^2+H_0Y^2$. If $X$ or $Y$ is a unit, the
entire parent has positive real residue and the strict-element lemma
applies immediately. Otherwise $X,Y$ are infinitesimal and one of
$W,Z$ is a unit.

The exact sign gauges make $u,b\geq0$. If $WZ\leq0$, the signed
core is a nonnegative pure-testing scalar by the full core Schur
identity SE9 and its literal homogeneous endpoints, so its source
multiple and all event squares are separated squares. Otherwise
choose either $W=1$, $z=Z$ bounded, or $Z=1$, $w=W$
infinitesimal. Both have $X,Y$ infinitesimal.

In the $W=1$ chart, set
$N_u=u-rbz$, $N_X=X-xbz^2$, $N_Y=Y-ybz$.
They are infinitesimal. The complete expansion of CR9 is its scalar
CR11, twice the three first normals CR16, and a polynomial quadratic
Gram in $N_u,N_X,N_Y$. Terms of normal degree at least three
have infinitesimal Gram entries. For weak $\beta$, or for unit
$\beta$ with $z$ infinitesimal, the hard residue is
\[
 C_0(N_u-N_Y)^2+(A_0+\alpha_0+1/2)N_Y^2
                         +(\beta_0+1/2)N_X^2\succ0.
 \tag{CR37}
\]
All first-normal lifts established above have infinitesimal
coefficients. The full hard-and-soft residue is therefore positive
definite. In the $Z=1$ chart use
$N_b=b-aw$, $N_X=X-xa$, $N_Y=Y-yaw$.
For weak $\beta$ its hard residue is
\[
 N_X^2/2+H_0N_Y^2-2\alpha_0N_bN_Y+\alpha_0N_b^2\succ0,
 \tag{CR38}
\]
since $H_0-\alpha_0=C_0+A_0+1/2>0$.
The same finite expansion handles every higher normal term. Thus
none of CR15, CR27 or CR34 is being used merely as a scalar
restriction.

It remains only unit $\beta$ and a testing state with $Z$ a unit
relative to $W$. Normalize $Z=1$ and $W=\omega$ bounded. Retain
the entire $m,D$ squares. The remainder belongs exactly to the
square of the ideal $(u,b,X,Y)$. Its quadratic real residue is
\[
\begin{aligned}
 &(\beta_0+1/2)\left(X+\frac{\beta_0}{\beta_0+1/2}u\right)^2
      +\frac{\beta_0}{1+2\beta_0}u^2\\
 &\quad+H_0\left(Y-\frac{C_0\omega u+\alpha_0 b}{H_0}\right)^2
      +a_0(\omega u-r_0b)^2,
 \qquad a_0=C_0-C_0^2/H_0>0.
\end{aligned}\tag{CR39}
\]
CR7 is precisely the critical relation needed for this identity.
The expression is positive definite in $(u,b,X,Y)$ for every
real $\omega$, including zero. Its complete bounded polynomial
matrix is consequently strict; every higher degree is an
infinitesimal correction. This completes all local states, including
the unit-beta case where $\mathcal F$ was signed.

\subsubsection*{Every other test and restoration}

The exact outer identities RC18--RC25 from the newly rederived
C2 reciprocal theorem apply here as polynomial identities, with
$n=0$ and the entire extra $\beta$ square retained. Their strictness
uses only $C,A,\alpha$ positive units and the critical residue
relations CR7; no C2 source membership is used. Explicitly:
\begin{itemize}
\item For bounded original $u,v$, the four rows $(uW,Y,X,Z)$
give RC18. After CR7 its real determinant is
$ (r_0+u^2v^2)D_u^2D_v^2(1+r_0-y_0)^2/
[16y_0(r_0-y_0)(1-y_0)]>0$.
\item For Alice reciprocal and Bob bounded, use the four rows
$(W,Y,vX,sX+Z)$ and the full core RC20. Its real residue RC21
is strict using $C,A>0$.
\item For Bob reciprocal and Alice a bounded unit, use
$(uW-Y,Y,X,bZ,\sqrt D\,Z)$ and the entire $D$ square as in
RC22--RC23. CR7 makes the three-by-three angular determinant
strict. Literal $D=0$ causes its last row to be zero, which is
allowed in a strict coefficient Gram.
\item For both angles reciprocal, RC24--RC25 give the five rows
$(W,X,Y,bZ,\sqrt D\,sZ)$ and a strictly positive real residue.
\end{itemize}
The only remaining angle chart is exactly $u,b=o(1)$, treated
above. These are whole-parent identities with every omitted event
square retained, and include all literal axes and infinities.
The exact source and testing sign gauges, positive pure-field
homogeneous normalization factors and real/imaginary state splitting
restore the stated projective tensor-SOS theorem.



\section{Exhaustion of the bounded source domain}
\label{component:pc0b2849f8d09}

\label{pc0b2849f8d09:sec:positive-scale-union-rederived}

\paragraph{Exact statement.}
Let \(K_c,K_t\) be labelled real closed extensions of \(\mathbb R\).
The original canonical coefficient vector is physical over \(K_c\),
has nonnegative coefficients, and has an actual finite noncommuting
zero contact. In C1 the exact omissions are \(\lambda=n=0\),
with \(m,C,A,\alpha>0\), \(D,\beta\geq0\), \(D+\beta>0\);
in C2 they are \(\lambda=\beta=0\), with
\(m,C,A,D,\alpha>0\), \(n\geq0\).
Use the actual native sign chart of this contact, with its explicitly
positive source gaps, as supplied by canonical entry.
Assume that normalization by \(d>0\) makes all coefficients bounded.
Equivalently the normalized total-scale coordinate
\(d/(d+\lambda+n+m+C+A+D+\alpha+\beta)\) is a positive unit.
Then the entire original homogeneous slack is a sum of squares
in \(K_c\otimes_{\mathbb R}K_t\) for every independent projective
testing pair and every real or complex testing state.
Chapter~\ref{ch:entry} supplies these contact coordinates for every
source in the two families.

Division by \(d\) is a positive pure-source normalization.
Hence take \(d=1\) below. Every nonnegative bounded source quantity
is either infinitesimal (including literal zero where permitted)
or a positive unit. These two alternatives give the source partition below.

We use the following two cases established above:
\[
\begin{array}{c|l|l}
 &\text{positive-unit coefficients}&\text{proof}\\\hline
\mathrm S&d,C,A,D,\alpha&\text{Section~\ref{component:pfb0cf454f653}}\\
\mathrm{CU}&d,n,m,C,A,D\quad(\alpha\text{ infinitesimal})
 &\text{Section~\ref{component:pf31e38136e03}}.
\end{array}
\]
Every other nonnegative coefficient may have arbitrary residue in
the stated canonical support. Both proofs certify the full state
polynomial at all independent tests. Their noncommuting case is
Theorem~\ref{thm:noncommuting-refined}; their finite-edge and
finite-ratio cases were proved at the start of this chapter.


\subsection*{Symmetries and the source partition}

Both transformations below preserve the core and nonnegative event
coefficients, actual physicality, finite nonzero contact, and pure
total-scale normalization. They act separately in source and testing
fields, so transport a full tensor certificate.

For C1 let
\[
 \sigma_1=(\lambda\ n)(m\ C)(A\ \alpha)(D\ \beta),\qquad
 G_v(W,X,Y,Z)=(vW-X,-W-vX,Z-vY,Y+vZ).
\]
The entire raw parent has the polynomial identity
\[
 u^2F_{\sigma_1c}(1/u,v;G_v\psi)
       =(1+v^2)F_c(u,v;\psi),\qquad
 G_v^{\mathsf T}G_v=(1+v^2)I .
 \tag{AU12}
\]
Homogenization interchanges the two Alice angle coordinates;
the Bob-homogeneous \(G\) has norm factor \(v_0^2+v_1^2>0\).
Thus literal endpoints are transported without dividing zero.

For C2 let
\[
 \sigma_2=(\lambda\ \beta)(m\ D)(C\ \alpha),\qquad
 J(W,X,Y,Z)=(Z,X,Y,W).
\]
The full identity is
\[
 u^2v^2F_{\sigma_2c}(1/v,1/u;J\psi)=F_c(u,v;\psi).
 \tag{AU13}
\]
It is the physical party interchange with both first observables
negated, followed by the local real \(X\) basis changes.
The homogeneous identity includes both projective endpoints.
In particular the source squared half-angles transform as
\((U,V)\mapsto(1/V,1/U)\).

The two actual forbidden weak faces are
\[
\begin{aligned}
 h|_{\lambda=\beta=A=0}
 &=-(1+2n)(2D+2C+1)(2\alpha+1+2m),\\
 h|_{\lambda=n=m=C=0}
 &=-(2D+1)(2\beta+1)(2A+2\alpha+1).
\end{aligned}\tag{AU14}
\]
In C2 with \(A\) weak, the only positive Laurent terms are
\(q_{01}/U+q_{02}V/U+q_{12}V\), which tend to zero along
\(U=t^2,V=1/t,\ t\to\infty\).
In C1 with both \(m,C\) weak the only positive term is \(q_{10}/V\),
which tends to zero as \(V\to\infty\).
The negative constants contradict physicality.
Thus C2 has \(A\) unit, and C1 has at least one of \(m,C\) unit.


\paragraph{C2.}
If \(\alpha\) is weak, WA1--WA6 force \(C,A,D\) to be units.
There are precisely three disjoint alternatives:
\[
\begin{array}{c|c}
 n,m\text{ units}&\mathrm{CU}\\
 n\text{ unit},\ m\text{ weak}&\mathrm{WA1\!-\!WA27}\\
 n\text{ weak},\ m\text{ arbitrary}&\mathrm{WN1\!-\!WN44}.
\end{array}\tag{AU15}
\]
If \(C\) is weak, apply \(\sigma_2\) and the same table.
Every unit requirement in the transformed table is a consequence
of the transformed physical source, including its required
\(D'=m\) bound; it is not assumed for an excluded original case.

It remains that \(C,A,\alpha\) are units.
The actual dictionary then has
\[
 q_{01}=ma_1,\ q_{02}=mDa_2,\ q_{10}=a_3,\
 q_{12}=Da_4,\ q_{20}=na_5,\ q_{21}=a_6 ,
\]
where every \(a_i\) is a positive bounded unit.
The contact equation bounds \(U,V^{-1},m/U,DV\).
The two angular equations give exactly
\[
 \begin{split}
 DV^2(a_4+a_2m/U)&=a_3+na_5U,\\
 U^2(a_6+na_5/V)&=m(a_1+a_2DV).
 \end{split}
\]
All parentheses are positive bounded units, so
\(DV^2\asymp1\), \(U^2\asymp m\).
Consequently the exhaustive partition is
\[
\begin{array}{c|c}
 D\text{ unit},\ m\text{ arbitrary}&\mathrm S\\
 D\text{ weak},\ m\text{ unit}&\sigma_2(\mathrm S)\\
 D,m\text{ weak}&\mathrm{RC1\!-\!RC25}.
\end{array}\tag{AU16}
\]
RC is the direct full-parent C2 reciprocal theorem; it permits
arbitrary \(n\geq0\). No unit-\(n\) condition has entered this
part of the partition.

\paragraph{C1.}
If \(\alpha\) is weak, WN applies directly, with all its common
source-entry implications proved in WA1--WA10.
If \(A\) is weak, apply \(\sigma_1\), making \(\alpha'=A\) weak,
and use WN on the transformed physical source.
Otherwise \(A,\alpha\) are units.
By AU14, orient with \(\sigma_1\), if necessary, so that \(C\)
is a unit. The exact dictionary becomes
\[
\begin{gathered}
 q_{01}=ma_1,\ q_{02}=mDa_2,\ q_{10}=a_3,\quad
 q_{12}=\beta ma_4+Da_5,\quad q_{21}=a_6,\quad q_{22}=\beta a_7,\\
 U^2(a_6+\beta a_7V)=m(a_1+a_2DV),\\
 a_3=V^2\{mDa_2/U+\beta ma_4+Da_5+\beta a_7U\}.
\end{gathered}\tag{AU17}
\]
The contact equation again bounds \(U,V^{-1},m/U,DV\).
If \(m\) is weak, \(U\) is infinitesimal.
If \(m\) is a unit, \(U\) is a positive unit.
Indeed, the contact equation also bounds \(\beta UV\); if
\(U\) were infinitesimal, then
\(U^2(a_6+\beta a_7V)\) would be infinitesimal, contradicting
the positive unit \(ma_1\) on the right of AU17.
For unit \(m\), the last braces are comparable with \(D+\beta\).
For weak \(m\), bounded \(V\) forces \(m/U\) weak in the
preceding equation, and then the last equation forces \(D\)
unit; conversely unit \(D\) forces bounded \(V\).
Thus the following source partition is exhaustive:
\[
\begin{array}{c|c}
 D\text{ unit},\ m\text{ arbitrary}&\mathrm S\\
 D\text{ weak},\ m\text{ unit},\ \beta\text{ unit}
       &\sigma_1(\mathrm S)\\
 m\text{ unit},\ D,\beta\text{ weak}
       &\mathrm{PP1\!-\!PP19}\\
 m,D\text{ weak},\ \beta\text{ arbitrary}
       &\mathrm{CR1\!-\!CR39}.
\end{array}\tag{AU18}
\]
In the second line all of \(m,C,A,\alpha\) are units before
reflection, so the reflected strong theorem has every required
unit. PP is the C1 projective-pitchfork theorem, with actual
contact \(U\) unit and \(V\) unbounded.
CR is the C1 reciprocal all-tests theorem, including unit
\(\beta\) and literal allowed \(D=0\) or \(\beta=0\).
The last line treats infinitesimal \(m,D\).
No boundary is promoted by a density argument.

AU15--AU18 cover every permitted bounded positive-\(d\) canonical
source. Each leaf proves the entire original parent at every
independent projective test, with the exact same source coefficients.
The full-parent symmetries, positive pure-source \(d\) scaling,
positive pure-testing homogeneous factors and complex-state split
restore the assertion in the opening paragraph. This completes
A5 at its canonical-entry interface.

The source estimates in each listed case use precisely its
physical contact equations and the displayed coefficient bounds.
The finite partition therefore proves the all-tests assertion for
every bounded canonical source.

\section{A three-event face and its stationary zeros}
\label{component:p23a094ff816d}

\label{p23a094ff816d:sec:rederived-three-event-contact}

The three-event face has an explicit matrix representation. Its
dual gives a separated sum of squares. We then use the stationary
equations to show that this face contains no source with a finite
noncommuting contact.

\subsubsection*{The selected full parent}

Write $D_u=1+u^2$, $D_v=1+v^2$, and retain the coefficient order
$(d,\lambda,n,m,C,A,D,\alpha,\beta)$.  The original complete parent is
\begin{align}
F={}&4\lambda D_u(W+vX)^2+4nD_uD_vX^2
 +4m(vW-X+uvY-uZ)^2\notag\\
&+4CD_v(uW-Y)^2+4AD_uD_vY^2
 +(4D+2d)(uW+uvX-Y-vZ)^2\notag\\
&+(4\alpha+2d)D_u(vY-Z)^2
 +(4\beta+2d)D_v(X+uZ)^2-2dD_uD_vZ^2.
\tag{CAP1}
\end{align}
All nine coefficients are nonnegative.  Let $F_*$ be the selected parent
with $\lambda=n=m=\alpha=\beta=0$, so that only the consecutive event
triple $(D,C,A)$ remains.  For $d,C,A,D\geq0$ define
\[
\Delta=8ADC-4ADd-d^2\{d+2(C+A+D)\}.
\tag{CAP2}
\]

\begin{proposition}[Exact physical wall and strict interior-angle positivity]
If $d>0$, the selected parent $F_*$ is positive semidefinite at every
projective measurement pair if and only if $\Delta\geq0$.  Whenever
this holds, its state matrix is positive definite at every finite pair
$u,v\ne0$.  In particular, the equality wall $\Delta=0$ has no finite
noncommuting contact.
\end{proposition}

\begin{proof}
Let $P_*(u,v)$ be the symmetric state matrix.  Direct expansion of the
full polynomial (CAP1), without a state or angle truncation, gives
\[
\det P_*=16D_u^2D_v^2u^2v^2
 \{d\Delta+4CDd(d+2A)v^2+4CAd(d+2D)u^2\}.
\tag{CAP3}
\]
Its restriction to $Z=0$ is the sum of the
nonnegative event squares and
\[
2d\{(uW+uvX-Y)^2+D_uv^2Y^2+D_vX^2\}.
\tag{CAP4}
\]
For $uv\ne0$ these three squared linear forms are independent.  Hence
the upper $3\times3$ state block is positive definite.  A Schur
complement taken in this one ordered field proves that $P_*\succeq0$
if and only if its determinant is nonnegative, and that a positive
determinant gives positive definiteness.  This is a statement about
ordinary field-valued source or testing matrices; it takes no inverse
in a tensor product.

Suppose $\Delta<0$.  Put $r=4CDd(d+2A)$ and
$s=4CAd(d+2D)$, both nonnegative.  Taking
\[
u^2=v^2=\frac{-d\Delta}{2(r+s+1)}>0
\]
makes the expression in braces in (CAP3) negative.  This is an explicit
ordered-field witness to nonphysicality, so no limiting argument is
needed.  Conversely, $\Delta\geq0$ makes all three terms in braces
nonnegative.  It also forces $C,A,D>0$: if one vanished, (CAP2) would
be strictly negative because $d>0$.  Thus $r,s>0$ and the determinant
is strictly positive whenever $uv\ne0$.

At an actual commuting angle, including a projective infinity, one
party's two reflections commute.  Simultaneous diagonalization makes
the CHSH operator on each local eigenspace either $\pm2$ times one
remote reflection.  The signed core $I-S/2$ is then positive
semidefinite.  Every event square is positive semidefinite too.  This
covers all literal endpoints directly.  The argument holds over every
real closed field, because it uses real diagonalization of symmetric
matrices and explicit ordered-field identities.  If $d=0$, every
parent in (CAP1) is its event-square sum separately.
\end{proof}

\subsubsection*{A small exact lift of the selected source cone}

Set
\[
p=4C+4D+2d,\qquad r=4D+2d,\qquad
B(d,C,A,D)=
\begin{pmatrix}
p+4A&4D&-p\\
4D&4D&-r\\
-p&-r&p
\end{pmatrix}.
\tag{CAP5}
\]
For nonnegative coefficients the selected source is physical exactly
when $B\succeq0$.  Here is a direct verification of this matrix
description.  When $d>0$, use the positive last diagonal entry $p$.
The Schur complement of that entry is
\[
\begin{pmatrix}4A&-2d\\-2d&4b/p\end{pmatrix},
\qquad b=4CD-2Dd-d^2,
\qquad \det B=8\Delta.
\tag{CAP6}
\]
Moreover
\[
\Delta=2Ab-d^2(d+2C+2D).
\tag{CAP7}
\]
If $\Delta\geq0$, then $A,b>0$, and (CAP6) is positive semidefinite.
If $B\succeq0$, its determinant is nonnegative, hence $\Delta\geq0$.
For $d=0$ one has the exact quadratic identity
\[
(y,z,x)B(0,C,A,D)(y,z,x)^T
 =4C(y-x)^2+4Ay^2+4D(y+z-x)^2.
\tag{CAP8}
\]
Thus the entire selected physical cone, including all zero faces, is
\[
\mathcal C=\{(d,C,A,D)\in\mathbb R_{\geq0}^4:B(d,C,A,D)\succeq0\}.
\tag{CAP9}
\]
It has a strict feasible point $(1,2,2,2)$: the leading principal
minors of $B$ there are $26,144,280$.

\begin{lemma}[Separated certificate for the selected wall]
For arbitrary labelled real closed fields $K_c,K_t\supset\mathbb R$,
every source $c\in\mathcal C_{K_c}$, and every projective angle pair
and real state in $K_t$, the complete selected parent is a finite sum
of squares in $K_c\otimes_{\mathbb R}K_t$.  The full parent (CAP1)
has the same property if its selected coefficients satisfy (CAP9).
\end{lemma}

\begin{proof}
We first establish the exact dual representation used in this lemma.
For a real linear functional $\ell$ nonnegative on $\mathcal C$,
there are $\mu\in\mathbb R_{\geq0}^4$ and a real symmetric
$Q\succeq0$ such that, for all $c=(d,C,A,D)$,
\[
\ell(c)=\mu\cdot c+\operatorname{tr}(B(c)Q).
\tag{CAP10}
\]
Indeed let $\mathcal D$ be the cone of functionals on the right.  It
is closed: if a sequence of these functionals converges, evaluation
at $c_0=(1,2,2,2)$ bounds both $\sum\mu_i$ and
$\operatorname{tr}Q$, because all coordinates of $c_0$ and the
smallest eigenvalue of $B(c_0)$ are strictly positive.  Positive
semidefinite matrices with bounded trace are bounded, so there is a
convergent subsequence of the representing pairs.  The dual cone
$\mathcal D^*$ is exactly (CAP9).  Finite-dimensional separation of
closed convex cones gives $\mathcal D=\mathcal C^*$ and proves
(CAP10).  This statement has a fixed finite first-order formulation
using the principal-minor inequalities for a $3\times3$ matrix.
Consequently it holds over every real closed field by transfer.

Fix an arbitrary testing tuple in $K_t$.  Its homogeneous selected
parent is a linear functional $\ell_t$ on the four source
coefficients.  The physical equivalence just proved says that
$\ell_t\geq0$ on $\mathcal C_{K_t}$.  Apply (CAP10) in $K_t$ to
obtain $\mu_t\geq0$, $Q_t\succeq0$.  Evaluating the resulting
linear coefficient identity in separated fields gives
\[
F_*(c,t)=\sum_{i=1}^4c_i\otimes\mu_{t,i}
 +\sum_{j,k=1}^3 B(c)_{jk}\otimes(Q_t)_{jk}.
\tag{CAP11}
\]
The first terms are squares by separate square roots.  Factor
$B(c)=UU^T$ within $K_c$ and $Q_t=VV^T$ within $K_t$.  The last
term in (CAP11) is exactly
\[
\sum_{a,b}\left(\sum_jU_{ja}\otimes V_{jb}\right)^2.
\]
This proves the tensor certificate with no mixed denominator.  The
five omitted event squares in (CAP1) are separate positive-source
factors times testing squares and can be restored exactly.  Every
literal projective endpoint is already included in the homogeneous
testing tuple.  Splitting a complex state into its real and imaginary
parts gives the same assertion for complex states.
\end{proof}

\subsubsection*{Every actual finite contact lies strictly below all eight walls}

The fixed event cycle is
\[
(a_0,\ldots,a_7)=(\lambda,D,C,A,\alpha,m,\beta,n).
\]
For indices modulo eight put
\[
\Delta_i=8a_i a_{i+1}a_{i+2}-4d a_i a_{i+2}
 -d^2\{d+2(a_i+a_{i+1}+a_{i+2})\}.
\tag{CAP12}
\]
The cyclic equivalence can be verified without an unspecified symmetry
import.  In the CHSH convention
$S=A_0B_0-A_0B_1+A_1B_0+A_1B_1$, the Bell relabeling
\[
(A'_0,A'_1,B'_0,B'_1)=(B_1,-B_0,-A_0,-A_1)
\tag{CAP13}
\]
preserves $S$ and maps each event in the displayed cycle to the next.
For example the transformed $\lambda=(+,+\mid0,1)$ is the original
$D=(-,+\mid1,1)$.  The complete eight-event permutation and the CHSH
identity are checked exactly.  Relabeling preserves whether each local
pair is noncommuting, so the strict assertion of the first proposition
applies to every cyclic triple.

\begin{theorem}[Automatic contact exclusion]
Let $d>0$, let every event coefficient be nonnegative, and suppose
the original parent (CAP1) has a zero on a nonzero source-field state
at an actual finite noncommuting source pair.  Then
\[
\Delta_i<0\qquad (i=0,\ldots,7).
\tag{CAP14}
\]
\end{theorem}

\begin{proof}
If some $\Delta_i\geq0$, the selected parent for that triple is
positive definite at the specified pair.  The full parent is this
selected parent plus five nonnegative event squares.  It is therefore
positive on every nonzero state at that pair, contradicting the
assumed zero.  This uses only ordinary ordered-field positivity at
the source contact.  It does not evaluate a source contact in an
independent testing field.
\end{proof}

More generally, suppose a proposed cap keeps the same signed coefficient
$d$ and removes nonnegative amounts $e_j$ from event coefficients.
At an actual zero $t_c$ one has the exact identity
\[
F_{\mathrm{cap}}(t_c)
 =-4\sum_j e_j\,\omega_j(t_c)R_j(t_c)^2,
\tag{CAP15}
\]
where $\omega_j$ is the positive angle factor multiplying the event
row in (CAP1).  Thus a globally physical cap can remove a strictly
positive amount only from an event whose row vanishes at that contact.
If one removed positive square has nonzero row, (CAP15) is already a
strictly negative physicality witness.  In the special three-event cap,
(CAP14) rules out even the case in which all removed rows vanish.
The same argument permits smaller retained event coefficients: add
their nonnegative excesses back as well.  It does not apply to a
modification of $d$, whose difference is a signed CHSH term rather
than a nonnegative event square.

\subsubsection*{Exact consequences for source normalization and routing}

For $d>0$ put
\[
\delta_{\rm tot}=\frac{d}{d+\sum_{i=0}^7a_i},
\qquad b_i=a_i/d.
\tag{CAP16}
\]
In a real closed extension, ``bounded'' means bounded in absolute
value by a real number, and ``infinitesimal'' means smaller in
absolute value than every positive real number.  Since all $b_i$
are nonnegative and the sum is finite,
\[
\begin{array}{ll}
\delta_{\rm tot}\text{ has positive real residue}
 &\Longleftrightarrow\text{every }b_i\text{ is bounded},\\
\delta_{\rm tot}\text{ is infinitesimal}
 &\Longleftrightarrow\text{at least one }b_i\text{ is unbounded}.
\end{array}
\tag{CAP17}
\]
These alternatives are disjoint and exhaustive because
$0<\delta_{\rm tot}\leq1$.  They are invariant under a positive
source rescaling and under Bell relabeling.  Normalizing $d=1$ is
legal in the coefficient field, and multiplying each tensor root by
the pure-source factor $\sqrt d$ restores the original scale.

There cannot be three consecutive unbounded $b_i$ at an actual finite
contact.  Indeed division of (CAP12) by its positive triple product
gives
\[
\frac{\Delta_i}{a_i a_{i+1}a_{i+2}}
 =8-\frac{4d}{a_{i+1}}
 -\frac{d^3}{a_i a_{i+1}a_{i+2}}
 -2d^2\left\{
 \frac1{a_i a_{i+1}}+\frac1{a_i a_{i+2}}
 +\frac1{a_{i+1}a_{i+2}}\right\}.
\tag{CAP18}
\]
If the three ratios $a_j/d$ were unbounded, its residue would be
$8$, contradicting (CAP14).  Equivalently, every consecutive triple
contains an event bounded relative to $d$.

If a later chart uses $C=2$, the rescaling factor is $2/C$, and its
new signed weight is $2d/C$; (CAP16) remains unchanged.  The implication
$d/C$ infinitesimal $\Rightarrow\delta_{\rm tot}$ infinitesimal is
valid, but the reverse implication does not follow from normalization
alone.  A different event may be the unbounded one.  An entry map that
claims more must prove its extra source inequalities explicitly.

For a more precise cap screen with $d>0$, (CAP7) yields
\[
b>0\quad\Longrightarrow\quad
\Delta<0\ \Longleftrightarrow\
A<\frac{d^2(d+2C+2D)}{2b}.
\tag{CAP19}
\]
If $b\leq0$, the strict inequality $\Delta<0$ is automatic.  Every
inverse in (CAP19) is within the source field, and the polynomial
condition (CAP14) remains the primary contact condition.



\section{\texorpdfstring{Weak-alpha source entry and the unit-\(n\) case}{Weak-alpha source entry and the unit-n case}}
\label{component:p3a3949869b9e}

\label{p3a3949869b9e:sec:new-weak-alpha-unit-n}

\paragraph{The source domain.}
All coefficients are bounded in the source field, \(d=1\), and the
original canonical source is physical with an actual finite noncommuting
zero contact. The C2 convention is \(\lambda=\beta=0\), and the C1
convention is \(\lambda=n=0\). The first subsection applies to either
family with \(\alpha\) infinitesimal. The all-tests theorem proved below
has the additional C2 hypotheses that \(n\) is a positive unit and \(m\)
is infinitesimal. It includes all independent testing states, signs and
literal projective endpoints. The weak-\(n\) and C1 cases are separate.

\subsubsection*{Source entry, proved before any shape quotient is used}

Physicality forces \(C,A,D\) to be positive units. Here are exact
residue-face tests proving this assertion. In C2 the negative central
coefficient at \(A=0\), and its vanishing positive Laurent terms along
\(U=t^2,V=1/t\), are the displayed forbidden-face proof in the actual
positive-\(d\) source partition. If also \(C=0\) in residue, the source
has only the four alternating events \(n,m,A,D\). At \(u=v=1\) the full
parent at state \((1,-\varepsilon,\varepsilon,1)\) is exactly
\[
 16\varepsilon\{(1+n+m+A+D)\varepsilon-1\}<0
 \tag{WA1}
\]
for a sufficiently small positive real \(\varepsilon\).
Thus \(C\) cannot be weak. The remaining central coefficient identities
are
\[
\begin{aligned}
h|_{\alpha=D=0}
 &=(2A+2C+1)\{8\beta mn-2\beta-4mn-2m-2n-1\},\\
h|_{\alpha=C=n=0}
 &=-(2A+1)(2D+1)(2\beta+2m+1),\\
h|_{\alpha=A=n=0}
 &=-(2C+2D+1)(2\beta+2m+1).
\end{aligned}\tag{WA2}
\]
In the first line \(n\beta=0\), so it is strictly negative. For C2
with \(\alpha=D=0\), the other Laurent terms are
\(q_{20}U/V+q_{21}U\), killed by \(U=t^2,V=t\).
For C1 with \(\alpha=D=0\), they are
\(q_{12}V+q_{21}U+q_{22}UV\), killed by \(U=V=t\).
For C1 with \(\alpha=C=0\), they are \(q_{02}V/U+q_{12}V\),
killed by \(U=1/t,V=t\). For C1 with \(\alpha=A=0\), they are
\(q_{02}V/U+q_{12}V+q_{22}UV\), killed by \(U=t,V=t^2\).
Here \(t\) tends to zero through positive reals. These contradictions
prove all the asserted unit bounds without a source-contact assumption.

Write \(U=u_c^2,V=v_c^2>0\). The actual determinant dictionary now has
\[
 q_{01}=\alpha m a_1,\quad q_{02}=ma_2,\quad
 q_{10}=\alpha a_3,\quad q_{12}=a_4,\quad
 q_{20}=na_5,\quad q_{21}=a_6,\quad q_{22}=\beta a_7,
 \tag{WA3}
\]
where the \(a_i\) are positive bounded units. The two angular equations
and their sum are
\[
\begin{aligned}
 U^2(na_5+a_6V+\beta a_7V^2)&=mV(\alpha a_1+a_2V),\\
 V^2(ma_2+a_4U+\beta a_7U^2)&=\alpha a_3U+na_5U^2,\\
 a_6U+a_4V+2\beta a_7UV
 &=\alpha(ma_1/U+a_3/V).
\end{aligned}\tag{WA4}
\]
Every \(t_i\) in the dictionary is a positive unit because it includes
\(d=1\). At contact all noncentral Laurent terms are nonnegative,
and their sum is bounded. In particular \(U,V,mV/U,nU/V\) are bounded.
Multiplying the last equation by \(V\) shows
\(\alpha/V=O(U+V)\). The first two equations then exclude either
\(U\) or \(V\) having nonzero residue: if \(U\) is a unit the first
forces \(V\) a unit, and conversely the second forces \(U\) a unit;
the last equation then contradicts weak \(\alpha\).
Both are infinitesimal. Consequently
\[
 \alpha/V=o(1),\qquad
 m\asymp (U/V)^2(n+V),\qquad
 \alpha\asymp V(U+V).
 \tag{WA5}
\]
For the last comparison, the denominator
\(a_3+ma_1V/U\) is a positive bounded unit. Actual finite contact
forces \(m,\alpha>0\) in WA4. If \(n\) is a unit and \(m\) weak,
\[
 U/V\asymp\sqrt m=o(1).
 \tag{WA6}
\]
If \(n\) is weak, then, whether \(U/V\) is bounded or unbounded,
WA5 gives \(mV/U,nU/V=o(1)\). These same two conclusions hold
under WA6. They are the common source bounds used in the shape argument.

Orient the actual signs positively, use \(W_c=1\), and put
\[
 p=v_c,\quad q=Z_c,\quad P=p^2,\quad Q=q^2,\qquad
 u_c=rpq,\quad X_c=xpq^2,\quad Y_c=ypq.
 \tag{WA7}
\]
The exact core sign identity SE9 ensures \(W_cZ_c>0\), so this
normalization is legitimate. To prove boundedness before using these
quotients, eliminate the \(X,Y\) block inside the source field. Set
\[
 H_X=n+m+\beta+\tfrac12,\quad K=C+D+\tfrac12,\quad H=K+A .
\]
The block has diagonal positive residue \(H_X,H\) and an infinitesimal
off-diagonal entry. The source Schur entries satisfy
\[
 S_{WW}/U\sim a=AK/H>0,\qquad
 -S_{WZ}/(u_cv_c)\sim b=\tfrac12+AD/H>0.
 \tag{WA8}
\]
Indeed the only additional leading contributions are multiples of
\(mV/U,nU/V,mn\), all infinitesimal on the domains just stated.
The kernel equation gives \(r\sim b/a\), a positive unit. The scaled
\(X,Y\) equations are then a bounded system with positive triangular
residue; its apparent \(m/Q\) term is bounded since
\(m/Q=r^2mV/U\). Thus \(x,y\) are bounded before taking residues.
Those equations and the two remaining kernel equations give
\[
 Ay\sim\tfrac12,\quad
 K(r-y)\sim D+\tfrac12,\quad
 D(1+y-r)\sim r/2,\qquad
 x\sim-\frac{m+\beta+1/2}{H_X}\,r.
 \tag{WA9}
\]
For the last kernel division the possible extra \(\alpha/V\),
\(nU/V\) and \(mV/U\) terms vanish by WA5. In particular \(y,r-y\)
and \(1+y-r\) are positive units, and the critical residue dictionary is
\[
 A_0=\frac1{2y_0},\qquad
 C_0=\frac{1+y_0}{2(r_0-y_0)},\qquad
 D_0=\frac{r_0}{2(1+y_0-r_0)}.
 \tag{WA10}
\]
In the unit-\(n\), weak-\(m\) domain, \(\beta=0\), and hence
\[
 P,Q=o(1),\qquad m\asymp Q^2,\qquad
 \alpha\asymp P^2,\qquad
 x_0=-r_0/(1+2n_0).
 \tag{WA11}
\]
No independent testing coordinate occurred in this entry proof.

\subsubsection*{The exact scalar and every first normal}

For reference the full finite quarter-parent is
\[
\begin{aligned}
F={}&nD_uD_vX^2+m(vW-X+uvY-uZ)^2
       +CD_v(uW-Y)^2+AD_uD_vY^2\\
 &+(D+\tfrac12)(uW+uvX-Y-vZ)^2
       +(\alpha+\tfrac12)D_u(vY-Z)^2\\
 &+(\beta+\tfrac12)D_v(X+uZ)^2-\tfrac12D_uD_vZ^2 .
\end{aligned}\tag{WA12}
\]
Here \(D_u=1+u^2,D_v=1+v^2\). We retain \(\beta\) symbolically to make
the exact common graph explicit, but the current case has \(\beta=0\).
In the testing chart \(W=1,Z=z\), use
\[
 u_g=rvz,\qquad X_g=xvz^2,\qquad Y_g=yvz,\qquad
 V_t=v^2,\quad T=z^2.
 \tag{WA13}
\]
In the rest of this subsection write \(V=V_t\) only. The full graph
restriction is
\[
\begin{aligned}
\Phi={}&mV+\alpha T+hVT+\mathcal C V^2T+EVT^2
   +\mathcal F V^2T^2+GV^3T^2+JV^2T^3+\mathcal K V^3T^3,\\
\mathcal C={}&Ay^2+C(r-y)^2+(\alpha+\tfrac12)y^2+2mry,\\
E={}&\alpha r^2+(m+\beta+\tfrac12)(r+x)^2+nx^2,\\
G={}&r^2y^2(A+\alpha+m+\tfrac12),\qquad
J=nr^2x^2,\qquad \mathcal K=(D+n+\tfrac12)r^2x^2,\\
\mathcal F={}&Ar^2y^2+(2D+1)rx(r-y-1)-(2\alpha+1)r^2y\\
 &-2mry(r+x)+\beta(r+x)^2+nx^2+rx+x^2/2 .
\end{aligned}\tag{WA14}
\]
The coefficient \(h\) is defined by this entire expansion; it is bounded.
Every displayed coefficient is bounded. In the current unit-\(n\)
domain \(\mathcal C,E\) are positive units, since
\[
 \mathcal C_0=r_0(1+y_0)/2,\qquad
 E_0=n_0r_0^2/(1+2n_0).
 \tag{WA15}
\]
The coefficient \(\mathcal F\) need not be positive.

The actual contact gives \(\Phi(P,Q)=\Phi_V(P,Q)=\Phi_T(P,Q)=0\).
Solving these three linear equations in \(m,\alpha,h\) is exact:
\[
\begin{aligned}
m={}&EQ^2+\mathcal F PQ^2+GP^2Q^2+2JPQ^3+2\mathcal K P^2Q^3,\\
\alpha={}&\mathcal C P^2+\mathcal F P^2Q+2GP^3Q
                         +JP^2Q^2+2\mathcal K P^3Q^2,\\
h={}&-2\mathcal CP-2EQ-3\mathcal FPQ-4GP^2Q
                         -4JPQ^2-5\mathcal K P^2Q^2 .
\end{aligned}\tag{WA16}
\]
Put
\[
 R=z(V-P),\quad S=v(T-Q),\quad
 \Psi=QR^2+PS^2+vzRS .
\]
With the exact balanced monomials
\[
\begin{aligned}
B_{32}&=2PQR^2+(vzR+PS)^2,\\
B_{23}&=2PQS^2+(QR+vzS)^2,\\
B_{33}&=(2PQ^2+VT^2)R^2+(2P^2Q+PVT)S^2
                              +(3PQ+PT)vzRS,
\end{aligned}\tag{WA17}
\]
the whole scalar, not a leading truncation, is
\[
 \Phi=\mathcal CR^2+ES^2+\mathcal F\Psi
                          +GB_{32}+JB_{23}+\mathcal K B_{33}.
 \tag{WA18}
\]
For \(v,z\) infinitesimal, all entries after the first two positive
unit pivots are infinitesimal, by \(P,Q=o(1)\). This is a strict
two-row Gram even when \(\mathcal F\) is signed.

Let \(L_u,L_X,L_Y\) be half the full first derivatives of WA12 at
WA13. The complete polynomials have the forms
\[
\begin{aligned}
 L_u&=vz A_u(V,T),\qquad L_Y=vz A_y(V,T),\\
 A_u&=a_{00}+a_{10}V+a_{01}T+a_{11}VT
                      +a_{12}VT^2+a_{21}V^2T+a_{22}V^2T^2,\\
 A_y&=e_{00}+e_{10}V+e_{11}VT+e_{21}V^2T,\\
 L_X&=v[-m+TB_X(V,T)],\qquad
 B_X=c_0+c_1V+c_2VT+c_3V^2T .
\end{aligned}\tag{WA19}
\]
All coefficients are bounded, and the expansion fixes every one
unambiguously. In particular
\[
 c_2=nr^2x,\qquad c_3=(D+n+\tfrac12)r^2x .
\]
Their actual source jets are \(A_u(P,Q)=A_y(P,Q)=0\) and
\(B_X(P,Q)=m/Q\). Finite monic difference division gives
\[
\begin{aligned}
 L_i&=v A_{i,V}(V,T)R+z B_i(T)S,\qquad i=u,Y,\\
 L_X&=vz B_{X,V}(V,T)R+
                \{T(Pc_2+P^2c_3)+m/Q\}S .
\end{aligned}\tag{WA20}
\]
Here \(A_{i,V}=[A_i(V,T)-A_i(P,T)]/(V-P)\), and \(B_{X,V}\)
is defined identically. The other divided polynomials are
\[
\begin{aligned}
B_u(T)&=a_{01}+Pa_{11}+Pa_{12}(T+Q)+P^2a_{21}
                                      +P^2a_{22}(T+Q),\\
B_Y(T)&=Pe_{11}+P^2e_{21}.
\end{aligned}
\]
These are literal polynomial identities with bounded coefficients,
including \(v=0\) or \(z=0\). Since \(m/Q=O(Q)=o(1)\),
every full first normal is an infinitesimal linear combination of \(R,S\).

Set \(N_u=u-rvz,N_X=X-xvz^2,N_Y=Y-yvz\). In the chart under
consideration these are infinitesimal. The full higher-normal
remainder is a polynomial in the three normals with no terms of
normal degree below two; all terms of degree at least three give
infinitesimal Gram entries. Its quadratic residue is
\[
 K_0(N_u-N_Y)^2+A_0N_Y^2+(n_0+\tfrac12)N_X^2\succ0.
 \tag{WA21}
\]
WA18 and WA20 therefore give a strict bounded Gram for the entire
parent on the three hard and two soft rows. The strict-Gram lemma
proves its tensor-SOS assertion, with no mixed inverse.

\subsubsection*{The complementary local states}

Normalize a maximal testing state coordinate. At \(u,v=o(1)\)
the state residue has positive pivots on \(X,Y\); if one is a unit,
the full parent has positive residue immediately. Otherwise one
of \(W,Z\) is a unit. The preceding proof covers \(W=1,Z=o(1)\).
Every remaining state can be normalized to \(Z=1,W=\omega\)
bounded. Retain the entire positive \(m,\alpha\) squares.
The remainder lies exactly in the square of \((u,v,X,Y)\).
Its quadratic real residue is
\[
\begin{aligned}
 &(n_0+\tfrac12)\left(X+\frac{u}{1+2n_0}\right)^2
     +\frac{n_0}{1+2n_0}u^2\\
 &+H_0\left(Y-\frac{K_0\omega u-D_0v}{H_0}\right)^2
     +a_0(\omega u-r_0v)^2,\qquad a_0=A_0K_0/H_0>0 .
\end{aligned}\tag{WA22}
\]
The critical identities WA10 prove the equality. This quadratic
form is positive definite for every real \(\omega\), including zero.
All higher normal terms are infinitesimal corrections on these
four rows, while \(\omega-\bar\omega\) is retained in coefficients.
Thus every state at the double-infinitesimal angle corner is covered.

\subsubsection*{All other testing angles, including the literal endpoints}

Again retain the entire \(m,\alpha\) squares and call the remainder
\(F_0\). Its source residue is the critical \(C,A,D,n\) source.
For bounded \(u,v\), \(F_0\) is exactly a bounded quadratic form
in the four rows \((uW,X,Y,Z)\). Its real \(Z=0\) block is positive:
the \(nX^2\) and \(AY^2\) squares and \(C(uW-Y)^2\) suffice.
The determinant is
\[
 \det \operatorname{Mat}_{(uW,X,Y,Z)}(\bar F_0)
 =D_u^2D_v^2\{q_{20}u^2+q_{12}v^4+q_{21}u^2v^2\}/16>0
 \tag{WA23}
\]
whenever at least one of the two real angle residues is nonzero.
Here \(q_{20},q_{12},q_{21}>0\) are the actual residue determinant
coefficients, and \(h=0\) follows from WA10. This proves strictness
of the full four-row Gram. Literal \(u=0\) is allowed.

For Bob reciprocal \(b=1/v\) infinitesimal and Alice bounded,
use the same rows \((uW,X,Y,Z)\) in \(b^2F_0(u,1/b)\).
At \(b=0\) its whole real expression is
\[
 C_0(uW-Y)^2+(A_0+\tfrac12)D_uY^2
             +D_0(uX-Z)^2+(n_0+\tfrac12)D_uX^2\succ0 .
 \tag{WA24}
\]
It is strict at every real \(u\), including zero.
For Alice reciprocal \(s=1/u\) infinitesimal and Bob bounded,
the original state rows suffice. The expression at \(s=0\) is
\[
 n_0D_vX^2+C_0D_vW^2+A_0D_vY^2
          +(D_0+\tfrac12)(W+vX)^2+\tfrac12(vY-Z)^2\succ0 .
 \tag{WA25}
\]
This includes \(v=0\), using the positive \(n_0\) pivot.

For both reciprocal angles \(s,b=o(1)\), use the five rows
\((W,X,Y,sZ,bZ)\). The whole reciprocal polynomial belongs
to their square ideal, and its real coefficient Gram is
\[
 C_0W^2+A_0Y^2+D_0(X-sZ)^2
          +(n_0+\tfrac12)X^2+\tfrac12(Y-bZ)^2\succ0 .
 \tag{WA26}
\]
Every remaining coefficient is infinitesimal. These exact finite
square-ideal allocations include \(s=0\), \(b=0\) and arbitrary
states; no projective endpoint is obtained by a limiting argument.

The four angle charts and the two local state charts are exhaustive.
Testing sign gauges, positive pure-field homogeneous normalization
factors and real/imaginary state splitting restore the original parent:
\[
 F_c(u,v;\psi)\ \in\ \Sigma(K_c\otimes_{\mathbb R}K_t)^2
 \quad\text{for every independent projective test and state.}
 \tag{WA27}
\]
This closes exactly the C2 unit-\(n\), weak-\(m,\alpha\) case.
The common source-entry subsection does not by itself close the
weak-\(n\) or C1 weak-\(\alpha\) branches.



\section{\texorpdfstring{Weak-alpha weak-\(n\) and C1: the complete parent}{Weak-alpha weak-n and C1: the complete parent}}
\label{component:p0db05df21ad1}

\label{p0db05df21ad1:sec:weak-alpha-weak-n-c1}

\paragraph{The source domain.}
The source is an actual physical canonical finite noncommuting contact,
all coefficients are bounded in the source field, and \(d=1\).
Either it is C2 with \(\lambda=\beta=0\) and \(n,\alpha\) infinitesimal,
or it is C1 with \(\lambda=n=0\), \(\alpha\) infinitesimal and arbitrary
bounded \(\beta\geq0\). All strict source gaps retain their original
meaning. We prove the tensor-SOS statement for the entire parent at
every independent projective testing angle and every state.
The common source-entry proof WA1--WA10 is used on precisely its
weak-\(n\) domain. In particular \(C,A,D,r,y,r-y,1+y-r\) are positive
units, \(x\) is bounded, and \(x_0=-r_0\).
Boundedness, infinitesimality and residue refer
to the separated bounded tensor ring and its real residue map.
The strict-Gram lemma is applied only after a finite explicit row
allocation; no mixed tensor is inverted.

Use the full quarter-parent WA12 and its exact source coordinates
\[
 p=v_c,\quad q=Z_c>0,\quad P=p^2,\quad Q=q^2,\qquad
 W_c=1,\quad u_c=rpq,\quad X_c=xpq^2,\quad Y_c=ypq.
 \tag{WN1}
\]
Thus \(P,PQ\) are positive infinitesimals. Put
\[
 H_0=m+\beta+\tfrac12,\quad H_X=H_0+n,\quad
 K=C+D+\tfrac12,\quad H_Y=K+A.
 \tag{WN2}
\]
These are positive units. The exact scalar coefficients
\(\mathcal C,E,\mathcal F,G,J,\mathcal K\), the contact balances, and
the complete first normal polynomials are those displayed in
WA14, WA16 and WA19; those identities retain \(\beta\) and \(n\).
Their relevant residues are
\[
 \mathcal C_0=\frac{r_0(1+y_0+4m_0y_0)}2>0,\qquad
 \mathcal F_0=\frac{r_0^2(1+y_0)}2>0,
 \qquad G,J,\mathcal K\geq0 .
 \tag{WN3}
\]
Both \(\mathcal C,\mathcal F\) are consequently positive units.
The exact source \(X\)-jet gives, with the WA19 coefficients,
\[
\begin{gathered}
 c_0=H_0(r+x)+nx,\quad c_2=nr^2x,\quad
 c_3=(D+n+\tfrac12)r^2x,\\
 c_0=m/Q-Pc_1-PQc_2-P^2Qc_3,\\
 E=\alpha r^2+c_0^2/H_X+nH_0r^2/H_X .
\end{gathered}\tag{WN4}
\]
The last equality is an exact completion inside the source field.

\subsubsection*{Two finite positive circuit identities}

For \(V=v^2,T=z^2\) put
\[
 R=z(V-P),\quad S=v(T-Q),\quad
 \Psi=V^2T^2+PQ^2V+P^2QT-3PQVT .
\]
Then \(\Psi=QR^2+PS^2+vzRS\). It is an explicit six-square
cubic circuit: if \(L_1^3=V^2T^2,L_2^3=PQ^2V,L_3^3=P^2QT\),
then \(L_1L_2L_3=PQVT\), and
\[
 \Psi=\sum_{i=1}^3 L_i
 \bigl\{(L_1-L_2)^2-(L_1-L_2)(L_1-L_3)+(L_1-L_3)^2\bigr\}.
 \tag{WN5}
\]
The real \(2\)-by-\(2\) Gram in each brace is strictly positive.
All radicals in this formula are elementary separated monomials:
source radicals and testing radicals are taken in their own real
closed fields and then tensored.

The balanced degree-\((3,3)\) term has the polynomial identity
\[
 a^{10}+2b^{10}+2c^{10}-5a^2b^4c^4
 =2(b^5-c^5)^2
 +(a^2-bc)^2(a^6+2a^4bc+3a^2b^2c^2+4b^3c^3).
 \tag{WN6}
\]
For nonnegative elementary monomials \(a,b,c\), the bracket is the
sum of squares of the four rows
\[
 a^3,\qquad \sqrt2a^2\sqrt{bc},\qquad
 \sqrt3abc,\qquad 2(bc)^{3/2}.
\]
This is a five-square separated certificate, including vanishing
variables. Applied with
\[
 a^{10}=V^3T^3,\quad b^{10}=P^2Q^3V,\quad c^{10}=P^3Q^2T,
\]
it proves positivity of the entire \(B_{33}\) in WA17.

A further exact identity supplies any fixed positive real pivot
on the row \(PS\):
\[
 2P\Psi+VT R^2-(PS)^2=(PS+vzR)^2+2PQR^2.
 \tag{WN7}
\]
Indeed, for every fixed \(b>0\),
\[
\begin{split}
 \mathcal C R^2+\mathcal F\Psi={}&
 (\mathcal C-bVT+2bPQ)R^2+(\mathcal F-2bP)\Psi\\
 &+b(PS)^2+b(PS+vzR)^2.
\end{split}\tag{WN8}
\]
On the charts where this is used, \(VT,P,PQ\) are infinitesimal,
so the original two unit residues are unchanged.

\subsubsection*{Bounded source ratio \(Q\)}

From WA5, WN4 and bounded \(Q\),
\[
 m\asymp Q^2(n+P),\quad \alpha\asymp P^2,\quad
 r+x=O(n+P),\quad E\asymp n+P^2.
 \tag{WN9}
\]
In particular \(m,n,E,\alpha\) are infinitesimal. None of these
estimates contains an independent testing coordinate.

First take \(W=1,Z=z\) bounded and \(u,v\) infinitesimal. Use the
entire graph \(u=rvz,X=xvz^2,Y=yvz\). The exact scalar is
\[
 \Phi=\mathcal CR^2+ES^2+\mathcal F\Psi
          +GB_{32}+JB_{23}+\mathcal K B_{33},
 \tag{WN10}
\]
where the last three terms retain their positive full certificates.
The strict soft rows are \(R\), the six cubic rows in WN5,
\(\sqrt E\,S\), and \(PS\); WN8 supplies any needed real pivot
on the last row.

The exact monic normal divisions WA20 apply unchanged.
Their \(R\)-coefficients contain \(v\) or \(vz\), hence are
infinitesimal. Their \(S\)-coefficients are \(O(E+P)\):
in particular
\[
 a_{01}=\alpha r+H_0(r+x)=O(E+P),\qquad
 a_{12}=nr x^2,\qquad m/Q=O(E+P).
 \tag{WN11}
\]
Every \(ES\) contribution has infinitesimal coefficient \(\sqrt E\)
on \(\sqrt E S\). Coefficients on \(PS\) may be bounded units.
There are only finitely many of them, so choose \(b\) in WN8
larger than their finite real Schur cost. This uses a real constant,
not an inverse of a mixed tensor.

The three hard normals
\(N_u=u-rvz,N_X=X-xvz^2,N_Y=Y-yvz\) are infinitesimal.
Their complete quadratic residue is
\[
 K_0(N_u-N_Y)^2+A_0N_Y^2+
          (\beta_0+\tfrac12)(N_X+\bar zN_u)^2\succ0.
 \tag{WN12}
\]
All higher-normal terms have degree at least three and bounded
coefficients, hence allocate as infinitesimal Gram entries on these
hard rows. WN10--WN12 prove the full parent in this chart.

For the complementary local states take \(Z=1,W=w\) infinitesimal,
set \(a=u/r\), and use
\[
 v=aw,\quad X=xa,\quad Y=yaw,\qquad U=a^2,\quad T=w^2.
 \tag{WN13}
\]
The entire scalar is
\[
\begin{split}
 \Phi'={}&\alpha+mUT^2+EU+hUT+\mathcal C U^2T^2
       +\mathcal F U^2T+GU^3T^2+JU^2+\mathcal K U^3T .
\end{split}\tag{WN14}
\]
Here \(a=(1/r)\otimes u\) is elementary; its roots remain separated.
The strict rows are \(P,a^2w,\sqrt E\,a\).
The term \(hUT\) is harmless: the contact balances give
\(h=O(P+E)\), and writing \(h=Ph_P+Eh_E\) with bounded source
coefficients gives the infinitesimal Gram terms
\(h_Pw\,P(a^2w)+h_Ew^2(\sqrt E a)^2\).
All other terms of WN14 are positive.

For completeness, half the entire \(v,Y,X\) derivatives have
the polynomial forms
\[
\begin{aligned}
 L_v'&=aw(d_0+mT+d_1U+d_2UT+d_3U^2+d_4U^2T),\\
 L_Y'&=aw(e_0+e_1U+e_2UT+e_3U^2T),\\
 L_X'&=a(c_0-mT+c_2U+c_1UT+c_3U^2T).
\end{aligned}\tag{WN15}
\]
Every coefficient is bounded. Source jets at
\(U=PQ,T=1/Q\) give
\[
\begin{split}
d_0&=-m/Q-d_1PQ-d_2P-d_3P^2Q^2-d_4P^2Q=O(n+P),\\
e_0&=-e_1PQ-e_2P-e_3P^2Q=O(P).
\end{split}
\]
Also \(c_0,m=O(n+P)\) and \(c_2=O(n)\).
Splitting \(n+P=O(E+P)\) in these displayed complete polynomials
allocates all first normals with infinitesimal coefficients to
\(P,a^2w,\sqrt E a\). For example \(awP\), \(awE\), and
\(a^3w\) have coefficients \(aw\), \(w\sqrt E\), and \(a\)
on those three rows; the other monomials have still more small
testing factors. The complete hard residue is
\[
 (\beta_0+\tfrac12)N_X^2+
 D_0N_v^2+2D_0N_vN_Y+H_{Y,0}N_Y^2\succ0,
 \tag{WN16}
\]
where \(N_v=v-aw,N_Y=Y-yaw,N_X=X-xa\).
The higher-normal argument is again exact polynomial allocation.
This proves all local states when \(Q\) is bounded.

\subsubsection*{Unbounded source ratio: the \(W\)-chart}

Suppose \(Q\) is unbounded. The source estimates and positive
contact balances give
\[
 PQ=o(1),\quad m\asymp Q^2(n+P),\quad nQ^2,PQ^2=O(1),
 \quad \alpha\asymp P^2Q .
 \tag{WN17}
\]
WN4 now gives
\[
 c_0\sim m/Q,\quad r+x\sim m/(H_XQ),\quad
 E\asymp n+(m/Q)^2,\quad h/m=o(1).
 \tag{WN18}
\]
To justify the first comparison before its use, divide the three
errors \(Pc_1,PQc_2,P^2Qc_3\) by \(m/Q\).
Their absolute values are \(O(1/Q),O(P),O(P)\), respectively;
\(n/(m/Q)=O(1/Q)\). The positive term
\(\alpha/(m/Q)^2=O(1/Q)\) follows from
\(m\geq cPQ^2\). The last comparison follows term by term from
WA16 using these same inequalities.

Take \(W=1,Z=z\) bounded, \(u,v\) infinitesimal. If \(m\) is a
unit, retain the entire \(n,\alpha\) squares. The remainder lies
in the square of \((u,v,X,Y)\). Its real quadratic form is
\[
\begin{split}
 &H_0\{X+zu-(m/H_0)v\}^2+
       \frac{m(\beta+1/2)}{H_0}v^2\\
 &\qquad+H_Y\{Y-(Ku-Dzv)/H_Y\}^2
       +a_*(u-rzv)^2,\qquad a_*=AK/H_Y>0,
\end{split}\tag{WN19}
\]
with coefficients interpreted in residue. WA10 verifies this
identity. It is strict for every bounded \(z\), so the complete
higher terms give a strict Gram.

If \(m\) is weak, use the graph WA13. Its entire scalar has strict
rows \(\sqrt m\,v,v^2z,\sqrt\alpha\,z\), since
\(mV+hVT=mV[1+(h/m)T]\), while \(\mathcal C,\mathcal F,E,G,J,\mathcal K\)
are nonnegative. The complete source jet equations yield
\[
 a_{00}=O(m),\quad e_{00}=O(PQ),\quad
 a_{01},c_0=O(m/Q).
 \tag{WN20}
\]
Thus \(a_{00}/\sqrt m\), \(e_{00}/\sqrt m\),
\(a_{01}/\sqrt m\), \(c_0/\sqrt m\) are infinitesimal.
For \(e_{00}\), use \(PQ/\sqrt m=O(\sqrt P)\).
Every other monomial of the complete normals WA19 has a factor
\(v^3z\) or \(v^3z^2\), so it has an infinitesimal coefficient
on \(v^2z\). The hard residue is WN12; all higher terms are
allocated exactly as there. This proves the whole \(W\)-chart.

\subsubsection*{Unbounded source ratio: the complementary scalar}

Use WN13 with \(a,w\) infinitesimal. Introduce only source quantities
\[
\begin{gathered}
 \epsilon=Q^{-1}=o(1),\quad \Lambda=PQ=o(1),\quad
 \tau=\Lambda/\epsilon=PQ^2=O(1),\\
 \nu=n/\epsilon^2=O(1),\qquad N=E/\epsilon^2=O(1).
\end{gathered}\tag{WN21}
\]
The quantities \(\epsilon,\Lambda,\tau,N\) are positive; \(\nu\geq0\).
Let
\[
 A_1=a(T-\epsilon),\quad K_1=UT-\Lambda\epsilon,\quad
 \mathcal M=U^2T+\tau UT^2+\Lambda^2\epsilon-3\Lambda UT .
 \tag{WN22}
\]
The whole scalar WN14, using exactly the three contact balances, is
\[
\begin{split}
 \Phi'={}&\mathcal C K_1^2+NA_1^2+\mathcal F\mathcal M\\
 &+G\{2\Lambda K_1^2+UT^2(U-\Lambda)^2\}\\
 &+J\{(U-\Lambda)^2+2\Lambda A_1^2/\epsilon^2\}
       +\mathcal K\mathcal B_{33},\\
 \mathcal B_{33}={}&U^3T+2(\Lambda^2/\epsilon)UT^2
                      +2\Lambda^3\epsilon-5\Lambda^2UT .
\end{split}\tag{WN23}
\]
The apparent \(1/\epsilon^2\) in the positive \(J\) square is
a source-field quantity, and \(J/\epsilon^2=\nu r^2x^2\) is bounded.
All positive terms are retained. WN6 applies to \(\mathcal B_{33}\)
with \(a^{10}=U^3T,b^{10}=(\Lambda^2/\epsilon)UT^2,
c^{10}=\Lambda^3\epsilon\).
Here \(\Lambda^2/\epsilon=\tau\Lambda\) is bounded.

Write \(a=t^3,w=s^3,\epsilon=h^6,\tau=c^6\), taking positive
elementary roots. Then \(t,s,h\) are infinitesimal and \(c\)
is bounded; it is allowed to be a unit. Put
\[
\begin{gathered}
 L_1=t^4s^2,\quad L_2=c^2t^2s^4,\quad L_3=h^6c^4,\\
 f_1=t^2s,\quad f_2=cts^2,\quad f_3=h^3c^2,\\
 R_{2i-1}=f_i(L_1-L_2),\qquad R_{2i}=f_i(L_1-L_3)
 \quad (i=1,2,3).
\end{gathered}\tag{WN24}
\]
Then \(\mathcal M=\sum_{i=1}^3(R_{2i-1}^2-R_{2i-1}R_{2i}+R_{2i}^2)\),
a strict six-row constant Gram. The following are exact identities:
\[
\begin{aligned}
 aw(U-\Lambda)&=t^3R_1+c^2ts^2R_2,\\
 a^2w(T-\epsilon)&=-h^6R_1+s^6R_2-h^3s^3R_5,\\
 \tau aw(T-\epsilon)&=c^2ts^2(R_2-R_1)-c^3s^3R_3,\\
 \tau\sqrt\epsilon\,A_1&=t^3R_6-t(t^2+c^2s^2)R_5 .
\end{aligned}\tag{WN25}
\]
Every coefficient on the right is infinitesimal, even when \(c\)
is a unit. These are finite linear conductors, not limiting
approximations. The further quadratic conductor is
\[
 \boxed{\ \Lambda(U-\Lambda)K_1
       =\Lambda\mathcal M-(\tau\sqrt\epsilon\,A_1)^2.\ }
 \tag{WN26}
\]
Its right side is an infinitesimal Gram on the six cubic rows.
It remains so after multiplication by any bounded signed coefficient.

If \(\tau\) is a unit, then \(K_1\) itself has an infinitesimal
linear representation on those rows:
\[
 K_1=\tau^{-1/3}
       \{\sqrt{L_1}(R_2-R_1)+\sqrt{L_3}R_6\}.
 \tag{WN27}
\]
This follows by expansion as
\(\tau^{-1/3}(L_1L_2-L_3^2)=UT-\Lambda\epsilon\).
Consequently any fixed real pivot on \(K_1\) can be supplied at
infinitesimal cost to the strict \(\mathcal M\) Gram in this case.

\subsubsection*{Every complementary first normal}

Let \(D_U,E_U\) denote the monic \(U\)-divided differences of
the bracket polynomials in the first two lines of WN15.
The source jets at \(U=\Lambda,T=\epsilon\) give the exact
identities
\[
\begin{aligned}
 L_v'={}&aw\{D_U(U,T)(U-\Lambda)
                 +(m+\Lambda d_2+\Lambda^2d_4)(T-\epsilon)\},\\
 L_Y'={}&aw\{E_U(U,T)(U-\Lambda)
                 +(\Lambda e_2+\Lambda^2e_3)(T-\epsilon)\},\\
 L_X'={}&-mA_1+
       a(c_2+c_3\Lambda\epsilon)(U-\Lambda)
                  +a(c_1+c_3U)K_1 .
\end{aligned}\tag{WN28}
\]
There are no omitted monomials. Define
\[
 d_*=(c_2+c_3\Lambda\epsilon)/\epsilon
       =\nu\epsilon r^2x+\Lambda c_3 .
\]
Since \(\epsilon(U-\Lambda)=K_1-U(T-\epsilon)\), the last line is
\[
 L_X'=-mA_1+P_X,\qquad
 P_X=a(c_1+c_3U+d_*)K_1-d_*UA_1 .
 \tag{WN29}
\]
The exact source equations also give
\[
\begin{aligned}
 m&=N+\tau(\mathcal F+G\Lambda\epsilon+2\mathcal K\Lambda)
                            +2\nu r^2x^2\Lambda,\\
 c_0/\epsilon&=m-\Lambda B,\qquad
 B=c_1+\Lambda c_3+\nu\epsilon r^2x,\\
 N&=\alpha r^2/\epsilon^2+
             (m-\Lambda B)^2/H_X+\nu H_0r^2/H_X .
\end{aligned}\tag{WN30}
\]
In particular
\[
 N,\nu,\tau=O(m),\quad \Lambda=O(m\epsilon),\quad
 N=m^2/H_X+\nu H_0r^2/H_X+O(m\Lambda).
 \tag{WN31}
\]
For the error term use \(\alpha/\epsilon^2=O(\tau\Lambda)\),
which is the exact positive alpha balance WN23/WA16.
All these comparisons are inside the source field.
We split by the bounded source ratio \(\nu/m\); \(m>0\).

\paragraph{Substantial \(\nu/m\).}
If \(\nu/m\) has positive residue, keep \(S_0=\sqrt m A_1\).
The \(N_X,S_0\) Gram, including the entire \(-mA_1\) first
normal, has residue from
\[
 \begin{pmatrix}H_X&-\sqrt m\\-\sqrt m&N/m\end{pmatrix}.
 \tag{WN32}
\]
Its determinant is
\((\nu/m)H_0r^2+O(\Lambda)\), a positive unit.
The \(v,Y\) first normals in WN28 use the first conductor WN25;
their remaining coefficients on \(S_0\) are
\(O(\sqrt m\,w)\), hence infinitesimal.
In WN29 the \(K_1\)-coefficient contains \(a\), and
\(d_*U/\sqrt m=O(\sqrt m\,\epsilon U)\) is infinitesimal.
Thus all other first normals have infinitesimal coefficients
on \(K_1\), the six cubic rows, and \(S_0\).
The \(v,Y\) hard block in WN16 is strict. All higher terms have
at least three infinitesimal hard normals and bounded coefficients.
This proves the full parent on this source subdomain.

\paragraph{Thin \(\nu/m\), including literal \(n=0\).}
Assume \(\nu/m=o(1)\). Equations WN30--WN31 imply
\(\tau/m\) is a positive unit: after division by \(m\),
\[
 1=N/m+(\tau/m)\mathcal F+o(1),\qquad
 N/m=m/H_X+o(1),
\]
and \(1-m/H_X=(n+\beta+1/2)/H_X\) has positive residue.
Write
\[
 N-m^2/H_X=\nu\xi+\eta,\qquad
 \xi=H_0r^2/H_X>0,\qquad \eta=O(m\Lambda).
\]
Complete only the source coefficient \(H_X\):
\[
\begin{split}
 H_XN_X^2-2mN_XA_1+NA_1^2
 &=H_X\widetilde N_X^2+\nu\xi A_1^2+\eta A_1^2,\\
 \widetilde N_X&=N_X-(m/H_X)A_1.
\end{split}\tag{WN33}
\]
Retain the strict row \(\sqrt\nu A_1\) when \(\nu>0\);
if \(\nu=0\), all terms allocated to it vanish literally.
The signed error is an infinitesimal cubic Gram, since
\[
 m\Lambda A_1^2=(m/\tau)(\tau\sqrt\epsilon\,A_1)^2.
 \tag{WN34}
\]
The first two normals in WN28 use WN25 and bounded \(m/\tau\).
For \(P_X\), its \(K_1\)-coefficient contains \(a\), while its
two other contributions have infinitesimal allocations
\[
 \nu\epsilon U A_1=(\sqrt\nu\,\epsilon U)(\sqrt\nu A_1),
 \qquad \Lambda U A_1=(\sqrt\epsilon\,U)
                         (\tau\sqrt\epsilon\,A_1).
 \tag{WN35}
\]
Therefore \(2P_X\widetilde N_X\) has infinitesimal Gram entries.

The shift in WN33 also creates the entire scalar term
\(2(m/H_X)A_1P_X\), which must not be dropped.
Its main factor is controlled by the exact identity
\[
 mU(T-\epsilon)K_1
   =mK_1^2-(m/\tau)\Lambda(U-\Lambda)K_1 .
 \tag{WN36}
\]
The second term is infinitesimal by WN26. If \(m\) is
infinitesimal, the first is an infinitesimal \(K_1\)-Gram.
If \(m\) is a unit, then \(\tau\) is a unit and WN27
makes it an infinitesimal cubic Gram. In that case a positive
real \(K_1\)-pivot may be reinserted by WN27 at infinitesimal
cubic cost. This allocation remains valid with each bounded
factor \(c_1+c_3U+d_*\).
The remaining scalar factors are multiples of
\(m\nu\epsilon U A_1^2\) and \(m\Lambda U A_1^2\);
they have infinitesimal Gram entries on \(\sqrt\nu A_1\)
and the fourth conductor in WN25, respectively.
This accounts for the whole scalar effect of the completion.

\subsubsection*{The entire higher-normal restoration in the thin case}

The full parent is quadratic in \(X\). Its coefficient beyond
the source constant \(H_X\) is exactly the nonnegative polynomial
\[
 D_X=(n+\beta+\tfrac12)v^2+nu^2+
                         (n+D+\tfrac12)u^2v^2 .
 \tag{WN37}
\]
Retain the entire square \(D_XN_X^2\), evaluated at the actual
independent \(u,v\), as separated squares. Each monomial coefficient
is a nonnegative source scalar times a testing square. It remains
a positive whole square after WN33.
The remaining polynomial has no quadratic \(N_X\) term except
\(H_XN_X^2\). At the complementary graph its half mixed
derivatives have the exact forms
\[
 [N_XN_v]=w[-m+a^2B_v],\qquad
 [N_XN_Y]=a^2wB_Y,
 \tag{WN38}
\]
where \(B_v,B_Y\) are bounded polynomials fixed by the full
parent expansion. After the shift, the new first-normal terms
are bounded multiples of \(\tau aw(T-\epsilon)\), using
\(m/\tau=O(1)\); WN25 gives infinitesimal coefficients.
Every other remaining monomial linear in \(N_X\) has degree at
least two in the other hard normals. Replacing \(N_X\) by
\(\widetilde N_X+(m/H_X)A_1\) gives an infinitesimal coefficient
on a product of at least two hard normals, since \(A_1=o(1)\).
Terms involving \(\widetilde N_X\) have total hard-normal
degree at least three and are handled in the same way.
Terms independent of \(N_X\) are unchanged. The full real hard block is
\[
 H_{X,0}\widetilde N_X^2+
       D_0N_v^2+2D_0N_vN_Y+H_{Y,0}N_Y^2\succ0.
 \tag{WN39}
\]
Equations WN23--WN39 therefore give a finite strict Gram for
the whole parent, including all first normals, every higher term,
all \(\beta\)-openings and literal \(n=0\).
This completes the double-infinitesimal testing corner for both
source ratios and all local states.

\subsubsection*{Every other angle, all states and literal endpoints}

Retain the entire \(m,n,\alpha,\beta\) squares and call the
remainder \(F_0\). Only its residue is the critical \(C,A,D\)
source in WA10; no global physicality of the subtracted finite
source is asserted. The following are complete polynomial row
allocations, so their strict residue Grams apply to the original
nearby coefficients.

For bounded \(u,v\) with \(v\) a unit use \((uW,X,Y,Z)\).
The \(Z=0\) principal block is positive definite and
\[
 \det\operatorname{Mat}_{(uW,X,Y,Z)}\bar F_0
 =D_u^2D_v^2(q_{12}v^4+q_{21}u^2v^2)/16>0.
 \tag{WN40}
\]
Both \(q_{12},q_{21}\) are positive units. If \(u\) is a unit
and \(v\) infinitesimal, use instead
\((uW-Y,Y,X+uZ,vZ)\). The entire parent is polynomial in
these rows with bounded coefficients; the only algebraic
coefficient divisions are by the pure testing unit \(u\).
Its determinant at \(v=0\) is
\(D_u^2q_{21}u^2/16>0\), and the first three-row principal block
is strict. This covers all bounded angle pairs outside the
double-infinitesimal corner.

For Bob reciprocal \(b=1/v=o(1)\), Alice bounded, the rows
\((uW,X,Y,Z)\) give the real form
\[
 C_0(uW-Y)^2+(A_0+\tfrac12)D_uY^2+
                 D_0(uX-Z)^2+\tfrac12D_uX^2\succ0.
 \tag{WN41}
\]
For Alice reciprocal \(s=1/u=o(1)\) and Bob a bounded unit,
the original state rows give
\[
 C_0D_vW^2+A_0D_vY^2+
       (D_0+\tfrac12)(W+vX)^2+\tfrac12(vY-Z)^2\succ0.
\]
If \(s,v\) are both infinitesimal, use
\((W,Y,vX,\Gamma)\), where \(\Gamma=sX+Z\).
The entire reciprocal parent is polynomial in these four
un-divided rows, and its real form is
\[
 C_0W^2+A_0Y^2+
                 (D_0+\tfrac12)(W+vX)^2+\tfrac12\Gamma^2\succ0.
 \tag{WN42}
\]
The coefficient-extraction substitution \(X=R/v,Z=\Gamma-sR/v\)
cancels all denominators; the asserted identity itself contains
no \(v^{-1}\), and includes literal \(v=0\).

For both reciprocal angles \(s,b=o(1)\), the full polynomial
belongs to the square ideal of \((W,X,Y,sZ,bZ)\), with real form
\[
 C_0W^2+A_0Y^2+D_0(X-sZ)^2+
                    \tfrac12X^2+\tfrac12(Y-bZ)^2\succ0.
 \tag{WN43}
\]
These finite identities include \(s=0\) and \(b=0\) exactly.

It remains to check that the local state charts were exhaustive.
Normalize a maximal state coordinate in the testing field.
At \(u,v=o(1)\), a unit \(X\) or \(Y\) gives positive full
residue immediately. Otherwise one of \(W,Z\) is a unit.
The exact sign gauges give \(u,v\geq0\).
If \(WZ\leq0\), the full-core sign identity SE9 supplies
separated nonnegative terms; all event squares are positive.
If \(WZ>0\), choose \(W=1,Z\) bounded, or \(Z=1,W=o(1)\).
These are exactly the states covered above; vanishing angles
or states are included by the same polynomial identities
(or by the nonpositive-sign case).
Positive pure-field homogeneous normalization and real/imaginary
state splitting now restore the original full parent:
\[
 F_c(\hbox{any independent projective test and state})
       \ \in\ \Sigma(K_c\otimes_{\mathbb R}K_t)^2 .
 \tag{WN44}
\]
This proves precisely the weak-\(n\) C2 and all-\(\beta\) C1
weak-\(\alpha\) source domains in the source domain stated above.
It supplies the missing local leaves usually called WB, WU
and CW. It does not, by itself, assert the global source
partition or the finite-lift conclusion.



\section{\texorpdfstring{A source-exact projective pitchfork for C1}{A source-exact projective pitchfork for C1}}
\label{component:p2cdf6197677b}

\label{p2cdf6197677b:sec:new-c1-projective-pitchfork}

\paragraph{Theorem and precise source domain.}
Take a bounded physical canonical C1 source with $d=1$,
$\lambda=n=0$, $m,C,A,\alpha$ positive bounded units and
$D,\beta\geq0$ infinitesimal. Assume the original actual finite
noncommuting stationary zero contact, and retain every original
source gap and intersection. Then the entire parent has a finite
sum-of-squares certificate in $K_c\otimes_{\mathbb R}K_t$ for
every independent projective test and every real or complex state.
This includes literal $D=0$ or $\beta=0$ whenever allowed by the
source domain. The two cannot vanish together at an actual contact.

\subsubsection*{Opening-safe source entry}

Write the positive source contact as $(u_c,w_c)=(c,1/v_c)$,
and $V=w_c^2$. The retained positive-$d$ balance proof gives
$c$ a positive bounded unit and $V$ a positive infinitesimal.
Its second stationarity equation, in the notation of that proof, is
\[
 a_3=V^{-2}\{mDa_2/c^2+\beta ma_4+Da_5+\beta a_7c^2\},
 \tag{PP1}
\]
where all $a_i$ are positive bounded units. Therefore
\[
 D=V^2\delta,\quad\beta=V^2\eta,\qquad
 \delta,\eta\geq0\text{ bounded},\qquad
 \delta+\eta\text{ a positive unit}.
 \tag{PP2}
\]
No inverse of either opening is used.

Normalize the source state by $Z_c=1$. Set
\[
 p=W_c/w_c,\qquad y=Y_c/w_c,\qquad x=X_c/V.
 \tag{PP3}
\]
All three are bounded source elements. To prove this without a
shape hypothesis, restrict the reciprocal full parent PP6 to
$(W,Y,X)$ at $w=0$, and then take residues. Its positive hard form is
\[
 H_u(W,Y,X)=m(W+uY)^2+C(uW-Y)^2
       +(A+\alpha+1/2)(1+u^2)Y^2+(1+u^2)X^2/2.
 \tag{PP4}
\]
It is positive definite for every bounded real $u$. Scale the two
$W,Y$ kernel equations at the source by $w_c$, and the $X$
equation by $V$. In the unknowns $(p,y,x)$ the resulting bounded
matrix has a block triangular residue: its $p,y$ block is the
positive $W,Y$ block of PP4, and its $x$ pivot is $(1+c^2)/2$.
The terms $(\beta-D)/V=V(\eta-\delta)$ are bounded. All right
hand sides are bounded, proving PP3 by an ordinary source-field
matrix inverse with unit determinant.

Let $q_{01},q_{10},q_{21},h$ be the original Laurent coefficients.
Their first three residues are positive units because $m,C,A,\alpha$
are units. The actual zero and angular stationarity equations,
using PP2, give
\[
 \bar q_{01}=\bar q_{21}\bar c^{\,4},\qquad
 \bar h=-2\bar q_{21}\bar c^{\,2}.
 \tag{PP5}
\]
Every other nonconstant Laurent contribution apart from $q_{10}w^2$
vanishes in this residue calculation. These are equations wholly in
the coefficient field followed by its real residue map.

\subsubsection*{A complete local normal and divided-difference argument}

It remains locally to treat testing $w$ infinitesimal or zero and
$u$ with the same positive real residue as $c$. The full reciprocal
quarter-parent is
\[
\begin{split}
 F={}&m(W-wX+uY-uwZ)^2+CD_w(uW-Y)^2+AD_uD_wY^2\\
 &+(D+1/2)(uwW+uX-wY-Z)^2
       +(\alpha+1/2)D_u(Y-wZ)^2\\
 &+(\beta+1/2)D_w(X+uZ)^2-D_uD_wZ^2/2,
 \qquad D_u=1+u^2,\quad D_w=1+w^2.
\end{split}\tag{PP6}
\]
Use three unrestricted normals and two scalar conductors
\[
 N_W=W-pwZ,\quad N_Y=Y-ywZ,\quad N_X=X-xw^2Z,
 \qquad Q=w(u-c)Z,\quad L=(w^2-V)Z.
 \tag{PP7}
\]
The centers agree with the actual source contact; no equation is
imposed on an independent state. Write the full normal expansion
\[
 F=N^TH(u,w)N+2Z\sum_iN_i\ell_i(u,w)+E(u,w^2)Z^2,
 \quad N=(N_W,N_Y,N_X)^T.
 \tag{PP8}
\]
It is an exact expansion of every term of PP6. The matrix $H$ has
bounded entries and residue PP4 at $u=\bar c$. Parity and direct
inspection of PP6 give bounded polynomials $a_W,a_Y,a_X$ such that
\[
 \ell_W=w a_W(u,b),\qquad \ell_Y=w a_Y(u,b),\qquad
 \ell_X=u(\beta-D)+b a_X(u,b),\quad b=w^2.
 \tag{PP9}
\]
Actual state stationarity gives $a_W(c,V)=a_Y(c,V)=0$ and
$\ell_X(c,w_c)=0$. Finite monic division thus gives bounded
polynomials $A_i,B_i$ with
\[
 a_i(u,b)=(u-c)A_i(u,b)+(b-V)B_i(b),\qquad i=W,Y.
\]
Consequently the entire first two normal couplings are
$Z\ell_i=A_iQ+wB_iL$.

For the third normal put
$r(u)=u(\beta-D)/V+a_X(u,V)$. This is a bounded source
polynomial by PP2, and $r(c)=0$. Hence $r=(u-c)R(u)$ for
a bounded polynomial $R$. Division in $b$ gives another bounded
polynomial $B_X$ with
\[
 Z\ell_X=\{B_X(u,b)-(u-c)R(u)\}L+wR(u)Q.
 \tag{PP10}
\]
Here the identity $V(u-c)Z=wQ-(u-c)L$ is exact. Moreover
$B_X(c,0)$ has zero residue: the source equation
$u(\beta-D)+Va_X(u,V)=0$ at $u=c$, divided by $V$,
gives $a_X(c,V)$ infinitesimal, and $B_X(c,0)$ has the same
residue. Thus the third normal has no residue coupling to $L$
or $Q$. The first two have no residue coupling to $L$.

For completeness, the scalar coefficient obeys
\[
 E(u,0)=D+\beta u^2,\qquad
 E(c,V)=E_u(c,V)=E_b(c,V)=0.
 \tag{PP11}
\]
The first assertion follows by setting $w=0$ in the entire PP6
after the graph substitution. The other three follow from the
actual full source contact and its two angle derivatives; neither
normal term is discarded in making this inference.
Define $d(u)=\delta+\eta u^2$ and the bounded polynomial
\[
 J(u,b)=\{E(u,b)-d(u)(b-V)^2\}/b.
 \tag{PP12}
\]
This notation means exact monic division: PP11--PP2 make the
constant term zero. No testing $b$ is inverted. At $(c,V)$
the value and both first derivatives of $J$ vanish. A second
finite monic Taylor division gives bounded polynomials $J_{20}$,
$J_{11},J_{02}$ with
\[
 J(u,b)=(u-c)^2J_{20}(u,b)
       +(u-c)(b-V)J_{11}(b)+(b-V)^2J_{02}(b).
 \tag{PP13}
\]
Substitution yields the complete exact scalar identity
\[
 E(u,w^2)Z^2=J_{20}(u,w^2)Q^2+wJ_{11}(w^2)QL
                     +\{d(u)+w^2J_{02}(w^2)\}L^2.
 \tag{PP14}
\]
All divisions preceded tensor evaluation, with monic divisors and
bounded source coefficients. These formulas hold literally at $w=0$.

Equations PP8--PP14 represent the entire parent by a bounded
Gram on the five rows $N_W,N_Y,N_X,Q,L$. Its $L$ pivot in
residue is $\bar\delta+\bar\eta\bar c^2>0$, and its $L$
crosses have zero residue. It remains to check the four-row hard
and angular residue; the following exact determinant computation
provides that check without taking a mixed inverse.

\subsubsection*{Strictness of the angular residue}

At $D=\beta=0$, replace $Z$ by $T/w$ in the whole reciprocal
parent and take its polynomial value at $w=0$. This gives a real
four-state quadratic form $B_u(W,Y,X,T)$, with hard block PP4.
The original determinant identity gives
\[
 \det B_u=\frac{(1+u^2)^2}{16}
                  (q_{01}+hu^2+q_{21}u^4).
 \tag{PP15}
\]
Using PP5, its real Schur coefficient at $T$ is
\[
 g(u)=\frac{(1+u^2)^2q_{21}}{16\det H_u}
                       (u^2-c^2)^2.
 \tag{PP16}
\]
Here all entries and inverses are real residues; this calculation
proves positivity of a real matrix and does not factor a matrix in
the tensor ring. The contact centers PP3 give the unique real
kernel of $B_c$, so the four-row residue in PP8--PP14 is precisely
the quadratic Taylor coefficient in hard normals and the change
$u-c$. Its angular Schur pivot is
\[
 g''(c)/2=\frac{(1+c^2)^2q_{21}c^2}{4\det H_c}>0.
 \tag{PP17}
\]
The hard block is PP4 and is positive definite. Thus this full
four-row residue is positive definite, as is the independent $L$
pivot. The separated strict bounded-Gram lemma proves that PP6
is tensor SOS in the complete local chart. The rows may be
unbounded or linearly dependent at a literal endpoint; only the
coefficient Gram needs a strict real residue.

\subsubsection*{Explicit testing complement and projective endpoints}

If $u,v$ are bounded original testing coordinates, use the ordinary
state matrix except near the double zero corner. At finite nonzero
real residues, the residue Laurent polynomial is
\[
 q_{21}(u^2-c^2)^2/u^2+q_{10}/v^2>0,
 \tag{PP18}
\]
so its state matrix is positive definite by the hard $Z=0$ block.
At $u=0$, $v\ne0$ the real quarter-parent is
\[
 m(vW-X)^2+(C+A+1/2)(1+v^2)Y^2
                       +\alpha(vY-Z)^2+(1+v^2)X^2/2\succ0.
\]
At $v=0$, $u\ne0$ it is
$m(X+uZ)^2+(C+1/2)(uW-Y)^2+A(1+u^2)Y^2
+\alpha(1+u^2)Z^2+(X+uZ)^2/2$, which is likewise strict.
These are residue identities; the bounded full matrices therefore
give tensor SOS on their respective charts.

When both $u,v$ are infinitesimal or zero, use the five rows
$(uW,vW,X,Y,Z)$ in the entire original finite parent. The exact
polynomial substitution is bounded and its residue is
\[
 m(vW-X)^2+(C+1/2)(uW-Y)^2+AY^2+\alpha Z^2+X^2/2\succ0
 \tag{PP19}
\]
as a coefficient Gram on those five rows. This treats that corner
with no state normalization.

For reciprocal Alice and bounded Bob, the complete core identity
RC20 applies unchanged with all $m,D,\alpha,\beta$ event squares
retained whole. Its residue RC21 is strict using only $C,A>0$.
For two reciprocal testing angles, retain those positive squares
and use RC24 for the four rows $W,X,Y,wZ$. The residue
$CW^2+AY^2+X^2/2+(Y-wZ)^2/2$ is strictly positive.
These identities allow the reciprocal coordinate to be literally zero.

For Bob reciprocal and bounded Alice, retain the whole $D,\beta$
squares, put $T=wZ$, and use the full four-row expression
\[
\begin{split}
 m(W-wX+uY-uT)^2+CD_w(uW-Y)^2
          +AD_uD_wY^2+\alpha D_u(Y-T)^2+K/2,
\end{split}
\]
where $K$ is the full core RC22 with $\rho_*=uW-Y$.
Its residue is exactly $B_u$. By PP15--PP16 it is strictly
positive whenever $u^2$ has residue different from $\bar c^2$,
including $u$ infinitesimal. The remaining matched residues are
PP7--PP17 after the exact Alice sign gauge if necessary.

Both signs of either angle are allowed by the exact simultaneous
state sign gauges of the original parent. Every normalization is
a positive square in a single source or testing field. Complex
states split into real and imaginary parts. The finite and
reciprocal homogeneous identities include all literal axes and
projective endpoints. This completes the all-tests theorem on the
stated positive-$d$ C1 projective source domain.



\section{\texorpdfstring{A direct certificate for the positive-$d$ C2 reciprocal corner}{A direct certificate for the positive-d C2 reciprocal corner}}
\label{component:p10686597065a}

\label{p10686597065a:sec:new-rc-direct}

\paragraph{The source domain.}
Normalize the original C2 source by $d=1$ and $\lambda=\beta=0$.
Assume that all source coefficients are bounded, that $C,A,\alpha$
are positive bounded units, that $m,D$ are positive infinitesimals,
and that $n\geq0$ is bounded. Retain the original physicality,
finite noncommuting stationary zero contact and positive native gaps. This section proves
the complete parent is tensor SOS at every independent projective
test and every state. The contact is used only inside the source
field. No mixed determinant is cancelled.

Equation~\eqref{eq:reciprocal-source-balances}, with $\beta=0$,
gives, for the reciprocal Bob contact coordinate $w_c=1/v_c$,
\[
 U=u_c^2\in\operatorname{Inf}_{>0},\qquad
 V=w_c^2\in\operatorname{Inf}_{>0},\qquad
 m\asymp U^2,\qquad D\asymp V^2.
 \tag{RC1}
\]
The positive-unit ratios follow from the final formula in
Section~\ref{lem:reciprocal-source-balance}.

\subsubsection*{The actual source has a bounded shape chart}

Choose the native positive contact chart with $Z_c=1$, and set
\[
 S=(u_c/w_c)W_c,\qquad y=Y_c/w_c,\qquad
 R=-X_c/(u_cV).
 \tag{RC2}
\]
Then $S,y,S-y,1-y,R$ are positive bounded units. Here is an entry
proof that does not assume those facts as extra source hypotheses.
In the reciprocal parent RC6 below, take the ordinary source-field
Schur complement of the $X,Y$ block. That block has residue
$\operatorname{diag}(n+1/2,H)$, where $H=C+A+\alpha+1/2$.
Its $W,W$ Schur entry divided by $U$ has residue
$a=C-C^2/H>0$, while its $W,Z$ entry divided by $u_cw_c$
has residue $-b=-(1/2+C\alpha/H)<0$. The balance
$m/U\in\operatorname{Inf}$ and $D/V\in\operatorname{Inf}$
justifies these divisions in the source field. The zero-state
equation therefore gives $S\sim b/a$, a positive bounded unit.

The two hard state equations, scaled in the source field by
$u_cV$ and $w_c$, are a bounded linear system in $R,y$. Its
residue is triangular with positive diagonal $n+1/2,H$.
Consequently both quantities are bounded. The second equation gives
$Hy=CS+\alpha$ in residue, so $y$ is a positive unit. The
$W$ equation divided by $u_cw_c$ gives
$C(S-y)=1/2$ in residue. The $Z$ equation divided by $V$
gives $\alpha(1-y)=S/2$ in residue: its apparently constant
$U$ terms cancel exactly before the division. Finally the first
hard equation gives $(n+1/2)R=(1+S-y)/2$ in residue. These
relations prove all five claimed positive-unit bounds. The exact
native gaps already give $0<y<1$ and $S-y>0$ before taking
residues.

The unchanged native source formulas are converted by the following
dictionary. Their letters are used only in this display:
\[
\begin{aligned}
 z=U,\quad a^2=1/(UV),\quad b=S/U,\quad t=(S-y)V,
 &\quad r=VR,\quad h=y/U,\\
 g=S/U+y-1+VR,&\qquad j=1+UVR-V(S-y).
\end{aligned}
 \tag{RC3}
\]
Substitution into the six exact native coefficients has denominators
with positive-unit leading parts after the displayed powers of
$U,V$ are removed. Thus $\mu=m/U^2$ and $\delta=D/V^2$
are bounded units, and the source residues are
\[
\begin{aligned}
 C_0&=\frac1{2(S-y)},& A_0&=\frac{1+S-y}{2y},&
 \alpha_0&=\frac{S}{2(1-y)},\\
 n_0&=\frac{1+S-y-R}{2R},&
 \mu_0&=\frac{1+S-y}{2S^2},&
 \delta_0&=\frac{S(1+S-y)}2.
\end{aligned}\tag{RC4}
\]
In residue formulas such as RC4 the shape letters mean their real
residues. In particular
\[
 C_0\{\alpha_0(A_0-1/2)-1/4\}
       =(A_0+\alpha_0+1/2)/4.
 \tag{RC5}
\]
These are source statements only. They do not identify any source
infinitesimal with an independent testing infinitesimal.

\subsubsection*{The complete local parent and both unrestricted normals}

For $u,w$ infinitesimal or zero in the independent testing field,
put $D_u=1+u^2$, $D_w=1+w^2$. The full quarter-parent after
the pure testing reciprocal Bob normalization is
\[
\begin{split}
 F={}&m(W-wX+uY-uwZ)^2+nD_uD_wX^2
       +CD_w(uW-Y)^2+AD_uD_wY^2\\
 &+(D+1/2)(uwW+uX-wY-Z)^2
       +(\alpha+1/2)D_u(Y-wZ)^2\\
 &+\tfrac12D_w(X+uZ)^2-\tfrac12D_uD_wZ^2.
\end{split}\tag{RC6}
\]
All four state coordinates are arbitrary. Introduce tensor rows
\[
 \begin{gathered}
 N_x=X+uRw^2Z,\qquad N_y=Y-ywZ,\\
 \rho=uW-SwZ,\qquad k=UW-SuwZ,\qquad \ell=(w^2-V)Z.
 \end{gathered}
 \tag{RC7}
\]
The testing-dependent $w^2$ in $N_x$ is essential; it agrees with
the source contact and removes the unbounded coefficient produced
by a contact-fixed center. The exact conductor identities are
\[
 U\rho-uk=Sw(u^2-U)Z,\qquad
 u\rho-k=(u^2-U)W.
 \tag{RC8}
\]

Expanding the entire RC6 gives
\[
 F=F_0+H_xN_x^2+H_yN_y^2+2H_{xy}N_xN_y
                         +2N_xL_x+2N_yL_y,
 \tag{RC9}
\]
where $F_0$ is only the constant coefficient of this unrestricted
normal expansion. The hard coefficients are exactly
\[
\begin{split}
 H_x&=mw^2+nD_uD_w+(D+1/2)u^2+D_w/2,\\
 H_y&=mu^2+CD_w+AD_uD_w+(D+1/2)w^2+(\alpha+1/2)D_u,\\
 H_{xy}&=-(m+D+1/2)uw.
\end{split}\tag{RC10}
\]
Set $J_x=n(1+w^2)+D+1/2$ and $K_x=n(1+U)+1/2+m$.
The complete first normal is
\[
 L_x=(D+1/2)uw\rho-(m/U)wk
 +u\{D/V-Rw^2K_x\}\ell
 -(R/S)uwJ_x(U\rho-uk).
 \tag{RC11}
\]
Every coefficient on $\rho,k,\ell$ is bounded infinitesimal.
For the other normal define
\[
 \begin{aligned}
 a_{10}&=Ay+(m+\alpha+1/2)(y-1),\\
 a_{01}&=-(C+D+1/2)S+(C+A+D+1/2)y,\\
 a_{11}&=R(m+D+1/2)+Ay.
 \end{aligned}
\]
Then the exact expression is
\[
\begin{split}
 L_y={}&\{m-C-(C+D+1/2)w^2+(U/S)(a_{10}+a_{11}w^2)\}\rho\\
 &-(u/S)(a_{10}+a_{11}w^2)k+w(a_{01}+a_{11}U)\ell.
\end{split}\tag{RC12}
\]
Its only nonzero residue coupling is $-C_0\rho$. RC11--RC12
follow by full expansion and the two actual source state equations;
they account for every term linear in the two normals in RC6.

\subsubsection*{A scalar certificate by finite monic divisions}

Put
\[
 c=C+(C+D+1/2)w^2,\quad b_2=(D+1/2)R,
 \quad b_1=(C+D+1/2)(y-S)-mR,\quad \gamma=b_1+b_2U.
\]
Define
\[
 F_{\rm base}=(c-2Ub_2w^2/S)\rho^2+\mu k^2
       +(2ub_2w^2/S)\rho k-2w\gamma\rho\ell.
 \tag{RC13}
\]
The source $W$ state equation removes the remaining $WZ$
coefficient. Direct expansion then gives
$F_0-F_{\rm base}=E(u^2,w^2)Z^2$, with $E(a,b)$ a
bounded-coefficient polynomial of degree at most two in $a$.
Its exact properties are
\[
\begin{gathered}
 E(a,0)=D,\qquad E(U,V)=E_a(U,V)=E_b(U,V)=0,\\
 E(a,b)=E_0(b)+(a-U)E_1(b)+(a-U)^2E_2(b),\\
 E_2(b)=R^2b^2\{n(1+b)+D+1/2\}.
\end{gathered}\tag{RC14}
\]
The three vanishing jets follow from the actual full contact:
every row in RC7 vanishes there, the complete scalar and both
angle derivatives vanish, and RC13 has zero first derivatives
there. They are also checked directly from the native source
coefficients. Since $E_1(0)=E_1(V)=0$ and $E_0$ has a double
zero at $V$, finite monic polynomial division gives bounded
polynomials $J_0,J_1$ such that
\[
 E_0(b)=(b-V)^2J_0(b),\qquad
 E_1(b)=b(b-V)J_1(b).
 \tag{RC15}
\]
These divisions take place in the source polynomial algebra before
tensor evaluation. They introduce no inverse of $V$ or of a mixed
quantity. In particular $V^2J_0(0)=D$, and the other coefficients
of $J_0$ remain bounded; hence $J_0(w^2)$ has residue $\delta_0$.

RC8 now gives the exact scalar identity
\[
\begin{split}
 F_0={}&F_{\rm base}
 +\frac{R^2w^2\{n(1+w^2)+D+1/2\}}{S^2}(U\rho-uk)^2\\
 &+\frac wS J_1(w^2)(U\rho-uk)\ell+J_0(w^2)\ell^2.
\end{split}\tag{RC16}
\]
This polynomial formula also holds at $u=0$ or $w=0$; neither
is inverted. Combining RC9--RC16 represents the entire parent
by a bounded Gram on $(N_x,N_y,\rho,k,\ell)$ with real residue
\[
 (n_0+1/2)N_x^2+H_0N_y^2-2C_0N_y\rho+C_0\rho^2
                              +\mu_0k^2+\delta_0\ell^2,
 \qquad H_0=C_0+A_0+\alpha_0+1/2.
 \tag{RC17}
\]
This is positive definite: the $N_y,\rho$ determinant is
$C_0(A_0+\alpha_0+1/2)>0$, and the other three pivots are
strictly positive. The separated strict bounded-Gram lemma
therefore gives a finite tensor SOS of the complete RC6. It
requires bounded matrix entries, not bounded rows or states.

\subsubsection*{All other testing charts}

The remaining measurement charts follow from three matrix
identities. In each one we retain the indicated event squares and
prove positivity of the remaining bounded Gram matrix.

\paragraph{Both original testing coordinates bounded.}
Put $p=uW$. Retain the whole $m,n,D$ squares and use the
remaining Gram on $(p,Y,X,Z)$:
\[
\begin{split}
 C D_v(p-Y)^2+A D_uD_vY^2
 &+\tfrac12(p+uvX-Y-vZ)^2+(\alpha+1/2)D_u(vY-Z)^2\\
 &+\tfrac12D_v(X+uZ)^2-\tfrac12D_uD_vZ^2.
\end{split}\tag{RC18}
\]
The $Z=0$ hard block is positive definite by its displayed
$C,A,X$ squares. After RC4 the determinant of the real residue is
\[
 \frac{(S+u^2v^2)(1+u^2)^2(1+v^2)^2(1+S-y)^2}
      {16y(S-y)(1-y)}>0.
 \tag{RC19}
\]
Here also $u,v$ mean real testing residues. Thus RC18 is a strict
bounded Gram, even when $u$ or $v$ has zero residue.

\paragraph{Alice reciprocal, Bob bounded.}
Let $s=1/u$ be infinitesimal or zero, set $D_s=1+s^2$, and use
rows $(W,Y,V_x,R_x)=(W,Y,vX,sX+Z)$. Retain the entire
$m,n,D,\alpha$ squares. The remaining expression is
$CD_v(W-sY)^2+AD_sD_vY^2+K/2$, where exactly
\[
\begin{split}
 K={}&\{W-sY+D_sV_x-svR_x\}^2+D_vR_x^2\\
 &+D_s\{v^2Y^2-2Y(vR_x-sV_x)-(vR_x-sV_x)^2\}.
\end{split}\tag{RC20}
\]
Its residue is
\[
 C_0D_vW^2+A_0D_vY^2
          +\tfrac12(W+V_x)^2+\tfrac12(vY-R_x)^2\succ0.
 \tag{RC21}
\]
This single formula includes all bounded $v$, including literal
zero and arbitrary competing small $s,v$. There is no $s/v$
coefficient.

\paragraph{Bob reciprocal, Alice a bounded unit.}
Let $w$ be infinitesimal or zero, and $u$ a bounded unit. Use
$\rho_*=uW-Y$, $T=wZ$, $S_D=\sqrt D\,Z$, and the five rows
$(\rho_*,Y,X,T,S_D)$. Retain $m,n$ whole. The $D$ square is
exactly $\{\sqrt D(uX+w\rho_*)-S_D\}^2$. The other terms are
\[
 CD_w\rho_*^2+AD_uD_wY^2+\alpha D_u(Y-T)^2+K/2,
\]
where the complete core is
\[
 K=(D_u+w^2)X^2+2uwX T+w^2\rho_*^2+2uw\rho_*X
                           -2\rho_*T-T^2+D_u(Y-T)^2.
 \tag{RC22}
\]
The $X,S_D$ residue pivots are $D_u/2,1$. On
$(\rho_*,Y,T)$ the residue matrix is
\[
 \begin{pmatrix}
 C_0&0&-1/2\\
 0&G D_u&-(\alpha_0+1/2)D_u\\
 -1/2&-(\alpha_0+1/2)D_u&\alpha_0D_u+u^2/2
 \end{pmatrix},\qquad G=A_0+\alpha_0+1/2.
 \tag{RC23}
\]
Its first two leading minors are $C_0,C_0GD_u$. By RC5 its
determinant is $GD_uu^2(1+2C_0)/4>0$. This proves strictness.
The alternative $u$ infinitesimal or zero is precisely RC7--RC17.

\paragraph{Both testing coordinates reciprocal.}
Let $s,w$ be infinitesimal or zero. Set $T=wZ$ and
$S_D=\sqrt D\,sZ$. Use the five rows $W,X,Y,T,S_D$, and retain the
entire $m,n,\alpha$ squares. The $D$ square is
$\{\sqrt D(wW+X-swY)-S_D\}^2$. The other terms are
$CD_w(W-sY)^2+AD_sD_wY^2+K/2$, where exactly
\[
\begin{split}
 K={}&(wW+X-swY)^2+D_s(Y-T)^2+s^2D_wX^2\\
 &-2sWT+2s^2YT+2swXT-s^2T^2.
\end{split}\tag{RC24}
\]
The full five-row residue is
\[
 C_0W^2+A_0Y^2+X^2/2+(Y-T)^2/2+S_D^2\succ0.
 \tag{RC25}
\]
This gives the final strict bounded Gram.

Every ordered-field testing coordinate is bounded or has
infinitesimal reciprocal. RC18--RC25 and the local proof therefore
cover both projective lines. The identities are polynomial in the
chosen homogeneous coordinates, so literal infinity is included by
$s=0$ or $w=0$. Both signs are allowed throughout; no state ratio
or state normalization is used. Positive pure testing square
normalizations restore the original homogeneous parent, and the
original source scale is restored by a pure source square root.
Real and imaginary parts of a complex state are handled separately.



\chapter{The C1 coefficient domain}
\label{ch:c1-short}

Let $K_c,K_t$ be separately labelled real closed extensions of
$\mathbb R$ and let $T=K_c\otimes_{\mathbb R}K_t$.
Source coefficients belong to $K_c$; independent testing coordinates
belong to $K_t$. A coefficient is bounded if it is a finite sum of
products of bounded elements of the two fields. Its residue is obtained
by taking the real residues in both factors. All sums of squares below
are finite sums in $T$.

We prove the complete C1 assertion under the native source hypotheses
of the contact-entry theorem: $d,m,C,A,\alpha>0$,
$\lambda_{\rm ev}=n=0$, $D,\beta\ge0$, $D+\beta>0$, the original
strict native gaps, and positive infinitesimal total scale. Normalize
$d=1$ by a pure-source square. The full target is the original
four-state quadratic parent in Chapter~\ref{ch:cover}, with its
homogeneous value used at projective testing endpoints.
No testing angle is identified with the source contact.

\section{Source coordinates and polynomial operations}
The contact equations determine the source coefficients. Their positive denominators and the lower bound for $E=1-x\eta h$ determine which rescalings remain bounded. We establish these facts before choosing the roots of the certificate.

\subsection{Finite polynomial operations used in the proof}
\label{sec:c1-finite-operations}

All coefficient arrays in this chapter have the following precise
meaning. If $\Pi$ is a displayed sum of retained positive squares,
$s$ is a vector of the remaining state coordinates and
$s=s_0+A\xi$ is the displayed change of state coordinates, put
\[
 \mathcal Q(\xi)=(F-\Pi)(s_0+A\xi),\qquad
 q_0=\mathcal Q(0),\quad q_i=[\xi_i]\mathcal Q,
 \quad H_{ij}={1\over2}\partial_i\partial_j\mathcal Q.
\]
Then exactly $\mathcal Q=q_0+\sum_i q_i\xi_i+\xi^TH\xi$.
Thus the scalar term, both linear state couplings and the entire state
quadratic form are computed from the same polynomial. Multiplying by the
displayed positive pure-source and pure-testing common denominators
makes these coefficient extractions polynomial. All later divisions
are divisions by the specified monic polynomials.

For a polynomial $f$ define its finite difference quotient by
\[
 \Delta_{z,a}f={f(z)-f(a)\over z-a}.
\]
It is the polynomial determined by
$(z^j-a^j)/(z-a)=\sum_{i=0}^{j-1}z^{j-1-i}a^i$.
More generally, a quotient/remainder by $z^r-a$ is determined uniquely
by repeatedly replacing $z^r$ by $a$ and collecting the difference.
Thus every quotient defined below is an explicit finite polynomial
operation, even where a source opening tends to zero. Weighted degree
is checked on the original polynomial before division; division by
a weighted-homogeneous monic polynomial preserves the filtration
with its stated degree shift. The support and weight bounds specified below ensure that the resulting
Gram coefficients remain bounded after rescaling.

\begin{lemma}[Auxiliary square allocation]
\label{lem:c1-square-allocation}
Let $r$ be roots with a bounded symmetric Gram whose real residue
$G_0$ is positive definite. Let $h=(h_1,\ldots,h_s)$ be additional
roots and suppose there are exact certificates
\[
 Q_j(r,h)-h_j^2\in\Sigma T^2,
\]
where every Gram coefficient of $Q_j$ is bounded with zero residue.
Suppose the remaining terms of the polynomial are a bounded quadratic
form in $(r,h)$ whose $r,r$ residue block is $G_0$. Then the complete
polynomial is a tensor sum of squares.
\end{lemma}
\begin{proof}
For a sufficiently large fixed real $M>0$, add
$M\sum_j(h_j^2-Q_j)$ to the remaining Gram. Its real residue is
\[
 \begin{pmatrix}G_0&B_0\\B_0^T&C_0+MI\end{pmatrix}\succ0.
\]
Existence of such an $M$ follows by a Schur complement of the finite
real matrix $G_0$. The strict positive-residue Gram lemma applies.
Restore $M\sum_j(Q_j-h_j^2)$ using the displayed exact certificates.
Only a finite real matrix was inverted in this proof.
\end{proof}

We also use the common product-relation allocation, the restricted
monic contact lift CC1--CC6, and the finite real two-block allocation
UA1--UA3. Their hypotheses are checked in the source cases below.
The arbitrary-scale noncommuting theorem is used only for the original
physical source, after its strong-event source conditions have been
proved. The generic LII and RR polynomial identities are used only
with the explicit C1 substitutions given below.

Whenever a proof below retains a nonnegative coefficient multiplying
a whole square, an unbounded positive excess stays outside its bounded
Gram. A strict leading block is obtained by clipping whole positive
squares at fixed real lower bounds, never by clipping individual
matrix entries. The real constrained-limit argument gives finite real
caps where a source opening makes a retained coefficient unbounded.
Every coefficient correction and both complete first state normals
are included before this argument is applied.

\subsection{The native source partition}
\label{sec:c1-source-partition}
\label{component:pfd0dd7bb10cf}

Use the native coordinates $x=k^2$, $0<k\le1$,
$\eta,\zeta>0$, $r=1-a\eta$, $\theta=t\zeta$, $h=1-\theta$,
$E=1-x\eta h>0$, $Q=1-x\zeta r>0$,
$M=1+xr-t>0$ and $J=1+xh-a>0$.
All satisfy $0<r,\theta<1$ and $0<a,t<2$.
Write I for positive infinitesimal, U for positive bounded with nonzero
residue and O for positive unbounded. Put $e=k\sqrt{\eta/\zeta}$.
The exhaustive source cases are
\[
\begin{array}{c|c|c|l}
k&\zeta&\text{other conditions}&\text{proof below}\\\hline
\mathrm U&\text{any}&\eta,\zeta:\mathrm I&\text{Section~\ref{sec:c1-uniform-projective-join}}\\
\mathrm U&\text{any}&\eta,\zeta\text{ not both I}&\text{Sections~\ref{component:pa0b854a04f49}, \ref{component:p42ee07a067c1}, \ref{component:pf7ffa921793c}}\\
\mathrm I&\mathrm U\text{ or }\mathrm O&\eta>0&\text{Section~\ref{sec:c1-L-partition}}\\
\mathrm I&\mathrm I&\eta:\mathrm I&\text{Section~\ref{component:pf5d2f33f694c}}\\
\mathrm I&\mathrm I&\eta:\mathrm O&\text{Section~\ref{component:p9afbe73899af}}\\
\mathrm I&\mathrm I&\eta:\mathrm U,e:\mathrm I&\text{Sections~\ref{component:pa7e744614f49}, \ref{component:p05a86be0046c}}\\
\mathrm I&\mathrm I&\eta,e,r:\mathrm U&\text{Section~\ref{component:p69995d829bb2}}\\
\mathrm I&\mathrm I&\eta,e:\mathrm U,r:\mathrm I&\text{Section~\ref{component:pccaa793fd0bf}}\\
\mathrm I&\mathrm I&\eta,r:\mathrm U,e:\mathrm O&\text{Section~\ref{component:p890eea5f538b}}\\
\mathrm I&\mathrm I&\eta:\mathrm U,e:\mathrm O,r:\mathrm I&\text{Section~\ref{component:pccaa793fd0bf}}
\end{array}
\]
The source identities below justify every listed implication, including
the ordinary UI/UH alternatives and the bound needed to split C1.1.
The last two centered-$A$ rows share the same bounded-reciprocal
construction. The source field alone determines this partition.

\subsubsection{The $L=x\zeta$ cases}
\label{sec:c1-L-partition}
For $k$ infinitesimal and $\zeta$ noninfinitesimal, put $L=x\zeta$.
When $L$ is unbounded, $Q>0$ forces $r$ infinitesimal and
$1-r=a\eta$, $a<2$ exclude infinitesimal $\eta$. Unit $\eta$
is treated in Section~\ref{component:pd4b92bf7de71}; unbounded
$\eta$ uses Theorem~\ref{thm:c1-unification-bounded-z-all-H}, with $z=1/L$ infinitesimal,
for $u,w$ infinitesimal. The unit-Alice results remain in
Sections~\ref{component:pb1104217f67b},~\ref{component:p909aa5765721}
and~\ref{component:pcc84202479d2}, and the complementary projective
tests remain in Section~\ref{component:p89fd3edb3edd}.
For unit $L$, infinitesimal and unit $\eta$ are covered in
Sections~\ref{component:p080aa62de16d} and~\ref{component:p686209b87229}.
Unbounded $\eta$ has $z$ a positive unit and uses
Theorem~\ref{thm:c1-unification-bounded-z-all-H}; its complementary
testing cases remain in Sections~\ref{component:pb0d9d28d0849}
and~\ref{component:pda4a21937f64}.

When $L$ is infinitesimal, unbounded $\eta$ gives
\[
 S=\sigma+L,\quad\sigma=1/\eta,\quad
 \kappa=x/S,\quad\nu=\sigma/S,\quad\ell=L/S,
 \quad \nu+\ell=1,\quad \zeta=\ell/\kappa.
\]
The identity $\kappa Q_s+\ell M+\nu E=1$, proved below, makes
$\kappa$ bounded and gives the five source cases
\[
\begin{array}{c|c|c|l}
\kappa&\nu&\ell&\text{proof}\\\hline
\mathrm U&\mathrm U&\mathrm U&\text{Sections~\ref{component:p3f02c51074be}, \ref{component:pd26b89cdb3f0}}\\
\mathrm U&\mathrm I&\mathrm U&\text{Section~\ref{component:p578719565e9d}}\\
\mathrm I&\mathrm U&\mathrm U&\text{Section~\ref{component:p2edb55cd1f69}}\\
\mathrm I&\mathrm U&\mathrm I&\text{Section~\ref{component:pb96c9bfcc12d}}\\
\mathrm I&\mathrm I&\mathrm U&\text{Sections~\ref{component:p80b0c5b9879f}, \ref{component:peaa56c1a3c7a}}
\end{array}
\]
In the last row $q_\nu=\nu/\ell$ and $\tau=h/q_\nu$ are
defined only there. BN proves every bounded $\tau$ and UN every
unbounded $\tau$. No separate restriction on $\beta/q_\nu^2$ is made.

If $L$ is infinitesimal and $\eta$ bounded, unbounded $\zeta$ with
unit $\eta$ is Section~\ref{component:p60eca5442075}. Small $\eta$
with unit or unbounded $\zeta$ uses
Sections~\ref{component:p024c821bfb59},~\ref{component:pc720c1e06aca}
and~\ref{component:pe3bbfd3a8151}.
The remaining $\eta,\zeta$ are units. Put
\[
 j=\eta(1+xh)-1+r=\eta J>0,\qquad j_0=\eta+r-1,
 \qquad j-j_0=\eta xh\in\operatorname{Inf}_c.
\]
Unit $j$ is Section~\ref{component:p3df6e6bdd8f4}; small $j$ with unit
$r$ is Section~\ref{component:p48da253a150f}; small $j,r$ is
Section~\ref{component:pbf6de1d7328a}.
Thus every positive source opening and each permitted literal
$D=0$ or $\beta=0$ has a case in the proof.


% Source formulas: p42ee07a067c1
\subsection{The native coefficient formulas and their uniform gap}

\label{component:p42ee07a067c1}

\label{p42ee07a067c1:sec:c1-global-native-cap}

For the physical C1 source, $\lambda=n=0$, $d>0$, $m,C,A,\alpha>0$,
$D,\beta\ge0$ with at most one zero, and the source has its actual
finite noncommuting stationary contact. Use the common global sign
entry and normalize $C=2$. The exact full-parent Alice reflection
permits $0<k\le1$. Put
\[
 x=k^2,\quad r=1-a\eta,\quad\theta=t\zeta,\quad
 h=1-\theta,\quad y=v_c^2=x\eta\zeta,
\]
\[
 K=1+x,\ R=1+xr,\ T=1+xh,\ G=h+\zeta,
 \quad M=R-t,\ J=T-a,\ E=1-x\eta h,\ Q=1-x\zeta r.
 \tag{NC1}
\]
The common exact signs give $0<r,\theta<1$, $E,Q,M,J>0$ and
$0<a,t<2$. There is no size assumption on $\eta,\zeta,e$ or a
positive lower bound on any opening.

\paragraph{Three short positive source identities.}
The original coquotient polynomials obey
\begin{align*}
 P_4&=TE+x(1-r)G,\\
 P_\alpha&=Rh(1+y)+\zeta ME,\\
 P_d&=p_0+\eta p_1,\qquad
 p_0=1+2xh+rx^2G+x\zeta,\quad
 p_1=x\{K\zeta+xh\theta\}.
 \tag{NC2}
\end{align*}
These are obtained by substituting $a\eta=1-r$ and $t\zeta=\theta$
in the unchanged seven native coquotients. In particular $P_d\ge1$.
The original nonnegative $D,\beta$ are precisely the inequalities
\[
 B_5:=K(1+y)(r+\eta)-(\eta T-1+r)P_d\ge0,
 \qquad B_\beta:=Kx(r+\eta)G-\theta P_d\ge0.
 \tag{NC3}
\]
The exact relative coefficients are
\begin{align*}
 \frac C d&=\frac{K(r+\eta)E}{2(1-r)P_d},&
 \frac A d&=\frac{\eta P_4}{2rP_d},\\
 \frac m d&=\frac{Kx(r+\eta)h(1+y)}{2M\zeta P_d},&
 \frac\alpha d&=\frac{x^2\zeta(r+\eta)P_\alpha}{2QP_d},\\
 \frac Dd&=\frac{B_5}{2J\eta P_d},&
 \frac\beta d&=\frac{B_\beta}{2\theta P_d}.
 \tag{NC4}
\end{align*}
The source inequalities make every denominator positive; their sizes
will enter the bounds below.

\begin{lemma}[A uniform native gap]
On every physical native C1 contact oriented by $0<k\le1$,
\[
                         E\ge\frac1{18}.             \tag{NC5}
\]
This includes the literal boundaries $D=0$ or $\beta=0$.
\end{lemma}
\begin{proof}
Define
\[
 \mathcal A=\theta T+K\zeta>0,\qquad
 \mathcal C=RG(1+rxh),
\]
\begin{align*}
 \mathcal B={}&h^2rx^2+hrx(2-h)+h^2x+hrx^2\zeta
             +hx\zeta+hx+h\\
             &+rx\zeta+x\zeta+2\zeta+1.
\end{align*}
Direct multiplication of (NC2)--(NC3), using $x\eta h=1-E$,
gives the polynomial identity
\[
                  hB_5=-\mathcal A E^2+\mathcal B E-\mathcal C.
 \tag{NC6}
\]
Every summand in $\mathcal B$ is nonnegative. Since
$0<x,r,h\le1$ and $h(2-h)\le1$, its six terms without $\zeta$
sum to at most $6$ and the coefficients of $\zeta$ sum to at most
$6$. Hence $\mathcal B\le6(1+\zeta)$. The native inequality
$\theta=t\zeta<2\zeta$ gives $1<h+2\zeta$, so
$1+\zeta<3G$. Also $\mathcal C\ge G>0$. From $B_5\ge0$ and
(NC6),
\[
 G\le\mathcal C\le\mathcal B E\le18GE.
\]
Cancel the positive pure source $G$. This proves (NC5) without
inverting an opening or assuming any of $\eta,\zeta$ bounded.
\end{proof}

\paragraph{An elementary division-free elimination.}
Keep $p_0,p_1$ from (NC2), and set
\begin{align*}
 a_2&=Kx\zeta-Tp_1,&
 a_1&=K(1+x\zeta r)-Tp_0+(1-r)p_1,&
 a_0&=Kr+(1-r)p_0,\\
 b_1&=KxG-\theta p_1,&
 b_0&=KxrG-\theta p_0.
\end{align*}
Then $B_5=a_2\eta^2+a_1\eta+a_0$ and
$B_\beta=b_1\eta+b_0$. Expanding these two expressions, or just
eliminating their quadratic and linear terms successively, gives
\begin{align*}
 &b_1^2B_5-a_2B_\beta^2-(b_1a_1-2a_2b_0)B_\beta\\
 &\qquad=b_1^2a_0-b_1a_1b_0+a_2b_0^2\\
 &\qquad=x^2h^2K^2RG^2
       \{1+2xh+rx^2h^2+(1-r)x\zeta\}.
 \tag{NC7}
\end{align*}
The last equality is a short polynomial identity in $x,r,h,\zeta$;
all five of its input coefficients are given above. No inverse of
$b_1$, no resultant premise, and no localization is used. For reference,
$b_1=xh\{K(1+\zeta)-x\theta^2\}>0$ as well.

\begin{theorem}[Actual unit/unit C1 contacts enter the five-event exit]
If both actual C1 half-angles $k,v_c$ are positive bounded units,
the original source has five consecutive events bounded below relative
to $d$. Consequently its entire original parent has the inherited
all-tests tensor-SOS certificate. All upper event scales and the
retained zero-coefficient boundaries are allowed.
\end{theorem}
\begin{proof}
After the exact Alice orientation, $x,y$ are bounded positive units.
Thus $\eta\zeta=y/x$ is a bounded positive unit. If $\eta$ were
unbounded, $\zeta$ would be infinitesimal. Then $\theta=t\zeta$
would be infinitesimal, contradicting $x\eta(1-\theta)<1$.
If $\zeta$ were unbounded, $\eta$ would be infinitesimal and
$r=1-a\eta$ would have residue one, contradicting $x\zeta r<1$.
Therefore $\eta,\zeta$ are both bounded positive units.

All polynomials in (NC2)--(NC3) now have bounded coefficients and
bounded arguments. If $h$ were infinitesimal, the exact specialization
\[
             B_\beta\big|_{\theta=1}
                    =-1-x\zeta(1-r)\le-1
\]
would contradict $B_\beta\ge0$. Thus $h$ is a positive bounded
unit. The uniform cap gives that $E$ is a unit. Equations (NC2) make
$P_4,P_\alpha,P_d$ positive bounded units: use $P_4\ge E$,
$P_\alpha\ge Rh(1+y)$, and $P_d\ge1$.

The first four equations of (NC4) now show that
$C,A,m,\alpha$ are all strong relative to $d$. Their denominators
are positive bounded quantities; they may be infinitesimal. Suppose
both $D/d$ and $\beta/d$ were infinitesimal, including a possible
literal zero. Their last two bounded denominator products in (NC4)
would make $B_5,B_\beta$ infinitesimal. Every coefficient on the
left of (NC7) is bounded, so that side would be infinitesimal. Its
right side is a positive unit, since $x,h,K,R,G$ are units and the
last bracket is at least one. This is a contradiction.

Thus $D$ or $\beta$ is strong. The original cycle
$(\lambda,D,C,A,\alpha,m,\beta,n)$ has respectively the five
consecutive strong events $(D,C,A,\alpha,m)$ or
$(C,A,\alpha,m,\beta)$. Apply the common five-consecutive-event
theorem, whose stated scope allows arbitrary upper weights and every
independent testing chart. No independent test is identified with
an actual source angle in this source-entry argument.
\end{proof}



% Source formulas: p51451360ff4a
\subsection{Consequences of the native source inequalities}

\label{component:p51451360ff4a}

\label{p51451360ff4a:sec:c1-residual-source-implications}

This section records source implications of (NC1)--(NC7), not further
all-tests certificates. All quantities are in the coefficient field.
Write $f\asymp g$ for two strictly positive quantities whose quotient
and reciprocal quotient are bounded. Set
\[
             L=x\zeta,\qquad S=1+L+y,\qquad \rho=r+\eta.
\]
The native signs and (NC5) imply the uniform estimates
\[
 S\le P_d\le6S,\qquad
 \frac{1+L}{54}\le P_4\le3(1+L),\qquad
 P_\alpha\asymp h(1+y)+\zeta M.                 \tag{NI1}
\]
Indeed $rx^2\zeta<x$ follows from $Q>0$, and
$\eta x^2h\theta=x\theta(1-E)\le1$, giving the $P_d$ bounds
from (NC2). For $P_4$, if $L\le2$ use $E\ge1/18$; if $L\ge2$,
then $r<1/L\le1/2$ and $P_4\ge L/2$. Its upper bound follows
immediately from (NC2). The last comparison uses $1\le R\le2$
and $1/18\le E<1$.
Consequently the unchanged original source satisfies
\begin{align*}
 C/d&\asymp\frac{\rho}{(1-r)S},&
 A/d&\asymp\frac{\eta(1+L)}{rS},\\
 m/d&\asymp\frac{x\rho h(1+y)}{M\zeta S},&
 \alpha/d&\asymp
 \frac{x^2\zeta\rho\{h(1+y)+\zeta M\}}{QS}.
 \tag{NI2}
\end{align*}
No opening has been bounded away from zero in these statements.
In particular, $E$ and $P_4$ do not supply additional hidden small
source denominators on C1.1, C1.4 or C1.5.

\paragraph{C1.1: two unavoidable large weights and its actual angle.}
Here $x,\zeta$ are infinitesimal and $\eta$ is unbounded.
The common original-$D$ cap gives $q=x\eta$ bounded with
residue at most $1/2$. Hence $y=q\zeta$ is infinitesimal:
the actual angle type is necessarily $(I,I)$, independently of $e$.
The exact formulas give
\[
 C/d\asymp\frac{\eta}{1-r},\qquad
 A/d\asymp\frac{\eta}{r},\qquad
 m/d\asymp\frac{e^2}{M},\qquad
 \alpha/d\asymp xq\zeta,\qquad
 \beta/d+\frac12\asymp\frac{e^2}{t}.             \tag{NI3}
\]
Thus $C,A$ are unbounded relative to $d$, and $\alpha$ is weak.
The common bounded polynomial formula for $B_5$ gives $D/d=O(1/\eta)$,
so $D$ is weak too, including its literal zero. In (NI3), $\beta\ge0$
forces $t=O(e^2)$. Therefore when $e$ is infinitesimal, $t$ is
infinitesimal, $M$ is a unit and $m$ is weak. When $e$ is a unit,
$m$ is strong; when $e$ is unbounded, both $m,\beta$ are unbounded.
These are restrictions on the original C1.1 row, not new certificates.
Its positive infinitesimal $\delta_{\rm tot}$ condition is automatic
because $C/d,A/d$ are unbounded.

\paragraph{C1.4: the natural scale is $L=x\zeta$.}
Now $x$ is infinitesimal and $\zeta$ is noninfinitesimal.
The following implications follow uniformly through all retained
coefficient boundaries.
\begin{enumerate}
\item If $L$ is unbounded, $Q>0$ forces $r$ infinitesimal, and
$1-r=a\eta<2\eta$ makes $\eta$ noninfinitesimal. Hence $y=L\eta$
is unbounded: the actual type is $(I,H)$. Here $\zeta$ is unbounded,
$t=\theta/\zeta$ is infinitesimal, and $M$ is a unit. Since
$h y/\zeta=x\eta h<1$, (NC2) gives $P_\alpha\asymp\zeta$.
Equations (NI1)--(NI2) then yield
\[
 C/d\asymp L^{-1},\qquad
 A/d\asymp r^{-1},\qquad
 \alpha/d\asymp L/Q,\qquad
 m/d\asymp (x\eta h)/\zeta\le O(\zeta^{-1}).
 \tag{NI4}
\]
In particular $C,m$ are weak, $A,\alpha$ unbounded, and
$(m/d)/(C/d)=O(x)$ is infinitesimal. The original total-scale
intersection is automatic.

\item If $L$ is a bounded unit, $\zeta$ is unbounded, $M$ is a
unit, $P_\alpha\asymp\zeta$, and $S\asymp1+\eta$. The identity
$1=r+a\eta$ gives $\rho/(1+\eta)$ a positive bounded unit.
Thus
\[
 C/d\asymp(1-r)^{-1},\qquad
 \alpha/d\asymp Q^{-1},\qquad
 A/d\asymp\frac{\eta}{r(1+\eta)},\qquad
 m/d=O(x).
 \tag{NI5}
\]
So $C,\alpha$ are strong and $m$ weak. Moreover $A$ is weak
exactly when $\eta$ is infinitesimal; that case has $r$ a unit and
$D/d$ unbounded, by $B_5\sim1$ in (NC3). The actual Bob angle is
infinitesimal, a unit, or unbounded exactly as $\eta$ is, because
$y=L\eta$. No restriction on $D,\beta$ is added in the other cases.

\item If $L$ is infinitesimal and $\eta$ is bounded, then $y$ is
infinitesimal. The $B_\beta\ge0$ equation bounds
$\theta\le Kx\rho G/P_d=O(x+L)$, so $h,M$ are units.
Now (NI2) gives $C$ strong, $m,\alpha$ weak, and $A$ weak exactly
when $\eta$ is infinitesimal. In that last case $D/d$ is unbounded,
as in the preceding item. The actual type is $(I,I)$.

\item If $L$ is infinitesimal and $\eta$ is unbounded, then
\[
 C/d\asymp\frac{\eta}{(1-r)(1+y)},\qquad
 A/d\asymp\frac{\eta}{r(1+y)}.                    \tag{NI6}
\]
Both are unbounded since $\eta/(1+y)=(\eta^{-1}+L)^{-1}$.
Also $Q$ is a unit and (NI2) bounds
\[
 \alpha/d\le O\!\left(x^2\zeta\eta h+
              \frac{x^2\zeta^2\eta}{1+y}\right)=O(L),
\]
so $\alpha$ is weak. Finally $J\sim1$ and
$P_d/(1+y)\sim1$ directly in (NC2); (NC4) gives
$D/d+1/2\sim1/2$, hence $D$ weak. This extends the
$C,A$-large/$D,\alpha$-weak source pattern of C1.1 into C1.4,
without claiming that its existing conductors apply. Here $y$ may
be infinitesimal, a unit, or unbounded. The total-scale intersection
is again automatic.
\end{enumerate}
These four alternatives exhaust C1.4, but they are necessary source
reductions only. In particular no entire C1.4 leaf is closed.

\paragraph{C1.5 after the actual unit/unit certificate.}
Here $k$ is a unit. The new theorem removes actual $(U,U)$; the
common both-$\eta,\zeta$-infinitesimal theorem is also subtracted.
Two useful remaining source types can be specified more tightly.

If $v_c$ is infinitesimal, $\eta\zeta$ is infinitesimal. Neither
$\eta$ nor $\zeta$ can be unbounded: $\zeta$ infinitesimal forces
$\theta$ infinitesimal and excludes unbounded $\eta$ by $E>0$;
$\eta$ infinitesimal gives $r\sim1$ and excludes unbounded $\zeta$
by $Q>0$. Thus exactly one of $\eta,\zeta$ is a unit in this residual
subcase. The exact stabilizer (AS5) exchanges these two alternatives.
It keeps $k$ a bounded unit, although a second orientation need not
keep $k\le1$. Choose the representative $\eta$ a unit and
$\zeta$ infinitesimal. Then
\[
 m/d,\beta/d\text{ are unbounded},\qquad
 C/d,A/d\text{ are strong},\qquad \alpha/d\text{ is weak}.
 \tag{NI7}
\]
Indeed $\theta=t\zeta$ is infinitesimal, $E$ is a unit,
$P_d,P_4,P_\alpha,r+\eta$ are units, and (NC4) gives
$m/d\asymp1/(M\zeta)$ and $\beta/d+1/2\asymp1/\theta$.
The reflected representative also has a unit $E$: its exact value
is $E'=h(1+y)/G$ in the original coordinates, with $h$ a unit by
the same $\theta=1$ exclusion used in the ordinary-angle proof.
There is no conclusion about the size of $D/d$ in (NI7).

If $v_c$ is unbounded, both $\eta,\zeta$ are noninfinitesimal and
at least one is unbounded. To see the first assertion, an infinitesimal
$\zeta$ would force $h\sim1$ and $\eta$ bounded by $E>0$;
an infinitesimal $\eta$ would force $r\sim1$ and $\zeta$ bounded
by $Q>0$. Two rate regimes are then useful.
\begin{itemize}
\item If $\eta$ is unbounded and $\zeta$ a unit, then $h$ is
infinitesimal. The exact alternative numerator identity
\[
 B_\beta=(x\eta h)\{K(1+\zeta)-x\theta^2\}
             +xrGT-\theta(1+2xh+x\zeta)             \tag{NI8}
\]
shows that $x\eta h$ is a unit: otherwise its residue at
$\theta=1$ would be $-1-x\zeta(1-r)<0$.
Thus $P_\alpha$ is a unit and $P_d\asymp\eta$.
Equations (NC4) give
\[
 C/d\asymp(1-r)^{-1},\quad A/d\asymp r^{-1},\quad
 m/d\asymp M^{-1},\quad\alpha/d\asymp Q^{-1},
 \quad D/d,\beta/d\text{ weak}.                    \tag{NI9}
\]
For the last assertion use $J\sim1$, $\theta\sim1$ and
$P_d/(Ky)\sim1$ in the original $D/d+1/2,\beta/d+1/2$
formulas. The total-scale condition $\delta_{\rm tot}$ infinitesimal
is therefore equivalent on this typed branch to at least one of
$r,1-r,M,Q$ being infinitesimal. If all four are units, the inherited
positive-total-scale theorem already applies.

\item If $\zeta$ is unbounded, $r$ is infinitesimal and $M$ a
unit. The noninfinitesimal $\eta$ obeys $\eta h\asymp1$.
For if $\theta$ is infinitesimal this follows from $E>0$ and
$h\sim1$; otherwise divide (NI8) by $\zeta$, use
$xrG/\zeta=O(r)$, and obtain that $x\eta h$ has a positive-unit
lower bound. Therefore (NC2) and (NC4) give
\[
 C/d\asymp\zeta^{-1},\quad m/d\asymp\zeta^{-1},\quad
 A/d\asymp r^{-1},\quad\alpha/d\asymp\zeta/Q.
 \tag{NI10}
\]
Thus $C,m$ are weak and $A,\alpha$ unbounded. No size restriction
on $D,\beta$ is inferred. This branch already has infinitesimal
$\delta_{\rm tot}$.
\end{itemize}
All implications retain the original physicality and exact source openings.



% Source formulas: pa133fae35e52
\subsection{The bound on $x\eta$}

\label{component:pa133fae35e52}

Throughout, $d$ is the original positive signed coefficient. The event
coefficient at the fourth cycle position is denoted $A_{\rm ev}$, and
$D_{\rm ev}=e_5$. These must not be confused with actual state coordinates.
The source has the canonical omissions $\lambda=n=0$ and is normalized
by $c=e_3=2$. We use the canonical actual-contact sign entry, including
its covered single-zero $D_{\rm ev}$ or $\beta$ boundaries. Explicitly,
the retained $m,c,A_{\rm ev},\alpha$ are positive, and $D_{\rm ev},\beta$
are positive or exactly one of them is zero. The noncoercive further-zero
faces are already separate exits. The sign and nonsingularity input is the native-entry theorem EN in Chapter~\ref{ch:entry}, proved for every finite noncommuting physical contact in this exact coefficient domain.

\paragraph{Exact transport from the actual contact.}
Write its positive-angle contact as
\[
 u_c=k,\quad v_c=\ell,\qquad
 \psi_c=(S,k(\theta-1),kSr,1),\quad
 S>0,\quad 0<r,\theta<1.
\]
The exact C1 reflection lets us orient $0<k\le1$. Define source-only
parameters
\[
 \eta=\frac{\ell}{kS},\quad \zeta=\frac{S\ell}{k},\quad
 f=\sqrt\eta,\quad g=\sqrt\zeta,\quad
 e=\frac{k}{S}=\frac{kf}{g},\quad
 a=\frac{1-r}{\eta},\quad t=\frac\theta\zeta.
                                                               \tag{CE1}
\]
Thus the classifier $\mathfrak T=k/(S\ell)$ is exactly $\zeta^{-1}$,
and its old state-ratio variable $B$ is exactly $\eta^{-1/2}$.
Put $\kappa=k^2$, $K=1+\kappa$. The original amplitude signs imply
\[
 R=1+\kappa r,\quad T=1+\kappa(1-\theta),\quad
 M=R-t>0,\quad J=T-a>0,
\]
\[
 E=1-\kappa\eta(1-\theta)>0,\qquad
 Q_c=1-\kappa\zeta r>0.                              \tag{CE2}
\]
Indeed the positive second amplitude is $k\zeta M$, the negative
fifth amplitude is $-\ell J$, the state gap is $SE$, and the remaining
state gap is $Q_c$. In particular $0<a,t<2$, without any shape-unit
hypothesis.

The seven native coquotients used by the uniform theorem are valid
under (CE1), not just asymptotically. Their constructed contact is
\[
 (W,X,Y,Z)=(\Delta k,-k(1-t\zeta),\Delta k^2(1-a\eta),1),
 \qquad \Delta=e^{-1}=S/k,
\]
which is exactly the original contact. The native contact identities
prove all six kernel and angular equations. To make the nonsingularity
premise explicit, put $Y_0=Sr$, $h_0=S+k\ell(\theta-1)$,
$R_2=\ell(S+k^2Y_0)-k\theta$, and
$R_5=k(S-Y_0)+k^2\ell(\theta-1)-\ell$. The common selected
six-column contact system, on $(d,m,A_{\rm ev},D_{\rm ev},\alpha,\beta)$
with $c=2$, has determinant
\[
 -8192\theta k^2\ell(1+k^2)^2(1+\ell^2)^2\mathcal D,
\]
\[
 \mathcal D=Y_0(1+k^2)(kY_0+\ell)(1-k\ell Y_0)h_0R_2(-R_5)>0.
\]
Each factor is strictly positive by the stated source sign theorem.
Thus the normalized coefficients equal the unique original source
coefficients in that common system. This
establishes full original-coefficient transport, including the exact
zero openings; it does not use an archived completion label.

\paragraph{A bounded polynomial consequence of the original fifth event.}
Assume now
\[
                 k\text{ infinitesimal},\qquad
                 \zeta=\mathfrak T^{-1}\text{ infinitesimal}.     \tag{CE3}
\]
Set $q=\kappa\eta$. Since $t<2$ and $E=1-q(1-t\zeta)>0$, $q$ is
bounded and $\bar q\le1$. The original fifth-event formula is
\[
 D_{\rm ev}=\frac{a}{Jc_3}(2c_5-J\eta c_d).
\]
Using $L_c=r$, $H_c=r+\eta$ and $V_c=q\zeta$, its sign is precisely
\[
 \mathcal B_5:=K(1+q\zeta)(r+\eta)-J\eta P_d\ge0.       \tag{CE4}
\]
Although $\eta$ need not be bounded, this whole expression has the
following bounded polynomial form. Define
\begin{align*}
 H={}&2+\kappa r+\zeta[K-(\kappa+2)t]
     +q\zeta(-\zeta t^2+t+1)
     +\kappa\zeta(1-r)(t-1),\\
 P_d={}&1+q\zeta+\kappa H.
\end{align*}
Then exact cancellation using $a\eta=1-r$ and $q=\kappa\eta$ gives
\[
 \mathcal B_5=Kr(1+q\zeta)+(1-r)P_d
       +q\{\zeta t(1+q\zeta)-(K-\kappa\zeta t)H\}.     \tag{CE5}
\]
Every variable on the right is bounded. Hence (CE4) gives
\[
       1-2\bar q\ge0,\qquad 0\le\bar q\le\tfrac12.     \tag{CE6}
\]
This is obtained before dividing by $1-q$ or by a coefficient opening.
In particular $E$ is a positive unit with residue at least $1/2$.

The fourth coquotient polynomial likewise has the bounded expression
\[
\begin{split}
 P_4={}&K(1-q)+\kappa(1-r)-\kappa\zeta t\\
 &+q\zeta(-\kappa\zeta t^2+2\kappa t+t)
   +\kappa\zeta(1-r)(1-t).
\end{split}                                                       \tag{CE7}
\]
Thus $\bar P_d=1$ and $\bar P_4=1-\bar q\ge1/2$.
These exact low-degree identities are the physicality input needed
for the entry, rather than an estimate on a truncated coefficient.



% Source formulas: p0acdb66bba2f
\section{Contact identities and complementary measurement charts}
The scalar contact polynomial supplies the first square roots. Polynomial division then expresses state couplings through these roots. The complementary measurement charts follow directly from the positive event terms and will be used throughout the remaining proof.

\subsection{Polynomial source coordinates and exact contact balances}

\label{component:p0acdb66bba2f}

\label{p0acdb66bba2f:app:c1-source-definitions}

The letter $d$ in the contact--conductor proof denotes the new parameter
$\delta$, not the Bell coefficient. Here the latter is denoted $d_{\rm B}$,
and the original event coefficient $e_5$ is denoted $D_{\rm e}$.
The testing factor is always $D=1+u^2$. The conductor monomial
$C=\delta b^2uw$ is distinct from the clearing multiplier
$\mathcal C=K^6Q^4T_n$.

\paragraph{Scope of the polynomial identities.} The seven coquotient formulas and the contact identities (SD9)--(SD11) are exact for arbitrary positive native parameters. The size assumptions below are needed only for the displayed bounded-Gram application.
\paragraph{Domain and seven native coquotients.}
The parameters $k,f,g>0$ belong to the source field; put $e=kf/g$ and $\Delta=e^{-1}$. The formulas are exact before a size condition is imposed. The shapes $a,t$ are positive bounded elements;
\[
 R=1+k^2(1-af^2),\qquad T=1+k^2-k^2g^2t,
 \qquad M=R-t>0,\quad J=T-a>0.
\]
All original event coefficients are assumed nonnegative. In particular,
the formulas below retain the requirements $D_{\rm e},\beta\ge0$.
No relation compares a testing rate with a coefficient rate.
Put
\[
 \kappa=k^2,\quad \eta=f^2,\quad \zeta=g^2,
 \quad V_c=\kappa\eta\zeta,\quad
 L_c=1-a\eta,\quad H_c=1+(1-a)\eta,\quad Q_c=1-\kappa\zeta+aV_c.
\]
Define
\begin{align*}
P_4={}&1+\kappa+\eta\kappa(a-\kappa-1)-\zeta\kappa t
 +V_c[-\zeta\kappa t^2+2\kappa t+a(1-t)+t],\\
P_\alpha={}&1+\kappa-a\eta\kappa
 +\zeta[(1+\kappa)-(2+\kappa)t]
 +V_c[-\zeta t^2+a(t-1)+t],\\
P_d={}&(1+\kappa)^2-a\eta\kappa^2
 +\zeta\kappa[(\kappa+1)-(\kappa+2)t]\\
 &+V_c[-\zeta\kappa t^2+\kappa a(t-1)+\kappa t+\kappa+1].
\end{align*}
The complete seven coquotients are
\begin{align*}
c_2&=2(1+\kappa)(1-t\zeta)L_c(1+V_c)H_cQ_c,\\
c_b&=2(1+\kappa)L_c[1+(1-t)\zeta]H_cQ_c,\\
c_3&=2(1+\kappa)L_cH_cQ_c(1-\kappa\eta+tV_c),\\
c_5&=2(1+\kappa)L_cH_cQ_c(1+V_c),\\
c_4&=2Q_cP_4,\qquad c_\alpha=2L_cH_cP_\alpha,
 \qquad c_d=4L_cQ_cP_d.
\end{align*}
At $\eta=\zeta=0$, the first six equal $2(1+k^2)$ and $c_d=4(1+k^2)^2$.
Here $c_b$ is the original $c_\beta$, and the subscript on $c_d$ is a
name rather than the conductor parameter.

\paragraph{Original source and full parent.}
The original event coefficients and multiplier are
\[
\begin{gathered}
\lambda=n=0,\qquad c=2,\qquad
 d_{\rm B}=\frac{2a\eta c_d}{c_3},\qquad
 m=\frac{2ae^2c_2}{Mc_3},\qquad A=\frac{2a\eta^2c_4}{c_3},\\
\alpha=\frac{2a\Delta^2k^6\eta^2c_\alpha}{c_3},\qquad
D_{\rm e}+\frac{d_{\rm B}}2=\frac{2ac_5}{Jc_3},\qquad
\beta+\frac{d_{\rm B}}2=\frac{2ae^2c_b}{tc_3},\qquad
\Lambda=\frac{RTc_3}{2a}.
\end{gathered}                                                    \tag{SD1}
\]
Let $F$ denote the original complete parent, including its signed core.
For a testing state $(W,X,Y,Z)$ define
\[
\begin{gathered}
 D=1+u^2,\quad D_v=1+v^2,\quad B=X+uZ,\quad Q_A=W+uY,\\
 r=uW-Y,\quad r_2=vQ_A-B,\quad r_5=r-v(Z-uX).
\end{gathered}
\]
With $H_2=4Tc_2$, $H_b=4Tc_b$, $H_3=4Rc_3$, $H_5=4Rc_5$,
the exact entire parent is
\begin{align*}
\Phi:=\Delta^2\Lambda F={}&H_2\frac RM r_2^2
 +H_b\frac RtD_v B^2+\Delta^2H_3\frac TaD_vr^2
 +\Delta^2H_5\frac TJr_5^2\\
 &+4RT\Delta^2\eta^2c_4DD_vY^2
 +4RT\Delta^4k^6\eta^2c_\alpha D(vY-Z)^2\\
 &+2RT\Delta^2\eta c_dD\{(vY-Z)^2-D_vZ^2\}.
                                                               \tag{SD2}
\end{align*}
Thus every original source event is accounted for.

\paragraph{Retained terms and state coordinates.}
In the remaining chart $Z=1$ with $u,v,W,X,Y$ infinitesimal, put
\[
 p=t/R,\quad h=a/T,\quad q_A=v(W+uY),\quad q_B=v(1-uX),
 \quad N=B-pq_A,\quad M_B=r-hq_B.
\]
Retain the four nonnegative terms
\[
 \Pi_1=H_2\frac tM N^2,\quad
 \Pi_2=H_b\frac MtD_vN^2,\quad
 \Pi_3=\Delta^2H_3\frac JaD_vM_B^2,\quad
 \Pi_4=\Delta^2H_5\frac aJM_B^2.                         \tag{SD3}
\]
Their subtraction is the exact identity
\begin{align*}
G=\Phi-\sum_j\Pi_j={}&H_2\{N^2-2q_AN+(1-p)q_A^2\}
 +H_bD_v\{N^2+2q_AN+pq_A^2\}\\
 &+\Delta^2H_3D_v\{M_B^2+2q_BM_B+hq_B^2\}
 +\Delta^2H_5\{M_B^2-2q_BM_B+(1-h)q_B^2\}
 +\mathcal W,\\
\mathcal W={}&4RT\Delta^2\eta^2c_4DD_vY^2
 +4RT\Delta^4k^6\eta^2c_\alpha D(vY-1)^2\\
 &+2RT\Delta^2\eta c_dD\{(vY-1)^2-D_v\}.
                                                               \tag{SD4}
\end{align*}
The testing coordinate $w=(W+uY)/D$ has the full inverse
\[
\begin{gathered}
 X=N+pDvw-u,\qquad q_B=v(D-uN-puDvw),\\
 W=w+\frac{u(M_B+hq_B)}D,\qquad
 Y=uw-\frac{M_B+hq_B}D.
\end{gathered}                                \tag{SD5}
\]
Its only testing denominator is the positive unit $D$.
On the negative-core branch, sign gauges give $u,v>0$, and $w>0$:
after eliminating the two positive core squares the remaining scalar is
\[
 \frac{uvDD_v}{\mathcal D}\{Duv w^2-2DD_vw+uv\},\qquad
 \mathcal D=u^2v^4+u^2v^2+u^2+v^4+2v^2+1>0,
\]
which is nonnegative if $w\le0$. The eliminated matrix is
$\left(\begin{smallmatrix}1+v^2+u^2v^2&uv\\uv&(D+v^2)/D\end{smallmatrix}\right)\succ0$.
The nonnegative-core branch and literal axes already give direct SOS.
Thus the radicals in the conductor proof separate between the two fields.

\paragraph{Short native proof of the exact contact identities.}
Write $V=1+V_c$ and $K=1+\kappa$ in this paragraph. The factored
coquotients above give immediately
\[
 c_b-c_2=\zeta c_3,\qquad c_3-c_5=-\kappa\eta c_b,
 \qquad c_2=(1-t\zeta)c_5.                            \tag{SD9}
\]
Hence $Vc_b-c_2=\zeta c_5$ and
$Vc_3-c_5=-\kappa\eta c_2$.
The remaining three balances are the following short polynomial
identities; substituting the displayed $P_4,P_\alpha,P_d$ proves them
by multiplication and collecting terms:
\begin{align*}
\mathfrak Y:={}&2\kappa c_2+2VL_cc_4
                 -2\kappa^2\zeta Q_cc_\alpha-Q_cc_d=0,\\
\mathfrak S:={}&\kappa L_cc_2+c_5+VL_c^2c_4+\kappa Q_c^2c_\alpha
                 +\tfrac12(\kappa\zeta L_c^2-2L_c-\eta)c_d=0,\\
\mathfrak V:={}&\kappa Rc_2+t\kappa V_cc_b+aV_cc_3+Tc_5
 +KV_cL_c^2c_4-\kappa^2\zeta KQ_cL_cc_\alpha
 -\tfrac K2(Q_cL_c+\eta)c_d=0.                        \tag{SD10}
\end{align*}
For example, after the common factor $4L_cQ_c$ is removed,
the first balance is just
\[
 \kappa K(1-t\zeta)VH_c+VP_4
       -\kappa^2\zeta H_cP_\alpha-Q_cP_d=0.
\]
The two other balances are equally direct identities in the same three
short polynomials; no expanded scalar or ideal calculation enters.

The exact contact, in the original coordinates of (SD2), is
\[
 u=k,\quad v=\Delta k^2\eta,\qquad
 (W,X,Y,Z)=(\Delta k,-k(1-t\zeta),\Delta k^2L_c,1).      \tag{SD11}
\]
At this point
\[
 r_2=Mk\zeta,\quad B=tk\zeta,\quad r=av,\quad
 r_5=-Jv,\quad vY-1=-Q_c,\quad D=K,\quad D_v=V.
\]
Set $\Psi=\Phi/(4RT)$ and let $E_s=\tfrac12\partial_s\Psi$,
with all other original testing coordinates held fixed.
Direct differentiation of the seven terms in (SD2), followed by (SD9),
gives at (SD11)
\begin{align*}
 \Psi&=K\zeta^2\mathfrak S,\\
 E_W&=kv\{\zeta c_2+\Delta^2(Vc_3-c_5)\}=0,\\
 E_X&=k\zeta\{Vc_b-c_2-\zeta c_5\}=0,\\
 E_Y&=\tfrac12K\Delta^3\kappa\eta^2\mathfrak Y,\\
 E_u&=k\zeta^2\{\mathfrak S-c_2+(1-t\zeta)c_5\},\\
 E_v&=v\Delta^2\mathfrak V.
\end{align*}
They all vanish by (SD9)--(SD10). Homogeneity in $(W,X,Y,Z)$ also gives
$E_Z=0$. This proves contact and full first-order stationarity directly
in the native source.



% Source formulas: p58894e12bcef
\subsection{A finite four-monomial circuit}

\label{component:p58894e12bcef}

Let
\[
a_1=\sqrt{wv},\quad a_2=\sqrt{\delta uv},\quad
a_3=b\sqrt{\delta kuw},\quad a_4=\delta k^{3/2}b.
\]
The radicals are separated coefficient/testing radicals. Put
\[
h_{13}=a_1^2-a_3^2,\quad h_{23}=a_2^2-a_3^2,\quad
h_{34}=a_3^2-a_4^2,
\qquad q_1=a_1a_2-a_3a_4,\quad q_2=a_1a_3-a_2a_4.
\]
Averaging the three-square arithmetic--geometric-mean identities for
pairings $12|34$ and $13|24$ gives exactly
\begin{equation}
\sum_{i=1}^4a_i^4-4a_1a_2a_3a_4
=h_{13}^2+h_{23}^2+h_{34}^2-h_{13}h_{23}+h_{23}h_{34}
 +q_1^2+q_2^2.\label{p58894e12bcef:eq:compact-circuit}
\end{equation}
The contrast block is
\[
G=\begin{pmatrix}1&-1/2&0\\-1/2&1&1/2\\0&1/2&1\end{pmatrix}\succ0;
\qquad \det G_{[1]}=1,\quad\det G_{[2]}=3/4,\quad\det G=1/2.
\]
Unlike a single pairing, the two pairings control all three independent
contrasts of the four squared monomials.

Write $K=1+k^2$, $D=1+u^2$ and
\[
g_0=\delta(v-b^2uw),\qquad g_4=uwv-\delta^2k^4b^2.
\]
The leading scalar used in the complete source is
\[
P=D(w^2+\delta^2)v^2+\delta^2Kb^4u^2w^2+
\delta^4Kk^6b^4-2\delta^2K^2b^2uwv
=g_0^2+g_4^2+\sum_i a_i^4-4\prod_i a_i.
\]
There is a second exact simplification:
\begin{equation}
g_4=u h_{13}+k h_{23}-ku g_0+k h_{34}.\label{p58894e12bcef:eq:g4-conductor}
\end{equation}
If $k,u$ are infinitesimal, the four conductor coefficients
are infinitesimal. Otherwise retain $g_4^2$ as a separate square.
The two-pairing contrast matrix displayed above is strict in either case.


% Source formulas: p466b31015532
\subsection{The restricted three-monomial contact lift}

\label{component:p466b31015532}

\label{p466b31015532:sec:three-circuit}

Work in the same tensor ring $T=K_c\otimes_{\mathbb R}K_t$ and bounded
subring $B$ as throughout this package. Every inverse and radical below
is of a pure-field or positive separated element. Let $k>0$ be a
pure-source infinitesimal and $s>0$ a positive pure-source bounded unit.
Let $u>0$ be a pure-testing infinitesimal and $t>0$ a positive
pure-testing bounded element (possibly infinitesimal). Define
\[
 A_1=u^2t,\qquad A_2=\frac{ku t^2}{s},\qquad A_3=k^2s,
 \qquad G=kut,\qquad d_i=A_i-G.
\]
Then $A_1A_2A_3=G^3$. Two exact conductor identities are
\begin{equation}\label{p466b31015532:T3:conductors}
 G(u-k)=k d_1,\qquad G(t-s)=s d_2.
\end{equation}
They permit a unit testing ratio $t$ without comparing a mixed tensor.

Put $Q=G\sum_i d_i$ and $L=\sum_i A_i^2-3G^2$.
For $a_i=A_i^{1/3}$, the elementary cubic identity gives
\begin{align}\label{p466b31015532:T3:sos}
 Q&=\frac12\sum_{1\le i<j\le3}\sum_{\ell=1}^3
       \{\sqrt{G a_\ell}(a_i-a_j)\}^2,\\
 L&=\sum_i d_i^2+2Q.\nonumber
\end{align}
Thus $L$ has a positive constant diagonal Gram on a finite list of
three difference roots and nine cubic-identity roots. No assertion
that every nonnegative tensor is SOS is used.

\paragraph{Exact finite lift.}
Let $\lambda\in B_c$ have positive real residue, and consider
\[
 F=\sum_{i=1}^3(\lambda+\epsilon_i)A_i^2
       -(3\lambda+\epsilon_g)G^2,
 \qquad h=\sum_i\epsilon_i-\epsilon_g.
\]
Write $\Delta_u=u-k$, $\Delta_t=t-s$. Assume explicit finite identities
\begin{align}
 \epsilon_i-\epsilon_{i0}&=K_{iu}\Delta_u+K_{it}\Delta_t,
       &\epsilon_i,\epsilon_{i0}&\in B_{\rm Inf},\label{p466b31015532:T3:coeff}\\
 h&=b_u\Delta_u+b_t\Delta_t+
          \sum_{j,\ell\in\{u,t\}}H_{j\ell}\Delta_j\Delta_\ell,
       &b_u,b_t&\in K_c,\label{p466b31015532:T3:taylor}\\
 2(\epsilon_{10}-\epsilon_{30})+k b_u&=0,
 &2(\epsilon_{20}-\epsilon_{30})+s b_t&=0.\label{p466b31015532:T3:jets}
\end{align}
Require $K_{iu},H_{uu},H_{ut},H_{tu}\in B$ and
$K_{it},H_{tt}\in B_{\rm Inf}$. The distinction is essential:
the second conductor has a unit coefficient, so boundedness alone of
the second-coordinate arrays does not suffice.

Then $F$ is tensor SOS. Indeed, expanding $A_i=G+d_i$ and applying
\eqref{p466b31015532:T3:conductors} and \eqref{p466b31015532:T3:jets} yields the exact identity
\begin{align}\label{p466b31015532:T3:error}
 F-\lambda L={}&\sum_i\epsilon_i d_i^2+2\epsilon_{30}Q
 +2\sum_i(kK_{iu}d_1d_i+sK_{it}d_2d_i)\\
 &+k^2H_{uu}d_1^2+ks(H_{ut}+H_{tu})d_1d_2
          +s^2H_{tt}d_2^2.\nonumber
\end{align}
In detail, the two bare-$G$ first-order terms are cancelled by
$G^2b_u\Delta_u=kb_uGd_1$ and
$G^2b_t\Delta_t=sb_tGd_2$. All remaining coefficients in
\eqref{p466b31015532:T3:error} are bounded infinitesimals. The strict diagonal
Gram for $\lambda L$ in \eqref{p466b31015532:T3:sos} therefore remains strictly
positive in real residue. The bounded-Gram lemma supplies a finite
tensor SOS. This argument uses no power series or mixed inverse.

\paragraph{Normal extension.}
Suppose a coupling to a normalized hard state normal has the exact form
$\ell=\sum_i c_iA_i$ with $c_i\in B_{\rm Inf}$,
$\sum_i c_{i0}=0$, and
\[
 \sum_i(c_i-c_{i0})=J_u\Delta_u+J_t\Delta_t,
 \qquad J_u\in B,\quad J_t\in B_{\rm Inf}.
\]
Then
\[
 \ell=\sum_i c_i d_i+kJ_u d_1+sJ_t d_2
\]
has infinitesimal coefficients on the same root list. It can be joined
to a hard normal block with strictly positive real residue. Extra
already-controlled scalar roots may be appended, provided their
coupling coefficients are likewise infinitesimal. The source-specific
arrays and hard-block bounds remain separate proof obligations.

\paragraph{Small-$L$ coordinates.}
For $x=k^2$, $s=\sqrt{\eta\zeta}$ and $t=v/u$, this gives
\[
 (A_1,A_2,A_3,G)=
 \left(uv,\sqrt{\frac{x}{\eta\zeta}}\frac{v^2}{u},
       \sqrt{\eta\zeta}\,x,\sqrt{x}\,v\right).
\]
The lemma applies only when the stated pure-field size assumptions and
all finite scalar and normal identities have been proved. A leading
three-term AM--GM form alone does not close a source domain.



% Source formulas: p89fd3edb3edd
\subsection{The complementary projective testing charts}

\label{component:p89fd3edb3edd}

\label{p89fd3edb3edd:sec:r6-c1-common-remote}
Normalize the original parent by its positive pure-source coefficient $d$.
Write $D_z=1+z^2$, and write the signed core as $2K$, where
\[
 K=(uW+uvX-Y-vZ)^2+D_u(vY-Z)^2+D_v(X+uZ)^2-D_uD_vZ^2.
 \tag{CE1}
\]
The event terms used below are
\[
4m(vW-X+uvY-uZ)^2,\quad4CD_v(uW-Y)^2,
\quad4AD_uD_vY^2,\quad4\alpha D_u(vY-Z)^2.
\]
Every omitted event and every nonnegative excess is retained as its exact
original square. Source and testing fields remain independent. ``Strong''
means bounded below by a positive real number; ``unbounded'' means larger
than every real number. Neither assertion constrains upper scales.

\begin{theorem}[Two adjacent unbounded events, with complete raw frames]
If $C,A$ are unbounded, the full original parent is tensor SOS at every
projective test for which $u$ is not infinitesimal, and every state.
No physicality assumption beyond nonnegative event weights is needed for
this particular assertion. By the exact cyclic relabeling below, if
$A,\alpha$ are unbounded, the conclusion holds whenever $v$ is not
projectively unbounded.
\end{theorem}
\begin{proof}
Select $C=A=\mu$, with a positive real $\mu$ chosen after the real testing
residues; unboundedness permits every such choice. The following are exact
full quadratic forms, before any restriction to a state ray or contact.
They therefore retain every first state normal and hard-normal coupling.

For $u$ a bounded nonzero unit and $v$ bounded, use the four raw roots
\[
 (W,Y,B,V),\qquad B=X+uZ,\quad V=vZ.
\]
The entire core equals
\[
 K=(uW-Y+uvB)^2+D_vB^2+D_uv^2Y^2+D_uu^2V^2
       -2D_uuV(W+vB).                         \tag{CE2}
\]
The two selected event forms remain $4\mu D_v(uW-Y)^2$ and
$4\mu D_uD_vY^2$. In their common real kernel $W=Y=0$, the core has
matrix on $(B,V)$
\[
 \begin{pmatrix}1+v^2D_u&-D_uuv\\-D_uuv&D_uu^2\end{pmatrix},
 \qquad \det=D_uu^2>0.                       \tag{CE3}
\]
Here and in residue matrices, bounded testing letters mean their real
residues. Thus a sufficiently large finite $\mu$ gives a strictly positive
residue on the entire four-root space, even when $v=0$ literally.

For $p=1/u$ infinitesimal and $v$ bounded, homogenize by $p^2$ and take
\[
 (W,Y,B,R),\qquad B=pX+Z,\quad R=vX.
\]
The exact full core is
\[
 K_p=B^2-2pvBW-2vBY+D_pR^2+2D_pRW
                  +(W-pY)^2+D_pv^2Y^2.       \tag{CE4}
\]
The selected terms are $4\mu D_v(W-pY)^2$ and
$4\mu D_pD_vY^2$. At $p=0$ their sum with $2K_p$ is
\[
 4\mu D_v(W^2+Y^2)+2(W+R)^2+2(B-vY)^2,       \tag{CE5}
\]
strict for every $\mu>0$. The formula allows literal $p=0$ and $v=0$;
the raw roots need not form independent state coordinates.

For projectively unbounded $v$, put $w=1/v$, homogenize by $w^2$, and
use $(W,X,Y,H)$ with $H=wZ$. The exact full core is
\begin{align*}
 K_w={}&u^2H^2-2u^2HY-2uHW+2uwHX+u^2w^2W^2\\
 &+2u^2wWX-2uw^2WY+(D_u+w^2)(X^2+Y^2)-2uwXY.
                                                        \tag{CE6}
\end{align*}
The selected terms are $4\mu D_w(uW-Y)^2$ and
$4\mu D_uD_wY^2$. If $u$ is a bounded nonzero unit, at $w=0$ their
common kernel has $W=Y=0$ and $K_w=D_uX^2+u^2H^2$ strictly positive.
A sufficiently large real $\mu$ again works. If also $p=1/u$ is
infinitesimal, homogenize (CE6) by $p^2$. The exact full core is
\begin{align*}
 K_{pw}={}&H^2-2pHW+2pwHX-2HY+w^2W^2+2wWX-2pw^2WY\\
 &+(D_p+p^2w^2)(X^2+Y^2)-2pwXY.              \tag{CE7}
\end{align*}
At $p=w=0$ the full selected residue is
\[
 4\mu(W^2+Y^2)+2X^2+2(H-Y)^2>0.             \tag{CE8}
\]
All these exact Gram coefficients are bounded. Strictness in (CE3) and
its counterpart following (CE6) implies a finite cap by the elementary
real Schur-complement argument: the event Gram is positive definite on
the complement of its kernel, and the core is positive definite on the
kernel. No source/testing rate comparison is involved. Apply the bounded
positive-residue Gram lemma, and restore every retained original square.

For the second assertion, compose the common Alice reflection with
the swap/reflection. The resulting one-step inverse cycle rotation is
\[
 (u',v')=(1/v,u),\quad
 (W',X',Y',Z')=(Y+vZ,-W-vX,Z-vY,vW-X),          \tag{CE9}
\]
and sends event index $i$ to $i-1$ in
$(\lambda,D,C,A,\alpha,m,\beta,n)$. In particular original $A,\alpha$
become new $C,A$. Its full-parent identity is
\[
 v^2F(\mathrm{rot}^{-1}c;1/v,u;\psi')=D_vF(c;u,v;\psi).
                                                        \tag{CE10}
\]
The state map has determinant of absolute value $D_v^2$, and its
homogeneous version at $v=\infty$ is invertible as well. Thus
$u'=1/v$ not infinitesimal is exactly $v$ not projectively unbounded.
Only pure testing factors are restored. Both signs, all literal endpoints,
and complex states (by real/imaginary splitting) are included.
\end{proof}

\begin{theorem}[Four consecutive strong events]
For an original physical parent with $C,A,\alpha,m$ strong, the full
parent is tensor SOS at every test except possibly $u$ a bounded nonzero
unit and $v$ projectively unbounded. All original upper coefficient
scales and all nonnegative unselected coefficients are allowed.
\end{theorem}
\begin{proof}
Cap these four events at any positive real $\mu$ smaller than all four
weights, and retain every excess. Noncommuting-residue tests use the
common physical-source theorem, applied to the original full parent.
For commuting-residue tests use the capped parent.

First suppose $v$ is bounded. On the three open commuting axes, the core
and the four caps have zero common kernel: at $u=0$, the core forces
$X=Y=0$, then $m$ and $\alpha$ force $W=Z=0$ when $v$ is a unit;
at $u=\infty$, $C,A$ force $W=Y=0$, then the core forces $X=Z=0$;
at $v=0$, $A,\alpha$ force $Y=Z=0$, then the core forces $W=X=0$
when $u$ is a unit. These are strictly positive bounded state Grams.

At $(u,v)=(0,0)$ the residue is
\[
 (2+4\mu)X^2+(2+8\mu)Y^2+4\mu Z^2.
\]
Only the $W$ ray remains. Normalize $W=1$ and take the five infinitesimal
roots $(u,v,X,Y,Z)$. The full polynomial has no constant or linear term;
its quadratic part is
\[
 2(u-Y)^2+2X^2+
 4\mu\{(v-X)^2+(u-Y)^2+Y^2+Z^2\},            \tag{CE11}
\]
which is strictly positive. At $(u,v)=(\infty,0)$ only the $X$ ray
remains. Normalize $X=1$ and use $(p,v,W,Y,Z)$. Its quadratic part is
\[
 2(W+v)^2+2(p+Z)^2+
 4\mu\{(p+Z)^2+W^2+Y^2+Z^2\},              \tag{CE12}
\]
also strictly positive. All higher monomials in each chart allocate to
products of two displayed roots with infinitesimal bounded coefficients.
This is an exact finite Gram, not a truncation certificate.

For $w=1/v$ infinitesimal use the full raw frame of (CE6). All four
selected terms, with every normal present, are exactly
\[
 4\mu\{(W-wX+uY-uH)^2+D_w(uW-Y)^2
              +D_uD_wY^2+D_u(Y-H)^2\}.        \tag{CE13}
\]
At $u=w=0$ their sum with $2K_w$ has residue
\[
 4\mu W^2+2X^2+(8\mu+2)Y^2+4\mu(Y-H)^2,    \tag{CE14}
\]
strict on all four roots. If $p=1/u$ is infinitesimal, homogenize again;
the exact terms become
\[
 4\mu\{(pW-pwX+Y-H)^2+D_w(W-pY)^2
              +D_pD_wY^2+D_p(Y-H)^2\}.        \tag{CE15}
\]
Together with (CE7), at $p=w=0$ the residue is
\[
 4\mu W^2+2X^2+4\mu Y^2+(8\mu+2)(Y-H)^2,    \tag{CE16}
\]
again strict. These frames require no state normalization, and no
division by either small testing coordinate. Literal $w=0$ is included.
The only chart not covered is ordinary $u$ and infinitesimal $w$.
Restore the original terms and all pure factors as above.
\end{proof}

\begin{lemma}[Matched weak outer events at the large-Bob corner]
Suppose $q>0$ is a pure-source infinitesimal and
\[
 C\ge c q^2,\quad m\ge m_*q^2,\quad
 A\ge a_*/q^2,\quad\alpha\ge b_*/q^2
                                                        \tag{CE17}
\]
for positive real $c,m_*,a_*,b_*$. At projectively unbounded $v$,
the complete parent is tensor SOS whenever $u$ is infinitesimal or
projectively unbounded. In the latter chart the lower bound on $m$
is unnecessary.
\end{lemma}
\begin{proof}
Keep $w=1/v$ and $H=wZ$ as in (CE6)--(CE7). If $u$ is infinitesimal,
use the five raw roots
\[
 E=qW,\quad X,\quad F=Y/q,\quad G=H/q,\quad V=uwW.
\]
The exact whole core (CE6) becomes
\begin{align*}
 K_w={}&V^2+(D_u+w^2)X^2
       +q^2\{u^2(G-F)^2+D_wF^2\}\\
 &-2uGE+2quwX(G-F)+2uVX-2uw^2EF.             \tag{CE18}
\end{align*}
Cap only $m,A,\alpha$ at the indicated weights. Their complete terms
are respectively
\[
 4m_*(E-qwX+uq^2(F-G))^2,\quad
 4a_*D_uD_wF^2,\quad4b_*D_u(F-G)^2.
\]
The selected residue is
\[
 2V^2+2X^2+4m_*E^2+4a_*F^2+4b_*(F-G)^2>0.   \tag{CE19}
\]
Every coefficient in (CE18) is bounded independently of all relative
source/testing rates.

If instead $p=1/u$ is infinitesimal, use
\[
 E=qW,\quad F=Y/q,\quad G=H/q,\quad
 R=X+wW,\quad S=pX.
\]
The exact whole core (CE7) is
\[
 K_{pw}=R^2+D_wS^2+q^2(G-F)^2+p^2D_wq^2F^2
       -2pE(G+w^2F)+2qwS(G-F).               \tag{CE20}
\]
Cap only $C,A,\alpha$. Their complete terms are
\[
 4cD_w(E-pq^2F)^2,\quad4a_*D_pD_wF^2,
 \quad4b_*D_p(F-G)^2.
\]
Together with $2K_{pw}$ their residue is
\[
 2R^2+2S^2+4cE^2+4a_*F^2+4b_*(F-G)^2>0.     \tag{CE21}
\]
Only the pure-source inverse $1/q$ occurs in a root. These exact raw
frames include every state normal and every literal endpoint. The
strict-residue lemma and restoration of excess terms prove the claim.
\end{proof}

\paragraph{Exact application to the original C1 source rows.}
The source premises here are precisely the already proved original
formulas (NC1)--(NC7), (NI3)--(NI10); no independent source family is
substituted. These common theorems give the following new quantified
test reductions, always intersected with prior certificates:
\begin{enumerate}
\item C1.1, every $e$ type: (NI3) gives $C,A$ unbounded;
only $u$ infinitesimal remains.
\item C1.4, $L$ infinitesimal and $\eta$ unbounded: (NI6)
again gives $C,A$ unbounded; only $u$ infinitesimal remains.
\item C1.4, $L$ unbounded: (NI4) and (CE20) leave only
$u$ bounded and $v$ unbounded.
\item C1.5 actual UH, $\eta$ unbounded and $\zeta$ unit:
(NI9) gives four strong events; only $u$ unit and $v$ unbounded remains.
\item C1.5 actual UH, $\zeta$ unbounded: (NI10) and (CE17)
leave the same sole strip, $u$ unit and $v$ unbounded.
\end{enumerate}
For the last line choose $q=\zeta^{-1/2}$. Formula (NI10) gives
$C,m\asymp q^2$, $\alpha\asymp1/(q^2Q)$ and $A\asymp1/r$;
the original $Q=1-x\zeta r>0$, with unit $x$, bounds $Aq^2$ below.
Thus (CE17) holds with positive real constants. The adjacent-event
theorem covers bounded $v$ and (CE18)--(CE21) remove both end types of
$u$ at unbounded $v$. Both actual UH source patterns now have the same
sole untreated testing strip, ordinary $u$ and unbounded $v$.

In the $L$-unbounded C1.4 line take $q=L^{-1/2}$.
Then (NI4) gives $C\asymp q^2$, $\alpha\asymp1/(q^2Q)$ and
$Aq^2\asymp1/(Lr)>1$. The $p$-chart (CE20), which needs no lower
bound on $m$, removes the projectively unbounded-$u$ tests there.

All bounds are relative to original $d$. The four-source-event result
does not require the additional NI9 total-scale condition that at least
one opening be infinitesimal. Every literal permitted $D=0$ or
$\beta=0$ is retained throughout, since neither coefficient was capped.
No claim of nonemptiness or full-row completion is made. In particular,
the $L$ unit and the $L$ infinitesimal/$\eta$ bounded parts of C1.4 are
unchanged, and these remote results do not supply the missing critical
contact conductors.



% Source formulas: p6cad7544d1b8
\subsection{Two exact complementary-current identities}

\label{component:p6cad7544d1b8}

\label{p6cad7544d1b8:sec:c1-uniform-all-tests}

Let $\mathcal H=2[(uW+uvX-Y-vZ)^2+D_u(vY-Z)^2+D_v(X+uZ)^2-D_uD_vZ^2]$ be the original testing core.

\begin{lemma}[Two strong currents]\label{lem:c1-two-current}
Suppose the original event coefficients are nonnegative, $d>0$, and
$d/c,d/D$ are infinitesimal, where $c=e_3$, $D=e_5$.
Then every testing chart with $1/v$ bounded is tensor-SOS, for every
other angle and every state.
\end{lemma}
\begin{proof}
Put $g=cD/(c+D)$, $h=D/(c+D)$, $R=uW-Y$,
$C=v(uX-Z)$, $B=X+uZ$, and $N=R+hC$.
Let $D_u=1+u^2$, $D_v=1+v^2$ and set
\[
 K_0=2\sqrt g\,C,\qquad
 B_0=\sqrt{2d}\sqrt{D_v/D_u}\,B,\quad
 c_u=D_u^{-1/2},\quad t_u=uD_u^{-1/2},\quad
 \eta_v=\sqrt{D_v}/v.
\]
The complete selected parent is exactly
\[
\begin{split}
4cD_vR^2+4D(R+C)^2+d\mathcal H
={}&4(c+D)N^2+4cv^2R^2+2d(R+C)^2\\
 &+2dD_u(vY-Z)^2+(K_0,B_0)Q(K_0,B_0)^T,
\end{split}
\]
where
\[
 Q=\begin{pmatrix}
 1-\dfrac{d}{2g}c_u^2\eta_v^2&
 \sqrt{\dfrac{d}{2g}}t_u\eta_v\\
 \sqrt{\dfrac{d}{2g}}t_u\eta_v&1
 \end{pmatrix}.
\]
To check the identity, use $Z=(uB-C/v)/D_u$ and complete the two current
squares. Since $d/g=d/c+d/D$ is infinitesimal and $c_u,t_u,\eta_v$
are bounded, $\bar Q=I_2$. Restore every unselected positive event.
All roots separate between fields. Clearing the original pure-testing
homogeneous angle factors gives the same identity at literal infinity:
$\eta_v^2=1+(1/v)^2$, and $c_u,t_u$ are the normalized homogeneous
Alice coordinates. Thus no infinite-angle endpoint is excluded.
\end{proof}

\begin{lemma}[The critical current coefficient]\label{lem:c1-critical-current}
Under the preceding strength assumptions, suppose
\[
 \gamma=\frac{A g}{d^2}\quad\hbox{is bounded},\qquad 4\bar\gamma>\bar c_u^2,
\]
where $A=e_4$ is the original event coefficient. Then every independent
test with $v$ infinitesimal satisfying this residue inequality is tensor-SOS, for every
state, including the homogeneous Alice-infinite endpoint.
\end{lemma}
\begin{proof}
Keep $R,C,B,N,g,h,c_u,t_u$ from the preceding proof, and put
$\epsilon=d/g$, $q=\sqrt\epsilon$ and
\[
 K_1=2\sqrt g\,C+\sqrt{d\epsilon}\,Y,\quad
 B_1=\sqrt d(B-uvY),\quad V_0=\sqrt d\,vY,
 \quad Y_0=\sqrt{d\epsilon}\sqrt{D_u}\,Y.
\]
Define
\[
 L_0=\frac{qc_u}{2}(K_1-c_uY_0)-t_uvB_1+c_uvV_0.
\]
Direct multiplication gives $L_0=\sqrt{dD_u}\,v(vY-Z)$ and the exact
full selected-parent identity
\[
\begin{split}
&4cD_vR^2+4D(R+C)^2+4AD_uD_vY^2+d\mathcal H\\
&\qquad=4(c+D)N^2+4cv^2R^2+2d(R+C)^2+\mathcal G,\\
\mathcal G={}&K_1^2+2D_v(B_1^2+V_0^2)
 +(4\gamma D_v-c_u^2)Y_0^2\\
&+2qv(K_1-c_uY_0)V_0-2L_0^2.
\end{split}
\]
The four-root Gram of $\mathcal G$ has residue
$\operatorname{diag}(1,2,2,4\bar\gamma-\bar c_u^2)$.
Its last entry is positive by the stated residue inequality. Every correction is polynomial in bounded
$q,v,c_u,t_u,\gamma$ and has zero residue. The strict-residue lemma and
restoration of the other events prove the claim. The displayed bounded
homogeneous factors also cover literal Alice infinity.
\end{proof}



\section{Bounded reciprocal contacts and a nonvanishing Alice scale}
For infinitesimal $\eta=f^2,\zeta=g^2$, one contact calculation covers bounded $k$ and bounded $\Delta=1/e$, including infinitesimal $k$. For unit $k$, a source reflection covers every positive ratio $g/f$. The remaining ordinary-Alice sources use the state completions derived below.

\subsection{One contact calculation for bounded $k$ and bounded reciprocal $e$}
\label{sec:c1-bounded-reciprocal}

Use the original C1 source in the normalization $C=2$, with the seven
native coquotients $c_2,c_b,c_3,c_5,c_4,c_\alpha,c_d$ and exact
source/parent SD1--SD2. Suppose $k>0$ is bounded, $f,g>0$ are
infinitesimal, and $\Delta=g/(kf)$ is bounded. Assume the original
positive shapes $a,t$ and positive gaps $M=R-t$, $J=T-a$.
They make $a,t$ bounded because $R,T\le1+k^2$. No bound $a,t<2$
is used; in particular the statement allows a bounded unit $k>1$.
Then the complete parent is tensor SOS on the testing chart with
$u$ bounded, $v$ infinitesimal, $Z$ a unit and all state coordinates
bounded. The proof allows bounded unit $u,W$ and both unit and
infinitesimal $k,\Delta$.

Normalize $Z=1$ in the testing field. Retain exactly the four positive
squares SD3 and use the full change of state coordinates SD5, with
$w=(W+uY)/(1+u^2)$. Put
\[
 K=1+k^2,\quad \delta=\Delta R/K,\quad b^2=\eta K/R,
 \quad Q=K+ak^2b^2,\quad T_n=K^2-t\delta^2k^4Qb^2.
\]
Here $\delta,k,a,t,u,w$ are bounded, $b$ and $v$ are infinitesimal,
and
\[
 \eta={Kb^2\over Q},\quad\Delta={\delta Q\over K},\quad
 \zeta={\delta^2k^2Qb^2\over K},\quad
 R={K^2\over Q},\quad T={T_n\over K}.\tag{SD6}
\]
Every denominator $K,Q,T_n,R,T$ is a positive pure-source unit;
its residue is respectively $K,K,K^2,K,K$.
The testing denominator $D=1+u^2$ is a positive bounded testing unit.
The positive source multiplier $\Delta^2\Lambda$ need not be a unit;
it is restored using pure-source square roots after the construction.

Define the polynomial pullback $C_i=K^2Q^3c_i$ by this substitution.
It is polynomial because a native monomial
$k^jf^{2r}g^{2s}a^\alpha t^\tau$ becomes
\[
 \delta^{2s}k^{j+2s}b^{2(r+s)}a^\alpha t^\tau
                         K^{2+r-s}Q^{3-r+s},\tag{SD7}
\]
and the seven native polynomials have $-2\le r-s\le3$.
Thus no shape inverse or negative exponent is introduced.
Let $\mathcal C=K^6Q^4T_n$ and define
\[
\begin{gathered}
 A=wv,\quad B=\delta v,\quad C_0=\delta b^2uw,
 \quad D_0=\delta^2k^3b^2,\quad D_v=1+v^2,\\
 V_n=K^2-tQuvw,\quad Y_n=KT_nuw-avV_n,\quad
 B_n=vY_n-KT_n,\quad A_c=C_5T_n+aK(C_3D_v-C_5).
\end{gathered}
\]
Put $G=\Delta^2\Lambda F-\sum_i\Pi_i$ exactly as in SD4.
The complete scalar of $\mathcal CG$ at the two state
normals $N=\delta M_B=0$ is the following five-term polynomial:
\[
\begin{split}
 F_0={}&4KQT_n^2[C_2(K^2-tQ)+C_bD_vtQ]D^2A^2
       +4Q^2A_cD^2V_n^2B^2\\
 &+4K^3C_4DD_v(KT_nC_0-ab^2BV_n)^2
       +4KQ^2C_\alpha DB_n^2D_0^2\\
 &+2K^2QC_dD\{b^2B^2(Y_n^2-K^2T_n^2+2aKT_nV_n)
                                     -2K^2T_n^2BC_0\}.
\end{split}
\]
This follows by the state substitution in SD4. In particular it lies
in $(A,B,C_0)^2+(D_0^2)$ before any contact division.
The complete two first-normal coefficients are
\[
\begin{split}
 S_B={}&-8aK^2b^4C_4D_vV_n
       +8\delta^2Q^2k^6b^4C_\alpha B_n+4KQb^2C_dB_n,\\
 S_C={}&8K^3T_nb^2C_4D_v,\\
 L_1={}&A[8K^3QT_n^2(C_bD_v-C_2)D]\\
       &+B[-8\delta K^2Q^2uDvV_nA_c+\delta K^3auvS_B]
                         +C_0[\delta K^3auvS_C],\\
 L_2={}&B[8K^2Q^2T_nDV_n(C_3D_v-C_5)-K^2T_nS_B]
                                      -C_0K^2T_nS_C.
\end{split}\tag{SD8}
\]
They include every first-order state term, and belong to
$(A,B,C_0)$ before contact division.

\paragraph{A single positive-weight estimate.}
Give $(\delta,k,b,a,t,u,w,v)$ weights $(0,0,1,0,0,0,0,2)$.
The zero-weight quantities are bounded; only $b,v$ require positive
weights. Each of $A,B,C_0,D_0$ has weight two. Directly from the
seven source polynomials,
\[
 C_i=2K^6+O_{\rm wt}(2)\ (i=2,b,3,5,4,\alpha),\quad
 C_d=4K^7+O_{\rm wt}(2),\quad Q=K+O_{\rm wt}(2),
\]
\[
 T_n=K^2+O_{\rm wt}(2),\quad V_n=K^2+O_{\rm wt}(2),
 \quad B_n=-K^3+O_{\rm wt}(2),\quad A_c=2K^8+O_{\rm wt}(2).
\]
The five scalar lines therefore have leading terms respectively
\[
 8K^{14}D^2A^2,\quad8K^{14}D^2B^2,\quad
 8K^{15}DC_0^2,\quad8K^{15}DD_0^2,
 \quad-16K^{16}DBC_0.
\]
These are identities modulo the coefficient ideal $(b^2,v)$ for
arbitrary bounded $k,\delta,u,w$, not identities only at contact.
The mixed term in the third line has a factor $b^2$; so does the
separate $B^2$ term in the fifth line. Consequently, with
\[
 P=D(A^2+B^2)+K(C_0^2+D_0^2)-2K^2BC_0,
 \qquad E_0=F_0-8K^2D\mathcal CP,
\]
one has $\operatorname{wt}(E_0)\ge6$.
For the complete normals use the exact differences
\[
 K(C_b-C_2)=\delta^2k^2Qb^2C_3,\qquad
 Q(C_3-C_5)=-k^2Kb^2C_b.
\]
Since $K,Q$ have positive unit constant terms in $b$, these give
weight at least two for both differences; also $S_B,S_C$ have
weight at least two. The displayed full normal expressions imply
$\operatorname{wt}(L_1),\operatorname{wt}(L_2)\ge4$.

The exact native contact is
$(u,w,v,N,M_B)=(k,\delta k,\delta k^2b^2,0,0)$.
The six native balances SD9--SD11 and the contact-zero retained
squares prove $E_0\in\mathfrak J^2$ and $L_i\in\mathfrak J$,
where $\mathfrak J=(u-k,w-\delta k,v-\delta k^2b^2)$.
Apply the restricted contact lift CC1--CC6.
For completeness, its coefficient filtration is particularly simple
here. The three contact substitutions preserve weight. Monic division
by $u-k$ and $w-\delta k$ loses no weight. A division by
$v-\delta k^2b^2$ loses two, but its quotient is always multiplied
by $b^2$ in the CC2 formula for $C_0S$.
Thus that formula preserves coefficient weight. The normal monomial
quotients have weight at least $4-2=2$; the scalar monomial quotients
have weight at least $6-4=2$. The successive CC reductions, including
the gradient subtraction $2k\rho_c(-2k,1,1,-1)$, preserve this bound.
Therefore every final scalar and normal correction coefficient has
weight at least two. It is a bounded tensor coefficient of zero
residue. No determinant inverse is a lifting coefficient.

\paragraph{The complete state-normal block.}
Before multiplication by $\mathcal C$, let
\[
 C_B=H_3D_vh_s+H_5(1-h_s),\quad h_s=a/T,
\]
\[
 B_Y=4RT\Delta^2\eta^2c_4DD_v
 +4RT\Delta^4k^6\eta^2c_\alpha Dv^2
 +2RT\Delta^2\eta c_dDv^2.
\]
The full matrix on $(N,\delta M_B)$ is defined equivalently by
coefficient extraction of SD4 after SD5. Its entries are
\[
\begin{split}
 A_{11}&=H_2+H_bD_v+\Delta^2C_Bu^2v^2
                                      +B_Yh_s^2u^2v^2/D^2,\\
 A_{12}&=-(\Delta^2/\delta)(H_3D_v-H_5)uv
                                      -(B_Y/\delta D^2)h_suv,\\
 A_{22}&=(\Delta^2/\delta^2)(H_3D_v+H_5)+B_Y/(\delta^2D^2).
\end{split}
\]
All negative powers cancel: $\Delta/\delta=Q/K$ is a source unit,
and $B_Y/\delta^2$ is an explicit bounded infinitesimal because each
summand contains $\eta^2$ or $\eta v^2$ after that cancellation.
Also $H_i\sim8K^2$, $C_B$ is bounded and
$H_3D_v-H_5=4R(c_3-c_5)+H_3v^2$ is infinitesimal.
Hence the block after multiplication by $\mathcal C$ has exact
residue $16\bar K^{14}I_2$, for arbitrary bounded $u,w$.

Finally use the six roots
\[
 x=B-C_0,\quad y=A-kC_0,\quad z=uB-kC_0,\quad
 h=kC_0-D_0,\quad q_1=a_1a_2-a_3a_4,\quad
 q_2=a_1a_3-a_2a_4,
\]
where $a_1=\sqrt{wv}$, $a_2=\sqrt{\delta uv}$,
$a_3=b\sqrt{\delta kuw}$ and $a_4=\delta k^{3/2}b$.
Writing $g_4=uwv-\delta^2k^4b^2$, the exact identities are
\[
\begin{split}
 P&=x^2+g_4^2+\sum_i a_i^4-4a_1a_2a_3a_4\\
  &=x^2+g_4^2+y^2+z^2+h^2-yz+zh+q_1^2+q_2^2.
\end{split}
\]
The $(y,z,h)$ matrix has successive principal minors
$1,3/4,1/2$; the other three coefficients are one.
Thus the six-root Gram is strictly positive and constant.
Retain $g_4^2$ separately; its coefficients
need not be infinitesimal when $k,u$ are units.
The full scalar and both normals are thus infinitesimal corrections
of this strict six-root Gram and the two strict state-normal roots.
The positive-residue Gram lemma applies. Restore $g_4^2$ and every
retained square, and undo the positive pure-source multiplier.
Literal axes and a nonnegative testing core have their direct original
event certificates. On the negative-core chart the exact SD core
calculation gives $u,v,w>0$, so all radicals above are separated
nonnegative radicals. This proves the stated local result for every
real state, and separately for both parts of a complex state.


\subsection{Every test in the bounded-reciprocal and ordinary-Alice sources}
\label{sec:c1-uniform-projective-join}
\label{component:p9f00f949fc8c}
\label{component:p64b054e815d0}
\label{component:p18a57c9c271d}
\label{component:p41b39e2f94e1}
\label{p41b39e2f94e1:thm:bounded-delta}
\label{p18a57c9c271d:thm:c1-ordinary-angle}

The preceding local calculation gives the complete projective result
when $k$ is positive bounded, $\eta=f^2,\zeta=g^2$ are infinitesimal
and $\Delta=g/(kf)$ is bounded. It also gives the entire source with
unit $k$ and arbitrary positive $g/f$, after the source reflection
below. We give the complementary testing and state cases explicitly.

\paragraph{Complementary testing angles.}
Let $g_c=CD/(C+D)$. The source coefficients SD1 give, with
$D_1=2c_5/c_3-J\eta c_d/c_3\sim2$,
\[
 {g_c\over a}={2D_1\over 2J+aD_1},\qquad
 2J+aD_1\sim2K,
\]
\[
 d/C,\ d/D\in\operatorname{Inf}_c,\qquad
 \gamma={Ag_c\over d^2}\sim{1\over4K^3}>0.
\]
Lemma~\ref{lem:c1-two-current} therefore applies when $1/v$ is bounded, including
Bob infinity, for every Alice angle. When $v$ is infinitesimal and
Alice is unbounded, Lemma~\ref{lem:c1-critical-current} has
residue diagonal
\[
 \operatorname{diag}(1,2,2,4\bar\gamma-\bar c_u^2),
 \qquad c_u=(1+u^2)^{-1/2}.
\]
Here $\bar c_u=0$ and $\bar\gamma>0$, so the matrix is strict.
Its homogeneous identity includes Alice infinity. Only bounded $u$
and infinitesimal $v$ remain. No upper bound on an event coefficient
was used; the complementary identities retain all original positive
excess squares.

\paragraph{The complementary state chart.}
Normalize a nonzero state by a coordinate of maximal absolute value.
Then all state coordinates are bounded. The testing signed core is a
pure-testing element. If it is nonnegative, it is a testing square
and every event term is already a separated square.

Suppose the core is negative. At $v=0$ it is
$2\{(X+uZ)^2+(uW-Y)^2\}$. Thus both displayed currents have zero
residue. If $Z$ is a unit, use
Section~\ref{sec:c1-bounded-reciprocal}; its coefficient bounds allow
bounded $u,w$ and every bounded $W,X,Y$ after $Z=1$ normalization.
If $Z$ is infinitesimal, then $X$ is infinitesimal and
$Y-uW$ infinitesimal. Some state coordinate is a unit, so $W$ is a
unit and $Q_A=W+uY$ is a unit with residue $(1+\bar u^2)\bar W$.

On this chart the exact five-root current identity uses
\[
 z=(r_2,B,\Delta r,\Delta r_5,\Delta\eta Y).
\]
Retain the four positive source-gap excesses, the whole alpha term
and the positive $v^2Y^2$ part of the core. The remaining diagonal is
\[
 H_2z_1^2+H_bD_vz_2^2+H_3D_vz_3^2+H_5z_4^2
                       +4RTc_4D_uD_vz_5^2,
\]
and its complete negative correction is
\[
 -4RTc_dD_u{\Delta Z\over Q_A}z_5(z_1+z_2)
 -2RTc_dD_u{\Delta^2\eta Z^2\over Q_A^2}(z_1+z_2)^2.
\]
The five positive coefficients are units and both correction
coefficients are infinitesimal because $\Delta$ is bounded and
$Z$ infinitesimal. This proves the state chart. The literal testing
axes have nonnegative core directly. The local and complementary
state arguments are therefore exhaustive.

\paragraph{Ordinary Alice with arbitrary state scale.}
Suppose $k$ is a positive bounded unit and $f,g$ are infinitesimal.
Only $\Delta$ unbounded remains. The exact complete-parent reflection
is
\[
 u'=1/u,\quad v'=v,\quad
 (W',X',Y',Z')=(vW-X,-W-vX,Z-vY,Y+vZ),\tag{AS3}
\]
\[
 u^2F(\operatorname{ref}(c);1/u,v;\psi')
                         =(1+v^2)F(c;u,v;\psi),\tag{AS4}
\]
with event permutation $(\lambda_{\rm ev}\ n)(m\ C)(A\ \alpha)
(D\ \beta)$. Put $H_f=1+(1-a)\eta$, $H_g=1+(1-t)\zeta$
and $V=1+k^2\eta\zeta$. At the actual source contact its native
parameters are
\[
\begin{gathered}
 k'=1/k,\quad
 \eta'=k^2\zeta H_f/H_g,\quad\zeta'=k^2\eta H_g/H_f,\\
 a'=M/(k^2H_f),\quad t'=J/(k^2H_g),\quad
 M'=aV/(k^2H_g),\quad J'=tV/(k^2H_f).
\end{gathered}\tag{AS5}
\]
These identities follow from $H_g-Q_c=\zeta M$ and
$H_f-E=\eta J$ and preserve all original signs and gaps.
Here $H_f,H_g$ are positive units, $k'$ is a unit,
$\eta',\zeta'$ are infinitesimal, and
\[
 \Delta'={\sqrt{\zeta'/\eta'}\over k'}
                         ={H_g\over\Delta H_f}
\]
is infinitesimal. Normalize the reflected coefficients by a positive
pure-source factor to $C'=2$ and apply the bounded-reciprocal result.
The homogeneous reflection restores the original projective endpoints
as well. This proves every ordinary-Alice source with small $f,g$,
without a separate comparable-ratio or shape-unit argument.


% Source formulas: p050c7ae6d9b8
\subsection{The positive Laurent form at the ordinary/small source}

\label{component:p050c7ae6d9b8}

\label{p050c7ae6d9b8:sec:c1-ui-leading}

This is the positive Laurent-model lemma used in the complete UI proof.
Write $x=k^2$, $r=1-a\eta$, $K=1+x$, $R=1+xr$ and
\[
 E_0=1-x\eta,\quad J=K-a,\quad
 P_4=KE_0+x(1-r),\quad P_d=1+2x+rx^2.
\]
At a real interior source point the leading coefficients of the
parent divided by $a$ are
\[
\begin{aligned}
 h&=\frac{2x\eta}{RE_0},& C&=\frac2a,&
 A&=\frac{2\eta^2P_4}{Kr(r+\eta)E_0},\\
 d&=\frac{4\eta P_d}{K(r+\eta)E_0},&
 D&=\frac2{JE_0}-\frac d2,&
 \alpha&=\frac{2x^2\eta R}{KE_0}.
\end{aligned}                                                     \tag{CL1}
\]
All are positive except $D\ge0$, including its literal zero. Here
$\alpha$ is the coefficient after extracting the weak $g^2$ scale.
For a test $(u,V)$ put $D_u=1+u^2$ and define
\begin{align*}
f={}&4hV^2(W+uY)^2+4C(uW-Y)^2+4AD_uY^2\\
 &+(4D+2d)(uW-Y-VD_uZ)^2+4\alpha D_uZ^2\\
 &-4dD_uVYZ-2dD_uV^2Z^2.                           \tag{CL2}
\end{align*}
Let $M_f$ be its symmetric three-state matrix and set $U=u^2$, $T=V^2$.
Direct collection gives
\[
 \det M_f=64D_u^2UT\,\mathcal L(U,T),                 \tag{CL3}
\]
where
\begin{align*}
\mathcal L={}&AC(D+d/2)U+
 \frac{h\alpha(A+C+D+d/2)}U+\frac{A\alpha(C+D+d/2)}T\\
 &+hD(A+C+d/2)\frac TU+hC(D+d/2)UT\\
 &+h(AD+2CD+Ad/2+Cd/2+Dd/2+d^2/4)T+L_0,\\
L_0={}&ACD-ADd/2+h\alpha(C+D+d/2)
                         -d^2(A+C+D)/4-d^3/8.       \tag{CL4}
\end{align*}
Every nonconstant coefficient is nonnegative. The coefficients of
$U,U^{-1},T,T^{-1}$ are strictly positive, also when $D=0$.
Thus in coordinates $(\log U,\log T)$ the Hessian is strictly
positive and the function is coercive. Substitution of (CL1) gives
\[
 \mathcal L(x,x\eta)=\partial_U\mathcal L(x,x\eta)
                   =\partial_T\mathcal L(x,x\eta)=0. \tag{CL5}
\]
Strict convexity proves that $\mathcal L$ is positive away from this
single contact. The $(W,Y)$ principal block of (CL2) is strictly
positive for $U,T>0$, because its $C$ and $A$ terms already have
trivial common kernel. Hence $f$ is positive semidefinite everywhere
and positive definite off the contact locus. Literal endpoints follow
by continuity. This proof uses no numerical determinant table.

There is a useful corrected formula at the large-$D$ boundary.
Impose its hard equation $uW-Y-VD_uZ=0$, eliminating $W$.
The resulting two-state determinant is $16D_u^2T\mathcal L_D$, where
\begin{align*}
\mathcal L_D={}&hT\{(C+A+d/2)/U+(2C+A+d/2)+CU\}\\
 &+h\alpha(1+1/U)+AC(U+1)+A\alpha/T-Ad/2-d^2/4.
                                                               \tag{CL6}
\end{align*}
In particular the coefficients of $T/U$ and $UT$ are
$h(C+A+d/2)$ and $hC$, respectively. Interchanging these coefficients
is incorrect. All three identities follow by polynomial collection in the displayed model.



% Source formulas: pa0b854a04f49
\subsection{The ordinary/small source: the complete state quadratic}

\label{component:pa0b854a04f49}

\label{pa0b854a04f49:sec:r6-c1-ui}
Use the unchanged native C1 source in the representative of (NI7):
$k,\eta$ are positive bounded units and $\zeta=g^2$ is positive
infinitesimal. All original openings and $D=0$ are retained.
The exact stabilizer recovers the other representative; no condition
$k\le1$ is imposed after that change. This section proves the missing
ordinary/small testing strip, and joins the remote theorem (CR1)--(CR13).

\paragraph{Exact bounded remainder.}
Normalize by the positive pure-source $a$. Write $x=k^2$,
$r=1-a\eta$, $K=1+x$, $R=1+xr$, $M=R-t$, and
\[
 J_s=K-xtg^2-a,\quad E=1-x\eta+tx\eta g^2,\quad
 Q_c=1-xg^2r,\quad y=x\eta g^2.
\]
Here $a,t,r,1-r,M,J_s$ are positive bounded openings; $E,Q_c,R,K,r+\eta$
are positive units. The exact polynomials $P_4,P_d,P_\alpha$ are (NC2).
The coefficients after division by $a$ are
\[
\begin{aligned}
 c&=2/a,& A_0&=\frac{2\eta^2P_4}{Kr(r+\eta)E},&
 d_0&=\frac{4\eta P_d}{K(r+\eta)E},\\
 D_0&=\frac{2(1+y)}{J_sE}-d_0/2,&
 \alpha_0&=\frac{2x^2\eta P_\alpha}{KQ_cE},&
 \alpha/a&=g^2\alpha_0.
\end{aligned}                                                     \tag{CU1}
\]
In particular $c,A_0$ are strong; $d_0,\alpha_0$ are units; and
$ca,A_0r,D_0J_s$ are bounded with continuous rational expressions
through every opening. Put
\[
 h_2=\frac{2x\eta(1-tg^2)(1+y)}E,\qquad
 h_b=\frac{2x\eta[1+(1-t)g^2]}E,
 \qquad h_b-h_2=2x\eta g^2.
\]
Both $h_2,h_b$ are positive units. The exact coefficients are
$m/a=h_2/(g^2M)$ and $\beta/a+d_0/2=h_b/(g^2t)$.

For the test put $D_u=1+u^2,D_v=1+v^2$ and
\[
 Q=W+uY,\quad B=X+uZ,\quad p=t/R,
 \quad N=(B-pvQ)/g,\quad L=vQ/g.
\]
The inverse $W=Q-uY$, $X=gN+pvQ-uZ$ is polynomial.
Retain the two positive squares
\[
 \frac{4h_2t}{MR}N^2,\qquad
 \frac{4h_bM}{tR}D_vN^2.                            \tag{CU2}
\]
Their remainder is exactly (CL7). Choose bounded caps
$c_*\le c$, $A_*\le A_0$, $D_*\le D_0$, with $c_*,A_*>0$,
$D_*\ge0$; bounded original weights can be kept unchanged.
Retain the three positive centered excesses
\[
\begin{aligned}
 &4(c-c_*)D_v(r_t-avZ)^2,\quad
 4(A_0-A_*)D_uD_v(Y-rvZ/\eta)^2,\\
 &4(D_0-D_*)(J_t+J_svZ)^2,
\end{aligned}                                                     \tag{CU3}
\]
where
\[
 r_t=uQ-D_uY,\quad J_t=r_t-D_uvZ+ugvN+upv^2Q.
\]
All five retained rows vanish at the actual contact
\[
 (u,v)=(k,k\sqrt\eta g),\quad
 (W,X,Y,Z)=(g/\sqrt\eta,k(tg^2-1),krg/\sqrt\eta,1).
                                                               \tag{CU4}
\]
Set $c_a=2-c_*a$, $A_r=A_0r-A_*r$, $D_J=D_0J_s-D_*J_s$.
After removing (CU2)--(CU3), the entire parent divided by $a$ is
\begin{align*}
G={}&\frac4R\{(h_2+h_bD_v)N^2+2(h_bD_v-h_2)LN
              +[h_2(1-p)+h_bpD_v]L^2\}\\
 &+4c_*D_vr_t^2+8c_aD_vvZr_t-4c_aaD_vv^2Z^2\\
 &+4A_*D_uD_vY^2+\frac{8A_r}{\eta}D_uD_vvZY
                   -\frac{4A_rr}{\eta^2}D_uD_vv^2Z^2\\
 &+4D_*J_t^2-8D_JvZJ_t-4D_JJ_sv^2Z^2\\
 &+4\alpha_0D_u(gvY-gZ)^2
       +2d_0\{J_t^2+D_u(v^2Y^2-2vYZ-v^2Z^2)\}.
                                                               \tag{CU5}
\end{align*}
Thus $G$ is an exact bounded Gram on $(N,Q,Y,vZ,gZ,L)$ whenever
$u$ is a bounded unit and $v$ infinitesimal. Every source denominator
in its coefficients is a unit. Formula (CU5), rather than a scalar
restriction, supplies all normal terms below.

\paragraph{Local contact lifting with all first normals.}
Put $\kappa=k\sqrt\eta$, $\varepsilon=\kappa g$,
$s=R/(k\eta)$, $\delta=u-k$, and define
\[
 E_Q=Q-svZ,\quad E_Y=Y-rvZ/\eta,\quad
 R_h=(vQ-s\varepsilon^2Z)/\varepsilon,
 \quad S=v\delta Z,\quad T=\varepsilon\delta Z.
\]
Let $c_L=4[h_2(1-p)+h_bpD_v]/R$, a positive bounded unit.
The sole potentially unbounded $L$ row can be replaced exactly:
\[
 L=\kappa R_h+kR gZ,\qquad
 L^2-\kappa^2R_h^2=2kR vQZ-k^2R^2g^2Z^2.           \tag{CU6}
\]
Consequently (CU5) equals $c_L\kappa^2R_h^2$
plus $8\kappa(h_bD_v-h_2)NR_h/R$ and a polynomial $G_0$.
Explicitly, in (CU5) replace its first line by
\[
 \frac4R(h_2+h_bD_v)N^2+8k(h_bD_v-h_2)gNZ
                 +c_L(2kRvQZ-k^2R^2g^2Z^2).        \tag{CU7}
\]
All coefficients of $G_0$ are bounded polynomials in $u,v,g$ with
pure-source unit denominators. The coefficient of $NR_h$ is
infinitesimal, since $h_bD_v-h_2=2x\eta g^2+h_bv^2$.

Here are exact prescriptions for every normal coefficient, avoiding
any implicit omission. Substitute $Q=E_Q+svZ$ and
$Y=E_Y+rvZ/\eta$ in $G_0$ and collect the quadratic polynomial
in $(N,E_Q,E_Y)$. The constant coefficient is $\phi(u,\xi)Z^2$,
where $\xi=v^2$ and $\phi\in(g^2,\xi)$. The linear coefficients
of $E_Q,E_Y$ are $vL_Q(u,\xi)Z,vL_Y(u,\xi)Z$; that of $N$ is
$gL_N(u,\xi)Z$, where $L_N\in(g^2,\xi)$. These parity and ideal
assertions follow term by term in (CU5),(CU7). The remaining quadratic
normal block is the full coefficient collection of the same expression.
For explicit pen-and-paper formulas, set
\[
 b=r/\eta,\quad F_0=us-D_ub,\quad
 G_0^{\rm sc}=F_0-D_u+up\xi s,\quad
 H_0=(8D_*+4d_0)G_0^{\rm sc}-8D_J.
\]
Then, writing $D_v=1+\xi$ in these polynomials,
\begin{align*}
 L_Q={}&2c_LkR+8D_vu(c_*F_0+c_a)+u(1+p\xi)H_0,\\
 L_Y={}&-8D_vD_u(c_*F_0+c_a)
       +8D_uD_v(A_*b+A_r/\eta)-D_uH_0\\
      &+8\alpha_0D_ug^2(b\xi-1)+4d_0D_u(b\xi-1),\\
 L_N={}&8k(h_bD_v-h_2)+u\xi H_0.
\end{align*}
The complete scalar coefficient is
\begin{align*}
 \phi={}&c_L(2kRs\xi-k^2R^2g^2)
   +4c_*D_vF_0^2\xi+8c_aD_vF_0\xi-4c_aaD_v\xi\\
 &+4A_*D_uD_vb^2\xi+8(A_r/\eta)D_uD_vb\xi
                    -4A_rrD_uD_v\xi/\eta^2\\
 &+4D_*(G_0^{\rm sc})^2\xi-8D_JG_0^{\rm sc}\xi-4D_JJ_s\xi\\
 &+4\alpha_0D_ug^2(b\xi-1)^2
   +2d_0\{(G_0^{\rm sc})^2\xi+D_u(b^2\xi^2-2b\xi-\xi)\}.
\end{align*}
These formulas display the two asserted ideals directly. They also
give a short explicit input for the monic divisions, without an
unprinted state-normal computation.

The source is an exact stationary zero and every removed square
vanishes there. Therefore
\[
 \phi=\phi_u=\phi_\xi=0,\qquad L_Q=L_Y=L_N=0
                       \quad\text{at }(u,\xi)=(k,\varepsilon^2).
\]
Successive monic polynomial division gives, with bounded coefficients,
\[
 \phi=\delta^2(\varepsilon^2B_0+\xi B_1)
       +\delta(\xi-\varepsilon^2)B_2
       +(\xi-\varepsilon^2)^2B_3,                  \tag{CU8}
\]
\[
 L_i=\delta L_{i0}+(\xi-\varepsilon^2)L_{i1} (i=Q,Y),\quad
 L_N=\delta(g^2L_{N0}+\xi L_{N1})
                         +(\xi-\varepsilon^2)L_{N2}.
\]
For clarity, each quotient is obtained by subtracting the value at
$u=k$ and dividing by the monic $u-k$, then by the monic
$\xi-\varepsilon^2$. The second $u$ quotient of $\phi$ stays in
$(g^2,\xi)$; its constant and first $u$ terms have the stated double
and simple zeros in $\xi$. The same argument gives the $L_N$ ideal.
No small source number is inverted in these divisions.

Put $R_s=(v^2-\varepsilon^2)Z/\varepsilon$. The exact root identities
\[
 \varepsilon R_s=(\varepsilon R_h-vE_Q)/s,\quad
 R_sT=(R_hT-SE_Q)/s,\quad
 S^2-T^2=\delta(R_hT-SE_Q)/s                       \tag{CU9}
\]
allocate (CU8) and every first normal on the six roots
$(N,E_Q,E_Y,R_h,S,T)$. For example $v\delta Z=S$,
$\delta(\xi-\varepsilon^2)Z^2=(R_hT-SE_Q)/s$, and
$(\xi-\varepsilon^2)^2Z^2=(\varepsilon R_h-vE_Q)^2/s^2$.
The $N$ normal has only infinitesimal couplings. All quadratic normal
coefficients are already bounded. This is the required exact full lift.

\paragraph{Strict matched Hessian and source boundary caps.}
First take finite interior source residues. The leading form is
(CL2). Put $V=k\sqrt\eta V_1$, $W=W_1/\sqrt\eta$,
$Y=kY_1/\sqrt\eta$, $Z=1$. At the contact
$(W_1,Y_1,V_1,u)=(1,r,1,k)$, the Schur complement of the two state
coordinates in half the Hessian, in coordinates $(V_1,u)$, is
\[
\begin{aligned}
 G_{11}&=\frac{32\eta x^2P_4}{KE_0},&
 G_{12}&=\frac{16\eta k^3(R^2-2KE_0)}{KE_0R},\\
 G_{22}&=\frac{32\eta x(P_4+xrP_d)}{R^2P_4},&
 \det G&=\frac{256\eta^2x^3(4K^2E_0-R^2)}{K^2E_0^2}.
\end{aligned}                                                     \tag{CU10}
\]
Here $E_0=1-x\eta$, $P_4=KE_0+x(1-r)$, $P_d=1+2x+rx^2$.
There are no opening denominators. The leading $D\ge0$ condition is
\[
 K(2+xr)E_0-R^2\ge0,
\]
so $4K^2E_0-R^2\ge K(4K-2-xr)E_0>0$.
The state block is uniformly bounded below since $c,A_0$ are strong.
Thus the matched Hessian is strict, uniformly through every permitted
source boundary after the appropriate hard directions are retained.

We give the boundary justification explicitly. Divergences can occur
only at $a=0$, $r=0$, or $J_s=0$, for $c,A_0,D_0$, respectively.
Since $a+J_s\sim K$, $c$ and $D_0$ cannot both diverge. If $a=0$,
then $r=1$, so $c$ is the only divergent weight. The only possible
divergent pair is $A_0,D_0$. Their hard-state rows are $E_Y$ and
$kE_Q-KE_Y$, with nonzero determinant $k$. All angular parts of
these centered rows are bounded. Solving these at most two real
linear equations therefore supplies a bounded competitor for the
minimizing state vector. The bounded finite part and (CU10) imply
that the state minimizer is bounded as well, by completing its
positive state square. The full matched Hessian hence has a uniform
positive lower bound on the Euclidean norm in a nearby finite-source
family. On a vector in the limiting hard kernel, adjust only the two
state coordinates to annihilate the nearby exact hard rows. This
adjustment tends to zero, and the divergent penalties vanish exactly.
The bounded finite part is consequently strictly positive on that
kernel. The elementary real finite-penalty lemma supplies finite caps.

Every allowed residue is approachable by such finite sources. For
fixed $x>0$, the real leading inequalities are
\[
 0\le r\le1,\qquad
 \eta_{\min}(r)=\frac{1-r}{K}\le\eta\le
 \eta_{\max}(r)=\frac{1-R^2/[K(2+xr)]}{x},
\]
with positive interval width $R/[xK(2+xr)]$. Perturb $r$ into
$(0,1)$ and its interval coordinate into $(0,1)$.
This makes $a,r,J$ and the $D$ numerator strictly positive while
converging to the given residue. Taking $t=R/2$ and then sufficiently
small real $g>0$ restores all original source signs, including
$\beta>0$. This construction, rather than substitution into a
singular leading denominator, justifies the boundary limit.

In the six-root residue Gram, $N$ is an isolated positive pivot.
The $(E_Q,E_Y)$ block is strict, $R_h$ has positive pivot
$c_L\kappa^2$, $S$ couples only to $E_Q,E_Y$, and $T$ only to $R_h$.
Setting $S=T$ is the matched Hessian under an invertible real change
of coordinates: $R_h/g$ has relative-angle derivative $2s\kappa\ne0$,
and $T/g=\kappa\delta$ supplies the ordinary-angle derivative.
After the hard Schur complements, the sum of the $S,T$ diagonals is
therefore strictly positive. Add a real multiple of the last zero
in (CU9) to redistribute this sum into two positive diagonals.
The extra coefficients contain $\delta$ and vanish in the residue.
The bounded strict-residue Gram lemma proves the full parent at
$u-k$ and $v$ infinitesimal, including every source opening.

\paragraph{Ordinary angles away from contact.}
Let now $\bar u^2\ne\bar k^2$, with $u$ a unit and $v$ infinitesimal.
Use the bounded frame $(N,r_t,Y,vZ,gZ,L)$ of (CU5), and the exact zero
\[
 L(gZ)-Q(vZ)=0,\qquad Q=(r_t+D_uY)/u.              \tag{CU11}
\]
For finite leading source weights, $N$ is an isolated positive pivot.
The remaining residue is
\[
 q^TPq+2Ub^Tq+c_2U^2+4hL^2+c_0T^2,
 \quad q=(r_t,Y)^T,\ U=vZ,\ T=gZ,
\]
where the source coefficients are (CL1) and
\[
\begin{gathered}
 s_D=4D+2d,\quad s_C=4C+s_D,\quad
 P=\operatorname{diag}(s_C,4AD_u),\\
 \ell=(1/u,D_u/u)^T,\quad b=(-s_DD_u,-2dD_u)^T,\\
 c_0=4\alpha D_u,\qquad c_2=s_DD_u^2-2dD_u.
\end{gathered}                                                     \tag{CU12}
\]
Apply the real rank-one allocation lemma (UA1)--(UA3), using (CU11).
We verify its hypotheses uniformly through all source boundaries.
Writing $d_a=\ell^TP^{-1}\ell$, $e_a=\ell^TP^{-1}b$,
$h_a=c_2-b^TP^{-1}b$, direct inversion yields
\[
\begin{aligned}
 d_a&=\frac{aJE_0}{8u^2(JE_0+a)}+\frac{D_u}{4Au^2},\\
 e_a&=-\frac{aD_u}{u(JE_0+a)}-\frac{dD_u}{2Au},\\
 h_a&=\frac{8D_u^2}{JE_0+a}-2dD_u-\frac{d^2D_u}{A},\\
 P^{-1}b&=(-aD_u/(JE_0+a),-d/(2A))^T.
\end{aligned}                                                     \tag{CU13}
\]
All are bounded continuous, since $JE_0+a$ is a positive unit and
$1/A$ has a continuous bounded formula proportional to $r$.
The same holds for $P^{-1}\ell$.

Normalize the Laurent polynomial (CL4) by $As_C$.
Its $U$ and $T^{-1}$ coefficients are, respectively,
$1/[2(JE_0+a)]$ and $\alpha/4$, positive units.
Its $T$ coefficient is
\[
 h\{(D+d/2)/s_C+2CD/(As_C)+d/(8A)\},
\]
also a positive unit: when $1/A=0$, $r=0$ forces $a$ a unit,
so $(D+d/2)/s_C=a/[4(JE_0+a)]$ is a unit; otherwise the last term
is positive. The $U^{-1}$ coefficient is
$h\alpha(1/s_C+1/(4A))$. If it vanishes, $r=0$ and $J=0$,
and the $T/U$ coefficient has the positive term $hD/s_C=h/4$.
All remaining normalized nonconstant coefficients are bounded and
nonnegative. The constant is bounded by the exact contact equation.
Thus strict convexity and coercivity in logarithmic coordinates
persist at every boundary. The normalized Laurent function retains
its unique stationary zero $(U,T)=(x,x\eta)$.

For fixed $u^2\ne x$, the rank-one determinant polynomial is
\[
 p(T)=c_0+(h_a+4hd_ac_0)T+4h(d_ah_a+e_a^2)T^2
      =\frac{16D_u}{As_C}T\mathcal L(u^2,T)>0.
\]
Its constant and leading coefficients are positive units by the
preceding coefficient bounds. The strict allocation interval
therefore survives at every source boundary. The two real Gram
blocks have a uniform positive lower bound; (CU13) controls their
state minimizers even if $P$ is unbounded. The centered divergent
rows have hard parts $r_t$, $Y$, $r_t$; the only simultaneous pair
is $Y,r_t$, whose determinant is one. Exact small state corrections
annihilate the large rows in the nearby finite family. The same
finite-penalty argument as above supplies caps in (CU3). The exact
bounded remainder (CU5), after (CU11), has strict residue.
This proves every ordinary/small test away from contact.

\paragraph{Exhaustive certificate.}
The two arguments cover the whole ordinary/small testing strip;
the exact gauges cover negative contact signs. The remote theorem
(CR1)--(CR13), including CR5a, covers its complement.
All transformations are homogeneous in the state; no state coordinate
was divided out in the tensor proof. Real and imaginary state parts
are treated separately. Hence the entire actual UI subdomain of C1.5,
and its exact stabilizer images, has an all-tests tensor-SOS proof.
The actual UH subdomain of C1.5 remains a separate obligation.



% Source formulas: p013c226a79b3
\subsection{The ordinary/small source: complementary testing charts}

\label{component:p013c226a79b3}

\label{p013c226a79b3:sec:c1-ui-remote}
The source has actual ordinary Alice angle and infinitesimal Bob
angle. Apply the exact native stabilizer (AS5) as in (NI7) and take
$\eta$ a positive bounded unit, $\zeta=g^2$ infinitesimal and $k$ a
bounded unit. The reflected orientation need not satisfy $k\le1$.
Its $E$ remains a positive unit by the exact source entry in (NI7).
Normalize by the positive pure-source $d$.
Equations (NC2)--(NC4) give
\[
 C/d,\ A/d,\quad g^2m/d,\quad g^2\beta/d,\quad
 \alpha/(dg^2)\quad\hbox{strong},\qquad D/d\ge0.       \tag{CR1}
\]
For example $P_d,P_\alpha,Q,r+\eta$ are units, whereas the
positive openings $M,t,r,1-r$ are bounded and may be infinitesimal.
The $-1/2$ in $\beta/d$ disappears after multiplication by $g^2$.
In this particular source subdomain $\beta=0$ is impossible; $D=0$
is allowed throughout. Choose positive real constants
$\mu,\nu,c,a,b$ below the five quantities in (CR1), respectively
in the order $g^2m/d,g^2\beta/d,C/d,A/d,\alpha/(dg^2)$.
Whenever needed we cap these coefficients at
\[
 m=\mu/g^2,\quad\beta=\nu/g^2,\quad C=c,\quad A=a,
 \quad\alpha=bg^2.                                  \tag{CR2}
\]
Every unused term, every positive excess, and the complete $D$
term is retained as a square. The signed core, divided by two, is
\[
 K=J^2+D_u(vY-Z)^2+D_v(X+uZ)^2-D_uD_vZ^2,
 \quad J=uW+uvX-Y-vZ.                              \tag{CR3}
\]
Only the core and the explicitly selected caps need a bounded Gram.
This observation keeps all upper coefficient scales out of the proof.

\begin{theorem}[C1 UI remote tests]
For every source satisfying (CR1), the complete original C1 parent
is tensor SOS at every independent testing pair except possibly
$u$ a bounded nonzero unit and $v$ infinitesimal. The conclusion
includes literal projective endpoints, both signs, all states,
and $D=0$.
\end{theorem}
\begin{proof}
All constants in residue forms below are real residues. Roots may be
unbounded, but all coefficients in their exact Gram are bounded.
First put $Q=W+uY$, $B=X+uZ$; then
\[
 M=vQ-B,\quad uW-Y=uQ-D_uY,\quad
 J=uQ-D_uY+uvB-vD_uZ.                               \tag{CR4}
\]
These identities account for every state normal.

If $u$ is a bounded unit and $v$ projectively infinite, put $w=1/v$
and homogenize by $w^2$. Use roots
$(E_Q,E_B,Y,Z)=(Q/g,B/g,Y,Z)$ and cap only $m,\beta$.
The exact homogenized core is
\[
 [w(ugE_Q-D_uY)+ugE_B-D_uZ]^2
 +D_u(Y-wZ)^2+D_wg^2E_B^2-D_uD_wZ^2.
\]
Its residue is $D_u(Y^2+u^2Z^2)$. Including the two caps gives
\[
 4\mu E_Q^2+4\nu E_B^2+2D_u(Y^2+u^2Z^2)>0.         \tag{CR5}
\]

If $u$ is infinitesimal and $v$ projectively infinite, the preceding
four-root frame loses the $Z$ pivot. Keep $w=1/v$ and instead use
$(E_Q,E_B,Y,U,T)$, where $U=uZ$ and $T=gwZ$. Put
$A_0=w(ugE_Q-D_uY)+ugE_B$. The exact homogenized core is
\[
 K_w=A_0^2+D_uY^2+D_wg^2E_B^2+D_uU^2
       -2D_ug(wE_Q+E_B)U+2wD_uuYU.
\]
All its coefficients are bounded, with residue $Y^2+U^2$.
Cap $m,\beta,\alpha$; their amplitudes after normalization are
$E_Q-wE_B$, $E_B$, and $gY-T$, with the positive factors
$D_w,D_u$ in the last two terms. The full selected residue is
\[
 4\mu E_Q^2+4\nu E_B^2+2Y^2+2U^2+4bT^2>0.       \tag{CR5a}
\]
This includes literal $u=0$ or $w=0$; no division by either was used.

If $u$ is infinitesimal and $v$ a bounded unit, use
$E_Q=Q/g,E_B=B/g,Y,U=uZ,Z_g=gZ$, and cap $m,\beta,\alpha$.
Substitute $W=gE_Q-uY$, $X=gE_B-uZ$ in (CR3).
All terms involving $Z$ collect into $uZ$; the resulting exact
core is a polynomial in $(u,v,g;E_Q,E_B,Y,U)$.
Its residue is $D_vY^2+v^2U^2$. The full selected residue is
\[
 4\mu(vE_Q-E_B)^2+4\nu D_vE_B^2
       +2D_vY^2+2v^2U^2+4bZ_g^2>0.                 \tag{CR6}
\]
The $\alpha$ amplitude is $gvY-Z_g$, so no inverse of $u$ was used
in constructing this Gram.

If both $u,v$ are infinitesimal, set
\[
 U=uQ,\quad V=vQ/g,\quad E_B=B/g,\quad Z_g=gZ,
 \quad S=uvZ,
\]
and include $Y$ as the sixth root. The exact core is
\begin{align*}
K={}&(U-D_uY+uvgE_B)^2+D_uv^2Y^2+D_vg^2E_B^2+D_uS^2\\
 &+2D_uuYS-2D_uuVZ_g-2D_uvgE_BS.                 \tag{CR7}
\end{align*}
In particular the apparently problematic term $UvZ$ is exactly
$uVZ_g$, with no mixed inverse. Cap all five weights in (CR2).
The entire selected residue is
\[
 4\mu(V-E_B)^2+4\nu E_B^2+(4c+2)(U-Y)^2
          +4aY^2+4bZ_g^2+2S^2>0.                  \tag{CR8}
\]

It remains to cover $u$ projectively infinite. Put $p=1/u$,
homogenize by $p^2$, and use $D_p=1+p^2$,
$Q=pW+Y$, $B=pX+Z$. Then
\[
 M=vQ-B,\quad r=W-pY,\quad
 J=D_p(W+vX)-p(Q+vB).                              \tag{CR9}
\]
For $v$ a bounded unit, use $(Q/g,B/g,W,X)$ and cap $m,\beta,C$.
The core after $Y=gE_Q-pW$, $Z=gE_B-pX$ is a bounded polynomial
with residue $(W+vX)^2$. The selected residue is
\[
 4\mu(vE_Q-E_B)^2+4\nu D_vE_B^2
           +4cD_vW^2+2(W+vX)^2>0.                 \tag{CR10}
\]

For infinitesimal $v$, use $(W,R,Y,V,E_B)$ with
$R=vX$, $V=vQ/g$, $E_B=B/g$. Cap $m,\beta,C,A$, retaining
$\alpha$ entirely. The exact core is
\begin{align*}
K_p={}&(W+D_pR-pY-pvgE_B)^2+D_vg^2E_B^2\\
 &+D_p\{v^2Y^2-2vgYE_B+2pYR-(vgE_B-pR)^2\}.
                                                               \tag{CR11}
\end{align*}
Every coefficient is bounded. The selected residue is
\[
 4\mu(V-E_B)^2+4\nu E_B^2+4cW^2+4aY^2
                    +2(W+R)^2>0.                 \tag{CR12}
\]

Finally, when both testing angles are infinite, also put $w=1/v$
and homogenize by $p^2w^2$. Use $(Q/g,B/g,W,X)$ and cap
$m,\beta,C$. Substitution $Y=gE_Q-pW$, $Z=gE_B-pX$ in the
homogenized core is polynomial and has residue $X^2$.
The selected residue is
\[
 4\mu E_Q^2+4\nu E_B^2+4cW^2+2X^2>0.               \tag{CR13}
\]
Each displayed positive form is strict on the complete root space.
The bounded positive-residue Gram lemma therefore supplies tensor
SOS. Restoring the retained terms proves the original parent.
The bounded unit/unit testing chart is already covered by the
common noncommuting-residue theorem. This leaves exactly the
ordinary/small testing strip stated in the theorem. The exact gauges
$(u,W,X)\mapsto(-u,-W,-X)$ and
$(v,W,Y)\mapsto(-v,-W,-Y)$ preserve every event square and the core.
All homogeneous formulas allow $p=0$ or $w=0$ literally.
Real and imaginary state components are treated separately.
\end{proof}

The ordinary/small testing strip is completed in Section~\ref{component:pa0b854a04f49}. Together with the remote charts proved here, it covers the entire actual UI source. The actual UH source is treated by the separate complete conductor in Section~\ref{component:pf7ffa921793c}.



% Source formulas: pf7ffa921793c
\subsection{The ordinary/unbounded source}

\label{component:pf7ffa921793c}

\label{pf7ffa921793c:sec:r7-c1-uh}
All coefficients below are divided by the original positive pure-source
$d$. We treat precisely (NI9) and (NI10), with no alteration of the native
source (NC1)--(NC7). The revision-6 common-event theorem already covers
all tests except positive unit $u$ and $w=1/v$ infinitesimal.
Put $x=k^2$, $K=1+x$, $\sigma=1/\eta$, $z=1/\zeta=q^2$,
$b=r/z$, $H=h/\sigma$, and retain $\sigma>0$ (it need not be a unit).
On (NI9), $\sigma$ is infinitesimal and $z$ is a unit; on (NI10),
$z$ is infinitesimal and $\sigma$ is bounded. Throughout, $k,H$ are
positive units, $b,z,\sigma$ are bounded, $Q=1-xb>0$, and
\[
 r=zb,\quad R=1+xzb,\quad M=R-z+\sigma zH,
 \quad J=1+\sigma(xH-1+zb),\quad\theta=1-\sigma H,
 \quad E_s=1-xH.
\]
The exact source polynomials, all with bounded arguments, are
\begin{align*}
 P={}&xK+\sigma x+\sigma z+\sigma zb x^2
 +x^2\sigma zH(1-\sigma H)+2x\sigma^2zH
 +\sigma^2z^2bx^2H,\\
 V={}&z(1+x\sigma H)E_s+x(1-zb)(1+\sigma zH),\\
 B_\alpha={}&R\sigma zH+RxH+ME_s.
\end{align*}
Writing $m=m_{\rm old}/z$, $c=C/z$, $a=Az$, $\ell=\alpha z$,
\[
 c={K(1+\sigma zb)E_s\over2(1-zb)P},\quad
 m={Kx(1+\sigma zb)H(\sigma z+x)\over2MP},\quad
 a={V\over2bP},\quad
 \ell={x^2(1+\sigma zb)B_\alpha\over2QP}.                 \tag{UH1}
\]
Define the polynomial quotients
\[
 N_D={K(\sigma z+x)(1+\sigma zb)-JP\over\sigma},\qquad
 N_\beta={Kx(1+\sigma zb)(1+\sigma zH)-\theta P\over\sigma}.
\]
Then $D=\sigma N_D/(2JP)$ and
$\beta=\sigma N_\beta/(2\theta P)$, exactly.

\paragraph{The bounded full parent.}
Let $D_u=1+u^2$, $D_w=1+w^2$. Its six event amplitudes are
\[
 M_w=W-wX+uY-uwZ,\ R_w=uW-Y,\ Y,\
 J_w=uwW+uX-wY-Z,\ Y-wZ,\ B_w=X+uZ.
\]
Thus the exact parent is
\[
 F=4mzM_w^2+4czD_wR_w^2+4(a/z)D_uD_wY^2
 +4DJ_w^2+4(\ell/z)D_u(Y-wZ)^2+4\beta D_wB_w^2+2K_w, \tag{UH2}
\]
where, without any restriction to a state ray,
\begin{align*}
 K_w={}&u^2w^2Z^2-2u^2wYZ-2uwWZ+2uw^2XZ
       +u^2w^2W^2+2u^2wWX-2uw^2WY\\
      &+(D_u+w^2)(X^2+Y^2)-2uwXY.
\end{align*}
For a potentially unbounded weight choose a finite cap, keeping bounded
weights unchanged. Subtract and retain the following nonnegative
centered excesses:
\begin{align*}
 &4(m-m_*)[qM_w-kM(wZ/q)]^2,\\
 &4(c-c_*)D_w[qR_w-x(1-r)(wZ/q)]^2,\\
 &4(a-a_*)D_uD_w[(Y-xbwZ)/q]^2,
 \quad4(\ell-\ell_*)D_u[(Y-xbwZ)/q]^2.          \tag{UH3}
\end{align*}
The first line uses $qM_w$ and the second $qR_w$; thus these are
exactly excesses of the scaled weights in (UH1). Their remainders
are bounded because $mM,c(1-r),ab,\ell Q$ are bounded.
When $\sigma$ is a unit, additionally retain
\[
 4(D-D_*)(J_w+JZ)^2,
 \qquad4(\beta-\beta_*)D_w(X+uhZ)^2.               \tag{UH4}
\]
Here $DJ$ and $\beta\theta$ are bounded. When $\sigma$ is
infinitesimal, no (UH4) cap is needed: $D,\beta=O(\sigma^2)$,
as proved below. Denote the complete remainder by $G$.
Every retained row vanishes at the actual contact
\[
 u=k,\quad w={q\sqrt\sigma\over k},\qquad
 (W,X,Y,Z)=({\sqrt\sigma\over q},-kh,kbq\sqrt\sigma,1).
                                                               \tag{UH5}
\]
No comparison between independent source and test rates occurs.

\paragraph{Exact homogeneous conductor, including every first normal.}
Put $h=\sigma H$, $\delta=u-k$, and use the six roots
\[
\begin{gathered}
 N_x=X+khZ,\quad N_y=(Y-kbq^2W)/q,\\
 U=wW-\sigma Z/k,\quad N=(wZ-q^2W/k)/q,\\
 S=q\delta W,\quad T=\sigma\delta Z.
\end{gathered}                                     \tag{UH6}
\]
Substitute $X=N_x-khZ$, $Y=qN_y+kbq^2W$ in the entire
bounded remainder $G$. Its complete quadratic normal block is the
coefficient collection in $(N_x,N_y)$. Its scalar and first-normal
coefficients have exactly the following supports:
\begin{align*}
 G_{00}&={A\over q^2}w^2Z^2+B w^2W^2+C\sigma^2 Z^2
                  +D_1 wWZ+E q^2W^2,\\
 G_{10}&=A_xwW+B_xZ+C_xw^2Z,\\
 G_{01}&=q(A_y+B_yw^2)W+C_ywZ/q.                    \tag{UH7}
\end{align*}
All eleven displayed coefficient families and the quadratic normal
coefficients are bounded polynomials in $u$, with bounded source
coefficients; $B_x$ is divisible by $\sigma$. In the unit-$\sigma$
case divisibility means a bounded quotient, not a new polynomial
identity. Formulas below specify every coefficient before caps, and
(UH3)--(UH4) specify their exact changes. No state normal is omitted.

Set $s=\sigma/k$, $a_0=q^2/k$ and
\[
 L=2A/k+2B\sigma/k+D_1,\quad
 J_0=C-B/k^2,\quad K_0=A/k^2+2B\sigma/k^2+D_1/k+E.
\]
The scalar is identically
\[
 A N^2+B U^2+qLWN+\sigma^2J_0Z^2+q^2K_0W^2.       \tag{UH8}
\]
Exact source contact and all stationary jets give
\[
 L(k)=J_0(k)=K_0(k)=0,\qquad K_0'(k)=-\sigma J_0'(k),
\]
\[
 B_x(k)+s[A_x(k)+a_0C_x(k)]=0,\qquad
 A_y(k)+C_y(k)/k+s a_0B_y(k)=0.                    \tag{UH9}
\]
These hold separately for every centered excess, hence for every cap.
Write the monic divisions
\[
 L=\delta L_1,\quad J_0=\delta j+\delta^2J_2,\quad
 K_0=-\sigma\delta j+\delta^2K_2,
\]
\[
 B_x+s(A_x+a_0C_x)=\sigma\delta H_x,\qquad
 A_y+C_y/k+s a_0B_y=\delta H_y.
\]
All quotients are bounded. The scalar and full first normals become
\begin{align*}
 G_{00}={}&A N^2+B U^2+(L_1+k\sigma j)SN-kjTU
                          +J_2T^2+K_2S^2,\\
 G_{10}={}&(A_x+a_0C_x)U+C_xqwN+H_xT,\\
 G_{01}={}&(C_y+s q^2B_y)N+qB_ywU+H_yS.            \tag{UH10}
\end{align*}
Together with the unchanged full quadratic normal block, this is an
exact bounded six-root Gram for $G$, not a scalar restriction.
The only root relation needed for strict completion is
\[
 T^2=\sigma S^2+k\sigma\delta SN-k\delta TU.        \tag{UH11}
\]
Its verification is one multiplication using (UH6). Adding any bounded
multiple of its zero changes no parent. No mixed inverse, root, order,
or source/test identification occurs in (UH6)--(UH11).

For the cap argument it is useful to print every centered row directly
in this same frame. In the order $m,c,(a,\ell),D,\beta$, they are
\begin{align*}
 &uzN_y-(kR+z\delta)N-qwN_x-(zQ/k)S,\\
 &S-zN_y-x(1-r)N,\qquad N_y-xbN,\\
 &(u-kbz)U+uN_x-qwN_y+(E_s/k)T,\qquad N_x+HT.
                                                               \tag{UH20}
\end{align*}
These are exact identities before taking any residue.

\paragraph{Explicit coefficient input.}
For (UH2) before caps put $j_u=u-kbz$, $p_u=1+ukbz$,
$t_u=kh-u$. The full scalar coefficients in (UH7) are
\begin{align*}
 A={}&4mz^2t_u^2+4\ell D_u+(4\beta+2)zt_u^2,\\
 B={}&(4cz+4D+2)j_u^2+4aD_uxb^2z,\\
 C={}&[4D(1+ukh)^2+4\beta(u-kh)^2+2D_uxh^2]/\sigma^2,\\
 D_1={}&8mzp_ut_u-8Dj_u(1+ukh)-8\ell D_ukb
          -4u^2k(h+bz)-4u+4uxhbz,\\
 E={}&4mp_u^2+4cj_u^2+4(a+\ell)D_uxb^2+2D_uxb^2z.
\end{align*}
The first normals are
\begin{align*}
 A_x={}&-8mzp_u+(8D+4)uj_u,\\
 B_x={}&-8Du(1+ukh)+8\beta(u-kh)-4D_ukh,\\
 C_x={}&-8mzt_u+(8\beta+4)(u-kh),\\
 A_y={}&8mzup_u-8czj_u+8(a+\ell)D_ukb+4D_ukbz,\\
 B_y={}&-8czj_u+8aD_ukb-(8D+4)j_u,\\
 C_y={}&8mz^2ut_u-8\ell D_u+8Dz(1+ukh)+4uzt_u.
\end{align*}
The coefficients of $N_x^2,N_y^2,N_xN_y$ are respectively
\[
 4mzw^2+4Du^2+4\beta D_w+2(D_u+w^2),
\]
\[
 4mz^2u^2+4cz^2D_w+4aD_uD_w+4\ell D_u+4Dzw^2
                            +2z(D_u+w^2),
 \qquad-8mzuqw-8Duqw-4uqw.                         \tag{UH12}
\]
Here products such as $Du$ mean $D\,u$, whereas $D_u=1+u^2$.

\paragraph{Infinitesimal $\sigma$: source collapse and strict residue.}
The exact polynomial $P_m=(N_D+xN_\beta)/\sigma$ is bounded.
At $\sigma=0$, $N_D=-xN_\beta$ and
\[
 N_\beta=x(K+z)H+bxz-x-z.
\]
Nonnegativity of $D,\beta$ forces
\[
 H=H_0+O(\sigma),\quad
 H_0={x+zQ\over x(K+z)},\qquad E_0=1-xH_0={R\over K+z}.
                                                               \tag{UH13}
\]
Indeed $JD+x\theta\beta=\sigma^2P_m/(2P)$ and $J,\theta$
are units. Thus $D,\beta=O(\sigma^2)$, including either literal zero.
Moreover
\[
 (D+x\beta)/\sigma^2\longrightarrow KH_0E_0/2>0.
\]
The other limiting coefficients are
\[
 c_0={R\over2x(1-zb)(K+z)},\quad m_0={xH_0\over2(1-zQ)},
 \quad a_0^{\rm wt}={H_0\over2b},\quad
 \ell_0={xR\over2Q(K+z)}.                           \tag{UH14}
\]
At a matched angle $u=k$, the residue Gram (UH10) splits into
$(N_y,N,S)$ and $(N_x,U,T)$, because their only off-block coefficients
contain $w$. Before caps, the $(N_y,N)$ block is positive definite;
its angular Schur complement is
\[
 \Gamma=8RH_0/K>0.                                \tag{UH15}
\]
The $(N_x,U)$ block is precisely
\[
 2K(N_x+xH_0U)^2+2KxH_0E_0U^2.                   \tag{UH16}
\]
The remaining $T$ couplings are bounded. As $\sigma,\delta$ have
zero residue, (UH11) permits an arbitrarily large positive $T^2$
pivot without changing any other residue coefficient. This proves a
strict full six-root residue Gram. For $u^2\ne x$, use the full raw
frame $(E_t,F_t,G_t,R_t,X,T_0)=(qW,Y/q,wZ/q,wW,X,\sigma Z)$.
Its leading first block is
\begin{align*}
 &4m_0[E_t+uz(F_t-G_t)]^2+4c_0(uE_t-zF_t)^2
       +4a_0^{\rm wt}D_uF_t^2+4\ell_0D_u(F_t-G_t)^2\\
 &\qquad+2z\{u^2(G_t-F_t)^2+F_t^2\}-4uE_tG_t.
\end{align*}
Put $v_0=u^2$ and $M_0=1-zQ$. In $(F_t,F_t-G_t)$ the hard block
is diagonal, with weights
\[
 B_h={2(x+zQ)[v_0(1-r)+R]\over bx(1-r)(K+z)},\qquad
 L_h={2R[v_0(x+zQ)+xM_0]\over Q(K+z)M_0}.
\]
The exact $E_t$ Schur complement is
\[
 {2(v_0-x)^2(1+v_0)RH_0\over
 [v_0(1-r)+R][v_0(x+zQ)+xM_0]}.                  \tag{UH17}
\]
It has a positive-unit limit through every opening when $v_0\ne x$.
The hard minimizers are
\[
 {F_t\over E_t}={ubK\over v_0(1-r)+R},\qquad
 {F_t-G_t\over E_t}={-uKQ\over v_0(x+zQ)+xM_0},
\]
also bounded through every opening. The other raw blocks are
$2\{u^2R_t^2+2u^2R_tX+D_uX^2\}$ and
$4(D+u^2\beta)T_0^2/\sigma^2$, both strictly positive.
This gives the full bounded away-angle lift, without a determinant-only
claim or any omitted first state normal.

\paragraph{Unit $\sigma$: angular Schur sign.}
Here (NI10) gives $z=0$ at residue. Put $h=\sigma H$,
$J=1+\sigma(xH-1)$, $\theta=1-\sigma H$. The source signs imply
\[
 0<\sigma<K,\qquad {1\over K+\sigma}\le H
       \le{x+\sigma\over x(K+\sigma)}.             \tag{UH18}
\]
These inequalities are stated on the interior; its closure includes
$\sigma=K$ when $J=\theta=0$.
The limits are
\[
 m={KxH\over2(K+\sigma)},\quad c={K(1-xH)\over2x(K+\sigma)},
 \quad a={1\over2b(K+\sigma)},\quad
 \ell={x\over2Q(K+\sigma)},
\]
\[
 D={K\over2J(K+\sigma)}-\tfrac12,
 \quad\beta={K\over2\theta(K+\sigma)}-\tfrac12.
\]
The two hard blocks in (UH10) are positive definite. Their angular
Schur coefficients $\gamma_S,\gamma_T$ have no $ST$ term and satisfy
\[
 \gamma_S+\sigma\gamma_T
 ={ -2p(H)\over\sigma(K+\sigma)}>0,\qquad
 p(H)=Kx(K+4\sigma)H^2-2K(x+2\sigma)H+x.           \tag{UH19}
\]
Convexity of $p$ and (UH18) prove strict negativity: its endpoint
values are respectively
\[
 {-\sigma(3\sigma x+4\sigma+4x+4)\over(K+\sigma)^2},\quad
 {-\sigma(4\sigma x+3\sigma+4x^2+4x)\over x(K+\sigma)^2}.
\]
Use (UH11) to distribute this positive total between $S^2,T^2$;
for example choose its multiple so both resulting Schur coefficients
are positive. This yields the required strict six-root Gram.

For completeness, angles $u^2\ne x$ are strictly positive as well.
Eliminating the two hard state variables from the leading form,
multiplying by the positive source factor $(K+\sigma)/2$, and
writing $U_0=u^2$, $P_0=(w/q)(Z/q)$, $T_0=(Z/q)^2$ gives a positive
Laurent function
\[
 A+B P_0^2+C P_0^2/T_0+D_0T_0-2E P_0,
\]
where
\[
 A=KxH+K(1-xH)U_0/x,\ B=x(1+U_0),\
 C={KU_0\over J+\theta U_0},
\]
\[
 D_0=(1+U_0)\left[{K(1+U_0)\over J+\theta U_0}-(K+\sigma)\right],
 \quad E={K\sqrt{U_0}(1+U_0)\over J+\theta U_0}.
\]
The strict criterion is $E<\sqrt{AB}+\sqrt{CD_0}$.
Squaring only to analyze these positive real residue quantities gives
\[
 (E^2-AB-CD_0)^2-4ABCD_0
 ={K^2(1+U_0)^2\over(J+\theta U_0)^2}(U_0-x)^2Q_2(U_0).
\]
The quadratic $Q_2$ has positive constant $H^2x^2J^2$, positive
leading coefficient $\theta^2(1-xH)^2$, and $Q_2(x)=xp(H)<0$.
Explicitly, with $E_s=1-xH$,
\[
 Q_2(U_0)=E_s^2\theta^2U_0^2
 +2E_s\{x\theta^2+E_s\theta J+(K+\sigma)\theta-2K\}U_0
 +x^2H^2J^2.
\]
The concave quadratic
$U_0(K+\sigma)-[x^2H+(1-xH)U_0](J+\theta U_0)$
is negative at zero and infinity and positive at $x$; at each of its
two zeros the preceding discriminant is negative. Consequently its
positive interval lies inside $Q_2<0$. This proves the criterion for
all $U_0\ne x$. Sign charts with $P_0\le0$ are immediate. These
square roots are an analysis of real residues only and are not tensor
roots or mixed localizations.

This inequality is directly a bounded Gram allocation, not a
comparison of mixed ratios. Before choosing projective coordinates,
put $E_t=qW$, $G_t=wZ/q$, $R_t=wW$. The exact relation
$E_tG_t=R_tZ$ permits distributing the cross term between the two
quadratic blocks $(E_t,G_t)$ and $(R_t,Z)$. The strict inequality
above is exactly the condition that both $2\times2$ blocks can be
made positive definite by a finite real allocation. It therefore
handles every independent rate of $q,w$.
At $J=0$ or $\theta=0$, the root-interval argument permits one endpoint
at zero or infinity. At the double face $J=\theta=0$, necessarily
$\sigma=K$, $H=1/K$. As $J+\theta U_0\to0$,
\[
 E-\sqrt{CD_0}\longrightarrow K\sqrt{U_0},\qquad
 AB-K^2U_0=(U_0-x)^2.
\]
Thus the constrained allocation remains strictly positive away from
$U_0=x$, including this double face without dividing by a vanishing
residue denominator.

\paragraph{Opening faces, tensor conclusion, and exact scope.}
The finite-cap argument is applied to the full Gram above. For
$\sigma$ infinitesimal the possible diverging rows have openings
$1-zb$, $1-zQ$, $b$, $Q$; at most two can diverge together. The pairs
$(c,m)$, $(m,a)$, $(c,\ell)$ have independent hard rows; other pairs
are excluded by $0<zb,xb<1$ and $1-zQ>0$. In the $(N_y,N)$ columns
of (UH20), their determinants have absolute values $kzK,k,x$,
respectively; each is a unit on its compatible face. The hard block
is uniformly coercive, since it contains positive-unit multiples
of $N_y^2$ and $(N_y-N)^2$. Its minimizing hard coordinates for
fixed angular $S$ are explicitly
\[
 {N_y\over S}={b(2bxz-x+1)\over K},\qquad
 {N\over S}={4bxz-x-2z+1\over xK}.
\]
They are bounded through every opening. Thus the limit of the
uncapped Schur complement is the Schur complement restricted to the
diverging centered rows' kernel. The quotients (UH15)--(UH16) stay
positive units there. The elementary finite-dimensional penalty
lemma now applies: a bounded symmetric form strictly positive on
the common kernel becomes strictly positive after sufficiently large
finite positive multiples of these rows' squares are added. This
proves the finite-cap claim, with the entire angular right-hand sides
in (UH20) retained. Away from $u^2=x$, (UH17) and its printed bounded
hard minimizers give the same constrained strictness directly.

For unit $\sigma$, only $b,Q,J,\theta$ can be openings.
The $b,Q$ rows coincide in the hard direction $N_y-xbN$ and cannot
both diverge. The $J,\theta$ rows on $(N_x,U)$ are independent even
if both vanish; their determinant is the unit $u$. Formula (UH19)
remains a positive unit through these faces. More explicitly, the
two hard blocks dominate positive-unit multiples of
$N_y^2+(N_y-N)^2$ and $N_x^2+(N_x+U)^2$. Their angular minimizers are
\[
 {N_y\over S}={b(1-HxK)\over K},\quad
 {N\over S}={1-HxK\over xK},
\]
\[
 {N_x\over T}={\sigma H^2xK-HxK-\sigma H(2x+1)+x\over\sigma K},
 \quad {U\over T}={HxK(K+\sigma)-xK-\sigma\over x\sigma K}.
\]
Their $D$ and $\beta$ row defects divided by $T$ are respectively
$k(HK-1)J/(\sigma K)$ and $-x\theta(HK-1)/(\sigma K)$.
Hence every relevant constrained limit exists with bounded
minimizers; the same penalty argument applies. On the angular
relation $T^2=\sigma S^2$, its strictly positive Schur value is
(UH19), and after taking finite caps the bounded zero (UH11)
distributes that value into two strict angular pivots. The away-angle
allocation and its single/double-face limits above give the same
conclusion there. All removed excesses remain exact positive
source/test squares. No positive lower bound on an original opening
has been imposed.

A strict bounded residue Gram lifts by the common tensor Gram
lemma. Every root in (UH6) is a finite sum of separated source and
testing factors. After restoring the retained squares this certifies
the complete original parent, including $D=0$ or $\beta=0$.
Together with revision-6 remote tests, this closes both (NI9) and
(NI10). Those are the entire actual UH portion of C1.5. The already
proved UU and UI portions and exact source stabilizers therefore close
the entire C1.5 row, subject to the exact checks and mathematical review
recorded alongside this section.



% Source formulas: p7d4e5164e84a
\section{Vanishing contact scales and the state normal directions}
When $k$ vanishes, some contact directions become singular. The following polynomial identities retain their exact factors before any rescaling. They control the scalar term and both state normal directions, including the complementary state charts.

\subsection{The endpoint polynomial source}

\label{component:p7d4e5164e84a}
\paragraph{Seven short source polynomials.}
The following definitions make the endpoint input reproducible without a stored expansion. Put
\[
K=k^2,\quad \eta=e^2b^2,\quad\zeta=k^2b^2,\quad V_c=K\eta\zeta.
\]
\[
L_c=1-a\eta,\qquad H_c=1+(1-a)\eta,\qquad Q_c=1-K\zeta+aV_c.
\]
Define
\begin{align*}
P_4={}&1+K+\eta K(a-K-1)-\zeta Kt
 +V_c[-\zeta Kt^2+2Kt+a(1-t)+t],\\
P_\alpha={}&1+K-a\eta K+\zeta[(1+K)-(2+K)t]
 +V_c[-\zeta t^2+a(t-1)+t],\\
P_d={}&(1+K)^2-a\eta K^2+\zeta K[(K+1)-(K+2)t]\\
&+V_c[-\zeta Kt^2+Ka(t-1)+Kt+K+1].
\end{align*}
The seven coquotients are
\begin{align*}
c_2&=2(1+K)(1-t\zeta)L_c(1+V_c)H_cQ_c,\\
c_\beta&=2(1+K)L_c[1+(1-t)\zeta]H_cQ_c,\\
c_3&=2(1+K)L_cH_cQ_c(1-K\eta+tV_c),\\
c_5&=2(1+K)L_cH_cQ_c(1+V_c),\\
c_4&=2Q_cP_4,\qquad c_\alpha=2L_cH_cP_\alpha,
\qquad c_d=4L_cQ_cP_d.
\end{align*}
Their residues are respectively $2,2,2,2,2,2,4$, so all are positive coefficient units on this domain. With $M=R_c-t$ and $J=T_0-a$, the unchanged source is
\begin{gather*}
\lambda=n=0,\quad c=2,\qquad
d=\frac{2a\eta c_d}{c_3},\quad m=\frac{2ae^2c_2}{Mc_3},\quad
A=\frac{2a\eta^2c_4}{c_3},\quad
\alpha=\frac{2aK^2\eta\zeta c_\alpha}{c_3},\\
D+\frac d2=\frac{2ac_5}{Jc_3},\qquad
\beta+\frac d2=\frac{2ae^2c_\beta}{tc_3},\qquad
\Lambda=\frac{R_cT_0c_3}{2a}>0.
\end{gather*}
The theorem retains the hypothesis $D,\beta\ge0$. Write
\[
r_2=v-X+uvY-uZ,\quad r_\beta=X+uZ,\quad r_3=u-Y,
\quad r_5=u+uvX-Y-vZ,\quad r_\alpha=vY-Z,
\]
and put $q_0=v_g$, $z_0=\eta u$. Then $G$ is exactly the sum of
\begin{align*}
G_2&=4T_0e^2c_2[(r_2+tq_0)^2-tR_cq_0^2],\\
G_\beta&=4T_0e^2c_\beta(1+v^2)[(r_\beta+Mq_0)^2-R_cMq_0^2],\\
G_3&=4R_cc_3(1+v^2)[(r_3+Jz_0)^2-T_0Jz_0^2],\\
G_5&=4R_cc_5[(r_5-az_0)^2-aT_0z_0^2],\\
G_4&=4R_cT_0\eta^2c_4(1+u^2)(1+v^2)Y^2,\\
G_\alpha&=4R_cT_0K^2\eta\zeta c_\alpha(1+u^2)r_\alpha^2,\\
G_d&=2R_cT_0\eta c_d(1+u^2)[r_\alpha^2-(1+v^2)Z^2].
\end{align*}
These seven factored expressions agree separately with the defining source. Their proof requires multiplication of short polynomials only.



% Source formulas: p935c9d6318ea
\subsection{Finite monic divisions at the endpoint}

\label{component:p935c9d6318ea}
\paragraph{Deriving the scalar and first-normal inputs by monic division.}
Here is a finite construction of every coefficient in the preceding table. It uses the seven factored terms, differentiation, and division by monic polynomials. Thus neither a stored expansion nor a rational-function ideal calculation is an input.

Let $\mathcal R=\Q[a,t,b,e,k]$ and use an indeterminate $U$ for $u^2$ in this paragraph only. Put $K=k^2$. Substitute $(X,Y,v)=(X_g,Y_g,v_g)$ in the peeled polynomial $E$. Directly from the seven terms, the resulting expression has the form
\[
\frac{E_g}{e^2b^4}=e^2\mathcal A(U)+ekZ\mathcal B(U)+k^2Z^2\mathcal C(U),
\qquad \mathcal A,\mathcal B,\mathcal C\in\mathcal R[U].
\]
The displayed division means coefficientwise polynomial extraction: the constant, linear and quadratic coefficients in $Z$ have the factors $e^2$, $ek$ and $k^2$, respectively. No evaluated tensor is divided. Substituting the graph into the factored terms gives the five short contact and axis identities
\begin{gather*}
(\mathcal A+\mathcal B+\mathcal C)(K)=0,\qquad
(\mathcal B+2\mathcal C)(K)=0,\qquad
(\mathcal A+\mathcal B+\mathcal C)'(K)=0,\\
\mathcal A(0)=\mathcal B(0)=0.
\end{gather*}
Consequently the following monic divisions have zero remainder:
\begin{align*}
\mathcal N&=\frac{\mathcal A+\mathcal B+\mathcal C}{(U-K)^2},&
\mathcal D&=\frac{\mathcal B+2\mathcal C}{U-K},\\
\mathcal R_0&=\frac{\mathcal C+K\mathcal D+K^2\mathcal N}{U},&
\mathcal J&=\frac{-\mathcal D-2K\mathcal N}{U}.
\end{align*}
For the last two divisions, evaluate the numerators at $U=0$ and use the first two division identities together with $\mathcal A(0)=\mathcal B(0)=0$. This is a polynomial-ring argument; the symbols $K$ and $U$ are not inverted after evaluation. Set
\[
A_{00}=\mathcal R_0(u^2),\qquad A_{01}=u\mathcal J(u^2),
\qquad A_{11}=\mathcal N(u^2).
\]
Comparing the three coefficients of $1,Z,Z^2$ proves exactly
\[
\frac{E_g}{e^2b^4}
=A_{00}\,[u(e-kZ)]^2+A_{01}\,u(e-kZ)(eu^2-k^3Z)
 +A_{11}(eu^2-k^3Z)^2.
\]
Multiplying by $e^2b^4$ gives the scalar identity used above.

For each $\xi\in\{X,Y,v\}$, let $R_\xi$ be the graph restriction of $\partial_\xi E$, divided coefficientwise by $e^2b^2,eb^2,e^2b^2$, respectively. The first two are linear in $Z$. The third has one quadratic coefficient, supplied solely by the signed term:
\[
[Z^2]R_v=-4R_cT_0c_d\,e k b^2u(1+u^2).
\]
Define
\[
C_X=C_Y=0,\qquad C_v=4R_cT_0c_d\,eb^2(1+u^2).
\]
Then $\widetilde R_\xi=R_\xi-C_\xi Z\,u(e-kZ)$ is linear in $Z$ and has the polynomial parity form
\[
\widetilde R_\xi=eu\,L_{\xi,0}(U)+kuZ\,L_{\xi,1}(U),
\qquad L_{\xi,i}\in\mathcal R[U].
\]
In particular, subtracting the quadratic term also leaves its corrected linear coefficient divisible by $ku$; this follows upon collecting the signed term with the two current terms. The contact identity is
$\big(L_{\xi,0}+L_{\xi,1}\big)(K)=0$. Form the monic quotient
\[
J_\xi=\frac{L_{\xi,0}+L_{\xi,1}}{U-K}.
\]
The promised full first-normal coefficients are now defined by the short formulas
\[
P_\xi=(L_{\xi,0}-UJ_\xi)(u^2)+C_\xi Z,
\qquad Q_\xi=uJ_\xi(u^2).
\]
Indeed direct multiplication gives
$R_\xi=P_\xi u(e-kZ)+Q_\xi(eu^2-k^3Z)$.
This proves the full first-normal identities, including the quadratic-in-$Z$ term; no first-normal term has been omitted.

For completeness, the weighted splits have an equally explicit finite rule. In a polynomial expanded in $(e,k,b,u)$, place a monomial with $k$-exponent at least $r$ in the first bucket and divide it by $k^r$; place every remaining monomial in the second bucket and divide it by $e^s b^s$. Use $(r,s)=(4,2)$ for $A_{00}$ and $(r,s)=(2,1)$ for
\[
A_{01},\quad P_X,\quad P_Y,\quad P_v-16kP_0.
\]
Substitution of the seven short terms into the monic formulas shows that every monomial left for the second bucket has the required $e^s b^s$ factor. Equivalently, the only polynomial remainders to check are
\[
A_{00}\equiv0\pmod{(k^4,e^2b^2)},\qquad
A_{01},P_X,P_Y,P_v-16kP_0\equiv0\pmod{(k^2,eb)}.
\]
These are ordinary finite polynomial truncations: discard the monomials divisible by the displayed generators. The two bucket quotients have respectively residues $(-8,-8)$ for $A_{00}$ and $(0,0)$ for each of the other four entries; $A_{11}$ has residue $8$, and all $Q_\xi$ have residue zero. To obtain these numbers, after the divisions set $e=k=u=Z=0$, keeping $a,t,b$ free. In particular the new term $C_vZ$ belongs to $eb\,\mathcal R[u,Z]$ and its second-bucket quotient contains the infinitesimal testing factor $Z$. No assumption $\bar b=0$ is used.

This is a reproducible coefficient construction: the source consists of seven short terms, the scalar construction of four monic quotients, and each first-normal construction of one monic quotient. The finite checks accompany these formulas as a convenience; the definitions and the multiplication identities, rather than expanded coefficient arrays, are the proof objects.



% Source formulas: p05180564ed05
\subsection{The endpoint product relation}

\label{component:p05180564ed05}
\label{p05180564ed05:sec:productrepair}
The new boundary mechanism is algebraic. A troublesome root is retained, and a second relation controls its product with an additional root. This supplies two positive pivots without comparing an infinitesimal in one field with an infinitesimal in the other.

\begin{lemma}[Product-relation completion]\label{p05180564ed05:lem:productrepair}
Let $z$ be a root vector over $T$, and $h,r\in T$. Suppose
\[
F=(z,h,r)^TQ(z,h,r),\qquad Q\in\Sym(B),\qquad
\bar Q=\operatorname{diag}(Q_0,-a,0),
\]
where $Q_0$ is a real positive-definite matrix and $a\in\R_{\ge0}$. Suppose $p\in B$ has nonzero residue, $A$ is a row with entries in $\mathfrak n$, and
\[
ph=Az+r,\qquad phr=(z,h,r)^TE(z,h,r),\quad E\in\Sym(\mathfrak n).
\]
Then $F\in\Sigma T^2$.
\end{lemma}
\begin{proof}
For a sufficiently large positive real $M$, add the exact zero
\[
M\big[p^2h^2+r^2-(Az)^2-2phr\big]=0.
\]
The residue Gram becomes $\operatorname{diag}(Q_0,M\bar p^2-a,M)\succ0$. The subtracted terms have zero residue by hypothesis. Lemma~\ref{p25af2c58da7e:lem:gram} applies. In particular $p$ is never inverted.
\end{proof}

\subsubsection*{Exact contact scope and full-polynomial input}
This application uses the canonical C1 contact model, with $\lambda=n=0$. Its source parameters are positive $e,k,b,a,t$ satisfying
\[
e,k\in\mathfrak m_c,\quad b\in B_c,\qquad
R_c=1+k^2(1-ae^2b^2),\quad T_0=1+k^2-tk^4b^2,
\]
\[
0<a<T_0,\qquad0<t<R_c.
\]
The signed weight satisfies $d>0$, and all eight event coefficients are nonnegative original physical C1 contact coefficients. This is a specified contact chart, not every coefficient vector with the same omissions. The independent testing chart is $W=1$ with $u,v,X,Y,Z$ infinitesimal. In this section there is no restriction on $k^3/e$ or $k/e$.

The complete seven-term contact polynomial is denoted $G$. Its short factored definition and unchanged original coefficient vector are given below. 

Define
\begin{gather*}
v_g=ekb^2u,\quad X_g=-uZ+tv_g,\quad Y_g=(1-ae^2b^2)u,\\
N_1=X-X_g,\quad N_2=Y_g-Y,\quad \delta v=v-v_g,\quad D_t=Z-uX,\\
U_0=1+u^2,\quad c_0=1-tk^2b^2,\quad P_0=T_0+u^2=U_0+k^2c_0,\\
S=u+uvX-Y-vZ+(T_0-a)e^2b^2u,\\
s_0=eb^2u(e-kZ),\qquad q=eb^2(eu^2-k^3Z).
\end{gather*}
\par\noindent\textit{Polynomial source. }The source is defined (Section~\ref{component:p7d4e5164e84a}).\par
Direct substitution gives the exact nonlinear relation
\[
L=P_0s_0-S+N_2-D_t\delta v+uv_gN_1-uc_0q=0. \tag{E1}
\]
Peel the two complete current squares:
\[
E=G-H_3(1+v^2)N_2^2-H_5S^2,\qquad
H_3=4R_cc_3,\quad H_5=4R_cc_5,\quad \bar H_3=\bar H_5=8.
\]
Each peeled term divided by $e^2$ is polynomial in bounded variables. The full graph scalar has the identity
\[
E_g=A_{00}s_0^2+A_{01}s_0q+A_{11}q^2,
\]
and the entire first normals, with $g_0=s_0/(eb^2)$ and $g_2=q/(eb^2)$ denoting their displayed polynomial expressions, satisfy
\[
E_X=e^2b^2(P_Xg_0+Q_Xg_2),\quad
E_Y=eb^2(P_Yg_0+Q_Yg_2),\quad
E_v=e^2b^2(P_vg_0+Q_vg_2).
\]
The exact polynomial splits and residues needed are
\begin{center}\small
\begin{tabular}{ccc}\toprule
Quantity&Polynomial split&Quotient residues\\\midrule
$A_{00}$&$k^4a_0+e^2b^2a_1$&$-8,-8$\\
$A_{01}$&$k^2c_1+ebc_2$&$0,0$\\
$P_X$&$k^2p_X+ebq_X$&$0,0$\\
$P_Y$&$k^2p_Y+ebq_Y$&$0,0$\\
$P_v-16kP_0$&$k^2p_v+ebq_v$&$0,0$\\\bottomrule
\end{tabular}
\end{center}
Also $\bar A_{11}=8$ and $\bar Q_X=\bar Q_Y=\bar Q_v=0$. These equalities remain valid with $b$ symbolic; boundedness, rather than infinitesimality, suffices for $b$.

\par\noindent\textit{Coefficient construction. }The finite coefficient construction is given (Section~\ref{component:p935c9d6318ea}).\par

Exact finite Taylor expansion in $(N_1,-N_2,\delta v)$ accounts for every higher-normal term. Because $E/e^2$ is polynomial, every term of normal degree at least two allocates to $x=eN_1$, $y=e\delta v$, $N_2$ with bounded coefficients. Terms of degree at least three have an extra infinitesimal normal. Terms involving $N_2$ retain an $e$ factor. The only additional nonzero residue is
\[
B_0=\begin{pmatrix}16&-8\\-8&8\end{pmatrix}
\quad\text{on }(x,y).
\]
It is visible without expansion of the full Hessian: at $u=Y=Z=0$, the quadratic residue of $E/e^2$ is $8(X-v)^2+8X^2$, with $b$ still free.

\subsubsection*{The new identity and its nine-root proof}
If the signed testing core is nonnegative, its pure testing square root and the original event terms already give SOS. Otherwise the sign gauges $u,v\ge0$ preserve every term, and the core Schur complement implies $u,v>0$ and $D_t>0$. To check the latter, after $Z=D_t+uX$ the Schur complement on $(W,D_t)$ is
\[
\frac{2(1+v^2)}{N_0}
\begin{pmatrix}
u^2v^2(1+u^2)&-uv(1+u^2)(1+v^2)\\
-uv(1+u^2)(1+v^2)&u^2v^2
\end{pmatrix},
\]
where $N_0=u^2v^4+u^2v^2+u^2+v^4+2v^2+1>0$. The eliminated block is positive definite. For $W=1,D_t\le0$, the remaining form is nonnegative, proving the claim.

Consequently $w=\sqrt{ek}\sqrt{D_t}\,\delta v$ is a legal separated root. Set
\begin{gather*}
h=k^2s_0-uq=e^2b^2u(k^2-u^2),\qquad z=ebs_0,\qquad r=k^2D_t\delta v,\\
A_0=k^2(S-N_2)-k^3b^2u^2x-uU_0q.
\end{gather*}
Multiplication of (E1) by $k^2$ gives $P_0h=A_0+r$. The additional exact identity is
\begin{align*}
P_0hr={}&ku^2(P_0-k^2)y(S-N_2)
-ku(P_0+c_0k^2u^2)yq\\
&+u^2(P_0-k^2)w^2-k^4b^2u^2T_0xy. \tag{E2}
\end{align*}
Every displayed coefficient has zero residue and is a polynomial in separately bounded factors. To derive (E2), substitute the two soft roots to obtain
\[
u(P_0-k^2)s_0-(1+c_0u^2)q
=ekb^2(k^2-u^2)(D_t+uN_1).
\]
Then use $D_t=U_0Z-uN_1-tu v_g$, $T_0=1+k^2c_0$, and (E1) to replace $S-N_2$ in (E2). The terms in $D_t(\delta v)^2$ cancel; the remaining terms give $P_0hr$. This is a whole-polynomial calculation.

Let $J_0$ denote the right-hand side of (E2). Thus
\[
Z_0=P_0^2h^2+r^2-A_0^2-2J_0=0.
\]
Add to the entire $G$ the exact zero expression
\[
-16kyL+16(ebL)^2+16Z_0.
\]
The first term cancels the exceptional $16kP_0$ first normal and supplies $16w^2$. Use $k^2s_0=h+uq$ in every other scalar and first-normal term. The scalar becomes
\[
a_0(h+uq)^2+a_1z^2+[c_1(h+uq)+c_2z]q+A_{11}q^2.
\]
Its only negative residue diagonals are $-8$ on $h,z$. The safe relation
\[
ebL=P_0z-ebS+ebN_2-bD_ty+ekb^3u^2x-ebu c_0q
\]
has only infinitesimal coefficients apart from $P_0$; hence $16(ebL)^2$ changes the $z$ diagonal to $8$. Every coefficient of $A_0$ is infinitesimal, as is every coefficient in (E2). Thus $16Z_0$ changes the $h$ diagonal to $8$ and supplies a new $r$ diagonal $16$, with no other residue change.

The full Gram on
\[
\boxed{(x,y,N_2,S,w,q,h,z,r)}
\]
therefore has residue
\[
\boxed{B_0\oplus\operatorname{diag}(8,8,16,8,8,8,16)\succ0.}
\]
Its first two LDL pivots are $16,4$; the other seven pivots are given above. The strict-residue lemma proves SOS of the complete seven-term polynomial for every coefficient ratio $k^3/e$.

\subsubsection*{Restoration and coverage}
The complete original-parent restoration is
\begin{align*}
\Lambda F={}&G
+4T_0e^2c_2\frac{t}{M}(r_2-Mq_0)^2
+4T_0e^2c_\beta\frac{M}{t}(1+v^2)(r_\beta-tq_0)^2\\
&+4R_cc_3\frac{J}{a}(1+v^2)(r_3-az_0)^2
+4R_cc_5\frac{a}{J}(r_5+Jz_0)^2.
\end{align*}
Every additional term is a nonnegative coefficient factor times a pure-testing factor times a square. This identity is four instances of
\[
S^2+2(A+B)qS+(A+B)Aq^2+\frac BA S^2
=\frac{A+B}{A}(S+Aq)^2.
\]
Use respectively $(A,B,q)=(M,t,q_0),(t,M,q_0),(a,J,z_0),(J,a,-z_0)$ and $S=r_2-Mq_0,r_\beta-tq_0,r_3-az_0,r_5+Jz_0$. Comparing the seven resulting coefficients with the displayed source proves the identity. Since $\Lambda$ is a positive pure coefficient factor, the pure-factor restoration in \hyperref[component:p9b00698ff36a]{RFA13} (p.~\pageref{component:p9b00698ff36a}) recovers the original parent. Normalizing the original state to $W=1$ and undoing that normalization use only testing factors.

The result applies to every positive bounded $b$, with no
condition on $k^3/e$. The unbounded-$b$ case is proved separately in
Section~\ref{component:pabbcd0905dcd}; the full source and state partition
is given in Section~\ref{component:p845b3cb8b114}.


% Source formulas: pa4f6d1e57f34
\subsection{The uniform small-$e$ source on the bounded $Z$ chart}

\label{component:pa4f6d1e57f34}

\label{pa4f6d1e57f34:sec:uniform-c1-z-new}

\paragraph{Domain and conclusion.}
Use the original physical uniform C1 source with $d=1$,
$\lambda=n=0$, $k,f,g,e=kf/g$ positive infinitesimals, and
all original nonnegative event weights. Retain
\[
 r=1-af^2,\quad R=1+k^2r,\quad
 T=1+k^2-tk^2g^2,\quad M=R-t>0,\quad J=T-a>0,
 \qquad a,t>0.                                                   \tag{EZ1}
\]
In particular $a,t$ are bounded, $r,R,T$ are positive units,
and $M,J$ may approach zero at arbitrary rates.
For every independent testing $u,v,X,Y$ infinitesimal and
$Z=1$, $W$ arbitrary bounded of either sign, the entire original
parent is a separated tensor SOS. This includes both bounded
and unbounded $b=g/k$, and every permitted $\beta=0$ face.
It does not claim the complementary $W=1,Z$-infinitesimal chart.

\subsubsection*{The exact source and a common set of bounded coefficients}

Use the seven short native coquotients
$c_2,c_\beta,c_3,c_5,c_4,c_\alpha,c_d$ in the retained
Section~\ref{component:p7d4e5164e84a}, now with $\eta=f^2,\zeta=g^2$.
Their residues are $2,2,2,2,2,2,4$.
In the present $d=1$ normalization the exact weights are
\[
\begin{gathered}
 C=\frac{c_3}{af^2c_d},\quad A=\frac{f^2c_4}{c_d},\quad
 m=\frac{k^2c_2}{g^2Mc_d},\quad
 \alpha=\frac{k^4g^2c_\alpha}{c_d},\\
 D+\tfrac12=\frac{c_5}{Jf^2c_d},\qquad
 \beta+\tfrac12=\frac{k^2c_\beta}{g^2tc_d}.
\end{gathered}                                                   \tag{EZ2}
\]
Every use below is of these actual weights. In particular
$\beta\geq0$ is retained, not replaced by a formal source condition.

The following two pure-source coordinate choices are exhaustive:
\[
\begin{array}{c|ccccc}
 &s&\omega&P&G&\kappa=P/G\\ \hline
 b\text{ bounded}&e&b^2&k^2&k^2&1\\
 b\text{ unbounded}&f&1&k^2&g^2&(k/g)^2
\end{array}                                                     \tag{EZ3}
\]
In both rows $f^2=s^2\omega$, $g^2=\omega G$,
$s,\sqrt P,\sqrt G$ are infinitesimal, and $\omega,\kappa$
are positive bounded. In the second row $\kappa$ is infinitesimal.
Here $\beta\geq0$ gives exactly
$t\leq2\kappa c_\beta/c_d$, so $t=\kappa\tau$ with
$\tau>0$ bounded; consequently $M$ is a unit in this row.

Define bounded source coefficients
\[
 c_*=\frac{c_3}{Tc_d},\quad d_*=\frac{c_5}{Tc_d},\quad
 a_*=\frac{c_4}{c_d},\quad \gamma_*=\frac{c_\alpha}{c_d},
 \qquad B_*=\omega\beta t
              =\frac{\kappa c_\beta}{c_d}-\frac{\omega t}{2}.
                                                               \tag{EZ4}
\]
The first four have residue $1/2$; $B_*\geq0$ may vanish.
Use the further table
\[
\begin{array}{c|ccc}
 &\mu&\chi&\psi\\ \hline
 b\text{ bounded}&c_2/(Rc_d)&t&tM\\
 b\text{ unbounded}&\kappa c_2/(Mc_d)&0&0 .
\end{array}                                                     \tag{EZ5}
\]
All entries are bounded. In the first row $\mu$ has residue
$1/2$; in the second $\mu$ is infinitesimal.

\subsubsection*{Exact removal of the source poles}

The original quarter-parent is
\[
\begin{split}
 F={}&mr_2^2+CD_vr_3^2+AD_uD_vY^2
 +(D+\tfrac12)r_5^2+(\alpha+\tfrac12)D_u(vY-1)^2\\
 &+(\beta+\tfrac12)D_v(X+u)^2-\tfrac12D_uD_v,
\end{split}                                                     \tag{EZ6}
\]
where $r_2=vW-X+uvY-u$, $r_3=uW-Y$,
$r_5=uW+uvX-Y-v$, $D_u=1+u^2$, $D_v=1+v^2$.
Put
\[
 z=sW,\qquad q_0=f^2uW^2=\omega uz^2,\qquad
 z_0=f^2uW=s\omega uz.
\]
Retain the following complete nonnegative terms of the original
polynomial:
\[
 \frac{CJ}{T}D_v(r_3-az_0)^2,\qquad
 \frac{(D+1/2)a}{T}(r_5+Jz_0)^2,\qquad
 \beta D_v(X+u-tq_0)^2.                                        \tag{EZ7}
\]
For bounded $b$, also retain
\[
                       \frac{mt}{R}(r_2-Mq_0)^2.              \tag{EZ8}
\]
Each vanishes at the actual source contact. They are literal
positive squares, even when $\beta=0$ or a source opening is small.
The first two removals use $a+J=T$ and the exact identity
\[
 \frac aT[(L+Jq)^2-TJq^2]+\frac JT(L-aq)^2=L^2.
\]
The last removal uses
$r_2^2-(t/R)(r_2-Mq_0)^2
=(M/R)(r_2^2+2tq_0r_2-tMq_0^2)$.
There is no mixed denominator.

Use all three unrestricted normal coordinates
\[
 N_x=X+u-tq_0,\qquad
 N_y=(Y-ruW)/s,\qquad N_v=(v-f^2uW)/s.
                                                               \tag{EZ9}
\]
Only $N_x$ is necessarily bounded. The exact bounds needed later are
$N_x,sN_y,sN_v\in\mathrm I$; they follow from the stated testing
chart and bounded $W$. No boundedness of $N_y,N_v$ is assumed.
Define
\[
 X_*=-u+tq_0+N_x,\qquad
 y_*=ruz+s^2N_y=sY,\qquad
 v_*=\omega uz+N_v=v/s,
\]
\[
 L_2=zv_*+uv_*y_*-X_*-u,\quad
 L_3=T\omega uz-N_y,\quad
 L_5=-N_y+v_*(uX_*-1).
\]
After multiplying the original quarter-parent by the positive
pure-source $\omega$ and removing precisely the scaled terms
EZ7--EZ8, the \emph{entire} remaining polynomial is
\[
\begin{split}
\mathcal F={}&
 \mu(L_2^2+2\chi q_0L_2-\psi q_0^2)
 +c_*D_v(L_3^2-TJ\omega^2u^2z^2)\\
 &+\omega^2a_*D_uD_vy_*^2
 +d_*(L_5^2-aT\omega^2u^2z^2)\\
 &+(\omega^2P^2G\gamma_*+\omega/2)
                         D_u(v_*y_*-1)^2
 +\frac{\omega}{2}D_v(X_*+u)^2\\
 &+B_*D_v\{2q_0(X_*+u)-tq_0^2\}
                         -\frac{\omega}{2}D_uD_v,
 \qquad D_v=1+s^2v_*^2 .
\end{split}                                                     \tag{EZ10}
\]
This is a short full-parent definition of every scalar and normal
coefficient used below. In particular no normal was set to zero
in obtaining it. All its coefficient arrays are bounded through
$a,J,M,t$; the unbounded original weights occur only in the
retained complete squares.

\subsubsection*{The entire scalar is a cubic circuit with finite corrections}

Set $U=u^2,V=z^2$. The polynomial obtained from EZ10 at
$N_x=N_y=N_v=0$ is exactly
\[
 \mathcal F_0=\omega^2\Phi(U,V),\qquad
 \Phi=\gamma_*P^2G(1+U)+UV\,H(U,V),                            \tag{EZ11}
\]
where $H$ has bounded finite polynomial coefficient arrays.
The apparent divisions by $\omega^2$ and $UV$ in this definition
are coefficientwise polynomial extractions: substitution in EZ10
shows those factors before any evaluation. No testing coordinate
is inverted.

At the actual source contact $u^2=P,z^2=G$ all three normal
coordinates and all retained terms vanish. The original contact
value and stationary equations give
\[
 \Phi(P,G)=\Phi_U(P,G)=\Phi_V(P,G)=0.                            \tag{EZ12}
\]
The exact native coquotients in EZ2 independently verify these
identities. Substituting EZ11, and cancelling nonzero pure-source
factors in the source field, gives
\[
 H(P,G)=-\gamma_*P(1+P),\quad
 H_U(P,G)=\gamma_*,\quad
 H_V(P,G)=\gamma_*\kappa(1+P).
\]
Write $\Delta=U-P$, $\Gamma=V-G$.
Finite monic division, first in $U-P$ and then in $V-G$,
therefore gives the complete scalar identity
\[
\begin{split}
 \Phi={}&\gamma_*\{\mathcal M+P\kappa U\Gamma^2\}\\
 &+UV\{\Delta^2H_{20}+\Delta\Gamma H_{11}
                                      +\Gamma^2H_{02}\},\\
 \mathcal M={}&U^2V+\kappa UV^2+P^2G-3PUV .
\end{split}                                                     \tag{EZ13}
\]
All quotients are finite bounded coefficient polynomials; no
inverse of $P,G$ is introduced by their monic construction.
The extra term $\gamma_*P\kappa U\Gamma^2$ is retained as a
nonnegative square.

In the inverse-$b$ row one further coefficient factor is needed:
\[
                         H_{02}=\kappa\widetilde H_{02},
 \qquad\widetilde H_{02}\text{ has bounded arrays}.              \tag{EZ14}
\]
Here is a finite source-division certificate, rather than an
estimate involving independent testing rates. Write
$P=\kappa G$, $t=\kappa\tau$ before substituting the native
coquotients into EZ10. Every denominator is among the pure-source
units $M,T,c_d$ and their factors. At formal $\kappa=0$,
\[
 P=t=\mu=B_*=0,\quad
 c_*=d_*=\frac{1+(1-a)s^2}{2},\quad
 a_*=\frac1{2(1-as^2)} .
\]
The formula gives $H(0,V)=0$ identically in $V$.
After clearing unit denominators, ordinary polynomial division
by $\kappa$ therefore has zero remainder. Apply the two monic
$V-G$ divisions to this quotient. This proves EZ14 with bounded
arrays. In the bounded-$b$ row $\kappa=1$, so it needs no such
factor.

\subsubsection*{Six explicit strict roots and their conductors}

The following identities apply to both rows of EZ3. Put
\[
 \mathsf h=G^{1/6},\quad\mathsf c=\kappa^{1/6},\quad
 \mathsf r=\sqrt[3]u,\quad\mathsf s=\sqrt[3]z,
\]
\[
 (L_1,L_2,L_3)=
 (\mathsf r^4\mathsf s^2,\
   \mathsf c^2\mathsf r^2\mathsf s^4,\
   \mathsf h^6\mathsf c^4),\qquad
 (f_1,f_2,f_3)=
 (\mathsf r^2\mathsf s,\
   \mathsf c\mathsf r\mathsf s^2,\
   \mathsf h^3\mathsf c^2).
\]
Define $R_{2i-1}=f_i(L_1-L_2)$ and
$R_{2i}=f_i(L_1-L_3)$ for $i=1,2,3$.
Then, exactly,
\[
 \mathcal M=\sum_{i=1}^3
       (R_{2i-1}^2-R_{2i-1}R_{2i}+R_{2i}^2).                  \tag{EZ15}
\]
Each pair has the fixed strictly positive Gram
$\left(\begin{smallmatrix}1&-1/2\\-1/2&1\end{smallmatrix}\right)$.
Real cube roots of signed testing quantities are permitted:
$z=sW$ has the separated cube root
$s^{1/3}\sqrt[3]W$. The $L_i=f_i^2$ identity is valid for
either sign of $u,z$. No mixed radical occurs.

For explicit finite allocation define
\[
 J_1=\mathsf r^3R_1+\mathsf c^2\mathsf r\mathsf s^2R_2,\qquad
 J_2=-\mathsf h^6R_1+\mathsf s^6R_2
                              -\mathsf h^3\mathsf s^3R_5.
\]
Direct multiplication gives
\[
\begin{aligned}
 uz\Delta&=J_1,&u^2z\Gamma&=J_2,\\
 \kappa uz\Gamma
   &=\mathsf c^2\mathsf r\mathsf s^2(R_2-R_1)
                                  -\mathsf c^3\mathsf s^3R_3,\\
 UV\Delta^2&=J_1^2,\\
 UV\Delta\Gamma
   &=\mathsf s^6R_2^2-\mathsf h^6R_1^2
             -\mathsf h^3\mathsf s^3R_1(R_5+R_6),\\
 \kappa UV\Gamma^2
   &=\mathsf s^6(2R_1^2-3R_1R_2+R_2^2+R_3^2)
                      +\mathsf h^3\mathsf s^3R_1(R_5-R_6).
\end{aligned}                                                   \tag{EZ16}
\]
All coefficients on the right are infinitesimal: $\mathsf r,
\mathsf s,\mathsf h$ are infinitesimal and $\mathsf c$ is
bounded. It is harmless that $\mathsf c=1$ in the first row.
Equations EZ13--EZ16 therefore put the entire scalar, apart
from its retained positive square, on the six rows $\omega R_i$
with a bounded Gram whose real residue is $\tfrac12$ times the
strict matrix in EZ15.

\subsubsection*{Every complete first normal and all higher terms}

Define $L_x,L_y,L_v$ as half the three derivatives of the
whole polynomial EZ10 at zero normals. Expanding its short
factored form gives exactly
\[
 L_x=\omega uz^2T_x(U,V),\qquad
 L_y=\omega uzT_y(U,V),\qquad
 L_v=\omega uzT_v(U,V),                                       \tag{EZ17}
\]
where each $T_i$ is a polynomial with bounded coefficient arrays.
This supplies a finite construction of every coefficient:
differentiate EZ10, set the three normals to zero, and extract
the displayed monomial. No derivative is dropped.
The exact actual-source contact equations give $T_i(P,G)=0$.
Consequently the monic divisions are
\[
 T_i=(U-P)A_i(U,V)+(V-G)B_i(V).                               \tag{EZ18}
\]
For inverse $b$, each $B_i=\kappa\widetilde B_i$ with
bounded arrays. Indeed the same source substitution used in
EZ14 gives $T_i(0,V)=0$ identically at $\kappa=0$.
Divide the cleared source polynomial by $\kappa$ and then by
the monic $V-G$. This is an exact coefficient construction;
it uses neither $1/u$, $1/z$, nor an evaluated mixed quotient.
For bounded $b$, $\kappa=1$ again suffices.

Combining EZ16--EZ18 proves that both complete state normals
and the complete testing-angle normal have infinitesimal
coefficients on the strict six-row list $\omega R_i$.
The additional factor $z$ in $L_x$ is bounded and infinitesimal.
This argument uses the full first-normal polynomial, not just
its value or linear approximation at contact.

The full normal expansion of EZ10 has a bounded quadratic Gram
whose real residue, in order $(N_x,N_y,N_v)$, is
\[
 \begin{pmatrix}
  \bar\mu+\bar\omega/2&0&0\\
 0&1&1/2\\
 0&1/2&1/2
 \end{pmatrix}\succ0.                                         \tag{EZ19}
\]
For bounded $b$, $\bar\mu=1/2$; for inverse $b$,
$\omega=1$ and $\bar\mu=0$. Both give a strictly positive
first pivot. The other block has pivots $1,1/4$.

Here is also a finite justification for all higher terms, since
the scaled normals $N_y,N_v$ may be unbounded. Expanding EZ10
in the three normals gives degree at most four. A monomial
of normal degree at least three containing $N_x$ has total
$N_y,N_v$ degree at most two. Assign those factors, and any
needed extra $N_x$ factor, to two Gram rows; the remaining
factor contains $N_x\in\mathrm I$. For a higher monomial
without $N_x$, each of the $N_y,N_v$ exponents is at most two,
and its coefficient has a factor $s^2$. If its degree is
three, the unassigned coefficient is a bounded multiple of
$s(sN_y)$ or $s(sN_v)$. If its degree is four, it is a
bounded multiple of $(sN_y)(sN_v)$ (or of one of their
squares). All are infinitesimal by EZ9. These support and
divisibility statements follow directly from the factored
quadratics in EZ10 and are checked on the entire polynomial.
Thus no unbounded scaled normal is hidden in a Gram coefficient.

The scalar, all first normals, and the complete hard block now
give a finite bounded Gram with strictly positive real residue.
The separated strict-Gram lemma supplies tensor SOS.
Restore the whole retained terms EZ7--EZ8 and the positive
square in EZ13, and divide by the pure-source positive $\omega$.
This proves the stated original-parent theorem on both source
rows, without a comparison of source and testing rates.

\paragraph{State and boundary scope.}
The signs of $W$ and of the testing angles were not restricted
in EZ10--EZ19. The polynomial and cube-root identities include
literal zero testing coordinates. If a bounded normalized
$Z=1$ state has $X$ or $Y$ a unit, the original pure-testing
core has positive residue at the double-small angle corner;
its positive source multiple is already a separated square,
and all original event terms are nonnegative squares.
Thus this proof supplies the entire $Z$-unit state chart.
Complex states split into their real and imaginary quadratic
evaluations. Section~\ref{component:p845b3cb8b114} combines this
certificate with the two strong-current outer testing identities
and the complementary $W$-unit chart to cover all tests.



% Source formulas: pabbcd0905dcd
\subsection{The uniform small-$e$ source at unbounded $b$}

\label{component:pabbcd0905dcd}

\label{pabbcd0905dcd:sec:c1-inverse-b-local-new}

\paragraph{Statement and certificate ring.}
Let $K_c,K_t$ be independent real closed extensions of $\mathbb R$,
with bounded rings $B_c,B_t$ and infinitesimal ideals
$\mathfrak m_c,\mathfrak m_t$. The certificate ring is
$\mathcal T=K_c\otimes_{\mathbb R}K_t$, with bounded coefficient ring
$\mathcal B=B_c\otimes_{\mathbb R}B_t$ and zero-residue ideal
$\mathfrak n=\mathfrak m_c\otimes B_t+B_c\otimes\mathfrak m_t$.
We use the finite strict-residue Gram lemma in this ring.
All radicals and inverses below are pure source, pure testing, or
products of separately specified such factors.

Use the actual uniform native C1 stationary source, with
$\lambda=n=0$, $d>0$, and the original nonnegative event coefficients.
Normalize $d=1$ by a positive pure-source scale. Suppose
\[
 k,f,g>0\text{ are infinitesimal},\qquad k/g\in\mathfrak m_c,
 \qquad r=1-af^2,\quad h=1-tg^2,
\]
\[
 R=1+k^2r,\quad T_0=1+k^2h,\quad
 0<a<T_0,\quad 0<t<R,\quad \beta\geq0.                 \tag{IW1}
\]
All original strict gaps remain imposed. Then the complete original
parent is tensor-SOS whenever the independent positive testing angles
$u,v$ are infinitesimal, for every real or complex testing state.
The proof below works in the rotated chart $w=1,z>0$ with $z$ arbitrary;
in particular it includes the previously missing original
$W$-unit, $Z$-infinitesimal chart. It uses no bound on $g/f$ and no
comparison between any testing and source infinitesimals.
The exact sign, zero-axis and state reductions are given at the end.

\subsubsection*{The actual C1 source and bounded source arrays}

For clarity, the six weights used here come from the following
unchanged native formulas. Put $x_n=k^2$, $\eta=f^2$, $\zeta=g^2$,
$y_n=x_n\eta\zeta$, $K=1+x_n$, $G=h+\zeta$, and
\[
 M=R-t,\quad J=T_0-a,\quad E_n=1-x_n\eta h,
 \quad Q_n=1-x_n\zeta r,\quad \theta=t\zeta.
\]
Define
\[
 P_d=1+2x_nh+rx_n^2G+x_n\zeta
       +\eta x_n(K\zeta+x_nh\theta),
\]
\[
 P_4=T_0E_n+x_n(1-r)G,\qquad
 P_\alpha=Rh(1+y_n)+\zeta M E_n.
\]
The actual relative coefficients, now with $d=1$, are
\[
\begin{array}{ll}
 m=\dfrac{Kx_n(r+\eta)h(1+y_n)}{2M\zeta P_d},&
 C=\dfrac{K(r+\eta)E_n}{2(1-r)P_d},\\[4pt]
 A=\dfrac{\eta P_4}{2rP_d},&
 D+\tfrac12=\dfrac{K(1+y_n)(r+\eta)}{2J\eta P_d},\\[4pt]
 \beta+\tfrac12=\dfrac{Kx_n(r+\eta)G}{2\theta P_d},&
 \alpha=\dfrac{x_n^2\zeta(r+\eta)P_\alpha}{2Q_nP_d}.
\end{array}                                                     \tag{IW2}
\]
No auxiliary C2 stationary source is used.
Since $a,t$ are bounded and $k,f,g$ are infinitesimal,
$r,h,P_d,R,T_0$ have residue one. The inequality
$\beta+\tfrac12\geq\tfrac12$ in (IW2) gives
$t=O(k^2/g^2)$. Thus $t$ and $p=t/R$ are infinitesimal, and
$M=R(1-p)$ is a positive unit. The two independent openings
$a$ and $J=T_0-a$ are not capped.

Introduce the exact compact coordinates
\[
 j=\frac{T_0}{g^2R},\quad B=\frac{k^2g^2R}{T_0},\quad
 T=\frac{f^2T_0}{R},\quad p=\frac tR,\quad H=\frac a{T_0},
 \qquad x=jB,\quad y=TB,\quad q_z=jT.                    \tag{IW3}
\]
Here $j$ is unbounded, $0<H<1$, and
\[
 T,x,B,p,\sigma:=1/j,\delta:=j^2B\in\mathfrak m_c,
 \qquad \delta\sim k^2/g^2,\qquad p=O(\delta).
                                                               \tag{IW4}
\]
The quantity $q_z$ can have any positive size. The exact inverse is
\[
\begin{gathered}
 R=\frac{1+x}{1+HTx},\quad T_0=\frac{1+x}{1+Bp},\quad
 r=\frac{1-HT}{1+HTx},\quad h=\frac{j-p}{j(1+Bp)},\\
 f^2=\frac{T(1+Bp)}{1+HTx},\quad
 g^2=\frac{1+HTx}{j(1+Bp)},\quad a=HT_0,\quad t=pR.
\end{gathered}                                                   \tag{IW5}
\]
In particular both strict gaps are retained exactly.

Here is a short rational dictionary sufficient to check every source
estimate. Write $D_+=D+\tfrac12$, $b_+=\beta+\tfrac12$,
$\mu=m(1-p)$, and set
\[
\begin{split}
 P={}&1+B(j+1-p)+BT(1+Hj)\\
 &+B^2T(Hj+jp-p^2+p)+B^2HT^2j(1+Bp),\\
 L={}&1+T(1-H+Bp),\qquad E=1+B(HT+p-1),\\
 F_\alpha={}&j+1-2p+BHTj(1-p)+BTp(j-p).
\end{split}
\]
Substitution of (IW5) in (IW2) gives exactly
\[
\begin{split}
 \mu&=\frac{Bj(j-p)(1+y)(1+Bp)L}{2(1+x)P},\\
 C&=\frac{(1+Bp)L[1+BTp-BTj(1-H)]}{2HT(1+x)P},\\
 A&=\frac{T(1+Bp)[1+BT(2Hj-Hp+H-j+p)+B^2H^2T^2j]}
              {2(1-HT)P},\\
 D_+&=\frac{(1+y)(1+Bp)(1+HTx)L}{2T(1-H)(1+x)P},\\
 b_+&=\frac{Bj(1+Bp)L(j+1-p+HTx)}{2p(1+x)P},\\
 \alpha&=\frac{B^2LF_\alpha}{2EP}.
\end{split}                                                     \tag{IW6}
\]
The denominators $P,E,1+x,1-HT$ are positive bounded units.
After $j=1/\sigma$, $B=\delta\sigma^2$, each of the following
normalized arrays is a rational function with polynomial numerator
and denominator in the bounded variables $\sigma,\delta,T,p,H$;
its denominator has a nonzero real constant at
$\sigma=\delta=T=p=0$, independently of $H$:
\[
 \frac{2\mu}{\delta},\quad\frac{2b_+p}{\delta},\quad
 2TCH,\quad2TD_+(1-H),\quad\frac{2A}{T},\quad
 \frac{2\alpha}{B^2j}\quad\hbox{all have residue }1.
                                                               \tag{IW7}
\]
This finite polynomial criterion proves boundedness and the stated
residues at every relative rate and at both $H$ openings. It also
shows $\beta p\leq b_+p=O(\delta)$, including $\beta=0$.
A useful exact cancellation is
\[
                  K_d:=CH-D_+(1-H)=-\frac{b_+p}{j}.    \tag{IW8}
\]

\subsubsection*{The entire parent and both state normals}

Put $D_u=1+u^2$, $D_v=1+v^2$. The quarter-parent is
\[
\begin{split}
 \mathcal P={}&m(vW-X+uvY-uZ)^2+CD_v(uW-Y)^2+AD_uD_vY^2\\
 &+D_+(uW+uvX-Y-vZ)^2
 +(\alpha+\tfrac12)D_u(vY-Z)^2\\
 &+b_+D_v(X+uZ)^2-\tfrac12D_uD_vZ^2.
\end{split}                                                     \tag{IW9}
\]
On the rotated chart
$w=(W+uY)/D_u=1$, $z=(Z-uX)/D_u>0$, use the two unrestricted normals
\[
 N=\frac{X+uZ-pD_uv}{D_u},\qquad
 M_0=\frac{uW-Y-HD_uvz}{D_u}.
\]
Their exact inverse is
\[
 X=N+pv-uz,\quad Z=z+uN+puv,\quad
 W=1+u(M_0+Hvz),\quad Y=u-M_0-Hvz.                     \tag{IW10}
\]
Only the pure-testing unit $D_u$ is inverted. Retain the positive
term $mpD_u^2N^2$, and put $G=(\mathcal P-mpD_u^2N^2)/D_u$.
Direct substitution, before any state specialization, gives
\[
\begin{split}
 G={}&\mu D_u[N^2-2vN+(1-p)v^2]
      +CD_uD_v(M_0+Hvz)^2+A D_v(u-M_0-Hvz)^2\\
 &+D_+D_u(M_0-(1-H)vz)^2\\
 &+(\alpha+\tfrac12)
       [u(1-p)v-(1+Hv^2)z-uN-vM_0]^2\\
 &+b_+D_uD_v(N+pv)^2-\tfrac12D_v(z+uN+puv)^2.
\end{split}                                                     \tag{IW11}
\]
In particular $G$ is exactly quadratic in the two normals:
\[
 G=F+2L_NN+2L_MM_0+H_NN^2+H_MM_0^2
                         +2(\alpha+\tfrac12)uvNM_0,
                                                               \tag{IW12}
\]
where $F=G|_{N=M_0=0}$ and
\[
\begin{split}
 h_0&=\mu+b_+\geq\tfrac12,\\
 H_N&=h_0D_u+b_+v^2+\beta u^2v^2+\alpha u^2,\\
 H_M&=D_u(CD_v+D_+)+AD_v+(\alpha+\tfrac12)v^2.
\end{split}                                                     \tag{IW13}
\]
The potentially unbounded $\beta$ is fully retained in (IW11)--(IW13).
There is no omitted higher-normal polynomial.

\subsubsection*{An exact separated scalar circuit}

Define the rows
\[
 U=u/\sqrt x,\quad V=v/\sqrt y,\quad Z_*=z/\sqrt{q_z},
 \quad Q_*=U-VZ_*,\quad B_*=V-UZ_*,\quad R_*=Z_*-UV,
\]
\[
 \Phi=U^2+V^2+Z_*^2+U^2V^2Z_*^2-4UVZ_*.
\]
None of $U,V,Z_*$ is assumed bounded. On the positive chart,
\[
 \Phi=(U-V)^2+Z_*^2(1-UV)^2+2UV(1-Z_*)^2             \tag{IW14}
\]
is a legal separated SOS; for example
$\sqrt{UV}=(xy)^{-1/4}\sqrt{uv}$.

For an explicit finite definition of the actual masses, put
$K_0=CH^2+D_+(1-H)^2$ and
\[
\begin{array}{ll}
 M_a=Ax,&M_b=[\mu(1-p)+b_+p^2]y,\\
 M_c=\alpha q_z,&
 M_{bc}=[K_0+AH^2+2H\alpha+H-\tfrac12]yq_z,\\
 M_{b^2}=b_+p^2y^2,&
 M_{b^2c}=H^2(A+C+\alpha+\tfrac12)y^2q_z,\\
 M_{ab^2c}=CH^2xy^2q_z,&M_{ab^2}=\beta p^2xy^2.
\end{array}
\]
Set
\[
\begin{gathered}
 f_c=M_b,\quad h_c=M_{ab^2},\quad
 k_c=(M_a-M_b+M_{ab^2c}-M_{bc}+h_c)/2,\\
 r_c=M_c-M_b-k_c,\quad q_c=M_a-f_c,\quad
 g_c=M_{ab^2c},\quad d_c=M_{b^2c}-k_c,
 \quad e_c=b_+p^2y^2-g_c.
\end{gathered}
\]
The complete scalar from (IW11) satisfies the exact identity
\[
\begin{split}
 F={}&q_cQ_*^2+r_cR_*^2+d_cV^2Q_*^2+e_cV^2B_*^2\\
 &+f_c\Phi+g_cV^2\Phi+k_cZ_*^2(V^2-1)^2+h_cV^2R_*^2.
\end{split}                                                     \tag{IW15}
\]
The last term is the extra nonnegative square supplied by the actual
$\beta$ event. 

The same finite bounded-array criterion as (IW7), with the displayed
mass definitions, gives
\[
 \frac{2f_c}{Tx^2}\sim1,\quad
 \frac{2q_c}{Tx}\sim1,\quad
 \frac{2r_c}{TjB^2}\sim1,\quad
 \frac{k_c}{T^2B^2j}\in\mathfrak m_c,\quad
 \frac{g_c}{f_c}\in\mathfrak m_c.                     \tag{IW16}
\]
Here $\sim1$ means residue one. In particular the three principal
masses $f_c,q_c,r_c$ are positive, and
$|k_c|/r_c$, $|k_c|/(q_cy)$ are infinitesimal.
All normalized mass denominators remain units after
$j=1/\sigma$, $B=\delta\sigma^2$; the companion finite-array
checks test their residues independently of $H$.

For the signed terms use the exact identities
\[
\begin{split}
 (1+U^2)\Phi&=R_*^2+B_*^2+U^2Q_*^2+U^2(1-UVZ_*)^2,\\
 (1+V^2)\Phi&=Q_*^2+R_*^2+V^2B_*^2+V^2(1-UVZ_*)^2,\\
 Z_*(V^2-1)&=-VQ_*-R_*.
\end{split}                                                     \tag{IW17}
\]
The $e_c,g_c$ combination in (IW15) is exactly
\[
 b_+p^2y^2V^2B_*^2-g_c\Phi
       +g_c[Q_*^2+R_*^2+V^2(1-UVZ_*)^2].
\]
A $k_c$ of either sign costs at most
$3|k_c|V^2Q_*^2+2|k_c|R_*^2$, including its occurrence in $d_c$.
The corresponding relative costs are
$3|k_c|v^2/(q_cy)$ and $2|k_c|/r_c$, both in $\mathfrak n$.
Every inequality here abbreviates a finite SOS identity, using
$2A^2+2B^2-(A+B)^2=(A-B)^2$ and nonnegative pure-source factors.
Consequently the scalar supplies strict fixed real fractions of
\[
                       q_cQ_*^2,\qquad r_cR_*^2,
                       \qquad f_c\Phi.                \tag{IW18}
\]
All other nonnegative terms remain available as whole squares.

\subsubsection*{Finite budgets for the complete normal rows}

Call a row free if any fixed positive real multiple of its square can
be subtracted at infinitesimal cost from (IW18), by a finite separated
SOS identity. This definition allows a bounded normal coupling to be
absorbed by a sufficiently large fixed real helper pivot.
Put
\[
 q=vz-Tu,\quad b=v-uz/j,\quad R_z=z-uv/B,\quad
 R_u=v(u^2-x),\quad R_v=u(v^2-y),
\]
\[
 J_u=u(u^2-x),\qquad J_v=v(v^2-y),\qquad
 \widehat q=q/\sqrt T.
\]
Since $Q_*=-q/(T\sqrt x)$, the first mass in (IW18) has a positive
unit coefficient on $\widehat q^2$.

First, $R_u=\sqrt y\,x(-B_*-UR_*)$. Applying the first identity
of (IW17) to $B_*^2$, and the actual $r_c$ pivot to $U^2R_*^2$,
shows that $R_u$ is free. The three relative costs are bounded by
constant multiples of
\[
                         B,\qquad u^2/j,\qquad u^2,
                                                               \tag{IW19}
\]
all infinitesimal. Second,
$R_v=\sqrt x\,y(-Q_*-VR_*)$, so $\sqrt T R_v$ is free;
its two costs are $O(y^2)$ and $O(Tv^2)$.

The same-angle rows have exact decompositions
\[
\begin{split}
 J_v&=\sqrt y\,v(v+\sqrt{y/x}\,u)(V-U)+(y/x)R_u,\\
 J_u&=\sqrt x\,u(u+\sqrt{x/y}\,v)(U-V)+(x/y)R_v.
\end{split}                                                     \tag{IW20}
\]
If $\rho=O(\delta)$ is a nonnegative source factor, then
$\rho J_v$ is free. Indeed $\rho\sqrt{y/f_c}=O(\sqrt\delta)$,
$\sqrt{y/x}=\sqrt{T/j}$ is bounded, and
$\rho y/x$ is infinitesimal. Thus every coefficient after division
by the strict $\sqrt{f_c}(V-U)$ pivot is infinitesimal.
Also $\sqrt T\,TxJ_u$ is free: the two coefficients of its
first summand relative to $\sqrt{f_c}(U-V)$ are bounded multiples of
\[
                  T\sqrt x\,u^2,\qquad\sqrt{T\delta}\,uv,
\]
and its second summand is exactly
$\delta(\sqrt T R_v)$, because $Tx^2/y=j^2B=\delta$.
This last use of $\delta\in\mathfrak m_c$ is the additional
source fact specific to the inverse-$b$ sector.

Finally,
\[
 j\alpha b=-(\alpha/B)R_u-\alpha uR_z.                \tag{IW21}
\]
Here $\alpha/B=O(x)$, and the square of $\alpha uR_z$ costs
$O(x^2u^2)$ relative to $r_cR_*^2$. Hence the entire left side
of (IW21) is free. These identities involve no bound on $z$.

\paragraph{The unrestricted first normal.}
Define the bounded source quantities
\[
 K_z=H(\alpha+\tfrac12)-\tfrac12,\qquad
 K_R=-\mu-\alpha(1-p)-\tfrac12+b_+p+TK_z,
\]
\[
 \rho=b_+p+\beta px\geq0.
\]
Direct differentiation of the entire (IW11), followed by the actual
source contact identity, gives
\[
 L_N=-j\alpha b+(K_R+\beta pv^2)R_u
                   +\rho J_v+K_zuvq.                 \tag{IW22}
\]
For completeness, before imposing the source, the remainder is
$v S_N$, with
\[
 S_N=j\alpha-\mu+b_+p+K_Rx+b_+py+\beta pxy=0
\]
identically under (IW6). Thus no constant or first-normal term is
silently dropped. Equations (IW7) imply
$K_R\sim-\tfrac12$, $\rho=O(\delta)$ and
$\beta p=O(\delta)$. The first three terms in (IW22) are free by
(IW19)--(IW21); the last has coefficient
$K_z\sqrt T\,uv\in\mathfrak n$ on $\widehat q$.
Dividing by $\sqrt{h_0}$ preserves all these bounds because
$h_0\geq\tfrac12$.

\paragraph{The unrestricted second normal.}
Use the exact $K_d$ from (IW8), and put
\[
\begin{split}
 C_q={}&K_dD_u+AH+\alpha+\tfrac12
        +v^2[CHD_u+AH+(\alpha+\tfrac12)H],\\
 c_2={}&T[CH+AH+(\alpha+\tfrac12)H]
                      -A-(\alpha+\tfrac12)(1-p).
\end{split}
\]
The complete second normal has the exact division
\[
 L_M=C_qq+(c_2+TCHu^2)R_v+T(K_d+CHy)J_u.             \tag{IW23}
\]
Before substituting the actual source its only remainder is $uS_M$,
where
\[
\begin{split}
 c_0&=T[K_d+AH+\alpha+\tfrac12]-A,\\
 S_M&=c_0+TK_dx+yc_2+TCHxy=0.
\end{split}
\]
The last equality is an exact identity from (IW6). Furthermore
$c_2$ is bounded (in fact infinitesimal), $TCH$ is a unit, and
$(K_d+CHy)/x$ is bounded with residue $-\tfrac12$.
It follows from the two free rows in (IW20) that
\[
 \sqrt{TH(1-H)}[(c_2+TCHu^2)R_v+T(K_d+CHy)J_u]
\]
is free with bounded coefficients. For the other term,
\[
 \sqrt{TH(1-H)}C_qq
     =T\sqrt{H(1-H)}C_q\widehat q.
\]
Its coefficient is in $\mathfrak n$: multiplying $C_q$ by $T$
clears the only $CH$ pole, and every term then contains either
$T$ or $v^2$. This statement is made on the full expression in
(IW23), before any relative source/testing rate is considered.

\subsubsection*{Strict full Gram and state restoration}

Normalize the hard rows by pure-source factors,
\[
 \widehat N=\sqrt{h_0}\,N,\qquad
 M_0=\sqrt{TH(1-H)}\,\widehat M.
\]
Their full hard matrix is bounded: $b_+/h_0\leq1$,
$\beta/h_0\leq1$, $\alpha/h_0=O(\alpha)$, and the two
source poles of $H_M$ are exactly cleared by $TH(1-H)$.
Equations (IW7), (IW13) give its real residue
\[
                           \operatorname{diag}(1,\tfrac12).
                                                               \tag{IW24}
\]
Its off-diagonal coefficient has the factor
$uv\sqrt{TH(1-H)}/\sqrt{h_0}$ and is infinitesimal.
The $\widehat q$ couplings in (IW22)--(IW23) are infinitesimal.
All remaining complete-normal couplings are bounded coefficients
of free rows. Buy sufficiently large fixed real pivots on the
finitely many free rows, use the elementary real Young square
completion, and apply the strict-residue Gram lemma to the
remaining finite matrix. The strict scalar pivots (IW18) lose only
infinitesimal amounts. This proves $G\in\Sigma\mathcal T^2$.
No state normal has been set to zero in this conclusion.
Restoring $mpD_u^2N^2$, the pure-testing factor $D_u$, and the
positive source scale $d$ proves the assertion for (IW9).

The proof has used $u,v>0$ and $w=1,z>0$, with arbitrary $z$.
The exact sign/state lemma SJ1--SJ10 applies to the same original
core: sign gauges make $u,v$ nonnegative while preserving every
term; a negative core at positive angles forces the rotated
$w,z$ to have the same nonzero sign. Dividing all state coordinates
by the nonzero pure-testing $w$ gives precisely this chart.
When the core is nonnegative it is itself a pure-testing square,
and every event term is already positive. Zero angles have the
direct core SOS exits. The zero state is immediate, and a complex
state splits into its real and imaginary quadratic evaluations.
Thus (IW1) is proved for all states at the double-infinitesimal
angles. The existing two-current outer-angle identities give the
separate all-tests join if needed; they are not required as a premise
of the local conductor proved here.



% Source formulas: p845b3cb8b114
\subsection{The uniform small-$e$ testing and state partition}

\label{component:p845b3cb8b114}

\label{p845b3cb8b114:sec:uniform-c1-complete-current-join}

\paragraph{Theorem and unchanged source scope.}
Use the original physical canonical C1 stationary source with
$\lambda=n=0$, $d>0$, every original nonnegative event coefficient,
and every original strict source gap. Suppose
\[
 k,f,g,e=kf/g>0\text{ are infinitesimal},\qquad
 r=1-af^2,\quad R=1+k^2r,\quad T_0=1+k^2-tk^2g^2,
\]
\[
                         a,t>0,\qquad J=T_0-a>0,\quad M=R-t>0.
                                                               \tag{UJ1}
\]
Then the entire original parent is a finite tensor sum of squares
for every independent projective testing pair and every real or
complex testing state. There is no restriction on $b=g/k$, $g/f$,
or any relative infinitesimal rate. In particular the original
$\beta=0$ face is retained.

This is a mathematical join of the actual bounded-$b$ endpoint
product-repair proof, the whole $Z$-chart proof EZ1--EZ19,
the inverse-$b$ local proof IW1--IW24, and the two full-parent
outer-current identities below. No unavailable EZ or IW
coefficient witness is used as a premise.

\subsubsection*{Exact original-source bounds for the two outer identities}

Normalize $d=1$ by a positive pure-source scale. The retained
Section~\ref{component:p7d4e5164e84a} prints seven original coquotients.
Their respective residues on (UJ1) are
\[
 (c_2,c_\beta,c_3,c_5,c_4,c_\alpha,c_d)
                    \sim(2,2,2,2,2,2,4).
\]
Their arguments are bounded, so they are positive source units.
The actual weights satisfy
\[
 C=\frac{c_3}{af^2c_d},\qquad
 D=\frac{B_c}{2Jf^2c_d},\qquad A=\frac{f^2c_4}{c_d},
 \qquad B_c=2c_5-Jf^2c_d\sim4.
\]
Consequently
\[
 C^{-1},D^{-1}\in\mathfrak m_c,\qquad
 g_B=\frac{CD}{C+D}>0,\qquad
 \gamma=A g_B=
 \frac{c_4c_3B_c}{c_d^2(aB_c+2Jc_3)}\sim\frac14.
                                                               \tag{UJ2}
\]
This is uniform in both openings $a,J$: the last denominator has
residue $16\cdot4(\bar a+\bar J)=64$, since $a+J=T_0\sim1$.
No positive lower bound on either opening is assumed.
These bounds use the actual original $C,D,A$, not selected
substitute weights. Any other positive event is restored whole.

\paragraph{Complementary testing angles.}
The actual source estimates UJ2 verify the hypotheses of
Lemmas~\ref{lem:c1-two-current} and~\ref{lem:c1-critical-current}.
For $v$ infinitesimal and $1/u$ bounded, $\bar\gamma=1/4$
and $\bar c_u^2<1$, so the latter real residue condition is strict.
Their whole homogeneous identities include both projective endpoints.
Only the two infinitesimal testing angles remain.
\subsubsection*{Every remaining source and state chart}

At $u,v$ infinitesimal normalize a nonzero testing state by a
coordinate of largest absolute value. This is a pure-testing
operation; all four coordinates are bounded and at least one
is a unit. If $X$ or $Y$ is a unit, the original pure-testing
core has positive real residue $2(X^2+Y^2)$. It is a positive
testing square, and every original event remains a positive
separated square. We can therefore assume $X,Y$ infinitesimal.

If $Z$ is a unit, normalize $Z=1$. The remaining $W$ is arbitrary
bounded and may have either sign. EZ1--EZ19 applies in each of
its two actual-source coordinate rows:
\[
\begin{array}{c|c|c}
 \text{source}&\text{exact coordinates}&\text{actual proof}\\ \hline
 b=g/k\text{ bounded}& f=eb,\quad g=kb&
                     \mathrm{EZ1\!-\!EZ19,\ bounded\ }b\\
 b=g/k\text{ unbounded}&
        (k,f,g),\quad k/g\in\mathfrak m_c&
                     \mathrm{EZ1\!-\!EZ19,\ inverse\ }b.
\end{array}                                                     \tag{UJ5}
\]
The extra inverse-$b$ source inequality used by EZ is derived from
the actual $\beta\geq0$: $tb^2\leq2c_\beta/c_d$. Thus it adds no
source hypothesis to (UJ1). Its full proof retains both state
normals and the testing-angle normal, including unbounded scaled
normal coordinates through their exact $sN_y,sN_v$ coefficients.

If $Z$ is infinitesimal, $W$ must be a unit. Normalize $W=1$;
then $u,v,X,Y,Z$ are all infinitesimal.
For bounded $b$, use the actual retained
Section~\ref{component:p05180564ed05} proof. Its source variables are exactly
$(e,k,b,a,t)$ with $f=eb$, $g=kb$; its gaps
$R_c=1+k^2(1-ae^2b^2)$ and
$T_0=1+k^2-tk^4b^2$ coincide with (UJ1).
It allows every bounded positive $b$, including infinitesimal $b$,
both strict-gap openings, and all original coefficient faces.
Its whole seven-term polynomial, complete first normals,
higher-normal support and nine-root product identity are printed
with Section~\ref{component:p7d4e5164e84a} and
Section~\ref{component:p935c9d6318ea}. Its final original-parent
restoration is an exact sum of complete positive terms.
There is no remaining condition on $k^3/e$ or $k/e$ in that proof.

For unbounded $b$, use IW1--IW24. Its exact source hypotheses
$k,f,g,k/g\in\mathfrak m_c$, physical $\beta\geq0$, and the
original gaps follow directly from (UJ1) and this ratio class.
It proves the complete original parent with both unrestricted
state normals for every positive rotated state ratio, so it
includes the present $W$-unit chart. The retained $\beta$ term,
including its zero face, and both source openings are part of
that proof. The state and sign reduction is the actual
SJ1--SJ10 identity, not a chart label.

These cases exhaust every source ratio and every state.
The zero state is immediate. A complex state splits into its
real and imaginary evaluations because the parent matrix is real.
Undo all pure-testing normalizations and the positive source
scale $d$. This proves the theorem on (UJ1) at all independent
projective tests.



% Source formulas: pf5d2f33f694c
\subsection{Every positive $e$ in the uniform source}

\label{component:pf5d2f33f694c}

\label{pf5d2f33f694c:sec:uniform-c1-all-e-rederived}

Use the original canonical C1 source and the exact native transport
EN13, with every original coefficient sign and strict gap retained.
Suppose
\[
 k,f,g>0\text{ are infinitesimal},\qquad e=kf/g>0 .
 \tag{UE1}
\]
There is no additional restriction on the size of \(e\), either
source ratio \(g/k,g/f\), or either separately labelled testing field.
Then the entire original parent is tensor SOS at every independent
projective testing pair and every real or complex testing state.

Here are the actual proof inputs and the complete join.  EN13 gives
\(0<a,t<2\), \(r=1-af^2\sim1\), and the original strict
\(M=R-t>0,J=T-a>0\).  Split the positive pure-source \(e\) into
\[
 e\text{ infinitesimal}\quad\hbox{or}\quad \Delta=1/e\text{ bounded}.
 \tag{UE2}
\]
In the first case UJ1--UJ5 is the full-parent theorem.  It
includes both source ratios \(g/k\) bounded and unbounded, both
state charts, both complete state normals and all angular/projective
joins; the actual bounded-ratio product-repair, EZ and IW proof bodies are its inputs.

In the second case use the actual retained
Section~\ref{component:p41b39e2f94e1}, Theorem \emph{Uniform C1, bounded
reciprocal \(e\)}.  Its hypothesis is exactly UE1 with
\(\Delta\) bounded, and its proof is complete in that body.
For clarity, the scope extension used there is substantive.
After its pure-source substitution
\[
 \delta=\Delta R/(1+k^2),\quad b^2=f^2(1+k^2)/R,
\]
both \(\delta\) and the pure-testing rotated state coordinate \(w\)
are allowed to be bounded units.  In its polynomial ring the weights
\[
 (\delta,k,b,a,t,u,w,v)=(0,2,6,0,0,2,0,16)
 \tag{UE3}
\]
make the full scalar error and both complete first normals have weights
at least \(42,32,26\).  The displayed unlocalized monic conductor,
restricted to \((A_m,B_m,C_m)^2+(D_0^2)\), gives a scalar Gram
whose every correction coefficient has weight at least \(2\);
the first-normal coefficient bounds are \(12,6\).
The exact hard block on its two normalized state normals has
residue \(16I_2\), including every quadratic term.  Its explicit
eleven-soft-root identity supplies a strict positive diagonal.
Thus the entire parent, with all four retained terms restored,
has a bounded strict-residue Gram.  This is a finite polynomial
division proof; neither \(\delta^{-1}\) nor a mixed element is a
lifting coefficient.

The same body proves the complementary \(W=1,Z\) infinitesimal
chart with five exact current roots.  Its two full correction
coefficients contain respectively \(\Delta Z\) and
\(\Delta^2 f^2Z^2\), so vanish in residue even when \(\Delta\)
is a unit.  The original source signs and both strict gap openings
remain intact.

For the outer angle joins, no conclusion from the old infinitesimal-\(e\)
assertion is imported.  Use the two exact whole-parent identities in
Lemmas~\ref{lem:c1-two-current} and~\ref{lem:c1-critical-current}.
Their source hypotheses hold on all of UE1.  Indeed the seven native
coquotients have residues \(2,2,2,2,2,2,4\), regardless of \(e\).
Writing \(B_c=2c_5-Jf^2c_d\sim4\), the original weights after
normalizing \(d=1\) satisfy
\[
 C^{-1},D^{-1}\text{ infinitesimal},\qquad
 A\frac{CD}{C+D}
 =\frac{c_4c_3B_c}{c_d^2(aB_c+2Jc_3)}\sim\frac14.
 \tag{UE4}
\]
Here \(a+J=T\sim1\), so the denominator has residue \(64\)
without an individual lower bound on either opening.
Those identities therefore cover every testing chart outside
\(u,v\) infinitesimal, including literal projective infinities.
Literal axes, the zero state, signs and complex states are covered
by their exact full-parent normalization identities.

The two alternatives UE2 exhaust every positive source \(e\).
All scalings are within one labelled field, and every subtracted
term is restored.  This proves UE1 at all tests, with literal
\(\beta=0\) and every allowed original coefficient boundary.
The two real closed extensions are arbitrarily chosen labelled
factors; no algebraic-disjointness assumption is imposed.



% Source formulas: p9afbe73899af
\subsection{The source with infinitesimal $k,\zeta$ and unbounded $\eta$}

\label{component:p9afbe73899af}

\label{p9afbe73899af:sec:r7-c1-small}
This section concerns the original C1.1 source, not an enlarged family.
All coefficients are divided by $d>0$. Put
\[
 b=\eta^{-1/2},\quad g=\sqrt\zeta,\quad q=k^2\eta,
 \qquad x=b^2q,\quad a=b^2(1-r).
 \tag{SC1}
\]
Here $b,g$ are positive infinitesimals, $q>0$ is bounded, and
$\overline q\le1/2$. The new bounded two-normal reduction below is
asserted for $q$ a positive unit. The exact polynomial identities
remain valid more generally, but their indicated boundedness argument
would not control $g/\sqrt q$ when $q$ is infinitesimal.
No all-tests closure is inferred from the reduction.

Keep $K,R,T,M,J,E,Q,P_4,P_d,P_\alpha$ from (NC1)--(NC2), and set
\[
\begin{aligned}
 c_b&=b^2C/d=\frac{K(1+b^2r)E}{2(1-r)P_d},&
 a_b&=b^2A/d=\frac{P_4}{2rP_d},\\
 D_b&=D/(db^2),&\alpha_b&=\alpha/(db^2g^2),\\
 h_2&=\frac{Kq(1+b^2r)(1-tg^2)(1+qg^2)}{2P_d},&
 h_b&=\frac{Kq(1+b^2r)[1+(1-t)g^2]}{2P_d}.
\end{aligned}
\]
Thus $m/d=h_2/(g^2M)$ and $\beta/d+1/2=h_b/(g^2t)$.
The exact relation
\[
 h_b-h_2=\frac{Kq(1+b^2r)E g^2}{2P_d}
 \tag{SC2}
\]
removes the apparent opening dependence from the first normal of the
paired $m,\beta$ currents. The quantities $D_b,\alpha_b,h_2,h_b$ are
bounded; $\alpha_b,h_2,h_b$ are units in the stated $q$-unit domain.
Their residues are respectively $(1-2q)/2,q^2/2,q/2,q/2$, with
bounded source letters on the right understood at residue.

\paragraph{Exact full-parent reduction.}
For bounded affine tests define
\[
 Q_b=(W+uY)/b,\quad Y_b=Y/b^2,\quad p=t/R,\quad
 N=\frac{(X+uZ)/b-pvQ_b}{g},
\]
\[
 H=uQ_b/b,\quad L=vQ_b/g,\quad V=uvZ/b,\quad
 B_b=gN+pvQ_b,\quad D_u=1+u^2,\quad D_v=1+v^2.
\]
All roots use only pure-source inverses. Retain the two nonnegative
original terms
\[
 \frac{4h_2t}{MR}N^2,
 \qquad \frac{4h_bM}{tR}D_vN^2.
\]
The complete original parent divided by $db^2$, minus these terms,
is exactly
\begin{align*}
 G={}&\frac4R\{(h_2+h_bD_v)N^2+2(h_bD_v-h_2)LN
              +[h_2(1-p)+h_bpD_v]L^2\}\\
 &+4D_v\{c_b(H-D_uY_b)^2+a_bD_uY_b^2\}
   +4D_bJ_b^2+4\alpha_b g^2D_u(vb^2Y_b-Z)^2+2K_b,
                                                        \tag{SC3}\\
 J_b={}&b^2(H-D_uY_b)+ubvB_b-D_uvZ,\\
 K_b={}&b^2(H-D_uY_b)^2+2buv(H-D_uY_b)B_b+u^2v^2B_b^2\\
 &+D_uv^2b^2Y_b^2+D_uV^2
       -2D_uV(Q_b+vB_b)+2D_uu^2vZY_b.
\end{align*}
These are entire quadratic polynomials in the state; no restriction
of a state normal has been made. In particular, (SC3) is a direct
input for collecting both complete first normals and their full hard
quadratic, rather than an identity solely on a contact graph.

\paragraph{A harmonic normal removes both $r$ openings.}
Put
\[
 s=c_b+a_b,\quad \rho=c_b/s,\quad h=c_ba_b/s,
 \quad \rho_0=Kr/R,\qquad
 N_y=\sqrt{s}\,(D_uY_b-\rho_0H).
\]
Then the complete $C,A$ line in (SC3), divided by $4D_v$, equals
\[
 \left[N_y+\sqrt{s}(\rho_0-\rho)H\right]^2+hH^2
 -\frac{a_bu^2}{D_u}
           \left(\frac{N_y}{\sqrt{s}}+\rho_0H\right)^2.
 \tag{SC4}
\]
The inverse $1/D_u$ is pure testing, and every coefficient in (SC4)
is bounded uniformly through either $r\to0$ or $r\to1$.
Indeed $s$ is bounded below, $h$ is a positive unit, and
$a_b/s$, $a_b\rho_0/\sqrt{s}$ and $a_b\rho_0^2$ are bounded.
There is no mixed harmonic denominator.

Here is the exact source cancellation that also centers the normal
without an opening loss. Define
\[
 S_0=Kr(1+b^2r)E+(1-r)P_4,
\]
\[
 T_0=q^2(1-tg^2)^2+qrg^2+2qtg^2-qg^2-2q+r.
\]
Then
\[
 \rho=\frac{Kr(1+b^2r)E}{S_0},\qquad
 Kr-\rho R=-\frac{b^2r(1-r)K T_0}{S_0},\qquad
 h=\frac{K(1+b^2r)E P_4}{2P_dS_0}.             \tag{SC5}
\]
Both $S_0$ and the other denominators are positive units. Also
$s\asymp[r(1-r)]^{-1}$, so $\sqrt{s}(\rho_0-\rho)$ is
infinitesimal without a lower bound on either opening.
At the actual contact,
\[
 (u,v)=(b\sqrt q,\sqrt q g),\quad
 Q_b=RgZ,\quad Y_b=r\sqrt q gZ,
\]
and both $N,N_y$ vanish exactly. Thus (SC3)--(SC5) remove every
unbounded source weight while retaining the actual contact and every
state normal.

\paragraph{The leading angular circuit.}
For $u,v$ infinitesimal, the leading scalar after the harmonic state
normal has been removed is twice
\[
 (1-q)H^2+qL^2+V^2+(1-2q)(vZ)^2+q^2(gZ)^2-2VQ_b.
\]
Here the leading relation $H(vZ)=VQ_b$ is exact for the displayed
roots. It splits as
\[
 (1-2q)(H-vZ)^2+
 qH^2+qL^2+V^2+q^2(gZ)^2-4qVQ_b.              \tag{SC6}
\]
The final four positive terms form one AM--GM circuit, since
\[
 (\sqrt qH)(\sqrt qL)V(qgZ)=(qVQ_b)^2.
\]
On the positive negative-core chart its radical factors separate
between the source and testing fields. This proves positivity of the
leading scalar, including $q=1/2$, where its separate binomial drops
out. The leading hard normal coefficients are $4q$ on $N$ and $4$
on $N_y$. This organizing form alone is not yet a certificate for
all the higher terms in (SC3).

A useful exact conductor can be stated without comparing any source
and testing rate. Normalize $Z=1$ on its pure-testing state chart,
put $c=\sqrt q$, $w=(W+uY)/R$, and now use the contact-exact roots
\[
 H_0=uw/b^2,\ L_0=vw/(bg),\ V_0=uv/b,\ D_0=c^2g,
 \quad A_0=cH_0,\ B_0=cL_0,\ C_0=V_0.
\]
Their contact is $(u,v,w)=(bc,cg,bg)$. Set
\[
 E_0=H_0-v,\quad
 S_v=A_0-D_0-cE_0,\quad
 S_u=C_0-D_0-S_v,\quad S_w=B_0-D_0-S_v.
\]
Direct multiplication gives
\[
\begin{aligned}
 H_0(u-bc)&=(u-bc)E_0+bS_u,\\
 H_0(v-cg)&=(v-cg)E_0+(v/c)S_v,\\
 H_0(w-bg)&=(w-bg)E_0+(bg/c)S_w.
\end{aligned}                                                     \tag{SC7}
\]
Only the positive-unit inverse $1/c$ occurs. These identities lift
contact-vanishing coefficients times $H_0$ to the displayed root
module with bounded coefficients on the small pure-testing chart.
They do not alone prove a full scalar square lift: a first-jet
syzygy among the redundant angular roots still has to be allocated,
particularly when the coefficient of $E_0^2$ vanishes at $q=1/2$.

\paragraph{Exact small-cost angular pivots.}
The four-term circuit has the exact symmetric decomposition
\[
 \sum_{i=1}^4A_i^2-4\sqrt{A_1A_2A_3A_4}
 =\frac13\sum_{i<j}(A_i-A_j)^2
  +\frac23\sum_{\{ij\}\mid\{kl\}}
       (\sqrt{A_iA_j}-\sqrt{A_kA_l})^2,
\]
where the second sum is over the three pair partitions. Use
$(A_1,A_2,A_3,A_4)=(A_0,B_0,C_0,D_0)$ above and assume the positive
chart $u,v,w>0$. Then the following identities are exact:
\[
\begin{aligned}
 c\sqrt{uw}\,E_0
 &=\sqrt{uw}(A_0-D_0)
      +b\sqrt{cg}(\sqrt{A_0D_0}-\sqrt{B_0C_0}),\\
 \sqrt{vw}(u-bc)
 &=\frac{b^{3/2}}{\sqrt c}
             (\sqrt{A_0C_0}-\sqrt{B_0D_0}),\\
 \sqrt{uw}(v-cg)
 &=b\sqrt{g/c}
             (\sqrt{B_0C_0}-\sqrt{A_0D_0}),\\
 \sqrt{uv}(w-bg)
 &=\frac{b^{3/2}\sqrt g}{c}
             (\sqrt{A_0B_0}-\sqrt{C_0D_0}).
\end{aligned}                                                     \tag{SC9}
\]
The coefficient factors on the right are infinitesimal; all radicals
separate between the source and testing fields. Consequently any fixed
positive real multiples of the squares of the four left-hand roots can
be inserted as additional positive Gram pivots at an infinitesimal cost
in the displayed strict circuit Gram. To see this without a limiting
assertion, write the four identities as $P=TZ$, with $Z$ the displayed
circuit roots and every entry of $T$ in the bounded maximal ideal. Then
$Z^TGZ=\tau P^TP+Z^T(G-\tau T^TT)Z$, and the second matrix retains
the strictly positive real residue of $G$ for every fixed real $\tau>0$.
This is a genuine full-rates Gram mechanism. It still requires showing
that the entire first-normal and scalar remainders of (SC3) allocate on
these weighted roots; the identities do not make that last membership
automatic.

\paragraph{A quantified remote improvement when $e$ is unbounded.}
The source identity $e^2=q/g^2$ and (NI3) imply that $m,\beta$ are
both unbounded when $e$ is unbounded. Three inverse-cycle rotations
send this adjacent pair to $C,A$ and transform tests as
$(u,v)\mapsto(v,1/u)$. Therefore the adjacent-event theorem covers
all tests with $v$ not infinitesimal. Combined with the previous
$C,A$ certificate, only $u,v$ infinitesimal remain in this $e$ type.
This includes every $q$-unit source considered above.

There is also a complete state-chart certificate there. Normalize the
pure-testing state by a coordinate of maximal absolute value. If
$W$ is a unit, cap $C,A,m,\beta$ at a sufficiently large positive
real $\mu$, retain all excesses, and leave every other event as its
original nonnegative square. At $u=v=X=Y=0$, the quadratic part in
$(u,v,X,Y)$ is
\[
 2(uW-Y)^2+2(X+uZ)^2-4uvWZ+
 4\mu\{(uW-Y)^2+Y^2+(vW-X-uZ)^2+(X+uZ)^2\}.
 \tag{SC8}
\]
The event part is strictly positive for every real $W\ne0$ and bounded
real $Z$: its zero equations give successively $Y=u=X=v=0$. A sufficiently
large finite $\mu$ therefore absorbs the entire signed quadratic core.
The exact remainder has degree at least three in these four small
roots, with bounded coefficients in $W,Z$, so it allocates with zero
residue. This proves the entire original parent on this chart.
If the core has negative sign and $W$ is infinitesimal, then $Z$ is
a unit, since otherwise the core residue forces $X=Y=0$, contrary
to maximal-coordinate normalization. Thus the remaining state chart
has $Z=1$, $W$ infinitesimal, and both testing half-angles infinitesimal.
This is a full quantified narrowing, not closure of C1.1.

\paragraph{Completion and exact outstanding obligation.}
The following section supplies a finite scalar square lift and both
complete first normals for the $q$-unit C1.1 subdomain, including
$D=0$. It uses a contact-centered $D$ term followed by the geometric
mean conductor; an unweighted $E_0^2$ perturbation by itself is not
controlled by the circuit uniformly at arbitrary relative rates.
For $q$ infinitesimal, the uniform normal and protected contact lift
in Sections~\ref{pc376284fc7ae:sec:r7-c1-q-inf-local}--\ref{pd2ecc6ac1891:sec:r7-c1-q-inf-remote}
resolve the additional $g/\sqrt q$ issue and all remote tests.
No assumption of bounded $e$ or $1/e$ is imposed. 



% Source formulas: p95350f69a353
\subsection{Positive-unit $q=x\eta$}

\label{component:p95350f69a353}

\label{p95350f69a353:sec:r7-c1-q-unit}
\begin{theorem}
On the original C1.1 source, suppose $q=k^2\eta$ is a positive bounded
unit. The entire original parent is tensor SOS at every independent
projective test and every real or complex state. All positive opening
rates and the permitted literal boundary $D=0$ are included.
\end{theorem}
The proof uses the exact full quadratic (SC3), not its leading scalar.
The $q$-infinitesimal part of C1.1 is not included in this theorem.

\paragraph{One more centered term.}
Use the variables of (SC1)--(SC5), and put $s_D=J/R$. Retain the
additional nonnegative square
\[
                  4D_b(J_b+s_DH)^2.                 \tag{SC10}
\]
It vanishes at actual contact, where $H=RvZ$ and $J_b=-JvZ$.
Its removal replaces $4D_bJ_b^2$ in (SC3) exactly by
\[
                  -8D_bs_DHJ_b-4D_bs_D^2H^2.
\]
In particular it removes the entire separate $(vZ)^2$ term.
This step allows $D_b=0$ literally and does not divide by $D_b$.

Set $Z=1$ on the remaining state chart. Put $c=\sqrt q$ and
$w=(W+uY)/R$, and use the four separated positive monomials
\[
 A=\frac{cuw}{b^2},\quad B=\frac{cvw}{bg},\quad
 C=\frac{uv}{b},\quad D=c^2g,\quad
 \Gamma=(ABCD)^{1/4}.
 \tag{SC11}
\]
Their exact contact is $(u,v,w)=(k,cg,bg)$, with $k=bc$, and
$A=B=C=D$ there. Notice that $H=RA/c$, $L=RB/c$, $V=C$,
$g=D/q$, and $Hv=R\Gamma^2/q$.

Here are explicit formulas for the entire scalar at $N=N_y=0$,
after (SC10). Write $a_t=1-\rho_0+pv^2$. Define
\begin{align*}
 P_H={}&4D_v\{c_b(1-\rho_0)^2+a_b\rho_0^2/D_u\}
       +2b^2\{a_t^2+v^2\rho_0^2/D_u\}\\
 &+4\alpha_bg^2b^4v^2\rho_0^2/D_u
       -8D_bs_Db^2a_t-4D_bs_D^2,\\
 P_L={}&4[h_2(1-p)+h_bpD_v]/R,\\
 P_{Hv}={}&-4[D_u(1+pv^2)-u^2\rho_0]
             -8\alpha_bg^2b^2\rho_0+8D_bs_DD_u.
\end{align*}
The exact scalar is
\[
 f=f_AA^2+f_BB^2+f_CC^2+f_DD^2-f_\Gamma\Gamma^2,
 \tag{SC12}
\]
where
\[
 f_A=R^2P_H/q,\quad f_B=R^2P_L/q,\quad
 f_C=2D_u,\quad f_D=4\alpha_bD_u/q^2,
 \quad f_\Gamma=-RP_{Hv}/q.
\]
These are rational functions of $u,v$ over bounded source coefficients,
with only positive-unit denominators. For example,
\[
 c_b(1-\rho_0)^2=
   \frac{K(1+b^2r)E(1-r)}{2P_dR^2},\qquad
 a_b\rho_0^2=\frac{P_4K^2r}{2P_dR^2}.
\]
Thus neither $r$ opening survives in their denominators.
For $u,v,b,g$ infinitesimal their residues are
\[
               \overline f_A=\overline f_B=
 \overline f_C=\overline f_D=2,\qquad
               \overline f_\Gamma=8.               \tag{SC13}
\]
This uses the exact source residues in (SC2), including
$\overline D_b=(1-2\overline q)/2$. It does not assume a positive
lower bound on $D_b$.

Here is an explicit finite denominator certificate used below.
Every function just displayed is $P(u,v)/(1+u^2)$ for a polynomial
$P$ with bounded pure-source coefficient array. For a monomial use
\[
 x^m-x_0^m=(x-x_0)\sum_{j=0}^{m-1}x^{m-1-j}x_0^j.
\]
Applying this identity to each monomial supplies finite first and
second difference arrays with bounded source coefficients. For the
denominator use only the exact identity
\[
 \frac1{1+u^2}-\frac1{1+k^2}
 =-\frac{(u-k)(u+k)}{(1+u^2)(1+k^2)}.
\]
Thus the first and second finite Taylor arrays about $(k,cg)$
are bounded and use only pure-testing $1+u^2$ and pure-source
$1+k^2$ denominators, both units. The same construction after
$u=U^2,v=V^2$ gives the half-variable arrays below. This is a
finite polynomial identity construction, not an analytic series.

\paragraph{Every first normal and the complete hard block.}
Let $\delta_\rho=\rho_0-\rho$. The full first-normal coefficients
of the remaining quadratic, with no normal suppressed, are
\begin{align*}
 \mathcal L_N={}&\frac8R(h_bD_v-h_2)L
       +4buvg(a_t-2D_bs_D)H-4D_uvgV,\\
 \mathcal L_y={}&8D_v\sqrt{s}\,\delta_\rho H\\
 &+\frac H{\sqrt s}\left[
 -\frac{8D_va_bu^2\rho_0}{D_u}+8D_bs_Db^2
 +\frac{8\alpha_bg^2b^4v^2\rho_0}{D_u}
 -4b^2a_t+\frac{4b^2v^2\rho_0}{D_u}\right]\\
 &+\frac{4bu}{\sqrt{s}}V
       -\frac{8\alpha_bgb^2v}{\sqrt{s}}(gZ).
                                                        \tag{SC14}
\end{align*}
Each is a bounded infinitesimal-coefficient linear combination of
$A,B,C,D$. At actual contact it is zero exactly, since the source
is a stationary zero and every removed term vanishes there.
The coefficients of $N^2,N_y^2,NN_y$ are respectively
\[
\begin{aligned}
 H_{NN}&=4(h_2+h_bD_v)/R+2g^2u^2v^2,\\
 H_{yy}&=4D_v\left(1-\frac{a_bu^2}{D_us}\right)
       +\frac{4\alpha_bg^2b^4v^2}{D_us}
       +\frac{2b^2}{s}\left(1+\frac{v^2}{D_u}\right),\\
 H_{Ny}&=-4buvg/\sqrt{s}.
\end{aligned}                                                     \tag{SC15}
\]
Thus the hard Gram has residue $\operatorname{diag}(4q,4)$.
In (SC15), $H_{Ny}$ denotes the polynomial coefficient of $NN_y$,
not the off-diagonal matrix entry.

\paragraph{A finite geometric-mean square lift.}
Here is a detailed lift for the particular raw module (SC12).
It is not a claim about arbitrary contact-vanishing polynomials.
Define $\alpha_i=A_i-\Gamma$, with
$(A_1,A_2,A_3,A_4)=(A,B,C,D)$, and put
\[
 T_4=\sum_iA_i^2-4\Gamma^2.
\]
A fixed positive Gram on the following finite roots represents $T_4$:
\[
 \alpha_i,\qquad
 Q_{ij|kl}=\sqrt{A_iA_j}-\sqrt{A_kA_l},\qquad
 \sqrt\Gamma(\sqrt{A_i}-\sqrt{A_j}),\qquad
 \sqrt\Gamma((A_iA_j)^{1/4}-(A_kA_l)^{1/4}).          \tag{SC16}
\]
Use all four $\alpha_i$, all three pair partitions, and all six
unordered differences in the third group. To verify strictness, average
positive multiples of the symmetric circuit decomposition in (SC9) and
\[
 T_4=\sum_i\alpha_i^2+2\Gamma\left(\sum_iA_i-4\Gamma\right).
\]

over the three decompositions
\[
 \sum_iA_i-4\Gamma=(\sqrt{A_i}-\sqrt{A_j})^2
 +(\sqrt{A_k}-\sqrt{A_l})^2
 +2((A_iA_j)^{1/4}-(A_kA_l)^{1/4})^2.
\]
The additional differences $A_i-A_j$ in the first decomposition are
$\alpha_i-\alpha_j$ and give nonnegative corrections to the already
strict $\alpha$ block.

Write $\epsilon_i=f_i-2$, $\epsilon_\Gamma=f_\Gamma-8$, and
$\sigma=\sum_i\epsilon_i-\epsilon_\Gamma$. Then
\[
 f-2T_4=\sum_i\epsilon_i\alpha_i^2
       +2\Gamma\sum_i\epsilon_i\alpha_i+\Gamma^2\sigma.
 \tag{SC17}
\]
All $\epsilon$ have zero residue. The exact contact value implies
$\sigma(k,cg)=0$. Clear the pure-testing denominator $D_u$ if needed;
monic division gives
$\sigma=(u-k)\sigma_u+(v-cg)\sigma_v$ with bounded coefficients.
For the monomials (SC11), the exact conductors are
\[
\begin{aligned}
 \Gamma(u-k)&=\sqrt{uk}\,Q_{13|24},\\
 \Gamma(v-cg)&=\sqrt{vcg}\,Q_{23|14},\\
 \Gamma(w-bg)&=\sqrt{wbg}\,Q_{12|34}.
\end{aligned}                                                     \tag{SC18}
\]
Thus the last term in (SC17) becomes $\Gamma$ times a linear
combination of the $Q$ roots with infinitesimal bounded coefficients.
The exact identity
\[
 Q_{ij|kl}-\tfrac12(A_i+A_j-A_k-A_l)
 =-\tfrac12(\sqrt{A_i}-\sqrt{A_j})^2
  +\tfrac12(\sqrt{A_k}-\sqrt{A_l})^2               \tag{SC19}
\]
changes these terms into infinitesimal quadratic corrections on (SC16)
and $\Gamma\sum_i\mu_i\alpha_i$, where each $\mu_i$ is infinitesimal
and is rational in $\sqrt u,\sqrt v$ with bounded coefficient arrays
and positive-unit denominators.

The exact three contact jets imply that the four source values
$\mu_i(k,cg)$ coincide. Indeed the differential vectors of
$\alpha_i/\Gamma$ at contact are the four rows
\[
 (1/2,-1/2,1/2),\quad(-1/2,1/2,1/2),\quad
 (1/2,1/2,-1/2),\quad(-1/2,-1/2,-1/2)
\]
in logarithmic coordinates $(u/k,v/(cg),w/(bg))$. Their sole linear
relation is the sum of all four rows. Cancelling the nonzero source
factors in this equality establishes coincidence; no inverse of a
small number enters a Gram coefficient.
Call the common infinitesimal value $\mu_*$. The common term is
\[
 \mu_*\Gamma\sum_i\alpha_i
   =\mu_*\Gamma\left(\sum_iA_i-4\Gamma\right),
\]
an infinitesimal correction on the explicitly displayed variance roots.
For each remaining $\mu_i-\mu_*$, monic division in $\sqrt u,\sqrt v$
has bounded coquotients. Apply the half-variable conductors
\[
\begin{aligned}
 \Gamma(\sqrt u-\sqrt k)
   &=(uk)^{1/4}\sqrt\Gamma
          ((AC)^{1/4}-(BD)^{1/4}),\\
 \Gamma(\sqrt v-\sqrt{cg})
   &=(vcg)^{1/4}\sqrt\Gamma
          ((BC)^{1/4}-(AD)^{1/4}).
\end{aligned}                                                     \tag{SC20}
\]
Their coefficient factors are infinitesimal. Multiplying by the
remaining $\alpha_i$ proves that the entire scalar correction in
(SC17), not merely its low-order jet, is a bounded zero-residue
quadratic form on (SC16). This uses only two finite monic divisions.
The boundedness of their coefficient arrays follows from the explicit
rational functions (SC12), whose test denominators are powers of
$D_u$, and from the displayed square-root substitutions; bounded
function values alone are not being used as a substitute for this fact.

The same argument for a first normal is shorter. If
$\mathcal L=\sum_i a_iA_i$, write
$\mathcal L=\sum_i a_i\alpha_i+\Gamma\sum_i a_i$.
Its exact contact zero gives $\sum_i a_i(k,cg)=0$; one monic division
and (SC18) allocate it on (SC16), with infinitesimal coefficients.
Apply this to both complete expressions (SC14).
Together with (SC15), this yields a bounded strictly positive-residue
Gram for the entire original remainder. Restore (SC10), the paired
terms, and all positive pure-source and pure-testing normalizations.

\paragraph{Exhaustion and legal radicals.}
For $q$ a unit, $e^2=q/g^2$ is unbounded. The exact adjacent-event
rotations and the $W$-unit state argument (SC8) therefore reduce every
uncovered negative-core test to $u,v,W$ infinitesimal and $Z$ a unit.
Literal axes are direct original-core exits. Use exact sign gauges and
normalize $Z=1$. The common pure-testing Schur identity forces
$W+uY>0$ on the remaining negative-core chart: with
$w_t=(W+uY)/D_u$, the eliminated scalar is
\[
 \frac{uvD_uD_v}{\mathcal D}
       (D_uuv w_t^2-2D_uD_vw_t+uv),\qquad
 \mathcal D=u^2v^4+u^2v^2+u^2+v^4+2v^2+1>0,
\]
which is nonnegative when $w_t\le0$; its eliminated two-square
matrix is strictly positive. Hence $w=(W+uY)/R>0$.
Every monomial in (SC11), and every product under a radical in
(SC16)--(SC20), factors into a positive pure-source quantity and a
positive pure-testing quantity. The occurrence of $R$ in $w$ is a
positive pure-source unit, not a mixed denominator. All test
coefficients remain bounded because $u,v,w$ are infinitesimal.
If the general toric contact lemma is used, its third coordinate
is literally the pure-testing $W+uY$, with pure-source contact
$Rbg$; this direction is inactive in every coefficient function.
Complex state evaluations split into real and imaginary parts.
This supplies every independent test and proves the theorem.



% Source formulas: p7d43f3fccb7d
\subsection{A uniform normal at infinitesimal $q$}

\label{component:p7d43f3fccb7d}

This intermediate lemma by itself is not an certificate of the full
$q$-infinitesimal C1.1 source. The following two completion sections
supply the scalar allocation and remote joins. Keep (SC1)--(SC3), but do
not discard the whole second paired term. Its addition to the
$N^2$ coefficient in (SC3) gives exactly
\[
 \frac4R(h_2+h_bD_v)+\frac{4h_bM}{tR}D_v
       =H_0+\frac{4h_b}{t}v^2,
 \qquad H_0=\frac{4h_2}{R}+\frac{4h_b}{t}.
\]
Retain the last summand as the positive square
$4h_b(vN)^2/t$, and use the pure-source normalization
\[
                         \widehat N=\sqrt{H_0}\,N.
\]
The exact original inequality $\beta\ge0$ gives
$h_b/t\ge g^2/2$. The explicit formula for $h_2$ gives
$h_2/q$ a positive bounded unit, uniformly through both openings.
Consequently $H_0\ge c_0(q+g^2)$ for some positive real $c_0$.
This argument also permits $\beta=0$ literally. It does not require
$e=\sqrt q/g$ or its reciprocal to be bounded.

The coefficient of $\widehat N^2$ from this source part is exactly
one. Its remaining contribution $2g^2u^2v^2/H_0$ is bounded and
infinitesimal on the small-angle testing chart. Moreover
\[
 \frac{g}{\sqrt{H_0}},\quad
 \frac{\sqrt q}{\sqrt{H_0}}
\]
are bounded. Since $h_b-h_2=O(qg^2)$, the complete first normal
(SC14), after this normalization, is an infinitesimal-coefficient
linear combination of $B,H,V$ for $u,v$ infinitesimal. The first
coefficient is bounded because
\[
 \frac{h_bD_v-h_2}{\sqrt{qH_0}}
            =O(g^2+v^2),
\]
and the other two coefficients are bounded by constant multiples of
$buv$ and $v$. The hard cross coefficient with $N_y$ is likewise
infinitesimal.

The issue left by this lemma is scalar and angular: $H$ is unweighted in this
statement, whereas the four-term circuit controls $\sqrt q H$.
One must retain a strict $D$-based contact normal and allocate every
unweighted $H$ coefficient on that normal and the circuit. The
naive full extraction (SC10) fails in this branch: its raw $P_H$
satisfies the exact formal specialization
\[
       P_H\big|_{q=0,\,v=0}=-\frac{2ru^2}{1+u^2},
\]
so $P_H/q$ need not be bounded. Thus the corrected hard-normal
normalization removes the former $g/\sqrt q$ obstruction. The later
protected-square lift establishes the missing scalar allocation.



% Source formulas: pc376284fc7ae
\subsection{Infinitesimal $q$ at the two small testing angles}

\label{component:pc376284fc7ae}

\label{pc376284fc7ae:sec:r7-c1-q-inf-local}
\begin{theorem}
On the original C1.1 source with $q=k^2\eta$ infinitesimal, the entire
original parent is tensor SOS whenever both independent testing
half-angles $u,v$ are infinitesimal. All original opening rates and
$\beta=0$ are retained. The boundary $D=0$ is empty here because
$D_b$ has residue $1/2$.
\end{theorem}
This statement includes every state; it does not yet assert the
ordinary or unbounded $v$ testing charts.

Use the preceding source notation and put $c=\sqrt q$, $e=c/g$,
$h=H/R$, and $E_*=h-v$. Work first with $Z=1$ and $u,v,h>0$.
Define
\[
 A=ch,\quad B=cL/R,\quad C=uv/b,\quad D=qg,\quad
 \Gamma=(ABCD)^{1/4},\quad d_i=A_i-\Gamma.
\]
Then
\[
 \Gamma^2=qhv,\qquad BC=e v^2h,\qquad
 u h=uE_*+bC.
 \tag{QI1}
\]
Here $h=u(W+uY)/(Rb^2)$ is separated into a positive pure-source
factor and a positive pure-testing factor. Every radical below is
therefore legal; no mixed inverse or order comparison is used.
The exact original inequality $\beta\ge0$ gives $t=O(e^2)$, while
$t$ is bounded. Hence $t/e$ and $\sqrt{t/e}$ are bounded.

\paragraph{Exact scalar, with the $D$ term restored.}
Write $U=u^2,V=v^2$, $D_u=1+U,D_v=1+V$, and
$a_t=1-\rho_0+(t/R)V$. At both state normals zero the entire original
scalar is
\[
 a h^2+b_*hv+c_*v^2+P_L L^2+2D_uC^2+4\alpha_b g^2D_u,
 \tag{QI2}
\]
where
\begin{align*}
 a={}&4D_vR^2\{c_b(1-\rho_0)^2+a_b\rho_0^2/D_u\}
 +2b^2R^2\{a_t^2+V\rho_0^2/D_u\}\\
 &+4\alpha_b g^2b^4VR^2\rho_0^2/D_u
 +4D_bb^4R^2a_t^2,\\
 b_*={}&-4R\{D_u(1+tV/R)-U\rho_0\}
 -8\alpha_b g^2b^2R\rho_0-8D_bb^2Ra_tD_u,\\
 c_*={}&4D_bD_u^2.
\end{align*}
All these rational functions have bounded source coefficient arrays;
only $D_u$ and positive pure-source units occur in their denominators.
The $r$ openings cancel exactly as in (SC12).
In particular
\[
 D_b=\frac{N_D}{2JP_d},\qquad
 N_D=\frac{K(1+qg^2)(1+b^2r)-JP_d}{b^2}
\]
has a polynomial numerator: the numerator before division vanishes
identically at $b^2=0$, where $P_d=1+qg^2$. This is a finite monic
division, so it introduces no unbounded source coefficient array.
Also $\alpha_b/q^2$, $h_2/q$, and $h_b/q$ have explicit bounded
arrays from their source formulas. These observations justify the
$q$ coquotients below uniformly through both source openings.

Use formal specialization at $q=0$ only to construct finite polynomial
coquotients; it is not a new physical source. Put $z=b^2$,
$J_0=1-z(1-r)$, and $S=1-r+tV$. The three specialized coefficients are
\begin{align*}
 a_0={}&2D_v\{(1+zr)(1-r)+r/D_u\}
       +2z\{S^2+Vr^2/D_u\}+2z^2S^2/J_0,\\
 b_0={}&-4\{D_u(1+tV)-Ur\}-4zSD_u/J_0,\\
 c_0={}&2D_u^2/J_0.
\end{align*}
Define $\widetilde b_0=b_0+4tD_uV$, and the bounded finite
coquotients
\[
 X=a_0+\widetilde b_0+c_0=U X_U+V X_V,
 \qquad Y=-\widetilde b_0-2c_0=U Y_U+V Y_V.
 \tag{QI3}
\]
For definiteness, first subtract the value at $U=0$, divide by $U$,
and then divide that value by $V$. Their two $V$ coquotients are
\[
 X_V=2+2zr(1-2t)+\frac{2z(1+zr)t^2V}{J_0},
 \qquad Y_V=\frac{4zt}{J_0}.
 \tag{QI4}
\]
Thus $X_V=2+\varepsilon_V$, $\varepsilon_V\in B_{\rm Inf}$;
$X_U,Y_U$ have bounded arrays. Likewise write
$a=a_0+qa_1$, $b_*=b_0+qb_1$, $c_*=c_0+qc_1$ by finite numerator
division. All $a_1,b_1,c_1$ have bounded arrays. Direct source
specialization gives
\[
 \overline{a_1-c_1}=2,\qquad
 \overline{b_1+2c_1}=-8,\qquad
 \overline{R^2P_L/q}=2,\qquad
 \overline{4\alpha_bD_u/q^2}=2.
 \tag{QI5}
\]

The following is the exact scalar identity to be allocated:
\begin{align*}
 (\mathrm{QI2})={}&(c_0+qc_1)E_*^2+2D_uC^2
 +X_U(uE_*+bC)^2+Y_Uu(uE_*+bC)E_*\\
 &+X_V(vh)^2+Y_Vv(vh)E_*
 -4(t v/e)D_uBC\\
 &+(a_1-c_1)A^2+(b_1+2c_1)\Gamma^2
 +(R^2P_L/q)B^2+(4\alpha_bD_u/q^2)D^2.
 \tag{QI6}
\end{align*}
It uses only (QI1) and $qv^2=2\Gamma^2-A^2+qE_*^2$.
In particular the apparently dangerous negative $t$ term is already
an infinitesimal-coefficient $BC$ term, since $v$ is infinitesimal.

\paragraph{Two extra contact squares and their finite first jets.}
Set
\[
 F_*=vh-gD,\quad
 R_F=\frac{\sqrt{g\Gamma}}{c}(\Gamma-D),\quad
 V_0=cv,\quad R_E=\sqrt\Gamma(\sqrt A-\sqrt{V_0}).
\]
All four roots vanish at actual contact. Exact multiplication gives
\begin{align*}
 (vh)^2&=F_*^2+2R_F^2+2gF_*d_4+g^2D^2-4g^2\Gamma d_4,\\
 c\Gamma E_*&=2\Gamma d_1-R_E^2,\\
 \Gamma F_*&=-2g\Gamma d_4+R_F^2/g.
 \tag{QI7}
\end{align*}
The root $R_E$ receives a strict real square pivot from $E_*^2$:
\[
 qE_*^2=(A-V_0)^2
       =4R_E^2+(\sqrt A-\sqrt{V_0})^4.
 \tag{QI8}
\]
Reserve this fraction $q$ of the strict $E_*$ pivot; the remaining
coefficient still has positive real residue. The leading $2(vh)^2$
in (QI6), by (QI7), gives strict $F_*$ and $R_F$ pivots, with only
infinitesimal quadratic corrections and the printed first-jet term.
The remaining leading angular form is $2(\sum A_i^2-4\Gamma^2)$,
with the strict finite toric square list from (SC16).

For precision, the only additional bare-$\Gamma$ normal couplings
left by (QI6), apart from ordinary toric first jets, have the form
$\Gamma E_*a_*+\Gamma F_*b_*'$ with
\[
 a_*=bu(2X_U+Y_U)+Y_Vvg,
 \qquad b_*'=2\varepsilon_Vg.
\]
At contact $u=bc,v=cg$ these satisfy
\[
 a_{*0}/c=b^2(2X_{U0}+Y_{U0})+Y_{V0}g^2\in B_{\rm Inf},
 \qquad b_{*0}'/g=2\varepsilon_{V0}\in B_{\rm Inf}.
\]
Their differences from contact have finite bounded monic Taylor
arrays. The conductors (SC18) allocate those differences against
$E_*,F_*$ with infinitesimal coefficients. Apply (QI7) to their
constant terms. What remains is a bounded infinitesimal toric
quadratic plus $\Gamma\sum\mu_i d_i+\Gamma^2\sigma$.
The exact scalar value and three original first jets imply its
contact value and first-jet balances. The same two finite divisions
as in (SC17)--(SC20), or the Hadamard contact-lift identity, allocate
it on the strict toric root list. No infinite series or unsupported
ideal-membership assertion is used: (QI3), the $q$ coquotients,
and the two displayed contact coefficient ratios are the finite
certificates needed for the additional roots.
Every coefficient function in these Taylor steps depends only on
$u,v$. Thus only the first two conductors in (SC18)--(SC20) are
used; the third state-coordinate jet supplies the inactive Hadamard
balance. No bound on $(W+uY)/R$ after $Z$ normalization is assumed.

\paragraph{Both complete state normals.}
Use $\widehat N=\sqrt{H_0}N$, $H_0=4h_2/R+4h_b/t$, as in the
preceding uniform-normal lemma. Add back the centered $D$ term
from (SC10), including its normals. Its scalar root is exactly
\[
 J_*=(T+b^2tV)h-D_uv
 =E_*+b^2c(1-tg^2)A-buC+b^2tvF_*+b^2tvgD.
 \tag{QI9}
\]
Its full root is $J_*+ubvgN-b^2N_y/\sqrt s$.
All additional hard coefficients are infinitesimal after the
$\widehat N$ normalization.

For the first $N$ normal in (SC14), its $H$ term is a bounded
infinitesimal multiple of $u h=uE_*+bC$; its other terms are
infinitesimal multiples of $B,C$. For the harmonic normal, multiply
its $H$ coefficient by $\sqrt s$. After the exact opening
cancellations, its value at $q=U=V=0$ is zero. Finite numerator
coquotients therefore give
\[
 \sqrt s\,[H\text{-coefficient in }\mathcal L_y]
       =qK_q+UK_U+VK_V
 \tag{QI10}
\]
with bounded arrays. Use
\[
 qH=RcA,\quad u^2H=Ru(uE_*+bC),\quad
 v^2H=Rv(F_*+gD).
\]
The two remaining terms of $\mathcal L_y$ have infinitesimal
coefficients on $C,D$ since $\alpha_b/q^2$ is bounded. Thus both
complete first normals, including (QI9), are infinitesimal
combinations of $E_*,F_*,A,B,C,D$. Their exact contact zeros make
the coefficient sum of $A,B,C,D$ zero at contact. One finite monic
division and (SC18) allocate these on the same strict square list.
The hard block has residue $\operatorname{diag}(1,4)$ and all its
other entries have zero residue. This proves a bounded positive
residue Gram for the entire original parent, after restoring all
retained positive terms.

\paragraph{States and signs.}
The restriction $Z=0$ is positive by the original pure-testing state
restriction. For $Z\ne0$, divide the state by the pure-testing $Z$;
its size is immaterial, since every coefficient array above depends
only on $u,v$. On the negative-core chart, exact sign gauges give
$u,v>0$ and the common pure-testing Schur identity forces
$W+uY>0$. Hence $h>0$, as required. The other core-sign chart and
literal testing axes are direct original exits. Complex states
split into real and imaginary parts. This completes the asserted
small-$u$, small-$v$ theorem.



% Source formulas: pd2ecc6ac1891
\subsection{Infinitesimal $q$ at the other Bob testing angles}

\label{component:pd2ecc6ac1891}

\label{pd2ecc6ac1891:sec:r7-c1-q-inf-remote}
\begin{theorem}
The original C1.1 source with $q=k^2\eta$ infinitesimal is tensor SOS
at every independent test. Together with the $q$-unit theorem this
closes the entire original C1.1 source.
\end{theorem}
The common $C,A$-large certificate already handles $u$ not
infinitesimal. The preceding local theorem handles $u,v$
infinitesimal. Here are the two remaining $v$ charts. Every positive
opening rate and the literal permitted coefficient boundary is kept.

\paragraph{Ordinary $v$: the full normal residue is strict.}
Let $v$ be a positive bounded unit, use $Z=1$ on the negative-core
chart, and retain the roots of (QI1). Normalize $N$ by $H_0$ as
before, but now restore the positive term $4h_b(vN)^2/t$.
Write
\[
 \tau=\operatorname{res}t,\qquad
 \rho=\operatorname{res}(t/e^2),\qquad
 a^2=\frac{\tau}{2(1+\tau)},\quad
 \gamma^2=\frac{\rho}{2(1+\tau)}.
\]
The original $\beta\ge0$ equation implies $0\le\rho\le1$;
$0\le\tau\le1$ follows from $M=R-t>0$ and $R\sim1$.
The entire residue on $(\widehat N,B,C)$ has diagonal entries
\[
 1+\frac{v^2}{1+\tau},\qquad 2(1+\tau v^2),\qquad 2,
\]
and polynomial cross coefficients
\[
 4av^2,\qquad -4\gamma v,\qquad -4\sqrt{\tau\rho}\,v.
\]
Its exact square decomposition is
\begin{align*}
 &2(C-\gamma v\widehat N-\sqrt{\tau\rho}\,vB)^2
       +\widehat N^2+2B^2\\
 &\qquad +(1-\rho)v^2
    \left(\frac{\widehat N}{\sqrt{1+\tau}}+
                         \sqrt{2\tau}\,B\right)^2.
 \tag{QR1}
\end{align*}
Thus it is strictly positive definite, also at $\rho=1$, the
possible $\beta=0$ endpoint, and at $\tau=0$ or $1$.
The other residue blocks are exactly
\[
                       2E_*^2+2(vh)^2+4D_vN_y^2.
\]
There are no nonzero leading cross entries with these blocks.

For completeness, the boundedness check uses the entire raw scalar
(QI2) and both full normal formulas (SC14), with (QI9) restored.
Terms containing $u h$ use $u h=uE_*+bC$; all their coefficients
are infinitesimal. The $N$--$B$ and $N$--$C$ coefficients with
nonzero residue are precisely those retained in (QR1). The harmonic
normal has infinitesimal coefficient on $h$ for bounded $v$:
after the opening cancellation its coefficient lies in
$(b^2,q,u^2)$ with bounded arrays on this chart. Its other
coefficients are infinitesimal as well.
All remaining source-constant terms and weak monomials enter the
same strict Gram through the exact bounded pure-testing identities
\[
                   h=(vh)/v,\qquad
                   1=((vh)-vE_*)/v^2.
\]
For example $A=ch$ and $D=qg$ are infinitesimal-coefficient
combinations of these two roots. Consequently the actual full Gram
has the positive residue just displayed. Restore all retained
terms. The same sign, $Z=0$, and complex-state joins as in the
local theorem apply; no restriction on the size of the normalized
state was used.

\paragraph{Unbounded $v$: a five-root endpoint frame.}
Put $w=1/v$, so $w$ is infinitesimal, and multiply the original
parent by $w^2$. The exact source formulas give positive real
constants $\mu,d_0$ with
\[
 C/d\ge\mu/b^2,\qquad A/d\ge\mu/b^2,\qquad
 D/d\ge d_0 b^2.
\]
For instance sufficiently small fixed constants below $1/4$ work,
since $b^2C/d,b^2A/d$ are uniformly bounded below and
$D/(db^2)\sim1/2$. Retain all excesses and the entire $m,\alpha,
\beta$ terms. No assumption on $e$ is needed.
Use the five roots
\[
              P=uW/b,\quad Q=Y/b,\quad X,\quad J_b'=bZ,
                     \quad V_b=uwZ.
\]
The homogenized signed core is exactly
\begin{align*}
 K={}&V_b^2-2buV_bQ-2wPJ_b'+2wV_bX+b^2w^2P^2
       +2buwPX-2b^2w^2PQ\\
 &+(D_u+w^2)(X^2+b^2Q^2)-2buwXQ.
 \tag{QR2}
\end{align*}
The three selected event terms are
\[
 4\mu D_w(P-Q)^2+4\mu D_uD_wQ^2,
 \qquad
 4d_0(b^2wP+buX-b^2wQ-J_b')^2.
 \tag{QR3}
\]
Combining these with $2K$, the bounded Gram has strictly positive
residue
\[
 4\mu\{(P-Q)^2+Q^2\}+4d_0(J_b')^2+2X^2+2V_b^2.
\]
All other displayed coefficients have zero residue. This is a full
state-quadratic certificate, not a restriction to a contact graph.
It includes the homogeneous endpoint $w=0$ directly and uses no
inverse involving independent source and testing variables.

The infinitesimal, unit, and unbounded $v$ charts exhaust all tests
remaining after the common $C,A$ certificate. Hence the entire
$q$-infinitesimal C1.1 source, and consequently C1.1 itself, is closed.



% Source formulas: p69995d829bb2
\subsection{Unit $\eta,e,r$}

\label{component:p69995d829bb2}

\begin{theorem}
Use the unchanged physical native C1 source, $c=2$, $\lambda=n=0$,
with all original event signs and positive openings $a,t,M,J$.
Suppose $k,\zeta$ are positive infinitesimals, and
\[
 \eta,\quad r=1-a\eta,\quad e=k\sqrt{\eta/\zeta}
\]
are positive bounded units. Then the complete original parent is tensor
SOS at every independent projective test and every real or complex state.
The openings may be arbitrarily small, and $\beta=0$ is included.
\end{theorem}

\paragraph{The double-infinity corner.}
With $p=1/u$, $q=1/v$, the exact homogeneous parent is
\begin{align*}
F_{\infty\infty}={}&4m(pW-pqX+Y-qZ)^2
 +4c(1+q^2)(W-pY)^2+4A(1+p^2)(1+q^2)Y^2\\
&+(4D+2d)(qW+X-pqY-pZ)^2
 +(4\alpha+2d)(1+p^2)(Y-qZ)^2\\
&+(4\beta+2d)(1+q^2)(pX+Z)^2
 -2d(1+p^2)(1+q^2)Z^2.
\end{align*}
At $Z=1$, after retaining $\alpha,\beta$, its complete quadratic part is
\[
 4cW^2+2dX^2+4D(X-p)^2+4AY^2+(4m+2d)(Y-q)^2.
 \tag{BE1}
\]
In particular the $(Y,q)$ determinant is $4A(4m+2d)>0$.
This is an exact full-parent expansion.

An even shorter full-parent certificate retains every event except
$c,A$. The remaining polynomial belongs to $(W,X,Y,qZ)^2$, with bounded
coefficients at $p,q=0$. Its residue Gram is the quadratic form
\[
 4cW^2+2dX^2+4AY^2+2d(Y-qZ)^2.                 \tag{BE2}
\]
It is strictly positive when $c/d,A/d$ are bounded below by positive
units. All cross terms of the complete polynomial are allocated on
these four roots, so (BE2) handles arbitrary states and both literal
projective endpoints, without a normal restriction.

\paragraph{Uniform lower bounds and outer tests.}
Write $H=r+\eta$. The short coquotients have residues $2rH$ for
$c_2,c_b,c_3,c_5,c_\alpha$, $2$ for $c_4$, and $4r$ for $c_d$.
The source identities imply
\[
 \frac md=\frac{e^2c_2}{M\eta c_d},\quad
 \frac cd=\frac{c_3}{a\eta c_d},\quad
 \frac Ad=\frac{\eta c_4}{c_d},\quad
 \frac Dd=\frac{2c_5-J\eta c_d}{2J\eta c_d}.
 \tag{BE3}
\]
Every ratio is bounded below by a positive real number. For the last,
$J+a=T\sim1$ and
\[
 \overline{2c_5-J\eta c_d}=4\bar r(\bar H-\bar J\bar\eta)
 =4\bar r>0.
\]
No reciprocal of an opening is required to be bounded. Choose positive
real lower bounds for the four ratios, cap these four coefficients at
the corresponding multiples of $d$, and retain their nonnegative
excess event squares. All selected-parent Gram coefficients divided
by $d$ are then bounded.

The physical-source theorem covers every noncommuting residue test.
For $u$ infinitesimal and $v$ a bounded unit, the selected residue form is
\[
 4m(vW-X)^2+2d(1+v^2)(X^2+Y^2)
 +4(c+A)(1+v^2)Y^2+4D(Y+vZ)^2\succ0.
\]
For $u=\infty$ in residue and $v$ a bounded unit, retain the signed
core at that commuting residue and use the four positive squares
$cW^2$, $AY^2$, $m(vY-Z)^2$, $D(W+vX)^2$; their common kernel is zero.
For $v=\infty$ in residue and $u$ a bounded unit, use the core
$2d(1+u^2)(X^2+Y^2)$ and the $c(uW-Y)^2,D(uX-Z)^2$ terms.
The $(0,\infty)$ corner has the strictly positive state residue
\[
 4mW^2+2dX^2+\{4(c+A)+2d\}Y^2+4DZ^2.
\]
The $(\infty,\infty)$ corner is (BE2).
At $(\infty,0)$ retain all events except $c,A$ and use the exact
four-root allocation on
\[
 (pX+Z,\,W,\,Y,\,vX),\qquad p=1/u.
\]
Its residue is
\[
 2d(pX+Z)^2+2d(W+vX)^2+4cW^2+4AY^2\succ0.
\]
The full selected polynomial lies in the square of this displayed
root ideal; the pure-testing change $Z=(pX+Z)-pX$ makes this immediate.
The other coefficients have zero residue. These homogeneous
identities include literal infinities. Literal zero axes have a
nonnegative pure-testing core and are immediate original-event SOS.
Only bounded $u$ and infinitesimal $v$ remain.

\paragraph{State entry and the $W$ endpoint.}
Normalize the state by a maximal coordinate. Set
$D_u=1+u^2$ and $w=(W+uY)/D_u$. On the negative pure-core branch,
if $w$ is infinitesimal then $Z$ is a unit: otherwise the residue core
$2\{(X+uZ)^2+(uW-Y)^2\}$ forces every normalized state coordinate
to have zero residue. Normalize $Z=1$; the exact pure-testing Schur
identity in (SD5) gives $w>0$ after the standard angle gauges.

If $w$ is a unit and $u$ is a unit, a zero of that core residue has
$Y=uw\ne0$. The positive lower bound for $A/d$ gives a strict state
residue after clipping the positive events. Thus only $u,v$
infinitesimal and $W\sim w$ a unit need an additional state check.
When $Z$ is infinitesimal, retain all events except clipped $m,A$ and
use the roots $(X,Y,uW,vW)$. Since $Z/W$ is a pure-testing
infinitesimal, the complete selected polynomial has residue
\[
 2d\{X^2+(uW-Y)^2\}+4m(vW-X)^2+4AY^2\succ0.
 \tag{BE4}
\]
Every remaining coefficient contains $u$, $v$ or $Z/W$; there is no
constant or first-normal term outside this frame. This includes
literal $Z=0$. If $Z$ is a unit, normalize $Z=1$ and use the next
raw scalar frame with $w$ a unit.

\paragraph{The exact remainder and its complete normals.}
Use exactly (SD3)--(SD8), their four positive terms, and the positive
clearing factor $\mathcal C$. Put
\[
 \delta=\Delta R/K,\quad E=b^2=\eta K/R,\quad
 A_0=wv,\ B_0=\delta v,\ C_0=\delta Euw,
 \ D_0=\delta^2k^3E,\quad x=B_0-C_0.
\]
Here $\delta,E$ are positive bounded units, and $a,t$ are bounded.
The factored five-term scalar is in
$(A_0,B_0,C_0)^2+(D_0^2)$, and the two complete first normals are
in $(A_0,B_0,C_0)$. These are the full formulas, without restricting
normals or replacing native coefficients by leading surrogates.

At $k=v=0$ and at the relevant small-$w$ or small-$u$ endpoint, put
$r_0=1-aE$ and $H_0=r_0+E$. Direct collection of the short formulas
gives the raw scalar residue
\[
 8r_0H_0D_u^2 A_0^2+
 8D_u\{x^2+u^2r_0H_0B_0^2\}+8r_0H_0D_uD_0^2.
 \tag{BE5}
\]
The hard block on $(N,\delta M_B)$ is
\[
 \operatorname{diag}(16r_0H_0,16r_0H_0+8E^2/D_u),
\]
and the complete first normals have residues $0,16Ex$.
The coupled $(\delta M_B,x)$ determinant is
$128D_ur_0H_0>0$.

If $u$ is a unit and $w$ infinitesimal, (BE5) is a strict Gram on
$(A_0,x,B_0,D_0,N,\delta M_B)$. If $u$ is infinitesimal and $w$ a
unit, use the exact relations
\[
 B_0=(\delta/w)A_0,\qquad C_0=(\delta/w)A_0-x.
\]
The factor $\delta/w$ is a product of a pure-coefficient unit and a
pure-testing unit; it is not a mixed inverse. The resulting residue
is the strict form
\[
 8r_0H_0A_0^2+8x^2+8r_0H_0D_0^2
 +16r_0H_0N^2+(16r_0H_0+8E^2)(\delta M_B)^2
 +16E(\delta M_B)x.                                  \tag{BE6}
\]
This supplies precisely the unit-$w$ chart for which the older
$\Delta$-small retained-current argument is not applicable.

For transparency, the only nontrivial raw coefficient collection in
(BE5) is
\[
 r_0H_0D_u+a^2E^2+r_0E(2a-1)=1+u^2r_0H_0,
 \qquad aE+r_0=1.
\]
The cross coefficient is $-16D_u(aE+r_0)=-16D_u$.
In the full normal formula $S_B$ has residue $-16E$ and $S_C$ has
residue $16E$. These prove the claimed full normal residues directly.

\paragraph{The sole contact chart.}
For $u,w,v$ infinitesimal and $Z=1$, use the unchanged monic
contact--conductor identities and the two-pairing circuit. Use weights
\[ (\delta,k,b,a,t,u,w,v):\quad(0,1,0,0,0,1,1,2). \]
The exact source collection is
\[
 (E_0)_{\le6}=8\rho T_4-8(\rho+E)k^2x^2,
 \quad (L_1)_{\le3}=0,\quad (L_2)_{\le3}=16Ex,
 \qquad \rho=r_0H_0-1.
\]
The source/contact substitutions preserve this grading. After the
displayed terms are removed, the scalar and normal lift coefficients
have positive weights, hence zero residues. The term $k^2x^2$ also
has a zero-residue coefficient. Thus the soft leading form is
$8x^2+8r_0H_0T_4$, with the same strict two-pairing circuit frame.
The hard block and its only first-normal coupling are those of (BE6),
and their determinant is $128r_0H_0>0$. Retain the old positive
$g_4^2$ square. The full eight-root remainder Gram is strict.
This invokes only the polynomial monic conductor already proved;
no opening or contact difference is inverted.

Finally restore the four complete positive terms, the retained
positive square, and all pure-coefficient and pure-testing factors.
The cases above exhaust both projective angle lines and every
projective state. Complex states split into their real and imaginary
evaluations. Neither this theorem nor the correction of (BE1) asserts
coverage when $r$ is infinitesimal, $\eta$ unbounded, or $e$
infinitesimal; those are separate source domains.



% Source formulas: p3763434080c2
\subsection{The centered-$A$ source with unit $e$}

\label{component:p3763434080c2}

\begin{theorem}[The centered-$A$ unit-$e$ slice]
Use the physical native C1 source with its unchanged original events,
$c=2$, $\lambda=n=0$, all positive openings $a,t,M,J$, and all original
nonnegative signs, including literal $\beta=0$.
Suppose $k,\zeta,r=1-a\eta$ are positive infinitesimals, while
$\eta$ and $e=k\sqrt{\eta/\zeta}$ are positive bounded units.
Then the full parent is tensor SOS at every independent projective test
and every real or complex state. No bound on $k^2/r$ is imposed.
\end{theorem}

The theorem is restricted to unit $e$. Its proof uses $\delta^{-1}$
as a bounded pure coefficient number; it does not assert the cases of
infinitesimal or unbounded $e$.

\paragraph{Centering the sole large event.}
The exact native contact has $Y_c=r u_cW_c$. Consequently
\[
 L_A=Y-r uW
\]
vanishes at contact with no asymptotic replacement. The native formulas
give $a$ a positive unit, $A$ unbounded and $Ar$ a positive unit of
residue $2$. Put $A_*=1/a$, a positive bounded pure coefficient unit.
The complete term
\[
             4(A-A_*)D_uD_v L_A^2                         \tag{CA1}
\]
is positive and has zero value and full first derivative at the exact
native contact. Retain it as a square.

It is essential to center this term before taking residues.
The remaining $A$ channel is exactly
\[
 4A_*D_uD_vL_A^2+
 4(Ar)D_uD_v\{2uWY-r u^2W^2\}.                         \tag{CA2}
\]
Only $Ar$ and $r$ occur in its coefficients. In particular this is a
bounded expression for every positive rate of $r$ relative to $k$.

\paragraph{A source chart with no hidden $k^2/r$ quotient.}
Keep the exact contact variables $\delta,k,E=b^2$ from (SD6), and put
\[
 K=1+k^2,\quad \ell=K-aE=Qr>0,\quad
 Q=K^2-\ell k^2,\quad a=(K-\ell)/E.
\]
Here $\delta,E,a$ are positive bounded units and $k,\ell$ are positive
infinitesimals. The exact source substitution is
\[
 \eta=KE/Q,\quad r=\ell/Q,\quad
 \Delta=\delta Q/K,\quad \zeta=\delta^2k^2QE/K,
 \quad R=K^2/Q,\quad T=K-k^2t\zeta.                    \tag{CA3}
\]
All denominators in (CA3) are positive pure coefficient units.
For $i=2,b,3,5,\alpha,d$ let $\widetilde c_i=c_i/\ell$,
defined by cancelling the explicit $L_c=\ell/Q$ factor in the short
native formulas. This cancellation gives rational expressions with
unit denominators; it does not leave an inverse of $\ell$.
For example, writing $H_c=r+\eta$ and $V_c=k^2\eta\zeta$,
\[
 \widetilde c_3=\frac{2KH_cQ_c(1-k^2\eta+tV_c)}Q,
 \qquad \widetilde c_d=\frac{4Q_cP_d}Q.
\]
Their residues are $2E$ for $i=2,b,3,5,\alpha$, and $4$ for $i=d$;
$c_4$ has residue $2$.

Define the following bounded pure coefficient quantities:
\[
 \begin{array}{ll}
 B_2=\dfrac{8a\widetilde c_2}{R\widetilde c_3\Delta^2},&
 B_b=\dfrac{8a\widetilde c_b}{R\widetilde c_3\Delta^2},\\[4pt]
 B_3=8a/T,&B_5=\dfrac{8a\widetilde c_5}{T\widetilde c_3},\\[4pt]
 L=Ar=\dfrac{2a\eta^2c_4}{Q\widetilde c_3},&
 d=\dfrac{2a\eta\widetilde c_d}{\widetilde c_3},\\[4pt]
 \alpha=\dfrac{2a\Delta^2k^6\eta^2\widetilde c_\alpha}
                    {\widetilde c_3}.&
 \end{array}                                                     \tag{CA4}
\]
These formulas explicitly account for the cancellation needed below.

\paragraph{The exact full remainder, before any collection.}
Retain the four original (SD3) terms, divided by the positive source
factor $\Delta^2\Lambda$, and retain (CA1). Use the unchanged pure
testing coordinate $w=(W+uY)/D_u$ and the hard normals $N,M_B$.
For $Z=1$ put $p=t/R$, $h=a/T$ and
\[
 q_A=vD_uw,\quad q_B=v(D_u-uN-puD_uv w),
\]
\[
 W=w+u(M_B+hq_B)/D_u,\qquad
 Y=uw-(M_B+hq_B)/D_u.
\]
The complete remainder is exactly
\begin{align*}
\mathscr G={}&B_2\{N^2-2q_AN+(1-p)q_A^2\}
 +B_bD_v\{N^2+2q_AN+pq_A^2\}\\
&+B_3D_v\{M_B^2+2q_BM_B+hq_B^2\}
 +B_5\{M_B^2-2q_BM_B+(1-h)q_B^2\}\\
&+4A_*D_uD_v(Y-r uW)^2
 +4LD_uD_v(2uWY-r u^2W^2)\\
&+4\alpha D_u(vY-1)^2
 +2dD_u\{(vY-1)^2-D_v\}.                         \tag{CA5}
\end{align*}
There are no normal terms above degree two. Every coefficient of
(CA5) is bounded. The only testing denominator is a power of the
positive unit $D_u$; every source denominator in (CA3)--(CA5) is a
positive unit. The possibly unbounded opening ratios remain exclusively
in the four retained positive terms.

\paragraph{The exact low-weight collection.}
The outer/state lemma in Section~\ref{component:p364e14238053} reduces the
independent tests to $Z=1$ and $u,w,v$ infinitesimal, with $w>0$ on
the negative pure-core branch. Use the unchanged circuit monomials
\[
 A_0=wv,\quad B_0=\delta v,\quad C_0=\delta Euw,
 \quad D_0=\delta^2k^3E,
\]
\[
 x=B_0-C_0,\quad y=A_0-kC_0,\quad
 z=uB_0-kC_0,\quad h_0=kC_0-D_0.
\]
Let $q_1,q_2$ be the same two separated radical pairings as in the
contact--conductor proof, and let
\[
 T_4=y^2+z^2+h_0^2-2yz+2q_1^2.
\]
Give $(\delta,k,E,\ell,t,u,w,v)$ weights $(0,1,0,1,0,1,1,2)$.
Then $x$ has weight two and $y,z,h_0,q_1,q_2$ have weight three.
Denote the scalar of (CA5) by $S$, and its two complete first normals
by $L_N,L_M$. Direct multiplication of (CA3)--(CA5) through the
indicated weights gives
\[
 \begin{aligned}
 S_{\le6}&=\frac{4P}{E^4\delta^2}x^2
 -\frac{16(k+u)}{E^2\delta^2}xz
 +\frac8{E\delta^2}T_4,\\
 (L_N)_{\le3}&=0,\\
 (L_M)_{\le3}&=
 \frac{8\{E^2+2E-2\ell(E+1)\}}{E^2\delta}\,x,
 \end{aligned}                                                   \tag{CA6}
\]
where the short coefficient polynomial is
\[
 \begin{split}
 P={}&E^2(E+4)-E(E^2+6E+4)\ell+2(E+1)(E+2)\ell^2\\
 &+E^2(8-3E)k^2+4E^2ku+E^2(E+4)u^2.
 \end{split}
\]
The full hard-normal block, from the quadratic part of (CA5), has
residue
\[
       \operatorname{diag}\left(\frac{16}{E\delta^2},
                                \frac{4(E^2+4)}E\right).          \tag{CA7}
\]
Here and in the residue Grams, $E,\delta$ denote their positive
residues. Formula (CA6), rather than a large expansion, is the finite
coefficient identity used by the proof. Its 37 raw scalar monomials
collect to the displayed three terms. 

\paragraph{Why all higher terms have harmless coefficients.}
The scalar of (CA5), after clearing the stated units, belongs to
$(A_0,B_0,C_0)^2+(D_0^2)$. This follows directly from $q_A=wvD_u$,
$q_B|_{N=M_B=0}\in(v)$, and $Y,W$ in (CA5): the centered scalar
$Y-r uW$ lies in $(v,uw)$; its square is in the same raw monomial
square ideal since $\delta,E$ are units. The sole constant state term
comes from $\alpha$, which has the factor $k^6$, and is a bounded
unit multiple of $D_0^2$. Both full first normals belong to
$(A_0,B_0,C_0)$ by the same displayed expression.

All retained terms vanish to first order at the exact contact
\[ (u,w,v,N,M_B)=(k,\delta k,\delta k^2E,0,0). \]
Thus the scalar belongs to the square of its contact ideal and the
first normals belong to that ideal. The unchanged monic
contact--conductor identities therefore apply, over the coefficient
ring in which only the displayed positive units have been inverted.
They never divide by $k$, $\ell$, $M$, $J$, $t$, or a contact difference;
the displayed $1/a$ is the inverse of a known positive unit.
The contact substitutions preserve the new weights. After (CA6) is
subtracted, scalar weights are at least seven and normal weights at
least four. Their lifted coefficients consequently have positive
weights, hence zero residues. Clearing the finitely many unit
denominators before this polynomial lift and restoring them afterward
changes no residue sign and introduces no mixed inverse.

\paragraph{Strictness, full normals, and restoration.}
The leading scalar is
\[
 \frac{4(E+4)}{E^2\delta^2}x^2+\frac8{E\delta^2}T_4.
\]
The standard two-pairing circuit supplies a strict Gram on
$(y,z,h_0,q_1,q_2)$. The only nonzero first-normal residue is
$8(E+2)x/(E\delta)$, the coefficient of the full normal $M_B$.
Accordingly the coupled $(M_B,x)$ Gram is
\[
 \begin{pmatrix}
 4(E^2+4)/E&4(E+2)/(E\delta)\\
 4(E+2)/(E\delta)&4(E+4)/(E^2\delta^2)
 \end{pmatrix},\qquad
 \det=\frac{256}{E^3\delta^2}>0.                      \tag{CA8}
\]
Together with the independent $N$ pivot in (CA7), this gives a
bounded strict eight-root Gram for the entire remainder (CA5).
Restore (CA1), all four original positive terms, and all separated
normalization factors. The outer/state lemma exhausts the remaining
tests; real and imaginary state evaluations give the complex case.

This proves precisely the stated unit-$e$, unit-$\eta$,
infinitesimal-$r$ slice. Combined with the bounded-unit-$r$ theorem,
it covers every positive $r$ when $\eta,e$ are bounded units and
$k,\zeta$ infinitesimal. It makes no assertion about unbounded
$\eta$ or either nonunit alternative for $e$.



% Source formulas: pccaa793fd0bf
\subsection{The centered-$A$ source with bounded reciprocal $e$}

\label{component:pccaa793fd0bf}

\begin{theorem}
Use the unchanged physical native C1 source with $k,\zeta,r=1-a\eta$
positive infinitesimals and $\eta$ a positive bounded unit. Suppose
$e=k\sqrt{\eta/\zeta}$ is bounded below by a positive real number;
it may be unbounded. Then the entire original parent is tensor SOS at
every independent projective test and every real or complex state.
Every positive opening rate, literal $\beta=0$, and every rate of
$k^2/r$ is included.
\end{theorem}

Use all exact definitions and the centered positive term (CA1) in
Section~\ref{component:p3763434080c2}. The sole change is that
$\delta=\Delta R/K$ is now any positive bounded coefficient, so it
must not be inverted in a bounded Gram. The source chart (CA3),
$\ell=Qr=K-aE$, and the exact cancelled coquotients remain unchanged.
Their only denominators are positive units, independent of the rate
of $\ell$ relative to $k$.

\paragraph{Scale the entire remainder and its full normal.}
Let $\mathscr G$ be the exact full remainder (CA5). Put
\[
 \widehat G=\delta^2\mathscr G,\qquad
 \widehat M=\delta M_B,\quad \widehat Y=\delta Y,
 \quad\widehat W=\delta W.
\]
The original testing variables have not been rescaled against a
coefficient rate; these displayed quantities are explicit separated
tensor roots. Their complete formulas are
\[
 q_A=vD_uw,\quad q_B=v(D_u-uN-puD_uv w),
\]
\[
 \widehat W=\delta w+u(\widehat M+\delta hq_B)/D_u,
 \quad\widehat Y=\delta uw-(\widehat M+\delta hq_B)/D_u.
\]
Set
\[
 \widehat B_2=\frac{8a\widetilde c_2K^2}{R\widetilde c_3Q^2},
 \qquad
 \widehat B_b=\frac{8a\widetilde c_bK^2}{R\widetilde c_3Q^2}.
\]
The entire exact scaled remainder is
\begin{align*}
\widehat G={}&\widehat B_2\{N^2-2q_AN+(1-p)q_A^2\}
 +\widehat B_bD_v\{N^2+2q_AN+pq_A^2\}\\
&+B_3D_v\{\widehat M^2+2\delta q_B\widehat M+\delta^2hq_B^2\}
 +B_5\{\widehat M^2-2\delta q_B\widehat M+\delta^2(1-h)q_B^2\}\\
&+4A_*D_uD_v(\widehat Y-r u\widehat W)^2
 +4LD_uD_v(2u\widehat W\widehat Y-r u^2\widehat W^2)\\
&+4\alpha D_u(v\widehat Y-\delta)^2
 +2dD_u\{(v\widehat Y-\delta)^2-\delta^2D_v\}.
 \tag{CS1}
\end{align*}
All coefficients are bounded; every occurrence of $\delta$ has
nonnegative exponent. There is no inverse of $\delta$, $\ell$, $k$,
$M$, $J$ or $t$. The inverse $1/a$ in $A_*$ is a positive unit.
Every first-normal and hard-normal term is retained in (CS1).

\paragraph{Uniform contact grading.}
Keep the exact roots $A_0,B_0,C_0,D_0,x,y,z,h_0,q_1,q_2$ and polynomial
$P$ in (CA6). Use the same weights
$\operatorname{wt}(\delta,k,E,\ell,t,u,w,v)=(0,1,0,1,0,1,1,2)$.
Weight zero means bounded here, not bounded away from zero.
The exact finite coefficient collection from (CS1) is
\[
 \begin{aligned}
 (\widehat G|_{N=\widehat M=0})_{\le6}
 &=\frac{4P}{E^4}x^2-\frac{16(k+u)}{E^2}xz+\frac8E T_4,\\
 (\widehat L_N)_{\le3}&=0,\\
 (\widehat L_M)_{\le3}
 &=\frac{8\{E^2+2E-2\ell(E+1)\}}{E^2}x.
 \end{aligned}                                                    \tag{CS2}
\]
The full hard block on $(N,\widehat M)$ has residue
\[
        \operatorname{diag}(16/E,4(E^2+4)/E).          \tag{CS3}
\]
These are precisely the old identities after the indicated scaling,
but their validity for an infinitesimal $\delta$ requires the
nonnegative-exponent assertion in (CS1), not cancellation of a
purported unit $\delta$.

The scalar is in
$(A_0,B_0,C_0)^2+(D_0^2)$ over the ring with $E$ and the displayed
positive units inverted. Both complete first normals are in
$(A_0,B_0,C_0)$. Direct monomial allocation in (CS1) proves these
memberships with nonnegative $\delta$ exponents. At zero normals,
$q_A=A_0D_u$, $\delta q_B=B_0D_u(1-puv w)$, and both
$\widehat Y,u\widehat W$ belong to $(B_0,C_0)$ with bounded-unit
coefficients. The sole constant from the $\alpha$ term satisfies
$\alpha\delta^2/D_0^2=2a\widetilde c_\alpha/\widetilde c_3$ exactly.
Thus these are memberships of the entire rational expression, not merely
its finite collection. The unchanged monic
contact conductor lifts them without a $\delta$ inverse. Its contact
substitutions preserve the weights, even if $\delta$ is infinitesimal.
After (CS2), all remaining scalar weights are at least seven and all
remaining normal weights at least four, hence the lifted correction
coefficients have zero residue. 

The leading soft form is
\[
            \frac{4(E+4)}{E^2}x^2+\frac8E T_4.
\]
Use the strict two-pairing circuit on $T_4$. Its only hard coupling
is the full first normal $8(E+2)x/E$ on $\widehat M$, and
\[
 \det\begin{pmatrix}4(E^2+4)/E&4(E+2)/E\\
                    4(E+2)/E&4(E+4)/E^2\end{pmatrix}
                   =256/E^3>0.                       \tag{CS4}
\]
Together with $16/E$ on $N$, this gives the strict eight-root Gram.

\paragraph{Exhaustion and restoration.}
The independent outer/state lemma Section~\ref{component:p364e14238053}
uses the unit-$e$ hypothesis only to infer a positive lower bound on
$m/d$ in (CAO2). When $e$ is unbounded, $m/d$ is instead unbounded,
so that same lower bound and every clipping argument hold. Its
$c/d,A/d,D/d$ bounds, exact centered root identity, ordinary-$u$
certificate, unit-$w$ certificate, and literal projective endpoint
arguments do not contain $e$ or $\delta$. Thus it reduces every
remaining independent test to exactly the contact chart settled above.
Restore $\delta^2$ times every retained original term, including
(CA1), and undo the positive pure coefficient factor $\delta^2$.
All factors and radicals remain separated between the two fields.
Complex states split into real and imaginary evaluations.

This closes the formerly missing infinitesimal-$r$, unbounded-$e$
slice, and incorporates the previous unit-$e$ theorem. Together with
the common unit-$r$ results, it covers all positive native $r$
when $\eta$ is a unit and $e$ is bounded below. It does not yet
claim the infinitesimal-$e$ slice.



% Source formulas: p364e14238053
\subsection{Complementary centered-$A$ testing and state charts}

\label{component:p364e14238053}

\label{p364e14238053:sec:centered-a-outer}

\begin{lemma}
Use the unchanged physical native C1 source, normalized by $c=2$ and
$\lambda=n=0$, with all original event signs and positive openings
$a,t,M,J$. Suppose $k,\zeta,r=1-a\eta$ are positive infinitesimals and
$\eta,e=k\sqrt{\eta/\zeta}$ are positive bounded units.
Then the outer-test and state-entry arguments reduce the all-tests
problem to precisely
\[
 u,v,w>0\text{ infinitesimal},\qquad Z=1,
 \qquad w=(W+uY)/(1+u^2).
 \tag{CAO1}
\]
All other independent projective tests and states admit a full-parent
tensor SOS. No comparison of $k^2$ with $r$ is made, and arbitrarily
small positive openings and literal $\beta=0$ are included.
\end{lemma}

\paragraph{Source ratios with the common factor cancelled.}
Here $r$ always denotes the native source parameter, not the event root
$uW-Y$. Put
\[
 K=1+k^2,\quad H_c=r+\eta,\quad V_c=k^2\eta\zeta,\quad
 S_3=1-k^2\eta+tV_c,
\]
The two required short native polynomials are
\begin{align*}
P_4={}&1+k^2+\eta k^2(a-k^2-1)-\zeta k^2t\\
 &+V_c[-\zeta k^2t^2+2k^2t+a(1-t)+t],\\
P_d={}&(1+k^2)^2-a\eta k^4
 +\zeta k^2[(k^2+1)-(k^2+2)t]\\
 &+V_c[-\zeta k^2t^2+k^2a(t-1)+k^2t+k^2+1].
\end{align*}
They are positive bounded units with residue one. The exact source
identities, after cancelling the common factor $r$ in the coquotients,
are
\begin{align*}
 d&=\frac{4a\eta P_d}{KH_cS_3},&
 Ar&=\frac{2a\eta^2P_4}{KH_cS_3},\\
 \frac md&=\frac{e^2K(1-t\zeta)(1+V_c)H_c}{2M\eta P_d},&
 \frac cd&=\frac{KH_cS_3}{2a\eta P_d},\\
 \frac Ad&=\frac{\eta P_4}{2rP_d},&
 \frac Dd&=\frac{KH_c(1+V_c)-J\eta P_d}{2J\eta P_d}.
 \tag{CAO2}
\end{align*}
The positive openings obey
$M+t=1+k^2r$ and $J+a=1+k^2-k^2\zeta t$; hence $a,t,M,J$ are
bounded. Since $a\eta=1-r$, one has
\[
 \bar\eta\ge1,\qquad \bar d=4/\bar\eta,\qquad Ar\sim2,
 \qquad A\text{ positive unbounded}.
 \tag{CAO3}
\]
The numerator of $D/d$ in (CAO2) has residue
$\bar H_c-\bar J\bar\eta=1$. Therefore $m/d,c/d,D/d,A/d$ all have
positive real lower bounds. This conclusion uses no inverse bound on
an opening. Moreover,
\[
 \begin{cases}
 \displaystyle \bar D=\frac{2}{\bar\eta(\bar\eta-1)},
       &\bar\eta>1,\\[4pt]
 D\text{ positive unbounded},&\bar\eta=1.
 \end{cases}
 \tag{CAO4}
\]

\paragraph{The previously proved outer routing still applies.}
Clip $m,c,A,D$ at positive real multiples of $d$, retaining all their
nonnegative excess original event squares. The resulting coefficients
divided by $d$ are bounded. The outer-test argument in
Section~\ref{component:p69995d829bb2}, through (BE2)--(BE4), uses only
these lower bounds and the independent physical-source certificate at
noncommuting residue tests. Its exact homogeneous identities and
four-root ideal memberships are unchanged. It therefore includes all
projective infinity charts, the $(0,\infty)$, $(\infty,0)$ and
$(\infty,\infty)$ corners, and the unit-$W$, infinitesimal-$Z$ state
chart. The original-event SOS on literal zero axes is unchanged as
well. This invokes none of the bounded-$r$ scalar or normal conclusions
(BE5)--(BE6). Only bounded $u$ and infinitesimal $v$ remain.

\paragraph{An exact centered root identity.}
Write $D_u=1+u^2$, $D_v=1+v^2$ and introduce
\[
 h=A^{-1/2},\quad L=Ar,\qquad
 H=\sqrt A(Y-ruW),\quad B=X+uZ,\quad U=uW,\quad V=vZ.
 \tag{CAO5}
\]
Thus $Y=rU+hH$, while $h,r$ are infinitesimal and $L\sim2$.
All coefficient radicals in (CAO5) lie in the coefficient field;
$B,U,V$ lie in the testing field. The root $H$ is their finite tensor
combination, with no mixed inverse or mixed radical.

For clarity, the complete original finite parent used here is
\begin{align*}
F={}&4m(vW+uvY-X-uZ)^2+8D_v(uW-Y)^2
       +4AD_uD_vY^2\\
 &+(4D+2d)(uW+uvX-Y-vZ)^2
   +(4\alpha+2d)D_u(vY-Z)^2\\
 &+(4\beta+2d)D_v(X+uZ)^2-2dD_uD_vZ^2.
\end{align*}

Choose $0<D_*\le D$ bounded, and let $m_*=0$ or choose a positive
real $m_*\le m$, as specified below. Retain the four nonnegative
original event contributions with weights $m-m_*,D-D_*,\alpha,\beta$.
For $C=4D_*+2d$, the exact remaining parent is
\begin{align*}
F_*={}&8D_v[(1-r)U-hH]^2
 +4D_uD_vH^2+4LrD_uD_vU^2+8LhD_uD_vUH\\
&+C[(1-r)U-hH+uvB-D_uV]^2+2dD_vB^2\\
&+2dD_u\{v^2(rU+hH)^2-2(rU+hH)V-V^2\}
 +4m_*R_m^2,                                      \tag{CAO6}\\
R_m={}&vW+uv(rU+hH)-B.
\end{align*}
In particular, the centered $A$ contribution is the exact identity
\[
 4AD_uD_vY^2
 =4D_uD_vH^2+4LrD_uD_vU^2+8LhD_uD_vUH.
\]
This retains its complete hard root. The last two coefficients have
zero residue. Formula (CAO6) follows directly from the original parent
and includes the complete signed core; it is not an eliminated scalar
or a restriction to state normals. Every displayed coefficient is
bounded. The term involving $Z$ in that core has already cancelled
to the displayed expression in $V=vZ$.

\paragraph{Every state when $u$ is a bounded unit.}
Set $m_*=0$, so the entire original $m$ event is retained. Formula
(CAO6) is an exact Gram allocation on $(H,B,U,V)$. At infinitesimal
$v$ its residue is
\[
 4D_uH^2+2dB^2+8U^2+C(U-D_uV)^2-2dD_uV^2.
 \tag{CAO7}
\]
Residues are understood in this and the following determinant
calculations. The $U,V$ determinant is
\[
 D_u\{8CD_u-2d(8+C)\}.
\]
If $\bar\eta>1$, take $D_*=D$, which is bounded by (CAO4).
Then $C=8/(\bar\eta-1)$ in residue and this determinant is
\[
 \frac{64u^2(1+u^2)}{\bar\eta-1}>0.
\]
If $\bar\eta=1$, choose a sufficiently large positive real $D_*$;
this is permitted because $D$ is unbounded. The determinant is then
\[
 32(1+u^2)(D_*u^2+2u^2-2)>0.
\]
The first diagonal is $8+C>0$, so the complete four-root residue
Gram is positive definite in both cases. The cap is a real constant
chosen from the nonzero testing residue of $u$; no division of tensor
elements is introduced. The bounded full Gram is therefore tensor
SOS, and restoring the retained events proves the claim for every
state in this chart.

\paragraph{The remaining unit-$w$ state chart.}
Suppose $u,v$ are infinitesimal and $w$ is a unit. Then $W$ is a unit
on the negative-core branch. When $Z$ is infinitesimal the inherited
(BE4) argument already applies. Otherwise normalize the state so that
$Z$ is a unit and put $q=W/Z$, a pure-testing unit. Choose a positive
real lower bound $m_*\le m$. Now
\[
 R_m=qV+uv(rU+hH)-B,
\]
so (CAO6) is again an exact bounded Gram on $(H,B,U,V)$, with residue
\[
 4H^2+2dB^2+4m_*(qV-B)^2
       +8U^2+C(U-V)^2-2dV^2.                    \tag{CAO8}
\]
The $B$ pivot is positive. Its Schur complement adds
\[
 \mu V^2,\qquad \mu=\frac{4dm_*q^2}{d+2m_*}>0
\]
to the $U,V$ block. This is just a residue determinant computation;
no mixed inverse is used to define a root or a full-parent coefficient.
If $\bar\eta>1$, take $D_*=D$ as before. The determinant becomes
\[
 \frac{8\bar\eta\mu}{\bar\eta-1}>0.
\]
If $\bar\eta=1$, choose a sufficiently large positive real cap
$D_*$ so that $\mu(C+8)-64>0$. The unboundedness of $D$ permits this
choice. Thus the complete residue Gram in (CAO8) is strict, including
all hard-soft terms, and restoring the retained events gives the
full-parent tensor SOS.

\paragraph{Completion of state entry.}
Normalize the state by a maximal coordinate and first remove the
nonnegative pure-testing-core branch, which has its direct original
SOS. On the negative-core branch with $v$ infinitesimal, the residue
core is
\[
 2\{(X+uZ)^2+(uW-Y)^2\}.
\]
Consequently these two roots have zero residue. If $w$ and $Z$ were
both infinitesimal, the exact inverse relations
\[
 W=w+\frac{u(uW-Y)}{D_u},\quad
 Y=uw-\frac{uW-Y}{D_u},\quad X=(X+uZ)-uZ
\]
would make every normalized state coordinate infinitesimal. Hence
small $w$ forces $Z$ to be a unit. The bounded-unit $u$ chart was
settled by (CAO7), and all unit-$w$ states by (CAO8) and (BE4).
Normalize the only remaining states by $Z=1$. The same exact
pure-testing Schur identity (SD5) and testing sign gauges give
$u,v,w>0$. This is exactly (CAO1). Literal axes and literal state
zeros have already been included. Complex states split into their
real and imaginary evaluations because the parent is real.



% Source formulas: p890eea5f538b
\subsection{Unit $\eta,r$ and unbounded $e$}

\label{component:p890eea5f538b}

\begin{theorem}[Unit $\eta$ and unbounded $e$]
Use the unchanged physical native C1 source, its original event signs,
and the positive shapes and gaps $a,t,M,J$. Suppose $k,\zeta$ are
positive infinitesimals, $\eta$ is a positive bounded unit,
$r=1-a\eta$ is a positive bounded unit, and $e=k\sqrt{\eta/\zeta}$
is unbounded. Then the full original parent is tensor SOS at every
independent projective test and every real or complex state.
No positive lower bound on $a,t,M,J$ is imposed.
\end{theorem}

Write $H=r+\eta$ and $\Delta=1/e$. The seven native coquotients
have residues $2rH$ for $c_2,c_b,c_3,c_5,c_\alpha$, $2$ for
$c_4$, and $4r$ for $c_d$. All are positive units. In particular the
four retained terms of the contact--conductor proof remain positive
and the remainder has bounded coefficients. We reuse their exact
definitions, not asymptotic substitutes for the original source.

\paragraph{Outer tests.}
The exact formulas give
\[
 \frac md=\frac{e^2c_2}{M\eta c_d},\qquad
 \frac\beta d=\frac{e^2c_b}{t\eta c_d}-\frac12.
\]
Both ratios are unbounded. Apply the exact full-parent Alice reflection
(AS3)--(AS4), which exchanges $m\leftrightarrow c$ and
$\beta\leftrightarrow D$ while leaving $v$ unchanged. The already
proved two-strong-current identity then covers every test with $1/v$
bounded, for every Alice angle and every state. Its homogeneous form
includes both literal Alice endpoints, so no inverse-angle point is lost.

For $v$ infinitesimal and Alice infinity, use the critical-current
identity with $g_c=cD/(c+D)$, $\epsilon=d/g_c$ and
$\gamma=Ag_c/d^2$. Its actual sufficient conditions on this chart
are bounded $\epsilon$ and positive-unit $\gamma$: in its exact
four-root identity the correction $L_0$ has factors $c_u$ or $v$,
where $c_u=(1+u^2)^{-1/2}$ is infinitesimal. Thus $\sqrt\epsilon$
need only be bounded. The residue Gram is
$\operatorname{diag}(1,2,2,4\bar\gamma)$.
Use the exact relation $J+a=T$ with $T-1$ infinitesimal, and the
cancelled original coefficient formula
\[
 \frac Da=\frac{2c_5/c_3-J\eta c_d/c_3}{J}.
\]
Its numerator has residue $2/\bar H$. The unchanged native formulas give
\[
 \frac d a\sim\frac{4\eta}{H},\quad
 \frac A a\sim\frac{2\eta^2}{rH},\quad
 \frac D a\sim\frac2{JH},\qquad
 \epsilon\sim\frac{2\eta(JH+a)}H,\quad
 \gamma\sim\frac{H}{4r(JH+a)}.
\]
The identity
$\overline{JH+a}=1+\bar\eta(1-\bar a)^2$ shows that the two required
conditions hold even when $a$ or $J$ is infinitesimal. The calculation
uses the cancelled harmonic coefficient $g_c/a$, so no lower bound
on an opening is inserted. Homogeneous angle normalization includes
literal Alice infinity. It remains to treat bounded $u$ and
infinitesimal $v$.

\paragraph{State entry.}
Normalize the state by a coordinate of maximal absolute value and put
$D_u=1+u^2$ and $w=(W+uY)/D_u$. If $w$ is a unit, the exact five-root
current identity (UW2)--(UW3) applies with $Q=W+uY=D_uw$ a unit.
Its first four diagonal residues are $8rH$ and its fifth is
$8D_u$; both negative corrections have factors $\Delta$ or
$\Delta^2\eta$. Thus this chart is strict.

If $w$ is infinitesimal, the nonnegative pure-testing core is a direct
exit. On the negative-core branch $Z$ is a unit. Indeed its residue
at $v=0$ is $2\{(X+uZ)^2+(uW-Y)^2\}$; if $Z,w$ and both currents
had zero residue, every normalized state coordinate would have zero
residue. Normalize $Z=1$. The exact core Schur identity gives $w>0$
after the usual positive-angle sign gauges. All transformations here
belong solely to the testing field.

\paragraph{The ordinary-$u$ chart needs no contact division.}
Use the exact source change
\[
 \delta=\Delta R/K,\quad b^2=\eta K/R,\quad K=1+k^2,
 \quad Q=K+ak^2b^2,\quad T_n=K^2-t\delta^2k^4Qb^2.
\]
All source denominators are positive units, $\delta,k$ are
infinitesimal and $b^2\sim\eta$. Let
\[
 A_0=wv,\quad B_0=\delta v,\quad C_0=\delta b^2uw,
 \quad D_0=\delta^2k^3b^2,\qquad x=B_0-C_0.
\]
The five-term exact scalar belongs to
$(A_0,B_0,C_0)^2+(D_0^2)$ and both complete first normals belong
to $(A_0,B_0,C_0)$. Collecting these displayed factors before taking
residues is legitimate over the bounded coefficient ring. At
$\delta=k=w=v=0$, keeping bounded $u$ free, the scalar Gram is
\[
 8rH D_u^2A_0^2+
 8D_u\{x^2+u^2rH B_0^2\}+8rH D_uD_0^2.
 \tag{UE1}
\]
The hard block on $(N,\delta M_B)$ is
$\operatorname{diag}(16rH,16rH+8\eta^2/D_u)$, and the complete
first normals have residues $0$ and $16\eta x$. The latter is the
coefficient of $\delta M_B$ in the full polynomial, so the symmetric
Gram off-diagonal entry between $\delta M_B$ and $x$ is $8\eta$.
The coupled $(\delta M_B,x)$ block has determinant
\[
 (16rH+8\eta^2/D_u)\,8D_u-(8\eta)^2=128D_urH>0.
 \tag{UE2}
\]
If $u$ is a unit, the remaining $B_0$ pivot in (UE1) is strict,
as are the $A_0,D_0,N$ pivots. Every omitted coefficient has zero
residue. This proves the chart by a bounded strict Gram on these
six roots, without any mixed inverse. Equation (UE1) follows directly
from $a\eta+r=1$: its $(B_0,C_0)$ part before collection is
$8D_u\{(1+u^2rH)B_0^2-2B_0C_0+C_0^2\}$.

\paragraph{The infinitesimal-$u$ chart uses the old conductor.}
Now $\delta,k,u,w,v$ are infinitesimal, while $b,a,t$ are bounded.
Use the unchanged roots
\[
 y=A_0-kC_0,\quad z=uB_0-kC_0,\quad h=kC_0-D_0,
 \quad q_1=a_1a_2-a_3a_4,
\]
and the second pairing $q_2=a_1a_3-a_2a_4$ of the four separated
radicals. Set
\[
 T_4=\sum_{i=1}^4a_i^4-4\prod_{i=1}^4a_i
     =y^2+z^2+h^2-2yz+2q_1^2,
 \qquad \rho=(1-ab^2)(1+(1-a)b^2)-1.
\]
Give $(\delta,k,b,a,t,u,w,v)$ weights $(1,1,0,0,0,1,2,3)$.
The root $x$ has weight four and the other soft roots have weight
five. Substitution at the exact contact preserves this grading.
Writing $E_0$ for the complete scalar correction relative to the
old leading polynomial, the five-term factorization and the two
complete normal formulas give
\[
 \begin{aligned}
 E_0&=8\rho T_4-8(\rho+b^2)k^2x^2+O_{\rm wt}(12),\\
 L_1&=O_{\rm wt}(6),\qquad
 L_2=16b^2x+O_{\rm wt}(6).
 \end{aligned}\tag{UE3}
\]
These identities require only multiplication of the short source
polynomials. In particular the residue correction is the single
factor $\rho$ multiplying the same circuit, rather than a new
polynomial positivity problem.

The unlocalized monic contact--conductor lift applies unchanged.
After the terms explicitly displayed in (UE3) are removed, scalar
coefficient weights are at least two and normal coefficient weights
are at least one. Therefore they all have zero residue. There is no
division by $k,b,\delta,a,t$ or a contact difference. The leading
strict soft Gram is now
\[
                  8x^2+8rH T_4.
\]
The two-pairing circuit identity supplies a strict Gram on
$(y,z,h,q_1,q_2)$, and $rH$ is a positive unit. The complete hard
block is $\operatorname{diag}(16rH,16rH+8\eta^2)$; its only
nonzero first-normal residue is $16\eta x$. Its coupling to the
$8x^2$ pivot has determinant $128rH>0$. Retain the old positive
$g_4^2$ square. The resulting eight-root remainder Gram is strict.

Finally restore all four positive terms, the positive square just
retained, and the positive pure-coefficient multiplier. The state and
angle cases above exhaust the independent tests. Complex states split
into their real and imaginary quadratic evaluations. The result
closes the entire stated source domain, rather than only an optimizing
or matching testing chart.



% Source formulas: pa7e744614f49
\subsection{Unit $\eta$ and infinitesimal $e$: local polynomial}

\label{component:pa7e744614f49}

\begin{proposition}[Local part of the small-$e$ theorem]
Use the unchanged physical native C1 source, with $k,\zeta,e$ positive
infinitesimals and $\eta$ a positive bounded unit. Every native positive
$r=1-a\eta$, and every positive opening rate, is allowed.
At every independent test with $u,v,w$ infinitesimal, $Z=1$, and $w>0$,
the entire original parent is tensor SOS. Literal $\beta=0$ is included.
\end{proposition}
This proposition is the contact part only. The independent outer/state
join is supplied separately; no assertion here identifies a testing rate
with a source rate.

\paragraph{The source entry forced by the original sign.}
Keep $K,Q,E,\ell,\delta$ from (CA3), with $E=b^2$ and
$\ell=Qr=K-aE$. Here $E$ is a positive unit, $0<\ell<K$, and no lower
bound on $a$ or $\ell$ is assumed. Put
\[
 \varepsilon=1/\delta,\qquad j=\delta k,\qquad t=\varepsilon^2\tau.
\]
Both $\varepsilon,j$ are infinitesimal: $e=\varepsilon K/Q$ and
$\zeta=j^2QE/K$. The original condition $\beta\ge0$ gives
\[
 \frac{t}{e^2}\le\frac{2c_b}{\eta c_d},
\]
whose right side is bounded after cancelling the common native $r$
factor. Thus $\tau$ is positive bounded. The exact chart is
\[
 k=\varepsilon j,\quad K=1+\varepsilon^2j^2,\quad
 Q=K^2-\ell\varepsilon^2j^2,\quad a=(K-\ell)/E,
\]
\begin{align*}
 \eta&=KE/Q,\quad r=\ell/Q,\quad H=r+\eta,\quad
 e^2=\varepsilon^2K^2/Q^2,\\
 \zeta&=j^2QE/K,\quad R=K^2/Q,\quad T=K-k^2t\zeta. \tag{SE1}
\end{align*}
All denominators displayed here are positive units. The cancelled
coquotients $\widetilde c_i=c_i/\ell$ are the same short exact formulas
as in (CA3); the cancellation introduces no inverse of $\ell$.
Their residues are $2(E+\ell)$ for $i=2,b,3,5,\alpha$, and $4$ for
$i=d$; $c_4$ has residue $2$. These formulas hold uniformly as
$\ell$ tends to zero or to one.

\paragraph{One bounded centered remainder, including both normals.}
Set
\[
 A_* =\frac{ArH}{1+rH},\qquad
 A-A_* =\frac{A}{1+rH}>0.                              \tag{SE2}
\]
Retain the entire square $4(A-A_*)D_uD_v(Y-r uW)^2$.
It and its full first derivative vanish at exact native contact.
Divide the remaining parent by the positive pure coefficient $a$;
every factor $a$ cancels before coefficients are used. Write
\[
 \begin{array}{ll}
 b_2=8\widetilde c_2e^2/(R\widetilde c_3),&
 b_b=8\widetilde c_be^2/(R\widetilde c_3),\\
 b_3=8/T,&b_5=8\widetilde c_5/(T\widetilde c_3),\\
 L=Ar/a=2\eta^2c_4/(Q\widetilde c_3),&
 d_0=d/a=2\eta\widetilde c_d/\widetilde c_3,\\
 A_0=A_*/a=LH/(1+rH),&
 \alpha_0=\alpha/a=2\varepsilon^4j^6Q^2\eta^2
       \widetilde c_\alpha/(K^2\widetilde c_3).
 \end{array}                                                     \tag{SE3}
\]
Every quantity in (SE3) is bounded with unit denominators, and
$b_2,b_b$ contain $\varepsilon^2$.
Use the original normals $N,M_B$, put $p=t/R$, $h=a/T$, and write
\begin{align*}
 q_A&=vD_uw,\qquad q_B=v(D_u-uN-puD_uv w),\\
 W&=w+u(M_B+hq_B)/D_u,\qquad Y=uw-(M_B+hq_B)/D_u.
\end{align*}
The entire bounded remainder is
\begin{align*}
G_s={}&b_2\{N^2-2q_AN+(1-p)q_A^2\}
 +b_bD_v\{N^2+2q_AN+pq_A^2\}\\
 &+b_3D_v\{M_B^2+2q_BM_B+hq_B^2\}
 +b_5\{M_B^2-2q_BM_B+(1-h)q_B^2\}\\
 &+4A_0D_uD_v(Y-r uW)^2
 +4LD_uD_v(2uWY-r u^2W^2)\\
 &+4\alpha_0D_u(vY-1)^2
 +2d_0D_u\{(vY-1)^2-D_v\}
 +\nu D_vN^2,\qquad \nu=d_0(R-t)/R.                 \tag{SE4}
\end{align*}
This is (SD4), after the indicated exact centered subtraction and
normalization, except that a positive part of the original $\Pi_2$ is
kept inside the remainder. Indeed the original coefficient of
$D_vN^2$ in $\Pi_2/a$ is
\[
 \frac{4\beta+2d}{a}\frac{R-t}{R}\ \ge\ 2\nu.
\]
Thus subtracting only $\Pi_2/a-\nu D_vN^2$ remains legitimate, even
when $\beta=0$. The other three original terms are retained unchanged.
All full first normals and every hard quadratic term are present in
(SE4); there are no terms of higher normal degree.

\paragraph{A conductor that never divides by $\varepsilon$.}
Define the separated monomials and differences
\[
 A=\varepsilon wv,\quad B=v,\quad C=Euw,\quad
 D_0=\varepsilon^2j^3E,\quad k=\varepsilon j,
\]
\[
 x=B-C,\quad y=A-kC,\quad z=uB-kC,\quad h_0=kC-D_0.
\]
The two standard radical pairings are obtained by dividing the original
four squared circuit monomials by $\delta$: they are
$A,uB,kC,D_0$, and their product is $(k^2BC)^2$.
Their strict circuit is
$T_4=y^2+z^2+h_0^2-2yz+2q_1^2$.
The exact contact is $(u,w,v)=(k,j,\varepsilon j^2E)$.
Besides the three basic conductor identities,
\[
 C(u-k)=z-ux,\quad
 \varepsilon C(w-j)=y-\varepsilon wx,\quad
 C(v-\varepsilon j^2E)=E(wz+jy-kwx),                \tag{SE5}
\]
use the additional, division-free identity
\[
 C v(w-j)=Ew(w-j)z+Ejw y-Ekw^2x.                    \tag{SE6}
\]
For a monic contact division
$S=(u-k)T_u+(w-j)T_w+(v-\varepsilon j^2E)T_v$,
whenever $T_w=\varepsilon U+vV$, the bounded conductor is
\[
 \begin{aligned}
 CS={}&x[-uT_u-\varepsilon wU-Ekw^2V-EkwT_v]\\
 &+y[U+EjwV+EjT_v]
 +z[T_u+Ew(w-j)V+EwT_v].
 \end{aligned}                                      \tag{SE7}
\]
No division by $\varepsilon,j,k,r,a,t$, or a contact difference occurs.

Here is the full rational membership argument required for applying
(SE7), not only a finite truncation. At zero normals let
$V_0=1-puv w$. Then
\[
 q_A=D_u wv,\quad q_B=D_uvV_0,\quad
 uW=C/E+u^2hvV_0,\quad Y=C/E-hvV_0.
\]
The scalar of (SE4) therefore has an exact raw allocation
\[
 f_AA^2+f_{BB}B^2+f_{BC}BC+f_{CC}C^2+\rho D_0^2.     \tag{SE8}
\]
Expand only the nonconstant part of $(vY-1)^2$ into $B,C$; its constant
has $\alpha_0/D_0^2=2\widetilde c_\alpha/\widetilde c_3$.
The coefficients in (SE8) have unit denominators independent of $w$.
Every dependence on $w$ in them comes through $V_0$, and hence has a
factor $\varepsilon^2$; modulo $\varepsilon$ they are independent of
$w$ and even in $u,v$. Both complete first normals are in $(A,B,C)$
with bounded coefficients by the same exact formulas.

Clear the finitely many unit denominators, chosen even in $u,v$ and
independent of $w$, and use the unchanged scalar square-lift algorithm.
For its first scalar conductor, the only $w$-dependent coefficient
$T_w$ is divisible by $\varepsilon$ (in fact by $\varepsilon^2$).
At $\varepsilon=u=0$, its residual polynomial $S$ is even in $v$,
and contact makes its constant zero; thus $S\in(v^2)$ there.
More explicitly, modulo $\varepsilon$ the first divided differences
satisfy $T_u\in(u)$, $T_w=0$, and $T_v\in(v)$ after the preceding
monic substitutions. It follows
that after this first conductor, the sole possible $w$ dependence
modulo $\varepsilon,u$ in the four second-stage remainders is of the
form $w v$ times a polynomial in $v$. Hence their $T_w$ all belong to
$(\varepsilon,v)$, exactly the hypothesis of (SE7). The complete
normal raw coefficients are independent of $w$ modulo $\varepsilon$,
so their one conductor has $T_w\in(\varepsilon)$. The retained contact
squares and the original native balances supply the exact contact
value and first-jet identities throughout. This proves the full
scalar-square and complete first-normal lifts over the bounded ring.

\paragraph{Finite residues of that full lift.}
Give $(\varepsilon,j,u,w,v,N,M_B)$ weights $(1,1,2,1,3,5,3)$;
$E,\ell,\tau$ have weight zero. Every positive-weight variable is
infinitesimal, while all weight-zero coefficients are bounded.
The contact substitutions preserve the grading; $x$ has weight three
and the other circuit roots weight five. Both (SE5) and (SE6) preserve
weights, so after the scalar collection through weight ten and full
first normals through weight five, every lifted remainder coefficient
has positive weight. Restore the unit denominators afterward.

Set $H_0=E+\ell$ and $V_*=1+\ell H_0$. Direct exact multiplication of
(SE1)--(SE4), followed by (SE7), gives scalar residue
\[
       \frac{8(E+2)}{H_0V_*}x^2+8T_4,                         \tag{SE9}
\]
and complete first-normal residues
\[
       L_N=0,\qquad L_M=\frac{16E(H_0+1)}{H_0V_*}\,x.
\]
The hard block has residue
\[
 \operatorname{diag}\left(\frac{4E}{H_0},
           \frac{8(E^2+2E\ell+2\ell^2+2)}{V_*}\right).       \tag{SE10}
\]
These statements include coefficient corrections of positive weight;
they are not asserted as literal equalities of the full rational scalar.

The only hard-soft coupled determinant is
\[
 \det G_{M,x}
 =\frac{64(E^2+2E\ell+4E+4\ell)}{H_0^2V_*}>0.              \tag{SE11}
\]
Together with $4E/H_0>0$ and the strict two-pairing circuit this gives
the full bounded strict eight-root Gram. All original positive terms,
the centered square, and the positive pure-coefficient factor $a$ are
then restored. All inverses and radicals remain separated between the
two independent real closed fields.



% Source formulas: p05a86be0046c
\subsection{Unit $\eta$ and infinitesimal $e$: complementary tests}

\label{component:p05a86be0046c}

\label{p05a86be0046c:sec:small-e-outer}

\begin{lemma}
Use the unchanged physical native C1 source with $c=2$,
$\lambda=n=0$, positive original openings, and nonnegative original
event coefficients. Suppose $k,\zeta,e=k\sqrt{\eta/\zeta}$ are positive
infinitesimals, $\eta$ is a positive bounded unit, and
$r=1-a\eta>0$. Then every independent projective test and state is
tensor SOS except possibly the chart
\[
 u,v,w>0\text{ infinitesimal},\qquad Z=1,
 \qquad w=(W+uY)/(1+u^2).
 \tag{SEO1}
\]
There is no restriction on the positive relative scales of $a,r,J$
or on $k^2/r$. This lemma supplies a join, not a certificate for the
remaining chart (SEO1).
\end{lemma}

\paragraph{The small weak-event scale.}
Put $\kappa=k^2$, $K=1+\kappa$, $H=r+\eta$,
$V_c=\kappa\eta\zeta$, $S=1-\kappa\eta+tV_c$,
$R=1+\kappa r$, $T=K-\kappa\zeta t$,
$M=R-t$, and $J=T-a$. Let $P_4,P_d$ be the short native
polynomials printed in the centered-A outer lemma. All of
$K,H,S,P_4,P_d$ are positive bounded units. Set
\[
 L=KHS,\quad N=KH(1+V_c)-J\eta P_d,\quad
 d_0=4\eta P_d,\quad A_0=2\eta^2P_4,\quad D_0=2N.
\]
After exact cancellation of the native $r$ factors, the source is
\[
 d=\frac{ad_0}{L},\qquad A=\frac{aA_0}{rL},\qquad
 D=\frac{aD_0}{JL},\qquad
 \frac{m}{ae^2}=
 \frac{2(1-t\zeta)(1+V_c)}{MS}.                 \tag{SEO2}
\]
Here $\bar N=1$. Native $\beta\ge0$ gives the exact bound
\[
 \frac{t}{e^2}\le\frac{2c_b}{\eta c_d}
 =\frac{K[1+(1-t)\zeta]H}{\eta P_d}.
\]
Thus $t$ is infinitesimal, $M\sim1$, and $m/(ae^2)\sim2$.
In particular
\[
 \frac{k^2}{e^2}=\frac\zeta\eta
 \quad\hbox{is infinitesimal}.                       \tag{SEO3}
\]
The ratios $d/a,A/a,D/a$ have positive real lower bounds, while
$d/a$ is a bounded unit. These statements allow arbitrary positive
opening rates and literal $\beta=0$.

\paragraph{Choose the critical parameter without a small denominator.}
There are two source cases. If $a$ is a bounded unit, keep $A,D,d$
and replace $c=2$ by the critical number
\[
 c_* =\frac{d(d+2A)(d+2D)}{2(4AD-d^2)}.
 \tag{SEO4}
\]
This is positive and bounded, and $c_*-2=O(k^2)$.
Here is an exact short verification of the potentially delicate bound.
Define
\[
 Q=4A_0D_0-rJd_0^2,
\]
so $\bar Q=16\bar\eta^2\bar a\bar H>0$. Then
\[
 c_*-2=
 \frac{ad_0(d_0r+2A_0)(d_0J+2D_0)-4QL}{2QL}.
 \tag{SEO5}
\]
After $r=1-a\eta$ is inserted, its polynomial numerator vanishes
identically at $\kappa=0$, with $\zeta$ still free. It is therefore
$\kappa$ times a polynomial in bounded native variables; its
denominator is a positive unit. No inverse of $r$ or $J$ occurs.

If $a$ is infinitesimal, then $r,J\sim1$. Divide the whole parent
by $a$ and replace the original $A$ by
\[
 A_* =\frac{d^2(4+d+2D)}{2(8D-2dD-d^2)}.       \tag{SEO6}
\]
The denominator is positive. In this case
$A_*/a,D/a,d/a$ are positive bounded units, $c/a=2/a$ is
unbounded, and $(A_*-A)/a=O(k^2)$. To verify this without a
rate comparison, put
\[
 B_0=8D_0L-a(2d_0D_0+d_0^2J).
\]
Its residue is $16\bar H>0$, and
\[
 \frac{A_*-A}{a}=
 \frac{r d_0^2\{4JL+a(d_0J+2D_0)\}-2A_0B_0}
 {2rLB_0}.                                           \tag{SEO7}
\]
Again the numerator vanishes identically at $\kappa=0$, so it has
an exact polynomial factor $\kappa$, with a unit denominator.

Use $\sigma=1$ in the first case and $\sigma=a$ in the second,
and henceforth normalize all parent coefficients by $\sigma$.
Let $\epsilon$ be the positive part of the critical increase in the
chosen normalized coefficient. From (SEO2)--(SEO7),
\[
 \epsilon\text{ is infinitesimal},\qquad
 \frac{\epsilon}{m/\sigma}=O(\zeta)\text{ is infinitesimal}.
 \tag{SEO8}
\]
If the chosen critical value is smaller than the actual coefficient,
retain the positive excess term instead. In either source case
the chosen critical parent obeys
\[
 2c(4AD-d^2)=d(d+2A)(d+2D),                  \tag{SEO9}
\]
with the letters in (SEO9) now denoting the normalized critical
coefficients. In the first case $c,d$ are units and $A,D$ are
bounded below by positive units; in the second $A,D,d$ are units
and $c$ may be unbounded.

\paragraph{A critical-wall frame permitting the large coefficients.}
Write $D_u=1+u^2$, $D_v=1+v^2$ and
\[
 B=X+uZ,\quad U=uW,\quad V=vZ,\quad J_t=U-Y,
 \quad C=4D+2d,
\]
and define the pure-coefficient numbers
\[
 y=\frac d{2A},\quad
 b=\frac{4AD-d^2}{AC},\quad s=1-b,\quad q=y+b.
\]
The critical identity gives the useful exact relations
\[
 b=\frac C{4c+C},\quad Cs=4cb,\quad
 cb=\frac{d(2A+d)}{4A}.                         \tag{SEO10}
\]
In both source cases, $y,b,s,q$ are bounded and $b,s>0$.
The entire signed core plus the three critical terms is exactly
\begin{align*}
 F_*={}&4cD_vJ_t^2+4AD_uD_vY^2
       +C(J_t+uvB-D_uV)^2+2dD_vB^2\\
       &+2dD_u\{v^2Y^2-2YV-V^2\}.              \tag{SEO11}
\end{align*}

In the first source case use
\[
 H_A=\sqrt A(Y-yV),\qquad
 H_D=\sqrt C(J_t+uvB-(b+u^2)V).
 \tag{SEO12}
\]
The inverse is $Y=yV+H_A/\sqrt A$ and
$J_t=(b+u^2)V-uvB+H_D/\sqrt C$.
All coefficients in (SEO11) after this substitution are bounded.
Indeed the large $A,C$ occur only in the normalized squares and in
the bounded combinations $Ay^2$, $\sqrt A\,y$, $Cs^2$,
$\sqrt C\,s$; use $Cs=4cb$ and bounded $c$ for the last two.
At $u=v=0$, the hard block on $(B,H_A,H_D)$ is
\[
 \operatorname{diag}(2d,4,1+4c/C)\succ0.
\]
On its exact zero section the scalar coefficient of $V^2$ is
\[
 g u^2+\gamma v^2+4cu^4+g(y+2)u^2v^2+4cu^4v^2,
 \quad g=\frac{d(2A+d)}A,\quad
 \gamma=\frac{Dd(2A+d)^2}{A^2(2D+d)}.             \tag{SEO13}
\]
Both $g,\gamma$ are positive bounded units. Every hard-$V$ entry
has angular degree at least two, with bounded coefficients; this is
an exact full-parent polynomial statement. Allocate it on $uV,vV$,
leaving a factor $u$ or $v$ in its coefficient. The resulting exact
five-root Gram on $(B,H_A,H_D,uV,vV)$ therefore has strict residue.

In the second source case set $L_c=4c+C$ and instead use
\[
 H_A=\sqrt A(Y-yV),\qquad
 H_c=\sqrt{L_c}(J_t-bV).                         \tag{SEO14}
\]
Its coefficients are bounded because $A,D,d$ are units and
$c/L_c$, $cb/\sqrt{L_c}$ and $cb^2$ are bounded by (SEO10).
Its hard block is $\operatorname{diag}(2d,4,1)$.
The exact zero-section scalar is
\[
 g u^2+\gamma v^2+C u^4+gy\,u^2v^2.
 \tag{SEO15}
\]
The same complete hard-$V$ allocation gives a strict five-root
Gram on $(B,H_A,H_c,uV,vV)$. Thus this is a full critical-wall
certificate in both source cases, even with the stated unbounded
coefficients. Every radical and denominator used in the roots is
pure coefficient; no mixed inverse is used.

\paragraph{All unit-$W$ states at the small-angle corner.}
Normalize $W=1$, with $Z,Y,X$ bounded, and take $u,v$ infinitesimal.
Retain the complete original $\alpha,\beta$ terms. Put
\[
 R_m=v-X+uvY-uZ=v(1+uY)-B,\qquad
 V=\frac{Z}{1+uY}(R_m+B).                        \tag{SEO16}
\]
The displayed denominator is a positive pure-testing unit.
In the first source case the only negative correction to the
critical parent is $-4\epsilon D_vJ_t^2$; (SEO12) and (SEO16)
express $\sqrt\epsilon J_t$ as a bounded combination of the five
critical roots and $\sqrt{m/\sigma}\,R_m$, with every coefficient
having zero residue. In the second case the only negative
correction is $-4\epsilon D_uD_vY^2$; use
$Y=yV+H_A/\sqrt A$ and (SEO16) for the identical conclusion.
The only possibly delicate coefficient is a bounded testing
multiple of $\sqrt{\epsilon/(m/\sigma)}$, and it is infinitesimal
by (SEO8). The extra original $m$ square contributes the strict
diagonal $4$. Hence the complete perturbed Gram remains strict.
This includes unit or infinitesimal $Z$, literal $Z=0$, and both
signs of the bounded testing state. It uses the entire original
weak $m$ event, without a false positive-unit lower bound on $m$.

\paragraph{Finite outer angles without the weak event.}
The same critical parent also gives the two finite outer charts.
In (SEO11) regard $B,Y,J_t,V$ as independent roots. Exactly,
\[
 F_*(0,v)=F_*(0,0)
       +v^2\{2dB^2+(4A+2d)Y^2+4cJ_t^2\}.       \tag{SEO17}
\]
The zero-angle form has three positive hard pivots and null ray
$(B,Y,J_t,V)=(0,y,b,1)V$. On this null ray the additional
coefficient is $\gamma v^2>0$. The normalized hard frames
(SEO12) or (SEO14) keep every coefficient bounded. Thus when
$u$ is infinitesimal and $v$ a bounded unit the full four-root
residue is strict. The infinitesimal critical-parameter deficit
is absorbed directly on these four roots; retain the entire
$m,\alpha,\beta$ terms.

For bounded-unit $u$ and infinitesimal $v$, set $V'=D_uV$.
The other exact identity is
\[
 F_*(u,0)=F_*(0,0;V')
        +u^2\{4AY^2+2d(V')^2/D_u\}.             \tag{SEO18}
\]
Use the same normalized critical roots with $V'$ in place of $V$.
In the first source case include $uvB$ in the $H_D$ root so that
the unbounded $D$ event remains isolated for nonzero small $v$.
The extra null mass in (SEO18) is positive in residue because
$d$ and $u$ are units. The bounded full four-root Gram is strict,
and the deficit is again a zero-residue correction. These two
identities include every state, not merely the normal restriction.

For completeness the same argument covers both finite angles of
nonzero residue, without a separate physical-source certificate.
The exact determinant of (SEO11) on $(B,Y,J_t,V)$ is
\[
 64cdD_u^2D_v^2
 \{A(2D+d)u^2+D(2A+d)v^2\}.                    \tag{SEO20}
\]
The determinant in the normalized frame (SEO12) is this number
divided by $AC$, namely
\[
 32cdD_u^2D_v^2
 \left\{u^2+\frac{D(2A+d)}{A(2D+d)}v^2\right\}.
\]
It is a positive unit whenever at least one finite angle has
nonzero residue. The hard restriction $V=0$ is explicitly
\[
 \left(4D_uD_v+\frac{2dD_uv^2}{A}\right)H_A^2
 +H_D^2+4cD_v(H_D/\sqrt C-uvB)^2+2dD_vB^2,
\]
which has a strict bounded residue. Hence its full residue Gram
is positive definite by (SEO20). In frame (SEO14), the determinant
is (SEO20) divided by $A(4c+C)$. Here $A,D,d$ are units and
$c/(4c+C)\sim1/4$, so it too is a positive unit. Its hard
restriction is
\[
 \left(4D_uD_v+\frac{2dD_uv^2}{A}\right)H_A^2
 +\frac{4cD_v}{4c+C}H_c^2
 +C(H_c/\sqrt{4c+C}+uvB)^2+2dD_vB^2,
\]
again strict. All entries of both full transformed Grams are
bounded as already proved. The infinitesimal deficit is a
zero-residue correction on these four roots. This settles every
finite angle pair outside the double-infinitesimal corner.

\paragraph{Projective infinity charts.}
For these charts normalize by $a$, clip the positive coefficients
$c/a,A/a,D/a$ at suitable positive real lower bounds, and retain
their positive original excess squares. Retain the whole original
$m,\alpha,\beta$ terms. The normalized $d/a$ is a unit.
All the infinity-infinity and infinity-zero four-root identities
in (BE2) and the subsequent outer routing apply without $m$.
For $u=\infty$ in residue and $v$ a bounded unit, the signed-core
residue itself retains
$2d\{(W+vX)^2+(vY-Z)^2\}$; together with $cW^2,AY^2$ it is
strict, so no $m(vY-Z)^2$ lower bound is needed.
For $v=\infty$ and $u$ a bounded unit the old signed-core and
$c,D$ root argument is unchanged.
At the remaining $(u,1/v)=(0,0)$ corner, replace the obsolete
state-$W$ mass argument by the exact root ideal
$(X,Y,U=uW,Z)^2$ after retaining $m$. In $q^2F(u,1/q)$ its
complete selected residue is
\[
 2dX^2+4c(U-Y)^2+(4A+2d)Y^2+4DZ^2\succ0.
 \tag{SEO19}
\]
Every remaining occurrence of $W$ carries $u$, so this is a full
bounded Gram including all states. Literal infinity is included
by these homogeneous identities. Literal zero axes have their
unchanged original-event SOS. Finite tests with both noncommuting
residues were included directly in (SEO20).

\paragraph{The residual state chart.}
Only $u,v$ infinitesimal remain. Normalize the state by a maximal
coordinate. On a negative pure-testing-core branch, $X,Y$ have
zero residue. The unit-$W$ chart was just proved for every bounded
$Z$. Otherwise $Z$ is a unit, so normalize $Z=1$ and $W$ is
infinitesimal. Then $w=(W+uY)/D_u$ is infinitesimal. The inherited
exact testing Schur identity and sign gauges give (SEO1).
Complex states split into their real and imaginary evaluations.



% Source formulas: pd4b92bf7de71
\section{Nonvanishing \texorpdfstring{$L=x\zeta$}{L=x zeta}}
For positive-unit or unbounded $L$, the reciprocal $z=1/L$ is bounded. When $\eta$ is unbounded, the uniform completion below combines two nonnegative source weights through their sum and covers every relative size of $h$ and the infinitesimal $\sigma=1/\eta$. The other source and testing domains retain their separate completions in complementary coordinates.

\subsection{Unbounded $L$ and unit $\eta$}

\label{component:pd4b92bf7de71}

\label{pd4b92bf7de71:sec:c14-large-l-unit}
This section proves this full original-parent subdomain, retaining every
state normal and every independent source/testing rate. It makes no claim
for infinitesimal $\sigma$. The common common remote theorem leaves
$u$ bounded and $w=1/v$ infinitesimal; both $u$ infinitesimal and $u$ a
bounded unit are covered below. All calculations use the exact NC source,
with $d=1$, $n=\lambda=0$, and its original nonnegative $D,\beta$ faces.

\paragraph{Source chart and bounded products.}
Put $x=k^2$, $z=q^2=L^{-1}$, $r=bz$, $Q=1-b$, $\theta=1-h$,
$K=1+x$, $R=1+xbz$, $M=R-xz\theta$, and
\[
 J=1+xh-\sigma(1-bz),\qquad E=1-xh/\sigma.
\]
Here $x,z$ are positive infinitesimals, $0<b<1$, and $\sigma$ is a positive
bounded unit. The NC source gives exactly
\begin{align*}
 P&=\sigma+K+\sigma z(1+bx+2xh)+\sigma bx^2hz^2+x^2h\theta z,\\
 V_0&=z(1+xh)E+(1-bz)(1+xhz),\\
 C&=cz,&c&=\frac{K(1+\sigma bz)E}{2(1-bz)P},\\
 m&=\mu x^2z,&\mu&=\frac{K(1+\sigma bz)h(1+\sigma z)}{2\sigma MP},\\
 A&=a/z,&a&=\frac{V_0}{2bP},\\
 \alpha&=\ell/z,&\ell&=\frac{(1+\sigma bz)[(x/\sigma)Rh(1+\sigma z)+ME]}{2QP},\\
 D&=\frac{K(1+\sigma z)(1+\sigma bz)}{2JP}-\frac12,&
 \beta&=\frac{K(1+\sigma bz)(1+xhz)}{2\theta P}-\frac12.
\end{align*}
In particular $c,\mu,ab,\ell Q,DJ,\beta\theta$ are bounded,
$c,\mu,ab,\ell Q,DJ$ are positive units, and $ab-\ell Q\in z\mathcal O$.
At $x=z=0$ their residues are
\[
\begin{gathered}
c_0=\frac1{2(1+\sigma)},\quad
\mu_0=\frac h{2\sigma(1+\sigma)},\quad
(ab)_0=(\ell Q)_0=\frac1{2(1+\sigma)},\\
(DJ)_0=\frac{\sigma^2}{2(1+\sigma)},\quad
(\beta\theta)_0=\frac{(1+\sigma)h-\sigma}{2(1+\sigma)}.
\end{gathered}
\]
Source nonnegativity gives $0<\sigma\leq1$ and
$h\geq\sigma/(1+\sigma)$ at residue. Thus $h$ is a positive unit.
The openings $b,Q,J,\theta$ remain independent; $\beta=0$ is retained.
The face $D=0$ cannot occur here since $DJ$ is a positive unit.

Write $D_u=1+u^2$, $D_w=1+w^2$ and
\begin{gather*}
 M_w=W-wX+uY-uwZ,\quad R_w=uW-Y,\quad
 J_w=uwW+uX-wY-Z,\quad B_w=X+uZ,\\
 K_w=J_w^2+D_u(Y-wZ)^2+D_wB_w^2-D_uD_wZ^2.
\end{gather*}
The complete reciprocal parent is
\[
 F_w=4mM_w^2+4czD_wR_w^2+4(a/z)D_uD_wY^2
 +4DJ_w^2+4(\ell/z)D_u(Y-wZ)^2+4\beta D_wB_w^2+2K_w.
\]
Choose bounded positive-unit caps $a_*\leq a$, $\ell_*\leq\ell$,
$D_*\leq D$, and a bounded nonnegative cap $\beta_*\leq\beta$.
Keep as nonnegative squares the four centered excesses
\begin{gather*}
4(a-a_*)D_uD_wN_y^2,\quad4(\ell-\ell_*)D_uN_y^2,\\
4(D-D_*)(J_w+JZ)^2,\quad4(\beta-\beta_*)D_wN_x^2,\\
N_x=X+uhZ,\qquad N_y=(Y-bwZ)/q.
\end{gather*}
These roots vanish at the exact source contact
\[
u=k,\quad w=q\sqrt\sigma,\quad
(W,X,Y,Z)=(\sqrt\sigma/(kq),-kh,bq\sqrt\sigma,1).
\]
Set $d_J=(D-D_*)J$ and $b_\theta=\beta\theta$.
Their coefficients and the centered excesses preserve the original parent.

\paragraph{Bounded full-normal remainder.}
Define $G=wZ/q$, $R_h=G-uqW$ and
\[
c_N=8D_u[abD_w-\ell Q],\qquad c_G=4D_u[ab^2D_w+\ell Q^2].
\]
The $A,\alpha$ remainder is
$4D_u(a_*D_w+\ell_*)N_y^2+c_NN_yG+c_GG^2$.
Here $c_N\in(z,w^2)\mathcal O$ and $c_G$ is a positive unit.
Keep $c_GR_h^2+c_NN_yR_h$ in the bounded Gram and use
$G^2-R_h^2=2uwWZ-u^2zW^2$. The rest is the bounded polynomial
\begin{align*}
G_0={}&4D_u(a_*D_w+\ell_*)N_y^2+c_NuqWN_y+c_G(2uwWZ-u^2zW^2)\\
&+4mM_w^2+4czD_wR_w^2+4D_*J_w^2-8d_JZJ_w-4d_JJZ^2\\
&+4\beta_*D_wB_w^2+8(b_\theta-\beta_*\theta)uD_wZB_w
 -4(b_\theta-\beta_*\theta)u^2\theta D_wZ^2+2K_w.
\end{align*}
All its coefficients are bounded. Put $N_z=Z-uwW/\sigma$ and substitute
$Z=N_z+uwW/\sigma$, $X=N_x-uhZ$, $Y=qN_y+bwZ$.
For $t=u^2$, $\xi=w^2$, define
\[
 T_0=1+th+b\xi,\quad M_0=1+t(h-Q)\xi/\sigma,\quad
 R_f=1-b\xi/\sigma,\quad J_f=1-T_0/\sigma.
\]
In the formulas below $c_N,c_G$ mean their even-polynomial expressions.
The exact expansion is
\[
G_0=\phi W^2+wL_xWN_x+uqL_yWN_y+uwL_zWN_z+H(N_x,N_y,N_z),
\]
where
\begin{align*}
\phi={}&4\mu x^2zM_0^2+4cz(1+\xi)tR_f^2+c_Gt(2\xi/\sigma-z)\\
 &+4D_*t\xi J_f^2-8d_Jt\xi J_f/\sigma-4d_JJt\xi/\sigma^2
 +4b_\theta\theta(1+\xi)t^2\xi/\sigma^2\\
 &+2[t\xi J_f^2+(1+t)tQ^2\xi^2/\sigma^2
 +(1+\xi)t^2\theta^2\xi/\sigma^2-(1+t)(1+\xi)t\xi/\sigma^2],\\
L_x={}&-8\mu x^2zM_0+t[(8D_*+4)J_f-8d_J/\sigma
 +(8b_\theta+4\theta)(1+\xi)/\sigma],\\
L_y={}&8\mu x^2zM_0-8cz(1+\xi)R_f+c_N-(8D_*+4)\xi J_f
 +8d_J\xi/\sigma-4(1+t)Q\xi/\sigma,\\
L_z={}&8\mu x^2z(h-Q)M_0-8cz(1+\xi)bR_f+2c_G
 -(8D_*+4)T_0J_f-8d_JJ_f+8d_J(T_0-J)/\sigma\\
 &+8b_\theta\theta(1+\xi)t/\sigma+4(1+t)Q^2\xi/\sigma
 +4(1+\xi)t\theta^2/\sigma-4(1+t)(1+\xi)/\sigma.
\end{align*}
For a compact explicit specification of the entire quadratic block, put
\[
\begin{gathered}
m_N=-wN_x+uqN_y+uw(h-Q)N_z,\quad
r_N=-qN_y-bwN_z,\\
j_N=uN_x-wqN_y-T_0N_z,\quad
b_N=N_x+u\theta N_z.
\end{gathered}
\]
Then every quadratic normal coefficient is retained by the finite formula
\begin{align*}
H={}&4D_u(a_*D_w+\ell_*)N_y^2+4mm_N^2+4czD_wr_N^2
 +(4D_*+2)j_N^2-8d_JN_zj_N-4d_JJN_z^2\\
&+(4\beta_*+2)D_wb_N^2+8(b_\theta-\beta_*\theta)uD_wN_zb_N
 -4(b_\theta-\beta_*\theta)u^2\theta D_wN_z^2\\
&+2D_u(qN_y-QwN_z)^2-2D_uD_wN_z^2.
\end{align*}
The exact source substitutions give, at $(t,\xi)=(x,\sigma z)$,
\[
\phi=\phi_t=\phi_\xi=L_x=L_y=L_z=0.
\]
Moreover $\phi(0,\xi)=4\mu x^2z$,
$L_x(0,\xi)=-8\mu x^2z$, and $\phi(t,\sigma z)\in z\mathcal O[t]$.

\paragraph{Exact six-root lift in the critical strip $u$ infinitesimal.}
Monic division, first by $\xi-\sigma z$ and then by $t-x$, gives bounded
polynomial coefficients satisfying
\begin{align*}
\phi&=z(t-x)^2 A+t(t-x)(\xi-\sigma z)B+t(\xi-\sigma z)^2C,\\
L_x&=(t-x)(tA_x+zB_x)+t(\xi-\sigma z)C_x,\\
L_y&=(t-x)A_y+(\xi-\sigma z)B_y,\\
L_z&=(t-x)A_z+(\xi-\sigma z)B_z.
\end{align*}
For example, the constant coefficient of
$L_x(t,\sigma z)/(t-x)$ equals $8\mu xz$, giving its $(t,z)$ support.
No small or mixed inverse enters these monic quotients. Also $A(0)=4\mu$.
Take the six finite separated roots
\[
 N_x,\quad N_y,\quad N_z,\quad R_h,\quad
 A_\theta=q(t-x)W,\quad V=uw(t-x)W.
\]
The complete scalar lift follows from
\begin{align*}
 z(t-x)^2W^2&=A_\theta^2,\\
 t(t-x)(\xi-\sigma z)W^2&=u\sigma A_\theta R_h-\sigma VN_z,\\
 t(\xi-\sigma z)^2W^2&=\sigma^2(qR_h-wN_z)^2.
\end{align*}
All first normals lift by the following finite list, which includes their
whole coefficient support:
\begin{align*}
wt(t-x)W&=uV,&wz(t-x)W&=qwA_\theta,\\
wt(\xi-\sigma z)W&=u\sigma qwR_h-u\sigma w^2N_z,\\
uq(t-x)W&=uA_\theta,&uq(\xi-\sigma z)W&=\sigma zR_h-\sigma qwN_z,\\
uw(t-x)W&=V,&uw(\xi-\sigma z)W&=\sigma qwR_h-\sigma w^2N_z.
\end{align*}
Together with the displayed $H$ these give an exact bounded six-root Gram,
not merely a restriction to the contact line. At residue its first five
diagonal pivots are respectively
\[
2+4\beta_*,\quad4(a_*+\ell_*),\quad
4[D_*+d_J(1+\sigma)],\quad \frac2{1+\sigma},\quad4\mu_0.
\]
They are positive units through all allowed openings. The only potentially
nonzero remaining cross entry is $N_zV$. The exact relation
\[
 V^2=\sigma tA_\theta^2+u\sigma(t-x)A_\theta R_h
       -\sigma(t-x)VN_z
\]
introduces an arbitrarily large finite positive $V^2$ pivot with only
infinitesimal changes to all other entries. Its real Schur complement
therefore makes the residue strictly positive definite. The exact bounded
Gram is positive definite by the standard finite-dimensional local
criterion. All excess squares are nonnegative with separated coefficients.

\paragraph{Away from the angular contact: $u$ a bounded unit.}
The same source chart has a positive-unit angular gap at every such real
residue. A convenient raw finite frame is
$E_t=qW,F_t=Y/q,G_t=wZ/q,R_t=wW,X,Z$; products obey $E_tG_t=R_tZ$.
The residue parent divided by $2/(1+\sigma)$ is exactly
\[
u^2E_t^2+D_u\{F_t^2/b+(F_t-G_t)^2/Q\}
 +\frac{(uR_t+uX-Z)^2}{J_0}
 +\frac{(X+uZ)^2}{\theta}-(1+\sigma)D_uZ^2,
\quad J_0=1-\sigma.
\]
Eliminate the positive $F_t$ and $X$ directions and allocate
the common product between the remaining two quadratic blocks. Up to a
common positive factor the two blocks have diagonal products $AB,CD_0$
and total off-diagonal magnitude $E_0$, with $U=u^2>0$ and
\[
\begin{gathered}
\delta=1-\sigma+\theta U,\qquad
A=U,\quad B=1+U,\quad C=U/\delta,\\
D_0=(1+U)[(1+U)/\delta-(1+\sigma)],\quad
E_0=\sqrt U(1+U)/\delta.
\end{gathered}
\]
Both diagonal products are strictly positive, since
$1+U-(1+\sigma)\delta=\sigma^2+U[(1+\sigma)h-\sigma]>0$.
All square roots in this sentence are real-residue quantities used to
choose real allocation constants, not mixed tensor radicals.
The required strict inequality is $E_0<\sqrt{AB}+\sqrt{CD_0}$.
Set $N=U(2\sigma-\theta U)$. Directly,
\[
E_0^2-AB-CD_0=(1+U)N/\delta,\quad
N^2-4U^2[1+U-(1+\sigma)\delta]
 =U^3(\theta^2U-4h).
\]
If $N\leq0$ strictness is immediate. If $N>0$, then
$\theta^2U<2\sigma\theta\leq2h<4h$, again giving strictness.
The single-opening faces follow by their constrained limits. At the only
double opening $J=\theta=0$, so $\sigma=h=1$, the exact limit is
$E_0-\sqrt{CD_0}\longrightarrow\sqrt U$, whereas $AB-U=U^2>0$.
Here is a full-frame square completion which also justifies those cap
limits. For clarity first take $u>0$ and select a finite real
$\rho$ in the nonempty interval
\[
\max\{-\sqrt{AB},E_0-\sqrt{CD_0}\}<\rho
 <\min\{\sqrt{AB},E_0+\sqrt{CD_0}\}.
\]
Its nonemptiness follows from the strict allocation inequality above.
Put
\[
r_\rho=D_u/u-\rho\delta/U,\qquad
s_\rho=D_u[2\rho/u-(1+\sigma)]-\rho^2\delta/U>0.
\]
The preceding complete residue form equals
\begin{align*}
&UE_t^2+D_uG_t^2-2\rho E_tG_t
 +\frac{D_u}{bQ}(F_t-bG_t)^2\\
&\quad+\frac{\delta}{J_0\theta}
 \left[X+\frac{U\theta R_t+u(J_0-\theta)Z}{\delta}\right]^2
 +\frac U\delta(R_t-r_\rho Z)^2+s_\rho Z^2.
\end{align*}
The equality uses only $E_tG_t=R_tZ$ and is checked exactly.
For $u<0$ take the correspondingly signed $\rho$; all sign choices are
real constants in this allocation. All triangular minimizer coefficients
are bounded through every opening because $U$ is a unit and
$\delta=J_0+U\theta$. Their positive hard weights stay bounded below.
At the double face choose $|u|<|\rho|<\sqrt{U(1+U)}$, with matching
sign; then $s_\rho$ has a positive-unit limit. This proves uniform
coercivity of the constrained residue, not just determinant positivity.

For completeness the exact centered cap rows in the raw frame are
\[
 F_t-bG_t,\quad
 uR_t+uX-qwF_t+(J-1)Z,\quad X+uhZ.
\]
Only one of $b,Q$ can tend to zero. The other two rows have hard parts
on $(R_t,X)$ with determinant $u$, a unit; each single row has a unit
hard entry. Their simultaneous kernel is solved exactly by
\[
F_t=bG_t,\qquad X=-uhZ,\qquad
R_t=\{(u^2h-J+1)Z+qwbG_t\}/u.
\]
More generally, prescribed changes $(e_F,e_D,e_\beta)$ in these three
rows are produced by
$\Delta F_t=e_F$, $\Delta X=e_\beta$,
$\Delta R_t=(e_D-u e_\beta+qw e_F)/u$.
This is a bounded separated right inverse: its only inversion is the
pure-testing unit $u$. For the finite-cap existence argument fix $u$
at its ordinary real residue, so the same inverse is an ordinary real
finite-dimensional right inverse. No order on the tensor product is
used. It lifts the limiting constrained kernel to the exact centered
kernel while keeping every opening independent. The
bounded residual form restricted to that kernel is the strictly positive
constrained form just computed. An elementary block Schur complement
then provides finite sufficiently large caps on every divergent row;
bounded rows retain their full coefficients. This is the finite-cap
lemma used in the UH proof. It yields a strictly positive bounded
full-frame Gram. Exact coefficients differ from their residues
infinitesimally, completing this strip.
All divisions by $\delta$ and radicals in this paragraph belong only
to the ordinary real residue calculation and its face limits. The exact
tensor Gram uses the polynomial raw frame, real allocation constants,
and finite source caps; it introduces no mixed inverse or radical.

Consequently this section closes the $L$ unbounded, $\sigma$ bounded-unit
subdomain of C1.4 under the original parent, including all state normals,
all independent opening rates, and $\beta=0$. It does not close the
$\sigma$ infinitesimal subdomain or claim a whole-row count change.



% Source formulas: p080aa62de16d
\subsection{Unit $L$ and infinitesimal $\eta$: C1 substitution in LII}

\label{component:p080aa62de16d}

\label{p080aa62de16d:sec:r7-c1-L-unit-eta-small}

\begin{theorem}
The entire original C1 source subdomain
\[
 x=k^2\in\Inf_{>0},\qquad L=x\zeta\text{ a positive bounded unit},
 \qquad \eta\in\Inf_{>0}
\]
has a full-parent tensor-SOS certificate for every independent
projective testing pair.  All original openings and the retained
boundaries $D=0$ or $\beta=0$ remain allowed whenever they occur
in this source subdomain.
\end{theorem}

The proof is an exact adapter to the strong-II local and remote
certificates.  Its additional event is retained with a centered
square; it is not deleted from the parent.

\paragraph{The native source in the shared contact geometry.}
Keep the native notation of (NC1)--(NC4), but in this section denote
its $r,h,a$ by $r_n,h_n,a_n$ to avoid collisions.  Normalize $d=1$.
The exact native contact with $Z_c=1$ is
\[
 (u_c,v_c;W_c,X_c,Y_c,Z_c)
 =\left(\sqrt x,\sqrt{x\eta\zeta};
 \sqrt{\zeta/\eta},-\sqrt x\,h_n,
 \sqrt x\,r_n\sqrt{\zeta/\eta},1\right).
\]
It therefore has the same contact geometry used in (LII2), with
\[
 z=x,\quad y=L\eta,\quad B=L,\quad H=Lr_n=B-a_ny,
 \quad r=h_n=1-\theta,\quad t=a_n,\quad
 j=1+zr-t=J_n.                                           \tag{LC1}
\]
In particular $z,y$ are positive infinitesimals.  Here $r_n\sim1$,
$0<Lr_n<1$, and hence $b_0:=\operatorname{res}B$ satisfies
$0<b_0\le1$.  We use $b_0$ only for this source residue.

Set, as in the local strong-II proof,
\[
 c=Cyt,\quad d_p=Dyj,\quad
 \mu=m/z^2,\quad a_A=A/y,\quad
 \sigma=\alpha(1-H),\quad \tau=\beta(1-r)=\beta\theta.
\]
These are bounded source quantities, and $c,d_p,\mu,a_A,\sigma$
are positive units.  Writing $r_0=\operatorname{res}r$, exact
substitution in (NC2)--(NC4), with opening factors canceled first,
gives
\[
\begin{gathered}
 c,d_p\sim\frac{b_0}{2(1+b_0)},\qquad
 a_A\sim\frac1{2b_0(1+b_0)},\\
 \mu\sim\frac{r_0}{2b_0(1+b_0)},\qquad
 \sigma\sim\frac{b_0^2}{2(1+b_0)},\\
 \operatorname{res}\tau=
 \frac{b_0}{2(1+b_0)}-\frac{1-r_0}{2}.                    
\end{gathered}
\tag{LC2}
\]
Since $\tau\ge0$, the last identity implies
$r_0\ge1/(1+b_0)>0$, proving the asserted unit lower bound for
$\mu$.  Also $t+j=1+zr\sim1$, and $c,d_p$ being units proves
$C,D$ unbounded uniformly through either opening $t=0$ or $j=0$
in residue.  Thus $D=0$ is in fact absent from this specific
subdomain by the exact source implications; it has not been excluded
by an additional hypothesis.  Literal $\beta=0$ is included.

The source-harmonic quantity $yCD/(C+D)$ and the cap behavior of
$\alpha$ have precisely the residues (RR1)--(RR2), with
$h=b_0$.  The remote proof (RR3)--(RR13) therefore applies unchanged,
retaining the extra positive term
$4\beta D_v(X+uZ)^2$ throughout.  This covers every test outside
the local region where $u,v$ are both infinitesimal, including the
original projective endpoints.

\paragraph{The additional event and its exact local remainder.}
For the local region use exactly the graph, roots, harmonic pair,
and caps (LII2)--(LII4), with the source quantities (LC1)--(LC2).
Put
\[
 R=v(U-z)W,\qquad
 E_\beta=\xi+u\zeta+R/B
          =(X+uZ)-\frac{z(1-r)v}{B}W.                    \tag{LC3}
\]
This root vanishes at the actual contact.  Retain the entire
positive term $4\beta D_vE_\beta^2$, in addition to the first,
second and fourth terms of (LII5); there is no $n$ term since
$n=0$.  The exact remainder from the $\beta$ event is
\[
 \begin{split}
 \Delta_\beta={}&
 \frac{8\tau z}{B}D_vvW(\xi+u\zeta)\\
 &+\frac{4\tau z}{B^2}D_vV\{2U-z(1+r)\}W^2.
 \end{split}                                             \tag{LC4}
\]
All its coefficients are bounded even if $\beta$ is unbounded:
only $\tau=\beta(1-r)$ occurs.  The exact full remainder is (LII6)
with its $n$ line deleted, plus (LC4).  Thus no state normal,
boundary coefficient, or first contact jet has been omitted.

The $R_hW$ gauge (LII7) is unchanged, because (LC4) contains no
$R_h$.  The coefficient representation (LII8), its exact contact
and first-normal vanishings (LII9)--(LII10), and its bounded monic
divisions (LII12)--(LII13) all remain valid.  Explicitly, (LC4)
adds $8\tau z(1+V)/B$ to $P_\xi$ and $P_\zeta$, adds zero to
$P_N$, and adds its printed scalar line to $\phi$.  Hence the
essential support $P_\xi\in(U,z)$ and $\deg_U\phi\le2$ remain
literal polynomial statements.

The only changed data in (LII11) are
\[
\begin{split}
 d_s&=\frac{4z^2}{B^2}
   \left\{\mu(B+zr)^2+\frac{r^2}{2}-\tau(1+r)\right\},\\
 e_s&=\frac{4z^2}{B^2}
   \left\{\frac{r^2}{2}-\tau(1+r)\right\}.
\end{split}                                             \tag{LC5}
\]
Both quotients by $z^2$ are bounded; $e_s$ need not be positive.
Its contribution to the final Gram has an infinitesimal coefficient
because it is multiplied by $V$.  The decisive exact residue
cancellation, using (LC2), is
\[
 \begin{split}
 \operatorname{res}\left\{
 \mu B^2+\frac{r^2}{2}-\tau(1+r)\right\}
 &=\frac{b_0r_0}{2(1+b_0)}+\frac{r_0^2}{2}
 -\left\{\frac{b_0}{2(1+b_0)}-\frac{1-r_0}{2}\right\}(1+r_0)\\
 &=\frac1{2(1+b_0)}.
 \end{split}                                             \tag{LC6}
\]
Consequently
\[
 \operatorname{res}(d_s/z^2)=\frac2{b_0^2(1+b_0)},\qquad
 \operatorname{res}A_\xi=\frac4{b_0(1+b_0)}.              \tag{LC7}
\]
The second equality is unchanged from (LII18): the new $P_\xi$
term is a bounded multiple of $z$ independent of $U$, so it has
zero $U$-quotient residue.  The four hard-normal residues (LII17)
are unchanged as well.

Apply the exact allocations (LII15) to all the coefficients,
including (LC4).  They give the same bounded six-root Gram
$(N,\xi,\zeta,R_h,R,T)$ as in (LII16).  Its finite $(\xi,R)$
block is exactly (LII19), with $h=b_0$, and has strictly positive
Schur complement $2/[b_0(1+b_0)^2]$.  The same exact conductor
(LII20) supplies an arbitrarily large real $T^2$ pivot while
changing every other residue entry by zero.  The complete bounded
Gram is strict.  Restoring the four retained terms, including
$4\beta D_vE_\beta^2$, proves the local tensor SOS for all
independent rates and literal affine axes.  The remote certificate
above completes the theorem.

\paragraph{Source inequalities.}
  Stationarity
itself is the unchanged native source identity, not an added
restriction.  The result removes the whole $L$-unit, $\eta$-infinitesimal
subdomain from C1.4.  It does not assert the other $L$-unit source
regimes or the entire C1.4 row.



% Source formulas: p686209b87229
\subsection{Unit $L,\eta$: C1 substitution in LII}

\label{component:p686209b87229}

\label{p686209b87229:sec:r7-c1-L-eta-unit}

\begin{theorem}
For the original C1 source with $x\in\Inf_{>0}$ and both
$L=x\zeta$ and $\eta$ positive bounded units, the entire original
parent is tensor-SOS for every independent projective testing pair.
All original positive openings and literal $\beta=0$ are included.
\end{theorem}

The exact source formulas force $D$ strong on this subdomain, so
$D=0$ is absent by source implication.  No opening is removed by
an extra hypothesis, and the proof does not invoke a table or atlas.

\paragraph{The source and the only testing strip left by four events.}
Use the exact shared geometry (LC1), and write $\rho=r_n$ for the
native shape parameter and $r=1-\theta$ for the geometric one.
Thus
\[
 B=L,\quad y=B\eta,\quad H=B\rho=B-yt,\quad
 t=(1-\rho)/\eta,\quad j=1+zr-t,
 \quad z=x\in\Inf_{>0}.
\]
Here $B,y,\eta$ are units.  For the boundary coefficient formulas
below set
\[
 P=1+B+y,\qquad P_4=1+B(1-\rho),\qquad S=H+y=P-P_4.
                                                               \tag{LU1}
\]
All three quantities are positive units, uniformly through
$\rho\to0$, $1-\rho\to0$, $j\to0$, and $1-H\to0$.
The native formulas imply that $D,C,A,\alpha$ are strong.
One forward cyclic relabeling sends these four consecutive events
to $C,A,\alpha,m$ and sends testing coordinates to
$(u',v')=(v,1/u)$.  The full-parent four-consecutive theorem
(CE11)--(CE16) therefore covers every test except
\[
                    u\in\Inf,\qquad v\text{ a bounded nonzero unit}.
                                                               \tag{LU2}
\]
The cycle is the inverse of (CE9)--(CE10), so its literal endpoint
scope is the original homogeneous parent's scope.

It remains to prove (LU2).  All roots and coefficients used below
are valid for both testing signs, including literal $u=0$.

\paragraph{A bounded full-parent remainder with every opening retained.}
Use the graph, harmonic roots and gauge (LII2)--(LII8), with
$k=\sqrt y$ now a source unit.  The bounded source data are
\[
 c=Cyt,\quad d_p=Dyj,\quad
 \mu=m/z^2,\quad \sigma=\alpha(1-H),\quad
 \tau=\beta(1-r),\quad A_H=AH.
\]
They are bounded; $c,d_p,\mu,\sigma,A_H$ are units.  In residue,
with the boundary symbols of (LU1), their exact values are
\[
 c=\frac S{2P},\quad d_p=\frac{BP_4}{2P},\quad
 \mu=\frac{Sr(1+y)}{2B^2P},\quad
 \sigma=\frac{BS}{2P},\quad A_H=\frac{yP_4}{2P},\quad
 \tau=\frac S{2P}-\frac{1-r}{2}.                         \tag{LU3}
\]
These are residue identities, not substitutions of source and
testing fields.  They are obtained after multiplying away every
possibly small opening in the exact native formulas.
Since $\tau\ge0$, $r\ge P_4/P$ in residue and is a positive unit.
Also $t+j\sim1$, so $\ell=y(C+D)$ has a positive lower bound,
$q=D/(C+D)$ is bounded in $[0,1]$, and
$g=yCD/(C+D)$ is a positive unit.

Choose positive bounded caps $\ell_*\le\ell$,
$A_*\le A$, and $\alpha_*\le\alpha$.  Choose a bounded cap
$0\le\beta_*\le\beta$ as specified below.  Retain the five
contact-centered positive terms
\[
\begin{gathered}
 4(\ell-\ell_*)N^2,\qquad 4CV\eta_t^2,\qquad
 4(A-A_*)D_uD_v\eta_t^2,\qquad
 4(\alpha-\alpha_*)D_u(v\eta_t-\zeta)^2,\\
 4(\beta-\beta_*)D_v(\xi+u\zeta+R/B)^2,
 \qquad \eta_t=k(N+qR_h),\quad R=v(U-z)W.
\end{gathered}
\tag{LU4}
\]
The exact remainder is (LII6), with its $n$ line deleted, its $A$
term replaced by
\[
 4A_*D_uD_v\eta_t^2
 +\frac{8A_H}B D_uD_vuW\eta_t
 +\frac{4A_HH}{B^2}D_uD_vUW^2,
                                                               \tag{LU5}
\]
and with (LC4) plus $4\beta_*D_v(\xi+u\zeta+R/B)^2$ added.
This specifies every coefficient of the full remainder.  Only
bounded source products occur, even at an unbounded original weight.

The same dependent-root gauge and exact contact equations give
(LII9)--(LII13).  The new constant scalar coefficients are
\[
 \frac{d_s+e_sV}{z^2}=
 \frac4{B^2}\left[
 \mu(B+zr)^2+
 \left\{\frac{r^2}{2}-\tau(1+r)+\beta_*\right\}(1+V)
 \right].                                                   \tag{LU6}
\]
The support $P_\xi\in(U,z)$ and scalar degree at most two in $U$
remain exact.  Apply all of (LII15).  As $v$ is a unit, eliminate
$T$ by the exact pure-factor identity
\[
                        T=(uk/v)R.                         \tag{LU7}
\]
Its coefficient is a bounded infinitesimal tensor.  Thus the entire
remainder has a bounded Gram on
$(N,\zeta,R_h,\xi,R)$.  No source/testing ratio or mixed inverse
has been introduced.  We now verify its complete residue.

\paragraph{A three-square identity supplies all three state/angle normals.}
The following boundary computation is an exact rational identity
before taking any opening limit.  Define
\[
\begin{aligned}
 C_0&=\frac{\rho+\eta}{2(1-\rho)P},&
 A_0&=\frac{\eta P_4}{2\rho P},\\
 D_0&=\frac{P_4}{2(\eta-1+\rho)P},&
 \alpha_0&=\frac{B^2(\rho+\eta)}{2(1-B\rho)P}.
\end{aligned}
\]
These may have unbounded openings; only the final capped identity
below is used as a bounded Gram.  For an indeterminate $p$, put
\[
 e=Y-\rho p,\qquad f=Z-vp/B,\qquad T_v=(V-y)p.
\]
The entire leading three-state form is exactly
\[
\begin{split}
 &4C_0D_v(p-Y)^2+4A_0D_vY^2
 +(4D_0+2)(p-Y-vZ)^2+(4\alpha_0+2)(vY-Z)^2-2D_vZ^2\\
 &\quad=\frac{4D_0}{B^2}(T_v+Be+Bvf)^2
       +4\alpha_0(f-ve)^2+4(C_0+A_0+\tfrac12)D_ve^2.
\end{split}                                                    \tag{LU8}
\]
This identity holds for every $v$, not only near the actual Bob
angle.  In the contact allocation at $z=u=0$,
$e=k(N+qR_h)$ and $T_v=-B(kR_h+v\zeta)$.
Subtract the residue contributions of the retained terms (LU4).
The harmonic identity then gives the complete three-normal block
\[
 \boxed{4\ell_*N^2+4gR_h^2
 +4(A_*+\tfrac12)D_vy(N+qR_h)^2
 +4\alpha_*\{\zeta-vk(N+qR_h)\}^2.}                     \tag{LU9}
\]
It has strictly positive residue on $(N,\zeta,R_h)$ because
$\ell_*,g,\alpha_*$ are positive units.  No normal has been
suppressed.  Every rational cancellation used to obtain (LU9)
has bounded source products on both sides, so the identity extends
continuously through every opening listed after (LU1).

\paragraph{The remaining two-root block and a finite cap.}
There is no finite residue coupling between these three roots and
$(\xi,R)$.  Such terms carry $u$, $z$, or the coefficient in (LU7).
For completeness, the exact first-normal support gives this block
without dropping a term.  To first order in $(U,z)$ write
$P_\xi=UA(V)+zZ(V)$; both coefficient functions are affine in $V$,
and
\[
 Z(V)=\frac{D_v}{B}(-4r+8\tau-8\beta_*).
\]
The exact contact equation gives $A(y)+Z(y)=0$ in residue.
After the first allocation of (LII15), the coefficient of $\xi R$
is therefore $A(V)-(V-y)(A'+Z')=-Z(V)$.
Using (LU3) and (LU6), the complete block is
\[
 M_{\xi R}=
 \begin{pmatrix}
 2D_v(1+2\beta_*)&\dfrac{2D_v}{B}(P_4/P+2\beta_*)\\[1mm]
 \dfrac{2D_v}{B}(P_4/P+2\beta_*)&
 \dfrac2{B^2}\{(P_4/P+2\beta_*)D_v-(S/P)r(V-y)\}
 \end{pmatrix}.                                           \tag{LU10}
\]
Here all letters denote real residues when positivity is assessed.

If $\beta$ is bounded, take $\beta_*=\beta$.  The native relation
$(1-r)(2\beta+1)=S/P$ follows from (LU3), and (LU10) has determinant
\[
 \frac{4S D_v}{B^2P}(1+2\beta)r(1+y)>0.                 \tag{LU11}
\]
Its first diagonal entry is positive, including $\beta=0$.
If $\beta$ is unbounded, take $\beta_*$ to be a sufficiently large
positive real cap.  Increasing it adds the positive square
$4\beta_*D_v(\xi+R/B)^2$.  On its kernel $\xi=-R/B$,
the uncapped block equals
\[
          \frac{2S}{B^2P}\{1+ry+(1-r)V\}R^2>0.           \tag{LU12}
\]
The elementary real two-by-two Schur argument therefore supplies
a finite cap making (LU10) strictly positive.  Unboundedness of the
original $\beta$ permits that cap.

Together (LU9)--(LU12) prove strict positive residue of the full
five-root Gram.  The bounded Gram lemma gives its tensor SOS;
restoring (LU4) proves (LU2).  The cyclic four-event complement
completes every testing chart and original projective endpoint.

\paragraph{Verification scope.}
  The exact source contact and
monic contact divisions are the unchanged full-parent argument of
(LII7)--(LII15), with the bounded modifications (LU4)--(LU6).
This theorem closes the $L$-unit, $\eta$-unit subdomain of C1.4;
the $\eta$-unbounded subdomain is not included.



% Source formulas: p10ee3d96a3f1
\subsection{Large $\eta$ with bounded reciprocal $x\zeta$: one proof for every $h/\sigma$}
\label{sec:c1-unification-bounded-z-all-H}

Let $K_c,K_t$ be separately labelled real closed extensions of $\mathbb R$,
$T=K_c\otimes_{\mathbb R}K_t$, and $\Sigma T^2$ the finite sums of
squares in $T$. Bounded coefficients belong to $B_c\otimes B_t$ and
their real residue is taken in both factors. We use the bounded-Gram
fact: a bounded symmetric matrix with positive-definite real residue
is a Gram matrix of a finite sum of squares in $T$. Roots in the Gram
need not be bounded. A positive source or testing factor can be absorbed
into a square using the real closedness of its own field.

\begin{theorem}\label{thm:c1-unification-bounded-z-all-H}
Use the stationary C1 source with $d=1$, $\lambda_{\rm ev}=n=0$,
$m,C,A,\alpha>0$, $D,\beta\ge0$, and the native strict gaps. Assume
explicitly that
\[
 x=k^2\in\operatorname{Inf}_{c,>0},\qquad
 \sigma=1/\eta\in\operatorname{Inf}_{c,>0},\qquad
 z=1/(x\zeta)>0\text{ is bounded}.
\]
There is no size restriction on $H=h/\sigma$. Then the complete
reciprocal parent is in $\Sigma T^2$ for every independent pair
$u,w\in\operatorname{Inf}_t$ and every real or complex state.
There is no state-ratio or testing-ratio restriction. The literal
faces $D=0$ and $\beta=0$ are included.

Equivalently, the theorem applies to the rational source below under
$0<b<1$, $0<r=bz<1$, $0<h<1$, $E,M,J>0$, $D,\beta\ge0$,
and the stated assumptions on $x,\sigma,z$.
\end{theorem}

\paragraph{The source and its uniform estimates.}
The rational identities and complete parent expansions below are
algebraic for every positive $z$ with the native gaps. Only their
boundedness estimates and the theorem require $z$ bounded.
Put
\[
 r=bz,\quad Q=1-b,\quad 0<b,r<1,\quad \theta=1-h,\quad
 e=xh/\sigma=xH,\quad E=1-e,\quad K=1+x,
\]
\[
 R_s=1+xr,\quad M=R_s-xz\theta>0,\quad
 J=1+xh-\sigma(1-r)>0,
\]
\[
 P=\sigma+K+\sigma z(1+bx+2xh)
              +\sigma bx^2hz^2+x^2h\theta z.
\]
Define the opening-cleared products
\[
\begin{split}
 p&=\frac{z(1+xh)E+(1-r)(1+xhz)}{2P},\\
 l&=\frac{(1+\sigma r)\{eR_s(1+\sigma z)+ME\}}{2P},\qquad
 \kappa=\frac{K(1+\sigma r)(1+\sigma z)}{2MP}.
\end{split}
\]
The event coefficients are
\[
\begin{gathered}
 m=\kappa x^2zH,\quad
 C=\frac{Kz(1+\sigma r)E}{2(1-r)P},\quad
 A=p/r,\quad \alpha=l/(zQ),\\
 D=\frac{K(1+\sigma z)(1+\sigma r)-JP}{2JP},\qquad
 \beta=\frac{K(1+\sigma r)(1+xhz)-\theta P}{2\theta P}.
\end{gathered}\tag{P1}\label{eq:c1-eta-large-source}
\]
Only the original nonnegative $D,\beta$ sources are considered.  Set
$a=Az=p/b$, $\ell=\alpha z=l/Q$, and
\[
 H_0=1+zQ,\quad B_s=P+xz(J-\sigma),
\]
\[
 \mathcal H=R_sH_0+z\{(1-r)xh(2+xr)
       -e[2+x(1+h+r\theta)]+Ke^2\}.
\]

Direct multiplication gives the two exact source identities
\begin{align}
 2\theta P\beta&=hB_s-\sigma H_0,\tag{BE3}\label{eq:c1-unification-source}\\
 2P(JD+x\theta\beta)&=\sigma^2\mathcal H.\tag{BE4}\label{eq:c1-unification-source-sphere}
\end{align}
They contain the entire scale estimate required below. Initially
$0<e<1$ follows from $E>0$, so $\mathcal H$ is bounded; also
$P,J,M,B_s\sim1$. Eliminating $\beta$ between the identities gives
\begin{equation}\label{eq:c1-unification-e-bound}
 eB_s=xH_0+\sigma\mathcal H-2PJD/\sigma
       \le xH_0+\sigma\mathcal H.
\end{equation}
Thus $e=O(x+\sigma)$ is infinitesimal. Consequently
\[
 \mathcal H-H_0\in\operatorname{Inf}_c,\quad
 D=O(\sigma^2),\quad \beta\theta=O(h+\sigma),\quad
 p-H_0/2,\ l-1/2,\ \kappa-1/2\in\operatorname{Inf}_c,
 \qquad H_0=1+zQ\ge1.
\]
The quantities $p,l,\kappa$ are bounded and positive, and
$m=\kappa xze$ is infinitesimal. The individual positive coefficients
$C,a=p/b,\ell=l/Q$ may be unbounded; they will stay in whole squares.

The following estimates do not assume that $H$ is bounded. Set
\[
 \nu=(2D-2\beta\theta+h)/\sigma,\quad
 \lambda=(8Dh+4\beta\theta^2+2h^2)/\sigma-4bp\sigma z,\quad
 K_A=4\kappa H-\lambda,
\]
and define
\[
 \lambda_0=\lambda-2H,\quad d_G=4D/\sigma^2,\quad
 \rho=\lambda x/\sigma,\quad L_0=\lambda\sigma.
\]
The source identities imply
\begin{align}
 \nu&=2D/\sigma+[H_0-ez(J-\sigma)]/P,\label{eq:c1-unification-nu}\\
 \lambda&=2H+2e\theta z(J-\sigma)/P
            -2\theta H_0/P+8Dh/\sigma-4bp\sigma z.
            \label{eq:c1-unification-lambda}
\end{align}
Hence $\nu,d_G,\lambda_0$ are bounded,
$\lambda_0+2\theta H_0$ is infinitesimal, and $L_0-2h$ is infinitesimal.
For $K_A=4\kappa H-\lambda$ one also has
\begin{equation}\label{eq:c1-unification-KA}
 K_A-\{2\theta H_0-2h(1-r)\}\in\operatorname{Inf}_c.
\end{equation}
Here is a finite polynomial derivation of the only possible pole in
this assertion. Put
\[
 N_\kappa=K(1+\sigma r)(1+\sigma z)-MP
          =xF_x+\sigma F_\sigma,
\]
where $F_x=(N_\kappa-N_\kappa|_{x=0})/x$ and
$F_\sigma=N_\kappa|_{x=0}/\sigma$ are polynomials in the bounded
source arguments. Direct substitution gives
$F_\sigma=r-1+\sigma rz$. Therefore
$H(2\kappa-1)=(eF_x+hF_\sigma)/(MP)$ has residue $-h(1-r)$,
which proves \eqref{eq:c1-unification-KA}. In particular $K_A$ is bounded.

For later use write $R_\gamma=R_0-8pu^2$, where
\[
 R_0=-4\{2(C+D)(r-1)+2p+r-1\}.
\]
There is the uniform estimate
\begin{equation}\label{eq:c1-unification-R0}
 R_0=O(x+\sigma+e),\qquad
 \lambda>0\ \Longrightarrow\ R_0\sqrt\lambda\in\operatorname{Inf}_c.
\end{equation}
To verify the first estimate without any source-pole division, temporarily
treat $e$ as an independent variable. The product $N=JPR_0$ is the
polynomial
\[
 -4\{-JKz(1+\sigma r)(1-e)
 +(r-1)K(1+\sigma z)(1+\sigma r)
 +J[z(1+xh)(1-e)+(1-r)(1+xhz)]\}.
\]
It vanishes identically at $x=\sigma=e=0$. Consecutive differences
in $x,\sigma,e$ express it as $xN_x+\sigma N_\sigma+eN_e$, with
bounded polynomial quotients. For the second estimate use
$\lambda=O(1+H)$ from \eqref{eq:c1-unification-lambda} and
$e=O(x+\sigma)$ from \eqref{eq:c1-unification-e-bound}. The bound
\[
 R_0^2\lambda=O((x+\sigma)^2(1+H)),\qquad
 (x+\sigma)^2H=xe+2\sigma e+\sigma h
\]
has an infinitesimal right side. This proves the square-root estimate
for every possible $H$.

The same source formulas also give the identities used by the
complementary testing charts. With $B=B_s$, put
\[
 d_0=\sqrt{2PJD/(\sigma^2\mathcal H)},\qquad
 b_0=\sqrt{2Px\theta\beta/(\sigma^2\mathcal H)}.
\]
The source identity gives
\[
 D=O(\sigma^2),\quad x\theta\beta=O(\sigma^2),\quad
 JD+x\theta\beta\asymp\sigma^2.\tag{BE5}
\]
Then exactly
\[
 d_0,b_0\ge0,\quad d_0^2+b_0^2=1,\quad
 D={\sigma^2\mathcal H\over2PJ}d_0^2,\quad
 \beta={\sigma^2\mathcal H\over2Px\theta}b_0^2.\tag{BE6}
\]
\[
 H={H_0\over B}+{\sigma\over x}{\mathcal H\over B}b_0^2.\tag{BE7}
\]
\[
 m=m_0+m_1b_0^2,\quad
 m_0=\kappa x^2zH_0/B\asymp x^2z,\quad
 m_1=\kappa x\sigma z\mathcal H/B\asymp x\sigma z.\tag{BE8}
\]
\[
 {C(1-r)\over z}\sim\tfrac12,\qquad
 Ar\sim H_0/2,\qquad \alpha Qz\sim\tfrac12.\tag{BE9}
\]
No inverse of a sphere coordinate is used. Another useful uniform
source identity is $p-l\in zB_c$: with $e$ temporarily independent,
$2P(p-l)$ is a polynomial in bounded $x,\sigma,z,b,h,e$, and
at $z=0$ its two source numerators are both one
($e+(1-e)=1$ for the second). Its finite quotient by $z$ is
therefore polynomial and bounded. Restoring $e=xh/\sigma$ proves
the assertion through both radial openings.

\paragraph{The complete state polynomial.}
Write $t=u^2$, $\xi=w^2$, $D_u=1+t$, $D_w=1+\xi$, $q=\sqrt z$,
$g=h-Q$.  The original reciprocal parent is
\[
\begin{split}
 F_w={}&4m(W-wX+uY-uwZ)^2+4CD_w(uW-Y)^2+4AD_uD_wY^2\\
 &+(4D+2)(uwW+uX-wY-Z)^2+(4\alpha+2)D_u(Y-wZ)^2\\
 &+(4\beta+2)D_w(X+uZ)^2-2D_uD_wZ^2.
\end{split}\tag{P4}\label{eq:c1-reciprocal-parent}
\]
Use the homogeneous roots
\[
\begin{gathered}
 N_x=X+uhZ,\quad N_y=(Y-bwZ)/q,\quad
 \Gamma=\sigma Z-uwW,\quad R=wZ/q-uqW,\\
 A_\theta=q(t-x)W,\quad P_\beta=uZ-(x/\sigma)wW,
 \quad X_e=X+ewW,\quad V=uw(t-x)W,\\
 J_C=N_y+bR,\qquad J_\ell=N_y-QR.
\end{gathered}
\]
The following definitions collect every remaining coefficient:
\[
\begin{gathered}
 \nu=(2D-2\beta\theta+h)/\sigma,\qquad
 \lambda=(8Dh+4\beta\theta^2+2h^2)/\sigma-4bp\sigma z,\\
 \lambda_0=\lambda-2H=4D(1+h)/\sigma-2\theta\nu-4bp\sigma z,\\
 K_A=4\kappa H-\lambda,\qquad B_x=8\kappa e,\\
 A_y=8(p-l)+4z\{(1+2m)g+2Dh+2p\sigma\},\\
 R_x=-4\{2Db+(1+2m)g-2\beta\theta\},\\
 R_\gamma=-4\{2(C+D)(r-1)+2p+r-1+2pt\},\\
 \chi=8(pb+lQ)-4h+4eh(2D+1)-8Dh\\
 \hspace{20mm}+4z\{4Dbh+4pb\sigma+(1+2m)g^2+2\beta\theta^2\},\\
 k_1=2(1+2m)g^2+8Dhb+8pb\sigma+4\beta\theta^2,
\end{gathered}\tag{SP3}\label{eq:c1-contact-coefficients}
\]
\[
\begin{split}
 E_{yy}&=2z\{(1+2m)t+(1+2D)\xi\},\\
 E_{RR}&=z\{8bD+8bp\sigma+4b\sigma(2C+2D+1)(r-1)
                                  +k_1t+2b^2(2D+1)\xi\},\\
 E_{yR}&=4z\{t[2Dh+2p\sigma+(1+2m)g]+b(2D+1)\xi
       +\sigma(2C+1)(r-1)+2D[1+\sigma(r-1)]+2p\sigma\}.
\end{split}
\]
Direct expansion of the original parent, using only the displayed source
identities, gives the full four-state identity
\begin{align}
 F_w={}&\mathcal E_++[2(1+\xi)+4m\xi]N_x^2
       +[4a+2z(2C+1)]J_C^2+4\ell J_\ell^2\nonumber\\
 &+E_{yy}N_y^2+E_{RR}R^2+E_{yR}N_yR
        +(4D/\sigma^2)\Gamma^2\nonumber\\
 &-4\nu N_x(u\Gamma)
       +R_\gamma(N_y+bR)(qw\Gamma)\nonumber\\
 &-4uqw(2m+2D+1)N_xN_y+uqwR_xN_xR
           +qwB_xN_xA_\theta+uA_yN_yA_\theta\nonumber\\
 &+K_AA_\theta^2+\chi uA_\theta R+\mathscr S,
 \label{eq:compression-full-parent}\tag{SP6}\\
 \mathscr S={}&\lambda\sigma P_\beta^2
       +(\lambda x/\sigma)(t-x)(\xi-\sigma z)W^2,\nonumber
\end{align}
where all terms in the following remainder are retained tensor squares:
\[
\begin{split}
 \mathcal E_+={}&4\beta D_wN_x^2+2(2D+1)(uX_e)^2
              +4bpzt\Gamma^2\\
 &+4a(D_uD_w-1)J_C^2+4\ell tJ_\ell^2+4Cz\xi J_C^2.
\end{split}
\]
This identity includes all four state coordinates and both complete Gamma normal terms.  


All $E$ coefficients are infinitesimal. The state-normal cross
coefficients outside $\mathscr S$, namely
$\nu,R_\gamma,R_x,B_x,A_y,\chi$, are bounded; so are $K_A,d_G$.
The possibly unbounded scalar coefficient $\lambda$ is handled below.
Indeed the only $C$ poles outside positive squares occur through
$C(r-1)$; $\beta$ occurs there through the bounded product
$\beta\theta$. Every non-Gamma cross coefficient has an explicit
$u$ or $w$ factor and is infinitesimal. These conclusions use the
uniform estimates just proved, not a bound on $H$.

Since $4a,4\ell>1$, clip these two whole positive radial squares at
one and retain their source excesses. The resulting radial form is
\[
 J_C^2+J_\ell^2=(N_y+bR)^2+(N_y-QR)^2,
\]
whose real residue is positive definite because $b+Q=1$.
The $N_x$ coefficient has residue two. The exact testing sign gauges
$(u,Y,Z)\mapsto(-u,-Y,-Z)$ and $(w,X,Z)\mapsto(-w,-X,-Z)$
allow $u,w\ge0$; literal zero values are allowed.

\paragraph{A single construction for $\lambda>0$.}
In this paragraph $\rho>0$. Direct elimination of $\beta$ in
\eqref{eq:c1-unification-source} gives the exact identity
\begin{equation}\label{eq:c1-unification-S}
\begin{split}
 S:=d_G+\rho={}&2\mathcal H/P
 +d_G(1-J+2hx)+2eH_0/P\\
 &-2e^2z(J-\sigma)/P-4bpxz.
\end{split}
\end{equation}
Thus $S-2H_0$ is infinitesimal. In particular $S$ is a positive
bounded source unit and $d_G,\rho$ are both bounded. Also
$q_0:=\lambda k$ is bounded, since $q_0^2=L_0\rho$, and
$\lambda x=\rho\sigma$ is infinitesimal.

Define the separated monomials, keeping the signs of $W,Z$,
\[
 B_1=\sqrt{\lambda\sigma}\,uZ,\quad
 B_2=\sqrt\rho\,uwW,\quad
 B_3=\sqrt\lambda\,x\sqrt zW,\quad
 B_4=\sqrt{\lambda x}\,wZ/\sqrt z.
\]
Exactly
\[
 \mathscr C:=\mathscr S+\lambda xR^2
        =\sum_i B_i^2-4\lambda xuwWZ.
\]
If $WZ>0$, change the overall state sign so $W,Z>0$, put
$G=(B_1B_2B_3B_4)^{1/4}$ and $d_i=B_i-G$. Then
\[
 \mathscr C=\sum_i d_i^2+\mathscr C_+,
\]
\[
 \mathscr C_+=2G\{(\sqrt{B_1}-\sqrt{B_2})^2
 +(\sqrt{B_3}-\sqrt{B_4})^2
 +2[(B_1B_2)^{1/4}-(B_3B_4)^{1/4}]^2\}\in\Sigma T^2.
\]
All radicals separate into factors in $K_c$ and $K_t$. If $WZ\le0$,
put $d_i=B_i$ and $\mathscr C_+=-4\lambda xuwWZ\in\Sigma T^2$.
In both cases the differences $d_i-d_j$ equal $B_i-B_j$.

Put $T_\theta=k\sqrt z(u-k)W$ and $\Gamma_*=\sqrt{d_G}\Gamma$.
The following exact conductors require no inverse of $D$ or sphere
coordinate:
\begin{align}
 \sqrt\rho\,u\Gamma={}&k(d_1-d_2)-(u-k)(d_2-d_3)
                      -\sqrt{\rho\sigma}\,T_\theta,\label{eq:c1-unification-conductor1}\\
 \sqrt\rho\sqrt z w\Gamma={}&\sqrt\rho\sigma zR
 -(\sqrt z w+z\sqrt\sigma)(d_2-d_3)\nonumber\\
 &+\sqrt\rho(\sigma z/k+\sqrt\sigma\sqrt z w/k)T_\theta.
 \label{eq:c1-unification-conductor2}
\end{align}
Use $1=(d_G+\rho)/S$ to write
\[
 u\Gamma={\sqrt{d_G}u\over S}\Gamma_*
              +{\sqrt\rho\over S}(\sqrt\rho\,u\Gamma).
\]
After \eqref{eq:c1-unification-conductor1}, every coefficient on
$(\Gamma_*,T_\theta,d_i)$ is infinitesimal. For the other complete
normal use
\[
 \sqrt z wR_\gamma\Gamma=R_0\sqrt z w\Gamma
                              -8p\sqrt z wu(u\Gamma),
\]
and apply the same identity $1=(d_G+\rho)/S$ only to the first term.
In \eqref{eq:c1-unification-conductor2} its only potentially singular
$T_\theta$ coefficients reduce exactly to
\[
 {R_0\rho\sigma z\over kS}={R_0q_0z\over S},\qquad
 {R_0\rho\sqrt\sigma\sqrt z w\over kS}
       ={(R_0\sqrt\lambda)\sqrt\rho\sqrt z w\over S}.
\]
They are infinitesimal by \eqref{eq:c1-unification-R0}; every other
coefficient visibly is too. Thus both entire Gamma normal terms
have infinitesimal coefficients on a single frame, including $D=0$.

It remains to handle a possibly signed $K_A$. The same construction
will work for every value of this source coefficient. Put
\[
 c=\max(0,1-K_A),\qquad a_0=k/\sigma,\qquad
 J_A=\sqrt z(u-k)W,\quad F_A=u(Z-a_0wW),
\]
\[
 E_u=\sqrt uJ_A,\quad E_k=\sqrt kJ_A,\qquad
 U_* =\sqrt{q_0}E_u,\quad K_* =\sqrt{q_0}E_k.
\]
One direct multiplication gives, for every state and all source sizes,
\begin{equation}\label{eq:c1-unification-US}
 \mathscr S=L_0F_A^2+2q_0E_u^2+q_0E_k^2
                                  +2q_0\sqrt uE_uR.
\end{equation}
Average this with $\mathscr S=\mathscr C-\lambda xR^2$:
\[
 \mathscr S=\tfrac12L_0F_A^2+U_*^2+\tfrac12K_*^2
 +\sqrt{q_0u}\,U_*R+\tfrac12\mathscr C-\tfrac12\lambda xR^2.
\]
The two cross/radial coefficients are infinitesimal. The term
$L_0F_A^2/2$ is retained as a positive square.

If $c>0$, \eqref{eq:c1-unification-KA} implies that $h$ is a positive
unit: in fact its residue is at least $1/4$. The inequality obtained
from \eqref{eq:c1-unification-e-bound},
\[
 a_0^2(hB_s-\sigma H_0)\le\mathcal H,
\]
then proves $a_0$ bounded. Also $L_0\sim2h$ is a positive unit,
so $\lambda$ is unbounded. The exact angular certificate proved
below is
\begin{equation}\label{eq:c1-unification-angular-budget}
 \Delta_A:=\varepsilon_C\mathscr C+u^2R^2
                         +\varepsilon_G\Gamma^2-A_\theta^2\in\Sigma T^2,
\end{equation}
where
\[
 \varepsilon_C=2/\lambda+u+k,\qquad
 \varepsilon_G=2a_0u+(u+k)/L_0.
\]
Both coefficients are infinitesimal whenever $c>0$. Subtract
$c(d_G/S)\Delta_A$ from the parent, keeping that nonnegative
quantity separately. The angular part of the remainder is
\[
 (K_A+c)A_\theta^2-c(\rho/S)A_\theta^2.
\]
The first coefficient is a positive unit. The second square is an
infinitesimal Gram correction on $(U_*,K_*)$, by the exact relation
\begin{equation}\label{eq:c1-unification-weighted-angular}
 \rho A_\theta^2=
       (\sqrt{a_0u}\,U_*+\sqrt{a_0k}\,K_*)^2.
\end{equation}
The other subtractions have coefficients
$c d_G\varepsilon_C/S$ on $\mathscr C$,
$c d_Gu^2/S$ on $R^2$, and $c\varepsilon_G/S$ on $\Gamma_*^2$;
all are infinitesimal. If $c=0$, none of these terms is used.
This verifies the same construction in both situations without a
new choice of roots or a source/test rate comparison.

Finally split the positive angular coefficient between $A_\theta$
and $T_\theta$, retaining the exact nonnegative difference
\[
 A_\theta^2-T_\theta^2=z(u-k)^2W^2(u^2+2uk)\in\Sigma T^2.
\]
The complete bounded Gram on
\[
 (N_x,N_y,R,\Gamma_*,A_\theta,T_\theta,d_1,\ldots,d_4,U_*,K_*)
\]
has positive-definite residue: its radial block is $J_C^2+J_\ell^2$,
the $N_x$ pivot is two, the Gamma pivot is one, the angular pivots
are positive, and the circuit and last two pivots are respectively
$1/2,\ldots,1/2,1,1/2$. All other entries have zero residue.
The remaining circuit coefficient
$1/2-cd_G\varepsilon_C/S$ has positive real residue $1/2$.
The one-dimensional bounded-Gram fact makes it a tensor SOS,
so its product with the retained $\mathscr C_+$ is also a tensor SOS.
Restore the retained positive squares and apply the bounded-Gram fact.

\paragraph{Proof of the angular certificate, including state signs.}
For $U,V,Z_*>0$ and positive separated $a_*$ set
\[
 \Psi=1+U^2V^2+U^2Z_*^2+V^2Z_*^2-4UVZ_*,\quad T_*=U^2-1,
\]
\[
 Q_*=VZ_*-U,\quad R_*=Z_*-UV,\quad G_*=U(V-Z_*).
\]
The polynomial identity
\begin{align*}
 &[2+a_*(U+1)]\Psi+U^2Q_*^2
       +[2U+(U+1)/a_*]R_*^2-T_*^2\\
 &=(T_*+UQ_*)^2+2Q_*^2
 +(U+1)(\sqrt{a_*}G_*-R_*/\sqrt{a_*})^2\\
 &\quad+a_*(U+1)(1-VZ_*)^2
       +2a_*(U+1)VZ_*(U-1)^2
\end{align*}
is verified by multiplication. Expand each $U+1$ factor before
taking square roots. The right side is a finite separated SOS.
For $W\ne0$ substitute
\[
 U=u/k,\quad V=w/(\sqrt\sigma\sqrt z),\quad
 Z_*={\sqrt\sigma\over k\sqrt z}{Z\over W},\quad
 a_*=\lambda k,\quad M_*=k^2\sqrt zW.
\]
Then $\mathscr C=\lambda M_*^2\Psi$,
$A_\theta=M_*T_*$, $R=k\sqrt zWQ_*$ and
$\Gamma=k\sqrt\sigma\sqrt zWR_*$. Multiplication by $M_*^2$
gives \eqref{eq:c1-unification-angular-budget}. These are legitimate
separated factors even when an intermediate ratio is unbounded.
If $Z_*<0$, apply the displayed identity first to $|Z_*|$.
Replacing $|Z_*|$ by $Z_*$ increases $\Psi$ by $8UV|Z_*|$ and
each of $Q_*^2,R_*^2$ by $4UV|Z_*|$. The three resulting additional
terms have positive separated coefficients and are SOS. Thus the
same certificate holds for either state sign. At $Z_*=0$ the original
identity applies directly. If $W=0$, $A_\theta=0$, so
\eqref{eq:c1-unification-angular-budget} follows immediately from
$\mathscr C\in\Sigma T^2$. No state coordinate is inverted there.

\paragraph{The case $\lambda\le0$.}
The definition of $\lambda$ gives
$H^2\le2bpz$ and $\beta\theta^2/\sigma^2\le bpz$.
Thus $H$ is bounded, $h$ is infinitesimal, $\theta\sim1$,
$M_\beta=\beta\theta/\sigma^2$ is bounded and
$H-H_0,d_G-2H_0$ are infinitesimal. In particular
\[
 L=\lambda/\sigma
   =8(D/\sigma^2)h+4M_\beta\theta+2H^2-4bpz
\]
is bounded. Direct multiplication gives
\[
 u^2\Gamma+V=\sigma uP_\beta,
\]
\[
 A_\theta^2+u\Gamma P_\beta+uA_\theta R
      =\sigma P_\beta^2+(x/\sigma)(u^2-x)(w^2-\sigma z)W^2,
\]
where $V=uw(u^2-x)W$. Consequently
\[
 K_AA_\theta^2+\mathscr S
 =4\kappa H A_\theta^2+Lu^2\Gamma^2+L\Gamma V
                                      +\lambda uA_\theta R.
\]
Add $1+L^2$ times the exact zero
\[
 V^2-\sigma u^2A_\theta^2-u\sigma(u^2-x)A_\theta R
                                      +(u^2-x)V\Gamma=0.
\]
The $(\Gamma,V)$ real residue block is
\[
 \begin{pmatrix}2\bar H_0&\bar L/2\\
                 \bar L/2&1+\bar L^2\end{pmatrix},\qquad
 \det=2\bar H_0+(2\bar H_0-1/4)\bar L^2>0.
\]
The $N_x$, radial and angular blocks are strict; every remaining
cross coefficient, including both complete Gamma normal terms,
is infinitesimal. The bounded-Gram fact proves this case, including
$\lambda=0$.

All constructions are literal on the testing axes. No inverse of
$D$ or $\beta$ occurs, so both permitted source faces are included.
The parent is real quadratic in the state; applying the proof to
the real and imaginary parts proves the complex-state assertion.
This completes the theorem.
\subsection{The six-root identity and unit-Alice angular relation}

\label{component:pb1104217f67b}

\label{pb1104217f67b:sec:r13-joint-radial-beta-zero}

\paragraph{Quantified domain and conclusion.}
Retain the original physical C1.4 source and all its residual-complement
conditions. Suppose $x=k^2$, $z=L^{-1}$ and $\sigma=\eta^{-1}$ are
positive infinitesimals, $0<b<1$, and the original coefficient $\beta=0$.
There are no further source hypotheses. In particular $h/\sigma\sim1$
and $D/\sigma^2\sim1/2$ follow from BE3--BE7, rather than being imposed.
The conclusion below covers every bounded independent real test $u$,
every infinitesimal or zero independent reciprocal test $w$, all their
signs, and real or complex states. The previously proved remote
complement supplies the other tests on this large-$L$ source chart.
Thus this proves the literal $\beta=0$ part of C14.1. The adjoining
Section~\ref{p909aa5765721:sec:r13-beta-positive-extension} extends the actual certificate
to all sources with $\beta\theta/\sigma^2$ bounded; the present result
is its beta-zero corollary and supplies the residue calculation used
there. C14.2 is not included.

The calculation joins the entire radial scalar quotient with every
radial first normal \emph{before} deciding boundedness. Their exact
cancellations remove the BR3 pole, including its separate
$8(ab)b/\sigma$ summand through $Q=1-b$ infinitesimal.

\paragraph{Notation and the original contact identities.}
Use the complete reciprocal parent \eqref{eq:c1-reciprocal-parent} and the large-$L$ source
notation. Thus $m=\mu x^2z$, $C=cz$, $A=a/z$, and
$\alpha=\ell/z$. In expressions below $p=ab$ and $l=\ell Q$ are
the two opening-cleared products; $p$ is not the source parameter $b$.
Set $t=u^2$, $\xi=w^2$, $\xi_0=\sigma z$, $\delta=\xi-\xi_0$,
$g=h-Q$, $D_u=1+t$, $D_w=1+\xi$, and $q=\sqrt z$.
The source identities established in the handoff give
\[
 \phi=\phi_t=\phi_\xi=L_x=L_y=L_z=0
 \quad\hbox{at }(t,\xi)=(x,\xi_0),
\]
as well as $\phi(0,\xi)=4m$ and $L_x(0,\xi)=-8m$.
Here $\phi,L_x,L_y,L_z$ are the \emph{complete} EN2 polynomials,
with $D_*=D$, $\beta_*=0$ and $\beta=0$.

Choose the fixed caps $a_* = \ell_* =1/4$. They are smaller than
the actual $a,\ell$, since $p,l\sim1/2$ and $0<b,Q<1$.
Retain the two exact nonnegative centered excesses
\[
 E_{\rm cap}=4(a-a_*)D_uD_wN_y^2+4(\ell-\ell_*)D_uN_y^2.
\]
Define six separated roots
\[
\begin{aligned}
N_x&=X+uhZ,&N_y&=(Y-bwZ)/q,\\
\Gamma&=\sigma Z-uwW,&R&=wZ/q-uqW,\\
A_\theta&=q(t-x)W,&V&=uw(t-x)W.
\end{aligned}\tag{JR1}
\]
No test/source quotient is used in their definition. Their inverses
and radicals are pure source.

\paragraph{Exact quotient specification.}
The contact identities give the exact divisions
\begin{align*}
 \phi&=z(t-x)^2A+t(t-x)\delta B+t\delta^2C_r,\\
 L_x&=(t-x)(tA_x+zB_x)+t\delta C_x,\\
 L_y&=(t-x)A_y+\delta B_y,\\
 L_z&=(t-x)A_z+\delta B_z.
\end{align*}
All quotients are monic polynomial divisions in one testing variable.
Their complete support can equivalently be specified by
\begin{gather*}
 A=4\mu+\frac{2h^2(2D+1)}\sigma t,\qquad
 A_x=-\frac{4h(2D+1)}\sigma,\qquad B_x=8\mu x,\\
 A_y=8(p-l)+4z\{(1+2m)g+2Dh+2p\sigma\},\\
 S:=A_z/\sigma-B=S_0+
       \frac{2h^2(2D+1)}{\sigma^2}(t-x).
\tag{JR2}
\end{gather*}
The same monic divisions give the much shorter exact expression
\[
 S_0=\frac{2\{2h^2(2D+1)x+4Dh-2bp\sigma^2z+h^2\}}{\sigma^2}.
\tag{JR2a}
\]
Thus $S_0\sim2$ follows immediately from $h/\sigma\sim1$ and
$D/\sigma^2\sim1/2$. This removes any dependency on the longer R10
reduction modulo the original beta numerator. The source contact jets
are still needed for the quotient remainders to vanish.
For a denominator-free quotient specification, $C_r$ is obtained by
dividing $\phi$ twice by $\xi-\xi_0$ and then dividing the quotient
by $t$; $B$ is the quotient of
$t^{-1}\phi_\xi(t,\xi_0)$ by $t-x$; and $A_z$ is the quotient of
$L_z(t,\xi_0)$ by $t-x$. These divisions introduce no mixed inverse.

The remaining exact coefficients are most compactly written as
\begin{align*}
 K_r:=\sigma^2C_r={}&k_0+k_1t+k_2\xi,\\
 k_0={}&2b\{(2C+1)b+4D+4p\sigma
             +2(2C+2D+1)\sigma(bz-1)\},\\
 k_1={}&2(1+2m)g^2+8Dhb+8pb\sigma,\\
 k_2={}&2b^2(2C+2D+1),\\
 R_x:=\sigma C_x={}&-4\{2Db+(1+2m)g\},\\
 R_y:=\sigma B_y={}&4\{t[2Dh+2p\sigma+(1+2m)g]
      +b(2C+2D+1)(\xi+\sigma z)\\
 &\hspace{23mm} +(2C+1)(b-\sigma)+2D(1-\sigma)+2p\sigma\},\\
 E_r={}&2(C+D)(bz-1)+2p+bz-1,\qquad
 R_\gamma=-4(E_r+2pt),\\
 Z_0={}&\frac{4D+2t(4Dh+h^2)+2t^2h^2(2D+1)}{\sigma^2},\\
 Z_x={}&-\frac{4u}{\sigma}\{2D+h+th(2D+1)\}.
\tag{JR3}
\end{align*}
In particular all of these coefficients are bounded when $u$ is bounded
and $w$ infinitesimal. The $\sigma^{-1}$ and $\sigma^{-2}$ in $Z_0,Z_x$
are controlled by the \emph{proved pure-source} bounds on $h,D$.

\paragraph{The joint cancellation, including all first normals.}
The exact full normal block EN3--EN4 obeys
\begin{align*}
 H_z/\sigma^2+\xi C_r-(\xi/\sigma)B_z&=Z_0,\\
 B_z-2\sigma C_r&=bR_\gamma,\\
 \frac4\sigma\{2DT_{\rm test}+bD_w(1+2C)+tg(1+2m)\}-B_y&=R_\gamma,\\
 \frac{4u}{\sigma}\{-2DT_{\rm test}-hD_u-(1+2m)\xi g\}
                         -u\xi C_x&=Z_x.
\tag{JR4}
\end{align*}
Here $T_{\rm test}=1+th+b\xi$. These are polynomial identities after clearing
pure-source powers of $\sigma$. No coefficient is discarded: $K_r$
contains the full BR3 term, and the last three lines keep its radial
couplings. The cancellation explains why a scalar-only extraction of
BR1 followed by boundedness checking cannot see this certificate.

Put
\begin{align*}
 H_x&=2(1+t+\xi)+4Dt+4m\xi,\\
 H_y&=4D_u(a_*D_w+\ell_*)+4mtz+4CzD_w+(4D+2)\xi z+2D_uz,\\
 c_G&=4D_u\{pbD_w+lQ\},\qquad
 c_N=8D_u\{pD_w-l\}.
\end{align*}
The \emph{entire original parent}, without restriction to a state
subspace, is exactly $F_w=E_{\rm cap}+G$, where
\begin{align*}
 G={}&H_xN_x^2+H_yN_y^2+Z_0\Gamma^2
        +(c_G+zK_r)R^2+AA_\theta^2\\
 &-4uqw(2m+2D+1)N_xN_y+Z_xN_x\Gamma
       +qwR_\gamma N_y\Gamma+qwbR_\gamma R\Gamma\\
 &+uqwR_xN_xR+(c_N+zR_y)N_yR
       +qwB_xN_xA_\theta+uA_yN_yA_\theta\\
 &+uA_xN_xV+u\sigma B A_\theta R+S\Gamma V.
\tag{JR5}
\end{align*}
For example, the scalar radial conductor is
$t\delta^2W^2=\sigma^2(qR-w\Gamma/\sigma)^2$;
the radial $Z$ first normal is
$uw\delta W(\Gamma/\sigma)=qwR\Gamma-(\xi/\sigma)\Gamma^2$.
Their joint use gives the first two identities of JR4. The other
conductors are those of the verified large-$L$ proof with
$N_z=\Gamma/\sigma$. These identities prove JR5 exactly and include
every first normal and every hard quadratic coupling.

\paragraph{One bounded positive Gram for every bounded Alice test.}
All entries of JR5 are bounded separated tensors. Indeed
\[
 h/\sigma\sim1,\quad D/\sigma^2\sim1/2,\quad
 \mu\sim1/2,\quad C=O(z),\quad p,l\sim1/2,
 \quad p-l\in z\mathcal O_{\rm source},
\]
and $\sigma B=A_z-\sigma S\sim4$ uniformly for bounded $u$.
The displayed formula for $A_z$ obtained from $L_z$ has only the
bounded terms $h^2/\sigma$, $Dh/\sigma$, $p,l$, and bounded multiples
of $z$; hence this last assertion does not order any mixed rate.

For completeness, the possibly needed opening-uniform divisibility
$p-l\in z\mathcal O_{\rm source}$ follows directly from the exact
source formulas, not from a residue subtraction. Write $H=h/\sigma$.
The two numerators over their common unit denominator $2P$ are
\begin{align*}
 U_A={}&z(1+x\sigma H)(1-xH)+(1-bz)(1+x\sigma Hz),\\
 U_\ell={}&(1+\sigma bz)\big[xH(1+xbz)(1+\sigma z)
       +\{1+xbz-xz(1-\sigma H)\}(1-xH)\big].
\end{align*}
They are bounded polynomials in the source letters and both equal
one at $z=0$. Consequently $(U_A-U_\ell)/z$ is a bounded polynomial.
Neither numerator has a $b$ or $Q$ denominator. This proves the claimed
divisibility through both openings and is the source argument in EN6's
adjoining radial paragraph.

The exact syzygy
\[
 V^2=\sigma t A_\theta^2+u\sigma(t-x)A_\theta R-(t-x)V\Gamma
\tag{JR6}
\]
allows the addition of zero times the pure-testing function
\[
 \lambda(t)=\frac{2(t^2+t+1)}{(1+t)^2}>0.
\]
Specifically add
$\lambda[V^2-\sigma tA_\theta^2-u\sigma(t-x)A_\theta R+(t-x)V\Gamma]$
to JR5. There is no source/test denominator in $\lambda$.

At any real residue $u_0$, write $t_0=u_0^2$. The $N_y$ pivot is
$2(1+t_0)$. The $(A_\theta,R)$ block is
\[
 \begin{pmatrix}2&2u_0\\2u_0&2(1+t_0)\end{pmatrix},
 \qquad \det=4.
\]
The remaining $(N_x,\Gamma,V)$ block is
\[
 \begin{pmatrix}
 2(1+t_0)&-2u_0(1+t_0)&-2u_0\\
 -2u_0(1+t_0)&2(1+t_0+t_0^2)&1+t_0+\lambda(t_0)t_0/2\\
 -2u_0&1+t_0+\lambda(t_0)t_0/2&\lambda(t_0)
 \end{pmatrix}.
\tag{JR7}
\]
Its first two leading pivots are positive, with leading $2\times2$
determinant $4(1+t_0)$. The last Schur complement is exactly
\[
 \lambda(1+t_0^2)+\frac{2t_0^2}{1+t_0}
       -\frac12(1+t_0+\lambda t_0/2)^2
   =\frac{4t_0+3}{2(1+t_0)^4}>0.
\tag{JR8}
\]
Thus the residue Gram is strictly positive definite for every finite
real $u_0$, including zero and both signs. Crucially it does not depend
on the residue of $b$: the result remains strict through $b\to0$
and $Q\to0$. The finite bounded positive-Gram criterion already used
by the project then gives the exact legal tensor-SOS.

\paragraph{Boundaries and scope.}
JR1--JR6 are identities, including $u=0$ and $w=0$; no root is inverted
at a zero. Their sign coverage is literal. Apply the real identity to
the real and imaginary parts of the state for complex states.
The exact remote join is the common-parent theorem in
Section~\ref{component:p89fd3edb3edd},
Section~\texttt{sec:r6-c1-common-remote}. First, $A,\alpha$ are
unbounded on this source, so its adjacent-event theorem transported by
(CE9)--(CE10) covers every $v$ that is not projectively unbounded.
For projectively unbounded $v$ and projectively unbounded $u$, put
$p_t=1/u$ and apply the complete raw frame (CE20), whose strict
residue is (CE21). Its exact source premises (CE17), without the
unneeded $m$ lower bound in this $p_t$ chart, hold with $q=\sqrt z$:
$C\sim z/2$, $A\ge1/(4z)$, and $\alpha\ge1/(4z)$.
The $C$ lower bound may likewise be chosen as $z/4$. The unselected
$D,\beta$ events are retained. The identities include literal
$p_t=0$, $w=0$ and all signs. These two remote cases and JR1--JR8
exhaust all projective tests on the printed source face. No source
domain with $L$ a unit, or $\beta>0$, is covered by this proof.

\paragraph{Corroboration.}
 The
original six contact identities are explicitly identified imported
results, not silently replaced by numerical residue checks.



% Source formulas: p909aa5765721
\subsection{Bounded $H$ at unit Alice testing angle}

\label{component:p909aa5765721}

\label{p909aa5765721:sec:r13-beta-positive-extension}

Retain the same original C14.1 source as JR1--JR8, with $x,z,\sigma$
positive infinitesimals, but allow $\beta\ge0$. Write
\[
 b_\theta=\beta\theta,\qquad H=h/\sigma,\qquad
 M_\beta=b_\theta/\sigma^2.
\]
This section proves two precise extensions. First, if $M_\beta$ is
bounded, the \emph{entire source domain is covered for every test}.
This includes the beta-zero theorem. Second, if $H$ is bounded but
$M_\beta$ is unbounded, every bounded nonzero-unit $u$ is covered
at the residual infinitesimal-$w$ tests. On this second source domain
only $u$ infinitesimal remains. Neither statement assumes $d_0$ a
unit: the second includes $D=0$. The $H$-unbounded domain and C14.2
are not covered here.

\paragraph{The generic joint radial identity.}
The exact six-root identity JR5 continues to hold if its coefficients
are given the following beta corrections:
\begin{gather*}
 k_1\ \mapsto\ k_1+4b_\theta\theta,\qquad
 R_x\ \mapsto\ R_x+8b_\theta,\\
 Z_0\ \mapsto\ Z_0+\frac{4b_\theta\theta t}{\sigma^2},\qquad
 Z_x\ \mapsto\ Z_x+8ub_\theta/\sigma,\\
 S\ \mapsto\ S+4b_\theta\theta/\sigma^2.
\tag{JB1}
\end{gather*}
Here the uncorrected formulas on the right are the polynomials in
JR2--JR3 evaluated at the actual $h,D,m,p,l$, not at a beta-zero
source. The scalar $B$ remains its full monic quotient. The additional
retained centered excess is $4\beta D_wN_x^2$. The complete generic
coefficient is, explicitly,
\[
 S=\frac{2\{h^2(2D+1)(t+x)+4Dh-2bp\sigma^2z
                       +2b_\theta\theta+h^2\}}{\sigma^2}.
\tag{JB2}
\]
All four cancellations JR4 are exact with these changes. They were
checked symbolically with $b_\theta$ free. Thus no radial pole remains
unallocated in this extension.

\paragraph{The source estimates needed by the extension.}
The source equation BE3 gives exactly
\[
 H=H_0/B+\frac{2P}{B}\frac{b_\theta}{\sigma}.
\tag{JB3}
\]
Consequently $M_\beta$ bounded implies $H\sim1$ and $\theta\sim1$.
The sphere budget then implies $D/\sigma^2\sim1/2$, since
$xb_\theta/\sigma^2=xM_\beta$ is infinitesimal. Also $\mu\sim1/2$.

More generally, suppose only $H$ bounded. Then $\res H\ge1$;
no exact inequality $H\ge1$ is imposed. One has $\theta\sim1$,
and $b_\theta/\sigma$ is bounded. Write
\[
 d=\res(D/\sigma^2),\qquad
 c=\res(xb_\theta\theta/\sigma^2).
\]
BE4--BE6 give $d,c\ge0$ and $d+c=1/2$. Formula JB3 gives
$\res(b_\theta/\sigma)=(\res H-1)/2$. The bounded source formulas
also give $\mu\sim H/2$, $p,l\sim1/2$, $C=O(z)$, and
$p-l\in z\mathcal O_{\rm source}$. These statements retain both
source-sphere endpoints and both $b,Q$ openings.

\paragraph{An exact square joining the beta pole to its angular reserve.}
Suppose first $u$ is a bounded nonzero unit, so that $1/u$ is an
allowed bounded pure-testing inverse. Put
\[
 a_\beta=b_\theta\theta,\qquad
 C_\beta=4a_\beta/\sigma^2,\qquad T_c=t\Gamma+V,
 \qquad P_\beta=uZ-(x/\sigma)wW.
\]
The identity $T_c=\sigma uP_\beta$ has no $u$ denominator.
Adding $(C_\beta/t)$ times JR6's zero, and completing the resulting
$T_c$ square, gives the exact entire-parent identity
\[
 G=4a_\beta P_\beta^2+\widehat G,
\tag{JB4}
\]
where $G$ is the complete generic JR5 and
\begin{align*}
 \widehat G=G
 &-C_\beta t\Gamma^2
 -(C_\beta+C_\beta x/t)\Gamma V
 -C_\beta\sigma A_\theta^2\\
 &-C_\beta\sigma(t-x)A_\theta R/u.
\tag{JB5}
\end{align*}
Every coefficient of $\widehat G$ is bounded when $H$ is bounded
and $u$ a unit. The only apparent angular reserve subtraction is
$4a_\beta/\sigma$. Its exact source comparison is
\[
 \mu-a_\beta/\sigma
 =\kappa H_0/B+
    \frac{b_\theta}{\sigma}\left(\frac{2\kappa P}{B}-\theta\right),
\tag{JB6}
\]
where $\kappa$ is the BE8 source factor. The bracket is infinitesimal,
so this difference has residue $1/2$. Thus the beta pole is paid for
by the actual shared $m$ term. No leading coefficient is substituted
in the exact identity JB4.

\paragraph{Bounded $M_\beta$: all residual tests.}
When $M_\beta$ is bounded and $u$ a unit, the residue of $\widehat G$
is exactly the beta-zero residue in JR7 before the finite $V$ pivot:
$H\sim1$, $d=1/2$, $c=0$. Add the same pure-testing pivot
$\lambda(t)=2(t^2+t+1)/(1+t)^2$ and apply JR8. Restore the positive
square in JB4 and every centered excess. This covers every unit $u$.

\paragraph{Unbounded $M_\beta$, bounded $H$: all unit-Alice tests.}
For any real nonzero residue $u_0$, set $t_0=u_0^2$. In the notation
of JB3, and writing $H$ now for its finite real residue, the only
nontrivial $\widehat G$ residue block is
\begin{align*}
 Q={}&2(1+t_0)N_x^2+[4d+2H^2t_0(1+t_0)]\Gamma^2\\
 &-4u_0(1+Ht_0)N_x\Gamma-4Hu_0N_xV\\
 &+[2H^2(1+t_0)-4c/t_0]\Gamma V.
\tag{JB7}
\end{align*}
Its angular $(A_\theta,R)$ block is the strict matrix from JR7 and
its $N_y$ pivot is $2(1+t_0)$. All other residue cross entries vanish.
The beta square in JB4 equals $(C_\beta/t)T_c^2$, with
$C_\beta$ unbounded. On its exact residue kernel $V=-t_0\Gamma$,
JB7 reduces to
\[
 2(1+t_0)N_x^2-4u_0N_x\Gamma+2\Gamma^2,
 \qquad \det=4.
\tag{JB8}
\]
Here the cancellation uses exactly $d+c=1/2$. It remains strict at
$d=0$ and $c=0$. The row $T_c=V+t\Gamma$ has a bounded separated
right inverse, obtained by changing $V$ alone. The elementary finite
penalty lemma therefore chooses a positive real $K$ such that
$\widehat G+(K/t)T_c^2$ has strictly positive residue. This is possible
because $t_0>0$. The exact retained excess is
$((C_\beta-K)/t)T_c^2$: $C_\beta-K$ is a positive pure-source
coefficient, and its square root separates from the pure-testing
$1/\sqrt t$. No ordering of $C_\beta/t$ minus a real number is used.
This proves the entire
parent for all bounded nonzero-unit $u$, with infinitesimal or zero
$w$, all signs and complex states.

The new exact remaining obstruction on the bounded-$H$ source branch
is therefore $M_\beta$ unbounded and $u$ infinitesimal. JB5 has
$x/u$ and $x/u^2$ terms there, so its bounded-Gram argument cannot be
imported. The square $P_\beta$ itself is still legal and has no
$u$ inverse; a joint allocation with the remaining angular terms is
required to remove those coefficient denominators without comparing
independent source and testing rates. No such allocation is claimed.

\paragraph{Denominators and coefficient bounds.}
Every source inverse in the roots or exact coefficients is a pure-source
inverse. On the large-$L$, bounded-$H$ domains, $P,J,B$ and the original
geometry coefficient $M=1+xbz-xz\theta$ are units; $1-r$ and $\theta$
are units as well. Possible $b,Q$ inverses occur only in the original
source weights and the retained positive centered excesses, while their
opening-cleared products $p,l$ are bounded. The normalized quantities
$H$, $D/\sigma^2$, $b_\theta/\sigma$, and
$xb_\theta\theta/\sigma^2$ are bounded by JB3 and the BE sphere.
The corrections in JB5 therefore leave bounded coefficients: the
remaining $\Gamma^2$ terms involve $D/\sigma^2$, $Dh/\sigma^2$ and
$H^2$; the $\Gamma V$ coefficient involves those same bounded source
products and $(xb_\theta\theta/\sigma^2)/t$; the angular-radial
coefficient involves $b_\theta\theta/\sigma$ times $(t-x)/u$.
The factors $1/u,1/t$ are used only when $u$ is a testing unit.
All $u=0$ and infinitesimal-$u$ assertions for bounded $M_\beta$
instead use the unaltered generic six-root identity and its bounded
pure-source pivot $(1+2M_\beta)^2$. The everywhere bounded testing
denominator $1+t=1+u^2$ in JR8 is also pure testing. No mixed inverse,
mixed radical, or source/testing rate comparison occurs in either
covered domain or its remote join.



% Source formulas: pcc84202479d2
\subsection{Unbounded $H$ at unit Alice testing angle}

\label{component:pcc84202479d2}

\label{pcc84202479d2:sec:r18-sphere-unit-u}

Retain the entire original C14.1b source intersection and the exact
SP6 parent from the common identity.  Assume $u$ is a bounded nonzero unit
and $w$ is infinitesimal.  The argument below proves the whole parent
for every testing state, including the remaining
$W/Z$-unbounded states.  It does not cover infinitesimal $u$.
All source gaps, $D=0$, both source sphere endpoints and the literal
coefficient boundaries remain unchanged.

\paragraph{No state chart is needed.}
All roots below are homogeneous polynomials in the original testing
state, with pure-source coefficients.  In particular the proof
uses no state-coordinate inverse, no bound on $W/Z$, and no sign
restriction.  The only testing inverse is that of the unit $u$.
Write $U=u^2$, $s=x/\sigma$, and preserve the exact SP/BS coefficients
\[
 L_h=\lambda\sigma/h,\quad \rho=\lambda x/\sigma,
 \quad d=4D/\sigma^2,\quad K_A=4\kappa h/\sigma-\lambda.
\]
The source identities give $s\in\mathrm{Inf}$,
$\sigma/\sqrt h\in\mathrm{Inf}$, $L_h\sim2$, $\nu\sim1$, and
\[
 d+\rho\sim2,\quad K_A-(2-4h)\in\mathrm{Inf},\quad
 \chi-(4-4h)\in\mathrm{Inf},\quad
 (c_G+zK_r)-2(1+U)\in\mathrm{Inf}.
 \tag{SU2}
\]
The last infinitesimal is a finite separated expression for bounded
$u$ and infinitesimal $w$.  For the only new source assertion in
SU2, the sphere identities give exactly
\[
 d=\frac{2\mathcal H}{PJ}d_0^2,\qquad
 \rho=\frac{2xH_0}{\sigma B}
       +\frac{2\mathcal H}{B}b_0^2+\lambda_0s.
 \tag{SU3}
\]
Thus $d+\rho\sim2$ follows from $d_0^2+b_0^2=1$, with no inverse
of either sphere coordinate.  All other bounds follow from the exact source cancellations in
equations~\eqref{eq:c1-unification-nu}--\eqref{eq:c1-unification-KA} and the source estimates,
including $p-l\in zB_c$, in Section~\ref{sec:c1-unification-bounded-z-all-H}.

\paragraph{Shift the troublesome angular product into the normal block.}
Set $A_0=qW$, $A=A_\theta=(U-x)A_0$, and retain $R$.
There is an exact conductor:
\[
 A(A+uR)=u(U-x)wWZ-x(U-x)A_0^2.
 \tag{SU4}
\]
Consequently, collecting the radial terms in \eqref{eq:compression-full-parent} gives
\[
\begin{split}
 K_AA^2+\chi uAR+(c_G+zK_r)R^2
 ={}&2(A+uR)^2+2R^2-4hu(U-x)wWZ\\
 &+4hx(U-x)A_0^2+\delta_A A^2+
   \delta_\chi uAR+\delta_RR^2,
\end{split} \tag{SU5}
\]
where the three $\delta$ coefficients are the infinitesimals in
SU2.  This step keeps the signed product exactly; it is not an
inequality based on the sign of $U-x$.

\paragraph{A complete strict normal block.}
Introduce the legal separated roots
\[
 H=\sqrt h\,uZ,\qquad S=uwW-hu^2Z,\qquad N=N_x,
 \qquad Z_g=uwW=S+u\sqrt h\,H.
\]
Their exact allocations are
\[
\begin{split}
 \Gamma&=-S+\left(\frac{\sigma}{u\sqrt h}-u\sqrt h\right)H,\\
 P_h:=\sqrt h P_\beta&=(1-sh)H-\frac{s\sqrt h}{u}S,\\
 X_e&=N-\sqrt h P_h.
\end{split} \tag{SU6}
\]
Every coefficient is bounded; the factors $u^{-1}$ are pure testing
units.  The full scalar term $\mathscr S$ in \eqref{eq:compression-full-parent} also satisfies exactly
\[
 \mathscr P=L_hP_h^2+
      \rho(1-x/U)Z_g^2-\rho\sigma(U-x)A_0^2.
 \tag{SU7}
\]
The apparently signed factor $1-x/U$ is simply a bounded Gram
coefficient of residue one.  No mixed inverse is used.

Combine the $wWZ$ term in SU5 with every $N,\Gamma,P_h,Z_g$ term:
\[
\begin{split}
 \mathscr N={}&[2(1+w^2)+4mw^2]N^2
     +2(2D+1)U X_e^2+d\Gamma^2-4u\nu N\Gamma\\
 &+L_hP_h^2+\rho(1-x/U)Z_g^2
       -4\sqrt h\,\frac{U-x}{u}\,Z_gH .
\end{split} \tag{SU8}
\]
This is an exact equality for the entire indicated normal block.
Its coefficients in $(N,S,H)$ are bounded separated expressions.
Using SU2 and SU6, its residue is exactly
\[
 2(1+u_0^2)N^2+4u_0NS+2S^2+2H^2.
 \tag{SU9}
\]
The $N,S$ determinant is $4$, so this residue is strictly positive.
The $H$ cross terms cancel: the $X_e$ square cancels the leading
$N\Gamma$ coupling to $H$, and the angular shift in SU5 cancels
the remaining $S H$ and $hH^2$ terms.  This is where the sphere
sum $d_0^2+b_0^2=1$ is essential.

The remaining angular/radial block is
\[
\begin{split}
 \mathscr A={}&2\{(U-x)A_0+uR\}^2+2R^2\\
 &+(4hx-\rho\sigma)(U-x)A_0^2
   +\delta_A(U-x)^2A_0^2
   +\delta_\chi u(U-x)A_0R+\delta_RR^2.
\end{split} \tag{SU10}
\]
Its residue on $(A_0,R)$ has determinant $4u_0^4>0$.

\paragraph{Every other original-parent term.}
The unchanged hard $N_y$ diagonal has residue $2(1+u_0^2)$.
All non-Gamma cross terms of \eqref{eq:compression-full-parent} already have
infinitesimal coefficients on $(N,N_y,A_0,R)$, by SP's complete
opening-uniform coefficient bounds and $A=(U-x)A_0$.
The two $\Gamma$ terms have a factor $qw$ and become infinitesimal
crosses after the bounded allocation SU6.  Thus the complete Gram
on $(N,N_y,S,H,A_0,R)$ has strictly positive residue, the direct
sum of SU9, the $N_y$ entry and SU10.

The remaining exact nonnegative rows are retained once: both
centered cap excesses, $4\beta D_wN^2$, and
$4bpzU\Gamma^2$.  The $d\Gamma^2$ part and
$2(2D+1)(uX_e)^2$ have already been used in SU8 and are not
retained a second time.  All scalar, hard and first-normal terms
of the original SP6 parent are therefore included.

The bounded separated-Gram criterion proves tensor SOS for all real
states and both testing signs, including literal $Z=0$ and $w=0$.
Applying the real identity separately to real and imaginary parts
gives all complex states.  Neither source sphere endpoint is inverted.
The precise remaining C14.1b testing domain has $u$ infinitesimal
(with the earlier state and sign certificates still retained).



\subsection{Unit $L$: noninfinitesimal Alice tests}

\label{component:pb0d9d28d0849}

\label{pb0d9d28d0849:sec:r18-c142-remote}

\begin{proposition}
Keep the unchanged physical C14.2 source: $d=1$, $x>0$
infinitesimal, $L=x\zeta$ a positive bounded unit, $\eta$ unbounded,
all original strict gaps, nonnegative event weights, total-scale and
earlier-exclusion intersections.  Then the complete original parent
is finite tensor SOS for every noninfinitesimal Alice angle $u$,
including $u=\infty$, every independent Bob angle, and all real or
complex states.  Literal $D=0$ and $\beta=0$ are retained.  Thus
only infinitesimal or zero $u$, with arbitrary Bob type, remains in
C14.2.  No entire-source certificate is asserted.
\end{proposition}

\paragraph{Opening-safe source input.}
Use the source estimates in
Section~\ref{sec:c1-unification-bounded-z-all-H}, including
the exact positive source inequality~\eqref{eq:c1-unification-e-bound}
and its sphere and opening-cleared products.
Put $e_r=1-r$, $Q=1-Lr>0$ and define positive pure-source weights
\[
 C_0=\frac1{2Le_r},\qquad
 A_0=\frac{1+Le_r}{2Lr},\qquad
 a_0=\frac L{2Q}.
                                                               \tag{LA1}
\]
Their actual-weight ratios satisfy
$C/C_0,A/A_0,\alpha/a_0\sim1$ uniformly through $r$, $e_r$,
and $Q$ openings.  Each leading and actual weight has a positive
real lower bound, since $L$ is a positive bounded unit and
$0<r,e_r,Q<1$.  No upper bound for the individual weights is assumed.
The BE source sphere and both its endpoints remain unchanged; no
coefficient $D$, $\beta$, $d_0$ or $b_0$ is divided by.

\paragraph{The reciprocal Bob strip with a unit Alice angle.}
Let $u$ be a bounded nonzero unit and $w=1/v$ infinitesimal or zero.
Use four legal raw roots
\[
 P=uW,\quad X,\quad Y,\quad T=wZ,
 \qquad U=u^2,\quad D_u=1+U,\quad D_w=1+w^2.
\]
The original signed core is given in Section~\ref{component:p89fd3edb3edd}.
Written in these roots it has the exact split
\[
 K_w=K_0+w^2\{(P-Y)^2+X^2\}
                      +2uwX(T+P-Y),                         \tag{LA2}
\]
\[
 K_0=D_uX^2+D_uY^2+UT^2-2UTY-2PT.
\]
This is the complete core, including its sign; no state graph is used.
The selected positive events are
$CD_w(P-Y)^2$, $AD_uD_wY^2$, and
$\alpha D_u(Y-T)^2$ in the quarter-parent.  The whole original
$m,D,\beta$ terms are retained separately.

Direct expansion using LA1 gives the following exact identity for
the leading selected form at $w=0$:
\[
\begin{split}
 Q_0={}&C_0(P-Y)^2+A_0D_uY^2+a_0D_u(Y-T)^2+K_0/2\\
 ={}&\frac{D_u}{2}X^2
       +C_0\{P-Y-Le_rT\}^2
       +(A_0+a_0+1/2)\{Y-LrT\}^2\\
 &+U\{A_0Y^2+(a_0+1/2)(Y-T)^2\}.
\end{split}                                                   \tag{LA3}
\]
Every displayed coefficient is nonnegative pure source or a positive
pure-testing factor.  The source coefficients in LA3 can be clipped
at positive real lower bounds, retaining their exact positive square
excesses.  The resulting bounded Gram on $(P,X,Y,T)$ has strictly
positive residue: $X$ is controlled directly, the two $U$ terms
control $Y,T$ because $u$ is a unit, and the first contact square
then controls $P$.  The coefficients $Le_r,Lr$ in these roots are
bounded pure source, including all original openings.

\paragraph{Exact common-factor accounting for the actual weights.}
To transfer LA3 to the actual parent, define solely in the source field
\[
 \delta=|C/C_0-1|+|A/A_0-1|+|\alpha/a_0-1|,
 \qquad f=1-\delta.
\]
Then $\delta$ is nonnegative infinitesimal, $f>0$ is a unit,
and each of $C-fC_0$, $A-fA_0$, $\alpha-fa_0$ is nonnegative.
Retain their original event terms as exact squares.  The complete
remaining selected polynomial is, identically,
\[
\begin{split}
 S={}&f\{C_0D_w(P-Y)^2+A_0D_uD_wY^2
                          +a_0D_u(Y-T)^2\}+K_w/2\\
  ={}&fQ_0+fw^2\{C_0(P-Y)^2+A_0D_uY^2\}\\
    &+(1-f)K_0/2+(K_w-K_0)/2.
\end{split} \tag{LA4}
\]
The second line's $w^2$ terms are retained positive squares with
their complete actual source scales.  The third line, including the
signed $(1-f)K_0/2$ term, is a bounded infinitesimal Gram on the
four raw roots, by LA2 and $1-f=\delta\in\operatorname{Inf}$.
Clip the completed source-square coefficients of $fQ_0$ at positive
real lower bounds as above.  Its retained bounded Gram remains
strict under this entire infinitesimal perturbation.  The separated
bounded-Gram lemma gives tensor SOS; restoring all square excesses
proves the whole original parent.  This is exact polynomial
accounting in the tensor ring, not an inequality that orders mixed
source and testing quantities.

\paragraph{The other Bob types and the unbounded Alice endpoint.}
If $u$ is a bounded nonzero unit and $v$ is a bounded nonzero unit,
use the common noncommuting-residue theorem on the unchanged
physical parent.  If $u$ is a bounded nonzero unit and $v$ is
infinitesimal or zero, clip the original $C,A,\alpha$ at positive
real lower bounds $c_*,a_*,a'_*$ and retain every excess and the
whole original $m,D,\beta$ terms.  The complete selected residue is
\[
 (c_*+1/2)(uW-Y)^2+a_*D_uY^2
          +a'_*D_uZ^2+\tfrac12(X+uZ)^2\succ0.     \tag{LA5}
\]
It is strict on all four original states when $u$ is a unit; every
finite-angle correction has bounded infinitesimal coefficients.
This covers all bounded Bob types without invoking LA4 outside its
stated reciprocal strip.

For $u$ unbounded, put $p=1/u$ and use the original homogeneous
parent.  The exact CE4--CE5 raw-frame proof, with any positive real
caps for $C,A$, covers every bounded $v$.  Its strict selected
residue is
\[
 4c_*(1+v^2)W^2+4a_*(1+v^2)Y^2
                      +2(W+R)^2+2(B-vY)^2,
 \quad B=pX+Z,\quad R=vX.                         \tag{LA6}
\]
For unbounded $v$, the exact twice-homogeneous CE7--CE8 proof has
strict residue
\[
 4c_*W^2+4a_*Y^2+2X^2+2(T-Y)^2,
 \qquad T=wZ.                                    \tag{LA7}
\]
These two homogeneous proofs need only positive real lower bounds
for $C,A$ at these endpoints.  They do not use the stronger
``$C,A$ unbounded'' hypothesis needed for the finite-unit-angle part
of the adjacent-event theorem.  All unselected original terms and
positive excesses are retained exactly.  Literal $p=0$ and $w=0$
are the polynomial values of these original homogeneous parents.

The finite unit-Alice and homogeneous unbounded-Alice cases exhaust
all noninfinitesimal $u$.  Exact sign gauges and real/imaginary splitting
give both signs and complex states.  The possible additional $D=0$
null direction at $u=w=0$ is untouched, since the proposition does
not include the infinitesimal-Alice chart.  The sole remaining C14.2
testing obligation is now $u$ infinitesimal or zero, with arbitrary
Bob types and the unchanged full source intersection.



% Source formulas: pda4a21937f64
\subsection{Unit $L$: bounded Bob tests}

\label{component:pda4a21937f64}

\label{pda4a21937f64:sec:r18-c142-bounded-bob}

\begin{proposition}
Retain the unchanged full C14.2 source and every original intersection
as in the preceding noninfinitesimal-Alice proposition.  For every
infinitesimal or zero Alice test $u$ and every bounded Bob test $v$,
the entire original parent is finite tensor SOS for all real or
complex states.  Every source opening and both $D=0$, $\beta=0$
are included.  Together with that proposition this leaves exactly
$u$ infinitesimal and $v$ projectively unbounded, equivalently
$u,w=1/v$ infinitesimal, as the remaining C14.2 testing corner.
\end{proposition}

\paragraph{The same source weights and a pure-source strict matrix.}
Keep the opening-safe positive source weights from LA1,
\[
 C_0=\frac1{2L(1-r)},\quad
 A_0=\frac{1+L(1-r)}{2Lr},\quad
 a_0=\frac L{2(1-Lr)}.
\]
Let $c=(1+v^2)^{-1/2}$, $s=vc$, so $c^2+s^2=1$.
Since $v$ is bounded, $c$ is a positive testing unit.  Write
$s_0=\operatorname{res}s\in(-1,1)$ and $c_0^2=1-s_0^2>0$.
These are real numbers, not source/testing comparisons.  Consider
the pure-source quadratic form in three placeholder variables
\[
 M_0(P,Y,T)=
  (C_0+c_0^2/2)(P-Y)^2+(A_0+s_0^2/2)Y^2
          +a_0(s_0Y-T)^2-s_0PT.                    \tag{LB1}
\]
Its Gram has first leading minor $C_1=C_0+c_0^2/2>0$,
second leading minor $C_1[A_0+s_0^2(a_0+1/2)]>0$, and
\[
 \det M_0=
 \frac{c_0^2\{1+L(1-r)\}^2}
                {8Lr(1-r)(1-Lr)}>0.               \tag{LB2}
\]
This is an exact source identity, not just a limiting determinant.
Moreover $M_0^{-1}$ has bounded entries through every opening.
Indeed every cofactor is a sum of bounded multiples of
$C_0A_0,C_0a_0,A_0a_0,C_0,A_0,a_0,1$; terms involving a square
of a single divergent weight cancel.  Multiplication by
$8L^2r(1-r)(1-Lr)$ makes every cofactor polynomial in bounded
$L,r,s_0$.  Division by LB2 therefore leaves only the source-unit
denominators $L$, $c_0^2$, and $1+L(1-r)$.
Explicitly, with $e=1-r$ and $B_L=1+Le$,
\[
 M_0^{-1}=\frac1{c_0^2B_L}
 \begin{pmatrix}
 2L&2Lr&2s_0\\
 2Lr&2Lr(1-s_0^2e)&2s_0r(1+Le c_0^2)\\
 2s_0&2s_0r(1+Le c_0^2)&
       2(1-Lr c_0^2)(1+Le c_0^2)/L
 \end{pmatrix}.                                      \tag{LB2a}
\]

Choose any exact positive pure-source congruence
$M_0=N_1^2+N_2^2+N_3^2$.  Its inverse allocation of $P,Y,T$
on these three roots is bounded: each inverse-row norm squared is
the corresponding diagonal of $M_0^{-1}$.  This uses only a
pure-source Gram, with the fixed real number $s_0$ as parameter;
no mixed matrix is inverted or square-rooted.

\paragraph{Freeze the only divergent Bob-dependent event exactly.}
Use the legal raw roots
\[
 P=uW,\quad X,\quad Y,\quad
 T=cZ-(s-s_0)Y,
 \qquad d_s=s-s_0\in\operatorname{Inf}_{\rm test}.
                                                               \tag{LB3}
\]
In particular the entire alpha event satisfies
$sY-cZ=s_0Y-T$ exactly.  Thus no product of an unbounded alpha
weight with $d_s$ is ever estimated.  No inverse of $c$, $u$, or
a state coordinate is needed in these root definitions.

Let $B_0=c_0^2(P-Y)^2+s_0^2Y^2-2s_0PT$.
The complete original half-core after Bob normalization is
$K_B/2$, where
\[
 K_B=\{c(P-Y)+usX-sZ\}^2
       +(1+u^2)(sY-cZ)^2+(X+uZ)^2-(1+u^2)Z^2.
\]
It has the following exact, division-free allocation on LB3:
\[
\begin{split}
 K_B={}&X^2+B_0
       +d_sP\{-(s+s_0)P+2s_0Y-2T\}\\
       &+2ucX(sP-s_0Y+T)
       +u^2s^2X^2+u^2(s_0Y-T)^2.
\end{split} \tag{LB4}
\]
This keeps the full signed core and all $u$-dependent hard normals.

\paragraph{Exact actual-parent reassembly.}
Use exactly LA4's pure-source factor
$f=1-(|C/C_0-1|+|A/A_0-1|+|\alpha/a_0-1|)>0$.
Retain the actual original event excesses
$C-fC_0$, $A-fA_0$, $\alpha-fa_0$, and the whole $m,D,\beta$
terms.  On the finite Bob chart the original quarter-parent divided
by $1+v^2$ is then the sum of those positive terms and exactly
\[
 S=f\{C_0(P-Y)^2+A_0(1+u^2)Y^2
                         +a_0(1+u^2)(s_0Y-T)^2\}+K_B/2.
\]
Substitution of LB1 and LB4 gives
\[
\begin{split}
 S={}&fM_0+\tfrac12X^2
       +fA_0u^2Y^2+fa_0u^2(s_0Y-T)^2\\
    &+\tfrac12u^2s^2X^2+\tfrac12u^2(s_0Y-T)^2\\
    &+\tfrac12(1-f)B_0
       +\tfrac12d_sP\{-(s+s_0)P+2s_0Y-2T\}\\
    &+ucX(sP-s_0Y+T).
\end{split} \tag{LB5}
\]
The first two lines consist of a strict source-normalized Gram on
$(N_1,N_2,N_3,X)$ and exact retained positive squares.  All terms
in the last two lines have infinitesimal bounded coefficients on
that same frame: $1-f$ is source infinitesimal, $d_s,u$ are testing
infinitesimals, $c,s,s_0$ bounded, and LB2 supplies a bounded inverse
allocation for all three raw roots.  Since $f\sim1$, the resulting
complete Gram has residue $\operatorname{diag}(1,1,1,1/2)$.
The separated bounded-Gram lemma and restoration of every retained
term prove the full original parent.  The normalization factor
$1+v^2$ is pure testing and positive.

Both coefficient-boundary terms were retained rather than capped
or inverted.  The argument works literally at $u=0$ or $v=0$,
allows signed testing angles, and extends to complex states by
real/imaginary splitting.  It gives no claim for $c$ infinitesimal:
there $c_0=0$, the determinant in LB2 vanishes, and the reciprocal
two-small-angle corner remains a separate obligation.

The new focused checks establish LB2, opening-regular cofactors,
the complete original-core identity LB4 and the entire selected-parent
identity LB5.  They do not rerun prior source identities or assert
the remaining C14.2 corner closed.



% Source formulas: p98e371f5b846
\section{Contact gaps when $L$ vanishes}
Here the coefficient gaps, rather than $L$ alone, determine the remaining quadratic form. The gap $j$ first separates a strict block from its singular limit. The completed $J$ relation controls the remaining state normal when that block degenerates.

\subsection{Small $L$: the source openings and complementary tests}

\label{component:p6d2f659c0fde}

\label{p6d2f659c0fde:sec:r8-c1-small-L-bounded-eta}

Work in the original C1 source, with $d=1$, $x\in\Inf_{>0}$,
$\zeta$ noninfinitesimal, $L=x\zeta\in\Inf_{>0}$, and $\eta$
bounded. Native $r=1-a\eta$, $h=1-\theta$ remain strictly positive.
This section admits every original projective test except possibly
$u,v$ both infinitesimal. It does not admit the remaining two-soft
local chart.

\paragraph{Exact source restrictions, including the total-scale intersection.}
Put $\rho=r+\eta$ and $j=\eta J=\eta(1+xh)-1+r>0$.
Since $a<1+xh$, $\rho$ is a positive bounded unit. The exact beta
numerator is
\[
 2\beta\theta P_d=K\rho(xh+L)-\theta P_d.       \tag{SB1}
\]
It implies $\theta=O(x+L)=O(L)$, since $x/L=1/\zeta$ is bounded.
Consequently
\[
 h,M,E,Q,P_d,P_4\sim1,\qquad B_5\sim1.           \tag{SB2}
\]
For the last assertion, use the exact expression
$B_5=K(1+L\eta)\rho-jP_d$: its residue is
$r+\eta-(\eta-1+r)=1$. Thus original $D=0$ cannot occur on
this source face; literal $\beta=0$ remains included.
The original formulas, with openings canceled before taking residues,
give uniformly
\[
\begin{gathered}
 C(1-r)\sim\rho/2,\quad Ar/\eta\sim1/2,
 \quad Dj\sim1/2,\\
 mL/x^2\sim\rho/2,\quad
 \alpha/[L(L+x)]\sim\rho/2,\quad
 (2\beta+1)\theta/(x+L)\sim\rho.
\end{gathered}                                         \tag{SB3}
\]
In particular $C,D$ are strong, $m,\alpha$ weak, and $A$ is strong
precisely when $\eta$ is a unit. The four exact large-weight
alternatives are
\[
\begin{array}{c|c}
\text{weight unbounded}&\text{equivalent opening}\\\hline
C&1-r\in\Inf\\
D&j\in\Inf\\
A&r\in\Inf\\
\beta&\theta/(x+L)\in\Inf.
\end{array}                                                   \tag{SB4}
\]
For the $A$ equivalence, $r\in\Inf$ forces $\eta$ to be a unit
because $j>0$; if $r$ is a unit, $\eta/r$ is bounded.
All coefficients are bounded exactly when the four quantities in
the right column of (SB4) are positive bounded units. That entire
interior is proved by the positive-total-scale theorem.
After intersection with $\delta_{\rm tot}\in\Inf_{>0}$, at least
one of (SB4) therefore occurs. These are exhaustive necessary source
alternatives, not assertions of their independent realizability.

\paragraph{The unit-$\eta$ remote tests: two complete three-state blocks.}
Suppose first $\eta$ is a unit, with real residue $e>0$, and write
$r_0=\operatorname{res}r\in[0,1]$. For an interior opening,
the source residues are
\[
 C_0=\frac{r_0+e}{2(1-r_0)},\quad
 A_0=\frac e{2r_0},\quad D_0=\frac1{2(e+r_0-1)}.       \tag{SB5}
\]
Retain the original $m,\alpha,\beta$ terms throughout. Bounded
selected weights retain their actual coefficients; unbounded ones
are capped at sufficiently large finite real numbers and their
excesses retained as exact squares.

For $u$ a bounded nonzero unit and $v$ infinitesimal, use the full
raw frame
\[
 p=uW,\quad Y,\quad R=vZ,\quad B_t=X+uZ.
\]
The exact core is (CE2), and the $D$ event is
$J=p-Y+uvB_t-D_uR$. Every Gram coefficient is bounded after
capping, and no state normal is omitted. Set $U=\operatorname{res}u^2>0$.
The residue, divided by four, is $B_t^2/2$ plus
\[
\begin{split}
 f_U={}&C_0(p-Y)^2+A_0(1+U)Y^2
       +D_0\{p-Y-(1+U)R\}^2\\
 &+\tfrac12\{(p-Y)^2+(1+U)UR^2-2(1+U)pR\}.
\end{split}                                            \tag{SB6}
\]
Its first two principal minors and determinant are, respectively,
\[
\begin{gathered}
\Delta_1=\frac{e(e+r_0)}{2(1-r_0)j_0},\qquad
\Delta_2=\frac{e^2(1+U)(e+r_0)}{4r_0(1-r_0)j_0},\\
\Delta_3=\frac{Ue(1+U)^2(e+r_0)^2}{8r_0(1-r_0)j_0}>0,
\qquad j_0=e+r_0-1.
\end{gathered}                                      \tag{SB7}
\]

For $u$ infinitesimal and $v$ a bounded nonzero unit, use the full
frame $(p=uW,X,Y,R=vZ)$. The same exact parent has residue,
divided by four, $(1+V)X^2/2$ plus
\[
\begin{split}
 f_V={}&C_0(1+V)(p-Y)^2+A_0(1+V)Y^2
       +D_0(p-Y-R)^2\\
       &+\tfrac12\{(p-Y)^2+VY^2-2pR\},
 \qquad V=\operatorname{res}v^2>0.
\end{split}                                            \tag{SB8}
\]
The leading two minors are positive, and
\[
 \det f_V=\frac{V(1+V)(e+r_0)^2}{8r_0(1-r_0)j_0}>0. \tag{SB9}
\]

These arguments include every opening in (SB4). To see explicitly
why the finite cap remains valid, impose the kernels of the divergent
events and inspect the remaining two-state Gram. The determinants
are
\[
\renewcommand{\arraystretch}{1.7}
\begin{array}{c|c|c}
\text{opening}&\det(f_U|_{\ker})&\det(f_V|_{\ker})\\\hline
r_0=0,\ e>1&\dfrac{Ue^2(1+U)}{4(e-1)}&\dfrac{Ve}{4(e-1)}\\
r_0=1&\dfrac{U(1+U)^2(e+1)}4&\dfrac{V(e+1)}{4e}\\
j_0=0,\ 0<r_0<1&
 \dfrac{U(1+U)^2}{4r_0}&\dfrac{V(1+V)}{4r_0(1-r_0)}.
\end{array}                                                   \tag{SB10}
\]
The respective kernels are $Y=0$, $p=Y$, and
$p-Y=(1+U)R$ for $f_U$ or $p-Y=R$ for $f_V$.
All first pivots are positive. At the only double opening,
$r_0=j_0=0,e=1$, the two kernels leave the positive scalar
$U(1+U)R^2/2$, respectively $VR^2/2$. There are no other
multiple openings since $e>0$. The elementary finite-penalty
Schur lemma therefore supplies finite real caps simultaneously
for every divergent selected event.

For completeness the remaining projective ends also use original
homogeneous degree. If $u$ is unbounded, (CE4)--(CE5) for bounded
$v$, and (CE7)--(CE8) for unbounded $v$, are strict with any positive
finite $C,A$ caps. If $v$ is unbounded and $u$ bounded, put
$w=1/v$ and use $(p=uW,X,Y,Z)$. The exact $D$ event is
$w(p-Y)+uX-Z$. At $w=0$ the core is
$(1+u^2)(X^2+Y^2)$, while the three caps give
$C(p-Y)^2+A(1+u^2)Y^2+D(uX-Z)^2$; these are strict on
all four roots, including $u=0$. Each term suppressed in this
residue has a bounded infinitesimal coefficient in this same
full frame. Literal $w=0$ is therefore included. The both-unit
tests are covered by the common noncommuting-residue theorem.
Together these statements exhaust the complement of $u,v\in\Inf$.

\paragraph{The infinitesimal-$\eta$ remote tests.}
Here $r\sim1$, $C,D$ are unbounded, and $y=L\eta$ is infinitesimal.
Define the pure-source harmonic coefficient
$\gamma=yCD/(C+D)$. Uniformly through the $a,J$ openings,
\[
 \gamma\sim L/2,\qquad A/y\sim1/(2L),\qquad
 \gamma A/y\sim1/4,\qquad y/\gamma\sim2\eta\in\Inf. \tag{SB11}
\]
The strong-II remote proof (RR3)--(RR6), which uses only $C,D$
unbounded, covers every noninfinitesimal $v$ and every $u$.
For $u$ noninfinitesimal and $v$ infinitesimal, use (RR7)--(RR9)
with zero alpha cap and $k=-1$. In its notation let
$a=u/\sqrt{1+u^2}$, $b=1/\sqrt{1+u^2}$,
$\ell=vY$, $x_t=bX$, $R_0=x_t+a(Z-\ell)$, and $q=\sqrt yY$.
Then
\[
 M_{-1}=\tfrac12R_0^2+\tfrac12(b\ell)^2,\qquad
 v(Z-\ell)=avR_0-bv(b\ell)-b\sqrt y\,V_h,       \tag{SB12}
\]
where $V_h=v(aX-bZ)/\sqrt y$. Normalize the harmonic pair as
$P_h=\sqrt\gamma V_h$, $Q_h=\sqrt{A/y}\,q$. Its residue
determinant is $1-b_0^2>0$ by (SB11), where $b_0=\operatorname{res}b<1$;
the off-diagonal entry is $b_0$. The sole negative term in (RR9),
$-v^2(Z-\ell)^2/2$, is an infinitesimal Gram perturbation of
$(R_0,b\ell,P_h)$ by (SB12), since $\sqrt{y/\gamma}$ is
infinitesimal. Thus the full four-root Gram is strict. Retain the
original positive $\alpha(Z-\ell)^2$, the C1 beta event, and every
other unused term. The homogeneous endpoint $u=\infty$ has
$a=\pm1,b=0$ and is included; literal $v=0$ also causes no change.



% Source formulas: pc15b7caff0a3
\subsection{The small-$L$ native scalar and complete state normals}

\label{component:pc15b7caff0a3}

\label{pc15b7caff0a3:sec:r8-small-L-native-normal}

Normalize $d=1$, use the native C1 source $\lambda=n=0$, and keep
all other original event weights, including literal $\beta=0$.
Put $U=u^2,V=v^2,D_u=1+U,D_v=1+V$, $\theta=1-h$, and
\[
             X=-huZ+\xi,\qquad Y=ruW+\nu.
\]
Here $\nu$ denotes this state normal only.  Direct expansion of the
original event terms gives
\[
 \frac{F_g}{4}=A_wW^2+A_zZ^2-2uv(b_0+b_1U)WZ,
\]
with the complete graph coefficients
\[
\begin{aligned}
 A_w={}&[(C+D+\tfrac12)(1-r)^2+Ar^2]U+Ar^2U^2\\
 &+\{m+[2mr+C(1-r)^2+(A+\alpha+\tfrac12)r^2]U
                  +r^2(m+A+\alpha+\tfrac12)U^2\}V,\\
 A_z={}&\alpha+[\alpha+(m+\beta+\tfrac12)\theta^2]U\\
 &+\{D+[2Dh+\tfrac12h^2+\beta\theta^2]U
                     +(D+\tfrac12)h^2U^2\}V,\\
 b_0={}&m\theta+D(1-r)+\alpha r+\tfrac12,\\
 b_1={}&mr\theta+(D+\tfrac12)h(1-r)+(\alpha+\tfrac12)r.
\end{aligned}
\]
These identities are homogeneous in $W,Z$ and need no state-ratio
bound.  The first-normal and hard-block formulas that follow are
obtained by the same substitution, or by differentiating the original
parent twice in $X,Y$.
The entire original parent is
exactly
\[
\begin{split}
 F/4={}&F_g/4+2\xi(vWf_w+uZf_z)+2\nu(uWk_w+vZk_z)\\
 &+H_{xx}\xi^2+2H_{xy}\xi\nu+H_{yy}\nu^2.
\end{split}                                               \tag{SN1}
\]
where
\[
\begin{aligned}
 f_w={}&-m(1+rU)+(D+\tfrac12)(1-r)U,\\
 f_z={}&m\theta-(D+\tfrac12)V(1+hU)
                    +(\beta+\tfrac12)D_v\theta,\\
 k_w={}&mV(1+rU)-CD_v(1-r)+AD_uD_vr
             -(D+\tfrac12)(1-r)+(\alpha+\tfrac12)D_urV,\\
 k_z={}&-m\theta U+(D+\tfrac12)(1+hU)
                              -(\alpha+\tfrac12)D_u,\\
 H_{xx}={}&m+\beta D_v+(D+\tfrac12)UV+D_v/2,\\
 H_{yy}={}&mUV+CD_v+AD_uD_v+D+\tfrac12
                              +(\alpha+\tfrac12)D_uV,\\
 H_{xy}={}&-(m+D+\tfrac12)uv.
\end{aligned}                                                   \tag{SN2}
\]
There is no discarded normal or excess term in this identity.
Its two exact first-normal contact equations are
\[
 \zeta f_w(x,y)+f_z(x,y)=0,\qquad
 k_w(x,y)+\eta k_z(x,y)=0,\qquad y=x\eta\zeta.             \tag{SN3}
\]
They follow by substituting (NC4), not by setting a testing angle
equal to a source angle in the desired tensor assertion.

On the small-$L$, bounded-$\eta$ source, with $x,L$ infinitesimal
and $\zeta$ noninfinitesimal, the beta-opening identity gives:
\[
 2\beta\theta P_d=K(r+\eta)(xh+L)-\theta P_d,
 \qquad \beta\theta=O(L),\quad \beta\theta^2=O(L^2).       \tag{SN4}
\]
In particular, after the pure-source normalization
$\widehat\xi=\sqrt{1+\beta}\,\xi$, its cross coefficient is
$\beta\theta/\sqrt{1+\beta}=O(L)$, including every unbounded
beta opening.  The complete normal allocations are the retained
RO/JO/HO and DG/LD/CJ terms.



% Source formulas: pe3ef590c01d6
\subsection{The small-$L$ scalar and finite Taylor quotients}

\label{component:pe3ef590c01d6}

\label{pe3ef590c01d6:sec:r9-c1-small-interior}

Use the normalized native C1 parent, $d=1$, with $x=k^2>0$
infinitesimal, and set
\[
 \rho=\eta+r,\quad y=x\eta\zeta,\quad L=x\zeta,\quad
 h=1-\theta,\quad s=\sqrt{\eta\zeta},\quad
 \lambda=\rho/(2\eta).
\]
The reference interior has
\[
 \eta,\zeta,r,1-r,j_0=\eta+r-1
       \quad\text{positive bounded units}.                    \tag{IC1}
\]
The formulas below are the exact scalar and full-normal dictionary.
The opening-uniform RO proof supplies their normalized bounds when
$r$ or $1-r$ tends to zero.  The displayed reference $X$ normal
uses $\sqrt\beta$ when $\beta>0$; RO replaces it throughout by
$\sqrt{\beta+1/2}$, which includes literal $\beta=0$.
The full scalar graph coefficients $A_w,A_z,b_0,b_1$ are printed
explicitly in the native dictionary SN1--SN2.
For IC4--IC21, first assume $Z\ne0$ and normalize $Z=1$ inside
the testing field.  Homogenizing the resulting state-linear roots
restores every $Z$; the literal $Z=0$ state has the direct core
certificate.  The reciprocal identities below are already homogeneous.
\paragraph{A bounded-ratio frame retaining all three state normals.}
First suppose the pure-testing ratio $t=v/u$ is bounded, possibly
infinitesimal. Define
\[
 \begin{gathered}
 E=\eta uW-v,\quad
 b={r(1-r)\over\eta(\eta+2r-1)},\quad
 N=Y-ruW-bE,\\
 \xi=X+hu,\qquad N_x=\sqrt\beta\,\xi,\quad h=1-\theta.
                                                               \end{gathered} \tag{IC4}
\]
The state change is exact and uses only the pure-testing inverse $u$:
$W=(v+E)/(\eta u)$, $Y=r(v+E)/\eta+bE+N$, $X=-hu+\xi$.
Every normal vanishes at the actual source contact. 

\paragraph{The exact scalar, with the negative defect absorbed.}
Set $U=u^2,V=v^2$, $s=\sqrt{\eta\zeta}$, $k=\sqrt x$, and introduce
the positive separated roots
\[
 A_t=uv,\quad B_t={k v^2\over s u},\quad
 C_t=xs,\quad G=kv,\qquad F_v=V-y.
                                                               \tag{IC6}
\]
They satisfy $A_tB_tC_t=G^3$ and coincide at actual contact
$(u,t)=(k,s)$. The pure-testing pole in $B_t$ is inside the root;
none of the coefficient arrays below has a negative power of $t$.

Use $A_w,A_z,b_0,b_1$ from (SN1)--(SN2), including the beta term.
The scalar at $E=N=\xi=0$ is exactly
\[
 \mathcal S={VA_w\over\eta^2U}+A_z
                  -{2V\over\eta}(b_0+b_1U)
       =c_FF_v^2+f_AA_t^2+f_BB_t^2+f_CC_t^2-f_GG^2.    \tag{IC7}
\]
Here the complete coefficients are
\[
\begin{aligned}
 c_F={}&{C(1-r)^2+Ar^2+(\alpha+\tfrac12)r^2+2mr\over\eta^2},\\
 a_{UV}={}&{Ar^2\over\eta^2}+{h^2\over2}+2Dh+\beta\theta^2
                                     -{2b_1\over\eta},\\
 a_{UUV}={}&h^2(D+\tfrac12),\qquad
 a_{UVV}={r^2(m+A+\alpha+\tfrac12)\over\eta^2},\\
 q_V={}&{C(1-r)^2+Ar^2+(D+\tfrac12)(1-r)^2\over\eta^2}
                                     +D-{2b_0\over\eta},\\
 q_U={}&\alpha+(m+\beta+\tfrac12)\theta^2,\\
 f_A={}&a_{UV}+a_{UUV}U+a_{UVV}V,\qquad
 f_B={m\zeta\over\eta x},\\
 f_C={}&{\alpha-c_Fy^2+q_UU\over\eta\zeta x^2},\qquad
 f_G=-{q_V+2c_Fy\over x}.                              \end{aligned} \tag{IC8}
\]

\paragraph{Finite scalar Taylor arrays.}
Let $\epsilon_i=f_i-\lambda$ and $\epsilon_G=f_G-3\lambda$,
$h_c=\epsilon_A+\epsilon_B+\epsilon_C-\epsilon_G$.
Equations (IC8) show that $h_c=h_{00}+aU+cV$ with bounded
pure-source $a,c$. Exact contact value gives
\[
 h_c=a(u^2-k^2)+c(u^2t^2-k^2s^2).                     \tag{IC11}
\]
For $d_u=u-k,d_t=t-s$, its finite second Taylor identity is
\[
\begin{split}
 h_c={}&2k(a+cs^2)d_u+2k^2sc\,d_t
 +(a+ct^2)d_u^2+2kc(t+s)d_ud_t+k^2c\,d_t^2.            \end{split} \tag{IC12}
\]
The coefficient of $d_t^2$ is infinitesimal. For any affine function
$f=f_{00}+a_fU+c_fV$ among (IC8), the exact first difference is
\[
 f-f(k,s)=(u-k)(u+k)(a_f+c_ft^2)
                    +(t-s)k^2(t+s)c_f.                \tag{IC13}
\]
These are the required bounded $u$ arrays and infinitesimal $t$ arrays;
no appeal to unspecified Taylor membership is needed.
The two exact first-jet balances are
\[
 2(\epsilon_{A0}-\epsilon_{C0})+k h_{c,u0}=0,
 \quad 2(\epsilon_{B0}-\epsilon_{C0})+s h_{c,t0}=0.      \tag{IC14}
\]
They follow from the original stationary contact, and are also checked
directly using (NC4). The extra square $c_FF_v^2$ has zero value and
first jets there. Apply the accompanying three-circuit contact-lift
lemma with the exact conductors
\[
 G(u-k)=k(A_t-G),\qquad G(t-s)=s(B_t-G).                \tag{IC15}
\]

\paragraph{Explicit finite allocations of every first normal.}
Use $f_w,f_z,k_w,k_z$ in (SN2), and define the polynomials
\[
\begin{aligned}
 P_Y(U,V)&=2(k_w/\eta+k_z),\\
 P_E(U,V)&={2\over\eta^2}(A_w/U-mV/U)
              -{2\over\eta}(b_0+b_1U)+bP_Y(U,V).
\end{aligned}                                                   \tag{IC16}
\]
The quotient $A_w/U-mV/U$ is polynomial, as its displayed numerator
shows. For either polynomial $P$ use the exact divisions
\[
 P=P_{00}+UP_U+VP_V,\quad
 P_U={P-P(0,V)\over U},\quad
 P_V={P(0,V)-P(0,0)\over V}.                           \tag{IC17}
\]
The normalized bounds on these exact arrays are supplied by RO. The complete first normals
of the original $F/4$ in the frame (IC4) are exactly
\[
\begin{aligned}
 \mathcal L_N&=vP_Y,\\
 \mathcal L_E&=2mv t^2/\eta^2+vP_E,\\
 \mathcal L_x&={1\over\sqrt\beta}\left[
 2u\theta(m+\beta+\tfrac12)
 -{2m\over\eta}ut^2(1+rU)+2uVZ_X\right],\\
 Z_X&=(\beta+\tfrac12)\theta-(D+\tfrac12)(1+hU)
                         +(D+\tfrac12)(1-r)/\eta .     \end{aligned} \tag{IC18}
\]
These identities follow by differentiating the entire original parent
after the full state substitution (IC4).

Their actual root allocations are
\[
\begin{aligned}
 vP&={P_{00}\over k}G+u(P_U+t^2P_V)A_t,\\
 {2m\over\eta^2}vt^2&={2ms t\over\eta^2k}B_t,\\
 \mathcal L_x&={2u\theta(m+\beta+\tfrac12)
                    \over\sqrt\beta\,xs}C_t
 -{2ms(1+rU)\over\eta k\sqrt\beta}B_t
 +{2vZ_X\over\sqrt\beta}A_t.                          \end{aligned} \tag{IC19}
\]

The exact original contact makes each complete normal zero at
$(u,t)=(k,s)$. For a normal $\ell=\sum b_iA_i+b_GG$, put
$B=\sum b_i+b_G$, and use the explicit monic divisions
\[
 B=(u-k){B(u,t)-B(k,t)\over u-k}
       +(t-s){B(k,t)-B(k,s)\over t-s}.                 \tag{IC21}
\]

\paragraph{Reciprocal testing-ratio identities.}
Return to homogeneous state coordinates and set
\[
 p=uW,\ q=vZ,\ E=\eta p-q,\quad
 Y=rp+bE+N,\quad B=X+uZ=\xi+\theta uZ.
\]
Use the complete raw root list
\[
 E,\ N,\ N_x=\sqrt\beta\xi,\quad
 P=\sqrt m\,vW,\quad Q=\sqrt\alpha Z,\quad
 R_t=u^2W,\quad S_t=uvW.                              \tag{IC22}
\]
\[
 \sqrt x\,p=\sqrt{x/m}\,(u/v)P,\quad
 Y=rp+bE+N,\quad q=\eta p-E,                          \tag{IC23}
\]
For completeness, with $B=0$ the exact higher-angle part of the
$C,A,D$ terms and half core, after the quadratic base is removed, is
\[
\begin{split}
 &U\{AY^2-2Dq(p-Y-q)+q^2/2-pq\}\\
 &\quad+V\{C(p-Y)^2+(A+\tfrac12)Y^2\}
 +(A+\tfrac12)UVY^2+(D+\tfrac12)U^2q^2.               \end{split} \tag{IC24}
\]
The terms involving $B$ are equally explicit. Relative to $B=0$, the
$D$ term adds $2DuvB(p-Y-D_uq)+DUVB^2$, and the half core adds
$uvB(p-Y-D_uq)+(D_v+UV)B^2/2$.
Here
\[
 B={N_x\over\sqrt\beta}+{\theta u\over\sqrt\alpha}Q,
 \quad uvp=uS_t,
\]
The opening-specific proof bounds their coefficients.  Write
$M=vW-X+uvY-uZ$. The three remaining
original terms keep their entire roots:
\[
\begin{aligned}
 \sqrt m\,M&=P+\sqrt m\,uvY-\sqrt m\,B,\\
 \sqrt\alpha\sqrt{D_u}(vY-Z)
  &=\sqrt{D_u}\{\sqrt\alpha(rS_t+v(bE+N))-Q\},\\
 \sqrt\beta\sqrt{D_v}B
  &=\sqrt{D_v}\{N_x+\sqrt{\beta\theta^2/\alpha}\,uQ\}.
\end{aligned}                                                   \tag{IC26}
\]



% Source formulas: p3df6e6bdd8f4
\subsection{Unit positive $j$: all $r$ openings}

\label{component:p3df6e6bdd8f4}

\label{p3df6e6bdd8f4:sec:r13-c1-r-openings}

\begin{theorem}
Use the original physical native C1 source, normalized by $d=1$, and
retain all original intersections. Suppose that $x=k^2$ is positive
infinitesimal and
\[
 \eta,\quad\zeta,\quad j_0=\eta+r-1
       \quad\hbox{are positive bounded units}.                 \tag{RO1}
\]
Then the complete original parent is tensor SOS at every independent
projective test and every state. The original $0<r<1$ may approach
either endpoint at any positive rate, and the literal $\beta=0$ face
is included. No lower positive real bound on $r$, $1-r$, or $\beta$
is imposed. This extends the covered IC1 domain; it does
not cover an infinitesimal $j_0$, infinitesimal $\eta$, or unbounded
$\zeta$.
\end{theorem}

All scalar and Gram identities in this section use the normalized
parent $F/4$; multiplication of every square root by $2$ restores $F$.

The proof reuses the unchanged exact scalar and full-normal identities
(IC7)--(IC8), (IC16)--(IC19), and the three-circuit lemma. Its new
content is a uniform weighting of the two divergent event openings,
and a different $X$ normal that retains bounded and zero beta.

\paragraph{Opening-cleared source arrays.}
Put $\rho=\eta+r$, $y=x\eta\zeta$, $h=1-\theta$,
$c=C(1-r)$ and $a=Ar$. The beta numerator (SB1) gives
$\theta=x\pi$ with $\pi>0$ bounded. We do not require $\pi$
infinitesimal. The original source formulas give, uniformly for
all $0<r<1$,
\[
\begin{gathered}
 c=\rho/2+O(x),\quad a=\eta/2+O(x),\quad
 D=1/(2j_0)+O(x),\quad h=1+O(x),\\
 m/x\sim\rho/(2\zeta),\qquad
 \alpha/x^2\sim\rho\zeta(1+\zeta)/2,\qquad
 B_x\theta/x={K\rho(h+\zeta)\over2P_d},\quad
 B_x=\beta+\tfrac12\ge\tfrac12.                      
\end{gathered} \tag{RO2}
\]
These are finite canceled numerator identities, followed by residue
evaluation. Every canceled denominator at $x=0$ is a product of
positive units from $\eta,\zeta,j_0,j_0+r$ and $\rho$, with no
factor $r$, $1-r$, or $\pi$. In particular $D>0$ on this domain;
the original $D=0$ face has empty intersection here, as already proved
in (SB2). Both $\theta/\sqrt\alpha$ and
$B_x\theta^2/\alpha$ are bounded, not necessarily infinitesimal.

The exact source formulas for the two cleared strong weights are
\[
 c={K\rho E\over2P_d},\qquad a={\eta P_4\over2P_d}.
\]
Their order-$x$ differences in (RO2), and the following arrays, have
opening-free finite canceled formulas:
\[
 {P_{Y,00}\over x},\quad {P_{E,00}\over x},\quad
 {\theta(m+B_x)\over x},\quad
 {\alpha+(m+B_x)\theta^2\over x^2},\quad {q_V\over x}.
                                                               \tag{RO3}
\]
Here $P_Y,P_E,q_V$ are precisely (IC16) and (IC8). To check why no
unbounded event remains in these arrays, replace every occurrence of
$C(1-r)^2$ by $c(1-r)$, every $Ar^2$ by $ar$, every
$C(1-r)$ by $c$, and every $Ar$ by $a$ before division. No other
occurrence of $C$ or $A$ survives in the scalar or first-normal
polynomials.  This is a finite symbolic
criterion for boundedness on the stated field domain, not a numerical
endpoint approximation.

\paragraph{The normalized hard frame.}
Define the pure-source coefficients
\[
 b={r(1-r)\over\eta(j_0+r)},\quad
 H_0={\eta(j_0+r)\over2r(1-r)j_0},\quad
 \gamma={\rho\over2\eta(j_0+r)}.
\]
Both $1/\sqrt{H_0}$ and $1/\sqrt{B_x}$ are bounded; $H_0$
itself may be unbounded. Put
\[
 E_t=\eta uW-vZ,\quad
 N_y=\sqrt{H_0}(Y-ruW-bE_t),\quad
 N_x=\sqrt{B_x}(X+huZ).                              \tag{RO4}
\]
All three normals vanish at the actual source contact. For the
comparison weights
\[
 C_0={\rho\over2(1-r)},\quad A_0={\eta\over2r},\quad
 D_0={1\over2j_0},
\]
the exact quadratic base is
\[
 C_0(p-Y)^2+A_0Y^2+D_0(p-Y-q)^2+\tfrac12(p-Y)^2-pq
       =H_0(Y-rp-b(\eta p-q))^2+\gamma(\eta p-q)^2.  \tag{RO5}
\]
Thus the divergent normal has already been normalized before taking
residue. The cancellation factors needed for its entire hard block are
\[
 {C_0\over H_0}={r\rho j_0\over\eta(j_0+r)},\quad
 {A_0\over H_0}={(1-r)j_0\over j_0+r},\quad
 {D_0\over H_0}={r(1-r)\over\eta(j_0+r)}.             \tag{RO6}
\]
For example, after $p=(v+E_t)/\eta$ and $Z=1$, the $E_t$
coefficients in $p-Y$ and $Y$ are respectively
\[
 {1-r\over\eta}-b={(1-r)j_0\over\eta(j_0+r)},\qquad
 {r\over\eta}+b={r\rho\over\eta(j_0+r)}.            \tag{RO7}
\]
The strong $C,A$ factors therefore cancel even on those complete
first-state directions. For instance an error
$\delta_c(p-Y)^2/(1-r)$, with $\delta_c=O(x)$, has a
$N_y^2$ coefficient $\delta_c/((1-r)H_0)=O(x)$,
an $E_t^2$ coefficient $O(x)(1-r)$, and an $E_tN_y$
coefficient $O(x)/\sqrt{H_0}$. The $A$ error has the same
properties using $rH_0$ and the second formula in (RO7).
The original higher-angle factors have the same cleared coefficients,
multiplied by $U=u^2$ or $V=v^2$.

Suppose first that $u,v>0$ are infinitesimal and the pure-testing
ratio $t=v/u$ is bounded. After the state substitution (RO4), the
complete hard quadratic on $(E_t,N_y,N_x)$ consequently has bounded
entries and residue
\[
                         \operatorname{diag}(\gamma,1,1).       \tag{RO8}
\]
Here $\gamma$ is a positive bounded unit even when either $r$
opening is infinitesimal. The additional $m t^2$ contribution is
infinitesimal. The beta and half-core coefficients jointly give the
third pivot because $m/B_x$ is infinitesimal. These estimates concern
the full hard block, not only its restriction to the contact graph.

\paragraph{The complete scalar and all first normals.}
Normalize $Z=1$, homogenizing at the end as in IC1. Use the unchanged
positive separated roots
\[
 A_t=uv,\quad B_t={kv^2\over s u},\quad C_t=xs,\quad
 G=kv,\quad F_v=V-y,\qquad s=\sqrt{\eta\zeta}.
\]
The scalar at all three normals zero is exactly (IC7). Its coefficient
residues are still
\[
 c_F\sim{\rho\over2\eta^2},\quad
 f_A,f_B,f_C\sim\lambda={\rho\over2\eta},\quad
 f_G\sim3\lambda.                                   \tag{RO9}
\]
They do not depend on the residue of $\pi$. Every coefficient of
$U,V$ in (IC8) is bounded after the clearing described above;
for the additional term in $f_C$ use
$B_x\theta^2/x^2=(B_x\theta/x)\pi$. The exact contact
value and the two first-jet identities (IC11)--(IC14) are unchanged.
Their finite arrays therefore obey precisely the hypotheses of the
three-circuit lemma, with $s$ a unit. It gives the same strict finite
circuit Gram and the strict $F_v$ pivot. In particular the negative
soft defect has been allocated by the entire scalar (IC7), not
dropped by comparing only its quadratic base.

The complete first normals are (IC18), with the $N$ line divided
by $\sqrt{H_0}$ and the $X$ line divided by $\sqrt{B_x}$
instead of $\sqrt\beta$. Their exact root allocations are (IC19)
with those same replacements. Every coefficient remains infinitesimal:
the $N_y$ factors are bounded source factors times the old cleared
polynomials; the $X$ coefficient on $C_t$ is
\[
 {2u\theta(m+B_x)\over\sqrt{B_x}\,xs},
\]
which is infinitesimal by (RO3), and its coefficients on $B_t,A_t$
are respectively $O(k)$ and $O(v)$. In the latter coefficient
$Z_X=B_x\theta-(D+\tfrac12)(1+hU)
 +(D+\tfrac12)(1-r)/\eta$ is bounded. The $E_t$ line uses
$m/x,P_{E,00}/x$ and bounded polynomial arrays, as in IC1.
There is no negative power of $t$ in any coefficient.

Each coefficient is polynomial in $u,t$ with bounded pure-source
arrays; every term depending on $t$ carries $u,k$ or $m/k=O(k)$.
The exact monic differences (IC21) hence have bounded $u$ quotient
and infinitesimal $t$ quotient after evaluation at $u=k$. The
unchanged exact source contact cancels the bare $G$ term in each
complete normal. The conductors
$G(u-k)=k(A_t-G)$ and $G(t-s)=s(B_t-G)$ therefore allocate
all normals with infinitesimal coefficients on the strict circuit
roots. Joining this allocation to (RO8) proves the entire parent
on the bounded-ratio chart.

\paragraph{The reciprocal ratio chart.}
For $v/u$ unbounded use the homogeneous root list
\[
 E_t,\quad N_y,\quad N_x,\quad P_m=\sqrt m\,vW,
 \quad Q_\alpha=\sqrt\alpha Z,\quad R_t=u^2W,
 \quad S_t=uvW.                                      \tag{RO10}
\]
The error of the actual quadratic base relative to (RO5) is an
$O(x)$ quadratic in $(p,E_t,N_y)$ after the opening cancellations
(RO6)--(RO7). Every $\sqrt x\,p$ factor is allocated exactly by
\[
 \sqrt x\,p=\sqrt{x/m}\,(u/v)P_m,
\]
with a bounded source unit $x/m$ and an infinitesimal pure-testing
ratio $u/v$. Thus the entire error has zero Gram residue on (RO10).
All coefficients in the exact higher-angle identity (IC24) remain
bounded after $Y=rp+bE_t+N_y/\sqrt{H_0}$. This follows term
by term from (RO6), $b=O(r(1-r))$, $Ar=a$, and $C(1-r)=c$.
Its only nonzero scalar residues are the two positive diagonal entries
\[
                 \kappa_R=\rho\eta/2,\qquad
                 \kappa_S=\rho/2.                  \tag{RO11}
\]
No opening has been substituted before cancellation.

It remains essential to combine beta with its half-core part. Set
$B_t^{\rm state}=X+uZ$ and $J_0=p-Y-D_uvZ$. Relative to
$B_t^{\rm state}=0$, the exact $D$, beta and half-core increment is
\[
 2(D+\tfrac12)uvB_t^{\rm state}J_0
 +(D+\tfrac12)UV(B_t^{\rm state})^2
 +B_xD_v(B_t^{\rm state})^2.                          \tag{RO12}
\]
This whole-polynomial identity is checked in  Use
\[
 B_t^{\rm state}={N_x\over\sqrt{B_x}}
               +{\theta u\over\sqrt\alpha}Q_\alpha.
\]
The first two terms of (RO12) have zero Gram residue on (RO10);
$D$ is bounded, $1/\sqrt{B_x}$ and $\theta/\sqrt\alpha$
are bounded, and every coefficient carries $u$ or $v$.
The three remaining positive terms have their entire roots
\[
\begin{aligned}
 \sqrt m\,R_m&=P_m+\sqrt m\,uvY-\sqrt m\,B_t^{\rm state},\\
 \sqrt\alpha\sqrt{D_u}(vY-Z)
   &=\sqrt{D_u}\{\sqrt\alpha[rS_t+v(bE_t+N_y/\sqrt{H_0})]
                         -Q_\alpha\},\\
 \sqrt{B_x}\sqrt{D_v}B_t^{\rm state}
   &=\sqrt{D_v}\{N_x+\sqrt{B_x\theta^2/\alpha}\,uQ_\alpha\}.
\end{aligned}                                                    \tag{RO13}
\]
Each is its leading root plus an infinitesimal-coefficient combination
of (RO10). The square root in the last line is pure source and bounded;
its factor $u$ supplies the needed infinitesimal even when beta is
zero or bounded. The full parent therefore has a bounded Gram with
strict diagonal residue
\[
                 (\gamma,1,1,1,1,\kappa_R,\kappa_S).
                                                               \tag{RO14}
\]
This proves the remaining ratio chart.

\paragraph{All-tests join and exact scope.}
The earlier small-$L$, bounded-$\eta$ remote theorem includes (RO1)
with every original opening, so it covers the complement of
$u,v$ both infinitesimal. Literal axes have the elementary full-core
identities recorded in IC1. Its exact sign gauges apply unchanged.
Homogenizing the $Z=1$ identity gives polynomial state identities,
including $Z=0$; applying them to real and imaginary parts covers
complex states. All inverse and radical factors above are pure
source or pure testing. This completes the claimed full-parent
all-tests certificate.

The newly removed source region is the portion of (RO1) outside IC1,
namely either $r$ or $1-r$ infinitesimal, with all total-scale and
native source inequalities retained. Remaining C14.3 source
configurations still include infinitesimal $\eta$, unbounded
$\zeta$, and nonunit or nonpositive $j_0$ while the actual
$j=\eta(1+xh)-1+r$ remains positive. No conclusion on those
regions or on C14.4 follows from this theorem.

\paragraph{A local algebraic extension below beta zero.}
The preceding two local ratio-chart arguments have a stronger precise
scope that is useful when another event is restored later. Assume the
exact native stationary coefficient formulas (NC1)--(NC4),
$x>0$ infinitesimal, (RO1), $0<r,h<1$, and
\[
              B_x=\beta+\tfrac12\ge b_*>0
              \quad\hbox{for some real }b_* .                 \tag{RO15}
\]
Original source physicality and $\beta\ge0$ are not required for
this local assertion. Then the full C1-form parent is tensor SOS at
all independent $u,v$ both infinitesimal, including their literal
zero axes. In particular $-1/2<\beta<0$ is allowed whenever
$\beta+1/2$ is a positive bounded unit.

Indeed $P_d\sim1$ before any $h$ estimate, since all its source
arguments are bounded. The exact formula
$B_x\theta/x=K\rho(h+\zeta)/(2P_d)$ has bounded right side,
and (RO15) therefore implies $\theta/x$ bounded. All of (RO2)--(RO14)
follows, with $1/\sqrt{B_x}\le1/\sqrt{b_*}$. In particular
$m,\alpha>0$, $C,A>0$, and $D\sim1/(2j_0)>0$ on this
domain. No remaining use of $\beta\ge0$ occurs in the local
scalar, first-normal or hard-block proof: its complete contribution
is combined with the half-core as (RO12). At $u=0$ the signed
core contributes $D_v(X^2+Y^2)/2$, so its $X^2$ part plus
the beta term is $B_xD_vX^2$. At $v=0$ the same combination
is $B_x(X+uZ)^2$ and the remaining half-core is
$(uW-Y)^2/2$. Thus both axes are also exact SOS under (RO15).
This proves the local and literal-axis assertions under (RO15).



% Source formulas: p48da253a150f
\subsection{The positive $j$ opening}

\label{component:p48da253a150f}

\label{p48da253a150f:sec:r13-c1-j-opening}

\begin{theorem}
Use the original physical native C1 source with $d=1$, $x=k^2>0$
infinitesimal, and
\[
 \eta,\ \zeta,\ r,\ 1-r\quad\hbox{positive bounded units},\qquad
 j=\eta(1+xh)-1+r\quad\hbox{positive infinitesimal}.       \tag{JO1}
\]
Then the complete original parent is tensor SOS for every independent
projective test and every state. Literal $\beta=0$ is retained.
There is no assumption that $j_0=\eta+r-1$ is positive: $j_0=0$
and a slightly negative $j_0$ are included. Simultaneous infinitesimal
$r$ and $j$ are outside this theorem.
\end{theorem}

All scalar and Gram identities in this section use the normalized
parent $F/4$; multiplication of every square root by $2$ restores $F$.

The divergence $D\asymp1/j$ must be centered by its full event,
including its testing-angle dependence. A normal centered only by its
quadratic limit would lose control when $j\ll x$. The following exact
frame removes that problem before any residue is taken.

\paragraph{Source bounds and contact-exact coordinates.}
Write $\theta=1-h=x\pi$, where $\pi>0$ is bounded by (SB1),
$B_x=\beta+1/2\ge1/2$, and $y=x\eta\zeta$. Put
\[
 R_0=1+j-\eta=r+\eta xh,\qquad
 g(U)=R_0-\eta hU,\qquad U=u^2,\ V=v^2.
                                                               \tag{JO2}
\]
Thus $g(x)=r$ exactly. The native source formulas, with
$r=1+j-\eta(1+xh)$, imply
\[
 \begin{gathered}
 C,A=O(1),\quad Dj=\tfrac12+O(x),\quad
 m/x\sim1/(2\zeta),\quad
 \alpha/x^2\sim\zeta(1+\zeta)/2,\\
 B_x\theta/x=K(\eta+r)(h+\zeta)/(2P_d),\qquad
 \theta/\sqrt\alpha,\quad B_x\theta^2/\alpha=O(1).
 \end{gathered}                                                \tag{JO3}
\]
In particular $D>0$, so its square root and inverse are pure source
and legitimate. Set
\[
 \begin{gathered}
 p=uW,\quad q=vZ,\quad E_t=\eta p-q,\qquad
 N_d=\sqrt D\,(J_t+jp),\\
 J_t=uW+uvX-Y-vZ,\qquad
 N_x=\sqrt{B_x}(X+huZ),\quad
 n=N_d/\sqrt D,\quad \xi=N_x/\sqrt{B_x}.
 \end{gathered}                                                \tag{JO4}
\]
The exact inverse state formulas are polynomial and homogeneous:
\[
 Y_0=(1+j)p-q-n,\qquad
 Y=Y_0-hUq+uv\xi,\qquad X=-huZ+\xi,\qquad
 J_t=n-jp.                                             \tag{JO5}
\]
All three normals vanish at the actual source contact. At their zero
graph, $q=\eta p$ and $Y=g(U)p$, so the original native contact
graph is recovered at $U=x$ without changing any source contact jet.

For later use introduce positive reference weights
\[
 \widetilde C={1+j\over2(\eta-j)},\quad
 \widetilde A={\eta\over2R_0},\quad
 \widetilde D={1\over2j}.
\]
Since $R_0,1-R_0,\eta$ are units, exact source numerator division gives
\[
 C-\widetilde C=O(x),\qquad A-\widetilde A=O(x),\qquad
 Dj-\tfrac12=O(x).                                    \tag{JO6}
\]
The last assertion is deliberately about $Dj$, not $D$. The canceled
denominators of (JO3),(JO6), at $x=j=0$, contain only the unit
factors $\eta,\zeta,1-\eta$ and constants. No comparison of $j/x$
has been made.

\paragraph{The complete bounded hard block.}
The reference quadratic base has the exact identity
\[
 \begin{aligned}
 &\widetilde C(p-Y_0)^2+\widetilde A Y_0^2+
 \widetilde D(p-Y_0-q)^2+\tfrac12(p-Y_0)^2-pq\\
 &=\widetilde\gamma E_t^2+
 \widetilde H\{(1-\widetilde b)E_t-n\}^2,               \\
 \widetilde H&={\eta(j+R_0)\over2R_0(1-R_0)j},\quad
 \widetilde b={R_0(1-R_0)\over\eta(j+R_0)},\quad
 \widetilde\gamma={1+j\over2\eta(j+R_0)} .
 \end{aligned} \tag{JO7}
\]
No $p$ term survives. Furthermore
\[
 1-\widetilde b={j(1+j)\over\eta(j+R_0)},\qquad
 \widetilde H/D\sim1.                                 \tag{JO8}
\]
On $(E_t,N_d)$ this is a bounded hard Gram with diagonal residue
$(1/(2\eta(1-\eta)),1)$ and vanishing off-diagonal residue.
The actual-base error has coefficients $O(x)$ in $(p,E_t,N_d)$.
For its apparently singular part use
\[
 (D-\widetilde D)(n-jp)^2,\quad
 {D-\widetilde D\over D}=O(x),\quad
 {(D-\widetilde D)j\over\sqrt D}=O(x\sqrt j),\quad
 (D-\widetilde D)j^2=O(xj).                           \tag{JO9}
\]
These are consequences of (JO3),(JO6) and hold for every relative
source rate. All other weights in the base are bounded.

For bounded $t=v/u$ at $u,v>0$ both infinitesimal, put $Z=1$.
Then $p=(v+E_t)/\eta$ in (JO5). The higher-angle terms have
bounded infinitesimal hard coefficients, and the $m$ hard coefficient
$m t^2$ is infinitesimal. Combining beta with its half-core part gives
the third hard pivot on $N_x$. Thus the full hard residue is
\[
                 \operatorname{diag}(1/(2\eta(1-\eta)),1,1).
                                                               \tag{JO10}
\]
The exact higher-angle identity given below also proves these hard
bounds directly; no state normal is suppressed here.

\paragraph{An exact scalar three-circuit with bounded polynomial coefficients.}
Use $s=\sqrt{\eta\zeta}$ and the same separated roots
\[
 A_t=uv,\quad B_t={kv^2\over su},\quad C_t=xs,\quad
 G=kv,\quad F_v=V-y.
\]
On the zero graph of (JO4), the exact scalar is
\[
 \mathcal S=c_FF_v^2+f_A(U,V)A_t^2+f_BB_t^2+
                       f_C(U)C_t^2-f_GG^2,            \tag{JO11}
\]
where, with $g=g(U)$,
\[
 \begin{aligned}
 c_F&={2mR_0+C(1-R_0)^2+(A+\alpha+\tfrac12)R_0^2\over\eta^2},\\
 q_U&=\alpha+(m+B_x)\theta^2,\\
 q_V&={C(1-R_0)^2+AR_0^2+(D+\tfrac12)j^2\over\eta^2}
       -{2m\theta\over\eta}-{(2\alpha+1)R_0\over\eta}-\tfrac12,\\
 f_B&={m\zeta\over\eta x},\quad
 f_C={\alpha-c_Fy^2+q_UU\over\eta\zeta x^2},\quad
 f_G=-{q_V+2c_Fy\over x},                              \\
 f_A&={mV(g^2-2\eta h)\over\eta^2}-{2m\theta g\over\eta}
       +C(1+V)\{2h(1-R_0)/\eta+h^2U\}\\
 &\quad+\{A(1+V)+(\alpha+\tfrac12)V\}
                   \{g^2/\eta^2-2R_0h/\eta+h^2U\}\\
 &\quad-(2\alpha+1)(g/\eta-h)+B_x\theta^2-\tfrac12 .
 \end{aligned} \tag{JO12}
\]
In particular $f_A$ is an explicit bounded polynomial of degree at
most two in $U$ and one in $V$; the divergent $D$ is absent.
The only $D$ factor elsewhere is $Dj^2=O(j)$. These formulas
follow by substituting $W=v/(\eta u)$, $X=-hu$, $Y=g(U)v/\eta$
in the entire parent; 

The source bounds and (JO12) give
\[
 c_F\sim1/(2\eta^2),\qquad
 f_A,f_B,f_C\sim\lambda=1/(2\eta).
\]
Exact scalar contact gives
$f_G=f_A(x,y)+f_B+f_C(x)$, proving that $f_G$ is bounded
and $f_G\sim3\lambda$ without estimating a singular quotient
$q_V/x$. Every nonconstant coefficient in $f_A,f_C$ is bounded;
for $q_U/x^2$ use (JO3).

For clarity the required finite Taylor step works for these higher
polynomial coefficients too. Let $H(u,t)=f_A(u^2,u^2t^2)+f_B+
f_C(u^2)-f_G$. It is a finite polynomial in $u,t$ with bounded
source arrays and $H(k,s)=0$. Write $\Delta_u=u-k$,
$\Delta_t=t-s$. A finite identity, obtained by monic polynomial
division, is
\[
 H=H_u(k,s)\Delta_u+H_t(k,s)\Delta_t+
       H_{uu}\Delta_u^2+H_{ut}\Delta_u\Delta_t+
       H_{tt}\Delta_t^2,                              \tag{JO13}
\]
where one may take
\[
 \begin{aligned}
 H_{uu}&={H(u,t)-H(k,t)-(u-k)H_u(k,t)\over(u-k)^2},\\
 H_{ut}&={H_u(k,t)-H_u(k,s)\over t-s},\\
 H_{tt}&={H(k,t)-H(k,s)-(t-s)H_t(k,s)\over(t-s)^2}.
 \end{aligned}
\]
These displayed divisions mean their finite polynomial quotients,
not mixed inverses. The first two arrays are bounded and $H_{tt}$
is infinitesimal: every $t$-dependent monomial of $H(k,t)$
contains $k^2$. The coefficient first differences similarly have a
bounded $u$ quotient and infinitesimal $t$ quotient. Exact first
angular stationarity supplies the two unchanged contact-jet balances
of the three-circuit lemma. The chain rule is legitimate here because
$g(x)=r$ and the original first state derivatives vanish at contact.
Consequently (JO11) has the strict finite three-circuit Gram and the
strict $F_v$ pivot, including the complete soft defect.

\paragraph{Every complete first normal.}
The following bounded polynomials are explicit:
\[
 \begin{aligned}
 P_Y/2={}&m\{V(1+gU)/\eta-\theta U\}
       -C D_v(1-g)/\eta+A D_uD_vg/\eta\\
       &+(D+\tfrac12)j/\eta+
                  (\alpha+\tfrac12)D_u(gV/\eta-1),\\
 P_E={}&{2mV\over\eta^2}\{g+1+j+g(1+j)U\}
       -{2m\theta\over\eta}\{1+(1+j)U\}\\
       &-{2CD_vj(1-g)\over\eta^2}
        +{2AD_uD_vg(1+j)\over\eta^2}
        +{2(D+\tfrac12)j^2\over\eta^2}\\
       &+{2(\alpha+\tfrac12)D_u(1+j)\over\eta}
                                      (gV/\eta-1),\\
 Z_X={}&B_x\theta+P_Y/2-(D+\tfrac12)j/\eta .
 \end{aligned}                                                \tag{JO14}
\]
Their only $D$ factors are $Dj$ and $Dj^2$; $D$
cancels entirely from $Z_X$. Direct full-parent differentiation gives
\[
 \begin{aligned}
 \mathcal L_d&=-vP_Y/\sqrt D,\\
 \mathcal L_E&=2mvt^2/\eta^2+vP_E,\\
 \mathcal L_x&={1\over\sqrt{B_x}}
 \{2u\theta(m+B_x)-2mut^2(1+gU)/\eta+2uVZ_X\}.
 \end{aligned}                                                \tag{JO15}
\]
No normal is missing. For $P=P_Y,P_E$ use the exact finite
division $P=P_{00}+UP_U+VP_V$. Contact implies respectively
\[
 \begin{aligned}
 P_{Y,00}/x&=-P_{Y,U}(x,y)-\eta\zeta P_{Y,V}(y),\\
 P_{E,00}/x&=-2(m/x)\eta\zeta/\eta^2
                  -P_{E,U}(x,y)-\eta\zeta P_{E,V}(y).
 \end{aligned}                                                \tag{JO16}
\]
Thus these coquotients are bounded by finite formulas with no $j$
inverse. Allocate $vP$ by
$vP=(P_{00}/k)G+u(P_U+t^2P_V)A_t$; allocate the additional
$2mvt^2/\eta^2$ on $B_t$ with coefficient
$2mst/(\eta^2k)$. The $X$ line allocates exactly on $(C_t,B_t,A_t)$
with coefficients
\[
 {2u\theta(m+B_x)\over\sqrt{B_x}xs},\quad
 -{2ms(1+gU)\over\eta k\sqrt{B_x}},\quad
 {2vZ_X\over\sqrt{B_x}}.                               \tag{JO17}
\]
Every displayed coefficient is infinitesimal; for the first use
$\theta(m+B_x)/x=O(1)$, and for the others use $m/x=O(1)$.
The $N_d$ row has the additional bounded factor $1/\sqrt D$.
Each sum of coefficient functions vanishes at contact. Its finite
$u$ difference is bounded and its $t$ difference at $u=k$ is
infinitesimal, because every $t$ dependence carries $u$ or $k$.
The two unchanged conductors therefore place every whole normal on the
strict circuit roots with infinitesimal coefficients. Joining (JO10)
proves the full parent on the bounded-ratio chart.

\paragraph{The reciprocal testing ratio: a complete polynomial remainder.}
For $v/u$ unbounded, return to homogeneous state coordinates and use
\[
 E_t,\ N_d,\ N_x,\ P_m=\sqrt m\,vW,\ Q_\alpha=\sqrt\alpha Z,
                  \ R_t=u^2W,\ S_t=uvW.              \tag{JO18}
\]
The entire base error from (JO9) has zero Gram residue on this list:
its $p^2$ coefficient is $O(x)$ and
$\sqrt x\,p=\sqrt{x/m}(u/v)P_m$. Every coefficient is a product
of a bounded source factor and a bounded testing factor; $u/v$ is
pure testing.

Here is the full higher-angle identity used to control the other terms.
Write $T_s=-hUq+uv\xi$, so $Y=Y_0+T_s$, with $Y_0$ in
(JO5). After removing the actual quadratic base and the three entire
positive terms $mR_m^2$, $\alpha D_u(vY-Z)^2$ and
$B_xD_v(X+uZ)^2$, the exact remainder is
\[
 \begin{aligned}
 \mathcal R={}&C\{V(p-Y_0)^2-2D_v(p-Y_0)T_s+D_vT_s^2\}\\
 &+A\{(D_uD_v-1)Y_0^2+2D_uD_vY_0T_s+D_uD_vT_s^2\}\\
 &+\tfrac12D_uVY^2-UqY+(h-\tfrac12)Uq^2-uvq\xi .
 \end{aligned}                                                \tag{JO19}
\]
This is a quadratic in the full state, with bounded coefficients
$C,A,h$ and no $D$. It is checked as a whole polynomial. Insert
$q=\eta p-E_t$, $n=N_d/\sqrt D$, $\xi=N_x/\sqrt{B_x}$.
Every term contains $U,V$ or $uv$. Cross terms with $p$ have
infinitesimal coefficients on (JO18), using $Up=uR_t$,
$Vp=vS_t$ and $uvp=uS_t$. Hard quadratic terms have zero
residue. Terms $U^aV^bp^2$ are allocated on $R_t^2$ when
$a>0$, and on $S_t^2$ when $a=0,b>0$. Their two only nonzero
residues, verified from (JO19), are
\[
                         \kappa_R=\eta/2,\qquad
                         \kappa_S=1/2.               \tag{JO20}
\]
Every remaining such coefficient contains another $U$ or $V$.
These identities explicitly exhaust the scalar and every normal
coupling of (JO19).

The three retained positive terms have leading roots
$P_m,-Q_\alpha,N_x$. Their remaining root coefficients are
infinitesimal because $m,\alpha$ are infinitesimal and
$\theta/\sqrt\alpha$, $B_x\theta^2/\alpha$ are bounded:
\[
 \sqrt m\,R_m=P_m+\sqrt m\,uvY-\sqrt m(\xi+\theta uZ),
 \quad
 \sqrt{B_xD_v}(X+uZ)=
 \sqrt{D_v}\{N_x+\sqrt{B_x\theta^2/\alpha}\,uQ_\alpha\}.
\]
The alpha term is $\sqrt{D_u}(\sqrt\alpha\,vY-Q_\alpha)$.
Thus the complete parent has a bounded Gram on (JO18) with strictly
positive diagonal residue
\[
             (1/(2\eta(1-\eta)),1,1,1,1,\eta/2,1/2).
                                                               \tag{JO21}
\]
This proves the reciprocal-ratio chart and completes the local proof.

\paragraph{All-tests join and receipt.}
The common small-$L$, bounded-$\eta$ remote theorem applies to
(JO1), including the exact original source openings and literal beta
zero. It covers the complement of both-infinitesimal testing angles.
The original exact sign gauges, literal axes, homogeneous state
extension and real/imaginary decomposition apply as in IC1 and RO.
All coefficients and radicals used above are pure-field quantities;
no mixed inverse or comparison of source and testing rates occurs.

 They verify
the original full-parent substitution, every complete first normal,
the scalar identity, all newly used source denominator factors,
the reference hard cancellation, the whole reciprocal higher-angle
identity and both strict soft pivots. The finite polynomial bounds
and full-rates circuit allocation are the proof above, rather than
an inference from those checks. This admits precisely (JO1) and its
exact symmetry images. It does not remove the simultaneous
$r,j$ opening or any unbounded-$\zeta$ or infinitesimal-$\eta$ source.



% Source formulas: pbf6de1d7328a
\subsection{The simultaneous positive $r,j$ opening}

\label{component:pbf6de1d7328a}

\label{pbf6de1d7328a:sec:r13-c1-simultaneous-rj}

\begin{theorem}
Use the original physical native C1 source, with $d=1$ and $x=k^2>0$
infinitesimal. Suppose $\eta,\zeta$ are positive bounded units and
$r$ and the actual $j=\eta(1+xh)-1+r$ are both positive
infinitesimals. Then the complete original parent is tensor SOS for
every independent projective test and every state. Literal $\beta=0$
is retained. No relative rate among $x,r,j$ or between source and
testing parameters is imposed.
\end{theorem}

All scalar and Gram identities in this section use the normalized
parent $F/4$; multiplication of every square root by $2$ restores $F$.

The essential change is to combine the whole positive $A,D$ events
before choosing the scalar graph. The weighted $uvX$ term is retained
inside a hard root. Solving for $p=uW$, rather than for $q=vZ$,
then gives a polynomial inverse with a pure-source unit denominator.

\paragraph{Source bounds and the harmonic frame.}
Write $U=u^2$, $V=v^2$, $D_u=1+U$, $D_v=1+V$,
$\theta=1-h=x\pi$, $B_x=\beta+1/2$, and
\[
 T=1+xh,\qquad e=\eta T=1-r+j,\qquad y=x\eta\zeta .
 \tag{HO1}
\]
The bounded-$\eta$ source bound gives $\pi=O(1)$.
Here $e,\eta,h,1-r\sim1$. Exact native numerator cancellation gives
\[
 \begin{gathered}
 C-C_* =O(x),\qquad Ar-e/2=O(x),\qquad Dj-1/2=O(x),\\
 C_*={1+j\over2(1-r)},\qquad A_*={e\over2r},\qquad D_*={1\over2j},\\
 m/x\sim1/(2\zeta),\qquad
 \alpha/x^2\sim\zeta(1+\zeta)/2,\\
 B_x\theta/x=O(1),\qquad
 q_U/x^2=O(1),\quad q_U=\alpha+(m+B_x)\theta^2 .
 \end{gathered}
 \tag{HO2}
\]
 This proves boundedness at arbitrary simultaneous
source rates, rather than just nonvanishing at a selected boundary.
In particular $A,D>0$ and both are unbounded.
Also $\theta/\sqrt\alpha$ and $B_x\theta^2/\alpha$ are bounded,
including $\beta=0$.

Define pure-source positive quantities
\[
 \mathsf H=A+D,\qquad \rho={D\over\mathsf H},\qquad
 \mathsf K={AD\over\mathsf H},\qquad
 \delta={Ar+Dj\over\mathsf H}.
 \tag{HO3}
\]
Both $\mathsf H$ and $\mathsf K$ are unbounded,
$0<\rho<1$, and
$\mathsf K^{-1}=A^{-1}+D^{-1}=O(r+j)$.
Put $p=uW$, $q=vZ$, $J_0=p+uvX-q$, and use the full-event roots
\[
 \begin{aligned}
 E&=\sqrt{\mathsf K}\{J_0-(r-j)p\},\\
 N&=\sqrt{\mathsf H}\{Y-\rho J_0-\delta p\},\\
 N_x&=\sqrt{B_x}(X+huZ),\qquad
 \ell=E/\sqrt{\mathsf K},\quad n=N/\sqrt{\mathsf H},
 \quad \xi=N_x/\sqrt{B_x}.
 \end{aligned}
 \tag{HO4}
\]
All three vanish at the actual source contact. Their exact inverse is
\[
 \begin{aligned}
 p&={(1+hU)q-uv\xi+\ell\over e},&
 Y&=rp+\rho\ell+n,&X&=-huZ+\xi,\\
 J_t&=-jp+(1-\rho)\ell-n,&J_t&=p+uvX-Y-q .
 \end{aligned}
 \tag{HO5}
\]
The denominator $e$ is pure source. In particular no inverse of
$1+hU$ occurs, even if the testing and source rates are unrelated.
The useful exact weighted identities are
\[
 {A\rho^2\over\mathsf K}=\rho,\qquad
 {A\over\mathsf H}=1-\rho,\qquad
 {D(1-\rho)^2\over\mathsf K}=1-\rho,\qquad
 {D\over\mathsf H}=\rho .
 \tag{HO6}
\]

\paragraph{The complete hard block and its reference cancellation.}
The $A,D$ base has the exact identity
\[
 \begin{split}
 AY^2+DJ_t^2={}&N^2+E^2
 +2p\left\{{Ar+Dj\over\sqrt{\mathsf H}}N
                 +\sqrt{\mathsf K}(r-j)E\right\}\\
 &+(Ar^2+Dj^2)p^2 .
 \end{split}
 \tag{HO7}
\]
Both displayed cross coefficients and the last coefficient are
infinitesimal: the squared $E$ cross coefficient is $O(r+j)$
and the squared $N$ cross coefficient is $O(\mathsf H^{-1})$.
For bounded testing ratio $t=v/u$, (HO5) has bounded coefficients
on the hard roots, and the only potentially large $m$ coefficient is
$mt^2/\mathsf K$, which is infinitesimal.
The higher $A$ term is controlled by (HO6); for example
\[
 {Ar^2\over\mathsf K}=r^2+{Ar\over Dj}rj=O((r+j)^2),
 \qquad {Ar\rho\over\mathsf K}=r .
 \tag{HO8}
\]
Consequently its hard coefficients carry an additional $U$ or $V$.
The bounded $C$ and half-core terms have infinitesimal coefficients
on $E,N$, and combining the beta term with its half-core part gives
the third hard pivot. The full hard residue on $(E,N,N_x)$ is
$\operatorname{diag}(1,1,1)$.

A stronger exact base statement is needed for unbounded $t$.
Set $\widetilde q=q-uvX=ep-\ell$, so that
$J_t=p-Y-\widetilde q$. The reference quadratic form
\[
 \mathcal B_*=
 C_*(p-Y)^2+A_*Y^2+D_*(p-Y-\widetilde q)^2
                    +\tfrac12(p-Y)^2-p\widetilde q
 \tag{HO9}
\]
has no soft $p$ coordinate after substituting
$Y=rp+\rho\ell+n$ and $\widetilde q=ep-\ell$:
\[
 \mathcal B_*=
 (A_*+C_*+\tfrac12)(\rho\ell+n)^2
                 +D_*\{(1-\rho)\ell-n\}^2 .
 \tag{HO10}
\]
This is an exact contact cancellation, checked before taking residues.
Since $A_*/A=1+O(x)$ and $D_*/D=1+O(x)$, (HO6) shows that
its bounded hard Gram on $(E,N)$ has residue $\operatorname{diag}(1,1)$.
Replacing the reference weights by the actual $C,A,D$ creates a
quadratic error with coefficients $O(x)$ in $(p,E,N)$.
For the singular weights this follows from
$(A-A_*)Y^2=O(x)\,AY^2$ and
$(D-D_*)J_t^2=O(x)\,DJ_t^2$, using (HO6),(HO7);
neither an $x/r$ nor an $x/j$ estimate is used.

\paragraph{The exact bounded-ratio scalar circuit.}
At $E=N=N_x=0$ and $Z=1$, let
\[
 w(U)={1+hU\over e},\qquad
 p=w(U)v,\quad Y=rw(U)v,\quad X=-hu .
 \tag{HO11}
\]
Since $w(x)=1/\eta$, this graph agrees exactly with the native
source contact. The source contact value and all first state and
angular jets are therefore common by the chain rule.

For compact scalar coefficients set
$a=C(1-r)^2$, $b=Ar^2$, $d=(D+1/2)j^2$,
$c=(\alpha+1/2)r^2$, and
\[
 L_A(U)=1+2h+(2h+h^2)U+h^2U^2 .
 \]
The complete scalar is exactly
\[
 \mathcal S=c_F(V-y)^2+f_A(U,V)A_t^2+f_BB_t^2+
                         f_C(U)C_t^2-f_GG^2 ,
 \tag{HO12}
\]
on the separated roots
$A_t=uv$, $B_t=kv^2/(su)$, $C_t=xs$, $G=kv$,
where $s=\sqrt{\eta\zeta}$. The coefficients are
\[
 \begin{aligned}
 c_F&={2m(h+r)+a+b+c\over e^2},\\
 q_V&={a+b+d\over e^2}-{2m\theta\over e}
                       -{(2\alpha+1)r\over e}-\tfrac12,\\
 f_B&={m\eta\zeta\over e^2x},\qquad
 f_C={\alpha-c_Fy^2+q_UU\over\eta\zeta x^2},\qquad
 f_G=-{q_V+2c_Fy\over x},\\
 f_A&={mV\over e^2}
 \{(h+r)^2+2hr+2hr(h+r)U+h^2r^2U^2\}
       -{2m\theta\over e}(h+r+hrU)\\
 &\quad+{a\over e^2}(2h+h^2U)(1+V)
       +{b\over e^2}L_A(U)(1+V)
       +{d\over e^2}(2h+h^2U)+{c\over e^2}L_A(U)V\\
 &\quad-{(2\alpha+1)r\over e}(1+h+hU)
                              +B_x\theta^2-\tfrac12 .
 \end{aligned}
 \tag{HO13}
\]
The only opening weights here are $Ar^2$ and $Dj^2$, both
infinitesimal. Thus $f_A$ is a bounded polynomial of degree at most
two in $U$ and one in $V$, and all nonconstant coefficients of
$f_C$ are bounded by (HO2). Direct source reduction gives
\[
 c_F\sim\tfrac12,\qquad f_A,f_B,f_C\sim\tfrac12 .
 \tag{HO14}
\]
Exact scalar contact gives
$f_G=f_A(x,y)+f_B+f_C(x)\sim3/2$, without separately estimating
$q_V/x$.

Apply the finite three-circuit lemma with $\lambda=1/2$.
The polynomial Taylor arrays are the monic quotients in (JO13),
now applied to
$f_A(u^2,u^2t^2)+f_B+f_C(u^2)-f_G$.
Their $u$ and mixed arrays are bounded; the $t,t$ array and the
coefficient $t$ first differences are infinitesimal, since every
$t$ dependence at $u=k$ contains $k^2$.
The two exact angular contact jets give the two circuit balances.
This supplies the strict three-circuit Gram and the strict
$V-y$ pivot for the entire scalar, including its negative soft term.

\paragraph{All complete first normals.}
The following formulas come from differentiating the entire original
parent after (HO5), with $Z=1$. Put
$a_E=r/e+\rho$, $b_E=(1-r)/e-\rho$,
$c_E=1-\rho-j/e$. The bounded polynomials are
\[
 \begin{aligned}
 P_N/2={}&m\{wV(1+rU)-\theta U\}
 -CD_v(1-r)w+AD_uD_vrw\\
 &+(D+\tfrac12)jw
             +(\alpha+\tfrac12)D_u(rwV-1),\\
 P_E/2={}&{mV\over e^2}(h+r+hrU)-{m\theta\over e}
 +ma_E\{wV(1+rU)-\theta U\}\\
 &+CD_v(1-r)wb_E+AD_uD_vrwa_E
 -(D+\tfrac12)jwc_E
            +(\alpha+\tfrac12)D_u(rwV-1)a_E,\\
 Z_X={}&{m\theta(1+rU)\over e}+B_x\theta\\
 &-{CD_v(1-r)^2w+AD_uD_vr^2w+(D+\tfrac12)j^2w
                 +(\alpha+\tfrac12)D_ur(rwV-1)\over e}.
 \end{aligned}
 \tag{HO15}
\]
Their only opening-weight factors are $Ar,Dj,Ar^2,Dj^2$.
Every coefficient is bounded uniformly in the source rates and
$0<\rho<1$. The three complete first normals are
\[
 \begin{aligned}
 \mathcal L_N&={vP_N\over\sqrt{\mathsf H}},\\
 \mathcal L_E&={2mvt^2/e^2+vP_E\over\sqrt{\mathsf K}},\\
 \mathcal L_x&={1\over\sqrt{B_x}}\left\{
 2u\theta(m+B_x)-2mwut^2(1+rU)
           \left(1+{V(1+rU)\over e}\right)+2uVZ_X\right\}.
 \end{aligned}
 \tag{HO16}
\]
For either $P=P_N,P_E$, use its exact finite monic decomposition
$P=P_{00}+UP_U+VP_V$.
State contact gives
\[
 \begin{aligned}
 P_{N,00}/x&=-P_{N,U}(x,y)-\eta\zeta P_{N,V}(y),\\
 P_{E,00}/x&=-2(m/x)\eta\zeta/e^2
                         -P_{E,U}(x,y)-\eta\zeta P_{E,V}(y).
 \end{aligned}
 \tag{HO17}
\]
Thus these coquotients are bounded without any new opening inverse.
Allocate $vP=(P_{00}/k)G+u(P_U+t^2P_V)A_t$.
The extra $E$ term allocates on $B_t$ with coefficient
$2mst/(e^2k\sqrt{\mathsf K})$.
The three $X$ terms allocate respectively on $C_t,B_t,A_t$,
with coefficients
\[
 {2u\theta(m+B_x)\over xs\sqrt{B_x}},\qquad
 {-2msw(1+rU)\{1+V(1+rU)/e\}\over k\sqrt{B_x}},\qquad
 {2vZ_X\over\sqrt{B_x}} .
 \tag{HO18}
\]
Every displayed allocation coefficient is infinitesimal.
The coefficient sums vanish at contact by the exact state jets;
their finite $u$ first differences are bounded and their $t$
first differences at $u=k$ are infinitesimal.
This follows directly from the bounded polynomial arrays in
(HO15): every $t$ dependence carries $u$ or $k$.
The two unchanged circuit conductors therefore place every whole
normal on the strict circuit roots with infinitesimal coefficients.
Together with the complete hard block, this proves all bounded $t$,
including infinitesimal $t$, without a missing state normal.

\paragraph{The reciprocal ratio: retain both soft directions.}
If $v/u$ is unbounded, work homogeneously and use seven roots
\[
 E,\quad N,\quad N_x,\quad
 P_m=\sqrt m\,vW,\quad Q_\alpha=\sqrt\alpha Z,\quad
 R_t=uvZ,\quad S_t=uvW .
 \tag{HO19}
\]
The base error following (HO10) has zero Gram residue on this list:
use the exact separated identity
$\sqrt x\,p=\sqrt{x/m}(u/v)P_m$.
This absorbs every $O(x)$ coefficient in $(p,E,N)$, since $x/m$
is a source unit and $u/v$ is a pure-testing infinitesimal.

After removing the actual quadratic base (HO9), with actual $C,A,D$,
and the three entire positive terms
$mM_t^2$, $\alpha D_u(vY-Z)^2$,
$B_xD_v(X+uZ)^2$, the full remaining quadratic is exactly
\[
 \begin{split}
 \mathcal R={}&CV(p-Y)^2+A(U+V+UV)Y^2
                      +\tfrac12D_uVY^2-\theta UqY\\
 &+(h-\tfrac12)Uq^2+\tfrac12h^2U^2q^2
       -uv\xi\{Y+(1+hU)q\}+\tfrac12UV\xi^2 .
 \end{split}
 \tag{HO20}
\]
This identity retains the whole $D$ event in the quadratic base,
so no coefficient containing a bare divergent $D$ remains.
It is checked as one complete homogeneous polynomial.

Here is a finite allocation of every term. The exact identity
\[
 up={(1+hU)R_t-u^2v\xi+uE/\sqrt{\mathsf K}\over e}
 \tag{HO21}
\]
has bounded coefficients and uses no mixed inverse.
The $C$ term has leading root
$\sqrt C(1-r)S_t$; its other coefficients are
$-\sqrt C\,\rho v/\sqrt{\mathsf K}$ and
$-\sqrt C\,v/\sqrt{\mathsf H}$ on $E,N$, both infinitesimal.
All three $A$ higher terms have zero residue. For example,
\[
 \sqrt A\,uY=\sqrt A\,r(up)+\sqrt\rho\,uE+
                               \sqrt{1-\rho}\,uN,\qquad
 \sqrt A\,vY=\sqrt A\,rS_t+\sqrt\rho\,vE+
                               \sqrt{1-\rho}\,vN .
 \tag{HO22}
\]
Here $\sqrt A\,r=O(\sqrt r)$ and (HO8),(HO21)
control every coefficient; the $uvY$ term is $v$ times the first.
The half-$VY^2$ term is smaller by the same identities.
The $-\theta UqY$ term is $-\theta R_t(uY)$ and has zero residue.
The next two terms are $(h-1/2)R_t^2$ and
$h^2U R_t^2/2$.
Finally, the cross term is
$-\xi v(uY)-\xi v(1+hU)R_t$ and all its coefficients,
as well as those of $UV\xi^2/2$, are infinitesimal.
Thus this allocation exhausts the complete quadratic (HO20);
its only nonzero soft residues are $1/2$ on each of $R_t,S_t$.

The three retained positive terms have leading roots
$P_m,-Q_\alpha,N_x$:
\[
 \begin{aligned}
 \sqrt m\,M_t&=P_m+\sqrt m\,uvY-\sqrt m\,\xi
                  -\sqrt{m/\alpha}\,\theta uQ_\alpha,\\
 \sqrt{\alpha D_u}(vY-Z)
             &=\sqrt{D_u}(\sqrt\alpha\,vY-Q_\alpha),\\
 \sqrt{B_xD_v}(X+uZ)
             &=\sqrt{D_v}\{N_x+
                      \sqrt{B_x\theta^2/\alpha}\,uQ_\alpha\}.
 \end{aligned}
 \tag{HO23}
\]
Every other coefficient here is infinitesimal by (HO2),(HO21),
because $\theta\sqrt{m/\alpha}=O(\sqrt x)$.
The complete bounded Gram on (HO19) therefore has strictly positive
diagonal residue
\[
             \operatorname{diag}(1,1,1,1,1,1/2,1/2).
 \tag{HO24}
\]
This proves the reciprocal chart for the full parent and every state.

\paragraph{All-tests join and receipt.}
The common small-$L$, bounded-$\eta$ remote theorem covers the
complement of both-infinitesimal testing angles on this original
physical source domain. Its exact source openings and literal beta
zero remain available. The original exact sign gauges, literal axes,
homogeneous state extension and real/imaginary decomposition join
as in IC1, RO and JO. All roots, radicals and inverses above are
pure-field or separated; every bound is a product of a source bound
and a testing bound.

Those checks support the finite bounds and circuit allocation proved
above; they are not being used as a substitute for them.
This theorem removes precisely the simultaneous positive
infinitesimal $r,j$ opening with unit $\eta,\zeta$.

\begin{corollary}[The entire unit-$\eta,\zeta$ small-$L$ region]
Within the original physical native small-$L$ C1.4 source, suppose
$x>0$ is infinitesimal and $\eta,\zeta$ are positive bounded units.
Then the full parent is tensor SOS at every independent projective
test and every state, with all original source intersections retained.
\end{corollary}
\begin{proof}
The positive actual gap $j=\eta(1+xh)-1+r$ is bounded.
If $j$ is a unit, then $j_0=j-\eta xh$ is a positive unit,
so RO applies, including either $r$ opening. If $j$ is infinitesimal,
then $1-r\sim\eta$ is a unit. A unit $r$ is exactly the JO domain;
an infinitesimal $r$ is exactly the HO domain. These alternatives
exhaust the stated source, require no comparison of $j$ with $x$,
and each supplies its complete local, remote and endpoint join.
\end{proof}



% Source formulas: pa73a17051258
\subsection{The simultaneous small-$\eta$ parent identity}

\label{component:pa73a17051258}

\paragraph{Polynomial identity.}
This is an independently reconstructed exact identity for the complete
original C1 parent, not a certificate admitting C14.3b. It retains the
same stationary source, $\Lambda=n=0$, $\beta\ge0$, all $j$ and $1-r$
openings, independent source/testing fields, and the remaining local
$v/u$-infinitesimal testing chart. No mixed inverse or comparison of
source and testing rates is used below.

Write $\mathcal P=F/4$, initially at $Z=1$. Set
\[
 a=1/\zeta,\quad t=\eta,\quad y=tx/a,\quad U=u^2,\quad V=v^2,
 \quad B=\beta+\tfrac12,\quad D_p=D+\tfrac12,
\]
\[
 \begin{gathered}
 E=tuW-v,\quad Q_t=Y-rv/t,\quad P=E/t-Q_t,\\
 \xi=X+hu,\quad R_D=v(U-x),\quad T_D=R_D-uv\xi/h.
 \end{gathered}
\]
Thus $Q_t=rE/t+\nu$ in the R15 notation. Put
\[
 q=\alpha+(m+B)(1-h)^2,\quad\mu=m-a^2q,
 \quad\lambda=\mu/(tax),\quad H=avW-u-a\xi,
\]
\[
 F_0=V-y,\quad F_b=F_0+vE,\quad F_t=F_0+(t/r)vQ_t,
 \quad T_0=A+\alpha+\tfrac12,\quad T_m=T_0+m,
 \quad e_4=r^2T_m/t^2.
\]
Here $F_t=\widetilde F$ and $F_b=\overline F$. Define pure-source
coefficients
\[
 b=Ar/t-m(1-h)-\alpha-\tfrac12,\qquad
 b'=b+e_4yt/(2r),
\]
\[
 \kappa=2[B(1-h)-D_pj/t-mr/t],\qquad
 \gamma=2b'/h+2D_pj/r-2tB(1-h)/r.
\]
They satisfy $b'=O(x)$, $t\kappa\sim-1$ and $\gamma\sim1$.

\paragraph{Complete parent, with no omitted hard row.}
Let
\[
 L=UV+(ax/t)V^2/U+tx^2/a-3xV,
\]
\[
 \mathscr C=\lambda L+\frac{\mu V}{t^2U}E^2
                +2\left(\frac{\mu v^3}{t^2U}-\lambda xv\right)E.
\]
The exact full-parent identity is
\begin{align*}
 \mathcal P={}&qH^2+\mathscr C
 +CP^2+AQ_t^2+D_p(P-hT_D)^2+2b'Q_tT_D\\
 &+\mathscr R+e_4UF_t^2-e_4yu(H+a\xi)F_t\\
 &+AUQ_t^2+B(1+V)\xi^2
                 -2(\mu/a)H\xi-\mu\xi^2\\
 &+\kappa uF_t\xi+\gamma uvQ_t\xi,                 \tag{JF1}
\end{align*}
where the entire joint radial row is
\[
 \mathscr R=\frac C{t^2}(F_b-rF_t)^2
       +\frac{r^2T_0}{t^2}F_t^2
       +\left(\frac{2mr}{t^2}+e_4x\right)F_bF_t.    \tag{JF2}
\]
In particular all $mVEQ_t/t$ terms are included in JF2; they are not
separate errors. The $mUVQ_t^2$ contribution is included in the complete
$e_4UF_t^2$ row, so it cannot also be retained independently.

A short independent way of certifying JF1 is useful. Without imposing
any source formulas, the difference between the original parent and
its displayed right side is exactly a linear combination of only
\[
                       vE,\ vQ_t,\ u\xi,\ UV,\ V,\ 1.
\]
The displayed expression has zero value and all five first derivatives
at the original contact $U=x,V=y,E=Q_t=\xi=0$.
Its six coefficients therefore vanish by the six original stationary
contact balances. This support calculation and the six contact identities
prove the identity for the complete original parent, including every state normal.

\paragraph{One strict radial pivot, with its exact companion.}
Set $\rho=m+rT_mx/2$. Then JF2 has the pure-source completion
\[
 \mathscr R=\frac C{t^2}
       [F_b-r(1-\rho/C)F_t]^2
       +\frac{r^2}{t^2}[T_0+2\rho-\rho^2/C]F_t^2.  \tag{JF3}
\]
Here $\rho=O(x)$ and $T_0+2\rho-\rho^2/C\sim1/2$;
all denominators are pure source and this remains uniform at both
source openings. Both $F_t/t$ and $F_b/t$ have bounded allocations
on JF3. This does not provide an independent strict $F_0/t$ pivot.
The non-small remaining coefficient is precisely $\gamma uvQ_t\xi$.
In particular a source split making $x/\sqrt t$ bounded does not
remove it, because $\gamma\sim1$.

\paragraph{The extra circuit row participates in the leading cancellation.}
Use $c^2=a/t$, $k^2=x$, $p=ckv/u$, and homogenize now in $Z$:
\[
 J_0=uv-x/c,\quad R_c=kv(tcW-Z),\quad
 R=\sqrt\lambda\,tk(cv-u)W,
 \quad \tau=tcW-Z.
\]
The generic exact circuit completion is
\[
 \mathscr C=\lambda J_0^2Z^2+R^2+
      \underbrace{\frac{2\lambda k}{p}R_c^2}_{\mathcal T^2}
      -\lambda xE^2.
\]
The full-$D$ conductor is
\[
 R_D=uJ_0Z+kR_c-(x/c)(H+a\xi),\qquad
 H+a\xi=avW-uZ.                                  \tag{JF4}
\]
The additional exact polarization is
\[
 E(H+a\xi)=uv\tau^2-t(u-cv)^2WZ.                  \tag{JF5}
\]
Consequently the extra circuit row and the principal cross left after
using JF4 satisfy
\[
 \mathcal T^2+\kappa h(x/c)E(H+a\xi)
 =\frac{x}{c}(2\lambda+\kappa h)uv\tau^2
      -\kappa ht\frac{x}{c}(u-cv)^2WZ.             \tag{JF6}
\]
This cancels the previously non-small $E$--angular Gram coupling.
JF6 is used as an exact identity, rather than as a separate positive
term. In particular, its first coefficient has not been allocated
to a strict positive row. The sign of $WZ$ must be retained. The second summand is a
literal SOS on the pure-testing sign chart $WZ\ge0$ because
$-\kappa>0$, but not on the opposite sign chart.

\paragraph{An exact allocation of the negative small extra-row fraction.}
There is nevertheless a useful conductor:
\[
 x\frac{k}{p}R_c^2
          =U R_c^2-kuJ_0R_c\tau.                   \tag{JF7}
\]
If the testing state is normalized with $Z$ a unit and $W$ bounded,
then $\tau/Z=tc(W/Z)-1$ is bounded. Thus the right side is an
infinitesimal Gram on $\sqrt\lambda R_c$ and
$\sqrt\lambda J_0Z$, with coefficients $U$ and
$ku(\tau/Z)$. Since
$\sqrt\lambda R_c=R+\sqrt\lambda kE$, its additional coefficient
on $E/\sqrt t$ is $\sqrt{\lambda t}\,k=O(k)$.
Therefore a negative $O(x)$ fraction of $\mathcal T^2$ has a bounded
infinitesimal allocation on this state chart. No bound on $p$ is used.
The companion identity
\[
 F_b/t=(W/Z)J_0Z+(x/a)\tau                          \tag{JF8}
\]
may supply a strict $xZ$ radial pivot in a simultaneous state-chart
reassembly, since $\tau/Z\sim-1$ there. All divisions in JF7--JF8
are pure testing. Neither identity is an authorization to assume a
strict $E/\sqrt t$ pivot after the leading source block has become
semidefinite. That final simultaneous allocation is still required.

\paragraph{A direct conductor on the normalized $Z$-unit state chart.}
There is a simpler route on this state chart that does not perform JF4--JF6.
Put $e_E=E/\sqrt t$ and retain the original strict source completion
of the first line of JF1. Since $Q_t=E/t-P$ and $vZ=tuW-E$,
\[
 uvQ_t\xi=\sqrt t\,U(W/Z)e_E\xi
             -u(\xi/Z)e_E^2-uvP\xi.                \tag{JF9}
\]
When the testing state is normalized with $Z$ a unit, $W/Z$ and
$X/Z$ are bounded and $\xi/Z=X/Z+hu$ is bounded. Thus all three
coefficients in JF9 are infinitesimal in the bounded tensor ring.
The last uses the bounded pure-source allocation of $P$ on the
completed source roots. The first line of JF1 has source Schur pivot
$tS\sim1/2$, and $E/\sqrt t$ has a bounded allocation there,
including after the full $D_p$ shift. Consequently JF9 removes the
remaining $\gamma$ coupling on this state chart without producing
a semidefinite source frame.
The accompanying complete state-chart certificate, if used, must still
list the retained positive rows and the strict angular block; JF9
alone does not cover states with normalized $Z$ infinitesimal.

\paragraph{Use of the full identity.}
The complete JF1 identity is the input to DG3/LD9 and the completed $J$-row conductor in Section~\ref{component:pe3bbfd3a8151}. Those arguments allocate the remaining normal and retain the full angular and radial terms.



% Source formulas: peea9e9c546c5
\subsection{The small-$\eta$ angular cancellation}

\label{component:peea9e9c546c5}

Use the small-$\eta$, unit-$\zeta$ source: $x,\eta>0$
infinitesimal, $\zeta$ a positive bounded unit, and
$\beta\ge0$.  Put
\[
 a=1/\zeta,\quad c^2=a/\eta,\quad k^2=x,\quad
 B_x=\beta+\tfrac12,\quad q=\alpha+(m+B_x)(1-h)^2,
 \quad\mu=m-a^2q,\quad\lambda=\mu/(\eta ax).
\]
For homogeneous testing states use the exact roots
\[
\begin{gathered}
 \xi=X+huZ,\quad E=\eta uW-vZ,\quad
 D_0=u-cv,\quad J=(uv-x/c)Z,\\
 R_D=v(u^2-x)Z,\quad R_c=kv(\eta cW-Z),\quad
 \mathcal Q=\sqrt q(avW-uZ-a\xi).
\end{gathered}
\]
Direct multiplication gives
\[
 R_D-uJ=(x/c)D_0Z,\qquad
 \mathcal Q/\sqrt q=-D_0Z+(c/k)R_c-a\xi.
\]
Hence the complete principal $\xi$ row cancels exactly to
\[
 -2\eta\lambda ckR_c\xi+2\eta\lambda c(R_D-uJ)\xi
 =-2\eta\lambda x\mathcal Q\xi/\sqrt q
                    -2\eta\lambda xa\xi^2.             \tag{SC5}
\]
Together with the singular rank's $+\mu\xi^2$, the last term
leaves $-\mu\xi^2$, because $\mu=\eta\lambda ax$.
The latter coefficient is infinitesimal relative to $B_x\xi^2$.
The complete angular block on $(\mathcal Q,\sqrt{B_x}\xi)$
has strictly positive residue: the original source expressions give
\[
 B_xq/x^2\sim
 \frac{B_x\zeta(1+\zeta)}2+\frac{(1+\zeta)^2}4
       >\frac14\sim(\eta\lambda)^2.
\]
For unbounded $B_x$ the normalized off-diagonal tends to zero.
For bounded $B_x\ge1/2$, the displayed difference is bounded
below by $\zeta(3+2\zeta)/4>0$.  The source substitutions are
the regular DG/LD dictionary; the direct derivative and margin
calculation is also printed in the refinement proof.
All inverses and radicals in these coefficient expressions are
pure source.  The full radial and remaining normal allocations are
the retained JF/DG/LD/CJ terms; SC5 supplies their angular block.



% Source formulas: p60eca5442075
\subsection{Small $L$, unbounded $\zeta$ and unit $\eta$}

\label{component:p60eca5442075}

\label{p60eca5442075:sec:r17-lz-unit-eta}

\begin{theorem}
Retain the original physical C1 source and every original contact,
strict-gap, nonnegative-weight, total-scale and earlier-exclusion
intersection. Suppose $x=k^2,L=x\zeta$ are positive infinitesimals,
$\zeta$ is unbounded and $\eta$ is a positive bounded unit.
Then the complete original parent is finite tensor SOS for every
independent projective test and every real or complex state.
No positivity or noninfinitesimal lower bound is imposed on
$j_0=\eta+r-1$: only the actual original gap
$j=\eta(1+xh)-1+r>0$ is used. Both $r$ openings, simultaneous
$r,j$ openings and literal $\beta=0$ are included.
\end{theorem}

This theorem treats the whole $\eta$-unit source portion of C14.3a.
Every identity below uses the normalized quarter-parent $\mathcal P=F/4$.
The complete parent identity JF1--JF3 from R16 is algebraic in the
original stationary source; none of its identities assumes unit
$\zeta$. We retain that identity and change its source normalization.

\paragraph{Opening-cleared source bounds.}
Put
\[
 a=1/\zeta,\quad x=aL,\quad y=\eta L,\quad
 \theta=1-h=L\pi,\quad B=\beta+\tfrac12,\quad D_p=D+\tfrac12,
 \quad \rho_0=r+\eta.
\]
The original beta numerator gives bounded $\pi>0$;
$h,M,E,Q,P_d,P_4\sim1$ as in SB1--SB3. In particular
$\rho_0$ is a positive bounded unit. All formulas here preserve
actual $j>0$, rather than replacing it by $j_0$.
Define, exactly as in JF1,
\[
\begin{gathered}
 q=\alpha+(m+B)\theta^2,\quad \mu=m-a^2q,\quad
 \lambda=\mu/(\eta a^2L),\\
 T_0=A+\alpha+\tfrac12,\quad T_m=T_0+m,\quad
 e_4=r^2T_m/\eta^2,\\
 b'=Ar/\eta-m\theta-\alpha-\tfrac12+rT_mL/2,\\
 \kappa=2[B\theta-D_pj/\eta-mr/\eta],\qquad
 \gamma=2b'/h+2D_pj/r-2\eta B\theta/r.
\end{gathered}                                                    \tag{EU1}
\]
Finite canceled source division gives
\[
\begin{gathered}
 C(1-r)\sim\rho_0/2,\quad Ar\sim\eta/2,\quad
 D_pj\sim\rho_0/2,\quad m/(a^2L)\sim\rho_0/2,\\
 \alpha/L^2\sim\rho_0/2,\quad B\theta/L\sim\rho_0/2,
 \quad q/L^2-\rho_0(1+\pi)/2\in\operatorname{Inf},\\
 \mu/(a^2L)\sim\rho_0/2,\quad rT_m\sim\rho_0/2,
 \quad b'/L\sim-\rho_0/4,\quad \eta\kappa\sim-\rho_0.
\end{gathered}                                                    \tag{EU2}
\]
The canceled denominators at $a=L=0$ are only $1,2,4$;
there is no $r$, $1-r$, $j$, $\theta$ or $j_0$ factor.
Thus these are uniform through every opening in the theorem.
In particular $C,A,D_p$ are bounded below by positive real constants,
although any of them may be unbounded. Also $q/L^2$ and $\lambda$
are positive bounded units, $\mu>0$, and $b'=O(L)$.
The positive lower bound for $D_p$ uses bounded positive actual $j$
and the third estimate, not positivity of $j_0$.

\paragraph{Use the complete homogeneous parent.}
For nonnegative testing angles $u,v$ put $U=u^2,V=v^2$ and define
\[
\begin{gathered}
 E_t=\eta uW-vZ,\quad Q_t=Y-rvZ/\eta,\quad
 P=E_t/\eta-Q_t,\quad \xi=X+huZ,\\
 R_D=v(U-x)Z,\quad T_D=R_D-uv\xi/h,\quad
 H=avW-uZ-a\xi,\\
 F_b=(V-y)Z+vE_t,\qquad
 F_t=(V-y)Z+(\eta/r)vQ_t.
\end{gathered}                                                    \tag{EU3}
\]
All inverses in these roots are pure source. The exact full parent is
\[
\begin{split}
\mathcal P={}&qH^2+\mathscr C
 +CP^2+AQ_t^2+D_p(P-hT_D)^2+2b'Q_tT_D+\mathscr R\\
 &+e_4UF_t^2-e_4yu(H+a\xi)F_t
 +AUQ_t^2+B(1+V)\xi^2-2(\mu/a)H\xi-\mu\xi^2\\
 &+\kappa uF_t\xi+\gamma uvQ_t\xi .
\end{split}                                                       \tag{EU4}
\]
This is the complete JF1 identity, including every hard row and every
first normal. Set
\[
 \rho=m+rT_mx/2=O(x),\qquad
 K_r=T_0+2\rho-\rho^2/C.
\]
Then $K_r/T_m\sim1$ and $K_r$ is bounded below by a positive
real constant. The complete radial row is exactly
\[
\mathscr R=R_b^2+R_t^2,\qquad
 R_b={\sqrt C\over\eta}[F_b-r(1-\rho/C)F_t],\qquad
 R_t={r\sqrt{K_r}\over\eta}F_t.                         \tag{EU5}
\]
It follows that
\[
 {F_b\over\eta}={R_b\over\sqrt C}
               +{1-\rho/C\over\sqrt{K_r}}R_t          \tag{EU6}
\]
is a bounded pure-source allocation even when $r$ is infinitesimal.
No independent bare $(V-y)Z$ pivot has been inserted.

\paragraph{A strict source frame without either gap inverse bound.}
Use the three original weighted triangle roots
\[
 N_C=\sqrt C\,P,\quad N_A=\sqrt A\,Q_t,\quad
 N_D=\sqrt{D_p}(P-hT_D).                               \tag{EU7}
\]
Their positive squares are exactly the three original source-triangle
terms in EU4. They give the bounded pure-source allocations
\[
 P={N_C\over\sqrt C},\quad Q_t={N_A\over\sqrt A},\quad
 T_D={N_C/\sqrt C-N_D/\sqrt{D_p}\over h},\quad
 E_t=\eta(N_C/\sqrt C+N_A/\sqrt A).                    \tag{EU8}
\]
Consequently the entire $2b'Q_tT_D$ term is an infinitesimal Gram
on this strict diagonal frame. This does not cap or remove a divergent
original event and does not assert boundedness of its coefficient.

\paragraph{The complete two-normal cancellation.}
The source definitions in EU1 give the exact identity
\[
             {\eta\kappa\over r}+\gamma=2b'/h-2m.
\]
Since $F_t=F_b-vE_t+(\eta/r)vQ_t$, both complete normal rows in EU4
therefore satisfy
\[
\boxed{\quad
 \kappa uF_t\xi+\gamma uvQ_t\xi
 =\kappa uF_b\xi
     +uv\{(2b'/h-2m)Q_t-\kappa E_t\}\xi .\quad}       \tag{EU9}
\]
This is the crucial simultaneous allocation. It cancels the potentially
unbounded $1/r$ coefficients before any estimate is made. It neither
estimates the two original rows separately nor uses their contact
vanishing as a substitute for the identity.

Set $N_x=\sqrt B\xi$ and $\mathcal Q=\sqrt q H$.
By EU2, EU6 and EU8, every coefficient of EU9 on
$(N_C,N_A,R_b,R_t,N_x)$ is infinitesimal: the first uses
$\eta\kappa=O(1)$ and an extra $u$; the second uses
$b',m\in\operatorname{Inf}$ and an extra $uv$; the third uses
$\eta\kappa=O(1)$ and an extra $uv$.
The factors $1/\sqrt B$, $1/\sqrt A$, $1/\sqrt C$ and
$1/\sqrt{K_r}$ are bounded. No testing-state normalization occurs.

\paragraph{Every remaining scalar and hard term.}
The complete angular block on $(\mathcal Q,N_x)$ is
\[
 \mathcal Q^2-{2\mu\over a\sqrt{qB}}\mathcal QN_x
                         +(1-\mu/B)N_x^2.
\]
Its off-diagonal coefficient is $O(a)$ and $\mu/B$ is infinitesimal:
\[
 {\mu\over a\sqrt{qB}}
   =a\,{\mu/(a^2L)\over\sqrt{q/L^2}\sqrt B}.
\]
Thus its residue is the identity, including literal $\beta=0$.
The radial-angular cross has the exact normalized expression
\[
 -e_4yu(H+a\xi)F_t
 =-{rT_mL\over\sqrt{K_r}}u
       \left({\mathcal Q\over\sqrt q}
                         +{aN_x\over\sqrt B}\right)R_t.          \tag{EU10}
\]
Both coefficients are infinitesimal because $rT_m=O(1)$,
$L/\sqrt q=O(1)$ and $u$ is infinitesimal.

The coupled circuit $\mathscr C$ is exactly the sum of the following
three whole positive squares, followed by $-\lambda xE_t^2$:
\[
 \sqrt\lambda(uv-x/c)Z,\qquad
 \sqrt\lambda\,\eta k(cv-u)W,\qquad
 \sqrt{2\lambda/c}\,k\sqrt{uv}(\eta cW-Z),
 \qquad c^2=a/\eta,\quad k^2=x.                       \tag{EU11}
\]
The last negative term is an infinitesimal Gram on EU8 since
$\lambda$ is bounded and $x$ infinitesimal. All radicals in EU11
are separated pure-source and pure-testing radicals.

Retain EU11's three squares, $e_4UF_t^2$, $AUQ_t^2$ and
$BV\xi^2$ as whole positive terms. Every other term of EU4 has now
been assigned explicitly to the seven-root frame
\[
                 (N_C,N_A,N_D,R_b,R_t,\mathcal Q,N_x).             \tag{EU12}
\]
The Gram on this frame has identity residue and bounded coefficients.
The bounded-Gram lemma proves finite tensor SOS for the complete
parent whenever both testing angles are infinitesimal. Every estimate
is a product of a source bound and a testing bound; there is no
source/testing order comparison and no inverse of a mixed tensor.

\paragraph{All tests, endpoints and states.}
The roots and the whole-parent identity are homogeneous in all four
state coordinates. In particular no $Z$ inverse or state restriction
was used. The positive circuit rows EU11 are defined on literal
$u=0$ or $v=0$ and their squared identity is polynomial there.
The common exact sign gauges give all local sign charts.
The unchanged small-$L$, bounded-$\eta$ remote theorem covers the
complement of both infinitesimal testing angles on this exact source
domain, retaining all its projective endpoints and openings.
Applying the real homogeneous identity to real and imaginary state
parts proves the complex-state assertion. All of the source unit-$\eta$
portion of C14.3a is therefore all-tests covered.

\paragraph{The complementary source.}
The remaining C14.3a sources have $\eta$ infinitesimal, hence $r\sim1$, with all original $1-r,j$ openings retained. Their complete parent is treated by the LD identities and the completed $J$-row conductor in Sections~\ref{component:pc720c1e06aca} and~\ref{component:pe3bbfd3a8151}.



% Source formulas: p024c821bfb59
\subsection{Small $\eta$ and unit $\zeta$: the damped positive block}

\label{component:p024c821bfb59}

\begin{theorem}
Normalize $d=1$ and retain the unchanged physical C14.3b source, with $x,\eta>0$
infinitesimal, $\zeta$ a positive bounded unit, every original $j$
and $1-r$ opening, $\beta\ge0$ including literal $\beta=0$, and all common intersections.
On the additional \emph{pure-source} subchart
\[
                         x/\eta\text{ bounded},              \tag{DG1}
\]
the entire original parent is finite tensor SOS for every independent
projective angle and state test. No relation between a source rate
and a testing rate is assumed. The proof below treats the whole local
corner; the common remote theorem gives its complement. The source
subchart $x/\eta$ unbounded is not covered by this theorem.
\end{theorem}

\paragraph{The sign reduction and the complete frame.}
Use the sign gauges to take $u,v\ge0$. For real states with $WZ\le0$
the independent half-core sign lemma proved in this revision gives a
direct tensor-SOS decomposition of the original signed half core;
all remaining original event terms are positive source multiples
of squares. Thus only $WZ>0$ needs the following construction. No
normalization or inverse of a state coordinate is used here.

Use all homogeneous JF1 definitions, and abbreviate $Q=Q_t$,
$T=T_D$, $B=B_x$, $a=1/\zeta$, $x=k^2$, $a=\eta c^2$. In particular
\[
 E=\eta uW-vZ=\eta(P+Q),\quad H=avW-uZ,
 \quad \tau=\eta cW-Z,\quad D_0=u-cv,
\]
\[
 J=(uv-x/c)Z,\quad
 R=-\sqrt\lambda\,\eta kD_0W,\quad
 f_t=F_t/\eta,\quad \mathcal Q=\sqrt q(H-a\xi),
 \quad N_x=\sqrt B\xi.
\]
The full radial term $\mathscr R$ is JF2, with the strict two-root
completion JF3. Its roots will be denoted $R_b,f_t$; both normalized
pivots have positive unit residue. The angular term is exactly
\[
 \mathscr A=\mathcal Q^2
       -\frac{2\mu}{a\sqrt{qB}}\mathcal QN_x
       +(1-\mu/B)N_x^2.                              \tag{DG2}
\]
Its normalized Schur pivot is strict by the unchanged R15/R16 source
margin, including $\beta=0$ and unbounded $B$. Those previously
proved angular and radial blocks are used once, as their complete
terms.

\paragraph{Damp the remaining normal and complete both circuit companions.}
Set the positive source unit
\[
 G_0=2\eta\lambda-2\eta\sim1,\qquad
 \epsilon=\gamma-G_0/h,
\]
\[
 \widehat R=R+\frac{G_0k}{2\sqrt\lambda}Q,
 \qquad \widehat J=\sqrt\lambda J+
                         \frac{G_0u}{2\sqrt\lambda}Q,
 \qquad c_U=A-\frac{G_0^2}{4\lambda}.
\]
Use the exact full-$D$ conductor and polarization displayed in Chapter~\ref{ch:mechanisms} to conduct only
$(G_0/h)uvQ\xi$, leaving $\epsilon uvQ\xi$ unaltered. Complete
\emph{both} $R$ and $J$, incorporating the entire native $AUQ^2$
row into the latter completion. Since
$G_0/(2\lambda)=\eta-\eta/\lambda$, the complete original parent
has the following exact identity:
\[
\begin{split}
 F/4={}&\mathscr A+\mathscr R+S_0
    -\lambda x(\eta P+\eta Q/\lambda)^2
    +\widehat R^2+\widehat J^2\\
 &+c_UUQ^2+e_4UF_t^2+BV\xi^2
    +\frac{2x}{c}uv\tau^2
    +\frac{G_0x}{c}D_0^2WZ\\
 &-e_4yuHF_t+\kappa uF_t\xi
                 +\frac{G_0x}{c}PH+\epsilon uvQ\xi,       
\end{split}\tag{DG3}
\]
where
\[
 S_0=CP^2+AQ^2+D_p(P-hT)^2+2b_0QT,
 \qquad b_0=b'-G_0/2.                                 \tag{DG4}
\]
This is a simultaneous completion of the entire remaining normal and
the coupled circuit. It retains every full-$D$, full-$Y$ and higher
normal of JF1. There is no second bare radial pivot and no extracted
uncentered $m$ event. The coefficient of the extra $uv\tau^2$ row
is exactly $2x/c>0$, because $2\lambda-G_0/\eta=2$.

\paragraph{Two strictly positive source deficits replace the failed soft pivot.}
Put $K=CD_p/(C+D_p)$, $s=x+\eta$, and
\[
                         \Delta_0=h^2AK-b_0^2.
\]
The original source formulas imply exact finite divisions
\[
 \Delta_0=x d_x+\eta d_\eta,\qquad
 \frac{c_U}{\eta}=x e_x+\eta e_\eta,                \tag{DG5}
\]
with bounded source coefficients having the following residues:
\[
 \begin{array}{c|c}
 d_x&\pi/2+1/4+1/(4a)\\
 d_\eta&\sigma/2+3/4\\
 e_x&\pi+1+1/(2a)\\
 e_\eta&\sigma+3/2
 \end{array}
 \qquad h=1-x\pi,\quad r=1-\eta\sigma.             \tag{DG6}
\]
All four residues are strictly positive. Here $\pi,\sigma\ge0$
are bounded by the unchanged physical source estimates. Thus
\[
 \Delta_0/s\text{ is a positive bounded unit},\qquad
 c_U/(\eta s)\text{ is a positive bounded unit}.     \tag{DG7}
\]
These assertions compare source quantities only and do not use DG1.

For clarity, the divisions in DG5 are opening-safe. Write
$C_n=(1-r)C$, $D_j=jD_p$ and $\Theta_B=B(1-h)$, all given by their
unchanged unit-denominator source formulas. The identity
$\eta\Theta_B-D_j=-C_n$ gives
$b'-\gamma h/2=-hC_n/r$. The gap-cleared harmonic coefficient is
\[
 \eta K=\frac{C_nD_j}
 {C_n(1+xh-\sigma)+D_j\sigma}.
\]
Its denominator is a positive unit, even at either opening. The
quantities $\eta\lambda$, $b'$, $A/\eta$ and $\gamma h$ also have
unit-denominator expressions after their displayed monomial factors
are canceled. Both functions on the left of DG5 vanish at
$x=\eta=0$. Define their $\eta$ coefficients by their exact $x=0$
values divided by $\eta$, and then divide the remainder by $x$.
This gives DG5 with only source-unit denominators; differentiation at
the origin gives DG6. No assumption on $x/\eta$ enters these exact
finite divisions.

\paragraph{The source roots control the complete soft direction.}
There is the exact pure-source completion
\[
 S_0=N_s^2+N_T^2+N_Q^2,
\]
\[
 N_s=\sqrt{C+D_p}\left(P-\frac{hD_p}{C+D_p}T\right),
 \quad N_T=h\sqrt K\left(T+\frac{b_0}{h^2K}Q\right),
 \quad N_Q=\sqrt{\frac{\Delta_0}{h^2K}}Q.            \tag{DG8}
\]
As $\eta K\sim1/2$, DG7 gives
$N_Q/\{\sqrt{\eta s}\,Q\}$ a positive source unit. The inverse
identities show
\[
 T=\frac{N_T}{h\sqrt K}-\frac{b_0}{h^2K}Q,
 \qquad
 P=\frac{N_s}{\sqrt{C+D_p}}
                      +\frac{hD_p}{C+D_p}T.         \tag{DG9}
\]
Consequently $P$ has bounded coefficients on $(N_s,N_T,N_Q)$:
the first two coefficients are $O(\sqrt\eta)$ and the last is
$O(\sqrt{\eta/s})$, which is bounded because $\eta\le s$.
This is the complete $P$ control after the full $D$ shift; it is not
an assumption that the old $E/\sqrt\eta$ pivot remains strict.

The negative square in DG3 is an infinitesimal Gram on these three
roots. Its root is
\[
 \sqrt{\lambda x}(\eta P+\eta Q/\lambda)
   =\eta\sqrt{\lambda x}\,P+\eta\sqrt{x/\lambda}\,Q.
\]
Since $\eta\lambda\sim1/2$, the first coefficient is
$O(\sqrt{\eta x})$ on the bounded allocation of $P$; the second,
after normalizing $Q$ by $N_Q$, is
$O(\eta\sqrt{x/s})$. Both are infinitesimal, without DG1.

\paragraph{Every remaining complete normal is small on this frame.}
The angular-radial term in DG3 is
\[
 -e_4y\eta u\left(\mathcal Q/\sqrt q
                             +aN_x/\sqrt B\right)f_t.
\]
Its coefficients are $O(u)$ and $O(xu)$, using
$e_4\eta^2\sim1/2$ and $q/x^2$ a positive bounded unit.
The term $\kappa uF_t\xi$ has coefficient
$(\eta\kappa)u/\sqrt B$, again infinitesimal.
The new $PH$ term has coefficients
\[
 \frac{G_0x}{c\sqrt q}\quad\hbox{and}\quad
 \frac{G_0ax}{c\sqrt B}
\]
on $P\mathcal Q$ and $PN_x$; these are respectively
$O(\sqrt\eta)$ and $O(x\sqrt\eta)$, so DG9 allocates the whole
term infinitesimally.

Finally $\epsilon=O(s)$ by the same gap-free source expressions.
The complete, still unconducted normal has coefficient
\[
 \frac{\epsilon uv}{\sqrt{\eta s}\sqrt B}
         =O\!\left(\sqrt{s/\eta}\,uv\right)          \tag{DG10}
\]
on $N_QN_x$. DG1 is used \emph{exactly here}: it makes $s/\eta$
bounded, and hence the coefficient infinitesimal for independent
infinitesimal $u,v$. No testing comparison is imposed.

Thus the remainder on
\[
                  (N_s,N_T,N_Q,R_b,f_t,\mathcal Q,N_x)
\]
has a bounded Gram with strictly positive residue. The bounded-Gram
lemma makes it finite tensor SOS. Add the whole positive rows
$\widehat R^2$, $\widehat J^2$, $c_UUQ^2$, $e_4UF_t^2$,
$BV\xi^2$, $(2x/c)uv\tau^2$ and
$(G_0x/c)D_0^2WZ$. The last is a literal tensor square on the current
pure-testing sign chart; all radicals are source-only or testing-only.
This proves the complete local parent, not merely its scalar
restriction.

\paragraph{Equality, endpoints and the global join.}
At original contact $P,Q,T,J,R,\tau,F_t,\xi$ vanish, so all newly
completed rows vanish too. DG3 preserves all contact normals exactly.
The row formulas have no angle ratio or state-coordinate inverse;
$\sqrt{uv}$ is defined as zero on a literal axis. Thus they include
both literal local axes. The common gauges cover every angle sign.
For a complex state, split the real-coefficient Hermitian parent into
its real and imaginary evaluations and apply either the $WZ\le0$
core lemma or this $WZ>0$ proof separately to each real state. There
is no assumption that both real and imaginary state components have
a unit $Z$ or $W$. Combine the local proof with the unchanged remote
all-state theorem to obtain the all-tests assertion on DG1.

The hypothesis DG1 is used only in the final allocation DG10. The exact DG3 identity and preceding source bounds also apply without it; the completed $J$-row conductor in Section~\ref{component:pe3bbfd3a8151} supplies the remaining allocation for arbitrary $x/\eta$.



% Source formulas: pc720c1e06aca
\subsection{Small $\eta$ and unbounded $\zeta$: the damped positive block}

\label{component:pc720c1e06aca}

\label{pc720c1e06aca:sec:r17-large-zeta-damped}

\begin{theorem}
Normalize $d=1$ and retain the unchanged physical C14.3a source,
its actual stationary contact, every strict source gap,
$\beta\ge0$ including literal $\beta=0$, the original total-scale
intersection and every earlier exclusion. Suppose $x,L=x\zeta,\eta$
are positive infinitesimals, $\zeta$ is unbounded, and
\[
                         L/\eta\text{ is bounded}.       \tag{LD1}
\]
Then the entire original parent is finite tensor SOS at every independent
projective angle and state test, real or complex. This is a whole-state,
all-tests source-subchart certificate. No condition is imposed on the
testing angle ratio or on the source ratio $a/\eta$, where $a=1/\zeta$.
The source subchart with $\eta$ infinitesimal and $L/\eta$ unbounded
is not covered here.
\end{theorem}

The proof uses the exact complete damped identity DG3, but its source
scale is $L+\eta$ rather than $x+\eta$. This change must be proved:
one cannot invoke the unit-$\zeta$ estimates by substituting an
infinitesimal $a$ into their displayed $1/a$ coefficients. All formulas
below use the actual physical source and the original contact.

\paragraph{Source variables and opening-cleared coefficients.}
Put
\[
 a=1/\zeta,\quad x=aL,\quad y=\eta L,\quad
 h=1-L\pi,\quad r=1-\eta\sigma,\quad B=\beta+\tfrac12.
\]
The original source estimates give bounded $\pi,\sigma\ge0$, and
actual $j>0$ says
$0<\sigma<1+xh$, $j=\eta(1+xh-\sigma)>0$.
Neither opening has a positive-residue lower bound.
The exact native source formulas displayed below give
\[
 \frac{m}{a^2L}\sim\frac12,\quad
 \frac{\alpha}{L^2}\sim\frac12,\quad
 \frac{q}{L^2}-\frac{1+\pi}{2}\in\operatorname{Inf},\quad
 \frac{\mu}{a^2L}\sim\frac12,
\]
\[
 \eta\lambda\sim\frac12,\quad
 \eta K\sim\frac12,\quad
 \eta\kappa\sim-1,\quad \eta^2e_4\sim\frac12,
 \quad b'=O(L),                                     \tag{LD2}
\]
where $q=\alpha+(m+B)(1-h)^2$, $\mu=m-a^2q$,
$\lambda=\mu/(\eta a^2L)$ and $K=CD_p/(C+D_p)$.

Here is a completely regular specification of the new source
expressions. Write
\[
 K_x=1+x,\ R_s=1+xr,\ T_s=1+xh,\ M_s=R_s-aL\pi,
 \ E_s=1-x\eta h,\ Q_s=1-Lr,
\]
\[
 P_d=1+2xh+rx^2h+rxL+L
                  +\eta\{K_xL+x^2hL\pi\},\qquad
 P_4=T_sE_s+(1-r)(xh+L).
\]
Define the opening-cleared quantities
\[
 C_n=\frac{K_x(r+\eta)E_s}{2P_d},\qquad
 D_j=\frac{K_x(1+y)(r+\eta)}{2P_d},\qquad
 A_\eta=\frac{P_4}{2rP_d},
\]
\[
 m_n=\frac{K_x(r+\eta)h(1+y)}{2M_sP_d},\quad
 \alpha_n=\frac{(r+\eta)\{aR_sh(1+y)+M_sE_s\}}
                         {2Q_sP_d},\quad
 B_n=\frac{K_x(r+\eta)(1+ah)}{2P_d}.                 \tag{LD3}
\]
Thus $C_n=(1-r)C$, $D_j=jD_p$, $A_\eta=A/\eta$,
$m_n=m/(a^2L)$, $\alpha_n=\alpha/L^2$ and
$B_n=B(1-h)/L$. In particular
\[
 q_n:=q/L^2=\alpha_n+m\pi^2+B_n\pi,\qquad
 \Lambda:=\eta\lambda=m_n-Lq_n,
\]
\[
 \eta K=\frac{C_nD_j}
                  {C_n(T_s-\sigma)+D_j\sigma}.       \tag{LD4}
\]
There is no $a^{-1}$ in (LD3)--(LD4). Their denominators at
$a=L=\eta=0$ have residues $1$, $2$ or $1/2$, independently
of $\pi,\sigma$. The remaining definitions are
\[
 T_m=A+\alpha+\tfrac12+m,\quad
 b'=A_\eta r-mL\pi-\alpha-\tfrac12+rT_mL/2,
\]
\[
 G_f:=\gamma h=2b'+2D_jh/r-2\eta LB_nh/r,
 \quad G_0=2\Lambda-2\eta\sim1,\quad b_0=b'-G_0/2.
                                                               \tag{LD5}
\]
These are the actual native-source expressions used in the completion.

\paragraph{Positive source deficits at the new scale.}
Set $s=L+\eta$ and retain DG's definitions
\[
 \Delta_0=h^2AK-b_0^2,
 \qquad c_U=A-\frac{G_0^2}{4\lambda},
 \qquad\epsilon=\gamma-G_0/h.
\]
The regular source formulas above give exact finite divisions
\[
 \Delta_0=Ld_L+\eta d_\eta,\qquad
 \frac{c_U}{\eta}=Le_L+\eta e_\eta,\qquad
 G_f-G_0=Lf_L+\eta f_\eta,                           \tag{LD6}
\]
with bounded coefficients and residues
\[
\begin{array}{c|cc}
 &L\text{ coefficient}&\eta\text{ coefficient}\\\hline
 \Delta_0&\pi/2+1/4&\sigma/2+3/4\\
 c_U/\eta&\pi+1/2&\sigma+3/2\\
 G_f-G_0&\pi+1/2&\sigma+2
\end{array}                                                     \tag{LD7}
\]
(as usual the bounded $\pi,\sigma$ on the right can be replaced
by their real residues). To justify the finite divisions explicitly,
each expression on the left of (LD6) vanishes at $L=\eta=0$
\emph{for arbitrary $a$}. Subtract its exact value at $L=0$, divide
that difference by $L$, and divide the $L=0$ value by $\eta$.
Clearing the source-unit denominators proves each divisibility in a
polynomial ring. Evaluation of these quotients at $a=L=\eta=0$
gives (LD7). There is consequently no hidden pure-$a$ remainder.

All four coefficients for the first two rows have strictly positive
residue. Hence
\[
 \Delta_0/s\text{ and }c_U/(\eta s)
       \text{ are positive bounded source units},\qquad
 \epsilon=O(s).                                      \tag{LD8}
\]
This source calculation itself does not use (LD1).

\paragraph{The same complete damped identity, without a state chart.}
Use DG's unshifted $H=avW-uZ$; the strict angular root is
$\mathcal Q=\sqrt q(H-a\xi)$, so this $H$ differs from the
shifted letter $H$ in EU3. Retain every other homogeneous JF/DG
definition. In particular
\[
 E=\eta(P+Q),\quad N_x=\sqrt B\xi,\quad
 \tau=\eta cW-Z,\quad D_0=u-cv,\quad c^2=a/\eta.
\]
After the pure-testing sign gauges, the half-core lemma supplies all
$WZ\le0$ states. On $WZ>0$, the exact identity DG3 is valid for
this source without any change:
\[
\begin{split}
 F/4={}&\mathscr A+\mathscr R+S_0
       -\lambda x(\eta P+\eta Q/\lambda)^2
       +\widehat R^2+\widehat J^2\\
 &+c_UUQ^2+e_4UF_t^2+BV\xi^2
       +\frac{2x}{c}uv\tau^2+\frac{G_0x}{c}D_0^2WZ\\
 &-e_4yuHF_t+\kappa uF_t\xi
       +\frac{G_0x}{c}PH+\epsilon uvQ\xi.             
\end{split}\tag{LD9}
\]
Here $S_0=CP^2+AQ^2+D_p(P-hT)^2+2b_0QT$,
$\widehat R=R+G_0\sqrt xQ/(2\sqrt\lambda)$ and
$\widehat J=\sqrt\lambda J+G_0uQ/(2\sqrt\lambda)$.
The identity conducts only $G_0uvQ\xi/h$ and completes both
circuit companions exactly, retaining every first normal and hard
coupling. Its derivation uses only $x=k^2$, $c^2=a/\eta$ and the
homogeneous JF conductor, not a bounded reciprocal of $a$ or $c$.

\paragraph{Strict source, radial and angular frames.}
Use the exact source roots
\[
 N_s=\sqrt{C+D_p}\left(P-\frac{hD_p}{C+D_p}T\right),
 \quad N_T=h\sqrt K\left(T+\frac{b_0}{h^2K}Q\right),
 \quad N_Q=\sqrt{\frac{\Delta_0}{h^2K}}Q.            \tag{LD10}
\]
Their squares sum to $S_0$, and
$N_Q/[\sqrt{\eta s}\,Q]$ is a positive source unit.
The inverse source identities give bounded coefficients for $P$
on this frame: the coefficients of $N_s,N_T$ are $O(\sqrt\eta)$
and that of $N_Q$ is $O(\sqrt{\eta/s})$. This uses only
$\eta\le s$, so no old $E/\sqrt\eta$ pivot is imported.

The complete radial term has its common strict roots $R_b,f_t$,
with $f_t=F_t/\eta$. Indeed $\rho=m+rT_mx/2=O(x)$ and
$K_r\sim1/2$. The complete angular block is
\[
 \mathscr A=\mathcal Q^2-
       \frac{2\mu}{a\sqrt{qB}}\mathcal QN_x
                  +(1-\mu/B)N_x^2.
\]
It has identity residue, since
\[
 \frac{\mu}{a\sqrt{qB}}
       =\frac{a[\mu/(a^2L)]}{\sqrt{q/L^2}\sqrt B}=O(a),
 \qquad\mu/B\in\operatorname{Inf}.                 \tag{LD11}
\]
This includes unbounded $B$ and literal $\beta=0$.

\paragraph{The entire remainder has infinitesimal coefficients.}
The negative-square root in (LD9) has, on the source frame,
the same coefficient bounds as DG:
\[
 \eta\sqrt{\lambda x}\,P+
           \eta\sqrt{x/\lambda}\,Q,
 \qquad O(\sqrt{\eta x}),\quad
           O(\eta\sqrt{x/s}).                       \tag{LD12}
\]
They are infinitesimal; here $x=aL$ and no comparison of $a$ with
$\eta$ is needed. For the radial-angular term use
$H=\mathcal Q/\sqrt q+aN_x/\sqrt B$ and $F_t=\eta f_t$.
Since $e_4y\eta=(e_4\eta^2)L=O(L)$ and $\sqrt q\asymp L$,
its two coefficients are $O(u)$ and $O(xu)$.
The $\kappa$ term has coefficient
$(\eta\kappa)u/\sqrt B$, which is also infinitesimal.

The new $PH$ term improves at this source end:
\[
 \frac{G_0x}{c\sqrt q}=O(\sqrt{a\eta}),\qquad
 \frac{G_0ax}{c\sqrt B}=O(x\sqrt{a\eta}).            \tag{LD13}
\]
Thus the bounded $P$ allocation makes both of its angular coefficients
infinitesimal. Finally, the complete unprocessed normal is
\[
 \frac{\epsilon uv}{\sqrt{\eta s}\sqrt B}N_QN_x
   \quad\hbox{up to a bounded positive source-unit factor}.
\]
Its coefficient is
\[
                  O\!\left(\sqrt{1+L/\eta}\,uv\right),          \tag{LD14}
\]
which is infinitesimal under the bounded-source-ratio
hypothesis (LD1). This is the only use of that hypothesis.

Consequently the complete remainder on
$(N_s,N_T,N_Q,R_b,f_t,\mathcal Q,N_x)$ has a bounded strict-residue
Gram. The bounded-Gram lemma gives its finite tensor SOS. Every other
term in (LD9) is a whole positive square: in particular $c_U>0$,
$G_0>0$, and $uv,WZ$ are nonnegative pure-testing factors on this
sign chart. No mixed inverse or radical is taken.

\paragraph{All tests and exact remaining source end.}
No state normalization, state-coordinate inverse or testing ratio
occurs in this proof. The division-free circuit rows and
$\sqrt{uv}=0$ include both literal local axes. The half-core exit
and the exact angle sign gauges cover every real state; splitting
the real-coefficient Hermitian form covers complex states. The
unchanged small-$L$, bounded-$\eta$ remote theorem covers all other
projective testing pairs, including literal infinities. This proves
the theorem with every original source opening retained.

The restriction LD1 is used only in LD14. The complete LD9 identity, source bounds and all preceding normal allocations are inputs to the completed $J$-row conductor in Section~\ref{component:pe3bbfd3a8151}, which treats the full infinitesimal-$\eta$ source at arbitrary $L/\eta$. The unit-$\eta$ source is the EU theorem of Section~\ref{component:p60eca5442075}.



% Source formulas: p8cfada778490
\subsection{The direct large-$X/Z$ state bound}

\label{component:p8cfada778490}

In any local chart with $u$ infinitesimal, suppose $Z\ne0$ and
the pure-testing ratio $|X/Z|$ is unbounded.  Then
\[
 \frac{1+v^2}{2}\{(X+uZ)^2-(1+u^2)Z^2\}>0.             \tag{DS3}
\]
Indeed, after division by $Z^2>0$, the square $(X/Z+u)^2$
is unbounded and $1+u^2$ is bounded.  The factor $1+v^2$
is positive.  This expression is the complete $X$ square and the
negative $Z$ term of the signed quarter-core.  Its other two
terms are $(uW+uvX-Y-vZ)^2/2$ and
$(1+u^2)(vY-Z)^2/2$, both nonnegative.  Thus the signed core
itself is a nonnegative pure-testing scalar and has a square root
in the testing field.  Every original event remains a square with
its nonnegative source coefficient.  This proves a direct separated
SOS certificate for the complete parent, with arbitrary $W,Y$.

The state chart $Z=0$ and the sign chart $WZ\le0$ have the
direct core certificates already stated in the geometric reduction.
Consequently the C14.3 local certificate only needs the remaining
case $Z\ne0$, $WZ>0$, $X/Z$ bounded.  The completed-J
conductor CJ treats that whole case without bounding $W/Z$ or
$Y/Z$.  Apply the appropriate real chart to each real and imaginary
part to include complex states.



% Source formulas: pe3bbfd3a8151
\subsection{The completed $J$ relation for the entire remaining normal}

\label{component:pe3bbfd3a8151}

\label{pe3bbfd3a8151:sec:r18-completed-j}

\paragraph{Scope and unchanged inputs.}
Retain the complete actual-parent identity DG3, equivalently LD9,
with $d=1$, every original source opening and residual intersection,
and $\beta\ge0$ including literal $\beta=0$.  On C14.3b use
$s=x+\eta$ and the unit $a=1/\zeta$; on C14.3a use $s=L+\eta$
with $x=aL$ and $a$ infinitesimal.  In both cases $s$ and $\eta$
are positive infinitesimals.  The exact DG5--DG9 and LD3--LD10 source identities give
\[
 \Lambda=\eta\lambda\sim\tfrac12,\qquad
 \epsilon=O(s),\qquad
 N_Q=\sqrt{\frac{\Delta_0}{h^2K}}Q,
 \qquad
 Q=\frac{\omega}{\sqrt{\eta s}}N_Q,                 \tag{CJ1}
\]
where $\omega$ is a bounded positive source unit.  The complete strict
frame is
\[
 \mathcal F=(N_s,N_T,N_Q,R_b,f_t,\mathcal Q,N_x),
 \qquad N_x=\sqrt B\,\xi,\quad B=\beta+\tfrac12.
                                                               \tag{CJ2}
\]
All DG3/LD9 residuals other than $\epsilon uvQ\xi$ already have
infinitesimal Gram coefficients on this frame without any bound on
$s/\eta$.  The two residuals in the exact split CJ6a below have that same
property whenever $Z\ne0$ and $X/Z$ is bounded.  Thus the only
expression to be allocated on that state chart is
\[
 N=\epsilon U(W/Z)E\xi,
 \qquad E=\eta(P+Q)=\eta uW-vZ.                    \tag{CJ3}
\]
There is no restriction on $Y/Z$ or $W/Z$ below.

\paragraph{An exact identity from the already retained completed row.}
Let $\phi=\xi/Z=X/Z+hu$, and retain exactly DG's completed square
\[
 \widehat J=\sqrt\lambda\,(uv-x/c)Z
       +\frac{G_0u}{2\sqrt\lambda}Q,
 \qquad G_0=2\eta\lambda-2\eta,\qquad c^2=a/\eta.
\]
The defining identity $vZ=\eta uW-\eta(P+Q)$ gives
\[
 \boxed{\quad
 \frac{\widehat J}{\sqrt\lambda}
 =\eta U W-\frac{x}{c}Z
                -\eta u\left(P+\frac Q\lambda\right).
 \quad}                                                       \tag{CJ4}
\]
Multiply CJ4 by $\epsilon\phi E/\eta$.  The complete remaining
normal has the exact conductor
\[
 \boxed{\quad
 N=\frac{\epsilon\phi}{\eta\sqrt\lambda}E\widehat J
       +\epsilon\phi uE\left(P+\frac Q\lambda\right)
       +\frac{\epsilon x}{\eta c}E\xi .
 \quad}                                                       \tag{CJ5}
\]
The only state inverse here is the pure-testing $Z^{-1}$.  No inverse
of $u$, $v$, $W$, or a mixed expression occurs.  No radial row, source
pivot, or completed square is duplicated.

\paragraph{Bounded allocations and infinitesimal coefficients.}
Define two dependent roots
\[
 \mathcal E=\sqrt{s/\eta}\,E,
 \qquad \mathcal P=P+Q/\lambda.
\]
Both have bounded allocations on the existing source roots
$(N_s,N_T,N_Q)$.  Indeed the exact DG9/LD10 inverse identities give
coefficients $O(\sqrt\eta)$, $O(\sqrt\eta)$ and
$O(\sqrt{\eta/s})$ for $P$.  Hence the coefficients for
$\sqrt{\eta s}\,P$ are bounded (in fact infinitesimal), whereas
\[
 \mathcal E=\sqrt{\eta s}\,P+\omega N_Q,
 \qquad
 Q/\lambda=\frac{\omega}{\Lambda}\sqrt{\eta/s}\,N_Q.
                                                               \tag{CJ6}
\]
Here $\eta/s\le1$ is a pure-source inequality.  These are bounded
allocations of dependent roots, not new positive pivots.
For completeness, the other two terms of the same normal have
the particularly simple allocations
\[
 \epsilon uvQ\xi-N
   =-\frac{\epsilon}{s}u\phi\mathcal E^2
     -\frac{\epsilon uv}{\sqrt B}PN_x.               \tag{CJ6a}
\]
Their coefficients are infinitesimal, since $u,v$ are local
infinitesimals and $P$ has a bounded source-root allocation.

In these roots CJ5 is exactly
\[
\begin{split}
 N={}&\frac{\epsilon}{s}\frac{\sqrt s}{\sqrt\Lambda}
                  \phi\mathcal E\widehat J\\
 &+\epsilon\sqrt{\eta/s}\,u\phi\mathcal E\mathcal P
   +\frac{\epsilon}{s}
       \frac{x\sqrt s}{\sqrt a\sqrt B}\mathcal E N_x .
\end{split}                                                    \tag{CJ7}
\]
All three coefficients are infinitesimal bounded tensor elements.
For the first two use $\epsilon/s$ bounded, $s$ infinitesimal,
$\eta/s\le1$, and bounded
$\phi=X/Z+hu$.  The last coefficient is also infinitesimal:
\[
 \frac{x\sqrt s}{\sqrt a}=
 \begin{cases}
 O(x\sqrt s),&a\text{ a bounded positive unit},\\
 \sqrt a\,L\sqrt s,&x=aL,\ a\text{ infinitesimal}.
 \end{cases}                                                   \tag{CJ8}
\]
Also $1/\sqrt B$ is bounded, even at literal $\beta=0$ and when
$B$ is unbounded.  The apparent inverse $\sqrt a$ is canceled in
the large-$\zeta$ case before taking bounds.  Every coefficient is
a finite separated expression; $\phi$ is expanded as $X/Z+hu$.
No source rate is compared with any testing rate.

\paragraph{Full local reassembly.}
Add the already present row $\widehat J^2$ to the strict frame
CJ2 and allocate CJ7 to it and the source roots.  This produces an
infinitesimal Gram perturbation of the previous strict-residue
matrix plus a unit $\widehat J$ diagonal entry.  The separated
bounded-Gram lemma therefore applies to the full remainder.
Retain each other positive row from DG3/LD9 exactly once:
\[
 \widehat R^2,\quad c_UUQ^2,\quad e_4UF_t^2,\quad BV\xi^2,
 \quad(2x/c)uv\tau^2,\quad(G_0x/c)D_0^2WZ.
\]
The last two are literal tensor squares on the unchanged sign chart
$u,v\ge0$, $WZ>0$.  The complete angular and radial terms remain
inside CJ2; all full-$D$, full-$Y$ and higher normals are those of
the exact original-parent identity, not truncated expansions.

Consequently the entire local original parent is tensor SOS when
$Z\ne0$ and $X/Z$ is bounded, with no condition on $s/\eta$,
$u^2W/Z$, or $Y/Z$.  The proof is valid on literal local angle axes
because it uses no angle inverse.  At original contact $E,P,Q,\xi$
and $\widehat J$ vanish, so this reassembly preserves equality and
all original contact normals exactly.

\paragraph{The local state and complementary testing cases.}
The reviewed pure-testing half-core exit covers $WZ\le0$, including
$Z=0$.  When $Z\ne0$ and $X/Z$ is unbounded, DS3 in Section~\ref{component:p8cfada778490} is a direct
original-parent tensor-SOS exit.  The above argument covers every
remaining real state on the local angle chart.  Applying these real
certificates to real and imaginary state parts gives the full complex
parent.  Thus the shared local C14.3a/b obligation is closed, with
all source gaps and $\beta=0$ retained.  The SB remote theorem in Section~\ref{component:p6d2f659c0fde} and the unit-$\eta$ theorem in Section~\ref{component:p60eca5442075} supply the other testing and source charts. The complete source partition is printed in Section~\ref{component:pfd0dd7bb10cf}.



% Source formulas: p66922bb7f235
\section{Vanishing $L$ with unbounded \texorpdfstring{$\eta$}{eta}}
Both $L$ and $\sigma=1/\eta$ now vanish. With $S=L+\sigma$, the three source ratios $\kappa=x/S$, $\nu=\sigma/S$ and $\ell=L/S$ satisfy $\nu+\ell=1$. Their exact relations organize the remaining contacts, including the endpoints at which two of these ratios vanish.

\subsection{The small-$L$, unbounded-$\eta$ source sphere}

\label{component:p66922bb7f235}

\label{p66922bb7f235:sec:r13-small-L-eta-large-bounds}

This is a source theorem, not an additional all-tests certificate.
Use the original physical C1 source with $x,L=x\zeta>0$
infinitesimal, $\zeta$ noninfinitesimal and $\eta$ unbounded.
Put $\sigma=1/\eta$, $S=\sigma+L$, and retain the original
$0<r,h<1$, $\theta=1-h$, $M>0$, and $D,\beta\ge0$.
The native gap gives $E=1-xh/\sigma\ge1/18$. Set
\[
 e={xh\over\sigma}=1-E,\quad
 P=\sigma P_d,\quad J=1+xh-\sigma(1-r),\quad
 H_0=1+L(1-r),\quad B=P+x(J-\sigma),
\]
\[
 \mathcal H=R H_0+(1-r)xh(2+xr)
 -e\{2+x(1+h+r\theta)\}+Ke^2,
 \quad R=1+xr,\quad K=1+x.                           \tag{ES1}
\]
No $E\sim1$ assumption is made here. The original source formulas
give the two exact identities
\[
 2\theta P\beta=hB-\sigma H_0,\qquad
 2P(JD+x\theta\beta)=\sigma^2\mathcal H.             \tag{ES2}
\]
They are the beta-shared identities with the reciprocal-$L$ factors
cleared. In particular their use does not require bounded $1/L$.

The positive denominator is exactly
\[
 P=\sigma(1+L+2xh+rxL+rx^2h)+KL+x^2h\theta.
                                                               \tag{ES3}
\]
Since $x/L=1/\zeta$ is bounded and $x^2h=x\sigma e$,
all terms beyond $\sigma+L$ in (ES3), divided by $S$, are
infinitesimal. Consequently
\[
 {P\over S}\sim1,\qquad J\sim1,\qquad
 {B\over S}-(1+x/S)\in\Inf,\qquad
 H_0\sim1,\qquad \mathcal H-E^2\in\Inf.             \tag{ES4}
\]
Here $x/S$ is bounded. Both $B/S$ and $\mathcal H$ are
positive bounded units, uniformly through all original openings;
for the latter use the actual bound $E\ge1/18$.

Define nonnegative pure-source sphere coordinates
\[
 d_0^2={2PJD\over\sigma^2\mathcal H},\qquad
 b_0^2={2Px\theta\beta\over\sigma^2\mathcal H}.
\]
Then
\[
 d_0^2+b_0^2=1,\qquad
 D={\sigma^2\mathcal H\over2PJ}d_0^2,\qquad
 x\theta\beta={\sigma^2\mathcal H\over2P}b_0^2.
                                                               \tag{ES5}
\]
Thus $D=O(\sigma^2/S)=O(\sigma)$ and
$x\theta\beta=O(\sigma^2/S)$, including either literal coefficient
boundary. The first identity of (ES2) becomes
\[
 eB=xH_0+\sigma\mathcal H b_0^2.                     \tag{ES6}
\]
Put $\kappa=x/S$ and $\nu=\sigma/S$. Multiplying
$e(B/S)=\kappa H_0+\nu\mathcal H b_0^2$ by $e$ and using
$e\nu=\kappa h$ gives
\[
 e^2(B/S)=\kappa\{eH_0+h\mathcal H b_0^2\}.
                                                               \tag{ES7}
\]
All factors in braces are bounded and nonnegative. Hence
$e^2=O(x/S)$, and (ES6) also gives $x/S=O(e)$.
In particular $\zeta$ unbounded implies $e$ infinitesimal,
since $x/S\le1/\zeta$. The unit-$\zeta$ cases may have
$e$ a positive bounded unit.

The other original coefficients have the opening-cleared bounds
\[
 {CS(1-r)\over E}\sim\tfrac12,\qquad
 {ASr\over E}\sim\tfrac12,\qquad
 m={K(1+\sigma r)eS\over2M\zeta P},                  \tag{ES8}
\]
\[
 \alpha={(1+\sigma r)\{LR eS+L^2ME\}\over2QP}
       =O(L),\qquad Q=1-Lr\sim1.                   \tag{ES9}
\]
For $A$ use $P_4=(1+xh)E+(1-r)(xh+L)\sim E$.
The bound on $\alpha$ follows because $0<M<R$ is bounded and
$L/S\le1$. Thus $C,A$ are unbounded while $D,\alpha$ are
infinitesimal on the entire stated source domain. No bound on $m$
is asserted through $M=0$, and no bound on beta is asserted through
$\theta=0$.

Finally put $\mu=K(1+\sigma r)S/(2M\zeta P)>0$. The
original $m$ term has the exact shared-sphere split
\[
 m=m_0+m_1b_0^2,\qquad
 m_0={\mu xH_0\over B},\qquad
 m_1={\mu\sigma\mathcal H\over B}.                  \tag{ES10}
\]
The remaining source opening $M$ is displayed explicitly. These
identities provide the weights and boundedness estimates for C14.4.
Section~\ref{component:pfd0dd7bb10cf} assembles the full-parent
certificates for the remaining infinitesimal-$u$ tests and the
complementary testing charts. The inequalities above use only the
order in each individual field.



% Source formulas: p8c30238a5f7d
\subsection{The unmatched Bob residue}

\label{component:p8c30238a5f7d}

\label{p8c30238a5f7d:sec:recovery-c144-small-e}

This is a partial-domain theorem. It retains the original physical C14.4
source, all strict gaps, nonnegative event weights, the original total-scale
and native source inequalities, and independent source and testing fields.
Its source has $x,L>0$ infinitesimal, $\zeta=L/x$ noninfinitesimal and
$\eta$ unbounded. Write $\sigma=1/\eta$, $S=\sigma+L$ and
\[
 e=\frac{xh}{\sigma},\quad E=1-e,\quad
 \kappa=\frac{x}{S},\quad \nu=\frac{\sigma}{S},\quad
 \ell=\frac{L}{S},\quad \nu+\ell=1.
 \tag{SE1}
\]
The letter $\kappa$ in this section denotes $x/S$, not the different
coefficient denoted by $\kappa$ in the R19 reciprocal-$L$ chart.
All source inequalities are taken in the coefficient field. Real residues
are compared only after taking residues separately in the two fields.

\paragraph{The exact source input.}
Use ES1--ES10 from the retained R13 source proof. In particular
$1/18\le E<1$, $0<h,r<1$, $\theta=1-h$, and
\[
\begin{gathered}
 J=1+xh-\sigma(1-r),\qquad M=1+xr-\theta/\zeta>0,\\
 Q=1-Lr,\quad H_0=1+L(1-r),\quad B=P+x(J-\sigma),\\
 P=\sigma(1+L+2xh+rxL+rx^2h)+(1+x)L+x^2h\theta,\\
 \mathcal H=(1+xr)H_0+(1-r)xh(2+xr)
       -e\{2+x(1+h+r\theta)\}+(1+x)e^2.
\end{gathered} \tag{SE2}
\]
The opening-safe sphere identities are
\[
\begin{gathered}
 2\theta P\beta=hB-\sigma H_0,\qquad
 d_0^2+b_0^2=1,\quad d_0,b_0\ge0,\\
 D=\frac{\sigma^2\mathcal H}{2PJ}d_0^2,\qquad
 x\theta\beta=\frac{\sigma^2\mathcal H}{2P}b_0^2,\\
 e\frac BS=\kappa H_0+\nu\mathcal H b_0^2,\qquad
 e^2\frac BS=\kappa(eH_0+h\mathcal H b_0^2).
\end{gathered}\tag{SE3}
\]
Here $P/S\sim1$, $J,Q,H_0\sim1$, $B/S-(1+\kappa)$ is
infinitesimal, and $\mathcal H-E^2$ is infinitesimal. In particular
$B/S$ and $\mathcal H$ are positive bounded units. These estimates
are the proved ES bounds; they do not use bounded $\eta$.

\paragraph{Source dichotomy.}
On the entire C14.4 source,
\[
 e\text{ is infinitesimal}\quad\Longleftrightarrow\quad
 \kappa=x/S\text{ is infinitesimal}.\tag{SE4}
\]
Indeed the first nonnegative sum in SE3 gives $\kappa=O(e)$,
and the second gives $e^2=O(\kappa)$. Both implications are
ordered-field statements entirely inside the source field.

Suppose $e$ is infinitesimal. If $\zeta$ is unbounded, then
$\theta/\zeta$ is infinitesimal and $M\sim1$ directly. If $\zeta$
is a positive bounded unit, then the exact relation
\[
                  \ell=\zeta\kappa\tag{SE5}
\]
gives $\ell$ infinitesimal and $\nu\sim1$. Consequently
$B/\sigma\sim1$. Since $\beta\ge0$, the first identity of SE3
implies $hB\ge\sigma H_0$, whence $h\sim1$ and $\theta$ is
infinitesimal. Thus $M\sim1$ in this case as well. We have proved,
uniformly on the whole small-$e$ source,
\[
                    E\sim1,\quad M\sim1,\quad m\in\operatorname{Inf}_{>0}.
                    \tag{SE6}
\]
For the last statement use the exact original coefficient
\[
 m=\frac{(1+x)(1+\sigma r)eS}{2M\zeta P};\tag{SE7}
\]
$1/\zeta$ is bounded and every other coefficient apart from $e$ is
bounded, including the reciprocal of $M$ just proved to be a unit.
No assumption on the ratio of a source infinitesimal to a testing
infinitesimal occurs.

The two nonnegative terms in the first sum in SE3 imply
$\nu b_0^2$ infinitesimal. Hence the exact sphere formula and the
exact formula for $\alpha$ yield
\[
 \frac DS-\frac{\nu^2}{2}\in\operatorname{Inf},\qquad
 \frac\alpha S-\frac{\ell^2}{2}\in\operatorname{Inf},\qquad
 SC(1-r)\sim\tfrac12,\qquad SAr\sim\tfrac12.\tag{SE8}
\]
For clarity the second formula is obtained without dividing by $\ell$:
\[
 \frac\alpha S=
 \frac{(1+\sigma r)\{\ell(1+xr)e+\ell^2ME\}}
      {2Q(P/S)}.\tag{SE9}
\]
These are additive bounded residue estimates, including $D=0$ whenever
that boundary has points in the stated domain. They do not assert a
relative estimate at a vanishing pivot.

\paragraph{Theorem: every unmatched Bob residue on the small-$e$ source.}
Let $e$ be infinitesimal. In the testing field choose normalized homogeneous
Bob coordinates $(c,s)$ with $c^2+s^2=1$. For every independent projective
test and every real or complex state satisfying
\[
                         \overline{s^2}\ne\overline\ell,\tag{SE10}
\]
the complete original C14.4 parent is a finite sum of squares in
$K_c\otimes_{\mathbb R}K_t$.

\paragraph{Proof with the complete parent retained.}
Noninfinitesimal Alice tests use the unchanged adjacent-$C,A$ theorem:
SE8 and $S$ infinitesimal imply that both $C$ and $A$ are unbounded.
For infinitesimal or zero $u$, apply the exact EB3--EB7 homogeneous
frame from the retained R18 proof. Those formulas are algebraic identities;
their only source-limit inputs to the bounded-Gram proof are precisely
$m$ infinitesimal and SE8, now established on the larger source domain.
To make the imported implication explicit, put
\[
 H=S(C+A),\quad \rho=\frac C{C+A},\quad a_h=\frac{SCA}{C+A},
 \qquad p=\frac{uW}{\sqrt S},\ q=\sqrt S Z,\ y=\frac Y{\sqrt S}.
\]
Then $1/\sqrt H$ is bounded, $0<\rho<1$, $a_h\sim1/2$, and
the harmonic normal $N_y=\sqrt H(y-\rho p)$ can be substituted
by $y=\rho p+N_y/\sqrt H$ with bounded source coefficients.
The remaining normals are the EB4 $N_m,N_b,N_r$, with only pure-field
radicals and the positive testing denominators $1+u^2$ and $1+u^2s^2$.
The exact full quarter-parent is EB7,
\[
 N_y^2+a_hp^2+ASu^2y^2+M_B^2+D_B^2+A_B^2+\beta(X+uZ)^2
                         +\tfrac12 K_B.\tag{SE11}
\]
All EB5 core terms and all EB6 event amplitudes are retained. In particular
no state graph is imposed and neither first-normal coupling is removed.
The two positive rows $ASu^2y^2$ and $\beta(X+uZ)^2$ are left whole.
On the other six roots $(N_y,N_m,N_b,N_r,p,q)$ the bounded Gram has
residue equal to $\operatorname{diag}(1,1,1/2,1/2)$ plus the scalar block
\[
 \frac12\begin{pmatrix}
 1&-\bar s\bar c\\
 -\bar s\bar c&\bar\nu^{\,2}\bar s^{\,2}
                            +\bar\ell^{\,2}\bar c^{\,2}
 \end{pmatrix}.
\tag{SE12}
\]
Its determinant is $(\overline{s^2}-\bar\ell)^2/4$, so SE10
makes this finite coefficient Gram strictly positive definite. The
separated bounded-Gram lemma supplies tensor SOS; restoring the two
whole rows gives the complete parent. The exact homogeneous formulas
allow $c=0$, $s=0$ and $u=0$. Signs and real/imaginary state parts are
handled exactly as in EB, and the original positive source and testing
normalizations are pure-field factors. This proves the stated domain.

\paragraph{New part, and exact retained complement.}
Unbounded $\zeta$ already forces $\kappa\le1/\zeta$ infinitesimal,
so the new theorem contains the entire R18 unmatched-residue certificate.
Its additional domain is $\zeta$ a positive bounded unit, $e$
infinitesimal, and Bob sine square with nonzero residue; this includes
Bob infinity. In this added domain $\ell$ has zero residue, so the
remaining Bob chart is exactly finite infinitesimal $v$ (including zero,
which may be handled by independent literal-axis exits).

After the unmatched-residue certificates, the complementary C14.4 domain is contained in the union of (i) $e$ infinitesimal and matched Bob residue, and (ii) $\zeta$ a positive bounded unit and $e$ a positive bounded unit. All original intersections remain in both. In case (ii), SE7 shows that $m$ has a positive real lower bound, while its upper bound may diverge through $M=0$. This is a source estimate. The complete certificates for both complementary domains, with their exact source partition, are joined in Section~\ref{component:pfd0dd7bb10cf}.



% Source formulas: p1a6cc124f5db
\subsection{The homogeneous unmatched-Bob quadratic}

\label{component:p1a6cc124f5db}

\label{p1a6cc124f5db:sec:r18-c144-bob-residue}

\paragraph{Precise certificate.}
Retain the entire original C14.4 source: $x,L>0$ infinitesimal,
$\eta$ unbounded, $\zeta=L/x$ noninfinitesimal, all original physical
intersections, every positive opening, and either permitted literal
$D=0$ or $\beta=0$. This section additionally assumes $\zeta$
unbounded. Put
\[
 \sigma=1/\eta,\qquad S=\sigma+L,\qquad
 \nu=\sigma/S,\qquad \ell=L/S=1-\nu.
\]
For a real projective Bob test choose its normalized homogeneous
coordinates $(c,s)$, so $c^2+s^2=1$; on the finite chart these are
$c=(1+v^2)^{-1/2}$, $s=vc$. The complete original parent is tensor SOS
at every independent test and state satisfying
\[
                         \overline{s^2}\ne\overline\ell.       \tag{EB1}
\]
Bars mean real residues of bounded elements in their respective fields.
This compares real residues only, not independent infinitesimal rates.
The adjacent-$C,A$ theorem already covers noninfinitesimal Alice
tests. Hence only $u$ infinitesimal is proved below. No all-tests
closure of C14.4 is asserted: the matched residue and bounded-unit
$\zeta$ source remain.

\paragraph{Source limits, including both sphere endpoints.}
Use ES1--ES10 without changing their source entry. Since
$x/S\le x/L=1/\zeta$ is infinitesimal, ES7 gives
$e=xh/\sigma$ infinitesimal, and hence $E\sim1$.
Moreover $M=1+xr-\theta/\zeta\sim1$. ES8 then gives $m$ infinitesimal.
The exact equation ES6, divided by $S$, is
\[
 e(B/S)=(x/S)H_0+\nu\mathcal H b_0^2.
\]
Both summands are nonnegative; $B/S$ and $\mathcal H$ tend to one.
Consequently $\nu b_0^2$ is infinitesimal. Therefore ES5 yields
\[
 \frac DS-\frac{\nu^2}{2}\in\operatorname{Inf},\qquad
 \frac\alpha S-\frac{\ell^2}{2}\in\operatorname{Inf},
 \qquad m\in\operatorname{Inf}_{>0}.                         \tag{EB2}
\]
For the second formula, divide the exact ES9 expression by $S$:
\[
 \frac\alpha S=
 \frac{(1+\sigma r)\{\ell R e+\ell^2ME\}}
 {2Q(P/S)}.
\]
All displayed denominator factors tend to one. These are additive
residue assertions, valid also when $\nu$ or $\ell$ is infinitesimal.
In particular no division by either sphere coordinate is used.

The opening-cleared ES8 bounds give
\[
 c_*=SC(1-r)\sim\tfrac12,\qquad
 a_*=SAr\sim\tfrac12.
\]
Define the pure-source harmonic quantities
\[
 H=S(C+A),\qquad \rho=C/(C+A),\qquad
 a_h=SCA/(C+A)=\frac{c_*a_*}{c_*r+a_*(1-r)}.
                                                               \tag{EB3}
\]
Then $1/\sqrt H$ is bounded, $0<\rho<1$, and $a_h\sim1/2$,
uniformly through both $r$ openings. No upper bound on $C$ or $A$
has been assumed.

\paragraph{An exact all-state frame for every Bob chart.}
Let $D_u=1+u^2$ and $A_u=1+u^2s^2$, both positive testing units.
Use the separated roots
\[
 p=\frac{uW}{\sqrt S},\quad q=\sqrt S Z,\quad
 y=\frac Y{\sqrt S},\quad N_y=\sqrt H(y-\rho p),
 \quad N_m=\sqrt m\,sW,
\]
\[
 N_b=\sqrt{A_u}\left(X+\frac{uc^2}{A_u}Z\right),
 \qquad N_r=\frac{\sqrt{D_u}\,usc}{\sqrt{A_u}}Z.
                                                               \tag{EB4}
\]
Only pure-source and pure-testing radicals or inverses occur.
In particular no inverse of $c$, $s$, $u$, or a state coordinate is
needed. The inverse harmonic allocation
$y=\rho p+N_y/\sqrt H$ has bounded coefficients.

Write $b=X+uZ$. The exact original signed core $K/D_v$ in finite
Bob coordinates equals the following homogeneous expression:
\[
\begin{split}
 K_B={}&N_b^2+N_r^2-2D_usc\,pq
       +Sc^2(p-y)^2+SD_us^2y^2\\
 &+\frac{2\sqrt S\,usc}{\sqrt{A_u}}(p-y)N_b
  +\frac{2\sqrt S\,us^2\sqrt{D_u}}{\sqrt{A_u}}(p-y)N_r.
\end{split}                                                    \tag{EB5}
\]
Here $D_us$ denotes $D_u\,s$, and similarly in the other products.
One derivation starts with CE2 divided by $D_v$, sets
$R=usZ$, and completes its entire $(b,R)$ block:
\[
 A_ub^2-2D_us\,bR+D_uR^2=N_b^2+N_r^2.
\]
The cross term from
$[\sqrt S\,c(p-y)+usb]^2$ is allocated by
$cb=cN_b/\sqrt{A_u}+\sqrt{D_u}sN_r/\sqrt{A_u}$.
This proves EB5 without deleting a first state normal.

The original positive event roots have the exact bounded expressions
\[
 \begin{split}
 M_B={}&N_m-\frac{\sqrt m\,c}{\sqrt{A_u}}N_b
       -\frac{\sqrt{mD_u}\,s}{\sqrt{A_u}}N_r
       +\sqrt{mS}\,us\,y,\\
 D_B={}&\sqrt{DS}\,c(p-y)
       +\frac{\sqrt D\,us}{\sqrt{A_u}}N_b
       -\sqrt{D/S}\,\frac{D_us}{A_u}q,\\
 A_B={}&\sqrt{\alpha D_uS}\,s y
       -\sqrt{\alpha D_u/S}\,c q.
 \end{split}                                                   \tag{EB6}
\]
They are, respectively, the normalized original $m,D,\alpha$ event
amplitudes. The formulas are valid at $D=0$ as well. In the $D$ row,
the identity
$usb-D_usZ=usN_b/\sqrt{A_u}-D_usZ/A_u$
avoids an apparent reciprocal-$c$ pole.

Consequently the full original quarter-parent, divided by $D_v$ on
the finite chart, is exactly
\[
 \mathcal P_B=N_y^2+a_hp^2+ASu^2y^2
              +M_B^2+D_B^2+A_B^2+\beta b^2+\tfrac12K_B.
                                                               \tag{EB7}
\]
The two explicit positive rows $ASu^2y^2$ and $\beta b^2$ can be
retained unchanged. All other rows in EB7 use EB4 once. The formula
is a full four-state identity and is not a restriction to a state ray.

\paragraph{Strictness precisely off the matched residue.}
Every coefficient of EB5 except the printed $pq$ term is either
its fixed positive pivot or infinitesimal after EB3. In EB6, the
$m$ row is $N_m$ plus infinitesimal coefficients; the $D$ and
$\alpha$ rows have only their $q$ coefficients at residue.
Use EB2, $A_u,D_u\sim1$, and EB3. The resulting residue Gram on
$(N_y,N_m,N_b,N_r,p,q)$ is the direct sum of
$\operatorname{diag}(1,1,1/2,1/2)$ and
\[
 \frac12
 \begin{pmatrix}
 1&-\bar s\bar c\\
 -\bar s\bar c&\bar\nu^2\bar s^2+\bar\ell^2\bar c^2
 \end{pmatrix}.                                                \tag{EB8}
\]
The determinant of this last block is
\[
             \frac14(\bar s^2-\bar\ell)^2.                    \tag{EB9}
\]
Thus EB1 makes the entire bounded finite Gram strictly positive.
The separated bounded-Gram lemma gives tensor SOS, and restoring the
two retained positive rows proves the complete parent. The roots are
allowed to be dependent or even identically zero at an endpoint; the
claim concerns the coefficient Gram, so this does not affect the
argument.

The same homogeneous formulas allow $c=0$ and $s=0$ literally.
They include every Bob projective endpoint, every sign, and every
local Alice axis. Nonlocal Alice charts use the common adjacent
$C,A$ theorem. The real and imaginary state parts are treated
separately. Every normalization is pure source or pure testing.

\paragraph{The exact remaining obstacle.}
The leading scalar has the useful two-square identity
\[
 \tfrac12p^2-sc\,pq
 +\tfrac12(\nu^2s^2+\ell^2c^2)q^2
 =\tfrac12(p-scq)^2+\tfrac12(s^2-\ell)^2q^2.                  \tag{EB10}
\]
It does not complete the matched-residue chart: there the second
pivot collapses, and the actual source corrections must be allocated
through exact contact conductors. Arbitrary small positive $D$ or
$\beta$ openings do not supply a missing fixed pivot. Unit $\zeta$
also requires new source limits because $e$ can have nonzero residue.
These two obligations are retained, not declared empty.



% Source formulas: pe55ec4890781
\subsection{The common finite matched-Bob angular relation}

\label{component:pe55ec4890781}

\label{pe55ec4890781:sec:new-c144-matched-budget}

This exact auxiliary identity supplies the angular helper budget in FM.
Work in the separated tensor ring $K_c\otimes_{\mathbb R}K_t$.
Let $S,\kappa,\nu,\ell$ be positive source quantities, with
$\nu+\ell=1$, and put
\[
 x=\kappa S,\qquad j=\sqrt{\nu\ell},\qquad
 a_B=\sqrt{\ell/\nu}.
 \tag{MB1}
\]
For independent testing variables $u,v,W,Z$ define
\[
\begin{gathered}
 T=\kappa W,\quad U=uZ,\quad P=uW/\sqrt S,\quad Q=\sqrt S Z,\\
 N=T-jU,\quad G=P-jQ,\quad \delta=v-a_B,\\
 B=\delta Q,\quad B_u=\delta U.
\end{gathered}
 \tag{MB2}
\]
The monomial relation $TU=\kappa PQ$ is exact. It gives
\[
 \delta^2N^2+2j\kappa\delta GB+2j^2\kappa B^2-j^2B_u^2
       =(\delta T)^2.
 \tag{MB3}
\]
In particular, if $\nu,\ell$ are positive bounded units,
$\kappa$ is infinitesimal and $\delta$ has zero real residue,
then any fixed real positive diagonal on $B_u$ can be allocationd
at an infinitesimal perturbation of the Gram matrix on $(N,G,B)$.
Indeed, for any fixed $c>0$ the exact SOS remainder is
\[
 \frac{c}{j^2}\delta^2N^2+
       \frac{2c\kappa}{j}\delta GB+2c\kappa B^2-cB_u^2
       =\frac{c}{j^2}(\delta T)^2.
 \tag{MB4}
\]
Every displayed coefficient is a finite sum of bounded separated
factors, with zero real residue except the constant $c$ in the
desired $B_u$ pivot. The only inverses are of positive source units.
The factor $\delta$ is a difference of a pure testing and a pure
source quantity; no inverse or radical of that difference is taken.
The right side has the displayed separated SOS root
$\sqrt c\,\delta T/j$.

Two exact radial conductors and one cross conductor are
\[
 \begin{split}
 (u^2-x)W/\sqrt S&=uG-\sqrt S N,\\
 (u^2-x)Z&=(\kappa\sqrt S/j)G-(u/j)N,\\
 (u^2-x)\delta Z^2&=(\kappa/j)GB-(1/j)NB_u.
 \end{split}
 \tag{MB5}
\]
These identities require no bound on the state coordinates or on
any ratio of source and testing infinitesimals. In the local matched
chart with $u$ infinitesimal their first two right sides have only
infinitesimal coefficients on $(N,G)$.



% Source formulas: p3f02c51074be
\subsection{Unit $\kappa,\nu,\ell$: complementary Bob charts}

\label{component:p3f02c51074be}

\label{p3f02c51074be:sec:new-c144-unit-kappa-all-tests}

\paragraph{Exact scope and conclusion.}
Retain the original physical C14.4 source, its strict source gaps,
and its native source inequalities. Assume
\[
 S=\sigma+L\in\operatorname{Inf}_{>0},\qquad
 \kappa=x/S,\quad\nu=\sigma/S,\quad\ell=L/S
 \quad\hbox{are positive bounded units},\qquad\nu+\ell=1.
 \tag{FU1}
\]
Together with the matched theorem UKR1--UKR21, the argument below
proves that the complete original parent is a finite separated tensor
SOS for every projective test and every state on FU1. It includes
arbitrary openings of $M,\theta,r,1-r$, literal $D=0$ and $\beta=0$,
and independent source and testing infinitesimals. It makes no claim
at the $\nu$- or $\ell$-infinitesimal source endpoints.

All normalizations are pure-source or pure-testing operations.
All real Schur complements below prove strictness of a real residue
matrix; no such inverse is used as a tensor-ring coefficient. The
strict bounded-Gram lemma is the one in the surviving
Lemma~\ref{p25af2c58da7e:lem:gram}.

\paragraph{The full parent and opening-safe source bounds.}
Write $D_u=1+u^2,D_v=1+v^2$, and divide the parent by $4$:
\[
\begin{split}
 F/4={}&m(vW-X+uvY-uZ)^2+CD_v(uW-Y)^2+AD_uD_vY^2\\
 &+(D+\tfrac12)(uW-Y+uvX-vZ)^2
 +(\alpha+\tfrac12)D_u(vY-Z)^2\\
 &+(\beta+\tfrac12)D_v(X+uZ)^2-\tfrac12D_uD_vZ^2.
\end{split}\tag{FU2}
\]
Keep the exact FM2--FM3 source formulas, with
$c_0=SC(1-r)$, $a_0=SAr$, $d_0=D/S$, and $a_1=\alpha/S$.
These four coefficients and $\beta\theta$ are bounded, with no
inverse of the possibly small $M$ or $\theta$ in their normalized
formulas. The $D/S$ numerator has its factor $S$ canceled exactly.

The following letters refer only to real residues:
\[
 A_* =\nu-\kappa h=\nu E,\quad R=1-A_*,\quad
 \nu=A_*+\kappa h,\quad\ell=1-A_*-\kappa h.
 \tag{FU3}
\]
The exact source yields
\[
\begin{gathered}
 E\ge1/18,\qquad h-A_*\ge0,\qquad
 A_*^2-\kappa(h-A_*)\ge0,\qquad R-\kappa\ge0,\\
 c_0=a_0=E/2,\qquad
 d_0=\{A_*^2-\kappa(h-A_*)\}/2,\\
 a_1=\frac{\ell\{\kappa h+(R-\kappa)A_*\}}{2\nu},
 \qquad 2\beta\theta=h-A_*.
\end{gathered}\tag{FU4}
\]
In particular $h\ge\nu/(1+\kappa)$, so $h,E,a_1,c_0,a_0$
have positive real lower bounds. The actual source has $m$ bounded
below by a positive real, while $C,A$ are unbounded. Here $m$
can be unbounded, and so can $\beta$. They cannot both be unbounded:
$m$ unbounded forces $R=\kappa$, whence
$\ell=\kappa\theta$ makes $\theta$ a unit; $\beta$ unbounded
forces $\theta=0$, whence $R-\kappa=\ell$ makes $M$ a unit.
For bounded $m,\beta$ their real values are
\[
 m=\frac{\kappa^2h}{2\nu(R-\kappa)},\qquad
 \beta=\frac{h-A_*}{2(1-h)}.
 \tag{FU5}
\]
These formulas are used only when the indicated real denominators
are positive. Every estimate below for the actual parent uses
bounded normalized source coefficients before taking residues.

\paragraph{A complete finite-Bob frame, including the state normals.}
Assume $u$ is infinitesimal and $v$ is finite. We may cap $m$ and
$\beta$ by finite pure-source values $m_c>0,\beta_c\ge0$ and
retain the complete excesses in their original event squares.
When a weight is bounded, keep it unchanged. When it is unbounded,
the sufficiently large real cap is chosen below. Define
\[
 B_0=X+uZ,\quad N_y=\frac{Y-ruW}{\sqrt{Sr(1-r)}},\qquad
 T=\kappa W,\ U=uZ,\ P=\frac{uW}{\sqrt S},\ Q=\sqrt S Z.
 \tag{FU6}
\]
All six rows are legal separated expressions, and $TU=\kappa PQ$
holds exactly. No restriction is placed on their values or on state
ratios. Set $a=u^2$, $s_y=\sqrt{Sr(1-r)}$, and
\[
\begin{split}
 M_0&=(v/\kappa)(1+ra)T-B_0,\\
 A_P&=D_v\{c_0(1-r)+a_0r(1+a)\}\\
 &\hspace{5mm}+S\{(D+\tfrac12)(1-r)^2
       +(\alpha+\tfrac12)(1+a)r^2v^2\}.
\end{split}\tag{FU7}
\]
The entire capped parent is $F_0+2L_yN_y+H_yN_y^2$, where
\[
\begin{split}
 F_0={}&m_cM_0^2+A_PP^2+(1+a)\{d_0(1+a)v^2+a_1\}Q^2
                   +\tfrac12(1+a)v^2U^2\\
 &-2(D+\tfrac12)(1+a)v^2B_0U
 -(1+a)v\{1+2D(1-r)+2\alpha r\}PQ\\
 &+\{(\beta_c+\tfrac12)D_v+(D+\tfrac12)av^2\}B_0^2
 +2(D+\tfrac12)(1-r)uv\sqrt S\,B_0P,
\end{split}\tag{FU8}
\]
\[
\begin{split}
 L_y={}&m_cuvs_yM_0
 +\sqrt{r(1-r)}D_v\{-c_0+a_0(1+a)\}P\\
 &+S\sqrt{r(1-r)}\{-(D+\tfrac12)(1-r)
          +(\alpha+\tfrac12)(1+a)rv^2\}P\\
 &-(D+\tfrac12)uvs_yB_0
       +\sqrt{r(1-r)}v(1+a)(D-\alpha)Q,\\
 H_y={}&rc_0D_v+(1-r)a_0(1+a)D_v\\
 &+Sr(1-r)\{m_cu^2v^2+D+\tfrac12
                      +(\alpha+\tfrac12)(1+a)v^2\}.
\end{split}\tag{FU9}
\]
These are exact identities in the two unrestricted original state
coordinates $X,Y$; $B_0$ is not set to zero. Their coefficients are
bounded across all source openings. In $L_y$ every coefficient has
zero residue, since $c_0,a_0$ have the same residue, and the other
terms contain $u$ or $S$. The $N_y$ pivot has residue $ED_v/2>0$.
Thus it suffices to make the remaining five-row real residue strict.

\paragraph{Finite nonzero unmatched Bob residue.}
Take $v>0$ with a positive finite residue different from
$b=\sqrt{\ell/\nu}$. The complete real residue of FU8--FU9 is
\[
\begin{split}
 &m_c\{(v/\kappa)T-B_0\}^2+qD_vB_0^2-v^2B_0U
       +\tfrac12v^2U^2-vPQ\\
 &\hspace{8mm}+\tfrac E2D_vP^2+(d_0v^2+a_1)Q^2
             +\tfrac E2D_vN_y^2,\qquad q=\beta_c+\tfrac12.
\end{split}\tag{FU10}
\]
For the moment use the uncapped finite values $m,\beta$ from FU5.
Eliminating $B_0$ in this real matrix, and using $TU=\kappa PQ$
only to name the adjusted cross coefficient, gives
\[
 a_TT^2+a_UU^2+a_PP^2+a_QQ^2-2fPQ,
 \tag{FU11}
\]
where, for $c=v^2$ and $H=m+q(1+c)$,
\[
\begin{gathered}
 a_T=\frac{mq(1+c)c}{\kappa^2H},\quad
 a_U=\frac{c(2H-c)}{4H},\quad a_P=E(1+c)/2,\quad
 a_Q=d_0c+a_1,\\
 f=\frac{v}{2}(1+mc/H).
\end{gathered}\tag{FU12}
\]
The replacement of the real Schur cross term $TU$ by $\kappa PQ$
is implemented in the actual bounded Gram by a real multiple of
the exact zero $TU-\kappa PQ$. The elimination of $B_0$ itself
is only a real strictness test. Every diagonal is positive,
including when $D$ or $\beta$ is zero.
The reduced determinant in a formal positive real variable $y$ is
\[
 \Delta(y,c)=(a_T\kappa^2+a_Py)(a_Uy+a_Q)-f^2y.
 \tag{FU13}
\]
It obeys $4H\Delta/(1+c)=\mathcal P(y,c)$ with
\[
 \mathcal P=p_{22}y^2c^2+p_{21}y^2c+p_{12}yc^2
              +p_{11}yc+p_{10}y+p_{02}c^2+p_{01}c,
 \tag{FU14}
\]
\[
\begin{array}{ll}
 p_{22}=E\beta,&p_{21}=E(m+q),\\
 p_{12}=2m\beta+2qEd_0,&
 p_{11}=-m-q+2E\{(m+q)d_0+qa_1\},\\
 p_{10}=2(m+q)Ea_1,&p_{02}=4mqd_0,\quad p_{01}=4mqa_1.
\end{array}\tag{FU15}
\]
All coefficients except possibly $p_{11}$ are nonnegative, and
$p_{21},p_{10},p_{01}$ are strictly positive. Direct substitution of
FU3--FU5 gives the exact three contact identities
\[
 \mathcal P(\kappa,\ell/\nu)=
 \mathcal P_y(\kappa,\ell/\nu)=
 \mathcal P_c(\kappa,\ell/\nu)=0.
 \tag{FU16}
\]
These finite rational identities can also be checked without
radicals. For positive real $y,c$,
\[
 \frac{\mathcal P(y,c)}{yc}
 =p_{22}yc+p_{21}y+p_{12}c+p_{11}
                     +p_{10}/c+p_{02}c/y+p_{01}/y.
 \tag{FU17}
\]
As a real function of $(\log y,\log c)$, this is strictly convex:
each nonconstant summand has nonnegative Hessian, and
$p_{21}y+p_{10}/c$ alone has positive definite Hessian. FU16 says
its value and gradient vanish at the contact. It is consequently
positive at every other positive point. This real argument is only
a way to determine a finite residue matrix, not a use of logarithms
or mixed inverses in the tensor ring.

Since $c\ne\ell/\nu$, FU13 is positive for all $y>0$.
Minimizing its quotient by $y$ gives
\[
 f<\kappa\sqrt{a_Ta_U}+\sqrt{a_Pa_Q}.
 \tag{FU18}
\]
Indeed the minimum of
$a_Pa_Uy+a_T\kappa^2a_Q/y$ is
$2\kappa\sqrt{a_Ta_Ua_Pa_Q}$, attained at a positive $y$;
equality in FU18 would give a zero of FU13. Choose a real
$\lambda$ with $|\lambda|<\sqrt{a_Ta_U}$ and
$|f-\kappa\lambda|<\sqrt{a_Pa_Q}$. Adding the exact zero
$-2\lambda(TU-\kappa PQ)$ makes both $2\times2$ blocks in
FU11 strictly positive. Restoring the positive real $B_0$ pivot
and the $N_y$ pivot proves strictness of the complete six-row
residue. All actual Gram entries were bounded in FU8--FU9, so
the strict bounded-Gram lemma gives the tensor SOS.

\paragraph{Unbounded $m$ or $\beta$ in the unmatched chart.}
The preceding strictness persists under a sufficiently large
finite cap in the one possible unbounded weight. This assertion
has an opening-safe direct verification. If $m\to\infty$,
divide FU14--FU15 by $m$ before taking the real limit. Then
\[
 (p_{22},p_{21},p_{12},p_{11},p_{10},p_{02},p_{01})/m
 \longrightarrow
 (0,E,2\beta,-1+2Ed_0,2Ea_1,4qd_0,4qa_1).
 \tag{FU19}
\]
At $R=\kappa$, the three contact identities FU16 remain exact
for this limit polynomial. The strict convexity proof FU17 still
applies because its $y$ and $1/c$ coefficients are positive.
The Schur coefficients tend to
\[
 a_T\to q(1+c)c/\kappa^2,\qquad
 a_U\to c/2,\qquad f\to v(1+c)/2,
 \tag{FU20}
\]
with $a_P,a_Q$ unchanged. Therefore FU18 is strict at the limit,
and continuity makes it strict for a sufficiently large finite
real $m_c$.

If $\beta\to\infty$, divide instead by $q$ before taking the
limit at $h=1$. The corresponding coefficient vector is
\[
 (E,E,2m+2Ed_0,-1+2E(d_0+a_1),
                         2Ea_1,4md_0,4ma_1),
 \tag{FU21}
\]
again with the three exact contact identities and the two strictly
positive convexity coefficients. Now
$a_T\to mc/\kappa^2$, $a_U\to c/2$, $f\to v/2$.
A sufficiently large finite $\beta_c$ therefore works.
All the other normalized actual coefficients have their real
residues with bounded unit denominators; no $O(S)$ term was
divided by $M$ or $\theta$. Restore the complete excesses
$(m-m_c)(vW-X+uvY-uZ)^2$ and
$(\beta-\beta_c)D_v(X+uZ)^2$. These caps are used only on the
unmatched chart and need not vanish at the actual contact.

\paragraph{Infinitesimal or zero Bob coordinate.}
If $v$ is infinitesimal, use FU8--FU9 with the six rows
\[
 (B_0,T_v,U_v,P,Q,N_y),\qquad T_v=vT,\quad U_v=vU.
 \tag{FU22}
\]
Cap $m$ at any sufficiently small fixed positive real $m_c\le m$,
and discard no excess. One may set $\beta_c=0$ and retain its
whole event. The exact coefficients remain bounded: the $B_0U$
coefficient is now a multiple of $vB_0U_v$, and $M_0$ becomes
$(1+ra)T_v/\kappa-B_0$. Their real residue is
\[
 m_c(T_v/\kappa-B_0)^2+\tfrac12B_0^2+\tfrac12U_v^2
                      +\tfrac E2(P^2+N_y^2)+a_1Q^2>0.
 \tag{FU23}
\]
This also covers literal $v=0$. The rows need not be independent
coordinates on the state space.

\paragraph{Infinitesimal reciprocal Bob coordinate.}
Put $w=1/v$ and use the homogeneous parent $w^2F$ at $w=0$ as
well as for nonzero infinitesimal $w$. The complete signed core
in the quarter-parent is $K_w/2$, with
\[
\begin{split}
 K_w={}&u^2H^2-2u^2HY-2uHW+2uwHX+u^2w^2W^2\\
 &+2u^2wWX-2uw^2WY+(1+u^2+w^2)(X^2+Y^2)-2uwXY,
 \qquad H=wZ.
\end{split}\tag{FU24}
\]
This is the exact CE6 full core, so all state normals are present.
FU4 gives the real alternative
\[
 d_0+\kappa\beta\theta=(\nu E)^2/2>0.
 \tag{FU25}
\]
Thus either $D/S$ has a positive real lower bound, or $\beta$
has one. Both cases allow independent rates for $u,w,S$.

If $\beta$ is strong, use the four rows $(W,X,Y,U)$, $U=uZ$,
in FU24 by replacing $uH$ with $wU$. Every coefficient is bounded
and the core residue is $X^2+Y^2$. The exact $m$ event is
$(W-wX+uY-wU)^2$, and the $\beta$ event is
$D_w(X+U)^2$. Positive finite caps give the strict residue
\[
 m_cW^2+\tfrac12(X^2+Y^2)+\beta_c(X+U)^2>0.
 \tag{FU26}
\]
All other event terms and both excesses remain whole.

If $D/S$ is strong, choose positive real $c,d,m_c$ with
$c\le c_0,a_0$, $d\le D/S$, and $m_c\le m$. Select exactly
\[
 C_c=c/[S(1-r)],\qquad A_c=c/(Sr),\qquad D_c=dS.
 \tag{FU27}
\]
Retain all excesses, including the whole $A_cu^2D_wY^2$.
Writing $Y=ruW+\sqrt{Sr(1-r)}N_y$, the other two selected
$C,A$ terms have the exact harmonic identity
\[
 C_cD_w(uW-Y)^2+A_cD_wY^2=cD_w(P^2+N_y^2).
 \tag{FU28}
\]
Use six rows $(W,V,X,P,N_y,Q)$ with $V=uwZ$ and the exact
monomial identity $VW=wPQ$. Replacing $uH$ by $V$ in FU24,
the potentially non-small cross term $-2VW$ is exactly $-2wPQ$.
Every other off-diagonal core coefficient is infinitesimal after
$Y=\sqrt S\{rP+\sqrt{r(1-r)}N_y\}$; its residue is $V^2+X^2$.
The selected $m$ and $D$ events are the complete squares
\[
 m_c(W-wX+uY-V)^2,\qquad
 d\{uw\sqrt S W+u\sqrt S X-w\sqrt S Y-Q\}^2.
 \tag{FU29}
\]
The full selected residue is strictly positive:
\[
 m_c(W-V)^2+\tfrac12(V^2+X^2)+c(P^2+N_y^2)+dQ^2>0.
 \tag{FU30}
\]
The strict bounded-Gram lemma applies to the exact complete
forms, and all unselected positive terms are restored. Both
reciprocal proofs include literal $w=0$ and $u=0$.

\paragraph{Exhaustion of tests and exact joins.}
If $u$ is noninfinitesimal, $C,A$ are unbounded by FU4. The
two-adjacent-events theorem CE1--CE10 in
Section~\ref{component:p89fd3edb3edd}
therefore applies to the original full parent and all states.
Its hypotheses require only nonnegative event weights and these
two unbounded coefficients; no small-$\kappa$ source statement
is imported.

For infinitesimal $u$, every positive projective $v$ is either
infinitesimal, a finite unit, or has infinitesimal reciprocal.
FU22--FU30 cover the two end cases. A finite unit has either the
matched residue, covered by UKR1--UKR21, or a different residue,
covered by FU10--FU21. The transformations
\[
 (u,Y,Z)\mapsto(-u,-Y,-Z),\qquad
 (v,X,Z)\mapsto(-v,-X,-Z)
 \tag{FU31}
\]
preserve FU2 exactly and handle both signs. Literal axes and
projective endpoints were included in the displayed raw frames.
Complex states split into real and imaginary parts. No proof in
this exhaustion bounds, normalizes, or relates the testing-state
ratios $W/Z$, $X/Z$, or $Y/Z$.

Consequently every testing chart is covered on exactly FU1,
with every original source predicate retained. The remaining
source endpoints are separate from this conclusion.



% Source formulas: pd26b89cdb3f0
\subsection{Unit $\kappa,\nu,\ell$: the matched Bob residue}

\label{component:pd26b89cdb3f0}

\label{pd26b89cdb3f0:sec:new-c144-unit-kappa-matched}

\paragraph{Scope.}
This theorem retains the entire original physical C14.4 source,
including its strict gaps and native source inequalities. Assume
$S\in\operatorname{Inf}_{>0}$ and $\kappa=x/S,\nu=\sigma/S,
\ell=L/S$ are positive bounded units, with $\nu+\ell=1$.
The testing chart is $u\in\operatorname{Inf}$ and finite $v$
whose absolute value has the same residue as
$b=\sqrt{\ell/\nu}$. All real or complex states are allowed.
The unmatched Bob tests on this source are proved in Section~\ref{component:p3f02c51074be}; the source endpoints are separate cells in Section~\ref{component:pfd0dd7bb10cf}. The present proof is confined to its stated matched-Bob chart and uses no mixed inverse.

The dictionary FM2--FM3 remains an exact source identity on this
larger domain, but its small-$\kappa$ estimates FM4 are replaced
below. Set $j=\sqrt{\nu\ell}$ and retain the rows
\[
 N_x=X+huZ,\quad N_y=(Y-ruW)/\sqrt{Sr(1-r)},\qquad
 T=\kappa W,\ U=uZ,\ P=uW/\sqrt S,\ Q=\sqrt S Z,
\]
\[
 N=T-jU,\quad G=P-jQ,\quad \delta=v-b,\quad
 B=\delta Q,\quad B_u=\delta U,\quad R_Z=(u^2-x)Z.
 \tag{UKR1}
\]
The radial and cross identities FM8 hold unchanged. Their
coefficients on $R_Z$ remain infinitesimal because $S,u$ are
infinitesimal and $\kappa,1/j$ are bounded.

\paragraph{A sharper exact angular budget.}
For any fixed real $0<\rho<1$, the monomial identity $TU=\kappa PQ$
gives
\[
 \frac{\delta^2}{j^2(1-\rho^2)}N^2+
 \frac{2\kappa\delta}{j(1+\rho)}GB+
 \frac{2\kappa}{1+\rho}B^2-B_u^2
   =\frac{(\delta T-\rho jB_u)^2}{j^2(1-\rho^2)}.
 \tag{UKR2}
\]
Its first two coefficients have zero residue. Thus a fixed
$cB_u^2$ pivot costs $q\kappa cB^2$ in residue, where
$q=2/(1+\rho)>1$ can be chosen arbitrarily close to one.
Every inverse in UKR2 is a pure source unit or a positive real.
There is no inverse of $\delta$.

\paragraph{The actual source through both coefficient openings.}
The source has $1/18\le E=1-\kappa h/\nu<1$ and
$\kappa$ is bounded because $1/\zeta$ is bounded. The exact
$\beta\theta$ identity gives a positive real lower bound on $h$.
Consequently $m$ has a positive real lower bound, although it may
be unbounded. Let
\[
 R_s=1+xr,\quad p=\frac{\kappa\theta}{\ell R_s}\in(0,1),\qquad
 \mu_c=m(1-p)=
 \frac{\kappa^2h(1+x)(1+\nu Sr)}{2\nu\ell P_0R_s}.
 \tag{UKR3}
\]
The quantity $\mu_c$ is a positive bounded unit through $M\to0$.
All of $SC(1-r),SAr,D/S,\alpha/S,\beta\theta$ are bounded
without inverses of $r,1-r,\theta,D,\beta$; this follows directly
from FM3 after the explicit factor $S$ in the $D/S$ numerator
is canceled. Their limits are evaluated below.

Only in the real residue calculation use the abbreviations
\[
 A=\nu-\kappa h=\nu E,\quad R=1-A,\quad d=h-A.
 \tag{UKR4}
\]
Here $A,R,h,\nu,\ell,\kappa$ are positive units, and the
original nonnegative weights and positive $M$ imply
\[
 d\ge0,\quad A^2-\kappa d\ge0,\quad R-\kappa\ge0,
 \qquad \nu=A+\kappa h,\quad\ell=1-A-\kappa h.
 \tag{UKR5}
\]
Indeed in residue $2\beta\theta=d$, $2D/S=A^2-\kappa d$,
and $\ell M=R-\kappa$. These are additive residue identities;
no relative error estimate is asserted at a vanishing opening.

\paragraph{Retain unbounded weights in actual contact-zero rows.}
Define
\[
 V=v(W+uY),\quad B_0=X+uZ,\quad N_m=B_0-pV.
 \tag{UKR6}
\]
At the actual source contact $u=\sqrt x$, $v=b$,
$W/Z=j\sqrt{S/\kappa}$, $X=-h\sqrt x Z$, $Y=r\sqrt x W$,
both $N_m$ and $N_x$ vanish. This is an identity at the source
contact, not an equation imposed on independent tests.

If $m$ is bounded, set $m_c=m$. If it is unbounded, choose a
positive real $m_c$ sufficiently large, as justified below.
Similarly keep $\beta_c=\beta$ when $\beta$ is bounded, and
choose a sufficiently large positive real $\beta_c$ when it is
unbounded. An unbounded positive field element exceeds either
fixed real cap. Retain the entire positive excesses
\[
 (m-m_c)N_m^2+(\beta-\beta_c)D_vN_x^2.
 \tag{UKR7}
\]
The remaining parent is exact, with its $m$ part replaced by
\[
 m_0 B_0^2-2m_1B_0V+m_2V^2,\qquad
 m_0=m_c,\quad m_1=\mu_c+m_cp,\quad
 m_2=\mu_c(1+p)+m_cp^2.
 \tag{UKR8}
\]
This follows from the identity
$m(V-B_0)^2-(m-m_c)(B_0-pV)^2$
and $\mu_c=m(1-p)$. Every coefficient of UKR8 is bounded.
The residual $\beta$ terms are exactly
\[
 \beta_cD_vN_x^2+2(\beta\theta)D_vN_xuZ
                         +(\beta\theta)\theta D_vu^2Z^2.
 \tag{UKR9}
\]
All coefficients of UKR9 are bounded. The complete cross terms
are retained, and both original excesses in UKR7 are restored
once at the end.

\paragraph{The complete bounded scalar and normal frame.}
Apply the FM9--FM19 polynomial construction to this capped
parent with the following explicit replacements: in $A_T$ replace
$m/\kappa^2$ by $m_2/\kappa^2$; in $A_U$ replace its
$m\theta^2$ term by $m_0\theta^2$; in $K$ replace its
$m\theta(1+ra)$ term by $m_1\theta(1+ra)$. In the complete
first-normal formulas replace
\[
\begin{array}{c|c}
 \text{entry}&\text{replacement}\\ \hline
 H_x\text{ first term}&m_0\\
 H_y\text{ term }mu^2v^2&m_2u^2v^2\\
 H_{xy}\text{ term }m&m_1\\
 L_x\text{ coefficient }-m/\kappa&-m_1/\kappa\\
 L_U\text{ terms }m\theta-(mj/\kappa)v(1+ra)&
                  m_0\theta-(m_1j/\kappa)v(1+ra)\\
 Y_P\text{ term }mv^2(1+ra)&m_2v^2(1+ra)\\
 Y_Q\text{ term }-m\theta a&-m_1\theta a.
\end{array}\tag{UKR10}
\]
Include $\beta_cD_v$ in $H_x$ and omit the separately retained
$\beta D_vN_x^2$ term of FM15. All other entries are unchanged.
These replacements follow by expanding UKR8--UKR9 in the two
unrestricted state normals. They are not a scalar-only operation.

The graph scalar and both complete first normals still vanish to
the orders required by FM11 and FM19: the subtracted rows in UKR7
vanish at the actual source contact. Their squares therefore
change neither the value nor any first derivative there.
The same finite monic divisions supply the exact bounded FM13
scalar and FM18--FM19 full normal conductors. The $Y$ normal has
zero-residue cross coefficients, since $SC(1-r)$ and $SAr$
both have residue $E/2$ and the remaining $Y$ coefficients
contain $S$ or $u^2$. Its diagonal residue is $E/(2\nu)>0$.

The only nontrivial real residue uses $(N_x,N,G,B,B_u)$.
Writing
\[
 H_* = h-\tfrac12+\tfrac\theta2(\ell+\kappa h),\qquad
 r_* = H_*-j^2m_2/\kappa^2,
 \tag{UKR11}
\]
its exact entries are
\[
\begin{gathered}
 H_c=\begin{pmatrix}
 m_0+(\beta_c+1/2)/\nu&-(m_1/\kappa)b\\
 -(m_1/\kappa)b&(m_2/\kappa^2)b^2
 \end{pmatrix},\qquad A_G=E/(2\nu),\\
 z=\begin{pmatrix}-j(m_1/\kappa+E)\\-r_*/\nu\end{pmatrix},\qquad
 g_*=E\ell-m_1\theta-\tfrac12+2m_2\ell/\kappa
                                        +\kappa r_*/\nu,\\
 H_B=\nu(m_1\theta+\tfrac12).
\end{gathered}\tag{UKR12}
\]
Thus the residue form is
$ (N_x,N)H_c(N_x,N)^t+A_GG^2+2g_*GB+H_BB^2
 +2z^t(N_x,N)^tB_u$ before the UKR2 pivot is added.
In particular the $\kappa r_*/\nu$ contribution to $g_*$ is
retained; it is not infinitesimal in the present source sector.

\paragraph{The exact real contact curvature and all source faces.}
First take an interior real source, $\theta>0$ and $R>\kappa$,
without capping, so
\[
 m=\frac{\kappa^2h}{2\nu(R-\kappa)},\qquad
 \beta=\frac{h-A}{2\theta},\qquad
 m_0=m_1=m_2=m.
 \tag{UKR13}
\]
The two hard pivots and the angular pivot obey the exact identities
\[
\begin{split}
 (N_x,N)H_c(N_x,N)^t
   &=\frac m{\kappa^2}(bN-\kappa N_x)^2
                                  +\frac R{2\nu\theta}N_x^2,\\
 J_*&=H_B-\frac{g_*^2}{A_G}-\kappa z^tH_c^{-1}z\\
 &=\frac{AR+\kappa(h-A)}{2AhR}
                  \{4A^2hR-\kappa(Rh-A)^2\}.
\end{split}\tag{UKR14}
\]
All inverses in UKR14 are used only in a real finite matrix
calculation to prove strictness. They are not inverses in the
tensor ring. The crucial bracket has the positive decomposition
\[
\begin{split}
4A^2hR-\kappa(Rh-A)^2={}&A^3R(4-A)+A^4(R-\kappa)\\
 &+A^2R(h-A)(3+A+2\kappa)\\
 &+R^2(h-A)\{A^2-\kappa(h-A)\}.
\end{split}\tag{UKR15}
\]
Its first term is a positive unit and all remaining terms are
nonnegative by UKR5. This includes the literal $D=0$ and
$\beta=0$ faces and both limiting openings.

For completeness the strictness is uniform toward those openings.
Put $\mu=m/\kappa^2$ and retain $H_*$ from UKR11 with uncapped
coefficients. Then the two entries of $H_c^{-1}z$ are
\[
 z_x=-\frac{2\nu\theta}{R}
                   \left(jE+\frac{\kappa H_*}{\nu b}\right),
 \qquad z_N=\frac\kappa b z_x+\nu-\frac{H_*}{\mu\ell}.
 \tag{UKR16}
\]
They stay bounded. Also
\[
 g_*=-\frac{2A^2+A h\kappa+A\kappa-A-h\kappa}{2\nu}
 \tag{UKR17}
\]
stays bounded, while $A_G$ and $J_*$ have positive real lower
bounds near any fixed allowed residue point. UKR14 gives a
positive real lower bound on $H_c$ because $\mu$ and
$R/(2\nu\theta)$ have positive lower bounds and $b$ is a unit.
Completing the real source quadratic form with the bounded
coefficients UKR16--UKR17 therefore gives a uniform positive
lower bound on
\[
 Q_{\rm phys}=(N_x,N)H_c(N_x,N)^t+A_GG^2+2g_*GB+H_BB^2
                         +2\sqrt\kappa\,z^t(N_x,N)^tB.
 \tag{UKR18}
\]
This is an auxiliary real matrix used to choose a diagonal budget.
It does not impose $B_u=\sqrt\kappa B$, or any relation between
source and testing angles, in the tensor certificate.

\paragraph{Why finite caps preserve a strict real residue.}
When $m,\beta$ are bounded no cap is needed and UKR18 is already
strict. If $m$ is unbounded, $R-\kappa=0$ in residue; since
$\ell$ is a unit and $\ell=R-\kappa h$, $\theta$ is then a unit,
so $\beta$ is bounded. If $\beta$ is unbounded, $\theta=0$ in
residue, and $R-\kappa=\ell>0$, so $m$ is bounded. Thus at
most one unbounded penalty is present.

For the $m$ opening, the real contact-zero linear form in UKR18 is
\[
 n_m=N_x-(pb/\kappa)N-(pj/\sqrt\kappa)B,
 \tag{UKR19}
\]
and $p\to1$. Its $N$ coefficient is a unit. UKR8 shows that
subtracting the whole $mn_m^2$ leaves a bounded coefficient
matrix with a well-defined limit at the opening. That limit is
strict on $\ker n_m$: approach the residue point by real
$\kappa<R$ with $A,h$ fixed, adjust any kernel vector continuously
to the exact kernel using the unit $N$ coefficient, and apply the
uniform lower bound from UKR18. The parameter inequalities UKR5
are preserved when $\kappa$ is decreased. A real symmetric matrix
strict on the kernel of a nonzero linear form becomes strictly
positive definite after adding a sufficiently large fixed real
multiple of the form's square; this follows by a real block Schur
complement. Choose $m_c$ that large in UKR7.

For the $\beta$ opening the same argument uses the exact penalty
$(\beta/\nu)N_x^2$ and its kernel $N_x=0$. Approach $h=1$ from
below with $A,\kappa$ fixed, which also preserves UKR5, and choose
$\beta_c$ sufficiently large. This is a calculation on real
residue matrices, not a mixed-field limiting certificate. The
actual capped coefficients have those residues because UKR3,
UKR8 and UKR9 have bounded denominators through the original
independently varying $S,M,\theta$ openings. In particular no
estimate divides an $O(S)$ error by $M$ or $\theta$.

We have proved that the actual caps can be chosen so that the
capped matrix UKR18 is strictly positive definite. It follows
that $H_c\succ0$, $A_G>0$, and
\[
 H_B-g_*^2/A_G-\kappa L>0,\qquad L=z^tH_c^{-1}z\ge0.
 \tag{UKR20}
\]

\paragraph{Complete separated SOS allocation.}
Choose $\rho<1$ sufficiently close to one and then a positive
real $c$ so that
\[
 L<c<\frac{H_B-g_*^2/A_G}{q\kappa},\qquad q=2/(1+\rho).
 \tag{UKR21}
\]
Such choices exist by UKR20. Use $c$ times UKR2 to introduce
the $cB_u^2$ pivot. The full real residue on
$(N_x,N,G,B,B_u)$ is strictly positive: eliminate $H_c$ to
leave $(c-L)B_u^2$, and eliminate $A_G$ to leave
$(H_B-q\kappa c-g_*^2/A_G)B^2$.
The $N_y$ pivot is $E/(2\nu)>0$ and its cross coefficients
have zero residue. Every other full normal and radial term has
already been allocated by the exact bounded FM constructions
and UKR10. The strict separated bounded-Gram lemma supplies a
finite tensor SOS. Restore the square on the right of UKR2 and
both complete excess terms UKR7. This is the full original
parent, not its restriction to a state graph.

The original sign transformations of FM5 handle negative $u,v$,
and the proof includes $u=0$. Complex states split into real and
imaginary parts. Arbitrary state ratios and independent source
and testing rates are allowed throughout. Unmatched unit-$\kappa$
Bob tests and the $\nu$- or $\ell$-infinitesimal source endpoints
remain outside this theorem.



% Source formulas: p578719565e9d
\subsection{Unit $\kappa,\ell$ and infinitesimal $\nu$}

\label{component:p578719565e9d}

\label{p578719565e9d:sec:new-c144-nu-endpoint}

\paragraph{Scope.}
Retain every original physical C14.4 source predicate, strict gap,
nonnegative event coefficient, and original native source inequalities.
This theorem assumes
\[
 S=\sigma+L\in\operatorname{Inf}_{>0},\quad
 \kappa=x/S\text{ a positive bounded unit},\quad
 \nu=\sigma/S\in\operatorname{Inf}_{>0},\quad\ell=L/S=1-\nu.
 \tag{NU1}
\]
It proves a finite separated tensor SOS for the complete original
parent at every projective test and every real or complex state.
The sector with both $\kappa,\nu$ infinitesimal is not included.
No relative source/testing rate, state-ratio bound, or mixed inverse
is used. The original literal $D=0$ and $\beta=0$ faces are included.

\paragraph{Opening-safe source parameters and the two sphere masses.}
Use the exact FM2--FM3 native dictionary, and write
\[
\begin{gathered}
 q=\nu/\ell,\quad \zeta=\ell/\kappa=L/x,\quad
 e=\kappa h/\nu,\quad E=1-e,\\
 c_0=SC(1-r),\quad a_0=SAr,\quad d_0=D/S,\quad a_1=\alpha/S.
\end{gathered}
 \tag{NU2}
\]
In NU1, $q$ is infinitesimal and $\zeta$ is a positive bounded
unit. The exact source sphere SE3 gives
$e(B/S)=\kappa H_0+\nu\mathcal H b_0^2$.
Since $B/S\sim1+\kappa$, $H_0\sim1$, and
$\mathcal H\sim E^2$, its real residues satisfy
\[
 e=\frac\kappa{1+\kappa},\quad E=\frac1{1+\kappa},\quad
 h\in\operatorname{Inf},\quad\theta\sim1,\quad
 c_0=a_0=a_1=E/2.
 \tag{NU3}
\]
The two equalities for $e,E$ in this display are residue
equalities; the definition $e=\kappa h/\nu$ remains exact.
In particular $e,E,c_0,a_0,a_1$ are positive bounded units.

Define actual nonnegative source coefficients
\[
 d_s=\frac{D}{Sq^2},\qquad n_s=\frac\beta{q^2}.
 \tag{NU4}
\]
These are bounded even when either is zero. More explicitly the
SE3 sphere identities give
\[
 d_s=\frac{\ell^2\mathcal H}{2P_0J}d_0^{\,\mathrm{sph}\,2},
 \qquad
 n_s=\frac{\ell^2\mathcal H}{2\kappa P_0\theta}
                                   b_0^{\,\mathrm{sph}\,2},
 \qquad
 Jd_s+\kappa\theta n_s=\frac{\ell^2\mathcal H}{2P_0}.
 \tag{NU5}
\]
Here $(d_0^{\mathrm{sph}})^2+(b_0^{\mathrm{sph}})^2=1$;
the superscript distinguishes sphere coordinates from $D/S$.
Every denominator in NU5 is a positive unit. At least one of
$d_s,n_s$ is a positive unit. No sphere coordinate is inverted.

The sole possibly unbounded event coefficient needed locally is
$m$. Remove its pole by its actual contact-zero event, not by
dividing an error by $M$. With $R_s=1+xr$ set
\[
 p=\frac{\kappa\theta}{\ell R_s}\in(0,1),\qquad
 \mu_c=m(1-p)=
 \frac{\kappa e(1+x)(1+\nu Sr)}{2\ell P_0R_s}.
 \tag{NU6}
\]
Thus $\mu_c$ is a positive bounded unit and $m\ge\mu_c$.
Choose a fixed positive real $\epsilon<m$. On the reciprocal
Bob chart define $V=W+uY$ and $B_c=w(X+uZ)$. Retain the whole
positive term $(m-\epsilon)(B_c-pV)^2$. The remaining $m$
quadratic is exactly
\[
 m_0B_c^2-2m_1B_cV+m_2V^2,\qquad
 m_0=\epsilon,\quad m_1=\mu_c+\epsilon p,\quad
 m_2=\mu_c(1+p)+\epsilon p^2.
 \tag{NU7}
\]
All three coefficients are bounded through $M\to0$ with arbitrary
rate. Their residues are
\[
 \mu_c=\frac{\kappa^2}{2(1+\kappa)},\quad p=\kappa,
 \qquad m_2=\kappa^2(\epsilon+\tfrac12)>0.
 \tag{NU8}
\]
The source $M>0$ also implies $0<\bar\kappa\le1$.

\paragraph{Local reciprocal chart and exact angular budgets.}
Assume $u,w$ are infinitesimal, including literal zero. Use
the homogeneous reciprocal parent $F_w/4$ from P4, with its
$m$ term replaced by NU7 and all other terms intact. Put
\[
\begin{gathered}
 X_e=X+ewW,\quad N_y=(Y-ruW)/\sqrt{Sr(1-r)},\quad
 a=u^2,\quad c=w^2,\quad \delta=c-q,\\
 P=uW/\sqrt S,\quad Q=\sqrt S\,wZ,\quad
 N=W-\zeta uwZ,\quad G=P-\ell Q,\\
 B=\sqrt S\,\delta Z,\quad B_u=u\delta Z,\quad
 R=w(a-x)Z=\frac{\sqrt S}{\zeta}G-\frac u\zeta N.
\end{gathered}\tag{NU9}
\]
Every row is a finite separated expression. No state coordinate
is divided by another. Since $W(uwZ)=PQ$, two exact identities are
\[
\begin{split}
 \delta^2N^2+2\zeta\delta wGB+2\zeta\ell w^2B^2
                  -\zeta^2(wB_u)^2&=(\delta W)^2,\\
 \delta^2G^2+2\ell\delta wNB_u+2\ell\zeta w^2B_u^2
                  -\ell^2(wB)^2&=(\delta P)^2.
\end{split}\tag{NU10}
\]
Consequently, from either positive sphere pivot in NU5, any fixed
positive diagonal budgets on both $wB$ and $wB_u$ can be obtained
at infinitesimal bounded-Gram cost. To see this without a hidden
division, if $d_s$ is a unit, replace $B$ by
$\sqrt{d_s}B/\sqrt{d_s}$ in the first identity, divide by
$\zeta^2$, and use $wB=w(\sqrt{d_s}B)/\sqrt{d_s}$ for the
other budget. All cost coefficients are infinitesimal. If $n_s$
is a unit, use the second identity divided by $\ell^2$ and
$wB_u=w(\sqrt{n_s}B_u)/\sqrt{n_s}$. The factor
$\zeta w^2$ is infinitesimal because $\zeta$ is a unit in
NU1. This is the precise restriction preventing a silent extension
to infinitesimal $\kappa$.

\paragraph{The entire scalar graph before normal restoration.}
Write $D=Sd_0$, $\alpha=Sa_1$, and define bounded coefficient
polynomials in the independent testing squares $a,c$:
\[
\begin{split}
 A_P={}&(1+c)\{c_0(1-r)+a_0r(1+a)\}\\
 &+S\{(D+\tfrac12)(1-r)^2c
                  +(\alpha+\tfrac12)(1+a)r^2\},\\
 A_N={}&m_2(1+ra)^2+2m_1ec(1+ra)+m_0e^2c^2
                  +(\beta+\tfrac12)e^2c(1+c)\\
 &-(D+\tfrac12)e\{2(1-r)-e\}ac,\\
 K={}&m_1(1+ra)+m_0ec+(D+\tfrac12)(1-r-e)\\
 &+(\alpha+\tfrac12)(1+a)r+(\beta+\tfrac12)e(1+c),\\
 F_Z={}&D+\alpha c+\beta a+(m_0+\alpha+\beta+\tfrac12)ac,\\
 J_1={}&K-\ell A_P,\qquad L_1=J_1-\zeta A_N,\\
 H={}&F_Z-\ell^2ScA_P+ac(\zeta^2A_N-2\zeta J_1).
\end{split}\tag{NU11}
\]
For $X=-ewW$, $Y=ruW$, the exact graph scalar is
\[
 F_0=A_NN^2+A_PG^2-2uwZL_1N+HZ^2.
 \tag{NU12}
\]
This graph is used to compute coefficients; the unrestricted
normals are restored explicitly below. The actual source contact
has $a=x,c=q,W=\zeta uwZ$, so $N=G=X_e=N_y=0$.
The retained row $B_c-pV$ also vanishes there. Direct source
substitution in NU11 yields
\[
 L_1(x,q)=0,\qquad H(x,q)=H_a(x,q)=H_c(x,q)=0,
 \qquad H(a,0)=D+\beta a.
 \tag{NU13}
\]
These are exact rational source identities; no independent
testing angle is identified with a source angle in the certificate.

All divisions in the next display are finite monic polynomial
divisions at the bounded pure-source contacts $x,q$. They define
bounded coefficient polynomials $L_a,L_c,H_{20},H_{11},H_{02}$:
\[
\begin{split}
 L_1(a,c)&=(a-x)L_a(a,c)+(c-q)L_c(c),\\
 H(a,c)&=H(x,c)+(a-x)H_a(x,c)+c(a-x)^2H_{20}(a,c),\\
 H_a(x,c)&=(c-q)\{-\beta/q+cH_{11}(c)\},\\
 H(x,c)&=S(c-q)^2\{d_s+\kappa n_s+cH_{02}(c)\}.
\end{split}\tag{NU14}
\]
Here the factor $c$ in the quadratic $a$ remainder follows from
$H(a,0)=D+\beta a$. The two constant terms follow by putting
$c=0$: $H_a(x,0)=\beta$ and $H(x,0)=D+\beta x$.
In particular the displayed apparent divisions by $q$ are
bounded: $\beta/q=n_sq$. Finally $H(x,c)/S$ is a bounded
coefficient polynomial directly from NU11, because
$x=\kappa S$, $D/S=d_0$, and $\alpha/S=a_1$ are bounded.
Thus dividing it by the monic $(c-q)^2$ introduces no unbounded
coefficient. This verifies all asserted coefficient bounds without
using relative orders of $S,q,u,w$.

The useful mixed identity is
\[
 c(a-x)(c-q)Z^2=\frac1\zeta\{G(wB)-N(wB_u)\}.
 \tag{NU15}
\]
Substituting NU14 into NU12 gives the complete exact scalar
allocation
\[
\begin{split}
 F_0={}&A_NN^2+A_PG^2+d_sB^2+n_sB_u^2
       +H_{02}(wB)^2+H_{20}R^2\\
 &+\frac{H_{11}-n_s}{\zeta}\{G(wB)-N(wB_u)\}
       -2uL_aNR-2L_cN(wB_u).
\end{split}\tag{NU16}
\]
The residues of $A_N,A_P$ are respectively
$\kappa^2(\epsilon+1/2)$ and $E/2$, both positive.
The row $R$ has infinitesimal coefficients on $(N,G)$ by NU9.
Every remaining signed term is a bounded quadratic term on
$N,G,wB,wB_u$. The angular rows are retained as the complete
nonnegative squares $(\sqrt{d_s}B)^2$ and
$(\sqrt{n_s}B_u)^2$; neither is prematurely discarded.

\paragraph{Both complete state normals.}
Put $s_y=\sqrt{Sr(1-r)}$ and $r_y=\sqrt{r(1-r)}$.
Define
\[
\begin{split}
 A_x={}&-m_0ec-m_1(1+ra)+(D+\tfrac12)a(1-r-e)
                          -e(\beta+\tfrac12)(1+c),\\
 T_x={}&\zeta cA_x+\beta-D+c(m_0+\beta+\tfrac12),\\
 Y_P={}&(1+c)\{-c_0+a_0(1+a)\}\\
 &+S\{m_2(1+ra)+m_1ec-(D+\tfrac12)c(1-r-e)
                             +(\alpha+\tfrac12)(1+a)r\},\\
 Y_Q={}&-m_1a+(D+\tfrac12)-(\alpha+\tfrac12)(1+a),\qquad
 Y_C=\ell Y_P+Y_Q.
\end{split}\tag{NU17}
\]
The entire capped quarter-parent, before any allocation, is
\[
 F_0+H_xX_e^2+H_yN_y^2+2H_{xy}X_eN_y+2L_xX_e+2L_yN_y,
 \tag{NU18}
\]
where
\[
\begin{split}
 H_x={}&m_0c+(D+\tfrac12)a+(\beta+\tfrac12)(1+c),\\
 H_y={}&rc_0(1+c)+(1-r)a_0(1+a)(1+c)\\
 &+Sr(1-r)\{m_2a+(D+\tfrac12)c+(\alpha+\tfrac12)(1+a)\},\\
 H_{xy}={}&-uws_y(m_1+D+\tfrac12),\\
 L_x={}&wA_xN+uZT_x,\qquad
 L_y=r_y\{Y_PG+Y_CQ\}.
\end{split}\tag{NU19}
\]
These formulas retain every quadratic state coefficient and both
first normals. Their real hard residues are
$H_x=1/2$, $H_y=E/2$, $H_{xy}=0$.

Exact source substitution gives $T_x(x,q)=Y_C(x,q)=0$.
Moreover $T_x(a,0)=\beta-D$, and its $a$ dependence has a
factor $c$. Thus bounded monic divisions give
\[
\begin{split}
 T_x(a,c)&=c(a-x)X_a+(c-q)\{(D-\beta)/q+cX_c(c)\},\\
 Y_C(a,c)&=(a-x)Y_a(a,c)+(c-q)Y_c(c).
\end{split}\tag{NU20}
\]
No coefficient here needs an inverse of a sphere coordinate.
Using NU4 and NU9, the full normal conductors become
\[
\begin{split}
 L_x={}&wA_xN+uwX_aR
       +d_s\sqrt S\,q u B-n_sqB_u+wX_c(wB_u),\\
 L_y={}&r_y\{Y_PG+\sqrt S\,Y_aR+Y_c(wB)\}.
\end{split}\tag{NU21}
\]
The two potentially weak terms are face-safe:
$d_s\sqrt S\,quB=(\sqrt{d_s}\sqrt S\,qu)(\sqrt{d_s}B)$
and $n_sqB_u=(\sqrt{n_s}q)(\sqrt{n_s}B_u)$.
Their coefficients are infinitesimal, even when either sphere
mass is zero. Also $Y_P$ has zero residue by NU3 and NU17.
All other terms in NU21 have infinitesimal coefficients on
$N,G,R$ or bounded coefficients on the two angular helper rows.

\paragraph{Strict complete Gram and restoration.}
Begin with the six rows
\[
 (N,G,X_e,N_y,\sqrt{d_s}B,\sqrt{n_s}B_u).
 \tag{NU22}
\]
Their real diagonal residue in NU16--NU21 is
\[
 \operatorname{diag}\bigl(
 \kappa^2(\epsilon+\tfrac12),E/2,1/2,E/2,1,1\bigr)>0.
 \tag{NU23}
\]
Every cross coefficient among these rows is infinitesimal after
the displayed $R$ substitution. Literal zero rows are harmless:
the coefficient matrix may be strict even when the chosen rows
are linearly dependent or vanish.

Use NU10 to add sufficiently large fixed real pivots for the
two further rows $wB,wB_u$, at infinitesimal cost to NU22.
The signed scalar terms in NU16 and the remaining normal terms
in NU21 have bounded coefficients on these helper rows.
A sufficiently large fixed pair of diagonals therefore gives
a strictly positive real eight-row Gram, by an ordinary real
Schur complement. All entries are finite sums of bounded separated
factors. The separated bounded-Gram lemma supplies a finite SOS.
Restore the exact positive square remainders in NU10 and the
whole excess $(m-\epsilon)(B_c-pV)^2$ in NU7. This gives the
entire original parent, with both unrestricted state normals.

\paragraph{All projective testing joins on NU1.}
Noninfinitesimal Alice tests use CE1--CE10 because $C,A$ are
unbounded. For infinitesimal $u$, the local proof above covers
every infinitesimal reciprocal $w$, including zero. It remains
to check finite Bob coordinates $v$.

For $v$ infinitesimal use the exact FU22--FU23 frame; its source
inputs here are only $m$ strong, $c_0,a_0$ with common positive
residue $E/2$, and $a_1$ positive, all proved in NU3--NU6.
For a finite positive unit $v$ use the complete FU6--FU12
bounded parent and real Schur calculation, capping the original
$m$ event at a finite real only if $m$ is unbounded. In the
present source residue $\beta=d_0=0$, $a_1=E/2$, and, if
$\kappa<1$,
\[
 m=\frac{\kappa^2}{2(1-\kappa^2)},\qquad E=\frac1{1+\kappa}.
 \tag{NU24}
\]
The determinant polynomial FU14 becomes exactly
\[
 \mathcal P(y,c)=p_{21}c(y-\kappa)^2+p_{10}y>0,
 \qquad p_{21}=E(m+\tfrac12),\quad
 p_{10}=(m+\tfrac12)E^2,\quad y,c>0.
 \tag{NU25}
\]
It gives the strict four-term circuit inequality FU18. If
$\bar\kappa=1$, then $m$ is unbounded. The normalized limit is
\[
 \mathcal P(y,c)/m\longrightarrow
              \tfrac12c(y-1)^2+\tfrac14y>0.
 \tag{NU26}
\]
The Schur coefficients have finite positive limits, so a
sufficiently large finite cap preserves strictness, exactly as
in FU19--FU20. The whole positive original event excess is
restored. These finite-Bob arguments use no lower bound on
$\nu$, and do not import the source scope of the FU theorem.

The exact sign transformations FU31 cover negative angles;
all displayed frames include literal axes and projective
endpoints. Complex states split into real and imaginary parts.
No chart normalizes a state coordinate or bounds a state ratio.
This completes the all-tests theorem on precisely NU1.



% Source formulas: p2edb55cd1f69
\subsection{Infinitesimal $\kappa$ and unit $\nu,\ell$}

\label{component:p2edb55cd1f69}

\label{p2edb55cd1f69:sec:new-c144-finite-matched}

\paragraph{Scope.}
This theorem treats the source sector stated in FM1: $S,\kappa$ are positive infinitesimals and $\nu,\ell$ are positive bounded units. The source endpoints $\nu\in\operatorname{Inf}$ and $\ell\in\operatorname{Inf}$, and the unit-$\kappa$ source, are separate cells in the complete partition of Section~\ref{component:pfd0dd7bb10cf}. All original physical C14.4 source conditions, strict gaps, nonnegative event coefficients, total-scale conditions and native source inequalities are retained.

\paragraph{Theorem.}
Let the original native C14.4 source have
\[
 S=\sigma+L>0,\qquad x=\kappa S,\qquad
 \nu=\sigma/S,\quad\ell=L/S,\quad\nu+\ell=1,
 \qquad S,\kappa\in\operatorname{Inf}_{>0},
 \tag{FM1}
\]
with $\nu,\ell$ positive bounded units. Then the complete original
parent is a finite SOS in $K_c\otimes_{\mathbb R}K_t$ for every
independent projective test and every real or complex state.
The argument below supplies the missing finite matched-Bob chart;
the other tests use the explicitly scoped SE theorem.

\paragraph{Exact opening-safe source dictionary.}
Write $\theta=1-h$, $0<r,h<1$, $e=\kappa h/\nu$, $E=1-e$,
and set
\[
\begin{split}
 J&=1+xh-\nu S(1-r),\quad
 M=1+xr-\kappa\theta/\ell,\quad Q_s=1-\ell Sr,\\
 P_0&=1+S\{\nu\ell+2\nu\kappa h+\kappa\ell+
                    \kappa^2h\theta\}
             +S^2\{\nu r\kappa\ell+\nu r\kappa^2h\}.
\end{split}\tag{FM2}
\]
These are the actual native source quantities, with $P=SP_0$.
The converted P1 coefficients, without applying P1's different
unit-$L$ asymptotic hypotheses, are exactly
\[
\begin{split}
 c_0=SC(1-r)&=\frac{(1+x)(1+\nu Sr)E}{2P_0},\\
 a_0=SAr&=\frac{(1+xh)E+(1-r)(\ell S+xh)}{2P_0},\\
 d_0=D/S&=\frac{(1+x)(1+\nu Sr)-JP_0}{2JSP_0},\\
 a_1=\alpha/S&=
 \frac{(1+\nu Sr)\{\ell(1+xr)e+\ell^2ME\}}{2Q_sP_0},\\
 b_1=\beta\theta&=
 \frac{(1+x)(1+\nu Sr)(\ell+\kappa h)-\theta P_0}{2P_0},\\
 \mu=m/\kappa^2&=
 \frac{(1+x)(1+\nu Sr)h}{2M\ell\nu P_0}.
\end{split}\tag{FM3}
\]
The numerator in the $d_0$ formula has a factor $S$, so its displayed
$S$ denominator cancels in the source polynomial ring. Thus all six
coefficients belong to the bounded source algebra generated by
$S,\kappa,h,r,\nu,\ell$ and inverses of $J,M,Q_s,P_0,\nu,\ell$.
None of these inverses is of an opening $r,1-r,\theta,D$ or $\beta$.
The original source gives
\[
\begin{gathered}
 J,M,Q_s,P_0,E\sim1,\qquad c_0,a_0\sim\tfrac12,\qquad
 d_0\sim\nu^2/2,\quad a_1\sim\ell^2/2,\\
 b_1-(h-\nu)/2\in\operatorname{Inf},\qquad
 \mu-h/(2\nu\ell)\in\operatorname{Inf},\qquad
 \bar h\ge\bar\nu>0.
\end{gathered}\tag{FM4}
\]
The last inequality follows from $b_1=\beta\theta\ge0$.
Hence $\mu$ is a positive bounded unit. These estimates are uniform
through $r,1-r,\theta\to0$; no comparison with a testing quantity
has been used. In this sector $D/S$ is a positive unit, so $D=0$
has empty intersection; $\beta=0$ is allowed.

\paragraph{Exact full parent and unrestricted normals.}
On the finite Bob chart write $D_u=1+u^2$, $D_v=1+v^2$.
The complete quarter-parent is
\[
\begin{split}
 F/4={}&m(vW-X+uvY-uZ)^2+C D_v(uW-Y)^2+A D_uD_vY^2\\
 &+(D+\tfrac12)(uW-Y+uvX-vZ)^2
 +(\alpha+\tfrac12)D_u(vY-Z)^2\\
 &+(\beta+\tfrac12)D_v(X+uZ)^2-\tfrac12D_uD_vZ^2.
\end{split}\tag{FM5}
\]
No state restriction is made. Put
\[
 N_x=X+huZ,\qquad N_y=\frac{Y-ruW}{\sqrt{Sr(1-r)}}.
 \tag{FM6}
\]
The source square root is nonzero even when its size is arbitrarily
small. It is used only to define a permitted tensor root; every
coefficient below is bounded without its inverse.

After the unmatched-residue and adjacent-$C,A$ exits it is
enough to take $u$ infinitesimal and $v$ a bounded positive unit with
the same residue as $b=\sqrt{\ell/\nu}$. Set
\[
\begin{gathered}
 a=u^2,\quad j=\sqrt{\nu\ell},\quad T=\kappa W,\quad U=uZ,\\
 P=uW/\sqrt S,\quad Q=\sqrt S Z,\quad N=T-jU,\quad G=P-jQ,\\
 \delta=v-b,\quad B=\delta Q,\quad B_u=\delta U.
\end{gathered}
 \tag{FM7}
\]
Every state coordinate remains arbitrary, including zero and
unbounded coordinates. The rows satisfy the exact MB identities
\[
\begin{gathered}
 R_Z=(a-x)Z=(\kappa\sqrt S/j)G-(u/j)N,\\
 (a-x)\delta Z^2=(\kappa/j)GB-(1/j)NB_u,\\
 \delta^2N^2+2j\kappa\delta GB+2j^2\kappa B^2-j^2B_u^2
                  =(\delta T)^2.
\end{gathered}\tag{FM8}
\]
These polynomial identities hold in the separated tensor ring.

\paragraph{The full graph scalar.}
Let $F_0$ denote FM5 with $N_x=N_y=0$; it is used only as the
constant coefficient of the full normal expansion. Define bounded
source/testing polynomials
\[
\begin{split}
 A_T={}&\mu v^2(1+ra)^2,\\
 A_P={}&(1+v^2)\{c_0(1-r)+a_0r(1+a)\}\\
       &+S\{(D+\tfrac12)(1-r)^2+
                   (\alpha+\tfrac12)(1+a)r^2v^2\},\\
 A_Q={}&d_0v^2+a_1,\\
 A_U={}&m\theta^2+\alpha+b_1\theta+\theta^2/2\\
       &+v^2\{2Dh+h-\tfrac12+b_1\theta+\theta^2/2\}
              +(D+\tfrac12)h^2v^2a,\\
 K={}&m\theta(1+ra)+(D+\tfrac12)(1-r)(1+ha)
                         +(\alpha+\tfrac12)(1+a)r,\\
 H={}&vK-\kappa jA_T,\quad R_U=A_U-j^2A_T,\\
 g={}&jA_P-H,\quad R_Q=A_Q+j^2A_P-2jH,\\
 J_Z={}&SR_Q+aR_U.
\end{split}\tag{FM9}
\]
Direct expansion of every term of FM5 gives
\[
 F_0=A_TN^2+A_PG^2+2gGQ+J_ZZ^2.
 \tag{FM10}
\]
For example the only relation used to obtain FM10 from the
five monomials $T^2,U^2,P^2,Q^2,PQ$ is $TU=\kappa PQ$.

The actual source stationarity identities, also verified directly
from FM2--FM3, are
\[
 g(x,b)=0,\qquad
 J_Z(x,b)=\partial_aJ_Z(x,b)=\partial_vJ_Z(x,b)=0.
 \tag{FM11}
\]
These are coefficient-field identities at the actual source
contact; they impose no equation on independent tests.

\paragraph{Finite polynomial division.}
All following quotients mean monic polynomial division in formal
$a$ or $v$, performed before evaluation in the tensor ring. No
mixed inverse is introduced. There are bounded coefficient
polynomials $g_a,g_v,H_{20},H_{11},H_{02}$ defined exactly by
\[
\begin{split}
 g(a,v)&=(a-x)g_a(a,v)+(v-b)g_v(v),\\
 J_Z(a,v)&=(a-x)^2H_{20}(a,v)+(a-x)(v-b)H_{11}(v)
                                      +S(v-b)^2H_{02}(v).
\end{split}\tag{FM12}
\]
For clarity, first divide $g(a,v)-g(x,v)$ by $a-x$ and then
$g(x,v)$ by $v-b$. For the second line, first subtract the
constant and linear Taylor terms in $a-x$ and divide the remainder
by $(a-x)^2$. Divide $\partial_aJ_Z(x,v)$ by $v-b$ for $H_{11}$.
Finally divide the bounded polynomial
$R_Q(x,v)+\kappa R_U(x,v)$ by $(v-b)^2$ for $H_{02}$.
FM11 makes every division exact. In particular $H_{02}$ is
bounded without dividing a bounded error by $S$.

Substituting FM8 into FM10 now yields the exact scalar identity
\[
\begin{split}
 F_0={}&A_TN^2+A_PG^2+2g_vGB+H_{02}B^2
          +H_{11}\{(\kappa/j)GB-(1/j)NB_u\}\\
       &+H_{20}R_Z^2+2\sqrt S\,g_aGR_Z.
\end{split}\tag{FM13}
\]
Thus every coefficient is bounded and all remaining radial
contributions are infinitesimal matrix entries on $(N,G)$.

Here and below a displayed residue formula uses the real residues
of its bounded source parameters. FM3 gives explicitly
\[
 \bar A_T=\frac h{2\nu^2},\quad
 \bar A_P=\frac1{2\nu},\quad
 \bar g_v=\frac{\ell-\nu}{2},\quad
 \bar H_{02}=\frac\nu2,\quad \bar H_{11}=-\theta j.
 \tag{FM14}
\]
For example at $S=\kappa=0$ and $a=0$ the three relevant
polynomials are
$g=j(1+v^2)/2-v/2$,
$R_Q=\nu(v-b)^2/2$, and
$R_U=\theta(\ell-\nu v^2)/2$.
These establish FM14 by the finite divisions just specified.

\paragraph{Both complete first-normal terms.}
Write the full expansion of FM5 as
\[
 F/4=F_0+H_xN_x^2+H_yN_y^2+2H_{xy}N_xN_y
                  +2N_xL_x+2N_yL_y+\beta D_vN_x^2.
 \tag{FM15}
\]
The exact bounded diagonal and cross coefficients are
\[
\begin{split}
 H_x={}&m+(D+\tfrac12)u^2v^2+D_v/2,\\
 H_y={}&r c_0D_v+(1-r)a_0D_uD_v\\
       &+Sr(1-r)\{mu^2v^2+D+\tfrac12
                            +(\alpha+\tfrac12)D_uv^2\},\\
 H_{xy}={}&-uv\sqrt{Sr(1-r)}(m+D+\tfrac12).
\end{split}\tag{FM16}
\]
In particular $\bar H_x=\bar H_y=1/(2\nu)$ and
$\bar H_{xy}=0$. The last term of FM15 is retained whole;
it allows $\beta$ to be arbitrarily large or zero.

Define
\[
\begin{split}
 L_U={}&m\theta-(mj/\kappa)v(1+ra)
 -(D+\tfrac12)v^2(1+ha)+(b_1+\theta/2)(1+v^2)\\
       &+(D+\tfrac12)(1-r)vjS,\\
 Y_P={}&(1+v^2)\{-c_0+a_0(1+a)\}
       +S\{mv^2(1+ra)-(D+\tfrac12)(1-r)
                        +(\alpha+\tfrac12)(1+a)rv^2\},\\
 Y_Q={}&-m\theta a+(D+\tfrac12)(1+ha)
                                  -(\alpha+\tfrac12)(1+a),\qquad
 Y_C=jY_P+vY_Q.
\end{split}\tag{FM17}
\]
The expression $mj/\kappa=\kappa\mu j$ is bounded and
infinitesimal. The complete original normal terms are exactly
\[
\begin{split}
 L_x={}&-\kappa\mu v(1+ra)N
       +(D+\tfrac12)(1-r)vu\sqrt S\,G+L_UU,\\
 L_y={}&\sqrt{r(1-r)}\{Y_PG+Y_CQ\}.
\end{split}\tag{FM18}
\]
The actual source gives $L_U(x,b)=Y_C(x,b)=0$. Perform the same
finite monic divisions to write
\[
 L_U=(a-x)l_a+(v-b)l_v,\qquad
 Y_C=(a-x)y_a+(v-b)y_v.
 \tag{FM19}
\]
Then $L_UU=u l_aR_Z+l_vB_u$ and
$Y_CQ=\sqrt S\,y_aR_Z+y_vB$. All their coefficients are
bounded. Moreover $\bar l_v=-j$, while $Y_P$ and $y_v$ have
zero residue. This follows directly from the source-limit
identities
$L_U=(\ell-\nu v^2)/2$ at $a=0$, and
$Y_P=Y_Q=0$ at $a=0$.
Consequently $L_x+jB_u$ is an infinitesimal bounded combination
of $(N,G,B_u)$ and $L_y$ is an infinitesimal bounded combination
of $(N,G,B)$. This accounts for both unrestricted first-normal
terms, including all $\beta\theta$ terms and the independent
$r,1-r,\theta$ openings.

\paragraph{A strict residue for the full six-row Gram.}
Use FM8 with a fixed real $c=4$ to introduce a $cB_u^2$ pivot
at infinitesimal Gram cost. More explicitly subtract
\[
 (c/j^2)\delta^2N^2+(2c\kappa/j)\delta GB+2c\kappa B^2
 \tag{FM20}
\]
from the bounded part of FM15, add $cB_u^2$ to that part, and
retain the exact positive square $(\sqrt c\,\delta T/j)^2$.
Their sum is unchanged by FM8. The residue on
$(N_x,N_y,N,G,B,B_u)$ is therefore
\[
 \frac{N_x^2+N_y^2}{2\nu}+\frac h{2\nu^2}N^2
 +\frac{G^2}{2\nu}+(\ell-\nu)GB+\frac\nu2B^2
       +\theta NB_u-2jN_xB_u+4B_u^2.
 \tag{FM21}
\]
The $(G,B)$ block has determinant $\nu\ell>0$ and positive
diagonal. The remaining block has positive diagonal pivots
$1/(2\nu),1/(2\nu),h/(2\nu^2)$, and its final scalar Schur
pivot is
\[
 4-2\nu j^2-\frac{\theta^2\nu^2}{2h}>0.
 \tag{FM22}
\]
Indeed $j^2\le1/4$, $0<\nu\le1$, $h\ge\nu$ and
$0\le\theta\le1$ in residue make the two subtracted terms
at most $1/2$ each. All entries before taking residues are finite
sums of bounded separated factors. The strict
bounded-Gram lemma therefore supplies a finite tensor SOS.
Restoring the single square from FM20 and the whole positive
$\beta D_vN_x^2$ row gives the entire FM5. No mixed matrix inverse
or mixed square root has been used.

\paragraph{Testing, sign, and boundary joins.}
The argument includes $u=0$ and every real state without a state
normalization or a sign condition. Negative $u$ is changed to
$-u$ together with $(Y,Z)\mapsto(-Y,-Z)$; negative $v$ is
changed to $-v$ together with $(X,Z)\mapsto(-X,-Z)$.
FM5 is exactly invariant under these transformations. On FM1,
$e=\kappa h/\nu$ is infinitesimal and $C,A$ are unbounded.
The SE10--SE12 theorem covers every unmatched
homogeneous Bob residue, including Bob zero and infinity,
and invokes the unchanged adjacent-$C,A$ theorem for
noninfinitesimal Alice tests. The only remaining Bob residue
has finite nonzero $v$ and $\bar v^2=\bar\ell/\bar\nu$,
which is the matched chart just proved. Thus the union covers
all projective tests on FM1. Real and imaginary state parts
are treated separately for complex states. All divisions and
normalizations used in these joins are pure source or pure
testing operations; all original source predicates remain.

\paragraph{Dependencies and limits.}
The analytical inputs are the native source formulas
P1 converted as FM2--FM3, the full original parent, the strict
separated bounded-Gram lemma in Lemma~\ref{p25af2c58da7e:lem:gram},
and Section~\ref{component:p8c30238a5f7d}
with its preserved EB and adjacent-$C,A$ dependencies.
The exact budget is independently printed in
Section~\ref{component:pe55ec4890781}. The source endpoint
and unit-$\kappa$ sectors require separate complete arguments;
this theorem makes no claim about them.



% Source formulas: pb96c9bfcc12d
\subsection{Infinitesimal $\kappa,\ell$ and unit $\nu$}

\label{component:pb96c9bfcc12d}

\label{pb96c9bfcc12d:sec:c144-ell-endpoint-rederived}

\paragraph{Theorem and exact scope.}
Retain the original native C14.4 physical source, all strict gaps,
nonnegative event weights, total-scale conditions and native source inequalities.  Let
\[
 S=\sigma+L,\quad x=\kappa S,\quad
 \nu=\sigma/S,\quad\ell=L/S,\quad\nu+\ell=1,
 \quad\zeta=L/x=\ell/\kappa.
\]
Suppose $S,\ell>0$ are infinitesimal and $\zeta$ is
noninfinitesimal.  Then the complete original parent is a finite
sum of squares in the separated ring
$K_c\otimes_{\mathbb R}K_t$ for every independently chosen
projective test and every real or complex state.  The proof includes
$\beta=0$ and unrestricted openings of $r,1-r,\theta$.
The intersection with $D=0$ is empty by the source estimate below.
No earlier missing R21 endpoint proof is used.

\paragraph{Source estimates at this endpoint.}
Use the exact actual-source formulas FM2--FM3 and SE1--SE3, with
$0<r,h<1$, $\theta=1-h$, $e=\kappa h/\nu$, $E=1-e\ge1/18$,
$U_s=(1+x)(1+\nu Sr)$, and
\[
\begin{gathered}
 J=1+xh-\nu S(1-r),\quad
 M=1+xr-\kappa\theta/\ell>0,\quad Q_s=1-\ell Sr,\\
 a_0=SAr,\qquad c_0=SC(1-r),\qquad P=SP_0.
\end{gathered}
\]
Put $g=\kappa/\ell=1/\zeta$, which is bounded.
Then $\nu\sim1$ and $\kappa\in\operatorname{Inf}$.
Since
\[
 0\le\beta\theta
  =\frac{U_s(\ell+\kappa h)-\theta P_0}{2P_0},
\]
we have $\theta=O(\ell+\kappa)=O(\ell)$ and
$\beta\theta=O(\ell)$, hence $h\sim1$,
$\beta\theta^2=O(\ell^2)$ and $M\sim1$.
The exact source formulas give
\[
\begin{gathered}
 J,Q_s,P_0,E\sim1,\quad c_0,a_0\sim\tfrac12,
 \quad D/S\sim\tfrac12,\quad m\in\operatorname{Inf}_{>0},\\
 \frac{qm}{\kappa^2}
   =\frac{U_s h}{2M\nu^2P_0}\sim\tfrac12,
 \qquad q=\ell/\nu,\\
 \frac{\alpha}{S\ell^2}
   =\frac{(1+\nu Sr)\{(1+xr)e/\ell+ME\}}{2Q_sP_0},
 \qquad
 \frac{\alpha}{S\ell^2}-\frac{1+g}{2}\in\operatorname{Inf}.
\end{gathered}                                                   \tag{EL0a}
\]
The last difference is uniform for bounded $g$, because
$e/\ell=gh/\nu$.  All statements are inside the source field;
none compares a testing rate to a source rate.  In particular
$D/S$ is a positive unit, while $\beta$ may be zero, bounded,
or unbounded.

\paragraph{The full parent and local testing chart.}
Noninfinitesimal Alice tests and unmatched Bob residues will be joined
at the end.  The remaining chart has $u,v\ge0$ infinitesimal,
including literal zero axes.  Write $a=u^2$, $c=v^2$,
$b=\sqrt q$, $j=\nu b$, and $d=D+1/2$.
The complete original quarter-parent is
\[
\begin{aligned}
 F/4={}&m(vW-X+uvY-uZ)^2+C(1+c)(uW-Y)^2\\
 &+A(1+a)(1+c)Y^2+d(uW-Y+uvX-vZ)^2\\
 &+(\alpha+\tfrac12)(1+a)(vY-Z)^2
 +(\beta+\tfrac12)(1+c)(X+uZ)^2\\
 &-\tfrac12(1+a)(1+c)Z^2.
\end{aligned}                                                     \tag{EL0b}
\]
All state coordinates are arbitrary elements of the testing field.
No state coordinate is divided by or bounded below.
Put $R_s=1+xr$, $p=\kappa\theta/(\ell R_s)$, $\mu=m(1-p)$, and
\[
 V=v(W+uY),\qquad N=X+uZ-pV,\qquad
 Y=ruW+\sqrt{Sr(1-r)}N_y.                         \tag{EL1}
\]
These are pure-source rescalings of polynomial state normals.
The original $m$ event is $(1-p)V-N$, the beta event is $N+pV$,
and the core event is
$uW-Y+puvV-(1+a)vZ+uvN$; thus no term has been dropped.
On $N=N_y=0$ define
\[
\begin{aligned}
 L_0&=\mu(1-p)+(\beta+\tfrac12)p^2,
 \qquad L_1=(\beta+\tfrac12)p^2,\\
 B_c&=1-r+pc(1+ra),\\
 \lambda(a,c)&=\frac q{\kappa^2}(1+ra)^2(L_0+L_1c),\\
 A_P&=(1+c)[c_0(1-r)+a_0r(1+a)]
       +S[dB_c^2+(\alpha+\tfrac12)(1+a)r^2c],\\
 K&=(1+a)[dB_c+(\alpha+\tfrac12)r],\\
 F_Z&=(1+a)[\alpha+Dc+dac].
\end{aligned}                                                    \tag{EL2}
\]
Here $c_0=SC(1-r)$ and $a_0=SAr$.  The complete graph scalar is
\[
 F_0=c(1+ra)^2(L_0+L_1c)W^2+A_PP^2-2vKPQ+F_ZZ^2,
 \quad P=uW/\sqrt S,\quad Q=\sqrt S Z.             \tag{EL3}
\]
Let $T=\kappa W,U=uZ$ and use the exact four-term circuit
\[
 \Psi=(\nu/\ell)cT^2+\nu^2cU^2+\kappa\nu\ell Q^2
       +\kappa P^2-4\kappa\nu vPQ.                \tag{EL4}
\]
Its four positive monomial square roots are
$A_1=(\kappa/b)vW$, $A_2=\nu vuZ$,
$A_3=\nu b\sqrt xZ$, and $A_4=\sqrt\kappa P$.
Their product is a square.  The usual two paired AM--GM identities
make $\Psi$ a separated SOS, including $WZ<0$ by sign excess.
The multiplier $\lambda$ itself is a sum of separated monomial
squares with positive pure-source coefficients; no mixed radical
of $\lambda$ is taken.

Set $\Gamma=P-\nu vQ$, $\Lambda=vP-\ell Q$, and
\[
 \widehat A=A_P-\kappa\lambda,\quad
 \widehat K=K-2\kappa\nu\lambda,\quad
 \widehat F=F_Z-\lambda(\nu^2ac+\kappa\nu\ell S).
\]
Then the following equality is purely polynomial and exact:
\[
 F_0-\lambda\Psi
   =\widehat A\Gamma^2-2vJ\Gamma Q+HZ^2,
 \quad J=\widehat K-\nu\widehat A,
 \quad H=\widehat F+\nu^2Sc\widehat A-2\nu Sc\widehat K.
                                                               \tag{EL5}
\]
The exact coefficient bounds are
$\lambda_0=qL_0/\kappa^2\sim1/2$ and
$\lambda_1=qL_1/\kappa^2=O(\ell)$.
Indeed the exact opening-safe formulas are
\[
 \lambda_1=\frac{\theta(2\beta\theta+\theta)}
                    {2\ell\nu R_s^2},\qquad
 \lambda_0=\frac{U_shM}{2\nu^2P_0R_s^2}+\lambda_1.
\]
Here $\theta/\ell$ and $(\beta\theta)/\ell$ are bounded, so
the formula for $\lambda_1/\ell$ has only bounded source factors.
No comparison between $S$ and $\ell$ is made.  The mixed term
arising from $-\nu^2\lambda_1c^2a$ is allocated in EL18 below
using an actual row of the same circuit.

\subsubsection*{Exact helper rows and the complete two-normal conductor}
The following exact budgets and complete normal conductors use a
positive real fraction of $\Psi$ and strict pivots on
$(\Gamma,\Lambda)$.  The scalar calculation below supplies precisely
these pivots; it does not use the state normals to obtain them.
Put $\delta=v-b$, $G_0=P-jQ$, $B_o=\delta Q$, and
\[
\begin{gathered}
 E_{12}=A_1-A_2=(v/b)(T-jU),\quad V_h=v\delta U,\\
 R_v=v(a-x)Z,\quad B=\sqrt S(c-q)Z,\quad
 R_0=\ell\sqrt S(a-x)Z.
\end{gathered}
                                                               \tag{EL6}
\]
The exact angular relations are
\[
 (v+b)G_0=\Lambda+b\Gamma,\qquad
 \nu(v+b)B_o=\Lambda-v\Gamma=\nu B.             \tag{EL7}
\]
In particular $B$ has a bounded linear expression in the two strict
angular rows.  Multiply MB3 by $v^2/j^2$ and use
$\eta=\kappa(v+b)/j\in\mathrm{Inf}$.  The exact SOS inequality is
\[
 V_h^2\preceq_{\rm SOS}
  \frac{\delta^2}{\nu^2}E_{12}^2
   +\eta(\Lambda+b\Gamma)^2
   +\frac{\eta+2\kappa}{\nu^2}(\Lambda-v\Gamma)^2.
                                                               \tag{EL8}
\]
For a direct verification, set $R_1=vG_0,R_2=vB_o$.
The mixed-term replacement uses the legal SOS
\[
 \eta(R_1^2+R_2^2)-2\kappa\delta R_1R_2/j
 =\frac{\kappa v}{j}(R_1-R_2)^2
  +\frac{\kappa b}{j}(R_1+R_2)^2.
\]
Replacing $v^2$ by $(v+b)^2$ leaves
$(2vb+b^2)G_0^2$ and $(2vb+b^2)B_o^2$, with positive
separated multipliers.  The remaining MB3 square is
$(v\delta T/j)^2$.  Every cost in EL8 is infinitesimal, since
$\kappa/j=(\kappa/\ell)b$ is infinitesimal.

Two further exact relations are
\[
 R_v=\frac{\kappa v\sqrt S}{\nu b}\Gamma
       -\frac u\nu E_{12}+\frac{xv(v-b)}bZ,          \tag{EL9}
\]
\[
 \frac{x^2}{qS}B^2-
       \left(\frac{xv(v-b)}bZ\right)^2
   =\frac{x^2}{q}(2vb+b^2)(v-b)^2Z^2.              \tag{EL10}
\]
Thus $R_v$ is free at infinitesimal Gram cost because
$x^2/(qS)=\kappa S(\kappa/q)$ is infinitesimal.  Finally
\[
 2\nu^2ScR_v^2+2\nu^2(a-x)^2B^2-R_0^2
       =[\nu\sqrt S(2c-q)(a-x)Z]^2               \tag{EL11}
\]
supplies $R_0$ with infinitesimal costs.  None of EL6--11 bounds or
divides a state coordinate or a mixed source/testing sum.

For the exact complete $N$ half derivative set
\[
\begin{aligned}
 B_0(c)&=[-\mu+(\beta+\tfrac12)p(1+c)]/\kappa,\\
 F_a(a,c)&=r\kappa B_0(c)+d(1-r)+dpc(1+ra),\\
 D_t(a,c)&=d(1+a)-\nu S F_a(a,c).
\end{aligned}                                                     \tag{EL12}
\]
All their coefficients are bounded, including at $\beta=0$ or
unbounded $\beta$: $\mu/\kappa$ and
$(\beta+1/2)p/\kappa=(\beta\theta+\theta/2)/(\ell R_s)$
are bounded.  Write the monic polynomial differences
$D_t-D_t(x,q)=(a-x)D_a+(c-q)D_c$ and
$B_0(c)-B_0(q)=(c-q)B_{01}$.  Source contact gives
$D_t(x,q)=\nu B_0(q)$.  The entire half derivative is
\[
\begin{aligned}
 L_N={}&bB_0(c)E_{12}+uv\sqrt S F_a\Gamma-uvD_aR_v\\
 &+[-D_t(x,q)+\nu bB_{01}(v+b)-v(v+b)D_c]V_h.
\end{aligned}                                                     \tag{EL13}
\]
Hence EL8--10 allocate every $N$ cross term without a state bound.
The normal $\widehat N=\sqrt{1+\beta}N$ has diagonal at least
$(\beta+1/2)/(1+\beta)\ge1/2$; dividing the cross coefficients
by $\sqrt{1+\beta}$ only improves the displayed bounds.

The complete other half derivative is
\[
 L_y=\sqrt{r(1-r)}[Y_P\Gamma+vY_KQ],             \tag{EL14}
\]
where
\[
\begin{aligned}
 Y_P={}&(1+c)[-c_0+a_0(1+a)]\\
 &+S\{c(1+ra)(L_0+L_1c)+d(-1+pac)B_c
                   +(\alpha+\tfrac12)(1+a)rc\},\\
 Y_K={}&\nu Y_P+(1+a)(D-\alpha-dpac).
\end{aligned}
\]
The polynomial coefficients are bounded, $Y_P$ is infinitesimal,
and actual source contact gives $Y_K(x,q)=0$.  Monic division gives
$Y_K=(a-x)Y_a+(c-q)Y_c$ with bounded coefficients, so
\[
 L_y=\sqrt{r(1-r)}[Y_P\Gamma+\sqrt S Y_aR_v+vY_cB].
                                                               \tag{EL15}
\]
The exact quadratic normal coefficients are
\[
\begin{aligned}
 H_N={}&m+dac+(\beta+\tfrac12)(1+c),\\
 H_y={}&rc_0(1+c)+(1-r)a_0(1+a)(1+c)\\
 &+Sr(1-r)[d(-1+pac)^2+(\alpha+\tfrac12)(1+a)c
                         +ac(L_0+L_1c)],\\
 H_{Ny}={}&\sqrt{Sr(1-r)}uv[\kappa B_0(c)-d+dpac].
\end{aligned}                                                     \tag{EL16}
\]
After the source-only normalization of $N$, the first two diagonal
residues are respectively in $[1/2,1]$ and $1/2$, and the cross
coefficient is infinitesimal.  Equations EL13--16 include all terms
of the full original parent; they require no $W/Z$ or testing/source
rate restriction.  The scalar allocation and exact strict-Gram assembly are given next.

\subsubsection*{Complete scalar allocation, including the mixed radial term}
Both $\widehat A$ and the coefficients of $J,H$ in EL5 are
bounded polynomials in $a,c$: their source coefficients are built
from bounded $p,d,\alpha,c_0,a_0,\lambda_0,\lambda_1,\kappa,S,\nu$.
In fact $p=O(\kappa)$,
$\mu/\kappa=O(\kappa/\ell)$, and
$(\beta+1/2)p^2=O(\kappa^2)$.
Consequently $\widehat A\sim1/2$ on the local chart.
Actual source contact, checked in the original parent, gives
\[
 J(x,q)=0,\qquad H(x,q)=H_a(x,q)=H_c(x,q)=0.        \tag{EL17}
\]
These are exact identities, not residue equalities.  They also follow
by evaluating the actual contact state in EL1: all four circuit
monomials agree, $\Gamma=0$, and both graph normals vanish.
The vanishing first derivatives of the original contact parent
then give the three scalar jets.  

The additional actual circuit row
\[
 E_{23}=A_2-A_3=\nu(uv-\sqrt{xq})Z
\]
resolves the otherwise unbounded mixed coefficient.  Direct expansion
using $x=\kappa S$ gives the exact identity
\[
 c(a-x)(c-q)Z^2
 =\frac{c-q}{\nu^2}E_{23}^2
  +\frac{2\sqrt{\kappa q}}{\nu}E_{23}B-\kappa B^2.
                                                               \tag{EL18}
\]
Every coefficient on the right is infinitesimal.  To see the equality
without division by a testing coordinate, use
\[
 (ac-xq)Z^2=E_{23}^2/\nu^2+2\sqrt{xq}\,ZE_{23}/\nu,
 \qquad c(a-x)=ac-xq-x(c-q).
\]
All square roots in this calculation are pure-source roots.

For completeness, the single circuit supplies positive fixed pivots
on both $E_{12}$ and $E_{23}$.  If $WZ\ge0$, use the two AM--GM
pairings having respectively these differences and average them;
then $\Psi-\tfrac14(E_{12}^2+E_{23}^2)$ is separated SOS.
If $W,Z$ are both negative, first change the common state sign;
the parent and all squared differences are unchanged.
If $WZ<0$, the negative product term of $\Psi$ is nonnegative,
and
\[
 \sum_{i=1}^4 A_i^2-\tfrac14[(A_1-A_2)^2+(A_2-A_3)^2]
\]
is a real positive-definite quadratic form in the $A_i$.
For example its leading $3\times3$ block has diagonal
$(3/4,1/2,3/4)$ and successive leading minors
$3/4,5/16,3/16$; the fourth diagonal is $1$.
This supplies a finite real Gram factorization.  The remaining
$-4\kappa\nu uvWZ$ is a separated square on this pure-testing sign
chart.  The zero-state boundaries use the same identities.
Since $\lambda_0\sim1/2$ and $\lambda_1\ge0$,
$\lambda-1/4$ is a separated SOS multiplier:
\[
 \lambda-\tfrac14=(\lambda_0-\tfrac14)
       +\lambda_0(2ra+r^2a^2)+\lambda_1c(1+ra)^2.
\]
Thus $\lambda\Psi$ leaves strict real diagonal $1/16$ on each of
$E_{12},E_{23}$ and an explicitly retained separated SOS remainder.
No mixed square root of $\lambda$ has been used.

Perform monic polynomial division in $a$ and then in $c$.
By EL17 there are bounded coefficient polynomials such that
\[
\begin{aligned}
 J&=(a-x)J_a+(c-q)J_c,\\
 H&=(a-x)^2H_{20}(a,c)
       +(a-x)(c-q)H_{11}(c)+(c-q)^2H_{02}(c).
\end{aligned}                                                     \tag{EL19}
\]
The division is by monic polynomials; it introduces no inverse of
$a-x$, $c-q$, or any mixed expression.  The following four particular
coefficient identities are exact:
\[
\begin{aligned}
 H_{20}(a,0)&=-\kappa\nu\ell S\lambda_0r^2,\\
 H_{11}(0)&=-\alpha/q+2\kappa\nu^2S\lambda_0r(1+rx),\\
 H_{02}(0)/S
   &=\nu^2\left[(1+x)\frac\alpha{S\ell^2}
                 -g\nu\lambda_0(1+rx)^2\right]\sim\tfrac12,\\
 H_{02}(c)&=H_{02}(0)+\nu^2S^2dp^2(1+rx)^2c.
\end{aligned}                                                     \tag{EL20}
\]
The first three also follow immediately by evaluating
$H(a,0)=\alpha(1+a)-\kappa\nu\ell S\lambda_0(1+ra)^2$
and using EL17; the last follows from the $c^3$ coefficient.
In particular $H_{11}(0)/(S\ell)$ is bounded and
$H_{02}(c)/S\sim1/2$.  The $g$ terms cancel in the real residue
of the third line by EL0a.  This cancellation is why the independent
ratio $S/\ell$ never enters the estimate.

Write $H_{20}=H_{20}(a,0)+c\widetilde H_{20}$ and
$H_{11}=H_{11}(0)+c\widetilde H_{11}$, with bounded coefficients.
Substitution into EL5 gives the exact scalar assembly
\[
\begin{aligned}
 F_0={}&\lambda\Psi+\widehat A\Gamma^2
                 +\frac{H_{02}(c)}S B^2\\
 &-2\sqrt S J_a\Gamma R_v-2vJ_c\Gamma B
       -g\nu\lambda_0r^2R_0^2+\widetilde H_{20}R_v^2\\
 &+\frac{H_{11}(0)}{S\ell}R_0B\\
 &+\widetilde H_{11}
       \left[\frac{c-q}{\nu^2}E_{23}^2
             +\frac{2\sqrt{\kappa q}}\nu E_{23}B-\kappa B^2\right].
\end{aligned}                                                     \tag{EL21}
\]
This includes every scalar term, including the one involving
$\lambda_1c^2a$.  It supplies strict real pivots on
$(E_{12},E_{23},\Gamma,B)$; the two radial rows have arbitrary
fixed positive pivots available at infinitesimal cost by EL9--11.
There is no discarded source-dependent signed term.

\subsubsection*{Strict bounded Gram and all-test completion}
Use the main rows
\[
 E_{12},\ E_{23},\ \Gamma,\ B,\ \widehat N=\sqrt{1+\beta}N,
 \ N_y
\]
and the three auxiliary rows $V_h,R_v,R_0$.
EL21 and EL13--16 express the entire EL0b as a bounded quadratic
form in these nine rows, together with the retained circuit SOS
remainder.  The real diagonal of the main block is bounded below
by $(1/16,1/16,1/2,1/2,1/2,1/2)$.
Its off-diagonal entries are infinitesimal except for couplings to
the auxiliary rows.  These are bounded, as are the possibly signed
auxiliary diagonal entries.

Choose a sufficiently large fixed real positive diagonal on each
auxiliary row, and allocation it using EL8--11.  This is an exact
SOS operation: if $Q-R^2$ is a displayed SOS budget, replace the
quadratic form by its difference with a real multiple of $Q-R^2$
and retain that multiple as a separate SOS.  EL8 and EL9--10 cost
only infinitesimal entries on the main rows; EL11 is allocationd
last and costs an infinitesimal entry on $R_v$ as well as on $B$.
The real auxiliary diagonals can be chosen large enough that the
real Schur complement is positive definite.  This argument uses
only a finite real matrix, never a mixed-field matrix inverse.
The resulting bounded coefficient Gram has strictly positive real
residue.  The common strict separated bounded-Gram lemma in
Lemma~\ref{p25af2c58da7e:lem:gram} makes the complete form a finite
SOS in $K_c\otimes_{\mathbb R}K_t$.  Restoring the budget and circuit
remainders proves the original full parent.

There is no restriction on $W,Z$, their ratio, or either normal;
$\beta$ was never discarded or divided by.  The actual retained
normal diagonal handles an arbitrarily large $\beta$ before any
cross-term estimate.  The only inverse involving an arbitrarily
small opening is the pure-source scaling used to define $N_y$;
its inverse never appears in a Gram coefficient.  All helper roots
are products of source and testing factors or finite sums thereof.

The local proof includes $u=0$ and $v=0$, since it never divides
by either.  Negative $u$ is reduced by
$(u,Y,Z)\mapsto(-u,-Y,-Z)$ and negative $v$ by
$(v,X,Z)\mapsto(-v,-X,-Z)$; EL0b is invariant under each map.
Real and imaginary state parts are handled separately.
By EL0a both $C,A$ are unbounded, so the adjacent-$C,A$
proof covers noninfinitesimal Alice tests.  Since $\kappa$ and $e$
are infinitesimal, the SE10--SE12 theorem covers every
homogeneous Bob sine-square residue different from $\bar\ell=0$,
including Bob infinity.  The only remaining Bob chart is exactly
finite infinitesimal $v$, proved above.  Therefore all projective
tests and all states on the original $\ell$-infinitesimal source
intersection are covered.  All original source predicates are
retained; no numerical sampling or completion marker
is an input to this conclusion.



% Source formulas: p1de5b12d1927
\subsection{Exact quotients at the doubly infinitesimal source}

\label{component:p1de5b12d1927}

\label{p1de5b12d1927:sec:new-c144-small-nu-bounded-beta}

Keep the original C14.4 source, its full native contact, and all
original coefficient signs.  The common doubly small endpoint has
\[
 S=\sigma+L\in\operatorname{Inf}_{>0},\quad\kappa=x/S\in\operatorname{Inf}_{>0},
 \quad\nu=\sigma/S\in\operatorname{Inf}_{>0},\quad
 \ell=L/S=1-\nu,\quad q=\nu/\ell,\quad n_s=\beta/q^2.
                                                               \tag{SN1}
\]
Here $n_s$ may be bounded or unbounded.  The formulas below are
exact identities; the complete BN and UN proofs supply their distinct
normalized bounds on the bounded and unbounded $\tau$ domains.
Set
\[
 c_0=SC(1-r),\quad a_0=SAr,\quad d_s=D/(Sq^2),\quad
 a_1=\alpha/S,\quad\mu=m/\kappa^2,\quad
 e=\kappa h/\nu,\quad\tau=h/q.                         \tag{SN2}
\]
The actual source-sphere identities are
\[
 d_s=\frac{\ell^2\mathcal H d_0^2}{2(P/S)J},\qquad
 n_s=\frac{\ell^2\mathcal H b_0^2}{2\kappa(P/S)\theta},\qquad
 Jd_s+\kappa\theta n_s=\frac{\ell^2\mathcal H}{2(P/S)}.   \tag{SN3}
\]
The source symbols in SN3 have the definitions in FM2--FM3 and
SE1--SE3.  In particular $e=\kappa\tau/\ell$ exactly.
\paragraph{The complete reciprocal parent.}
For infinitesimal Alice coordinate $u$ and reciprocal Bob coordinate
$w$, including either literal zero, use $D_u=1+u^2,D_w=1+w^2$ and
\[
\begin{split}
 F_w/4={}&m(W-wX+uY-uwZ)^2+CD_w(uW-Y)^2+AD_uD_wY^2\\
 &+(D+\tfrac12)(uwW+uX-wY-Z)^2
 +(\alpha+\tfrac12)D_u(Y-wZ)^2\\
 &+(\beta+\tfrac12)D_w(X+uZ)^2-\tfrac12D_uD_wZ^2.
\end{split}\tag{SN5}
\]
Keep both unrestricted state normals
\[
 X_e=X+ewW,\qquad N_y=(Y-ruW)/\sqrt{Sr(1-r)}.
 \tag{SN6}
\]
Every inverse and radical in these expressions is pure-source or
pure-testing. Write $a=u^2,c=w^2,x=\kappa S,\delta=c-q$ and set
\[
\begin{gathered}
 N=\kappa W-\ell uwZ,\quad G=uW/\sqrt S-\ell\sqrt S\,wZ,\\
 B=\sqrt S\,\delta Z,\quad B_u=u\delta Z,\quad
 R=w(a-x)Z=\frac{\kappa\sqrt S\,G-uN}{\ell}.
\end{gathered}\tag{SN7}
\]
The last equality is the exact homogeneous radial conductor.

\paragraph{Two free helper pivots.}
The exact polynomial identity
\[
 \delta^2N^2+2\kappa\ell\delta wGB
       +2\kappa\ell^2w^2B^2-\ell^2(wB_u)^2
                  =(\kappa\delta W)^2
 \tag{SN8}
\]

is an exact SOS identity.  Whenever a current strict frame controls
$N,G,B$, it supplies a prescribed real pivot on $wB_u$ at
infinitesimal Gram cost.  BN and UN give the required ordered helper
constructions even when the original $d_sB^2$ pivot is absent.

\paragraph{Exact scalar dictionary and the $1/\kappa$ cancellations.}
Put $d=D+1/2$ and define the following polynomials in $a,c$:
\[
\begin{split}
 A_P={}&(1+c)\{c_0(1-r)+a_0r(1+a)\}
       +S\{d(1-r-e)^2c+(\alpha+\tfrac12)(1+a)r^2\},\\
 A_N={}&m(1+ra+ec)^2+(\beta+\tfrac12)e^2c(1+c),\\
 K={}&m(1+ra+ec)+d(1-r-e)
       +(\alpha+\tfrac12)(1+a)r+(\beta+\tfrac12)e(1+c),\\
 F_Z={}&D+\alpha c+\beta a+(m+\alpha+\beta+\tfrac12)ac,\\
 J_1={}&K-\ell A_P,\quad L_1=J_1-(\ell/\kappa)A_N,\\
 H={}&F_Z-\ell^2ScA_P
                 +ac\{(\ell/\kappa)^2A_N-2(\ell/\kappa)J_1\}.
\end{split}\tag{SN9}
\]
These are the exact graph coefficients at $X=-ewW,Y=ruW$;
the negative $dac,e\{2(1-r)-e\}$ contribution has been placed
in $A_P$, not left in $A_N$. The displayed combination is used intact before taking any
normalized bounds.

The actual source satisfies
$L_1(x,q)=H(x,q)=H_a(x,q)=H_c(x,q)=0$ and
$H(a,0)=D+\beta a$. Finite monic division therefore gives
\[
\begin{gathered}
 L_1=(a-x)L_a+\delta L_c,\\
 H=S\delta^2(d_s+\kappa n_s+cH_{02})
        +(a-x)\delta(-qn_s+cH_{11})+c(a-x)^2H_{20}.
\end{gathered}\tag{SN10}
\]
The quotients mean literal polynomial differences at $a=x,c=q$,
not division by a mixed tensor. For example $L_a$ is
$(L_1(a,c)-L_1(x,c))/(a-x)$, and $L_c$ is the quotient of
$L_1(x,c)-L_1(x,q)$ by $c-q$.

Individually $H_{20},H_{11},L_c/\kappa$ need not be bounded.
The complete scalar uses the following combined coefficients:
\[
 C_R=H_{20}+2\ell L_a/\kappa,\quad
 C_G=\kappa(H_{11}-n_s)/\ell,\quad
 C_N=(n_s-H_{11})/\ell-2L_c/\kappa.
 \tag{SN11}
\]
Substitution of $e=\kappa\tau/\ell$ gives the exact cancellations
\[
\begin{split}
 C_R={}&-\ell\mu r(a\ell r+2c\kappa\tau+2\ell),\\
 C_N={}&n_s/\ell-Sa_0\ell r
       +(\beta+m+\tfrac12)\tau^2(c/\ell+1/\ell^2)
                   +2\kappa\mu\tau-m\tau^2/\ell,\\
 H_{02}={}&\kappa\tau^2(\beta+m+\tfrac12).
\end{split}\tag{SN12}
\]
For further pole cancellation, write
$A_P=P_0(a)+cP_1(a)$. Direct differentiation in SN9 gives
\[
\begin{split}
 \kappa H_{11}={}&2\ell^2P_1(x)+\ell^2x a_0r
 +2\ell^2\kappa\mu e(1+2rx)
 +\frac{\ell^2(\beta+1/2)e^2}{\kappa}
 -2\ell(m+\beta+\tfrac12)e\\
 &+(c+q)\frac{\ell^2(m+\beta+1/2)e^2}{\kappa}.
\end{split}\tag{SN13}
\]
These are identities of the complete coefficient combinations.
Their respective normalized estimates are proved in BN and UN.

The entire scalar graph is now
\[
\begin{split}
 F_0={}&(A_N/\kappa^2)N^2+A_PG^2+d_sB^2+n_sB_u^2
       +H_{02}(wB)^2+C_RR^2-2\sqrt S\,L_aGR\\
 &+C_GG(wB)+C_NN(wB_u).
\end{split}\tag{SN14}
\]
For verification the mixed identity is
$c(a-x)\delta Z^2=\{\kappa G(wB)-N(wB_u)\}/\ell$.
Combining this with $uN=\kappa\sqrt S\,G-\ell R$ proves
SN14 directly from SN9--SN10.

\paragraph{Both complete state normals.}
Define
\[
\begin{split}
 A_X={}&-m(1+ra+ec)+da(1-r-e)-e(\beta+\tfrac12)(1+c),\\
 T_X={}&(\ell/\kappa)cA_X+\beta-D+c(m+\beta+\tfrac12),\\
 Y_P={}&(1+c)\{-c_0+a_0(1+a)\}
 +S\{m(1+ra+ec)-dc(1-r-e)
                    +(\alpha+\tfrac12)(1+a)r\},\\
 Y_C={}&\ell Y_P-ma+d-(\alpha+\tfrac12)(1+a).
\end{split}\tag{SN15}
\]
Their actual contact values satisfy $T_X(x,q)=Y_C(x,q)=0$.
Use finite polynomial quotients
\[
 T_X=c(a-x)X_a+\delta\{(D-\beta)/q+cX_c\},\qquad
 Y_C=(a-x)Y_a+\delta Y_c.
 \tag{SN16}
\]
Although $X_a,A_X/\kappa$ separately may be unbounded, their
complete conductor has
\[
\begin{gathered}
 \widehat A_X=A_X/\kappa-aX_a/\ell
 =-\kappa\mu-\frac\tau\ell
                  \{(\beta+\tfrac12)(1+c)+mc\},\\
 \widehat X_a=\kappa X_a=\ell\{d(1-r-e)-mr\},\qquad
 X_c=-\tau(m+\beta+\tfrac12).
\end{gathered}\tag{SN17}
\]
The exact quarter-parent is $F_w/4=F_0+H_XX_e^2+H_YN_y^2+
2H_{XY}X_eN_y+2L_XX_e+2L_YN_y$, with
\[
\begin{split}
 H_X={}&mc+da+(\beta+\tfrac12)(1+c),\\
 H_Y={}&rc_0(1+c)+(1-r)a_0(1+a)(1+c)
       +Sr(1-r)\{ma+dc+(\alpha+\tfrac12)(1+a)\},\\
 H_{XY}={}&-uw\sqrt{Sr(1-r)}(m+d),\\
 L_X={}&w\widehat A_XN+
        \frac{uw\sqrt S}{\ell}\widehat X_aG
        +d_s\sqrt S\,quB-n_sqB_u+wX_c(wB_u),\\
 L_Y={}&\sqrt{r(1-r)}\{Y_PG+\sqrt S\,Y_aR+Y_c(wB)\}.
\end{split}\tag{SN18}
\]

Every formula retains the original $r$ and $1-r$ openings and both
complete state normals.  In particular, the term $n_sqB_u$ may be
written $q\sqrt{n_s}(\sqrt{n_s}B_u)$ when that row is used;
literal $n_s=0$ needs no division.  BN and UN allocate these actual
rows, with their own weights, instead of treating separate uncanceled
coefficients as bounded.

\paragraph{Explicit finite divisions.}
All quotient symbols above can be defined directly in the polynomial
ring in $a,c$ over the source field.  In addition to $L_a,L_c$
specified after SN10, use
\[
\begin{aligned}
 H_{20}&=\frac{H(a,c)-H(x,c)-(a-x)H_a(x,c)}{c(a-x)^2},\\
 H_{11}&=\frac{H_a(x,c)/(c-q)+qn_s}{c},\\
 H_{02}&=\frac{H(x,c)/(S(c-q)^2)-d_s-\kappa n_s}{c},\\
 X_a&=\frac{T_X(a,c)-T_X(x,c)}{c(a-x)},\qquad
 X_c=\frac{T_X(x,c)/(c-q)-(D-\beta)/q}{c},\\
 Y_a&=\frac{Y_C(a,c)-Y_C(x,c)}{a-x},\qquad
 Y_c=\frac{Y_C(x,c)-Y_C(x,q)}{c-q}.
\end{aligned}
\]
These displays specify polynomial division, not evaluation of an
inverse of $a-x$, $c-q$ or $c$.  Divisibility follows from
$H(x,q)=H_a(x,q)=H_c(x,q)=0$, $H(a,0)=D+\beta a$,
$T_X(x,q)=Y_C(x,q)=0$, and the explicitly affine dependence of
$T_X$ on $a$.  Thus all exact coefficients are defined without an
external program.  Pure-source divisions such as $1/S$ or $1/q$
are retained until their cancellations in the printed combinations.



% Source formulas: p80b0c5b9879f
\subsection{The doubly infinitesimal source with bounded $\tau$}

\label{component:p80b0c5b9879f}

\label{p80b0c5b9879f:sec:new-c144-bounded-h-ratio}

\paragraph{Source scope and result.}
Retain the original entered physical C14.4 source, all strict gaps and
all native source inequalities. Assume
\[
 S=\sigma+L,\ \kappa=x/S,\ \nu=\sigma/S
       \in\operatorname{Inf}_{>0},\qquad
 \ell=1-\nu,\quad q=\nu/\ell,\quad
 \tau=h/q\ \hbox{is bounded}.
 \tag{BN1}
\]
The complete original parent is a finite separated tensor SOS for
every test and every real or complex state on BN1. The quantity
$\beta/q^2$ is unrestricted, including zero and unbounded values.
The proof below is a circuit completion of SN, not an import of
a missing endpoint proof.

\paragraph{Exact source estimates.}
Use the definitions of SN2, except that $n_s=\beta/q^2$ need not be
bounded. The exact sphere SN3 remains valid. Put $b=\beta/q$ and
$z=\kappa\beta/q^2$. Then
\[
\begin{gathered}
 h=q\tau\in\operatorname{Inf},\quad\theta\sim1,\quad
 e=\kappa\tau/\ell\in\operatorname{Inf},\quad M\sim1,\\
 b=\frac{\tau(B/S)-\ell H_0}{2\theta(P/S)},\qquad
 \mu\ell^2=\frac{\tau(1+x)(1+\sigma r)}{2M(P/S)},\\
 Jd_s+\theta z=\frac{\ell^2\mathcal H}{2(P/S)}\sim\tfrac12.
\end{gathered}\tag{BN2}
\]
In particular $b,z,d_s,\mu,\tau$ are bounded, $\tau$ and $\mu$
have positive real lower bounds, and at least one of $d_s,z$ is a
positive bounded unit. The exact inequality $\beta\ge0$ supplies
the lower bound on $\tau$. Also
$c_0,a_0,a_1\sim1/2$, $D,\alpha,m,\beta$ are infinitesimal.
Although the two terms below need not have small difference before
subtraction, their displayed bounded formulas in BN2 give
\[
                  A_{00}:=\mu\ell^2-b\sim\tfrac12.
 \tag{BN3}
\]
This argument takes residues only of bounded source coefficients and
never uses an estimate divided by $q$. It allows arbitrary relative
rates of $S,\kappa,q$.

\paragraph{A retained circuit and two actual positive rows.}
Use the full reciprocal parent SN5, the two unrestricted normals SN6,
and the shifted exact scalar dictionary SN9. Let $\zeta=\ell/\kappa$,
$a=u^2,c=w^2$, $x=\kappa S$, and assume $u,w\ge0$ by the original
sign gauges. Both may be zero. For every state define
\[
\begin{array}{ll}
 A_1=\sqrt\beta\,uZ,&A_2=\sqrt\beta\,W/(\zeta\sqrt q),\\
 A_3=\sqrt\beta\sqrt{x/q}\,wZ,
 &A_4=\sqrt\beta\,uwW/(\zeta q\sqrt x),
\end{array}\tag{BN4}
\]
\[
 \Psi=\sum_{i=1}^4 A_i^2-\frac{4\beta}{\zeta q}uwWZ,
 \qquad E_1=A_1-A_2,\quad E_2=A_3-A_4.
 \tag{BN5}
\]
Every monomial and its radical is a product of a pure-source and a
pure-testing element. This remains true if $x/q$ is unbounded.
The exact inequality meant here is a separated SOS certificate:
\[
                    \Psi-\tfrac12(E_1^2+E_2^2)\in\Sigma T^2.
 \tag{BN6}
\]
Indeed, if $WZ\ge0$, then $A_1A_2,A_3A_4\ge0$ and
\[
 \Psi=E_1^2+E_2^2+
              2(\sqrt{A_1A_2}-\sqrt{A_3A_4})^2.
\]
The radicals displayed here split into pure-source and pure-testing
radicals. If $WZ<0$, use instead
\[
 \Psi-\tfrac12(E_1^2+E_2^2)
 =\tfrac12(A_1+A_2)^2+\tfrac12(A_3+A_4)^2
                  +\frac{4\beta}{\zeta q}|uwWZ|.
\]
The last positive separated monomial is a square. Thus BN6 is valid
for all states, including $WZ=0$. At the actual source contact every
$A_i$ agrees, so $\Psi$ and its first contact derivatives vanish.

\paragraph{The exact scalar after circuit subtraction.}
Write $N_0=W/\zeta-uwZ=N/\ell$, where $N$ is SN7, and keep
$G,B,B_u,R$ from SN7. Thus
\[
 R=\frac{\sqrt S}{\zeta}G-uN_0.
 \tag{BN7}
\]
Subtract the entire circuit BN5 from the exact scalar graph in SN9.
Define
\[
\begin{gathered}
 A=\zeta^2 A_N-b,\qquad A_P'=A_P-zc/\ell^2,\\
 K'=K-2\beta/(\zeta q),\quad J'=K'-\ell A_P',
 \quad L'=J'-A/\zeta,\\
 H'=F_Z-\beta a-\beta xc/q-\ell^2ScA_P'
                          +ac(A-2\zeta J').
\end{gathered}\tag{BN8}
\]
Then $F_0-\Psi=A N_0^2+A_P'G^2-2\zeta uwZL'N_0+H'Z^2$.
All coefficients of $A$ and $A_P'$ are bounded:
\[
 A=\mu\ell^2(1+ra+ec)^2
       +(\beta+\tfrac12)\tau^2c(1+c)-b,
 \qquad \overline A=\overline{A_P'}=\tfrac12.
 \tag{BN9}
\]
Here and below the residue is taken at infinitesimal $u,w$.

The circuit subtraction has especially simple exact effects on the
SN10 quotient dictionary:
\[
\begin{gathered}
 L'_a=L_a,\qquad L'_c=L_c+\kappa n_s/\ell,\qquad
 H'_{20}=H_{20},\quad H'_{02}=H_{02},\quad
 H'_{11}=H_{11}-2n_s,\\
 H'=S\delta^2(d_s+cH_{02})
               +c(a-x)\delta H'_{11}+c(a-x)^2H_{20}.
\end{gathered}\tag{BN10}
\]
In particular no coefficient $n_s$ remains on its own in the final
conductors. With $C_R$ and $L_a$ from SN11--SN13 put
\[
 C_G'=\kappa(H_{11}-2n_s)/\ell,\qquad
 C_N'=-H_{11}-2\ell L_c/\kappa.
 \tag{BN11}
\]
These are bounded. The cancellation for the second coefficient is
explicit:
\[
 C_N'=-Sa_0\ell^2r+(\beta+m+\tfrac12)\tau^2(c+1/\ell)
                 +2\ell\kappa\mu\tau-m\tau^2.
 \tag{BN12}
\]
For $C_G'$ use the bounded formula SN13 and $\kappa n_s=z$.
The exact full scalar is consequently
\[
\begin{split}
 F_0={}&\Psi+A N_0^2+A_P'G^2+d_sB^2+H_{02}(wB)^2
       +C_RR^2-2\sqrt S\,L_aGR\\
 &+C_G'G(wB)+C_N'N_0(wB_u).
\end{split}\tag{BN13}
\]
This is an identity in unrestricted state coordinates, not a
restriction to a contact graph in the full parent.

\paragraph{Free helper rows across the whole source sphere.}
First reserve the positive rows $E_1,E_2$ from BN6. If $d_s$ is a
unit, $wB$ has infinitesimal cost on $\sqrt{d_s}B$. Otherwise
$z$ is a positive bounded unit. The exact identity
\[
 E_2=-\sqrt z\,\{(w/\ell)G+
                          \sqrt S\,w(w-\sqrt q)Z\}
\]
then gives
\[
 wB=-\frac{w+\sqrt q}{\sqrt z}E_2
                   -\frac{w(w+\sqrt q)}\ell G.
 \tag{BN14}
\]
Every coefficient in BN14 is infinitesimal, so any fixed real pivot
on $wB$ is free at infinitesimal Gram cost. Only the chosen pure-source
unit $\sqrt z$ is inverted. Literal $\beta=0$ uses the first branch.

After this pivot has been allocationd, the normalized SN8 identity is
\[
 \delta^2N_0^2+\frac{2\kappa\delta}{\ell}G(wB)
          +2\kappa(wB)^2-(wB_u)^2
                         =(\delta W/\zeta)^2.
 \tag{BN15}
\]
It supplies any fixed pivot on $wB_u$ at infinitesimal cost. The order of
these two allocations is deliberate; no weak $B$ diagonal is needed
in the second step.

\paragraph{Both full normals and their only additional row.}
Use the complete $H_X,H_Y,H_{XY},Y_P,Y_a,Y_c$ from SN18 and its
bounded $\widehat A_X,\widehat X_a,X_c$ from SN17. These formulas
only require $\tau,\mu,\beta$ bounded, not $n_s$ bounded.
The entire first-normal part is
\[
\begin{split}
 L_X={}&\ell w\widehat A_XN_0
       +\frac{uw\sqrt S}{\ell}\widehat X_aG
       +d_s\sqrt S\,quB-bB_u+wX_c(wB_u),\\
 L_Y={}&\sqrt{r(1-r)}\{Y_PG+\sqrt S\,Y_aR+Y_c(wB)\}.
\end{split}\tag{BN16}
\]
The apparent extra $B_u$ coupling is paid by the actual circuit row
$E_1$, using the exact, zero-face-safe identity
\[
 bB_u=-(w+\sqrt q)\sqrt b\,E_1-b(w+\sqrt q)N_0.
 \tag{BN17}
\]
All coefficients on the right have zero residue because $b$ is
bounded. No inverse of $b$ or $\beta$ is used. In the $B$ term in
$L_X$, write $d_sB=\sqrt{d_s}(\sqrt{d_s}B)$; this is likewise
face-safe when $D=0$. All direct normal coefficients are
infinitesimal, and $H_X,H_Y\sim1/2$, $H_{XY}$ is infinitesimal.
There is no missing $X$ or $Y$ normal and no imposed state ratio.

\paragraph{Allocation and all-tests conclusion.}
Keep a fixed positive fraction of each circuit row in BN6.
Use the strict base rows $(N_0,G,X_e,N_y,\sqrt{d_s}B,E_1,E_2)$,
omitting any zero row. Their main four diagonal residues are $1/2$;
the last three have fixed positive pivots wherever their rows are
nonzero. Terms containing $R$ are infinitesimal matrix perturbations
by BN7. Purchase sufficiently large fixed pivots on $wB,wB_u$ by
BN14--BN15, and use them for the bounded cross terms in BN13 and
BN16. Formula BN17 pays the remaining first-normal row at
infinitesimal cost. Every allocation is taken from an explicitly
retained positive square or from the exact SOS remainder BN6.
The resulting bounded Gram has strictly positive definite real
residue. The strict bounded-Gram lemma therefore gives
finite tensor SOS for SN5 with both unrestricted normals restored.

This proves all infinitesimal $u,w$, including the literal axes.
For noninfinitesimal Alice use the unchanged adjacent-$C,A$ theorem
CE1--CE10; $C,A$ are unbounded on BN1. For the other Bob charts use
the small-$e$ unmatched theorem SE10--SE12, since
$\overline\ell=1$. The sign gauges and real/imaginary state split
are the original exact ones. The local argument already handles
either sign of $WZ$, so no state normalization or state-chart
exclusion is required. Both literal $D=0$ and $\beta=0$ are covered
whenever they intersect BN1.

\paragraph{The other source cells.}
The unbounded-$h/q$ source and the $\ell$-infinitesimal endpoint are proved in Sections~\ref{component:peaa56c1a3c7a} and~\ref{component:pb96c9bfcc12d}. Their source partition with the present bounded-$h/q$ theorem is printed in Section~\ref{component:pfd0dd7bb10cf}.



% Source formulas: peaa56c1a3c7a
\subsection{The doubly infinitesimal source with unbounded $\tau$}

\label{component:peaa56c1a3c7a}

\label{peaa56c1a3c7a:sec:new-c144-unbounded-h-ratio}

\paragraph{Exact scope.}
Retain the full original entered C14.4 source, all its strict gaps and
native source inequalities. Suppose $S,\kappa,\nu$ are positive
infinitesimals, $\ell=1-\nu$, $q=\nu/\ell$, and
\[
 \tau=h/q\ \hbox{is unbounded},\qquad
 \theta=1-h>0,\quad e=\kappa h/\nu\in\operatorname{Inf},\quad
 \zeta=\ell/\kappa.
 \tag{UN1}
\]
This section proves finite separated tensor SOS for the complete
parent at every projective test and every state on UN1. The opening
$\theta\in\operatorname{Inf}$ is included. There is no state-ratio
restriction or source/testing relative-rate assumption.

\paragraph{Normalized source estimates and a canceled radial mass.}
Keep $c_0=SC(1-r)$, $a_0=SAr$, $d_s=D/(Sq^2)$,
$a_1=\alpha/S$, $P_0=P/S$, and the exact SE2--SE3 source sphere.
Define
\[
 d=D+\tfrac12,\quad
 n=\beta\theta^2+2Dh+\tfrac12h^2>0,\quad
 z_n=\kappa n/q^2,\quad A_{00}=m\zeta^2-n/q.
 \tag{UN2}
\]
Here $n$ is a scalar coefficient, not an omitted event weight.
It is bounded even when $\beta$ is unbounded. The exact source gives
\[
\begin{gathered}
 m\zeta^2=\frac{\tau(1+x)(1+\sigma r)}{2MP_0},\qquad
 m\zeta\in\operatorname{Inf},\quad D,\alpha,m\in\operatorname{Inf},\quad
 c_0,a_0,a_1\sim\tfrac12,\\
 e\tau=\frac{\ell\{eH_0+h\mathcal H b_0^2\}}{B/S}\ \hbox{is bounded},\\
 z_n=\frac{\ell^2\mathcal H\theta b_0^2}{2P_0}
 +\frac{\ell^2(eH_0+h\mathcal H b_0^2)}{2(B/S)}
 +2\kappa S d_s h.
\end{gathered}\tag{UN3}
\]
Consequently $d_s,z_n$ are bounded and
$d_s+z_n\sim1/2$; at least one is a positive bounded unit. This
includes literal $D=0$. Literal $\beta=0$ has no points on UN1,
because $\beta\ge0$ has the exact boundary equation
$\tau(B/S)=\ell H_0$ when $\beta=0$.

For completeness, the following finite source division proves the
bound on $A_{00}$ without dividing an error by $\theta$ or $q$.
Put $U=(1+x)(1+\sigma r)$ and
$M_*=\ell(1+xr)-\kappa\theta=\ell M$, and form the bounded
source polynomial
\[
 T=\ell JU+JM_*P_0-JM_*\kappa\theta(J-\sigma)-2M_*U.
\]
In this expression $P_0,J,U,M_*$ are their exact native polynomial
formulas in $\kappa,\nu,S,h,r$, with $\ell=1-\nu$. Thus
\[
\begin{gathered}
 T=\kappa T_\kappa+\ell\nu S T_\nu,\qquad
 T_\nu=r-1-\nu-\nu S(1-r)(r+\ell),\\
 T_\kappa=\{T-T|_{\kappa=0}\}/\kappa
             \quad\hbox{is a bounded source polynomial},\\
 A_{00}=\frac{\theta\ell H_0}{2P_0}
            +\frac{eT_\kappa+h\ell S T_\nu}{2JMP_0}.
\end{gathered}\tag{UN4}
\]
The quotient is a finite monic polynomial division in a pure source
variable. All denominator factors are positive bounded units.
Hence $A_{00}$ is bounded and $A_{00}-\theta/2$ is infinitesimal.
Moreover $n/q=m\zeta^2-A_{00}\sim\tau/2$ is positive unbounded.
If $\theta$ is infinitesimal, then $n\sim1/2$ and
$\kappa/q=e\ell/h$ is infinitesimal. These facts supply the two
source subcases used below.

\paragraph{The actual circuit and three reserved positive rows.}
Use the complete original reciprocal quarter-parent SN5. For
infinitesimal $u,w\ge0$, including zero, retain the unrestricted
normals
\[
 N_x=X+huZ,\qquad N_y=(Y-ruW)/\sqrt{Sr(1-r)}.
 \tag{UN5}
\]
Put $a=u^2,c=w^2,x=\kappa S,\delta=c-q$ and
\[
\begin{gathered}
 N_0=W/\zeta-uwZ,\quad G=uW/\sqrt S-\ell\sqrt S\,wZ,\\
 B=\sqrt S\delta Z,\quad B_u=u\delta Z,\quad
 R=w(a-x)Z=\sqrt S\,G/\zeta-uN_0,\quad
 G_h=h(a-xc/q)Z.
\end{gathered}\tag{UN6}
\]
Define $A_1,A_2,A_3,A_4$ by BN4 with $\beta$ replaced by the
positive $n$, and set
\[
 \Psi_n=\sum_iA_i^2-4nuwWZ/(\zeta q),\qquad
 E_1=A_1-A_2,\quad E_2=A_3-A_4,\quad E_3=A_1-A_3.
 \tag{UN7}
\]
The explicit separated AM--GM square decompositions in BN6 give
\[
             \Psi_n-\tfrac14(E_1^2+E_2^2+E_3^2)\in\Sigma T^2.
 \tag{UN8}
\]
To verify the simultaneous reserve, average BN6 with the pairing
$\Psi_n-(A_1-A_3)^2-(A_2-A_4)^2\in\Sigma T^2$.
For the latter pairing the products $A_1A_3,A_2A_4$ are nonnegative;
if $WZ<0$, the additional signed cross term is a positive separated
monomial. Thus UN8 holds for both signs of $WZ$, as well as its
zero face. No radical of a general mixed sum is used.

\paragraph{Exact scalar subtraction and radial cancellation.}
The complete scalar graph at $N_x=N_y=0$ is
\[
 F_0=m(1+ra)^2W^2+A_PP^2+F_ZZ^2-2KuwWZ,
 \quad P=uW/\sqrt S,
 \tag{UN9}
\]
where
\[
\begin{split}
 A_P={}&(1+c)\{c_0(1-r)+a_0r(1+a)\}
       +S\{d(1-r)^2c+(\alpha+\tfrac12)(1+a)r^2\},\\
 K={}&m\theta(1+ra)+d(1-r)(1+ha)
                         +(\alpha+\tfrac12)(1+a)r,\\
 F_Z={}&D+\alpha c+na+
          \{m\theta^2+\alpha+(\beta+\tfrac12)\theta^2\}ac
                                                 +dh^2a^2.
\end{split}\tag{UN10}
\]
Subtract $\Psi_n$, put
\[
\begin{gathered}
 A=m\zeta^2(1+ra)^2-n/q,\quad A_P'=A_P-z_nc/\ell^2,\\
 K'=K-2n/(\zeta q),\quad J'=K'-\ell A_P',\quad
 L'=J'-A/\zeta,\\
 H'=F_Z-na-nxc/q-\ell^2ScA_P'+ac(A-2\zeta J').
\end{gathered}\tag{UN11}
\]
The exact residual graph is $A N_0^2+A_P'G^2-
2\zeta uwZL'N_0+H'Z^2$. At the actual source contact
$L'(x,q)=H'(x,q)=H'_a(x,q)=H'_c(x,q)=0$.
Also $H'(a,0)=D+dh^2a^2$. Finite monic contact division gives
\[
\begin{gathered}
 L'=(a-x)L_a+\delta L_c,\\
 H'=(D+dh^2x^2)\delta^2/q^2
       +(a-x)\delta(-2dh^2x/q+cH_{11})
       +(a-x)^2(dh^2+cH_{20}).
\end{gathered}\tag{UN12}
\]
The otherwise possible third-$c$ quotient is exactly zero.
Define $P_1(a)$ to be the coefficient of $c$ in $A_P$. The complete
combined coefficients have short bounded formulas:
\[
\begin{split}
 L_a={}&m\theta r+dh(1-r)+(\alpha+\tfrac12)r
       -\ell\{(1+c)a_0r+S(\alpha+\tfrac12)r^2\}\\
       &-m\zeta r\{2+r(a+x)\},\\
 C_R={}&H_{20}+2\zeta L_a=-m\zeta^2r(2+ra),\\
 C_G={}&H_{11}/\zeta=2\ell P_1(x)+\ell x a_0r-2z_n/\ell,\\
 C_N={}&-H_{11}-2\zeta L_c=-Sa_0\ell^2r.
\end{split}\tag{UN13}
\]
Only $C_R$ may be unbounded separately. Its entire contribution
must be combined before estimating: $A+aC_R=A_{00}$ exactly.
Using $R=\sqrt S\,G/\zeta-uN_0$ therefore gives
\[
\begin{split}
 F_0={}&\Psi_n+A_{00}N_0^2+\widetilde A_PG^2
             +2u\sqrt S\,(L_a-C_R/\zeta)N_0G\\
 &+d_sB^2+dG_h^2+C_GG(wB)+C_NN_0(wB_u),\\
 \widetilde A_P={}&A_P'+SC_R/\zeta^2-2SL_a/\zeta\sim\tfrac12.
\end{split}\tag{UN14}
\]
Every coefficient in UN14 is bounded after this combination:
$C_R/\zeta=-m\zeta r(2+ra)$ is infinitesimal and
$C_R/\zeta^2=-mr(2+ra)$ is bounded. The $N_0G$ coefficient
is infinitesimal. No estimate of an unbounded isolated radial
coefficient is used.

\paragraph{Complete $X$ conductor and its exact positive completion.}
The full $N_x$ hard coefficient is
$H_X=mc+da+(\beta+1/2)(1+c)$. Its complete half-linear
coefficient, after using the actual contact identity, is
\[
\begin{split}
 L_X={}&-m\zeta wN_0+d\tau a(wN_0+B_u)
        +uw\sqrt S\{d(1-r-e)-mr\}G+t_0B_u,\\
 t_0={}&\frac{D+h/2-\beta\theta}{q}
       =D/q+\frac{\ell H_0-\kappa\tau(J-\sigma)}{2P_0}
          \sim\tfrac12.
\end{split}\tag{UN15}
\]
The exact source identity in the second line cancels the potential
$1/q$ pole. Here $D/q=S q\,d_s$ is infinitesimal. The single
unbounded $\tau$ term is removed with the actual square $dG_h^2$,
not estimated. Namely
\[
 G_h=-\tau u(wN_0+B_u)+ew\sqrt S\,G,\quad
 N_x+\tau(wN_0+B_u)=X+ewW,
\]
\[
\begin{split}
 &dG_h^2+daN_x^2+2d\tau a(wN_0+B_u)N_x\\
 &\quad=da(X+ewW)^2
       -2de\tau uw\sqrt S\,(wN_0+B_u)G+de^2Sw^2G^2.
\end{split}\tag{UN16}
\]
The remainder has bounded infinitesimal coefficients on $N_0,G,wB_u$
because $e\tau$ is bounded. Retain the entire positive square on the
right and also the whole term $\beta(1+c)N_x^2$. The remaining
$N_x$ diagonal is $mc+(1+c)/2\sim1/2$. The remaining $L_X$ is
UN15 with its $d\tau a$ term removed; all its direct coefficients
are infinitesimal except the bounded coefficient $t_0$ on $B_u$.
The construction thus includes arbitrarily large $\beta$.

\paragraph{The unrestricted $Y$ normal and cross coefficient.}
Put
\[
\begin{split}
 Y_P={}&(1+c)\{-c_0+a_0(1+a)\}
 +S\{m(1+ra)-d(1-r)c+(\alpha+\tfrac12)(1+a)r\},\\
 Y_C={}&\ell Y_P-m\theta a+d(1+ha)-(\alpha+\tfrac12)(1+a).
\end{split}\tag{UN17}
\]
The actual source has $Y_C(x,q)=0$, so finite monic division gives
$Y_C=(a-x)Y_a+\delta Y_c$ with bounded polynomial coefficients.
The complete remaining entries are
\[
\begin{split}
 L_Y={}&\sqrt{r(1-r)}\{Y_PG+\sqrt S\,Y_aR+Y_c(wB)\},\\
 H_Y={}&rc_0(1+c)+(1-r)a_0(1+a)(1+c)\\
 &+Sr(1-r)\{ma+dc+(\alpha+\tfrac12)(1+a)\}\sim\tfrac12,\\
 H_{XY}={}&-uw\sqrt{Sr(1-r)}(m+d)\in\operatorname{Inf}.
\end{split}\tag{UN18}
\]
In particular $Y_P$ is infinitesimal. Every coefficient is bounded
through $r=0$ and $r=1$ openings; no inverse of $r$ or $1-r$ remains
after the displayed pure-source normal scaling. UN14--UN18 account
for all terms of SN5, including both independent state normals.

\paragraph{Free pivots, including the $\theta$-infinitesimal opening.}
Reserve the three actual circuit rows in UN8. In either sphere
chart, $wB$ is free at infinitesimal cost: use $\sqrt{d_s}B$ when
$d_s$ is a unit, and otherwise use the exact identity
\[
 wB=-\frac{w+\sqrt q}{\sqrt{z_n}}E_2
                 -\frac{w(w+\sqrt q)}\ell G.
 \tag{UN19}
\]
If $\theta$ is a unit, $A_{00}$ already supplies a strict $N_0$
pivot. Since $q/n$ is infinitesimal, the exact identity
\[
 B_u=-(w+\sqrt q)\sqrt{q/n}\,E_1-(w+\sqrt q)N_0
 \tag{UN20}
\]
supplies any fixed $B_u$ pivot at infinitesimal cost. Then $wB_u$ is
free by multiplication by the pure-testing infinitesimal $w$.

If $\theta$ is infinitesimal, first supply the missing $N_0$ pivot
from the actual circuit, using
\[
 N_0=\sqrt{q/n}(E_2-E_1)
       -\frac{w-\sqrt q}{\sqrt n}E_3
       +\frac{\sqrt{\kappa/q}}\ell\,wG.
 \tag{UN21}
\]
All coefficients are infinitesimal: $n\sim1/2$ and $\kappa/q$
is infinitesimal by UN4. Identity UN21 holds without requiring
$z_n$ to be a unit; the apparent source factor in UN19 has
canceled exactly before estimating. After this allocation, UN20
supplies $B_u$, and multiplication by $w$ supplies $wB_u$. Thus the
order of allocations involves no circular use of a missing pivot.
The bounded, possibly signed $A_{00}$ in this source subcase
has zero residue and is absorbed after the strict $N_0$ pivot is
obtained. No inverse of $\theta$ is introduced.

\paragraph{Final allocation and unchanged testing joins.}
Use the strict rows $G,N_x,N_y$, the applicable $N_0$ pivot, and
fixed retained fractions of $E_1,E_2,E_3,\sqrt{d_s}B$ wherever
nonzero. Purchase sufficiently large real pivots on $B_u,wB,wB_u$
by UN19--UN21. Every allocation has infinitesimal Gram cost on
the preceding strict rows. All signed terms in UN14--UN18 are
then either infinitesimal matrix perturbations or bounded
couplings to those freely allocationd helper pivots. A fixed
sufficiently large choice makes the complete residue matrix
strictly positive definite. The separated strict bounded-Gram
lemma proves finite SOS. Restore the positive terms retained
in UN8 and UN16, including $\beta(1+c)N_x^2$; nothing from the
original parent is omitted.

This proves all infinitesimal $u,w$, their literal zero faces and
arbitrary states. Noninfinitesimal Alice tests use CE1--CE10 because
$C,A$ are unbounded. The other Bob tests use the SE10--SE12
unmatched certificate, since $e$ is infinitesimal and
$\overline\ell=1$. The original sign gauges and real/imaginary state
split are exact; the circuit argument treats both signs of $WZ$.
These joins give the claimed all-tests theorem on UN1.

The theorem just proved has precisely the source domain UN1. Its join with the bounded-$h/q$ theorem and the other C14.4 source cells is printed in Section~\ref{component:pfd0dd7bb10cf}.



\section{Completion of the C1 certificate}
The source partition in Section~\ref{sec:c1-source-partition} is
exhausted. For unbounded $\eta$ and bounded $z=1/(x\zeta)$,
Theorem~\ref{thm:c1-unification-bounded-z-all-H} gives the local
certificate for every $H=h/\sigma$; the unit-Alice and complementary
projective charts cover the other tests. When $L$ is infinitesimal,
the positive-$j$ completions, the completed $J$ relation and the five
normalized source domains give the remaining certificates. Each
represents the full state quadratic form, with its positive event
terms restored.

The real statement applies separately to the real and imaginary state
parts because the parent is a real quadratic form. Source and testing
homogenization factors are pure-field squares. Restoring the original
positive $d$ therefore proves the complete projective C1 assertion.

\chapter{The C2 cases}
\label{ch:c2-short}
All source coefficients lie in the real closed field $K_c$, and all
testing measurements and states lie in the independently labelled real
closed field $K_t$. We use $K_c\otimes_{\mathbb R}K_t$ without an
algebraic-disjointness assumption. A retained positive square may have
unbounded pure-source coefficients. Only the matrix coefficients to
which the bounded strict-Gram lemma is applied must be bounded.
The native physical entry, bounded-coefficient theorem, and arbitrary-
scale noncommuting-residue theorem are the common preceding results.
Normalize $d=1$ by a pure-source square and write $P=F/4$. The common
four-state parent is
\begin{equation}
\begin{split}
F=4P={}&4nD_uD_vX^2+4mM^2+4CD_v(uW-Y)^2+4AD_uD_vY^2\\
&+4DJ^2+4\alpha D_u(vY-Z)^2\\
&+2\{J^2+D_u(vY-Z)^2+D_v(X+uZ)^2-D_uD_vZ^2\},\\
M={}&vW-X+uvY-uZ,\quad J=uW+uvX-Y-vZ,\quad
D_u=1+u^2,\quad D_v=1+v^2.
\end{split}\label{eq:c2-master-parent}
\end{equation}
The finite coefficient-extraction and monic-division operations are
those in Section~\ref{sec:c1-finite-operations}; they are field-independent.
For every displayed state substitution $s=s_0+A\xi$, every scalar,
complete first normal and hard coefficient below means respectively
$[1]\mathcal Q$, $[\xi_i]\mathcal Q$, and
$\frac12\partial_i\partial_j\mathcal Q$, where
$\mathcal Q=(F-\Pi)(s_0+A\xi)$ and $\Pi$ is the displayed sum of retained
positive squares. Quadratic state degree makes this the exact whole
polynomial. Coefficients need not be expanded to define them. All
support, divisibility, source-bound and residue verifications remain
part of the relevant calculation below.

The proof follows the source geometry. We first determine the stationary
coefficients and treat two weak complementary events by a weighted radial
completion. The remaining families are organized by the scales of the actual
angles. For one small angle the state parameter separates bounded and
unbounded source scales; for two small angles the source ratios define
R1--R10. Within each family the scalar calculation is completed by the
state-normal and complementary-testing arguments.

\section{Stationary sources and their classification}
\subsection{Source domains}
\label{component:pe351eb06b28b}
\label{pe351eb06b28b:sec:c2-complete-source-routing}
The native entry theorem supplies an actual physical stationary source
with $d=1$, $\lambda_{\rm event}=\beta=0$, $n\ge0$, and positive
$m,C,A,D,\alpha$. In the native coordinates $u_c=c$, $v_c=ac$ its
strict gaps are
\[
 c,a,b,t>0,\quad 0<r<1,\quad h=b-a^2t>0,\quad
 g=(1+c^2)b-(1-r+a^2c^2t)>0,
\]
\[
 j=1+c^2r-t>0,\qquad1-c^2h>0,\qquad b-a^2c^2r>0.
\]
These are the original source equations. If
$(1+n+m+C+A+D+\alpha)^{-1}$ is a unit, the bounded-coefficient
theorem applies. Henceforth that reciprocal is infinitesimal.
Call an event weak when its coefficient is infinitesimal and strong
otherwise; strong coefficients may be unbounded.

The exact reflection
\begin{equation}
 (u,v;W,X,Y,Z;n,m,C,A,D,\alpha)
 \longmapsto(1/v,1/u;Z,X,Y,W;n,D,\alpha,A,m,C)
 \label{eq:c2-short-reflection}
\end{equation}
satisfies $P=u^2v^2P'$ and is a homogeneous polynomial identity at
the literal endpoints. With $z=c^2$ its source map is
\[
 c'=(ac)^{-1},\ a'=a,\ b'=a^2/b,\ r'=a^2zr/b,
 \quad t'=(1-zh)/b,
\]
\[
 h'=a^2zh/b,\quad g'=a^2j/b,\quad j'=g/b,
 \quad\Delta'=\Delta/(b^2z^2),\quad N'=a^2N/b^2.
\]
The remaining three strict gaps are positive by the original gaps.
This map preserves the total coefficient scale. In particular it
orients an exactly-one-weak source to $C$ weak and $\alpha$ strong.

Write $\mathrm I,\mathrm U,\mathrm O$ for positive infinitesimal,
bounded-unit, and unbounded finite actual angles. The complete cases are
\[
\begin{array}{c|c|c}
 (u_c,v_c)&C,\alpha\text{ strong}&C\text{ weak},\alpha\text{ strong}\\\hline
 \mathrm{II}&\text{II calculation}&\varnothing\\
 \mathrm{IU}&\text{one small angle}&\varnothing\\
 \mathrm{IO}&\text{strong IH}&\text{weak IH}\\
 \mathrm{UI}&\text{strong UI}&\varnothing\\
 \mathrm{UU}&\text{five positive events}&\varnothing\\
 \mathrm{UO}&\text{reflected IU}&\text{reflected IU}\\
 \mathrm{OI}&\text{strong HI}&\varnothing\\
 \mathrm{OU}&\text{reflected UI}&\varnothing\\
 \mathrm{OO}&\text{reflected II}&\text{reflected weak II}.
\end{array}
\]
Each source domain is treated below for all independent testing angles and
states. The following bounds establish that these domains exhaust the source.

If $C,\alpha$ are both weak, the four event-cycle physicality tests
force $n,m,A,D$ unbounded. One test is
$4D\varepsilon^2-2\varepsilon+4(2+\varepsilon)C
+8(1+\varepsilon)\alpha$: a sufficiently small positive real
$\varepsilon$ excludes bounded $D$. The matrix-square core bound gives
$H\ge-2u_cv_c\|\psi_c\|^2$. Thus each matching event amplitude divided
by $\sqrt{u_cv_c}$ is infinitesimal. Compact source angles and the
reflection imply $c$ bounded, $a,b,t\sim1$, and
$r,h,g,j\in\mathrm I$. For $c\in\mathrm I$ the relative matching
equations give $v_c/u_c\sim1$ and $W_c/Z_c\sim1$ before scaled
stationarity is used. The radial calculations give
$N\asymp r+h+g+j$ and $4\alpha/(c^4N)\sim1$; they retain all four
state amplitudes and both state charts. For $c$ a unit,
$\alpha/C\asymp1$. The outer calculations in both cases use
$1/(Cn),1/(Cm),1/(CA),1/(CD)=O(1)$, with their reflected counterparts
when $c$ is a unit. This covers all two-weak sources.

The empty cells follow from the original source inequalities.
For II, $L<y(1+zr)$ and the positive source numerator show that weak
$C$ forces weak $\alpha$. For UI the same numerator and $n\ge0$
would force $r$ infinitesimal and $h$ a unit, contradicting weak $C$.
For IU, bounded $b$ obeys
\[
 r_0\{b_0+k_0^2(b_0+1)\}\le b_0,\quad
 r_0\ge1-b_0,\quad b_0\ge{k_0\over\sqrt{1+k_0^2}}>0,
 \qquad k=v_c,
\]
and gives strong $C$. For unbounded $b$, $C\asymp L^{-1}$ with
$L$ bounded, again giving strong $C$. For UU the positive numerator
and unit source denominator make $C,\alpha$ weak or strong together.
Finally, for OI or OU put $q=u_c^{-2}\in\mathrm I$ and $y=v_c^2$
bounded. The exact $n$ numerator is $-qP/(2rQ)$, $Q>0$, and
\[
 P\ge r^2(1+y)(H+yr)-Tqr(y+2)-q^2(T^2+T),\quad0<T<H+y.
\]
If $q/r$ were infinitesimal, $P>0$ would contradict $n\ge0$.
Thus $r/q$ is bounded, $H\sim1$, and the canceled $C$ numerator
is a unit with bounded positive denominator. These arguments include
equality $n=0$.

For IU, bounded $b$ is a unit by the preceding bound. Its exhaustive
opening cases are $g$ a unit; $g$ infinitesimal and $r$ a unit; and
$g,r$ both infinitesimal. The three one-angle calculations retain,
respectively, the ordinary hard normal, the full $m$ normal, and both
centered $n,m$ squares. If $b$ is unbounded, put
$x=u_c^2$, $B=xb$, $H=xh$, $L=x(b-h)$. The gaps imply
$0<H<1$, $0<L<k^2(1+xr)$ and $B=H+L$ bounded. Exact cancellation gives
\[
 1-r\asymp1,\quad m\asymp b^{-2},\quad C\asymp L^{-1},
 \quad A\asymp H^{-1},\quad D\asymp B/j,\quad
 \alpha\asymp B/(1-H).
\]
For infinitesimal $B$ use the homogeneous one-angle calculation; for
unit $B$ use the strong-event calculation, which includes unbounded
$C$ or $\alpha$. The reflection supplies every UO source.

On strong II the source equations give $B$ a unit, $H,1-r$ units, and
\[
 m/z^2,\quad A/y,\quad CL,\quad Dyj,\quad\alpha(1-H)\asymp1,
 \qquad nr=O(1).
\]
These are the coefficient assumptions of the shared II local and
outer calculations. For IO the weak IH source has $B$ unbounded and
$\rho=B/y$ with residue in $[0,1]$; its interior, two endpoints and
open-axis calculations exhaust that interval. The strong IH source
instead has $B$ a unit, $r=O(1/y)$ and $Dy^2\asymp1$ and uses the
same polynomial relations with its own coefficient bounds. Strong UI
has bounded $nr,m\mu,\alpha\pi,C\rho k^2,D\sigma k^2,Ct-Dj$,
which are the inputs of its three-normal calculation. On OI, direct
source cancellation gives $R=r/q$ bounded, $H\sim1$ and
\[
 mqK,\quad Cyt,\quad Dyj,\quad\alpha qE,\quad A/(qy)\asymp1.
\]
The HI circuit calculation includes every $n\ge0$, including $R_0=0$;
its multiplier $2(1+R_0)^2/(3+R_0)$ lies in $[2/3,2]$. On strong UU,
the source equations force $m,A,D$ strong, so the five-event proof
uses the positive real lower bound of $\min(1,m,C,A,D,\alpha)$.
The remaining reflected entries follow from
\eqref{eq:c2-short-reflection}. The directed weak OO entry is the
following weak-II calculation.

\subsubsection{The directed weak II source}
\label{component:pec6c617d3a3d}
\label{pec6c617d3a3d:sec:weak-c2-complete-join-rederived}
Use $z=u_c^2$, $y=v_c^2$, $H=v_cY_c$,
$L=v_c(u_cW_c-Y_c)$, $B=H+L$, $r=-X_c/u_c$.
The source equations are \eqref{eq:c2-ba-source}, with
$\delta=B\Delta$ and the positive numerator
\[
 M_s=B^2M=(1+z)(B+zr)H+yG+yzj+zrL(1+y),
 \qquad G=B+z(H+r-1).
\]
Since $L<y(1+zr)$, $L$ is
infinitesimal and $B$ bounded. If $B$ is infinitesimal, the exact bound
\[
 0<\alpha\le{(B+zr)(1+z)(H+y)\over2(1-H)}\in\mathrm I
\]
uses $\delta\ge B+z$. If $B$ is a unit, $H=B-L$, $\delta$ and the
positive numerator are units, giving strong $\alpha$. Thus weak
$\alpha$ is exactly $B\in\mathrm I$ in this chart.

Set $a=z/B$, $t=y/B$, $h=H/B$, $l=L/B=1-h$, $e=1-r$.
The exact gaps give
\[
 g=1-ae+hz>0,\quad j=1+zr-l/t>0,\quad ae<1+hz.
\]
Hence unit $e$ forces bounded $a$; unbounded $a$ forces infinitesimal
$e$; infinitesimal $a$ gives $g\sim1$; and infinitesimal $g$ gives
$ae\sim1$. The source formula gives
\[
 r(2n+1)\asymp1,\quad AH/y\asymp1,\qquad
 m={ (1+y)(1+zr)S_c\over2}{z\over\delta}
                       {1+zH/B-g\over g}.
\]
Thus $m$ is bounded for unit $g$, and $m\asymp1/g$ for
infinitesimal $g$. The positive source numerator also gives the
unbounded-coefficient equivalences
\[
\begin{array}{c|ccccc}
 \text{coefficient}&n&m&A&C&D\\\hline
 \text{condition}&r\in\mathrm I&g\in\mathrm I&H/y\in\mathrm I&
 L/(H+y)\in\mathrm I&t\in\mathrm I\text{ or }j\in\mathrm I.
\end{array}
\]
For $t$ unbounded, $j\sim1$ and $D=O(1/t)$.
These estimates concern the original source.

Split $t$ first. Infinitesimal $t$ is R1 and unbounded $t$ is R2,
proved below with their complete scalar and both-normal calculations.
For unit $t$ and bounded $a$, Proposition~\ref{prop:c2-bounded-a}
proves all cases at once.
For unit $t$ and unbounded $a$, necessarily $e$ is infinitesimal,
$r\sim1$. The remaining exhaustive cases are
\[
\begin{array}{c|l}
 h\text{ unit},\ l\in\mathrm I&\text{R3, the EN calculation}\\
 h,l\text{ units},\ j\in\mathrm I&\text{R4, the EN calculation}\\
 h,l,j\text{ units},\ g\in\mathrm I&\text{R6, RI or RB}\\
 h,g\in\mathrm I&\text{R7, RU}\\
 h,j\in\mathrm I,\ g\text{ unit}&\text{R8, the joint-D calculation}\\
 h\in\mathrm I,\ g,j\text{ units}&\text{R10, the weighted/thin calculation}.
\end{array}
\]
Indeed if $h,l,j,g$ are all units, all original coefficients are
bounded by the displayed identities, contradicting the assumed total
scale. Infinitesimal $r$ is impossible when $a$ is unbounded. No
other ordered alternative remains.

The R6 source inequality forces $b=a^2B$ bounded, and its two cases
$b$ unit or infinitesimal are RB and RI. RU treats both $b$ cases for
R7. In R8 and R10, split $n/e^2$ into noninfinitesimal, positive
infinitesimal, and zero. The first uses the weighted source Schur
bound; the second uses the explicitly paired positive coefficients and
the full contact-normal correction; equality uses the exact OP4
native C1 map. OP4 preserves the physical parent and the actual
contact, so only the native C1 small-$e$ or centered-$A$ theorem is
used at $n=0$. The positive-thin calculation never declares its
formal negative-$\beta$ tuple physical. These source splits and the
finite and homogeneous outer calculations below exhaust every
independent testing pair and state.

\subsection{Stationary coefficients and contact}\label{component:p6e5d41ed80b3}\label{app:c2-source}
In this appendix, $c$ denotes the actual half-angle and
$C$ the event weight previously denoted $c$. Set
\[
 z=c^2,\quad A_c=1+zr,\quad
 h=b-a^2t,\quad g=(1+z)b-(1-r+a^2zt),\quad j=1+zr-t,
\]
\begin{align*}
\Delta={}&(1+z)zb^2+2z(1+z)br+(1+z)(1+a^2z)b\\
 &+z^2(1-a^2)r^2+2z(1+a^2z)r+1+a^2z,\\
N={}&(1+z)(b+r)(b+a^2(1-t))-a^2(1+a^2z)(1-r)t.
\end{align*}
The unchanged $d=1$, $\lambda=\beta=0$ source is
\begin{equation}\begin{aligned}
n&=\frac{(b+r)A_c(A_c+a^2zt)-r\Delta}{2r\Delta},&
m&=\frac{(1+a^2z)(1-r)A_c(A_c+a^2zt)}{2g\Delta},\\
C&=\frac{A_cN}{2a^2t\Delta},&
A&=\frac{a^2A_c(b+r)(A_c+a^2zt)}{2h\Delta},\\
D&=\frac{(1+z)(b+r)(b-a^2zr)(A_c+a^2zt)}{2a^2j\Delta},&
\alpha&=\frac{z^2(b+r)N}{2(1-zh)\Delta}.
\end{aligned}\label{eq:c2-native-source}\end{equation}
Its exact contact is
\begin{equation}
 (u_c,v_c;W_c,X_c,Y_c,Z_c)
   =\left(c,ac;\frac ba,-cr,c\left(\frac ba-at\right),1\right).
\label{eq:c2-native-contact}\end{equation}
The positive chart requires $a,c,b,t>0$, $0<r<1$, $h,g,j>0$,
$1-zh>0$, $b-a^2zr>0$, and nonnegative original event weights.
The source-entry proof uses this positive chart without a shape-unit
assumption.

The numerator has the useful positive identity
\[
N=(1+c^2)(b+r)h+a^2\{g+j+rt(1+a^2c^2)\}>0.
\]
It follows by substituting the three source-gap definitions; every term
is positive in the domain used by the contact-entry theorem.
For the determinant calculation use the aliases
\[
 R_2=R_m,\quad R_3=uW-Y,\quad R_5=R_D,\quad R_\alpha=vY-Z,
 \qquad(e_1,e_2,e_3,e_4,e_5)=(n,m,C,A,D).
\]

\subsubsection{Uniqueness of the stationary coefficients}\label{component:p61209c7390b0}
\label{sec:c2-contact-determinant}

The six stationary coefficients are unique, as follows.
Let $M$ have rows indexed by $(W,X,Y,Z,u,v)$ and columns equal to the
derivatives of
\[
D_uD_vX^2,\quad R_2^2,\quad D_vR_3^2,\quad
D_uD_vY^2,\quad R_5^2,\quad D_uR_\alpha^2,
\]
evaluated at $(u,v;W,X,Y,Z)=(c,ac;b/a,-cr,c(b/a-at),1)$.
Write $z=c^2$, $D=1+z$, $E=1+a^2z$, $h=b-a^2t$,
$g=Db-(1-r+a^2zt)$, and $j=1+zr-t$.
Factoring the six columns by
\[
2cr,\quad2cg,\quad2act,\quad-2ch/a,\quad2acj,
\quad2(1-zh)
\]
and the first row by $c$ leaves the matrix
\[
B=\begin{pmatrix}
0&a&E&0&-1&0\\
-DE&-1&0&0&-az&0\\
0&az&-E&-DE&1&-acD\\
0&-c&0&0&ac&D\\
zrE&zh-1&(b/a)E&-zhE/a&azr-b/a&c(1-zh)\\
azrD&b/a+zh/a&a^2zt&-zDh&1+zr&-cDh/a
\end{pmatrix}.
\]
Use the pivots in row $2$, column $1$; row $1$, column $3$; and
row $4$, column $6$.  Eliminate those three columns, multiply the
remaining second and third rows by $D$ and $E$, and factor $-DE$
from their middle column.  The determinant reduces to
$\det B=D^2E^2\det T$, where
\[
T=\begin{pmatrix}
a&1&a^2z\\
-1-b-z(r+a^2t)&zh/a&az(r-1+zh)\\
b/a+az(h-r)&zh&E(1+zh)+z(r+a^2t)
\end{pmatrix}.
\]
Subtract $(zh/a)$ times the first row from the second and $zh$ times
the first row from the third.  Expanding along the middle column gives
\[
\det T=[1+b+z(b+r)][1+z(a^2+b+r)]
-z(1-r)(b-a^2zr)=\Delta.
\]
The last equality follows by expanding the two products and is exactly
the source definition of $\Delta$.  Restoring the factors proves
\begin{equation}
\det M=-64a c^6rgt h j(1-zh)D^2E^2\Delta.
\label{eq:c2-entry-determinant}
\end{equation}
Each source coefficient occurs with a factor $4$ in the parent, so the
stationary linear system in $(e_1,e_2,e_3,e_4,e_5,\alpha)$ has determinant
$4^6\det M$.  On the positive source domain this is nonzero.  Since the
six source formulas above satisfy all six stationary equations,
they are the unique stationary coefficients.

\section{Two weak complementary coefficients}
\subsection{Entry from two weak coefficients}
\label{component:p534f57e1da6f}

\label{p534f57e1da6f:sec:chart-free-c2}

Here weak means infinitesimal after division by the original positive
signed coefficient \(d\). All source calculations take place in one real
closed coefficient field. Testing fields remain independent.

\begin{theorem}[C2 entry without a contact-chart hypothesis]
\label{p534f57e1da6f:thm:chart-free-c2}
Let the physical canonical parent have \(\lambda=\beta=0\), \(d>0\),
and nonnegative original event coefficients. Suppose it has an actual
finite noncommuting stationary zero contact, and
\[
                    C/d,\ \alpha/d\quad\text{are infinitesimal}.
\]
Then its full parent is tensor SOS at every independent real or complex
projective test. No hypothesis on the other four event weights, on the
contact-angle ratio, or on the signs of the source gaps is imposed.

More precisely, normalize \(d=1\). After local real sign gauges and at
most the exact C2 reflection,
the actual contact enters the positive rational C2 source with
\[
 c\text{ bounded},\qquad a,b,t\longrightarrow1,\qquad
 r,h,g,j>0\text{ infinitesimal}.
\]
In particular all four matching weights are unbounded and both
complementary weights are strictly positive on this finite-contact domain.
\end{theorem}

The proof has three ingredients: a finite physical witness makes the
matching weights large; the quadratic CHSH excess controls the contact
amplitudes at the scale of the actual angles; stationarity fixes their
signs. This last step establishes the source chart rather than assuming it.

\subsubsection*{The physical bound at the correct angular scale}
For \(u,v>0\), put \(D_u=1+u^2,D_v=1+v^2\), and let \(M(u,v)\) be the
real symmetric matrix of the signed core
\[
H=2\{R_D^2+D_u(vY-Z)^2+D_v(X+uZ)^2-D_uD_vZ^2\},
\]
where
\[
R_m=vW-X+uvY-uZ,\qquad R_D=uW+uvX-Y-vZ.
\]
Write \(L=D_uD_v\). Direct matrix multiplication gives
\begin{equation}
(M-LI)^2=L^2I+4uvLJ,\qquad
J=\begin{pmatrix}0&0&0&1\\0&0&-1&0\\0&-1&0&0\\1&0&0&0\end{pmatrix},
\qquad J^2=I.                                      \label{p534f57e1da6f:eq:cf-core-square}
\end{equation}
For a unit real state \(\psi\), Cauchy--Schwarz therefore gives
\[
\psi^TM\psi\ge L-\sqrt{L^2+4uvL}\ge-2uv.
\]
The last inequality follows by squaring the positive quantities
\(L+2uv\) and \(\sqrt{L^2+4uvL}\); their squared difference is
\(4u^2v^2\). Homogeneity proves
\begin{equation}
                         H\ge-2uv\,\|\psi\|^2.        \label{p534f57e1da6f:eq:cf-core-bound}
\end{equation}
This is the quadratic CHSH bound in the present normalization.

At a normalized zero contact, every event term is nonnegative. Define
the four matching amplitudes, in the fixed order \((n,A,m,D)\), by
\[
\mathcal L=(\sqrt{D_uD_v}X,\sqrt{D_uD_v}Y,R_m,R_D).
\]
Equation~\eqref{p534f57e1da6f:eq:cf-core-bound} gives, for each matching event,
\begin{equation}
                   e_i\mathcal L_i^2\le uv/2.         \label{p534f57e1da6f:eq:cf-amplitude-bound}
\end{equation}
All square roots and inequalities in this paragraph are coefficient-side.

\subsubsection*{Matching weights and the compact contact geometry}
First every matching weight is unbounded. To see the mechanism for
\(D\), use the coefficient-field experiment
\[
 u=1,\quad v=\sqrt{1+\delta},\quad
 (W,X,Y,Z)=(1,0,0,\sqrt{1+\delta}),\qquad\delta>0.
\]
The \(n,A,m\) positive square terms vanish, and the full eight-event parent is exactly
\[
4D\delta^2-2\delta+8\lambda+4(2+\delta)C
 +8(1+\delta)\alpha+4(2+\delta)(1+\delta)\beta.
\]
If \(D\le B\) for a positive real \(B\), choose the positive real
\(\delta=1/[4(B+1)]\). The first two terms have a strictly negative
real upper bound, whereas the complementary terms are infinitesimal.
This contradicts physicality. The exact even rotations of the event
cycle give the same argument for \(n,A,m\), preserving the sum of the
four complementary weights.  Hence \(1/n,1/A,1/m,1/D\) are infinitesimal.

Normalize the actual state to norm one and use positive actual
half-angles, obtainable by local real basis gauges. The weaker
bound \(H/(D_uD_v)\ge-2\) shows that all normalized matching
probabilities vanish in residue. In particular \(X,Y\to0\).
Put
\[
c_u=(1+u^2)^{-1/2},\ s_u=uc_u,\qquad
c_v=(1+v^2)^{-1/2},\ s_v=vc_v.
\]
If \(\sigma=\operatorname{res}(c_us_v)\) and
\(\tau=\operatorname{res}(s_uc_v)\), the two remaining matching
equations read
\[
\sigma\overline W-\tau\overline Z=0,\qquad
\tau\overline W-\sigma\overline Z=0.
\]
At least one of \(\overline W,\overline Z\) is nonzero. Thus
\(\sigma=\tau\ge0\). If positive, the two actual half-angles are
bounded units with the same residue. If zero, both are infinitesimal
or both unbounded, by the two unit-circle identities. The exact C2
reflection sends \((u,v)\) to \((1/v,1/u)\) and exchanges
\((C,\alpha)\) and \((m,D)\), preserving the hypotheses. We may
therefore assume either ordinary bounded-unit angles or two
infinitesimal angles.

\subsubsection*{Relative coalescence when both angles vanish}
Suppose \(u,v\) are infinitesimal. From
\eqref{p534f57e1da6f:eq:cf-amplitude-bound}, all four
\(\mathcal L_i/\sqrt{uv}\) are infinitesimal. In particular
\(X/\sqrt{uv},Y/\sqrt{uv}\to0\). Set \(a=v/u>0\).
Dividing the two other matching amplitudes by \(\sqrt{uv}\) gives
\[
\sqrt a\,W-\frac Z{\sqrt a}\longrightarrow0,\qquad
\frac W{\sqrt a}-\sqrt a\,Z\longrightarrow0.
\]
The discarded terms are \(X/\sqrt{uv},Y/\sqrt{uv}\) and
\(\sqrt{uv}X,\sqrt{uv}Y\), all infinitesimal.
If \(a\) were unbounded, these equations would force both \(W,Z\)
infinitesimal; the same follows if \(a\) is infinitesimal.
Both contradict normalization. Thus \(a\) is a bounded unit.
Taking residues now yields \(a\to1\) and \(W/Z\to1\), with
both \(W,Z\) units. Consequently
\[
                  X/u,\ Y/u,\ R_m/u,\ R_D/u\longrightarrow0.
\]

Normalize \(W=1\) by a pure coefficient factor. Write
\[
x=X/u,\quad y=Y/u,\quad q=v/u,\qquad
(u,x,y,Z,q)\longrightarrow(0,0,0,1,1).
\]
The core divided by \(u^2\) is the polynomial
\begin{align}
\mathcal H:=H/u^2=2\{&
(1-y-qZ+u^2qx)^2+(1+u^2q^2)(x+Z)^2 \notag\\
&+(1+u^2)(u^2q^2y^2-2qyZ-q^2Z^2)\}. \label{p534f57e1da6f:eq:cf-scaled-core}
\end{align}
It has zero residue at this point. The nonnegative alpha positive square term and
the zero-contact equation give
\[
0\le4\frac{\alpha}{u^2}(1+u^2)(u^2qy-Z)^2
                   \le-\mathcal H.
\]
Its square factor is a positive unit. Hence
\begin{equation}
                         \alpha/u^2\longrightarrow0. \label{p534f57e1da6f:eq:cf-alpha-scale}
\end{equation}
This estimate is derived before the scaled stationary equations are used.

\subsubsection*{Stationarity fixes the source signs}
In the small-angle case introduce the four exact local coordinates
\[
\ell=\left(\sqrt{D_uD_v}\,x,\sqrt{D_uD_v}\,y,
 q-x+u^2qy-Z,\ 1+u^2qx-y-qZ\right).
\]
Their Jacobian with respect to \((x,y,Z,q)\) has residue
\[
J_0=\begin{pmatrix}
1&0&0&0\\0&1&0&0\\-1&0&-1&1\\0&-1&-1&-1
\end{pmatrix},\qquad \det J_0=2.
\]
Thus its inverse and the transformed derivatives are bounded at the
actual contact. Only this pointwise coefficient-field inverse is used.
Direct differentiation of \eqref{p534f57e1da6f:eq:cf-scaled-core} gives
\[
\operatorname{res}(\nabla_\ell\mathcal H)=2(1,-1,-1,1).
\]
The full parent divided by \(u^2\) is exactly
\[
4\sum_i e_i\ell_i^2
 +4C(1+u^2q^2)(1-y)^2
 +4(\alpha/u^2)(1+u^2)(u^2qy-Z)^2+\mathcal H.
\]
The two complementary-gradient terms have zero residue, by weak \(C\)
and \eqref{p534f57e1da6f:eq:cf-alpha-scale}. The original stationary equations imply
stationarity in \(\ell\), with \(u\) fixed. These equations first show
that every \(e_i\ell_i\) is bounded, so its residue is defined. Consequently
\begin{equation}
\operatorname{res}(n\ell_n,A\ell_A,m\ell_m,D\ell_D)
                         =\tfrac14(-1,1,1,-1).        \label{p534f57e1da6f:eq:cf-sign-residue}
\end{equation}
Every matching weight is positive, so \(X<0,Y>0,R_m>0,R_D<0\)
as exact source inequalities.

For ordinary angles the argument is shorter. Normalize \(W=1\);
the actual contact has residue
\((X,Y,Z,v)=(0,0,1,u)\), where \(u\) is a bounded unit.
Use the unscaled \(\mathcal L\) coordinates. Their Jacobian determinant
at that residue is \(2u(1+u^2)^2\), and
\[
\nabla_{\mathcal L}H=2u(1,-1,-1,1)
\]
there. All complementary gradients are bounded times \(C,\alpha\),
hence infinitesimal. Stationarity gives
\(\operatorname{res}(e_i\mathcal L_i)
 =\overline u\,(-1,1,1,-1)/4\), with the same exact signs.

\subsubsection*{The original source and all independent tests}
In either case \(W/Z\to1\), \(v/u\to1\), and \(X/u,Y/u\to0\).
Set, using the original actual contact,
\[
c=u,\quad a=v/u,\quad b=aW/Z,\quad
r=-X/(uZ),\quad t=(uW-Y)/(vZ).
\]
Then \(a,b,t\to1\) and
\[
h=b-a^2t=\frac{aY}{uZ}>0,\qquad
g=\frac{R_m}{uZ}>0,\qquad j=-\frac{R_D}{vZ}>0.
\]
All four \(r,h,g,j\) are infinitesimal, so \(0<r<1\).
The remaining chart boundaries are positive units:
\[
1-c^2h=\frac{Z-vY}{Z}>0,\qquad
b-a^2c^2r=\frac{a(W+vX)}Z>0.
\]
These definitions give exactly
\[
h=b-a^2t,\qquad
g=(1+c^2)b-(1-r+a^2c^2t),\qquad j=1+c^2r-t.
\]

Finally the stationary coefficient matrix has the nonzero determinant
\[
4^6[-64ac^6rgt\,hj(1-c^2h)(1+c^2)^2(1+a^2c^2)^2\Delta].
\]
The preceding three-by-three reduction gives this determinant.
For \(z=c^2\), its positive expansion is
\[
\Delta=1+(1+z)b+z(1+z)b^2+2z(1+z)br+z^2r^2+2zr
       +a^2z\{(1+z)b+1+zr(2-r)\}\ge1.
\]
The original coefficients
therefore equal the unique positive rational C2 source formulas.
In particular their positive numerator identity forces \(C,\alpha>0\).
If \(c\) is infinitesimal, apply the small-angle radial theorem and its
outer-test corollary; if \(c\) is a unit, apply the ordinary-angle radial
corollary. Reflect back when necessary and restore \(d\) by a pure-source
square root. This proves Theorem~\ref{p534f57e1da6f:thm:chart-free-c2}.

This covers the finite-contact two-weak branch. Weak \(C\) alone need
not imply weak \(\alpha\); the one-weak family below supplies the
complementary case.

\subsection{Matching the weak coefficients to the strong events}
\label{component:p8b53b68d8af6}

\label{p8b53b68d8af6:sec:c2-matching-strength}

A coefficient condition symmetric in the two weak positive square terms closes both
commuting angle corners at once.  It also gives a second exact radial
source family whose contact angles have ordinary positive residue.

\begin{theorem}[Two matched weak positive square terms]
\label{p8b53b68d8af6:thm:c2-matching-strength}
Let the original physical C2 parent have $d=1$, $\lambda=\beta=0$,
and event weights $(n,m,C,A,D,\alpha)$.  Suppose $C,\alpha>0$ are
infinitesimal, all four other weights are positive, and each of
\[
 Cn,\ Cm,\ CA,\ CD,\qquad
 \alpha n,\ \alpha m,\ \alpha A,\ \alpha D
\]
has bounded reciprocal in the coefficient field.  Then the full parent
is tensor SOS at every independent projective test and every real or
complex state.
\end{theorem}
\begin{proof}
Lemma~\ref{pc7dacaa3258e:lem:c2-four-strong-outer} applies directly to every commuting
residue test except the finite zero-zero corner.
For that corner use the exact reflection
\begin{equation}
 (u,v;W,X,Y,Z;n,m,C,A,D,\alpha)
 \longmapsto
 (1/v,1/u;Z,X,Y,W;n,D,\alpha,A,m,C).
\label{p8b53b68d8af6:eq:c2ms-reflection}
\end{equation}
More explicitly, if $F$ denotes the full parent with its indicated
arguments, then
\begin{equation}
 F(u,v;W,X,Y,Z;n,m,C,A,D,\alpha)
 =u^2v^2F(1/v,1/u;Z,X,Y,W;n,D,\alpha,A,m,C).
\label{p8b53b68d8af6:eq:c2ms-reflection-identity}
\end{equation}
This is a homogeneous polynomial identity, so its interpretation in
reciprocal charts includes literal zero and infinity endpoints.

For a short verification, the reflected $R_2,R_5$ amplitudes, after
multiplication by $uv$, are $-R_5,-R_2$.  The reflected $C$ and
$\alpha$ positive square terms exchange, while the $n,A$ positive square terms are fixed.
If $\mathcal H$ is the signed core, its common expanded form is
\begin{align*}
 \mathcal H/2={}&u^2(W^2+Z^2)+(1+v^2+u^2v^2)(X^2+Y^2)\\
 &+2u^2v(WX-YZ)+2u(XZ-WY)-2uv(WZ+XY).
\end{align*}
Multiplying its reflected expression by $u^2v^2$ fixes the first,
second and last terms, and exchanges the middle two.  This proves the
whole-parent identity, including its signed part.

The reflected weak weight is $\alpha$, its four strong weights are
$n,D,A,m$, and its other nonnegative weak positive square term is $C$.
Thus all hypotheses of Lemma~\ref{pc7dacaa3258e:lem:c2-four-strong-outer} hold for
the reflected parent.  Its infinity-infinity chart is precisely the
original zero-zero chart.  Restore the homogeneous pure-testing
factors to obtain the certificate there as well.

For noncommuting residue tests apply the physical-source theorem, using
positive pure-source normalization when a coefficient is unbounded.
Restore that factor and split real and imaginary state parts.
\end{proof}

\begin{corollary}[Radial C2 source with ordinary contact angles]
\label{p8b53b68d8af6:cor:c2-ordinary-angle-radial}
Use the same rational source \eqref{eq:c2-native-source}, now with source
half-angle $c$ a positive bounded unit.  Suppose $h=b-a^2t,r,g,j$ are
positive infinitesimals and $a,1-r,t$ are positive bounded units.
Then the full parent is tensor SOS at every independent projective test
and every real or complex state.
\end{corollary}
\begin{proof}
Put $z=c^2$ and continue to write $C=e_3$ for the weak event weight.
The exact source identities
\[
 b-1=-r+g-zh,\qquad t-1=zr-j,\qquad
 a-1=\frac{g+j-(1+z)(r+h)}{t(a+1)}
\]
remain valid for this ordinary source angle.  Thus $a,b,t$ have residue
one.  The positive source identity
\[
 N=(1+z)(b+r)h+a^2\{g+j+rt(1+a^2z)\}
\]
shows that $N>0$ is infinitesimal and $h/N,r/N,g/N,j/N$ are bounded:
each of the four positive coefficients in this sum is a bounded unit.

Here is the entire normalized coefficient calculation.  Write
$z_0=\overline z>0$.  Substitution of $a=b=t=1,r=0$ in the bounded
source numerators and unit denominators gives
\begin{gather}
 \overline\Delta=2(1+z_0)^2,\qquad
 \overline{C/N}=\frac1{4(1+z_0)^2},\qquad
 \overline{\alpha/C}=z_0^2,\notag\\
 \overline{rn}=\overline{hA}=\frac1{4(1+z_0)},\qquad
 \overline{gm}=\overline{jD}=\frac14.
\label{p8b53b68d8af6:eq:c2ms-ordinary-residues}
\end{gather}
For example, the exact quotients used in the first line are
\[
 \frac CN=\frac{1+zr}{2a^2t\Delta},\qquad
 \frac\alpha C=\frac{z^2(b+r)a^2t}{(1-zh)(1+zr)}.
\]
The four strong normalized products use the numerators of
\eqref{eq:c2-native-source} after multiplication by $r,h,g,j$, respectively;
no vanishing source gap is inverted in this calculation.
All six weights are therefore positive, $C,\alpha$ are infinitesimal,
and their normalized quotients displayed above are bounded units.
For instance
\[
 \frac1{Cn}=\frac{r/N}{(C/N)(rn)},\qquad
 \frac1{\alpha n}=\frac{C/\alpha}{Cn};
\]
the other six product bounds follow in the same way from $g/N,h/N,j/N$.
These are exactly the hypotheses of Theorem~\ref{p8b53b68d8af6:thm:c2-matching-strength}.

For completeness, the source is physical by the unchanged exact contact
argument.  The rational stationary identities
\eqref{eq:c2-native-contact} are valid at this value of $c$ as well.
At that noncommuting contact the state matrix has a nonzero kernel
vector.  Its restriction to $Z=0$ is positive definite: the signed core
is nonnegative there and the positive event positive square terms have zero common
kernel.  Hence the other three eigenvalues are positive.  The exact
contact and angular stationarity give the stationary zero of the
Laurent determinant; its constant-plus-positive-exponentials form makes
this its global minimum.  The determinant criterion and homogeneous
continuity prove physicality, exactly as for the small-angle source.
Theorem~\ref{p8b53b68d8af6:thm:c2-matching-strength} now applies.

The corollary applies to the stated radial subfamily of the ordinary-angle
cell.
\end{proof}

\subsection{Weighted radial completion}
\label{component:p9f1235511111}
\label{component:p53752031d16c}
\label{component:p163e0ebf59f6}
\label{p9f1235511111:sec:c2-unified}
\label{p53752031d16c:sec:c2-y-current}
\label{p163e0ebf59f6:sec:c2-radial-corner}
Normalize the original positive $d$ by a pure-source square. The
two-weak source reduction is applied after this normalization; it has
already given $c\in\mathrm I$,
$a,b,t\sim1$ and $r,h,g,j\in\mathrm I$, where $h=b-a^2t$.
Use the exact native coefficients in Section~\ref{app:c2-source}, with
$(e_1,e_2,e_3,e_4,e_5)=(n,m,C,A,D)$ and
\[
 (u_c,v_c;W_c,X_c,Y_c,Z_c)
 =(c,ac;b/a,-cr,c(b/a-at),1).
\]
The two source identities needed for all relative opening rates are
\begin{equation}
 N=(1+c^2)(b+r)h+a^2\{g+j+rt(1+a^2c^2)\}>0,
 \qquad N\asymp h+r+g+j.
 \label{p53752031d16c:eq:c2y-N-positive}
\end{equation}
Consequently
\[
 U_1=re_1,\quad U_2=ge_2,\quad U_4=he_4,\quad U_5=je_5
 \sim\tfrac14,\qquad e_3/N\sim\tfrac14,
 \qquad {4\alpha\over c^4N}\sim1.
\]
These follow directly from the native fractions after their respective
opening factors are canceled. In particular all four $e_i$ above are
unbounded, and $h/N,r/N,g/N,j/N$ are bounded. They also imply
$1/(Cn),1/(Cm),1/(CA),1/(CD)=O(1)$, the assumptions of the complete
radial outer calculation.

The same source reduction forces $n,m,A,D$ unbounded, so $n=0$ is
impossible in this case. No zero-coefficient face is being discarded.
\begin{proposition}
\label{prop:c2-weighted-radial}
\label{p163e0ebf59f6:thm:c2-radial-corner}
On this source domain the full parent is tensor SOS at every independent
projective testing pair and every real or complex state.
\end{proposition}
First let the independent tests $u,v$ be infinitesimal or zero and take
the pure-testing state chart $Z=1$, $W\ge0$ bounded. Define
\[
 R_m=vW-X+uvY-u,\qquad J=uW+uvX-Y-v,
\]
\[
 H_1=X+ru,\quad H_2=R_m-gu,\quad
 H_4=Y-(h/a)u,\quad H_5=J+jv.
\]
Retain the four whole positive squares
\[
 4e_1D_uD_vH_1^2+4e_2H_2^2+4e_4D_uD_vH_4^2+4e_5H_5^2.
\]
Their differences from the original terms are, respectively,
\[
 -8U_1D_uD_vuH_1+4rU_1D_uD_vu^2,
 \quad8U_2uH_2+4gU_2u^2,
\]
\[
 {8U_4\over a}D_uD_vuH_4+{4hU_4\over a^2}D_uD_vu^2,
 \quad-8U_5vH_5+4jU_5v^2.
\]
Thus the entire remaining polynomial has bounded coefficients,
including all its state normals. With $n_1=H_1$, $n_5=H_5$ the exact
inverse is
\[
 X=-ru+n_1,\qquad
 Y=uW-(1-j+ru^2)v+uvn_1-n_5.
\]
Define its scalar $f$, linear coefficients $L_1,L_5$, and quadratic
normal matrix by coefficient extraction in $(n_1,n_5)$. Quadratic
state degree proves that this is the whole remainder.

Put
\[
 R=au-v,\qquad S=u(W-W_c),\qquad T=u^2-c^2.
\]
Then $f\in(R,S,T)^2$ and $L_1,L_5\in(R,S,T)$ with bounded
coefficients. Here is the finite division proof. Parity makes
$f(u,au,W)=f_0(u^2,W)$ even and $f_v(u,au,W),L_i(u,au,W)$ odd.
At zero angles the value is $4\alpha$, independent of $W$, so both
nonconstant $W$ coefficients of $f_0$ are divisible by $u^2$.
Write $f_0=A(\xi)(W-W_c)^2+B(\xi)(W-W_c)+C(\xi)$.
The exact source contact gives
\[
 A=\xi A_1,\quad B=\xi(\xi-c^2)B_1,
 \quad C=(\xi-c^2)^2C_1.
\]
Monic division gives these polynomials without a factor $1/c$.
Likewise $f_v(u,au,W)=S\ell_1+uT\ell_2$. Finite Taylor division in
$v-au$ proves the scalar assertion. The identical odd-parity argument
proves the linear-normal assertions. All higher normal terms are
already quadratic by construction.

The exact current identity removes $S$. Put $K=1-j+ru^2$ and
\[
 Q(\xi)=aW_c-a^2(1-j)-a^2r(\xi+c^2).
\]
The identities $H_2=G-(1-u^2v^2)H_1-uvH_5$ and
$G=-[D_uW-Ku(v+au)]R+aD_uS+uTQ(u^2)$ give
\begin{equation}
 S={H_2+(1-u^2v^2)H_1+uvH_5
       +[D_uW-Ku(v+au)]R-uTQ(u^2)\over aD_u}.
 \label{p9f1235511111:eq:c2u-eliminate-S}
\end{equation}
Its denominator is a product of pure-source and pure-testing units.
After restoring the four retained squares this is an exact bounded
Gram on $K_i=\sqrt{e_i}H_i$ for $i=1,2,4,5$, together with $R,T$.

There is one further exact linear relation:
\begin{equation}
 \begin{split}
 \zeta={}&aD_uH_4-H_2-(1-u^2v^2+aD_uuv)H_1
 +(aD_u-uv)H_5\\
 &-[D_uW+(a-uv)K]R+u[Q(u^2)+a^2D_ur]T=0.
 \end{split}
 \label{p53752031d16c:eq:c2y-zero-relation}
\end{equation}
Add the exact zero
\begin{equation}
 {8U_4\over a^3}\zeta R+\zeta^2=0.
 \label{p53752031d16c:eq:c2y-zero-correction}
\end{equation}
The hard and $R$ residue can be computed directly on the present
all-small domain. Modulo $(h,r,g,j)$, followed by the joint real
residue $c=u=v=0$, the degree-two scalar is
\[
 2(u-v)(uW-v)=2R(R+S).
\]
Equation~\eqref{p9f1235511111:eq:c2u-eliminate-S} gives $S=WR$ on
that residue, while $\zeta=-(1+W)R$ and $U_4=1/4$.
The displayed zero changes $2(1+W)R^2$ into $(1+W)^2R^2$.
The four scaled hard diagonals are $4$ and their other couplings
vanish.

It remains to divide the radial column by $\sqrt N$. Let
$I=(h,r,g,j)$ in the bounded source coefficient ring. The exact
identities
\[
 b-1=-r+g-c^2h,\quad t-1=c^2r-j,\quad
 a-1={g+j-(1+c^2)(r+h)\over t(a+1)}
\]
show that subtraction of the $I=0$ specialization of any coefficient
leaves $I$ with bounded coquotients. Its denominators are units after
the four opening factors have been absorbed in $U_i$. Finite monic
division and the displayed current substitution preserve this fact.

The exact simplifications $Q(u^2)=h-a^2ru^2$ and
$Q(u^2)+a^2D_ur=h+a^2r$ belong to $I$. The only radial cross
coefficient requiring a further correction is the $RT$ coefficient.
For an explicit calculation before taking the testing residue, set
$D_c=1+c^2$, $Y_*=uW-v$, $G_*=D_uvW-D_vu$. At $I=0$ the whole
scalar remainder is
\[
 f_*=2uG_*+{2\over D_c}D_uD_vuY_*
       +2\{D_u(vY_*-1)^2+D_vu^2-D_uD_v\}.
\]
It satisfies
\[
 f_*(u,u,1)=0,\quad
 (f_*)_W(u,u,1)={2u^2(u^2-c^2)D_u\over D_c},\quad
 (f_*)_v(u,u,1)={-2u(u^2-c^2)D_u\over D_c}.
\]
The monic division order just specified therefore gives, after
eliminating $S$, the $RT$ coefficient
$2u(D_uW+1-uv)/D_c$. The earlier zero adds none at this
specialization. The $R$ coefficient of $\zeta$ is
$-(D_uW+1-uv)$, so add the further exact zero
\[
                         (2u/D_c)\zeta T=0.
\]
The resulting $RT$ coefficient belongs to $uI$.

Now use $T_N=\sqrt N T$. Every hard--$T_N$ coefficient is $u$
times bounded factors divided by $\sqrt{Ne_i}$, which is bounded
because $1/\sqrt{Ne_i}=\sqrt{(r,g,h,j)_i/N}/\sqrt{U_i}$.
Its residue is zero. The corrected $RT_N$ coefficient belongs to
$uI/\sqrt N$, again bounded with zero residue. The $T^2$ coefficient
belongs to $I$ because $f_*(u,u,1)=0$; every contribution other than
its value at $u=0$ has a further factor $u^2$. Its value at $u=0$ is
exactly $4\alpha/c^4$. Therefore its quotient by $N$ is bounded and
has residue $1$. The complete weighted residue is
\[
 4K_1^2+4K_2^2+4K_4^2+4K_5^2+(1+W_0)^2R^2+T_N^2\succ0.
\]
The separated strict-Gram lemma proves the full chart.

For $W=1$, $Z$ infinitesimal, retain the original four squares in
$X,R_m,Y,J$ and the whole $\alpha$ square. The exact testing identities
\[
 v={R_m+(1-u^2v^2)X+uD_vZ+uvJ\over D_u},
\]
\[
 u={D_uY+[Z(1-u^2v^2)-D_uuv]X+(D_u+uvZ)J+ZR_m
       \over D_u-D_vZ^2}
\]
have pure-testing unit denominators. They express the entire remaining
quadratic form on the four scaled retained amplitudes. To check this
without discarding a state term, the half-core has the exact completion
\[
 H=2\{C_*^2+D_vQ_*^2+D_uv^2Y^2
       +D_uu^2v^2(1-u^2v^2)Z^2/D_v\}-4D_uuvZ,
\]
where $C_*=u+uv(X+uZ)-Y$ and
$Q_*=X+uZ-D_uuv^2Z/D_v$. Substitute the two identities in every
term. All coefficients are bounded and the residue on the four scaled
amplitudes is $4I$; thus the same lemma applies.

The testing sign gauges allow $u,v\ge0$. The $X,Y$ block of half
the core has eigenvalues $L\pm uv>0$, $L=1+v^2+u^2v^2$, and Schur
complement
\[
 {uvD_uD_v\{uv(W^2+Z^2)-2LWZ\}\over L^2-u^2v^2}.
\]
For $WZ\le0$ the core is a nonnegative pure-testing value and all
events are already tensor squares. For $WZ>0$, the two preceding
state charts are exhaustive. Homogeneous restoration and real/imaginary
splitting prove every state. Section~\ref{component:pc7dacaa3258e}
then proves all other independent projective testing pairs.

\subsection{Complementary testing angles}
\label{component:pc7dacaa3258e}

\label{pc7dacaa3258e:sec:c2-radial-all-tests}

The radial corner has four strong event positive square terms and one weak positive square term.
The weak positive square term must be retained at its exact scale at the opposite
commuting corner.  An elementary conductor identity makes that scale
sufficient; it does not compare coefficient and testing infinitesimals.
Throughout this section $C$ denotes the event weight $e_3$, whereas $c$
continues to denote the source half-angle in \eqref{eq:c2-native-source}.

\begin{lemma}[Four strong positive square terms and a weak opposite-corner pivot]
\label{pc7dacaa3258e:lem:c2-four-strong-outer}
Consider the full parent \eqref{eq:c2-master-parent}, with
\[
 (e_1,e_2,e_3,e_4,e_5)=(n,m,C,A,D),\qquad \alpha\geq0.
\]
Suppose $C>0$ is infinitesimal and the four positive pure coefficient
quantities $Cn,Cm,CA,CD$ have bounded reciprocals.  Then the full parent
is tensor SOS at every commuting residue test outside the finite
zero-zero corner, for every real or complex projective state.
The four strong coefficients and the roots in the proof may be unbounded.
\end{lemma}

\begin{proof}
The hypotheses imply that $1/n,1/m,1/A,1/D$ are infinitesimal.
First consider the outer axes other than the infinity-infinity corner.
Choose $\mu=\min(1,n,m,A,D)$, a positive bounded unit.  Retain the
nonnegative excess of each strong positive square term, together with the whole
$C$ and $\alpha$ positive square terms.  The remaining parent consists of the signed
core and the four positive square terms with common weight $\mu$.

Here are the four kernel calculations, with positive factors
$1+u^2$ or $1+v^2$ omitted only from the positive square term amplitudes.
At $u=0$, finite nonzero $v$, the four amplitudes are
$(X,vW-X,Y,-Y-vZ)$; at $v=0$, finite nonzero $u$, they are
$(X,-X-uZ,Y,uW-Y)$.  Their common kernels are zero, and half the
signed core is, respectively,
\[
 D_v(X^2+Y^2),\qquad (uW-Y)^2+(X+uZ)^2.
\]
At $u=\infty$, finite $v$, the four homogenized amplitudes are
$(X,vY-Z,Y,W+vX)$; at $v=\infty$, finite $u$, they are
$(X,W+uY,Y,uX-Z)$.  These common kernels are zero even when the
finite angle is zero.  Half the homogenized signed core is, respectively,
\[
 (W+vX)^2+(vY-Z)^2,\qquad D_u(X^2+Y^2).
\]
Thus the selected state Gram has strict positive residue.  Its
coefficients are bounded in each indicated local testing chart, so the
strict-residue Gram lemma proves the whole chart.  This includes the
mixed zero-infinity endpoints and either sign of every finite angle.

It remains to treat $u=v=\infty$ at residue.  Put $p=1/u$, $q=1/v$,
using the homogeneous reciprocal coordinates at exact infinity, and
write $D_p=1+p^2$, $D_q=1+q^2$.  Thus $p,q$ are infinitesimal or zero.
The homogeneous full parent in this chart is exactly
\begin{align}
\widehat F={}&4nD_pD_qX^2+4mM^2+4CD_q(W-pY)^2
       +4AD_pD_qY^2+4DJ^2+4\alpha D_p(Y-qZ)^2+H,\notag\\
 M={}&pW-pqX+Y-qZ,\qquad
 J=qW+X-pqY-pZ,\notag\\
 H={}&2\{J^2+D_p(Y-qZ)^2+D_q(pX+Z)^2-D_pD_qZ^2\}.
\label{pc7dacaa3258e:eq:c2ro-parent}
\end{align}
For finite nonzero $u,v$ this is $p^2q^2F$; homogeneous interpretation
also gives the exact endpoint polynomial.

Set $L=1+p^2+p^2q^2$.  The needed conductor identity is
\begin{equation}
\frac H2=qWJ-qZM+L(X^2+Y^2)-2pqXY
             +q(W+pqZ)X-q(pqW+Z)Y.
\label{pc7dacaa3258e:eq:c2ro-core-conductor}
\end{equation}
For a direct check, put $M_0=pW-qZ$, $J_0=qW-pZ$.  The scalar part
of $H/2$ is
\[
 J_0^2+(q^2-p^2)Z^2=qWJ_0-qZM_0.
\]
Its linear $X,Y$ part is
$2q(W+pqZ)X-2q(pqW+Z)Y$, and its quadratic part is
$L(X^2+Y^2)-2pqXY$.  Substitution of
$M=M_0-pqX+Y$, $J=J_0+X-pqY$ gives
\eqref{pc7dacaa3258e:eq:c2ro-core-conductor}.

First normalize $W=1$ and assume $Z$ bounded.  No bound is required
on $X$ or $Y$.  Retain the $\alpha$ positive square term and use the five roots
\[
 K_1=\sqrt nX,\quad K_2=\sqrt mM,\quad
 K_4=\sqrt AY,\quad K_5=\sqrt DJ,\quad T_C=\sqrt C\,W.
\]
Here $T_C=\sqrt C$ in this state chart.  Equations
\eqref{pc7dacaa3258e:eq:c2ro-parent}--\eqref{pc7dacaa3258e:eq:c2ro-core-conductor} give the following
\emph{exact} Gram expression for the remainder:
\begin{align}
\widehat F-4\alpha D_p(Y-qZ)^2={}&
 4D_pD_q(K_1^2+K_4^2)+4K_2^2+4K_5^2
 +4D_q\left(T_C-p\sqrt{C/A}\,K_4\right)^2\notag\\
&+2L\left(\frac{K_1^2}{n}+\frac{K_4^2}{A}\right)
 -\frac{4pq}{\sqrt{nA}}K_1K_4\notag\\
&+2qT_C\left\{
 \frac{K_5}{\sqrt{CD}}-\frac{ZK_2}{\sqrt{Cm}}
 +\frac{(1+pqZ)K_1}{\sqrt{Cn}}
 -\frac{(pq+Z)K_4}{\sqrt{CA}}\right\}.
\label{pc7dacaa3258e:eq:c2ro-Wunit-Gram}
\end{align}
Every denominator is pure coefficient.  The four inverse square roots
$1/\sqrt{Ce_i}$ are bounded by hypothesis, and
$\sqrt{C/A}=C/\sqrt{CA}$ is infinitesimal.
All coefficients of this Gram are therefore in the bounded tensor
ring.  Every displayed correction has zero residue, whereas its five
positive square term pivots have residue $4$.  Thus its residue is exactly
\[
             4(K_1^2+K_2^2+K_4^2+K_5^2+T_C^2)\succ0.
\]
The strict-residue Gram lemma applies without any bound on the roots.
In particular it makes no comparison between $C$ and $p$ or $q$.

The other state chart is $Z=1$, $W$ infinitesimal or zero.  Here the
exact pure-testing conductor is a unit:
\begin{equation}
 q=\Gamma:=\frac{M+WJ+(pq-W)X+(Wpq-1)Y}{W^2-1}.
\label{pc7dacaa3258e:eq:c2ro-Zunit-conductor}
\end{equation}
Indeed $M+WJ=q(W^2-1)+(W-pq)X+(1-Wpq)Y$.
Substitution in \eqref{pc7dacaa3258e:eq:c2ro-core-conductor} gives
\begin{equation}
 \frac H2=L(X^2+Y^2)-2pqXY
 +\Gamma\{WJ-M+(W+pq)X-(pqW+1)Y\}.
\label{pc7dacaa3258e:eq:c2ro-Zunit-core}
\end{equation}
This is an exact quadratic form on $X,M,Y,J$, with bounded pure-testing
coefficients.  Retain the whole $C$ and $\alpha$ positive square terms separately.
Scale the four remaining roots as before.  Every coefficient of the
scaled core has zero residue, because each contains two inverse strong
positive square term square roots.  Adding the four strong positive square terms consequently
produces the strict residue
$4(K_1^2+K_2^2+K_4^2+K_5^2)$.  Again only the Gram coefficients,
not $X,Y$ or the positive square term roots, needed boundedness.

These two state charts are exhaustive when $(W,Z)\ne(0,0)$: if
$W\ne0$ and $Z/W$ is bounded use the first; otherwise $Z\ne0$ and
$W/Z$ is infinitesimal, giving the second.  Normalization and restoration
use pure-testing factors.  When $W=Z=0$,
\eqref{pc7dacaa3258e:eq:c2ro-core-conductor} is just the bounded quadratic form
$H/2=L(X^2+Y^2)-2pqXY$; after the same four-root scaling its
coefficients have zero residue, so the four retained strong positive square terms
again give a strict Gram.  This includes the zero state.  Splitting a
complex state into real and imaginary parts completes the proof.
\end{proof}

\begin{corollary}[The radial source has no missing outer commuting chart]
\label{pc7dacaa3258e:cor:c2-radial-outer}
Use the unchanged exact source \eqref{eq:c2-native-source} under the hypotheses
of Theorem~\ref{p163e0ebf59f6:thm:c2-radial-corner}: $h,r,g,j$ are positive
infinitesimals and $a,1-r,t$ are positive bounded units.
Then the full parent is tensor SOS at every commuting residue test and
every real or complex state.  Together with the physical-source
noncommuting-residue theorem, this proves all tests for this exact
radial source family.
\end{corollary}
\begin{proof}
The source identities and positive decomposition of $N$ proved in
Section~\ref{p163e0ebf59f6:sec:c2-radial-corner} give
\[
 N\asymp h+r+g+j,\qquad
 U_1=re_1,\ U_2=ge_2,\ U_4=he_4,\ U_5=je_5>0
 \quad\hbox{bounded units},
\]
where $N/(h+r+g+j)$ has residue $1$.  Moreover, directly from
\eqref{eq:c2-native-source},
\[
 \frac CN=\frac{A_c}{2a^2t\Delta}>0
 \quad\hbox{is a bounded unit}.
\]
For example,
\[
 \frac1{Ce_1}=\frac{r/N}{(C/N)U_1},
\]
and the three analogous formulas use $g/N,h/N,j/N$.
Thus all four product reciprocals in
Lemma~\ref{pc7dacaa3258e:lem:c2-four-strong-outer} are bounded, while $C$ is positive
infinitesimal.  That lemma covers every outer commuting chart, and
Theorem~\ref{p163e0ebf59f6:thm:c2-radial-corner} covers the zero-zero chart.

The final noncommuting join uses only the already-established
physicality of the exact source and the physical-source
noncommuting-residue theorem.  In particular the exact contact and
Laurent-determinant physicality argument used for the unit-$N$ source
requires $N>0$, not that $N$ have positive residue.  All five event
weights are positive.  On the three-dimensional state subspace $Z=0$
the signed core is nonnegative, and the positive event positive square terms have
zero common kernel at the noncommuting contact; hence the state
restriction is positive definite.  The exact stationary state is a
nonzero kernel vector, so the other three state eigenvalues are positive.
Determinant stationarity remains exact, and the same Laurent-determinant
argument proves physicality.
No coefficient-entry claim is made here for C2 sources outside the
stated radial family.
\end{proof}

\section{Strong sources with two small actual angles}
The following full-parent polynomial identities and normal divisions are
also used by the C1 coefficient domains. Their use there requires the
C1 source bounds and coefficient identities proved in the C1 chapter;
it does not use a C2 source conclusion with altered coefficients.

\subsection{Local strong-II Gram}
\label{component:pe16499ccc4c6}

\label{pe16499ccc4c6:sec:r7-local-strong-II}

\begin{theorem}[Strong-II local allocation]
Let the positive C2 source satisfy $z,y\in\Inf_{>0}$ and let
$B=H+L$ be a positive unit.  Keep the exact source coefficients of
(GE1), with $n\geq0$ and the remaining displayed weights positive.
Then the original C2 parent belongs to $\sum(K_c\otimes_{\mathbb R}
K_t)^2$ whenever both independent affine testing coordinates $u,v$
are infinitesimal, including either literal testing axis.
There is no comparison hypothesis between $u^2,z$, between $v^2,y$,
or between the two source rates.  All openings of $r$, $L/y$,
$1+zr-L/y$, and $1-H$ are retained, as is literal $n=0$.
\end{theorem}

This is the local full-parent calculation. Section~\ref{component:p590112f102bd} proves the other testing charts and the literal projective endpoints.
All identities below are polynomial in testing coordinates.  Inverses
are taken only in the pure source field, and the final step uses the
bounded Gram strict-real-residue lemma already used in the completed
local UI certificate.  In particular no ordering or inverse of the
mixed tensor is used.

\paragraph{Source units and notation.}
Write $U=u^2,V=v^2,k=\sqrt y$ and
\[
 a_c=1+zr,\quad t=L/y,\quad j=a_c-t,\quad s=a_c/B,
 \quad c=CL,\quad d=Dyj.
\]
The source entry gives bounded $t,j,r$, positive units $B,c,d,m/z^2,A/y$,
bounded $nr$ and $\alpha(1-H)$, and strong $\alpha$.
These facts include the arbitrary openings $t\to0$, $j\to0$,
$r\to0$, and $1-H\to0$.  They follow directly from the positive
source formulas and were proved in (DI7).  Let $h=\operatorname{st}B$;
then $0<h\leq1$ and, with $r_0=\operatorname{st}r$,
\[
\begin{gathered}
 c_0=d_0^{\rm pair}=\frac{h}{2(1+h)},\quad
 \operatorname{st}(m/z^2)=\frac{1-r_0}{2h^2(1+h)},\quad
 \operatorname{st}(A/y)=\frac1{2h(1+h)},\\
 \operatorname{st}(nr)=\frac1{2(1+h)}-\frac{r_0}{2},\qquad
 \operatorname{st}\{\alpha(1-H)\}=\frac{h^2}{2(1+h)}.
\end{gathered}
\tag{LII1}
\]
These residues follow directly from (GE1): $L\to0$, $H\to h$,
$a_c\to1$, $G\to h$, $\delta\to h(1+h)$, and the source
numerator $M\to h^2$.  Substitution in the six rational coefficient
formulas gives (LII1), including the identity
$\operatorname{st}(nr)=1/[2(1+h)]-r_0/2$ without division by $n$.
Define the source quantities (with $q,g$ bounded)
\[
 \ell=y(C+D)=c/t+d/j,\qquad
 q=\frac{D}{C+D}=\frac{dt}{cj+dt},\qquad
 g=\frac{yCD}{C+D}=\frac{cd}{cj+dt}.
\]
Here $0<q<1$, $g$ is a positive unit with residue $h/[2(1+h)]$,
and $\ell$ has a positive real lower bound, although it may be
unbounded.  Indeed $cj+dt$ is a unit because $t+j=a_c$ is a unit.
Choose positive real $\ell_*\leq\ell$ and
$\alpha_*\leq\alpha$.

\paragraph{The state roots and the exact C,D completion.}
Use the source-fixed parity graph
\[
 X=-\frac{zr v}{B}W+\xi,\qquad
 Y=\frac{Hu}{B}W+\eta,\qquad
 Z=\frac{uv}{B}W+\zeta.                         \tag{LII2}
\]
It agrees with the actual kernel at $U=z,V=y$ and is a change of all
four homogeneous state coordinates with determinant one.  Set
\[
 Q_n=u\xi-\zeta,\quad K=ukW,\quad
 R_h=\frac{vQ_n-su(V-y)W}{k},\qquad
 N=\frac\eta k-qR_h.                              \tag{LII3}
\]
Thus $\eta=k(N+qR_h)$.  Every root in (LII3) vanishes at the actual
contact.  If $Q=uW-Y$ and $J=Q+v(uX-Z)$, then
\[
 Q/k=(t/B)K-N-qR_h,\qquad
 J/k=(1-q)R_h-N-(j/B)K.
\]
Consequently the exact weighted completion is
\[
\begin{split}
 4CQ^2+4DJ^2={}&4\ell N^2+4gR_h^2
 +\frac{8(d-c)}BKN-8gsKR_h
 +\frac{4(ct+dj)}{B^2}K^2,\\
 4CVQ^2={}&4CV\eta^2
 +\frac{4ct}{B^2}VK^2-\frac{8c}BVK(N+qR_h).
                                                        \end{split}\tag{LII4}
\]
There is no estimate in these identities.  Retain the following
four nonnegative, contact-centered positive square terms outside the remainder:
\[
 4(\ell-\ell_*)N^2,\qquad 4CV\eta^2,\qquad
 4n(1+U)(1+V)\xi^2,\qquad
 4(\alpha-\alpha_*)(1+U)(v\eta-\zeta)^2.          \tag{LII5}
\]
Call the exact remainder $G$.  Its coefficients are bounded after
writing $\eta=k(N+qR_h)$ and using (LII4).  For completeness, its
entire formula is
\[
\begin{split}
G={}&4\ell_*N^2+4gR_h^2+\frac{8(d-c)}BKN-8gsKR_h\\
 &+\frac{4(ct+dj)}{B^2}K^2
 +\frac{4ct}{B^2}VK^2-\frac{8c}BVK(N+qR_h)\\
 &+4n(1+U)(1+V)
   \left\{\frac{z^2r^2V}{B^2}W^2-\frac{2zr v}B W\xi\right\}\\
 &+4\alpha_*(1+U)(v\eta-\zeta)^2
  +4\alpha(1+U)
   \left\{\frac{(1-H)^2UV}{B^2}W^2
        -\frac{2(1-H)uv}B W(v\eta-\zeta)\right\}\\
 &+4mM^2+4A(1+U)(1+V)Y^2\\
 &+2\{J^2+(1+U)(vY-Z)^2+(1+V)(X+uZ)^2
                      -(1+U)(1+V)Z^2\},        \end{split}\tag{LII6}
\]
where $M=vW-X+uvY-uZ$ and (LII2) is used in the last line.
Every potentially unbounded original coefficient occurs here only
through $nr$, $\alpha(1-H)$, or the bounded quantities in (LII1).
No division by a testing coordinate has been introduced.

\paragraph{One exact gauge removes the dependent-root first normal.}
In the formal bounded polynomial (LII6), the coefficient of $R_hW$
is $uk\rho(U,V)$ for a bounded source polynomial $\rho$.  Replace
this entire term using the exact identity
\[
 KR_h=uvWQ_n-sU(V-y)W^2.                         \tag{LII7}
\]
Call the resulting polynomial $\widetilde G$.  It is equal to $G$
on the roots (LII3), and has no $R_hW$ term.  This gauge is useful
because $R_h$ is a dependent root: it permits the state and angular
contact equations to be read as formal coefficient vanishings.

More precisely, by parity write uniquely
\[
 \widetilde G=Q_0(N,\xi,\zeta,R_h)
 +W\{vP_\xi(U,V)\xi+uvP_\zeta(U,V)\zeta
                         +ukP_N(U,V)N\}
 +\phi(U,V)W^2,                                  \tag{LII8}
\]
where all coefficients are bounded source polynomials and $Q_0$ is
quadratic in the four displayed roots.  Formula (LII6), followed by
the one replacement (LII7), is an exact coefficient-extraction
specification for every entry in (LII8).

The retained positive square terms (LII5) vanish together with their first jets at
the actual contact.  The exact original source state equations thus
give
\[
 P_\xi(z,y)=P_\zeta(z,y)=P_N(z,y)=0.              \tag{LII9}
\]
Indeed the physical $\eta$ derivative first forces the $N$ coefficient
to vanish; the physical $\xi,\zeta$ derivatives then force their
formal coefficients to vanish because the $R_hW$ coefficient was
removed.  The exact original contact and the two angular derivatives
now give
\[
 \phi(z,y)=\phi_U(z,y)=\phi_V(z,y)=0.              \tag{LII10}
\]
This uses the six exact original source substitutions already checked
in R6.  It imposes no new relation between source and testing fields.

\paragraph{Bounded division and the complete angular allocation.}
There are two additional elementary support facts in (LII8):
\[
 P_\xi\in(U,z),\qquad \deg_U\phi\leq2,\qquad
 \phi(0,V)=V(d_s+e_sV),                           \tag{LII11}
\]
where
\[
 d_s=\frac4{B^2}\{m(B+zr)^2+(n+\tfrac12)z^2r^2\},\qquad
 e_s=\frac4{B^2}(n+\tfrac12)z^2r^2.
\]
Both $d_s/z^2$ and $e_s/z^2$ are bounded.  The first support assertion
is literal bounded-coefficient membership: setting $U=z=0$ in the
extracted polynomial gives zero.  It is also clear before expansion:
all $\xi W$ terms not supplied by the gauge contain $U$, $m$, $nzr$,
or $zr$; those supplied by the gauge contain $U$.

Monic division using (LII9) therefore gives bounded polynomials
\[
\begin{split}
 P_\xi&=(U-z)A_\xi+z(V-y)B_\xi,\\
 P_\zeta&=(U-z)A_\zeta+(V-y)B_\zeta,\\
 P_N&=(U-z)A_N+(V-y)B_N.                         \end{split}\tag{LII12}
\]
For the first line, evaluate $P_\xi$ at $U=z$: the support assertion
supplies the factor $z$ with bounded quotient, and its value at $V=y$
is zero.  All other divisions are by the monic polynomials $U-z$
or $V-y$.  In particular none creates an inverse of a small source
or a testing difference.

The scalar has the exact bounded decomposition
\[
 \phi=\frac{d_s+e_sV}{z^2}V(U-z)^2
       +U(U-z)(V-y)A_s(V)+U(V-y)^2C_s(V).        \tag{LII13}
\]
Here is a constructive proof, including boundedness.  Write
$\phi=a(V)U^2+b(V)U+V(d_s+e_sV)$ and subtract the first term of
(LII13).  The remaining $U^2,U$ coefficients $a_1,b_1$ vanish at
$V=y$ by the value and $U$ derivative in (LII10).  Divide them by
$V-y$ to obtain $A_s,B_s$.  The $V$ derivative gives
$zA_s(y)+B_s(y)=0$, so set $C_s=(B_s+zA_s)/(V-y)$.
All three divisions are monic divisions of bounded polynomials.
This proves (LII13) exactly, not asymptotically.

Introduce the two legal angular roots
\[
 R=v(U-z)W,\qquad T=uk(U-z)W.                    \tag{LII14}
\]
The following identities allocate every term in (LII12)--(LII13):
\[
\begin{split}
 zv(V-y)W&=\frac{uVQ_n-uvkR_h}{s}-(V-y)R,\\
 uv(V-y)W&=\frac{VQ_n-vkR_h}{s},\\
 uk(V-y)W&=\frac{vkQ_n-yR_h}{s},\\
 U(U-z)(V-y)W^2&=-\frac{R_hT-uRQ_n}{s},\\
 U(V-y)^2W^2&=\frac{VQ_n^2-2vkR_hQ_n+yR_h^2}{s^2}.
                                                        \end{split}\tag{LII15}
\]
Together with $v(U-z)W=R$, $uv(U-z)W=uR$, and
$uk(U-z)W=T$, these are an explicit complete allocation of the
full remainder as a bounded symmetric Gram in
\[
                    (N,\xi,\zeta,R_h,R,T).       \tag{LII16}
\]
The original first state normals are all present.  The scalar
restriction has not been used as a substitute for them.

\paragraph{The finite block and the missing pivot.}
The residue of $Q_0$ on the first four roots is
\[
             \operatorname{diag}(4\ell_*,2,4\alpha_*,4g_0).
                                                        \tag{LII17}
\]
For example the signed core at $W=0$ is exactly
$2(\xi+u\zeta-uv\eta)^2+2(1+U)V\xi^2+2(1+V)\eta^2$;
all its other entries vanish in the present residue.  The extra
terms in the first three lines of (LII15) have infinitesimal Gram
coefficients except for $A_\xi\xi R$ and the allowed $A_NNT$.
Likewise the last line has only infinitesimal coefficients, and
the fourth line has one potentially finite $R_hT$ entry.  Thus,
apart from entries in the last row and column $T$, the sole finite
off-diagonal entry is between $\xi$ and $R$.

The exact coefficient extraction yields
\[
 \operatorname{st}A_\xi=\frac4{h(1+h)},\qquad
 \operatorname{st}(d_s/z^2)=\frac2{h^2(1+h)}.     \tag{LII18}
\]
For a short independent check of the first value, the unscaled
original first-normal coefficient on the $V=y$ slice is
\[
 -4nzr(1+y)+4m(1-H)-(4D+2)yj+2(1+y)
 \longrightarrow 2-4d_0^{\rm pair}=\frac2{1+h}.
\]
Its full multiplier is $2/B$.  Passing from the original normal
to (LII8) changes this $U$-quotient only by a bounded multiple of
$U$ or $z$, so its residue is unchanged.  The second value follows
from (LII1):
\[
 \operatorname{st}\{m(B+zr)^2/z^2+(n+\tfrac12)r^2\}
 =\frac{1-r_0}{2(1+h)}+\frac{r_0}{2(1+h)}
 =\frac1{2(1+h)}.
\]
Hence the finite $(\xi,R)$ Gram block is exactly
\[
 \begin{pmatrix}
  2&\dfrac2{h(1+h)}\\[2mm]
  \dfrac2{h(1+h)}&\dfrac2{h^2(1+h)}
 \end{pmatrix},\qquad
 \text{Schur complement}=\frac2{h(1+h)^2}>0.      \tag{LII19}
\]
This calculation retains $r_0=0$, $n=0$, and $h=1$.

Finally the exact toric conductor is the zero identity
\[
 \boxed{\quad
 T^2-UR^2-\frac{U-z}{s}R_hT
             +\frac{u(U-z)}s RQ_n=0.\quad}       \tag{LII20}
\]
Add any positive real multiple $\Lambda$ of this zero identity to
the Gram.  Its only nonzero residue contribution is $\Lambda T^2$.
The five-by-five block excluding $T$ is strictly positive by
(LII17)--(LII19); all entries incident to $T$ are bounded.  Choose
a sufficiently large real $\Lambda$ by the ordinary real Schur
complement.  The entire six-by-six residue Gram is then strictly
positive.  The bounded Gram lemma gives a tensor SOS for $G$, and
adding the four positive square terms (LII5) proves the theorem.

\paragraph{Quantifiers and endpoint discipline.}
No proof step divides by $u$, $v$, $U-z$, or $V-y$: parity quotients
and monic polynomial divisions are coefficient extractions whose
final identities are polynomial.  Thus the two literal affine axes
are included.  The source inverses $1/B$, $1/k$, and the source
positive square term square roots are legal even when their values are large;
all coefficients of the final Gram are bounded.  The source ratios
$t,j,r,1-H$ have not been bounded away from zero.  The new roots
raise testing degree, so the local certificate is not homogenized
to assert the original literal projective endpoints.  Those are
joined by the separate direct-parent remote certificate.

\paragraph{Source input.}
Source contact and stationarity follow from the native-entry identification in Chapter~\ref{ch:entry} and the exact source formulas in Section~\ref{component:p6e5d41ed80b3}. The scaled bounds used above are the II specialization of GE1--GE2 in Section~\ref{component:pf5726218aec8}, recorded there as DI7.

\subsection{Complementary strong-II tests}
\label{component:p590112f102bd}

\label{p590112f102bd:sec:r7-c2-remote}

Use the positive C2 source with $z,y$ infinitesimal and $B=H+L$ a
positive unit. We prove the full parent for every testing pair with at
least one noninfinitesimal angle, including literal axes and projective
infinities. The double-small testing chart was treated above.

\paragraph{Source constants, including all openings.}
Put $t=L/y$, $j=1+zr-t$, and write $h=\operatorname{res}H$.
Then $0<h\leq1$, $t+j=1+zr\sim1$, and the exact source formulas give
\[
 Cyt\sim Dyj\sim\frac{h}{2(1+h)},\qquad
 \frac Ay\sim\frac1{2h(1+h)}.
\]
Consequently $C,D$ are unbounded, and the positive pure-source quantity
\[
 \gamma=\frac{yCD}{C+D}
 =\frac{(Cyt)(Dyj)}{(Cyt)j+(Dyj)t}
 \sim\frac{h}{2(1+h)}.                                      \tag{RR1}
\]
These statements are uniform when either $t$ or $j$ tends to zero.
If $h<1$, then
\[
 \alpha\sim\frac{h^2}{2(1-h^2)};
\]
if $h=1$, then $\alpha$ is unbounded.  Thus, for every real
$0<h_0<h$, the positive real number
\[
 \beta_0=\frac{h_0^2}{2(1-h_0^2)}<\alpha                     \tag{RR2}
\]
is an admissible cap.  No lower bound on $r,j,t,1-H$ is imposed,
and $n=0$ is retained.

\paragraph{Homogeneous testing coordinates.}
For finite $u,v$ define
\[
 a=\frac{u}{\sqrt{1+u^2}},\quad b=\frac1{\sqrt{1+u^2}},\qquad
 c=\frac{v}{\sqrt{1+v^2}},\quad d=\frac1{\sqrt{1+v^2}}.
\]
Use the same homogeneous formulas at infinity, with $(a,b)=(1,0)$
or $(c,d)=(1,0)$.  These are the original parent's homogeneous
endpoint values.  In particular no raised-degree graph is being
specialized at an endpoint.  Put
\[
 E=aW-bY,\quad K=aX-bZ,\quad L=bX+aZ,
 \quad J=dE+cK,\quad Q=cY-dZ.
\]
The original parent divided by $4(1+u^2)(1+v^2)$ is exactly
\[
 nX^2+m\{bcW-bdX+acY-adZ\}^2+CE^2+AY^2
 +(D+\tfrac12)J^2+(\alpha+\tfrac12)Q^2
 +\tfrac12(L^2-Z^2).                                    \tag{RR3}
\]
At projective endpoints this equation means the original homogeneous
form, not a division by an infinite number.  All radicals and
normalizations used here belong to a single testing field.

\begin{proposition}[All tests with $v$ not infinitesimal]
For every such $v$, including $v=\infty$, and every $u$, including
$u=0,\infty$, the original full parent is tensor-SOS.
\end{proposition}
\begin{proof}
Here $c$ is a unit.  Set $P=bL$, $R_1=\sqrt C E$, and
$R_2=\sqrt D J$.  Orthogonality gives
\[
 L^2-Z^2=P^2+2aKP-b^2K^2,
 \qquad
 K=\frac{R_2}{c\sqrt D}-\frac{dR_1}{c\sqrt C}.            \tag{RR4}
\]
Both coefficients on the right are infinitesimal bounded tensors:
$1/c$ is bounded in the testing field, while $1/\sqrt C$ and
$1/\sqrt D$ are infinitesimal in the source field.  Therefore
\[
 R_1^2+R_2^2+\tfrac12P^2+aKP-\tfrac12b^2K^2
\]
has a bounded Gram with residue $\operatorname{diag}(1,1,1/2)$.
It is tensor-SOS by the strict bounded-Gram lemma.  Every other term
of (RR3), namely the retained $n,m,A,\alpha$ positive square terms and
$J^2/2+Q^2/2$, is already SOS.  This uses all four state coordinates.
\end{proof}

\paragraph{An exact harmonic identity for the remaining remote strip.}
Now let $v$ be infinitesimal, possibly zero, and let $u$ not be
infinitesimal, possibly infinity.  Thus $a$ is a unit.
Write $D_v=1+v^2$, and put
\[
 s=C+D,\quad\theta=D/s,\quad
 N=\sqrt{s}(E+\theta vK),\quad
 V=\frac{vK}{\sqrt y},\quad q=\sqrt yY.
\]
No comparison between $v$ and $\sqrt y$ is made.  The roots may be
unbounded; their Gram coefficients below are bounded.  The exact
pure-source harmonic identity is
\[
 CD_vE^2+D(E+vK)^2=N^2+\gamma V^2+Cv^2E^2.               \tag{RR5}
\]
Also set $x=bX$, $z_0=Z$, and $\ell=vY$.  The single toric relation
needed below is the polynomial identity
\[
 a x\ell-b^2z_0\ell=bVq.                                \tag{RR6}
\]
For a real cap $0<\beta_0<\alpha$, the original parent divided
by $4(1+u^2)$ is the sum of the explicitly positive square terms
\[
 \begin{split}
 &nD_vX^2+m\{b(vW-X)+a(vY-Z)\}^2+N^2+Cv^2E^2\\
 &\hspace{15mm}+\tfrac12(E+vK)^2+(\alpha-\beta_0)(\ell-z_0)^2
 \end{split}                                             \tag{RR7}
\]
and the following exact bounded quadratic form:
\[
 G=\gamma V^2+\frac AyD_vq^2
 +(\beta_0+\tfrac12)(\ell-z_0)^2
 +\frac{D_v}{2}(x^2+2axz_0-b^2z_0^2).                   \tag{RR8}
\]
This decomposition is valid also at $u=\infty$, by the original
homogeneous convention.  It has used no state graph or state inverse.

\begin{proposition}[All tests with $u$ not infinitesimal and $v$ infinitesimal]
The full original parent is tensor-SOS for all these tests, including
$v=0$ and $u=\infty$.
\end{proposition}
\begin{proof}
For a real constant $k$ define
\[
 M_k=(\beta_0+\tfrac12)(\ell-z_0)^2
 +\tfrac12(x^2+2axz_0-b^2z_0^2)
 +k(ax\ell-b^2z_0\ell).
\]
By (RR6), the exact form (RR8) equals
\[
 \boxed{G=D_vM_k+
 \gamma V^2+\frac AyD_vq^2-kD_vbVq
 -v^2(\beta_0+\tfrac12)(\ell-z_0)^2.}                   \tag{RR9}
\]
This is a full-parent identity; every state-normal term is present.

First suppose $b$ has positive residue $b_0$.  Since $a$ is a unit,
$0<b_0<1$.  Choose $0<h_0<h$ sufficiently close to $h$ that
\[
 b_0(1+h-h_0)<1,
\]
and take the admissible cap (RR2).  There exists a positive real
$\kappa$ satisfying
\[
 \frac1{1+b_0h_0}<\kappa<
 \min\left\{\frac1{1-b_0h_0},\frac1{b_0(1+h)}\right\}.
                                                               \tag{RR10}
\]
Set $k=-\kappa$.  The Gram of $M_k$ on $(x,z_0,\ell)$ has
first pivot $1/2$, second Schur pivot $\beta_0$, and determinant
\[
 \det M_k=\frac{2\beta_0 b^2k^2-(2\beta_0+1)(k+1)^2}{8}.
                                                               \tag{RR11}
\]
Its residue is strictly positive by the first two inequalities in
(RR10), because $h_0^2=2\beta_0/(2\beta_0+1)$.  The Gram on $(V,q)$
in (RR9) has positive diagonal residues and determinant residue
\[
 \frac14\left\{\frac1{(1+h)^2}-k^2b_0^2\right\}>0.       \tag{RR12}
\]
Thus both blocks have strictly positive residue.  The final term in
(RR9) is a bounded infinitesimal perturbation on $(x,z_0,\ell)$,
so the combined bounded Gram remains strict and is tensor-SOS.

If $b$ is infinitesimal, including $b=0$, choose any positive real
cap $\beta_0<\alpha$ and take $k=-1$.  There is now an exact
three-square identity
\[
 M_{-1}=\tfrac12\{x+a(z_0-\ell)\}^2
       +\beta_0(z_0-\ell)^2+\tfrac12(b\ell)^2.           \tag{RR13}
\]
On these three roots and $(V,q)$, (RR9) has positive diagonal
residue; the only off-block term is $D_vbVq$, whose coefficient
is infinitesimal.  The negative perturbation changes only the
$(z_0-\ell)^2$ coefficient by an infinitesimal.  Its Gram is again
strict.  Formula (RR7) restores all retained positive square terms.
The argument does not divide by $u,v,b$, or any state coordinate,
so the asserted literal boundaries are included.
\end{proof}

\paragraph{Scope and join.}
The two propositions cover precisely the complement of the region
where both testing angles are infinitesimal, for every strong-II
source with $B$ a unit.  They preserve independent source and testing
fields, arbitrary source openings, testing signs, and literal $n=0$.
Combined with a separately verified full-parent case of that
local region, they close the strong II/HH orbit by the exact C2
reflection.  They do not close the $B$-infinitesimal weak face.

\section{One small actual angle: bounded state scale}
\subsection{A lower bound on the state parameter}
\label{component:p7294aae40ee2}

\label{p7294aae40ee2:sec:c2-bounded-state-entry}

This short reduction concerns the same exact rational source, before
assuming that $b$ is a unit.  Suppose $x>0$ is infinitesimal, $k$ is
a positive bounded unit, $b>0$ is bounded, $0<r<1$, $0<h<b$,
$g=b-1+r+xh>0$, and $n\geq0$.  No further source physicality is
needed for the following necessary bound.

All quantities in
\[
             (b+r)A_cS_c\geq r\Delta
\]
are bounded, so taking residues gives
\[
 r_0\{b_0+k_0^2(b_0+1)\}\leq b_0.
\]
The brace is strictly positive even if $b_0=0$.  On the other hand
$g>0$ gives $r_0\geq1-b_0$.  If $b_0\leq1$, combining the two
inequalities yields
\[
 (1-b_0)\{b_0+k_0^2(b_0+1)\}\leq b_0,
 \quad\text{or equivalently}\quad
              b_0^2\geq\frac{k_0^2}{1+k_0^2}.
\]
For $b_0\geq1$ this lower bound is automatic.  Hence in every case
\[
                b_0\geq\frac{k_0}{\sqrt{1+k_0^2}}>0.
\]
Thus bounded $b$ is necessarily a unit.  In particular there is no
additional $b$-infinitesimal source domain in this actual-contact
scaling.  Combining the two new single-angle certificates, the
uncaptured part of this bounded-$b$ entry domain has exactly the
joint necessary degeneration $r_0=g_0=0$, hence $b_0=1$; the other
remaining scale is unbounded $b$.  This is a necessary source
reduction, not a certificate for the joint degeneration.

\subsection{A bounded-shape family at all tests}
\label{component:pec59340d644f}

\label{pec59340d644f:sec:c2-general-single-angle}

This calculation gives a source-domain extension of the fixed rational
family.  Its scalar uniqueness proof is an exact six-term Laurent
identity and a three-cycle arithmetic--geometric mean inequality.
The full-normal and outer-chart arguments include literal $n=0$.

Put $x>0$ infinitesimal in the coefficient field and $\rho=\sqrt{x}$.
Choose positive bounded coefficient units $k,b,r,h$ such that
\[
 1-r,\quad b-h,\quad g=b-1+r+xh,\quad b-k^2r
\]
are positive bounded units.  Set
\[
 c=\rho,\quad a=k/\rho,\quad t=x(b-h)/k^2,
 \quad j=1+xr-t,\quad a_2=k^2/x,
 \quad A_c=1+xr,\quad S_c=A_c+a_2xt.
\]
Thus $h=b-a_2t$, $g=(1+x)b-(1-r+a_2xt)$, and $j=1+xr-t$.
Use the exact original rational source
\begin{align*}
 \Delta={}&(1+x)xb^2+2x(1+x)br+(1+x)(1+a_2x)b
       +x^2(1-a_2)r^2+2x(1+a_2x)r+1+a_2x,\\
 N={}&(1+x)(b+r)(b+a_2(1-t))-a_2(1+a_2x)(1-r)t,\\
 n={}&\frac{(b+r)A_cS_c-r\Delta}{2r\Delta},\\
 m={}&\frac{(1+a_2x)(1-r)A_cS_c}{2g\Delta},\\
 C={}&\frac{A_cN}{2a_2t\Delta},\\
 A={}&\frac{a_2A_c(b+r)S_c}{2h\Delta},\\
 D={}&\frac{(1+x)(b+r)(b-a_2xr)S_c}{2a_2j\Delta},\\
 \alpha={}&\frac{x^2(b+r)N}{2(1-xh)\Delta}.
\end{align*}
Assume $n\geq0$ and physicality of this canonical source.  The latter
is the same inherited source-physicality hypothesis used for the
general noncommuting-residue exit; no conclusion below is predicated
on a new source enumeration.  Literal $n=0$ is allowed.

The parent, with $d=1$ and $\lambda=\beta=0$, is
\begin{align*}
 F={}&4nD_uD_vX^2+4mM^2+4CD_v(uW-Y)^2+4AD_uD_vY^2
                +4DJ^2+4\alpha D_u(vY-Z)^2+H,\\
 M={}&vW-X+uvY-uZ,\qquad J=uW+uvX-Y-vZ,\\
 H/2={}&J^2+D_u(vY-Z)^2+D_v(X+uZ)^2-D_uD_vZ^2,
 \qquad D_u=1+u^2,\quad D_v=1+v^2.
\end{align*}
Its exact stationary contact is
\[
 (u,v;W,X,Y,Z)=(\rho,k;\rho b/k,-\rho r,xh/k,1).
\]
Writing $\kappa=b/k$ and using a subscript zero for residues gives
\begin{align*}
 \Delta_0&=(1+k_0^2)(b_0+1),\qquad
 n_0=\frac{b_0+r_0-r_0\Delta_0}{2r_0\Delta_0}\geq0,\\
 m_0&=\frac{1-r_0}{2(b_0-1+r_0)(b_0+1)}>0,\\
 C_*:=(xC)_0&=\frac{k_0^2(b_0+r_0)}{2(b_0-h_0)\Delta_0},\qquad
 A_*:=(xA)_0=\frac{k_0^2(b_0+r_0)}{2h_0\Delta_0},\\
 D_*:=(D/x)_0&=\frac{(b_0+r_0)(b_0-k_0^2r_0)}{2k_0^2\Delta_0},\qquad
 E_*:=(\alpha/x)_0=\frac{k_0^2(b_0+r_0)^2}{2\Delta_0}.
\end{align*}
All starred quantities are positive.  Thus $n$ is bounded, $m$ is a
positive unit, $C,A$ have scale $x^{-1}$, and $D,\alpha$ have scale $x$.

\begin{theorem}[General bounded-shape single-angle source]
\label{pec59340d644f:thm:c2-general-single-angle}
Every source in the displayed domain has tensor SOS for the full
original parent at all independent real or complex projective testing
states and all projective testing angles.  No comparison between a
coefficient infinitesimal and a testing infinitesimal is needed.
The conclusion includes $n=0$.  Applying the exact C2 reflection gives
the corresponding original orientation with weak $C$ and strong
$\alpha$; this is a sufficient source domain, not a global one-weak
C2 classification.
\end{theorem}

\paragraph{The finite $u=0$ axis and its full normals.}
Set $a_0=4m_0$, $b_0'=a_0+2$, and
\[
 H_0(v)=(4n_0+2)(1+v^2)+a_0,\quad \eta=h/b,
 \quad \gamma=-r-\frac{a_0b-b_0'}{H_0(k)}.
\]
Here $\gamma$ is infinitesimal and $\eta_0=C_* /(C_*+A_*)$.
Use the exact contact graph
\[
 X_g=\frac{a_0vW-b_0'uZ}{H_0(v)}+\gamma uZ,
 \qquad Y_g=\eta uW.
\]
The denominator $H_0(v)$ is a pure testing unit and $H_0(k)$ is a
pure coefficient unit.  Let $N_X=X-X_g$, $N_Y=(Y-Y_g)/\rho$.
Writing $f=F|_g=f_{ww}W^2+f_{wz}WZ+f_{zz}Z^2$, exact monic divisions give
\[
 f=aP^2+b_s B^2+c_s V^2+eE^2+g_sBV,
 \qquad(P,B,V,E)=(uW/\rho,uZ,vW,\rho Z),
\]
where
\[
 \begin{gathered}
 a=x(f_{ww}-f_{ww}|_{u=0})/u^2,\qquad
 b_s=(f_{zz}-f_{zz}|_{u=0})/u^2,\\
 c_s=f_{ww}|_{u=0}/v^2,\qquad e=f_{zz}|_{u=0}/x,\qquad
 g_s=f_{wz}/(uv).
 \end{gathered}
\]
These displayed quotients are polynomial divisions, not inverses of
testing infinitesimals.  Parity gives the divisions by $u^2,v^2,uv$;
$f_{zz}|_{u=0}=4Dv^2+4\alpha$ gives the $x$ factor.  Only $C,A$ can
produce coefficient poles, and they occur in $f_{ww}$ with a factor
$u^2$; its prefactor $x$ removes them.

Put $Q=4C_*A_* /(C_*+A_*)>0$.  The bounded coefficient residues are
\begin{align*}
 q(v)=a_0^{\rm sc}&=Q D_v,&
 b(v)&=b_0'\frac{H_0-a_0-2}{H_0},\\
 \bar c(v)&=a_0\frac{H_0-a_0}{H_0},&
 g(v)&=-2b_0'\frac{H_0-a_0}{H_0},&
 e(v)&=4D_*v^2+4E_*.
\end{align*}
Thus $c(v):=v^2\bar c(v)$ and $d(v):=-vg(v)$ give the weighted
leading scalar
\[
 q(v)u^2W^2/x+b(v)u^2Z^2+c(v)W^2-d(v)uWZ+e(v)xZ^2.
\]
The full expansion is
\[
 F=f+N_XL_X+N_YL_Y+H_{XX}N_X^2+2H_{XY}N_XN_Y+H_{YY}N_Y^2.
\]
Its hard residue is $\operatorname{diag}(H_0(v),4D_v(C_*+A_*))$.
Direct differentiation gives
\[
 L_X=V\ell_w+B\ell_z,\qquad
 L_Y=P(B_2+xe_w)+B(\rho u e_z)+E e_*,
\]
where
\[
 B_2=8D_v\{(xA)\eta D_u-(xC)(1-\eta)\},\qquad
 e_*=8v(D-\alpha).
\]
All $\ell_w,\ell_z,e_w,e_z$ are bounded; the residues of
$\ell_w,\ell_z,B_2,e_*$ vanish identically in finite $v$.
For the second identity, after subtracting $uWB_2/x$ from $F_Y|_g$,
the remaining $W$ coefficient is $ue_w$ and the remaining $Z$
coefficient is $u^2e_z+e_*$.  The constant signed-core derivatives
cancel.  This proves full normal restoration wherever the scalar
Gram on these four roots is strict.

\paragraph{Global scalar strictness by an exact cycle identity.}
Let $\nu=2n_0+1\geq1$, $\mu=2m_0+1>1$, $z=v^2$,
$\mathcal H=\nu(1+z)+\mu-1$, $e_2=4D_*$ and $e_0=4E_*$.  Then
\[
 b=\frac{2\mu(\nu(1+z)-1)}{\mathcal H},\quad
 c=\frac{2(\mu-1)z\nu(1+z)}{\mathcal H},\quad
 d=\frac{4\mu v\nu(1+z)}{\mathcal H},\quad
 e=e_2z+e_0.
\]
For $s,z>0$ define
$\mathcal D(s,z)=(q s+c)(b s+e)-d^2s/4$.
The following exact identity is the key:
\begin{align*}
 \frac{\mathcal H\mathcal D(s,z)}{(1+z)sz}
 ={}&2Q\mu(\nu-1)s/z+2Q\mu\nu s
 +2(\mu-1)\nu e_2z/s+2(\mu-1)\nu e_0/s\\
 &+Qe_0(\nu+\mu-1)/z+Qe_2\nu z\\
 &+Q\{e_2(\nu+\mu-1)+e_0\nu\}-4\mu\nu.
\end{align*}
Every nonconstant coefficient is positive, except the first may be
zero when $n_0=0$.  Exact contact and stationarity, followed by taking
residues, give a stationary zero at $(s,z)=(1,k_0^2)$.  Here is an
algebraic proof that it is the unique zero on the whole positive
quadrant.  Rescale $z=k_0^2t$ and write the Laurent expression as
\[
 \mathcal L=a_\ell s/t+b_\ell s+c_\ell t/s+d_\ell/s
                      +e_\ell/t+f_\ell t-\Sigma,
 \qquad \Sigma=a_\ell+b_\ell+c_\ell+d_\ell+e_\ell+f_\ell.
\]
The constant term has this form by its zero value at $(1,1)$.
Stationarity gives
$a_\ell-c_\ell=d_\ell-b_\ell=f_\ell-e_\ell=:\tau$.
Put $p_\ell=\min(a_\ell,c_\ell)\geq0$,
$q_\ell=\min(b_\ell,d_\ell)>0$,
$r_\ell=\min(e_\ell,f_\ell)>0$.  Exactly,
\begin{align*}
 \mathcal L={}&p_\ell(s/t+t/s-2)+q_\ell(s+1/s-2)
                              +r_\ell(t+1/t-2)\\
 &+\begin{cases}
 \tau(s/t+1/s+t-3),&\tau\geq0,\\
 (-\tau)(t/s+s+1/t-3),&\tau\leq0.
 \end{cases}
\end{align*}
Each bracket is nonnegative by the two- or three-term
arithmetic--geometric mean inequality.  The two strictly positive
pair terms force $s=t=1$ at equality.  They also give a positive
definite Hessian there: the other brackets are nonnegative and
vanish at that point.  Thus both remote uniqueness and local
nondegeneracy are proved by the same finite algebraic identity, over
any real closed residue field.  No inference from a merely local
Hessian is being used.

Consequently, for finite nonzero $v$ with $v^2\ne k_0^2$,
\[
 |d(v)|<2\sqrt{b(v)c(v)}+2\sqrt{q(v)e(v)}.
\]
Indeed the minimum of $\mathcal D(s,z)/s$ over $s>0$ is
$(\sqrt{bc}+\sqrt{qe})^2-d^2/4$ and is attained, so positivity is
strict.  Split the exact cross term using $BV=vPE$ and the pure
testing function
\[
 \theta=\frac{|v|\sqrt{b(v)\bar c(v)}}
 {|v|\sqrt{b(v)\bar c(v)}+\sqrt{q(v)e(v)}}.
\]
The resulting residue blocks on $(B,V)$ and $(P,E)$ are strictly
positive.  The denominator is a pure testing unit.  At $v_0=0$ the
same conclusion holds when $n_0>0$, with zero cross residues.
The case $n_0=0,v_0=0$ is treated separately below.

\paragraph{The matched finite-angle residue.}
For $v_0=k_0$, normalize $Z=1$ when $Z\ne0$.  If $W$ is
infinitesimal, put
\[
 R=(uW-\kappa x)/\rho,\quad S=u-\kappa^{-1}W,
 \quad U=\rho(v-k),\quad T=u(v-k).
\]
The exact conductors are
\[
 u^2-x=uS+\kappa^{-1}\rho R,\qquad
 W^2-\kappa^2x=-\kappa WS+\kappa\rho R,
 \qquad uW-\kappa x=\rho R.
\]
Write $\xi=u^2,\omega=W^2,p=uW$ and
\[
 f=p^2Q_b/x+B_b,\qquad
 Q_b=4D_v\{(xC)(1-\eta)^2+(xA)\eta^2D_u\}.
\]
Here $B_b$ is bounded and polynomial in these three invariants with
only the separate unit denominators already specified.  Put
$G=\kappa^2xQ_b+B_b$ and take successive monic differences
\[
 G=G_0(v)+(\xi-x)G_\xi+(\omega-\kappa^2x)G_\omega
                    +(p-\kappa x)G_p.
\]
Exactly,
\[
 f=R^2Q_b+\kappa^2S^2G_\omega+\rho R\mathcal A+uS\mathcal B+G_0,
\]
where $\mathcal A=2\kappa Q_b+\kappa^{-1}G_\xi+
\kappa G_\omega+G_p$ and $\mathcal B=G_\xi-\kappa^2G_\omega$.
The contact Jacobian of $(R,S)$ with respect to $(u,W)$ has
determinant $-2$.  Exact stationarity therefore gives
$\mathcal A_c=\mathcal B_c=0$.  Cancelling $\rho$ here proves
vanishing and does not introduce a Gram coefficient inverse.
Monic contact division of these two functions and the conductors
above give a bounded Gram on $(R,S,U,T)$.
Also $G_0=xg_0(v)$ with bounded $g_0$, and exact contact gives a
double factor $(v-k)^2$, producing its bounded $U^2$ coefficient.

The full normals have zero-residue lifts on these roots.  For $L_X$,
write $L_X=uA_1+WB_1=-\kappa B_1S+u(A_1+\kappa B_1)$.
Both coefficient functions have identically zero leading part for
every $v$; exact normal contact makes $A_1+\kappa B_1$ vanish at
contact.  Monic division gives coefficients with factors $u\rho$,
$u^2$, or $uW$, and a $(v-k)$ quotient of an identically zero
leading function.  All have zero residue.  For $L_Y$, write
$L_Y=RB_2+\rho E_2$.  Its bounded even $E_2$ vanishes exactly at
contact, and both $B_2,E_2$ have identically zero leading part.
Monic division gives factors $x,\rho u,\rho W$, or the same
zero-leading-function quotient.  Thus the entire normal restoration
is bounded and has zero residue.

The soft residue has the form
\[
 q(k_0)R^2+\kappa_0^2c(k_0)S^2+A_1'RU+B_1'ST+E_1'U^2.
\]
It satisfies
\[
 E_1'>\frac{(A_1')^2}{4q(k_0)}+
          \frac{(B_1')^2}{4\kappa_0^2c(k_0)}.
\]
To prove this strict inequality, the leading scalar at
$u=\rho s,W=\rho w$ is
$\Phi=q s^2w^2+b s^2+c w^2-dsw+e$.
Complete its $w$ square:
\[
 \Phi=(qs^2+c)\left(w-\frac{ds}{2(qs^2+c)}\right)^2
           +\frac{\mathcal D(s^2,v^2)}{qs^2+c}.
\]
The global Laurent calculation gives a positive definite Hessian of
the second term in $(s,v)$ at $(1,k_0)$; both coordinate derivatives
are nonzero there.  The first square supplies the missing $w$
direction.  Hence the full three-variable Hessian is positive
definite.  On first variations $T=U$; the displayed inequality is
exactly its Schur complement after the invertible $(R,S)$ change.
Choose a real $\delta$ strictly between the second summand and
$E_1'-(A_1')^2/(4q(k_0))$.  Add $\delta$ times the exact zero
\[
 T^2-U^2-ST(v-k)-\kappa^{-1}RU(v-k)=0.
\]
The two blocks $(R,U)$ and $(S,T)$ are now strictly positive, and so
is the full six-root residue Gram.

If the ratio $W/Z$ is not infinitesimal, instead normalize $W=1$;
then $Z$ is bounded.  The earlier scalar is $aP^2$ plus a term of
positive residue $c(k_0)$; its normal lifts have zero residue on
$(P,1)$.  This is strict.  The case $Z=0$ is a literal sum of the
event squares and the three positive core squares.  The simultaneous
sign change $(u,v,X,Y)\mapsto(-u,-v,-X,-Y)$ covers $v_0=-k_0$.

\paragraph{The double-zero corner when $n_0=0$.}
Use instead the pure coefficient quantities
\[
 H_c=4n+4m+2,\quad \xi_c=4m/H_c,\quad
 \zeta=(4m+2)/H_c,\quad \eta_c=C/(C+A),
\]
and graph $X_g=\xi_cvW-\zeta uZ$, $Y_g=\eta_cuW$.
Put $\epsilon=4\alpha+4n(4m+2)/H_c>0$.
Then $\epsilon$ is infinitesimal, while $n/\epsilon$ and
$x/\epsilon$ are bounded.  Use the hard roots as before and scalar
roots
\[
 P=uW/\rho,\quad V=vW,\quad E=\rho Z,\quad
 B_\epsilon=\sqrt\epsilon\,uZ,\quad T=uvZ.
\]
Direct positive square term expansion gives exactly
\[
 f_{zz}=4\alpha+\epsilon u^2+4n\zeta^2u^4+4Dv^2
                         +u^2v^2 B_b,
\]
where $B_b$ is bounded with residue $2$.  Its first four terms
give coefficients $4\alpha/x$, $1+(4n\zeta^2/\epsilon)u^2$,
and $(4D/x)v^2$ on $E^2,B_\epsilon^2,E^2$ respectively.  The
$W^2$ contribution has bounded coefficients on $P^2,V^2$ with
positive residues $Q$ and $4m_0/(2m_0+1)$.  The $WZ$ contribution
is $uvWZ$ times a bounded function; use $uvWZ=vPE$ to give it zero
cross residue.  Thus the five scalar roots have strict diagonal
residue $(Q,4m_0/(2m_0+1),4E_*,1,2)$.

For completeness, $L_X$ has a zero-residue bounded coefficient on
$V$.  At $v=0$ its $Z$ coefficient is exactly $-8n\zeta u^3$;
on $B_\epsilon$ this is $-8n\zeta u^2/\sqrt\epsilon$, which is
infinitesimal because $n/\epsilon$ is bounded.  The remaining
$Z$ coefficient has a factor $uv^2$, giving a zero-residue multiple
of $T$.  For $L_Y$, the large-positive square term coefficient is now exactly
$B_2=8D_v(xA)\eta_cu^2$; the remaining terms lift as
\[
 L_Y=P(B_2+xe_w)+T(\rho u e_z)+E\,8v(D-\alpha),
\]
with bounded $e_w,e_z$.  All coefficients have zero residue.  The
hard residue is $\operatorname{diag}(4m_0+2,4(C_*+A_*))$.
This includes literal $n=0$ without deleting a coordinate.

\paragraph{The remaining commuting charts.}
The following exact core identities retain all unused event positive square terms.
Every retained coefficient is nonnegative; in particular the $n$
positive square term can be retained even when $n=0$.  They prove strict bounded
Grams without dividing by a testing infinitesimal.

If $u$ is a finite unit and $v\to0$, take
\[
 R=(uW-Y)/\rho,\quad L=Y/\rho,\quad B=X+uZ,\quad V=vZ.
\]
Then $M=-B+(\rho v/u)(R+D_uL)$, and exactly
\begin{align*}
 H/2={}&(1+D_uv^2)B^2+u^2D_uV^2-2uvD_uBV\\
 &+2\rho R(uvB-D_uV)+xR^2-2D_u\rho LV+D_uxv^2L^2.
\end{align*}
Retain $n,D,\alpha$.  The residue on $(B,V,R,L)$ is diagonal
with entries $(4m_0+2,2u_0^2D_{u,0},4C_*,4A_*D_{u,0})$.

If $u\to\infty$ and $v$ is finite, put $p=1/u$ and use $p^2F$.
Take $R=(W-pY)/\rho$, $L=Y/\rho$, $B=Z+pX$, $V=vX$.
The homogenized $m$ amplitude is
$-B+\rho pvR+\rho vD_pL$, and the exact core is
\begin{align*}
 \widehat H/2={}&B^2+D_pV^2+2\rho D_pRV-2\rho pvRB+xR^2\\
 &+2D_p\rho pLV-2D_p\rho vLB+D_pxv^2L^2.
\end{align*}
Retain $n,D,\alpha$.  The residue on $(B,V,R,L)$ is diagonal
with entries $(4m_0+2,2,4C_*D_{v,0},4A_*D_{v,0})$.

If $v\to\infty$ and $u$ is bounded, put $q=1/v$ and use $q^2F$.
Take $R=(uW-Y)/\rho$, $L=Y/\rho$, $T=uqZ$, $N_1=\rho Z$.
Exactly,
\begin{align*}
 \widehat H/2={}&(D_u+q^2)X^2+D_uxL^2+T^2+2qXT
       +q^2xR^2+2uq\rho XR-2qRN_1-2D_uqLN_1,\\
 4D\widehat J^2={}&4(D/x)(-N_1+\rho uX+xqR)^2.
\end{align*}
Retain $n,m,\alpha$.  The residue on $(X,R,L,T,N_1)$ is diagonal
with entries
\[ (2D_{u,0},4C_*,4A_*D_{u,0},2,4D_*). \]

Finally if both angles are infinite, use $p=1/u,q=1/v$ and
$p^2q^2F$.  Set $R=(W-pY)/\rho$, $L=Y/\rho$, $T=qZ$.  Then
\[
 \widehat H/2=(1+p^2D_q)X^2+xq^2R^2+xD_pL^2+T^2
       +2\rho qXR+2pqXT-2\rho pRT-2\rho D_pLT.
\]
Retain $n,m,D,\alpha$.  The residue on $(X,R,L,T)$ is
diagonal with entries $(2,4C_*,4A_*,2)$.

\paragraph{Exhaustion and independence.}
These charts cover every commuting angle residue: the whole finite
$u=0$ axis, the finite-unit $v=0$ axis, either infinite axis, and
their intersections.  When both angle residues are noncommuting,
apply the inherited physical noncommuting-residue theorem at arbitrary
coefficient scales.  Complex states require no additional chart:
if $Z\ne0$, multiply by a phase and normalize $Z=1$; the real parent
is the sum of its real-state value and its imaginary-state value,
the latter having $Z=0$ and hence being a literal square sum.
If $Z=0$ this is already true of both parts.  Homogeneity covers all
projective state charts.  Roots may be unbounded throughout.
All inverses or radicals used in coefficients are pure coefficient
or pure testing operations on units, except $\sqrt\epsilon$ in
roots and its controlled coefficient ratio $n/\sqrt\epsilon$.
No mixed inverse and no comparison of independent infinitesimals
occurs.  Applying the strict bounded-tensor Gram lemma completes
the proof.

\subsection{Three arbitrarily small source openings}
\label{component:p81893ad08bed}

\label{p81893ad08bed:sec:c2-three-openings}

This theorem removes three unit assumptions from the source of
Theorem~\ref{pec59340d644f:thm:c2-general-single-angle}.  Both positive openings $h$
and $b-h$ may approach zero, and $r$ may approach zero, making $n$
arbitrarily large.  The conclusions concern the entire original parent
and every independent test.  Source physicality remains an explicit
hypothesis; no new source classification is claimed.

\begin{theorem}[Three-opening single-angle source]
\label{p81893ad08bed:thm:c2-three-openings}
Use exactly the rational C2 source and contact of
Section~\ref{pec59340d644f:sec:c2-general-single-angle}.  Suppose $x>0$ is
infinitesimal, $k,b$ are positive bounded units,
\[
             0<r<1,\qquad 0<h<b,
\]
and $g=b-1+r+xh$ is a positive bounded unit.  Assume $n\geq0$
and physicality of the original source.  There is no lower bound on
$r,h,b-h$.  Then the entire original parent has tensor SOS at all
independent projective measurement pairs and all real or complex
projective states.  Literal $n=0$ is included.
\end{theorem}

Write $\rho=\sqrt x$, $\eta=h/b$, and retain the source quantities
$A_c,S_c,\Delta$ from that section.

\paragraph{Entry and the two complementary openings.}
First the former assumptions that $1-r$ and $b-k^2r$ are units follow
from $n\geq0$.  Indeed all quantities in
\[
                 (b+r)A_cS_c\geq r\Delta
\]
are bounded, and $\Delta_0=(1+k_0^2)(b_0+1)$.  Taking residues yields
\[
 r_0\leq\frac{b_0}{b_0+k_0^2(b_0+1)}<1,
 \qquad
 b_0-k_0^2r_0\geq
 \frac{b_0^2(1+k_0^2)}{b_0+k_0^2(b_0+1)}>0.
\]
Thus $m$ is a positive bounded unit, and $D/x,\alpha/x$ are positive
bounded units, regardless of whether $r$ has zero residue.

Set
\[
 T_C=xC(b-h),\qquad T_A=xAh,\qquad H=x(C+A).
\]
The original source formulas, with their apparent opening denominators
cancelled, show that $T_C,T_A$ are bounded positive units and have the
same residue
\[
 T_0=\frac{k_0^2(b_0+r_0)}{2\Delta_0}>0.
\]
Consequently $H$ is bounded below by a positive real, although it may
be unbounded.  The four identities
\begin{align*}
 xC(1-\eta)&=T_C/b,&xA\eta&=T_A/b,\\
 xC(1-\eta)^2&=T_C(b-h)/b^2,&
 xA\eta^2&=T_Ah/b^2                                      \tag{CO1}
\end{align*}
are the reason no new scalar singularity appears when either opening
collapses.  In particular their last two terms add to a bounded unit,
with residue $T_0/b_0$.  The effective scalar coefficient is
\[
                         Q=4T_0/b_0>0.                  \tag{CO2}
\]
It is independent of the limiting value of $h$.

\paragraph{A bounded $n$: restoration with the normalized hard root.}
If $r$ is a unit, the exact formula for $n$ makes $n$ bounded.  Use the
contact graph $X_g,Y_g=\eta uW$ of the prior theorem.  All of its scalar
coefficients remain bounded: the only potentially unbounded weights
are $C,A$, and on the graph they occur solely through the last two
products in (CO1), divided by $x$ and multiplied by $u^2W^2$.
Its leading scalar is exactly the scalar in that theorem with $Q$
replaced by (CO2).  The quantities $n_0,m_0,D_*,E_*$ have their former
values.  Hence the entire six-term Laurent identity, its balanced
AM--GM decomposition, global equality characterization, and positive
contact Hessian apply without change.

Replace the old hard root $N_Y=(Y-Y_g)/\rho$ by
\[
        \widehat N_Y=\sqrt H\,N_Y=\sqrt{C+A}(Y-Y_g).
                                                               \tag{CO3}
\]
Its exact hard coefficient from the $C,A$ positive square terms is
\[
             4D_v\frac{C+AD_u}{C+A}
       =4D_v\left(1+\frac{A}{C+A}u^2\right),             \tag{CO4}
\]
which is bounded and has residue $4D_{v,0}$ on the finite $u=0$
residue axis.  The other positive square term contributions to this coefficient
and to the mixed hard entry are bounded: $1/(C+A)=x/H$ and
$1/\sqrt{C+A}=\rho/\sqrt H$ are bounded and infinitesimal.

For clarity the complete first-normal verification is also retained.
The previous unnormalized $Y$ derivative has exactly the form
\[
 L_Y=P(B_2+xe_w)+B(\rho u e_z)+E e_*,\qquad
 B_2=8D_v\{(xA)\eta D_u-(xC)(1-\eta)\},
 \quad e_*=8v(D-\alpha).                                \tag{CO5}
\]
Here $P=uW/\rho$, $B=uZ$, $E=\rho Z$.  The functions $e_w,e_z$
come only from the bounded $n,m,D,\alpha$ positive square terms and the signed
core, so remain bounded.  Equation (CO1) makes $B_2$ bounded with
identically zero residue for every finite $v$ on $u=0$; $e_*$ has
the same property.  The coefficient of $\widehat N_Y$ is exactly
$L_Y/\sqrt H$, so its bounded zero-residue lifts are preserved.
The $X$ derivative contains neither $C$ nor $A$ and is unchanged.
Thus all four-root remote certificates restore both full normals.

For matched finite-angle residues the scalar has exactly the old
form $f=p^2Q_b/x+B_b$, with
\[
 Q_b=4D_v\{(xC)(1-\eta)^2+(xA)\eta^2D_u\}.
\]
Both $Q_b$ and $B_b$ have bounded coefficients by (CO1).  The monic
divisions in $u^2-x$, $W^2-\kappa^2x$, $uW-\kappa x$, and $v-k$,
where $\kappa=b/k$, therefore give precisely the same bounded Gram
on $R,S,U,T$.  Its residue is unchanged.  The old first-normal
contact divisions use only bounded $B_2,e_w,e_z$ and exact stationary
contact; multiplication by the pure coefficient $1/\sqrt H$
preserves their zero-residue lifts.  Hence the complete six-root
Gram has strictly positive residue.

If $n_0=0$ at the double-zero test, use the previous exact graph
with $\eta_c=C/(C+A)$.  Now
\[
 xC(1-\eta_c)=xA\eta_c=\frac{xCA}{C+A},
 \qquad
 \left(\frac{xCA}{C+A}\right)_0=T_0/b_0.                \tag{CO6}
\]
These quantities are bounded units even when $H$ is unbounded.
Thus the exact $\epsilon$ reserve and five scalar roots in that
proof remain valid.  Its $Y$ normal is again divided by $\sqrt H$,
and its hard $Y$ residue is now $4$ rather than $4(C_*+A_*)$.
This proves the literal $n=0$ case as well as its infinitesimal
neighborhood without deleting a state coordinate.

\paragraph{An infinitesimal $r$: the opposite event can be unbounded.}
Suppose now $r$ is infinitesimal.  Let $N=4n+2$.  The exact source
formula gives
\[
             Nr=\frac{2(b+r)A_cS_c}{\Delta}.             \tag{CO7}
\]
Thus $Nr$ is a bounded positive unit, $N$ is unbounded, and both
$nr$ and $nr^2$ are bounded.  Use the exact-contact graph
\[
 X_g=-ruZ,\qquad Y_g=\eta uW,
 \qquad
 \widehat N_X=\sqrt N(X+ruZ),\quad
 \widehat N_Y=\sqrt{C+A}(Y-\eta uW).                    \tag{CO8}
\]
At the original contact $u=\rho,v=k,W=\rho b/k,Z=1$ this graph
has exactly the required $X,Y$ values.  The potentially large $n$
positive square term contributes $4nr^2D_uD_vu^2Z^2$ to the graph scalar, and
$-8nrD_uD_vuZ$ to its unscaled $X$ derivative.  Equation (CO7)
therefore controls the entire scalar and every full first normal,
not only their values at contact.

The hard residue on $(\widehat N_X,\widehat N_Y)$ is
\[
                  \operatorname{diag}(D_{v,0},4D_{v,0}). \tag{CO9}
\]
Indeed $H_{XX}/N=D_uD_v+O(1/N)$, since the remaining source
weights that enter $H_{XX}$ are $m,D$; (CO4) controls $H_{YY}$.
The mixed entry is bounded times
$1/(\sqrt N\sqrt{C+A})$ and has zero residue.

The five exact scalar monic divisions on $P,B,V,E$ from the previous
proof still apply.  Their residues, with $z=v^2$, are now
\[
 q=Q(1+z),\quad b_s=2\mu,\quad c_s=2(\mu-1)z,
 \quad d_s=4\mu v,\quad e=E_2z+E_0,                    \tag{CO10}
\]
where $\mu=2m_0+1>1$, $E_2=4(D/x)_0>0$,
$E_0=4(\alpha/x)_0>0$, and $Q$ is given by (CO2).
Define
$\mathcal D(s,z)=(qs+c_s)(b_ss+e)-d_s^2s/4$.
Exactly,
\begin{align*}
 \frac{\mathcal D(s,z)}{sz}={}&
 2Q\mu s/z+2Q\mu s
 +2(\mu-1)E_2z/s+2(\mu-1)E_0/s\\
 &+QE_0/z+QE_2z+Q(E_0+E_2)-4\mu.                    \tag{CO11}
\end{align*}
All six nonconstant Laurent coefficients are positive.  Exact
stationary contact, applied after the bounded graph substitution
$u=\rho s_1,W=\rho w_1$, gives a stationary zero at
$(s,z)=(1,k_0^2)$.  Alternatively this follows by direct substitution
of
\[
 m_0=\frac1{2(b_0-1)(b_0+1)},\quad
 Q=\frac{2k_0^2}{\Delta_0},\quad
 E_2=\frac{2b_0^2}{k_0^2\Delta_0},\quad
 E_0=\frac{2k_0^2b_0^2}{\Delta_0}.
\]
Here $b_0>1$ follows from $g_0=b_0-1>0$.
Rescale $z=k_0^2t$ and use exactly the balanced reciprocal-pair and
three-cycle identity of the previous theorem.  Its $s+1/s-2$ and
$t+1/t-2$ coefficients are strictly positive.  Thus (CO11) has a
unique zero in the positive quadrant and a positive definite Hessian
there.  This proves remote scalar strictness and the matched-contact
Schur inequality, with no appeal to numerical evidence.

For the full $X$ normal, direct differentiation and parity give
\[
 \left.F_X\right|_g=V A_X+B B_X,
\]
where the complete coefficients are
\begin{align*}
 A_X={}&-8m(1+\eta u^2)+(8D+4)(1-\eta)u^2,\\
 B_X={}&-8nrD_uD_v+8m(1-r)
       -(8D+4)v^2(1+ru^2)+4D_v(1-r).
\end{align*}
These are bounded polynomials in the bounded testing coordinates.
Consequently the coefficient of $\widehat N_X$ has bounded
zero-residue coefficients on $V,B$ because $1/\sqrt N$ is
infinitesimal.  The full $Y$ normal is (CO5), with the bounded
graph (CO8); (CO1) gives all of its required bounds and zero residues.
Its complete residual coefficients are
\begin{align*}
 e_w={}&8mv^2(1+\eta u^2)-(8D+4)(1-\eta)
                     +(8\alpha+4)D_uv^2\eta,\\
 e_z={}&v\{-8m(1-r)+(8D+4)r-(8\alpha+4)\}.
\end{align*}
This restores the remote scalar certificates, including $v_0=0$,
where (CO10) has strictly positive coefficients on $P,B,V,E$ and
zero cross residues.  No special $n_0=0$ corner occurs in this case.

At a matched residue write, exactly,
\[
 \frac{F_X|_g}{\sqrt N}=uA_1+WB_1
   =-\kappa B_1S+u(A_1+\kappa B_1).
\]
Both $A_1,B_1$ have identically zero residue for every finite $v$;
their finite monic coquotients also retain the factor $1/\sqrt N$.
Exact normal contact makes $A_1+\kappa B_1$ vanish at the source
contact.  The same four conductors $R,S,U,T$ therefore give bounded
zero-residue coefficients, exactly as in the bounded-$n$ argument.
For the $Y$ normal the previous division
$L_Y=RB_2+\rho E_2'$ has bounded $B_2,E_2'$ with identically zero
leading functions; division by $\sqrt H$ preserves those properties.
All quadratic normal coefficients are already bounded by
(CO4),(CO7),(CO9).  Thus the matched-contact Gram for the entire
parent is strictly positive in residue.  The other state charts and
the simultaneous sign copy use precisely the previous proof.

\paragraph{Every other commuting test.}
Choose the pure coefficient
\[
                         \sigma=\min(1,xC,xA).
\]
It is a positive bounded unit, since $T_C,T_A$ are units while
$h,b-h$ are bounded.  Retain the nonnegative excess positive square terms
\[
 4(C-\sigma/x)D_v(uW-Y)^2,
 \qquad4(A-\sigma/x)D_uD_vY^2,
\]
and use the four outer core identities in the previous theorem with
$C=A=\sigma/x$.  These identities retain the entire $n$ positive square term;
they never use boundedness of $n$.  Their displayed hard entries
$4C_*,4A_*$ are simply replaced by $4\sigma_0$, while $m_0>0$
and $(D/x)_0>0$ are unchanged.  Every residue Gram stays strictly
positive.  This covers the finite-unit $u$ axis, both reciprocal
axes and their double corners, including literal projective ends.

The finite $u=0$ axis, those outer axes, and the inherited theorem for
noncommuting-residue tests exhaust all independent measurements.
Restore the retained positive square terms and positive homogeneous factors, and
split complex states into real and imaginary parts.  This proves
Theorem~\ref{p81893ad08bed:thm:c2-three-openings}.  In particular the newly proved
source scales permit either $A$ or $C$ to be arbitrarily larger than
$x^{-1}$, independently of an arbitrarily large $n$.  They cannot
both have unbounded ratios to this scale because $h+(b-h)=b$ is a unit.  There is no
prescribed comparison between these coefficient scales and any
testing infinitesimal.

\subsection{The remaining opening with bounded opposite event}
\label{component:pd7bd7f0d7fb6}

\label{pd7bd7f0d7fb6:sec:c2-last-opening}

\begin{theorem}[The unbounded-$m$ opening with bounded $n$]
\label{pd7bd7f0d7fb6:thm:c2-no-opening-bounds}
In Theorem~\ref{p81893ad08bed:thm:c2-three-openings}, replace the assumption that
$g=b-1+r+xh$ is a positive bounded unit by $g>0$ and the assumption that at least one of $r,g$ is a bounded
unit.  The all-tests tensor-SOS conclusion still holds.
\end{theorem}

The only new case is infinitesimal $g$; it is automatically bounded.
Our hypothesis now makes $r$ a unit.  Since $b_0+r_0=1$ and $b$ is
a positive unit, $r_0<1$.  Thus $n$ is bounded,
$0<b_0<1$, and $1-r$ is a unit.  The source formulas imply
\[
               mg=\frac{(1+k^2)(1-r)A_cS_c}{2\Delta},   \tag{CG1}
\]
which is a positive bounded unit.  Thus $m$ is unbounded.  The
two-opening bounds (CO1), positivity of $D/x,\alpha/x$, and the
consequence that $b-k^2r$ is a unit remain valid.

\paragraph{Use the original strong amplitude as the hard normal.}
Put $\eta=h/b$ and use the exact contact graph
\[
 Y_g=\eta uW,\qquad
 X_g=v(1+\eta u^2)W-(1+g)uZ.                            \tag{CG2}
\]
At contact this gives $X_g=\rho[b+xh-1-g]=-\rho r$.
The original $m$ amplitude satisfies exactly $M|_g=guZ$.
Normalize the full hard directions by
\[
 E_M=\sqrt m\{X-X_g-uv(Y-Y_g)\},\qquad
 E_Y=\sqrt{C+A}(Y-Y_g).                                 \tag{CG3}
\]
In these variables
\[
 M=guZ-E_M/\sqrt m,
 \quad
 X=X_g+E_M/\sqrt m+uvE_Y/\sqrt{C+A}.
\]
Consequently the unbounded positive square term has no $E_Y$ contribution at all.
This triangular change is essential: scaling $X$ alone would leave
the uncontrolled coefficient $\sqrt{m/(C+A)}uv$.
The entire scalar contribution of that positive square term is $4mg^2u^2Z^2$,
which is bounded by (CG1).  Its full $X$ derivative contributes
$-8mg\,uZ$, also bounded, before division by $\sqrt m$.

All other source weights entering the $X$ row are $n,D$ and the
unit signed weight, so the full $X$ derivative on (CG2) has bounded
coefficients on $V=vW,B=uZ$.  Hence the $E_M$ first normal has
zero-residue bounded coefficients on those roots.  The $E_Y$ first
normal is
\[
       \frac{(F_Y+uvF_X)|_g}{\sqrt{C+A}}.               \tag{CG4}
\]
The $m$ positive square term cancels identically in $F_Y+uvF_X$, before any graph
substitution.  The remaining large terms are exactly the $C,A$
terms controlled by (CO1); all other coefficients are bounded.
In particular the old (CO5) form on $P,B,E$ holds with bounded
coquotients and identically zero leading corrections.
The complete quadratic hard block is bounded and has residue
\[
                         \operatorname{diag}(4,4D_{v,0})  \tag{CG5}
\]
on the finite $u=0$ axis.  The $m$ positive square term gives precisely its first
diagonal entry, the $C,A$ positive square terms give the second, and all mixed
entries from bounded positive square terms contain $1/(\sqrt m\sqrt{C+A})$.

\paragraph{The remaining five Laurent terms already force strictness.}
The exact scalar divisions on $P,B,V,E$ are bounded for the same
reason: $C,A$ enter through (CO1), and $m$ only through $mg^2$.
Their leading weighted scalar is
\[
 q\,u^2W^2/x+b_su^2Z^2+c_sW^2-d_suWZ+exZ^2,
\]
where, for $z=v^2$ and $\nu=2n_0+1\geq1$,
\[
 \begin{gathered}
 q=Q(1+z),\quad b_s=2\nu(1+z)-2,\quad
 c_s=2z\nu(1+z),\quad d_s=4v\nu(1+z),\\
 e=E_2z+E_0,\qquad Q=4T_0/b_0,
 \quad E_2=4(D/x)_0>0,\quad E_0=4(\alpha/x)_0>0.
 \end{gathered}                                         \tag{CG6}
\]
Set $\mathcal D(s,z)=(qs+c_s)(b_ss+e)-d_s^2s/4$.
Its exact Laurent identity is
\begin{align*}
 \frac{\mathcal D(s,z)}{(1+z)sz}={}&
 2Q(\nu-1)s/z+2Q\nu s
 +2\nu E_2z/s+2\nu E_0/s\\
 &+QE_0/z+QE_2-4\nu.                                  \tag{CG7}
\end{align*}
Contact and stationarity give a zero at $(s,z)=(1,k_0^2)$.
After $z=k_0^2t$, write the five coefficients in order as
$a,b,c,d,e$.  Here $a\geq0$ and $b,c,d,e>0$.  Stationarity gives
$c=a+e$ and $b=d+e$.  The zero value fixes the constant, so exactly
\[
 \mathcal L=
 a(s/t+t/s-2)+d(s+1/s-2)+e(t/s+s+1/t-3).                 \tag{CG8}
\]
Every summand is nonnegative by AM--GM.  The last bracket alone
forces $s=t=1$ at equality: its three positive terms have product
one, and equality makes all three equal.  Its exponent vectors
$(-1,1),(1,0),(0,-1)$ span two dimensions; equivalently its Hessian
in $\log s,\log t$ at $(1,1)$ is
$\left(\begin{smallmatrix}2&-1\\-1&2\end{smallmatrix}\right)$.
Hence it already has a positive definite contact Hessian.
This proves global uniqueness and local nondegeneracy even when
$n_0=0$ and the first pair is absent.

For remote finite nonzero $v$ the same scalar cross-term split as
before gives a strict Gram.  The case $v_0=0,n_0>0$ is also strict
on $P,B,V,E$.  The full normals are restored by (CG3)--(CG5).
At matched finite-angle residues the whole scalar still has the
form $p^2Q_b/x+B_b$ with bounded $Q_b,B_b$.  The four monic
conductors $R,S,U,T$ of the prior theorem therefore give a bounded
Gram, and (CG8) proves its strict Schur inequality.

The full first normals have the required same conductor lifts.
For $E_M$, their coefficient is $(uA_1+WB_1)/\sqrt m$ with bounded
polynomial $A_1,B_1$.  Monic contact division retains the factor
$1/\sqrt m$ in every coquotient.  Exact contact makes the residual
coefficient $A_1+\kappa B_1$ vanish, giving the old four-root lift
with zero residue.  For $E_Y$, the cancellation in (CG4) leaves
precisely the old large-positive square term term $B_2$ plus bounded residual
terms; its exact contact and identically zero leading part give
the same $R B_2+\rho E_2'$ division, then the bounded factor
$1/\sqrt H$.  These are full normal identities, not just Hessian
checks.  Quadratic state degree and (CG5) finish the complete
matched-contact certificate.

\paragraph{The double-zero test when $n_0=0$.}
The graph (CG2) is replaced in this remote corner.  Put
\[
 H_c=4m+4n+2,\quad \xi=4m/H_c,\quad
 \zeta=(4m+2)/H_c,\quad \eta_c=C/(C+A).
\]
Use
\[
 Y_g=\eta_cuW,\quad
 X_g=\xi v(1+\eta_cu^2)W-\zeta uZ,
\]
and the triangular normals
\[
 E_X=\sqrt{H_c}\{X-X_g-uv(Y-Y_g)\},\qquad
 E_Y=\sqrt{C+A}(Y-Y_g).                                \tag{CG9}
\]
The entire graph amplitude is
\[
 M|_g=(1-\xi)v(1+\eta_cu^2)W+(\zeta-1)uZ.
\]
Both $m(1-\xi)$ and $m(\zeta-1)$ are bounded because $n$ is
bounded.  Thus the scalar and full first normals have bounded
coquotients after the same $x$ scaling as before.

Set $\epsilon=4\alpha+4n\zeta>0$.  It is infinitesimal, with
$n/\epsilon,x/\epsilon$ bounded.  The complete scalar $Z^2$
coefficient is exactly
\[
 f_{zz}=4\alpha+\epsilon u^2+4n\zeta^2u^4+4Dv^2
                            +u^2v^2 B_b,               \tag{CG10}
\]
where $B_b$ is bounded with residue $2$.  This identity follows
directly from the positive square terms and the exact relation
$4n\zeta^2+(4m+2)(1-\zeta)^2=4n\zeta$.
The $W^2$ coefficients on $P=uW/\rho,V=vW$ have positive residues
$Q,2$.  The cross term is $uvWZ$ times a bounded function and
can be written as a zero-residue multiple of $PE$ using
$uvWZ=vPE$.  Thus the five scalar roots
\[
           P,\quad V,\quad E=\rho Z,\quad
           B_\epsilon=\sqrt\epsilon\,uZ,\quad T=uvZ
\]
have strict diagonal residue $(Q,2,E_0,1,2)$.

For the full $X$ derivative, at $v=0$ the $Z$ coefficient is exactly
$-8n\zeta u^3$.  After division by $\sqrt{H_c}$ its coefficient
on $B_\epsilon$ is infinitesimal because
$n/\sqrt\epsilon=(n/\epsilon)\sqrt\epsilon$ is infinitesimal.
The remaining $Z$ coefficient has the factor $uv^2$ and bounded
coquotient; its coefficient on $T$ has zero residue.  The $W$
coefficient has a factor $v$ and bounded coquotient, multiplied by
$1/\sqrt{H_c}$.  For $E_Y$ the exact $m$ cancellation in
$F_Y+uvF_X$ again applies; (CO6) gives
$B_2=8D_v(xA)\eta_c u^2$, and the remaining lift is
\[
 P(B_2+xe_w)+T(\rho u e_z)+E\,8v(D-\alpha),
\]
divided by $\sqrt H$, with bounded $e_w,e_z$.
Every coefficient has zero residue.  The hard residue of (CG9) is
$\operatorname{diag}(1,4)$, and all mixed hard entries vanish in
residue.  This proves the full-parent corner including literal $n=0$.

\paragraph{All other tests.}
For the remaining commuting axes use the earlier outer identities,
cap $C,A$ at $\sigma/x$ as before, and also cap $m$ at
$\min(m,1)$.  Retain the entire nonnegative excess $m$ positive square term.
The selected $m$ weight is a unit and all previous displayed
strict residues apply.  The noncommuting-residue theorem covers
the other tests.  Restore all positive factors and split complex
states as before.  This proves the theorem for arbitrary positive
$g$, including the new unbounded-$m$ regime.

The remaining assumptions in this source family are now only that
$k,b$ are units, $x$ is positive infinitesimal, all original contact
openings are positive, at least one of $r,g$ is a unit, $n\geq0$,
and the source is physical.  The joint degeneration $r,g\to0$,
where both $n,m$ can be unbounded, is not covered.  The
result does not yet treat unbounded $b$ or a degenerate residue of
the finite actual angle $k$.  Section~\ref{p7294aae40ee2:sec:c2-bounded-state-entry}
excludes infinitesimal $b$ under these entry hypotheses.

\subsection{Joint opening of the two event coefficients}
\label{component:p8f912583feaa}

\label{p8f912583feaa:sec:c2-joint-centered}

\begin{theorem}[Both event openings may close]
Use the exact rational C2 source and original parent of
Section~\ref{component:pec59340d644f}.  Suppose $x>0$ is infinitesimal,
$k,b$ are positive bounded units, $0<r<1$, $0<h<b$, and
$g=b-1+r+xh>0$.  Suppose $r,g$ are infinitesimal, $n\geq0$, and
the source is physical.  Then the entire original parent is tensor
SOS at every independent real or complex projective testing state
and every independent projective pair of testing angles.  No lower
bound on either $h$ or $b-h$ is required.
\end{theorem}

Together with the three-opening and last-opening theorems this covers
the bounded-$k,b$ rational source family without requiring either $r$
or $g$ to be a unit. Nonunit $k$ and unbounded $b$ remain separate
source domains.

\paragraph{Source bounds, including both complementary openings.}
Write $\rho=\sqrt{x}$, $\eta=h/b$, $\kappa=b/k$, $K=k_0>0$,
and $D_K=1+K^2$.  Since $r,g$ are infinitesimal, $b_0=1$.
The source formulas give
\[
 (nr)_0=\frac1{4D_K},\qquad (mg)_0=\frac14,
 \qquad \Delta_0=2D_K.
\]
Thus $n,m$ are positive unbounded pure coefficient quantities,
$nr,mg$ are bounded units, and $nr^2,mg^2$ are infinitesimal.
In particular $m-1>0$.  The existing opening cancellation identities,
with $T_C=xC(b-h)$ and $T_A=xAh$, give bounded positive units
\[
 (T_C)_0=(T_A)_0=\frac{K^2}{4D_K},\qquad
 (D/x)_0=\frac1{4K^2D_K},\quad
 (\alpha/x)_0=\frac{K^2}{4D_K}.
\]
They hold even if $h$ or $b-h$ is infinitesimal.  Specifically,
\[
 xC(1-\eta)=T_C/b,\quad xA\eta=T_A/b,\quad
 xC(1-\eta)^2=T_C(b-h)/b^2,\quad
 xA\eta^2=T_Ah/b^2.                                  \tag{JC1}
\]
Put $H_c=x(C+A)$.  This quantity is bounded below by a positive
real, although it can be unbounded.  Its reciprocal and reciprocal
square root are bounded pure coefficient elements.

\paragraph{Center the excess positive square terms before taking normals.}
Retain exactly the original amplitudes
\[
 M=vW-X+uvY-uZ,\qquad J=uW+uvX-Y-vZ,
 \qquad D_u=1+u^2,\quad D_v=1+v^2.
\]
Let $F$ denote the full original parent, including its signed core
$H$.  Define $G$ by the following exact polynomial identity:
\[
 F=G+4nD_uD_v(X+ruZ)^2
             +4(m-1)(M-guZ)^2.                       \tag{JC2}
\]
An explicit expression, requiring no elimination or mixed inverse, is
\begin{align*}
G={}&4M^2+4CD_v(uW-Y)^2+4AD_uD_vY^2
       +4DJ^2+4\alpha D_u(vY-Z)^2+H\\
 &-8nrD_uD_vuZX-4nr^2D_uD_vu^2Z^2\\
 &+8(m-1)g\,uZM-4(m-1)g^2u^2Z^2.                  \tag{JC3}
\end{align*}
Every coefficient in (JC3) except $C,A$ is bounded.  At the exact
original contact
\[
 (u,v;W,X,Y,Z)=(\rho,k;\kappa\rho,-r\rho,xh/k,1)
\]
both subtracted squares and all of their first derivatives vanish.
Consequently $G$ has exactly the same zero and full stationary
contact as $F$.  This statement is essential: merely capping the
uncentered original positive square terms would not preserve the contact.

On the finite $u=0$ residue axis use
\begin{gather*}
 X_g=-ruZ,\qquad Y_g=\eta uW,\qquad E_X=\sqrt n(X+ruZ),\\
 E_Y=\sqrt{C+A}(Y-\eta uW),\qquad
 E_M=\sqrt{m-1}(M-guZ).                              \tag{JC4}
\end{gather*}
Expand only $G$ in $E_X,E_Y$ and retain $4E_M^2$ as written.
The root $E_M$ is allowed to be redundant.  No substitution for
$W$ is made.  In particular, the positive $C$ positive square term introduces
no coefficient of the form $Cu^2/n$ or $Cu^2/m$.

\paragraph{The complete graph scalar and its global leading inequality.}
Write $f=G|_g=f_{ww}W^2+f_{wz}WZ+f_{zz}Z^2$.  The same exact monic
divisions as in the prior single-angle proof give
\[
 f=aP^2+b_sB^2+c_sV^2+eE^2+g_sBV,
 \qquad (P,B,V,E)=(uW/\rho,uZ,vW,\rho Z),              \tag{JC5}
\]
where
\[
 a=x\frac{f_{ww}-f_{ww}|_{u=0}}{u^2},\quad
 b_s=\frac{f_{zz}-f_{zz}|_{u=0}}{u^2},\quad
 c_s=\frac{f_{ww}|_{u=0}}{v^2},\quad
 e=\frac{f_{zz}|_{u=0}}x,\quad g_s=\frac{f_{wz}}{uv}.
\]
These are polynomial coquotients, not inverses of testing
infinitesimals.  Their coefficients are bounded by (JC1)--(JC3).
The exact constant $f_{zz}|_{u=0}=4Dv^2+4\alpha$ supplies the $x$
factor.  For finite $v$ their residues are
\[
 q(v)=\frac{K^2(1+v^2)}{D_K},\qquad
 b_s=4,\quad c_s=4,\quad g_s=-10,\qquad
 e(v)=\frac{v^2/K^2+K^2}{D_K}.                        \tag{JC6}
\]
Thus the leading scalar, using the exact identity $BV=vPE$, is
\[
                  qP^2-2vPE+eE^2+4(V-B)^2.           \tag{JC7}
\]
The exact identity on the bounded full state graph is
\[
                   q(v)e(v)-v^2
                    =\frac{(v^2-K^2)^2}{D_K^2}.      \tag{JC8}
\]
If $v$ is a finite nonzero residue with $v^2\ne K^2$, then
$10|v|<8|v|+2\sqrt{qe}$.  Splitting the exact $BV=vPE$ cross
term with the pure testing function
\[
             \theta=\frac{4|v|}{4|v|+\sqrt{q(v)e(v)}}
\]
gives strictly positive residue blocks on $(B,V)$ and $(P,E)$.
At $v_0=0$, allocate the entire cross term using $BV=vPE$;
(JC5) then has strictly positive diagonal residue
$(q(0),4,4,e(0))$ and zero cross residues.  There is no exceptional
$n_0=0$ corner in the joint regime.

Here is also a global equality and Hessian check needed at the
matched angles.  For $z=v^2>0$ and $s>0$, set
\[
 \mathcal D(s,z)=(q s+4z)(4s+e)-25zs.
\]
Exactly,
\begin{align*}
 \frac{\mathcal D(s,z)}s
 &=4qs+4ez/s+qe-9z\\
 &=4\bigl(\sqrt{qs}-\sqrt{ez/s}\bigr)^2
       +(\sqrt{qe}-\sqrt z)(\sqrt{qe}+9\sqrt z).    \tag{JC9}
\end{align*}
The second term is nonnegative by (JC8); its equality forces
$z=K^2$.  The first then forces $s=1$.  Moreover the second term is
\[
 \frac{(z-K^2)^2}{D_K^2}
       \frac{\sqrt{qe}+9\sqrt z}{\sqrt{qe}+\sqrt z},
\]
whose quadratic coefficient at contact is strictly positive.
The first square supplies a nonzero $s$ derivative there.
Hence $\mathcal D$ has a positive definite Hessian at $(1,K^2)$.
These radicals are used only in the residue-field inequality; no
radical of an element mixing coefficient and testing fields is used.

\paragraph{Restore every normal, with bounded coefficients.}
From (JC2)--(JC4) the entire parent is exactly
\[
 F=f+E_XL_X+E_YL_Y+
 H_{XX}E_X^2+2H_{XY}E_XE_Y+H_{YY}E_Y^2+4E_M^2.        \tag{JC10}
\]
The first normals are $L_X=G_X|_g/\sqrt n$ and
$L_Y=G_Y|_g/\sqrt{C+A}$.  The only unbounded weights $C,A$ do not
enter the $X$ row.  Thus
\[
 L_X=V\ell_w+B\ell_z
\]
with bounded coefficient functions carrying the infinitesimal
pure coefficient factor $1/\sqrt n$.  Both have identically zero
residue on the entire finite $u=0$ axis, including $v=0$.

For the second normal the exact form is
\[
 L_Y=\frac1{\sqrt{H_c}}\left\{
 P(B_2+xe_w)+B(\rho u e_z)+E e_*\right\},            \tag{JC11}
\]
where
\[
 B_2=8D_v\{(xA)\eta D_u-(xC)(1-\eta)\},\qquad
 e_*=8v(D-\alpha).
\]
All $e_w,e_z$ are bounded: all new terms in (JC3) have bounded
coefficients, and their $Y$ derivative has an additional $u^2Z$
factor.  Equation (JC1) makes $B_2$ bounded with identically zero
residue, even at either closing complementary opening.  Thus every
coefficient in (JC11) has zero residue.

The entire hard block is bounded and has residue
\[
             \operatorname{diag}(4D_{v,0},4D_{v,0},4)
                   \quad\hbox{on }(E_X,E_Y,E_M).      \tag{JC12}
\]
The first entry comes from the retained centered $n$ positive square term;
the bounded $G$ contribution is divided by $n$.  The second is
\[
 4D_v\frac{C+AD_u}{C+A}
   =4D_v\left(1+\frac A{C+A}u^2\right)
\]
plus bounded terms divided by $C+A$.  The mixed entry contains
$1/(\sqrt n\sqrt{C+A})$ times a bounded function.  The third
entry is the unexpanded centered $m$ positive square term.  This proves the
full-parent assertion at all remote finite-axis tests, by the
strict Gram in (JC5) and the zero-residue first normals.

\paragraph{The matched finite-angle residues.}
At $v_0=K$, normalize $Z=1$ and suppose $W$ is infinitesimal.
Use exactly the prior contact conductors
\[
 R=(uW-\kappa x)/\rho,\quad S=u-\kappa^{-1}W,
 \quad U=\rho(v-k),\quad T=u(v-k),
\]
with exact identities
\[
 u^2-x=uS+\kappa^{-1}\rho R,\quad
 W^2-\kappa^2x=-\kappa WS+\kappa\rho R,
 \quad uW-\kappa x=\rho R.                           \tag{JC13}
\]
Here $f=p^2Q_b/x+B_b$, with bounded $B_b$ and
\[
 Q_b=4D_v\{(xC)(1-\eta)^2+(xA)\eta^2D_u\}.
\]
Boundedness follows from (JC1); the remaining terms use only
the bounded coefficients of (JC3).  The full exact stationary
contact of $G$ proved after (JC3) supplies every vanishing required
by the prior monic-contact-division argument.  Successive monic
differences in $u^2-x,W^2-\kappa^2x,uW-\kappa x,v-k$ therefore
give a bounded Gram on $R,S,U,T$.

For precision, its strict residue follows from the same explicit
Schur-complement step in that argument, now with
\[
 \Phi=q(v)s^2w^2+4s^2+4v^2w^2-10vsw+e(v).
\]
Completing the $w$ square gives
\[
 \Phi=(qs^2+4v^2)
       \left(w-\frac{5vs}{qs^2+4v^2}\right)^2
       +\frac{\mathcal D(s^2,v^2)}{qs^2+4v^2}.
\]
Equation (JC9) proves a positive definite three-variable Hessian
at $(s,w,v)=(1,1/K,K)$.  Consequently the residue before the
conductor-zero correction has form
\[
 q(K)R^2+4S^2+A_1RU+B_1ST+E_1U^2,
 \qquad E_1>A_1^2/(4q(K))+B_1^2/16.
\]
Choose a real $\delta$ strictly between $B_1^2/16$ and
$E_1-A_1^2/(4q(K))$, and add $\delta$ times the exact zero
\[
 T^2-U^2-ST(v-k)-\kappa^{-1}RU(v-k)=0.
\]
The two residue blocks $(R,U)$ and $(S,T)$ are strictly positive.

An explicit equivalent residue check can be used in place of this
Schur choice.  The capped current has the exact conductor lift
\[
 (M-gu)|_g=-v\kappa S+\kappa T
                  +\eta uv\rho R+\eta u\kappa\rho U.
\]
Retaining its $T^2$ term, rather than first moving that term to $U^2$,
gives the scalar residue
\[
 K^2R^2+\frac{2(K^2-1)}{D_K}RU+\frac{U^2}{K^2}
                  +4(S-T/K)^2.
\]
Adding $2/D_K^2$ times the same conductor zero leaves the $(R,U)$
Schur margin $2/D_K^2>0$, and leaves
$4(S-T/K)^2+(2/D_K^2)T^2$ on $(S,T)$, also strictly positive.

The full-normal lifts are also the ones actually required for
(JC10), rather than scalar-only lifts.  Write
$L_X=uA_1'+WB_1'=-\kappa B_1'S+u(A_1'+\kappa B_1')$.
Both functions and all their bounded monic coquotients retain
$1/\sqrt n$.  Exact $G_X=0$ at contact supplies vanishing of
$A_1'+\kappa B_1'$.  Using (JC13) gives zero-residue coefficients
on $R,S,U,T$.  For (JC11), write its numerator as $RB_2+\rho E_2$.
Both $B_2,E_2$ have identically zero leading functions and $E_2$
vanishes at contact.  The same monic divisions give factors
$x,\rho u,\rho W$, or a zero-leading-function $(v-k)$ coquotient;
division by $\sqrt{H_c}$ preserves boundedness and zero residue.
Thus the complete seven-root Gram on
$(R,S,U,T,E_X,E_Y,E_M)$ is strict in residue.

If $W/Z$ is not infinitesimal, normalize $W=1$ and keep $Z$
bounded.  Then the scalar has a strict residue on $(P,1)$:
its positive constant contribution is $4K^2$.  All normalized
first normals have zero residue as above.  For $Z=0$, the original
parent is already the literal sum of its event positive square terms and three
positive core squares.  The simultaneous sign copy gives $v_0=-K$.

\paragraph{Every other independent test.}
The outer-chart proof need not use the centered decomposition.
Use the original parent, cap $C,A$ at $\sigma/x$, where
$\sigma=\min(1,xC,xA)$ is a positive bounded unit, and cap $m$ at
$1$.  Retain the original nonnegative excess positive square terms and the
entire original $n$ positive square term.  The four exact outer-core identities
in the prior single-angle theorem apply without change: their
strict residues require only $\sigma_0>0$, the selected $m=1$,
and $(D/x)_0>0$.  They never bound $n$.  This covers the finite
unit $u$ axis, both reciprocal axes, all their intersections, and
literal projective ends.  The inherited physical
noncommuting-residue theorem covers the remaining tests.

These charts and the finite $u=0$ axis exhaust the projective measurement
pairs. Restore the retained squares and homogeneous factors, then apply
the real identities to the real and imaginary state parts.

\section{One small actual angle: unbounded state scale}
\subsection{Source reduction at unbounded state scale}
\label{component:pb60046a37191}

\label{pb60046a37191:sec:c2-unbounded-b-reduction}

Use the original positive rational source with $d=1$, actual
$u_c=\sqrt x$ positive infinitesimal and $v_c=k$ a positive bounded
unit. The global entry theorem supplies its positive chart inequalities.
The preceding rational theorems, including the joint-opening theorem,
now cover every bounded $b$: the inherited $n\ge0$ entry inequality
makes bounded $b$ a unit. We classify the remaining $b$ unbounded
without changing any original coefficient.

Put
\[
 B=xb,\qquad H=xh,\qquad L=x(b-h)=B-H,\qquad
 a^2=k^2/x,\quad t=L/k^2.
\]
The original gaps give
\[
 0<H<1,\qquad0<L<k^2(1+xr),\qquad B=H+L.
\]
Thus $B,H,L$ are bounded and $x/B=1/b$ is infinitesimal.
Set $A_c=1+xr$ and $S=A_c+L$, both positive bounded units. Also
\[
 g/b=1+(r-1)/b+H/b\longrightarrow1,
 \qquad j=1+xr-L/k^2>0,\quad j\le1+x.
\]
The following exact rearrangements of the original short polynomials
control every event scale:
\begin{align*}
 x^2N={}&(1+x)(B+xr)(H+k^2)-x(1+k^2)(1-r)L,\\
 \Delta/b={}&(1+x)B+2x(1+x)r+(1+x)(1+k^2)\\
 &+(x^2-k^2x)r^2/b+2x(1+k^2)r/b+(1+k^2)/b.
\end{align*}
Consequently $xN/b=x^2N/B$ and $\Delta/b$ are positive bounded
units, with respective residues
\[
 H_0+k_0^2,\qquad B_0+1+k_0^2.
\]
Here $H,L\le B$ controls the divided corrections, and $x/B\to0$.

The exact physical condition $n\ge0$ is
\[
 r(\Delta/b)\le(1+r/b)A_cS.
\]
It implies
\[
 r_0\le\frac{1+L_0}{1+k_0^2+B_0},\qquad
 1-r_0\ge\frac{k_0^2+H_0}{1+k_0^2+B_0}>0.
\]
Thus $1-r$ is a positive bounded unit without any lower bound on $r$.
The original formulas now give the exact comparabilities
\[
 m\asymp b^{-2},\qquad C\asymp L^{-1},\qquad A\asymp H^{-1},
 \qquad D\asymp B/j,\qquad\alpha\asymp B/(1-H).
 \tag{SR1}
\]
Every omitted factor in $\asymp$ is a positive bounded unit with
bounded reciprocal. For instance,
\[
 CL=\frac{A_c(xN/b)}{2\Delta/b},\qquad
 AH=\frac{k^2A_c(1+r/b)S}{2\Delta/b}
\]
are positive units; the other identities follow using $g/b\to1$
and $(b-k^2r)/b\to1$.

There are exactly two unbounded-$b$ scale leaves:
\begin{enumerate}
\item If $B=xb$ is infinitesimal, so are $H,L$, while $j\to1$
and $1-H\to1$. Thus $C,A$ are unbounded and $m,D,\alpha$ are
weak, with $m\asymp b^{-2}$ and $D\asymp\alpha\asymp xb$.
This is the one-weak-complementary-event coefficient domain, with weak
$\alpha$ and strong $C$ in this actual-angle orientation.
\item If $B$ is a positive bounded unit, then $C,A,D,\alpha$
are all strong and $m$ is weak. The new strong-complement pitchfork
applies to this entire leaf at all upper event scales. In particular
it includes $C$ unbounded when $L$ is infinitesimal and $\alpha$
unbounded when $1-H$ is infinitesimal.
\end{enumerate}
Unbounded $B$ is impossible by the original chart inequalities.
The homogeneous all-tests theorem of Section~\ref{pab6222b899ef:sec:c2-unbounded-b-small-B}
closes the remaining infinitesimal-$B$ leaf without assuming a bounded
actual state ratio. Thus both unbounded-$b$ leaves, as well as all
bounded-$b$ openings, are now covered. Literal $n=0$ is preserved.

\subsection{Homogeneous completion of the scalar graph}
\label{component:pbb29c73cc124}

\label{pbb29c73cc124:sec:unbounded-b-homogeneous-scalar}

This lemma is a scalar lifting statement.  It neither supplies the full
state normals nor proves positivity of its coefficient Gram by itself.
Its purpose is to avoid any bound on the contact state ratio.

Let $B,\sigma>0$ be pure coefficient infinitesimals, let $k>0$ be a
pure coefficient bounded unit, and put $x=\sigma B$, $\xi=u^2$.
Suppose
\[
 f=a(\xi,v)W^2+2u d(\xi,v)WZ+c(\xi,v)Z^2,
 \quad
 a=\frac\xi B Q+\sigma^2 a_0,\qquad
 c=B e+\xi b,                                      \tag{HL1}
\]
where $Q,a_0,d,e,b$ have bounded polynomial coefficients.  Bounded
pure testing unit denominators may be cleared beforehand.
Assume the value and all first derivatives of $f$ vanish at
\[
 u=\sqrt{x},\quad v=k,\quad
 W=\frac1k\sqrt{B/\sigma},\quad Z=1.                \tag{HL2}
\]
No bound on $B/\sigma$ is assumed.

Define the homogeneous roots
\[
 R=\frac{uW-(B/k)Z}{\sqrt B},\qquad
 S=uZ-\sigma kW,\qquad
 U=\sqrt B(v-k)Z,\qquad T=u(v-k)Z.                 \tag{HL3}
\]
They satisfy the exact conductors
\[
 (\xi-x)W=u\sqrt B R+(B/k)S,
 \qquad
 (\xi-x)Z=uS+\sigma k\sqrt B R.                    \tag{HL4}
\]
Then $f$ has an exact symmetric Gram on $(R,S,U,T)$ whose entries
are bounded polynomials.  In particular no inverse of $\sigma$,
$B$, a testing infinitesimal, or a mixed expression is required in
these Gram entries.  The roots themselves need not be bounded.

\paragraph{First set $v=k$.}
In this paragraph every scalar coefficient is evaluated at $v=k$.
State stationarity and $u$ stationarity give
\begin{gather*}
 a_c=-\sigma k d_c,\qquad c_c=-B d_c/k,\\
 B(a_\xi)_c+kd_c+2kx(d_\xi)_c
                      +\sigma k^2(c_\xi)_c=0,     \tag{HL5}
\end{gather*}
where $c$ on a subscript means evaluation at $\xi=x$.
These are identities in the pure coefficient field; deriving them
does not insert an inverse into a tensor Gram.

Set, by monic polynomial division in $\xi$,
\[
 s=\frac{c+Bd/k}{\xi-x},\qquad
 q=-kd-\sigma k^2s,\qquad
 h_2=\frac{Ba+k\xi d+\sigma k^2c+\sigma Bkd}
                    {(\xi-x)^2}.                 \tag{HL6}
\]
All three quotients have bounded polynomial coefficients.  The first
numerator vanishes at $x$ by (HL5), and the last numerator and its
$\xi$ derivative vanish there by the same three identities.
Moreover $Ba=\xi Q+B\sigma^2a_0$ is bounded.  Division by a monic
polynomial preserves bounded coefficients and is not evaluation of
an inverse of $\xi-x$.

Direct multiplication gives the exact identity
\[
 f(\xi,k)=qR^2+sS^2
                  +h_2\left(uR+\frac{\sqrt B}{k}S\right)^2.
                                                               \tag{HL7}
\]
For later residue computations, if $E_0=e+d/k$, then
\[
 s=b+B\frac{E_0(\xi)-E_0(x)}{\xi-x}
       +x\frac{b(\xi)-b(x)}{\xi-x}.                \tag{HL8}
\]
Thus $s_0=b_0$, $q_0=-k_0d_0$, and the $h_2$ correction has zero
residue in every Gram entry.  Here subscripts $0$ denote residues,
not evaluation at the contact polynomial variable.

\paragraph{The complete angular coquotient.}
Write $\tau=v-k$ and, for each of $Q,a_0,d,e,b$, form its bounded
monic difference in $v$, denoted respectively by
$Q_1,a_{01},d_1,e_1,b_1$:
\[
 Q(\xi,v)-Q(\xi,k)=\tau Q_1(\xi,v),
\]
and similarly for the other four functions.  Put
\[
 H_1=Q_1/k^2+2d_1/k+e_1,\qquad
 L_1=a_{01}/k^2+b_1,\qquad J_1=H_1+\sigma L_1.
\]
Angular stationarity at (HL2) gives exactly $J_1(x,k)=0$.
Consequently the successive monic differences
\[
 J_1=(\xi-x)J_\xi+\tau J_v,
 \quad
 J_\xi=\frac{J_1(\xi,v)-J_1(x,v)}{\xi-x},\qquad
 J_v=\frac{J_1(x,v)}{v-k}                            \tag{HL9}
\]
are bounded polynomials.  Let $D_1=L_1+BJ_\xi$.  The full exact
angular correction is
\begin{align*}
 f(\xi,v)-f(\xi,k)={}&
 \tau Q_1R^2+\frac{\tau a_{01}}{k^2}S^2\\
 &+\left\{2(Q_1/k+d_1)+\sigma kD_1\right\}RU\\
 &+\left\{D_1-2a_{01}/k^2\right\}ST+J_vU^2.       \tag{HL10}
\end{align*}
To check this without any division, use
$uW/\sqrt B=R+\sqrt B Z/k$ and
$\sigma W=(uZ-S)/k$.  The only term not already among the displayed
root products becomes
\[
 \tau Z^2\{BH_1+\xi L_1\}
 =D_1\tau(\xi-x)Z^2+BJ_v\tau^2Z^2
 =D_1(ST+\sigma kRU)+J_vU^2,
\]
by (HL4).  Equations (HL7) and (HL10) prove the asserted homogeneous
Gram lift for every state $W,Z$ simultaneously.  There is no state
normalization and no discarded endpoint $Z=0$.

The usual exact zero, now with the new source scale, is
\[
 T^2-\sigma U^2-ST(v-k)-\sigma kRU(v-k)=0.          \tag{HL11}
\]
This is the appropriate identity for redistributing the angular
diagonal.  Its $U^2$ coefficient is $\sigma$, which is infinitesimal;
replacing it by $1$ would erase the smaller radial scale.  Any
positivity argument must use the actual residue coefficients of
(HL7)--(HL10), rather than importing the bounded-state Schur test
without this scale correction.

\paragraph{Strict scalar corollary for the small-$B$ source.}
Suppose additionally that the coefficientwise leading functions on
the finite $u=0$ axis are
\[
 Q_0(v)=\frac{2K^2(1+v^2)}{1+K^2},\qquad
 d_0(v)=-2v,\qquad
 e_0(v)=\frac{2(v^2/K^2+K^2)}{1+K^2},
 \qquad b_0(K)>0,
\]
where $K=k_0$.  Then the $(R,U)$ residue block produced above is
\[
 2K^2R^2+\frac{4(K^2-1)}{1+K^2}RU+\frac2{K^2}U^2,
\]
with Schur margin $8/(1+K^2)^2>0$.  The $S^2$ coefficient has
residue $b_0(K)>0$, the $RS$ residue is zero, and the only other
possible residue entry is a bounded $ST$ coefficient, say $\ell$.
Add a positive real multiple $\delta$ of (HL11), with
$\delta>\ell^2/(4b_0(K))$.  Its correction to the $U^2$ entry has
zero residue because $\sigma$ is infinitesimal, whereas it supplies
the positive $T^2$ diagonal $\delta$.  Both residue blocks are then
strictly positive.  This proves the scalar strict bounded Gram for
all homogeneous states in the matched-angle chart.  Full normals
and the other testing charts remain separate requirements.

\subsection{The case $b$ unbounded and $xb$ infinitesimal}
\label{component:pab6222b899ef}

\label{pab6222b899ef:sec:c2-unbounded-b-small-B}

\begin{theorem}[The remaining unbounded-state rational C2 source]
Use the exact original rational C2 source.  Suppose $x>0$ is
infinitesimal, $k$ is a positive bounded unit, $b$ is positive
unbounded, $B=xb$ is infinitesimal, $0<r<1$, $0<h<b$, and all
original positive-chart inequalities hold.  Assume $n\geq0$ and
physicality of the source.  Then the full original parent has tensor
SOS for every independent real or complex projective state and
every independent projective pair of measurement angles.
\end{theorem}

The conclusion includes $n=0$, arbitrarily large $n$, arbitrary
complementary openings, and every value of the pure coefficient
ratio $B/\sigma$, where $\sigma=1/b$.  In particular, the actual
contact ratio $W/Z=k^{-1}\sqrt{B/\sigma}$ need not be bounded.
This proof uses the homogeneous scalar lemma in
Section~\ref{pbb29c73cc124:sec:unbounded-b-homogeneous-scalar}; it restores both
full state normals and exhausts the other testing charts.

\paragraph{Exact source scales.}
Put
\[
 \sigma=1/b,\quad x=\sigma B,\quad \eta=h/b,\quad
 y=B\eta,\quad z=B(1-\eta),\quad
 A_c=1+\sigma Br,\quad S_c=A_c+z.
\]
The positive chart gives $0<\eta<1$; both $y,z$ are infinitesimal.
Let $\widehat\Delta=\sigma\Delta$,
$\widehat N=\sigma^2 B N_{\rm src}$,
$\widehat g=\sigma g=1+\sigma(r-1+B\eta)$ and
$j=A_c-z/k^2$.  The cancelled exact formulas are
\begin{align*}
 nr&=\frac{(1+\sigma r)A_cS_c-r\widehat\Delta}
                       {2\widehat\Delta},\\
 m&=\frac{\sigma^2(1+k^2)(1-r)A_cS_c}
                       {2\widehat g\widehat\Delta},\\
 C&=\frac{A_c\widehat N}{2z\widehat\Delta},\qquad
 A=\frac{k^2A_c(1+\sigma r)S_c}{2y\widehat\Delta},\\
 D&=\frac{B(1+\sigma B)(1+\sigma r)(1-\sigma k^2r)S_c}
                       {2k^2j\widehat\Delta},\\
 \alpha&=\frac{B(1+\sigma r)\widehat N}
                       {2(1-y)\widehat\Delta}.
                                                               \tag{UB1}
\end{align*}
Here $\widehat\Delta,\widehat N,\widehat g,j,1-y$ are bounded
positive units.  Write $K=k_0>0$ and $D_K=1+K^2$.  Their relevant
residues imply
\begin{gather*}
 (nr)_0=\tfrac12(D_K^{-1}-r_0),\quad
 (m/\sigma^2)_0=(1-r_0)/2,\\
 (Cz)_0=(Ay)_0=K^2/(2D_K),\quad
 (D/B)_0=1/(2K^2D_K),\quad
 (\alpha/B)_0=K^2/(2D_K).                            \tag{UB2}
\end{gather*}
Since $n\geq0$, $r_0\leq D_K^{-1}<1$, so $1-r$ is a unit.
Set $N=4n+2\geq2$.  Exactly,
\[
 Nr=\frac{2(1+\sigma r)A_cS_c}{\widehat\Delta},      \tag{UB3}
\]
which is a bounded positive unit.  Thus $nr,nr^2$ are bounded even
when $n$ is unbounded, while $1/N$ and $1/\sqrt N$ are bounded.

All opening cancellations needed below are uniform:
\[
 BC(1-\eta)=Cz,\quad BA\eta=Ay,\quad
 BC(1-\eta)^2=Cz(1-\eta),\quad BA\eta^2=Ay\eta.
                                                               \tag{UB4}
\]
The quantity $H_B=B(C+A)$ is bounded below by a positive real;
its reciprocal and reciprocal square root are bounded.  No
assumption that $\eta$ or $1-\eta$ is a unit is made.

\paragraph{The whole finite $u=0$ axis.}
Use the original parent and the exact-contact graph
\[
 X_g=-ruZ,\qquad Y_g=\eta uW,
 \qquad E_X=\sqrt N(X+ruZ),\quad
 E_Y=\sqrt{C+A}(Y-\eta uW).                          \tag{UB5}
\]
The actual contact is
\[
 u=\sqrt{\sigma B},\quad v=k,\quad
 W=k^{-1}\sqrt{B/\sigma},\quad Z=1.
\]
Let $\xi=u^2$ and $f=F|_g$.  Exact polynomial coquotients give
\[
 f=a(\xi,v)W^2+2u d(\xi,v)WZ+c(\xi,v)Z^2,
 \quad a=\frac\xi B Q+\sigma^2a_0,\quad c=B e+\xi b.
                                                               \tag{UB6}
\]
Every coefficient of $Q,a_0,d,e,b$ is bounded by (UB3)--(UB4).
For example $a_0=4(m/\sigma^2)v^2$,
$e=4(D/B)v^2+4\alpha/B$, and
$Q=B(a-a|_{u=0})/u^2$.  These are monic polynomial divisions.
The full leading functions at $u=0$ are
\begin{align*}
 q(v)&=2K^2(1+v^2)/D_K,& d_0(v)&=-2v,\\
 e_0(v)&=2(v^2/K^2+K^2)/D_K,&
 (a_0)_0(v)&=2(1-r_0)v^2,\\
 b_0(v)&=2r_0(1+v^2)/D_K-4r_0+2.                  \tag{UB7}
\end{align*}

For remote tests use
\[
 P=uW/\sqrt B,\quad B_u=uZ,\quad
 V=\sigma vW,\quad E=\sqrt B Z.
\]
Then (UB6) has the exact bounded Gram
\[
 QP^2+b B_u^2+(a_0/v^2)V^2+eE^2+2dPE,             \tag{UB8}
\]
where $a_0/v^2$ is a polynomial coquotient.  The $(P,E)$ residue
block is strictly positive whenever $v_0^2\ne K^2$, since
\[
 q(v)e_0(v)-4v^2=4(v^2-K^2)^2/D_K^2.               \tag{UB9}
\]
This is a global identity for every finite residue $v$, not a
local-contact positivity inference.

\paragraph{Both complete state normals.}
The exact quadratic normal expansion is
\[
 F=f+E_X L_X+E_Y L_Y+H_{XX}E_X^2
                       +2H_{XY}E_XE_Y+H_{YY}E_Y^2.
                                                               \tag{UB10}
\]
Its entire hard block is bounded with residue
$\operatorname{diag}(D_{v,0},4D_{v,0})$.
Indeed the $X$ pivot is $ND_uD_v+O(1)$ before division by $N$,
with the remaining terms either infinitesimal in $u$ or of order
$m$; the $Y$ pivot from the strong positive square terms is exactly
$4D_v(1+[A/(C+A)]u^2)$.  All mixed entries have zero residue and
contain no $C,A$ contribution.

The complete unnormalized $X$ derivative has the form
\[
 F_X|_g=uZ A_X+W(\xi B_1+\sigma^2 B_0),             \tag{UB11}
\]
with bounded polynomial coefficients, where
\begin{align*}
 B_0&=-8(m/\sigma^2)v(1+\eta\xi),\\
 B_1&=v(8D+4)(1-\eta),\\
 A_X&=-8nrD_uD_v+8m(1-r)
       -(8D+4)v^2(1+r\xi)+4D_v(1-r).
\end{align*}
Thus in the remote roots its $P,V$ coefficients have zero residue;
its only possible nonzero residue coupling is to $B_u$.
The exact residue Schur expression for the $(B_u,E_X)$ block is
\[
 (N r^2D_v-4r+2)
       -\frac{(4-2NrD_v)^2}{4ND_v}
                 =2-\frac4{ND_v}.                  \tag{UB12}
\]
It is strictly positive for every finite nonzero residue $v$,
because $N\geq2$.  At $v_0=0$ it is strict if $n$ is unbounded
or has positive residue.  The sole zero case $n_0=v_0=0$ is
treated below.

For the other complete first normal,
\[
 L_Y=\frac1{\sqrt{H_B}}
 \{P(B_2+B e_w)+B_u(\sqrt B\,u e_z)+E e_*\},       \tag{UB13}
\]
where $e_w,e_z$ are bounded,
\[
 B_2=8D_v\{(BA)\eta D_u-(BC)(1-\eta)\},\qquad
 e_*=8v(D-\alpha).
\]
All coefficients have zero residue, by (UB4), $B\to0$, and
$u\to0$.  This verifies both full normals and proves the full
parent at every remote finite-axis test covered by (UB12).

\paragraph{Matched residues without any state-ratio restriction.}
For $v_0=K$, apply the homogeneous scalar lemma to (UB6).  Its
contact and full stationarity hypotheses hold exactly because the
graph (UB5) passes through the original stationary contact.
Use the homogeneous roots
\[
 R=\frac{uW-(B/k)Z}{\sqrt B},\quad
 S=uZ-\sigma kW,\quad
 U=\sqrt B(v-k)Z,\quad T=u(v-k)Z.                  \tag{UB14}
\]
They satisfy
\[
 (\xi-x)W=u\sqrt B R+(B/k)S,\qquad
 (\xi-x)Z=uS+\sigma k\sqrt B R.                    \tag{UB15}
\]
The lemma constructs an exact bounded scalar Gram, valid for
every $W,Z$ simultaneously.  Its $(R,U)$ residue is
\[
 2K^2R^2+\frac{4(K^2-1)}{D_K}RU+\frac2{K^2}U^2,
\]
whose Schur margin is $8/D_K^2>0$.  Its radial $S^2$ residue is
$s_0=2(1-r_0)>0$; its only other possible residue coupling is
$\ell ST$, with a bounded real residue $\ell$.

To lift the full $X$ normal, set
\[
 L=A_X+(B/k)B_1+(\sigma/k)B_0.
\]
Equations (UB11) and (UB15) give exactly
\[
 F_X|_g=u\sqrt B B_1R-(\sigma/k)B_0S+uZL.
\]
At the source contact $L(x,k)=0$, because $F_X=0$ there.
Write its bounded monic differences as
$L=(\xi-x)L_\xi+(v-k)L_v$.  Then
\[
 F_X|_g=
 u\sqrt B(B_1+\sigma kL_\xi)R
 +\{\xi L_\xi-(\sigma/k)B_0\}S+L_vT.             \tag{UB16}
\]
After division by $\sqrt N$, all coefficients are bounded.
The $R,S$ coefficients have zero residue; the $T$ coefficient
may have a finite residue $\gamma$ (it is
$-8K/(D_K\sqrt{N_0})$ when $N$ is bounded).
This nonzero full-normal coupling is retained, not silently
discarded.

For (UB13), put $A_Y=B_2+B e_w$ and
$L_Y'=A_Y/k+\xi e_z+e_*$.  Its numerator is exactly
$RA_Y+\sqrt B ZL_Y'$.  Exact normal contact gives
$L_Y'(x,k)=0$.  Monic differences and (UB15) yield bounded
coefficients on $R,S,U$, all of zero residue: the $R,S$ corrections
carry factors $\sigma B$ or $u\sqrt B$, and the $U$ coefficient
is the angular coquotient of an identically zero leading function.
Division by $\sqrt{H_B}$ preserves these bounds.

Finally add $\delta$ times the exact zero
\[
       T^2-\sigma U^2-ST(v-k)-\sigma kRU(v-k)=0.    \tag{UB17}
\]
Choose a real $\delta>\ell^2/(4s_0)+\gamma^2/(4D_K)$.
This supplies a $T^2$ residue $\delta$ at zero residue cost to
the $U^2$ coefficient, since $\sigma$ is infinitesimal.  The block
on $(S,E_X,T)$ has strictly positive Schur complement; the $(R,U)$
block and $E_Y$ diagonal are already strict.  Thus the complete
six-root Gram on $(R,S,U,T,E_X,E_Y)$ is strictly positive in
residue.  The calculation is homogeneous in all four state
coordinates and needs no bound on $W/Z$, no comparison of $B$ with
$\sigma$, and no state endpoint exclusion.  The exact simultaneous
sign copy treats $v_0=-K$.

\paragraph{The double-zero test, including literal $n=0$.}
Suppose now $n_0=v_0=0$.  Use the pure coefficient quantities
\[
 H_c=4n+4m+2,\quad \zeta=(4m+2)/H_c,\quad
 \xi_c=4m/H_c,\quad \eta_c=C/(C+A),
\]
and graph $X_g=\xi_cvW-\zeta uZ$, $Y_g=\eta_cuW$.
Normalize its $X$ normal by $\sqrt{H_c}$ and its $Y$ normal by
$\sqrt{C+A}$.  Put
\[
 \epsilon=4\alpha+4n(4m+2)/H_c>0.
\]
Then $\epsilon$ is infinitesimal, $n/\epsilon$ and $B/\epsilon$
are bounded, and the harmonic product $BCA/(C+A)$ is a unit.
Use the five scalar roots
\[
 P=uW/\sqrt B,\quad V=\sigma vW,\quad E=\sqrt B Z,
 \quad B_\epsilon=\sqrt\epsilon\,uZ,\quad T_0=uvZ.
\]
The exact inherited identity, with its original coefficients,
is
\[
 f_{zz}=4\alpha+\epsilon u^2+4n\zeta^2u^4
                       +4Dv^2+u^2v^2 B_b,
 \qquad (B_b)_0=2.
\]
The $W^2$ coefficient at $u=0$ is exactly
\[
 \frac{4m(2n+1)}{2n+2m+1}v^2
              +(4n+2)\xi_c^2v^4.
\]
Therefore its coefficient on $V^2$ has positive residue
$4(m/\sigma^2)_0=2(1-r_0)$.  The scalar residue on the five roots
is strictly diagonal, with entries
\[
 \left(\frac{2K^2}{D_K},\ 2(1-r_0),\
        \frac{2K^2}{D_K},\ 1,\ 2\right).
\]
The cross term is allocated using $uvWZ=vPE$, so it has zero residue.

Both full normals are restored, including when $n=0$ exactly.
At $v=0$ the $X$ derivative is $-8n\zeta u^3Z$, whose coefficient
on $B_\epsilon$ is infinitesimal since $n/\epsilon$ is bounded.
Its remaining $Z$ coefficient is $uv^2$ times a bounded function,
so it has zero residue on $T_0$.  Its $W$ coefficient is a sum of
$u^2vW$ times a bounded function and $mvW$ times a bounded
function divisible by $v^2$ at $u=0$.  Allocate the former to
$P$ with coefficient $u\sqrt B\,v$, and the latter to $V$ with
coefficient $m/\sigma$; both have zero residue.  In particular one
must not divide the former by $\sigma$.
The $Y$ derivative is the inherited harmonic-opening identity:
$B_2=8D_v(BA)\eta_cu^2$, with bounded residual coquotients.
Its coefficients on $P,T_0,E$ all have zero residue.  The hard
residue is $\operatorname{diag}(1,4)$.  This proves the full-parent
corner without deleting a state coordinate.

\paragraph{All remaining projective tests.}
For the other commuting charts use the four original outer-core
identities with $\rho$ replaced by $\sqrt B$, and cap $C,A$ at
$s/B$, where $s=\min(1,BC,BA)$ is a positive bounded unit.
Retain both original nonnegative excess positive square terms and the entire
original $n$ positive square term.  No positive lower bound on $m$ is needed:
the first two outer pivots are $4m_0+2=2$, while the other two
charts retain the whole $m$ positive square term and do not use it for strictness.
The reciprocal-$v$ chart uses $(D/B)_0>0$ exactly as before.
Thus the finite unit $u$ axis, both reciprocal axes, all double
corners, and literal projective endpoints have strict bounded
Grams.  The inherited physical noncommuting-residue theorem covers
all remaining testing pairs.

The finite-axis identities are homogeneous in the state and include
unbounded actual state ratios. Every mixed quotient used above is a
monic polynomial quotient; coefficient inverses and radicals remain in
one labelled field. Restore the retained squares and apply the strict
bounded-Gram lemma to each real and imaginary state part.

\section{One small actual angle: strong complementary events}
\subsection{Two unbounded retained events}
\label{component:p0c7a15f800d9}

\label{p0c7a15f800d9:sec:c2-two-large}

\begin{theorem}[The $n,D$-unbounded C2 pitchfork]
\label{p0c7a15f800d9:thm:c2-two-large}
Let the physical canonical C2 source over a real closed coefficient
field have $\lambda=\beta=0$, with $d,C,\alpha$ positive bounded
units, $A\geq0$ bounded, and $n,m,D\geq0$ arbitrary. Suppose there is
an actual finite noncommuting stationary zero with $u_c=k>0$
infinitesimal and $v_c=a>0$ a bounded unit. Then the entire original
parent is tensor SOS at every independent projective measurement
pair and every real or complex projective state.
There is no bound on $n,D$, or on their ratio. In fact physicality and
contact force $A$ to be a positive unit, $D$ to be bounded below by
a positive real, and $m/k^4$ to be a positive bounded unit.
\end{theorem}

This extends the preceding pitchfork by removing the upper bound on
$D$. A bound on $m$ need not be assumed: it follows from physicality.
The original $D$ positive square term supplies a state normal, so an unbounded
weight is not put into a bounded Gram coefficient.

Write the unchanged parent as
\[
\begin{split}
F={}&dH+4nD_uD_vX^2+4mR_2^2+4CD_vr^2
       +4AD_uD_vY^2+4DJ^2+4\alpha D_up^2,\\
H={}&2\{J^2+D_up^2+D_vB^2-D_uD_vZ^2\},\\
R_2={}&vW-X+uvY-uZ,\quad r=uW-Y,\quad
J=uW+uvX-Y-vZ,\\
p={}&vY-Z,\quad B=X+uZ,\quad D_u=1+u^2,\quad D_v=1+v^2.
\end{split}                                                    \tag{TD1}
\]

\paragraph{Entry and bounded contact coordinates.}
At $u=0,v=a$, the state block on $(X,Y,Z)$ of
$dH+4CD_vr^2+4\alpha D_up^2$ has positive definite residue.
Its $W$ column is $O(k)$ and its $W^2$ entry is $O(k^2)$.
Completing this ordinary bounded hard block gives, for positive
real constants $b,B_0$, the coefficient-field inequality
\[
 dH+4CD_vr^2+4\alpha D_up^2
 \geq b(X^2+Y^2+Z^2)-B_0 k^2W^2
\]
at the actual angles. All omitted positive square terms are nonnegative.
Thus an actual zero has $W\ne0$ and, on setting $W=1$,
$X,Y,Z=O(k)$. Since $R_2=a+O(k)$ is a unit, the same inequality
and its retained $m$ positive square term give $m=O(k^2)$.

The state equation $F_Y=0$ gives $DJ=O(k)$: every other term is
$O(k)$, using bounded $A,C,\alpha,d$, $m=O(k^2)$ and the bounds
just obtained. The equation $F_X=0$ then gives
$(d+2n)X=O(k^2)$, since the $J$ contribution is $O(kDJ)$ and the
$m$ contribution is $O(m)$, while $Z=O(k)$. Consequently, with
\[
 N=d+2n,\qquad x=X_c/k^2,\quad z=Z_c/k,\quad j=J_c/k,
\]
we have
\[
 Nx,\ z,\ j,\ Dj\quad\hbox{bounded},\qquad N^{-1}\ \hbox{bounded}.
                                                               \tag{TD2}
\]
The bound on $j$ follows directly from $X,Y,Z=O(k)$, even before a
lower bound for $D$ has been obtained. No inverse of $D$ has yet
been used.

For the scale entry, put $U=k^2,V=a^2$, $L=d+2D$,
$\theta=n/N$ and use
\[
\begin{gathered}
t_2=d+2(\alpha+m+A),\quad t_3=d+2(C+D),\quad
 t_4=d+2(\alpha+m),\\
t_6=d+2(\alpha+C+A),\quad t_8=d+2(C+A+D).
\end{gathered}
\]
The original Laurent-determinant stationary equations, divided only
by $N$, are
\[
\begin{split}
U^2\{4\theta A t_3t_4+4CA(1+2D/N)t_4V\}
 &=mV\{4\alpha t_8+4Dt_6V\},\\
4\alpha A(1+2m/N)t_3+4\theta A t_3t_4U
 &=\{4mDt_6/U+4CDt_2\}V^2.
\end{split}                       \tag{TD3}
\]
Here $t_2,t_4,t_6,t_3/L,t_8/L$ are positive bounded units.
After division by $L$, the braces in the first equation are
$A$ times a positive unit on the left and a positive unit on the
right. To check the lower bound uniformly when both $n,D$ diverge,
observe
\[
 \theta+\frac{1+2D/N}{L}
 =\frac{n+1}{N}+\frac{2n}{NL},                    \tag{TD4}
\]
which is a positive bounded unit; all its neighboring factors in
(TD3) have positive upper and lower real bounds. Thus
$m/U^2$ is $A$ times a positive unit. The second equation, again
divided by $L$, has the form
\[
 A\cdot(\hbox{positive unit})+O(AU)
    =(D/L)\cdot(\hbox{positive unit})+O(AU).       \tag{TD5}
\]
If $A$ were infinitesimal, then $D/L$ and hence $D$ would be
infinitesimal, as would $m$. This is incompatible with physicality:
on the test state $(W,X,Y,Z)=(1,0,u,zu)$ the $n$ positive square term vanishes
exactly, and at the residue source $m=A=D=0$ choose
\[
 z=\frac{v(2\bar\alpha+\bar d)}{2\bar\alpha}.
\]
The coefficient of $u^2$ in the full polynomial is then
$-\bar d v^2(2\bar\alpha+\bar d)/\bar\alpha<0$ for real $v\ne0$.
The bounded residue polynomial cannot be nonnegative for all real
small $u$. Therefore $A$ is a unit. Equations (TD3)--(TD5) now prove
that $D/L$ and $m/k^4$ are positive bounded units. In particular
$D>0$ and $D^{-1}$ is bounded. Set $\gamma=4ma^2/k^4$.

\paragraph{A graph that keeps both large positive square terms bounded.}
For arbitrary independent testing angles use the exact graph
\[
 X_g=xu^2W,\qquad Z_g=zuW,\qquad J_g=juW,
\]
which fixes
\[
 Y_g=u(1-j-vz+xu^2v)W.                            \tag{TD6}
\]
It matches the entire actual contact, including its $Y$ coordinate.
The three exact state normals are
\[
 E_X=\sqrt N(X-xu^2W),\qquad
 E_J=\sqrt D(J-juW),\qquad E_Z=Z-zuW.             \tag{TD7}
\]
The inverse uses only pure coefficient inverses:
\[
 X=xu^2W+E_X/\sqrt N,\quad Z=zuW+E_Z,\quad
 Y=u(1-j)W+uvX-vZ-E_J/\sqrt D.                   \tag{TD8}
\]
The $n$ positive square term on the graph is
$4nx^2u^4D_uD_vW^2$, and the $D$ positive square term is
$4Dj^2u^2W^2$. Both have bounded coefficients by (TD2).
Their linear coefficients are bounded too, since $nx$ and
$\sqrt D\,j=(Dj)/\sqrt D$ are bounded. The exact quadratic
normal block is bounded: its only large positive square terms become
$4(n/N)D_uD_vE_X^2$ and $4E_J^2$.

Set $W=1$, $K=k^2$, and
\[
 R=u^2-(K/a)v,\qquad T=u(v-a).
\]
The entire graph scalar has the exact shape
\[
 F|_g=4mv^2+u^2H_0(u^2,v),\qquad H_0\in B_c[u^2,v].             \tag{TD9}
\]
The three contact scalar jets give exactly the same monic division
as for the preceding pitchfork:
\[
 F|_g=\gamma R^2+u^2\{(u^2-K)^2H_{20}
                   +(u^2-K)(v-a)H_{11}+(v-a)^2H_{02}\},         \tag{TD10}
\]
with bounded polynomial coquotients. After
$u^2-K=R+(K/a)(v-a)$ this is a bounded Gram on $(R,T)$.
Its $R$ residue pivot is $\bar\gamma>0$.

The full first normals, before the indicated scale divisions, have
exact shapes
\[
 (F_X+uvF_Y)|_g=-8mv+u^2L_X(u^2,v),\quad
 (-F_Y)|_g=uL_J(u^2,v),\quad
 (F_Z-vF_Y)|_g=uL_Z(u^2,v).                       \tag{TD11}
\]
The coefficient of $E_X$ is the first expression divided by
$\sqrt N$; that of $E_J$ is the second divided by $\sqrt D$.
These scaled polynomials have bounded coefficients, as shown above.
Contact makes $L_X(K,a)=8ma/K$, and makes the scaled $J,Z$ functions
vanish at $(K,a)$. Monic division by $u^2-K,v-a$ therefore puts
every first normal in $(R,T)$ with bounded coquotients.
Every coefficient coupling $E_X$ to $R,T$ is infinitesimal; the
$R$ coefficients in the $E_J,E_Z$ rows are infinitesimal.
For example, the $X$ row is exactly
\[
\frac1{\sqrt N}\left[
 (8ma/K+u^2L_{X1})R
 +u\{(K/a)L_{X1}+L_{X2}\}T\right],              \tag{TD12}
\]
and the other two rows have form
$uL_{i1}R+\{(K/a)L_{i1}+L_{i2}\}T$ after their scale division.
State degree is two, so there are no omitted higher-normal terms.
Equations (TD7)--(TD12) give an exact bounded full-parent Gram on
$(E_X,E_J,E_Z,R,T)$.

\paragraph{The residue Gram is strict.}
The $X$ pivot has residue $2(1+\bar a^2)$ and the $R$ pivot has
residue $\gamma$. For the remaining hard pair the unscaled state
form is
\[
 B(v)[Y,Z]=\{4(C+A)+2d\}(1+v^2)Y^2
           +4D(Y+vZ)^2+4\alpha(vY-Z)^2.           \tag{TD13}
\]
In the scaled coordinates $Y=-vE_Z-E_J/\sqrt D$, $Z=E_Z$ its
coefficients are bounded and its residue is positive definite.
If $D$ is bounded, this follows from the invertible residue change
and the positive form (TD13). If $D$ is unbounded, its residue is
\[
 4E_J^2+\bigl(\{4(\bar C+\bar A)+2\bar d\}(1+v^2)v^2
                  +4\bar\alpha(1+v^2)^2\bigr)E_Z^2\succ0.
\]
The exact minimized boundary coefficient of $u^2$, before residue,
is
\[
 f(v)=\frac{8(1+v^2)}{\det B(v)/L}
       \left(\frac{Q_{10}}L+\frac{Q_0}L v^2+\frac{Q_{12}}L v^4\right),
                                                               \tag{TD14}
\]
where
\[
 Q_{10}=4\alpha A[d+2(C+D)],\qquad
 Q_{12}=4CD[d+2(\alpha+A)].
\]
The middle coefficient $Q_0$ is defined by the exact Schur numerator;
its explicit value is unnecessary. All quantities in (TD14) are
bounded, $\det B/L$ is a positive unit for bounded real $v$, and
$Q_{10}/L,Q_{12}/L$ are positive units. One may verify boundedness
of $Q_0/L$ directly from the determinant expansion, which is affine
in $D$. This is only a residue calculation of a real hard Schur
complement, not inversion of a mixed tensor element.

The full contact normal and scalar equations imply that the residue
of (TD6) at $v=\bar a$ minimizes the hard form, and that
$\bar f(\bar a)=\bar f'(\bar a)=0$. Thus the real quartic bracket
in (TD14) equals
\[
 \overline{Q_{12}/L}\,(v^2-\bar a^2)^2.
\]
Its second derivative at $\bar a$ is strictly positive. This is
exactly the Schur complement on $T$ of the remaining residue block
$(E_J,E_Z,T)$. Hence the whole five-root residue Gram is positive
definite. The bounded strict-residue Gram lemma proves the local
certificate at $\bar u=0,\bar v=\bar a$. The Bob sign copy covers
$\bar v=-\bar a$.

\paragraph{All other independent tests.}
For $u$ infinitesimal and $v$ bounded with
$\bar v\ne\pm\bar a$, retain the original $m,n$ positive square terms. Set
$K_J=\sqrt D J$ and substitute
\[
 Y=uW+uvX-vZ-K_J/\sqrt D.
\]
The exact remainder is a bounded quadratic Gram on
$(X,Z,K_J,uW)$. Its hard residue is the block just considered with
$X$ pivot $2\bar d(1+\bar v^2)$; its soft Schur complement is
$\bar f(\bar v)>0$. This includes $v=0$ and all infinitesimal $v$.
For $u,q=1/v$ infinitesimal use the homogeneous parent $q^2F(u,1/q)$.
Here $J_h=quW+uX-qY-Z$. Retain $m,n$ and substitute
$Z=quW+uX-qY-K_J/\sqrt D$ with $K_J=\sqrt D J_h$.
The remainder is a bounded Gram on $(X,Y,K_J,uW)$. Its soft residue
is $\lim_{q\to0}q^2\bar f(1/q)>0$, since
$\overline{Q_{12}/L}>0$ and
$\overline{\det B_\infty/L}>0$.
These are polynomial homogeneous expressions, so literal projective
endpoints are included.

At every other commuting test the capping proof from the preceding
all-tests corollary applies verbatim: choose
$\mu=\min(1,C/d,A/d,D/d,\alpha/d)$, a positive unit, retain every
excess and the whole $m,n$ positive square terms, and use its displayed open-axis
state forms and two corner quadratic forms. That proof used lower
bounds, not upper bounds, on these four source weights. Noncommuting
residue tests use the retained arbitrary-physical-source theorem.
Finally homogenize (TD6)--(TD12) in $W$, replacing $R,T$ by $RW,TW$;
restore pure testing projective factors and split complex states into
real and imaginary parts. This proves the theorem for all states,
including $W=0$, without comparing $k,n,D$ with independent testing
infinitesimals.

\paragraph{Exact scope.}
This is an all-tests theorem for the stated actual-contact domain.
It does not bound or remove the residual regimes where $C/d$ or
$\alpha/d$ is unbounded, $A/d$ is unbounded, or the actual contact
has a different angle-residue type. The C2 reflection
$(m,C,D,\alpha;u_c,v_c)\mapsto(D,\alpha,m,C;1/v_c,1/u_c)$ supplies
its corresponding reflected domain. No exhaustive C2 join follows
from these two domains alone.

\subsection{Three arbitrary event weights}
\label{component:p6c62a53b1454}

\label{p6c62a53b1454:sec:c2-three-large}

\begin{theorem}[All scales of the three retained events]
\label{p6c62a53b1454:thm:c2-three-large}
Let $K_c,K_t$ be independent real closed extensions of $\mathbb R$.
Let a physical canonical C2 source over $K_c$ have $\lambda=\beta=0$,
with $d,C,\alpha$ positive bounded units and
$n,m,A,D\geq0$ arbitrary. Suppose an actual finite noncommuting
stationary zero has $u_c=k>0$ infinitesimal and $v_c=a>0$ a bounded
unit. Then the entire original parent is tensor SOS at every
independent projective measurement pair and every real or complex
projective state. No ratios among $n,A,D$ are restricted.
\end{theorem}

The bounded-$A$ result is a special case. The mechanism is to keep
all three possibly large event amplitudes as state coordinates.
The exact parent and its amplitudes are (TD1); in particular
$J=uW+uvX-Y-vZ$.

\paragraph{Source entry with three large weights.}
The initial coercivity estimate in the preceding proof uses only
$d,C,\alpha$, not $A$. Thus after $W_c=1$,
\[
 X_c,Y_c,Z_c=O(k),\qquad m=O(k^2).
\]
Now use $F_Z=0$ first. The $A$ and $n$ positive square terms do not enter this
state equation. Every term other than $-8DvJ$ is $O(k)$, because
$v_c=a$ is a unit and the just displayed bounds hold. Therefore
$DJ_c=O(k)$. The $Y$ equation then gives $AY_c=O(k)$, and the $X$
equation gives $(d+2n)X_c=O(k^2)$. Write
\[
 N=d+2n,\quad T_A=d+2A,\quad L=d+2D,\qquad
 x=X_c/k^2,\quad y=Y_c/k,\quad j=J_c/k.
\]
Consequently
\[
 x,y,j,Nx,Ay,Dj\ \hbox{are bounded},\qquad N^{-1}\ \hbox{bounded}.
                                                               \tag{TA1}
\]
These estimates do not yet invert $A$ or $D$.

Use the exact stationary equations (TD3) and divide both by
$T_A L$. Now $t_2/T_A,t_6/T_A,t_3/L,t_4$ are positive bounded
units. In the first equation the coefficient multiplying $U^2$ is
$(A/T_A)$ times a positive bounded unit: its remaining bracket is
\[
 4t_4\left\{\theta\frac{t_3}{L}
            +CV\frac{1+2D/N}{L}\right\},
\]
whose lower bound follows from (TD4). The coefficient multiplying
$m$ is also a positive bounded unit:
\[
 4V\left\{\alpha\frac{t_8}{T_A L}
              +\frac{D}{L}\frac{t_6}{T_A}V\right\}.             \tag{TA2}
\]
Indeed, if $D/L$ is bounded below, the second summand is a unit.
Otherwise $D$ is infinitesimal and $L$ is a unit, while
$t_8/T_A=1+2(C+D)/T_A$ is a unit. All displayed factors have
bounded upper bounds in either case. Hence
\[
 \frac{m}{U^2}=\frac{A}{T_A}\cdot(\hbox{positive bounded unit}). \tag{TA3}
\]
The second stationary equation becomes
\[
 \frac{A}{T_A}\cdot(\hbox{positive unit})
       +O\left(\frac{A}{T_A}U\right)
 =\frac{D}{L}\cdot(\hbox{positive unit})
       +O\left(\frac{A}{T_A}U\right).                          \tag{TA4}
\]
Thus $A/T_A$ and $D/L$ are comparable. If they were infinitesimal,
$A,D,m$ would all be infinitesimal, contradicting exactly the
physicality test printed after (TD5). Therefore
\[
 A/T_A,\ D/L,\ m/k^4\quad\hbox{are positive bounded units}.       \tag{TA5}
\]
In particular $A,D$ have bounded reciprocals, but may be arbitrarily
large. Set $\gamma=4ma^2/k^4$, also a positive bounded unit.

\paragraph{The three-positive square term graph.}
On the local testing domain $\bar v=\bar a\ne0$, put
$b=1-j-y$ and use
\[
 X_g=xu^2W,\qquad Y_g=yuW,\qquad J_g=juW,\qquad
 Z_g=u\left(\frac b v+xu^2\right)W.                            \tag{TA6}
\]
This graph matches the whole actual contact: the definition of $J_c$
gives $Z_c/k=b/a+xk^2$. Its three hard normals are
\[
 E_X=\sqrt N(X-xu^2W),\quad
 E_Y=\sqrt A(Y-yuW),\quad E_J=\sqrt D(J-juW).                    \tag{TA7}
\]
The exact inverse is
\[
 X=xu^2W+E_X/\sqrt N,\quad Y=yuW+E_Y/\sqrt A,
\quad Z=\frac{uW+uvX-Y-juW-E_J/\sqrt D}{v}.                    \tag{TA8}
\]
Every denominator is either pure coefficient or the pure testing
unit $v$. All graph coefficients and all scaled normal coefficients
are bounded. For the three large positive square terms this follows explicitly
from
\[
 n x^2,\ nx,\ n/N,\quad A y^2,\ \sqrt A\,y,\quad
 D j^2,\ \sqrt D\,j\quad\hbox{bounded},                       \tag{TA9}
\]
using (TA1),(TA5). Their quadratic normal contributions are exactly
$4(n/N)D_uD_v E_X^2$, $4D_uD_vE_Y^2$, and $4E_J^2$.
Every other original source coefficient is bounded, including $m$.

For clarity, the finite Taylor division used below works in the
bounded Laurent polynomial ring in $v$: for positive integer $r$,
\[
 v^{-r}-a^{-r}
 =-(v-a)\frac{\sum_{i=0}^{r-1}v^i a^{r-1-i}}{a^rv^r}.
\]
Thus a Laurent polynomial vanishing at $(U,v)=(K,a)$ belongs to
$(U-K,v-a)$ with bounded Laurent coquotients; repeating the same
identity gives the second-order division. There is no localization
at a difference of the two fields.

Set $K=k^2$, $R=u^2-(K/a)v$ and $T=u(v-a)$. The entire graph
scalar is again $4mv^2+u^2H_0(u^2,v)$, now with bounded Laurent
coefficients in $v$. Its three scalar contact jets give (TD10),
and hence an exact bounded Gram on $(R,T)$ with $R$ residue pivot
$\bar\gamma$. The full first normals before their pure source
scale divisions are
\[
\begin{split}
 (F_X+uF_Z)|_g&=-8mv+u^2L_X(u^2,v),\\
 (F_Y-F_Z/v)|_g&=uL_Y(u^2,v),\\
 (-F_Z/v)|_g&=uL_J(u^2,v).                                     \end{split} \tag{TA10}
\]
Divide the rows by $\sqrt N,\sqrt A,\sqrt D$, respectively.
The scaled coefficients are bounded by (TA9). At contact they
vanish. Laurent Taylor division gives exactly the root membership
and infinitesimal $R$ couplings proved in (TD11)--(TD12).
The $X$ row has infinitesimal coefficients on both $R,T$ because
$m/K=O(K)$; the other rows have infinitesimal coefficients on $R$.
The exact quadratic state remainder is bounded by (TA8)--(TA9),
so the entire original parent has a bounded five-root Gram on
$(E_X,E_Y,E_J,R,T)$.

\paragraph{Uniform strictness without a largest-event choice.}
The $X$ and $R$ residue pivots are $2(1+\bar a^2)$ and
$\bar\gamma$. The remaining hard state form is still (TD13).
At $u=0$, (TA8) gives
\[
 Y=E_Y/\sqrt A,\qquad
 Z=-\frac{E_Y}{v\sqrt A}-\frac{E_J}{v\sqrt D}.
\]
In these coordinates (TD13) is the sum of the two explicit pivots
\[
 4(1+v^2)E_Y^2+4E_J^2
\]
and the nonnegative bounded form
$(4C+2d)(1+v^2)Y^2+4\alpha(vY-Z)^2$.
Its residue is therefore positive definite for every combination of
bounded or unbounded $A,D$.

The boundary Schur function is the exact expression
\[
 f(v)=\frac{8(1+v^2)}{\det B(v)/(T_A L)}
 \left(\frac{Q_{10}}{T_A L}+
       \frac{Q_0}{T_A L}v^2+\frac{Q_{12}}{T_A L}v^4\right),     \tag{TA11}
\]
with $Q_{10},Q_{12}$ as in (TD14). For $v$ a bounded unit all its
coefficients and denominator are bounded, and
\[
 \frac{Q_{10}}{T_A L}=4\alpha\frac A{T_A}\frac{t_3}{L},\qquad
 \frac{Q_{12}}{T_A L}=4C\frac D L\frac{t_2-2m}{T_A}
\]
are positive units. The determinant has the positive term
$16AD(1+v^2)v^2$, so its quotient by $T_A L$ has a positive lower
bound. Its other terms are bounded after that division.
The exact middle coefficient is polynomial of degree at most one
in each of $A,D$, so $Q_0/(T_A L)$ is bounded too.

The contact normal equations and scalar jets give
$\bar f(\bar a)=\bar f'(\bar a)=0$ as before. Therefore the real
quartic bracket equals
$\overline{Q_{12}/(T_A L)}(v^2-\bar a^2)^2$.
Its strict second derivative gives the positive soft Schur pivot on
$T$. The whole five-root residue Gram is strict. This proves the
local all-state certificate; the Bob sign identity gives its negative
sign copy. Homogenization in $W$ includes $W=0$.

\paragraph{Outer tests when both $A$ and $D$ are large.}
For $u$ infinitesimal and $v$ a bounded unit with
$\bar v\ne\pm\bar a$, retain the entire $m,n$ positive square terms and use the
four roots
\[
 X,\quad K_A=\sqrt A Y,\quad K_D=\sqrt D J,\quad uW.
\]
Substitute $Y=K_A/\sqrt A$ and
$Z=(uW+uvX-K_A/\sqrt A-K_D/\sqrt D)/v$.
The exact remainder is a bounded quadratic Gram. Its hard residue
has the positive $X$ pivot $2\bar d(1+\bar v^2)$ and the two
explicit pivots just displayed; its soft Schur complement is
$\bar f(\bar v)>0$ by (TA11).

At $u,v$ both infinitesimal, and at $u,1/v$ both infinitesimal,
one need not divide by $v$. Cap $C,A,D,\alpha$ at the positive
unit $\mu=\min(1,C/d,A/d,D/d,\alpha/d)$ after division by $d$,
and retain all excesses and the full $m,n$ positive square terms. The bounded
selected parent has an exact Gram on $(X,Y,Z,uW)$ in the finite
chart and on the same four roots in the homogeneous reciprocal
chart. At $(u,v)=(0,0)$ its soft Schur complement is
$8Q_{10,\mu}/\det B_\mu(0)>0$; at $(u,1/v)=(0,0)$ it is
$8Q_{12,\mu}/\det B_{\mu,\infty}>0$. The hard blocks are positive
by (TD13) and its homogeneous limit. Hence both entire endpoint
neighborhoods have strict-residue Grams, including literal zeros
and infinities. These expressions use only the bounded selected
parent; no residue of an unbounded $A$ or $D$ is taken.

For every other commuting test, the open-axis and corner forms of the
preceding proof apply with the four lower bounds (TA5). Noncommuting
residue tests use the physical-source theorem. Restore the positive pure
factors and split complex states into real and imaginary parts.

\paragraph{Coverage gained and remaining restriction.}
The theorem removes every relative-scale restriction among $n,A,D$
from the finite pitchfork. Its reflected form covers the corresponding
ordinary--infinite actual contacts. It still assumes $C/d$ and
$\alpha/d$ are bounded units and this actual-angle residue type.
Unbounded $C/d$ or $\alpha/d$, other actual-angle types, and the
exhaustive C2 join remain outside this theorem.

\subsection{Arbitrary complementary event scales}
\label{component:pf7943b34a7dc}

\label{pf7943b34a7dc:sec:strong-complement-pitchfork}

\begin{theorem}[No upper bound on the two complementary events]
\label{pf7943b34a7dc:thm:strong-complement-pitchfork}
Let $K_c,K_t\supset\mathbb R$ be independent real closed fields.
Let the original canonical C2 parent be physical over $K_c$, with
$\lambda=\beta=0$, $d$ a positive bounded unit, $C,\alpha>0$
having bounded reciprocals, and $n,m,A,D\geq0$ arbitrary.
Suppose it has an actual finite noncommuting stationary zero with
$u_c=k>0$ infinitesimal and $v_c=a>0$ a bounded unit.
Then the entire original parent is tensor SOS at every independent
projective measurement pair and every real or complex projective state.
There is no upper bound or relative-scale condition on any event weight.
\end{theorem}

The case of bounded $C,\alpha$ is the existing three-large-event
theorem. We prove both remaining cases, without modifying the parent
or identifying the source and testing fields. We use
\[
\begin{gathered}
 D_u=1+u^2,\quad D_v=1+v^2,\quad r=uW-Y,\quad p=vY-Z,\\
 J=uW+uvX-Y-vZ,\quad R_2=vW-X+up,\quad B=X+uZ,\\
 H=2\{J^2+D_up^2+D_vB^2-D_uD_vZ^2\},\\
 F=dH+4nD_uD_vX^2+4mR_2^2+4CD_vr^2
       +4AD_uD_vY^2+4DJ^2+4\alpha D_up^2.              
\end{gathered}
\tag{SC1}
\]
Every square in this formula retains its original coefficient.

\paragraph{Necessary reductions from the full parent.}
Cap only $C,\alpha$ at positive bounded units in the initial hard
state estimate. At the actual angles the bounded selected form is
bounded below by
$b(X^2+Y^2+Z^2)-B_0k^2W^2$, with positive real $b,B_0$.
Every retained positive square term is nonnegative. Hence the contact has $W_c\ne0$;
on setting $W_c=1$,
\[
 X_c,Y_c,Z_c=O(k),\qquad m=O(k^2).
\]
Each positive event positive square term is $O(k^2)$ as well. Thus, writing
$y=Y_c/k$, $z=Z_c/k$, $x_1=X_c/k$, $j=J_c/k$,
\[
 Ay^2,\ Dj^2,\ C(1-y)^2,\ \alpha(ay-z)^2=O(1),\qquad
 j=1-y-az+aX_c.                                      \tag{SC2}
\]
If $C$ is unbounded then $\bar y=1$, so $A$ is bounded. If $D$
were also unbounded, $j,z$ would be infinitesimal, and the residue of
\[
 H_c/(2k^2)=j^2+(1+k^2)(ay-z)^2
            +(1+a^2)(x_1+Z_c)^2-(1+k^2)(1+a^2)z^2
\]
would be $\bar a^2+(1+\bar a^2)\bar x_1^2>0$.
This contradicts $H_c\leq0$ at the zero of (SC1). Therefore
$C$ unbounded forces both $A,D$ bounded. If $\alpha$ is unbounded
and $A$ were unbounded, (SC2) would make $y,z$ infinitesimal and
$\bar j=1$; the same residue would instead be
$1+(1+\bar a^2)\bar x_1^2>0$. Thus $\alpha$ unbounded forces $A$
bounded. It remains to prove precisely
\[
 \text{(I) } C\text{ unbounded},\ A,D\text{ bounded};\qquad
 \text{(II) } C\text{ bounded},\ \alpha\text{ unbounded},\ A\text{ bounded}.
                                                               \tag{SC3}
\]

\paragraph{Uniform stationary scale entry.}
Set $U=k^2$, $V=a^2$, $N=d+2n$, $\theta=n/N$, and
\[
\begin{gathered}
 t_2=d+2(\alpha+m+A),\quad T=t_3=d+2(C+D),\quad
 t_4=d+2(\alpha+m),\\
 t_6=d+2(\alpha+C+A),\quad t_8=d+2(C+A+D).
\end{gathered}
\]
The exact original Laurent stationary equations are
\[
\begin{split}
 U^2 A t_4\{\theta T+C(1+2D/N)V\}
     &=mV\{\alpha t_8+Dt_6V\},\\
 \alpha A(1+2m/N)T+\theta A Tt_4U
     &=\{mDt_6/U+CDt_2\}V^2.                         
\end{split}\tag{SC4}
\]
Define only pure coefficient quantities
\[
 \delta=CD/T,\quad e=(D/T)(t_6/\alpha),\quad b=t_2/\alpha,
 \quad \mathcal B=\theta+(C/T)(1+2D/N)V,
 \quad \mathcal G=t_8/T+eV.
\]
Because one of $C,D$ is bounded in either case (SC3),
$\delta\leq\min(C,D)/2$ is bounded. The identities
\[
 e=(D/T)\{2+(d+2A)/\alpha\}+2\delta/\alpha,
 \quad \mathcal G=1+2A/T+eV,
 \quad \mathcal B=\theta+(C/T)V+(2\delta/N)V
\]
show boundedness of all these quantities. The quantities
$b,\mathcal G,t_4/\alpha$ are positive bounded units. So is
$\mathcal B$: since $C,N$ have bounded reciprocals and $d$ is
bounded,
\[
 \frac{T}{C(N+2D)}\leq\frac d{CN}+\frac2N+\frac1C=O(1),
\]
and consequently
\[
 \theta+(C/T)(1+2D/N)
 =\frac{n+C(N+2D)/T}{d+2n}
 \geq\frac{n+c_0}{d+2n}
\]
for a positive real $c_0$. The bounded-unit factor $V$ preserves
this lower bound. Dividing (SC4) by $\alpha T$ gives exactly
\[
 \frac m{U^2}=A\rho,\qquad
 A\mathcal H=\delta bV^2,\qquad
 \rho=\frac{(t_4/\alpha)\mathcal B}{V\mathcal G},\quad
 \mathcal H=1+2m/N+\theta(t_4/\alpha)U-U\rho eV^2.      \tag{SC5}
\]
Here $\rho$ is a positive bounded unit and $\mathcal H=1+O(U)$.
These identities were obtained before inverting $A$ or $D$.
They show $A\asymp\delta$. Since $C$ has a bounded reciprocal,
$\delta$ infinitesimal implies $D$ infinitesimal.
If $A$ were infinitesimal, therefore, so would $D,m$ be.
But the exact physical test
\[
 u=v=1,\qquad (W,X,Y,Z)=(1,0,1,1)
\]
has $X=r=p=0$, $R_2=1$, $J=-1$, $H=-2$, and hence
$F=-2d+4m+16A+4D$. Its residue would be negative. We conclude
\[
 A\asymp1,\qquad D^{-1}=O(1),\qquad
 m/k^4\asymp1,\qquad \gamma=4ma^2/k^4\asymp1.        \tag{SC6}
\]
In case (I), $D$ is a bounded unit; in case (II), it may be
arbitrarily large. In both cases $\delta$ is a positive bounded unit.

\paragraph{Full state equations and bounded graph data.}
For later use, the three entire state derivatives of (SC1) are
\[
\begin{split}
 F_X/4={}&D_v[d+2nD_u]X+(d+2D)uvJ+d uD_vZ-2mR_2,\\
 F_Y/4={}&-(d+2D)J+(d+2\alpha)vD_up-2CD_vr
                  +2AD_uD_vY+2muvR_2,\\
 F_Z/4={}&-(d+2D)vJ-(d+2\alpha)D_up+d uD_vB
                  -dD_uD_vZ-2muR_2.                
\end{split}\tag{SC7}
\]
In case (I), the $Z$ equation first yields $\alpha p_c=O(k)$,
because $D$ is bounded and every core coordinate is $O(k)$.
Then the $Y$ equation yields $Cr_c=O(k)$. In case (II), $C,A$
are bounded, so the $Y,Z$ equations are a system with bounded right
side $O(k)$ and unit coefficient matrix
\[
 \begin{pmatrix}-1&a\\-a&-1\end{pmatrix}
\]
on $((d+2D)J_c,(d+2\alpha)(1+k^2)p_c)$.
They yield $DJ_c,\alpha p_c=O(k)$. Finally the $X$ equation in
either case yields $NX_c=O(k^2)$. Thus put $x=X_c/k^2$ and
\[
\begin{array}{ll}
 \text{(I)}&r_0=r_c/k,\ p_0=p_c/k,
       \qquad Nx,Cr_0,\alpha p_0=O(1),\\
 \text{(II)}&j_0=J_c/k,\ p_0=p_c/k,
       \qquad Nx,Dj_0,\alpha p_0=O(1).
\end{array}                                                    \tag{SC8}
\]
All unweighted graph constants are bounded as well. The reciprocals
of $N,C,D,\alpha$ used below are bounded pure coefficient elements.

\paragraph{Two exact graphs retaining the large positive square terms.}
For independent finite testing angles, case (I) uses
\[
 X_g=xu^2W,\quad r_g=r_0uW,\quad p_g=p_0uW,
 \quad Y_g=(1-r_0)uW,\quad Z_g=\{v(1-r_0)-p_0\}uW.    \tag{SC9}
\]
Its exact normals and inverse are
\[
\begin{gathered}
 E_X=\sqrt N(X-xu^2W),\quad E_r=\sqrt C(r-r_0uW),
 \quad E_p=\sqrt\alpha(p-p_0uW),\\
 X=xu^2W+E_X/\sqrt N,\quad
 Y=(1-r_0)uW-E_r/\sqrt C,\quad
 Z=vY-p_0uW-E_p/\sqrt\alpha.                         
\end{gathered}
\tag{SC10}
\]
Case (II) uses
\[
 X_g=xu^2W,\quad J_g=j_0uW,\quad p_g=p_0uW,
 \quad Y_g=\frac{u(1-j_0+vp_0+xu^2v)}{D_v}W,
 \quad Z_g=vY_g-p_0uW.                              \tag{SC11}
\]
Here the exact normals and inverse are
\[
\begin{gathered}
 E_X=\sqrt N(X-xu^2W),\quad E_J=\sqrt D(J-j_0uW),
 \quad E_p=\sqrt\alpha(p-p_0uW),\\
 X=xu^2W+E_X/\sqrt N,\quad J=j_0uW+E_J/\sqrt D,
 \quad p=p_0uW+E_p/\sqrt\alpha,\\
 Y=(uW+uvX+vp-J)/D_v,\qquad Z=vY-p.                 
\end{gathered}
\tag{SC12}
\]
Both graphs match the entire actual contact. All denominators in
these formulas are pure coefficient factors or the pure testing
unit $D_v$. No inverse of a mixed tensor element occurs.
The possibly unbounded positive square terms become respectively
\[
\begin{array}{ll}
 \text{(I)}&4(n/N)D_uD_vE_X^2,quad4D_vE_r^2,quad4D_uE_p^2,\\
 \text{(II)}&4(n/N)D_uD_vE_X^2,quad4E_J^2,quad4D_uE_p^2
\end{array}
\]
in the quadratic normal remainder. Their graph and first-normal
coefficients are bounded by (SC8), for example
$nx,nx^2,Cr_0^2,\sqrt C r_0,\alpha p_0^2,\sqrt\alpha p_0$
in case (I), and $Dj_0^2,\sqrt D j_0$ in case (II).
Every remaining source coefficient is bounded.

\paragraph{Entire scalar and entire first-normal division.}
Set $W=1$, $K=k^2$, and
$R=u^2-(K/a)v$, $T_s=u(v-a)$.
Since $R_2|_g=v+(p_0-x)u^2$ for both graphs, the entire scalar has
the exact shape
\[
 \Phi=F|_g=4mv^2+u^2H_0(u^2,v),                    \tag{SC13}
\]
with bounded polynomial coefficients in case (I), and bounded
coefficients in $B_c[v,D_v^{-1}]$ in case (II).
The three scalar contact jets give
$H_0(K,a)=-\gamma K$, $H_{0,U}(K,a)=\gamma$, and
$H_{0,v}(K,a)=-2\gamma K/a$. In particular the explicit expression
\[
 \Psi=H_0-\gamma u^2+2\gamma(K/a)v
\]
satisfies $\Phi-\gamma R^2=u^2\Psi$, and its value and both first
derivatives vanish at $(u^2,v)=(K,a)$. Finite Taylor division of
$\Psi$ therefore gives
\[
 \Phi=\gamma R^2+u^2\{(u^2-K)^2H_{20}
       +(u^2-K)(v-a)H_{11}+(v-a)^2H_{02}\},           \tag{SC14}
\]
with bounded coquotients. This is an exact bounded Gram on
$(R,T_s)$, with separate residue pivot $\bar\gamma$ on $R$.
For the rational case, Taylor division is legitimate because
\[
 D_v^{-r}-D_a^{-r}
 =-(v-a)(v+a)
   \frac{\sum_{i=0}^{r-1}D_v^iD_a^{r-1-i}}{D_a^rD_v^r}.
\]
Iterating this identity gives the second-order division; only pure
factors have been inverted.

Here are the full first normals, rather than initial forms. In
case (I) they are
\[
 \mathcal N_X=F_X|_g,\quad
 \mathcal N_r=-(F_Y+vF_Z)|_g,\quad
 \mathcal N_p=-F_Z|_g,                              \tag{SC15}
\]
and in case (II) they are
\[
 \mathcal N_X=\{F_X+(uv/D_v)(F_Y+vF_Z)\}|_g,\quad
 \mathcal N_J=-(F_Y+vF_Z)|_g/D_v,\quad
 \mathcal N_p=(vF_Y-F_Z)|_g/D_v.                     \tag{SC16}
\]
The derivatives in (SC7) specify every term. In particular the
following exact positive square term expressions expose the cancellations:
\[
\begin{split}
 -(F_Y+vF_Z)/(4D_v)
   &=(d+2D)J+2Cr-2AD_uY-duvB+dvD_uZ,\\
 (vF_Y-F_Z)/(4D_v)
   &=(d+2\alpha)D_up-2Cvr+2AvD_uY+2muR_2-duB+dD_uZ.
\end{split}\tag{SC17}
\]
The actual coefficients of the normal roots are obtained by dividing
these rows by $\sqrt N,\sqrt C,\sqrt\alpha$ in case (I), and by
$\sqrt N,\sqrt D,\sqrt\alpha$ in case (II).
Using (SC8), their exact forms after those divisions are
\[
 \mathcal N_X/\sqrt N=-8mv/\sqrt N+u^2L_X(u^2,v),
 \qquad \mathcal N_i/\sqrt{w_i}=uL_i(u^2,v),          \tag{SC18}
\]
where all $L_i$ have bounded coefficients in the indicated ring.
The contact equations give
$L_X(K,a)=8ma/(K\sqrt N)$ and $L_i(K,a)=0$.
Writing $L_X=L_X(K,a)+(u^2-K)L_{X1}+(v-a)L_{X2}$, and similarly
for $L_i$, gives the full exact row identities
\[
\begin{split}
 \mathcal N_X/\sqrt N
  &=(8ma/(K\sqrt N)+u^2L_{X1})R
       +u\{(K/a)L_{X1}+L_{X2}\}T_s,\\
 \mathcal N_i/\sqrt{w_i}
  &=uL_{i1}R+\{(K/a)L_{i1}+L_{i2}\}T_s.             
\end{split}\tag{SC19}
\]
Thus the $X$ row has infinitesimal couplings to both soft roots,
and the other rows have infinitesimal couplings to $R$.
The exact state degree is two, so the bounded quadratic remainder
already accounts for every higher-normal term. Equations
(SC10)--(SC19) give an exact full-parent bounded five-root Gram.

\paragraph{A common exact real Schur certificate.}
The $X$ and $R$ residue pivots are $2(1+\bar a^2)$ and
$\bar\gamma$. The unscaled boundary hard form on $Y,Z$ is
\[
 \mathsf B(v)[Y,Z]
 =sD_vY^2+4D(Y+vZ)^2+4\alpha(vY-Z)^2,
 \qquad s=4(C+A)+2d.                               \tag{SC20}
\]
In case (I), $Y=-E_r/\sqrt C$,
$Z=-vE_r/\sqrt C-E_p/\sqrt\alpha$ at $u=0$.
The form has the explicit pivots $4D_vE_r^2+4E_p^2$ plus a
bounded nonnegative form. In case (II),
$Y=(vE_p/\sqrt\alpha-E_J/\sqrt D)/D_v$ and
$Z=(-E_p/\sqrt\alpha-vE_J/\sqrt D)/D_v$.
The explicit pivots are $4E_J^2+4E_p^2$, again with a bounded
nonnegative remainder. The hard residue is strict in both cases.

The determinant and the exact boundary Schur numerator are
\[
\begin{gathered}
 \det\mathsf B(v)=4(1+v^2)\{s(D\,v^2+\alpha)+4D\alpha(1+v^2)\},\\
 Q_{10}=4\alpha A[d+2(C+D)],\qquad
 Q_{12}=4CD[d+2(\alpha+A)],\\
 Q_0=8ACD+8AC\alpha-4ADd-2Ad^2-2Cd^2-4Cd\alpha
       -2Dd^2-4Dd\alpha-d^3-2d^2\alpha.             
\end{gathered}
\tag{SC21}
\]
Put $S=\alpha T$. Every coefficient of $\det\mathsf B/S$ and of
$Q_{10}/S,Q_0/S,Q_{12}/S$ is bounded. This follows directly from
(SC21), bounded $A$, bounded $C/T,D/T,\delta$, and bounded
$1/\alpha$. Moreover
\[
 Q_{10}/S=4A\asymp1,\qquad
 Q_{12}/S=4\delta[d+2(\alpha+A)]/\alpha\asymp1.       \tag{SC22}
\]
The normalized determinant has a positive lower bound for every
bounded unit real $v$. Indeed it equals $C/T$ times the determinant
of the scaled hard form in case (I), and equals
$(D/T)D_v^2$ times that determinant in case (II). The factors
$C/T$ and $D/T$, respectively, are positive units in their cases.

We now take residues of these bounded coefficients and use an
ordinary real variable $v$. The exact real hard Schur function is
\[
 \bar f(v)=
 \frac{8(1+v^2)}{\overline{\det\mathsf B/S}(v)}
 \left\{\overline{Q_{10}/S}+\overline{Q_0/S}v^2
                      +\overline{Q_{12}/S}v^4\right\}.          \tag{SC23}
\]
This formula is an ordinary real Schur calculation, never an
inverse of $\det\mathsf B(v)$ in $K_c\otimes K_t$.
The contact first-normal identities identify the graph at $\bar a$
with the unique minimizer of the scaled real hard form. The scalar
contact jets give $\bar f(\bar a)=\bar f'(\bar a)=0$.
Because $\bar a\ne0$, the quartic bracket in (SC23) is exactly
\[
 \overline{Q_{12}/S}(v^2-\bar a^2)^2.
\]
Its strict second derivative supplies the soft Schur pivot
\[
 \tfrac12\bar f''(\bar a)=
 \frac{32(1+\bar a^2)\bar a^2\overline{Q_{12}/S}}
      {\overline{\det\mathsf B/S}(\bar a)}>0.
\]
The entire five-root Gram therefore has strictly positive residue
when $u$ is infinitesimal or zero and $\bar v=\bar a$.
The strict-residue Gram lemma proves tensor SOS. The exact identity
\[
 F(u,-v;W,-X,Y,-Z)=F(u,v;W,X,Y,Z)
\]
supplies $\bar v=-\bar a$.

\paragraph{All other independent tests, including literal endpoints.}
If $u$ is infinitesimal and $v$ is a bounded unit with
$\bar v\ne\pm\bar a$, retain the full original $m,n$ positive square terms.
In case (I) use the roots
\[
 X,\quad K_C=\sqrt C r,\quad K_\alpha=\sqrt\alpha p,\quad uW
\]
and the exact inverse
$Y=uW-K_C/\sqrt C$, $Z=vY-K_\alpha/\sqrt\alpha$.
In case (II) use
\[
 X,\quad K_D=\sqrt D J,\quad K_\alpha=\sqrt\alpha p,\quad uW
\]
and the inverse in (SC12), with unshifted $J,p$.
The entire remaining parent is a bounded quadratic Gram on these
four roots. Its separate $X$ residue pivot is
$2\bar d(1+\bar v^2)$; its hard pair is the strict form (SC20)
in the scaled coordinates just displayed, and its soft Schur
complement is $\bar f(\bar v)>0$ by (SC23).

For the two endpoint neighborhoods $u,v\sim0$ and
$u,q=1/v\sim0$, divide by $d$, choose
\[
 \mu=\min(1,C/d,A/d,D/d,\alpha/d)>0,
\]
a bounded unit by (SC6), and retain all excesses together with
the full $m,n$ positive square terms. The selected bounded parent has signed
weight one, $C=A=D=\alpha=\mu$, and $m=n=0$.
In the finite chart it is exactly a Gram on $(X,Y,Z,uW)$.
At $(u,v)=(0,0)$ its hard block is positive and its soft Schur
complement is $8Q_{10,\mu}/\det\mathsf B_\mu(0)>0$.
For the reciprocal chart use the polynomial homogeneous parent
$q^2F_\mu(u,1/q)$, on the same four roots. At $(u,q)=(0,0)$
its hard block is
$\operatorname{diag}(2,4(C+A+\alpha)+2,4D)$ in the
$(X,Y,Z)$ variables, evaluated at the capped weights, and its
soft Schur complement is
$8Q_{12,\mu}/\det\mathsf B_{\mu,\infty}>0$.
These are bounded polynomial Grams, so they include the full
infinitesimal neighborhoods and literal $v=0$ or $q=0$.

For completeness, the other open commuting axes use the same cap;
here the other angle has a nonzero finite residue.
At $v=0$ the core is $2\{(uW-Y)^2+(X+uZ)^2\}$ and the
$A,\alpha$ positive square terms force $Y=Z=0$; finite nonzero $u$ then forces
$W=X=0$. At $u=\infty$ the homogeneous core is
$2\{(W+vX)^2+(vY-Z)^2\}$; the $C,A$ positive square terms force $W=Y=0$,
and finite nonzero $v$ forces $X=Z=0$. At $v=\infty$ the core
is $2(1+u^2)(X^2+Y^2)$; the $C,D$ positive square terms force $W=Z=0$
for finite nonzero $u$. Thus these complete state matrices are strict.

Only the two corners with $u=\infty$ remain. Put $p_u=1/u$.
At $(p_u,v)=(0,0)$ the selected state form is
$(2+8\mu)W^2+4\mu Y^2+(2+4\mu)Z^2$.
Its sole zero ray is $X$. In the chart $X=1$, the five small
variables $(p_u,v,W,Y,Z)$ have exact quadratic part
\[
 (2+4\mu)(W+v)^2+2(Z+p_u)^2+4\mu(W^2+Y^2+Z^2)\succ0.
\]
At $(p_u,q)=(0,0)$ the selected state form is
$4\mu W^2+(2+4\mu)X^2+(2+8\mu)Y^2$.
In its sole zero-ray chart $Z=1$, the variables
$(p_u,q,W,X,Y)$ have exact quadratic part
\[
 2X^2+(2+4\mu)(Y-q)^2+4\mu\{W^2+Y^2+(X-p_u)^2\}\succ0.
\]
There are no lower-order terms. Allocating every higher monomial
to two coordinate roots leaves an infinitesimal factor, giving a
strict-residue Gram for the entire selected polynomial at each
corner. Retain the original excess positive square terms throughout.
After pure testing projective normalization, states outside each
displayed zero ray have positive residue value and are covered by
the scalar strict-residue lemma.

Noncommuting-residue tests use the inherited arbitrary-physical-source
theorem. This exhausts all independent projective tests. Finally
homogenize the local graph identities in $W$, replacing $R,T_s$
by $RW,T_sW$; they then include $W=0$. Restore all pure positive
coefficient and testing factors, and split complex states into real
and imaginary parts. No source infinitesimal was compared with an
independent testing infinitesimal anywhere in the argument.

\paragraph{Domain.}
This proves the finite pitchfork when $C/d,\alpha/d$ have positive real
lower bounds, with arbitrary upper event scales, and its C2 reflection.
Other actual-angle types and vanishing complementary lower bounds remain
separate cases.

\subsection{Source bounds from an infinitesimal first angle}
\label{component:pf5726218aec8}

\label{pf5726218aec8:sec:strong-angle-entry}
Let $z=u_c^2$ be infinitesimal, $y=v_c^2>0$ arbitrary, and use the exact
positive C2 source. Set
\[
 \begin{gathered}
 H=zh\in(0,1),\quad L=yt>0,\quad B=zb=H+L,\\
 A_c=1+zr,\quad S=A_c+L,\quad j=A_c-L/y>0.
 \end{gathered}
\]
Here $0<r<1$. Define
\begin{align*}
 \delta={}&z+(1+z)B+(1+z)B^2+2z(1+z)Br+z^3r^2+2z^2r\\
 &+y\{(1+z)B+z+z^2r(2-r)\},\\
 G={}&B+z(H+r-1),\\
 M={}&(1+z)(B+zr)(H+y)-z(1+y)(1-r)L\\
 ={}&(1+z)(B+zr)H+yG+yzj+zrL(1+y)>0.
\end{align*}
The unchanged original coefficient formulas are
\begin{align*}
 n&=\frac{(B+zr)A_cS-r\delta}{2r\delta},&
 m&=\frac{z^2(1+y)(1-r)A_cS}{2G\delta},\\
 C&=\frac{A_cM}{2L\delta},&
 A&=\frac{yA_c(B+zr)S}{2H\delta},\\
 D&=\frac{(1+z)(B+zr)(B-zyr)S}{2yj\delta},&
 \alpha&=\frac{(B+zr)M}{2(1-H)\delta}.
\tag{GE1}
\end{align*}
Assume $C$ and $\alpha$ are strong, meaning each has a positive real
lower bound in the coefficient field. Then $B$ is a bounded positive
unit. Indeed
\[
 \delta\ge(1+y)(B+z),\qquad
 M\le(1+z)(B+z)(1+y)
\]
give $CL\le A_c(1+z)/2$, so $L$ and $B$ are bounded.
If $B$ were infinitesimal then so would $H$, while
$\alpha(1-H)\le(B+zr)(1+z)/2$ would force $\alpha$ infinitesimal,
a contradiction.

Next, $n\ge0$ gives $r\delta\le(B+zr)A_cS$ and therefore
$yr\le A_cS$, since $\delta\ge y(B+zr)$.
Thus $G$, $B-zyr$ and $\delta/(1+y)$ are positive units.
Since $L<yA_c$, the ratio $(H+y)/(1+y)$ is a positive unit: if $y$
is infinitesimal then $H=B-L$ is a unit, and otherwise $y/(1+y)$
is a unit. The positive second expression for $M$ now shows that
$M/(1+y)$ is a positive unit.

Finally $1-r$ is a unit. To see this, let
$E(r)=(B+zr)A_cS-r\delta$, holding $B,H,y,z$ fixed.
Its coefficients divided by $1+y$ are bounded, while direct expansion
at $r=1$ gives
\[
 E(1)=-(H+y)(B+z)(1+z).
\]
The right side divided by $1+y$ is negative with nonzero residue.
If $1-r$ were infinitesimal, $E(r)/(1+y)$ would have the same negative
residue, contradicting $n\ge0$.

Consequently the following exact scaled quantities are positive units:
\[
 \frac m{z^2},\quad CL,\quad\alpha(1-H),\quad
 \frac{AH(1+y)}y,\quad Dyj(1+y).
\tag{GE2}
\]
For the IH source subcase $y$ is unbounded, hence $r=O(1/y)$ is
infinitesimal, $j$ is a unit, and $Dy^2$ is a positive unit.
Put $\rho=B/y$, $s=r/\rho$, and $c_s=(1+L)/B$.
Then $s$ is bounded and, with $N=4n+2$ and $\nu=Nr$,
\[
 \frac\nu\rho\longrightarrow2c_s,\quad
 \frac{nr}\rho\longrightarrow\frac{c_s-s}{2},\quad
 \frac D{\rho^2}\longrightarrow\frac{c_s}{2},\quad
 CL\longrightarrow\frac12,\quad
 \frac{Bm}{z^2}\longrightarrow\frac{c_s}{2},
\]
\[
 AH\longrightarrow\frac{1+L}{2},\qquad
 \alpha(1-H)\longrightarrow\frac B2.
\tag{GE3}
\]
Here arrows mean equality of residues of bounded expressions, with
$B,H,L,s$ replaced on the right by their residues. The expressions
remain valid at $H_0=0$, $H_0=1$, or $L_0=0$ whenever allowed by $B_0>0$.
The reason is that the scaled rational functions have denominators
bounded away from zero after the displayed opening factors have
been absorbed; no unscaled $1/H$, $1/(1-H)$ or $1/L$ is asserted bounded.

\label{component:pfb3a83f3849a}
\label{pfb3a83f3849a:sec:common-ii}
\paragraph{The II source corollary.}
\label{sec:ii-source-bounds-refined}
For the original positive C2 source of GE1, suppose $z,y$ are
infinitesimal, without an initial strength assumption on $C$ or $\alpha$. Then $\alpha$ is strong if and only if $B=H+L$
is a positive unit.  On this face $C$ is unbounded, $H$ and
$1-r$ are positive units, and
\[
 \frac{m}{z^2},\quad\frac{A}{y},\quad CL,\quad Dyj,
 \quad\alpha(1-H)\asymp1,\qquad nr=O(1). \tag{DI7}
\]
Here every displayed comparison is in the coefficient field.
Indeed $0<L<yA_c$ makes $L$ infinitesimal and $B$ bounded.
If $B$ is infinitesimal, then $H$ is infinitesimal and GE1 gives
\[
 0<\alpha\le
 \frac{(B+zr)(1+z)(H+y)}{2(1-H)}\in\operatorname{Inf}_c,
\]
using $\delta\ge B+z$ and
$M\le(1+z)(B+zr)(H+y)$.  If $B$ is a unit, then
$H=B-L$ is a unit.  The positive expressions in GE1 show that
$G,\delta,M$ are units and hence that $CL$ and
$\alpha(1-H)$ are positive units.  Thus $C$ is unbounded
and $\alpha$ is strong, so the hypotheses used to prove GE2
hold.  That proof also gives $1-r$ a positive unit, from the
strictly negative value of $E(1)$ and $E(r)\ge0$.
Now GE2 yields the five bounds in DI7 because $H$ and $1+y$
are units.  Finally the exact GE1 identity
\[
 nr=\frac{(B+zr)A_cS-r\delta}{2\delta}
\]
is bounded.  No lower bound on $r$, $j$, or $1-H$ was imposed,
and the argument includes $n=0$.

\section{One small and one unbounded actual angle}
\subsection{Source bounds for weak $C$}
\label{component:p93b05ea6ad1e}

\label{p93b05ea6ad1e:sec:c2-weak-ih-frontier}

Write $z=u_c^2$, $y=v_c^2$, and use the exact contact quantities
\[
 H=v_cY_c=zh,\quad L=v_c(u_cW_c-Y_c)=yt,\quad
 B=H+L=u_cv_cW_c=zb,\quad R=-u_cX_c=zr.
\]
Thus $0<H<1$, $0<R<z$, $L>0$, and
\[
 L<y(1+R),\quad B>yR,\quad
 G=B+R-z(1-H)>0.
\]
Set $S=1+R+L$ and
\begin{align*}
 \widetilde\Delta={}&(1+z)B^2+2(1+z)BR+(1+z)(1+y)B\\
 &+(z-y)R^2+2z(1+y)R+z(1+y),\\
 \widehat N={}&(1+z)(B+R)(H+y)-(1+y)(z-R)L.
\end{align*}
These are exactly $z\Delta,z^2N_{\rm src}$. In particular,
\[
 \widehat N=(1+z)(B+R)H+yG+z\{y(1+R)-L\}+RL(1+y)>0.
\]
The original coefficients, without an approximation, are
\begin{align*}
 C&=\frac{(1+R)\widehat N}{2L\widetilde\Delta},&
 \alpha&=\frac{(B+R)\widehat N}{2(1-H)\widetilde\Delta},\\
 A&=\frac{y(1+R)(B+R)S}{2H\widetilde\Delta},&
 D&=\frac{(1+z)(B+R)(B-yR)S}
 {2\{y(1+R)-L\}\widetilde\Delta},\\
 m&=\frac{z(1+y)(z-R)(1+R)S}{2G\widetilde\Delta},&
 n+\frac12&=\frac{z(B+R)(1+R)S}{2R\widetilde\Delta}.
                                                        \tag{WI1}
\end{align*}

\paragraph{Uniform source reduction for IH.}
Suppose $z$ is positive infinitesimal, $y$ is positive unbounded,
and $n\ge0$. Then $C$ is weak if and only if $B$ is unbounded.
To prove the nontrivial direction, suppose $B$ is bounded. Since
$R<z$ and $G>0$, one has $z\le4(B+R)$: use $B\ge1/2$
or, otherwise, $B+R>z(1-H)>z/2$. Direct division in (WI1)
gives bounded $\widetilde\Delta/[y(B+R)]$. Moreover,
\[
 \frac{\widehat N}{y(B+R)}
 =(1+z)(1+H/y)
 -(1+1/y)\frac{(z-R)L}{B+R}
\]
has residue $1$, since $L\le B$ and $z$ is infinitesimal.
Consequently $C\ge c/L$ for a positive real $c$, which is strong
because $L$ is bounded.

Now let $B$ be unbounded. Put
\[
 \rho=B/y,\qquad e=1+R-L/y>0.
\]
Here $0<\rho<1+R+H/y$, so $\rho$ is bounded and its residue
lies in $[0,1]$. The following exact normalized quantities have
residue $1$:
\[
 \frac{\widetilde\Delta}{B(B+y)},\qquad
 \frac{\widehat N}{yB},\qquad
 \frac{L}{B},\qquad \frac{S}{B},\qquad\frac{G}{B}.
                                                        \tag{WI2}
\]
They give
\[
 n+\tfrac12=\frac{\rho}{2r(1+\rho)}(1+o(1)).
\]
The inequality $n\ge0$ therefore gives $r/\rho$ bounded and
$r\le\rho/[1+\rho](1+o(1))\le1/2+o(1)$.
In particular $1-r$ is a positive unit, and
$R/\rho=z(r/\rho)$ is infinitesimal. The normalized leading
relations, uniform at arbitrary openings $H,1-H,e$, are
\begin{align*}
 C&=\frac{1+o(1)}{2B(1+\rho)},&
 A&=\frac{B(1+o(1))}{2H(1+\rho)},\\
 \alpha&=\frac{B(1+o(1))}{2(1-H)(1+\rho)},&
 m&=\frac{z^2(1-r)(1+o(1))}{2B(1+\rho)},\\
 D&=\frac{\rho^2(1+o(1))}{2e(1+\rho)}.
                                                        \tag{WI3}
\end{align*}
Here each $1+o(1)$ means a positive coefficient unit of residue one;
it never denotes a comparison with a testing infinitesimal. In
particular $A,\alpha$ are unbounded, $m,C$ are infinitesimal, and
$Cm$ is infinitesimal for every weak-$C$ IH source. Thus the four
reciprocal-product exit can never admit this entire directed leaf.
If $\rho$ is infinitesimal then $e$ is a unit and $D$ is weak;
if $\rho$ is a unit then $D$ is strong. Literal $n=0$ remains
allowed throughout the reduction. No lower bound on $r$ was used.

\subsection{Interior double corner}
\label{component:pc00b84929b12}

\label{pc00b84929b12:sec:c2-weak-ih-interior-corner}

This proof covers the stated local testing chart; the remote commuting
charts are treated below.

\begin{theorem}[Full-parent IH corner, interior shape]
Use the exact positive C2 source with $d=1$, the omitted original
$\lambda,\beta$ coefficients zero, and $n\ge0$. Suppose its actual
angles have type IH, and $C$ is weak. Use the source notation of
\S\ref{p93b05ea6ad1e:sec:c2-weak-ih-frontier}; thus $z$ is infinitesimal and
$B$ unbounded. Write $\rho=B/y$ (a shape parameter, not an event
coefficient). Assume
\[
 \rho,\quad1-\rho
 \quad\hbox{are positive bounded units}.                 \tag{IC1}
\]
At every independent testing pair for which $u$ and the homogeneous
reciprocal coordinate $w=1/v$ are infinitesimal or zero, the entire
original parent has a separated-field tensor-SOS certificate, for
all real or complex state coordinates. The source openings $H,1-H$ and the positive radial parameter $r$
are arbitrary; $n$ may be unbounded. Literal
testing $u=0$ and $v=\infty$, either testing sign, and literal source
$n=0$ are included.
No state coordinate or certificate root is required to be bounded.
\end{theorem}

\begin{proof}
Put $N=4n+2$ and $\nu=Nr$. Source reduction (WI1)--(WI3)
shows that $\nu$ and $1-r$ are bounded positive units, even when
$N$ is unbounded. In residue,
\[
 \nu_0=\frac{2\rho_0}{1+\rho_0},\qquad 2r_0\le\nu_0,
 \qquad D_0=\frac{\rho_0^2}{2(1-\rho_0^2)}.             \tag{IC2}
\]
All assumptions and divisions in this sentence concern the coefficient
field alone.
The uniform coefficient estimates needed below are
\[
 BC\asymp1,\quad \frac{Bm}{z^2}\asymp1,\quad
 \frac{AH}B\asymp1,\quad\frac{\alpha(1-H)}B\asymp1,
 \quad D\asymp1,\quad nr\le\nu/4.
\]
In particular $C,m$ are infinitesimal, $nr,nr^2$ are bounded,
$mB/z\asymp z$, and $m\sqrt B/z\asymp z/\sqrt B$.
These estimates are uniform in $r,H,1-H$ and do not assert that
$A/B$ or $\alpha/B$ is bounded at an opening.

The homogeneous testing parent is exactly
\begin{align*}
 \widehat F={}&4nD_uD_wX^2+4mM^2+4CD_w(uW-Y)^2
     +4AD_uD_wY^2+4DJ^2+4\alpha D_u(Y-wZ)^2\\
 &+2\{J^2+D_u(Y-wZ)^2+D_w(X+uZ)^2-D_uD_wZ^2\},\\
 M={}&W-wX+uY-uwZ,\qquad
 J=uwW+uX-wY-Z,
 \qquad D_u=1+u^2,\quad D_w=1+w^2.
                                                               \tag{IC3}
\end{align*}
For $w\ne0$ it is $w^2F(u,1/w)$; (IC3) defines the literal
projective endpoint as well. Let $G$ be (IC3) restricted to the graph
\[
                         X=-ruZ,\qquad Y=HwZ.
\tag{IC4}
\]
Use the following five scalar roots, all defined using only pure
coefficient radicals and inverses:
\begin{align*}
 R_0&=(uW-BwZ)/\sqrt B,&
 R_1&=(zW-BuwZ)/\sqrt B,\\
 R_2&=uwW-\rho Z,&
 T&=(u^2-z)Z,&
 E&=\sqrt{BH}\,(w^2-1/y)Z.                               \tag{IC5}
\end{align*}
They vanish at every signed exact source contact. Define the pure
coefficient quantities
\begin{align*}
 K={}&4AH^2+8DHr+4\alpha(1-H)^2
      +4m(H+r-1)^2+2(H+r-1)^2+4nr^2,\\
 k={}&K-4B^2m/z^2+8znr^2+8AH^2/y.                      \tag{IC6}
\end{align*}
The exact source formulas and state/angular stationarity give the
following finite polynomial Gram identity:
\begin{align*}
G={}&\left(4BC-\frac{kz}{B}\right)R_0^2
     +\frac{4Bm}{z^2}R_1^2+2(2C+2D+1)R_2^2\\
 &+r^2\{4nw^2+2(2D+2n+1)\}T^2
  +\frac H B\{4Au^2+2(2A+2C+2D+1)\}E^2\\
 &+\frac{ku}{B}R_0R_1
  +\left\{-4r(2D+1)+\frac{k}{B}\right\}R_2T
  -\frac{4\sqrt H(2C+2D+1)}{\sqrt B}R_2E.                   \tag{IC7}
\end{align*}
For clarity, the only non-obvious algebraic step is the exact toric
conductor
\[
 B(u^2-z)(w^2-1/y)Z^2
 =(u^2-z)ZR_2-zR_0^2+uR_0R_1.                         \tag{IC8}
\]
Before (IC8), the last mixed radial term is
$k(u^2-z)(w^2-1/y)Z^2$. Expanding (IC7), its difference from
$G$ has only the four monomials $uwWZ,u^2Z^2,w^2Z^2,Z^2$;
the unchanged original source sets all four coefficients exactly to
zero. 

All coefficients of (IC7) are bounded: $BC$, $Bm/z^2$, $AH/B$,
$\alpha(1-H)/B$, $nr,D$ and $k/B$ are bounded. No inverse of $r$
is used, and $1/\sqrt N$ is bounded because $N\ge2$. Moreover,
\[
 (k/B)_0=\frac{2r_0}{1+\rho_0}.
\]
Add to the right side of (IC7) the following exact zero, obtained
by one more toric conductor:
\[
 \rho T^2+(u^2-z)TR_2
       -u(zR_0-uR_1)(uR_0-R_1)=0.                    \tag{IC8a}
\]
The new $T^2$ term has positive unit residue $\rho_0$. Every other
coefficient introduced by (IC8a) is infinitesimal. This step is what
makes the proof uniform when $r$ is infinitesimal and $n$ unbounded.

Suppress residue bars temporarily. The $R_0,R_1,E$ diagonal
residues are respectively
\[
 \frac2{1+\rho},\qquad
 \frac{2(1-r)}{1+\rho},\qquad
 \frac2{1+\rho},
\]
which are positive. The only surviving off-diagonal block is
\[
 (R_2,T):\qquad
 \begin{pmatrix}
 \dfrac2{1-\rho^2}&-\dfrac r{1-\rho}\\[3pt]
 -\dfrac r{1-\rho}&\rho+r\nu+4Dr^2
 \end{pmatrix}.                                      \tag{IC9}
\]
Using (IC2), its determinant is exactly
\[
 \frac{2\rho}{1-\rho^2}
  +\frac{r\{4\rho-r(3\rho+1)\}}
       {(1-\rho)(1+\rho)^2}>0.
\]
Indeed $0\le r\le\rho/(1+\rho)$, and consequently
$4\rho-r(3\rho+1)\ge\rho(3+\rho)/(1+\rho)>0$.
The first determinant summand is already strictly positive, including
$r_0=0$. Thus (IC7) after the exact zero (IC8a) has strictly positive
definite residue, including literal $n=0$ and unbounded $n$.

It remains to restore the full state normals; the graph certificate
alone is not the claimed conclusion. Put
\[
 E_X=\sqrt N(X+ruZ),\qquad E_Y=\sqrt{B/[H(1-H)]}(Y-HwZ).
\]
We give exact factorizations of both linear normal terms. Let
$U=u^2,V=w^2,s=1/y$ and define $L_X,L_Y$ by
$4L_X=\partial_X\widehat F|_{\rm graph}$ and
$4L_Y=\partial_Y\widehat F|_{\rm graph}$. Direct expansion gives
\begin{align*}
 L_X={}&(2D+1)uR_2-\frac{2m\sqrt B}{z}wR_1+uZP_X(U,V),\\
 P_X={}&(2D+1)\rho-\frac{2mB}{z}V
    -2nr(1+U)(1+V)-r(2D+1)U-2D-r\\
 &+V\{1-H-r-2DH-2m(H+r-1)\}.                         \tag{IC10}
\end{align*}
Every coefficient of $P_X$ is bounded, and source stationarity gives
$P_X(z,s)=0$. Write its exact monic division as
\[
 P_X=(U-z)f_X(V)+(V-s)g_X.
\]
Both $f_X,g_X$ have bounded coefficients. Explicitly,
\begin{align*}
 f_X(V)&=-2nr(1+V)-r(2D+1),\\
 g_X&=-2mB/z-2nr(1+z)+1-H-r-2DH-2m(H+r-1).
\end{align*}
Thus no inverse of $r$ or of either opening is introduced. We obtain
\[
 L_X=\left(2D+1+\frac{g_X}B\right)uR_2
       -\frac{2m\sqrt B}{z}wR_1
       +u f_X(V)T
       -\frac{uwg_X}{\sqrt B}R_0.                       \tag{IC11}
\]
Here we used the exact identity
$B(V-s)Z=R_2-\sqrt B\,wR_0$; this avoids dividing by $\sqrt H$
when $H$ is infinitesimal. Every displayed multiplier is infinitesimal
in the bounded tensor ring; in particular $m\sqrt B/z\asymp z/\sqrt B$.

For the other normal, put
\begin{align*}
 f(V)={}&2m-2C(1+V)-(2D+1)V,\\
 P_Y(U,V)={}&Bf(V)+2AH(1+U)(1+V)+2CH(1+V)
       +2DHV+2DrU+2D\\
 &-2\alpha(1-H)(1+U)+2m(H+r-1)U
       +(H+r-1)U+H(V+1).
\end{align*}
Then
\[
 L_Y=\sqrt B f(V)R_0+wZP_Y(U,V),\qquad P_Y(z,s)=0.
\]
The coefficients of $P_Y/B$ are bounded. Write
$P_Y=(U-z)f_Y(V)+(V-s)g_Y$. The two quotient coefficients are explicitly
\begin{align*}
 f_Y(V)&=2AH(1+V)+2Dr-2\alpha(1-H)
                   +(2m+1)(H+r-1),\\
 g_Y&=-B(2C+2D+1)+2AH(1+z)+2CH+2DH+H.
\end{align*}
The displayed source estimates show directly that $f_Y/B$ and
$g_Y/B$ have bounded coefficients, including both opening limits.
The elementary exact identity
\[
          (U-z)\sqrt B\,wZ=zR_0-uR_1
\]
therefore gives
\[
 \frac{L_Y}{\sqrt{B/[H(1-H)]}}
 =\sqrt{H(1-H)}\left\{f(V)+\frac{z f_Y(V)}B\right\}R_0
       -\frac{\sqrt{H(1-H)}u f_Y(V)}B R_1
       +\frac{\sqrt{1-H}w g_Y}B E.          \tag{IC12}
\]
Again every displayed multiplier is infinitesimal. This argument
retains the full $X,Y$ linear terms; none is estimated on the graph
and then discarded.

Finally the normal quadratic block, in $(E_X,E_Y)$, is bounded.
Its residue is diagonal with entries
\[
             1,\qquad \frac2{1+\rho_0}>0.
\]
Indeed the $X^2$ coefficient before scaling is
$4nD_uD_w+4mw^2+(4D+2)u^2+2D_w$; the $Y^2$ coefficient is
$4CD_w+4AD_uD_w+(4D+2)w^2+(4\alpha+2)D_u+4mu^2$.
Their mixed coefficient is $-4uw(2m+2D+1)$.
For the $Y$ normalization, if $L=B/[H(1-H)]$, then $A/L$ and
$\alpha/L$ are bounded and
\[
 \frac{4(A+\alpha)}L
 =4\left\{\frac{AH}B(1-H)
               +\frac{\alpha(1-H)}B H\right\}
 \longrightarrow\frac2{1+\rho_0}.
\]
Thus the stated normal residue remains strict even when $H_0=0$
or $H_0=1$; no denominator involving a vanishing opening was dropped.
Equations (IC11)--(IC12) show that all cross terms from these normals
to the five scalar roots have zero residue. Thus the exact seven-root
Gram for the entire (IC3) is bounded and has positive definite
residue. The bounded positive-residue Gram lemma yields tensor SOS.
Its roots may be unbounded; it uses no mixed inverse or radical.
The homogeneous chart and real/imaginary splitting give all endpoints
and states stated in the theorem.
\end{proof}

This calculation proves the interior $0<\rho_0<1$ testing corner. Section~\ref{component:p18dbedfff317} treats the source endpoints; Sections~\ref{component:p0fb32cd2f1bc} and~\ref{component:pf95b1c5fc32f} prove the other testing charts.

\subsection{The two source endpoints}
\label{component:p18dbedfff317}

\label{p18dbedfff317:sec:weak-ih-endpoints}
We use the exact parent, graph, roots, coefficients $K,k$, and both
normal factorizations (IC3)--(IC12) in the preserved interior-corner
proof.  Those identities are exact for every positive canonical C2
source; their validity does not require the interior-shape assumption.
Here $z$ and testing $u,w$ are infinitesimal (zero testing coordinates
are allowed), $B$ is unbounded, $H\in(0,1)$, $\rho=B/y$, and
$N=4n+2$, $\nu=Nr$. All source and testing rates are independent.

\paragraph{The endpoint $\rho_0=0$.}
Put $s=r/\rho$ and $\tau=1/B$. Nonnegativity of $n$ and the exact
source formulas imply that $s$ is bounded. The following normalized
source quantities have the indicated residues:
\[
 \nu/\rho\to2,\quad nr/\rho\to(1-s)/2,\quad
 D/\rho^2\to1/2,\quad BC,Bm/z^2,AH/B,\alpha(1-H)/B\to1/2,
\]
\[
 k/(B\rho)\to2s,\qquad 0\le s_0\le1.
\tag{WIE1}
\]
These are uniform opening identities: after substituting
$B=1/\tau$, $y=1/(\tau\rho)$, $r=s\rho$, their normalized rational
functions have a nonzero denominator at $z=\tau=\rho=0$, including
$H=0,1$. 
Thus they concern simultaneous, not iterated, infinitesimal limits.

Replace $T$ by $T_\rho=\rho T$ in (IC7). Its coefficients are bounded;
for example $r^2n/\rho^2=s(nr/\rho)$ is bounded. Add the exact zero
\[
 T_\rho^2+(u^2-z)T_\rho R_2
 -\rho u(zR_0-uR_1)(uR_0-R_1)=0.
\tag{WIE2}
\]
This is (IC8a) multiplied by $\rho$, and uses no mixed inverse.
The residue Gram has $R_0,R_1,E$ diagonal entries $2,2,2$, and
\[
 (R_2,T_\rho):\quad
 \begin{pmatrix}2&-s_0\\-s_0&1+2s_0\end{pmatrix}.
\]
Its determinant $2+4s_0-s_0^2\ge2$ is strictly positive on $[0,1]$.

The entire normal restoration survives this rescaling. In (IC11),
replace $f_XT$ by $(f_X/\rho)T_\rho$, where
\[
 f_X/\rho=-2(nr/\rho)(1+w^2)-s(2D+1)
\]
is bounded. Every multiplier still contains testing $u$ or $w$, and
$1/\sqrt N$ is bounded. The remaining coefficients in (IC11) are
unchanged and infinitesimal. Formula (IC12) uses the root $R_1$ rather
than $T$ and remains bounded and infinitesimal unchanged. The normal
quadratic residue is $\operatorname{diag}(1,2)$. Hence the full
seven-root Gram is strict. This includes arbitrary positive $r,H,1-H$
and literal $n=0$.

\paragraph{The endpoint $\rho_0=1$.}
Set $e=1+zr-\rho+H/y=j>0$, an infinitesimal, and write
\[
 \rho=\frac{1+zr-e}{1-H\tau},\qquad \tau=1/B.
\]
The exact source formulas give the simultaneous normalized residues
\[
 De\to1/4,\quad\nu\to1,\quad nr\to(1-2r)/4,\quad
 BC\to1/4,\quad Bm/z^2\to(1-r)/4,
\]
\[
 AH/B,\alpha(1-H)/B\to1/4,\quad
 0\le r_0\le1/2.
\tag{WIE3}
\]
Put $\kappa=k-8DHr$; then $\kappa/B\to r$.
The scaled rational functions have regular nonzero denominators at
$z=\tau=e=0$ after the displayed powers of $e$ are absorbed.

The large coefficient $D$ must be retained in a square, not treated
as a bounded Gram coefficient. Introduce the full-state root
\[
 F_D=\sqrt D\,(J+eZ),\qquad a=\sqrt{H/B},\qquad
 Q=R_2-rT-aE.
\]
On the graph, $Q=J+eZ$. Subtracting its square from (IC7) gives the
exact identity
\begin{align*}
 G-4DQ^2={}&(4BC-\kappa z/B)R_0^2+4BmR_1^2/z^2
       +2(2C+1)R_2^2\\
 &+r^2(4nw^2+4n+2)T^2
 +\frac HB\{4A(1+u^2)+4C+2\}E^2\\
 &+\kappa uR_0R_1/B+(-4r+\kappa/B)R_2T
       -4a(2C+1)R_2E.
\tag{WIE4}
\end{align*}
Add the exact zero (IC8a). Restore the full normals
\[
 E_X=\sqrt N(X+ruZ),\qquad
 E_Y=\sqrt{B/[H(1-H)]}(Y-HwZ).
\]
Use the exact relation
\[
 R_2=F_D/\sqrt D+rT+aE-uE_X/\sqrt N
                         +wE_Y/\sqrt{B/[H(1-H)]}.
\tag{WIE5}
\]
The inverse $1/\sqrt D$ is purely in the source field and infinitesimal.
All coefficients after this substitution are bounded.

For completeness, the first-normal quotients of
$\widehat F-4F_D^2$ are precisely the D-removed versions of
(IC11)--(IC12). Explicitly,
\begin{align*}
 \bar f_X&=-2nr(1+V)-r,\\
 \bar g_X&=-2mB/z-2nr(1+z)+1-H-r-2m(H+r-1),\\
 \bar f&=2m-2C(1+V)-V,\\
 \bar f_Y&=2AH(1+V)-2\alpha(1-H)+(2m+1)(H+r-1),\\
 \bar g_Y&=-B(2C+1)+2AH(1+z)+2CH+H.
\end{align*}
The constant term in the $Y$ monic polynomial acquires $2De$, but
its monic quotient coefficients are unchanged; the $X$ polynomial
has the constant correction $-2De$. Exact source stationarity
still sets the constant remainders to zero because $J+eZ$ vanishes
at the contact. Thus the displayed quotient formulas reconstruct
both linear normals. Their normalized coefficients are infinitesimal
by (WIE3); in particular $De$ is bounded and no bare $D$ survives.
The normal quadratic block of the D-removed expression has residue
$\operatorname{diag}(1,1)$.

In the order $(R_0,R_1,T,E,F_D,E_X,E_Y)$, the complete residue Gram
is diagonal:
\[
 \operatorname{diag}\{1,1-r_0,1+r_0-r_0^2,1,4,1,1\}.
\tag{WIE6}
\]
For instance the $T^2$ entry comes from substituting $R_2=r_0T$
in (WIE4): it is $r_0\nu_0-r_0^2$, plus the unit contribution from
(IC8a). Every entry is positive for $0\le r_0\le1/2$.
This proves full-parent tensor SOS at the corner also when
$1-\rho$ is infinitesimal. Along with the preserved interior proof
and the $\rho_0=0$ argument, every allowed source shape is covered
at the independent testing corner $u,w$ infinitesimal.

\subsection{The open testing axis}
\label{component:pf95b1c5fc32f}

\label{pf95b1c5fc32f:sec:weak-ih-open-axis}
\paragraph{Domain and source input.}
The source and testing real closed fields are independent.  Let $B_c,B_t$
be their rings of bounded elements and use $B=B_c\otimes_{\mathbb R}B_t$
for Gram coefficients.  Source square roots and testing square roots are
always taken in their respective fields.  In the weak-$C$ IH source chart,
write $z=u_c^2$ infinitesimal, $y=v_c^2$ unbounded,
$H=zh\in(0,1)$, $B_*=zb$ unbounded and $\rho=B_*/y$ bounded.
The source entry calculation gives
\[
 0\leq\rho_0\leq1,\qquad
 (B_*C)_0=\left(\frac{AH}{B_*}\right)_0
 =\left(\frac{\alpha(1-H)}{B_*}\right)_0
 =c:=\frac1{2(1+\rho_0)}.
\tag{OA1}
\]
In particular $C$ is infinitesimal and $A,\alpha$ are unbounded.
For $0<\rho_0<1$,
\[
 D_0=\frac{\rho_0^2}{2(1-\rho_0^2)};
\tag{OA2}
\]
for $\rho_0=0$, $D$ is positive infinitesimal; for $\rho_0=1$, $D$ is
unbounded.  These statements follow from
$D\sim\rho^2/[2e(1+\rho)]$ and $e_0=1-\rho_0$.
No lower bound on $n\geq0$ is imposed.

\paragraph{Statement.}
For every testing $u$ with nonzero finite residue and every infinitesimal
$w=1/v$, including $w=0$, the complete reciprocal-$v$ canonical C2 parent
is a tensor sum of squares for every state.  The same assertion holds for
complex states by applying the real quadratic form to real and imaginary
parts separately.

\paragraph{Exact parent and retained squares.}
Put $D_u=1+u^2$, $D_w=1+w^2$ and
\[
 M=W-wX+uY-uwZ,\qquad J=uwW+uX-wY-Z.
\]
The parent is
\begin{align*}
\widehat F={}&4nD_uD_wX^2+4mM^2+4CD_w(uW-Y)^2
 +4AD_uD_wY^2+4DJ^2+4\alpha D_u(Y-wZ)^2\\
&+2\{J^2+D_u(Y-wZ)^2+D_w(X+uZ)^2-D_uD_wZ^2\}.
\end{align*}
Retain the entire $n,m$ positive square terms. When $\rho_0=0$, retain the entire
$D$ positive square term as well and set $d=0$. When $0<\rho_0<1$, set $d=D$.
When $\rho_0=1$, choose a sufficiently large positive real $d$ below $D$
and retain $4(D-d)J^2$.  Denote the remainder by $F_*$.  All retained
terms are tensor squares or sums of tensor squares.

Use the graph $X=0$, $Y=HwZ$, and define
\[
 P=\frac{uW}{\sqrt{B_*}},\quad E=\sqrt{B_*}\,wZ,
 \quad V=uwW,\quad
 E_Y=\sqrt{\frac{B_*}{H(1-H)}}(Y-HwZ).
\tag{OA3}
\]
Crucially, the exact root relation is
\[
 PE=VZ.\tag{OA4}
\]
On the graph, the $C$, $A,\alpha$, $d$ positive square terms are respectively
\begin{align*}
 &4B_*CD_w(P-HE/B_*)^2,\\
 &\frac{4D_u}{B_*}\{AH^2D_w+\alpha(1-H)^2\}E^2,\\
 &4d\{V-(1+Hw^2)Z\}^2.
\end{align*}
The remaining core is exactly
\[
 2V^2-4(1+Hw^2)PE
 +\frac{2}{B_*}\{H^2+u^2(1-H)^2+H^2w^2\}E^2.
\tag{OA5}
\]
Thus every graph coefficient is bounded.  The last $E^2$ coefficient is
infinitesimal, uniformly for all $0<H<1$.

\paragraph{Restoring the $Y$ normal.}
Let $L=B_*/[H(1-H)]$ and $Y_*=HwZ$.
The first normal, divided by four, is
\begin{align*}
 L_Y:=\frac14\partial_YF_*\big|_{X=0,Y=Y_*}
={}&2CD_w(Y_*-uW)+2D_u\{AH D_w-\alpha(1-H)\}wZ\\
&-2dw\{V-(1+Hw^2)Z\}
 -w\{V-(1+Hw^2)Z\}+D_u(H-1)wZ.
\tag{OA6}
\end{align*}
Every coefficient of $L_Y/\sqrt L$ in $(P,E,V,Z)$ is infinitesimal
when $d$ is bounded, including $d=0$.  Indeed the only potentially finite contribution is
\[
 2D_u\frac{AH-\alpha(1-H)}{B_*}\sqrt{H(1-H)}E,
\]
and it has zero residue by (OA1); the additional $w^2$ contribution is
infinitesimal.  The other terms contain $C$, $w/\sqrt L$ or $1/B_*$.
The quadratic $E_Y$ coefficient has residue $4cD_{u,0}>0$, and the
$X E_Y$ coefficient has zero residue.  These statements remain valid
when $H$ or $1-H$ is infinitesimal.  No inverse of $d$, $H$ or $1-H$ enters a Gram coefficient
unaccompanied by its compensating factor.

\paragraph{The strict residue when $0<\rho_0<1$.}
Write $s=\sqrt{1+u_0^2}>1$ and $r=\rho_0$.
Using (OA4), split the core cross term as
$-4PE=-4aPE-4(1-a)VZ$, where $a$ is any real number satisfying
\[
 \frac{s}{s+r}<a<\min\left\{\frac{s}{1+r},\frac{s}{s-r}\right\}.
\tag{OA7}
\]
The interval is nonempty: both upper endpoints exceed $s/(s+r)$.
The residue Gram splits into the positive $E_Y$ pivot, a $(P,E)$ block
\[
 G_{PE}=\begin{pmatrix}2/(1+r)&-2a\\-2a&2s^2/(1+r)\end{pmatrix},
\]
and a $(V,Z,X)$ block
\[
 G_{VZX}=\begin{pmatrix}
 2+4d_0&-4d_0-2(1-a)&2u_0(1+2d_0)\\
 -4d_0-2(1-a)&4d_0&-4d_0u_0\\
 2u_0(1+2d_0)&-4d_0u_0&2s^2+4d_0u_0^2
 \end{pmatrix}.
\tag{OA8}
\]
The $(V,X)$ principal block is positive definite.  Its Schur complement
is positive exactly when
\[
 \frac{s}{s+r}<a<\frac{s}{s-r},
\]
as direct substitution of $d_0=r^2/[2(1-r^2)]$ shows.  Hence (OA7)
makes both blocks strictly positive.  This checks the entire $X$ normal,
not merely the scalar graph.

\paragraph{The endpoint $r=1$.}
For finite $d>0$, positivity of the $(V,Z,X)$ block is equivalent to
$a_-<a<a_+$, where
\[
 a_\pm(d)=
 \frac{(2d+1)s^2\pm s\sqrt{2d(2d+1)}}{(2d+1)u_0^2+1}.
\]
As $d\to\infty$, the lower endpoint tends to $s/(s+1)$ and the upper
endpoint to $s/(s-1)$.  Choose first a real $a$ strictly between
$s/(s+1)$ and $\min\{s/2,s/(s-1)\}$, then a sufficiently large finite
$d$.  The interval is nonempty because $s>1$.  Since source $D$ is
unbounded this cap is admissible, and the full residue is strict.

\paragraph{The endpoint $r=0$.}
Take $a=1$, $d=0$, having retained the whole $D$ positive square term, and omit
$Z$ from the root list. The residue is the direct sum of
\[
 2(P-E)^2+2u_0^2E^2,
 \qquad 2(V+u_0X)^2+2X^2,
 \qquad 2D_{u,0}E_Y^2.
\tag{OA9}
\]
It is strictly positive. The exact rewriting has bounded coefficients:
(OA5) uses only $P,E,V$, and the first $X$ normal from the core is
$uV+u(1-H)w^2Z=uV+u(1-H)wE/\sqrt{B_*}$.
Thus no unscaled $Z$ or inverse of $D$ remains.

\paragraph{Conclusion.}
In all three cases the original remainder has a bounded Gram matrix with
strictly positive real residue.  The bounded-Gram lifting lemma therefore
makes it a tensor sum of squares.  Adding the retained source-positive
positive square terms proves the assertion.  The argument is valid for arbitrary $n\geq0$,
including $n=0$, and literal projective $w=0$; no comparison between source
and testing infinitesimals was used.

\subsection{Complementary tests for weak $C$}
\label{component:p0fb32cd2f1bc}

\label{p0fb32cd2f1bc:sec:weak-ih-remote}
Assume the weak-$C$ IH source entry, with no restriction on
$\rho_0\in[0,1]$, on positive $r,H,1-H$, or on $n\ge0$.
The uniform source inputs are $B$ unbounded, $C,m$ infinitesimal,
$BC,AH/B,\alpha(1-H)/B$ positive units with common residue
$c=1/[2(1+\rho_0)]$, and $A,\alpha$ unbounded.
Every square retained below has nonnegative pure-source coefficient;
products $1+u^2$ and $1+v^2$ are themselves testing sums of squares.

\paragraph{Finite $v$, infinitesimal $u$.}
Write $Y_A=\sqrt A Y$ and $Z_\alpha=\sqrt\alpha(vY-Z)$, so
$Y=Y_A/\sqrt A$ and $Z=vY_A/\sqrt A-Z_\alpha/\sqrt\alpha$.
Retain the complete $n,C,D$ positive square terms. Use the five roots
\[
 U=uW,\quad X,\quad V=\sqrt m\,vW,\quad Y_A,\quad Z_\alpha.
\]
The $m$ positive square term is exactly
$4(V-\sqrt m X+u\sqrt{m/\alpha}Z_\alpha)^2$.
All coefficients in the remaining core are bounded after substituting
$Y,Z$, and all coefficients involving $Y$ or $Z$ are infinitesimal.
The full residue Gram is diagonal:
\[
 2U^2+2(1+v_0^2)X^2+4V^2
       +4(1+v_0^2)Y_A^2+4Z_\alpha^2.
\tag{IR1}
\]
This includes literal $u=0$ or $v=0$; redundant zero roots cause no
problem for a positive Gram representation.

\paragraph{Finite $v$, unit $u$.}
Retain $n,m,C,D$. Use $(W,X,Y_A,Z_\alpha)$.
The residue of the core and selected $A,\alpha$ positive square terms is
\[
 2\{u_0^2(W+v_0X)^2+(1+v_0^2)X^2\}
       +4(1+u_0^2)(1+v_0^2)Y_A^2
       +4(1+u_0^2)Z_\alpha^2,
\tag{IR2}
\]
which is strictly positive because $u_0\ne0$.

\paragraph{Finite $v$, unbounded $u$.}
Use the homogeneous coordinate $p=1/u$, infinitesimal or zero.
Retain $n,m,D$ and use
\[
 P=pX,\quad Q=W+vX,\quad S=\sqrt C W,\quad Y_A,\quad Z_\alpha.
\]
The $C$ positive square term becomes
$4(1+v^2)(S-p\sqrt{C/A}Y_A)^2$.
The entire selected expression has bounded coefficients, and residue
\[
 2Q^2+2(1+v_0^2)P^2+4(1+v_0^2)S^2
       +4(1+v_0^2)Y_A^2+4Z_\alpha^2.
\tag{IR3}
\]
Thus every finite-$v$ chart is complete, independently of $n$.

\paragraph{Both testing angles unbounded.}
Let $p=1/u$, $w=1/v$ be infinitesimal or zero and put
$D_p=1+p^2$, $D_w=1+w^2$.
Retain the entire $n,m$ positive square terms. Set $d=\min(1,D)>0$ in the source
field and retain $4(D-d)J^2$. The remainder is exactly
\begin{align*}
 F_*={}&4CD_w(W-pY)^2+4AD_pD_wY^2+4dJ^2
       +4\alpha D_p(Y-wZ)^2\\
 &+2\{J^2+D_p(Y-wZ)^2+D_w(pX+Z)^2-D_pD_wZ^2\},\\
 J={}&wW+X-pwY-pZ.
\end{align*}
Use the graph $X=-wW$, $Y=HwZ$, the normals
$E_X=X+wW$, $E_Y=\sqrt L(Y-HwZ)$ with
$L=B/[H(1-H)]$, and scalar roots
\[
 S=\sqrt C W,\quad E=\sqrt B\,wZ,\quad
 U=pwW,\quad P=\sqrt d\,pZ.
\tag{IR4}
\]
On the graph the exact expression is
\begin{align*}
 F_*|_{\rm graph}={}&4D_w(S-pH\sqrt{C/B}E)^2
 +\frac{4D_p}{B}\{AH^2D_w+\alpha(1-H)^2\}E^2\\
 &+4(1+Hw^2)^2P^2+2D_wU^2
 -\frac{4pD_w}{\sqrt{BC}}SE
 +\frac2B\{(1-H)^2+H^2p^2D_w\}E^2.
\tag{IR5}
\end{align*}
All coefficients are bounded and its residue is
$4S^2+4cE^2+2U^2+4P^2$.

The first $X$ normal, divided by four, is exactly
\[
 L_X=-pD_wU+\frac{p(1-H)w}{\sqrt B}E
               -2\sqrt d(1+Hw^2)P.
\tag{IR6}
\]
The first $Y$ normal is
\begin{align*}
 L_Y={}&-2p\sqrt C D_w S+\frac{E}{\sqrt B}
 \{2Cp^2HD_w+2D_p(AHD_w-\alpha(1-H))\\
 &\hspace{38mm}+(2d+1)p^2(1+Hw^2)-D_p(1-H)\}.
\tag{IR7}
\end{align*}
After dividing by $\sqrt L$, all its coefficients are infinitesimal;
the only possible nonzero term is proportional to
$(AH-\alpha(1-H))/B$, whose residue is zero.
The quadratic normal coefficients have residues
$2+4d_0$ and $4c$, and their mixed coefficient has zero residue.
The only nonzero scalar-normal coupling is the $(P,E_X)$ block
\[
 \begin{pmatrix}4&-4\sqrt{d_0}\\-4\sqrt{d_0}&2+4d_0\end{pmatrix},
\tag{IR8}
\]
whose determinant is exactly $8$. Thus the full Gram is strict even
when $D$ is infinitesimal and even when $n=0$. The normalization uses
$\sqrt d$ but never $1/\sqrt d$ in a Gram coefficient.

Each displayed Gram is a representation of the complete selected
state quadratic form, not just its graph restriction. The bounded
positive-residue lemma and the retained squares prove tensor SOS.
These charts, the open-axis theorem, and the three local source-shape
proofs cover every independent testing pair.

\subsection{Local completion for strong $C$}
\label{component:pf3e0ca8bf054}

\label{pf3e0ca8bf054:sec:strong-ih-local}
Assume the strong IH source entry of \S\ref{pf5726218aec8:sec:strong-angle-entry}:
$z$ is infinitesimal, $y$ unbounded, $B=H+L$ a bounded positive unit,
$0<H<1$, $L>0$, and $C,\alpha$ strong. The source parameter
$\rho=B/y$ is infinitesimal, $s=r/\rho$ bounded, and $c_s=(1+L)/B$
is a positive unit. Testing $u,w=1/v$ are infinitesimal or zero.
We use the exact graph, five roots and both full normal identities
(IC3)--(IC12) of the preserved corner calculation; none of those
algebraic identities assumes $C$ weak.

The source quantities obey (GE3), and direct rational simplification
gives additionally
\[
 \frac{k}{B\rho}\longrightarrow
 \frac{2(B^2s+BHs+BH-Bs-H^2+H)}{B^2}.
\tag{SL1}
\]
No opening denominator occurs in this normalized expression.

Set $T_\rho=\rho T$ and
$\eta=\sqrt{H(1-H)/B}$, so $Y=HwZ+\eta E_Y$.
The full Gram's entire explicit dependence on $C$ is exactly
\[
 4C(U_1^2+U_2^2),\quad
 U_1=uW-Y-LwZ=\sqrt B R_0-\eta E_Y,
\]
\[
 U_2=w(uW-Y)-(L/y)Z
       =R_2-\sqrt{H/B}E-w\eta E_Y.
\tag{SL2}
\]
The scalar and normal identities verifying this assertion use
\[
 \sqrt H R_2-\sqrt{BH}\,wR_0-\sqrt B E=0.
\tag{SL3}
\]
Choose a positive real cap $c_*\le C$ and retain
$4(C-c_*)(U_1^2+U_2^2)$. In the remaining exact Gram replace explicit
$C$ by $c_*$, leaving $k$ and all other original source quantities
unchanged. The distinction matters: this is an exact square
subtraction from the complete Gram, not a perturbation of the source.
Add $\lambda$ times the exact zero (WIE2), where $\lambda>0$ is a
real number to be chosen. All resulting coefficients are bounded,
using (GE3), (SL1), and the opening-normalized roots.

We display the entire residue. The roots $R_1$ and $E_X$ have isolated
positive pivots $2c_s$ and $1$. The $(R_0,E_Y)$ block is
\[
 \begin{pmatrix}
 4Bc_*&-4c_*\sqrt{H(1-H)}\\
 -4c_*\sqrt{H(1-H)}&
 (4c_*+4A+4\alpha+2)H(1-H)/B
 \end{pmatrix}.
\tag{SL4}
\]
Here and below all expressions in a residue matrix mean their bounded
residues; products such as $AH$ are evaluated by (GE3), even when
$A$ itself is unbounded. Its determinant is exactly $8c_*(1+L)>0$.
Put $a_*=2(2c_*+1)$. The $(R_2,E)$ block is
\[
 \begin{pmatrix}
 a_*&-a_*\sqrt{H/B}\\
 -a_*\sqrt{H/B}&H(4A+a_*)/B
 \end{pmatrix},
\tag{SL5}
\]
with determinant $2a_*c_s>0$. The only remaining coupling is between
$R_2$ and $T_\rho$, whose Gram entry is
$c_s(H/B-s)$. The $T_\rho$ diagonal is $\lambda+2c_s s$.
Choose $\lambda$ larger than the finite Schur-complement threshold
of the positive block (SL5). This produces a strictly positive
residue matrix. Since this choice is made after taking a real
residue, it is a real scalar and introduces no mixed inverse.

We spell out why no normal term is hidden in this block description.
In (IC11), $f_X/\rho$ is bounded by (GE3) and every scalar-root
multiplier contains $u$ or $w$. In (IC12), the only nonzero limiting
term is the explicit $C$ contribution to $R_0E_Y$, already included
in (SL2)--(SL4). After the cap, all other normal multipliers contain
$z$, $u$, or $w$. The unbounded factors $A$ and $\alpha$ always occur
as $AH$ and $\alpha(1-H)$ after normalization, hence remain bounded
also at $H_0=0,1$. The normal $X$ coefficient divided by $N=4n+2$
has residue one, and its coupling to $E_Y$ vanishes. The remaining
$Y$ coefficient is exactly the lower-right entry of (SL4).
These facts follow directly from the displayed exact normal
factorizations, not from scalar-graph positivity.

The bounded positive-residue Gram lemma now gives tensor SOS for the
complete original parent, after restoring the retained positive
$C-c_*$ squares. This holds for all states and all independent local
tests, for both opening endpoints, for arbitrary $n\ge0$, and at
literal $u=0$ or $w=0$.

\subsection{Complementary tests for strong $C$}
\label{component:p7405dbf89341}

\label{p7405dbf89341:sec:strong-ih-remote}
Use the source entry of \S\ref{pf5726218aec8:sec:strong-angle-entry}. In particular
$B=H+L$ is a bounded positive unit, $m,D$ are infinitesimal,
$C,A,\alpha$ are strong, and the bounded products satisfy
\[
 CL\to\tfrac12,\qquad AH\to(1+L)/2,
 \qquad\alpha(1-H)\to B/2.
\tag{SR1}
\]
Whenever an event coefficient is unbounded, a positive finite real cap
may be used and its excess retained in its original positive square.
When a coefficient is bounded it may be used unchanged.

\paragraph{Finite-penalty principle.}
Suppose a real symmetric quadratic form is strictly positive on the
common kernel of finitely many real linear forms. Adding sufficiently
large positive multiples of their squares makes the complete form
strictly positive. To verify this without a limiting assertion, choose
an orthogonal splitting into their common kernel and its complement.
The upper block has a positive lower eigenvalue $c$, and the penalty
matrix on the complementary block has a positive lower eigenvalue
$e$. If the remaining lower block has norm at most $b$ and the cross
block norm is at most $a$, any penalty $M$ with $Me>b+a^2/c$ makes the
Schur complement strictly positive. All choices are real scalars.
This principle makes the opening cases below uniform; no unbounded
coefficient is inserted into a bounded Gram.

\paragraph{Testing $u$ infinitesimal, $v$ finite.}
Retain the $n,D$ positive square terms and excess caps of $C,A,\alpha$.
Write their finite selected coefficients as $c,a,b>0$. Use roots
$U=uW,Y,Z,X,V=\sqrt m\,vW$. The $m$ positive square term is exactly
\[
 4(V-\sqrt m X+\sqrt m\,uvY-\sqrt m\,uZ)^2.
\]
Together with the selected positive square terms and the original core this is an
exact bounded Gram expression. Its residue has isolated $X,V$ pivots
$2(1+v^2),4$ and the following $(U,Y,Z)$ block (residue notation suppressed):
\[
 \begin{pmatrix}
 2+4cD_v&-2-4cD_v&-2v\\
 -2-4cD_v&2D_v+4cD_v+4aD_v+4bv^2&-4bv\\
 -2v&-4bv&4b
 \end{pmatrix}.
\tag{SR2}
\]
In the interior, use $c=1/(2L)$, $a=(1+L)/(2H)$ and
$b=B/[2(1-H)]$. Its leading principal minors are
\[
 \frac{2(L+D_v)}L,\quad
 \frac{4(L+1)(D_v-H)(L+D_v)}{HL(1-H)},\quad
 \frac{8B(L+1)^2D_v}{HL(1-H)},
\tag{SR3}
\]
all positive. At an opening, apply the finite-penalty principle:
\begin{itemize}
\item $H_0=0$: the large $A$ positive square term forces $Y=0$; the restricted
      two-dimensional determinant is $4(L+1)>0$.
\item $H_0=1$: the large $\alpha$ positive square term forces $Z=vY$; the restricted
      determinant is $4(L+1)^2D_v/L>0$.
\item $L_0=0$: the large $C$ positive square term forces $U=Y$; the restricted
      determinant is $4/(1-H)>0$.
\item $L_0=0$, $H_0=1$: both constraints leave the strictly positive
      form $2Y^2$.
\end{itemize}
The case $L_0=H_0=0$ is excluded by $B_0>0$. Thus (SR2) can always be
made strictly positive with bounded selected coefficients.

\paragraph{Testing $u$ a unit, $w=1/v$ infinitesimal.}
Retain $n,m$ and the excess caps of $C,A,\alpha$; retain the full
$D$ positive square term using its natural root. Put
$U=uW$, $E=wZ$, $P_Z=\sqrt D Z$. In roots $(U,Y,E,X,P_Z)$ the exact
selected expression is
\begin{align*}
 &4cD_w(U-Y)^2+4aD_uD_wY^2+4bD_u(Y-E)^2\\
 &\quad+4\{\sqrt D(wU+uX-wY)-P_Z\}^2\\
 &\quad+2\{(wU+uX-wY)^2+D_u(Y-E)^2+D_wX^2-E^2
                   -2UE+2YE+2uwXE\}.
\tag{SR4}
\end{align*}
Its isolated $X,P_Z$ residue pivots are $2D_u,4$. The $(U,Y,E)$ block is
\[
 \begin{pmatrix}
 4c&-4c&-2\\
 -4c&4c+(4a+4b+2)D_u&-4bD_u-2u^2\\
 -2&-4bD_u-2u^2&4bD_u+2u^2
 \end{pmatrix}.
\tag{SR5}
\]
For interior openings the leading minors are
\[
 2/L,\qquad4(1+L)D_u/[HL(1-H)],\qquad
 8u^2(1+L)^2D_u/[HL(1-H)],
\tag{SR6}
\]
strict because $u_0\ne0$. The opening constraints and restricted
positive determinants are respectively
\[
 \begin{array}{c|c|c}
 H_0=0&Y=0&4u^2(1+1/L)\\
 H_0=1&Y=E&4u^2(1+L)/L\\
 L_0=0&U=Y&4u^2D_u/[H(1-H)]
 \end{array}
\]
and at $L_0=0,H_0=1$ the remaining form is $2u^2Y^2$.
Finite penalties therefore give a strict full Gram at every opening.

\paragraph{Both testing angles unbounded.}
Use (IR4)--(IR7), replacing $C$ by a finite positive cap $c_*$ and
setting $d=D$ (which is infinitesimal in the strong IH source).
Retain the excess $C$ positive square term and the entire $n,m$ positive square terms.
The roots are
\[
 \begin{gathered}
 S=\sqrt{c_*}W,\quad E=\sqrt B\,wZ,\quad U=pwW,\quad P=\sqrt D\,pZ,\\
 E_X=X+wW,\qquad E_Y=\sqrt{B/[H(1-H)]}(Y-HwZ).
 \end{gathered}
\]
The exact graph identity (IR5) remains valid. The $E^2$ coefficient
has residue $2c_s$, $c_s=(1+L)/B$, because its core contribution
$2(1-H)^2/B$ must be retained when $B$ is a unit.
In (IR7), the potentially nonzero first-normal coefficient is now
\[
 2\{AH-\alpha(1-H)\}-(1-H),
\]
whose residue is zero by (SR1). The normal $E_Y$ pivot is also $2c_s$.
All scalar-normal cross coefficients are infinitesimal, since
$\sqrt D$ is infinitesimal. Thus the full diagonal residue is
\[
 \operatorname{diag}(4,2c_s,2,4,2,2c_s)>0.
\tag{SR7}
\]

\paragraph{The remaining remote charts.}
When $u$ is unbounded and $v$ is a unit, the homogeneous core has
residue $2\{(W+vX)^2+(vY-Z)^2\}$. Positive finite caps on the $C$
and $A$ positive square terms supply $W^2$ and $Y^2$, respectively, so the residue
is strict on $(W,X,Y,Z)$. When $u$ is unbounded and $v$ infinitesimal,
put $p=1/u$ and use the five roots
$(P,Q,W,Y,Z)=(pX,W+vX,W,Y,Z)$. Positive caps $c,a,b$ give the residue
\[
 2Q^2+2(P+Z)^2+4cW^2+4aY^2+4bZ^2>0.
\tag{SR8}
\]
When $u$ is a unit and $v$ infinitesimal, the core has residue
$2\{(uW-Y)^2+(X+uZ)^2\}$; the positive $A$ and $\alpha$ caps give
$Y^2$ and $Z^2$, making the residue strict. In these expressions all
excess event coefficients are retained in their positive squares.
When both testing angles are units, the inherited noncommuting-test
theorem applies independently of source scales.

\paragraph{Exhaustive join.}
Classify each testing angle in the projective line as infinitesimal,
a bounded unit, or unbounded, with literal zero and infinity included
in the corresponding homogeneous chart. The infinitesimal/unbounded
cell is the local theorem of \S\ref{pf3e0ca8bf054:sec:strong-ih-local}.
The infinitesimal/finite cells follow from (SR2); the unit/unbounded
cell from (SR4); the unbounded/unbounded cell from (SR7); the
unbounded/finite and unit/infinitesimal cells from (SR8) and its
preceding argument; and the unit/unit cell from the inherited theorem.
These nine cells are exhaustive. Both exact sign gauges
$(u,W,X)\mapsto(-u,-W,-X)$ and
$(v,W,Y)\mapsto(-v,-W,-Y)$ preserve the full parent, so every
signed testing chart is included. Every certificate retains the full
state normals and the omitted original $\lambda,\beta$ remain zero.
Therefore the entire strong IH source leaf has all-tests tensor-SOS
case, including $n=0$ and every projective endpoint.

\section{One ordinary and one small actual angle}
\subsection{Source reduction at an ordinary--small contact}
\label{component:pc3daa35e9a9b}

\label{pc3daa35e9a9b:sec:ui-full-parent-reduction}

This section records the original source, every parent positive square term, and the
bounded products used in the following reciprocal and local certificates.
It is an exact algebraic reduction, not an all-tests case. The
coefficient and testing fields remain independent real closed extensions
of $\mathbb R$.

\paragraph{Source definitions and physical domain.}
Take the actual source angles $u_c=a$ and $v_c=k$, where $a>0$ is a
bounded unit and $k>0$ is infinitesimal. Set $z=a^2$ and
$\varepsilon=k^2/z$. In the original canonical C2 chart put
\[
 h=b-\varepsilon t,\quad A_c=1+zr,\quad j=A_c-t,\quad
 g=(1+z)b-1+r-\varepsilon zt,\quad S=A_c+\varepsilon zt.
 \tag{UR1}\label{pc3daa35e9a9b:eq:ui-source-gaps}
\]
The native positive gaps are $r,1-r,t,h,g,j,1-zh$; assume all original
event signs, including $n\ge0$, and assume $C,\alpha$ are bounded below
by positive real numbers. No positive lower bound on a native opening
is imposed. Define the two complete source polynomials
\begin{align*}
 \Delta={}&1+(1+z)b+z(1+z)b^2+2z(1+z)br+z^2r^2+2zr\\
 &+\varepsilon z\{(1+z)b+1+zr(2-r)\},\\
 N={}&(1+z)(b+r)\{b+\varepsilon(1-t)\}
          -\varepsilon(1+\varepsilon z)(1-r)t.
 \tag{UR2}\label{pc3daa35e9a9b:eq:ui-source-polynomials}
\end{align*}
With $d=1$ and $\lambda=\beta=0$, the unchanged original coefficients
are
\[
 \begin{aligned}
 n&=\frac{(b+r)A_cS-r\Delta}{2r\Delta},&
 m&=\frac{(1+k^2)(1-r)A_cS}{2g\Delta},\\
 C&=\frac{A_cN}{2\varepsilon t\Delta},&
 A&=\frac{\varepsilon A_c(b+r)S}{2h\Delta},\\
 D&=\frac{(1+z)(b+r)(b-k^2r)S}{2\varepsilon j\Delta},&
 \alpha&=\frac{z^2(b+r)N}{2(1-zh)\Delta}.
 \end{aligned}                                                   \tag{UR3}
 \label{pc3daa35e9a9b:eq:ui-original-source}
\]
These are the original rational C2 source formulas, with its angle-square
parameter replaced by $z$ and its squared angle-ratio parameter by
$\varepsilon$; all quantities formerly supplied by a source-file
reference are displayed in (UR1)--(UR3).

The source bounds used below follow directly from these definitions.
First $t<A_c$, $h<1/z$, and $b=h+\varepsilon t$ give bounded $t,h,b$.
The inequality $(1+z)b+r>1$ gives a positive lower bound on $b+r$;
(UR2) makes $\Delta$ a positive bounded unit. The exact positive identity
\[
 N=(1+z)(b+r)h+
       \varepsilon\{g+j+rt(1+\varepsilon z)\}                 \tag{UR4}
\]
shows that infinitesimal $h$ would make $b,N$ infinitesimal and
$1-zh$ a unit, contradicting strong $\alpha$. Thus $h,b,N$ are units.
The numerator of $2rn\Delta$ evaluated at $r=1$ is exactly
\[
 -z(\varepsilon+h)(b+1)(z+1).
\]
It has negative unit residue, so $n\ge0$ forces $1-r$ to be a unit.
Consequently
\[
 mg,\quad A/\varepsilon,\quad \varepsilon Ct,\quad
 \varepsilon Dj,\quad \alpha(1-zh)
 \quad\hbox{are positive bounded units},\qquad nr\text{ is bounded}.
 \tag{UR5}\label{pc3daa35e9a9b:eq:ui-source-bounds}
\]
In particular $m$ is strong, $A$ is infinitesimal, and $C,D$ may have
arbitrarily large positive sizes. Literal $n=0$ is retained.

\paragraph{The entire parent and exact contact.}
For independent testing angles and a homogeneous state define
\[
 \begin{gathered}
 D_u=1+u^2,\quad D_v=1+v^2,\quad
 r_{\rm amp}=uW-Y,\quad Q=uX-Z,\\
 J=r_{\rm amp}+vQ=uW+uvX-Y-vZ,\\
 M=vW-X+uvY-uZ,\quad p=vY-Z,\quad B=X+uZ.
 \end{gathered}
\]
The signed core and the full original parent are
\[
 \mathcal K=2\{J^2+D_up^2+D_vB^2-D_uD_vZ^2\},
\]
\[
 F=4nD_uD_vX^2+4mM^2+4CD_vr_{\rm amp}^2
       +4AD_uD_vY^2+4DJ^2+4\alpha D_up^2+\mathcal K.
 \tag{UR6}\label{pc3daa35e9a9b:eq:ui-full-parent}
\]
Normalize the actual contact by $W_c=1$. Introduce the bounded pure-source
quantities
\[
 x=-r/b,\quad \rho=t/(ab),\quad \sigma=-j/(ab),\quad
 s=\sigma-\rho=-A_c/(ab),\quad
 \mu=g/b,\quad \pi=(zh-1)/(ab).                       \tag{UR7}
 \label{pc3daa35e9a9b:eq:ui-contact-amplitudes}
\]
Here $s$ is a nonzero bounded unit. The exact stationary contact is
\[
 (u_c,v_c;W_c,X_c,Y_c,Z_c)
 =\left(a,k;1,kx,a-\rho k^2,\frac{k}{ab}\right),       \tag{UR8}
\]
and its amplitudes are
\[
 r_{{\rm amp},c}=\rho k^2,\quad J_c=\sigma k^2,\quad
 M_c=\mu k,\quad p_c=\pi k,\quad Q_c=sk.
\]
Substitution of (UR3) gives $F=0$ and all four state and both angle
first derivatives zero at (UR8). These are the unchanged native
contact identities; no equality of independent testing rates is used.

\paragraph{Products with all possible opening poles cancelled.}
Define two positive bounded units
\[
 c_1=\frac{zA_cN}{2\Delta},\qquad
 c_2=\frac{z(1+z)(b+r)(b-k^2r)S}{2\Delta}.
\]
The exact products required by the local cap argument are
\[
 \begin{aligned}
 nx&=-\frac{(b+r)A_cS-r\Delta}{2b\Delta},&
 m\mu&=\frac{(1+k^2)(1-r)A_cS}{2b\Delta},\\
 \alpha\pi&=-\frac{z^2(b+r)N}{2ab\Delta},&
 C\rho k^2&=\frac{c_1}{ab},\qquad
 D\sigma k^2=-\frac{c_2}{ab}.
 \end{aligned}                                                   \tag{UR9}
 \label{pc3daa35e9a9b:eq:ui-cancelled-products}
\]
Their products with the bounded factors $x,\mu,\pi$ also show
$nx^2,m\mu^2,\alpha\pi^2$ bounded. Crucially, cancellation before
estimating gives the further exact identity
\[
 Ct-Dj=
 \frac{(1+z)(b+r)\{A_cj-zt(b-\varepsilon zr)\}
       -A_c(1+\varepsilon z)(1-r)t}{2\Delta}.          \tag{UR10}
 \label{pc3daa35e9a9b:eq:ui-cancelled-CD-difference}
\]
Thus $C\rho+D\sigma=(Ct-Dj)/(ab)$ is bounded for every positive rate
of $t,j$, including infinitesimal openings.

Set
\[
 L=C+D,\qquad q=D/L,\qquad H=CD/L,\qquad
 \tau=(C\rho+D\sigma)/L,\qquad \gamma=k^2H.
\]
Here $q,1/L,\tau,L\tau$ are bounded, and
\[
 \gamma=\frac{c_1c_2}{c_1j+c_2t}\asymp1,
 \qquad t+j=A_c\asymp1.                              \tag{UR11}
 \label{pc3daa35e9a9b:eq:ui-harmonic-unit}
\]
The denominator in (UR11) is a positive unit even if either opening
is infinitesimal. At finite source residues it gives precisely $H_0$
in (RC3). Equations (UR9)--(UR11) also show that the bounded finite
part extends continuously to the allowed boundary residues: no
inverse of $r,g,1-zh,t,j$ remains in these products.

\begin{proposition}[Exact bounded UI remainder]
\label{pc3daa35e9a9b:prop:ui-bounded-remainder}
For bounded source caps $0\le n_*\le n$ and positive bounded-unit
caps $m_*\le m$, $\alpha_*\le\alpha$, the parent (UR6) is the sum
of explicit positive centered positive square terms and a bounded polynomial
remainder $G$. Every centered positive square term vanishes with its full first jet
at (UR8), so $G$ has exact zero value and all first derivatives there.
\end{proposition}
\begin{proof}
The exact original $C,D$ combination is
\[
 CD_vr_{\rm amp}^2+DJ^2
   =L(r_{\rm amp}+qvQ)^2+Hv^2Q^2+Cv^2r_{\rm amp}^2.
\]
Retain the three nonnegative positive square terms
\[
 \begin{aligned}
 P_L&=4L(r_{\rm amp}+qvQ-\tau k^2W)^2,\\
 P_H&=4H(vQ-sk^2W)^2,\\
 P_C&=4Cv^2(r_{\rm amp}-\rho k^2W)^2.
 \end{aligned}                                                   \tag{UR12}
\]
Their subtraction leaves exactly
\begin{align*}
 G_{CD}={}&8L\tau k^2W(r_{\rm amp}+qvQ)-4L\tau^2k^4W^2\\
 &+8Hs k^2vWQ-4Hs^2k^4W^2\\
 &+8C\rho k^2v^2Wr_{\rm amp}-4C\rho^2k^4v^2W^2.
 \tag{UR13}
\end{align*}
All its coefficients are bounded by (UR9)--(UR11), including
$L\tau^2=(L\tau)^2/L$. Centering at $k^2W$ is essential to this
boundedness assertion.

Also retain
\[
 \begin{aligned}
 P_n&=4(n-n_*)D_uD_v(X-xvW)^2,\\
 P_m&=4(m-m_*)(M-\mu vW)^2,\\
 P_\alpha&=4(\alpha-\alpha_*)D_u(p-\pi vW)^2.
 \end{aligned}                                                   \tag{UR14}
\]
The exact corresponding remainders are
\[
 \begin{aligned}
 G_n={}&4n_*D_uD_vX^2+8(n-n_*)xvWD_uD_vX
                    -4(n-n_*)x^2v^2W^2D_uD_v,\\
 G_m={}&4m_*M^2+8(m-m_*)\mu vWM
                    -4(m-m_*)\mu^2v^2W^2,\\
 G_\alpha={}&4\alpha_*D_up^2+8(\alpha-\alpha_*)\pi vWD_up
                    -4(\alpha-\alpha_*)\pi^2v^2W^2D_u.
 \end{aligned}                                                   \tag{UR15}
\]
These coefficients are bounded by (UR9) and the bounded caps. With
\[
 G=G_{CD}+G_n+G_m+G_\alpha+4AD_uD_vY^2+\mathcal K,
\]
the entire exact identity is
\[
 F=P_L+P_H+P_C+P_n+P_m+P_\alpha+G.                 \tag{UR16}
 \label{pc3daa35e9a9b:eq:ui-entire-reduction}
\]
It includes the signed core and every original state-normal term.
All six retained roots vanish at (UR8), proving the contact-jet
assertion. This proposition does not claim positivity of $G$.
\end{proof}

For the local theorem retain $P_L,P_C,P_n,P_m,P_\alpha$ outside the
bounded Gram and combine $G$ with
\[
 P_H=4\gamma R_h^2,\qquad R_h=(vQ-sk^2W)/k.
 \tag{UR17}
\]
The factor $1/k$ is pure coefficient. The exact graph normals are
\[
 E_X=X-xvW,\quad E_Z=Z-(ux-s)vW,\quad
 E_r=r_{\rm amp}-\rho k^2W.
\]
Their inverse is polynomial:
\[
 X=E_X+xvW,\quad Z=E_Z+(ux-s)vW,\quad
 Y=uW-\rho k^2W-E_r,\quad Q=svW+uE_X-E_Z.
\]
Thus no state is discarded, including $W=0$. The following reciprocal
lemma identifies the leading relative-angle completion; the local
certificate supplies the complete normal lift, cap choice, and strict
Gram. The outer-angle and other commuting-axis join remains a separate
obligation, so this reduction adds no all-tests source case.

\subsection{Reciprocal completion of a singular angle}
\label{component:pa142a027a282}

\label{pa142a027a282:sec:ui-reciprocal-contact}

The following algebraic identity applies at the indicated singular angle.

\begin{lemma}[Reciprocal contact completion]
Let $A_0,A_2$ be symmetric matrices over a real closed field, let
$q_0$ be a column vector, and consider
\[
 f(q,V,W)=q^T(A_0+V^2A_2)q+2VW b^Tq+(c_0+c_2V^2)W^2.
\]
Suppose $(q,V,W)=(q_0,1,1)$ is a zero with all first derivatives
zero. Then
\[
 b=-(A_0+A_2)q_0,\qquad
 c_0=q_0^TA_2q_0,\qquad c_2=q_0^TA_0q_0,
\]
and the entire polynomial has the exact decomposition
\[
 \boxed{f=(q-Vq_0W)^TA_0(q-Vq_0W)
             +(Vq-q_0W)^TA_2(Vq-q_0W).}             \tag{RC1}
\]
Consequently $A_0,A_2\succeq0$ imply a sum of squares. If
$A_0\succ0$ and $q_0^TA_2q_0>0$, its nonzero zero states have
$V=\pm1$ and $q=Vq_0W$.
\end{lemma}
\begin{proof}
The $q$ equation gives the displayed value of $b$. The $V$
equation then gives $c_2=q_0^TA_0q_0$, and the $W$ equation gives
$c_0=q_0^TA_2q_0$. Expanding the right side of (RC1) proves the
identity. The zero statement follows first from the strictly positive
$A_0$ term and then from $(V^2-1)^2W^2q_0^TA_2q_0$.
\end{proof}

\paragraph{Homogeneous form without a mixed inverse.}
For a positive pure-coefficient $k$, substitute $q=Q/k$ and
$V=v/k$, and multiply (RC1) by $k^2$. It becomes
\[
 (Q-vq_0W)^TA_0(Q-vq_0W)
 +\frac1{k^2}(vQ-k^2q_0W)^TA_2(vQ-k^2q_0W).       \tag{RC2}
\]
Thus its roots use only the pure-coefficient factor $1/k$.
There is no inverse or radical of a mixed tensor element and no
comparison between a testing rate and $k$.

\paragraph{The bounded-source C2 UI leading model.}
Here $u_c=a$ is a positive bounded unit and $v_c=k$ is
infinitesimal, with $d=1$. To state a residue model, first restrict
to the subdomain where the original $n,m,\alpha$ are bounded and
the source parameters $r,g,1-a^2b$ have nonzero residues. In the
following formulas $a,b,r$ denote those residues, $z=a^2$, and
\[
 A_c=1+zr,\quad g=(1+z)b-1+r,\quad
 \Delta=1+(1+z)b+z(1+z)b^2+2z(1+z)br+z^2r^2+2zr.
\]
The source formulas give
\[
\begin{gathered}
 n=\frac{(b+r)A_c^2-r\Delta}{2r\Delta},\qquad
 m=\frac{(1-r)A_c^2}{2g\Delta},\qquad
 \alpha=\frac{z^2(1+z)b(b+r)^2}{2(1-zb)\Delta},\\
 H_0=\overline{k^2CD/(C+D)}
       =\frac{z(1+z)(b+r)b}{2\Delta},\qquad
 A_0^{\rm weak}=\overline{A/k^2}
       =\frac{A_c^2(b+r)}{2zb\Delta}.
\end{gathered}                                                   \tag{RC3}
\]
Normalize the actual contact by $W=1$. With $v=kV$, $X=k\xi$,
$Z=k\zeta$ and the leading constraint $Y=aW$, the leading
polynomial of $F/k^2$ is
\begin{align*}
 f={}&4n(1+a^2)\xi^2
 +4m\{V(1+a^2)W-\xi-a\zeta\}^2
 +4H_0V^2(a\xi-\zeta)^2\\
 &+4A_0^{\rm weak}(1+a^2)a^2W^2
 +(4\alpha+2)(1+a^2)(VaW-\zeta)^2\\
 &+2(\xi+a\zeta)^2-2(1+a^2)\zeta^2.               \tag{RC4}
\end{align*}
It has the form of the lemma, with $q=(\xi,\zeta)^T$,
\[
 q_0=(-r/b,1/(ab))^T,
\]
\[
 q^TA_0q=4n(1+a^2)\xi^2+(4m+2)(\xi+a\zeta)^2
                      +4\alpha(1+a^2)\zeta^2,
 \qquad q^TA_2q=4H_0(a\xi-\zeta)^2.
\]
Here $A_0\succ0$ and $H_0>0$. Put
$s_0=aq_{0,1}-q_{0,2}=-(1+a^2r)/(ab)\ne0$.
The exact source identities (RC3) verify all three coefficient
relations in the lemma. Hence the apparently problematic relative
normal and scale difference combine into the single root
\[
             \frac{v(aX-Z)-s_0k^2W}{k}.             \tag{RC5}
\]
The ordinary centered normals $(X,Z)-vq_0W$ supply the remaining
positive directions. Requiring separate positive pivots for the
relative normal and the scale difference would obscure this
completion.

This proves an exact leading-model identity on its stated bounded
subdomain. It does not discard the additional large source scales:
those are handled algebraically by the centered-positive square term reduction in
Section~\ref{pc3daa35e9a9b:sec:ui-full-parent-reduction}. The complete local lifted
Gram is proved in Section~\ref{p76545700d90b:sec:local-ui-certificate}; the outer
projective join remains a separate obligation.

\subsection{Local contact certificate}
\label{component:p76545700d90b}

\label{p76545700d90b:sec:local-ui-certificate}

This section uses the exact full-parent reduction and source notation
in Section~\ref{pc3daa35e9a9b:sec:ui-full-parent-reduction}. It proves a local independent-test
statement. The other commuting axis and the outer projective join
are not included.

\begin{theorem}[Local ordinary--small C2 certificate]
Use the unchanged physical C2 source with $d=1$, actual contact
$u_c=a$ a positive bounded unit and $v_c=k$ a positive infinitesimal,
and $C,\alpha$ bounded below by positive real numbers. Then the full
parent is tensor SOS for every independent test with $v$ infinitesimal
and $u-a$ infinitesimal, and every projective state. The original
$n=0$ boundary and all positive source opening rates are included.
No upper bound on $n,m,C,D,\alpha$ is imposed.
\end{theorem}

Here $u-a$ is the difference of the corresponding tensor factors.
The residue condition means that their real residues agree; it does
not identify the source and testing fields or compare their rates.
The exact sign gauges give the corresponding negative-angle charts.

\paragraph{Bounded full remainder and its exact graph.}
The positive square term reduction leaves a bounded polynomial remainder $G$ and
the positive relative square
\[
 4\gamma R_h^2,\qquad
 \gamma=k^2CD/(C+D)\asymp1,\qquad
 R_h=\frac{v(uX-Z)-s k^2W}{k}.
\]
Every other retained positive square term is an exact positive square vanishing
at the full actual contact. For positive bounded caps $m_*,\alpha_*$
and $0\le n_*\le n$, set
\[
 E_X=X-xvW,\quad E_Z=Z-(ux-s)vW,\quad
 E_r=uW-Y-\rho k^2W,
 \qquad Q_N=uE_X-E_Z.
\]
The full inverse is polynomial. All terms of $G$, including its
complete first-normal rows and quadratic normal block, are retained.
There are no higher state-normal terms.

Put $w=u-a$, $\xi=v^2$, and use the additional separated roots
\[
 S=vwW,\qquad T=kwW,\qquad R=(v^2-k^2)W/k.
\]
The exact identities
\[
 kR=\frac{kR_h-vQ_N}{s},\qquad
 RT=\frac{R_hT-SQ_N}{s},\qquad
 S^2-T^2=\frac w s(R_hT-SQ_N)                      \tag{UI1}
\]
use only the bounded pure-source inverse $1/s$.

\paragraph{Complete scalar and normal lifting.}
On the graph $E_X=E_Z=E_r=0$, write $G=\phi(u,\xi)W^2$.
Its bounded polynomial coefficients and the original parity imply
$\phi\in(k^2,\xi)$. The exact contact identities give
$\phi=\phi_u=\phi_\xi=0$ at $(u,\xi)=(a,k^2)$.
Successive monic Taylor division gives
\[
 \phi=w^2(k^2A_0+\xi A_1)
             +w(\xi-k^2)B+(\xi-k^2)^2C_0,          \tag{UI2}
\]
with bounded coefficients. To see the asserted first factor directly,
divide first in $u-a$; the second $u$ coquotient stays in
$(k^2,\xi)$. The constant and first $u$ terms are then divided in
$\xi-k^2$. No source infinitesimal is inverted by these operations.
Equations (UI1) allocate the three terms of (UI2) on
$(S,T,R_h,E_X,E_Z)$.

The full first $E_X,E_Z$ coefficients are $vL_X(u,\xi)$ and
$vL_Z(u,\xi)$. Exact contact gives $L_i(a,k^2)=0$, so they are
\[
 v\{wL_{i0}+(\xi-k^2)L_{i1}\}.
\]
The first summand is allocated on $S$, and the second uses
$vkR=v(kR_h-vQ_N)/s$. The entire $E_r$ coefficient belongs to
$(k^2,\xi)$ and vanishes at contact. The same division writes it as
\[
 w(k^2L_{r0}+\xi L_{r1})+(\xi-k^2)L_{r2}.
\]
It consequently uses $kT,vS,kR$. All these $E_r$ couplings have
zero residue. Thus (UI1)--(UI2) produce an exact bounded Gram on
\[
             (E_r,E_X,E_Z,R_h,S,T).                 \tag{UI3}
\]
This is a full-parent lift, not a restriction to a state kernel.

\paragraph{The strict matched Hessian.}
First take finite positive residue source weights. The leading model
before setting the ordinary test angle equal to $a$ is (RC4), with
every testing occurrence of $a$ replaced by $u$; its source
coefficients remain those of (RC3). Let $\mathsf G$ denote one half
of its angular Hessian after eliminating its two state directions,
in variables $(V,u)$ at the exact contact. Direct differentiation and
two-by-two Schur elimination give
\[
 \mathsf G_{uu}=
 \frac{8(1-r)(1+zr)^2}{zb\Delta}>0,                 \tag{UI4}
\]
\[
 \det\mathsf G=
 \frac{16(b+r)^2(1+zr)^3
                \{3z+4-(z+2)^2r\}}{b^2\Delta^2}.  \tag{UI5}
\]
The other two entries, making boundedness explicit, are
\begin{align*}
 \mathsf G_{Vu}&=
 \frac{4(1+zr)^2\{zb+zr+2r-2\}}{ab\Delta},\\
 \mathsf G_{VV}&=
 \frac{8(1+zr)\{z^2b^2+z^2br+zb^2+2zbr-zb-r+1\}}
      {b\Delta}.
\end{align*}
For a short check of strictness, the exact leading $n$ numerator is
$bQ_n$, where
\[
 Q_n=1+r(z-1)-z(z+2)r^2-z(z+1)br
     =(1-(z+1)r)(1+zr)-zr g.
\]
Since $Q_n\ge0$ and $g\ge0$, one has $r\le1/(z+1)$.
Consequently
\[
 3z+4-(z+2)^2r\ \ge\ \frac{z(2z+3)}{z+1}>0.       \tag{UI6}
\]
All denominators in (UI4)--(UI6) are positive bounded units.
These formulas therefore give a strict, uniformly bounded-away-from-zero
angular Schur form on the source residue domain, including $n=0$.
They contain no inverse of $r$, $g$, or $1-zb$.

\paragraph{Large weights and finite real caps.}
We use the elementary finite-dimensional penalty lemma: if a real
symmetric form $B$ is strictly positive on the common kernel of the
rows of a matrix $V$, then $B+M V^TV$ is positive definite for some
finite real $M>0$. A proof is to split into $\ker V$ and its
orthogonal complement. The first diagonal block is strict; taking
$M$ sufficiently large makes its Schur complement strict too.

The bounded finite part obtained by the centered positive square terms is
continuous at all source boundary residues. This assertion is
algebraic: its only large-weight products are
\[
 nx,\ nx^2,\ m\mu,\ m\mu^2,\ \alpha\pi,\ \alpha\pi^2,
 \quad C\rho k^2,\ D\sigma k^2,\ k^2H,
\]
and their bounded cap analogues. Substitution of the original source
formulas cancels the possible $r,g,1-zh,t,j$ poles in those products;
their denominators are the displayed positive units. For example,
\[
 m\mu=\frac{(1+k^2)(1-r)A_cS}{2b\Delta},\qquad
 \alpha\pi=-\frac{z^2(b+r)N}{2ab\Delta}.
\]
The other cancellations are
\[
 nx=-\frac{(b+r)A_cS-r\Delta}{2b\Delta},\quad
 nx^2=x(nx),\quad m\mu^2=\mu(m\mu),\quad
 \alpha\pi^2=\pi(\alpha\pi),
\]
\[
 C\rho k^2=\frac{c_1}{ab},\qquad
 D\sigma k^2=-\frac{c_2}{ab},\qquad
 c_1=\frac{zA_cN}{2\Delta},\quad
 c_2=\frac{z(1+z)(b+r)(b-k^2r)S}{2\Delta}.
\]
Finally $C\rho+D\sigma=(Ct-Dj)/(ab)$, using the exact canceled
formula in the full-parent reduction. Thus even this difference has
no hidden small-opening denominator.
For $k^2H$, write $k^2C=c_1/t$, $k^2D=c_2/j$, with
$c_1,c_2$ positive bounded units. Then
$k^2H=c_1c_2/(c_1j+c_2t)$, whose denominator is a unit because
$t+j=A_c$ is a unit. Its residue is exactly $H_0$ in (RC3).

At most two of $n,m,\alpha$ can diverge. Indeed $m$ can diverge
only at $g=0$, while $\alpha$ can diverge only at $1-zb=0$;
these two equalities would give $g=1/z+r>0$. The hard-state parts
of the three possible diverging centered rows are, up to nonzero
unit factors,
\[
                   E_X,\qquad E_X+aE_Z,\qquad E_Z.
\]
Every possible pair is independent with a unit determinant.

Here is why their constrained limiting finite form is strict.
Approach the boundary by the positive leading-source domain. The
state block is bounded below by a positive real form because $m$ and
$\alpha$ are strong. Its minimizing hard vector for a fixed angular
vector is bounded: solve the at most two diverging positive square term equations
in the hard variables, using the unit determinants just displayed;
this supplies a bounded competitor, and all remaining coefficients
stay bounded. If $h_0$ is that competitor and $h_{\min}$ the state
minimizer for angular vector $q$, the exact identity
\[
 (h_0-h_{\min})^TP(h_0-h_{\min})
       =F(h_0,q)-q^T\mathsf Gq
\]
and $P\succeq cI$ make boundedness of $h_{\min}$ explicit.
Together with (UI4)--(UI6), the state Schur factorization
then bounds the full matched Hessian below by a positive real multiple
of the Euclidean norm. For a vector in the common limiting positive square term
kernel, adjust only its hard coordinates to lie in the exact nearby
positive square term kernels. The adjustment tends to zero and annihilates the
large weights exactly. Taking the limit proves strict positivity of
the bounded finite part on that common kernel. The penalty lemma now
chooses finite real caps for every diverging weight. Bounded source
weights can be retained in full. All caps are below their original
weights, so the remaining centered excesses are positive squares.

This limiting argument uses only real finite-dimensional forms to
choose caps. No inverse of a tensor matrix and no mixed root are
used in the full-parent identity. The boundary residues are limits
of the positive leading-source domain: $g$ and $1-zb$ specify the
interval for $b$, and decreasing a positive $r$ if necessary makes
$Q_n\ge0$ strict. Thus the argument also covers a literal original
$n=0$ and all combinations of the allowed boundary residues.

\paragraph{Separate angular roots and full strictness.}
The residue of the hard block in (UI3) is
\[
 2E_r^2+4n_*(1+a^2)E_X^2
 +(4m_*+2)(E_X+aE_Z)^2+4\alpha_*(1+a^2)E_Z^2
 +4\gamma R_h^2,
\]
which is strict. In the residue Gram, $S$ couples only to the two
ordinary hard normals, and $T$ only to $R_h$; the $E_r$ couplings
vanish. Setting $S=T$ gives the matched leading Hessian under an
invertible bounded change of coordinates. Indeed at $V=1$ the
ordinary normals are its two state deviations, $R_h$ supplies the
relative-angle deviation with nonzero coefficient $2s$, and
$S=T$ supplies the ordinary-angle deviation. The chosen caps therefore
make the sum of the two angular Schur diagonals strictly positive.

Add an appropriate real multiple of the last zero identity in (UI1).
It redistributes mass between $S^2$ and $T^2$; all accompanying
coefficients contain $w$ and have zero residue. Since their Schur
diagonal sum is positive, choose this multiple so that both Schur
diagonals are positive. The entire six-root residue Gram is now
positive definite. The strict-residue Gram lemma proves the exact
bounded remainder plus relative positive square term is tensor SOS.

Restore the centered excesses and the other C/D squares. The construction
is homogeneous in $W$ with a polynomial state inverse, so it includes
$W=0$ and every state. Real and imaginary parts give complex states.
The complementary angle charts are proved next.

\subsection{Complementary angle charts}

\label{p6b22e944bd7c:sec:ui-all-tests-recovered}

\begin{theorem}[Strong C2 UI, all independent tests]
Use the original C2 source (UR1)--(UR3), with $d=1$, actual source
angles $u_c=a>0$ a bounded unit and $v_c=k>0$ infinitesimal, and
$C,\alpha$ strong. Then the entire parent (UR6) is a sum of squares
in the tensor product of the independent source and testing fields,
for every testing angle, including projective endpoints, and every
homogeneous state. The statement retains the literal boundary $n=0$
and arbitrary positive rates of every source opening.
\end{theorem}

Section~\ref{p76545700d90b:sec:local-ui-certificate} covers $v$ infinitesimal
and $u$ with the same residue as $a$. Its complement follows from a
rank-one allocation and positive frames at the testing endpoints.

\paragraph{A rank-one allocation lemma.}
Let $P$ be a positive definite real symmetric matrix, let $\ell,b$ be
real columns, and let $H,c_0>0$. Set
\[
 d=\ell^TP^{-1}\ell,\qquad e=\ell^TP^{-1}b,\qquad
 h=c_2-b^TP^{-1}b.
\]
Suppose
\[
 p(T)=c_0+(h+4Hdc_0)T+4H(dh+e^2)T^2>0\quad(T\ge0),
 \qquad dh+e^2>0.                                      \tag{UA1}
\]
There exists a real $\lambda$ such that
\[
 |\lambda|<\sqrt{4Hc_0},\qquad
 h+2e\lambda-d\lambda^2>0.                             \tag{UA2}
\]
Indeed put $s=\sqrt{4Hc_0}$. If $d>0$ and $|e|<ds$, take
$\lambda=e/d$. If $|e|\ge ds$, the value at the appropriate endpoint
is $h+2|e|s-ds^2$. Were this nonpositive, writing
$x=|e|-ds\ge0$ would give
$h+ds^2\le-2sx$ and $dh+e^2\le x^2$; the quadratic in (UA1)
would then have a nonpositive value at a positive argument. Thus a
nearby interior $\lambda$ works. This includes equality $|e|=ds$.
The scalar interval argument also holds in a limiting case $d=0$;
then $e\ne0$ follows from $dh+e^2>0$. (For a finite positive $P$
itself, $d=0$ would force $\ell=e=0$, so this case occurs only
when using the continuous scalar extension at a boundary.)

Consequently the two real matrices
\[
 \begin{pmatrix}P&b-\lambda\ell\\(b-\lambda\ell)^T&c_2\end{pmatrix},
 \qquad
 \begin{pmatrix}4H&\lambda\\\lambda&c_0\end{pmatrix}       \tag{UA3}
\]
are positive definite. This is an elementary real lemma. Its inverse
and square root are never taken in the tensor ring.

\paragraph{The finite ordinary angle away from contact.}
Assume $v$ infinitesimal and $u$ a bounded unit whose real residue has
square different from $a_0^2$. Keep the centered positive square terms
$P_L,P_C,P_n,P_m,P_\alpha$ of (UR12), (UR14) outside the Gram, and
combine $P_H$ with the exact remainder. Use the six separated roots
\[
 r_0=uW-Y,\quad X,\quad Z,\quad V_0=vW,\quad K_0=kW,
 \quad R_0=\frac{v(uX-Z)}{k}.                           \tag{UA4}
\]
They satisfy the exact zero relation
\[
 R_0K_0-(uX-Z)V_0=0.                                  \tag{UA5}
\]
All divisions in (UA4) are pure source divisions.

For completeness the boundedness of the $C,D$ remainder in this frame
is explicit. Write $\ell_0=L\tau$, $c_R=C\rho k^2$, and
$Q=uX-Z$. Then (UR13) gives exactly
\begin{align*}
 G_{CD}+P_H={}&4\gamma R_0^2
 +8\ell_0kK_0r_0+8\ell_0qk^2V_0Q-4\ell_0\tau k^2K_0^2\\
 &+8c_RvV_0r_0-4c_R\rho k^2V_0^2.                  \tag{UA6}
\end{align*}
All displayed coefficients are bounded by (UR9)--(UR11); all terms
except the first have zero residue. The $A$ positive square term is exactly
$4(A/k^2)D_uD_v(uK_0-kr_0)^2$. Substituting
$Y=uW-r_0$ in (UR15) and the signed core gives a bounded Gram on
(UA4), with no discarded normal term. Its residue has an isolated
$2r_0^2$ term. First suppose the leading source weights are finite.
The remaining residue is
\[
 q_s^TPq_s+2V_0b_*^Tq_s+c_2V_0^2+4H R_0^2+c_0K_0^2,
 \quad q_s=(X,Z)^T,                                  \tag{UA7}
\]
where $D_u=1+u^2$, $H=H_0$, $A_w=A_0^{\rm weak}$ are (RC3), and
\[
\begin{gathered}
 P=\begin{pmatrix}
 4nD_u+4m+2&(4m+2)u\\
 (4m+2)u&(4m+2)u^2+4\alpha D_u
 \end{pmatrix},\quad \ell=(u,-1)^T,\\
 b_*=\binom{-4mD_u}{-D_uu(4m+4\alpha+2)},\quad
 c_0=4A_wD_uu^2,\quad
 c_2=4mD_u^2+(4\alpha+2)D_uu^2.                      
\end{gathered}
\tag{UA8}
\]
The matrix $P$ is strictly positive because $m,\alpha$ are strong.
The last sign-changing part of (UA7) is handled by adding
$2\lambda$ times (UA5). This gives precisely the two blocks (UA3).

Here is a direct proof of (UA1) for every ordinary angle away from
contact. Put $N_*=2n+1$, $T_*=2m+2\alpha+1$, $U=u^2$, and
$T=V^2$. The leading state polynomial (RC4), with its testing $a$
replaced by $u$ but its source unchanged, has determinant
\[
 \det f=8D_u^2u^2V^2\,\mathcal L(U,T),                \tag{UA9}
\]
where
\begin{align*}
\mathcal L={}&\frac{4\alpha mN_*}{U}
 +4A_wHT_*U+\frac{4A_w\alpha(2m+N_*)}{T}
 +\frac{4A_wnT_*U}{T}\\
 &+2HN_*T_*T+\frac{4HmN_*T}{U}
 +4A_wH(2m+N_*)-N_*T_*.
\end{align*}
After division by $N_*T_*$ all possible opening poles cancel. Define
\[
 i_N=N_*^{-1},\quad t_m=m/T_*,\quad t_{am}=\alpha m/T_*,
 \quad q_\alpha=\frac{\alpha(2m+N_*)}{N_*T_*},\quad i_\alpha=1/\alpha.
\]
Writing $z=a^2$, $A_c=1+zr$, and using the leading $\Delta$ of (RC3),
these quantities are exactly
\[
\begin{aligned}
 i_N&=\frac{r\Delta}{(b+r)A_c^2},&
 t_m&=\frac{(1-r)(1-zb)}{2b(b+r)(z+1)},\\
 t_{am}&=\frac{z^2(b+r)(1-r)}{4\Delta},&
 q_\alpha&=\frac{bz^2\{(z+1)b+(z+2)r-1\}}{2A_c^2},\\
 i_\alpha&=\frac{2(1-zb)\Delta}{z^2(z+1)b(b+r)^2}.
\end{aligned}\tag{UA10}
\]
All denominators here are units. The normalized Laurent polynomial is
\[
 \frac{\mathcal L}{N_*T_*}
 =\frac{4t_{am}}U+4A_wHi_NU+\frac{4A_wq_\alpha}T
  +2A_w(1-i_N)\frac UT+2HT+4Ht_m\frac TU
  +4A_wHq_\alpha i_\alpha-1.                         \tag{UA11}
\]
Its value and its two first derivatives vanish at $(U,T)=(z,1)$,
as direct substitution verifies. In logarithmic coordinates the
nonconstant terms have nonnegative coefficients. The coefficients of
$U^{-1}$ and $T$ are positive units. If $i_N>0$, the coefficients of
$U$ and $T^{-1}$ are positive; if $i_N=0$, that of $U/T$ is positive.
The function is therefore strictly convex in logarithmic coordinates
and has its unique minimum at $(z,1)$. This proves strict positivity
at fixed $U\ne z$ for every $T>0$. The constant and leading
coefficients of the equivalent quadratic (UA1) are strictly positive:
$c_0>0$, and the positive $T$ coefficient in (UA11) gives the second
assertion. Thus (UA2)--(UA3) prove the required strict residue Gram.

\paragraph{Opening boundaries and divergent source weights.}
The preceding proof extends to all allowed opening residues, rather
than assuming $n,m,\alpha$ bounded. The following formulas make the
needed uniformity checkable. Put
\[
 J_0=q_\alpha+\frac{u^2(1-i_N)}2.
\]
Then $J_0$ is a positive unit, and direct two-by-two inversion gives
\[
\begin{aligned}
 \frac{\det P}{N_*T_*}&=8D_uJ_0,\\
 d&=\frac{u^2i_N+q_\alpha i_\alpha}{4J_0},&
 e&=\frac{uD_u}{2J_0},&
 h&=\frac{D_u(4t_{am}-u^2)}{J_0},\\
 (P^{-1}b_*)_1&=\frac{i_N(u^2-4t_{am})}{2J_0},&
 (P^{-1}b_*)_2&=-\frac{u\{4t_{am}i_N+(1-i_N)u^2+1\}}{2J_0}.
\end{aligned}\tag{UA12}
\]
To check the unit assertion, if $i_N=0$ then $J_0=u^2/2>0$.
If $i_N>0$, the weight $N_*$ is finite and the strong lower bounds
on $m,\alpha$ give $q_\alpha>0$: this follows directly from
$\alpha(2m+N_*)/[N_*(2m+2\alpha+1)]$, whether both weights are
finite or one diverges. Thus (UA12) has only unit denominators on
the entire source residue domain.

The strictly positive leading coefficient is also explicit:
\[
 dh+e^2=\frac{D_u(u^2+2t_m)}{2J_0}>0.
\]
Equations (UA10)--(UA12) prove continuity of $d,e,h$ and of the
minimizing hard vector. The convexity proof for (UA11) remains strict
at every allowed boundary. Choose $\lambda$ from the strict interval
(UA2) at that boundary; the same strict inequalities hold nearby.
Moreover $P$ is bounded below by a positive real matrix, while
$P^{-1}b_*$ and $P^{-1}\ell$ are bounded. Its Schur factorization
therefore bounds the first block in (UA3) below by a positive real
multiple of the identity, uniformly in the nearby source. The second
block is a strictly positive real matrix as well.

The finite leading sources used in this limiting argument are valid
real interior sources. This point must be established separately:
the actual source has $b=h+(k^2/z)t$, and it is not legitimate to
substitute that actual $b$ into a leading formula with denominator
$1-zb$ when this denominator has an infinitesimal sign change.
For the limiting real parameters, put
\[
 f(b,r)=1+r(z-1)-brz(z+1)-r^2z(z+2).
\]
Their domain is $b>0$, $0\le r<1$, $g=(1+z)b-1+r\ge0$,
$1-zb\ge0$, and $f(b,r)\ge0$. For $0<s<1$, let
\[
 b_s=sb+(1-s)/(1+z),\qquad r_s=sr.
\]
Direct expansion gives
\[
 g_s=sg,\qquad 1-zb_s>0,\qquad
 f(b_s,r_s)=(1-s)(1+s-sr)+s^2f(b,r)>0.
\]
Increase $b_s$ by a sufficiently small $\epsilon>0$. This makes
$g_s$ strict, decreases the last two strict margins by respectively
$z\epsilon$ and $r_sz(z+1)\epsilon$, and therefore preserves them.
If $r=0$, then increase $r_s$ by a sufficiently small positive
amount; $g$ increases and $f$ remains positive by continuity.
Letting these perturbations tend to zero as $s\to1$ supplies valid
finite real leading sources approaching every allowed residue.
Their $t$ may be chosen as $(1+zr)/2$, since the leading harmonic
weight is independent of $t$. The normalized leading identity
$2rn\Delta=bf(b,r)$ verifies the nonnegative-$n$ requirement.
The same construction supplies the real-interior approximation
used in the finite-cap step of the local UI certificate.

At most two of $n,m,\alpha$ diverge: $m$ can diverge only at $g=0$,
and $\alpha$ only at $1-zb=0$, which cannot occur together. Their
centered excess rows have hard-state parts proportional to
\[
 X,\qquad X+uZ,\qquad Z.
\]
Any possible pair is independent because $u$ is a unit. For a vector
in their common limiting kernel, change only $(X,Z)$ by a vanishing
amount to annihilate the corresponding nearby exact rows. Their
large penalties then vanish exactly. Taking the limit in the uniform
strict lower bound proves that the bounded finite part is strictly
positive on that limiting kernel. The elementary finite-dimensional
penalty lemma used in the local theorem now supplies finite real caps
$n_*,m_*,\alpha_*$. Their centered excesses are precisely the retained
positive square terms (UR14). The exact bounded remainder with those caps, after
adding (UA5), has a strictly positive residue Gram on (UA4).
This proves the finite ordinary-angle complement for all source
openings and literal $n=0$.

\paragraph{Endpoint frames.}
Choose positive real $\mu,m_*,\alpha_*$ below, respectively,
$k^2C,k^2D$, $m$, and $\alpha$. The products $k^2C,k^2D$ are strong
by (UR5), and $A_w=A/k^2$ is a positive unit. Retain $n$ and all
positive excesses in full. The following formulas concern the parent
with $C=D=\mu/k^2$, $m=m_*$, $\alpha=\alpha_*$ and $A=A_wk^2$.
They therefore suffice. The root coordinates need not themselves be
bounded; only their displayed Gram coefficients must be bounded.

At $u,v$ infinitesimal use
\[
 E=(uW-Y)/k,\quad R=v(uX-Z)/k,\quad V=vW,\quad K=kuW,
 \quad X,Z.
\]
The exact amplitudes are
\[
 J=k(E+R),\quad M=D_uV-X-uZ-uvkE,\quad
 vY-Z=uV-Z-vkE,\quad kY=K-k^2E.
\]
Their full residue form is
\[
 4\mu\{E^2+(E+R)^2\}+4m_*(V-X)^2+4\alpha_*Z^2
       +4\overline{A_w}K^2+2X^2,                    \tag{UA13}
\]
which is strict. Literal zero testing angles are included without
any division by a testing angle.

For bounded $u$ and finite unit $v$, put
$E=(uW-Y)/k$, $F=J/k$. Then
\[
 Y=uW-kE,\quad Z=uX+k(E-F)/v.
\]
The roots $(E,F,W,X)$ have strict residue
\[
 4\mu D_vE^2+4\mu F^2+
 \{4m_*D_u^2+(4\alpha_*+2)D_uu^2\}(vW-X)^2
 +2D_vD_uX^2.                                      \tag{UA14}
\]
For bounded $u$ and $v=\infty$ write $q=1/v$ and homogenize by $q^2$.
The same roots are defined using the homogeneous $J$, and
\[
 Y=uW-kE,\qquad Z=uX+kqE-kF.
\]
The strict residue is (UA14) with $D_v=1$ and $(vW-X)^2$ replaced
by $W^2$. This includes the literal endpoint $q=0$.

For $u=\infty$, put $p=1/u$, $D_p=1+p^2$, and homogenize by $p^2$.
Write $r_p=W-pY$, $Q_p=X-pZ$, and $P=vY-Z$. Then
\[
 J=r_p+vQ_p,\qquad M=D_pP-pQ_p+pv r_p.
\]
For $v$ infinitesimal use the six roots
\[
 E=r_p/k,\quad F=vQ_p/k,\quad P,\quad R=pQ_p,
 \quad K=kY,\quad S=pZ.
\]
The signed core is exactly
\[
 2k^2(E+F)^2+2D_pP^2+2D_vR^2-4D_pD_vRP
    +4D_pD_vpFK+2D_pD_vS^2.
\]
Its cross-term reduction uses $RvY=pFK$, an exact identity.
Adding every capped positive square term gives the strict residue
\[
 4\mu\{E^2+(E+F)^2\}+(4m_*+2)(P-R)^2
 +4\alpha_*P^2+4\overline{A_w}K^2+2S^2.             \tag{UA15}
\]
For finite unit $v$, instead use $E=r_p/k$, $F=J/k$, $P,K,S$.
Since $Q_p=k(F-E)/v$, the strict residue is
\[
 4\mu D_vE^2+4\mu F^2+(4m_*+4\alpha_*+2)P^2
 +4\overline{A_w}D_vK^2+2D_vS^2.                   \tag{UA16}
\]
Finally, at both infinite endpoints put $q=1/v$ and homogenize by
$p^2q^2$. Use
\[
 E=(W-pY)/k,\quad F=(qW+X-pqY-pZ)/k,\quad
 P=Y-qZ,\quad K=kY,\quad S=pZ.
\]
Here $M=D_pP+pkD_qE-pqkF$ and the entire signed core is
\[
 2k^2F^2+2D_pP^2+2D_qp^2k^2(F-qE)^2
 +4D_qD_pk(F-qE)S+2D_qD_pS^2.
\]
The strict residue is (UA16) with $D_v=1$.

\paragraph{Exhaustion and restoration.}
A testing angle is infinitesimal, a bounded unit, or projectively
infinite. Equations (UA13)--(UA16) cover every pair except ordinary
$u$ with infinitesimal $v$. The local theorem covers the equal-residue
contact, and (UA4)--(UA12) cover its complement. Exact sign gauges are
\[
 (u,W,X)\mapsto(-u,-W,-X),\qquad
 (v,W,Y)\mapsto(-v,-W,-Y);
\]
all unspecified variables are fixed. Both preserve the full parent,
so negative ordinary contact residues are included. All endpoint
frames have been written after homogenization, including literal
projective endpoints. Each proof has an exact bounded Gram with
strict real residue; the strict-residue Gram lemma gives tensor SOS.
Restoring all retained positive square terms completes the original parent.
Real and imaginary parts treat complex states. This closes the
entire strong C2 UI case and its exact symmetry image HU.

\section{An unbounded first angle and a small second angle}
\subsection{Source entry at the HI corner}
\label{component:p54b8f20ad2dd}

Use the original positive rational C2 source with $d=1$, actual
$u_c$ unbounded and $v_c$ infinitesimal. Put
\[
 q=u_c^{-2},\quad y=v_c^2,\quad H=u_c^2h\in(0,1),
 \quad R=r/q,\quad B=b/q,\quad L=yt=B-H,\quad T=B-yR.
\]
The already proved source identities (WE1)--(WE3) show that $R$ is
bounded using only $n\geq0$, $q$ infinitesimal and $y$ bounded.
Here is their short contradiction argument, to make this use
explicit. The positive source gaps imply
\[
 0<T<H+y,\quad g=H+(1+y)r-1+qT>0,
\]
and $n=-qP/(2rQ)$ with $Q>0$ and
\[
 P\ge r^2(1+y)(H+yr)-Tqr(y+2)-q^2(T^2+T).
\]
If $q/r$ were infinitesimal, then $T/(H+yr)$ and
$(T^2+T)/(H+yr)$ would be bounded. For unit $r$ this follows
from $T<H+y\le(H+yr)/r$; for infinitesimal $r$, the inequality
$g>0$ forces $H$ to have residue one. Division by $r^2(H+yr)$
would give $P>0$, contrary to $n\geq0$. Thus $R$ is bounded.
(The exact positive remainder in this lower bound is printed in
(WE3), so this argument uses no asymptotic sign substitution.)

Now $r=qR$ is infinitesimal and $g>0$ gives $H_0=1$.
Define the positive opening
\[
 E=(1-H)/q.
\]
The same gap implies
$E<(1+y)R+T$, so $E$ is bounded. Since $j=1+R-t>0$,
$t$ is bounded; consequently
\[
 H=1-qE,\quad B=H+yt\sim1,\quad
 j=1+R-t,\quad K=B+R-E=g/q>0.
\]
All of $E,t,j,K$ are positive bounded, while
$t+j=1+R$ and $E+K=B+R$ are positive units.

Let $\Delta$ and $N$ denote the unchanged original source
polynomials, and set $\mathcal N=N/q$. Their exact expressions are
\begin{align*}
\Delta={}&(1+q)\{B^2+2BR+(1+y)B\}
 +(1-yq)R^2+2(1+y)R+1+y,\\
\mathcal N={}&(1+q)(B+R)(H+y)-(1+y)(1-qR)yt.
\end{align*}
Thus $\Delta$ is a positive bounded unit and
\[
 \bar\Delta=(1+\bar R)(3+\bar R),\qquad
 \overline{\mathcal N}=1+\bar R>0.
\]
Set $S=1+R+yt$, also a positive unit. Direct substitution into
the original coefficients gives the following exact equalities:
\[
\begin{aligned}
 mqK&=\frac{(1+y)(1-qR)(1+R)S}{2\Delta},&
 Cyt&=\frac{(1+R)\mathcal N}{2\Delta},\\
 \frac{A}{qy}&=\frac{(1+R)(B+R)S}{2H\Delta},&
 Dyj&=\frac{(1+q)(B+R)(B-yR)S}{2\Delta},\\
 \alpha qE&=\frac{(B+R)\mathcal N}{2\Delta},&
 n&=\frac{(B+R)(1+R)S-R\Delta}{2R\Delta}.
\end{aligned}                                                   \tag{HI-entry}
\]
Every one of the first five right-hand sides is a positive bounded
unit; in particular $B-yR$ has residue one. The last numerator has
residue $(1+\bar R)(1-\bar R)$, so $n\geq0$ implies
$0\le\bar R\le1$. This does not imply $R\le1$ as a source
inequality, and no such stronger claim is used.

Since $q,y$ are infinitesimal and $E,t,j,K$ bounded, the positive
numbers $qK,yt,yj,qE$ are infinitesimal. Dividing their unit
products in (HI-entry) proves
\[
 m,C,D,\alpha\quad\hbox{are all unbounded},\qquad
 A\asymp qy.
\]
This conclusion follows from the original source formulas and
$n\geq0$; no additional strong-weight assumption is needed for
an actual HI source.

\subsection{The HI contact for every $n\ge0$}
\label{component:p888d9f79ba3b}
\label{component:p4363feef60dd}
\label{p888d9f79ba3b:sec:r7-hi-n0-retained}
\label{p4363feef60dd:sec:r7-hi-full-retained}
Use the HI source estimates of Section~\ref{component:p54b8f20ad2dd}.
Let $q=u_c^{-2}$ and $y=v_c^2$ be positive infinitesimals, and put
$p=1/u$ for the testing angle. The full target here is
$\widehat F=p^2F(1/p,v)$, with its polynomial value at $p=0$.
Here the independent testing $p,v$ are infinitesimal or zero.
The source numbers $R=r/q$, $H$, and the opening quantities $t,j,K,E$
are those of the source calculation; $R$ is bounded, $H\sim1$, and
\[
 S_c=1+R+yt=j+t(1+y)=K+E(1+q).
\]
Set $d_p=1+p^2$, $d_v=1+v^2$, $\delta_p=p^2-q$,
$\delta_v=v^2-y$, and use the bounded positive source products
\[
 c=Cyt,\quad d=Dyj,\quad \mu=mqK,\quad a_*=\alpha qE.
\]
The original $n$ is included from the start. Define
\[
 \nu=2nR={(B_{\rm src}+R)(1+R)S_c-R\Delta\over\Delta},
 \qquad \zeta=R+\nu,\quad h_*=H-\nu.
\]
The exact source identities are
\[
 {2AH\over qy}=\zeta,\qquad
 c-d={y\over2}\{\zeta(1+q)-t\},\qquad
 \mu-a_*={q\over2}\{E-\zeta(1+y)\}.
\]
They cancel the apparent opening denominators before any normal is
estimated. At $R_0=r_0\in[0,1]$ they give
\[
 \nu_0={1-r_0\over3+r_0},\qquad
 c_0=d_0=\mu_0=(a_*)_0={1+r_0\over2(3+r_0)}.
\]
Thus these products remain units at $r_0=0$, even when $n$ is
unbounded. Literal $n=0$ is retained in every formula.

Define the state coordinates
\[
 A_0=W-pY,\quad B_0=vY-Z,\quad
 S=X-vW-pZ+pvY,\quad s=-S/S_c,\quad Q=\sqrt{qy}\,Y/H,
\]
\[
 P=p/\sqrt q,\quad V=v/\sqrt y,\quad T=PV,
 \quad N_c=(A_0-tv s)/\sqrt y,\quad
 N_m=(B_0+Ep s)/\sqrt q,
\]
\[
 N_D=d_vN_c+t\delta_vVs,\qquad
 N_M=d_pN_m-E\delta_pPs.
\]
The exact source has $P=V=1$, $s=Q=\sqrt q$ and all four $N$'s zero.
Retain the positive squares
\[
 \mathcal E={4cj d_v\over tS_c}N_c^2
 +{4dt(1+y)\over jS_c}N_D^2
 +{4\mu E(1+q)\over KS_c}N_M^2
 +{4a_*K d_p\over ES_c}N_m^2
\]
and $\mathcal N=4nd_pd_vx^2$, where
\[
 x=X+(R/H)pvY=-g s+(H+R)TQ+v\sqrt yN_c-p\sqrt qN_m,
 \qquad g=S_c-tv^2-Ep^2.
\]
The complete difference of the original $n$ term and this square is
\[
 \begin{split}
 4nd_pd_v(X^2-x^2)={}&4\nu d_pd_vg\,TsQ
 -2\nu(2H+R)d_pd_vT^2Q^2\\
 &-4\nu d_pd_vTQ(v\sqrt yN_c-p\sqrt qN_m).
 \end{split}
\]
All its coefficients are bounded, regardless of the size of $n$.

The exact quadratic-state expansion is
\[
 \begin{split}
 \widehat F-\mathcal E-\mathcal N={}&h_cN_c^2+h_mN_m^2
 -4\sqrt{qy}\,pv d_pd_vN_cN_m\\
 &+N_c(C_sVs+\beta_cTQ)+N_m(M_sPs+\beta_mTQ)+\mathcal G,
 \end{split}
\]
where
\[
 h_c={4d_v\over S_c}\{c(1+y)+dd_v\}+2yd_v(1+v^2+p^2v^2),
\quad
 h_m={4d_p\over S_c}\{\mu d_p+a_*(1+q)\}+2qd_p(1+p^2+p^2v^2),
\]
\[
 C_s=4d_v[y\{\zeta(1+q)-d_pg\}+2dt\delta_v/S_c],
 \quad M_s=4d_p[q\{d_vg-\zeta(1+y)\}-2\mu E\delta_p/S_c],
\]
\[
 \beta_c=4h_*\sqrt y\,v d_pd_v,
 \qquad\beta_m=-4h_*\sqrt q\,p d_pd_v,
\]
and
\[
 \mathcal G=f_AP^2s^2+f_BV^2s^2+f_CT^2Q^2+f_DQ^2-eTsQ,
\]
\[
 \begin{split}
 f_A={}&4[\mu K+a_*Ed_p-2\mu E\delta_p
                   +\mu E^2\delta_p^2/S_c]
       +2qd_v(E\delta_p-K+tv^2)^2-2E^2qv^2d_p,\\
 f_B={}&4[ctd_v+dj-2dt\delta_v+dt^2\delta_v^2/S_c]
                   +2y(t\delta_v-j)^2,\\
 f_C={}&2\{H^2-\nu(2H+R)\}d_pd_v,\quad
 f_D=2H(R+\nu)d_pd_v,\quad e=4h_*d_pd_vg.
 \end{split}
\]
This identity follows by the exact quadratic-state extraction rule:
substitute the displayed inverse coordinates, subtract the five
retained squares, collect $N_c,N_m$, and use the three source identities
above. It specifies the full scalar, both complete linear normals and
the entire hard block. Its coefficients are polynomials in $p,v$ with
bounded source coefficients. In particular $C_s,M_s,\beta_c,\beta_m$
are infinitesimal. The hard diagonal residues are $4/(3+r_0)>0$ and
\[
 (f_A)_0=(f_B)_0=(f_C)_0=(f_D)_0
 =\lambda(r_0)={2(1+r_0)^2\over3+r_0},\qquad e_0=4\lambda(r_0).
\]

For completeness the finite radical division is given here once. A
star means $p=\sqrt q$, $v=\sqrt y$ and $\sigma=\sum f_i-e$.
Because every subtracted square vanishes at the actual contact, the
original state and angular contact equations give exactly
\[
 f_A^*+f_B^*=f_C^*+f_D^*=e^*/2,
 \quad q^{1/4}\sigma_{\sqrt p}^*=-4(f_A^*-f_D^*),
 \quad y^{1/4}\sigma_{\sqrt v}^*=-4(f_B^*-f_D^*),
\]
\[
 C_s^*+\beta_c^*=M_s^*+\beta_m^*=0.
\]
If $sQ>0$, a pure-testing sign choice makes $s,Q>0$.
Put $(A_1,A_2,A_3,A_4)=(Ps,Vs,TQ,Q)$,
$G=(A_1A_2A_3A_4)^{1/4}$, $E_i=A_i-G$. Each radical separates into
pure-source and pure-testing radicals. The exact circuit is
\[
 L=\sum A_i^2-4G^2=\sum E_i^2+2G\sum E_i.
\]
For every pairing $ij|k\ell$,
\[
 G\sum E_i=G(\sqrt{A_i}-\sqrt{A_j})^2
 +G(\sqrt{A_k}-\sqrt{A_\ell})^2
 +2G\{(A_iA_j)^{1/4}-(A_kA_\ell)^{1/4}\}^2.
\]
Also $L=(A_1-A_3)^2+(A_2-A_4)^2
+2(\sqrt{A_1A_3}-\sqrt{A_2A_4})^2$, and for $rr'=G^2$,
$(r-r')^2=(r-G)^2+(r'-G)^2+2G(\sqrt r-\sqrt{r'})^2$.
A finite average supplies positive pivots simultaneously on the $E_i$
and on $F_p=\sqrt{A_1A_3}-G$, $F_v=\sqrt{A_2A_3}-G$.
The exact relations
\[
 G(\sqrt p-q^{1/4})=q^{1/4}F_p,
 \qquad G(\sqrt v-y^{1/4})=y^{1/4}F_v
\]
turn finite monic Taylor division in $\sqrt p-q^{1/4}$ and
$\sqrt v-y^{1/4}$ into bounded infinitesimal corrections to these
positive pivots. More explicitly put $z=\sqrt p$, $w=\sqrt v$,
$z_0=q^{1/4}$, $w_0=y^{1/4}$ and $k=(f_C^*+f_D^*)/2$.
Each monomial has the finite expansion
\[
 z^iw^j=\sum_{a=0}^i\sum_{b=0}^j
 {i\choose a}{j\choose b}z_0^{i-a}w_0^{j-b}
 (z-z_0)^a(w-w_0)^b.
\]
All its coefficients are bounded. In $\mathcal G-kL$, terms with two
differences multiplying $G^2$ become bounded infinitesimal multiples
of $F_p^2,F_pF_v,F_v^2$, because each replacement contributes a
source factor $z_0$ or $w_0$. Terms with one difference multiplying
$GE_i$ become infinitesimal multiples of $F_pE_i,F_vE_i$.
Terms multiplying $E_i^2$ have infinitesimal coefficients since
$f_i-k$ has zero residue. The only possible first-jet expression is
\[
 G\{2\sum_i(f_i^*-k)E_i+z_0\sigma_z^*F_p+w_0\sigma_w^*F_v\}.
\]
Using the displayed contact jets, it equals exactly
\[
 \begin{split}
 2(f_D^*-k)G\sum_iE_i
 &-{z_0\sigma_z^*\over2}
       \{\sqrt G(\sqrt{A_1}-\sqrt{A_3})\}^2\\
 &-{w_0\sigma_w^*\over2}
       \{\sqrt G(\sqrt{A_2}-\sqrt{A_3})\}^2.
 \end{split}
\]
All three source coefficients are infinitesimal, and the finite circuit
average already supplies strictly positive pivots on their roots.
For each normal, its two infinitesimal coefficients have zero sum at
contact. Write their two raw amplitudes as $G+E_i$; finite monic
division of that sum and the same two relations express the whole
normal coupling by infinitesimal entries on the already positive
roots. Thus the hard block and circuit block have strict residue.
For $sQ<0$ the raw diagonal and $\sqrt{T|sQ|}$ already have positive
residue coefficients $\lambda,\lambda,\lambda,\lambda,4\lambda$;
for $sQ=0$ omit the last root. The strict-Gram lemma proves all three
cases. Restore the five positive squares. Section~\ref{component:pcee1372511ed}
then proves the other measurements, including the literal axes and
infinities, for every $n\ge0$.

\subsection{Tests outside the HI critical corner}
\label{component:pcee1372511ed}

\label{pcee1372511ed:sec:hi-remote}
\begin{theorem}[All tests outside the HI critical corner]
Let $K_c,K_t$ be independent real closed extensions of $\mathbb R$.
Consider the original C2 parent with $d=1$, $\lambda=\beta=0$,
$n,A\geq0$, and $m,C,D,\alpha$ positive unbounded source numbers.
Then at every projective testing pair except possibly
$1/u$ and $v$ both nonzero infinitesimal, and every real or complex testing
state, the
original parent is a sum of squares in $K_c\otimes_{\mathbb R}K_t$.
No comparison of source and testing rates is made.
\end{theorem}

\begin{proof}
For any positive real $\mu$, retain the four nonnegative event
positive square terms with coefficients $m-\mu,C-\mu,D-\mu,\alpha-\mu$, and
retain the $n$ and $A$ positive square terms in their entirety. It suffices to
prove the capped parent obtained by setting
$m=C=D=\alpha=\mu$ and $n=A=0$ is tensor SOS at the given test.
The real constant $\mu$ may be chosen after the real testing
residue is fixed, since each of the original four source
coefficients is larger than every real number.

Write $D_u=1+u^2$, $D_v=1+v^2$.
First suppose $v$ is infinitesimal, including zero. Use the roots
\[
 S=vW,\qquad E=Y-uW,\qquad X,\quad Z.
\]
All four capped event amplitudes are bounded linear forms in
these roots:
\[
 M=D_uS-X-uZ+uvE,\quad uW-Y=-E,\quad
 J=-E+uvX-vZ,\quad vY-Z=uS+vE-Z.
\]
The remaining core depends only on the same four roots, so the
exact capped parent has a bounded four-root Gram.
Let $a=\bar u$ and $d_a=1+a^2$. Its residue is
\[
 (8\mu+2)E^2+
 4\mu(d_aS-X-aZ)^2+(4\mu+2)d_a(aS-Z)^2
       +2(X+aZ)^2-2d_aZ^2.
\]
Put $T=X+aZ$ and $R=aS-Z$; this is an invertible real change of
coordinates. The last three-variable form is
\[
 4\mu(d_aS-T)^2+4\mu d_aR^2+2T^2
       -2d_a a^2S^2+4d_a aSR.
\]
The $T$ and $R$ pivots are $4\mu+2$ and $4\mu d_a$.
The remaining Schur pivot is exactly
\[
 \frac{d_a[4\mu^2-a^2(4\mu+1)]}{\mu(2\mu+1)}.
\]
Taking $\mu=1+a^2$ makes its numerator positive, because
$4\mu^2-a^2(4\mu+1)=3a^2+4$.
Thus this bounded Gram has strictly positive real residue and
is tensor SOS by the inherited strict-residue Gram lemma.

Next suppose $v$ is a bounded unit. At the real residues
$(a,b)=(\bar u,\bar v)$ the four event forms have common kernel
\[
 (W,X,Y,Z)=W(1,b,a,ab).
\]
The core restricted to this line is
$2(1+a^2)(1+b^2)b^2W^2$, which is strictly positive away from zero.
Consequently a sufficiently large positive real $\mu$ makes the
whole capped residue form positive definite. For completeness,
this elementary assertion has a direct quantitative proof.
Let $E_0\succeq0$ be the event Gram, split the state space into
$\ker E_0$ and its orthogonal complement, and let $\sigma>0$ be
its smallest eigenvalue on that complement. Let $\kappa>0$ be
the core's eigenvalue on the one-dimensional kernel and bound
the norms of the two remaining core blocks by $M$.
Then $4\mu\sigma>M+M^2/\kappa$ gives a positive Schur
complement. Every constant here belongs to $\mathbb R$.
The exact state Gram is bounded and has this positive residue.

Finally suppose $v$ is unbounded, including the projective
endpoint. Set $w=1/v$, homogenizing before setting $w=0$.
The four event forms become
\[
 W-wX+uY-uwZ,\quad uW-Y,\quad
 uwW+uX-wY-Z,\quad Y-wZ,
\]
with bounded positive factors $D_w$ or $D_u$ where appropriate.
At $w=0$ their common kernel is
$(W,X,Y,Z)=(0,X,0,aX)$, on which the homogenized core is
$2(1+a^2)X^2$. The same real Schur argument gives a positive
residue Gram for a sufficiently large $\mu$.
For finite unbounded $v$, undo homogenization by the pure
positive testing square $v^2$; for the projective endpoint the
homogeneous expression is already the required parent.
These three cases exhaust every projective Bob test when $u$ is bounded.

It remains to consider $p=1/u$ infinitesimal, after homogenizing
by $p^2$. If $v$ is a bounded unit, use the four roots
\[
 S=pY,\qquad T=vY-Z,\qquad W,\quad X.
\]
The exact event amplitudes are
\[
 W-S,\quad J=W+vX-D_vS+pvT,\quad
 M=pvW-pX+T,\quad T.
\]
The signed core also has bounded coefficients in these roots:
\[
 2J^2+2D_pT^2+2D_v
 [p^2X^2+2vSX-2pXT-v^2S^2+2pvST-p^2T^2].
\]
At $p=0$, set $b=\bar v\ne0$, $A_0=W-S$, and
$B_0=X-bS$. The full capped residue is
\[
 (8\mu+2)T^2+4\mu(1+b^2)A_0^2
 +(4\mu+2)(A_0+bB_0)^2
 +4(1+b^2)bSB_0+2(1+b^2)b^2S^2.
\]
As $\mu$ grows, the first three positive penalties constrain
$T=A_0=B_0=0$, and the residual kernel has positive coefficient
$2(1+b^2)b^2S^2$. The quantitative real Schur argument above
therefore gives a finite $\mu$ with a strictly positive residue.

If $u,v$ are both projectively unbounded, put $w=1/v$ and
homogenize by $p^2w^2$. Use $P=Y-wZ$, $S=pZ$, $W,X$.
Now the exact event amplitudes are
\[
 M=pW-pwX+P,\quad W-pP-wS,\quad
 J=wW+X-pwP-D_wS,\quad P.
\]
The exact core is
\[
 2J^2+2D_pP^2+2D_w[p^2X^2+2SX-S^2].
\]
At $p=w=0$ the capped residue is exactly
\[
 4\mu W^2+(8\mu+2)P^2+4\mu(X-S)^2+2X^2,
\]
which is strict for every $\mu>0$. The literal projective
endpoints cause no division. The literal critical axes also have direct exact certificates:
\begin{align*}
 F|_{p=0}={}&4nD_vX^2+(4m+4\alpha+2)(vY-Z)^2
 +4CD_vW^2+4AD_vY^2+(4D+2)(W+vX)^2,\\
 F|_{v=0}={}&4nD_pX^2+(4m+2)(pX+Z)^2
 +(4C+4D+2)(W-pY)^2+4AD_pY^2+4\alpha D_pZ^2.
\end{align*}
Thus only $p,v$ both nonzero infinitesimal remain untreated.
Applying the real-state result separately to real and imaginary
state components proves the complex-state claim.
\end{proof}

For a strong C2 HI source, the exact source-entry lemma below proves
that all four coefficients in the theorem are unbounded. Hence the
theorem applies without any additional source-shape assumption.
Its precise uncovered complement is the critical testing corner
$1/u$ and $v$ both nonzero infinitesimal. Exact sign gauges may
make both positive.
The original HI source leaf is therefore narrowed but not closed.

\subsection{Further source reductions}
\label{component:pe56fb394b522}

\label{pe56fb394b522:sec:c2-angle-reductions}

All angles in this section are actual source contacts. They are strictly
positive finite field elements; ``infinite'' below means unbounded in
the coefficient field, not a literal projective endpoint. Write
$z=c^2=u_c^2$, $y=a^2z=v_c^2$, and
$A_c=1+zr$, $S=A_c+a^2zt$. The exact source identities used are
\[
 n\ge0\ \Longleftrightarrow\ (b+r)A_cS-r\Delta\ge0,
 \qquad
 N=(1+z)(b+r)h+a^2\{g+j+rt(1+y)\}>0.
\]

\begin{theorem}[The ordinary--ordinary source orbit is already closed]
If both actual C2 half-angles are positive bounded units, then the
source belongs either to the two-weak-complementary-event all-tests
exit or to the five-consecutive-strong-event all-tests exit. No upper
bound on an original event coefficient is required.
\end{theorem}
\begin{proof}
Here $a,z$ are units. The positive gaps give $0<h<1/z$ and
$0<t<1+zr$, so $h,t,b=h+a^2t$ are bounded. From
$g=b-1+r+zh>0$ we obtain $1<b+r+zh\le(1+z)b+r$;
thus $b+r$ is a positive unit. The positive expansion of $\Delta$
makes it a bounded unit too. The original $A$ formula has unit
numerator and bounded denominator $h$, so $A$ is strong.

The exact numerator for $n$, after substituting $t=(b-h)/a^2$,
satisfies the polynomial specialization
\[
 \left[(b+r)A_cS-r\Delta\right]_{r=1}
     =-z(a^2+h)(b+1)(z+1).
\]
If $1-r$ were infinitesimal, taking residues would contradict
$n\ge0$: the right side is strictly negative. Hence $1-r$ is a
unit. The original $m$ numerator is now a unit and $g\Delta$ is
bounded, so $m$ is strong. Apply the exact C2 reflection; it preserves
ordinary actual angles and exchanges $m,D$. Thus $D$ is strong.

If either $C$ or $\alpha$ is weak, its original rational formula
implies $N$ infinitesimal: all the multiplying denominator factors
are bounded and $z,a,b+r,\Delta$ have bounded inverses. The positive
identity for $N$ then forces $h,j$ infinitesimal. Consequently
$t=A_c-j$ and $1-zh$ are units. Both original complementary formulas
therefore give $C,\alpha$ infinitesimal. The two-weak theorem applies.
Otherwise $C,\alpha$ are both strong; with $m,A,D$ strong, the
five-consecutive-event theorem applies. These cases exhaust the source.
\end{proof}

\paragraph{Some directed one-weak angle types are empty.}
Orient an exactly-one-weak source so that $C$ is weak and $\alpha$
strong. In addition to the ordinary--ordinary exclusion above, its
actual angle types cannot be $(\mathrm{unit},\mathrm{Inf})$,
$(\mathrm{Inf},\mathrm{unit})$ or $(\mathrm{Inf},\mathrm{Inf})$.
Here $\mathrm{Inf}$ denotes positive infinitesimal.

For $(\mathrm{unit},\mathrm{Inf})$, $z$ is a unit and $a^2$ is
infinitesimal. Again $t,h,b,\Delta$ are bounded, $\Delta$ is a unit,
and $b+r$ is a unit. Weak $C$ implies $N/a^2$ infinitesimal.
The positive identity forces $j$ and $rt$ infinitesimal. Since
$t=A_c-j$ is a unit, $r$ is infinitesimal. Then $g>0$ makes
$b+zh$ a unit, and $b=h+a^2t$ makes $h$ a unit. The first positive
summand of $N$ is consequently a unit, contradicting $N/a^2$
infinitesimal.

For $(\mathrm{Inf},\mathrm{unit})$, bounded $b$ is a unit by the
inherited source-entry inequality. Explicitly,
\[
 xC(b-h)=\frac{A_c(xN)}{2\Delta},\qquad
 xN=(1+x)(b+r)(xh+k^2)-x(1+k^2)(1-r)(b-h).
\]
The last numerator has residue $k_0^2(b_0+r_0)>0$, independently
of any lower bound on $g$. Thus $xC(b-h)$ is a unit and $C$ is strong. For unbounded
$b$, (SR1) gives $C\asymp1/L$ with bounded $L$, again strong.
This exhausts $b$ and excludes weak $C$ in this angle orientation.

For $(\mathrm{Inf},\mathrm{Inf})$, both $z,y$ are infinitesimal
and $t<A_c$ is bounded. Put $H=zh$, $L=z(b-h)=yt$ and $B=zb=H+L$.
Here $0<H<1$, $L$ is infinitesimal, and $B$ is bounded. If $H$ is
a unit then $b$ is unbounded; direct division of the original
$\Delta$ by $b$ gives a positive unit of residue $1+B_0$.
The positive first summand of $N$ consequently gives
$C\ge c_0H/L$ for some positive real $c_0$, contradicting weak $C$.
Thus weak $C$ forces $H,B$ infinitesimal. The exact ratio is then
\[
 \frac\alpha C=
 \frac{z^2(b+r)(b-h)}{A_c(1-H)}
 =\frac{(B+zr)L}{A_c(1-H)},
\]
which is infinitesimal. Therefore $\alpha$ is weak as well, contrary
to the selected one-weak orientation.

After these exclusions, an exactly-one-weak source oriented with weak
$C$ has only the following candidate actual types:
\begin{gather*}
 (\mathrm{Inf},\mathrm{unbounded}),\quad
 (\mathrm{unit},\mathrm{unbounded}),\quad
 (\mathrm{unbounded},\mathrm{Inf}),\\
 (\mathrm{unbounded},\mathrm{unit}),\quad
 (\mathrm{unbounded},\mathrm{unbounded}).
\end{gather*}
The second type is the reflection of the single-angle domain and
is now fully covered by the homogeneous unbounded-$b$ theorem and the
other single-angle cases. The four other listed directed types
remain conditional source leaves; no nonemptiness is inferred from
their listing.

\paragraph{Other strong-event angle types are not all empty.}
For $x>0$ infinitesimal, the exact positive chart
\[
 c=\sqrt x,\quad a=1,\quad b=1/x,\quad r=1/4,\quad t=1/2
\]
has $h=1/x-1/2$, $g=1/x+1/4-x/2$, $j=1/2+x/4$ and all
remaining chart gaps positive. Its original coefficients are
\begin{align*}
 n&=\frac{8-5x^2}{16(1+x)^2},&
 m&=\frac{3x^2(3x+4)}{16(1+x)(4+x-2x^2)},\\
 C=\alpha&=\frac{4-x}{8x},&
 A&=\frac{x(x+4)(3x+4)}{32(2-x)(1+x)^2},\\
 D&=\frac{(2-x)(3x+4)}{16x(1+x)}.
\end{align*}
Every coefficient is positive. The unchanged exact contact is
\[
 (u_c,v_c;W,X,Y,Z)
 =(\sqrt x,\sqrt x;1/x,-\sqrt x/4,1/\sqrt x-\sqrt x/2,1).
\]
Original source stationarity, positive transverse state inertia and
Laurent convexity prove physicality, including projective limits.
Here $C,\alpha,D$ are unbounded, $m\sim3x^2/16$, $A\sim x/4$,
and $n\to1/2$. Thus a genuine strong-complementary source can have
two infinitesimal actual angles and weak $m,A$. This is an obstruction
to deleting that angle orbit, not a counterexample to tensor SOS.

\subsection{A strong $C$ event from an unbounded first angle}
\label{component:pb7189b9c1f6e}

\label{pb7189b9c1f6e:sec:c2-weak-hi-hu-empty}

\begin{theorem}[The directed weak-$C$ types HI and HU are empty]
Use the exact positive rational C2 source, with $d=1$ and $n\geq0$.
If the first actual half-angle $u_c$ is unbounded and the second
actual half-angle $v_c$ is bounded, then $1/C$ is bounded.
Consequently no exactly-one-weak source oriented with $C$ weak can
have directed actual type HI or HU.  Literal $n=0$ is included.
\end{theorem}

This is an original-source impossibility theorem.  It makes no
claim that a scalar contact identity by itself is a tensor-SOS
certificate.  All independent testing variables and state coordinates
remain unrestricted; the two source leaves are empty before any
testing restriction is imposed.

\begin{proof}
Write
\[
 q=u_c^{-2}>0,\qquad y=v_c^2>0,
 \qquad H=u_c^2h=v_cY_c\in(0,1),
\]
and retain the original $r\in(0,1)$.  Here $q$ is infinitesimal
and $y$ is bounded; $y$ may itself be infinitesimal.  Put
\[
 R=r/q=-u_cX_c,\quad B=u_cv_cW_c,\quad
 L=v_c(u_cW_c-Y_c),\quad T=B-yR.
\]
The exact positive gaps give
\[
 B=H+L=yR+T,\qquad 0<T<H+y,\qquad L>0.
\]
In particular $H,T,y,r$ are bounded.  The original source coordinates
become exactly
\[
 z=1/q,\quad a^2=yq,\quad b=yr+qT,\quad h=qH,
 \quad t=r/q+(T-H)/y,
\]
and the two gaps satisfy
\[
             g=H+(1+y)r-1+qT>0,
 \qquad j=(H+y-T)/y>0.                             \tag{WE1}
\]
The displayed expressions are pure coefficient identities; no
inverse of a mixed source/testing element is used.

We first show that $R$ is bounded.  Set $Q=q^2\Delta>0$.
Direct expansion of the unchanged original formula for $n$ gives
\[
                         n=-\frac{qP}{2rQ},        \tag{WE2}
\]
where the following exact decomposition of $P$ is sufficient:
\begin{align*}
P={}&r^2(1+y)(H+yr)-Tqr(y+2)-q^2(T^2+T)+P_+,\\
P_+={}&HTq^2+HTqr+Hqr(y+1)+T^2q^2r\\
     &+Tq^2r(y+1)+2Tqr^2(y+1)+qr^2y(y+1)>0.        \tag{WE3}
\end{align*}
Suppose for contradiction that $R$ is unbounded, and put
$\kappa=q/r=1/R$, which is infinitesimal.  Let $S=H+yr>0$.
The ratio $T/S$ is bounded.  Indeed, if $r$ is a unit, then
\[
                T<H+y\leq S/r.
\]
If $r$ is infinitesimal, (WE1) and boundedness of $y,T$ imply
$H_0=1$, so $S$ is a unit and $T/S$ is again bounded.
Since $T$ is bounded, $(T^2+T)/S$ is bounded as well.  Dividing
the lower bound in (WE3) by the positive $r^2S$ gives
\[
 \frac{P}{r^2S}\geq (1+y)
       -\kappa\frac{T(y+2)}S
       -\kappa^2\frac{T^2+T}S.
\]
The last two terms are infinitesimal, while $1+y\geq1$.
Thus $P>0$.  Equation (WE2) contradicts $n\geq0$, including
the equality case $n=0$.  Therefore $R$ is bounded.

Now $r=qR$ is infinitesimal.  Equation (WE1) becomes
\[
             g=H-1+q\{(1+y)R+T\}>0,
\]
so $H_0=1$.  Both $B=yR+T$ and $L=B-H$ are bounded.
The original denominator has the exact form
\begin{align*}
 \Delta={}&(1+q)\{B^2+2BR+(1+y)B\}
           +(1-yq)R^2+2(1+y)R+1+y.                \tag{WE4}
\end{align*}
It is a positive bounded unit: all quantities are bounded, and
$1-yq$ is a positive unit, with constant contribution $1+y\geq1$.
The unchanged numerator $N_{\rm src}$ satisfies
\[
 \frac{N_{\rm src}}q
  =(1+q)(B+R)(H+y)-(1+y)(1-qR)L.                  \tag{WE5}
\]
Since $H_0=1$ and $L=B-H$, its residue is
\[
       (N_{\rm src}/q)_0=(1+y_0)(1+R_0)>0.
\]
Thus $N_{\rm src}/q$ is a positive bounded unit.  Finally,
$A_c=1+R$ and $a^2t=qL$, so the original coefficient is exactly
\[
          C=\frac{(1+R)(N_{\rm src}/q)}{2L\Delta}.
                                                               \tag{WE6}
\]
The numerator has bounded positive inverse, and $L\Delta$ is
bounded.  Therefore $1/C$ is bounded, as asserted.  No lower bound
on $L$ was required; if $L$ is infinitesimal the conclusion is
stronger.
\end{proof}

The exact C2 reflection preserves directed type HI and exchanges
$C,\alpha$.  Applying the theorem in both orientations shows that
every physical HI source has both complementary positive square terms strong.
For type HU the reflected type is UI; the reflected consequence is
that $\alpha$ is automatically strong at every UI actual contact.
These observations respect the directed weak-$C$ orientation; they
do not identify IH with HI.

\section{Sources with five lower-bounded events}
\subsection{Five consecutive lower-bounded events}
\label{component:p80253fe9bc1b}

\label{p80253fe9bc1b:sec:five-retained-events}

The following argument permits every retained event coefficient to be
unbounded. It requires no canonical omission, actual-contact coordinates,
or source-shape entry. Five consecutive events in the eight-cycle suffice.

\begin{theorem}[Five consecutive lower bounds]
\label{p80253fe9bc1b:thm:five-retained-events}
Let the original canonical source be physical, with $d>0$ and all eight
event coefficients nonnegative. Suppose
\[
 d/m,\ d/c,\ d/A,\ d/D,\ d/\alpha
 \quad\hbox{are positive bounded coefficient elements}.
\]
Then the entire original parent is tensor SOS at every independent
projective measurement pair and every real or complex projective state.
There is no upper bound on any event coefficient divided by $d$.
By exact CHSH relabeling the same conclusion holds if any five
consecutive event coefficients in the cycle
$(\lambda,D,c,A,\alpha,m,\beta,n)$ have these lower bounds.
\end{theorem}

\begin{proof}
Divide by the positive pure coefficient $d$, so $d=1$.
Let $\mu=\min(1,m,c,A,D,\alpha)$. It is a positive bounded unit.
Select the parent $F_\mu$ with $m=c=A=D=\alpha=\mu$ and
$\lambda=n=\beta=0$. Keep every nonnegative excess and all three
unselected coefficients as their original event squares. This uses only
pure coefficient square roots. It makes no assertion that $F_\mu$ is
globally physical.

At noncommuting residue tests apply
the inherited noncommuting-residue theorem to the full physical source. It
performs its own positive coefficient normalization if the event
weights are unbounded.
It remains to certify $F_\mu$ near every commuting residue test. Its
coefficients are bounded, so a positive-definite real residue Gram is
sufficient by the bounded positive-residue Gram lemma.

Write $\mathcal H$ for the signed core with $d=1$ and
$D_u=1+u^2,D_v=1+v^2$. On the four open commuting axes its exact
homogeneous restrictions are
\[
\begin{aligned}
\mathcal H|_{u=0}&=2D_v(X^2+Y^2),\\
\mathcal H|_{v=0}&=2\{(uW-Y)^2+(X+uZ)^2\},\\
\mathcal H|_{u=\infty}&=2\{(W+vX)^2+(vY-Z)^2\},\\
\mathcal H|_{v=\infty}&=2D_u(X^2+Y^2).
\end{aligned}                                                   \tag{NW1}
\]
The remaining finite angle has nonzero residue. At $u=0$, the retained
$m,D$ positive square terms force $W=Z=0$ on the core kernel. At $v=0$, the retained
$A,\alpha$ positive square terms force $Y=Z=0$, and then $W=X=0$. At $u=\infty$,
the retained $c,A$ positive square terms force $W=Y=0$, and then $X=Z=0$. At
$v=\infty$, the retained $c,D$ positive square terms force $W=Z=0$ after the core
has forced $X=Y=0$. Thus the
selected residue state matrix is positive definite on every open axis,
for either sign of the other finite angle.

Normalize a nonzero testing state by a coordinate of largest absolute
value. Its coordinates are then bounded and one is a unit. At the four
double corners, write $F_{ab}$ for the homogeneous restriction of
$F_\mu$ to $(u,v)=(a,b)$. Directly,
\[
\begin{aligned}
F_{00}&=(2+4\mu)X^2+(2+12\mu)Y^2+4\mu Z^2,\\
F_{\infty0}&=(2+8\mu)(W^2+Z^2)+4\mu Y^2,\\
F_{0\infty}&=4\mu(W^2+Z^2)+2X^2+(2+12\mu)Y^2,\\
F_{\infty\infty}&=4\mu W^2+(2+4\mu)X^2+(2+12\mu)Y^2.
\end{aligned}                                                    \tag{NW2}
\]
Here infinity means the polynomial homogeneous chart, obtained by
dividing the finite parent by the appropriate pure-testing squared
angle. Every form is strict except its explicitly missing state
coordinate. Consequently the only nontrivial residue states are the
$W$ ray at $(0,0)$, the $X$ ray at $(\infty,0)$,
and the $Z$ ray at $(\infty,\infty)$.

At $(0,0)$ normalize $W=1$. The five local coordinates
$(u,v,X,Y,Z)$ are infinitesimal. The complete polynomial $F_\mu$
has no constant or linear term, and its degree-two part is
\[
 (2+8\mu)(u-Y)^2+2X^2
       +4\mu\{(v-X)^2+Y^2+Z^2\}.                    \tag{NW3}
\]
Its residue is positive definite: vanishing successively gives
$X=Y=Z=v=u=0$.

At the corner $(\infty,0)$ normalize $X=1$ and put $p=1/u$.
The five infinitesimal coordinates are $(p,v,W,Y,Z)$. The complete
homogeneous polynomial again has no terms below degree two, and its
quadratic part is
\[
 (2+4\mu)(W+v)^2+(2+4\mu)(p+Z)^2
                         +4\mu(W^2+Y^2+Z^2).        \tag{NW4}
\]
This is positive definite. At the corner $(\infty,\infty)$,
normalize $Z=1$ and put $p=1/u,q=1/v$. On $(p,q,W,X,Y)$ the exact
quadratic part is
\[
 (2+8\mu)(q-Y)^2+2X^2
                         +4\mu\{(p-X)^2+Y^2+W^2\}. \tag{NW5}
\]
It is positive definite as well. One can obtain (NW5) directly from
(NW3) by the exact symmetry
$(u,v;W,X,Y,Z)\mapsto(1/v,1/u;Z,X,Y,W)$, which permutes the five
equal retained weights.

For completeness, these quadratic calculations control the entire
polynomial. In each local chart, allocate every higher monomial to a
product of two of the five coordinate roots. The remaining monomial
has positive degree, hence is infinitesimal; its coefficient is bounded
because it is a real polynomial in the bounded unit $\mu$. Symmetrizing
gives an exact finite five-root Gram whose residue is (NW3), (NW4), or
(NW5). No inversion, root, or comparison mixing the two fields occurs.
The bounded positive-residue Gram lemma therefore certifies the full
selected polynomial, including all higher terms.

Restore the excess event squares, the positive pure-testing homogeneous
and projective state factors, and the positive coefficient $d$. The
zero state is immediate. Splitting a complex state into real and
imaginary parts proves the complex assertion.

Finally the exact full-parent Alice reflection acts on the displayed
cycle indices as $i\mapsto7-i$, and the swap/reflection
$(u,v;W,X,Y,Z)\mapsto(1/v,1/u;Z,X,Y,W)$ acts as $i\mapsto6-i$.
Their composition gives a one-step cycle rotation. Their full-parent
identities have only positive pure-testing homogeneous factors, so they
transport the five-positive square term conclusion to every cyclic translate.
\end{proof}

This theorem is a sufficient source domain, not an exhaustive zero-$d$
source cover. It requires five selected events to be positive and
noninfinitesimal relative to $d$. Its all-tests conclusion is independent
of actual contact angles and of all original event hierarchies. In
particular, it contains both C1 and C2 sources with all six retained
events bounded below relative to $d$.

\section{Two small actual angles: common identities}
\subsection{Source coefficients and normal arrays}
\label{component:pd709027c7805}

The following are the exact native scalar and state-normal coefficient arrays.
Put $r=1-e$, $z=aB$, $y=tB$, and
\[
 c_a=1+aBr,\quad S_c=1+B(1-h+ar),\quad
 g=1+a(hB+r-1),\quad j=c_a-(1-h)/t,
\]
\begin{align*}
\Delta={}&a+(1+aB)+(1+aB)B+2aB(1+aB)r+a^3B^2r^2+2a^2Br\\
 &+tB\{1+a+aB+a^2Br(2-r)\},\\
M={}&(1+aB)(1+ar)(h+t)-a(1+tB)e(1-h).
\end{align*}
The physical stationary event coefficients are
\[
\begin{aligned}
n&=\frac{(1+ar)c_aS_c-r\Delta}{2r\Delta},&
m&=\frac{a^2(1+tB)e c_aS_c}{2g\Delta},\\
C&=\frac{c_aM}{2(1-h)\Delta},&
Ah&=\frac{t c_a(1+ar)S_c}{2\Delta},\\
D&=\frac{(1+aB)(1+ar)(1-atBr)S_c}{2tj\Delta},&
\alpha&=\frac{B^2(1+ar)M}{2(1-hB)\Delta}.
\end{aligned}
\]
The following scalar and normal arrays are exact rational functions of
the native stationary source. Each case certificate proves the source
bounds it needs at its own independent openings; the arrays themselves
have no bounded-$a$ or old R9-domain restriction.

\paragraph{Scalar arrays.}
The notation $Ah$ in these formulas means the product of $A$ and $h$.
\[
\begin{aligned}
q_s&=nr^2+(m+\tfrac12)e^2,\qquad q=q_s+\alpha,\\
a_w&=(C+D+\tfrac12)(1-h)^2+Ah\,h,\\
b_2&=2mh+C(1-h)^2+Ah\,h+(\alpha+\tfrac12)h^2,\\
b_4&=h^2(m+\alpha+\tfrac12)+Ah\,h,\\
d_2&=(n+\tfrac12)r^2+2Dr,\qquad
d_4=(n+D+\tfrac12)r^2,\\
e_0&=me+D(1-h)+\alpha h+\tfrac12,\\
e_1&=mhe+r(D+\tfrac12)(1-h)+(\alpha+\tfrac12)h,\\
c_F&=(b_2-a^2nr^2)/t^2,\qquad
\lambda=(m-a^2q)/(ta^2B),\\
f_{A0}&=Ah\,h/t^2+d_2-2e_1/t+2anr^2/t,\\
f_{C0}&=(\alpha-c_Ft^2B^2)/(taB^2),\\
f_G'&=-\{a_w/t^2+D-2e_0/t+2aq/t+2c_FtB\}/(aB).
\end{aligned}
\]
The potentially singular $a,B$ factors cancel in the combinations
used in RC1. One direct proof is to clear the denominators in the displayed source
coefficients and the printed monomials, and then expand; all three resulting numerator polynomials are
zero. Equivalently, the original value and two first derivatives at
contact give precisely RC1. Thus RC1 is an exact identity, including
arbitrarily small $a$ and literal $n=0$.

\paragraph{All first-normal arrays.}
For the angular weight $q=q_s+\alpha$ define
\[
\begin{aligned}
K_X'&=-nr+(m+\tfrac12)e-aq,\\
K_E'&=a_w/t^2-e_0/t+aq/t,\\
K_Y&=\{-(C+D+\tfrac12)(1-h)+Ah\}/t+D-\alpha,\\
x_3&=2\{(D+\tfrac12)[(1-h)/t-1]+e/2-nr(1+a/t)-mh/t\},\\
x_5&=-2r(n+D+\tfrac12),\\
a_E&=2Ah\,h/t^2-2e_1/t,\qquad b_E=2b_2/t^2,\qquad c_E=2b_4/t^2,\\
a_Y&=2\{Ah/t-me+r(D+\tfrac12)-\alpha-\tfrac12\},\\
b_Y&=2\{-C(1-h)+Ah+m+(\alpha+\tfrac12)h\}/t,\\
c_Y&=2\{Ah+h(m+\alpha+\tfrac12)\}/t.
\end{aligned}
\]
\paragraph{The six exact contact balances.}
With $x=aB$, $y=tB$, the scalar balances are
\[
 f_{A0}=\lambda-2d_4x-b_4y/t^2,\qquad
 f_{C0}=\lambda+b_4y/t^2,\qquad
 f_G'=3\lambda-d_4x+b_4y/t^2.\tag{RC1}
\]
The complete first-normal balances are
\[
 2K_E'+a_Ex+b_Ey+c_Exy+2\lambda x=0,
\]
\[
 2K_X'+x_3y+x_5xy-2t\lambda x=0,\qquad
 2K_Y+a_Yx+b_Yy+c_Yxy=0.\tag{RC6}
\]
Substitution of the native weights and multiplication by the displayed
source denominators gives six zero polynomial numerators. Equivalently
these are the original contact value and its five first derivatives.
They underlie the generic whole-parent identity BA7 (also printed as
R10.9). BA6 explicitly gives its entire hard quadratic. Those identities
hold before any source size assumption; the bounded-$a$, R8 and R10
certificates separately establish their coefficient bounds and pivots.

\subsection{The five-root identity}
\label{component:p9ec3d6bbd717}
Set $\kappa=1/t$, $\ell=1/a$. The letters $e_0,e_1$ are the native
scalar coefficients above. In this subsection $a_g$ denotes the scalar
coefficient called $a$ in the five-root table, and not the source ratio.
\paragraph{Five homogeneous contact roots.}
Put $U=u^2$, $V=v^2$, and define
\[
 \xi=X+ruZ,\qquad q=\sqrt A\,(Y-huW),\qquad
 R_1=uW-\kappa vZ,\quad R_2=vW-\ell uZ,\quad
 T=uvW-BZ.                                             \tag{WA3}
\]
Every root vanishes at the actual source contact. Inverses and
radicals in their coefficients are pure source. The large $A$
positive square term is retained in $q$; no positive square term is subtracted from
the parent and no physicality hypothesis is transferred to a
modified remainder.

Use the following source coefficients for the exact scalar graph:
\begin{align*}
 a_g&=(C+D+\tfrac12)(1-h)^2+Ah^2,&a_4&=Ah^2,\\
 b_2&=2mh+C(1-h)^2+(A+\alpha+\tfrac12)h^2,
 &b_4&=h^2(m+A+\alpha+\tfrac12),\\
 c&=\alpha,&c_2&=nr^2+(m+\tfrac12)(1-r)^2+\alpha,\\
 c_4&=nr^2,&d_2&=(n+\tfrac12)r^2+2Dr,\\
 d_4&=(n+D+\tfrac12)r^2.
\end{align*}
Thus the graph on $\xi=q=0$ is
\[
 \begin{split}
 F_g/4={}&[a_gU+a_4U^2+(m+b_2U+b_4U^2)V]W^2\\
 &+[c+c_2U+c_4U^2+(D+d_2U+d_4U^2)V]Z^2
       -2uv(e_0+e_1U)WZ,
 \end{split}
\]
where $e_0,e_1$ are exactly the original coefficients in (WS3).
Set
\[
 \begin{gathered}
 \beta_s=\frac{Bd_4+\kappa c_4}{\kappa\ell^2},\qquad
 \gamma_s=b_4+\frac{d_4}{\kappa\ell},\\
 j_0=\frac{Bd_4-a_4\kappa^2\ell+d_2\ell}
                   {2\kappa\ell^2},\qquad
 j_1=\frac{d_4}{2\kappa\ell},\qquad J_s(U)=j_0+j_1U,\\
 f_s=\frac12\left(b_2-\frac{c_4}{\ell^2}
                -\frac c{B^2}-\frac{Bd_4}{\kappa\ell^2}\right).
 \end{gathered}                                        \tag{WA4}
\]
The four original scalar contact equations (value and the three
logarithmic derivatives in $u,v,W/Z$) give the exact identity
\[
 F_g/4=(R_1,R_2,T)\,\mathsf S\,(R_1,R_2,T)^T,\qquad
 \mathsf S=
 \begin{pmatrix}
 a_g+a_4U&uvJ_s&-vJ_s\\
 uvJ_s&m+\beta_sU&u(f_s-j_1U)\\
 -vJ_s&u(f_s-j_1U)&c/B^2+\gamma_sU
 \end{pmatrix}.                                       \tag{WA5}
\]
This is a different scalar identity from the excluded polynomial
two-block splitter (JG7). In particular, the last diagonal retains
the full radial mass $\alpha/B^2$ and the Gram depends on the
independent testing variables. The source bounds for this exact identity are established in each coefficient domain below.

\paragraph{Both full first normals and the complete hard block.}
For the $\xi$ normal define
\[
 \begin{gathered}
 p_0=\frac{D+r(n+1/2)}\kappa,\qquad
 p_1=\frac{r(n+D+1/2)}\kappa,\qquad
 q_1=\frac{nr+Bp_1}\ell,\\
 f_x=-mh+(D+\tfrac12)(1-h)-p_0-q_1,\\
 \mathsf L_x=\bigl(uv(p_0+p_1U),\ -m+q_1U,\
                                u(f_x-p_1U)\bigr).
 \end{gathered}                                        \tag{WA6}
\]
For the unnormalized $\eta=Y-huW$ normal set
\[
 \begin{gathered}
 p_y=-(C+D+\tfrac12)(1-h)+Ah,\qquad p_y'=Ah,\\
 k_v=-C(1-h)+Ah+m+(\alpha+\tfrac12)h,\qquad
 k_{uv}=h(A+m+\alpha+\tfrac12),\\
 k_z=-m(1-r)+r(D+\tfrac12)-\alpha-\tfrac12,\\
 q_y=\frac{-k_z-\kappa p_y'-Bk_{uv}}\ell,\qquad
 f_y=k_v-q_y,\\
 \mathsf L_y=\bigl(p_y+p_y'U,\ uvq_y,\
                                      v(f_y+k_{uv}U)\bigr).
 \end{gathered}                                        \tag{WA7}
\]
The two remaining original source stationarity equations, in $X,Y$,
give these exact representations of the complete first normals.
No first-normal term is discarded. Define also
\[
 \begin{split}
 h_{xx}&=n(1+U)(1+V)+m+(D+\tfrac12)UV+\tfrac12(1+V),\\
 h_{xy}&=-uv(m+D+\tfrac12),\\
 h_{yy}&=A(1+U)(1+V)+C(1+V)+D+\tfrac12
                   +mUV+(\alpha+\tfrac12)(1+U)V.
 \end{split}
\]
Then the \emph{entire original four-state parent} is exactly
\[
 F/4=\mathcal R^T\mathsf G(u,v)\mathcal R,\qquad
 \mathcal R=(\xi,q,R_1,R_2,T)^T,
\]
\[
 \mathsf G=
 \begin{pmatrix}
 h_{xx}&h_{xy}/\sqrt A&\mathsf L_x\\
 h_{xy}/\sqrt A&h_{yy}/A&\mathsf L_y/\sqrt A\\
 \mathsf L_x^T&\mathsf L_y^T/\sqrt A&\mathsf S
 \end{pmatrix}.                                       \tag{WA8}
\]
Block notation in (WA8) has sizes $1,1,3$. Every coefficient of this finite polynomial matrix is pure source.
 It then verifies (WA5)--(WA7) from the corresponding
four plus two contact equations. Thus (WA8) is not merely a graph
or Hessian certificate.

\subsection{Division at the source-ratio zero}
\label{component:p296e824cb97a}
Use the native coefficient dictionary and put $\mu=m/a^2$,
$\alpha_*=\alpha/B^2$.

The only new potentially unbounded local coefficient is
$(b_2-\alpha_*)/a$, where
\[
 b_2=2a^2\mu h+C(1-h)^2+(Ah)h+(B^2\alpha_*+1/2)h^2.
\]
At $a=0$, with $d_*=1+B(1+t)$, the exact values are
\[
 C=\frac{h+t}{2(1-h)d_*},\quad
 Ah=\frac{t[1+B(1-h)]}{2d_*},\quad
 \alpha_* =\frac{h+t}{2(1-Bh)d_*}.
\]
Substitution gives $b_2=\alpha_*$ identically for arbitrary
$B,h,t$. Clearing the unit denominators in (WE1), the numerator
of $b_2-\alpha_*$ is therefore divisible by $a$ in the polynomial
ring. Its quotient has a unit denominator at $B=a=h=0$.
Hence
\[
 Q_*=(b_2-\alpha_*)/a\quad\hbox{is bounded}.       \tag{WE3}
\]
This is a joint regularity argument, so it permits arbitrary
$h/a$, including an unbounded ratio. 

\paragraph{The complete local five-root Gram survives.}
Keep $\kappa=B/y$ and every coefficient and normal in
(WA3)--(WA8), but set $\ell=1/a$ and replace $R_2$ by
\[
 \widetilde R_2=aR_2=avW-uZ.
\]
The entire full-parent matrix is the exact congruence
\[
 \widetilde G=EGE,\qquad E=\operatorname{diag}(1,1,1,a^{-1},1),
 \qquad \widetilde{\mathcal R}=(\xi,q,R_1,\widetilde R_2,T).
                                                               \tag{WE4}
\]
Its newly rescaled coefficients are
\[
\begin{gathered}
 m/a^2=\mu,\qquad
 \beta_s/a^2=(Bd_4+\kappa c_4)/\kappa,\\
 J_s(U)/a=\frac{aBd_4-a_4\kappa^2+d_2+d_4U}{2\kappa},\qquad
 f_s/a=\tfrac12[Q_*-ac_4-aBd_4/\kappa],\\
 (-m+q_1U)/a=-a\mu+(nr+Bp_1)U,\\
 q_y/a=-k_z-\kappa p_y'-Bk_{uv}.                  \end{gathered}\tag{WE5}
\]

\subsection{The angular circuit}
\label{sec:c2-short-angular-identity}
For $\lambda=(m-a^2q)/(ta^2B)>0$ and $q=nr^2+(m+1/2)e^2+\alpha$,
the following identity holds before any source-size assumption.
Set $U=u^2$, $V=v^2$, $k=\sqrt{x}$, and $c=\sqrt{a/t}$.
Use the state-linear forms
\[
 E=tuW-vZ,\quad \eta=Y-huW=\sqrt h\,N_y,\quad
 \xi=X+ruZ,\quad F_v=(V-y)Z,\quad R_D=v(U-x)Z,
\]
\[
 J=(uv-x/c)Z,\quad Q=\sqrt q(avW-uZ-a\xi),\quad
 N=\sqrt n\{r[(a/t)V-U]Z+u\xi\},
\]
\[
 R=\sqrt\lambda\,tk(cv-u)W,\qquad
 T=\sqrt{2\lambda/c}\,k\sqrt{uv}(tcW-Z).              \tag{BA4}
\]
The square root in $T$ is pure testing. Initially take $u,v>0$;
literal axes are addressed below.

For a compact coefficient table put
\[
 \begin{aligned}
 b_2&=2mh+Cl^2+(A+\alpha+\tfrac12)h^2,&
 b_4&=h^2(m+A+\alpha+\tfrac12),\\
 d_4&=(n+D+\tfrac12)r^2,&
 e_1&=mhe+r(D+\tfrac12)l+(\alpha+\tfrac12)h,\\
 c_F&=(b_2-a^2nr^2)/t^2,&
 a_E&=2Ah^2/t^2-2e_1/t,\\
 b_E&=2b_2/t^2,&c_E&=2b_4/t^2,\\
 a_Y&=2[Ah/t-me+r(D+\tfrac12)-\alpha-\tfrac12],\\
 b_Y&=2[-Cl+Ah+m+(\alpha+\tfrac12)h]/t,\\
 c_Y&=2[Ah+h(m+\alpha+\tfrac12)]/t,\\
 x_3&=2[(D+\tfrac12)(l/t-1)+e/2-nr(1+a/t)-mh/t],&
 x_5&=-2r(n+D+\tfrac12).
 \end{aligned}                                                     \tag{BA5}
\]
Let $Y_*=hE/t+\eta$ and define the entire hard quadratic explicitly by
\[
 \begin{split}
 H_0={}&CD_v(lE/t-\eta)^2+AD_uD_vY_*^2
 +(D+\tfrac12)(lE/t-\eta+uv\xi)^2\\
 &+(\alpha+\tfrac12)D_uVY_*^2
 +[n(1+V+UV)+D_v/2]\xi^2\\
 &+2m(VE/t-uv\xi)Y_*+mUVY_*^2.
 \end{split}                                                       \tag{BA6}
\]
In particular BA6 has no inverse testing angle and no $a$ denominator.
Upon replacing $\eta$ by $\sqrt hN_y$, every occurrence of $A$
is through $Ah$, $\sqrt hAh$, or $hAh$.

With $\sigma=t\lambda x/\sqrt q$, the exact identity is
\[
 \begin{split}
 P={}&Q^2+N^2+T^2+R^2+\lambda J^2+d_4R_D^2
 +(c_F+b_4x/t^2)F_v^2+(b_4/t^2)vR_DF_v\\
 &+H_0-\lambda xE^2-\mu\xi^2-2\sigma Q\xi\\
 &+E[(a_E+c_EV)R_D+(b_E+c_Ex)vF_v]\\
 &+\sqrt hN_y[(a_Y+c_YV)R_D+(b_Y+c_Yx)vF_v]\\
 &+\xi[x_5uvR_D+u(x_3+x_5x)F_v].
 \end{split}\tag{BA7}
\]
This is the generic RC/R10 identity, not a scalar restriction.
One pen-and-paper verification is to substitute BA4 into BA7 and use
\[
 R^2+T^2+\lambda J^2
 =\lambda\{t^2x[U+(a/t)V]W^2-4txuvWZ+(UV+xy)Z^2\}.
\]
The hard quadratic is then exactly BA6. The remaining six scalar
and first-normal coefficient identities follow by substituting the native source fractions in Section~\ref{component:pd709027c7805};
a common denominator leaves zero polynomial numerators. Equivalently
these six identities are the original value and all five first-contact
conditions. Both the full identity and those six numerator equalities
have also been independently expanded during verification.

\subsection{Homogeneous identities for unit $t>1$}
The following polynomial identities hold before any source-size
specialization. Their strict residues in this subsection use
$t_0=\tau>1$, $C\sim c_0=\tau/2$, $D\sim d_0=1/[2(\tau-1)]$,
$A^{-1/2}$ infinitesimal and $n,m,\alpha$ infinitesimal. Each source
calculation establishes its own coefficient bounds; in particular R10.6
below supplies the exact modification when $m$ is merely bounded.
Put $\widehat Y=\sqrt A Y$, so $Y=A^{-1/2}\widehat Y$.

First suppose $u$ is a bounded unit or unbounded and $v$ is
infinitesimal. Put $p=1/u$, $D_p=1+p^2$ and use
\[
 N_\infty=\sqrt n X,\quad V_X=vX,\quad K=pX+Z,\quad
 R_\infty=\sqrt\alpha pX,\quad W,\quad\widehat Y.
\]
Let $J_\infty=W+D_pV_X-pY-pvK$. The entire homogeneous parent is
\[
 \begin{split}
 p^2P(1/p,v)={}&D_pD_vN_\infty^2+m(pvW+vY-K)^2
 +CD_v(W-pY)^2+D_pD_v\widehat Y^2\\
 &+DJ_\infty^2+D_p[R_\infty+\sqrt\alpha(vY-K)]^2\\
 &+\tfrac12\{J_\infty^2+(1-p^2v^2)K^2\\
 &\hspace{9mm}+D_p[v^2Y^2-2vYK+2pYV_X+2pvKV_X-p^2V_X^2]\}.
 \end{split}                                                       \tag{BA11}
\]
It follows directly by substituting $Z=K-pX$ in the quarter-parent $P=F/4$ defined by \eqref{eq:c2-master-parent}. Its bounded
six-root residue, with $d_p=1+p_0^2$, is
\[
 d_p(N_\infty^2+R_\infty^2+\widehat Y^2)+K^2/2
 +c_0W^2+(d_0+\tfrac12)(W+d_pV_X)^2-d_pp_0^2V_X^2/2.
\]
The $W,V_X$ determinant is $d_pc_0(d_0+1/2)>0$, using
$c_0(d_0+1/2)=(c_0+d_0+1/2)/2$. All other pivots are positive.
No inverse of $n$ or $\alpha$ occurs, and either root may vanish.
This covers literal $u=\infty$ by $p=0$ as well.

Next suppose $u$ is infinitesimal and $v$ is a bounded unit.
Use the five roots
\[
 (S,R_u,X,\widehat Y,Z)=(\sqrt mW,uW,X,\sqrt A Y,Z).
\]
The exact full polynomial is obtained from the quarter-parent $P=F/4$ defined by \eqref{eq:c2-master-parent} by the two replacements
\[
 m(vW-X+uvY-uZ)^2=(vS-\sqrt mX+\sqrt m\,uvY-\sqrt m\,uZ)^2,
 \qquad uW=R_u
\]
in all other terms. This operation uses no inverse of $m$ or $u$.
Its separate residue pivots on $S,X,\widehat Y$ are
$V,(1+V)/2,1+V$, where $V=v_0^2>0$. The remaining $R_u,Z$ block is
\[
 \begin{pmatrix}
 c_0(1+V)+d_0+1/2&-(d_0+1/2)v_0\\
 -(d_0+1/2)v_0&d_0V
 \end{pmatrix},\qquad \det=c_0d_0V^2>0.             \tag{BA12}
\]
Here $c_0d_0-d_0/2-1/4=0$. If instead $v$ is unbounded, use
$q_v=1/v$ and the same five roots. Multiplying the same exact
polynomial by $q_v^2$ gives, at $u=q_v=0$, the diagonal residue
\[
 (S,R_u,X,\widehat Y,Z):\quad (1,c_0,1/2,1,d_0)>0.   \tag{BA13}
\]
This homogeneous substitution is polynomial and also applies at
literal $q_v=0$.

If both finite testing residues are nonzero, the hard
$\widehat Y$ pivot is $(1+U)(1+V)$. The complete remaining
three-state residue on $(W,X,Z)$ is
\[
 c_0(1+V)UW^2+(d_0+\tfrac12)(u_0W+u_0v_0X-v_0Z)^2
 +\tfrac12(1+V)(X+u_0Z)^2-\tfrac12(1+U)VZ^2,
\]
where $U=u_0^2,V=v_0^2$. Its determinant is
\[
 (1+U)(1+V)UV\left\{\frac{c_0(2d_0+1)}4U
                              +\frac{c_0d_0}2V\right\}>0.        \tag{BA14}
\]
The $Z=0$ restriction is strictly positive, so the full matrix is
positive definite. Every coefficient is bounded after the hard-state
normalization.

The one-infinity cases not already covered by BA11 or BA13 are
strict directly. At $u=\infty$ and a bounded unit $v$, the residue is
\[
 c_0(1+V)W^2+(d_0+\tfrac12)(W+v_0X)^2+Z^2/2
 +(1+V)\widehat Y^2\succ0.
\]
At $v=\infty$ and a bounded unit $u$, it is
\[
 c_0UW^2+(1+U)X^2/2+d_0(u_0X-Z)^2
 +(1+U)\widehat Y^2\succ0.
\]
These are the degree-two homogeneous values of the quarter-parent $P=F/4$ defined by \eqref{eq:c2-master-parent} in the respective
angles; nearby reciprocal coordinates are pure-testing infinitesimals.

For two infinite angles put $p=1/u$, $q_v=1/v$ and introduce
$P_\infty=pZ$, $Q_\infty=q_vZ$. The entire exact polynomial is
\[
 \begin{split}
 p^2q_v^2P(1/p,1/q_v)={}&n(1+p^2)(1+q_v^2)X^2
 +m(pW-pq_vX+Y-Q_\infty)^2\\
 &+C(1+q_v^2)(W-pY)^2
 +(1+p^2)(1+q_v^2)\widehat Y^2\\
 &+D(q_vW+X-pq_vY-P_\infty)^2
 +\alpha(1+p^2)(Y-Q_\infty)^2\\
 &+\tfrac12\{(q_vW+X-pq_vY-P_\infty)^2
 +(1+p^2)(Y-Q_\infty)^2\\
 &\hspace{9mm}+(1+q_v^2)[p^2X^2+2XP_\infty-P_\infty^2]\}.
 \end{split}                                                       \tag{BA15}
\]
Its residue on $(W,X,\widehat Y,P_\infty,Q_\infty)$ is
\[
 c_0W^2+(d_0+\tfrac12)X^2-2d_0XP_\infty
 +d_0P_\infty^2+\widehat Y^2+Q_\infty^2/2\succ0.
\]
The nontrivial determinant is $d_0/2>0$. At $p=q_v=0$ the two
extra roots vanish, respecting the original null state rather than
inventing a positive $Z$ pivot.

\subsection{Bounded source ratio}
\label{sec:c2-bounded-a-new}

Let $K_c,K_t$ be independently labelled real closed extensions of
$\mathbb R$. All source letters below belong to $K_c$, and all testing
letters belong to $K_t$. No algebraic-disjointness hypothesis is imposed.
Write $\mathrm I$ for positive infinitesimals and call a positive bounded
element a unit when its reciprocal is bounded.

We use the actual stationary C2 source, normalized by $d=1$, with
$\lambda_{\rm event}=\beta=0$. Put
\[
 x=aB,\quad y=tB,\quad l=1-h,\quad e=1-r,\quad
 g=1-ae+haB,\quad j=1+aBr-l/t.
\]
The source assumptions are
\begin{equation}
 \begin{gathered}
 B,x,y\in\mathrm I,\qquad a>0\text{ bounded},\qquad
 t\text{ a positive bounded unit},\\
 0<h,r<1,\qquad g,j,1-txr>0,\qquad n\ge0,\qquad
 C\ge c_*>0\quad(c_*\in\mathbb R).
 \end{gathered}
 \label{eq:c2-ba-domain}
\end{equation}
The last strict gap is exactly the original $B-xyr>0$ divided by $B$.
The source formulas, which also fix our normalization, are
\[
 c_a=1+aBr,\quad S=1+B(l+ar),
\]
\[
 \begin{split}
 \Delta={}&a+(1+aB)+(1+aB)B+2aB(1+aB)r+a^3B^2r^2+2a^2Br\\
 &+tB\{1+a+aB+a^2Br(2-r)\},\\
 M={}&(1+aB)(1+ar)(h+t)-a(1+tB)el,
 \end{split}
\]
\begin{equation}
 \begin{aligned}
 n&={ (1+ar)c_aS-r\Delta\over2r\Delta},&
 m&={a^2(1+tB)e c_aS\over2g\Delta},\\
 C&={c_aM\over2l\Delta},&
 Ah&={t c_a(1+ar)S\over2\Delta},\\
 D&={(1+aB)(1+ar)(1-taBr)S\over2t j\Delta},&
 \alpha&={B^2(1+ar)M\over2(1-hB)\Delta}.
 \end{aligned}
 \label{eq:c2-ba-source}
\end{equation}
The original source is physical; no physicality assertion about a
polynomial obtained by a formal coefficient substitution will be used.

\begin{proposition}
\label{prop:c2-bounded-a}
Under \eqref{eq:c2-ba-domain}, the entire C2 parent is a finite sum of
squares in $K_c\otimes_{\mathbb R}K_t$ for every independent projective
testing pair and every state. The conclusion includes $n=0$, every
relative rate of $g,j,h,l,r,e$, and $a$ infinitesimal in residue.
\end{proposition}

For finite tests set $D_u=1+u^2$, $D_v=1+v^2$ and define the quarter-parent
\begin{equation}
 \begin{split}
 P={}&m(vW-X+uvY-uZ)^2+nD_uD_vX^2+CD_v(uW-Y)^2+AD_uD_vY^2\\
 &+(D+\tfrac12)(uW+uvX-Y-vZ)^2
 +(\alpha+\tfrac12)D_u(vY-Z)^2\\
 &+\tfrac12D_v(X+uZ)^2-\tfrac12D_uD_vZ^2.
 \end{split}
 \label{eq:c2-ba-parent}
\end{equation}
All identities below are identities of this four-state polynomial.
\paragraph{Gram completion and real caps.}
Apply Lemma~\ref{p25af2c58da7e:lem:gram} to the bounded remainder after
retaining the displayed nonnegative squares. The roots may be unbounded;
only its matrix and positive-definite real residue enter this lemma.

For real caps, let a real quadratic form be positive definite on the
common kernel of finitely many squared rows. A sufficiently large real
multiple of those squares makes the form positive definite: decompose
into the kernel and its orthogonal complement and take the Schur
complement. Each unbounded source weight admits that cap; a finite weight
admits a sufficiently close smaller cap when the finite-residue form is
strict. Retain every nonnegative excess.

\paragraph{Source bounds when $e$ is a unit.}
The denominator has residue $1+a_0>0$. The physical numerator has the
exact identity
\begin{equation}
 2r\Delta n=eQ_s-BK_s\ge0,
 \quad
 \begin{aligned}
 Q_s={}&1+B(1+t+ah+2ar+at)\\
       &+B^2\{a(h+t)+a^2[h(r+1)+t(1+r-r^2)]\},\\
 K_s={}&(1+a)(h+t)(1+aB).
 \end{aligned}
 \label{eq:c2-ba-physical}
\end{equation}
Consequently $(nr)_0=e_0/[2(1+a_0)]>0$. In particular $n$ has a
positive real lower bound and $n=0$ is impossible on this branch.
The displayed source formulas give positive bounded units
\begin{equation}
 nr,\qquad mg/a^2,\qquad Ah,\qquad Cl,\qquad Dj,
 \qquad \alpha_*:=\alpha/B^2.
 \label{eq:c2-ba-unit-products}
\end{equation}
For $Cl$ and $\alpha_*$ one needs the actual strong-$C$ assumption:
if $l$ is a unit it makes $M$ a unit; if $l$ is infinitesimal then
$h_0=1$ and $M_0=(1+a_0r_0)(1+t_0)>0$. This explicitly excludes the
weak-$C$ corner where $M_0=0$. Thus $A,C,D,n$ have positive real
lower bounds. If $g$ is a unit, $m/a^2$ is a positive unit; if $g$ is
infinitesimal, $a$ is a unit and $m$ is unbounded.

\paragraph{One exact five-root matrix for the unit-$e$ local chart.}
Use the complete WA matrix, including both normal rows, in the roots
\[
 z=(\xi,\eta,R_1,R_2',T)^T,
 \quad \xi=X+ruZ,\quad\eta=Y-huW,
\]
\[
 R_1=uW-vZ/t,\qquad R_2'=avW-uZ,\qquad T=uvW-BZ,
 \qquad P=z^TGz.
\]
The precise generic input here is WA3--WA8: replace its normalized
$q=\sqrt A\eta$ by $\eta$, and its $R_2$ by $R_2'/a$.
It is an exact full-parent identity derived from the six original
contact equations, not an assertion about a Hessian.

The full divergent $D$ channel has the exact identity
\begin{equation}
 P=(D-d_0)L_D^2+
 z^T\{G|_{D=d_0}+(D-d_0)jH_D\}z,
 \label{eq:c2-ba-D-five}
\end{equation}
where $0<d_0<D$ is real,
\[
 L_D=uv\xi-\eta+lR_1+ruvR_2'-arvT,
\]
and the only nonzero symmetric entries of $H_D$ are
\[
 (H_D)_{13}=tuv,\quad (H_D)_{15}=-tu,\quad
 (H_D)_{34}=rtuv,\quad (H_D)_{35}=-artv.
\]
For completeness the coefficient identity proving this is
\[
 \partial_DG=dd^T+\tfrac12(sk^T+ks^T)+jH_D,
\]
\[
 \begin{aligned}
 d&=(uv,-1,l,ruv,-arv)^T,\\
 s&=(0,0,(aB-u^2)v,-(B-v^2/t)u,u^2-av^2/t)^T,\\
 k&=(-2rtu,0,0,-r^2tu,ar^2t)^T.
 \end{aligned}
\]
The root syzygy $s^Tz=0$ and $d^Tz=L_D$ prove
\eqref{eq:c2-ba-D-five} including every normal entry. Since $Dj$ is
bounded, the correction is bounded and vanishes in local residue.

If $g$ is infinitesimal, also retain $(m-K)R_{mc}^2$, where $K>0$ is
real and
\[
 R_{mc}=-\xi+uv\eta+{1-g\over a}R_2'+huT,
 \qquad \tau=(m-K)g.
\]
The remaining exact matrix is $G|_{m=K}+\tau H_m$, before the
independent $D$ operation. In unscaled $R_2=R_2'/a$ coordinates the
only nonzero symmetric entries are
\[
 (H_m)_{14}=-1,\quad (H_m)_{25}=v,\quad
 (H_m)_{44}=2-g,\quad(H_m)_{45}=hu.
\]
This is the full centered-$m$ identity R5.2--R5.5. It commutes with
\eqref{eq:c2-ba-D-five}; neither correction depends on the other
coefficient. Here $a$ is a unit and $\tau$ is bounded.

Put $N=n+1/2$, $L=A+C+d_0+1/2$ and normalize
$\widehat\xi=\sqrt N\xi$, $\widehat\eta=\sqrt L\eta$.
Every coefficient of the resulting five-root matrix is bounded.
Here are the complete bounds needed from the WA table. Its scalar
entries involve $nr^2$, $Cl^2$, $Ah^2$, $\alpha_*$, the capped $m$,
and bounded $D$; its first normals involve $nr$, $Cl$, $Ah$ and these
same bounded products; the hard block uses only $n/N$, $A/L$, $C/L$
and bounded entries divided by $\sqrt N$ or $\sqrt L$. All are bounded
by \eqref{eq:c2-ba-unit-products}. If $a$ is infinitesimal, then $g$
is a unit and the only apparently singular scaled entry is
\begin{equation}
 {b_2-\alpha_*\over a},\qquad
 b_2=2mh+Cl^2+(Ah)h+(B^2\alpha_*+1/2)h^2.
 \label{eq:c2-ba-a-divisibility}
\end{equation}
The WE3 numerator vanishes exactly at $a=0$, for arbitrary
$B,h,r,t$. Cancel that factor in the actual source formula. Its remaining
denominator uses only $g,\Delta,t,1-hB$, all units here; the $h,l,r,j$
denominators have already canceled in the displayed products. The WE5
rewrites cover both normal rows. Thus this is joint regularity, including
all simultaneous openings, rather than separate one-variable limits.

At $u=v=0$ the residue splits into two $2\times2$ blocks and a scalar.
The $(\widehat\eta,R_1)$ block has determinant
\begin{equation}
 {A(C+d_0+1/2)\over A+C+d_0+1/2},\qquad
 {1\over\det}={1\over A}+{1\over C+d_0+1/2}.
 \label{eq:c2-ba-eta-determinant}
\end{equation}
Its upper-left entry is $1$, its other entries are bounded, and the
determinant is a positive unit. If $g$ is a unit, the
$(\widehat\xi,R_2')$ block is
\[
 \begin{pmatrix}
 1+m/N&-m/(a\sqrt N)\\-m/(a\sqrt N)&m/a^2
 \end{pmatrix},\qquad \det=m/a^2>0.
\]
If $g$ is infinitesimal, its determinant is
\[
 {K+2\tau_0-\tau_0^2\nu\over a_0^2}>0,
\]
where $\nu=\operatorname{st}(N^{-1})$, also when $N$ is unbounded.
Choose $K>2\tau_0^2+1$, since $0\le\nu\le2$.
The remaining $T$ pivot is $\alpha_*>0$. Therefore the complete bounded
matrix has strict residue. This proves every local independent test on
the unit-$e$ branch without any contact conductor.

\paragraph{The small-$e$ branch: source limits and physical Schur bound.}
Now $e$ is infinitesimal and $r_0=1$. The actual source formulas imply
\begin{equation}
 n,m,\alpha\in\mathrm I\cup\{0\},\quad g_0=1,\quad
 (Ah)_0=t_0/2,\quad (Cl)_0=(h_0+t_0)/2,\quad
 (Dj)_0=1/(2t_0),\quad
 (\alpha_*)_0=(1+a_0)(h_0+t_0)/2.
 \label{eq:c2-ba-small-e-products}
\end{equation}
In particular $A,C,D$ have positive real lower bounds, even if the
strong-$C$ hypothesis had not been imposed. From $j>0$ we have $t_0\ge l_0$.
Define
\[
 q=nr^2+(m+1/2)e^2+\alpha,\quad
 \mu=m-a^2q,\quad \lambda={\mu\over ta^2B},\quad
 \sigma={t\lambda aB\over\sqrt q}.
\]
The $a^2B$ factor in $\mu$ cancels exactly in the native rational
formula. Its canceled denominator at $B=0,r=1$, with $h$ free, is
$2t(a+1)$; hence
\begin{equation}
 \lambda_0={h_0+t_0\over2t_0}>0,\qquad \mu\in\mathrm I.
 \label{eq:c2-ba-free-h-lambda}
\end{equation}
This includes $a_0=0$ and does not assume $h$ infinitesimal.

The exact inequality \eqref{eq:c2-ba-physical} gives
$e/B\ge K_s/Q_s$, where $Q_s\sim1$ and
$K_s\sim(1+a)(h+t)$. Since $q\ge e^2/2+\alpha$, it follows that
\begin{equation}
 \operatorname{st}(\sigma^2)
 \le {a_0^2(h_0+t_0)\over
 2(1+a_0)\{1+(1+a_0)(h_0+t_0)\}}.
 \label{eq:c2-ba-free-h-sigma}
\end{equation}
If $e/B$ is unbounded the left side is zero. In all cases
\begin{equation}
 \operatorname{st}(1/2-\mu-\sigma^2)
 \ge {1+a_0+(h_0+t_0)(1+2a_0)\over
 2(1+a_0)\{1+(1+a_0)(h_0+t_0)\}}>0.
 \label{eq:c2-ba-free-h-schur}
\end{equation}
No division by $n$ occurs, so $n=0$ is retained.

\paragraph{The exact small-$e$ full-parent decomposition.}
Use the generic BA4--BA7 identity and its complete displayed $H_0$.
In particular take
\[
 E=tuW-vZ,\quad\eta=Y-huW,\quad\xi=X+ruZ,
 \quad F_v=(v^2-y)Z,\quad R_D=v(u^2-x)Z,
\]
\[
 J=(uv-\sqrt{xy})Z,\quad
 Q=\sqrt q(avW-uZ-a\xi),\quad
 N=\sqrt n\{r[(a/t)v^2-u^2]Z+u\xi\}.
\]
With $k=\sqrt x$ and $c=\sqrt{a/t}$ retain the whole nonnegative positive square term
\[
 T^2={2\lambda\over c}k^2uv(tcW-Z)^2,
 \qquad R=\sqrt\lambda\,tk(cv-u)W.
\]
Initially $u,v>0$; the literal axes and the other signs are restored below.
The generic identity has the exact circuit
\[
 R^2+T^2+\lambda J^2
 =\lambda\{t^2x[u^2+(a/t)v^2]W^2-4txuvWZ
                     +(u^2v^2+xy)Z^2\}.
\]
Let $\mathscr P=P-N^2-T^2$. Its entire $D$ channel is
\begin{equation}
 \mathscr P=D L_D^2+\mathscr P|_{D=0}-2Dj\,uF_v\xi,
 \qquad L_D=lE/t-rR_D-\eta+uv\xi.
 \label{eq:c2-ba-D-circuit}
\end{equation}
This identity holds before any source specialization. Indeed
the $D$ terms are
\[
 D(lE/t-\eta+uv\xi)^2+Dr^2R_D^2
 -2Dr(lE/t-\eta+uv\xi)R_D
 +2D(l/t-1-rx)uF_v\xi,
\]
whose sum is \eqref{eq:c2-ba-D-circuit} by $j=1+xr-l/t$.

Choose positive real caps $A_c<A$, $C_c<C$, $d_0<D$. Subtract the
following exact nonnegative positive square terms:
\begin{equation}
 N^2+T^2+(D-d_0)L_D^2
 +(A-A_c)D_uD_v\eta^2+(C-C_c)D_v\eta^2.
 \label{eq:c2-ba-excesses}
\end{equation}
The remainder is the full BA quadratic with $D$ formally set to $d_0$,
minus $2(D-d_0)j\,uF_v\xi$, and with precisely the two displayed
centered-$\eta$ squares removed. Every coefficient is bounded. The
apparent $A,C$ poles are now only $Ah,Cl$; for example
\[
 \begin{split}
 AD_uD_v(hE/t+\eta)^2-(A-A_c)D_uD_v\eta^2
 &=D_uD_v\{A_c\eta^2+2(Ah)E\eta/t+(Ah)hE^2/t^2\},\\
 CD_v(lE/t-\eta)^2-(C-C_c)D_v\eta^2
 &=D_v\{C_c\eta^2-2(Cl)E\eta/t+(Cl)lE^2/t^2\}.
 \end{split}
\]
Every other entry of BA5--BA7 uses the bounded products in
\eqref{eq:c2-ba-small-e-products}, $a,t,t^{-1}$, or the bounded
$\lambda,\sigma$. The correction in \eqref{eq:c2-ba-D-circuit} is
bounded because $Dj$ is bounded. There is no $a$, $h$, $l$, or $j$
inverse left in the matrix coefficients. Pure-source inverse factors
inside the separately retained positive square term $T^2$ do not need to
be bounded.

At $u=v=0$ write $d_+=d_0+1/2$ and use residues of the bounded source
letters on the right. The entire nontrivial hard block is exactly
\begin{equation}
 d_+(R_D-lE/t+\eta)^2
 +{t+hl\over2t^2}E^2-{h\over t}E\eta+(A_c+C_c)\eta^2.
 \label{eq:c2-ba-hard-completion}
\end{equation}
The remaining blocks are
\[
 Q^2-2\sigma_0Q\xi+\tfrac12\xi^2,
 \quad {h_0+t_0\over2t_0^2}F_v^2,
 \quad \lambda_0J^2,\quad R^2.
\]
The first is strict by \eqref{eq:c2-ba-free-h-schur}. To make
\eqref{eq:c2-ba-hard-completion} strict it suffices to choose the caps
so that $2(A_c+C_c)(t_0+h_0l_0)>h_0^2$. If $h_0,l_0>0$, the actual
finite source residues satisfy
\[
 A_0={t_0\over2h_0},\quad C_0={h_0+t_0\over2l_0},\quad
 2(A_0+C_0)(t_0+h_0l_0)-h_0^2
 ={t_0(t_0+h_0)\over h_0l_0}>0.
\]
Thus sufficiently close smaller real caps work. If $h_0=0$ or
$l_0=0$, respectively $A$ or $C$ is unbounded, so a sufficiently large
real cap works. This proves strictness also when $j$ is infinitesimal
or opens simultaneously with either endpoint. The full bounded-matrix
completion principle proves the local small-$e$ claim.

\paragraph{A single critical family supplies every weak-$m$ outer join.}
When $e$ is infinitesimal, or when $e$ is a unit and $a$ infinitesimal,
$m$ is infinitesimal. The finite coefficients $A,C,D$ or their
unbounded limits have the common critical source residues
\begin{equation}
 A_0={t_0\over2h_0},\qquad C_0={h_0+t_0\over2l_0},\qquad
 D_0={1\over2(t_0-l_0)},\qquad h_0+l_0=1,\quad t_0\ge l_0.
 \label{eq:c2-ba-critical-family}
\end{equation}
Each quotient is used only when its denominator is positive. Its zero
denominator means the corresponding coefficient is unbounded. We retain
the whole $\alpha$ positive square term. We retain the $n$ positive square term too, or cap it at
a positive real lower bound on the unit-$e$ branch. In both cases the
remaining $X$ mass from the signed core is at least $D_v/2$.

For $u$ infinitesimal and $v$ a bounded unit put $S_m=\sqrt m W$ and
$R_u=uW$. Rewrite the entire $m$ positive square term as
\[
 (vS_m-\sqrt mX+\sqrt m\,uvY-\sqrt m\,uZ)^2
\]
and replace every other $uW$ by $R_u$. These operations introduce no
inverse. The $S_m$ pivot is $v_0^2>0$ and the $X$ pivot is at least
$D_v/2>0$. The remaining $(R_u,Y,Z)$ residue is
\begin{equation}
 \begin{split}
 F_v^{\rm crit}={}&C_0(1+V)(R_u-Y)^2+A_0(1+V)Y^2
 +D_0(R_u-Y-v_0Z)^2\\
 &+\tfrac12\{(R_u-Y)^2+VY^2-2v_0R_uZ\},\qquad V=v_0^2>0.
 \end{split}
 \label{eq:c2-ba-outer-v}
\end{equation}
In the finite interior its $Z=0$ restriction is strict and
\[
 \det F_v^{\rm crit}
 ={V^2(h_0+t_0)^2(1+V)\over8h_0l_0(t_0-l_0)}>0.
\]
At $h_0=0<t_0-1$, the unbounded $A$ row imposes $Y=0$ and the
remaining determinant is $C_0D_0V^2>0$. At $l_0=0$, the unbounded
$C$ row imposes $R_u=Y$ and the determinant is
$V^2(t_0+1)/(4t_0)>0$. At $t_0=l_0$, $h_0>0$, the unbounded $D$
row imposes $R_u=Y+v_0Z$ and the determinant is
$V^2(1+V)/(4h_0l_0)>0$. At the only simultaneous endpoint
$h_0=0$, $t_0=l_0=1$, both $A,D$ are unbounded; on their common
kernel $Y=0$, $R_u=v_0Z$ the remaining form is $V^2Z^2/2>0$.
These exhaust the boundary of \eqref{eq:c2-ba-critical-family} because
$t_0>0$. The cap principle now supplies a bounded strict full matrix
in every case, with all positive original excesses restored.

For $u$ a bounded unit, or unbounded, and $v$ infinitesimal set
$p=1/u$, $K=pX+Z$, $B_v=vX$, and $D_p=1+p^2$. The complete identity is
\begin{equation}
 \begin{split}
 p^2P(1/p,v)={}&nD_pD_vX^2+m(pvW+vY-K)^2
 +CD_v(W-pY)^2+AD_pD_vY^2\\
 &+D J_p^2+\alpha D_p(pX+vY-K)^2\\
 &+\tfrac12\{J_p^2+(1-p^2v^2)K^2\\
 &\hspace{10mm}+D_p[v^2Y^2-2vYK+2pYB_v+2pvKB_v-p^2B_v^2]\},\\
 J_p={}&W+D_pB_v-pY-pvK.
 \end{split}
 \label{eq:c2-ba-outer-u-full}
\end{equation}
The $n,m,\alpha$ positive square terms are retained exactly. On writing
$L=W-pY$, the remaining residue is $K^2/2$ plus
\begin{equation}
 C_0L^2+D_pA_0Y^2+(D_0+1/2)(L+D_pB_v)^2
 +D_ppYB_v-\tfrac12D_pp^2B_v^2.
 \label{eq:c2-ba-outer-u}
\end{equation}
In the finite interior its $L,Y$ diagonal block is strict and its
$B_v$ Schur complement is $D_p(h_0+t_0)/(2t_0)>0$, since
\[
 {C_0(D_0+1/2)\over C_0+D_0+1/2}
 =\tfrac12+{1\over4A_0}={h_0+t_0\over2t_0}.
\]
The same four boundary kernels give the following positive transverse
Schur coefficients: $D_p/2$ when $h_0=0<t_0-1$;
$D_p(D_0+1/2)$ when $l_0=0$;
$D_pC_0$ when $t_0=l_0$, $h_0>0$;
and $D_p/2$ at $h_0=0$, $t_0=l_0=1$.
Thus the cap principle proves every such chart, including literal $p=0$.
There is no inverse of $v$ or of a source opening in these roots.

If $e$ is a unit and $a$ is a unit, $m$ has a positive real lower
bound instead. For finite tests with only one nonzero residue, cap
$n,m,A,C,D$ at positive real lower bounds and restore their exact
excess positive square terms. At $v_0=0$, $u_0\ne0$ a strict selected residue is
\[
 n_cD_uX^2+A_cD_uY^2+(C_c+D_c+1/2)(u_0W-Y)^2
                         +\tfrac12(X+u_0Z)^2.
\]
At $u_0=0$, $v_0\ne0$ it is
\[
 (n_c+1/2)D_vX^2+m_c(v_0W-X)^2
 +(A_c+C_c+1/2)D_vY^2+D_c(Y+v_0Z)^2.
\]
For both finite testing residues nonzero, apply the common
all-noncommuting-residues theorem to the unchanged physical parent.
Its arbitrary-coefficient-scale version is the only remote physical
theorem imported here; the cap arguments above do not assert that
an arbitrary clipped source is physical.

\paragraph{All remaining projective charts.}
On the unit-$e$ branch, $n,A,C,D$ have positive real lower bounds.
After capping them and retaining all excesses and the whole $m,\alpha$
positive square terms, the three exact homogeneous templates R4E14--R4E16 have strict
residues
\[
 n_cD_vX^2+C_cD_vW^2+A_cD_vY^2
 +(D_c+1/2)(W+v_0X)^2+\tfrac12(v_0Y-Z)^2
 \quad(u=\infty),
\]
\[
 (n_c+1/2)D_uX^2+C_c(R_u-Y)^2+(A_c+1/2)D_uY^2
 +D_c(u_0X-Z)^2\quad(v=\infty),
\]
\[
 (n_c+1/2)X^2+C_cW^2+A_cY^2+D_c(X-P_\infty)^2
 +\tfrac12(Y-Q_\infty)^2\quad(u=v=\infty),
\]
where $R_u=uW$, $P_\infty=Z/u$, $Q_\infty=Z/v$ are legal roots;
the displayed matrices are strict on their respective root lists.
Their entire homogeneous polynomials have bounded coefficients.

On the small-$e$ branch, $m$ is infinitesimal and $n$ is retained whole.
The $u$-infinite, $v$-small chart is already
\eqref{eq:c2-ba-outer-u-full} at $p_0=0$. For $u$ infinite and $v$ a
bounded unit, the selected residue is
\[
 C_cD_vW^2+A_cD_vY^2+(D_c+1/2)(W+v_0X)^2
 +\tfrac12(v_0Y-Z)^2\succ0.
\]
For $v$ infinite and $u$ bounded use $S_m=\sqrt mW$ and $R_u=uW$.
The complete homogeneous polynomial obtained from
\eqref{eq:c2-ba-parent} has strict selected residue
\[
 S_m^2+\tfrac12D_uX^2+C_c(R_u-Y)^2
 +(A_c+1/2)D_uY^2+D_c(u_0X-Z)^2\succ0.
\]
For two infinite tests put $p=1/u$, $q_v=1/v$,
$P_\infty=pZ$, $Q_\infty=q_vZ$. The exact half-core is
\[
 \tfrac12\{(q_vW+X-pq_vY-P_\infty)^2
 +(1+p^2)(Y-Q_\infty)^2
 +(1+q_v^2)(p^2X^2+2XP_\infty-P_\infty^2)\}.
\]
Together with the capped $A,C,D$ positive square terms its residue is
\[
 C_cW^2+(D_c+1/2)X^2-2D_cXP_\infty+D_cP_\infty^2
 +A_cY^2+\tfrac12(Y-Q_\infty)^2\succ0;
\]
the $X,P_\infty$ determinant is $D_c/2>0$. All omitted original
positive square terms are exact nonnegative squares. Reciprocal deviations are pure
testing infinitesimals, so the same completion principle applies. The
literal endpoints are the homogeneous polynomial values, not a division
by a zero coordinate.

Finally, the literal finite axes obey the complete identities
\begin{equation}
 P|_{u=0}=(n+1/2)D_vX^2+m(vW-X)^2
 +(C+A+1/2)D_vY^2+D(Y+vZ)^2+\alpha(vY-Z)^2,
\label{eq:c2-literal-u-axis}
\end{equation}
\begin{equation}
 P|_{v=0}=nD_uX^2+(m+1/2)(X+uZ)^2
 +(C+D+1/2)(uW-Y)^2+AD_uY^2+\alpha D_uZ^2.
\label{eq:c2-literal-v-axis}
\end{equation}
The sign gauges $(u,Y,Z)\mapsto(-u,-Y,-Z)$ and
$(v,X,Z)\mapsto(-v,-X,-Z)$ preserve the parent. Applying the real
identity to the real and imaginary parts proves the complex-state
statement. The two local branches, the two finite one-unit-angle
branches, the noncommuting finite branch, and the displayed homogeneous
charts exhaust every independent source and test type. This proves the
proposition.

\subsection{State signs and projective endpoints}
\label{component:pf77de02d80c4}

\label{pf77de02d80c4:sec:rederived-state-join}

We reduce the complete weak-C2 parent to a positive state-ratio chart
and evaluate its projective measurement endpoints. The subsequent local
certificates treat the remaining positive chart.

Let $K_c,K_t\supset\mathbb R$ be real closed fields, occupying separately
labelled tensor factors.  All quantities $u,v,W,X,Y,Z$ below belong to $K_t$;
all six coefficients $n,m,A,C,D,\alpha$ belong to $K_c$ and are nonnegative.
Write $D_u=1+u^2$, $D_v=1+v^2$, and define the complete parent
\begin{align}
 F={}&nD_uD_vX^2+m\{v(W+uY)-X-uZ\}^2
       +AD_uD_vY^2+CD_v(uW-Y)^2\\
    &+D(uW+uvX-Y-vZ)^2+\alpha D_u(vY-Z)^2+K,\tag{SJ1}\\
 K={}&\tfrac12(uW+uvX-Y-vZ)^2
       +\tfrac12D_u(vY-Z)^2
       +\tfrac12D_v(X+uZ)^2-\tfrac12D_uD_vZ^2.\tag{SJ2}
\end{align}
Products involving source and testing quantities mean tensor products.
For example, $nD_uD_vX^2$ means
$n\otimes(D_uD_vX^2)$; a testing-only term means $1\otimes K$.
The positive square terms in (SJ1) are sums of squares in
$K_c\otimes_{\mathbb R}K_t$, since each coefficient and each testing
multiplier has a square root in its own real closed field.

\paragraph{Sign gauges preserve the entire parent.}
For $\varepsilon,\eta\in\{1,-1\}$, the substitution
\[
 (u,v,W,X,Y,Z)\longmapsto
 (\varepsilon u,\eta v,W,\eta X,\varepsilon Y,
  \varepsilon\eta Z)                                      \tag{SJ3}
\]
preserves every summand of (SJ1), including (SJ2).  Indeed the roots of the
$m,C,D,\alpha$ positive square terms acquire, respectively, the signs
$\eta,\varepsilon,\varepsilon,\varepsilon\eta$, and the other roots also
change only by a sign.  Thus $u,v$ may both be taken nonnegative without
altering the source coefficients or the full expression.

\paragraph{An exact positive hard block.}
Set
\[
 t=\frac{X+uZ}{D_u},\qquad r=\frac{Y-uW}{D_u},\qquad
 w=\frac{W+uY}{D_u},\qquad z=\frac{Z-uX}{D_u}.             \tag{SJ4}
\]
The inverse transformation is
$X=t-uz$, $Z=z+ut$, $Y=r+uw$, $W=w-ur$.
Let $H=1+u^2+v^2$ and $d=H^2-u^2v^2$.  Both
$H-uv=1+(u-v)^2+uv$ and $H+uv$ are positive for $u,v\geq0$;
in fact $H>|uv|$ for arbitrary $u,v$.  Consequently $d>0$.
Define
\[
 \mathsf A=\frac12\begin{pmatrix}H&-uv\\-uv&H\end{pmatrix},\qquad
 \mathsf B=\frac12\begin{pmatrix}-u^2v&-uv^2\\uv^2&u^2v\end{pmatrix}.
\]
The matrix $\mathsf A$ is positive definite.  Direct expansion and square
completion give the exact identity
\begin{align}
 \frac K{D_u}
 ={}&\left(\binom tr+\mathsf A^{-1}\mathsf B\binom wz\right)^{\!T}
       \mathsf A
       \left(\binom tr+\mathsf A^{-1}\mathsf B\binom wz\right)\\
    &+\frac{D_uD_v}{2d}
       \{u^2v^2(w^2+z^2)-2uvH wz\}.                    \tag{SJ5}
\end{align}
The rational matrix used in the shift is explicitly
\[
 \mathsf A^{-1}\mathsf B=
 \frac1d\begin{pmatrix}
 -u^2v(1+u^2)&-uv^2(1+v^2)\\
 uv^2(1+v^2)&u^2v(1+u^2)
 \end{pmatrix}.                                         \tag{SJ6}
\]
For any $a,b\in K_t$, the hard quadratic form is
\[
 \binom{a}{b}^{\!T}\mathsf A\binom{a}{b}
   =\frac{H-uv}{4}(a+b)^2+\frac{H+uv}{4}(a-b)^2.         \tag{SJ7}
\]
This also gives an explicit square decomposition of the first term of
(SJ5), with square roots taken only in $K_t$.

If $u,v\geq0$ and $wz\leq0$, the second term of (SJ5) is nonnegative.
It too is a square in $K_t$ (or can be split into its three nonnegative
terms).  Hence $K\geq0$, and (SJ1) is a tensor sum of squares immediately.
In particular the axes $w=0$ and $z=0$ are included, without division by
either state coordinate.  The same identity, without sign gauges, proves
the sharper implication $K<0\Longrightarrow uvwz>0$.

\paragraph{Reduction of the remaining states.}
The only remaining case after (SJ3) is $K<0$, $u,v>0$, $wz>0$.
Then $w\ne0$.  Divide all four state coordinates by $w$.  Equation (SJ4)
now has $w'=1$ and $z'=z/w>0$.  The ratio $z'$ is arbitrary: it may be
infinitesimal, bounded, or unbounded in $K_t$.  No estimate in this
reduction replaces $z'$ by a bounded coefficient.

The parent is homogeneous of degree two in the state coordinates, so
\[
 F(u,v;W,X,Y,Z)=w^2F(u,v;W/w,X/w,Y/w,Z/w).               \tag{SJ8}
\]
If the parent on the right has a certificate $\sum_i h_i^2$ in
$K_c\otimes_{\mathbb R}K_t$, then the parent on the left has certificate
$\sum_i((1\otimes w)h_i)^2$.  This proves that certificates for the full
parent with $u,v>0$, $w=1$, and arbitrary positive $z$ suffice for all
finite affine testing parameters and all real states.

\paragraph{Literal projective endpoints.}
For clarity, the projective endpoints can be checked without a limiting
argument.  Homogenize $u=U_1/U_0$, $v=V_1/V_0$, and put
$D_U=U_0^2+U_1^2$, $D_V=V_0^2+V_1^2$, with neither pair identically zero.
Multiplication by $U_0^2V_0^2$ turns (SJ2) into
\begin{align}
 K_h={}&\tfrac12(U_1V_0W+U_1V_1X-U_0V_0Y-U_0V_1Z)^2\\
       &+\tfrac12D_U(V_1Y-V_0Z)^2
        +\tfrac12D_V(U_0X+U_1Z)^2-\tfrac12D_UD_VZ^2.     \tag{SJ9}
\end{align}
At the four axes this is, exactly,
\begin{align}
 U_0=0:\quad &K_h=\tfrac12U_1^2\{(V_0W+V_1X)^2+(V_1Y-V_0Z)^2\},\\
 V_0=0:\quad &K_h=\tfrac12V_1^2D_U(X^2+Y^2),\\
 U_1=0:\quad &K_h=\tfrac12U_0^2D_V(X^2+Y^2),\\
 V_1=0:\quad &K_h=\tfrac12V_0^2\{(U_1W-U_0Y)^2+(U_0X+U_1Z)^2\}.
                                                               \tag{SJ10}
\end{align}
All other homogenized positive square terms in (SJ1) retain nonnegative coefficients
and are sums of squares.  Thus the endpoints are proved directly,
including the intersections of these axes.  Equations (SJ3)--(SJ10)
establish the stated full-state, sign, and endpoint reduction.  They do
not establish the still-required positive-ratio chart certificate.

\subsection{Source coordinates and the two-square decomposition}
\label{component:p667905169990}

\label{p667905169990:sec:r2-reconstructed-source-entry}

\paragraph{Actual source and conventions.}
The native source formulas give the following coordinate identities
and boundedness estimates. They provide the source parameters used
in the subsequent R2 contact and state-normal certificates.

Use the original weak-C2 native source $(a,B,t,r,h)$, with
$0<r,h<1$, $a,B,t>0$, the original strict source gaps, and $n\geq0$.
In R2, $t$ is positive unbounded while
\[
 B,\qquad z=aB,\qquad y=tB
\]
are positive infinitesimals.  There is no bound on $a$ and no lower
bound on $r$, $h$, $1-h$, or $e=1-r$.  Put $l=1-h$ and $S=1/t$.
Define the actual native arrays
\[
 c=1+zr,\quad s_c=1+B(l+ar),\quad
 R=1+hz,\quad g=R-ae>0,\quad j_D=c-l/t>0,
\]
\begin{align*}
 \Delta={}&a+(1+z)+(1+z)B+2z(1+z)r
          +a^3B^2r^2+2a^2Br\\
          &+tB\{1+a+z+a^2Br(2-r)\},\\
 M={}&(1+z)(1+ar)(h+t)-a(1+y)e\,l.
\end{align*}
The six weights used throughout this section are exactly
\[
\begin{aligned}
 n&=\frac{(1+ar)c s_c-r\Delta}{2r\Delta},&
 m&=\frac{a^2(1+y)e c s_c}{2g\Delta},\\
 C&=\frac{cM}{2l\Delta},&
 A&=\frac{tc(1+ar)s_c}{2h\Delta},\\
 D&=\frac{(1+z)(1+ar)(1-atBr)s_c}{2t j_D\Delta},&
 \alpha&=\frac{B^2(1+ar)M}{2(1-hB)\Delta}.
\end{aligned}                                                     \tag{RER1}
\]
All denominators are pure source.  The original positivity assumptions
are retained without adding a rate comparison.

\paragraph{Physicality controls the apparently unbounded source scale.}
The exact numerator identity is
\[
 (1+ar)(1+zr)\{1+B(l+ar)\}-r\Delta=eQ-BK,                \tag{RER2}
\]
where
\[
 K=(a+1)(h+t)(1+z),
\]
\[
 Q=1+B+y+z(h+t+2r)
       +z^2\{t(1+er)+h(1+r)\}+Bz(h+t).                 \tag{RER3}
\]
Every term of $Q$ is nonnegative.  Since $r\Delta>0$, original
physicality is exactly $eQ\geq BK$, including equality when $n=0$.
Put $\kappa=ay=tz$.  Before assuming that $\kappa$ is bounded, direct
termwise inspection of (RER3), using $B,y,z$ infinitesimal, gives
$Q\leq C_0(1+\kappa)$ for a fixed positive real $C_0$.
The strict gap $g>0$ gives $ae<R\sim1$, so multiplication of the
physicality inequality by $a$ yields
\[
 \kappa(a+1)(1+hS)(1+z)\leq aeQ<RQ
                                      \leq C_1(1+\kappa)           \tag{RER4}
\]
for a fixed positive real $C_1$.

If $a$ is bounded, $\kappa=ay$ is infinitesimal.  If $a$ is unbounded,
move $C_1\kappa$ to the left in (RER4).  Since
$a+1-C_1\geq(a+1)/2$, one obtains
$\kappa\leq2C_1/(a+1)$, again infinitesimal.  In either case,
(RER4) also bounds $(a+1)\kappa$.  Returning to (RER3) and then to
$eQ\geq BK$ proves
\[
 ay\in\mathrm I,\qquad a(a+1)y=O(1),\qquad Q\sim1,
 \qquad \frac{(a+1)y}{e}=O(1).                         \tag{RER5}
\]
No division by $n$ occurred; the argument includes the face $n=0$.

\paragraph{Uniform source normalizations.}
Put
\[
 k=\frac{1+ar}{1+a},\qquad p=\frac{ae}{R}\in(0,1),
 \qquad \mu=m(1-p)=\frac{mg}{R},\qquad\delta_n=nr.
\]
The quantity $k$ is a positive bounded unit.  Indeed $k\leq1$; if
$a\leq1$, then $k\geq1/(1+a)\geq1/2$.  If $a>1$, use
$ae<R=1+hz$ to get $k>(a-hz)/(1+a)$, which has a positive real lower
bound.  Direct termwise division of the displayed polynomial
$\Delta$ by $1+a$ gives $\Delta/(1+a)\sim1$.  Also
\[
 \frac{M}{t(1+ar)}\sim1:
\]
the first term has ratio $(1+z)(1+hS)\sim1$, while the absolute
ratio of the second is at most $R(1+y)S/(1+ar)$ and hence is
infinitesimal.  Since $c,s_c,j_D,1-hB\sim1$, (RER1) now gives
\[
 \frac{SC\,l}{k}\sim\frac12,\quad
 \frac{SA\,h}{k}\sim\frac12,\quad
 \frac{D}{Sk}\sim\frac12,\quad
 \frac{\alpha}{Sy^2(1+a)k^2}\sim\frac12.               \tag{RER6}
\]
For the $D$ relation use $atB=ay=\kappa\in\mathrm I$ from
(RER5).  Thus all original $h$ and $l$ openings are present.

Two exact bounded products remove the remaining nonuniform source
poles:
\[
 mg=\frac{a^2e(1+y)c s_c}{2\Delta},\qquad
 \delta_n=\frac{eQ-BK}{2\Delta}.                       \tag{RER7}
\]
The quantities $mg$, $mg/a$, $mg/a^2$, $\mu$, $\delta_n$, and
$nr^2$ are bounded.  For example
$a^2e/(1+a)\leq R a/(1+a)=O(1)$, while
$0\leq\delta_n\leq eQ/(2\Delta)=O(1)$.
These estimates do not imply that $m$ or $n$ themselves are bounded.

\paragraph{Exact extraction from the entire parent.}
Let testing $u,v,W,X,Y,Z$ be arbitrary, independently of the source;
put $D_u=1+u^2$, $D_v=1+v^2$.  The complete original quarter-parent is
\begin{align*}
 \mathcal P={}&m(vW-X+uvY-uZ)^2+nD_uD_vX^2
     +CD_v(uW-Y)^2+AD_uD_vY^2\\
 &+(D+\tfrac12)(uW+uvX-Y-vZ)^2
     +(\alpha+\tfrac12)D_u(vY-Z)^2\\
 &+\tfrac12D_v(X+uZ)^2-\tfrac12D_uD_vZ^2.
\end{align*}
Define the mixed rows
\[
 q_a=v(W+uY),\qquad N=X+uZ-pq_a,\qquad \xi=X+ruZ.
\]
Then the following identity retains both actual positive square terms,
without setting either normal to zero:
\[
 \mathcal P=mpN^2+nD_uD_v\xi^2+G,                       \tag{RER8}
\]
where
\begin{align*}
 G={}&\mu\{N^2-2q_aN+(1-p)q_a^2\}
        +\delta_nD_uD_v(-2uZX-ru^2Z^2)\\
 &+CD_v(uW-Y)^2+AD_uD_vY^2\\
 &+(D+\tfrac12)(uW+uvX-Y-vZ)^2
        +(\alpha+\tfrac12)D_u(vY-Z)^2\\
 &+\tfrac12D_v(X+uZ)^2-\tfrac12D_uD_vZ^2.
\end{align*}
It follows simply by expanding
$m\{(1-p)q_a-N\}^2-mpN^2$ and
$n(X^2-\xi^2)=nr(-2uZX-ru^2Z^2)$.
Both retained positive square terms are finite separated SOS: their source weights
are nonnegative, $D_uD_v$ is pure testing, and $N,\xi$ are allowed
mixed rows.  When $n=0$ the second positive square term vanishes literally.
The exact identity and all estimates above are the exact
structural conclusion; no positivity conclusion about $G$ is asserted.

\subsection{The C1--C2 boundary identity}
\label{component:pc024f5300388}

\label{pc024f5300388:sec:r13-source-overlap}

This section records an exact source-only identity. It does not by itself
admit any new independent testing domain. Its purpose is to identify the
same physical boundary in the two remaining source charts and to retain
the nonnegative slack when that boundary opens.

Use the original weak-C2 variables $B=H+L$, $z,y>0$ and $0<r<1$,
with the unchanged source gaps and nonnegative $n$. Write $e=1-r$,
$K=1+z$, $T_c=B+zr$, and let $\delta$ be the original positive
C2 denominator. Define the positive source polynomial
\[
 P_n=B+B^2+2Bzr+BHz+yB(1+z)+Hrz^2+yz^2re.
 \tag{OP1}
\]
Then the following identity holds on the entire original source:
\[
 \boxed{\quad eP_n=(1+z)(H+y)(B+zr)+2r\delta n.\quad}
 \tag{OP2}
\]
Every factor in the first summand is positive, and the second summand
is nonnegative, including literal $n=0$. In particular the pure-source
weights
\[
 q={ (1+z)(H+y)(B+zr)\over eP_n}>0,
 \qquad s_n={2r\delta n\over eP_n}\ge0
 \quad\hbox{satisfy}\quad q+s_n=1.                   \tag{OP3}
\]
Both square roots are legitimate source-field elements. Neither weight
is assumed to have nonzero residue.

\begin{proof}
The unchanged formula for $n$ is
\[
 2r\delta n=(B+zr)(1+zr)(1+zr+B-H)-r\delta.
\]
Substitution of the printed polynomial for $\delta$, followed by
collecting terms, gives OP2. This is a polynomial identity, so no
source boundary is obtained by continuity and no opening has been
inverted except the displayed positive pure-source factors.
\end{proof}

\paragraph{Exact correspondence of the source contacts.}
Use a subscript $1$ for native C1 notation in this paragraph. The map
\[
 x_1=z,\qquad \eta_1={y\over B},\qquad
 \zeta_1={B\over z},\qquad r_1={H\over B},\qquad
 \theta_1=e,
 \tag{OP4}
\]
with $k_1=\sqrt z$ and the native parameter
$\Delta_1=B/(z\sqrt y)$, gives
\[
 (u_c,v_c;W_c,X_c,Y_c,Z_c)
 =\left(\sqrt z,\sqrt y;{B\over\sqrt{zy}},
                       -r\sqrt z,{H\over\sqrt y},1\right)
\]
from the original native C1 contact formula. Thus it is the same
contact, including all state coordinates. Direct substitution in the
native positive polynomial $P_d$ gives $BP_d=P_n$.
The native C1 beta coefficient at this formal contact therefore obeys
\[
 \beta_1=-{r\delta n\over eP_n}=-{s_n\over2},
 \qquad \beta_1+\tfrac12={q\over2}>0.                 \tag{OP5}
\]
For $n=0$, the original C2 parent has
$\lambda=n=\beta=0$. It is literally an original C1 parent,
up to an optional positive normalization of $C$. The native source
entry and uniqueness at stationary contact identify its coefficients
with OP4. For $n>0$, OP5 makes the formal C1 beta negative. Hence
OP4 does \emph{not} map that source to the nonnegative C1 event cone,
and no physical C1 all-tests theorem may be imported there.

\section{Two small actual angles: the R1 sector}
\subsection{Scalar completion at unbounded source ratio}
\label{component:p9e725a065583}

\label{p9e725a065583:sec:reconstructed-r1-qo-source-scalar}

\paragraph{Scope of the dictionary and budgets.}
The identities below use the actual native weak-C2 stationary coefficients,
with $d=1$, $\lambda_{\rm event}=\beta=0$, all strict native gaps and
$n\ge0$. Source quantities lie in $K_c$ and testing quantities in a
separately labelled real closed extension $K_t$; all inverses and radicals
are pure source or pure testing. The rational identities have no source-size
restriction. Their R1 estimates assume $B,T,x,y$ infinitesimal; the later
RL and EN certificates establish their own replacement estimates.

For the native source write $a_s,B_s,t_s,r,h$, and put
$c_s=1+a_sB_sr$, $R_s=1+ha_sB_s$, $K_s=1+a_sB_s$.
The formerly used auxiliary ratio is simply
\[
 q_{\rm old}=R_s^2/(a_st_sK_s^2).\tag{RSR1}
\]
No size restriction on $q_{\rm old}$ or $q_z$ is imposed in this
refined dictionary. The actual full-R1 estimates, including arbitrary
positive $q_z=jT$, are justified by R1N1--R1N11.

\subsubsection*{The exact compact source and its uniform bounds}

Define pure-source quantities
\[
\begin{gathered}
 j=\frac{a_sc_s}{R_s},\quad B=\frac{B_sR_s}{c_s},\quad
 T=\frac{t_sc_s}{R_s},\\
 p=\frac{a_s(1-r)}{R_s},\quad H=\frac{1-h}{t_sc_s},\\
 x=jB,\quad y=TB,\quad q_z=jT .
\end{gathered}
 \tag{RSR2}
\]
The native gaps give $0<p<\min(1,j)$ and $0<H<1$.
In the R1 estimates $B,T,x,y$ are positive infinitesimals, while
$j$ and $q_z=jT$ are arbitrary positive. The exact dictionary itself
also applies at the different source scales used in RL and EN. Direct inversion gives
\[
 r=\frac{j-p}{j+px},\quad h=\frac{1-HT}{1+HTx},\quad
 a_s=\frac{j+px}{1+HTx},\quad
 B_s=\frac{B(1+HTx)}{1+pB},\quad
 t_s=\frac{T(1+pB)}{1+HTx}.
 \tag{RSR3}
\]
In particular the compact variables describe the actual source,
including its $r\to0$, $p\to1$, $H\to0$, and $H\to1$ openings.
Set
\[
 K=j+1-p,\qquad E=1+B(HT+p-1),\qquad
 W_0=(1-H)j+1-p.
 \tag{RSR4}
\]
Then $\tfrac12(j+1)<K<j+1$, $E\sim1$, and $W_0>0$.
The following polynomials contain the entire stationary dictionary:
\begin{align*}
\mathcal P={}&j+1+B\{(j+1)^2-p\}\\
 &+TB\{Hj+j+1+B[Hj(2j+2-p)+p(2j+1-p)]\}\\
 &+HjT^2B^2(1+Bp+BHj),\\
\mathcal L={}&K+T(W_0-H)+BT[pK+Hj(1-p)]+BHT^2(W_0-1),\\
\mathcal N={}&p+B(-j^2+3jp-j-p^2+p)\\
 &+BT(Hj^2+Hj-j^2+jp-j+p)\\
 &+B^2T(2Hj^2p-Hj^2-Hjp^2+2Hjp
             -2j^2p+3jp^2-jp-p^3+p^2)\\
 &+B^2HT^2j(Hj-j+p)(1+Bp).
 \tag{RSR5}
\end{align*}
With $\mu=m(1-p)$, the exact six original stationary weights are
\[
\begin{aligned}
 \mu&=\frac{jp(1+y)(1+x)(1+Hy)}{2\mathcal P},&
 n&=\frac{\mathcal N}{2(j-p)\mathcal P},\\
 C&=\frac{\mathcal L}{2HT\mathcal P},&
 A&=\frac{T(1+x)(1+Hy)(K+Hjy)}{2(1-HT)\mathcal P},\\
 D&=\frac{(1+Hy)(K+Hjy)\{1-y[(1-H)j-p]\}}
             {2T(1-H)\mathcal P},&
 \alpha&=\frac{B^2(K+Hjy)\mathcal L}{2E\mathcal P}.
\end{aligned}
 \tag{RSR6}
\]
For example, the original stationary denominator and its source mass
satisfy
\[
 \Delta=\frac{(1+x)\mathcal P}{(1+Bp)(1+HTx)},\qquad
 M=\frac{(1+x)\mathcal L}{(1+HTx)^2}.
\]
The two native strict gaps become
\[
 g=\frac{(1-p)(1+x)}{1+HTx},\qquad
 j_{\rm gap}=\frac{(1-H)(1+x)}{1+Bp}.
\]
These equations verify the source origin of RSR6 directly; no new
stationary solution has been substituted for the original one.

Every term after $j+1$ in $\mathcal P$, divided by $j+1$, is
infinitesimal: use $B(j+1)=B+x\in\mathrm I$, $T\in\mathrm I$,
$p<1$, and $H<1$. The same termwise estimate, using
$K>(j+1)/2$, applies to $\mathcal L-K$. Thus
\[
 \mathcal P\sim j+1,\qquad \mathcal L\sim K.
 \tag{RSR7}
\]
These are relative equivalences, valid also when $j$ is infinitesimal
or unbounded. They show, in particular, that $A/T$, $THC$, and
$T(1-H)D$ are positive bounded units, $\mu/p$ is bounded, and
$h_0=\mu+1/2$ is a positive bounded unit. Moreover $\alpha/B$ is
infinitesimal. Only the displayed source openings can make $m,n,C,D$
unbounded.

The original physicality condition gives a useful uniform restriction.
Regroup the exact polynomial $\mathcal N$ as
\[
\mathcal N=p\{1+B(3j+1-p)+BT(j+1)\}
 -Bj(j+1)\{1+(1-H)T\}+\mathcal R.
 \tag{RSR8}
\]
Here $\mathcal R$ is precisely the last two lines of $\mathcal N$
in RSR5. Since $0<p<\min(1,j)$, each coefficient in its $B^2T$
line is bounded in absolute value by a fixed multiple of $j(j+1)$.
The $B^2HT^2$ line has the same bound, multiplied by the bounded
factor $T(1+Bp)$. Consequently
\[
 |\mathcal R|\leq C_0BT\,Bj(j+1)
\]
for a fixed real $C_0$. All corrections multiplying $p$ and
$Bj(j+1)$ in RSR8 are infinitesimal. As $n\geq0$ is exactly
$\mathcal N\geq0$, it follows that
\[
 \eta_*:=\frac{Bj(j+1)}p\ \hbox{is bounded},\qquad
 \eta_*\in\mathrm I\ \Longrightarrow\
 n\sim\frac{p}{2(j-p)(j+1)}>0.
 \tag{RSR9}
\]
Indeed $\mathcal N=p(1+\varepsilon_1)-Bj(j+1)(1+\varepsilon_2)$
with $\varepsilon_i\in\mathrm I$. This proves both assertions without
dividing by $n$ and includes literal $n=0$ in the bounded-$\eta_*$
alternative. It also gives $nr$ bounded.

\subsubsection*{The complete graph and its actual contact circuit}

The quarter-parent, before any positive square term is removed, is
\begin{align*}
\mathscr P={}&m\{vW+uvY-X-uZ\}^2
 +CD_v(uW-Y)^2+AD_uD_vY^2\\
&+(D+\tfrac12)(uW+uvX-Y-vZ)^2
 +(\alpha+\tfrac12)D_u(vY-Z)^2+nD_uD_vX^2\\
&+\tfrac12D_v(X+uZ)^2-\tfrac12D_uD_vZ^2,
\qquad D_u=1+u^2,\ D_v=1+v^2.
 \tag{RSR10}
\end{align*}
The source-free rotated state coordinate is $w=(W+uY)/D_u$.
On its nonzero chart set $w=1$ and $z=(Z-uX)/D_u$. Define
the two complete state normals
\[
 N=X+uZ-pD_uv,\qquad M=uW-Y-HD_uvz.
\]
Their exact inverse is
\[
\begin{aligned}
 X&=N/D_u+pv-uz,& Z&=z+uN/D_u+puv,\\
 Y&=u-M/D_u-Hvz,& W&=1+uM/D_u+uHvz.
\end{aligned}
 \tag{RSR11}
\]
The scalar studied here is exactly
$F_{\rm full}=(\mathscr P/D_u)|_{N=M=0}$.
No $n$ positive square term, $m$ positive square term, or state-normal term is silently removed.
Substitution of RSR11 into RSR10 gives the complete full-parent
normal expression RFN7 in Section~\ref{component:p261f71af685d}, with
$N_0=N/D_u$, $M_0=M/D_u$ and the retained positive square term
$mpD_u^2N_0^2$ restored.

Use the separated rows
\[
 U=u/\sqrt x,\quad V=v/\sqrt y,\quad Z_*=z/\sqrt{q_z},\qquad
 Q_*=U-VZ_*,\quad B_*=V-UZ_*,\quad R_*=Z_*-UV,
\]
\[
 \Psi=U^2+V^2+Z_*^2+U^2V^2Z_*^2-4UVZ_*.
 \tag{RSR12}
\]
These are rows; $U,V,Z_*$ are not assumed bounded coefficients.
The actual contact is $(U,V,Z_*;N,M)=(1,1,1;0,0)$.
The exact native formulas make the complete parent value and all five
first derivatives vanish there. On the positive core,
\[
 \Psi=(U-V)^2+Z_*^2(1-UV)^2+2UV(1-Z_*)^2
 \tag{RSR13}
\]
is a legal separated sum of squares: the last root is
$\sqrt2\,(uv)^{1/2}(xy)^{-1/4}(1-Z_*)$, with its radicals in
the two separate real closed fields.

Let the ten positive monomial coefficients and the seven circuit
masses be exactly RB6--RB8. In particular
\[
 M_a=Ax,\quad M_b=\{\mu(1-p)+(n+\tfrac12)p^2\}y,\quad
 M_c=\alpha q_z,\quad M_{ac}=nj^2y,
\]
\[
 f_c=M_b-M_{ac},\quad q_c=M_a-f_c,\quad
 e_c=(n+\tfrac12)p^2y^2-g_c,\quad d_c=M_{b^2c}-k_c.
\]
The original contact balances give the exact scalar identity
\begin{align*}
 F_{\rm full}={}&q_cQ_*^2+M_{ac}B_*^2+r_cR_*^2
  +d_cV^2Q_*^2+e_cV^2B_*^2\\
 &+f_c\Psi+g_cV^2\Psi+k_cZ_*^2(V^2-1)^2.
 \tag{RSR14}
\end{align*}
The identities RB7 and the two odd coefficient balances are independently
checked on the actual source. Thus RSR14 is a finite polynomial identity,
not a Taylor approximation to the contact.

Here is an explicit source-only formula for its radial mass. Put
$\mathcal F=K+T\mathcal F_1+T^2\mathcal F_2$, where
\begin{align*}
\mathcal F_1={}&-2BHjp+BHj+2Bjp-2Bp^2+2Bp\\
 &-Hj-Hp-H+j+1-p^2/j+p/j,\\
\mathcal F_2={}&-B^2H^2jp+B^2Hjp-B^2Hp^2+B^2Hp\\
 &-BH^2j-BH^2p+BHj-BHp^2/j+BHp/j.
 \tag{RSR15}
\end{align*}
Every displayed term divided by $K$ is bounded, using $p/j<1$;
terms containing $Bj$ are even infinitesimal. Hence $\mathcal F\sim K$.
Direct source substitution yields
\[
\begin{aligned}
 f_c&=\frac{Tx^2\mathcal F}{2\mathcal P}>0,&
 q_c&=Ax-f_c\sim\frac{TxK}{2(j+1)}>0,\\
 r_c&=\frac{B^2jT(1-p)(1+x)(1+Hy)\mathcal L}
                {2E\mathcal P}>0,&
 g_c&=\frac{HT^2j^2B^3\mathcal L}{2\mathcal P}>0.
\end{aligned}
 \tag{RSR16}
\]
In particular
\[
 f_c\sim\frac{Tx^2K}{2(j+1)},\qquad
 \frac{r_c}{f_c}\sim\frac{1-p}{j},\qquad
 \frac{g_c}{f_c}=Hy\frac{\mathcal L}{\mathcal F}.
 \tag{RSR17}
\]

For the remaining signed mass define
\begin{align*}
\mathcal K={}&pW_0+
 B(1-p)\{j[H(p+1)-p]-(1-p)(p-H)\}\\
 &+TBH\{(p+1)(1-H)j+(1-p)(BHj+p+1-H)\}\\
 &+T^2B^2H^2j(1-H).
 \tag{RSR18}
\end{align*}
Then
\[
 k_c=-\frac{T^2B^2j(1+x)\mathcal K}{2E\mathcal P}.
\]
The identity
\begin{align*}
\mathcal K-pW_0+Bp(1-p)(j+1)
={}&B(1-p)\{Hj(p+1)+H(1-p)+p^2\}\\
 &+TBH\{(p+1)(1-H)j+(1-p)(BHj+p+1-H)\}\\
 &+T^2B^2H^2j(1-H)
 \tag{RSR19}
\end{align*}
has a nonnegative right side. Since $W_0\geq1-p$ and
$B(j+1)\in\mathrm I$, it proves $\mathcal K>0$ and $k_c<0$.
The displayed finite arrays also give
\[
 \mathcal K/K=O(p+B),\qquad
 |k_c|/(q_cy)=O(p+B).
 \tag{RSR20}
\]
All constants in these estimates are real and uniform over the source
openings; neither $j-p$ nor $1-H$ has been bounded below.

\subsubsection*{Paying the signed masses, with the full current reserved}

Put
\[
 h_0=\mu+\tfrac12,\qquad
 M_e=\frac{nh_0j^2}{n+h_0},\qquad J_c=M_ey,
 \qquad e_0=\tfrac12p^2y^2-g_c.
 \tag{RSR21}
\]
These are pure-source quantities and $0\leq J_c\leq M_{ac}$.
Writing $e_c=np^2y^2+e_0$ preserves the entire second correlated
current positive square term. The conclusion is the following precise scalar
statement: after subtracting the reserved positive square terms
\[
 (M_{ac}-J_c)B_*^2+np^2y^2V^2B_*^2,
 \tag{RSR22}
\]
the remainder is a separated SOS and retains, for example, one half
of $q_cQ_*^2+r_cR_*^2+f_c\Psi+J_cB_*^2$.
When $n=0$ the two reserved positive square terms and $J_c$ vanish literally.
This assertion supplies precisely the effective current mass needed
when the full $N$ normal is later completed against the original $n$
positive square term. It is not a completion of that normal itself.

A budget has infinitesimal cost when its Gram coefficients are
bounded and have zero residue. A finite number of such corrections
leaves any fixed fraction smaller than one of each protected square.
The remaining Gram matrix has positive-definite real residue, so
Lemma~\ref{p25af2c58da7e:lem:gram} gives its exact sum of squares.
Restoring the budget identities and the retained positive source
terms proves the required identity.

First $e_0^-\leq g_c$. The exact budget
\[
 (1+V^2)\Psi-V^2B_*^2
 =Q_*^2+R_*^2+V^2(1-UVZ_*)^2
 \tag{RSR23}
\]
pays its negative part from the circuit at cost
\[
 \frac{e_0^-}{f_c}+\frac{e_0^-}{f_cy}v^2
 \ \leq\ Hy\frac{\mathcal L}{\mathcal F}
       +H\frac{\mathcal L}{\mathcal F}v^2\in\mathrm I.
\]
Its positive part is retained as a square. Also $d_c>0$ because
$M_{b^2c}>0$ and $k_c<0$; the $d_c$ and $g_c$ terms cause no loss.

For $k_c$, use the exact identity
\[
 Z_*(V^2-1)=-VQ_*-R_*.
 \tag{RSR24}
\]
The $VQ_*$ part is paid from $q_cQ_*^2$ with coefficient
$2|k_c|v^2/(q_cy)$, which is infinitesimal by RSR20. The remaining
$R_*$ part is paid in the following exhaustive three cases. Fixed
real multipliers such as the factor two do not change any cost claim.

\paragraph{Case 1: $1-p$ has a positive real lower bound.}
The exact ratio is
\[
 \frac{|k_c|}{r_c}
 =\frac{T\mathcal K}{(1-p)(1+Hy)\mathcal L}\in\mathrm I.
 \tag{RSR25}
\]
This uses only $\mathcal K/\mathcal L=O(p+B)$ and $T\in\mathrm I$.
Neither current mass nor a bound on $j$ is required.

\paragraph{Case 2: $1-p\in\mathrm I$ and $\eta_*$ has a positive real
lower bound.}
Here $p\sim1$, so $Bj(j+1)$ is a bounded positive unit. Since $B$ is
infinitesimal, $j$ is unbounded and $j^2B$ is a positive bounded unit.
Equations RSR16--RSR20 give
\[
 \frac{|k_c|}{f_c}\in\mathrm I,\qquad
 \frac{|k_c|}{f_cx}
 =O\!\left(\frac{T(p+B)}{j^2B}\right)\in\mathrm I.
 \tag{RSR26}
\]
The exact circuit budget
\[
 (1+U^2)\Psi=R_*^2+B_*^2+U^2Q_*^2+U^2(1-UVZ_*)^2
 \tag{RSR27}
\]
therefore pays $|k_c|R_*^2$ at cost
$|k_c|/f_c+(|k_c|/(f_cx))u^2\in\mathrm I$.
This case includes any intersection with literal $n=0$.

\paragraph{Case 3: $1-p,\eta_*\in\mathrm I$.}
Now $p\sim1$, $j>p$, and RSR9 gives $n>0$ together with its exact
relative estimate. The source mass $M_e$ is a positive bounded unit.
For the lower bound use
\[
 \frac{p^2}{M_e}
 =\frac{p^2}{h_0j^2}+\frac{p^2}{nj^2}=O(1).
\]
For the upper bound use $M_e\leq h_0j^2$ when $j$ is bounded,
and $M_e\leq nj^2\sim p/2$ when $j$ is unbounded. In particular,
\[
 \frac{|k_c|}{f_c}\in\mathrm I,\qquad
 \frac{|k_c|}{q_cy}=O(1),\qquad
 \frac{|k_c|}{J_cx}
 =O\!\left(\frac{T(p+B)(1+1/n)}{j^2}\right)=O(1).
 \tag{RSR28}
\]
The last bound follows directly from
$1/n\sim2(j-p)(j+1)/p$, with $p\sim1$ and $j>p$.
There is an exact separated SOS budget
\[
 2\{\Psi+V^2Q_*^2+U^2B_*^2\}-R_*^2\quad\hbox{is SOS}.
 \tag{RSR29}
\]
For completeness, set $t_*=UV$, $\delta_*=U-V$,
$E_*=Z_*(1-t_*)$, $T_*=t_*(1-Z_*)$, and
$J_*=T_*-\delta_*^2Z_*$. Its SOS decomposition is
\[
 (E_*+T_*)^2+2\delta_*^2+4t_*(1-Z_*)^2
                       +2J_*^2+4t_*\delta_*^2Z_*^2.
\]
All roots are separated as in RSR13. Applying RSR29 costs a multiple
of $|k_c|/f_c$ on the circuit, $|k_c|v^2/(q_cy)$ on $Q_*$,
and $|k_c|u^2/(J_cx)$ on $B_*$. Each is infinitesimal by RSR28.

The three cases exhaust the source because $0<1-p<1$ and $\eta_*$
is bounded. A finite number of the displayed allocations therefore
proves the scalar assertion following RSR22 with a strict retained
fraction. All coefficient bounds depend only on the source and
$u,v$; no bound on the independent state coordinate $z$ was used.
The proof allows arbitrary relative rates of all source and testing
infinitesimals. Literal $n=0$ is covered by Cases 1 and 2; Case 3
excludes it by the original physicality equation itself.

\subsection{Completion of both state normals}
\label{component:p261f71af685d}

\label{p261f71af685d:sec:fn-reconstructed}

\paragraph{Use of the complete normal calculation.}
The parent identities below are exact for the compact stationary dictionary.
The source estimates and square budgets give the full local R1 certificate
under the full-R1 source hypotheses verified in R1N1--R1N11, with arbitrary
positive $jT$. The later RL and EN bodies reuse only the printed exact
identities and supply their own source estimates.

\paragraph{Fields, source scope, and notation.}
Let $K_c,K_t\supset\mathbb R$ be real closed fields, with bounded rings
$B_c,B_t$ and infinitesimal ideals $\mathfrak m_c,\mathfrak m_t$.
The ring of certificate expressions is
$\mathcal T=K_c\otimes_{\mathbb R}K_t$, and its bounded coefficient ring
is $\mathcal B=B_c\otimes_{\mathbb R}B_t$.  A bounded coefficient with
zero real residue belongs to
$\mathfrak n=\mathfrak m_c\otimes B_t+B_c\otimes\mathfrak m_t$.
Every radical or inverse used below is pure source, pure testing, or
an explicitly separated product of such factors.  No order on
$\mathcal T$ and no inverse of a mixed expression is used.

Use the positive compact source coordinates
\[
 x=jB,\qquad y=TB,\qquad q_z=jT,\qquad
 0<p<\min\{1,j\},\qquad 0<H<1,                         \tag{RFN1}
\]
where $B,T,x,y\in\mathfrak m_c$ and $q_z>0$ is arbitrary. These are
the compact coordinates of the complete original R1 sector; the absence
of any $q_z$ restriction is verified in R1N1--R1N11.  The original
strict source gaps remain in force and $n\geq0$, including $n=0$.
For reference, their native substitution is
\[
 r=\frac{j-p}{j+px},\quad
 h=\frac{1-HT}{1+HTx},\quad
 a_s=\frac{j+px}{1+HTx},\quad
 B_s=\frac{B(1+HTx)}{1+pB},\quad
 t_s=\frac{T(1+pB)}{1+HTx}.                              \tag{RFN2}
\]
The six weights $n,m,C,A,D,\alpha$ are always the original source
weights at (RFN2), rather than weights of an auxiliary stationary
source.  Put $\mu=m(1-p)$, $h_0=\mu+\tfrac12$,
$k=(j+1-p)/(j+1)$, and $\eta_*=Bj(j+1)/p$.
In this source sector RS proves
\begin{gather}
 \tfrac12<k<1,\quad \eta_*=O(1),\quad
 \mu=O(p),\quad h_0\asymp1,\quad A=O(T),\quad
 \alpha\asymp B^2(j+1),                                  \tag{RFN3a}\\
 TCH\sim k/2,\qquad T(D+\tfrac12)(1-H)\sim k/2,\qquad
 \frac{n(j-p)}j=O(1),                                   \tag{RFN3b}\\
 f_c\asymp Tx^2,\qquad q_c\asymp Tx,\qquad
 r_c\asymp TjB^2(1-p),                                  \tag{RFN3c}\\
 \eta_*\in\mathfrak m_c\quad\Longrightarrow\quad
 n\sim\frac{p}{2(j-p)(j+1)}>0.                           \tag{RFN3d}
\end{gather}
Here $\asymp$ means that the ratio is a positive bounded unit, and
$\sim$ means that the ratio has residue one.  All estimates are uniform
at the independent openings $H$, $1-H$, $j-p$, and $1-p$.
In particular $B(j+1)=B+x$ and $\alpha$ are infinitesimal.
No comparison between source and testing infinitesimals is allowed.

\subsubsection*{The complete parent and the correlated source mass}

Write $D_u=1+u^2$, $D_v=1+v^2$.  The original quarter-parent is
\begin{align}
 \mathcal P={}&m(vW-X+uvY-uZ)^2+nD_uD_vX^2
  +CD_v(uW-Y)^2+AD_uD_vY^2\\
 &+(D+\tfrac12)(uW+uvX-Y-vZ)^2
  +(\alpha+\tfrac12)D_u(vY-Z)^2\\
 &+\tfrac12D_v(X+uZ)^2-\tfrac12D_uD_vZ^2.               \tag{RFN4}
\end{align}
Use the two mixed normal rows
\[
 N_0=\frac{X+uZ-pD_uv}{D_u},\qquad
 M_0=\frac{uW-Y-HD_uvz}{D_u}.                            \tag{RFN5}
\]
Their exact inverse, with $w=1$, is
\[
 \begin{aligned}
 X&=N_0+pv-uz,&Z&=z+uN_0+puv,\\
 W&=1+u(M_0+Hvz),&Y&=u-M_0-Hvz.
 \end{aligned}                                           \tag{RFN6}
\]
Only the pure testing unit $D_u$ is inverted.  Retain the nonnegative
positive square term $mpD_u^2N_0^2$ and define
$G=(\mathcal P-mpD_u^2N_0^2)/D_u$.  Direct substitution gives the
\emph{entire} expression
\begin{align}
 G={}&\mu D_u\{N_0^2-2vN_0+(1-p)v^2\}
       +nD_v(N_0+pv-uz)^2\\
 &+CD_uD_v(M_0+Hvz)^2+A D_v(u-M_0-Hvz)^2\\
 &+(D+\tfrac12)D_u(M_0-(1-H)vz)^2\\
 &+(\alpha+\tfrac12)
   \{u(1-p)v-(1+Hv^2)z-uN_0-vM_0\}^2\\
 &+\tfrac12D_uD_v(N_0+pv)^2
       -\tfrac12D_v(z+uN_0+puv)^2.                      \tag{RFN7}
\end{align}
Thus $G$ is exactly quadratic in $(N_0,M_0)$:
\[
 G=F+2L_NN_0+2L_MM_0+H_NN_0^2+H_MM_0^2
                  +2(\alpha+\tfrac12)uvN_0M_0,          \tag{RFN8}
\]
where $F=G|_{N_0=M_0=0}$ and
\[
 \begin{split}
 H_N&=h_0D_u+\tfrac12v^2+\alpha u^2+nD_v,\\
 H_M&=D_u(CD_v+D+\tfrac12)+AD_v+(\alpha+\tfrac12)v^2.
 \end{split}                                           \tag{RFN9}
\]
There are no omitted higher state-normal terms.

\paragraph{Exact scalar interface, before spending the $n$-mass.}
Set
\[
 U=\frac u{\sqrt x},\quad V=\frac v{\sqrt y},\quad
 Z_* =\frac z{\sqrt{q_z}},\qquad
 Q_*=U-VZ_*,\quad B_*=V-UZ_*,\quad R_*=Z_*-UV,
\]
\[
 \Psi=U^2+V^2+Z_*^2+U^2V^2Z_*^2-4UVZ_*.                 \tag{RFN10}
\]
All three normalized testing variables are positive, but none is
assumed bounded.  A legal separated certificate is
\[
 \Psi=(U-V)^2+Z_*^2(1-UV)^2+2UV(1-Z_*)^2.               \tag{RFN11}
\]
For example, $\sqrt{UV}$ is the product of the source factor
$(xy)^{-1/4}$ and the testing factor $\sqrt{uv}$.

The exact RS/RB scalar identity for the actual source is
\begin{align}
 F={}&q_cQ_*^2+nj^2y B_*^2+r_cR_*^2+d_cV^2Q_*^2
   +(np^2y^2+e_0)V^2B_*^2\\
   &+f_c\Psi+g_cV^2\Psi+k_cZ_*^2(V^2-1)^2,
 \qquad e_0=\tfrac12p^2y^2-g_c.                        \tag{RFN12}
\end{align}
For an entirely algebraic definition of the masses, put
$K_0=CH^2+(D+\tfrac12)(1-H)^2$ and
\[
\begin{array}{ll}
 M_a=Ax,&M_b=\{\mu(1-p)+(n+\tfrac12)p^2\}y,\\
 M_c=\alpha q_z,&
 M_{ab}=\{A+\alpha(1-p)^2+\mu(1-p)+\tfrac12(1-p)^2\}xy,\\
 M_{ac}=nxq_z=nj^2y,&
 M_{bc}=\{K_0+AH^2+2H\alpha+H-\tfrac12\}yq_z,\\
 M_{b^2}=(n+\tfrac12)p^2y^2,&M_{abc}=(K_0+n)xyq_z,\\
 M_{b^2c}=H^2(A+C+\alpha+\tfrac12)y^2q_z,&
 M_{ab^2c}=CH^2xy^2q_z.
\end{array}
\]
Then
\[
\begin{split}
 k_c&=(M_a-M_b+M_{ac}+M_{ab^2c}-M_{bc})/2,\\
 r_c&=M_c-M_b+M_{ac}-k_c,\quad f_c=M_b-M_{ac},\quad q_c=M_a-f_c,\\
 d_c&=M_{b^2c}-k_c,\qquad g_c=M_{ab^2c}.
\end{split}
\]
The source contact identities make (RFN12) exact, including both odd
monomials of $F$.  It is checked coefficientwise against (RFN7), not
only at the contact point.

Define the unnormalized current $b=v-uz/j=\sqrt y B_*$ and the
pure source quantities
\[
 c_b=\frac{nj}{n+h_0},\qquad M_e=\frac{nh_0j^2}{n+h_0}.
                                                               \tag{RFN13}
\]
The whole scalar pieces $nj^2b^2+np^2v^2b^2$ in (RFN12) must be kept
until their correlations with $N_0$ have been completed.  The full
$N_0$-linear term below contains $njb+npv^2b$, so these pieces and the
corresponding hard terms give
\[
 n(N_0+jb)^2+n v^2(N_0+pb)^2.
\]
Complete only the first square, using a pure source inverse:
\[
 h_0N_0^2+n(N_0+jb)^2
   =(n+h_0)(N_0+c_bb)^2+M_e b^2.                       \tag{RFN14}
\]
Keep $nv^2(N_0+pb)^2$ whole.  Put
$\widehat N=\sqrt{n+h_0}(N_0+c_bb)$.
The other hard terms in (RFN9) are the nonnegative whole squares
\[
 h_0(uN_0)^2+\tfrac12(vN_0)^2+\alpha(uN_0)^2.             \tag{RFN15}
\]
They are still kept whole after the shift; no term $nu^2$ or $nv^2$
is treated as a bounded perturbation.

For $n>0$ one has the exact bounded-shift identity
\[
 \frac{c_b^2}{M_e}=\frac{n}{h_0(n+h_0)}=O(1).
                                                               \tag{RFN16}
\]
At $n=0$, set $c_b=M_e=0$ and omit the effective current row.  Formula
(RFN14) remains literal, and no division by $n$ or $\sqrt{M_e}$ is
made in that case.

RS pays the signed $e_0$ and $k_c$ terms in (RFN12) at infinitesimal
cost after (RFN14), with the effective current mass $M_e y$, and leaves
fixed positive real fractions of
\[
 q_cQ_*^2,\qquad M_e b^2,\qquad r_cR_*^2,\qquad f_c\Psi. \tag{RFN17}
\]
The second term is omitted when $n=0$.  This is the precise scalar
allocation needed here; an allocation that spends the raw $nj^2y$
mass before (RFN14) would not suffice.

\subsubsection*{Finite exact budgets for the full normal rows}

In the rest of the proof, a row is called free if any prescribed fixed
real multiple of its square can be purchased at infinitesimal cost
from (RFN17).  This has an exact meaning: a finite SOS identity writes
its square below a finite sum of those retained squares, with bounded
separated coefficients in $\mathfrak n$.  Subtracting any fixed real
multiple changes their strict Gram diagonals only infinitesimally.

\paragraph{Three elementary circuit budgets.}
Direct expansion gives
\begin{align}
 (1+U^2)\Psi-R_*^2
     &=B_*^2+U^2Q_*^2+U^2(1-UVZ_*)^2,                  \tag{RFN18a}\\
 (1+V^2)\Psi-V^2B_*^2
     &=Q_*^2+R_*^2+V^2(1-UVZ_*)^2.                   \tag{RFN18b}
\end{align}
There is also a useful opposite-current budget.  Put
$t_*=UV$, $\delta_*=U-V$, $E_*=Z_*(1-t_*)$,
$T_*=t_*(1-Z_*)$, and $J_*=T_*-\delta_*^2Z_*$.  Then
\begin{align}
 2(\Psi+V^2Q_*^2+U^2B_*^2)-R_*^2
 ={}&(E_*+T_*)^2+2\delta_*^2+4t_*(1-Z_*)^2\\
   &+2J_*^2+4t_*\delta_*^2Z_*^2.                      \tag{RFN18c}
\end{align}
Every term on the right is a legal separated square.  Thus a
nonnegative multiplier $a$ provides $aR_*^2$ at infinitesimal cost whenever
\[
 a/f_c\in\mathfrak n,\qquad a/(q_cy)=O(1),\qquad
 a/(M_e yx)=O(1),                                      \tag{RFN18d}
\]
since the last two diagonal costs have the additional testing factors
$v^2$ and $u^2$.  This criterion is used only when $M_e>0$.

\paragraph{An exact radial budget.}
The row $R_u=v(u^2-x)$ is free in all source subcases, including the
weak opening $1-p\in\mathfrak m_c$.  Here is an explicit finite
certificate, so that this assertion does not rest on pointwise
nonnegativity.  Write $D_*=V(U^2-1)$ and
\[
\begin{split}
 t_7={}&(Z_*-U^2V,\ V-U^2V,\ U-U^2V,\ UZ_*-U^2V,\\
       &\hspace{25mm}UV-U^2V,\ UVZ_*-U^2V,\ U^2-U^2V)^T,\\
 t_3={}&(1-UV,\ V-UV,\ U-UV)^T,
\end{split}
\]
\[
 G_7=\begin{pmatrix}
 4&-2&-3/2&0&-1/2&0&1/2\\
 -2&3&0&-3/2&0&-1&2\\
 -3/2&0&4&3/2&-2&-3&0\\
 0&-3/2&3/2&4&-3/2&0&-3/2\\
 -1/2&0&-2&-3/2&3&2&-1\\
 0&-1&-3&0&2&8&-4\\
 1/2&2&0&-3/2&-1&-4&4
 \end{pmatrix},\qquad
 G_3=\begin{pmatrix}3&-2&-2\\-2&2&2\\-2&2&3\end{pmatrix}.
\]
The exact identity is
\[
 4\Psi+4U^2Q_*^2+4U^2R_*^2-D_*^2
   =t_7^TG_7t_7+4VZ_*(1-U)^2+UZ_*t_3^TG_3t_3.         \tag{RFN19}
\]
The leading principal minors of $G_7$ are
$4,8,101/4,1049/16,907/16,3279/16,743/16$; those of $G_3$ are
$3,2,2$.  Both are positive definite real matrices.  Their real
Cholesky factors and the separated square roots of $VZ_*$ and $UZ_*$
therefore make the right side a finite legal SOS.  Denote it by
$\mathcal W(U,V,Z_*)$.

Since $R_u=x\sqrt y D_*$, the exhaustive source cases are as follows.

\emph{Case 1: $1-p$ is a positive unit.}
Multiply (RFN19) by $x^2y$.  The costs relative to
$f_c\Psi,q_cQ_*^2,r_cR_*^2$ are, respectively,
\[
 O(B),\qquad O(Bu^2),\qquad O(u^2/(1-p)).
\]
All belong to $\mathfrak n$.

\emph{Case 2: $p\sim1$ and $\eta_*$ is a positive unit.}
Then $j$ has a positive real lower bound, and $j^2B$ is a positive
unit.  Adding $4U^2$ times (RFN18a) to (RFN19) gives
\[
 D_*^2\ \preceq_{\rm sos}\
 4(1+U^2+U^4)\Psi+4U^2Q_*^2.
\]
After multiplication by $x^2y$, the costs are
\[
 O\!\left(B+u^2/j+u^4/(j^2B)\right),\qquad O(Bu^2),
\]
again infinitesimal.  The notation $a\preceq_{\rm sos}b$ means
$b-a\in\Sigma\mathcal T^2$ and is only shorthand for the displayed
exact identities.

\emph{Case 3: $p\sim1$ and $\eta_*\in\mathfrak m_c$.}
Formula (RFN3d) gives $n>0$ and $1/M_e=O(1)$; it also gives
$p^2/M_e=O(1)$.  Put $E^{\rm rad}=Z_*(U^2-1)$.
Permuting $V$ and $Z_*$ in (RFN19) yields
\[
 \mathcal W(U,Z_*,V)
   =4\Psi+4U^2Q_*^2+4U^2B_*^2-(E^{\rm rad})^2.
\]
The identity $D_*=UE^{\rm rad}+(U^2-1)B_*$ then gives, exactly,
\begin{align}
 &8U^2\Psi+8U^4Q_*^2+
   \{8U^4+2(U^2-1)^2\}B_*^2-D_*^2\\
 &\hspace{10mm}=2U^2\mathcal W(U,Z_*,V)
                 +\{UE^{\rm rad}-(U^2-1)B_*\}^2.        \tag{RFN20}
\end{align}
After multiplication by $x^2y$, the three costs are
\[
 O(u^2/j),\qquad O(u^4/j),\qquad
 \frac{8u^4+2(u^2-x)^2}{M_e}.
\]
They are infinitesimal.  These three cases prove the asserted free
radial row without assuming a rate for any testing variable.

\paragraph{The smaller current helpers.}
If $\eta_*\in\mathfrak m_c$, (RFN3d) implies
\[
 \frac{p^2}{M_e}
 =\frac{p^2}{h_0j^2}+\frac{p^2}{nj^2}=O(1),
\]
because $p/j<1$ and
$p(j-p)(j+1)/j^2\leq p+p/j$.
Thus $(pv^2b)^2$ costs only $(p^2/M_e)v^4$ of the effective current.
If $\eta_*$ is a positive unit instead, then
\[
 \frac{p^2y}{f_c}=O\!\left(\frac{p+p/j}{\eta_*}\right)=O(1).
\]
Multiply (RFN18b) by $p^2y^2v^2$.  This provides $(pv^2b)^2$ with circuit
cost
\[
 \frac{p^2y}{f_c}(yv^2+v^4)\in\mathfrak n.
\]
The same identity multiplied by $Tp^2u^2y^2$ shows that
$\sqrt T\,puvb$ is free: its cost is either
$(p^2/M_e)Tu^2v^2$ or
$(p^2y/f_c)Tu^2(y+v^2)$.

\subsubsection*{The entire two-normal conductor}

\paragraph{First normal: exact source cancellation and allocation.}
Put
\[
 q=vz-Tu,\qquad R=z-uv/B,\qquad R_u=v(u^2-x),
\]
\[
 K_c=\frac p2+j\{H(\alpha+\tfrac12)-\tfrac12\}+n(p-j),
 \qquad
 K_R=-\mu-(\alpha+\tfrac12)(1-p)+\frac TjK_c.
\]
Both $K_c/j$ and $K_R$ are bounded by (RFN3).  The complete first normal
is exactly
\[
 L_N=\{j(n-\alpha)+p(n+\tfrac12)v^2\}b
                  +K_RR_u+\frac{K_c}{j}uvq.            \tag{RFN21}
\]
In particular the sign of the $\mu$ term in $K_R$ is negative.
For a direct source check, the difference between this formula and
the derivative of (RFN7) is $vS_N$, where
\[
 S_N=np-\mu+\tfrac p2+
       \{-\mu-(\alpha+\tfrac12)(1-p)\}x
       +j(\alpha-n)+yK_c=0.
\]
The last equality is an exact rational identity of the actual native
source under (RFN2), not an asymptotic contact equation.

After the two correlated $n$-squares are formed, the remaining linear
coefficient is
\[
 L_N^{\rm rem}=-j\alpha b+\tfrac12pv^2b
                       +K_RR_u+\frac{K_c}{j}uvq.
\]
The $q$ row is strict because
$q_cQ_*^2=(q_c/(Tx))(q^2/T)$, with $q_c/(Tx)$ a positive unit.
Its coefficient in $L_N^{\rm rem}$ relative to $q/\sqrt T$ is
$(K_c/j)uv\sqrt T\in\mathfrak n$.
The $pv^2b$ helper is free by the preceding paragraph.
For the remaining small row, the exact identity
\[
 j\alpha b=-\frac\alpha B R_u-\alpha uR
\]
has $\alpha/B=O(B+x)\in\mathfrak m_c$.
When $1-p$ is a unit, the square of $\alpha uR$ costs
$O(B^2(j+1)^2u^2/(1-p))\in\mathfrak n$ against $r_cR_*^2$.
When $p\sim1$, use $B_*^2\preceq_{\rm sos}(1+U^2)\Psi$ from
(RFN18a).  The circuit cost of $(j\alpha b)^2$ is
\[
 O\!\left(B^3(j+1)^2+B^2(j+1)^2u^2/j\right)\in\mathfrak n,
\]
since $j$ is bounded below and $B(j+1)=B+x$ is infinitesimal.
Finally, $K_RR_u$ has bounded coefficient; provide an arbitrarily large
fixed real $R_u^2$ pivot using (RFN19)--(RFN20).

The shift $N_0=\widehat N/\sqrt{n+h_0}-c_bb$ introduces no new
unbounded Gram coefficient: $1/\sqrt{n+h_0}=O(1)$, and for $n>0$
$c_b/\sqrt{M_e}=O(1)$ by (RFN16).  The only bounded, potentially
non-infinitesimal couplings are therefore those of $R_u$ to
$\widehat N$ and $\sqrt{M_e}b$.  A sufficiently large fixed real radial
pivot absorbs both while preserving a strict positive real residue.
For $n=0$ the latter coupling is exactly absent.

\paragraph{Second normal: cancellation of both source openings.}
Define
\[
 K_d=CH-(D+\tfrac12)(1-H),\qquad
 C_q=K_d+AH+\tfrac12
       +v^2\{CH(1+u^2)+AH+(\alpha+\tfrac12)H\}.
\]
The complete second normal is
\[
 L_M=C_q q+\mu uvb+v(\alpha-BTCHu^2)R.                  \tag{RFN22}
\]
The actual source satisfies the three exact identities
\[
\begin{split}
 S_0&=T(K_d+AH+\tfrac12)-A=0,\\
 S_1&=jK_d+\mu+xTCH=0,\\
 S_2&=jK_d+xTCH+T\{CH+AH+(\alpha+\tfrac12)H\}
          -A-(\alpha+\tfrac12)(1-p)+\alpha/B=0.
\end{split}
\]
Before imposing these identities, the complete derivative minus
(RFN22) is exactly
$uS_0+u^2vzS_1/j+uv^2(S_2-S_1)$.
This records all state terms and rules out a hidden $j$-weighted
radial error.  In particular $S_1$ gives
$K_d=-\mu/j-BTCH=O(1)$ even when $C$ or $D$ is unbounded.
Consequently $TC_q=O(T+v^2)\in\mathfrak n$.

Normalize
\[
 M_0=\sqrt{TH(1-H)}\,\widehat M.
\]
The full diagonal (RFN9), not only its zero-angle value, obeys
\[
 TH(1-H)H_M\sim k/2>0.                                 \tag{RFN23}
\]
Indeed the $C$ and $D$ terms contribute
$(k/2)D_u\{(1-H)D_v+H\}$ up to infinitesimal bounded coefficients,
and all remaining terms are infinitesimal after this normalization.
This is uniform at both independent openings.
The normalized $C_qq$ coupling has coefficient
$T\sqrt{H(1-H)}C_q\in\mathfrak n$ on the strict row $q/\sqrt T$.
The normalized $\mu uvb$ row is free because $\mu=O(p)$ and
$\sqrt T\,puvb$ has already been purchased.

It remains to allocate the entire opposite row, including its
higher testing terms.  Put
\[
 J=B(j+1)+u^2\in\mathfrak n,\qquad
 a_R=Tq_zv^2(\alpha-BTCHu^2)^2.
\]
The factor $H(1-H)$ may be discarded from the cost bound, since it is
bounded.  Since $\alpha-BTCHu^2=O(BJ)$, the exact squared multiplier
$a_R$ satisfies
\[
 \frac{a_R}{f_c}=O((T/j)v^2J^2),\qquad
 \frac{a_R}{q_cy}=O(v^2J^2).                            \tag{RFN24a}
\]
If $1-p$ is a unit, its direct cost against $r_cR_*^2$ is
$O(Tv^2J^2/(1-p))\in\mathfrak n$.
If $p\sim1$ and $\eta_*$ is a unit, apply (RFN18a); the cost is
$(a_R/f_c)(1+u^2/x)\in\mathfrak n$, using that $j^2B$ is a unit.
If $p\sim1$ and $\eta_*\in\mathfrak m_c$, (RFN3d) gives
\[
 \frac{a_R}{M_e yx}
     =O\!\left(\frac{T(1+1/n)}{j^2}v^2J^2\right)=O(1),  \tag{RFN24b}
\]
because $(1+1/n)/j^2=O(1)$ in this source case.  The other two
conditions in (RFN18d) follow from (RFN24a).  The exact budget
(RFN18c) therefore provides this row at infinitesimal cost as well.
These estimates treat the whole square
$(\alpha-BTCHu^2)^2$; there is no truncation at a testing degree.

\paragraph{Full Gram, legal roots, and restoration.}
After the normalizations, the entire remaining mixed hard term is
\[
 2(\alpha+\tfrac12)uv\sqrt{TH(1-H)}
    \left(\frac{\widehat N}{\sqrt{n+h_0}}-c_bb\right)\widehat M.
\]
Its coefficients relative to $\widehat N$, $\widehat M$, and, when
present, $\sqrt{M_e}b$ are infinitesimal; (RFN16) controls the shift.
All other remaining normal couplings were allocated above.  Choose a
fixed sufficiently large radial pivot, and then purchase all the
finitely many free helpers.  Their scalar costs lie in $\mathfrak n$,
so fixed positive fractions of (RFN17) survive.  The complete bounded
Gram has strict positive real residue: the $\widehat N$ pivot is one,
the $\widehat M$ pivot has residue $\overline{k}/2$, the retained scalar pivots
are positive, and the only bounded radial off-diagonals have been
absorbed by the chosen real radial pivot.

The separated strict bounded-Gram lemma now gives a finite sum of
squares in $\mathcal T$.  Its row entries need not be bounded.  All
budget remainders are explicit finite separated sums of squares, all
source-only completions use positive pure source denominators, and
all state-coordinate inverses are pure testing.  Restore the whole
squares (RFN15), the retained $nv^2(N_0+pb)^2$ square, the scalar
budget remainders, the $mpD_u^2N_0^2$ positive square term, and the pure testing
factor $D_u$.  This certifies the complete original parent (RFN4).
Every bounded coefficient used in the argument depends only on the
source and on $u,v$.  The arbitrary state ratio $z$ remains inside
retained rows throughout.  Hence the proof has no bounded-state
restriction, and it includes literal $n=0$ without any limiting
argument.

\subsection{All source ratios and testing angles}
\label{component:pbaab204a8ab1}

\label{pbaab204a8ab1:sec:new-r1-all-tests}

\paragraph{Statement and exact source.}
Use the original common-II source, with $d=1$, $\lambda=\beta=0$,
the six native coefficients of the actual
Section~\ref{component:pd709027c7805}, every original strict gap,
and $n\geq0$.  In the original R1 sector,
\[
 t_s,\quad B_s,\quad a_sB_s,\quad t_sB_s>0
 \quad\hbox{are infinitesimal}.                         \tag{R1N1}
\]
There is no bound on $a_s$, on any original positive opening, or on
the source ratio previously called $q_{\rm old}$.  The conclusion is:
the entire original parent is a finite separated tensor sum of
squares at every independent real projective measurement pair and
every real or complex state.  This proof treats R1 only.

Source and testing quantities belong to separate real closed fields
$K_c,K_t\supset\mathbb R$.  All SOS assertions are in
$K_c\otimes_{\mathbb R}K_t$.  We use the actual actual strict
bounded-Gram lemma: a bounded finite Gram with positive definite
real residue gives tensor SOS, with no boundedness requirement on
its row entries.  No order or mixed inverse is imposed on the
tensor ring.

Put $c_s=1+a_sB_sr$, $R_s=1+ha_sB_s$, and define
\[
\begin{aligned}
 j&=\frac{a_sc_s}{R_s},\quad B=\frac{B_sR_s}{c_s},\quad
 T=\frac{t_sc_s}{R_s},\quad p=\frac{a_s(1-r)}{R_s},\\
 H&=\frac{1-h}{t_sc_s},\qquad x=jB,\quad y=TB,\quad q_z=jT.
\end{aligned}                                          \tag{R1N2}
\]
The native strict gaps give $0<p<\min(1,j)$ and $0<H<1$.
They also give $1-h<t_sc_s$, so $h\sim1$.
Consequently $B,T,x,y$ are infinitesimal, while
\emph{$q_z$ is an arbitrary positive source quantity}.
In fact $x=a_sB_s$ and $y=t_sB_s$ exactly.
The inverse source substitution is
\[
 r=\frac{j-p}{j+px},\quad h=\frac{1-HT}{1+HTx},\quad
 a_s=\frac{j+px}{1+HTx},\quad
 B_s=\frac{B(1+HTx)}{1+pB},\quad
 t_s=\frac{T(1+pB)}{1+HTx}.                             \tag{R1N3}
\]
All arguments below therefore concern the actual original source.

The source bounds are RSR7--RSR9: their proofs use only $B,T,Bj$ infinitesimal and $p,H<1$, and impose no condition on $jT$.

\subsubsection*{The complete small-angle theorem extends to all R1 sources}

The exact proof Section~\ref{component:p261f71af685d}, RFN4--RFN24,
is a full-parent proof, not just a restriction to a scalar graph.
We now justify its application under the larger source domain
(R1N1), rather than importing its former domain label.

First, all of its exact identities hold for the rational dictionary
(R1N3), without any size assumption on $q_z$.  The parent, its two
normals, their quadratic block, the scalar circuit, the radial Gram
and the three source cancellation identities were all proved as
identities of this dictionary.  Second, every uniform source bound
used in that proof is exactly (R1N4)--(R1N7) or the following scalar
bound:
\[
 f_c\asymp Tx^2,\quad q_c\asymp Tx,\quad
 r_c\asymp TjB^2(1-p),\quad
 g_c/f_c=O(Hy),\quad |k_c|/(q_cy)=O(p+B).                \tag{R1N9}
\]
To check the latter bounds without $jT\in\mathrm I$, use the explicit
finite polynomials RSR15 and RSR18: each term of
$\mathcal F_1/K,\mathcal F_2/K$ is bounded using $p/j<1$
and $Bj=x$, and $\mathcal K/K=O(p+B)$ by the same inequalities.
Thus $\mathcal F\sim K$ and all statements (R1N9) follow from
RSR16--RSR20.  No $jT$ factor occurs in these estimates.

All source and normal estimates in RFN18--RFN24 use precisely these
source bounds. In particular the apparently unbounded factor $q_z=jT$
cancels in the exact ratios RFN24a--RFN24b; their proofs use $Bj=x$,
$B,T,x,y$ infinitesimal, and no restriction on $jT$. The effective
$n$ square is completed before the signed scalar terms are allocated.
Consequently that full scalar-and-two-normal calculation proves every
R1 source in the local testing chart, including $n=0$.

\subsubsection*{Outer testing angles: two exact current identities}

We include the outer mechanism and its precise positivity condition.
Let
\[
 R=uW-Y,\quad C_t=v(uX-Z),\quad B_t=X+uZ,\quad
 N=R+h_*C_t,\quad D_u=1+u^2,\quad D_v=1+v^2,
\]
and let $K$ be the signed quarter-core
\[
 K=\tfrac12(R+C_t)^2+\tfrac12D_u(vY-Z)^2
             +\tfrac12D_v B_t^2-\tfrac12D_uD_vZ^2.       \tag{R1N12}
\]
All unselected original nonnegative events are kept as whole squares.

If $1/v$ is bounded, use
\[
 K_0=2\sqrt{g_*}C_t,\quad
 B_0=\sqrt{2D_v/D_u}\,B_t,\quad
 c_u=D_u^{-1/2},\quad t_u=uD_u^{-1/2},\quad
 \eta_v=\sqrt{D_v}/v .
\]
The exact full selected-parent identity is
\begin{align*}
 4CD_vR^2+4D(R+C_t)^2+4K
 ={}&4(C+D)N^2+4Cv^2R^2+2(R+C_t)^2\\
 &+2D_u(vY-Z)^2+(K_0,B_0)Q(K_0,B_0)^T,\\
 Q={}&
 \begin{pmatrix}
 1-c_u^2\eta_v^2/(2g_*)&t_u\eta_v/\sqrt{2g_*}\\
 t_u\eta_v/\sqrt{2g_*}&1
 \end{pmatrix}.
\tag{R1N13}
\end{align*}
Since $1/g_*$ is infinitesimal and all the testing factors are bounded,
$\overline Q=I$.  This treats every other angle and every state,
including Bob infinity by homogeneous clearing.

For $v$ infinitesimal put $\epsilon=1/g_*$, $q=\sqrt\epsilon$,
$\gamma=Ag_*$, and
\[
 K_1=2\sqrt{g_*}C_t+\sqrt\epsilon Y,\quad
 B_1=B_t-uvY,\quad V_0=vY,\quad Y_1=\sqrt{\epsilon D_u}\,Y,
\]
\[
 L_0=\tfrac12qc_u(K_1-c_uY_1)-t_uvB_1+c_uvV_0.
\]
Then $L_0=\sqrt{D_u}\,v(vY-Z)$ exactly, and
\begin{align*}
 &4CD_vR^2+4D(R+C_t)^2+4AD_uD_vY^2+4K\\
 &\quad=4(C+D)N^2+4Cv^2R^2+2(R+C_t)^2+\mathcal G,\\
 \mathcal G={}&K_1^2+2D_v(B_1^2+V_0^2)
 +(4\gamma D_v-c_u^2)Y_1^2\\
 &+2qv(K_1-c_uY_1)V_0-2L_0^2 .
\tag{R1N14}
\end{align*}
The complete four-root Gram has residue
\[
 \operatorname{diag}(1,2,2,\overline{k}^{\,2}
                                      -\overline{c_u}^{\,2}).    \tag{R1N15}
\]
Thus it is strict whenever the last entry is positive.  In particular
it covers every unbounded $u$, and covers every noninfinitesimal
$u$ whenever $\overline k=1$.  All coefficients in the identity
are bounded functions of the displayed normalized homogeneous Alice
coordinates, so it includes the literal Alice-infinite endpoint.

\subsubsection*{The remaining bounded nonzero Alice chart}

The only outer chart left by (R1N13)--(R1N15) has $v$ infinitesimal,
$u$ a bounded positive unit, and $\overline k<1$.  Since
$1-k=p/(j+1)$, this implies that $p,j$ are positive units and $j$
is bounded.  Therefore $\eta_*\in\mathrm I$ and (R1N6) gives the
exact relative critical limits
\[
 m\sim\frac{jp}{2(1-p)(j+1)},\qquad
 n\sim\frac{p}{2(j-p)(j+1)}.                              \tag{R1N16}
\]
These statements allow either or both coefficients to be unbounded.

Choose bounded positive real caps $m_c\le m$, $n_c\le n$ as specified
below and keep their nonnegative excesses.  Retain all of $\alpha$
as its original square.  Define six rows
\[
 \widehat N=\sqrt{C+D}(R+h_*C_t),\quad
 P=C_t/\sqrt T,\quad Y_0=\sqrt T\,Y,\quad
 X,\quad Z,\quad V=vY.                                  \tag{R1N17}
\]
They need not be independent or bounded.  Put
$\delta=(C+D)^{-1/2}=O(\sqrt T)$, so
$R=\delta\widehat N-h_*\sqrt T\,P$.
The selected \emph{entire} parent is exactly
\begin{align*}
 \mathcal P_{\rm sel}={}&
 \widehat N^2+Cv^2R^2+\tfrac12(R+C_t)^2
       +(Tg_*)P^2+(A/T)D_uD_vY_0^2\\
 &+m_c\{(v/u)R+(D_u/u)V-X-uZ\}^2+n_cD_uD_vX^2\\
 &+\tfrac12D_u(V-Z)^2+\tfrac12D_v(X+uZ)^2
                                      -\tfrac12D_uD_vZ^2 .
\tag{R1N18}
\end{align*}
This is obtained by the exact completion
$CD_vR^2+D(R+C_t)^2=(C+D)N^2+g_*C_t^2+Cv^2R^2$.
It keeps every state-normal term.

The row relation
\[
 PY_0=V(uX-Z)                                            \tag{R1N19}
\]
allows an arbitrary real zero shift
$2b(PY_0-V(uX-Z))$.  Keep the two explicitly nonnegative squares
$Cv^2R^2$ and $(R+C_t)^2/2$ whole.  The remaining complete Gram
on (R1N17) is bounded: $u,1/u$ are bounded, $Tg_*,A/T$ are
units, and the $R$ contribution in the $m_c$ square has coefficients
$\delta v/u$ and $h_*\sqrt T v/u$, both infinitesimal.
Its residue is the direct sum of a unit pivot and the two blocks
\[
 \begin{pmatrix}k_0/2&b\\b&k_0(1+u_0^2)/2\end{pmatrix},\qquad
 Q_3(m_c,n_c;b),                                        \tag{R1N20}
\]
where subscripts denote real residues and $Q_3$ is the quadratic
form in $(X,Z,V)$
\begin{align*}
 Q_3={}&m_c\{(D_u/u)V-X-uZ\}^2+n_cD_uX^2\\
       &+\tfrac12D_uV^2-D_uVZ+\tfrac12(X+uZ)^2
                                            -2bV(uX-Z).
\tag{R1N21}
\end{align*}
In (R1N21), $u=u_0>0$ is real.  The choices of $b,m_c,n_c$
that make both blocks strict are now explicit.

\paragraph{Interior openings.}
Suppose $p_0<1$ and $j_0>p_0$.  Use initially
$m_0=j_0p_0/[2(1-p_0)(j_0+1)]$ and
$n_0=p_0/[2(j_0-p_0)(j_0+1)]$, and take $b=k_0/2$.
The first leading principal minor of $Q_3$ is
$m_0+n_0D_u+1/2>0$; the other two are
\[
 \frac{p_0u_0^2D_uK_0}
      {4(j_0-p_0)(j_0+1)^2(1-p_0)},\qquad
 \frac{p_0u_0^2D_uK_0^2}{8(j_0-p_0)(j_0+1)^3},          \tag{R1N22}
\]
where $K_0=j_0+1-p_0$.  They are positive.
The first block in (R1N20) has determinant
$k_0^2u_0^2/4>0$.  Positive definiteness is open, so fixed
real caps slightly below $m_0,n_0$ preserve strictness and are
smaller than the actual source coefficients.

\paragraph{Only the \(n\) opening is infinite.}
If $j_0=p_0<1$, then $n$ is unbounded and
$m_0=j_0^2/[2(1-j_0)(j_0+1)]$ is finite positive.
At $b=k_0/2=1/[2(j_0+1)]$, restrict (R1N21) to the kernel
$X=0$ of the $n$ positive square term.  Its first diagonal on $(Z,V)$ is
$u_0^2(m_0+1/2)>0$, and its determinant is
$u_0^2k_0^2/4>0$.  Choose $m_c<m_0$ sufficiently close and
then a sufficiently large fixed real $n_c$.  The whole block is
strict, with both caps below the original weights.

\paragraph{Only the \(m\) opening is infinite.}
If $p_0=1<j_0$, then $m$ is unbounded and
$n_0=1/[2(j_0^2-1)]$ is finite positive.
Take a real $b$ just above $k_0/2=j_0/[2(j_0+1)]$, with
\[
 b<\min\{k_0\sqrt{D_u}/2,\ j_0/[2(j_0-1)]\}.
\]
This interval is nonempty.  The first block in (R1N20) is strict.
On the kernel $(D_u/u)V-X-uZ=0$ of the $m$ positive square term, use $(X,V)$
coordinates.  The first diagonal of (R1N21) is
$D_u/[2(j_0^2-1)]>0$, and its determinant is
\[
 -\frac{D_u^2[\,2b(j_0-1)-j_0\,][\,2b(j_0+1)-j_0\,]}
        {4u_0^2(j_0^2-1)}>0.                            \tag{R1N23}
\]
Choose $n_c<n_0$ sufficiently close and then a sufficiently large
real $m_c$.  This gives the required full strict block.

\paragraph{Both openings are infinite.}
If $p_0=j_0=1$, both $m,n$ are unbounded and $k_0=1/2$.
Choose
$1/4<b<\sqrt{D_u}/4$.
The common kernel of the two event positive square terms is
$X=0$, $Z=D_uV/u_0^2$.  On it the rest of (R1N21) is exactly
\[
 \frac{D_u(4b-1)}{2u_0^2}V^2>0.                        \tag{R1N24}
\]
A sufficiently large real cap on both positive square terms therefore gives
a strict block.

The elementary finite-cap fact used in these cases is as follows:
if a real positive semidefinite event Gram is positive definite on
the complement of its kernel, and the remaining real quadratic
form is positive definite on that kernel, a sufficiently large
real multiple of the event Gram makes their sum positive definite.
It follows by a block Schur complement, and works for the sum of
two event Grams as well.  Thus the choices above prove strictness
of the \emph{complete} bounded Gram (R1N18), without restricting
any normal or state coordinate.

\subsubsection*{Exhaustive testing and state join}

The sign gauges SJ3 preserve every original source positive square term and the
whole core, so affine $u,v$ may be taken nonnegative.
If $v$ is not infinitesimal, (R1N13) applies.
If $v$ is infinitesimal and $u$ is not infinitesimal, either
(R1N15) applies or the complete six-root argument
(R1N16)--(R1N24) applies.  These cases exhaust the outer angles.
If $u,v>0$ are both infinitesimal, the enlarged local FN theorem
proved above applies after SJ's full-state reduction to $w=1,z>0$;
SJ gives the other state signs and the axes directly.
Literal homogeneous measurement endpoints are given by SJ10 or
the homogeneous current identities, without a limiting argument.
The zero state is immediate.  For a complex state, the real
parent matrix splits its value into its real and imaginary
quadratic evaluations, to which the same proof applies.

Restore every capped excess, the original $n,m,\alpha$ positive square terms
not used in a particular chart, and all cleared positive pure
testing homogeneous factors.  All normalizations, radicals and
inverses are pure source or pure testing.  This proves the
complete original R1 source row at all independent tests and
states, without a source-ratio split or a comparison of source
and testing rates.

\section{Two small actual angles: the R2 sector}
\subsection{Bounded $a$ or an infinitesimal paired-event parameter}
\label{component:p393d3faa269a}

\label{p393d3faa269a:sec:r2-shifted-local}

\paragraph{Exact statement.}
Use the actual original weak-C2 source and all its original strict
gaps, with \(d=1\), \(\lambda=\beta=0\), and \(n\geq0\).
Suppose \(t_s\) is positive unbounded and
\(B_s,a_sB_s,t_sB_s\) are positive infinitesimals.
Put \(e=1-r\), \(R_s=1+ha_sB_s\).  On the source subset
\[
       a_s=O(1)\quad\hbox{or}\quad
       p=a_s e/R_s\in\mathrm I_c,                     \tag{RL1}
\]
the entire parent is separated tensor SOS for every independent
projective testing pair and every real or complex state.
Literal \(n=0\) and every opening of \(r,h,1-h,g\) are retained.
For the all-testing conclusion, the local proof below is joined
explicitly to R2O1--R2O24.

Use the field and bounded-Gram conventions of EN1--EN4.

\subsubsection*{Actual source entry and bounds}

Use the exact compact source dictionary RSR2--RSR8, with
\[
 j=a_sc_s/R_s,\quad B=B_sR_s/c_s,\quad T=t_sc_s/R_s,\quad
 p=a_se/R_s,\quad H=(1-h)/(t_sc_s),\quad c_s=1+a_sB_sr.
\]
Set
\[
 S=1/T,\quad \ell=HT,\quad \bar h=1-\ell,\quad
 x=jB,\quad y=TB,\quad B=Sy,\quad q_z=j/S .
                                                               \tag{RL2}
\]
Thus \(S,x,y\) are positive infinitesimals,
\(0<p<\min(1,j)\), \(0<\ell<1\), and \(j\asymp a_s\).
Both \(\ell\) and \(\bar h\) may be arbitrarily small.
The original physicality inequality RER2--RER7 gives \(jy\in\mathrm I_c\).
In terms of the actual finite polynomials \(\mathcal P,\mathcal L,
\mathcal N,\mathcal F,\mathcal K\) of RSR, put
\[
 K=j+1-p,\quad k=K/(j+1),\quad
 F_b=j+1+p(1-p)/j,\quad E=1+B(\ell+p-1).
                                                               \tag{RL3}
\]
Here \(1/2<k<1\) and \(1\le F_b/(j+1)<2\).  Direct substitution
\(B=Sy,T=1/S,H=S\ell\) in those finite polynomials gives
\[
 \mathcal P/(j+1)\sim1,\qquad
 S\mathcal L/K\sim1,\qquad S\mathcal F/F_b\sim1.         \tag{RL4}
\]
For clarity, the terms after division contain only bounded
\(p/j,j/(j+1),K/(j+1)\), multiplied by \(S,y\), or \(jy\).
No inverse of \(\ell\), \(\bar h\), or \(j-p\) is involved.

The same substitution in the exact numerator of \(n\) yields
\[
 \mathcal N=p(1+\epsilon_p)
                -yj(j+1)(1+\epsilon_y),\qquad
 \epsilon_p,\epsilon_y\in\mathrm I_c.                  \tag{RL5}
\]
For an explicit finite division certificate, the zero-\(p\) term is
\[
 \mathcal N|_{p=0}
 =-yj(j+1)\{1+S(1-\ell)\}
             \{1+S\ell yj/(j+1)\}.
\]
The quotient \((\mathcal N-\mathcal N|_{p=0})/p\) is a
finite polynomial.  This grouping can be read termwise from RSR5:
the first
coefficient is \(1+y(j+1)+O(Sy(j+1)+Sy^2(j+1))\);
the second coefficient is \(1+O(S+y+Syj)\).
Using \(\mathcal N\geq0\), therefore,
\[
 \eta_2=\frac{yj(j+1)}p=O(1),\qquad
 \eta_2\in\mathrm I_c\ \Longrightarrow\
 n\sim\frac{p}{2(j-p)(j+1)}>0 .                        \tag{RL6}
\]
The estimate includes the face \(\mathcal N=0\), which can occur
only in the positive-unit \(\eta_2\) case.

Let \(\mu=m(1-p)\) and \(h_0=\mu+1/2\).
The complete source formulas now imply
\[
\begin{gathered}
 \mu\sim\frac{jp}{2(j+1)},\quad h_0\asymp1,\quad
 CH\sim k/2,\quad SA\bar h\sim k/2,\quad D/S\sim k/2,\\
 \alpha\sim\frac{Sy^2K^2}{2(j+1)},\quad
 f_c\sim\frac{j^2y^2F_b}{2(j+1)},\quad
 r_c\sim\frac{jy^2(1-p)K}{2(j+1)} .
\end{gathered}                                                   \tag{RL7}
\]
These use the exact masses in RFN12 and the closed RSR formulas.
All displayed positive masses are strictly positive.
Write \(\kappa=-k_c>0\).  The exact positive-surplus formula
RSR19 proves this sign.  Its substituted finite polynomial gives
\[
 \frac{\kappa}{r_c}=O(p+Sy)\quad(p\in\mathrm I_c),\qquad
 \frac{\kappa}{jy^2}=O(1).                              \tag{RL8}
\]
For the remaining case in (RL1), \(p\) and \(j\) are positive
bounded units, and
\[
 \frac{\kappa}{f_c}\sim\frac{pK}{jF_b}\leq\frac12,
 \qquad jF_b-2pK=(j-p)K\geq0.                          \tag{RL9}
\]
The inequality is a bound on real residues when combined with
the preceding asymptotic equality.  On this case \(\eta_2\) is
infinitesimal, and the effective current defined below is a
positive bounded unit.

\subsubsection*{The entire parent and its harmonic shift}

Retain the complete original quarter-parent RFN4.  Its exact inverse
state change at \(w=(W+uY)/(1+u^2)=1\) is
\[
 X=N_0+pv-uz,\quad Z=z+uN_0+puv,\quad
 W=1+u(M_0+Hvz),\quad Y=u-M_0-Hvz.                    \tag{RL10}
\]
Here \(z>0\) is arbitrary.  Retain \(mpD_u^2N_0^2\) and divide
the remaining whole parent by \(D_u=1+u^2\).
The exact expression is RFN7, with both complete state normals,
hard entries RFN9, and scalar RFN12.

Use the additional exact change
\[
 q=vz-Tu,\qquad L=M_0+Hq,\qquad
 M_0=L-Hq,\qquad L=\sqrt{S\ell\bar h}\,\widehat L .
                                                               \tag{RL11}
\]
All divisions are pure source or pure testing.  In particular
this change does not invert the harmonic mixed unit
\(1+\rho u^2\).
The complete diagonal of \(\widehat L\) is
\[
 S\ell\bar h\{D_u(CD_v+D+\tfrac12)+AD_v+
                   (\alpha+\tfrac12)v^2\},
                                                               \tag{RL12}
\]
with positive real residue \(\overline k/2\) and bounded
coefficient array through both openings.

Use the circuit coordinates \(U,V,Z_*,Q_*,B_*,R_*,\Psi\)
of SH1--SH2, and \(b=v-uz/j=\sqrt y B_*\).
Before spending any scalar \(n\)-mass, complete the exact correlated
block, as in RFN13--RFN16:
\[
 h_0N_0^2+n(N_0+jb)^2
 =(n+h_0)(N_0+c_bb)^2+M_e b^2,\quad
 c_b=\frac{nj}{n+h_0},\quad M_e=\frac{nh_0j^2}{n+h_0}.  \tag{RL13}
\]
Keep \(nv^2(N_0+pb)^2\) and all three additional positive hard
squares RFN15 whole, including after the change
\(\widehat N=\sqrt{n+h_0}(N_0+c_bb)\).
At \(n=0\), put \(c_b=M_e=0\) and omit the current row.
For \(n>0\), \(c_b^2/M_e=n/[h_0(n+h_0)]\) is bounded.

\paragraph{The exact shifted scalar.}
Let \(F\) be RFN12.  The shifted scalar is
\(F-2HqL_M+H^2q^2H_M\).
Use the whole \(d_cV^2Q_*^2\) term in its \(q^2\) diagonal.
Using the exact actual-source identity \(S_0=0\) in RFN22,
the resulting coefficient on \((\sqrt S q)^2\) is exactly
\[
 D_q=S\{A\bar h^2+\ell^2(C+D+\tfrac12)D_u\}
       -\frac{f_c}{Tx}-TH^2C\,u^2v^2
       -\frac{k_cv^2}{yTx}.                            \tag{RL14}
\]
It has bounded array and positive real residue \(\overline k/2\).
Indeed \(f_c/(Tx)=O(jy)\) is infinitesimal, \(TH^2C=\ell CH\)
is bounded, and \(\kappa/(yTx)=\kappa/(jy^2)\) is bounded.
Its equivalent \(Q_*^2\) mass has order \(jy\).
The only extra scalar cross terms are
\[
 -2Hq\{\mu uvb+v(\alpha-yCHu^2)R\},\qquad
 R=z-uv/(Sy).                                        \tag{RL15}
\]
The budgets below allocate their entire expressions.

The remaining signed scalar terms are exactly
\(e_0V^2B_*^2+k_cZ_*^2(V^2-1)^2\), where
\(e_0=p^2y^2/2-g_c\).
The negative part of \(e_0\) is at most \(g_c\).
SH2 purchases it at relative circuit cost
\[
 \frac{g_c}{f_c}(1+V^2)
   =O\bigl(S(y+v^2)\bigr)\in\mathrm I.                \tag{RL16}
\]
Use \(Z_*(V^2-1)=-VQ_*-R_*\) for the other signed term.
If \(p\) is infinitesimal, (RL8) pays its opposite part
directly from \(r_cR_*^2\), and its \(VQ_*\) cost is
\(O(v^2)\).
Otherwise \(p,j,M_e\) are positive units.  Young's exact
square identity with parameter \(1/6\), followed by SH1 at
\(a=3/2\), gives the legal SOS domination
\[
 \kappa Z_*^2(V^2-1)^2
 \preceq_{\rm sos}
 \frac74\kappa\Psi+
 \frac{21}{2}\kappa V^2Q_*^2+
 \frac72\kappa V^2B_*^2 .                             \tag{RL17}
\]
The circuit cost is at most \(7/8+o(1)\), by (RL9).
The two other relative costs are \(O(v^2)\).
Thus fixed positive real fractions of \(D_q(\sqrt S q)^2\),
\(r_cR_*^2\), \(f_c\Psi\), and \(M_eb^2\) remain;
the last is absent only at \(n=0\).
The opposite mass \(r_c\) was not spent in the second case.

\subsubsection*{All helper rows and both complete normals}

As usual, a free row permits any fixed real square pivot at
infinitesimal cost from the preceding retained scalar forms.
SH2 proves that both \(E_u=u(uvz-jy)\) and \(E_v=v(uvz-jy)\)
are free.  The row \(R_u=v(u^2-x)\) is free as well.
For \(p\) infinitesimal, use the actual RFN19 rational Gram:
its three costs are \(O(S^2y)\), \(O(Syu^2)\),
and \(O(Su^2/(1-p))\).
When \(p,j,M_e\) are units, apply SH1 to the \(R_*^2\)
term in RFN19.  The resulting costs are
\[
 O(S^2y+Su^2/j),\quad O(Syu^2+Su^2v^2),\quad
 O(jSu^2v^2/M_e),                                    \tag{RL18}
\]
all infinitesimal.

The rows \(pv^2b\) and \(\sqrt S\,\mu uvb\) are free.
If \(\eta_2\) is infinitesimal, (RL6) and (RL13) imply
\[
 \frac{p^2}{M_e}
  =\frac{p^2}{h_0j^2}+\frac{p^2}{nj^2}=O(1).
                                                               \tag{RL19}
\]
They then cost \(O(v^4)\) and \(O(Su^2v^2)\).
If \(\eta_2\) is a positive unit, the actual source gives
\(p^2y/f_c=O(1)\) and \(\mu=O(p)\).
The second identity in SH2 provides the same rows with costs
\[
 O\bigl(v^2(y+v^2)\bigr),\qquad
 O\bigl(Su^2(y+v^2)\bigr),                             \tag{RL20}
\]
including \(n=0\).

The full opposite term \(\sqrt S\,v(\alpha-yCHu^2)R\)
is free by SH4--SH8.  Also \(j\alpha b\) is free:
\[
 j\alpha b=-(\alpha/B)R_u-\alpha uR,\qquad
 \alpha/B=O(y(j+1))\in\mathrm I_c.
                                                               \tag{RL21}
\]
For \(p\) infinitesimal, the direct \(r_c\) cost for
\(\alpha uR\) is \(O(Sy^2(j+1)^2u^2)\).
In the unit \(p,j,M_e\) case use SH1; the three costs
are \(O(Sy^2u^2)\), \(O(Sy^2u^2v^2)\), and
\(O(Sjy^2u^2v^2/M_e)\).  All are infinitesimal.
These statements account for the whole source \(\alpha\), not
only its leading coefficient.

\paragraph{The first normal.}
The exact RFN21 source cancellations give, after (RL11),
\[
 L_N'=[j(n-\alpha)+p(n+\tfrac12)v^2]b
             +K_sR_u+\kappa_NE_v,                     \tag{RL22}
\]
\[
 K_s=-\mu+(\alpha+\tfrac12)(\ell+p-1),\qquad
 \kappa_N=-(n+\tfrac12)(1-p/j).
\]
Both coefficients are bounded.  For \(\kappa_N\), use the
exact numerator of \(n\) and
\(n(j-p)/j=O(p/[j(j+1)])\); this cancels the \(r\) opening.
After the entire correlated completion (RL13), the remaining
normal is
\(-j\alpha b+\tfrac12pv^2b+K_sR_u+\kappa_NE_v\).
Every row is free by (RL18)--(RL21).
Division by \(\sqrt{n+h_0}\) is bounded.
The shift generates couplings to \(\sqrt{M_e}b\) with bounded
coefficients \(c_b/\sqrt{M_e}\); they are absent at \(n=0\).
Large fixed real helper pivots absorb these couplings by exact
Young identities.

\paragraph{The second normal and the entire hard block.}
The complete RFN22 normal after (RL11) is
\[
 L_L=(C_q-HH_M)q+\mu uvb+v(\alpha-yCHu^2)R,\qquad
 C_q-HH_M=-D-H(C+D+\tfrac12)u^2 .                      \tag{RL23}
\]
The coefficient of \(\sqrt S q\) after the normalization
of \(L\) is infinitesimal; \(D\) is infinitesimal and
\(H(C+D+1/2)\) is bounded.
The other two terms are free by (RL20) and SH6--SH8.
Exactly the same rows pay (RL15), whose first factor is
\(\sqrt S q\) with a bounded source multiplier.

The remaining hard off-diagonal is
\((\alpha+1/2)uv\sqrt{S\ell\bar h}/\sqrt{n+h_0}\),
which is infinitesimal.  Its \(c_b b\) correction has the
same bound after division by \(\sqrt{M_e}\).
The complete finite Gram therefore has a positive definite
real residue, with principal normal pivots
\[
                   1,\quad\overline k/2,\quad\overline k/2
                                                               \tag{RL24}
\]
on \(\widehat N,\widehat L,\sqrt S q\), together with the
retained strict scalar and purchased helper pivots.
All extra positive hard squares and both original retained
current positive square terms are restored whole.  The strict bounded-Gram
lemma proves tensor SOS for the entire parent.

\subsubsection*{State and outer-chart exhaustion}

The fresh full-state sign lemma SJ1--SJ10 permits positive
\(u,v\), and on the remaining negative-core chart gives
\(wz>0\).  Divide by the nonzero pure-testing \(w\) to obtain
the normalization used above with arbitrary \(z>0\).
No estimate in this proof uses its size.  Literal axes and
homogeneous endpoints have the direct original-core SOS exit.
Real and imaginary state parts separate.
The exact outer theorem R2O1--R2O24 covers every other testing
pair on all original R2 sources.  This proves the stated
all-tests conclusion on (RL1).

The complementary source \(a_s\) unbounded and \(p\) a positive
unit is not inferred from this proof.  It requires the separate
TU source comparison and additional-current conductor.

\subsection{Circuit budgets for the opposite normal}
\label{component:p09b2107cc05e}

Use the compact actual-source dictionary of RSR1--RSR8.  In the R2
source put \(S=1/T\), \(x=jSy\), \(B=Sy\), and \(H=S\ell\), where
\(0<\ell<1\).  Both testing angles \(u,v\) are positive
infinitesimals.  Normalize the rotated testing coordinate \(w=1\);
the other rotated coordinate \(z>0\) is arbitrary, including
unbounded values.  No coefficient estimate below involves \(z\).

Set
\[
 U=\frac{u}{\sqrt x},\quad V=\frac v{\sqrt y},\quad
 Z_*=\frac z{\sqrt{j/S}},\quad
 Q_*=U-VZ_*,\quad B_*=V-UZ_*,\quad R_*=Z_*-UV,
\]
\[
 \Psi=U^2+V^2+Z_*^2+U^2V^2Z_*^2-4UVZ_*.
\]
All monomials and radicals separate between the two fields.

\paragraph{A parameterized opposite-square identity.}
For every positive real \(a>1\), let \(K=a/(a-1)\),
\(d=U-V\), \(t=UV\), \(E=Z_*(1-t)\), \(T_*=t(Z_*-1)\).
Exact multiplication gives
\[
\begin{split}
 a\Psi+K V^2(Q_*^2+B_*^2)-R_*^2
 ={}&a d^2+(a-1)(E-T_*/(a-1))^2\\
 &+(2at+KV^4)(1-Z_*)^2+2KZ_*(t-V^2)^2 .
\end{split}                                                     \tag{SH1}
\]
The right side is a finite legal separated SOS.  For example
\(a=3/2\) gives \(K=3\).  This is a different budget from the
older one with \(U^2B_*^2\); neither budget dominates the other
uniformly at independent source and testing rates.

The two symmetric current budgets are
\[
\begin{split}
 (1+U^2)\Psi&=R_*^2+B_*^2+U^2Q_*^2
                         +U^2(1-UVZ_*)^2,\\
 (1+V^2)\Psi&=Q_*^2+R_*^2+V^2B_*^2
                         +V^2(1-UVZ_*)^2 .
\end{split}                                                     \tag{SH2}
\]
In particular the two actual rows
\[
 E_u=u(uvz-jy),\qquad E_v=v(uvz-jy)                       \tag{SH3}
\]
are free from a fixed positive fraction of \(f_c\Psi\) whenever
\(f_c\asymp j^2y^2\).  Their respective relative circuit costs
are \(O(x+u^2)\) and \(O(y+v^2)\), both infinitesimal.  Here
\emph{free} means that an arbitrarily large fixed real multiple of
the row square can be inserted at infinitesimal cost in the
strict circuit Gram.  It is not a pointwise bound without an
SOS witness: multiply (SH2) by \(j^2y^2x\) or \(j^2y^3\).

\paragraph{The full opposite row after the harmonic shift.}
Let
\[
 \begin{gathered}
 q=vz-u/S,\quad R=z-\frac{uv}{Sy},\quad R_v=u(v^2-y),\\
 d_s=\frac{BCH(1-p)(1+x)}{E},\qquad E=1+B(HT+p-1).
 \end{gathered}
\]
The exact actual source identity is
\[
                 \frac{\alpha}{y}-CHx=d_s.             \tag{SH4}
\]
Indeed it is equivalent to
\((K+Hjy)-jE=(1-p)(1+x)\), after canceling the displayed
positive source factors in the RSR dictionary.
The two exact testing identities
\[
 u(u^2-x)=S E_u-Su^2q,\qquad
 v(u^2-x)=S E_v-Suvq                                  \tag{SH5}
\]
now give the complete identity
\[
\begin{split}
 \sqrt S\,v(\alpha-yCHu^2)R
 ={}&[\alpha-yCHu^2-CHu^2(v^2-y)]\,(\sqrt S q)\\
 &+\sqrt S\,CH(v^2-y)E_u-\frac{d_s}{\sqrt S}R_v .
\end{split}                                                     \tag{SH6}
\]
There is no restriction on \(z\), and no normal has been set to
zero.  The first coefficient in (SH6) is infinitesimal, because
\(CH=O(1)\), \(y,u,v\) are infinitesimal, and
\(\alpha=O(Sy^2(j+1))\) is infinitesimal by \(jy\in\mathrm I\).
The second coefficient is infinitesimal and (SH3) is free.

For the last term write
\[
 R_v=-\sqrt x\,y(Q_*+VR_*),\qquad
 \frac{d_s}{\sqrt S}
   =\sqrt S\,y\,CH(1-p)(1+x)/E.                         \tag{SH7}
\]
Suppose the scalar has retained strict pivots of respective orders
\(jy\) and \(jy^2(1-p)\) on \(Q_*^2,R_*^2\), as in the actual
R2 source.  Twice Young's identity and (SH7) provide this row at
relative costs
\[
 O(S^2y^4(1-p)^2),\qquad
 O(S^2y^2(1-p)v^2),                                   \tag{SH8}
\]
respectively.  Both are infinitesimal even when \(1-p\) has
arbitrarily small positive order.  No inverse of that opening
remains.  Multiplication by the additional bounded source
factor \(\sqrt{\ell(1-\ell)}\), used to normalize the complete
second normal, preserves the estimate.

This controls the opposite normal uniformly through all R2 source ratios
once the stated scalar pivots are retained. The signed scalar allocation
remains a separate step.

\subsection{Unbounded source ratio and bounded relative current}
\label{component:pb88ccd68a4bf}

\label{pb88ccd68a4bf:sec:r2-thin-new}

\paragraph{Statement and domain.}
Use the actual original weak-C2 source RER1--RER8, normalized by
$d=1$, with $\lambda=\beta=0$ and every original strict gap.
Suppose $t,a$ are positive unbounded, while $B,aB,tB$ are
positive infinitesimals, and put $e=1-r$.  Assume
\[
             \lambda_n=n/(tB)=O(1),\qquad n\geq0.       \tag{TU1}
\]
The entire original parent is tensor SOS for every independent
projective test and every state.  Literal $n=0$ is included.
No comparison of a testing quantity with a source quantity is made.
This theorem treats precisely this source sector; it is not a
completion of the other R2 sectors or the global finite-lift proof.

The bounded-Gram and finite toric allocation lemmas used below are
the actual retained SC16--SC20 and QI1--QI10 proofs.  We use their
printed finite algebra and prove its source bounds on the present
domain.  We do not apply the physical C1 theorem to a formal source
whose event coefficient could be negative.

\subsubsection*{An exact common contact and bounded source comparison}

Set
\[
 S=1/t,\quad z=1/a,\quad q=a(tB),\quad w=ae,\qquad
 b=\sqrt S,\quad g=\sqrt z,\quad c=\sqrt q.
                                                               \tag{TU2}
\]
Thus $B=Szq$, $e=zw$, the actual squared angles are
$u_c^2=Sq$ and $v_c^2=zq$, and $S,z,q$ are positive
infinitesimals.  Keep $h$ as the original C2 source parameter.
The strict $m$ gap gives
\[
 R=1+Sqh,\quad M=R-w>0,
 \qquad 0<w<R\sim1.                                      \tag{TU3}
\]
Both original openings $h,1-h$ remain arbitrary positive quantities.

The exact RER2 identity, divided by $e>0$, gives
\[
 Q_{\rm RER}
 =\tau\{(1+z)(1+Sh)(1+Sq)+2r(z\Delta)\lambda_n\},
 \qquad \tau=q/(zw),\quad \lambda_n=n/(qz).
\]
Here $Q_{\rm RER}\sim1$, $z\Delta\sim1$ and $r\sim1$.
The bracket is a positive bounded unit, since $\lambda_n\geq0$
is bounded and its first term is a positive unit.  Consequently
\[
 \tau\text{ is a positive bounded unit},\qquad
 \overline\tau=\frac1{1+2\overline\lambda_n},\qquad
 q=zw\tau.
                                                               \tag{TU4}
\]
No division by $n$ has occurred.  This sector includes the earlier
thin condition $\rho_n=n/e^2\in\mathrm I$: in that case the same
RER identity gives $\tau\sim1$ and
$\lambda_n=\rho_nw/\tau\in\mathrm I$ directly, without first
assuming (TU1).

For the comparison only, substitute into the native C1 rational
formulas NC1--NC4
\[
 k^2=Sq,\quad \eta=1/S,\quad \zeta=z,\quad
 r_{\rm C1}=h,\quad t_{\rm C1}=w,\quad
 a_{\rm C1}=S(1-h).                                     \tag{TU5}
\]
Their $\theta$ is $e=zw$ and their $R,M$ are exactly (TU3).
Write a superscript $0$ for these formal C1 coefficients.
All their source denominators have the requisite positive signs:
\[
\begin{gathered}
 K=1+Sq,\quad J=1+Sq(1-e)-S(1-h),\quad
 E=1-q(1-e),\quad Q=1-Sqzh,
\end{gathered}
\]
and $J,E,Q$ are positive bounded units.  Their common polynomial is
\[
 P=1+qz+Sq(2+z-2e)+Sq^2(z+e-e^2)+S^2hq^2(1+z-e).
                                                               \tag{TU6}
\]
In particular $P\sim1$.  We make no sign claim for $\beta^0$.

The two sources have exactly the same state contact.  With $Z_c=1$,
\[
 u_c=bc,\quad v_c=cg,\quad W_c=bg,\quad
 Y_c=hu_cW_c,\quad X_c=-(1-e)u_c.                       \tag{TU7}
\]
This follows either by substitution in the two native dictionaries,
or from $Y_c=hu_cW_c$, $X_c=-ru_cZ_c$ and
$Z_c/W_c=\sqrt{at}$ in the actual C2 source.  In particular
the comparison does not move the contact.

Use the actual, not formal, normalized coefficients
\[
 c_b=SC,\ a_b=SA,\ D_b=D/S,\ \alpha_b=\alpha/(Sz),\quad
 h_2=mzM,\quad h_b=zw/2.
                                                               \tag{TU8}
\]
They satisfy $m=h_2/(zM)$ and $\beta+1/2=h_b/(zw)$.
Use the same normalizations with superscript $0$, and let $\delta$
mean actual minus formal.  The following explicit polynomials make
all needed cancellations finite and opening-safe.  Put $G=1+z-e$,
and define
\begin{align*}
 P_C={}&1+qz+S\{1+q(2+2z-e)+q^2z\}\\
 &+S^2\{q[h(1+z)+1-e]
          +q^2[hez-e(1+z)+2z+1]\}+S^3hq^2G,\\
 P_A={}&1+qz+S\{q(2+z-2e)+q^2z(1-e)\}
                 +S^2q^2G(G-hz),\\
 P_D={}&1+q(z+e)+S\{q[h(1+z)+1-e]+q^2(z+e-e^2)\}
                 +S^2hq^2G,\\
 P_m={}&1+qz+S\{q(2+z-2e)+q^2z\}
                 +S^2q^2\{(1-e)^2+z(1-e)+hez\},\\
 P_b={}&(1+z)(1+qz)+Sq(1+z)(2+z-2e)\\
       &+Sq^2z(1+z-e^2)+S^2q^2G^2.
\end{align*}
Then the exact rational identities are
\[
\begin{aligned}
 (1-h)\delta c_b&=-\frac{n(1-e)P_C}{zP},&
 h\delta a_b&=-\frac{n(1-e)P_A}{zP},\\
 \delta D_b&=-\frac{n(1-e)(1+qz)P_D}{zJP},&
 \delta h_2&=\frac{n(1-e)(1+qz)P_m}{zP},\\
 \delta h_b&=\frac{n(1-e)P_b}{zP},&
 \delta\alpha_b&=
 \frac{n(1-e)q^2(1+qz)[(1+S)e-S(1-h)(1+z)]}{zQP}.
\end{aligned}                                                   \tag{TU9}
\]
These identities are proved simply by substituting the six printed
RER1 weights and the six NC4 weights, clearing the indicated pure
source denominators, and multiplying.  They include $n=0$; the
right sides define no quotient by $n$.

After substituting $e=zw$, the polynomial
\[
       P_\Delta=\frac{P_b-(1+qz)P_m}{z}                 \tag{TU10}
\]
is a genuine finite polynomial in $S,z,q,w,h$.  Indeed the numerator
specialized at $z=0$ is zero identically, so ordinary monic division
by $z$ has zero remainder.  Consequently
\[
 \delta(h_b-h_2)=n(1-e)P_\Delta/P.
                                                               \tag{TU11}
\]
Every polynomial in (TU6), (TU9), (TU10) has bounded coefficients
and bounded arguments.  From (TU4) and (TU9)--(TU11) we obtain
\[
\begin{gathered}
 (1-h)\delta c_b,\ h\delta a_b,\ \delta D_b
                     =q\lambda_n\,O(1),\qquad
 \delta h_2,\delta h_b=q\lambda_n\,O(1),\\
 \delta\alpha_b/q^2=\lambda_n q\,O(S+e),\qquad
 \delta(h_b-h_2)/(qz)=\lambda_n O(1).
\end{gathered}                                                   \tag{TU12}
\]
Here $O(1)$ denotes an explicit bounded rational source array with
denominators among the units $P,J,Q,R$; it does not hide a state
or testing factor.  These conclusions require no lower bound on
$h$, $1-h$, or $M$.

In particular $c_b\asymp1/(1-h)$, $a_b\asymp1/h$,
$D_b\sim1/2$, $\alpha_b/q^2\sim1/2$, and
$h_2/q,h_b/q$ positive bounded units with common residue
$1/2+\overline\lambda_n$.  We also have
\[
 h_b/w=z/2,\qquad h_b-h_2=O(qz),\qquad
 H_0:=4h_2/R+4h_b/w\geq2z,
 \quad H_0\asymp q+z.                                  \tag{TU13}
\]
The last statements use the formal exact SC2 difference together
with (TU11), and hold also when $M$ is arbitrarily small.

\subsubsection*{The complete additional current, including its normal}

Let the testing $u,v$ be positive infinitesimals and first take
$Z=1$ on its pure-testing state chart.  Use the exact SC coordinates
\[
 Q_b=(W+uY)/b,\quad Y_b=Y/S,\quad p=w/R,\quad
 N=((X+u)/b-pvQ_b)/g,
\]
\[
 H=uQ_b/b,\quad L=vQ_b/g,\quad
 \sigma=c_b+a_b,\quad \rho_0=Kh/R,\quad
 N_y=\sqrt\sigma\{D_uY_b-\rho_0H\}.
                                                               \tag{TU14}
\]
Both $N,N_y$ vanish at (TU7).  The exact SC3 full-parent identity
is an algebraic identity for arbitrary $c_b,a_b,D_b,\alpha_b,h_2,h_b$;
no C1 coefficient relation is used in deriving it.  With (TU8),
it represents all the actual C2 parent except its additional term
$4nD_uD_vX^2/S$.  Retain its original nonnegative $m$ positive square term
$4h_2wN^2/(MR)$, incorporate the second paired positive square term into
$H_0N^2$, and retain $4h_b(vN)^2/w$ exactly as in the uniform
QI normal lemma.  Define $\widehat N=\sqrt{H_0}N$.

Put $h_t=H/R$, and introduce the same four positive separated
monomials and their geometric mean as in QI:
\[
 A_*=ch_t,\quad B_*=cL/R,\quad C_*=uv/b,\quad
 D_*=qg,\qquad \Gamma=(A_*B_*C_*D_*)^{1/4}.
                                                               \tag{TU15}
\]
Their exact relations include
$\Gamma^2=qh_tv$ and $B_*C_*=(c/g)v^2h_t$.
The extra monomial
\[
                   U_*=cgu/b
\]
has the two exact identities
\[
       U_*B_*=\Gamma^2,\qquad
       cgX/b=czN+eB_*-U_*.
                                                               \tag{TU16}
\]
The entire additional current therefore has the following square
completion, with no state normal set to zero:
\begin{align*}
 \frac{4n}{S}D_uD_vX^2
 ={}&4\lambda_nD_uD_v(U_*-B_*-czN)^2\\
 &+4\lambda_nD_uD_v
   \{2(1-e)\Gamma^2-(1-e^2)B_*^2
                         -2(1-e)czB_*N\}.
                                                               \tag{TU17}
\end{align*}
The first line is a nonnegative separated SOS and vanishes at the
actual contact, including its unrestricted state-normal coordinate.
Its coefficient is literally zero if $n=0$.  The remaining scalar
changes have bounded coefficients on $B_*^2$ and $\Gamma^2$.
Their possibly nonzero real residues will cancel exactly with
the source comparison below.  The full additional first-normal coefficient, on
$B_*\widehat N$, is infinitesimal because
\[
       \lambda_n cz/\sqrt{H_0}=O(\lambda_n c\sqrt z)
                                            \in\mathrm I.
                                                               \tag{TU18}
\]
In particular (TU17) does not estimate or discard $u/b$, which can
be arbitrarily large.  The new positive square contains that entire
row, while the exact product relation in (TU16) places its cross
term on the old angular circuit.

\subsubsection*{Why the full QI allocation is stable here}

We verify the actual finite hypotheses of QI1--QI10 rather than
an abstract continuity assertion.  Its coefficient arrays have
only the pure-testing denominators $D_u$ and positive pure-source
units after cancellation of the two $h$ openings.  The latter
cancellation is preserved exactly because
\[
 1-\rho_0=(1-h)/R,\qquad \rho_0=Kh/R.
                                                               \tag{TU19}
\]
For example all perturbations of
$c_b(1-\rho_0)^2$, $a_b\rho_0^2$,
$c_b(1-\rho_0)$ and $a_b\rho_0$ are $q\lambda_n$
times bounded rational arrays, by (TU9).  Moreover
\[
 \delta(\sigma\rho_0-c_b)
       =\rho_0\delta a_b-(1-\rho_0)\delta c_b
                             =q\lambda_n O(1).
                                                               \tag{TU20}
\]
This is the complete potentially singular harmonic first-normal
coefficient after multiplication by $\sqrt\sigma$.
Both $\sigma$ and its formal comparison are comparable to
$[h(1-h)]^{-1}$, and their ratio has residue one.
Their square-root ratios and inverses are consequently bounded
pure-source units.  Thus this change of hard normal loses no
opening bound.

In the exact generic scalar QI2, the changes of $a,b_*,c_*$ are
$q\lambda_n$ times bounded arrays in $u,v$.  This follows term by
term from (TU12), (TU19): its only $C,A$ occurrences are the first
two expressions in (TU19); its $D_b$ occurrences have bounded
factors; its $\alpha_b$ occurrences have bounded factors and the
additional small factors printed in QI2.  Consequently the
formal specialization $a_0,b_0,c_0$ and its exact $U,V$ divisions
in QI3--QI4 can be used unchanged, while the remaining $q$
coquotients change by bounded arrays.  Write
$\lambda_0=\overline\lambda_n$.  Their complete leading changes
are respectively $0,-8\lambda_0,4\lambda_0,0$: substituting
$S=z=q=u=v=0$ in the bounded coquotients (TU9) gives
$\delta a/q=-4\lambda_0$, $\delta b_*/q=0$,
$\delta c_*/q=-4\lambda_0$, and
$\delta(R^2P_L)/q=4\lambda_0$.  Thus before the extra current
the four residues are
\[
 \overline{a_1-c_1}=2,\qquad
 \overline{b_1+2c_1}=-8-8\lambda_0,\qquad
 \overline{R^2P_L/q}=2+4\lambda_0,\qquad
 \overline{4\alpha_bD_u/q^2}=2.
                                                               \tag{TU21}
\]
Equation (TU17) adds $-4\lambda_0$ to the $B_*^2$ residue
and $+8\lambda_0$ to the $\Gamma^2$ residue.  Consequently
the \emph{combined} full scalar has exactly the original four
residues $(2,-8,2,2)$ on $(A_*^2,\Gamma^2,B_*^2,D_*^2)$.
All differences from those residues are infinitesimal bounded
arrays, because $\lambda_n$ itself is bounded.  This cancellation
is why a thinness assumption is unnecessary on (TU1).
The one additional bound QI uses from physical $\beta\geq0$ is
also available directly, without any sign for $\beta^0$:
\[
  w/(c/g)=\sqrt{w/\tau}=O(1),\qquad w=O(1).
                                                               \tag{TU22}
\]
Thus the exact dangerous cross term $-4(wv/(c/g))D_uB_*C_*$
still has infinitesimal coefficient.

For the complete first $N$ normal, (TU13) gives exactly the two
bounds used in the uniform QI normal lemma:
$g/\sqrt{H_0}$ and $c/\sqrt{H_0}$ are bounded, and
\[
 \frac{h_bD_v-h_2}{\sqrt{qH_0}}=O(z+v^2).
                                                               \tag{TU23}
\]
The $H$ term and all other terms in SC14 therefore have the same
infinitesimal coefficients on the QI rows; (TU18) handles the whole
extra term.  For the harmonic normal, (TU20) adds a
$q\lambda_n$ bounded array to the left side of QI10.  It can be
absorbed into its existing $qK_q$ coquotient.  All other terms are
the same bounded rational arrays multiplied by the opening-safe
quantities in (TU12), (TU19).  The exact $D_b$ root and its two
normal terms in QI9 remain unchanged as algebraic identities.
The complete hard block therefore retains the real residue
$\operatorname{diag}(1,4)$ in $(\widehat N,N_y)$.

For clarity, here is the finite allocation being invoked.  Use
the QI roots
\[
 E_*=h_t-v,\quad F_*=vh_t-gD_*,\quad
 R_F=\frac{\sqrt{g\Gamma}}c(\Gamma-D_*),\quad
 R_E=\sqrt\Gamma(\sqrt{A_*}-\sqrt{cv}),
\]
together with the complete finite toric root list SC16.
The exact identities QI7--QI8 give strict real square pivots on
$F_*,R_F,R_E$; QI6 gives the strict $E_*$ pivot and twice the
four-monomial circuit.  Every remaining coefficient is an
infinitesimal bounded rational array in the testing $u,v$, with
the same finite denominators and divisions just verified.
All removed positive square terms, including (TU17), vanish at (TU7).
The actual source contact value, its three scalar first jets,
and both full first-normal contact values consequently vanish
exactly for this remaining expression.  Applying the two finite
monic divisions SC17--SC20 and the displayed QI7 first-jet
identities therefore puts the entire scalar correction and
both complete first normals on this strict root list with
infinitesimal coefficients.  The coefficient functions are
independent of the state ratio $Q_b/Z$; no boundedness of that
ratio is assumed.  This is precisely the constructive QI
allocation, with its source inequalities checked in
(TU12)--(TU23), not a conclusion from a scalar restriction alone.

The resulting finite Gram is bounded and has positive definite
real residue.  The strict bounded-Gram lemma supplies tensor SOS.
Restore the original positive paired positive square terms and the new full
square in (TU17).  This proves the local assertion for the entire
parent on the positive small-angle, negative-core chart.

\subsubsection*{States, signs, and all remaining tests}

Every inverse and radical used above is pure source, pure testing,
or the explicitly separated positive monomial radical in (TU15).
For $Z\ne0$ the normalization $Z=1$ is pure testing; multiplying
back by $Z^2$ restores homogeneity.  At $Z=0$ the original
pure-testing core restriction is nonnegative.  The exact core
Schur identity and sign gauges used in the retained QI proof
force $W+uY>0$ on the remaining negative-core chart, and hence
all four monomials are positive.  Otherwise the core and all
original event positive square terms already give the direct exit.  Literal
zero-angle charts have that same original exit.  Complex states
split into their real and imaginary parts, since the real
measurement matrices and full quadratic parent are unchanged.

The new R2 outer theorem R2O1--R2O24 treats every projective
testing pair outside the double-infinitesimal original chart,
with no condition on $a$ or $n/e^2$.  Its hypotheses are exactly
the original RER source conditions already retained here.
Combining that result with this local proof gives the claimed
all-tests, all-states theorem on (TU1).  No other source branch
is proved by this conclusion.

\subsection{Unbounded source ratio and relative current}
\label{component:pef021eabacde}

\label{pef021eabacde:sec:r2-unbounded-current-new}

\paragraph{Precise assertion.}
Retain the original RER source, all its strict gaps, and the notation
TU2--TU16 of the preceding bounded-relative-current proof.
Assume the native $a,t$ are positive unbounded and
\[
             \lambda=\frac{n}{tB}=\frac{n}{qz}
                   \text{ is positive unbounded}.                 \tag{UL1}
\]
Then the whole original weak-C2 parent is tensor SOS for every
independent projective test and every state. The argument below
complements TU1; it never assigns a sign to the formal C1 coefficient
$\beta^0$, and never replaces the original physical source by that
formal source.

\subsubsection*{Source bounds without a thinness assumption}

The exact RER physicality equation already used in TU4 is
\[
 Q_{\rm RER}=\tau\{A+2r(z\Delta)\lambda\},\qquad
 \tau=\frac q{zw},\qquad
 A=(1+z)(1+Sh)(1+Sq).
\]
Here $A,Q_{\rm RER},r,z\Delta$ are positive units with residue one.
Since $\lambda$ is unbounded, this proves, with no cancellation of
mixed factors,
\[
 \tau\in\mathrm I,\qquad \lambda\tau\sim\frac12,\qquad
 \lambda q\sim\frac{zw}{2},\qquad
 \frac n{z^2w}\sim\frac12.                                     \tag{UL2}
\]
All four divisions are pure source. Recall $0<w<R\sim1$.
The exact TU9 formulas therefore imply
\[
\begin{gathered}
 (1-h)\delta c_b,\ h\delta a_b,\ \delta D_b,\
            \delta h_2,\delta h_b=\lambda q\,O(1),\\
 \delta\alpha_b/q^2=\lambda q\,O(S+zw),\qquad
 h_2\asymp zw,\quad h_b=zw/2,\\
 H_0=4h_2/R+2z\asymp z,\qquad H_0\geq2z,\\
 \frac g{\sqrt{H_0}}=O(1),\quad
 \frac c{\sqrt{H_0}}=O(1),\quad
 \frac{\lambda q}{H_0}=O(1),\quad \frac n{z^2}=O(1).
\end{gathered}                                                   \tag{UL3}
\]
In particular $\lambda q$ is infinitesimal. The displayed coefficient
arrays have denominators only among the pure-source units
$P,J,Q,R$ after the original $h,1-h$ openings cancel.
The actual and formal $\sigma=c_b+a_b$ differ by a relative
infinitesimal. Their square-root ratio is a bounded unit.
No lower bound on $w,h,1-h$ or the strict gap $M$ is assumed.

Use the complete generic SC3 identity with the actual coefficients.
Retain its positive $m$ positive square term and use $\widehat N=\sqrt{H_0}N$.
The formal comparison is used only for finite scalar and normal
identities. Its $q$-coquotients and the QI coefficient arrays are
bounded here: their derivation requires $S,z,q$ infinitesimal and
$w$ bounded, but not $\beta^0\geq0$. The sole QI bound that is no
longer available, $w/(c/g)=O(1)$, is replaced explicitly below.

\subsubsection*{The entire extra current supplies a new angular pivot}

Use the exact positive-chart rows TU15--TU16, denoted here by
$A_*,B_*,C_*,D_*,U_*$ and $\Gamma$. Define
\[
 J_3=\sqrt\lambda(U_*-B_*),\qquad
 E_*=h_t-v,\qquad F_*=vh_t-gD_*.
                                                               \tag{UL4}
\]
The first line of the exact whole-current identity TU17 is
\[
 4D_uD_v\left(J_3-d_N\widehat N\right)^2,\qquad
 d_N=\frac{\sqrt\lambda\,cz}{\sqrt{H_0}}
       =\sqrt{\frac{nz}{H_0}}\in\mathrm I.                       \tag{UL5}
\]
Thus it supplies a strict real pivot on $J_3$, with only
infinitesimal changes to the hard normal Gram. Its scalar
remainder is still exactly
$8\lambda D_uD_v(1-e)\Gamma^2
-4\lambda D_uD_v(1-e^2)B_*^2$.
Its complete first-normal remainder, retained below, is
$-8\lambda D_uD_v(1-e)czB_*N$.

Write $\widetilde c,\widetilde a,\widetilde d,\widetilde\alpha,
\widetilde m,\widetilde b$ for the six differences
$\delta c_b,\delta a_b,\delta D_b,\delta\alpha_b,
\delta h_2,\delta h_b$ divided by $\lambda q$.
They are defined explicitly by the six right sides in TU9, using
$n=\lambda qz$; their apparent opening poles cancel in every
combination below. Put $U=u^2,V=v^2$, and retain
$\rho_0=Kh/R$, $a_t=1-\rho_0+(w/R)V$.
Define the exact rational arrays
\[
\begin{split}
\mathsf a={}&4D_vR^2\{\widetilde c(1-\rho_0)^2
                      +\widetilde a\rho_0^2/D_u\}
 +4\widetilde\alpha zS^2VR^2\rho_0^2/D_u
 +4\widetilde d S^2R^2a_t^2,\\
\mathsf b={}&-8\widetilde\alpha zSR\rho_0
                         -8\widetilde d SRa_tD_u,\\
\mathsf c={}&4\widetilde dD_u^2,\\
\mathsf l={}&4R\{\widetilde m(1-w/R)
                              +\widetilde b(w/R)D_v\}.
\end{split}                                                     \tag{UL6}
\]
For example $\delta a/q=\lambda\mathsf a$ in QI2,
and $\delta(R^2P_L/q)=\lambda\mathsf l$.
Direct specialization of TU9 gives, keeping $S,z,w,h$ exact,
\[
\begin{split}
 (\mathsf a-\mathsf c)|_{q=U=V=0}&=0,\\
 (\mathsf b+2\mathsf c+8D_uD_v(1-e))|_{q=U=V=0}&=0,\\
 (\mathsf l-4D_uD_v(1-e^2))|_{q=U=V=0}&=0.
\end{split}                                                     \tag{UL7}
\]
These are stronger than residue equalities. Clear the unit
denominators and divide successively by the monic variables
$q,U,V$. Each of the three expressions is consequently
$qF_q+UF_U+VF_V$, with finite bounded rational coefficient arrays.
The division does not invert $q,U,V$ in the tensor ring.

The combined full scalar correction to QI6 therefore consists of
bounded-array multiples of $\lambda q,\lambda U,\lambda V$
on $A_*^2,\Gamma^2,B_*^2$, an infinitesimal coefficient on
$E_*^2$, and the infinitesimal $\alpha$ correction on $D_*^2$.
The following exact identities allocate every apparently unbounded
testing term:
\[
\begin{aligned}
\lambda U A_*^2&=\lambda q(uE_*+bC_*)^2,\\
\lambda V A_*^2&=\lambda q(F_*+gD_*)^2,\\
\lambda U\Gamma^2&=\lambda q\,bC_*(uE_*+bC_*),\\
\lambda V\Gamma^2&=\lambda q(F_*+gD_*)
                                    (F_*+gD_*-vE_*),\\
\sqrt\lambda\,uB_*&=\sqrt{nS/z^2}\,(F_*+gD_*),\\
\sqrt\lambda\,vB_*&=\sqrt n\,C_*-vJ_3 .
\end{aligned}                                                   \tag{UL8}
\]
Every coefficient introduced on squaring the last two identities
is infinitesimal by UL2--UL3 and $S,u,v\in\mathrm I$.
Finally the formerly dangerous QI6 term has the exact replacement
\[
 \frac{wg}{c}\,vB_*C_*
       =wz\,C_*^2-\frac{wg}{c\sqrt\lambda}\,vC_*J_3,
 \qquad
 \frac{wg}{c\sqrt\lambda}
          =w\sqrt{\frac{z^2}{n}}=O(1).                          \tag{UL9}
\]
Both resulting coefficients are infinitesimal. This retains the
entire term for arbitrary independent source and testing rates.

\subsubsection*{Both unrestricted state-normal corrections}

The new complete first $N$ coefficient, including the remainder
of TU17, differs from the formal generic SC coefficient by
\[
\begin{split}
\delta\mathcal L_N={}&
 8\lambda c(1-e)\mathsf F_N B_*
 +8\lambda q\widetilde d\,ubvg(Sa_tH-D_uv),\\
\mathsf F_N={}&
 \frac{P_bD_v-(1+qz)P_m}{P}-zD_uD_v
                 =V-zUD_v+qF_{Nq}.
\end{split}                                                     \tag{UL10}
\]
The last equality is a finite monic $q$ division; its coquotient
has bounded arrays because the numerator at $q=0$ is exactly
$V-zUD_v$. After division by $\sqrt{H_0}$, its $q$ part is
infinitesimal on $B_*$ by UL3. The other parts use
\[
\begin{split}
\lambda cV B_*&=\lambda qg\,vC_*
                            -\sqrt{\lambda q}\,v^2J_3,\\
\lambda czU B_*&=\lambda qbg\,v(uE_*+bC_*).
\end{split}                                                     \tag{UL11}
\]
Their coefficients on $\widehat N$ are infinitesimal since
$g/\sqrt{H_0}$ and $\sqrt{\lambda q/H_0}$ are bounded.
The $\widetilde d$ term in UL10 uses
$uH=R(uE_*+bC_*)$ and $uv=bC_*$; it has the same
infinitesimal coefficient bounds. Thus this is a lift of the
complete first normal, with no state component put to zero.

For the second normal it is convenient to compare the coefficient
after multiplication by $\sqrt\sigma$. Differentiating the full
generic SC3 identity, including the restored $D$ positive square term, gives
\[
 \sqrt\sigma\,\delta\mathcal L_y
       =8\lambda q\{\mathsf F_HH+\mathsf F_vv\},                 \tag{UL12}
\]
where
\[
\begin{split}
\mathsf F_H={}&D_v\{\rho_0\widetilde a
                    -(1-\rho_0)\widetilde c\}
 -D_v\widetilde aU\rho_0/D_u
 +\widetilde\alpha zS^2V\rho_0/D_u
 -\widetilde dS^2a_t,\\
\mathsf F_v={}&\widetilde dSD_u-\widetilde\alpha zS.
\end{split}
\]
These arrays are bounded through both source openings. Exactly,
\[
 \mathsf F_H|_{q=U=V=0}
       =\frac{(1-e)S}{1-S(1-h)},\qquad
 \mathsf F_v|_{q=U=V=0}
       =-\frac{(1-e)S}{1-S(1-h)}.                              \tag{UL13}
\]
As $R=1$ there, finite monic division gives
$R\mathsf F_H+\mathsf F_v=qF_q+UF_U+VF_V$ with bounded
arrays. Since $H=R(E_*+v)$, the whole UL12 coefficient is
therefore an infinitesimal combination of $E_*,F_*,A_*,C_*,D_*$:
use only
\[
 qv=cA_*-qE_*,\qquad Uv=buC_*,\qquad
 Vv=v(F_*+gD_*-vE_*).                                         \tag{UL14}
\]
The factor $1/\sqrt\sigma$ is bounded.
Changing the formal harmonic-normal normalization to the actual
one multiplies its previously bounded arrays by a pure-source
unit with residue one.

For completeness, the unchanged formal first $N$ coefficient is
bounded after the actual $\sqrt{H_0}$ normalization:
$h_2^0,h_b^0=O(q)$ and $h_b^0-h_2^0=O(qz)$, while
$g/\sqrt{H_0}$ and $c/\sqrt{H_0}$ are bounded. These are exactly
the two estimates used in the QI normal allocation.
All remaining hard coefficients from the original $C,A,D,\alpha$
lines have the same residues as before; the relative $C,A$
changes and the absolute $D,\alpha$ changes are infinitesimal.
Together with UL5 the complete hard block on
$\widehat N,N_y,J_3$ has a positive definite real residue.

\subsubsection*{A finite contact lift for the additional row}

It remains to allocate the full scalar correction on roots that
vanish at contact. Merely bounding a coefficient on $A_*,B_*,C_*,D_*$
would not suffice. Keep the complete strict toric root list SC16,
the QI rows $E_*,F_*,R_F,R_E$, and add
\[
 V_3=\sqrt{\lambda\Gamma}(\sqrt{U_*}-\sqrt{B_*}).
\]
Since $U_*B_*=\Gamma^2$, the exact identities are
\[
\begin{split}
 J_3^2&=4V_3^2+\lambda(\sqrt{U_*}-\sqrt{B_*})^4,\\
 \Gamma J_3&=\frac{V_3^2}{\sqrt\lambda}
                      -2\sqrt\lambda\,\Gamma(B_*-\Gamma).
\end{split}                                                     \tag{UL15}
\]
Reserve half the $J_3^2$ coefficient in UL5 for the first
identity. It supplies strict real pivots on both $J_3,V_3$,
with the quartic square retained nonnegative. The normal
coupling in UL5 remains an infinitesimal hard Gram correction.

Here are the additional first-jet coefficient checks needed for
the second identity, so that its factor $\sqrt\lambda$ causes
no hidden loss. The only new scalar terms linear in $J_3$ in
UL8--UL9 have coefficients, up to bounded arrays, of the forms
\[
 \epsilon_B=-2\sqrt n\,v,\qquad
 \epsilon_T=\frac{wg}{c\sqrt\lambda}\,v .
                                                               \tag{UL16}
\]
Both have bounded finite testing difference arrays and are
infinitesimal. At the actual pure-source contact $v_c=cg$,
\[
 \sqrt\lambda\,\epsilon_B(u_c,v_c)=-2n\in\mathrm I,\qquad
 \sqrt\lambda\,\epsilon_T(u_c,v_c)=wz=e\in\mathrm I.             \tag{UL17}
\]
Multiplication by any of the bounded arrays produced in UL7
preserves these conclusions. No other scalar term introduces
$J_3$ linearly: the other five identities in UL8, and the
$\lambda q$ scalar corrections, use only the old QI rows.

Write a term $\epsilon J_3A_i$ as
$\epsilon J_3(A_i-\Gamma)+\epsilon\Gamma J_3$.
The first summand is already an infinitesimal quadratic on the
strict list. Divide $\epsilon-\epsilon(u_c,v_c)$ finitely by
$u-u_c,v-v_c$ and apply the exact SC18 conductors to
$\Gamma(u-u_c)$ and $\Gamma(v-v_c)$. This allocates the
nonconstant part as infinitesimal quadratic corrections involving
$J_3$. For the constant part use UL15. Its $V_3^2$
coefficient is infinitesimal, and its remaining coefficient on
$\Gamma(B_*-\Gamma)$ is infinitesimal by UL17.

There are also new bare-$\Gamma$ $E_*,F_*$ terms when the
first five identities in UL8 are expanded and their $C_*,D_*$
factors are written as $(A_i-\Gamma)+\Gamma$. Put
$L=\lambda q$ and $M_1=nS/z^2$, both infinitesimal.
The additional $E_*$ coefficients are bounded-array multiples
of $Lbu$ and $Lgv$; the additional $F_*$ coefficients are
bounded-array multiples of $Lg$ and $M_1g$.
At $u_c=bc,v_c=cg$ their required contact ratios are
\[
 \frac{Lbu_c}{c}=LS,\qquad
 \frac{Lgv_c}{c}=Lz,\qquad
 \frac{Lg}{g}=L,\qquad
 \frac{M_1g}{g}=M_1 .
                                                               \tag{UL18}
\]
All are infinitesimal, with bounded finite difference arrays
before these contact substitutions. The exact QI7 identities
therefore allocate these additional terms by the same finite
division as the old $E_*,F_*$ terms. This check is necessary:
the two new normal rows cannot be ignored when applying the
contact-jet argument.

After these explicit additional steps, the remainder has exactly
the QI form: a bounded infinitesimal quadratic on the listed
roots, the now-verified bare-$\Gamma$ $E_*,F_*$ couplings, and
a toric expression
$\Gamma\sum_i\mu_i(A_i-\Gamma)+\Gamma^2\sigma$.
The old $E_*,F_*$ contact coefficient ratios remain the bounded
ones in QI7--QI8: their unchanged $q=0$ arrays depend on bounded
$w$ and do not require $w/(c/g)$ bounded.
Every retained positive square term, including UL5 and UL15, vanishes at the
actual contact. The complete actual scalar consequently has zero
contact value and all three zero first jets. Apply the two
explicit finite divisions SC17--SC20. Their three-jet balance
makes the constant toric first-jet coefficients equal; SC19 and
the variance roots then allocate this common term. This puts
the entire scalar on a bounded Gram with strictly positive
real residue.

The full normal lifts UL10--UL14 have exact contact zero as well.
Their $J_3,E_*,F_*$ terms already vanish there. The sum of their
remaining four monomial coefficients is zero at contact, so one
finite SC18 division allocates them on the same strict list.
Every coefficient function in these steps depends only on the
testing $u,v$. No bound on the state ratio $Q_b/Z$ is used, and
the third contact jet is used only for its exact toric balance.

The strict bounded-Gram lemma now proves tensor SOS for the whole
original parent on the small-angle positive negative-core chart.
All its radicals are radicals of separated positive monomials;
every inverse is pure source or pure testing.

\subsubsection*{Projective completion and scope}

The state, sign, zero-axis and real/imaginary-state arguments are
the same exact original-core arguments at the end of TU.
In particular $Z=0$ is an original positive restriction, and the
negative-core $Z=1$ chart forces $W+uY>0$, permitting the
separated radicals above. The theorem R2O covers every testing
angle chart outside $u,v$ both infinitesimal without any
condition on $\lambda$ or the native $a$. This proves UL1 for
all tests and states.

Together with TU this covers native $a$ unbounded: the nonnegative
$\lambda=n/(tB)$ is bounded, including zero, or unbounded. The
bounded-$a$ theorem supplies the complementary source domain.

\subsection{Complementary testing angles}
\label{component:p8f8780350b13}

\label{p8f8780350b13:sec:r2-rederived-outer-join}

This proof treats every independent testing pair outside the positive
double-infinitesimal angle chart of the original R2 source.  It uses an
explicit full-parent Bell reflection and the two-current and six-root
identities that were freshly proved for R1.  Only those universal algebraic
identities and their real matrix arguments are reused.  No assertion that
the reflected source is an original R1 source is made.

\subsubsection*{Actual source and exact reflected parent}

The source is the original weak-C2 tuple, with $d=1$, $\lambda=\beta=0$,
the six coefficients in (RER1), all original strict gaps, $n\geq0$,
$a,B,t_s>0$, and $0<r,h<1$.  R2 means
\[
t_s\text{ unbounded},\qquad B, z=aB, y=t_sB\text{ infinitesimal}.
\tag{R2O1}
\]
Put $S=1/t_s$, $l=1-h$, $e=1-r$, $R_s=1+hz$,
$g_s=R_s-ae>0$, and $k=(1+ar)/(1+a)$.  The rederived exact source
lemma (RER1)--(RER7) gives
\[
k\text{ a positive bounded unit},\qquad k<1,\qquad
SC l\sim SAh\sim D/S\sim k/2,
\qquad \Delta/(1+a)\sim1.
\tag{R2O2}
\]
It also gives $ay$ infinitesimal, $Q\sim1$, and the exact expressions
\[
n=\frac{eQ-BK_s}{2r\Delta},\qquad
m=\frac{a^2e(1+y)c s_c}{2g_s\Delta},\qquad
BK_s=(1+a)y(1+hS)(1+z),
\tag{R2O3}
\]
where $c=1+zr\sim1$ and $s_c=1+B(l+ar)\sim1$.
There is no positive lower bound assumed on $r,h,l$, or $g_s$.

For arbitrary finite testing angles define the unnormalized testing
state transformation
\[
\begin{split}
G\psi=(&uvW-uX-vY+Z,\ -vW+X-uvY+uZ,\\
       &-uW-uvX+Y+vZ,\ W+vX+uY+uvZ).
\end{split}
\tag{R2O4}
\]
It has $G^TG=D_uD_v I$.  The Bell transformation on observables is
\[
(A'_0,A'_1,B'_0,B'_1)=(-B_1,-B_0,-A_1,-A_0).
\]
After the displayed real local basis gauges it sends the angles to
$(u',v')=(v,u)$ and the source coefficients to
\[
(\lambda',n',m',C',A',D',\alpha',\beta')
 =(\alpha,m,n,C,D,A,\lambda,\beta).
\tag{R2O5}
\]
Direct substitution in the whole parent gives
\[
F_{c'}(v,u;G\psi)=D_uD_v F_c(u,v;\psi).
\tag{R2O6}
\]
  The transformed
$\lambda'=\alpha$ event is kept as its whole positive square.
Because $D_uD_v>0$ is pure testing, a certificate on the left restores
the right by dividing each root by $\sqrt{D_uD_v}$ within $K_t$.
No source or testing contact is identified in this operation.

In the next two subsections write the transformed testing angles again
as $u,v$, and denote the reflected coefficients by
\[
\mathsf C=C,\quad\mathsf D=A,\quad\mathsf A=D,
\quad\mathsf m=n,\quad\mathsf n=m.
\]
Thus the symbols on the right refer to the original R2 source.
Define the pure-source quantities
\[
g_* =\frac{\mathsf C\mathsf D}{\mathsf C+\mathsf D},\qquad
h_* =\frac{\mathsf D}{\mathsf C+\mathsf D},\qquad
\gamma=\mathsf A g_*.
\tag{R2O7}
\]
If $U_C=SC l$ and $U_A=SAh$, then exactly
\[
Sg_* =\frac{U_CU_A}{hU_C+lU_A}\sim k/2,
\qquad \gamma\sim k^2/4.
\tag{R2O8}
\]
The denominator is a positive bounded unit even when $h$ or $l$ is
infinitesimal.  In particular $1/\mathsf C$, $1/\mathsf D$,
$1/g_*$ are infinitesimal,
$(\mathsf C+\mathsf D)^{-1/2}=O(\sqrt S)$, and $0<h_*<1$.
These are the complete source bounds needed for the first two outer
current identities.

\subsubsection*{Universal outer current identities}

Set
\[
R_t=uW-Y,\quad C_t=v(uX-Z),\quad B_t=X+uZ,\quad
J=R_t+C_t,\quad N=R_t+h_*C_t,
\]
and let $K$ be the exact signed quarter-core
\[
K=\tfrac12J^2+\tfrac12D_u(vY-Z)^2
 +\tfrac12D_vB_t^2-\tfrac12D_uD_vZ^2.
\tag{R2O9}
\]
All unselected original events, including $\lambda'$, are retained
as whole nonnegative positive square terms.

If $1/v$ is bounded, put
\[
K_0=2\sqrt{g_*}C_t,\quad B_0=\sqrt{2D_v/D_u}\,B_t,
\quad c_u=D_u^{-1/2},\quad t_u=uD_u^{-1/2},\quad
\eta_v=\sqrt{D_v}/v.
\]
The entire selected-parent identity is
\begin{align*}
4\mathsf C D_vR_t^2+4\mathsf D J^2+4K
={}&4(\mathsf C+\mathsf D)N^2+4\mathsf C v^2R_t^2
 +2J^2+2D_u(vY-Z)^2\\
&+(K_0,B_0)
\begin{pmatrix}
1-c_u^2\eta_v^2/(2g_*)&t_u\eta_v/\sqrt{2g_*}\\
t_u\eta_v/\sqrt{2g_*}&1
\end{pmatrix}
\binom{K_0}{B_0}.
\tag{R2O10}
\end{align*}
The finite Gram is bounded and has real residue $I$, because $c_u,t_u$
and $\eta_v$ are bounded while $1/g_*$ is infinitesimal.  The actual
separated bounded-Gram lemma makes it tensor SOS.  The roots themselves
need not be bounded.  This gives every state and every other finite angle
when transformed Bob is noninfinitesimal, equivalently when the original
Alice angle is noninfinitesimal.

For positive infinitesimal transformed $v$, put
$q=g_*^{-1/2}$ and define
\[
K_1=2\sqrt{g_*}C_t+qY,\quad
B_1=B_t-uvY,\quad V_0=vY,\quad Y_1=q\sqrt{D_u}Y,
\]
\[
L_0=\tfrac12q c_u(K_1-c_uY_1)-t_uvB_1+c_uvV_0.
\]
Then $L_0=\sqrt{D_u}\,v(vY-Z)$ exactly, and
\begin{align*}
&4\mathsf C D_vR_t^2+4\mathsf D J^2
       +4\mathsf A D_uD_vY^2+4K\\
&\quad=4(\mathsf C+\mathsf D)N^2+4\mathsf C v^2R_t^2
 +2J^2+\mathcal G,\\
\mathcal G={}&K_1^2+2D_v(B_1^2+V_0^2)
 +(4\gamma D_v-c_u^2)Y_1^2\\
&+2qv(K_1-c_uY_1)V_0-2L_0^2.
\tag{R2O11}
\end{align*}
This is the complete four-root Gram, including all normal terms.  Its
residue in the ordered frame $(K_1,B_1,V_0,Y_1)$ is
\[
\operatorname{diag}(1,2,2,k_0^2-c_{u,0}^2).
\tag{R2O12}
\]
Subscript $0$ denotes real residue.  Hence it is strict whenever its
last entry is positive.  This includes every unbounded transformed
$u$ and every bounded positive-unit transformed $u$ when $k_0=1$.
All inverse and square-root factors in these identities belong to one
field at a time.

\subsubsection*{The remaining unit angle and both coefficient openings}

It remains to treat transformed $v$ infinitesimal, transformed $u$ a
bounded positive unit, and $k_0<1$.  The source identity
\[
1-k=\frac{ae}{1+a},\qquad 0<ae<R_s\sim1
\tag{R2O13}
\]
shows that $a,e$ are positive bounded units: if $a$ were infinitesimal
or unbounded, or if $e$ were infinitesimal while $a$ is bounded, then
$1-k$ would be infinitesimal.  Thus in the current chart $a$ is a
positive unit, $e$ is a positive unit, and both $1+a$ and $\Delta$
are units.

Use the opening-safe source coordinates
\[
j=R_s/a,\qquad p=e,\qquad j-p=g_s/a>0.
\tag{R2O14}
\]
They satisfy $0<p<\min(1,j)$, $j,p$ are positive bounded units, and
the residue of $(j+1-p)/(j+1)$ equals $k_0$.  The exact relative
critical coefficient identities are
\[
\frac{\mathsf m}{jp/[2(1-p)(j+1)]}
 =\left(Q-\frac{BK_s}{e}\right)\frac{R_s+a}{R_s\Delta},
\qquad
\frac{\mathsf n}{p/[2(j-p)(j+1)]}
 =(1+y)c s_c\frac{R_s+a}{\Delta}.
\tag{R2O15}
\]
Their right sides have residue one.  Indeed $BK_s/e$ is infinitesimal
here by (R2O3), and all other normalized factors have the limits stated
in (RER).  No division by $r$, $g_s$, $n$, or an infinitesimal gap
was used to assert a bounded error.  The displayed quotients cancel
those openings exactly before the limit is taken.  In particular the
critical chart itself forces original $n>0$, even though the theorem
as a whole retains $n=0$.

Choose positive real caps $m_c\leq\mathsf m$ and
$n_c\leq\mathsf n$, with the choices specified below.  Keep both
nonnegative excesses as the original whole positive square terms.  Define six rows
\[
\widehat N=\sqrt{\mathsf C+\mathsf D}\,N,\quad
P=C_t/\sqrt S,\quad Y_0=\sqrt S\,Y,
\quad X,\quad Z,\quad V=vY.
\tag{R2O16}
\]
They can be unbounded or linearly dependent.  If
$\delta=(\mathsf C+\mathsf D)^{-1/2}$, then
$R_t=\delta\widehat N-h_*\sqrt S P$.  The selected entire
quarter-parent is exactly
\begin{align*}
\mathcal P_{\rm sel}={}&\widehat N^2+\mathsf C v^2R_t^2
 +\tfrac12J^2+(Sg_*)P^2
 +(\mathsf A/S)D_uD_vY_0^2\\
&+m_c\{(v/u)R_t+(D_u/u)V-X-uZ\}^2
 +n_cD_uD_vX^2\\
&+\tfrac12D_u(V-Z)^2+\tfrac12D_v(X+uZ)^2
 -\tfrac12D_uD_vZ^2.
\tag{R2O17}
\end{align*}
The full-state row relation is
\[
PY_0=V(uX-Z).
\tag{R2O18}
\]
Add $2b(PY_0-V(uX-Z))=0$ for a real $b$.  Keep the two displayed
positive squares $\mathsf C v^2R_t^2$ and $J^2/2$ whole.  The
remaining Gram is bounded: $u,1/u$ are bounded, $Sg_*$ and
$\mathsf A/S$ are units, and the $R_t$ contributions in the
$m_c$ square have coefficients $\delta v/u$ and
$h_*\sqrt S v/u$, both infinitesimal.  Its residue is the direct sum
of a unit pivot, the matrix
\[
\begin{pmatrix}k_0/2&b\\b&k_0(1+u_0^2)/2\end{pmatrix},
\tag{R2O19}
\]
and the following real quadratic form on $(X,Z,V)$:
\[
Q_3=m_c\{(D_u/u)V-X-uZ\}^2+n_cD_uX^2
 +\tfrac12D_uV^2-D_uVZ+\tfrac12(X+uZ)^2
 -2bV(uX-Z).
\tag{R2O20}
\]
In this formula $u=u_0>0$ is real.  We now give all finite-cap choices;
no opening of either source coefficient is suppressed.

If $p_0<1$ and $j_0>p_0$, use initially
\[
m_0=\frac{j_0p_0}{2(1-p_0)(j_0+1)},\quad
n_0=\frac{p_0}{2(j_0-p_0)(j_0+1)},\quad b=k_0/2.
\]
The three leading principal minors of $Q_3$ are
\[
m_0+n_0D_u+\tfrac12,\quad
\frac{p_0u_0^2D_u(j_0+1-p_0)}
 {4(j_0-p_0)(j_0+1)^2(1-p_0)},\quad
\frac{p_0u_0^2D_u(j_0+1-p_0)^2}
 {8(j_0-p_0)(j_0+1)^3},
\tag{R2O21}
\]
all strictly positive.  Matrix (R2O19) has determinant
$k_0^2u_0^2/4>0$.  Choose real caps slightly below $m_0,n_0$;
openness of real positive definiteness and (R2O15) ensure they remain
strict and are below the actual coefficients.

If $j_0=p_0<1$, the coefficient $\mathsf n$ is unbounded and
$m_0=j_0^2/[2(1-j_0)(j_0+1)]$ is finite positive.  Take
$b=1/[2(j_0+1)]$.  On the kernel $X=0$ of the $\mathsf n$ positive square term,
the first diagonal and determinant of $Q_3$ on $(Z,V)$ are
\[
u_0^2(m_0+1/2)>0,\qquad
u_0^2/[4(j_0+1)^2]>0.
\tag{R2O22}
\]
Choose $m_c<m_0$ sufficiently close and then a sufficiently large
real $n_c$.  A real block Schur complement makes the full matrix strict.

If $p_0=1<j_0$, the coefficient $\mathsf m$ is unbounded and
$n_0=1/[2(j_0^2-1)]$ is finite positive.  Choose a real $b$ with
\[
\frac{j_0}{2(j_0+1)}<b<
\min\left\{\frac{k_0\sqrt{D_u}}2,
                 \frac{j_0}{2(j_0-1)}\right\}.
\]
The interval is nonempty because $u_0>0$.  On the kernel
$(D_u/u)V-X-uZ=0$ of the $\mathsf m$ positive square term, the first diagonal
and determinant in $(X,V)$ are
\[
\frac{D_u}{2(j_0^2-1)}>0,\qquad
-\frac{D_u^2[2b(j_0-1)-j_0][2b(j_0+1)-j_0]}
       {4u_0^2(j_0^2-1)}>0.
\tag{R2O23}
\]
Choose $n_c<n_0$ sufficiently close, then a sufficiently large real
$m_c$.  Matrix (R2O19) and the full $Q_3$ are strict.

If $p_0=j_0=1$, both coefficients are unbounded and $k_0=1/2$.
Choose $1/4<b<\sqrt{D_u}/4$.  On the common event kernel
$X=0$, $Z=D_uV/u_0^2$, the remaining form is
\[
\frac{D_u(4b-1)}{2u_0^2}V^2>0.
\tag{R2O24}
\]
Sufficiently large real caps on both events make the full form strict.
The finite-cap step in each case is the elementary real fact that a
quadratic form positive definite on the kernel of a nonnegative event
Gram becomes positive definite after a sufficiently large multiple of
that event Gram is added; a block Schur complement proves it.

These four cases exhaust $0<p_0\leq\min(1,j_0)$ with positive $j_0,p_0$.
The separated bounded-Gram lemma now proves the entire selected parent
(R2O17), with every state normal retained.  Restore the original excesses
and all unused positive square terms.  The real caps are used only on this testing
chart: no globally physical capped source is asserted, so this argument
does not conflict with the finite-contact cap exclusion (CAP15).

\subsubsection*{Exact outer and projective join}

Return to the original testing angles.  The full-parent sign gauges
(SJ3) permit $u,v\geq0$ without changing source coefficients.  Literal
zeros and projective infinities are the exact signed-core squares
(SJ10), plus the untouched event squares.  For finite positive angles:
\begin{itemize}
\item if the original $u$ is noninfinitesimal, use (R2O10);
\item if the original $u$ is infinitesimal and $v$ is unbounded,
      use (R2O11)--(R2O12);
\item if the original $u$ is infinitesimal and $v$ is a positive
      bounded unit, use (R2O11) when $k_0=1$, and the complete
      six-root argument (R2O13)--(R2O24) when $k_0<1$.
\end{itemize}
These are exactly all finite positive angle pairs outside the case in
which both are infinitesimal.  Every argument is a whole-parent state
Gram, so no additional state chart or state-ratio restriction is present.
Equation (R2O6) restores the original frame and positive pure-testing
factor.  The zero state and the real/imaginary split for complex states
are immediate.  No source-to-testing rate comparison was made.

\section{Two small actual angles: the R3 and R4 sectors}
\subsection{All ratios at infinitesimal $e$}
\label{component:p1de8b5ffab69}

\label{p1de8b5ffab69:sec:new-r34-all-tests}

\paragraph{Statement and exact domain.}
For the present source partition it suffices to impose additionally $a$ unbounded; the identities below remain valid on this smaller domain.

Use the original common-II stationary source, its six native weights
$n,m,C,A,D,\alpha$, all its strict source gaps, $d=1$,
$\lambda=\beta=0$, and $n\geq0$. Suppose $B_s,a_sB_s,t_sB_s$
are positive infinitesimals, $t_s,h$ are positive bounded units,
and $e=1-r$ is positive infinitesimal. In addition impose either
\[
 \begin{array}{ll}
 \mathrm{R3}:&1-h\text{ is infinitesimal},\\
 \mathrm{R4}:&1-h\text{ is a positive unit and }
       1+a_sB_sr-(1-h)/t_s\text{ is infinitesimal}.
 \end{array}                                                    \tag{EN1}
\]
There is no condition on the size of $a_s$, $a_s^2B_s$,
$a_s e$, or their ratios. The conclusion is finite separated tensor
SOS for the entire original parent at every independent projective
testing pair and every state, including signs, axes, and literal $n=0$.
The bounded-$a$ parts of R3 and R4 are covered by Proposition~\ref{prop:c2-bounded-a}; the present calculation supplies their unbounded-$a$ parts.

The source and testing fields are independent real closed extensions
$K_c,K_t\supset\mathbb R$. A bounded coefficient means an element
of $B_c\otimes B_t$; its residue is the homomorphism to $\mathbb R$.
The strict bounded-Gram lemma is used only for matrices with positive
definite real residue. Row entries may be arbitrarily large.
Every inverse and radical below is pure source, pure testing, or a
product of such factors.

\subsubsection*{One exact compact source for both ends}

Use the exact compact variables of the exact actual-source
dictionary RSR2--RSR6:
\[
\begin{aligned}
c_s&=1+a_sB_sr,\quad R_s=1+ha_sB_s,\quad
j=a_sc_s/R_s,\quad B=B_sR_s/c_s,\quad T=t_sc_s/R_s,\\
p&=a_s e/R_s,\quad H=(1-h)/(t_sc_s),\qquad
x=jB=a_sB_s,\quad y=TB=t_sB_s,\quad q_z=jT .
\end{aligned}                                                    \tag{EN2}
\]
Thus $0<p<\min(1,j)$, $0<H<1$, and
\[
B,x,y\in\mathrm I_c,\quad T\in\mathrm U_c,\quad
1-HT\in\mathrm U_c,\quad p/j\in\mathrm I_c.              \tag{EN3}
\]
The last assertion is exact source entry:
$e=p(1+x)/(j+px)\sim p/j$.
In R3, $H$ is infinitesimal. In R4, $1-H$ is infinitesimal
and $0<\overline T<1$. In particular
$H(1-H)\in\mathrm I_c$ in both rows.
The native inverse is
\[
r=\frac{j-p}{j+px},\quad h=\frac{1-HT}{1+HTx},\quad
a_s=\frac{j+px}{1+HTx},\quad
B_s=\frac{B(1+HTx)}{1+pB},\quad
t_s=\frac{T(1+pB)}{1+HTx}.                              \tag{EN4}
\]
The exact six weights are RSR6 at this substitution.
These rational identities, including the contact first jets, hold at the displayed substitution.

Put $\sigma=1+(1-H)T$, $K=j+1-p$, $E=1+B(HT+p-1)$, and
use the explicit polynomials $\mathcal P,\mathcal L,\mathcal N,
\mathcal F,\mathcal K$ in RSR5, RSR15, RSR18. Then
\[
\frac{\mathcal P}{j+1}\sim1,\qquad
\frac{\mathcal L}{j+1}\sim\sigma,\qquad
\frac{\mathcal F}{j+1}\sim\sigma.                       \tag{EN5}
\]
Here $T,\sigma,(1-HT)^{-1}$ are bounded positive units.
To prove these uniform statements, divide the printed finite
polynomials by $j+1$. The corrections in $\mathcal P$ contain
$B+x$ or $TB$; those in $\mathcal L$ after
$(j+1)\sigma$ contain $p/(j+1)$ or $B+x$.
In $\mathcal F$ the only additional quotient is
$p(1-p)/(j(j+1))$, infinitesimal by (EN3).
Every other term contains $B$ or $Bj/(j+1)$.
Thus no restriction on $j$ or $jT$ is needed.

The exact physicality numerator has the grouping
\[
\mathcal N=p\{1+B(3j+1-p)+BT(j+1)\}
-Bj(j+1)\{1+(1-H)T\}+\mathcal R,\quad
|\mathcal R|\le C_0BT\,Bj(j+1).                         \tag{EN6}
\]
The bound follows termwise from RSR5 using $p<\min(1,j)$ and
bounded $T$. Since $n=\mathcal N/[2(j-p)\mathcal P]\ge0$,
\[
\begin{aligned}
&\eta_*=\frac{Bj(j+1)}p=O(1),\qquad
0\le n=O\!\left(\frac{p}{j(j+1)}\right)\in\mathrm I_c,\\
&\eta_*\in\mathrm I_c\Longrightarrow
n\sim\frac{p}{2(j-p)(j+1)}>0 .
\end{aligned}                                                    \tag{EN7}
\]
This includes $\mathcal N=0$ literally.

Write $\mu=m(1-p)$, $h_0=\mu+1/2$. The full source formulas give
\[
\begin{aligned}
&\mu=O(p),\quad h_0\asymp1,\quad
\alpha/[B^2(j+1)]\sim\sigma/2,\quad \alpha/B\in\mathrm I_c,\\
&A\sim\frac{T}{2(1-HT)},\qquad
CH\sim\frac{\sigma}{2T},\qquad
D(1-H)\sim\frac1{2T}.
\end{aligned}                                                    \tag{EN8}
\]
Consequently $A$ and the finite one of $C,D$ have positive real
lower bounds. The other coefficient is unbounded. The exact
normalized source limits are
\[
\begin{array}{c|c|c|c}
 & C&A&D\\ \hline
R3&+\infty&T/2&1/(2T)\\
R4&1/(2T)&T/[2(1-T)]&+\infty .
\end{array}                                                    \tag{EN9}
\]
Only real residues of the bounded entries are meant in this table.
The coefficient $m$ may be infinitesimal, a unit, or unbounded.

\subsubsection*{The scalar allocation in the complete two-normal frame}

Use the entire parent RFN4 and the exact full-state change RFN5--RFN9,
with $w=(W+uY)/(1+u^2)=1$ and arbitrary $z=(Z-uX)/(1+u^2)>0$.
The identities are rational identities of (EN4), so they retain
both complete normals and every state term in the present source.
The tests in this subsection satisfy $u,v>0$ infinitesimal.
For clarity, the inverse is
\[
X=N_0+pv-uz,\quad Z=z+uN_0+puv,\quad
W=1+u(M_0+Hvz),\quad Y=u-M_0-Hvz.                       \tag{EN10}
\]
Retain the actual positive square term $mpD_u^2N_0^2$, and divide the
remaining entire parent by the pure testing factor $D_u=1+u^2$.
Its exact expansion is RFN7--RFN9. In particular the hard entries are
\[
\begin{aligned}
&H_N=h_0D_u+\tfrac12v^2+\alpha u^2+nD_v,\\
&H_M=D_u(CD_v+D+\tfrac12)+AD_v+(\alpha+\tfrac12)v^2,\qquad
H_{NM}=(\alpha+\tfrac12)uv .
\end{aligned}                                                    \tag{EN11}
\]

Keep the separated rows
$U=u/\sqrt{x}$, $V=v/\sqrt{y}$, $Z_*=z/\sqrt{q_z}$,
$Q_*=U-VZ_*$, $B_*=V-UZ_*$, $R_*=Z_*-UV$, and
\[
\Psi=U^2+V^2+Z_*^2+U^2V^2Z_*^2-4UVZ_*
=(U-V)^2+Z_*^2(1-UV)^2+2UV(1-Z_*)^2.                 \tag{EN12}
\]
The last radical is separated, for example
$\sqrt{UV}=(xy)^{-1/4}\sqrt{uv}$.
The exact scalar identity is RFN12 with the same actual circuit
masses. From (EN5), the closed formulas RSR16--RSR20 now give
\[
\begin{aligned}
&f_c\sim Tx^2\sigma/2,\quad
q_c\sim\frac{Tx}{2(1-HT)},\quad
r_c\sim TjB^2(1-p)\sigma/2,\\
&g_c/f_c=Hy\,\mathcal L/\mathcal F\in\mathrm I_c,\quad
k_c<0,\quad |k_c|/(q_cy)=O(p+B).
\end{aligned}                                                    \tag{EN13}
\]
In particular all three displayed positive masses are positive.
The sign of $k_c$ follows from the same exact positive-surplus
identity RSR19: its error bound uses $B(j+1)\in\mathrm I_c$,
not $T$ infinitesimal.

Before spending any current mass, perform the exact source-only
completion
\[
\begin{aligned}
h_0N_0^2+n(N_0+jb)^2
&=(n+h_0)(N_0+c_bb)^2+M_eb^2,\\
b&=v-uz/j,\quad c_b=\frac{nj}{n+h_0},\quad
M_e=\frac{nh_0j^2}{n+h_0}.
\end{aligned}                                                    \tag{EN14}
\]
Keep $nv^2(N_0+pb)^2$ and the positive extra hard terms RFN15
as whole squares. This reserves exactly the correlated $n$ positive square terms
of the complete parent. If $n=0$, set $c_b=M_e=0$ and omit that
row. For $n>0$, $c_b^2/M_e=n/[h_0(n+h_0)]$ is bounded.

The signed scalar allocation differs from R1 in one source split.
The $e_0$ cost is unchanged: its negative part is at most $g_c$,
and RFN18b provides it at cost
$H(\mathcal L/\mathcal F)(y+v^2)\in\mathrm I$.
Also $d_c>0$. Use $Z_*(V^2-1)=-VQ_*-R_*$ for $k_c$.
The $VQ_*$ cost is $O((p+B)v^2)$, infinitesimal.
For its opposite-row part the exhaustive cases are:
\[
\begin{array}{c|l}
p\in\mathrm I_c&
 |k_c|/r_c=
 T\mathcal K/[(1-p)(1+Hy)\mathcal L]=O(p+B)\in\mathrm I_c\\
p\in\mathrm U_c,\ \eta_*\in\mathrm U_c&
 j\text{ unbounded},\ j^2B\in\mathrm U_c,\ 
 |k_c|/f_c\in\mathrm I_c,\ |k_c|/(f_cx)=O(1)\\
p\in\mathrm U_c,\ \eta_*\in\mathrm I_c&
 j\text{ unbounded},\ M_e\in\mathrm U_c,\ 
 |k_c|/f_c\in\mathrm I_c,\\
& |k_c|/(q_cy)=O(1),\ |k_c|/(M_eyx)=O(1).
\end{array}                                                    \tag{EN15}
\]
In the second case $j$ is unbounded because $p/j$ is infinitesimal;
$j^2B\asymp p\eta_*$. In the third,
$nj^2\sim p/2$ and $n$ is infinitesimal, so $M_e\sim p/2$.
The closed formulas show
$|k_c|/f_c=O((p+B)/j)$ in the last two cases and
$|k_c|/(M_eyx)=O((p+B)(1+1/n)/j^2)$ in the third.
These bounds are uniform through $p\to1$.

Case two uses RFN18a:
$R_*^2\preceq_{\rm sos}(1+U^2)\Psi$.
Case three uses RFN18c:
$R_*^2\preceq_{\rm sos}2(\Psi+V^2Q_*^2+U^2B_*^2)$.
The non-infinitesimal source ratios in (EN15) acquire the testing
factors $u^2$ or $v^2$. Thus all costs are infinitesimal.
Literal $n=0$ lies in the first or second case; it is not discarded.
After this finite allocation, fixed positive fractions of
\[
q_cQ_*^2,\quad r_cR_*^2,\quad f_c\Psi,\quad M_eb^2       \tag{EN16}
\]
remain, with the last row absent exactly at $n=0$.

\subsubsection*{Both complete normal arrays, with the new endpoint weights}

The exact derivatives RFN21--RFN22, and their exact source cancellations
$S_N=S_0=S_1=S_2=0$, hold unchanged. Let
$q=vz-Tu$, $R=z-uv/B$, $R_u=v(u^2-x)$. In particular
\[
\begin{aligned}
L_N^{\rm rem}
&=-j\alpha b+\tfrac12pv^2b+K_RR_u+(K_c/j)uvq,\\
L_M&=C_qq+\mu uvb+v(\alpha-BTCHu^2)R,\\
K_c&=p/2+j\{H(\alpha+1/2)-1/2\}+n(p-j),\\
K_R&=-\mu-(\alpha+1/2)(1-p)+(T/j)K_c,\\
K_d&=CH-(D+1/2)(1-H)=-\mu/j-BTCH,\\
C_q&=K_d+AH+1/2+
v^2\{CH(1+u^2)+AH+(\alpha+1/2)H\}.
\end{aligned}                                                    \tag{EN17}
\]
Every displayed coefficient other than a retained row is bounded:
$p/j$ and $n$ are infinitesimal, $\mu/j=O(p/(j+1))$ is bounded,
and $A,CH,D(1-H),T$ are bounded by (EN8).
In particular $K_R$, $K_c/j$, and $C_q$ are bounded.

The complete bounds for these costs are as follows; they use the actual
$T\in\mathrm I_c$ hypothesis of RFN.
The exact radial Gram RFN19--RFN20 purchases $R_u^2$ at infinitesimal
cost from (EN16). Its three cases are still $1-p$ a unit;
$p\sim1,\eta_*$ a unit; and $p\sim1,\eta_*$ infinitesimal.
Their costs are respectively
\[
\begin{array}{c|l}
1-p\in\mathrm U_c&O(B),\ O(Bu^2),\ O(u^2/(1-p))\\
p\sim1,\eta_*\in\mathrm U_c&
 O(B+u^2/j+u^4/(j^2B)),\ O(Bu^2)\\
p\sim1,\eta_*\in\mathrm I_c&
 O(u^2/j),\ O(u^4/j),\
 [8u^4+2(u^2-x)^2]/M_e .
\end{array}                                                    \tag{EN18}
\]
They require only $T$ bounded above and below, because $T$ cancels
between $y=TB$ and the masses (EN13).
The small-current budgets NSR3--NSR5 likewise apply unchanged:
both $pv^2b$ and $\sqrt T\,\mu uvb$ are free.
For $\eta_*$ infinitesimal their coefficients are
$O(v^4)$ and $O(Tu^2v^2)$, using $p^2/M_e=O(1)$.
For $\eta_*$ a unit the coefficients are
$O(v^2(y+v^2))$ and $O(Tu^2(y+v^2))$, using
$p^2y/f_c=O(1)$. Each is infinitesimal.

The row $q/\sqrt T$ is strict by (EN13).
Its $N$ coupling has coefficient $(K_c/j)uv\sqrt T$ infinitesimal.
The identity $j\alpha b=-(\alpha/B)R_u-\alpha uR$ handles
the remaining $N$ term. For $1-p$ a unit the latter square costs
$O(B^2(j+1)^2u^2/(1-p))$ against $r_cR_*^2$.
For $p\sim1$, RFN18a bounds the entire $j\alpha b$ square by
$O(B^3(j+1)^2+B^2(j+1)^2u^2/j)$ of the circuit.
These are infinitesimal. A sufficiently large fixed real radial
pivot pays the bounded $K_R$ coupling, including the bounded shift
$c_b/\sqrt{M_e}$ from (EN14).

Normalize the complete second normal by
$M_0=\sqrt{TH(1-H)}\,\widehat M$. The full hard diagonal (EN11)
then has residue
\[
\overline{TH(1-H)H_M}=
\begin{cases}
(1+\overline T)/2,&R3,\\
1/2,&R4 .
\end{cases}                                                    \tag{EN19}
\]
Every term of the complete diagonal is bounded: the only divergent
factors occur in $CH$ and $D(1-H)$.
The normalized $q$ coupling is
$T\sqrt{H(1-H)}C_q$, now infinitesimal because
$H(1-H)$ is infinitesimal. This is the replacement for the
$T$-infinitesimal argument in R1.
The normalized current coupling is a bounded multiple of the
already purchased $\sqrt T\,\mu uvb$.

For the entire remaining opposite row put $J=B(j+1)+u^2$ and
$a_R=Tq_zv^2(\alpha-BTCHu^2)^2$.
Since $\alpha-BTCHu^2=O(BJ)$, its costs are
\[
\begin{aligned}
a_R/f_c&=O((T/j)v^2J^2),&
a_R/(q_cy)&=O(v^2J^2),\\
a_R/r_c&=O(Tv^2J^2/(1-p)),&
a_R/(M_eyx)&=O(T(1+1/n)v^2J^2/j^2).
\end{aligned}                                                    \tag{EN20}
\]
The direct $r_c$ cost applies for $1-p$ a unit.
For $p\sim1,\eta_*$ a unit, RFN18a applies using $j^2B$ a unit.
For $p\sim1,\eta_*$ infinitesimal, RFN18c applies using
$(1+1/n)/j^2=O(1)$. Thus the entire square, including every
higher testing degree, is free. No state ratio occurs in a cost.

Finally the complete mixed hard term is
$2(\alpha+1/2)uv\sqrt{TH(1-H)}
(\widehat N/\sqrt{n+h_0}-c_bb)\widehat M$.
Its bounded coefficients have zero residue.
The strict bounded-Gram argument therefore applies exactly as in
the final paragraph of RFN: choose a large fixed radial pivot, provide
the finitely many free rows, and retain positive fractions of
(EN16). The two hard pivots are $1$ and the positive number
(EN19). Restore all whole positive square terms in (EN14), RFN15, and
$mpD_u^2N_0^2$, and the pure testing factor $D_u$.
This proves the full local parent for all infinitesimal $u,v>0$
and arbitrary $z>0$, with both unrestricted state normals.
The exact SJ state/sign reduction supplies every state and the
literal axes in this local chart.

\subsubsection*{Every other testing angle, with weak \(m\) retained}

This outer join uses only (EN9), positivity of the six native
weights, and exact full-parent identities. Set
$\widetilde m=m/(1+m)$, a bounded positive source coefficient below
$m$, and retain the difference as the whole original $m$ square.
Keep the entire $n$ and $\alpha$ events throughout.
For the single unbounded coefficient $C$ or $D$, retain its excess
over a sufficiently large positive real cap. Other selected
coefficients have the actual bounded positive residues (EN9).

\paragraph{Finite Bob with nonzero residue and bounded Alice.}
Let $v$ have nonzero finite residue, $u$ be bounded, and use the
five full-state rows
\[
P=uW,\quad Y,\quad X,\quad Z,\quad S=\sqrt{\widetilde m}\,W.
                                                               \tag{EN21}
\]
The selected $m$ event is the entire square
$[vS-\sqrt{\widetilde m}(X-uvY+uZ)]^2$.
All other terms use $P,Y,X,Z$ with bounded coefficients.
In R3 the strong row is $P-Y$; on its kernel the selected remaining
form on $(Y,X,Z)$ is
\[
\frac T2D_uD_vY^2+\frac{1+T}{2T}v^2(uX-Z)^2+K_0 .
\]
In R4 the strong row is $P+uvX-Y-vZ$; on its kernel it is
\[
\frac{D_vv^2}{2T}(Z-uX)^2+
\frac{T}{2(1-T)}D_uD_vY^2+K_0 ,
\quad
K_0=\frac{D_u(vY-Z)^2+D_v(X+uZ)^2-D_uD_vZ^2}{2}.
\]
The leading principal minors in the order $(Y,X,Z)$ are
\[
\begin{array}{c|c|c}
 &R3&R4\\ \hline
1&
\dfrac{D_u(TD_v+v^2)}2&
\dfrac{D_u(T+v^2)}{2(1-T)}\\[4pt]
2&
\dfrac{D_u(TD_v+v^2)(Tu^2v^2+Tv^2+T+u^2v^2)}{4T}&
\dfrac{D_uD_v(T+v^2)(T+u^2v^2)}{4T(1-T)}\\[4pt]
3&
\dfrac{v^2(1+T)D_u^2D_v(Tu^2+v^2)}{8T}&
\dfrac{v^2D_u^2D_v^2(Tu^2+v^2)}{8T(1-T)} .
\end{array}                                                    \tag{EN22}
\]
All are strictly positive at the real residues, even $u=0$.
The $S$ direction has the positive coefficient $v^2$ in the
selected $m$ event. A sufficiently large real cap on the strong
row makes the entire five-row residue positive definite.
This is the real finite-cap Schur lemma: if a bounded real
quadratic form is positive on the kernel of a linear map, adding
a sufficiently large real multiple of the squared map makes it
positive definite. The unbounded actual coefficient permits
every such cap. This argument includes $m$ infinitesimal.

\paragraph{Infinitesimal Bob and Alice a finite positive unit.}
Use $R=uW-Y$, $B_t=X+uZ$, $V_Z=vZ$, and $Y$; the testing inverse
$1/u$ is bounded. The full selected parent is a bounded quadratic
on these four rows because
$v(W+uY)=(v/u)(R+D_uY)$ and
$J=R+uvB_t-D_uV_Z$.
Its residue splits off $(\widetilde m+1/2)B_t^2$ and the form
\[
CR^2+AD_uY^2+(D+1/2)(R-D_uV_Z)^2
-D_uYV_Z-\tfrac12D_uV_Z^2.                            \tag{EN23}
\]
In R3 restrict to the strong-row kernel $R=0$.
After removing a positive factor $D_u$, the matrix on $(Y,V_Z)$ is
$\left(\begin{smallmatrix}T/2&-1/2\\-1/2&
(1+T)D_u/(2T)-1/2\end{smallmatrix}\right)$;
its determinant is $(1+T)u^2/4>0$.
In R4 the strong kernel is $R=D_uV_Z$.
On $(R,Y)$ its first diagonal is $1/(2T)-1/(2D_u)>0$,
and its determinant is $u^2/[4(1-T)]>0$.
A large real cap again makes (EN23) strict.

\paragraph{Infinite Bob and bounded Alice.}
Put $w=1/v$, divide the homogeneous parent by $v^2$, and use
the five rows (EN21). Select any fixed positive real lower caps
$c_0,a_0,d_0$ below $C,A,D$, and the actual $\widetilde m$ event.
The full quadratic is bounded; its residue is
\[
[S+\sqrt{\widetilde m}\,uY]^2+
c_0(P-Y)^2+(a_0+\tfrac12)D_uY^2+
d_0(uX-Z)^2+\tfrac12D_uX^2.                           \tag{EN24}
\]
This is positive definite on the five rows, for every bounded
residue of $u$. It includes $u=0$ and $m$ infinitesimal.

\paragraph{Infinite Alice.}
Let $s=1/u$ and keep the entire $m,n,\alpha$ events.
For finite $v$ of nonzero residue, after division by $u^2$
the selected residue on $(W,Y,X,Z)$ is
\[
c_0D_vW^2+a_0D_vY^2+
(d_0+\tfrac12)(W+vX)^2+\tfrac12(vY-Z)^2,
\]
which is positive definite. For $v$ infinitesimal use instead
$(N,W,Y,V)=(sX+Z,W,Y,vX)$.
The complete signed core after division by $u^2$ is
\[
\frac12\{W+(1+s^2)V-sY-svN\}^2
+\frac12(1+s^2)v^2Y^2
-(1+s^2)Y(vN-sV)-\frac12(1+s^2)(vN-sV)^2
+\frac12D_vN^2 .
\]
Adding the selected $C,A,D$ positive square terms gives strict residue
$c_0W^2+a_0Y^2+(d_0+1/2)(W+V)^2+N^2/2$.
Every full coefficient is bounded; no $s/v$ is used.
If also Bob is infinite, divide by $u^2v^2$ and use
$(W,X,Y,S_Z,T_Z)=(W,X,Y,sZ,wZ)$.
The complete selected residue is
\[
c_0W^2+a_0Y^2+d_0(X-S_Z)^2+\tfrac12X^2
+\tfrac12(Y-T_Z)^2,                                   \tag{EN25}
\]
again positive definite. The full quadratic follows directly from
$J/(uv)=wW+X-swY-S_Z$ and the exact signed core;
its coefficients are bounded polynomials in $s,w$.

These charts exhaust all nonnegative projective angles:
both infinitesimal are the complete local theorem, and every other
pair is one of (EN21)--(EN25) or (EN23).
All strict-matrix coefficients use only bounded pure source or
pure testing factors. Restore every selected excess, the untouched
$m,n,\alpha$ positive square terms, and the homogeneous testing squares.
The exact SJ sign gauges, literal projective endpoints and
real/imaginary state split finish the asserted all-tests theorem.

\section{Two small actual angles: the R6 sector}
\subsection{Positive residue of $a^2B$}
\label{component:pb60eba754641}

\label{pb60eba754641:sec:r6-b-unit}

\paragraph{Scope.}
All notation and all physical source constraints are the original weak-C2
ones. In particular $d=1$, $\lambda=\beta=0$, $n\geq0$, and the actual
finite stationary contact, every strict gap, the total-scale condition,
and the complement of earlier cases are retained. On R6,
$t,h,l=1-h,j,r$ are positive bounded units, $B,x=aB,y=tB,g$ are
positive infinitesimals, and here $a$ is unbounded. Define
\[
 s=a^{-1},\quad w_s=ae,\quad b=a^2B,\quad e=1-r,
 \qquad B=s^2b,\quad x=sb,\quad y=ts^2b.
 \tag{RB1}
\]
The existing exact source identity gives
\[
 2r\Delta_*n/s^2=w_sQ_*-bK_*,\qquad
 Q_*\sim\Delta_*\sim1,\quad K_*\sim h+t.
 \tag{RB2}
\]
The gap $g=1-w_s+hsb$ gives $w_s\sim1$. Hence $b$ is bounded and
\[
 \rho=n/e^2\sim\frac{1-b_0(h_0+t_0)}2,\qquad
 0\leq b_0(h_0+t_0)\leq1.                         \tag{RB3}
\]
Here and below a subscript zero means the real residue; a displayed
residue formula with bounded source letters has the same meaning.
No testing variable occurs in these source estimates.

\begin{theorem}
On the entire original R6 source intersection, assume in addition that
$b=a^2B$ is a positive bounded unit. Then the complete original parent
is a finite sum of squares in the separated source--testing tensor ring
at every independent projective testing pair and every real or complex
state, with either angle sign and every literal projective endpoint.
All positive $s,g$ opening rates are independent. Literal $n=0$ is
included. The theorem makes no assertion about $b$ infinitesimal.
\end{theorem}

\paragraph{Source bounds needed by the proof.}
Put $R=1+hx$, $K=1+x$ and $T=1+x(1-e)$. The original rational source
formulas, after the pure-source substitutions (RB1), have only unit
denominators after $mg$, $n/s^2$ and $\alpha/(s^3b^2)$ are formed.
Their joint residues are
\[
 mg\sim\tfrac12,\quad
 C\sim\frac{h+t}{2l},\quad A\sim\frac t{2h},\quad
 D\sim\frac1{2(t-l)},\quad
 \frac\alpha{s^3b^2}\sim\frac{h+t}2.               \tag{RB4}
\]
Thus $m$ is unbounded, $C,A,D$ are positive units, and all four have
positive real lower bounds. There is no $s/g$ or $g/s$ assumption.
These are rational cancellations in the actual physical source;
no formal negative-beta source is imported.

\paragraph{Two contact-zero positive square terms remove the two troublesome terms.}
Write the complete parent in the normalization $P=F/4$:
\[
\begin{split}
 P={}&nD_uD_vX^2+mM^2+C D_v(uW-Y)^2+A D_uD_vY^2\\
 &+(D+\tfrac12)J^2+(\alpha+\tfrac12)D_u(vY-Z)^2
       +\tfrac12D_v(X+uZ)^2-\tfrac12D_uD_vZ^2,\\
 M={}&vW-X+uvY-uZ,\qquad J=uW+uvX-Y-vZ,
 \quad D_u=1+u^2,\quad D_v=1+v^2.
\end{split}                                                     \tag{RB5}
\]
In the finite state chart $Z=1$, use the pure-testing coordinate
$\omega=(W+uY)/D_u$ and the exact normals
\[
\begin{gathered}
 p=\frac{w_s}{R}=1-\frac gR,\quad c=\frac{l}{tT},\quad
 q_A=D_uv\omega,\quad q_B=v(1-uX),\\
 \qquad N=X+u-pq_A,\quad M_B=uW-Y-cq_B.
\end{gathered}                                    \tag{RB6}
\]
Their full inverse is
\[
\begin{gathered}
 X=N+pD_uv\omega-u,\quad q_B=v(D_u-uN-puD_uv\omega),\\
 \quad W=\omega+\frac{u(M_B+cq_B)}{D_u},\quad
 Y=u\omega-\frac{M_B+cq_B}{D_u}.
\end{gathered}                                    \tag{RB7}
\]
The only testing denominator is the positive pure-testing unit $D_u$.
Set
\[
 \delta=\frac{R}{K\sqrt{bt}},\quad E=\frac{tK}{R},\quad
 k=\sqrt x,\quad K_c=\delta^2Ek^2=\frac{sR}{K},\quad
 \gamma=\frac r{K_c}=\frac{rK}{sR}.                \tag{RB8}
\]
All of $\delta,E,b$ are positive bounded units in this theorem.
The actual contact in these coordinates is
\[
 (u,\omega,v;N,M_B)=(k,\delta k,\delta k^2E;0,0).
 \tag{RB9}
\]
For direct verification, its original state coordinates with $Z=1$ are
\[
 u_c=\sqrt{sb},\quad v_c=s\sqrt{bt},\quad
 W_c=\sqrt{s/t},\quad X_c=-r\sqrt{sb},\quad
 Y_c=hs\sqrt{b/t},\quad
 \omega_c=\sqrt{s/t}\,R/K.
 \tag{RB9a}
\]
They give $u_cW_c-Y_c=ls\sqrt{b/t}=c\,v_c(1-u_cX_c)$
and $X_c+u_c=e\sqrt{sb}=pD_{u_c}v_c\omega_c$.
Also $\gamma v_c\omega_c=r\sqrt{sb}=-X_c$.
Thus both normals and the retained $n$ row vanish exactly, before
any residue calculation. The same contact substitution also gives
$\omega_c=\delta k$, $v_c=\delta k^2E$, and $K_c=sR/K$.
For any fixed positive real $K_m$, retain the two whole squares
\[
 (m-K_m)N^2,
 \qquad nD_uD_v(X+\gamma v\omega)^2.                \tag{RB10}
\]
The first coefficient is positive because $m$ is unbounded. The second
vanishes identically when $n=0$. Both rows vanish at (RB9), with every
first jet of their squares. No inverse of $n$ is used.

The complete remaining polynomial is exactly
\[
\begin{split}
 G={}&K_mN^2-2\frac{mg}{R}q_AN+\frac{mg^2}{R^2}q_A^2\\
 &-2n\gamma D_uD_vX(v\omega)-n\gamma^2D_uD_v(v\omega)^2\\
 &+CD_v(uW-Y)^2+AD_uD_vY^2+(D+\tfrac12)J^2\\
 &+(\alpha+\tfrac12)D_u(vY-1)^2
       +\tfrac12D_v(X+u)^2-\tfrac12D_uD_v.
\end{split}                                                     \tag{RB11}
\]
This is the entire parent minus precisely (RB10), including every
first normal and every hard quadratic term. The exact cancellation
$M=(g/R)q_A-N$ removes the whole $g^{-1}$ pole before an estimate.
Moreover $n\gamma=O(s)$ and $n\gamma^2$ is bounded. All coefficients
in (RB11), after (RB7), are therefore bounded. The only nonconstant
state-normal degree is one or two.

\paragraph{Raw support before the conductor.}
Define
\[
 A_0=\omega v,\quad B_0=\delta v,\quad C_0=\delta E u\omega,
 \quad D_0=\delta^2k^3E,\quad X_0=B_0-C_0.           \tag{RB12}
\]
At $N=M_B=0$ let $V_0=1-puv\omega$. Then
\[
 q_B=D_uvV_0,\quad uW=u\omega+u^2cvV_0,\quad
 Y=u\omega-cvV_0,\quad J=-(1-c)D_uvV_0,
 \quad X+u=pD_uA_0.                               \tag{RB13}
\]
Thus every term of the scalar of (RB11) lies, with bounded coefficients,
in $(A_0,B_0,C_0)^2+(D_0^2)$. The only constant is from $\alpha$,
and $\alpha/D_0^2$ is bounded with residue $b(h+t)/2$.
For the new negative $n$ contribution use exactly
\[
\begin{split}
 -2n\gamma D_uD_vX A_0-n\gamma^2D_uD_vA_0^2
 ={}&2n\gamma D_uD_vuA_0\\
 &-(2n\gamma pD_u+n\gamma^2)D_uD_vA_0^2
 -2n\gamma D_uD_vA_0N,
\end{split}                                                     \tag{RB14}
\]
where $uA_0=B_0C_0/(\delta^2E)$. Both complete first normals lie in
$(A_0,B_0,C_0)$ with bounded coefficients. Every remaining hard entry
is bounded. These are full rational support statements, not a scalar
contact-value test. Clearing $D_u^2$ makes them polynomial; its inverse
is a pure-testing square factor when restored.

\paragraph{The contact Gram and its exact leading calculation.}
Let
\[
 Y_0=A_0-kC_0,\quad Z_0=uB_0-kC_0,\quad H_0=kC_0-D_0.
\]
Use the separated nonnegative monomials $a_i^2=(A_0,uB_0,kC_0,D_0)$
and the two pairings $Q_1=a_1a_2-a_3a_4$,
$Q_2=a_1a_3-a_2a_4$. Their exact circuit is
\[
\begin{split}
 T_4&=A_0^2+(uB_0)^2+(kC_0)^2+D_0^2-4k^2B_0C_0\\
 &=\tfrac12\{(Y_0-Z_0)^2+H_0^2+Y_0^2+(Z_0+H_0)^2\}
                   +Q_1^2+Q_2^2.
\end{split}                                      \tag{RB15}
\]
Its matrix on $(Y_0,Z_0,H_0)$ is
$\left(\begin{smallmatrix}1&-1/2&0\\-1/2&1&1/2\\0&1/2&1\end{smallmatrix}\right)$,
which is positive definite with determinant $1/2$. All radicals
separate: $a_1$ is pure testing, $a_2,a_3$ are products of pure-source
and pure-testing radicals, and $a_4$ is pure source.

Here is enough explicit algebra to reproduce the leading collection.
At $g=0$, the canceled source expansions through the needed orders are
\[
\begin{aligned}
 n/s^2&=\tfrac12\{1-b(h+t)\}+O(s),&
 mg&=\tfrac12+\tfrac12(bh-1)s+O(s^2),\\
 C&=\frac{h+t}{2l}-\frac{t+1}{2l}s+O(s^2),&
 A&=\frac t{2h}-\frac t{2h}s+O(s^2),\\
 D&=\frac1{2(t-l)}-
       \frac{bt(h+t)+t-l}{2(t-l)^2}s+O(s^2),&
 \alpha&=\tfrac12 b^2(h+t)s^3+O(s^4).
\end{aligned}                                                     \tag{RB16}
\]
Here $O(s^j)$ means exact divisibility by $s^j$ in a rational ring with
unit denominator, not an asymptotic substitution in the proof identity.
Give $k,u,\omega$ weight one and $v$ weight two. Put $q=b(h+t)$.
Using $s=k^2/b$ and
$\delta=(1-lk^2+O(k^4))/\sqrt{bt}$, $\delta E=\sqrt{t/b}$,
the scalar $G_0=G|_{N=M_B=0}$ satisfies
\[
 G_0=\frac b{2h}(1+u^2)X_0^2+\frac q2T_4
       -\frac{k^2[1+bh(h+t)]}{2h}X_0^2+O_{\rm wt}(8)
       \quad\pmod g.                               \tag{RB17}
\]
The complete first-normal coefficients obey
\[
 L_N=O_{\rm wt}(5),\qquad
 L_{M_B}=\frac{X_0}{h\delta}+O_{\rm wt}(4)\quad\pmod g,
 \tag{RB18}
\]
and the hard block has residue
\[
 \operatorname{diag}\left(K_m+\tfrac12,\,
           A_*+C_*+D_*+\tfrac12\right),
 \quad A_*={t\over2h},\ C_*={h+t\over2l},\ D_*={1\over2(t-l)}.
 \tag{RB19}
\]

Apply the inherited monic contact--conductor identities CC1--CC6 to
the exact difference in (RB17), and to the exact first-normal
differences in (RB18). Their hypotheses are exactly the raw support in
(RB13)--(RB14), together with zero contact value and full first jet.
Both follow from the original stationary parent and (RB10); the model
$X_0^2,T_4$ also has these zeros. No contact difference is inverted.
The conductor maps preserve this grading: $X_0$ has weight two,
the five circuit roots have weight three. The scalar difference of
weight at least eight and normal differences of the displayed weights
therefore have strictly positive-weight Gram coefficients. Every such
coefficient is infinitesimal. Terms divisible by $g$ also lift with
bounded coefficients and remain infinitesimal: the contact (RB9),
$\delta$ and $E$ are independent of $g$, and the canceled coefficient
ring has unit denominators at $g=0$. Clearing and restoring its finite
unit denominators changes none of these conclusions.

Consequently the residue on the eight roots
$(X_0,Y_0,Z_0,H_0,Q_1,Q_2,N,\delta M_B)$ is the strict circuit with
coefficient $q/2>0$, the $N$ pivot $K_m+1/2>0$, and the coupled block
\[
 \begin{pmatrix}
 b/(2h)&bt/(2h)\\
 bt/(2h)&bt(A_*+C_*+D_*+1/2)
 \end{pmatrix}.                                    \tag{RB20}
\]
Its determinant is
\[
 \frac{b^2t}{2h}(C_*+D_*+1/2)>0.                    \tag{RB21}
\]
All pivots are positive bounded units. The separated strict-residue
Gram lemma therefore gives finite tensor SOS for the exact remainder.
Restore (RB10) and the pure-source/pure-testing clearing factors.
This proves the whole local contact parent, for every independent rate
and every normal direction.

\paragraph{All remaining finite state and angle charts.}
The following details make the join an all-tests result.
For bounded $u$ with nonzero residue and infinitesimal $v$, retain
$nD_uD_vX^2$ and $\alpha D_u(vY-Z)^2$, and the centered $m$ excess
$(m-K_m)N^2$ from (RB6), now homogeneous in $Z$. Put
$V_Z=vZ$ and $R_t=uW-Y$. After $X=N-uZ+pv(W+uY)$, the whole
remainder is a bounded quadratic on $(N,R_t,Y,V_Z)$. Its residue
splits off $K_m+1/2$ on $N$, with the other form
\[
 C_*R_t^2+A_*(1+u^2)Y^2
 +(D_*+1/2)\{R_t-(1+u^2)V_Z\}^2
 -(1+u^2)YV_Z-\tfrac12(1+u^2)V_Z^2.               \tag{RB22}
\]
Its first two leading principal minors are positive, and its determinant
is exactly
\[
 (1+u^2)^2A_*C_*(D_*+1/2)u^2>0.                    \tag{RB23}
\]
This direct four-root proof covers every state, including $Z=0$, with
no negative-beta or small-$\rho$ argument.

It remains to consider both angles infinitesimal. The exact testing-core
sign reduction from SD5 handles nonnegative core and literal axes by
original-event squares. On its negative-core branch one may use the
standard exact gauges $u,v>0$. Normalizing the original state by a
maximal coordinate, the case $Z$ infinitesimal has $W$ a unit and is
covered by BE4, with positive real caps of $m,A$ supplied by (RB4).
Retain the actual nonnegative $n,\alpha$ positive square terms. That exact selected
full-parent frame is $(X,Y,uW,vW)$ and has strict residue; no coefficient
bound involving $b$ enters its proof.

If $Z$ is a unit normalize $Z=1$. The same core identity gives
$\omega>0$, and $\omega$ is bounded. If $\omega$ is infinitesimal,
use (RB17)--(RB21). If $\omega$ is a positive unit, use the exact
relations $B_0=(\delta/\omega)A_0$ and $C_0=B_0-X_0$ in the raw
support (RB13)--(RB14). The inverse is pure testing. Its complete
residue on $(A_0,X_0,D_0,N,\delta M_B)$ consists of the positive
$q/2$ pivots on $A_0,D_0$, the $K_m+1/2$ normal pivot, and precisely
(RB20). Direct substitution in (RB11) gives these residues; all other
coefficients contain $k,g,u,v$ after the cancellations already recorded.
Thus this whole state chart is strict as well.

\paragraph{Every other projective test.}
Use the exact centered-event identity R5.5 for the finite remote pairs
$v_0\ne0$. All source residues are (RB4) and $n_0=\alpha_0=0$.
On the kernel of the limiting centered $m$ row, the determinant is
\[
 \frac{(1+U)^2(1+V)^2(h+t)}{8hl(t-l)}
       \{V^2+t(h+t)U^2V+hUV^2\}>0.                 \tag{RB24}
\]
The restriction to $Z=0$ is positive by the original $A,C$ positive square terms
and the nonnegative signed-core restriction, exactly as in RT10.
A sufficiently large positive real cap $K_m$ makes the full bounded
residue strict. Its choice depends on the real testing residues only;
unbounded $m$ permits every such real cap.

For unbounded $u$ and bounded unit $v$, R6.5 at $n_0=0$ is strict
because $C_*W^2$, $A_*Y^2$ and $(D_*+1/2)(W+vX)^2$ have common
kernel zero. For unbounded $u$ and infinitesimal $v$, the exact
BE $(\infty,0)$ frame $(pX+Z,W,Y,vX)$ retains every event except
positive caps of $C,A$, and has strict residue. For unbounded $v$
and bounded $u$, R6.6 at $n_0=0$ is strict; the centered correction
vanishes on its $m$-row kernel. For both angles unbounded, R6.7 at
$n_0=0$ has block determinants $D_*/2$ and
$A_*(K_m+1/2)$, both positive. These are exact homogeneous full-parent
identities. Their literal reciprocal zeros include the actual projective
endpoints; undoing a nonzero homogeneous factor uses a pure-testing
square only.

The charts exhaust the stated $b$-unit source domain. Restore all retained
squares and normalizations and apply the real identities to the real and
imaginary parts of a complex state.

\paragraph{Imported exact identities.}
The BE4 four-root state identity and the unnumbered $(\infty,0)$
identity in Section~\ref{component:p69995d829bb2} supply the two
selected full-parent frames used above. The latter is the paragraph
beginning ``At $(\infty,0)$ retain all events except $c,A$''.
The remote kernel determinant and its $Z=0$ restriction are printed
in (RB24) and the paragraph following it. The centered-event identity
is R5.5 in Section~\ref{component:p8da9ea6ed9e5}. The homogeneous
kernel and root matrices R6.5--R6.7 are printed in
Section~\ref{component:pd5fa361b4623}; they are used at $n_0=0$
with the source residues (RB4). Finally SD5 and its following sign
Schur identity in Section~\ref{component:p0acdb66bba2f} give the
pure-testing inverse and the exact sign reduction. Only these
algebraic identities are used with the present source estimates;
no C1 physical-source or negative-beta perturbation theorem is imported.

\subsection{Infinitesimal $a^2B$}
\label{component:p75b395956985}

\label{p75b395956985:sec:r6-b-inf}

\paragraph{Source scope and relation to the preceding theorem.}
Retain the complete original weak-C2 source: $d=1$, $\lambda=\beta=0$,
$n\geq0$, the actual stationary contact, every strict gap, total scale,
and the complement of prior cases. On R6, $t,h,l=1-h,r$ and
the original $j$ gap are positive bounded units, while $g,B,x=aB,y=tB$
are positive infinitesimals. This section treats $a$ unbounded and
$b=a^2B$ positive infinitesimal. The notation $s=1/a$, $w_s=ae$,
$B=s^2b$, $x=sb$, $e=s(1+hsb-g)$ is unchanged from (RB1).
The exact physical identity (RB2), equivalently RT2, remains in force.
It bounds $b$ on all unbounded-$a$ R6 and gives $\rho_0=1/2$ here,
where $\rho=n/e^2$. In particular this subdomain has $n>0$; the
literal $n=0$ source face is covered by the preceding $b$-unit theorem.

\begin{theorem}
On this entire original R6 source intersection with $b=a^2B$
infinitesimal, the complete original parent is a finite separated
source--testing tensor sum of squares at every independent projective
testing pair and every real or complex state. Both angle signs, all
literal projective endpoints, and all relative opening rates of
$s,b,g$ are included. No mixed inverse or mixed square root is used.
\end{theorem}

Together with the $b$-unit theorem and the inherited bounded-$a$
R6 theorem, this proves the entire R6 row. It does not by itself
admit another weak-C2 row or the global finite-lift theorem.

\paragraph{A regular source chart.}
Keep $R=1+hx$, $K=1+x$, $T=1+xr$, $k=\sqrt{x}$ and the contact
parameters $\delta=R/(K\sqrt{bt})$, $E=tK/R$ from (RB8). Define
\[
 \varepsilon=\delta^{-1},\qquad \jmath=\delta k,
 \qquad k=\varepsilon\jmath.
 \tag{RI1}
\]
Both $\varepsilon$ and $\jmath$ are positive infinitesimals. In this
section $\jmath$ denotes the contact scale, not the original source
gap called $j$. The exact source substitution is
\[
 \begin{gathered}
 x=\varepsilon^2\jmath^2,\quad K=1+x,\quad R=1+hx,\quad
 E=tK/R,\\
 b=\frac{\varepsilon^2R^2}{K^2t},\qquad
 s=\frac{t\jmath^2K^2}{R^2}.
 \end{gathered}                                                   \tag{RI2}
\]
All denominators in (RI2) are positive source units. There is no
comparison of a testing variable with either source infinitesimal.

Use precisely the complete contact-centered remainder $G$ in (RB11),
with the full coordinate inverse (RB7). Thus the two retained positive square terms
on the contact chart are $(m-K_m)N^2$ and
$nD_uD_v(X+\gamma v\omega)^2$, where
\[
 p=(1+hx-g)/R,\quad c=l/(tT),\quad
 \gamma=rK/(sR),\quad D_u=1+u^2,\quad D_v=1+v^2.
\]
Its exact contact is now
\[
 (u,\omega,v;N,M_B)
   =(\varepsilon\jmath,\jmath,\varepsilon\jmath^2E;0,0).
 \tag{RI3}
\]
Substitution of (RI3) in the complete contact formulas
(RB7)--(RB9a) makes both normals and the retained $n$ row vanish.

\paragraph{Two exact source cancellations.}
The quantities $mg$, $n\gamma$, $n\gamma^2$ and
$\alpha/(s^3b^2)$ have canceled denominators that are units at
$s=b=g=0$. Introduce
\[
 \mu_1=mg/R,\quad \mu_2=(mg)g/R^2,\quad
 F_0=p^2/2-n\gamma^2-2n\gamma p+\mu_2,\qquad
 A_N=p-2\mu_1-2n\gamma.
 \tag{RI4}
\]
Both $F_0$ and $A_N$ are divisible by $b$ in this rational unit ring.
Here is a short exact verification, independent of a series expansion.
At $b=0$, put $q=1-g$. Direct cancellation in the original source gives
\[
 \begin{gathered}
 p=q,\quad mg=\frac{q}{2(1+s)},\quad
 n\gamma=\frac{sq}{2(1+s)},\quad
 n\gamma^2=\frac{q(1-sq)}{2(1+s)}.
 \end{gathered}                                                    \tag{RI5}
\]
Substitution makes both expressions in (RI4) identically zero.
The denominators are units before this specialization; hence the
polynomial numerators are divisible by $b$, without introducing a
pole. Their needed residues are
\[
 \frac{F_0}{b}\sim\frac{h+t}{2},\qquad
 \left.\frac{A_N}{b}\right|_{s=g=0}=0,\qquad
 \left.\frac{A_N}{bs}\right|_{g=0}\sim h+t.
 \tag{RI6}
\]
The last quotient also has a unit canceled denominator at $s=b=0$.
Every statement in (RI5)--(RI6) is separately checked from the original
source formulas. In particular the source estimate is uniform in
$g/s$, $s/b$ and all other opening ratios.

\paragraph{Complete bounded raw support, including both first normals.}
Use the separated raw monomials
\[
 \mathsf A=\varepsilon\omega v,\quad
 \mathsf B=v,\quad \mathsf C=Eu\omega,\quad
 \mathsf D=\varepsilon^2\jmath^3E,
 \qquad \mathsf x=\mathsf B-\mathsf C.
 \tag{RI7}
\]
Set $H=\alpha+1/2$, $L=AD_v+Hv^2$, and $V_0=1-puv\omega$.
At zero normals, (RB7) gives
\[
 q_B=D_uvV_0,\quad uW-Y=cD_uvV_0,\quad
 Y=u\omega-cvV_0,\quad
 J=-(1-c)D_uvV_0,
\]
and the entire scalar has the following exact allocation:
\[
 G_0=f_A\mathsf A^2+f_{BB}\mathsf B^2+
 f_{BC}\mathsf B\mathsf C+f_{CC}\mathsf C^2+
 \varrho\mathsf D^2,                               \tag{RI8}
\]
where
\[
\begin{aligned}
 f_A={}&D_u^2D_vF_0/\varepsilon^2,\\
 f_{BB}={}&[Cc^2D_u^2D_v+(D+1/2)(1-c)^2D_u^2
                  +D_uLc^2]V_0^2\\
          &+2HD_ucV_0-D_u/2-D_u^2\mu_2v^2\omega^2,\\
 f_{BC}={}&[-2D_uLcV_0-2HD_u+2n\gamma D_uD_v]/E,\\
 f_{CC}={}&[D_uL+D_uD_vn\gamma^2v^2]/E^2,\\
 \varrho={}&\alpha D_u/\mathsf D^2.
\end{aligned}                                                      \tag{RI9}
\]
These coefficients are bounded. In particular
\[
 \frac{F_0}{\varepsilon^2}=
       \frac{R^2}{K^2t}\frac{F_0}{b},\qquad
 \frac{\alpha}{\mathsf D^2}=
       \frac1t\frac{\alpha}{s^3b^2}.
 \tag{RI10}
\]
The term that would give $u^2/\varepsilon^2$ in the first raw
coefficient has instead been placed in $f_{CC}$, using the exact
identity $u^2(v\omega)^2=(v\mathsf C/E)^2$.
The final term of $f_{BB}$ likewise allocates the whole
$-D_u^2\mu_2v^4\omega^2$ contribution without dividing by
$\varepsilon$. These are identities for the full scalar.

Write the exact normal expansion as
$G=G_0+L_NN+L_MM_B+G_{\rm hard}$, with $G_{\rm hard}$ quadratic.
The two complete linear coefficients are
\[
 L_N=R_A\mathsf A+R_B\mathsf B+R_C\mathsf C,
 \qquad L_M=S_B\mathsf B+S_C\mathsf C,               \tag{RI11}
\]
where
\[
 R_A=D_uA_N/\varepsilon,\qquad R_C=2cuvL/E,\qquad S_C=-2L/E,
\]
\[
 R_B=\Delta_{v,0}(L_N-R_A\mathsf A-R_C\mathsf C),\qquad
 S_B=\Delta_{v,0}(L_M-S_C\mathsf C).\tag{RI12}
\]
Here $L_N,L_M$ are the complete coefficient extractions from the bounded
full parent in the stated state coordinates. The displayed source
cancellations make both quotient numerators vanish identically at
$v=0$. Thus these are finite monic polynomial divisions, with the pure
testing denominator $D_u^2$ cleared first. They give all normal
coefficients, including the higher testing terms. Their coefficients
are bounded; division shifts testing degree by one and preserves the
weighted coefficient estimates used below.

Here $A_N/\varepsilon=
\varepsilon R^2(A_N/b)/(K^2t)$ is bounded. Thus no first-normal
pole has been discarded. Every hard coefficient is bounded as well:
the only testing inverse in the full change of state is $D_u^{-1}$,
and all original source coefficients in (RB11) are bounded after the
displayed cancellation. At $u=v=0$ the complete hard matrix is exactly
\[
 \operatorname{diag}(K_m+1/2,A+C+D+1/2).             \tag{RI13}
\]

\paragraph{A bounded conductor for these complete coefficients.}
Define the contact differences
\[
 \mathsf y=\mathsf A-k\mathsf C,\quad
 \mathsf z=u\mathsf B-k\mathsf C,\quad
 \mathsf h=k\mathsf C-\mathsf D.
\]
The scaled conductor identities are
\[
\begin{aligned}
 \mathsf C(u-k)&=\mathsf z-u\mathsf x,\\
 \varepsilon\mathsf C(\omega-\jmath)
       &=\mathsf y-\varepsilon\omega\mathsf x,\\
 \mathsf C(v-\varepsilon\jmath^2E)
       &=E(\omega\mathsf z+\jmath\mathsf y-k\omega\mathsf x),\\
 \mathsf Cv(\omega-\jmath)
       &=E\omega(\omega-\jmath)\mathsf z+
         E\jmath\omega\mathsf y-Ek\omega^2\mathsf x.
\end{aligned}                                                       \tag{RI14}
\]
These are SE5--SE6 with $w=\omega$, also checked directly here.
For a monic contact division
\[
 S=(u-k)T_u+(\omega-\jmath)T_\omega+
           (v-\varepsilon\jmath^2E)T_v,
 \qquad T_\omega=\varepsilon U+vV,
\]
the exact bounded lift is
\[
\begin{aligned}
 \mathsf CS={}&\mathsf x[-uT_u-\varepsilon\omega U
                    -Ek\omega^2V-Ek\omega T_v]\\
 &+\mathsf y[U+E\jmath\omega V+E\jmath T_v]\\
 &+\mathsf z[T_u+E\omega(\omega-\jmath)V+E\omega T_v].
\end{aligned}                                                       \tag{RI15}
\]
No inverse of a small source parameter or a contact difference occurs.

We verify the hypothesis of (RI15) for the full coefficients, rather
than import the different small-$e$ source assumptions of SE8.
Clear the finite source unit denominators and $D_u^2$. They can be
chosen independent of $\omega$ and even in $u,v$. In (RI9), all
$\omega$ dependence is through $uv\omega$ and $v^2\omega^2$;
the other testing dependence is even in $u,v$. The same property
holds after subtracting any model of the form used below.

For the scalar square-lift algorithm CC4, substitute
$\mathsf A=\mathsf y+k\mathsf C$,
$\mathsf B=\mathsf x+\mathsf C$,
$\mathsf D=k\mathsf C-\mathsf h$ into (RI8).
The coefficient requiring its first conductor is
\[
 S=k^2f_A+f_{BB}+f_{BC}+f_{CC}+k^2\varrho.
\]
Its zero contact value follows by domain cancellation, as in CC4.
Divide monically in the order $u$, then $\omega$, then $v$.
After $u=k=\varepsilon\jmath$, the $\omega$ divided difference
belongs to $(\varepsilon v,v^2)$. Accordingly choose
$T_\omega=\varepsilon U+vV$ with $U,V\in(v)$.
Modulo $\varepsilon$, the remaining $v$ polynomial is even, so
$T_v\in(\varepsilon,v)$; also
$T_u|_{\varepsilon=u=v=0}=0$.

After this first conductor the coefficient of $\mathsf h$ remains
$-2k\varrho$. The sole gradient relation is
$(-2k,1,1,-1)$ on
$(\mathsf x,\mathsf y,\mathsf z,\mathsf h)$; the determinant of
the first three contact gradients is
$-\varepsilon^3\jmath^4E^2\ne0$. The adjugate cancellation of
CC5 therefore removes the contact values without putting an inverse
of this determinant into any coefficient. Subtract the contact
multiple $2k\varrho_c(-2k,1,1,-1)$.
Each of the four resulting coefficient remainders $J_i$ vanishes at
contact. Its restriction at $\varepsilon=u=v=0$ is independent of
$\omega$: for the $\mathsf x$ coefficient this follows from
$2f_{BB}+f_{BC}$; the other three restrictions are zero by the
$U,V,T_u,T_v$ properties just established and the factor $k$.
Consequently the next monic $\omega$ divided difference of every
$J_i$ belongs to $(\varepsilon,v)$. Formula (RI15) applies again.

The last scalar repair uses the exact circuit relation
\[
 k\mathsf C(\mathsf y+\mathsf z-\mathsf h-2k\mathsf x)
       =Q_1^2-\mathsf y\mathsf z,
\]
where $Q_1$ is the first of the separated radical pairings below.
This completes the same restricted scalar square lift as CC6.
For (RI11), one conductor suffices: after replacing
$\mathsf A,\mathsf B$ as above, the coefficient multiplying
$\mathsf C$, restricted at $\varepsilon=u=v=0$, is independent of
$\omega$ directly from (RI12). Its divided difference therefore
lies in $(\varepsilon,v)$ and (RI15) applies. The full scalar and
both full normals have thus been lifted over the bounded coefficient
ring. The restricted raw support $\mathsf A^2,\mathsf B^2,
\mathsf B\mathsf C,\mathsf C^2,\mathsf D^2$ is essential;
an unrestricted ideal claim is neither needed nor asserted.

The original stationary parent and the two retained contact-zero
positive square terms give the exact scalar value and full first-jet identities
used here. The contact (RI3) is independent of $g$. Hence taking
the exact coefficient difference divisible by $g$ preserves these
identities and every raw-support condition.

\paragraph{Strict residues of the complete contact lift.}
Suppose now that $u,v,\omega$ are all positive infinitesimals, as
in the remaining contact state chart. Give
\[
 (\varepsilon,\jmath,u,\omega,v)
       \text{ weights }(1,1,2,1,3).
\]
Then $\mathsf x$ has weight three and the five circuit roots have
weight five. Every contact substitution and (RI14)--(RI15) preserves
the weights. Put $\kappa=(h+t)/(2t)$ and
\[
 \Lambda=\frac{1-s+s^2+u^2-[2l+h(h+t)]k^2}{2ht},\qquad
 \mu=\frac1{h(1+s)}.
\]
Direct exact collection of (RI9) through weight ten gives
\[
 G_0=\Lambda\mathsf x^2+\kappa T_4+O_{\rm wt}(12)
                   \quad\pmod g,                 \tag{RI16}
\]
where
\[
 T_4=\mathsf A^2+(u\mathsf B)^2+(k\mathsf C)^2+
          \mathsf D^2-4k^2\mathsf B\mathsf C.
\]
The two complete normal collections are
\[
 L_N=O_{\rm wt}(8),\qquad
 L_M=\mu\mathsf x+O_{\rm wt}(7)\quad\pmod g.
 \tag{RI17}
\]
For transparency, the only source terms needed for (RI16)--(RI17),
besides (RI6), are the following exact Taylor coefficients at $g=0$:
\[
\begin{aligned}
 A&=\frac t{2h}(1-s+s^2)+O_{\deg(s,b)}(3),\\
 C&=\frac{h+t}{2l}-\frac{t+1}{2l}s+
                      \frac{t+1}{2l}s^2+O_{\deg(s,b)}(3),\\
 D&=\frac{1-s+s^2}{2(t-l)}-
        \frac{sbt(h+t)}{2(t-l)^2}+O_{\deg(s,b)}(3),\\
 n\gamma&=\frac s2-\frac{s^2}2-\frac{sb(h+t)}2+
                         O_{\deg(s,b)}(3).
\end{aligned}                                                       \tag{RI18}
\]
Here the error denotes exact degree filtration in a rational ring
with unit denominators. Also $s=t\jmath^2+O_{\rm wt}(6)$,
$c=(l/t)(1-k^2)+O_{\rm wt}(6)$ and
$E=t(1+lk^2)+O_{\rm wt}(8)$. The weight eight bound for $L_N$
follows term by term from (RI6), (RI12).
At $g=0$, after the displayed unit denominators are cleared, all
scalar weighted monomials have even weight, and all $M_B$ linear
monomials have odd weight. Hence the next possible degrees are
twelve and seven, respectively.

The stronger coefficient-wise statement also holds for the displayed
raw representation: $f_A-\kappa$ and $\varrho-\kappa$ have weight at
least two, while
\[
 f_{BB}-\Lambda-\kappa u^2,\qquad
 f_{BC}+2\Lambda+4\kappa k^2,\qquad
 f_{CC}-\Lambda-\kappa k^2
\]
each has weight at least six. Similarly $S_B-\mu,S_C+\mu$ have
weight at least four separately, and $(R_A,R_B,R_C)$ have weights
at least $(3,5,5)$. Thus the chosen raw allocation of each exact
difference already has the claimed filtration; no low-weight
Gram coefficients are hidden by cancellations between raw products.

Apply the bounded conductor just proved to the exact differences in
(RI16)--(RI17). The scalar correction Gram coefficients have weight
at least two; the $M_B$ normal corrections have weight at least two;
the $N$ normal corrections have weight at least three. They are all
infinitesimal. The coefficients divisible by $g$ lift boundedly by the
same argument and are infinitesimal independently of every other
opening. Restoring the finite unit denominators preserves this fact.

For strictness use the two pairings on the separated monomials
\[
 a_i^2=(\mathsf A,u\mathsf B,k\mathsf C,\mathsf D),\qquad
 Q_1=a_1a_2-a_3a_4,\quad Q_2=a_1a_3-a_2a_4.
\]
Their product relation gives the exact two-pairing identity (RB15),
with the scaled roots. Its three-dimensional difference block has
determinant $1/2$, and both pairing pivots are positive. The complete
residue on
$(\mathsf x,\mathsf y,\mathsf z,\mathsf h,Q_1,Q_2,N,M_B)$
therefore consists of this strict circuit with coefficient
$\kappa>0$, the $N$ pivot $K_m+1/2>0$, and
\[
 \begin{pmatrix}
 1/(2ht)&1/(2h)\\
 1/(2h)&A_*+C_*+D_*+1/2
 \end{pmatrix},
 \quad A_*=t/(2h),\ C_*=(h+t)/(2l),\ D_*=1/[2(t-l)].
 \tag{RI19}
\]
Its determinant is exactly
\[
 \frac{C_*+D_*+1/2}{2ht}>0.                         \tag{RI20}
\]
All pivots are positive bounded units. The strict separated Gram
lemma WA10 yields a finite tensor-SOS identity for the complete
remainder. Restore the two actual nonnegative retained positive square terms and
all separated clearing factors. This proves the entire local contact
chart, including every normal direction and arbitrary relative rates.

\paragraph{The $\omega$-unit state chart uses the original $n$ positive square term.}
When $u,v$ are infinitesimal, $Z=1$, and $\omega$ is a positive unit,
retain $nD_uD_vX^2$ itself, together with $(m-K_m)N^2$.
The complete remainder is (RB11) with its two negative $n\gamma$
and $n\gamma^2$ contributions absent. This choice is legitimate
directly from $n\geq0$ and does not require a contact lift away from
the contact state scale.

Use the raw roots
\[
 (A_o,\mathsf x,\mathsf D,N,M_B),\qquad A_o=\omega v.
\]
The only new inverse is $1/\omega$, which is pure testing and
bounded. The relations $\mathsf B=A_o/\omega$ and
$\mathsf C=A_o/\omega-\mathsf x$ put the entire scalar and both
first normals in this frame, by the same full formulas (RI9),
(RI12) with the absent $n$ terms deleted. All coefficients are
bounded; the constant term is still $\alpha D_u=\varrho\mathsf D^2$.
Explicitly $\mathsf A=\varepsilon A_o$, so the first scalar term
becomes $D_u^2D_v(p^2/2+\mu_2)A_o^2$ and the first $N$-normal
term becomes $D_u(p-2\mu_1)A_o$. The printed $\varepsilon$
denominators therefore cancel before any coefficient estimate.
Its exact residue is the $A_o$ pivot $1/2$, the $\mathsf D$ pivot
$\kappa$, the $N$ pivot $K_m+1/2$, and the coupled block (RI19).
The first $N$ normal has zero residue and the first $M_B$ normal
has residue $\mathsf x/h$. Thus this entire state chart
also has a strict finite separated Gram.

\paragraph{All other independent tests and states.}
Every source bound used by the noncontact joins of the preceding
$b$-unit theorem is (RB4), which is valid here without a lower bound
on $b$. For bounded unit $u$ and infinitesimal $v$, retain the original
$n,\alpha$ events and use the exact full-state form (RB22), whose
determinant (RB23) is positive. For both angles infinitesimal, the
SD5 exact sign reduction disposes of the nonnegative-core branch
and literal zero axes by original-event squares. On the negative
branch choose the exact gauges $u,v>0$. If the normalized state has
$Z$ infinitesimal, BE4 gives the strict full-state frame
$(X,Y,uW,vW)$ using positive real caps of $m,A$.
Otherwise normalize $Z=1$; $\omega>0$ is bounded, and its two
possibilities are exactly the contact chart and the unit chart just
proved. No source/testing ratio split is used.

For finite remote $v_0\ne0$, use the actual centered-event identity
R5.5, the strict kernel determinant RT10 (printed as RB24), and a
sufficiently large positive real cap below unbounded $m$.
For unbounded $u$ and bounded unit $v$, use R6.5 specialized to
$n_0=0$ and the source residues (RB4). For unbounded $u$ and
infinitesimal $v$, use the exact BE $(\infty,0)$ full frame
$(p_uX+Z,W,Y,vX)$, $p_u=1/u$, with positive caps of $C,A$.
For unbounded $v$ and bounded $u$, R6.6 at $n_0=0$ is strict.
For both angles unbounded, R6.7 at $n_0=0$ has determinants
$D_*/2$ and $A_*(K_m+1/2)$ on its two nontrivial blocks.
These uses import the displayed homogeneous identities with the
source bounds (RB4), without importing a source theorem from another
stratum.
Literal reciprocal zeros are polynomial values of those charts.
All signs and complex states follow by the same exact gauges and
the real/imaginary split. This exhausts all independent tests.

\paragraph{The contact identities.}
Equations (RI8)--(RI15) establish the rational-support hypotheses of
SE5--SE7 in Section~\ref{component:pa7e744614f49} for the present
coefficients, with the bounds proved here. We then apply the monic
square-lift and gradient repair CC4--CC6 in
Section~\ref{component:pfabc29e28ab3}.

\section{Two small actual angles: the R7 sector}
\subsection{Unbounded $a$, including the physical zero face}
\label{component:p405d96527133}

\label{p405d96527133:sec:r7-unbounded-a}

\paragraph{Quantified source and conclusion.}
Keep the entire original weak-C2 source: $d=1$, $\lambda=\beta=0$,
the actual rational stationary coefficients, $n\ge0$, every strict gap,
the original finite stationary contact, total scale, and all prior
exclusions. In R7, $B,z,y$ are positive infinitesimals,
$a=z/B$, $t=y/B$ is a positive bounded unit, and $h,g$ are positive
infinitesimals. The original exclusion of simultaneously infinitesimal
$r$ and the source gap $j_{\rm gap}$ is retained.

\begin{theorem}
On the whole original R7 source intersection with $a$ unbounded,
the complete original parent is a finite separated source--testing
tensor sum of squares for every independent projective testing pair
and every real or complex state. Both angle signs, all literal zero
and infinite angles, literal $n=0$, and every relative opening rate
of the source infinitesimals are included. No inverse or radical of
a mixed source--testing expression is used.
\end{theorem}

Together with the bounded-$a$ theorem this covers R7.

\paragraph{Actual source scale and the physical inequality.}
Use
\[
 \begin{gathered}
 s=1/a,\quad b=a^2B,\quad x=sb=z,\quad B=s^2b,\\
 w_s=ae=1+hx-g,\quad e=sw_s,\quad r=1-e,\quad l=1-h.
 \end{gathered}
 \tag{RU1}
\]
Set $K=1+x$, $R=1+hx$, $T=1+xr$ and
$j_{\rm gap}=T-l/t$. The exact original source identity is
\[
\begin{gathered}
 2r\Delta n=eQ-BK_{\rm phys},\qquad
 2r\Delta_*n/s^2=w_sQ-bK_*,\qquad \Delta_*=s\Delta,\\
 Q=1+B(1+t+ah+2ar+at)
    +B^2\{a(h+t)+a^2[h(r+1)+t(1+r-r^2)]\},\\
 K_{\rm phys}=(a+1)(h+t)(1+x),\qquad
 K_*=(1+s)(h+t)(1+x).
\end{gathered}                                                     \tag{RU2}
\]
Before bounding $b$, $Q\sim\Delta_*\sim1$ since all their nonconstant
terms contain $s$ or $x$. Here $w_s\sim1$ and $K_*\sim t$ is a
positive unit. The original inequality $eQ\ge BK_{\rm phys}$
therefore gives $b$ bounded and $n=O(s^2)$, including its literal
equality case. In particular $r\sim1$, so the original R7 exclusion
is respected. The source gap gives $t_0\ge1$; $j_{\rm gap}$ is a
unit exactly when $t_0>1$, and otherwise $t_0=1$.

The original rational formulas, after cancellation, give jointly
\[
 mg\sim\tfrac12,\quad C\sim\tfrac t2,\quad
 A_h:=Ah\sim\tfrac t2,\quad
 \beta_D:=D(1-c)=Dj_{\rm gap}/T\sim\frac1{2t},
 \quad \frac\alpha{s^3b^2}\sim\tfrac t2,
 \qquad c=\frac l{tT}.
 \tag{RU3}
\]
Thus $m,A$ are unbounded; $C$ is a positive unit. If $t_0>1$,
$D\sim1/[2(t-1)]$ is a positive unit. If $t_0=1$, $D$ is unbounded.
The canceled coefficients $mg,C,A_h,\beta_D,n\gamma,n\gamma^2$
below have unit denominators through both small-gap ends. No ratio
$h/g$, $j_{\rm gap}/h$, or comparison with a testing rate is used.

\paragraph{Four whole positive square terms and a bounded full-parent remainder.}
Use the complete parent $P=F/4$:
\[
\begin{aligned}
 P={}&nD_uD_vX^2+mM^2+CD_v(uW-Y)^2+AD_uD_vY^2\\
 &+(D+1/2)J^2+(\alpha+1/2)D_u(vY-Z)^2
                   +D_v(X+uZ)^2/2-D_uD_vZ^2/2,\\
 M={}&vW-X+uvY-uZ,\qquad J=uW+uvX-Y-vZ.
\end{aligned}
\]
The exact coordinate inverse (RB6)--(RB7) is
\[
\begin{gathered}
 D_u=1+u^2,\quad D_v=1+v^2,\quad Z=1,\quad
 p=w_s/R,\quad q_A=D_uv\omega,\\
 N=X+u-pq_A,\qquad q_B=v(1-uX),\qquad M_B=uW-Y-cq_B,\\
 X=N+pD_uv\omega-u,\quad q_B=v(D_u-uN-puD_uv\omega),\\
 W=\omega+u(M_B+cq_B)/D_u,\qquad
 Y=u\omega-(M_B+cq_B)/D_u.
\end{gathered}                                                     \tag{RU4}
\]
The sole testing inverse is $D_u^{-1}$. Put
\[
 \begin{gathered}
 \gamma=\frac{rK}{sR},\quad \eta=Y-huW,\quad q=1-c=j_{\rm gap}/T,\\
 \mu_1=mg/R,\quad \mu_2=(mg)g/R^2,\quad n_1=n\gamma,\quad n_2=n\gamma^2.
 \end{gathered}
\]
Retain the four entire nonnegative squares
\[
 mN^2,\qquad nD_uD_v(X+\gamma v\omega)^2,
 \qquad AD_uD_v\eta^2,\qquad DM_B^2.
 \tag{RU5}
\]
All four rows vanish at the actual contact; (RB9a) verifies the first
two, and directly $Y_c=hu_cW_c$, $M_{B,c}=0$.
The $n$ positive square term vanishes identically on the literal $n=0$ face.

The full $A,D$ cancellations are the exact identities
\[
\begin{aligned}
 AD_uD_v(Y^2-\eta^2)&=2A_hD_uD_vuWY-A_hhD_uD_v(uW)^2,\\
 D(J^2-M_B^2)&=-2\beta_DM_Bq_B+\beta_Dq q_B^2,
 \qquad J=M_B-q q_B.
\end{aligned}                                                     \tag{RU6}
\]
Since $M=(g/R)q_A-N$, subtraction of precisely (RU5) leaves
\[
\begin{aligned}
 G={}&-2\mu_1q_AN+\mu_2q_A^2
       -2n_1D_uD_vX(v\omega)-n_2D_uD_v(v\omega)^2\\
    &+CD_v(uW-Y)^2+2A_hD_uD_vuWY-A_hhD_uD_v(uW)^2\\
    &-2\beta_DM_Bq_B+\beta_Dq q_B^2+J^2/2\\
    &+(\alpha+1/2)D_u(vY-1)^2+D_v(X+u)^2/2-D_uD_v/2.
\end{aligned}                                                     \tag{RU7}
\]
This is a whole-parent identity, including every scalar, normal and
hard coefficient. Equations (RU3), (RU6) bound all of them before
any contact restriction. In particular neither $h^{-1}$ nor
$j_{\rm gap}^{-1}$ survives.

\paragraph{One regular contact frame for the two $b$ alternatives.}
Use the pure-source parameters
\[
 \varepsilon=\frac{K\sqrt{bt}}R,\quad
 \jmath=\sqrt{s/t}\,\frac RK,\quad k=\varepsilon\jmath=\sqrt x,
 \quad E=tK/R.
 \tag{RU8}
\]
The symbol $\jmath$ is a contact scale, not $j_{\rm gap}$.
Here $\jmath$ is infinitesimal; $\varepsilon$ is a unit or
infinitesimal exactly as $b$ is. The exact inverse source chart is
\[
 x=\varepsilon^2\jmath^2,\qquad
 b=\frac{\varepsilon^2R^2}{K^2t},\qquad
 s=\frac{t\jmath^2K^2}{R^2}.
 \tag{RU9}
\]
Only source units are inverted. The contact is
$(u,\omega,v;N,M_B)=(k,\jmath,\varepsilon\jmath^2E;0,0)$.

Define
\[
\begin{gathered}
 \mathsf A=\varepsilon\omega v,\quad\mathsf B=v,
 \quad\mathsf C=Eu\omega,\quad\mathsf D=\varepsilon^2\jmath^3E,\\
 \mathsf x=\mathsf B-\mathsf C,\quad
 \mathsf y=\mathsf A-k\mathsf C,\quad
 \mathsf z=u\mathsf B-k\mathsf C,\quad
 \mathsf h=k\mathsf C-\mathsf D.
\end{gathered}                                                     \tag{RU10}
\]
The two potentially singular source combinations are
\[
 F_0=p^2/2-n_2-2n_1p+\mu_2,\qquad A_N=p-2\mu_1-2n_1.
\]
Their exact canceled quotients satisfy
\[
 F_0/b\sim t/2,\qquad A_N/(bs)\sim t.
 \tag{RU11}
\]
Both denominators are units through $s=b=h=g=0$. These assertions
follow from the original rational source, not from the inequality
$n\ge0$. In particular
\[
 F_0/\varepsilon^2=\frac{R^2}{K^2t}\frac{F_0}b,
 \qquad \alpha/\mathsf D^2=\frac1t\frac\alpha{s^3b^2},
 \qquad A_N/\varepsilon=\frac{\varepsilon R^2s}{K^2t}
                               \frac{A_N}{bs}
 \tag{RU12}
\]
are bounded without losing the zero-$n$ face.

\paragraph{Exact full scalar and both first normals.}
Put $H=\alpha+1/2$, $V_0=1-puv\omega$. At $N=M_B=0$,
$uW=\mathsf C/E+u^2c\mathsf BV_0$ and
$Y=\mathsf C/E-c\mathsf BV_0$. The complete scalar is
\[
 G_0=f_A\mathsf A^2+f_{BB}\mathsf B^2+
       f_{BC}\mathsf B\mathsf C+f_{CC}\mathsf C^2+
       \varrho\mathsf D^2,                                      \tag{RU13}
\]
with the bounded coefficients
\[
\begin{aligned}
 f_A={}&D_u^2D_vF_0/\varepsilon^2,\qquad
 \varrho=\alpha D_u/\mathsf D^2,\\
 f_{BB}={}&[Cc^2D_u^2D_v+(\beta_Dq+q^2/2)D_u^2+HD_uv^2c^2
             -A_hD_uD_v(2u^2+hu^4)c^2]V_0^2\\
          &+2HD_ucV_0-D_u/2-D_u^2\mu_2v^2\omega^2,\\
 f_{BC}={}&[2A_hD_uD_v(lu^2-1)cV_0-2HD_uv^2cV_0
                   -2HD_u+2n_1D_uD_v]/E,\\
 f_{CC}={}&[A_h(2-h)D_uD_v+HD_uv^2+D_uD_vn_2v^2]/E^2.
\end{aligned}                                                     \tag{RU14}
\]
For the full normal expansion
$G=G_0+L_NN+L_MM_B+G_{\rm hard}$, the exact coefficients are
\[
 L_N=R_A\mathsf A+R_B\mathsf B+R_C\mathsf C,
 \qquad L_M=S_B\mathsf B+S_C\mathsf C,                           \tag{RU15}
\]
where
\[
 R_A=D_uA_N/\varepsilon,\qquad R_C=2cuv\{A_hD_v(1-lu^2)+Hv^2\}/E,\qquad S_C=2\{A_hD_v(lu^2-1)-Hv^2\}/E,
\]
\[
 R_B=\Delta_{v,0}(L_N-R_A\mathsf A-R_C\mathsf C),\qquad
 S_B=\Delta_{v,0}(L_M-S_C\mathsf C).\tag{RU16}
\]
Here $L_N,L_M$ are the complete coefficient extractions from the bounded
full parent in the stated state coordinates. The displayed source
cancellations make both quotient numerators vanish identically at
$v=0$. Thus these are finite monic polynomial divisions, with the pure
testing denominator $D_u^2$ cleared first. They give all normal
coefficients, including the higher testing terms. Their coefficients
are bounded; division shifts testing degree by one and preserves the
weighted coefficient estimates used below.

At $u=v=0$ the complete hard matrix on $(N,M_B)$ is exactly
$\operatorname{diag}(1/2,C+1/2)$.

\paragraph{Bounded monic conductor, including the small-$\varepsilon$ end.}
The identities (RI14)--(RI15) apply to (RU10) without a source-shape
assumption. Their crucial inputs are checked here on the new complete
coefficients. After finite source unit denominators and $D_u^2$ are
cleared, every $\omega$ occurrence of (RU14) is through $uv\omega$
or $v^2\omega^2$, and all remaining testing powers are even in $u,v$.
The scalar raw support is exactly the five terms in (RU13).

For completeness, first replace $\mathsf A=\mathsf y+k\mathsf C$,
$\mathsf B=\mathsf x+\mathsf C$, and
$\mathsf D=k\mathsf C-\mathsf h$. The coefficient requiring a
first conductor is $S=k^2f_A+f_{BB}+f_{BC}+f_{CC}+k^2\varrho$.
Divide its contact-zero difference monically in order $u,\omega,v$.
At $u=k=\varepsilon\jmath$, the $\omega$ divided difference is in
$(\varepsilon v,v^2)$; write it as $\varepsilon U+vV$ with
$U,V\in(v)$. The last $v$ divided difference is in
$(\varepsilon,v)$ because, modulo $\varepsilon$, the remaining
$v$ dependence is even. The first $u$ divided difference restricts
to zero at $\varepsilon=u=v=0$.

These facts permit the exact bounded lift (RI15), which uses the
identity for $\mathsf Cv(\omega-\jmath)$ and never divides by
$\varepsilon$. Subtract the gradient multiple
$2k\varrho_c(-2k,1,1,-1)$. The first three contact gradients have
determinant $-\varepsilon^3\jmath^4E^2\ne0$; the CC5 adjugate
argument cancels a zero identity, and introduces no inverse of that
determinant. Each second-stage coefficient at
$\varepsilon=u=v=0$ is independent of $\omega$: the $\mathsf x$
one restricts to $2f_{BB}+f_{BC}$, and the other three to zero.
Hence every second $\omega$ divided difference is in
$(\varepsilon,v)$ and has the same bounded lift. The final repair is
\[
 k\mathsf C(\mathsf y+\mathsf z-\mathsf h-2k\mathsf x)
      =Q_1^2-\mathsf y\mathsf z.
 \tag{RU17}
\]
The pairings $Q_i$ are defined below. For each complete normal in
(RU15), one conductor suffices: after replacing $\mathsf A,\mathsf B$,
its $\mathsf C$ coefficient restricted at
$\varepsilon=u=v=0$ is independent of $\omega$, directly from
(RU16). This is precisely the normal criterion in the same proof.
When $b$ is a unit, $\varepsilon^{-1}$ is bounded and the ordinary
CC1--CC6 conductor is also available.

All contact value and first-jet hypotheses follow from the stationary
original parent and the four exact contact-zero positive square terms (RU5).
The algorithm is a bounded, source-linear map on the stated raw
coefficient class and preserves the filtrations used next. In
particular it preserves the source ideal $(h,g)$; no separate
contact-zero assertion for its individual coefficient summands is
needed. The source substitution in each monic step uses the actual
$E$, which is a bounded unit through $h=0$.

\paragraph{The complete filtered leading model.}
Set
\[
 \Lambda=\frac1{t(1+s)},\qquad \mu=\frac1{1+s},\qquad
 T_4=\mathsf A^2+(u\mathsf B)^2+(k\mathsf C)^2+
                \mathsf D^2-4k^2\mathsf B\mathsf C.
\]
The model scalar is
\[
\begin{aligned}
 G_{\rm mod}
 &=\Lambda\mathsf x^2+T_4/2-(1/t+1/2)k^2\mathsf x^2
                      +(k^2-u^2)\mathsf C\mathsf x/t\\
 &=\left[\Lambda+\frac{u(u+k)}t-(1/t+1/2)k^2\right]\mathsf x^2
                  +T_4/2-\frac{u+k}t\mathsf x\mathsf z.
\end{aligned}                                                     \tag{RU18}
\]
The second equality is the exact conductor identity
$\mathsf C(u-k)=\mathsf z-u\mathsf x$. Use the first expression
for the raw coefficient filtration, since it retains the even-$u$
class. Its angular raw coefficients are
\[
\begin{aligned}
 F_{BB}&=\Lambda+u^2/2-(1/t+1/2)k^2,\\
 F_{BC}&=-2\Lambda-u^2/t+(3/t-1)k^2,\\
 F_{CC}&=\Lambda+u^2/t-2k^2/t.
\end{aligned}                                                     \tag{RU19}
\]
At $h=g=0$, exact source collection gives the following two
coefficientwise filtrations:
\[
\begin{array}{c|c|c}
 &b\text{ a unit}&b\text{ infinitesimal}\\ \hline
 \text{variable weights}&(k,u,\omega,v)=(1,1,1,2)
     &(\varepsilon,\jmath,u,\omega,v)=(1,1,2,1,3)\\
 \operatorname{wt}(\mathsf x)&2&3\\
 \operatorname{wt}(\mathsf y,\mathsf z,\mathsf h,Q_1,Q_2)&3&5\\
 f_A-1/2,\ \varrho-1/2&\ge2&\ge2\\
 f_{BB}-F_{BB},\ f_{BC}-F_{BC},\ f_{CC}-F_{CC}&\ge4&\ge6\\
 S_B-\mu,\ S_C+\mu&\ge2&\ge4\\
 (R_A,R_B,R_C)&\ge(2,3,3)&\ge(3,5,5).
\end{array}                                                       \tag{RU20}
\]
These statements concern the chosen raw coefficients separately;
polynomial cancellation between raw products is not used to infer
them. Consequently
\[
\begin{array}{c|c|c}
 &b\text{ a unit}&b\text{ infinitesimal}\\ \hline
 G_0-G_{\rm mod}&O_{\rm wt}(8)&O_{\rm wt}(12)\\
 L_M-\mu\mathsf x&O_{\rm wt}(4)&O_{\rm wt}(7)\\
 L_N&O_{\rm wt}(5)&O_{\rm wt}(8).
\end{array}                                                       \tag{RU21}
\]
Here each equality is modulo the source ideal $(h,g)$.

To reproduce the collection, the necessary $b$-infinitesimal source
jets at $h=g=0$ are
\[
\begin{aligned}
 C&=t/2-(t+1)s/2+(t+1)s^2/2+O_{\deg(s,b)}(3),\\
 A_h&=t(1-s+s^2)/2+O_{\deg(s,b)}(3),\\
 \beta_D&=[1-s+s^2-(t+1)sb]/(2t)+O_{\deg(s,b)}(3),\\
 n_1&=(s-s^2-tsb)/2+O_{\deg(s,b)}(3).
\end{aligned}                                                     \tag{RU22}
\]
Also $c=(1-k^2)/t+O_{\rm wt}(6)$,
$E=t(1+k^2)$ at $h=0$, and
$s=t\jmath^2+O_{\rm wt}(6)$. For the unit-$b$ filtration use
$s=k^2/b$ and the respective coefficients of $s$:
$-(t+1)/2,-t/2,-[1+(t+1)b]/(2t),(1-bt)/2$.
The source cancellations (RU11)--(RU12) supply the first and last
rows of (RU20). Scalar weights are even in both
charts. The $M_B$ weights are even in the first and odd in the
second, giving exactly the next weights in (RU21).

Apply the bounded conductor to the exact differences in (RU21).
Their coefficientwise raw filtration makes every correction Gram
coefficient have strictly positive weight. Terms in $(h,g)$ lift
boundedly and remain infinitesimal at arbitrary relative rates.
The whole hard block differs infinitesimally from
$\operatorname{diag}(1/2,(t+1)/2)$, since all its coefficients are
bounded and $u,v,h,g,s$ are infinitesimal.

On the negative-core contact chart take $u,v,\omega>0$. The separated
nonnegative monomials
\[
 a_i^2=(\mathsf A,u\mathsf B,k\mathsf C,\mathsf D),\qquad
 Q_1=a_1a_2-a_3a_4,\quad Q_2=a_1a_3-a_2a_4
\]
have the exact product relation yielding
\[
 T_4=\tfrac12\{(\mathsf y-\mathsf z)^2+\mathsf h^2+
                 \mathsf y^2+(\mathsf z+\mathsf h)^2\}
                  +Q_1^2+Q_2^2.                                \tag{RU23}
\]
Its three-dimensional difference matrix has determinant $1/2>0$.
Each radical separates into source and testing factors.
The full eight-root residue on
$(\mathsf x,\mathsf y,\mathsf z,\mathsf h,Q_1,Q_2,N,M_B)$
therefore consists of half this strict circuit, the $N$ pivot $1/2$,
and the complete coupled block
\[
 \begin{pmatrix}1/t&1/2\\1/2&(t+1)/2\end{pmatrix},\qquad
 \det=\frac{t+2}{4t}>0.                          \tag{RU24}
\]
The additional cross term in the second line of (RU18) has
infinitesimal coefficient. WA10 now gives a finite separated SOS
for the exact remainder. Restore all four actual squares (RU5).
This proves the whole local contact chart in both $b$ alternatives.

\paragraph{The $\omega$-unit state chart.}
If $u,v$ are infinitesimal, $Z=1$, and $\omega$ is a positive unit,
replace the retained centered $n$ positive square term by $nD_uD_vX^2$ itself.
The complete remainder is (RU7) with the two $n_1,n_2$ terms removed.
Use the frame $(A_o,\mathsf x,\mathsf D,N,M_B)$, $A_o=\omega v$.
The exact identities
$\mathsf A=\varepsilon A_o$, $\mathsf B=A_o/\omega$,
$\mathsf C=A_o/\omega-\mathsf x$ use only a pure-testing unit
inverse. In particular the first scalar coefficient becomes
$D_u^2D_v(p^2/2+\mu_2)A_o^2$, and the first $N$ coefficient
becomes $D_u(p-2\mu_1)A_o$; the printed $\varepsilon$ denominators
cancel completely. The full residue has the $A_o,\mathsf D,N$
pivots $1/2,1/2,1/2$ and the same block (RU24).
The complete normal residues are $0,\mathsf x$. Hence this chart
also has a strict finite separated Gram.

\paragraph{A remote cap lemma and complete bounded parents.}
Away from the contact chart there is no need for centered excesses.
Retain the original nonnegative $n,\alpha$ positive square terms and the positive
excesses of $m,A$ above chosen positive real caps $m_c,A_c$.
If $D$ is unbounded, retain its original excess above a real cap
$D_c$ too; otherwise leave $D$ at its actual bounded value.
The remainder is the full original parent with these three
coefficients capped, and all its coefficients are bounded.

We use the elementary finite-dimensional fact that a real quadratic
form positive definite on the common kernel of finitely many real
linear rows becomes positive definite after sufficiently large
positive real multiples of the squared rows are added. Choose
coordinates on the common kernel and its orthogonal complement;
the assertion follows by a Schur complement. Unbounded actual
$m,A,D$, when used, exceed each chosen real cap. Thus every cap
is an exact nonnegative positive square term subtraction in the original parent,
and a strict residue gives WA10 for the entire bounded remainder.
The real caps depend only on real testing residues, never on a
source/testing rate comparison.

\paragraph{Finite remote $v$ with nonzero residue.}
Write $U=u^2,V=v^2$ for real testing residues, with $V>0$.
On the common kernel $Y=M=0$, $X=vW-uZ$, and the remaining
two-state form is exactly
\[
 CD_vU W^2+(D+1/2)(uD_vW-vD_uZ)^2
                  +VD_vW^2/2-D_uVZ^2/2.                         \tag{RU25}
\]
Its determinant is
\[
 VD_uD_v\{DV/2+U[CD-(D+1/2)/2]+C(D+1/2)U^2\}.
 \tag{RU26}
\]
For $t_0>1$, insert $C=t/2,D=1/[2(t-1)]$; it equals
$D_uD_v[V^2+t^2U^2V]/[4(t-1)]>0$.
Its $W^2$ coefficient is positive since $V>0$.
For $t_0=1$, $C=1/2$ and one may choose a real $D$ cap: the braces
are $D(V+U^2)/2+(U^2-U)/4$, positive for sufficiently large $D$.
The cap lemma then chooses $m_c,A_c$ for a strict full state matrix.

\paragraph{Finite unit $u$ with infinitesimal $v$, including zero.}
Use $N=-M=X+uZ-v(W+uY)$, $R_t=uW-Y$, and $V_Z=vZ$.
After the original $n,\alpha$ positive square terms and allowed positive excesses
are retained, the entire remainder lies in the square of the four
roots $(N,R_t,Y,V_Z)$. Only the pure-testing unit $u$ is inverted
to write $W=(R_t+Y)/u$; no inverse of $v$ is needed. Its exact
residue is
\[
\begin{aligned}
 (m_c+1/2)N^2+CR_t^2+A_cD_uY^2
 +(D+1/2)(R_t-D_uV_Z)^2-D_uYV_Z-D_uV_Z^2/2.
\end{aligned}                                                     \tag{RU27}
\]
For $t_0>1$, its restriction to $Y=0$ has determinant
$t^2D_u(D_u-1)/[4(t-1)]>0$ on $(R_t,V_Z)$.
Choose $A_c$ sufficiently large. For $t_0=1$, the common kernel
$Y=0,R_t-D_uV_Z=0$ has value $D_uu^2V_Z^2/2>0$;
large real $D_c,A_c$ give strictness. The normal pivot is positive.
The full four-root support and the residue are checked directly;
literal $v=0$ is their polynomial value.

\paragraph{Both finite angles infinitesimal, and all state charts.}
The exact SD5 pure-core sign identity deals with the nonnegative-core
branch and literal zero axes by original-event squares. On the
negative branch use its standard gauges $u,v>0$ and normalize the
state by a maximal coordinate. If $Z$ is infinitesimal, $W$ is a
unit and the exact BE4 selected four-root identity applies to
$(X,Y,uW,vW)$, with positive caps of $m,A$ provided by (RU3).
All other original events remain as nonnegative positive square terms.
If $Z$ is a unit normalize $Z=1$; the same sign identity gives
bounded positive $\omega$. Its infinitesimal and unit alternatives
are exactly (RU18)--(RU24) and the subsequent $\omega$-unit chart.

\paragraph{Every reciprocal chart and literal projective infinity.}
For unbounded $u$, put $p_u=1/u$ and use the entire homogeneous
polynomial $p_u^2P$. If bounded $v$ has nonzero residue, the full
capped four-state residue is
\[
 CD_vW^2+A_cD_vY^2+(D_c+1/2)(W+vX)^2
                          +(m_c+1/2)(vY-Z)^2\succ0.             \tag{RU28}
\]
Here $D_c$ denotes either a positive real cap or the actual positive
unit residue. The four square rows have common kernel zero because
$v$ is a unit. If instead $v$ is infinitesimal, including zero,
use the exact BE $(\infty,0)$ selected full frame
$(p_uX+Z,W,Y,vX)$ with positive caps of $C,A$. Its strict residue
is the sum of positive multiples of
$(p_uX+Z)^2,(W+vX)^2,W^2,Y^2$; the other original positive square terms are
retained. This handles the zero-$n$ degeneration at that corner.

For unbounded $v$ and bounded $u$, put $q_v=1/v$. The full capped
residue of $q_v^2P$ is
\[
 m_c(W+uY)^2+C(uW-Y)^2+(A_c+1/2)D_uY^2
                         +D_c(uX-Z)^2+D_uX^2/2\succ0.           \tag{RU29}
\]
Its common square kernel is zero for every bounded real $u$,
including zero. Finally, if both angles are unbounded, use
$p_u^2q_v^2P$ and the full frame
$(X,W,Y,P_Z=p_uZ,Q_Z=q_vZ)$. The exact residue is
\[
 CW^2+A_cY^2+D_c(X-P_Z)^2+X^2/2
                           +(m_c+1/2)(Y-Q_Z)^2\succ0.           \tag{RU30}
\]
The two nontrivial block determinants are $D_c/2$ and
$A_c(m_c+1/2)$, both positive. The entire polynomial has this
five-root support, with no inverse of $p_u$ or $q_v$ after
homogenization. All three formulas are checked from the original
parent. Their literal reciprocal zeros are polynomial values;
restoring a nonzero homogeneous factor uses a pure-testing square.

The charts exhaust independent testing pairs and states. Apply each
real identity to real and imaginary parts of a complex state, restore
every retained actual positive square term and separated clearing factor, and use
the exact sign gauges for both angle signs. This completes the stated
all-tests theorem without a restriction on any opening ratio.

\section{Two small actual angles: the R8 and R10 sectors}
\subsection{The opening $D$ gap}
\label{component:pe07041e033da}
\label{pe07041e033da:sec:r17-R8-joint-D}
Assume $t,g$ units, $h,j$ infinitesimal, and $a$ unbounded. Then
$e$ is infinitesimal, $r\sim1$, and $t\sim1$. Use the generic BA
identity above and the exact D identity in
\eqref{eq:c2-ba-D-circuit}. The following normalization and residue
formula hold before the thick/thin source alternatives are used.
Put $L_D=lE/t-rR_D-\sqrt hN_y+uv\xi$ and $M_D=\sqrt D L_D$.
Use the polynomial inverse
\[
 E={t\over l}\{D^{-1/2}M_D+rR_D+\sqrt hN_y-uv\xi\}. \tag{JD5}\label{eq:c2-jd-inverse}
\]
On R8, $t/l$ is a unit, $D^{-1/2}$ is infinitesimal, and $Dj$ is
bounded.  Every coefficient of the new remainder is bounded whenever
the remaining source arrays of R10.9 are bounded.  In particular no
ratio $h/j$, $x/j$, or testing angle divided by a source opening occurs.
The full divergent rank-one block has been normalized jointly.

\paragraph{The exact residue after the joint normalization.}
Put $\sigma=t\lambda x/\sqrt q$ and $\mu=m-a^2q$ as in R10.9.
At the R8 source/testing corner, \eqref{eq:c2-ba-D-circuit} and \eqref{eq:c2-jd-inverse} have the complete residue
\[
 \begin{split}
 M_D^2+r_0^2(C_0+n_0)R_D^2+(Ah)_0N_y^2
  +(n_0+\tfrac12-\mu_0)\xi^2
  +(c_F)_0F_v^2+\lambda_0J^2+R^2+Q^2-2\sigma_0Q\xi.
 \end{split}                                                   \tag{JD6}
\]
This formula applies when $r,n$ are bounded and $r$ a unit.  Its
derivation can retain $n,C,Ah,\mu,c_F,\lambda,\sigma$ as free source
letters: specialize $u=v=\sqrt h=D^{-1/2}=x=\alpha=0$ and $t=1$
in \eqref{eq:c2-ba-D-circuit} and \eqref{eq:c2-jd-inverse}.  In particular the $E,R_D,N_y$ terms become
$r^2(C+n)R_D^2+AhN_y^2$, with no missing $D$ contribution.
The separate $N^2,T^2$ remain intact.

\paragraph{Unbounded $a$: the thick layer.}
Put $s=1/a$, $w=ae$, $b=a^2B$.  The exact R10 source identity
still reads
\[
 n=s^2{wQ_*-bK_*\over2r\Delta_*},\qquad
 Q_*\sim1,\ K_*\sim t,\ \Delta_*\sim1.             \tag{JD11}
\]
Thus $b/w$ is bounded, $b$ is bounded, $g_0=1-w_0>0$, and
\[
\begin{gathered}
 (m/w)_0={1\over2(1-w_0)},\quad C_0=(Ah)_0=1/2,
 \quad (Dj)_0=1/2,\\ \lambda_0=1/2,
 \quad(c_F)_0=(1-w_0+b_0)/2>0.
\end{gathered}                       \tag{JD12}
\]
The source formulas for these quantities have no $j$ denominator;
only $D$ itself is unbounded, and \eqref{eq:c2-ba-D-circuit} removes that pole.
For $\rho=n/e^2$ positive noninfinitesimal, retain the exact R10.10
identities and R10.11 Schur pivot:
\[
 \bar q=a^2q=w^2\{\rho r^2+H\}+a^2\alpha,
 \quad H=m+\tfrac12,\quad \chi=m/w,
 \quad \chi-H\in\operatorname{Inf},
\]
\[
 S=\tfrac12-\mu-{(t\lambda x)^2\over q}
 =H-{m^2\over\bar q}
 \ge {H\rho r^2\over\rho r^2+H}
       +{H^2-\chi^2\over\rho r^2+H}.                 \tag{JD13}
\]
The first term has positive residue and the second is infinitesimal,
also for unbounded $\rho$.  Moreover $x/\sqrt q=b/\sqrt{\bar q}$
is bounded by JD11 and $\bar q\ge Hw^2$.  JD6 consequently has
strict residue on $(M_D,R_D,N_y,\xi,F_v,J,R,Q)$, with the only
possible nonzero cross residue in its strictly positive $Q,\xi$
block.  Restore $N^2,T^2$.  This proves the full local parent, not
just its contact-state restriction.

\paragraph{Unbounded $a$: the thin layer and the exact paired-$D$ repair.}
Assume $0<\rho=n/e^2$ infinitesimal.  JD11 gives
$b/w\sim1/t$, so $\psi=w/b$ is a unit.  The formal C1 contact is
exactly OP4, with native parameters
\[
 x_1=x,\quad\eta_1=t,\quad\zeta_1=s,\quad r_1=h,
 \quad\theta_1=e,\quad t_1=w.
\]
TL2 still gives $\beta_1=-s_n/2$ infinitesimal and
$\kappa=(1+2n)/(1-s_n)\sim1$.  The native gaps satisfy
$M_1=g$ a unit and $J_1=j>0$.  The canceled CA/CS/SE remainders
have no inverse $J_1$; their divergent paired-$D$ mass is retained
in SD3.  The formal source has $C_1,A_1h$ positive units,
$D_1j$ a positive unit, positive $m_1,\alpha_1$, and positive
$\beta_1+1/2$.  These facts follow directly from TL3 and JD12.
Its bounded opening ratio $t_1/e_{1,\mathrm{scale}}^2=w/(tb)$
is a unit.  Thus the explicit CA/CS/SE strict base and their finite
conductors apply exactly as in the R16 thin proof; the reservation
$(2+4\beta_1)\nu-\nu$ is positive because $\beta_1$ is
infinitesimal.  This is an algebraic use of those formulas, not an
application of a physical C1 theorem with negative beta.

Here is the one change needed to the R16 perturbation.  In native
notation put
\[
 T_1=c_a=1+xr,\quad a_1=l/t,\quad
 b_D={J_1\over T_1}={j\over c_a},\quad
 h_1={a_1\over T_1}=1-b_D,
\]
\[
 q_B=v(Z-uX),\quad M_B=(uW-Y)-h_1q_B,
 \quad R_{D,\mathrm{event}}=M_B-b_Dq_B.               \tag{JD14}
\]
The old retained paired-$D$ positive square term in $\kappa F_1/4$ is
$\kappa(D_1+1/2)h_1M_B^2$.  In the actual parent retain instead
the positive contact-zero positive square term $(D_2+1/2)h_1M_B^2$.
For $\delta_D=(D_2+1/2)-\kappa(D_1+1/2)$ the entire resulting
difference is
\[
 \boxed{\quad
 \delta_D\{R_{D,\mathrm{event}}^2-h_1M_B^2\}
 =\eta_D\{M_B^2-2M_Bq_B+b_Dq_B^2\},\quad
 \eta_D=\delta_D b_D.\quad}                          \tag{JD15}
\]
By the exact TL3 coefficients, with $\Theta=n/B$,
\[
 \eta_D={\Theta\{S_cP_D/t-jP_\kappa\}
                    \over K(h+t)c_a},\qquad
       \eta_D/\rho=O(w).                            \tag{JD16}
\]
Every denominator is a unit.  In particular the previously unbounded
$\rho w/j$ correction is absent; no relative $w/j$ bound is needed.

Keep the R16 positive contact-zero $n$ row TL4, the actual centered
$A$ positive square term, and every other unchanged retained positive square term.  The exact
full stress is TL6 with its $\delta_D R_{D,\mathrm{event}}^2$
line replaced by the right side of JD15.  Both parents and all
retained positive square terms have zero contact value and full first jet, so this
new entire stress has the same exact vanishing.

For completeness its new term satisfies every finite-support premise
of the inherited conductor.  The exact inverse state formulas give
\[
 q_B=v\{D_u-uN-puD_uv\mathfrak w\},\qquad p=t_1/R_1.
                                                               \tag{JD17}
\]
At zero normals $q_B=vD_uV_0$.  Hence the new scalar term is in
$(v)^2$, both complete first normals are in $(v)$, and the remaining
hard entries are bounded.  There are no terms above normal degree
two.  Every dependence on $\mathfrak w$ has a factor $p$.
In the CA chart JD16 divided by $\rho$ has bounded coefficients
on the same raw roots as TL8.  In the SE chart
\[
 \Theta/\rho=\varepsilon^2\tau^2QE/K
\]
gives the same extra $\varepsilon^2$ factor to every new coefficient.
Modulo $\varepsilon$, scalar coefficients are even in $u,v$ and
the normal coefficients are independent of $\mathfrak w$.
Thus the exact support, full-jet, and finite divisibility hypotheses
of CC3--CC6 and SE5--SE8 are unchanged.  The complete stress divided
by $\rho$ has a bounded Gram on the same eight roots.  Multiplication
by $\rho$ makes it an infinitesimal perturbation of the strict formal
C1 base.  Restoring the positive square terms proves the actual C2 parent
on the entire thin contact chart, including both full state normals.

At literal $n=0$, OP4 is a physical C1 identification with
$k_1,\zeta_1,r_1$ infinitesimal, $\eta_1=t$ a unit and
$e_{1,\mathrm{scale}}^2=tb$ bounded.  The physical C1.2/C1.3
all-tests theorems allow arbitrary positive $J_1$ and literal beta
zero, so they cover this boundary exactly.

\paragraph{The unit-$W$ entry on the unbounded-$a$ thin layer.}
The actual physical C2 critical coefficient is
\[
 C_*={(h+2Ah)(1+2D)\over2(4AhD-h)}.
\]
Use the polynomial UW3 numerator/denominator with $b=wz_*$,
$z_*=b/w$ bounded.  Its numerator is divisible by $sw$ because it
vanishes identically at $s=0$ and $w=0$.  The denominator after
canceling the $m$ numerator has residue $16t_0=16$, also when
$j_0=0$.  Consequently
\[
                   (C_*-C)/m=s\times O(1).            \tag{JD18}
\]
There is no $j$ inverse in this estimate.  SEO12 permits both $A,D$
unbounded: its entire hard block remains strict and its two radial
masses tend to $2,2$.  With $W=1$, the exact testing identity
$vZ=Z(R_m+X+uZ)/(1+uY)$ absorbs the full negative critical correction
by JD18 and the retained original $m$ square.  All denominators are
pure testing units.  This proves every local normalized unit-$W$
state, including $Z=0$; negative-core state normalization leaves
exactly the $Z=1$, $u,v,\mathfrak w>0$ infinitesimal chart proved
above.  This is the same actual-C2 state entry as UW, with its
source bound checked at the new $t_0=1$ boundary.

\paragraph{Finite remote tests when $e$ is infinitesimal.}
Here $C_*-C$ is infinitesimal by JD18 and bounded $m$.
The full critical-wall identities SEO17, SEO18, SEO20, in the
normalized SEO12 frame, are strict whenever at least one finite
angle has nonzero residue.  They explicitly allow $A,D$ unbounded.
The entire correction $-4(C_*-C)_+D_v(uW-Y)^2$ has infinitesimal
bounded entries in that frame.  Retain original $n,m,\alpha$ and
any positive excess.  This proves every finite remote test, with
all states, for the entire $e$-infinitesimal R8 domain.

\paragraph{Projective endpoints, signs, and the complete join.}
On all of R8, $C$ is a positive unit and $A,D$ are unbounded.
For infinity charts retain the whole original $m,n,\alpha$ positive square terms,
clip $C,A,D$ at suitable positive real lower bounds, and retain their
positive excesses.  The exact homogeneous infinity identities used
in SEO19 and its accompanying BE2 routing require only these lower
bounds and $d=1$; they do not require $D$ bounded or $m$ bounded
below.  In particular the delicate $(u,1/v)=(0,0)$ corner has the
strict selected residue
\[
 2X^2+4c(uW-Y)^2+(4A_*+2)Y^2+4D_*Z^2
\]
on $(X,Y,uW,Z)$ after retaining $m$.  All leftover $W$ occurrences
carry $u$.  The remaining single and double infinities are exactly
the same homogeneous event identities.  Literal endpoint values are
taken in those original homogeneous formulas; no higher-degree
contact graph is substituted at infinity.

Here $a$ is unbounded, so $e$ is infinitesimal and $r\sim1$.
The thick, positive-thin, and literal-zero alternatives for $n/e^2$
are exhaustive. The finite outer calculation above uses the exact
coefficient repair JD18 and the strict forms SEO17, SEO18, SEO20 in
Section~\ref{component:p05a86be0046c}. The displayed homogeneous
forms prove every infinite test. The sign gauges, exact axis identities,
state homogenization and real/imaginary splitting therefore give all
independent tests and states on this unbounded-$a$ R8 domain.
The bounded-$a$ part of R8 is Proposition~\ref{prop:c2-bounded-a}.

\subsection{Source scale, remote tests, and weighted local completion}
\label{component:p9ee4ba7326a7}

Use the native coefficient dictionary, the exact full-parent identity
BA4--BA7, the homogeneous identities BA11--BA15, and the physical
source identity \eqref{eq:c2-ba-physical}. The following bounds are
derived on the present unbounded-$a$ coefficient domain.

\paragraph{Scope and result.}
Retain the original weak-C2 R10 physical stationary source, every strict
source gap, $n\ge0$, the positive-infinitesimal total-scale intersection,
and every inherited exclusion. Thus $t,g,j$ are positive bounded units,
$B,h,e$ are infinitesimal, $a$ is unbounded, and $x=aB$ is infinitesimal.
The new all-tests case proved here is the source subdomain
\[
                         n/e^2\notin\operatorname{Inf}.
\tag{R10.1}
\]
Every nonlocal independent testing chart is also proved on all of R10.
The positive-$n$, infinitesimal-$n/e^2$ local complement is handled by the following thin-layer and state-entry sections; it is not covered by this weighted pivot.
The literal $n=0$ boundary is already an inherited C1 case, by the
exact map specified below. These statements retain all independent rates.

\paragraph{The decisive exact source scale.}
Use the pure-source variables
\[
 s=1/a,\qquad w=ae,\qquad b=a^2B=x/s,
 \quad B=s^2b,\quad x=sb,\quad e=sw,\quad r=1-sw.
\]
Initially $b$ is not assumed bounded. The gap $g=1-w+hx$ makes $w$
bounded and $1-w$ a positive unit in residue. In the exact positive
polynomials $Q_s,K_s$ in \eqref{eq:c2-ba-physical} put
\[
 Q_*=1+B(1+t+ah+2ar+at)
 +B^2\{a(h+t)+a^2[h(r+1)+t(1+r-r^2)]\},
\]
\[
 K_*=(1+s)(h+t)(1+x),\qquad \Delta_*=s\Delta.
\]
Before imposing a bound on $b$, $Q_*\sim1$ and $\Delta_*\sim1$:
write them using bounded $w$, infinitesimal $s,x,h$; their only
nonconstant monomials contain $s$ or $x$. Equation~\eqref{eq:c2-ba-physical} gives exactly
\[
 \boxed{\quad 2r\Delta_*\,n/s^2=wQ_*-bK_* .\quad}
\tag{R10.2}
\]
Consequently
\[
 0<b/w\le Q_*/K_*,\qquad
 0\le n/s^2\le wQ_*/(2r\Delta_*).
\tag{R10.3}
\]
In particular $b=O(w)$ is bounded, $n=O(s^2w)$, $an$ is
infinitesimal, and $ck=\sqrt{b/t}$ is bounded for
$c=\sqrt{a/t}$, $k=\sqrt{x}$. This is an original-source bound,
including equality $n=0$, not a comparison with either testing angle.

After this bound every displayed source denominator is a unit in
$(s,b,w,h,t)$ except the explicit positive monomials already canceled.
Direct finite rational cancellation in the source formulas gives
\[
 \begin{gathered}
 m/w\sim\frac1{2(1-w)},\quad
 C\sim Ah\sim t/2,\quad D\sim\frac1{2(t-1)},\quad
 \alpha/(s^3b^2)\sim t/2,\\
 \lambda=\frac{m-a^2q}{ta^2B}\sim\frac12,\qquad
 c_F\sim\frac{t-w+tb}{2t^2}>0,\qquad
 q=nr^2+(m+\tfrac12)e^2+\alpha.
 \end{gathered}
\tag{R10.4}
\]
The last strictness follows from $t_0>1$ and $w_0<1$.
These are joint residue identities in bounded letters; they do not
set $w$ or $b$ to a positive unit. In particular they hold when either
is infinitesimal. The cancellation defining $\lambda$ leaves no $b$
or $s$ denominator; the new checker verifies this explicitly.
We also have $m$ bounded, $n,\alpha,A^{-1/2}$ infinitesimal, and
\[
 d_4\sim\frac{t}{2(t-1)},\quad b_4\sim0,\quad
 a_E\sim-\frac1{t-1},\quad x_3\sim-1.
\tag{R10.5}
\]
All first-normal arrays printed in
Section~\ref{component:pd709027c7805} are bounded:
the only newly unbounded letter $a$ occurs there through $an$ or
$aq$, both infinitesimal. The factors $a^2n$ in $c_F$ are bounded.

\paragraph{All remote tests on all of R10.}
The exact full-parent identities printed with BA11--BA15 and the exact identities \eqref{eq:c2-literal-u-axis}--\eqref{eq:c2-literal-v-axis} apply algebraically. The
bounded-$a$ proof used $m$ infinitesimal; under R10.4 its entire
nonnegative event positive square term is retained. In BA11 this adds a nonnegative
rank-one form to the strict six-root residue displayed immediately
after that identity; all its coefficients are bounded. For the chart
$u$ infinitesimal, $v$ a bounded unit, keep $S=\sqrt m W$.
The complete $(S,X)$ residue is now
\[
 (v_0S-\sqrt{m_0}X)^2+\frac{1+v_0^2}{2}X^2,
 \qquad \det=\frac{v_0^2(1+v_0^2)}2>0.
\tag{R10.6}
\]
The $(uW,Z)$ block remains BA12 and the normalized $Y$ pivot remains
strict; the extra $uvY,uZ$ terms in the $m$ positive square term are infinitesimal.
For unbounded $v$ its reciprocal identity has the same diagonal
residue BA13 because the potential $\sqrt m X$ term carries $1/v$.
On the remaining remote charts retain $mR_m^2$ as a full positive
square and apply the unchanged strict residue to the remainder.
The original homogeneous identities include either literal infinity,
and the exact literal-axis identities \eqref{eq:c2-literal-u-axis}--\eqref{eq:c2-literal-v-axis} retain $n=0$ without a division. The original
sign gauges and real/imaginary split are unchanged. Hence every test
outside both angles infinitesimal is covered for the entire R10 source.

\paragraph{The exact local full-parent identity.}
Use the generic contact balances RC1, RC6 in
Section~\ref{component:pd709027c7805} and the complete identity
BA4--BA7. These are rational identities of the original source and
never require bounded $a$. Set $U=u^2,V=v^2$, $y=tB$,
\[
 E=tuW-vZ,\quad N_y=(Y-huW)/\sqrt h,\quad \xi=X+ruZ,
\]
\[
 F_v=(V-y)Z,\quad R_D=v(U-x)Z,\quad J=(uv-x/c)Z,
\]
\[
 Q=\sqrt q(avW-uZ-a\xi),\qquad
 N=\sqrt n\{r[(a/t)V-U]Z+u\xi\},
\]
\[
 R=\sqrt\lambda\,tk(cv-u)W,\qquad
 T=\sqrt{2\lambda/c}\,k\sqrt{uv}(tcW-Z),\quad
 \mu=m-a^2q=ta\lambda x.
\]
At positive nonzero angles these are exactly the state-linear roots
in BA4, with no testing-ratio inverse. The entire hard quadratic
$\mathcal H_0$ is the explicit expression BA6. It has no inverse
testing angle or $a$ denominator; after $\eta=\sqrt hN_y$, every
occurrence of $A$ is through $Ah$, $\sqrt hAh$, or $hAh$.
Its R10 local residue is
\[
 \operatorname{diag}\left(\frac1{2(t_0-1)},\frac{t_0}2,\frac12\right)
 \quad\hbox{on }(E,N_y,\xi).
\tag{R10.7}
\]
This permits bounded $m$; every occurrence of $m$ in the remaining
hard entries carries a positive testing degree or $h$.

The principal angular--$\xi$ coupling is resolved as a whole.
Writing $R_c=kv(tcW-Z)$, direct multiplication gives
\[
 -2t\lambda ck R_c\xi+2t\lambda c(R_D-uJ)\xi
 =-2t\lambda x Q\xi/\sqrt q-2\mu\xi^2.
\tag{R10.8}
\]
Thus neither $cu$ nor $cR_D$ is left as a Gram coefficient. The
complete generic identity BA7, with the source arrays just specified,
gives the entire original $F/4$ exactly as
\[
\begin{split}
 &Q^2+N^2+T^2+R^2+\lambda J^2+d_4R_D^2
 +(c_F+b_4x/t^2)F_v^2+(b_4/t^2)vR_DF_v\\
 &\quad+\mathcal H_0-\lambda xE^2-\mu\xi^2
                  -2(t\lambda x/\sqrt q)Q\xi\\
 &\quad+E\{(a_E+c_EV)R_D+(b_E+c_Ex)vF_v\}\\
 &\quad+\sqrt hN_y\{(a_Y+c_YV)R_D+(b_Y+c_Yx)vF_v\}\\
 &\quad+\xi\{x_5uvR_D+u(x_3+x_5x)F_v\}.
\end{split}
\tag{R10.9}
\]
All coefficients are the exact source arrays in
Section~\ref{component:pd709027c7805} and BA5.
This identity retains every scalar, first-normal and hard-normal term.

\paragraph{A strict weighted pivot when $n/e^2$ is noninfinitesimal.}
Put $\rho=n/e^2$, $H=m+1/2$, $\chi=m/w$, and
$\overline q=a^2q$. Then
\[
 \overline q=w^2(\rho r^2+H)+a^2\alpha,
 \qquad \chi-H\in\operatorname{Inf}.
\tag{R10.10}
\]
The $Q,\xi$ source block in R10.9 has Schur pivot
\[
 S=\frac12-\mu-\frac{(t\lambda x)^2}{q}
   =H-\frac{m^2}{\overline q}.
\tag{R10.11}
\]
The equality uses $a^2q+\mu=m$ exactly. Since $\alpha\ge0$,
\[
 S\ge H-\frac{\chi^2}{\rho r^2+H}
 =\frac{H\rho r^2}{\rho r^2+H}
    +\frac{H^2-\chi^2}{\rho r^2+H}.
\tag{R10.12}
\]
Here $H$ is a positive bounded unit. Under R10.1 the first summand
has positive residue and the second is infinitesimal, also when
$\rho$ is unbounded. Hence $S$ is a positive bounded unit.
Moreover $b/w$ is bounded and $\overline q\ge Hw^2$, so
$x/\sqrt q=b/\sqrt{\overline q}$ is bounded. All entries of this
strict source block are therefore bounded.

In the rest of R10.9 the only nonzero off-diagonal residue is
$a_E$ between $E$ and $R_D$, giving
\[
 \frac1{2(t_0-1)}
 \begin{pmatrix}1&-1\\-1&t_0\end{pmatrix}\succ0.
\]
The $N_y,F_v,J,R$ pivots have positive residues
$t_0/2,(t_0-w_0+t_0b_0)/(2t_0^2),1/2,1$.
Every other entry is infinitesimal by R10.4--R10.5 and the explicit
positive testing degrees in R10.9. Retain $N^2,T^2$ as exact positive
squares. The bounded-Gram lemma supplies finite tensor SOS for the
entire parent on R10.1. No source/testing rate comparison occurs.
All roots homogenize without $Z$ or testing-ratio inverses, so literal
axes and $Z=0$ are included directly. The original signs, complex-state
split and the remote theorem complete the all-tests case.

\paragraph{Literal $n=0$ and the exact remaining complement.}
At $n=0$, OP4 is a physical C1 identification with
\[
 x_1=x,\quad\eta_1=t,\quad\zeta_1=s,\quad r_1=h,
 \quad\theta_1=e,\quad e_{1,\mathrm{scale}}=\sqrt{tb}.
\]
Thus $k_1,\zeta_1,r_1$ are infinitesimal, $\eta_1$ a unit, and
$e_{1,\mathrm{scale}}$ bounded. Its infinitesimal type is covered by
the local calculation in Section~\ref{component:pa7e744614f49} together with the outer calculation in Section~\ref{component:p05a86be0046c};
its positive-unit type is covered by the neighboring
Section~\ref{component:pccaa793fd0bf}. Both retain $\beta=0$.
This is C1.2/C1.3 routing, not the unbounded-$\eta$ C1.1 theorem.
For $n>0$ OP5 gives formal $\beta_1<0$, so this physical C1 theorem
cannot be extended by continuity or imported into the neighborhood.
The complement left for the companion thin-layer proof is exactly
\[
       0<n/e^2\in\operatorname{Inf},\qquad u,v\in\operatorname{Inf},
\]
with every original physical and total-scale intersection retained.
In this complement R10.11 can have zero residue: at a unit $w$ and
$b_0=w_0/t_0$, its residue is zero. Thus R10.9 is a useful full-parent
identity there, but its stated strict-Gram argument does not prove that
complement. The companion thin-layer proof and direct state join complete this complement; the final join states the whole-R10 theorem.

\subsection{The positive-$n$ thin layer}
\label{component:pef37475d831e}

\paragraph{Domain.}
All polynomial identities below use the convention $F/4$, as in
R10.9; in particular the original $n$ positive square term is $nD_uD_vX^2$.
All parent and positive square term identities in this section use the common normalization $F/4$.
Keep every original R10 condition and assume
$\rho=n/e^2>0$ infinitesimal. Use $s,w,b,x,B$ from R10.2.
Dividing that exact identity by $w$ gives
\[
 Q_*-(b/w)K_*=2r\Delta_*\rho w,
 \qquad b/w\sim1/t.
\tag{TL1}
\]
Thus $\psi=w/b$ is a positive bounded unit. No testing rate occurs.
The formal native C1 map OP4 is
\[
 x_1=x,\quad\eta_1=t,\quad\zeta_1=s,\quad r_1=h,\quad
 \theta_1=e,\quad t_1=\theta_1/\zeta_1=w.
\]
Write $P_d=P_n/B$. OP3--OP5 give the exact slack
\[
 s_n=2r\Delta_*\rho w/P_d,\quad
 \beta_1=-s_n/2,\quad \pi=1-s_n=1+2\beta_1>0,\quad
 \kappa=(1+2n)/\pi.
\tag{TL2}
\]
Hence $s_n/w=O(\rho)$, $\beta_1$ is infinitesimal, and
$\kappa\sim1$. We substitute these coefficients into the C1 polynomial identities to obtain the comparison below.

\paragraph{Exact coefficient comparison.}
Let $F_1$ be the native C1 parent at this same contact, normalized by
$d_1=1$, and $F_2$ the actual C2 parent, with $d_2=1$.
For a coefficient $L$ put $\delta L=L_2-\kappa L_1$.
Set $K=1+x$, $c_a=1+xr$, $S_c=1+B(1-h)+xr$, $Q_c=1-hB$,
and $\Theta=n/B=\rho w\psi$. The following short polynomials are
bounded with the displayed unit denominators:
\[
 P_\kappa=c_a^2+sx(1+t)+sx^2(r+he+t),
\]
\[
 P_D=tx^2r^2+x\{t-1+e(1-h-2t)\}-1,\quad
 P_\alpha=1+s-2e+x\{r^2+s(r+he)\}.
\]
Exact cancellation of the original C2 and native C1 formulas gives
\[
\begin{array}{c|l}
 \delta m&0\\
 \delta C&-\Theta c_a/(1-h)\\
 \delta(Ah)&-\Theta t c_aS_c/[K(h+t)]\\
 \delta D&\Theta S_cP_D/[K(h+t)tj]\\
 \delta\alpha&-nxP_\alpha/(KQ_c)\\
 \delta d=1-\kappa&-2\Theta P_\kappa/[K(h+t)]\\
 \delta\beta+\delta d/2&-n.
\end{array}
\tag{TL3}
\]
In particular $\delta C,\delta D,\delta d,\delta(Ah)=O(\rho w)$,
while $\delta\alpha=O(nx)$. There is no surviving $1/h$ after the
large $A$ event is centered.

\paragraph{The formal C1 local identities still supply a strict base.}
The native $\eta_1$ is a unit, $k_1,\zeta_1,r_1$ are infinitesimal,
and $e_{1,\mathrm{scale}}^2=tb$ is bounded. Separate its positive-unit
and infinitesimal types. Use respectively the exact centered-$A$
CA5/CA8 construction (or its CS1 scaling) and the SE4/SE9--SE11
construction. This use concerns the printed formulas and conductor
lemmas; the physical-source theorem is not being invoked at $\beta_1<0$.
Here are all changes to their source premises:

The formulas and native stationary identities are rational identities,
valid at the formal OP4 contact. All native gaps remain positive; $M=g$
and $J=j$ are units. The other native event coefficients are positive,
with $C_1,D_1,A_1h$ positive units and $A_1$ unbounded. The opening
ratio $t_1/e_{1,\mathrm{scale}}^2=w/(tb)$ is a positive bounded unit
by TL1, which supplies the only scale bound previously obtained from
$\beta_1\ge0$. All four retained positive square terms SD3 involve the positive
paired coefficient $\beta_1+1/2=\pi/2$, never a separate
$\beta_1$ square. In SE4, the coefficient of the reserved normal positive square term
is $(2+4\beta_1)\nu$, so its subtraction of $\nu$ is still positive,
since $\beta_1$ is infinitesimal. The remaining source bounds, finite
coquotients and strict residue matrices are unchanged. A positive-unit
normalization to $C_1=2$, when used by those formulas, changes none of
these bounds. Consequently their exact local construction yields the
same strict bounded eight-root Gram and positive retained positive square terms for
$\kappa F_1$ on its contact chart.

\paragraph{A new positive contact-zero $n$ row.}
Use the inherited pure-testing coordinate and normals
\[
 \mathfrak w=(W+uY)/D_u,\quad A_t=v\mathfrak w,\quad
 X=N+pD_uA_t-uZ,\quad p=t_1/R_1,
\]
with $D_u=1+u^2,D_v=1+v^2$. In the exact CA chart the contact is
\[
 (u,\mathfrak w,v)=(k,\delta k,\delta k^2E),\quad
 K_c=\delta^2Ek^2=\zeta_1K/Q.
\]
Here $Q,E$ denote the inherited CA source letters, not the angular
weight of R10. The SE chart has $k=\varepsilon j$,
$\delta=1/\varepsilon$, so $K_c=j^2E$.
Put $\gamma=r/K_c$ and retain the entire state-linear square
\[
 \boxed{\quad \mathcal N^2=nD_uD_v(X+\gamma A_t)^2.\quad}
\tag{TL4}
\]
At the actual contact $A_t=K_ck$ and $X=-rk$, so its value and
all first derivatives vanish exactly. The square root of $D_uD_v$ is
pure testing; every other radical or inverse is pure source or pure
testing. The following entire-parent identity includes its full normal:
\[
\begin{split}
 nD_uD_vX^2-\mathcal N^2
 ={}&2n\gamma D_uD_vuA_tZ
 -(2n\gamma pD_u+n\gamma^2)D_uD_vA_t^2\\
 &-2n\gamma D_uD_vA_tN.
\end{split}
\tag{TL5}
\]
There is no remaining hard quadratic from this subtraction. This row
is necessary: without it the scalar $nu^2$ lies outside the old raw
square ideal, even if $n/k^4$ is infinitesimal.

\paragraph{Exact bounded stress after the positive square terms.}
Let $A_{*,1}$ be the inherited bounded centered-$A$ cap. In the actual
parent retain $(A_2-\kappa A_{*,1})D_uD_v(Y-huW)^2$, which is positive
because $A_2$ is unbounded and $\kappa A_{*,1}$ bounded. Comparing with
the corresponding retained positive square term of $\kappa F_1$ replaces the whole
$A$ difference by
\[
 \delta(Ah)D_uD_v(2uWY-hu^2W^2).
\]
After retaining TL4, the full difference of the bounded remainders is
exactly
\[
\begin{split}
 \mathscr E={}&\bigl[nD_uD_vX^2-\mathcal N^2\bigr]
   -nD_v(X+uZ)^2+\delta C D_v(uW-Y)^2\\
 &+(\delta D+\delta d/2)R_{D,\mathrm{event}}^2
   +\delta\alpha D_u(vY-Z)^2\\
 &+(\delta d/2)D_u\{(vY-Z)^2-D_vZ^2\}\\
 &+\delta(Ah)D_uD_v(2uWY-hu^2W^2),
\end{split}
\tag{TL6}
\]
where $R_{D,\mathrm{event}}=uW+uvX-Y-vZ$ is the original event
row. This follows from TL3; in particular the paired beta difference
is exactly $-n$. Every scalar, first-normal and hard-normal term is
present. The full expression and its first derivatives vanish at the
unchanged source contact, because both full parents and every retained
square do.

\paragraph{Finite support and bounds for the unchanged conductor.}
Work on the local chart $Z=1$, with $u,v,\mathfrak w$ positive
infinitesimal. At zero normals set
\[
 V_0=1-puv\mathfrak w,\quad q_B=D_uvV_0,\quad
 uW=u\mathfrak w+u^2h_1vV_0,\quad
 Y=u\mathfrak w-h_1vV_0,\quad h_1=a_1/T_1.
\tag{TL7}
\]
These are exact inherited state inverses. The rows $uW-Y$ and
$R_{D,\mathrm{event}}$ are in $(v)$ at zero normals;
$X+u=pD_uA_t$. The centered $A$ expression in TL6 is a quadratic
in $(v,u\mathfrak w)$ with bounded coefficients. The signed-core
line has no constant: $(vY-1)^2-D_v$ is explicitly
$(vY)^2-2vY-v^2$. The sole state-constant term is
$\delta\alpha D_u$. Finally TL5 is a combination of
$uA_t=(v)(u\mathfrak w)$ and $A_t^2$, with first normal in $(A_t)$.
Thus this is an explicit raw allocation, not an ideal-membership guess.

For the unit-$e_{1,\mathrm{scale}}$ chart use
\[
 (A,B_c,C_c,D_0)
 =(\mathfrak wv,\delta v,\delta E u\mathfrak w,\delta^2k^3E).
\]
Every scalar coefficient of $\mathscr E/\rho$ in
$(A,B_c,C_c)^2+(D_0^2)$ is bounded with unit denominators. Both complete
first normals belong to $(A,B_c,C_c)$ with bounded coefficients, and
the hard-normal entries are bounded. The decisive source coquotients
are
\[
 \Theta/\rho=\delta^2QEw^2/K,\quad
 n\gamma^2/\rho=w^2r^2Q^2/K^2,\quad
 nx/(\rho D_0^2)=Q^2w^2/K^2.
\tag{TL8}
\]
The remaining $n\gamma$ coquotient has an additional $\delta^2k^2E$.
All apparent negative $k$ powers have therefore canceled before the
conductor is used.

For the infinitesimal-$e_{1,\mathrm{scale}}$ chart use the exact SE
parameters $w=t_1=\varepsilon^2\tau$,
$\zeta_1=j^2QE/K$, $k=\varepsilon j$, and
\[
 (A,B_c,C_c,D_0)
 =(\varepsilon\mathfrak wv,v,Eu\mathfrak w,\varepsilon^2j^3E).
\]
The full raw scalar allocation of $\mathscr E/\rho$ is
\[
 f_AA^2+f_{BB}B_c^2+f_{BC}B_cC_c+f_{CC}C_c^2+f_0D_0^2.
\tag{TL9}
\]
All five coefficients are bounded; after division by $\rho$ each has a
factor $\varepsilon^2$. The $A^2$ coefficient has this factor by TL10;
all its other terms carry $p=\varepsilon^2\tau/R_1$, because
$\delta m=0$. The coefficients of $B_c^2,B_cC_c,C_c^2$ have it
from $\Theta/\rho$, and the constant coefficient from the final
identity in TL10. In particular
\[
 \Theta/\rho=\varepsilon^2\tau^2QE/K,
\]
\[
 \frac{n\gamma^2}{\rho\varepsilon^2}
 =\frac{\varepsilon^2\tau^2r^2Q^2}{K^2},\quad
 \frac{n\gamma}{\rho\varepsilon}
 =\frac{\varepsilon^3j^2E\tau^2rQ^2}{K^2},\quad
 \frac{nx}{\rho D_0^2}
 =\frac{\varepsilon^2\tau^2Q^2}{K^2}.
\tag{TL10}
\]
The first-normal coefficients are in $(A,B_c,C_c)$ with bounded
arrays, also divisible by $\varepsilon^2$ after division by $\rho$;
the new $A$ coefficient has the stronger $\varepsilon^3j^2$ factor
printed in TL10. The extra term $-nD_v(X+u)^2$ has scalar coefficient
$-np^2D_u^2D_v$ on $A_t^2$, first normal
$-2npD_uD_vA_tN$, and hard term $-nD_vN^2$; TL10 and
$p=\varepsilon^2\tau/R_1$ bound each normalized coefficient.

Exactly as in SE8, all dependence of the raw coefficients on
$\mathfrak w$ comes through $V_0$ and has a factor
$p=O(\varepsilon^2)$. The additional TL5 coefficients are independent
of $\mathfrak w$. Modulo $\varepsilon$, scalar raw coefficients are
even in $u,v$; the normal raw coefficients are independent of
$\mathfrak w$. Clearing powers of $D_u$ and the source unit
denominators preserves these properties. These are precisely the
finite divisibility premises of SE7: the first scalar conductor has
$T_{\mathfrak w}\in(\varepsilon)$, the second-stage remainders have
$T_{\mathfrak w}\in(\varepsilon,v)$, and each full normal has
$T_{\mathfrak w}\in(\varepsilon)$.

For either chart, TL6 has exact contact value and first-jet zero.
Apply the inherited monic contact lift to $\mathscr E/\rho$:
CC3--CC6 in the first chart and the division-free SE7 lift in the
second. The raw support restriction retains only $D_0^2$, not arbitrary
mixed $D_0$ products. The proofs use only finite monic division and
bounded unit denominators; extending their coefficient ring by the
bounded source letters in TL8--TL10 changes nothing. They give a
bounded Gram matrix on the same eight inherited roots for
$\mathscr E/\rho$. Hence the exact stress $\mathscr E$ has an
infinitesimal Gram matrix on that list. No inverse of $\rho$ occurs in
its final coefficients: division by $\rho$ here records the common
positive source factor before a linear finite lift.

Adding this infinitesimal matrix to the strict formal-C1 base Gram
preserves its positive residue. Restore all its retained positive square terms,
the modified positive centered-$A$ positive square term, and TL4. This proves the
full actual C2 parent on the thin-layer contact chart, with both
complete normals, all hard couplings and every independent rate.
The direct actual-C2 outer/state argument and the all-R10 remote
argument are required for the all-tests join; no negative formal-beta
positive square term is used in either join.

\subsection{The critical wall for unit $W$}
\label{component:pf1aac1dff7c7}

\paragraph{Statement.}
On the entire original R10 source, every independent test with $u,v$
infinitesimal and normalized $W$ a testing unit is tensor SOS. This
includes bounded $Z$ of every sign, infinitesimal $Z$, and $Z=0$.
The theorem uses the actual physical C2 coefficients. No continuation
of a negative-beta C1 source is involved.

Use the R10 source scale $s=1/a$, $w=ae$, $b=a^2B$, and put
$z=b/w$. The source constraint proves $z$ bounded. Write $A_h=Ah$.
The source estimates give $A_h,C,D$ positive units, $h$ infinitesimal,
$m/w$ a positive unit, and $t_0>1$, $w_0<1$.
Define the positive pure-source critical coefficient
\[
 C_*={ (h+2A_h)(1+2D)\over 2(4A_hD-h)}
     ={(1+2A)(1+2D)\over2(4AD-1)}.
 \tag{UW1}
\]
Its denominator after the first normalization is a positive unit.
We prove the opening-uniform estimate
\[
                 {C_*-C\over m}=s\,R,
                 \qquad R\text{ bounded}.             \tag{UW2}
\]

\paragraph{A short polynomial proof of the source estimate.}
All the following letters are polynomials in bounded source variables
$(s,w,z,h,t)$ after $b=wz$, with $x=sb$, $B=s^2b$, $r=1-sw$:
\[
 c_a=1+xr,\quad S_c=1+B(1-h)+sbr,\quad
 g=1-w+hx,\quad J=t c_a-(1-h),\quad \Delta_*=s\Delta,
\]
\[
 M_*=(1+x)(s+r)(h+t)-sw(1+tB)(1-h).
\]
Use the exact source numerators and denominators
\[
 \begin{array}{ll}
 N_C=c_aM_*,&D_C=2(1-h)\Delta_*,\\
 N_A=t c_a(s+r)S_c,&D_A=2\Delta_*,\\
 N_D=(1+x)(s+r)(1-txr)S_c,&D_D=2J\Delta_*,\\
 N_m=w(1+tB)c_aS_c,&D_m=2g\Delta_*.
 \end{array}
\]
Thus $C=N_C/D_C$, $A_h=N_A/D_A$, $D=N_D/D_D$,
and $m=N_m/D_m$. Clear the denominator in UW1:
\[
 \begin{split}
 P={}&(hD_A+2N_A)(D_D+2N_D)D_C\\
     &-2N_C(4N_AN_D-hD_AD_D),\\
 Q={}&2(4N_AN_D-hD_AD_D)D_C.
 \end{split}                                                  \tag{UW3}
\]
Then $C_*-C=P/Q$. At $s=0$ one has
\[
 C={t+h\over2(1-h)},\quad A_h=t/2,\quad
 D={1\over2(t-1+h)},
\]
and substitution gives $C_*=C$ identically. At $w=0$ (hence
$b=0$), the same three formulas hold for every $s$, because all
factors $1+s$ cancel. Therefore the polynomial $P$ is divisible by
both $s$ and $w$, hence $P=swP_1$. Since $N_m=wN_{m,1}$,
\[
 {C_*-C\over m}=s\,{P_1D_m\over QN_{m,1}}.
\]
The denominator $QN_{m,1}$ at $s=h=0$ is $16t>0$.
This proves UW2 by finite polynomial cancellation, without comparing
$w$ to $s$ or $h$.  Its displayed residue also gives
$R\sim t(1-w)(tz+2)$, though the proof only needs boundedness.

\paragraph{Complete parent and state allocation.}
Set $\epsilon=(C_*-C)_+$. By UW2, $\epsilon/m$ is
infinitesimal. If $C>C_*$, retain the nonnegative excess $C-C_*$.
Retain the entire original $n$ and $\alpha$ positive square terms; beta is zero.
The remaining signed core and $C_*,A,D$ positive square terms are exactly the
critical parent SEO11, with $c=C_*$ and $d=1$, because
\[
                  2C_*(4AD-1)=(1+2A)(1+2D).
\]
The exact critical-wall identities SEO12--SEO15 give a bounded strict
Gram on
\[
 B_t=X+uZ,\quad H_A=\sqrt A(Y-yV),\quad
 H_D=\sqrt{4D+2}\{uW-Y+uvB_t-(b_*+u^2)V\},
 \quad uV,\ vV,
\]
where $V=vZ$, $y=1/(2A)$ and
$b_*=(4AD-1)/(A(4D+2))$. Here $y,b_*$ are bounded;
$C_*,D$ are units and $A$ may be unbounded. The hard residue is
strict and the two radial masses have positive residues $2$ and
$2/t_0$. This is an application of those full polynomial identities,
not of their original C1 source theorem.

Normalize $W=1$ by a pure-testing factor. The entire $m$ event obeys
\[
 R_m=v(1+uY)-B_t,\qquad
 V={Z\over1+uY}(R_m+B_t).                         \tag{UW4}
\]
The denominator is a positive testing unit. From the displayed
$H_D$ root,
\[
 uW-Y=(b_*+u^2)V-uvB_t+H_D/\sqrt{4D+2}.
\]
Thus $\sqrt\epsilon(uW-Y)$ is a linear combination of the strict
critical roots and $\sqrt mR_m$ with infinitesimal coefficients.
The only potentially small denominator is accounted for by
$\sqrt{\epsilon/m}$, infinitesimal by UW2. Consequently the complete
negative correction $-4\epsilon(1+v^2)(uW-Y)^2$ has an
infinitesimal bounded Gram. The $m$ positive square term supplies a strict additional
pivot. Restore all retained original positive square terms to obtain the entire
C2 parent, including every hard state direction.

\paragraph{State entry for the remaining local work.}
Normalize a real testing state by a maximal coordinate. On the
negative pure-testing-core branch at $u,v$ infinitesimal, $X,Y$
are infinitesimal. If $W$ is a unit the theorem applies. Otherwise
$Z$ is a unit; set $Z=1$, and then
$w_{\rm test}=(W+uY)/(1+u^2)$ is infinitesimal. The inherited exact
Schur sign identity gives the positive contact chart. Literal axes
have the original event-square exits; complex states split into real
and imaginary parts. Together with the R10 remote theorem, the only
remaining task is the contact chart $u,v,w_{\rm test}>0$
infinitesimal, $Z=1$, on any source not already proved locally.

\label{component:p78dfeedee062}
\label{p78dfeedee062:sec:r17-r4-eunit}
\label{component:p1c9902e9a6f2}
\label{p1c9902e9a6f2:sec:r18-r3-eunit}
\label{component:p8da9ea6ed9e5}
\label{p8da9ea6ed9e5:sec:r17-r5}
\label{component:pd5fa361b4623}
\label{pd5fa361b4623:sec:r17-r6-bounded-a}
\label{component:p9a223c96c696}
\label{p9a223c96c696:sec:r17-r7-bounded-a}
\label{component:p482fcad1a2e2}
\label{p482fcad1a2e2:sec:r17-r8-small-r}
\label{component:p0d4eedd6c7a4}
\label{p0d4eedd6c7a4:sec:r9-weak-A-remote}
\label{component:p99372d145f17}
\label{p99372d145f17:sec:r11-small-r-local}
\label{component:p34609a6519a6}
\label{p34609a6519a6:sec:r11-small-r-remote}


\backmatter
\chapter{Reference}
\begingroup
\let\chapter\section
\begin{thebibliography}{9}
\bibitem{scheiderer}
C. Scheiderer, \emph{Second-order cone representation for convex sets
in the plane}, SIAM J. Appl. Algebra Geom. \textbf{5}, 114--139 (2021),
arXiv:2004.04196v2.
The general tensor-evaluation criterion used here is Corollary 3.19,
p.~15; the quantitative tensor-rank criterion is Theorem 3.10.
\url{https://arxiv.org/abs/2004.04196}.
\end{thebibliography}
\endgroup
\end{document}

